\documentclass[11pt,reqno]{amsart}
\usepackage{iftex}
\ifPDFTeX
  \usepackage[T1]{fontenc}
  \usepackage[utf8]{inputenc}
  \usepackage{lmodern}
\else
  \usepackage[no-math]{fontspec}
\fi
\usepackage{amsmath}
\usepackage{amssymb}
\usepackage{amsfonts}
\usepackage{mathrsfs}
\usepackage{amsthm}
\usepackage{mathtools}
\usepackage{geometry}
\usepackage{enumitem}
\usepackage{longtable}
\usepackage{booktabs}
\usepackage{array}
\usepackage{microtype}
\usepackage{graphicx}
\graphicspath{{../figures/}{figures/}}
\usepackage[dvipsnames]{xcolor}
\usepackage{listings}
\lstset{basicstyle=\ttfamily\scriptsize,breaklines=true,columns=fullflexible,keepspaces=true,
  frame=single,framerule=0.3pt,showstringspaces=false,keywordstyle=\color{NavyBlue},
  commentstyle=\color{OliveGreen},stringstyle=\color{BrickRed},upquote=true}
\lstdefinestyle{out}{basicstyle=\ttfamily\scriptsize,language={},keywordstyle={},commentstyle={},stringstyle={}}
\usepackage{float}
\usepackage[pagebackref=true]{hyperref}
\usepackage{bookmark}
\usepackage{pdfcomment}

% Hover previews for internal references (no layout change).
\makeatletter
\let\CG@origref\ref
\let\CG@origeqref\eqref
\newcommand{\CG@gettip}[1]{\@ifundefined{CGTip@#1}{Reference: #1}{\csname CGTip@#1\endcsname}}
\DeclareRobustCommand{\CGref}[1]{\pdftooltip{\CG@origref{#1}}{\CG@gettip{#1}}}
\DeclareRobustCommand{\CGeqref}[1]{\pdftooltip{\CG@origeqref{#1}}{\CG@gettip{#1}}}
\makeatother
% Tooltip texts shown when hovering over internal references.
\makeatletter
\expandafter\gdef\csname CGTip@sec:notation\endcsname{Section : Notation (p. 6)}
\expandafter\gdef\csname CGTip@sec:introduction\endcsname{Section 1: Introduction (p. 10)}
\expandafter\gdef\csname CGTip@subsec:background\endcsname{Section 1.1: Background (p. 10)}
\expandafter\gdef\csname CGTip@eq:mertens-def\endcsname{Eq. (1.1), p. 11: M(x):=sum\_n<= xmu(n).}
\expandafter\gdef\csname CGTip@subsec:known\endcsname{Section 1.2: Known results on the Riemann--von Mangoldt law, the ordinates and the clock (p. 11)}
\expandafter\gdef\csname CGTip@thm:known-RvM\endcsname{Theorem 1.1 (Riemann--von Mangoldt; [ref], [ref, Thms. 9.3--9.4]) (p. 11): For T>0 not an ordinate, N\_zeta(T)=1+(1)/(pi)vartheta(T)+S(T), with 1+(1)/(pi)vartheta(T)=(T)/(2pi)log(T)/(2pi e)+(7)/(8)+O(T\textasciicircum{}-1) and S(T)=O(log T).}
\expandafter\gdef\csname CGTip@thm:known-S\endcsname{Theorem 1.2 (Littlewood; Selberg; [ref], [ref, Thm. 9.9], [ref]) (p. 11): int\_0\textasciicircum{}TS(t) dt=O(log T), and int\_0\textasciicircum{}TS(t)\textasciicircum{}2 dt=(T)/(2pi\textasciicircum{}2)log log T+O(T(log log T)\textasciicircum{}1/2).}
\expandafter\gdef\csname CGTip@thm:known-landau\endcsname{Theorem 1.3 (Landau's formula; [ref], uniform version [ref]) (p. 11): For fixed x>1, sum\_0<gamma <= Tx\textasciicircum{}rho=-(T)/(2pi)Lambda(x)+O\_x(log T), where Lambda(x)=0 unless x is a prime power.}
\expandafter\gdef\csname CGTip@thm:known-hlawka\endcsname{Theorem 1.4 (Hlawka; [ref], see also [ref]) (p. 11): For every real alpha !=0 the sequence (alpha gamma\_k)\_k>=1 is uniformly distributed modulo 1; in particular (1)/(N)sum\_k<= Ne\textasciicircum{}iugamma\_k->0 for every u!=0.}
\expandafter\gdef\csname CGTip@thm:known-cramer\endcsname{Theorem 1.5 (Cramer; [ref]) (p. 11): The function V\_+(z)=sum\_Im rho>0e\textasciicircum{}rho z, holomorphic in Im z>0, has a meromorphic continuation to the plane cut along the non-positive imaginary axis; its poles are simple and lie at the points +-log p\textasciicircum{}m.}
\expandafter\gdef\csname CGTip@subsec:midpoint-system\endcsname{Section 1.3: The midpoint RVM system (p. 12)}
\expandafter\gdef\csname CGTip@eq:intro-clock\endcsname{Eq. (1.2), p. 12: N\_RVM(T)=(T)/(2pi)log(T)/(2pi e)+(7)/(8),}
\expandafter\gdef\csname CGTip@eq:intro-lambert\endcsname{Eq. (1.3), p. 12: tau\_k=(2pi(k-11/8))/(W\_0((k-11/8)/e)) .}
\expandafter\gdef\csname CGTip@eq:intro-zeta-rvm\endcsname{Eq. (1.4), p. 12: P\_RVM\textasciicircum{}crit(s)=prod\_k>=1(1+(s(s-1))/((1)/(4)+tau\_k\textasciicircum{}2)), zeta\_RVM(s)=(P\_RVM\textasciicircum{}crit(s))/(G(s)), G(s)=s(s-1)Gamma(s/2)pi\textasciicircum{}-s/2.}
\expandafter\gdef\csname CGTip@eq:intro-pi\endcsname{Eq. (1.5), p. 13: pi\_RVM(x)=sum\_1<= m<=log\_2x(mu(m))/(m) J\_RVM\textasciicircum{}*(x\textasciicircum{}1/m), J\_RVM\textasciicircum{}*(x)=sum\_1<= m<=log\_2x(1)/(m) pi\_RVM(x\textasciicircum{}1/m).}
\expandafter\gdef\csname CGTip@eq:intro-real-primes\endcsname{Eq. (1.6), p. 13: p\_j\textasciicircum{}RVM:=pi\_RVM\textasciicircum{}-1(j), pi\_RVM(p\_j\textasciicircum{}RVM)=j,}
\expandafter\gdef\csname CGTip@eq:intro-transition\endcsname{Eq. (1.7), p. 13: zeta\_RVM(s)=varpi\_RVM(1+w(s))(1+O(1/|t|)), w(s)=exp -2pi i N\_RVM (-i(s-(1)/(2))),}
\expandafter\gdef\csname CGTip@eq:intro-Mbound\endcsname{Eq. (1.8), p. 13: M\_RVM\textasciicircum{}crit(x)=L+O\_eps(x\textasciicircum{}1/2+eps)}
\expandafter\gdef\csname CGTip@eq:intro-muh\endcsname{Eq. (1.9), p. 14: mu\_h(n)=mu(n)\textasciicircum{}2varsigma\_n,varsigma\_nin -1,+1, M\_h(n):=sum\_m<= nmu\_h(m),}
\expandafter\gdef\csname CGTip@eq:intro-quantizer\endcsname{Eq. (1.10), p. 14: varsigma\_n= +1, M\_h(n-1)< M(n),; -1, M\_h(n-1)>= M(n),}
\expandafter\gdef\csname CGTip@subsec:core\endcsname{Section 1.4: The core of the method (p. 14)}
\expandafter\gdef\csname CGTip@eq:intro-family\endcsname{Eq. (1.11), p. 14: gamma\_k(lambda)=(1-lambda)gamma\_k+lambda tau\_k, zeta\_lambda(s)=(1)/(G(s))prod\_k>=1(1+(s(s-1))/((1)/(4)+gamma\_k(lambda)\textasciicircum{}2)),}
\expandafter\gdef\csname CGTip@eq:intro-hamburger\endcsname{Eq. (1.12), p. 14: zeta\_Z1=zeta <=> (D): zeta\_Z1 is an absolutely convergent Dirichlet series in Re s>1 .}
\expandafter\gdef\csname CGTip@eq:intro-varpi-Z1\endcsname{Eq. (1.13), p. 15: varpi\_Z1=lim\_sigma ->+infzeta\_Z1(sigma)=varpi\_RVMlim\_K->infprod\_k<= K((1)/(4)+tau\_k\textasciicircum{}2)/((1)/(4)+gamma\_k\textasciicircum{}2) >=1,}
\expandafter\gdef\csname CGTip@eq:intro-step3\endcsname{Eq. (1.14), p. 15: zeta\_Z1=zeta => M(x)=O\_eps (x\textasciicircum{}1/2+eps)(eps>0),}
\expandafter\gdef\csname CGTip@eq:intro-prime-chain\endcsname{Eq. (1.15), p. 15: tau\_k->zeta\_RVM-> J\_RVMlongleftrightarrow pi\_RVM->(dpi\_RVM, p\_j\textasciicircum{}RVM, c\_n\textasciicircum{}RVM),}
\expandafter\gdef\csname CGTip@eq:intro-M-chain\endcsname{Eq. (1.16), p. 15: F\_RVM=zeta\_RVM\textasciicircum{}-1-> M\_RVM\textasciicircum{}crit->(M,mu\_h,M\_h),}
\expandafter\gdef\csname CGTip@eq:intro-family-chain\endcsname{Eq. (1.17), p. 15: zeta\_RVM=zeta\_1-> lambda decreasing to 0 zeta\_0=zeta\_Z1?=zeta => M(x)=O\_eps(x\textasciicircum{}1/2+eps).}
\expandafter\gdef\csname CGTip@subsec:intro-expsums\endcsname{Section 1.5: Exponential sums over the ordinates and the Riemann hypothesis (p. 15)}
\expandafter\gdef\csname CGTip@eq:intro-centroid\endcsname{Eq. (1.18), p. 15: C\textasciicircum{}pathc(N;u)=(1)/(N)sum\_n<= Nsum\_k<= ne\textasciicircum{}iuo\_k=sum\_k<= N(1-(k-1)/(N))e\textasciicircum{}iuo\_k}
\expandafter\gdef\csname CGTip@eq:intro-imaginary\endcsname{Eq. (1.19), p. 16: RH<=>sum\_k>=1e\textasciicircum{}-y\_0gamma\_k=e\textasciicircum{}-iy\_0/2V\_+(iy\_0)}
\expandafter\gdef\csname CGTip@eq:intro-window\endcsname{Eq. (1.20), p. 16: RH<=>sum\_k>=1hat f(gamma\_k)=int\_J f(u) e\textasciicircum{}-u/2V\_+(u) dufor every fin C\_c\textasciicircum{}inf(J),}
\expandafter\gdef\csname CGTip@eq:intro-centroid-value\endcsname{Eq. (1.21), p. 16: e\textasciicircum{}-1/2V\_+(1)=(1)/(2sqrt(e))(e-(1)/(e\textasciicircum{}2-1))+i 0.76785566886...=0.77689426079...+0.76785566886... i}
\expandafter\gdef\csname CGTip@subsec:status\endcsname{Section 1.6: Originality, scope and status (p. 16)}
\expandafter\gdef\csname CGTip@subsec:organization\endcsname{Section 1.7: Organization (p. 20)}
\expandafter\gdef\csname CGTip@sec:clock\endcsname{Section 2: The midpoint Riemann--von Mangoldt clock (p. 21)}
\expandafter\gdef\csname CGTip@eq:zero-count\endcsname{Eq. (2.1), p. 21: N\_zeta(T):=no.rho: 0<Im rho <= T=no.k:gamma\_k<= T.}
\expandafter\gdef\csname CGTip@eq:RvM-full\endcsname{Eq. (2.2), p. 21: N\_zeta(T)=(T)/(2pi)log(T)/(2pi e)+(7)/(8)+S(T)+O((1)/(T)), S(T):=(1)/(pi)arg zeta ((1)/(2)+iT)=O(log T),}
\expandafter\gdef\csname CGTip@eq:clock\endcsname{Eq. (2.3), p. 21: N\_RVM(T):=(T)/(2pi)log(T)/(2pi e)+(7)/(8) .}
\expandafter\gdef\csname CGTip@eq:midpoint-condition\endcsname{Eq. (2.4), p. 21: N\_RVM(tau\_k)=k-(1)/(2), k>=1,}
\expandafter\gdef\csname CGTip@prop:lambert\endcsname{Proposition 2.1 (Lambert inversion) (p. 21): For every k>=1, tau\_k=(2pi(k-11/8))/(W\_0((k-11/8)/e)), The quotient is positive also for k=1.}
\expandafter\gdef\csname CGTip@eq:lambert\endcsname{Eq. (2.5), p. 21: tau\_k=(2pi(k-11/8))/(W\_0((k-11/8)/e)),}
\expandafter\gdef\csname CGTip@eq:lambert-proof-y\endcsname{Eq. (2.6), p. 21: y(log y-1)=k-(11)/(8).}
\expandafter\gdef\csname CGTip@eq:clock-derivative\endcsname{Eq. (2.7), p. 21: N\_RVM'(T)=(1)/(2pi)log(T)/(2pi),}
\expandafter\gdef\csname CGTip@eq:clock-growth\endcsname{Eq. (2.8), p. 22: tau\_k\textasciitilde{}(k)/(log k)(k>=2).}
\expandafter\gdef\csname CGTip@eq:RvM-formula\endcsname{Eq. (2.9), p. 22: N\_zeta(T)=N\_RVM(T)+S(T)+O(T\textasciicircum{}-1), S(T)=O(log T)(T>=2).}
\expandafter\gdef\csname CGTip@eq:symmetric-count\endcsname{Eq. (2.10), p. 22: N\_zeta\textasciicircum{}*(T):=(1)/(2)(N\_zeta(T+0)+N\_zeta(T-0))}
\expandafter\gdef\csname CGTip@eq:eta-def\endcsname{Eq. (2.11), p. 22: eta\_k:=tau\_k-gamma\_k .}
\expandafter\gdef\csname CGTip@lem:bounded-displacement\endcsname{Lemma 2.2 (Bounded displacement) (p. 22): There is a constant C\_eta such that |eta\_k|<= C\_eta (k>=1).}
\expandafter\gdef\csname CGTip@eq:bounded-displacement\endcsname{Eq. (2.12), p. 22: |eta\_k|<= C\_eta (k>=1).}
\expandafter\gdef\csname CGTip@eq:block-count\endcsname{Eq. (2.13), p. 22: |N\_zeta\textasciicircum{}*(gamma\_k)-(k-(1)/(2))|<=(m(gamma\_k))/(2) .}
\expandafter\gdef\csname CGTip@eq:clock-discrepancy\endcsname{Eq. (2.14), p. 22: |N\_RVM(gamma\_k)-N\_RVM(tau\_k)|=O(log gamma\_k).}
\expandafter\gdef\csname CGTip@sec:zeta-rvm\endcsname{Section 3: The pure critical RVM zeta function (p. 22)}
\expandafter\gdef\csname CGTip@eq:q-def\endcsname{Eq. (3.1), p. 22: q(s):=s(s-1),}
\expandafter\gdef\csname CGTip@eq:Pcrit\endcsname{Eq. (3.2), p. 22: P\_RVM\textasciicircum{}crit(s):=prod\_k>=1(1+(q(s))/((1)/(4)+tau\_k\textasciicircum{}2)),}
\expandafter\gdef\csname CGTip@eq:Gdef\endcsname{Eq. (3.3), p. 22: G(s):=s(s-1)Gamma(s/2)pi\textasciicircum{}-s/2,}
\expandafter\gdef\csname CGTip@eq:zetaRVM\endcsname{Eq. (3.4), p. 22: xi\_RVM(s):=(1)/(2)P\_RVM\textasciicircum{}crit(s), zeta\_RVM(s):=(P\_RVM\textasciicircum{}crit(s))/(G(s)).}
\expandafter\gdef\csname CGTip@thm:canonical\endcsname{Theorem 3.1 (Canonical characterization of zeta\_RVM) (p. 22): The product (3.2) converges absolutely and locally uniformly in sin C; as a function of the symmetric variable q=q(s) it is a canonical product of genus zero. It defines an entire function with P\_RVM\textasciicircum{}crit(1-s)=P\_RVM\textasciicircum{}crit(s), P\_RVM\textasciicircum{}crit(conj s)=conj P\_RVM\textasciicircum{}crit(s), P\_RVM\textasciicircum{}crit(0)=P\_RVM\textasciicircum{}crit(1)=1 . Its zeros are exactly...}
\expandafter\gdef\csname CGTip@eq:P-symmetry\endcsname{Eq. (3.5), p. 22: P\_RVM\textasciicircum{}crit(1-s)=P\_RVM\textasciicircum{}crit(s), P\_RVM\textasciicircum{}crit(conj s)=conj P\_RVM\textasciicircum{}crit(s), P\_RVM\textasciicircum{}crit(0)=P\_RVM\textasciicircum{}crit(1)=1 .}
\expandafter\gdef\csname CGTip@eq:synthetic-divisor\endcsname{Eq. (3.6), p. 22: rho\_k\textasciicircum{}RVM=(1)/(2)+itau\_k, conj rho\_k\textasciicircum{}RVM=(1)/(2)-itau\_k, k>=1 .}
\expandafter\gdef\csname CGTip@eq:canonical-normal\endcsname{Eq. (3.7), p. 23: sup\_sin K|(q(s))/((1)/(4)+tau\_k\textasciicircum{}2)|<<\_Ktau\_k\textasciicircum{}-2,}
\expandafter\gdef\csname CGTip@eq:factor-zero\endcsname{Eq. (3.8), p. 23: (s-(1)/(2))\textasciicircum{}2+tau\_k\textasciicircum{}2=0 .}
\expandafter\gdef\csname CGTip@sec:exp-sums\endcsname{Section 4: RVM exponential sums, centroids, and the height transform (p. 23)}
\expandafter\gdef\csname CGTip@eq:S-RVM\endcsname{Eq. (4.1), p. 23: S\_RVM(N;u):=sum\_k<= Ne\textasciicircum{}iutau\_k, C\_RVM\textasciicircum{}pts(N;u):=(S\_RVM(N;u))/(N),}
\expandafter\gdef\csname CGTip@eq:path-centroid\endcsname{Eq. (4.2), p. 23: C\_RVM\textasciicircum{}pathc(N;u):=(1)/(N)sum\_n<= NS\_RVM(n;u) =sum\_k<= N(N-k+1)/(N)e\textasciicircum{}iutau\_k.}
\expandafter\gdef\csname CGTip@thm:exp-sum\endcsname{Theorem 4.1 (Uniform logarithmic exponential-sum bound) (p. 23): Let 0<u\_1<= u\_2<inf. Then sup\_uin[u\_1,u\_2]|S\_RVM(N;u)|<<\_u\_1,u\_21+log N .}
\expandafter\gdef\csname CGTip@eq:exp-sum-bound\endcsname{Eq. (4.3), p. 23: sup\_uin[u\_1,u\_2]|S\_RVM(N;u)|<<\_u\_1,u\_21+log N .}
\expandafter\gdef\csname CGTip@eq:inverse-clock\endcsname{Eq. (4.4), p. 23: T'(y)=(2pi)/(log(T(y)/2pi)), T''(y)=-(4pi\textasciicircum{}2)/(T(y)log\textasciicircum{}3(T(y)/2pi)) .}
\expandafter\gdef\csname CGTip@eq:rvm-gap-bounds\endcsname{Eq. (4.5), p. 23: (2pi)/(l\_k+1)<=tau\_k+1-tau\_k<=(2pi)/(l\_k),}
\expandafter\gdef\csname CGTip@eq:rvm-second-difference\endcsname{Eq. (4.6), p. 23: |tau\_k+2-2tau\_k+1+tau\_k|<=(4pi\textasciicircum{}2)/(tau\_kl\_k\textasciicircum{}3).}
\expandafter\gdef\csname CGTip@eq:tau-ell-k\endcsname{Eq. (4.7), p. 23: tau\_kl\_k\textasciitilde{} k(k>=2).}
\expandafter\gdef\csname CGTip@eq:Ek-identity\endcsname{Eq. (4.8), p. 23: E\_k= b\_k\textasciicircum{}rm sbp(E\_k+1-E\_k).}
\expandafter\gdef\csname CGTip@eq:summation-by-parts\endcsname{Eq. (4.9), p. 24: sum\_k=k\_0\textasciicircum{}NE\_k= b\_N\textasciicircum{}rm sbpE\_N+1- b\_k\_0\textasciicircum{}rm sbpE\_k\_0+sum\_k=k\_0+1\textasciicircum{}N(b\_k-1\textasciicircum{}rm sbp- b\_k\textasciicircum{}rm sbp)E\_k .}
\expandafter\gdef\csname CGTip@eq:kappa-bounds\endcsname{Eq. (4.10), p. 24: | b\_k\textasciicircum{}rm sbp|<=(pi)/(2Delta\_k)<=(l\_k+1)/(4u\_1), | b\_k-1\textasciicircum{}rm sbp- b\_k\textasciicircum{}rm sbp|<=|Delta\_k-Delta\_k-1| | b\_k-1\textasciicircum{}rm sbp|| b\_k\textasciicircum{}rm sbp| <=(pi\textasciicircum{}2u\_2)/(4u\_1\textasciicircum{}2) (l\_k+1\textasciicircum{}2)/(tau\_k-1l\_k-1\textasciicircum{}3)<<\_u\_1,u\_2(1)/(tau\_kl\_k)<<(1)/(k),}
\expandafter\gdef\csname CGTip@cor:point-centroid\endcsname{Corollary 4.2 (Point centroid) (p. 24): Uniformly for uin[u\_1,u\_2]in(0,inf), C\_RVM\textasciicircum{}pts(N;u)=O\_u\_1,u\_2 ((1+log N)/(N))->0 . By direct averaging, Theorem 4.1 gives only C\_RVM\textasciicircum{}pathc(N;u)=O(log N) for the path centroid (4.2), not even boundedness; the path centroid converges by a second summation by parts (Theorem 5.13).}
\expandafter\gdef\csname CGTip@eq:point-centroid\endcsname{Eq. (4.11), p. 24: C\_RVM\textasciicircum{}pts(N;u)=O\_u\_1,u\_2 ((1+log N)/(N))->0 .}
\expandafter\gdef\csname CGTip@eq:n-RVM\endcsname{Eq. (4.12), p. 24: n\_RVM(H):=no.k:tau\_k<= H.}
\expandafter\gdef\csname CGTip@eq:MH-RVM\endcsname{Eq. (4.13), p. 24: M\_H\textasciicircum{}RVM(x):=(1)/(n\_RVM(H))sum\_tau\_k<= Hx\textasciicircum{}itau\_k,}
\expandafter\gdef\csname CGTip@eq:MH-genuine\endcsname{Eq. (4.14), p. 24: M\_H(x):=(1)/(N\_zeta(H))sum\_gamma\_k<= Hx\textasciicircum{}igamma\_k.}
\expandafter\gdef\csname CGTip@cor:MH-decay\endcsname{Corollary 4.3 (Synthetic height-transform decay) (p. 24): For every fixed x>1, with N=n\_RVM(H), M\_H\textasciicircum{}RVM(x)=O\_x ((1+log N)/(N))->0(H->inf).}
\expandafter\gdef\csname CGTip@eq:MH-decay\endcsname{Eq. (4.15), p. 24: M\_H\textasciicircum{}RVM(x)=O\_x ((1+log N)/(N))->0(H->inf).}
\expandafter\gdef\csname CGTip@rem:MH-not-mertens\endcsname{Remark 4.4 (p. 24): Neither (4.13) nor (4.14) is the Mertens function; both are normalized phase averages.}
\expandafter\gdef\csname CGTip@sec:ordinate-sums\endcsname{Section 5: Exponential sums over the ordinates of zeta: centroids, Cramer's function and the window form of the Riemann hypothesis (p. 24)}
\expandafter\gdef\csname CGTip@eq:f-hat\endcsname{Eq. (5.1), p. 25: hat f(z):=int\_R f(u) e\textasciicircum{}iuz du(zin C);}
\expandafter\gdef\csname CGTip@eq:f-hat-decay\endcsname{Eq. (5.2), p. 25: |hat f(t+iy)|<= C\_ f,M e\textasciicircum{}u\_2|y|(1+|t|)\textasciicircum{}-M(t,yin R, M>=0).}
\expandafter\gdef\csname CGTip@subsec:cramer\endcsname{Section 5.1: The ordinate transform and Cramer's function (p. 25)}
\expandafter\gdef\csname CGTip@eq:ordinate-transforms\endcsname{Eq. (5.3), p. 25: Z\_gamma(w):=sum\_k>=1e\textasciicircum{}iwgamma\_k, Z\_tau(w):=sum\_k>=1e\textasciicircum{}iwtau\_k, Z\_xi(w):=sum\_k>=1e\textasciicircum{}w(rho\_k-(1)/(2))=e\textasciicircum{}-w/2 V\_+(w),}
\expandafter\gdef\csname CGTip@eq:cramer-V\endcsname{Eq. (5.4), p. 25: V\_+(z):=sum\_Im rho>0e\textasciicircum{}rho z(Im z>0)}
\expandafter\gdef\csname CGTip@eq:Z-gamma-growth\endcsname{Eq. (5.5), p. 25: | Z\_gamma(u+iy)|<=sum\_ke\textasciicircum{}-ygamma\_k<= C y\textasciicircum{}-2(uin R, 0<y<=1).}
\expandafter\gdef\csname CGTip@eq:ordinate-defect\endcsname{Eq. (5.6), p. 25: X(w):= Z\_gamma(w)- Z\_xi(w)(Im w>0).}
\expandafter\gdef\csname CGTip@prop:pair-decomposition\endcsname{Proposition 5.1 (Pair decomposition of the defect) (p. 25): For Im w>0, X(w)=-sum\_k>=1(cosh(wd\_k)-1)e\textasciicircum{}iwgamma\_k, the series converging absolutely; here k runs over both members of each pair (1)/(2)+- d+igamma, with multiplicity, and only the zeros with d\_k!=0 contribute. In particular X==0 under RH, and for y>0 X(iy)=sum\_k>=1(1-cos(yd\_k))e\textasciicircum{}-ygamma\_k >=0 .}
\expandafter\gdef\csname CGTip@eq:pair-decomposition\endcsname{Eq. (5.7), p. 25: X(w)=-sum\_k>=1(cosh(wd\_k)-1)e\textasciicircum{}iwgamma\_k,}
\expandafter\gdef\csname CGTip@eq:defect-imaginary\endcsname{Eq. (5.8), p. 25: X(iy)=sum\_k>=1(1-cos(yd\_k))e\textasciicircum{}-ygamma\_k >=0 .}
\expandafter\gdef\csname CGTip@prop:imaginary-frequency\endcsname{Proposition 5.2 (RH at one imaginary frequency) (p. 25): For every y\_0in(0,4pi], RH<=>sum\_k>=1e\textasciicircum{}-y\_0gamma\_k=e\textasciicircum{}-iy\_0/2 V\_+(iy\_0) .}
\expandafter\gdef\csname CGTip@eq:imaginary-frequency\endcsname{Eq. (5.9), p. 25: RH<=>sum\_k>=1e\textasciicircum{}-y\_0gamma\_k=e\textasciicircum{}-iy\_0/2 V\_+(iy\_0) .}
\expandafter\gdef\csname CGTip@prop:cramer-zero-free\endcsname{Proposition 5.3 (Zero-free representation of Cramer's function) (p. 26): Let a>0 and 0<theta<(pi)/(2), and put g(s):=(1)/(s)+(1)/(s-1)-(1)/(2)log pi+(1)/(2) (Gamma')/(Gamma)((s)/(2)),so that (xi')/(xi)(s)=g(s)+(zeta')/(zeta)(s). For Im z>0, 2pi i V\_+(z)= V\_rm b(z)+ V\_rm R(z)+ V\_rm L(z)+ V\_rm p(z), V\_rm b(z) =int\_-a\textasciicircum{}1+ae\textasciicircum{}sigma z (xi')/(xi)(sigma) dsigma, V\_rm R(z)=i e\textasciicircum{}(1+a)zint\_0\textasciicircum{}inf e\textasciicircum{}it...}
\expandafter\gdef\csname CGTip@eq:g-arch\endcsname{Eq. (5.10), p. 26: g(s):=(1)/(s)+(1)/(s-1)-(1)/(2)log pi+(1)/(2) (Gamma')/(Gamma)((s)/(2)),so that (xi')/(xi)(s)=g(s)+(zeta')/(zeta)(s).}
\expandafter\gdef\csname CGTip@eq:cramer-zero-free\endcsname{Eq. (5.11), p. 26: 2pi i V\_+(z)= V\_rm b(z)+ V\_rm R(z)+ V\_rm L(z)+ V\_rm p(z),}
\expandafter\gdef\csname CGTip@eq:cramer-pieces\endcsname{Eq. (5.12), p. 26: V\_rm b(z) =int\_-a\textasciicircum{}1+ae\textasciicircum{}sigma z (xi')/(xi)(sigma) dsigma, V\_rm R(z)=i e\textasciicircum{}(1+a)zint\_0\textasciicircum{}inf e\textasciicircum{}ithetae\textasciicircum{}itz g(1+a+it) dt,; V\_rm L(z) =i e\textasciicircum{}-azint\_0\textasciicircum{}inf e\textasciicircum{}ithetae\textasciicircum{}itz g(1+a-it) dt, V\_rm p(z)=sum\_n>=2Lambda(n)((e\textasciicircum{}(1+a)(z-log n))/(z-log n)+(n\textasciicircum{}-1-ae\textasciicircum{}-az)/(z+log n)),}
\expandafter\gdef\csname CGTip@eq:cramer-real-window\endcsname{Eq. (5.13), p. 26: 2Re V\_+(u)=e\textasciicircum{}u-(1)/(e\textasciicircum{}2u-1)(log 2<u<log 3).}
\expandafter\gdef\csname CGTip@subsec:window\endcsname{Section 5.2: The window form of the Riemann hypothesis (p. 27)}
\expandafter\gdef\csname CGTip@lem:edge\endcsname{Lemma 5.4 (Vanishing boundary values) (p. 27): Let Jin R be an open interval, h>0, and F holomorphic on Jx(0,h) with |F(u+iy)|<= Cy\textasciicircum{}-m there for some m>=0. If lim\_ydecreasing to 0int F(u+iy) f(u) du=0 for every fin C\_c\textasciicircum{}inf(J), then F==0.}
\expandafter\gdef\csname CGTip@thm:window-RH\endcsname{Theorem 5.5 (The window form of the Riemann hypothesis) (p. 27): Let Jin(0,inf) be an open interval containing no point log p\textasciicircum{}m, for instance J=(log 2,log 3). Without any hypothesis, for every fin C\_c\textasciicircum{}inf(J), sum\_k>=1hat f(gamma\_k)-int\_J f(u) e\textasciicircum{}-u/2V\_+(u) du =sum\_k>=1(hat f(gamma\_k)-(1)/(2)hat f(gamma\_k-id\_k)-(1)/(2)hat f(gamma\_k+id\_k)), both series converging absolutely; only th...}
\expandafter\gdef\csname CGTip@eq:window-residual\endcsname{Eq. (5.14), p. 27: sum\_k>=1hat f(gamma\_k)-int\_J f(u) e\textasciicircum{}-u/2V\_+(u) du =sum\_k>=1(hat f(gamma\_k)-(1)/(2)hat f(gamma\_k-id\_k)-(1)/(2)hat f(gamma\_k+id\_k)),}
\expandafter\gdef\csname CGTip@rem:window\endcsname{Remark 5.6 (What the window form needs and what it does not) (p. 28): item After inserting (5.11), both sides of (ii) are expressed without the abscissas of the zeros: the left side uses the ordinates only, the right side is given in terms of Lambda, Gamma and the values of zeta on a real segment. Identity (ii) is the explicit formula of Guinand and Weil [ref] written for the measure...}
\expandafter\gdef\csname CGTip@subsec:one-window\endcsname{Section 5.3: One window is one cell: the cosine and sine identities (p. 29)}
\expandafter\gdef\csname CGTip@eq:Y-offline\endcsname{Eq. (5.15), p. 29: < Y, f>:=sum\_k>=1int\_0\textasciicircum{}inf f(u)(cosh(d\_ku)-1)cos(gamma\_ku) du(fin C\_c\textasciicircum{}inf(0,inf)),}
\expandafter\gdef\csname CGTip@thm:rung-zero-sum\endcsname{Theorem 5.7 (The rung functional of the endpoint is a sum over the zeros off the line) (p. 29): For every real fin C\_c\textasciicircum{}inf(0,inf), R\_ f(0)-(1)/(2)sum\_n>=2(Lambda(n))/(sqrt(n)) f(log n)=< Y, f> . Equivalently, in D'(1,inf), (log x)((dJ\textasciicircum{}rm M\_Z1)/(dx)-(dJ)/(dx))=(2)/(sqrt(x)) Y(log x). Here the derivatives are distributional derivatives on (1,inf) (dJ/dx is the measure sum\_p\textasciicircum{}mm\textasciicircum{}-1delta\_p\textasciicircum{}m), and Y(log x) denotes...}
\expandafter\gdef\csname CGTip@eq:rung-zero-sum\endcsname{Eq. (5.16), p. 29: R\_ f(0)-(1)/(2)sum\_n>=2(Lambda(n))/(sqrt(n)) f(log n)=< Y, f> .}
\expandafter\gdef\csname CGTip@eq:cell-defect-identity\endcsname{Eq. (5.17), p. 29: (log x)((dJ\textasciicircum{}rm M\_Z1)/(dx)-(dJ)/(dx))=(2)/(sqrt(x)) Y(log x).}
\expandafter\gdef\csname CGTip@eq:cosine-difference\endcsname{Eq. (5.18), p. 29: S\_ f(0)=sum\_k>=1(c\_ f(gamma\_k)- c\_ f(tau\_k))= R\_ f(1)-< Y, f> .}
\expandafter\gdef\csname CGTip@cor:one-window\endcsname{Corollary 5.8 (One window is one cell) (p. 29): Let J=(log 2,log 3). item Re bv X=0 on J <=> < Y, f>=0 for every fin C\_c\textasciicircum{}inf(J) <=> R\_ f(0)=0 for every fin C\_c\textasciicircum{}inf(J) <=> J\textasciicircum{}rm M\_Z1 is constant on (2,3), the rung N=3 of Theorem 15.36. item RH holds if and only if J\textasciicircum{}rm M\_Z1 is constant on (2,3) and, for every real fin C\_c\textasciicircum{}inf(J), the sine identity sum\_k>=1int\_J f(u...}
\expandafter\gdef\csname CGTip@eq:sine-identity\endcsname{Eq. (5.19), p. 29: sum\_k>=1int\_J f(u)sin(gamma\_ku) du=int\_J f(u) e\textasciicircum{}-u/2 Im V\_+(u) du}
\expandafter\gdef\csname CGTip@rem:conjugate-nonlocal\endcsname{Remark 5.9 (The conjugate part is not local) (p. 30): Put Q(u):=2 Y(u)=u e\textasciicircum{}u/2((dJ\textasciicircum{}rm M\_Z1)/(dx)-(dJ)/(dx))(e\textasciicircum{}u). For a holomorphic function on the upper half-plane with suitable decay, the boundary values satisfy bv=-(1)/(2)(Q+i HQ) with H the Hilbert transform; X grows like e\textasciicircum{}max|d\_k| |Re w| and the relation holds only after a regularization, so the following is a he...}
\expandafter\gdef\csname CGTip@prop:finite-offline\endcsname{Proposition 5.10 (Finitely many zeros off the line) (p. 30): Suppose that the set of zeros with d\_k!=0 is finite and non-empty. Then for every open interval Iin(0,inf) there is a real fin C\_c\textasciicircum{}inf(I) with < Y, f>>0. Consequently, if J\textasciicircum{}rm M\_Z1 is constant on one interval (n,n+1), n>=1, then either RH holds or infinitely many zeros lie off the critical line.}
\expandafter\gdef\csname CGTip@thm:frullani\endcsname{Theorem 5.11 (Frullani form of the value at infinity) (p. 30): For K>=1 put Q\_K:=2int\_0\textasciicircum{}inf(e\textasciicircum{}-u/2)/(u)sum\_k<= K(cos(gamma\_ku)-cos(tau\_ku)) du . Then Q\_K=sum\_k<= Klog((1)/(4)+tau\_k\textasciicircum{}2)/((1)/(4)+gamma\_k\textasciicircum{}2), the limit Q\_inf=lim\_K->inf Q\_K exists, and Q\_inf=log L+log varpi\_Z1=log L+sum\_pairslog c\_gamma >= log L>0.0206 , with c\_gamma the pair factor of (13.11). Hence RH holds if and...}
\expandafter\gdef\csname CGTip@eq:frullani-QK\endcsname{Eq. (5.20), p. 30: Q\_K:=2int\_0\textasciicircum{}inf(e\textasciicircum{}-u/2)/(u)sum\_k<= K(cos(gamma\_ku)-cos(tau\_ku)) du .}
\expandafter\gdef\csname CGTip@eq:frullani-limit\endcsname{Eq. (5.21), p. 30: Q\_inf=log L+log varpi\_Z1=log L+sum\_pairslog c\_gamma >= log L>0.0206 ,}
\expandafter\gdef\csname CGTip@rem:frullani-arithmetic\endcsname{Remark 5.12 (The arithmetic of both sides and the size of the defect) (p. 31): item Inserting the explicit formula, sum\_kcos(gamma\_ku)= k(u)-(1)/(2)sum\_nLambda(n)n\textasciicircum{}-1/2delta\_log n- Y(u) (proof of Theorem 5.7), into (5.20) returns the same identity: by Frullani, int\_0\textasciicircum{}inf e\textasciicircum{}-u/2(cosh(du)-1)cos(gamma u)(du)/(u)=(1)/(4)log(((1)/(4)+gamma\textasciicircum{}2)\textasciicircum{}2)/(|rho|\textasciicircum{}2|1-rho|\textasciicircum{}2)=-(1)/(4)log c\_gamma for each zero...}
\expandafter\gdef\csname CGTip@subsec:centroids\endcsname{Section 5.4: Centroids of the two spectra (p. 31)}
\expandafter\gdef\csname CGTip@eq:triangular-centroid\endcsname{Eq. (5.22), p. 31: C\textasciicircum{}pathc(N;u)=(1)/(N)sum\_n<= Nsum\_k<= ne\textasciicircum{}iuo\_k=sum\_k<= N(1-(k-1)/(N))e\textasciicircum{}iuo\_k;}
\expandafter\gdef\csname CGTip@thm:clock-centroid\endcsname{Theorem 5.13 (The clock centroid converges) (p. 31): Let u>0, E\_k=e\textasciicircum{}iutau\_k, S\_n=sum\_k<= nE\_k, b\_k=(e\textasciicircum{}iu(tau\_k+1-tau\_k)-1)\textasciicircum{}-1 (this is b\textasciicircum{}rm sbp\_k of the proof of Theorem 4.1), and let k\_0=k\_0(u) be as in that proof. The series C\textasciicircum{}pathc\_RVM(inf;u):=sum\_k<k\_0E\_k-b\_k\_0E\_k\_0+sum\_k>k\_0(b\_k-1-b\_k)E\_k converges, S\_n-b\_nE\_n+1=C\textasciicircum{}pathc\_RVM(inf;u)+O\_u(log n/n), and C\textasciicircum{}pathc\_RVM(N;...}
\expandafter\gdef\csname CGTip@eq:clock-centroid-limit\endcsname{Eq. (5.23), p. 31: C\textasciicircum{}pathc\_RVM(inf;u):=sum\_k<k\_0E\_k-b\_k\_0E\_k\_0+sum\_k>k\_0(b\_k-1-b\_k)E\_k}
\expandafter\gdef\csname CGTip@eq:clock-centroid-rate\endcsname{Eq. (5.24), p. 31: C\textasciicircum{}pathc\_RVM(N;u)=C\textasciicircum{}pathc\_RVM(inf;u)+O\_u(((log N)\textasciicircum{}2)/(N)).}
\expandafter\gdef\csname CGTip@lem:local-tauberian\endcsname{Lemma 5.14 (A local Tauberian lemma) (p. 32): Let F:[0,inf)->C be locally absolutely continuous, F=0 on [0,x\_0], with |F(x)|<= Clog(2+x) and |F'(x)|<= Clog(2+x)/(1+x) almost everywhere. Suppose that for Re s>0 int\_0\textasciicircum{}inf F(x)e\textasciicircum{}-sx dx=(c)/(s)+c'Log s+Phi(s), with constants c,c', Log the principal branch on Re s>0, and Phi holomorphic in a disc |s|<r. Then F(x)->...}
\expandafter\gdef\csname CGTip@eq:tauberian-hyp\endcsname{Eq. (5.25), p. 32: int\_0\textasciicircum{}inf F(x)e\textasciicircum{}-sx dx=(c)/(s)+c'Log s+Phi(s),}
\expandafter\gdef\csname CGTip@thm:centroid-RH\endcsname{Theorem 5.15 (The ordinate centroid under RH) (p. 32): Assume RH. For every uin(0,inf)log p\textasciicircum{}m, lim\_N->infC\textasciicircum{}pathc\_zeta(N;u)=lim\_N->infsum\_k<= N(1-(k-1)/(N))e\textasciicircum{}iugamma\_k=e\textasciicircum{}-u/2 V\_+(u). For u=1 the limit is e\textasciicircum{}-1/2V\_+(1), whose real part is (1)/(2sqrt(e))(e-(1)/(e\textasciicircum{}2-1))=0.7768942607932... by (5.13), and whose numerical value is 0.77689426079325...+0.76785566886... i (Experim...}
\expandafter\gdef\csname CGTip@eq:centroid-RH\endcsname{Eq. (5.26), p. 32: lim\_N->infC\textasciicircum{}pathc\_zeta(N;u)=lim\_N->infsum\_k<= N(1-(k-1)/(N))e\textasciicircum{}iugamma\_k=e\textasciicircum{}-u/2 V\_+(u).}
\expandafter\gdef\csname CGTip@rem:centroid-status\endcsname{Remark 5.16 (Status of the centroid value) (p. 33): Without RH neither the existence of lim\_NC\textasciicircum{}pathc\_zeta(N;1) nor its value is proved: Step 3 uses RH to identify Z\_gamma with the zero-free function Z\_xi, and Step 1 uses B\_u(T)=O(log T). The equality (5.26) at u=1 is implied by RH; the converse does not follow from the argument of Proposition 5.2 (Remark 5.6(b)). The...}
\expandafter\gdef\csname CGTip@sec:explicit-formula\endcsname{Section 6: The RVM explicit formula, J\_RVM, and pi\_RVM (p. 34)}
\expandafter\gdef\csname CGTip@eq:Ei-def\endcsname{Eq. (6.1), p. 34: Ei(w):=int\_-inf\textasciicircum{}a(e\textasciicircum{}v+ib)/(v+ib) dv, Ei(conj w):=conj Ei(w) .}
\expandafter\gdef\csname CGTip@eq:J-RVM\endcsname{Eq. (6.2), p. 34: J\_RVM(x):=li(x)-2Re sum\_k>=1Ei (((1)/(2)+itau\_k)log x)-log 2+int\_x\textasciicircum{}inf(dy)/(y(y\textasciicircum{}2-1)log y),}
\expandafter\gdef\csname CGTip@lem:Ei\endcsname{Lemma 6.1 (Asymptotics of Ei off the real axis) (p. 34): For w=a+ib, b>0, Ei(w)=(e\textasciicircum{}w)/(w)+varrho(w),|varrho(w)|<=(e\textasciicircum{}a)/(b\textasciicircum{}2).}
\expandafter\gdef\csname CGTip@eq:Ei-remainder\endcsname{Eq. (6.3), p. 34: Ei(w)=(e\textasciicircum{}w)/(w)+varrho(w),|varrho(w)|<=(e\textasciicircum{}a)/(b\textasciicircum{}2).}
\expandafter\gdef\csname CGTip@thm:J-exists\endcsname{Theorem 6.2 (Compact-uniform existence of the RVM zero sum) (p. 34): The series sum\_kRe Ei (((1)/(2)+itau\_k)u) converges uniformly for u in every compact interval [u\_1,u\_2]in(0,inf). Hence J\_RVM is a continuous function on (1,inf).}
\expandafter\gdef\csname CGTip@eq:Ei-asymptotic\endcsname{Eq. (6.4), p. 34: Re Ei (((1)/(2)+itau\_k)u)=(e\textasciicircum{}u/2)/(u) (sin(utau\_k))/(tau\_k)+O\_u\_1,u\_2(tau\_k\textasciicircum{}-2),}
\expandafter\gdef\csname CGTip@eq:abel-weights\endcsname{Eq. (6.5), p. 34: (1)/(tau\_K)-(1)/(tau\_K+1)<=(tau\_K+1-tau\_K)/(tau\_K\textasciicircum{}2)<=(2pi)/(l\_Ktau\_K\textasciicircum{}2)<<(log K)/(K\textasciicircum{}2)}
\expandafter\gdef\csname CGTip@eq:J-star\endcsname{Eq. (6.6), p. 34: J\_RVM\textasciicircum{}*(x):= J\_RVM(x), x>=2,; 0, 0<x<2,}
\expandafter\gdef\csname CGTip@eq:pi-RVM\endcsname{Eq. (6.7), p. 34: pi\_RVM(x):=sum\_1<= m<=log\_2x(mu(m))/(m) J\_RVM\textasciicircum{}*(x\textasciicircum{}1/m).}
\expandafter\gdef\csname CGTip@prop:finite-inversion\endcsname{Proposition 6.3 (Exact finite prime-power inversion) (p. 34): For every x>=2, J\_RVM\textasciicircum{}*(x)=sum\_1<= m<=log\_2x(1)/(m) pi\_RVM(x\textasciicircum{}1/m).}
\expandafter\gdef\csname CGTip@eq:J-pi-inverse\endcsname{Eq. (6.8), p. 34: J\_RVM\textasciicircum{}*(x)=sum\_1<= m<=log\_2x(1)/(m) pi\_RVM(x\textasciicircum{}1/m).}
\expandafter\gdef\csname CGTip@eq:mobius-coefficient\endcsname{Eq. (6.9), p. 35: (1)/(r)sum\_d| rmu(d),}
\expandafter\gdef\csname CGTip@subsec:square-root\endcsname{Section 6.1: The square-root estimate (p. 35)}
\expandafter\gdef\csname CGTip@eq:A-RVM\endcsname{Eq. (6.10), p. 35: A\_RVM(u):=sum\_k>=1(e\textasciicircum{}iutau\_k)/(tau\_k),}
\expandafter\gdef\csname CGTip@lem:A-RVM\endcsname{Lemma 6.4 (The weighted clock sum) (p. 35): For every u>=log 2 the series (6.10) converges, and | A\_RVM(u)|<=(log\textasciicircum{}2(2+2u))/(4pi)+(log(2+2u))/(2pi)+C\_ A, with an absolute constant C\_ A.}
\expandafter\gdef\csname CGTip@eq:A-RVM-bound\endcsname{Eq. (6.11), p. 35: | A\_RVM(u)|<=(log\textasciicircum{}2(2+2u))/(4pi)+(log(2+2u))/(2pi)+C\_ A,}
\expandafter\gdef\csname CGTip@eq:A-head\endcsname{Eq. (6.12), p. 35: sum\_tau\_k<= Y\_0(1)/(tau\_k)=(n\_RVM(Y\_0))/(Y\_0)+int\_tau\_1\textasciicircum{}Y\_0(n\_RVM(t))/(t\textasciicircum{}2) dt<=(l(Y\_0))/(2pi)+int\_tau\_1\textasciicircum{}Y\_0(l(t))/(2pi t) dt<=(l(Y\_0))/(2pi)+(l(Y\_0)\textasciicircum{}2)/(4pi),}
\expandafter\gdef\csname CGTip@eq:phi-second\endcsname{Eq. (6.13), p. 35: phi''(y)=-(2pi u)/(T l(T)\textasciicircum{}3), T= T(y-(1)/(2)).}
\expandafter\gdef\csname CGTip@eq:block-vdc\endcsname{Eq. (6.14), p. 35: |sum\_a<k<= be\textasciicircum{}iutau\_k|<= A\_rm vdC(256((Y\_jl(Y\_j+1))/(2pi)+1)varkappa\_j\textasciicircum{}1/2+varkappa\_j\textasciicircum{}-1/2)=: P\_j,}
\expandafter\gdef\csname CGTip@eq:block-weighted\endcsname{Eq. (6.15), p. 35: (P\_j)/(Y\_j)<= A\_rm vdC((16)/(2pi)((pi u)/(Y\_jl(Y\_j+1)))\textasciicircum{}1/2+(16)/(Y\_j)+((l(Y\_j+1)\textasciicircum{}3)/(pi uY\_j))\textasciicircum{}1/2).}
\expandafter\gdef\csname CGTip@thm:Lambert-estimate\endcsname{Theorem 6.5 (Square-root estimate for J\_RVM and pi\_RVM) (p. 36): For x>=8, J\_RVM(x)-li(x)=O ((sqrt(x)(log log x)\textasciicircum{}2)/(log x)), pi\_RVM(x)-li(x)=O ((sqrt(x)(log log x)\textasciicircum{}2)/(log x)), with absolute implied constants. No hypothesis on zeta is used.}
\expandafter\gdef\csname CGTip@eq:J-baseline\endcsname{Eq. (6.16), p. 36: J\_RVM(x)-li(x)=O ((sqrt(x)(log log x)\textasciicircum{}2)/(log x)),}
\expandafter\gdef\csname CGTip@eq:pi-baseline\endcsname{Eq. (6.17), p. 36: pi\_RVM(x)-li(x)=O ((sqrt(x)(log log x)\textasciicircum{}2)/(log x)),}
\expandafter\gdef\csname CGTip@eq:Ei-split\endcsname{Eq. (6.18), p. 36: Ei(rho\_k\textasciicircum{}RVMu)=(e\textasciicircum{}rho\_k\textasciicircum{}RVMu)/(rho\_k\textasciicircum{}RVMu)+varrho\_k,|varrho\_k|<=(e\textasciicircum{}u/2)/(tau\_k\textasciicircum{}2u\textasciicircum{}2).}
\expandafter\gdef\csname CGTip@eq:Ei-sum-A\endcsname{Eq. (6.19), p. 36: sum\_k>=1Ei(rho\_k\textasciicircum{}RVMu)=(e\textasciicircum{}u/2)/(u)(-i A\_RVM(u)+sum\_kr\_ke\textasciicircum{}iutau\_k)+sum\_kvarrho\_k =(e\textasciicircum{}u/2)/(u)(-i A\_RVM(u)+O(1)).}
\expandafter\gdef\csname CGTip@eq:pi-minus-li\endcsname{Eq. (6.20), p. 36: pi\_RVM(x)-li(x)=(J\_RVM(x)-li(x))+sum\_2<= m<=log\_2x(mu(m))/(m) J\_RVM\textasciicircum{}*(x\textasciicircum{}1/m).}
\expandafter\gdef\csname CGTip@cor:three-limits\endcsname{Corollary 6.6 (The three limits) (p. 36): As x->inf, (li(x))/(pi\_RVM(x))->1, (pi(x))/(li(x))->1, (pi(x))/(pi\_RVM(x))->1 . More precisely, for some absolute c\_0>0, pi(x)-pi\_RVM(x)=O (x e\textasciicircum{}-c\_0sqrt(log x)).}
\expandafter\gdef\csname CGTip@eq:three-limits\endcsname{Eq. (6.21), p. 36: (li(x))/(pi\_RVM(x))->1, (pi(x))/(li(x))->1, (pi(x))/(pi\_RVM(x))->1 .}
\expandafter\gdef\csname CGTip@eq:pi-minus-piRVM\endcsname{Eq. (6.22), p. 36: pi(x)-pi\_RVM(x)=O (x e\textasciicircum{}-c\_0sqrt(log x)).}
\expandafter\gdef\csname CGTip@sec:integer-weights\endcsname{Section 7: Integer weights and the continuous-to-atomic problem (p. 37)}
\expandafter\gdef\csname CGTip@eq:cn\endcsname{Eq. (7.1), p. 37: c\_n\textasciicircum{}RVM:=pi\_RVM(n)-pi\_RVM(n-1), n>=2,}
\expandafter\gdef\csname CGTip@eq:weighted-zeta\endcsname{Eq. (7.2), p. 37: zeta\_Z,RVM(s):=prod\_n>=2(1-n\textasciicircum{}-s)\textasciicircum{}-c\_n\textasciicircum{}RVM.}
\expandafter\gdef\csname CGTip@eq:JZ\endcsname{Eq. (7.3), p. 37: J\_Z,RVM(x):=sum\_1<= m<=log\_2x(1)/(m) pi\_RVM(floor(x\textasciicircum{}1/m)),}
\expandafter\gdef\csname CGTip@eq:J-star-Z\endcsname{Eq. (7.4), p. 37: J\_RVM\textasciicircum{}*,Z(x):=J\_RVM\textasciicircum{}*(floor(x)).}
\expandafter\gdef\csname CGTip@thm:unit-cell\endcsname{Theorem 7.1 (Lipschitz control of J\_RVM) (p. 37): There is an absolute constant C\_J such that |J\_RVM(x+theta)-J\_RVM(x)|<= C\_J theta (x>=2, 0<theta <=1).}
\expandafter\gdef\csname CGTip@eq:unit-cell\endcsname{Eq. (7.5), p. 37: |J\_RVM(x+theta)-J\_RVM(x)|<= C\_J theta (x>=2, 0<theta <=1).}
\expandafter\gdef\csname CGTip@eq:Dk-def\endcsname{Eq. (7.6), p. 37: D\_k:=Ei(rho\_k\textasciicircum{}RVMu\_2)-Ei(rho\_k\textasciicircum{}RVMu\_1)=int\_u\_1\textasciicircum{}u\_2(e\textasciicircum{}v/2e\textasciicircum{}ivtau\_k)/(v) dv .}
\expandafter\gdef\csname CGTip@eq:damped-sum\endcsname{Eq. (7.7), p. 37: sum\_ke\textasciicircum{}-alpha tau\_kD\_k=int\_u\_1\textasciicircum{}u\_2(e\textasciicircum{}v/2)/(v) S\_alpha(v) dv, S\_alpha(v):=sum\_k>=1g(tau\_k), g(tau):=e\textasciicircum{}(iv-alpha)tau.}
\expandafter\gdef\csname CGTip@eq:S-eps-split\endcsname{Eq. (7.8), p. 37: S\_alpha(v)=int\_tau\_0\textasciicircum{}inf g dN\_RVM+g(tau\_0)N\_RVM(tau\_0)+int\_tau\_0\textasciicircum{}inf psi (N\_RVM(tau)+(1)/(2))g'(tau) dtau .}
\expandafter\gdef\csname CGTip@eq:K-harmonics\endcsname{Eq. (7.9), p. 37: int\_tau\_0\textasciicircum{}inf psi (N\_RVM+(1)/(2))g' dtau=-sum\_m>=1((-1)\textasciicircum{}m)/(2pi i m)(K\_m\textasciicircum{}+-K\_m\textasciicircum{}-), K\_m\textasciicircum{}+-:=int\_tau\_0\textasciicircum{}inf(iv-alpha)e\textasciicircum{}-alpha taue\textasciicircum{}iPsi\_m\textasciicircum{}+-(tau) dtau,}
\expandafter\gdef\csname CGTip@eq:K-minus-bound\endcsname{Eq. (7.10), p. 38: |K\_m\textasciicircum{}-|<=2(v+1)((8)/(mlog 2)+8sqrt((2T\_m)/(m))) =2(v+1)((8)/(mlog 2)+16sqrt(pi) (e\textasciicircum{}v/(2m))/(sqrt(m))).}
\expandafter\gdef\csname CGTip@eq:S-eps-bound\endcsname{Eq. (7.11), p. 38: |S\_alpha(v)|<= C (v+1) e\textasciicircum{}v/2(v>=log 2, 0<alpha <=1),}
\expandafter\gdef\csname CGTip@eq:Dk-sum-bound\endcsname{Eq. (7.12), p. 38: |sum\_ke\textasciicircum{}-alpha tau\_kD\_k|<= C'(u\_2-u\_1) x<= C'theta,}
\expandafter\gdef\csname CGTip@cor:cn-bounded\endcsname{Corollary 7.2 (Bounded cell masses) (p. 38): pi\_RVM(y)-pi\_RVM(floor(y))=O(1)(y>=2), c\_n\textasciicircum{}RVM=O(1)(n>=2), with absolute implied constants, uniform in y>=2 and n>=2.}
\expandafter\gdef\csname CGTip@eq:unit-cell-consequence\endcsname{Eq. (7.13), p. 38: pi\_RVM(y)-pi\_RVM(floor(y))=O(1)(y>=2), c\_n\textasciicircum{}RVM=O(1)(n>=2),}
\expandafter\gdef\csname CGTip@cor:ZRVM-Dirichlet\endcsname{Corollary 7.3 (The integer-weighted zeta is a Dirichlet series) (p. 38): The product (7.2) converges absolutely for Re s>1 and zeta\_Z,RVM(s)=sum\_N>=1(a\_N\textasciicircum{}Z)/(N\textasciicircum{}s), a\_1\textasciicircum{}Z=1, sum\_N>=1(|a\_N\textasciicircum{}Z|)/(N\textasciicircum{}sigma)<inf (sigma>1). The inverse is also an absolutely convergent Dirichlet series on the whole half-plane: (1)/(zeta\_Z,RVM(s))=sum\_N>=1(conj a\_N\textasciicircum{}Z)/(N\textasciicircum{}s), conj a\_1\textasciicircum{}Z=1, sum\_N>=1(|conj a\_N\textasciicircum{}Z|)/(N...}
\expandafter\gdef\csname CGTip@eq:ZRVM-Dirichlet\endcsname{Eq. (7.14), p. 38: zeta\_Z,RVM(s)=sum\_N>=1(a\_N\textasciicircum{}Z)/(N\textasciicircum{}s), a\_1\textasciicircum{}Z=1, sum\_N>=1(|a\_N\textasciicircum{}Z|)/(N\textasciicircum{}sigma)<inf (sigma>1).}
\expandafter\gdef\csname CGTip@eq:ZRVM-inverse-Dirichlet\endcsname{Eq. (7.15), p. 38: (1)/(zeta\_Z,RVM(s))=sum\_N>=1(conj a\_N\textasciicircum{}Z)/(N\textasciicircum{}s), conj a\_1\textasciicircum{}Z=1, sum\_N>=1(|conj a\_N\textasciicircum{}Z|)/(N\textasciicircum{}sigma)<inf (sigma>1).}
\expandafter\gdef\csname CGTip@eq:log-ZRVM\endcsname{Eq. (7.16), p. 38: log zeta\_Z,RVM(s)=sum\_n>=2c\_n\textasciicircum{}RVMsum\_m>=1(n\textasciicircum{}-ms)/(m)=sum\_N>=2(b\_N)/(N\textasciicircum{}s), b\_N:=sum\_n\textasciicircum{}m=N(c\_n\textasciicircum{}RVM)/(m),}
\expandafter\gdef\csname CGTip@thm:MS\endcsname{Theorem 7.4 (Mellin--Stieltjes identification) (p. 39): There is a function J\_RVM\textasciicircum{}rm M on (1,inf) with J\_RVM\textasciicircum{}rm M(x)x\textasciicircum{}-cin L\textasciicircum{}2((1,inf),dx/x) for every c>1 such that log zeta\_RVM(s)=sint\_1\textasciicircum{}inf J\_RVM\textasciicircum{}rm M(x) x\textasciicircum{}-s-1 dx(Re s>1), and J\_RVM\textasciicircum{}rm M-J\_RVM is constant on (1,inf). In particular J\_RVM\textasciicircum{}rm M=J\_RVM\textasciicircum{}*+C\_RVM on [2,inf) for a constant C\_RVM.}
\expandafter\gdef\csname CGTip@eq:MS-RVM\endcsname{Eq. (7.17), p. 39: log zeta\_RVM(s)=sint\_1\textasciicircum{}inf J\_RVM\textasciicircum{}rm M(x) x\textasciicircum{}-s-1 dx(Re s>1),}
\expandafter\gdef\csname CGTip@cor:weighted-factorization\endcsname{Corollary 7.5 (Weighted factorization) (p. 39): J\_Z,RVM(x)-J\_RVM\textasciicircum{}*(x)=O(1+log log x), and zeta\_Z,RVM(s)=zeta\_RVM(s) e\textasciicircum{}h\_RVM(s), where h\_RVM(s):=sint\_1\textasciicircum{}inf (J\_Z,RVM(x)-J\_RVM\textasciicircum{}rm M(x))x\textasciicircum{}-s-1 dx is holomorphic for Re s>0. The two functions have the same zeros in Re s>0, and zeta\_Z,RVM inherits through e\textasciicircum{}h\_RVM the pole of zeta\_RVM at s=1; compare the modified Euler pr...}
\expandafter\gdef\csname CGTip@eq:JZ-diff\endcsname{Eq. (7.18), p. 39: J\_Z,RVM(x)-J\_RVM\textasciicircum{}*(x)=O(1+log log x),}
\expandafter\gdef\csname CGTip@eq:weighted-factorization\endcsname{Eq. (7.19), p. 39: zeta\_Z,RVM(s)=zeta\_RVM(s) e\textasciicircum{}h\_RVM(s),}
\expandafter\gdef\csname CGTip@eq:hRVM\endcsname{Eq. (7.20), p. 39: h\_RVM(s):=sint\_1\textasciicircum{}inf (J\_Z,RVM(x)-J\_RVM\textasciicircum{}rm M(x))x\textasciicircum{}-s-1 dx}
\expandafter\gdef\csname CGTip@lem:cell-cancel\endcsname{Lemma 7.6 (Exact first-harmonic cell cancellation) (p. 39): For every integer n, int\_n-1\textasciicircum{}n2cos(2pi x) dx=0 .}
\expandafter\gdef\csname CGTip@eq:cell-cancel\endcsname{Eq. (7.21), p. 39: int\_n-1\textasciicircum{}n2cos(2pi x) dx=0 .}
\expandafter\gdef\csname CGTip@hyp:cell-remainder\endcsname{Hypothesis 7.7 (Uniform cell remainder away from the exceptional powers of 2; proved in Corollary 15.19) (p. 39): For n>=3 with nnot in 2\textasciicircum{}j:mu(j)!=0, c\_n\textasciicircum{}RVM-int\_n-1\textasciicircum{}n((1)/(log x)+2cos(2pi x))dx=o ((1)/(log n)).}
\expandafter\gdef\csname CGTip@eq:cell-remainder\endcsname{Eq. (7.22), p. 39: c\_n\textasciicircum{}RVM-int\_n-1\textasciicircum{}n((1)/(log x)+2cos(2pi x))dx=o ((1)/(log n)).}
\expandafter\gdef\csname CGTip@prop:eventual-positive\endcsname{Proposition 7.8 (Eventual positivity) (p. 40): Hypothesis 7.7 implies c\_n\textasciicircum{}RVM\textasciitilde{}(1)/(log n), hence 0<c\_n\textasciicircum{}RVM<1 for all sufficiently large n in that set.}
\expandafter\gdef\csname CGTip@eq:eventual-positive\endcsname{Eq. (7.23), p. 40: c\_n\textasciicircum{}RVM\textasciitilde{}(1)/(log n),}
\expandafter\gdef\csname CGTip@sec:mertens-transform\endcsname{Section 8: The critical RVM Mertens transform (p. 40)}
\expandafter\gdef\csname CGTip@eq:F-RVM\endcsname{Eq. (8.1), p. 40: F\_RVM(s)=hat M\_RVM\textasciicircum{}crit(s):=(1)/(zeta\_RVM(s))=(G(s))/(P\_RVM\textasciicircum{}crit(s)) .}
\expandafter\gdef\csname CGTip@eq:Mcrit-formal\endcsname{Eq. (8.2), p. 40: M\_RVM\textasciicircum{}crit(x)=(1)/(2pi i)int\_(c)F\_RVM(s)(x\textasciicircum{}s)/(s) ds ;}
\expandafter\gdef\csname CGTip@prop:Mcrit-poles\endcsname{Proposition 8.1 (Synthetic pole abscissa) (p. 40): F\_RVM=G/P\_RVM\textasciicircum{}crit is meromorphic on C and holomorphic in Re s>(1)/(2). The poles of this meromorphic function are the simple poles (1)/(2)+- itau\_k(k>=1)and s=-2,-4,-6,... and s=1 is a simple zero of F\_RVM.}
\expandafter\gdef\csname CGTip@eq:Mcrit-poles\endcsname{Eq. (8.3), p. 40: (1)/(2)+- itau\_k(k>=1)and s=-2,-4,-6,...}
\expandafter\gdef\csname CGTip@prop:Gaussian\endcsname{Proposition 8.2 (Gaussian square-root bound under vertical growth) (p. 40): Fix c>1 and 0<eps<c-(1)/(2). Assume, uniformly for (1)/(2)+eps <=sigma <= c, F\_RVM(sigma+it)=exp(o(t\textasciicircum{}2))(|t|->inf). Define M\_RVM,G\textasciicircum{}crit(x):=(1)/(2pi i)int\_(c)e\textasciicircum{}s\textasciicircum{}2F\_RVM(s)(x\textasciicircum{}s)/(s) ds . Then M\_RVM,G\textasciicircum{}crit(x)=O\_eps(x\textasciicircum{}1/2+eps).}
\expandafter\gdef\csname CGTip@eq:subgaussian\endcsname{Eq. (8.4), p. 40: F\_RVM(sigma+it)=exp(o(t\textasciicircum{}2))(|t|->inf).}
\expandafter\gdef\csname CGTip@eq:MG\endcsname{Eq. (8.5), p. 40: M\_RVM,G\textasciicircum{}crit(x):=(1)/(2pi i)int\_(c)e\textasciicircum{}s\textasciicircum{}2F\_RVM(s)(x\textasciicircum{}s)/(s) ds .}
\expandafter\gdef\csname CGTip@eq:MG-bound\endcsname{Eq. (8.6), p. 40: M\_RVM,G\textasciicircum{}crit(x)=O\_eps(x\textasciicircum{}1/2+eps).}
\expandafter\gdef\csname CGTip@sec:transition\endcsname{Section 9: The transition theorem and the unsmoothed synthetic Mertens bound (p. 41)}
\expandafter\gdef\csname CGTip@eq:chi\endcsname{Eq. (9.1), p. 41: chi(s):=2\textasciicircum{}spi\textasciicircum{}s-1sin ((pi s)/(2))Gamma(1-s)=pi\textasciicircum{}s-1/2(Gamma((1-s)/2))/(Gamma(s/2)),}
\expandafter\gdef\csname CGTip@eq:transition-zw\endcsname{Eq. (9.2), p. 41: s=(1)/(2)+z, z=sigma'+it, tau\_*:=-iz=t-isigma', w(s):=exp -2pi i N\_RVM(tau\_*).}
\expandafter\gdef\csname CGTip@eq:chi-size\endcsname{Eq. (9.3), p. 41: |chi(sigma+it)|=((|t|)/(2pi))\textasciicircum{}1/2-sigma(1+O(|t|\textasciicircum{}-1)).}
\expandafter\gdef\csname CGTip@thm:transition\endcsname{Theorem 9.1 (Transition evaluation of the midpoint product) (p. 41): There is a constant varpi\_RVM>0 with the following properties. item For every B>0 there are t\_0(B) and C\_B such that, for 0<sigma'<= B and t>= t\_0(B), zeta\_RVM(s)=varpi\_RVM(1+w(s))(1+ E\_0(s)),| E\_0(s)|<=(C\_B)/(t) . item For every 0<sigma'\_1<sigma'\_2<(1)/(2) there are t\_0 and C such that, for sigma'\_1<=sigma'<=sigma'...}
\expandafter\gdef\csname CGTip@eq:transition-w\endcsname{Eq. (9.4), p. 41: zeta\_RVM(s)=varpi\_RVM(1+w(s))(1+ E\_0(s)),| E\_0(s)|<=(C\_B)/(t) .}
\expandafter\gdef\csname CGTip@eq:transition-chi\endcsname{Eq. (9.5), p. 41: zeta\_RVM(s)=varpi\_RVM(1+chi(s))(1+ E(s)),| E(s)|<=(C)/(t) .}
\expandafter\gdef\csname CGTip@subsec:transition-proof\endcsname{Section 9.1: Proof of Theorem 9.1 (p. 41)}
\expandafter\gdef\csname CGTip@eq:transition-product\endcsname{Eq. (9.6), p. 41: (P\_RVM\textasciicircum{}crit(s))/(P\_RVM\textasciicircum{}crit((1)/(2)))=prod\_k>=1(1+(z\textasciicircum{}2)/(tau\_k\textasciicircum{}2))=exp Lambda(z), Lambda(z):=sum\_k>=1f(tau\_k), f(tau):=Log (1+(z\textasciicircum{}2)/(tau\textasciicircum{}2)).}
\expandafter\gdef\csname CGTip@eq:fprime\endcsname{Eq. (9.7), p. 41: f'(tau)=-(2z\textasciicircum{}2)/(tau(tau\textasciicircum{}2+z\textasciicircum{}2))=-(2)/(tau)+(1)/(tau-iz)+(1)/(tau+iz).}
\expandafter\gdef\csname CGTip@eq:sawtooth-count\endcsname{Eq. (9.8), p. 41: n\_RVM(tau)=floor(N\_RVM(tau)+(1)/(2))=N\_RVM(tau)-psi (N\_RVM(tau)+(1)/(2)), n\_RVM(tau\_0)=0 .}
\expandafter\gdef\csname CGTip@eq:transition-split\endcsname{Eq. (9.9), p. 41: Lambda(z)=Lambda\_rm sm(z)+Lambda\_rm saw(z), Lambda\_rm sm(z):=int\_tau\_0\textasciicircum{}inf f dN\_RVM+f(tau\_0)N\_RVM(tau\_0),; Lambda\_rm saw(z):=int\_tau\_0\textasciicircum{}infpsi (N\_RVM(tau)+(1)/(2))f'(tau) dtau .}
\expandafter\gdef\csname CGTip@lem:TR-smooth\endcsname{Lemma 9.2 (Smooth part) (p. 42): With log G(s):=Log(z\textasciicircum{}2-(1)/(4))+Log Gamma((s)/(2))-(s)/(2)log pi, Lambda\_rm sm(z)-log G(s)=nu\_3+O(t\textasciicircum{}-1),nu\_3in R independent of z .}
\expandafter\gdef\csname CGTip@eq:smooth-difference\endcsname{Eq. (9.10), p. 42: Lambda\_rm sm(z)-log G(s)=nu\_3+O(t\textasciicircum{}-1),nu\_3in R independent of z .}
\expandafter\gdef\csname CGTip@eq:smooth-J-ibp\endcsname{Eq. (9.11), p. 42: Lambda\_rm sm(z)=-int\_tau\_0\textasciicircum{}inf N\_RVM f' dtau=int\_tau\_0\textasciicircum{}inf N\_RVM(tau)((2)/(tau)-(2tau)/(tau\textasciicircum{}2+z\textasciicircum{}2))dtau .}
\expandafter\gdef\csname CGTip@eq:J2\endcsname{Eq. (9.12), p. 42: (7)/(8)Log (1+(z\textasciicircum{}2)/(tau\_0\textasciicircum{}2))=(7)/(4)Log z-(7)/(4)log tau\_0+O(t\textasciicircum{}-2).}
\expandafter\gdef\csname CGTip@eq:I-def\endcsname{Eq. (9.13), p. 42: I(z):=(1)/(pi)int\_0\textasciicircum{}inf log(tau)/(2pi e) (z\textasciicircum{}2)/(tau\textasciicircum{}2+z\textasciicircum{}2) dtau .}
\expandafter\gdef\csname CGTip@eq:J-asymp\endcsname{Eq. (9.14), p. 42: Lambda\_rm sm(z)=(z)/(2)Log(z)/(2pi e)+(7)/(4)Log z+nu\_1+O(t\textasciicircum{}-2), nu\_1=-(7)/(4)log tau\_0-(1)/(pi)int\_0\textasciicircum{}tau\_0log(tau)/(2pi e) dtau .}
\expandafter\gdef\csname CGTip@eq:stirling-loggamma\endcsname{Eq. (9.15), p. 42: Log Gamma((s)/(2))=(z)/(2)Log(z)/(2)-(1)/(4)Log z+(1)/(4)log 2-(z)/(2)+(1)/(2)log 2pi+O(t\textasciicircum{}-1),}
\expandafter\gdef\csname CGTip@eq:G-asymp\endcsname{Eq. (9.16), p. 42: log G(s)=(z)/(2)Log(z)/(2pi e)+(7)/(4)Log z+nu\_2+O(t\textasciicircum{}-1), nu\_2=(1)/(4)log 2+(1)/(2)log 2pi-(1)/(4)log pi .}
\expandafter\gdef\csname CGTip@lem:TR-clock\endcsname{Lemma 9.3 (The clock as a holomorphic phase) (p. 42): N\_RVM(tau)=(tau)/(2pi)(Log(tau)/(2pi)-1)+(7)/(8) is holomorphic on Re tau>0, N\_RVM(conj tau)=conj N\_RVM(tau), and for Re tau >=tau\_0, 2pi Im N\_RVM(tau)>=(log 2) Im tau (Im tau >=0), 2pi Im N\_RVM(tau)<=(log 2) Im tau (Im tau <=0).}
\expandafter\gdef\csname CGTip@eq:clock-halfplane-bound\endcsname{Eq. (9.17), p. 42: 2pi Im N\_RVM(tau)>=(log 2) Im tau (Im tau >=0), 2pi Im N\_RVM(tau)<=(log 2) Im tau (Im tau <=0).}
\expandafter\gdef\csname CGTip@eq:clock-im-identity\endcsname{Eq. (9.18), p. 42: 2pi Im N\_RVM(tau)=Ylog(|tau|)/(2pi e)+Xarg tau =Y[log(X)/(2pi e)+(1)/(2)log(1+v\textasciicircum{}2)+(arctan v)/(v)]>= Y[log 2+varphi(v)],}
\expandafter\gdef\csname CGTip@lem:TR-harmonic\endcsname{Lemma 9.4 (Abel-regularized harmonic sum) (p. 42): Lambda\_rm saw(z)=Log (1+w(s))+nu\_4+O(t\textasciicircum{}-2),nu\_4in R independent of z .}
\expandafter\gdef\csname CGTip@eq:Psi-asymp\endcsname{Eq. (9.19), p. 42: Lambda\_rm saw(z)=Log (1+w(s))+nu\_4+O(t\textasciicircum{}-2),nu\_4in R independent of z .}
\expandafter\gdef\csname CGTip@eq:psi-r\endcsname{Eq. (9.20), p. 43: psi\_r(y):=-sum\_m>=1r\textasciicircum{}m(sin 2pi my)/(pi m)=(1)/(pi)Im Log (1-re\textasciicircum{}2pi iy).}
\expandafter\gdef\csname CGTip@eq:Psi-r-limit\endcsname{Eq. (9.21), p. 43: Lambda\_rm saw(z)=lim\_rincreasing to 1Lambda\_r(z), Lambda\_r(z):=int\_tau\_0\textasciicircum{}inf psi\_r(N\_RVM+(1)/(2))f' dtau .}
\expandafter\gdef\csname CGTip@eq:Psi-r-series\endcsname{Eq. (9.22), p. 43: Lambda\_r(z)=-sum\_m>=1((-r)\textasciicircum{}m)/(2pi i m)(I\_m\textasciicircum{}+-I\_m\textasciicircum{}-), I\_m\textasciicircum{}+-:=int\_tau\_0\textasciicircum{}inf e\textasciicircum{}+-2pi imN\_RVM(tau)f'(tau) dtau .}
\expandafter\gdef\csname CGTip@eq:ray-integrals\endcsname{Eq. (9.23), p. 43: R\_m\textasciicircum{}+:=iint\_0\textasciicircum{}inf e\textasciicircum{}2pi imN\_RVM(tau\_0+iy)f'(tau\_0+iy) dy, R\_m\textasciicircum{}-:=-iint\_0\textasciicircum{}inf e\textasciicircum{}-2pi imN\_RVM(tau\_0-iy)f'(tau\_0-iy) dy .}
\expandafter\gdef\csname CGTip@eq:lower-contour\endcsname{Eq. (9.24), p. 43: R\_m\textasciicircum{}--I\_m\textasciicircum{}-=2pi i e\textasciicircum{}-2pi imN\_RVM(tau\_*)=2pi i w\textasciicircum{}m .}
\expandafter\gdef\csname CGTip@eq:Psi-r-closed\endcsname{Eq. (9.25), p. 43: Lambda\_r(z)=-sum\_m>=1((-rw)\textasciicircum{}m)/(m)+Phi\_r(z)=Log(1+rw)+Phi\_r(z), Phi\_r(z):=-sum\_m>=1((-r)\textasciicircum{}m)/(2pi i m)(R\_m\textasciicircum{}+-R\_m\textasciicircum{}-).}
\expandafter\gdef\csname CGTip@eq:ray-bound\endcsname{Eq. (9.26), p. 43: |R\textasciicircum{}+-,(z)\_m|<=(4)/(t\textasciicircum{}2)int\_0\textasciicircum{}inf 2\textasciicircum{}-my(tau\_0+y) dy=(4)/(t\textasciicircum{}2)((tau\_0)/(mlog 2)+(1)/((mlog 2)\textasciicircum{}2)).}
\expandafter\gdef\csname CGTip@eq:assembly\endcsname{Eq. (9.27), p. 44: zeta\_RVM(s) =(P\_RVM\textasciicircum{}crit(s))/(G(s))=P\_RVM\textasciicircum{}crit((1)/(2))exp (Lambda\_rm sm(z)-log G(s)+Lambda\_rm saw(z)); =P\_RVM\textasciicircum{}crit((1)/(2)) e\textasciicircum{}nu\_3+nu\_4(1+w)(1+O(t\textasciicircum{}-1)),}
\expandafter\gdef\csname CGTip@eq:varpi-def\endcsname{Eq. (9.28), p. 44: varpi\_RVM:=P\_RVM\textasciicircum{}crit((1)/(2)) e\textasciicircum{}nu\_3+nu\_4>0,}
\expandafter\gdef\csname CGTip@eq:theta-def\endcsname{Eq. (9.29), p. 44: vartheta(v):=(1)/(2i)(Log Gamma ((1)/(4)+(iv)/(2))-Log Gamma ((1)/(4)-(iv)/(2)))-(v)/(2)log pi,}
\expandafter\gdef\csname CGTip@eq:clock-theta\endcsname{Eq. (9.30), p. 44: N\_RVM(v)=(vartheta(v))/(pi)+1+O(|v|\textasciicircum{}-1).}
\expandafter\gdef\csname CGTip@eq:w-chi\endcsname{Eq. (9.31), p. 44: w=e\textasciicircum{}-2ivartheta(tau\_*)-2pi i+O(1/t)=chi(s)(1+O(t\textasciicircum{}-1)).}
\expandafter\gdef\csname CGTip@lem:TR-boundary\endcsname{Lemma 9.5 (Extension to the critical line) (p. 44): The function E\_0=zeta\_RVM/(varpi\_RVM(1+w))-1 extends holomorphically to |sigma'|<(1)/(2), t>t\_0, where t\_0=t\_0((1)/(4)) is the threshold of Theorem 9.1(a) with B=(1)/(4), and | E\_0|<= C/t for 0<=sigma'<=(1)/(4).}
\expandafter\gdef\csname CGTip@cor:derivative-zeros\endcsname{Corollary 9.6 (Derivative at the zeros) (p. 44): For k with tau\_k>t\_0, zeta\_RVM'(rho\_k\textasciicircum{}RVM)=varpi\_RVMlog(tau\_k)/(2pi)(1+O((1)/(tau\_k))).}
\expandafter\gdef\csname CGTip@eq:derivative-zeros\endcsname{Eq. (9.32), p. 44: zeta\_RVM'(rho\_k\textasciicircum{}RVM)=varpi\_RVMlog(tau\_k)/(2pi)(1+O((1)/(tau\_k))).}
\expandafter\gdef\csname CGTip@rem:TR-audit\endcsname{Remark 9.7 (Structure of the proof) (p. 44): The delicate points are handled as follows: the smooth part is computed with principal branches by the identity theorem (Lemma 9.2); the clock estimate is quantified (Lemma 9.3); the sawtooth series is replaced by its Abel means, termwise integration is justified for r<1, and rincreasing to 1 is taken after the cont...}
\expandafter\gdef\csname CGTip@cor:vertical-limit\endcsname{Corollary 9.8 (Vertical constant limit) (p. 45): For every fixed sigma>(1)/(2), zeta\_RVM(sigma+it)->varpi\_RVM, F\_RVM(sigma+it)-> L:=(1)/(varpi\_RVM)(|t|->inf). More precisely, with sigma'=sigma-(1)/(2), |w(sigma+it)|<= (|t|/2pi)\textasciicircum{}-sigma'(1+O\_B(|t|\textasciicircum{}-1))(0<sigma'<= B), and, with alpha\_c:=min(c-(1)/(2),1) for c>(1)/(2), F\_RVM(c+it)-L=O\_c(|t|\textasciicircum{}-alpha\_c).}
\expandafter\gdef\csname CGTip@eq:vertical-limit\endcsname{Eq. (9.33), p. 45: zeta\_RVM(sigma+it)->varpi\_RVM, F\_RVM(sigma+it)-> L:=(1)/(varpi\_RVM)(|t|->inf).}
\expandafter\gdef\csname CGTip@eq:w-decay\endcsname{Eq. (9.34), p. 45: |w(sigma+it)|<= (|t|/2pi)\textasciicircum{}-sigma'(1+O\_B(|t|\textasciicircum{}-1))(0<sigma'<= B),}
\expandafter\gdef\csname CGTip@eq:FminusL-all\endcsname{Eq. (9.35), p. 45: F\_RVM(c+it)-L=O\_c(|t|\textasciicircum{}-alpha\_c).}
\expandafter\gdef\csname CGTip@subsec:halfplane\endcsname{Section 9.2: Bounds on the whole half-plane (p. 45)}
\expandafter\gdef\csname CGTip@eq:saw-primitive\endcsname{Eq. (9.36), p. 45: s\_1(tau):=int\_tau\_0\textasciicircum{}taupsi (N\_RVM(v)+(1)/(2)) dv(tau >=tau\_0).}
\expandafter\gdef\csname CGTip@lem:saw-primitive\endcsname{Lemma 9.9 (Bounded primitive of the midpoint sawtooth) (p. 45): | s\_1(tau)|<=(pi)/(4log 2) for all tau >=tau\_0.}
\expandafter\gdef\csname CGTip@thm:halfplane-RVM\endcsname{Theorem 9.10 (Two-sided bound and limit at +inf) (p. 45): There is an absolute constant R\_0 such that, for s=(1)/(2)+z, z=sigma'+it, sigma'>(1)/(2) and |z|>= R\_0, zeta\_RVM(s)=varpi\_RVM e\textasciicircum{}theta\_RVM(s), |theta\_RVM(s)|<=(pi\textasciicircum{}2)/(2log 2)*(1)/(sigma')+(C)/(|z|). Consequently, for every sigma\_0>1 there is C(sigma\_0) with C(sigma\_0)\textasciicircum{}-1<=|zeta\_RVM(s)|<= C(sigma\_0)(Re s>=sigma\_0), a...}
\expandafter\gdef\csname CGTip@eq:halfplane-RVM\endcsname{Eq. (9.37), p. 45: zeta\_RVM(s)=varpi\_RVM e\textasciicircum{}theta\_RVM(s), |theta\_RVM(s)|<=(pi\textasciicircum{}2)/(2log 2)*(1)/(sigma')+(C)/(|z|).}
\expandafter\gdef\csname CGTip@eq:halfplane-two-sided\endcsname{Eq. (9.38), p. 45: C(sigma\_0)\textasciicircum{}-1<=|zeta\_RVM(s)|<= C(sigma\_0)(Re s>=sigma\_0),}
\expandafter\gdef\csname CGTip@eq:nu4-closed\endcsname{Eq. (9.39), p. 45: nu\_4=-2int\_tau\_0\textasciicircum{}inf(s\_1(tau))/(tau\textasciicircum{}2) dtau .}
\expandafter\gdef\csname CGTip@eq:smooth-exact\endcsname{Eq. (9.40), p. 45: Lambda\_rm sm(z)=(z)/(2)Log(z)/(2pi e)+(7)/(8)Log (1+(z\textasciicircum{}2)/(tau\_0\textasciicircum{}2))-(1)/(pi)int\_0\textasciicircum{}tau\_0log(tau)/(2pi e) (z\textasciicircum{}2)/(tau\textasciicircum{}2+z\textasciicircum{}2) dtau .}
\expandafter\gdef\csname CGTip@eq:saw-ibp\endcsname{Eq. (9.41), p. 45: Lambda\_rm saw(z)=-int\_tau\_0\textasciicircum{}inf s\_1f'' dtau, f''(tau)=(2)/(tau\textasciicircum{}2)-(1)/((tau-iz)\textasciicircum{}2)-(1)/((tau+iz)\textasciicircum{}2).}
\expandafter\gdef\csname CGTip@eq:saw-bound\endcsname{Eq. (9.42), p. 46: |Lambda\_rm saw(z)+2int\_tau\_0\textasciicircum{}inf(s\_1(tau))/(tau\textasciicircum{}2) dtau |<=(pi)/(4log 2)*(2pi)/(sigma') .}
\expandafter\gdef\csname CGTip@thm:varpi-less-one\endcsname{Theorem 9.11 (The transition constant is less than one) (p. 46): The smooth constant of Lemma 9.2 is nu\_3=8-(33)/(4)log 2-2log pi=-0.0079240113..., the sawtooth constant of Lemma 9.4 satisfies nu\_4<-0.012691, and varpi\_RVM=P\textasciicircum{}crit\_RVM((1)/(2)) e\textasciicircum{}nu\_3+nu\_4<P\textasciicircum{}crit\_RVM((1)/(2)) e\textasciicircum{}-0.020615<0.97960 . In particular varpi\_RVM<1 and L=1/varpi\_RVM>1.}
\expandafter\gdef\csname CGTip@eq:nu3-closed\endcsname{Eq. (9.43), p. 46: nu\_3=8-(33)/(4)log 2-2log pi=-0.0079240113...,}
\expandafter\gdef\csname CGTip@eq:varpi-less-one\endcsname{Eq. (9.44), p. 46: varpi\_RVM=P\textasciicircum{}crit\_RVM((1)/(2)) e\textasciicircum{}nu\_3+nu\_4<P\textasciicircum{}crit\_RVM((1)/(2)) e\textasciicircum{}-0.020615<0.97960 .}
\expandafter\gdef\csname CGTip@eq:nu4-limit\endcsname{Eq. (9.45), p. 46: nu\_4=-2lim\_X->infint\_tau\_0\textasciicircum{}Xpsi (N\_RVM(tau)+(1)/(2))(dtau)/(tau) ,}
\expandafter\gdef\csname CGTip@eq:nu4-I0\endcsname{Eq. (9.46), p. 46: I\_0(tau\_1):=int\_tau\_0\textasciicircum{}tau\_1N\_RVM(tau)(dtau)/(tau)=(tau\_1l\_1-2tau\_1-tau\_0log 2+2tau\_0)/(2pi)+(7)/(8)log(tau\_1)/(tau\_0), l\_1:=log(tau\_1)/(2pi).}
\expandafter\gdef\csname CGTip@eq:nu4-weight\endcsname{Eq. (9.47), p. 46: (dtau)/(tau)= W(u) du, W(u):=(2pi)/(T log(T/2pi)), T= T(u-(1)/(2)),}
\expandafter\gdef\csname CGTip@eq:nu4-bound\endcsname{Eq. (9.48), p. 47: nu\_4=-2(I\_0(tau\_1)+sum\_n>=1 I\_n)<=-2 M(tau\_1), M(t):= I\_0(t)-(pi)/(6tl)-(sqrt(3))/(216)*(4pi\textasciicircum{}2(l+1))/(t\textasciicircum{}2l\textasciicircum{}3), l=log(t)/(2pi).}
\expandafter\gdef\csname CGTip@rem:varpi-values\endcsname{Remark 9.12 (High-precision values; what the proof uses) (p. 47): The proof uses only P\textasciicircum{}crit\_RVM((1)/(2))<1 and the two upper bounds nu\_3<0, nu\_4<=-2 M(tau\_1). The floating-point values computed by the same program (Part B of its output, not part of the proof) are nu\_4=-0.0251569083 (period-by-period quadrature of I\_n for n<=2000 and the Euler--Maclaurin tail), e\textasciicircum{}nu\_3+nu\_4=0.96746...}
\expandafter\gdef\csname CGTip@thm:no-Dirichlet-RVM\endcsname{Theorem 9.13 (No absolutely convergent Dirichlet expansion) (p. 47): Neither zeta\_RVM nor F\_RVM=1/zeta\_RVM is represented by a Dirichlet series sum\_na\_nn\textasciicircum{}-s converging absolutely in some half-plane Re s>sigma\_0.}
\expandafter\gdef\csname CGTip@eq:mean-square\endcsname{Eq. (9.49), p. 47: (1)/(2T)int\_-T\textasciicircum{}T|F\_RVM(sigma+it)-L|\textasciicircum{}2dt =sum\_m,ntilde a\_mconj tilde a\_n(mn)\textasciicircum{}-sigma (1)/(2T)int\_-T\textasciicircum{}T((n)/(m))\textasciicircum{}itdt ,}
\expandafter\gdef\csname CGTip@eq:mean-square-limit\endcsname{Eq. (9.50), p. 48: lim\_T->inf(1)/(2T)int\_-T\textasciicircum{}T|F\_RVM(sigma+it)-L|\textasciicircum{}2dt=sum\_n>=1|tilde a\_n|\textasciicircum{}2n\textasciicircum{}-2sigma.}
\expandafter\gdef\csname CGTip@thm:unsmoothed-RVM\endcsname{Theorem 9.14 (Unsmoothed synthetic Mertens bound) (p. 48): For every x>1 and every c>(1)/(2) the symmetric Perron integral M\_RVM\textasciicircum{}crit(x):=lim\_U->inf(1)/(2pi i)int\_c-iU\textasciicircum{}c+iUF\_RVM(s)(x\textasciicircum{}s)/(s) ds exists, is independent of c>(1)/(2), and equals the absolutely convergent expression M\_RVM\textasciicircum{}crit(x)=L+(1)/(2pi i)int\_(c)(F\_RVM(s)-L)(x\textasciicircum{}s)/(s) ds . For every eps in(0,(1)/(2)), M\_RVM\textasciicircum{}cr...}
\expandafter\gdef\csname CGTip@eq:Mcrit-symmetric\endcsname{Eq. (9.51), p. 48: M\_RVM\textasciicircum{}crit(x):=lim\_U->inf(1)/(2pi i)int\_c-iU\textasciicircum{}c+iUF\_RVM(s)(x\textasciicircum{}s)/(s) ds}
\expandafter\gdef\csname CGTip@eq:Mcrit-renorm\endcsname{Eq. (9.52), p. 48: M\_RVM\textasciicircum{}crit(x)=L+(1)/(2pi i)int\_(c)(F\_RVM(s)-L)(x\textasciicircum{}s)/(s) ds .}
\expandafter\gdef\csname CGTip@eq:Mcrit-unsmoothed-bound\endcsname{Eq. (9.53), p. 48: M\_RVM\textasciicircum{}crit(x)=L+O\_eps(x\textasciicircum{}1/2+eps).}
\expandafter\gdef\csname CGTip@eq:bromwich\endcsname{Eq. (9.54), p. 48: lim\_U->inf(1)/(2pi i)int\_c-iU\textasciicircum{}c+iU(x\textasciicircum{}s)/(s) ds=1 .}
\expandafter\gdef\csname CGTip@eq:Mcrit-final-estimate\endcsname{Eq. (9.55), p. 48: |M\_RVM\textasciicircum{}crit(x)-L|<=(x\textasciicircum{}1/2+eps)/(2pi)int\_-inf\textasciicircum{}inf(|F\_RVM((1)/(2)+eps+it)-L|)/(|(1)/(2)+eps+it|) dt,}
\expandafter\gdef\csname CGTip@rem:unsmoothed-scope\endcsname{Remark 9.15 (Scope) (p. 48): Theorem 9.14 concerns 1/zeta\_RVM. The transfer to M goes through the endpoint of the lambda-family and is Theorem 12.8.}
\expandafter\gdef\csname CGTip@sec:comparison\endcsname{Section 10: The exact Hadamard product and the comparison with zeta (p. 48)}
\expandafter\gdef\csname CGTip@eq:Pzeta\endcsname{Eq. (10.1), p. 48: P\_zeta(s):=G(s)zeta(s)=2xi(s), xi(s)=(1)/(2)s(s-1)pi\textasciicircum{}-s/2Gamma(s/2)zeta(s).}
\expandafter\gdef\csname CGTip@thm:Hadamard\endcsname{Theorem 10.1 (Exact Hadamard product) (p. 48): Let rho run over the nontrivial zeros of zeta with Im rho>0, with multiplicity, ordered by Im rho. The product P\_zeta(s)=prod\_Im rho>0(1+(q(s))/(rho(1-rho))) converges absolutely and locally uniformly in sin C, and the identity holds with no exponential or constant factor.}
\expandafter\gdef\csname CGTip@eq:Hadamard-exact\endcsname{Eq. (10.2), p. 48: P\_zeta(s)=prod\_Im rho>0(1+(q(s))/(rho(1-rho)))}
\expandafter\gdef\csname CGTip@eq:hadamard-general\endcsname{Eq. (10.3), p. 49: 2xi(s)=e\textasciicircum{}Bsprod\_rho(1-(s)/(rho))e\textasciicircum{}s/rho,}
\expandafter\gdef\csname CGTip@eq:pair-factor\endcsname{Eq. (10.4), p. 49: (1-(s)/(rho))(1-(s)/(1-rho))e\textasciicircum{}s/rho+s/(1-rho) =(1+(q(s))/(rho(1-rho)))e\textasciicircum{}s/(rho(1-rho)).}
\expandafter\gdef\csname CGTip@eq:hadamard-exp\endcsname{Eq. (10.5), p. 49: 2xi(s)=e\textasciicircum{}B's Pi\_zeta(s), Pi\_zeta(s):=prod\_Im rho>0(1+(q(s))/(rho(1-rho))), B'=B+sum\_Im rho>0(1)/(rho(1-rho)) .}
\expandafter\gdef\csname CGTip@eq:Q-RVM\endcsname{Eq. (10.6), p. 49: Q\_RVM(s):=(zeta(s))/(zeta\_RVM(s))=(P\_zeta(s))/(P\_RVM\textasciicircum{}crit(s)).}
\expandafter\gdef\csname CGTip@prop:M-transform-relation\endcsname{Proposition 10.2 (Exact transform relation) (p. 49): For Re s>1, and by meromorphic continuation, F\_RVM(s)= Q\_RVM(s) (1)/(zeta(s)) . If the inverse Mellin transform Q\_RVM\textasciicircum{}v of Q\_RVM and the inverse defining M\_RVM\textasciicircum{}crit exist in a common distribution space with Mellin convolution, then M\_RVM\textasciicircum{}crit= Q\_RVM\textasciicircum{}v*Mel M .}
\expandafter\gdef\csname CGTip@eq:M-transform-relation\endcsname{Eq. (10.7), p. 49: F\_RVM(s)= Q\_RVM(s) (1)/(zeta(s)) .}
\expandafter\gdef\csname CGTip@eq:M-convolution\endcsname{Eq. (10.8), p. 49: M\_RVM\textasciicircum{}crit= Q\_RVM\textasciicircum{}v*Mel M .}
\expandafter\gdef\csname CGTip@cor:Q-control\endcsname{Corollary 10.3 (All growth of zeta to the right of (1)/(2) lies in Q\_RVM) (p. 49): For every c'>(1)/(2) and 0<eps<c'-(1)/(2) there are t\_0=t\_0(eps,c') and 0<c\_-(eps,c')<= c\_+(eps,c') such that c\_-<=|zeta\_RVM(sigma+it)|<= c\_+((1)/(2)+eps <=sigma <= c', |t|>= t\_0), and therefore, in the same region, c\_-| Q\_RVM(s)|<=|zeta(s)|<= c\_+| Q\_RVM(s)|, (1)/(c\_+| Q\_RVM(s)|)<=|(1)/(zeta(s))|<=(1)/(c\_-| Q\_RVM(s)...}
\expandafter\gdef\csname CGTip@eq:zetaRVM-two-sided\endcsname{Eq. (10.9), p. 49: c\_-<=|zeta\_RVM(sigma+it)|<= c\_+((1)/(2)+eps <=sigma <= c', |t|>= t\_0),}
\expandafter\gdef\csname CGTip@eq:zeta-asymp-Q\endcsname{Eq. (10.10), p. 49: c\_-| Q\_RVM(s)|<=|zeta(s)|<= c\_+| Q\_RVM(s)|, (1)/(c\_+| Q\_RVM(s)|)<=|(1)/(zeta(s))|<=(1)/(c\_-| Q\_RVM(s)|).}
\expandafter\gdef\csname CGTip@eq:prime-defect\endcsname{Eq. (10.11), p. 50: d\_n\textasciicircum{}RVM:= 1\_n prime-c\_n\textasciicircum{}RVM.}
\expandafter\gdef\csname CGTip@eq:prime-defect-cumulative\endcsname{Eq. (10.12), p. 50: sum\_2<= n<= N d\_n\textasciicircum{}RVM=pi(N)-pi\_RVM(N).}
\expandafter\gdef\csname CGTip@sec:lambda-family\endcsname{Section 11: The lambda-family from zeta\_RVM to zeta\_Z1 (p. 50)}
\expandafter\gdef\csname CGTip@eq:gamma-lambda\endcsname{Eq. (11.1), p. 50: gamma\_k(lambda):=(1-lambda)gamma\_k+lambda tau\_k=gamma\_k+lambda eta\_k,}
\expandafter\gdef\csname CGTip@eq:P-lambda\endcsname{Eq. (11.2), p. 50: P\_lambda(s):=prod\_k>=1(1+(q(s))/((1)/(4)+gamma\_k(lambda)\textasciicircum{}2)),}
\expandafter\gdef\csname CGTip@eq:zeta-lambda\endcsname{Eq. (11.3), p. 50: zeta\_lambda(s):=(P\_lambda(s))/(G(s)), F\_lambda(s):=(1)/(zeta\_lambda(s)),}
\expandafter\gdef\csname CGTip@eq:family-endpoints\endcsname{Eq. (11.4), p. 50: P\_Z1:=P\_0, zeta\_Z1:=zeta\_0, zeta\_1=zeta\_RVM.}
\expandafter\gdef\csname CGTip@thm:path\endcsname{Theorem 11.1 (The lambda-family) (p. 50): The product (11.2) converges absolutely and locally uniformly jointly on [0,1]x C. For every lambda, P\_lambda is entire of order exactly one, P\_lambda(1-s)=P\_lambda(s), P\_lambda(conj s)=conj P\_lambda(s), P\_lambda(0)=P\_lambda(1)=1, and its zeros, with multiplicity, are exactly (1)/(2)+- igamma\_k(lambda), k>=1 . For 0...}
\expandafter\gdef\csname CGTip@eq:family-divisor\endcsname{Eq. (11.5), p. 50: (1)/(2)+- igamma\_k(lambda), k>=1 .}
\expandafter\gdef\csname CGTip@eq:lambda-normal\endcsname{Eq. (11.6), p. 50: sum\_k>=1sup\_sin K, 0<=lambda <=1|(q(s))/((1)/(4)+gamma\_k(lambda)\textasciicircum{}2)|<inf,}
\expandafter\gdef\csname CGTip@eq:order-one\endcsname{Eq. (11.7), p. 50: log max\_|s|<= R\_rm fam|P\_lambda(s)|<=sum\_k>=1log (1+(R\_rm fam\textasciicircum{}2+R\_rm fam)/((1)/(4)+gamma\_k(lambda)\textasciicircum{}2))=O(R\_rm famlog R\_rm fam),}
\expandafter\gdef\csname CGTip@eq:ordinate-gaps\endcsname{Eq. (11.8), p. 50: gamma\_k+1(lambda)-gamma\_k(lambda)=(1-lambda)(gamma\_k+1-gamma\_k)+lambda(tau\_k+1-tau\_k)>=lambda(tau\_k+1-tau\_k)>0,}
\expandafter\gdef\csname CGTip@rem:only-gamma-tau\endcsname{Remark 11.2 (All differences lie in gamma\_k-tau\_k) (p. 51): For lambda,lambda'in[0,1] and s off the divisors, (zeta\_lambda(s))/(zeta\_lambda'(s))=prod\_k>=1(((1)/(4)+(gamma\_k+lambda'eta\_k)\textasciicircum{}2)((1)/(4)+(gamma\_k+lambda eta\_k)\textasciicircum{}2+q(s)))/(((1)/(4)+(gamma\_k+lambda eta\_k)\textasciicircum{}2)((1)/(4)+(gamma\_k+lambda'eta\_k)\textasciicircum{}2+q(s))) . The gamma factor cancels; the members differ only through eta\_k=tau\_k...}
\expandafter\gdef\csname CGTip@eq:ratio-members\endcsname{Eq. (11.9), p. 51: (zeta\_lambda(s))/(zeta\_lambda'(s))=prod\_k>=1(((1)/(4)+(gamma\_k+lambda'eta\_k)\textasciicircum{}2)((1)/(4)+(gamma\_k+lambda eta\_k)\textasciicircum{}2+q(s)))/(((1)/(4)+(gamma\_k+lambda eta\_k)\textasciicircum{}2)((1)/(4)+(gamma\_k+lambda'eta\_k)\textasciicircum{}2+q(s))) .}
\expandafter\gdef\csname CGTip@thm:generator\endcsname{Theorem 11.3 (Infinitesimal generator) (p. 51): Put A\_k,lambda:=(1)/(4)+gamma\_k(lambda)\textasciicircum{}2. Off the divisor of P\_lambda, V\_lambda(s):=d\_lambda log P\_lambda(s)=-2q(s)sum\_k>=1(gamma\_k(lambda) eta\_k)/(A\_k,lambda(A\_k,lambda+q(s))), the series converging locally uniformly on compact sets separated from the divisor. Since G does not depend on lambda, d\_lambda log zeta\_l...}
\expandafter\gdef\csname CGTip@eq:generator\endcsname{Eq. (11.10), p. 51: V\_lambda(s):=d\_lambda log P\_lambda(s)=-2q(s)sum\_k>=1(gamma\_k(lambda) eta\_k)/(A\_k,lambda(A\_k,lambda+q(s))),}
\expandafter\gdef\csname CGTip@eq:generator-zeta\endcsname{Eq. (11.11), p. 51: d\_lambda log zeta\_lambda= V\_lambda, d\_lambda log F\_lambda=- V\_lambda .}
\expandafter\gdef\csname CGTip@eq:mode-derivative\endcsname{Eq. (11.12), p. 51: d\_lambda log (1+(q)/(A\_k,lambda))=-(2q gamma\_k(lambda)eta\_k)/(A\_k,lambda(A\_k,lambda+q)) .}
\expandafter\gdef\csname CGTip@prop:integrated-multiplier\endcsname{Proposition 11.4 (Integrated multiplier) (p. 51): On every component of C bigcup\_lambda'in[0,1](1)/(2)+- igamma\_k(lambda'):k>=1, zeta\_lambda(s)=zeta\_Z1(s)exp (int\_0\textasciicircum{}lambda V\_lambda'(s) dlambda'), F\_lambda(s)=F\_Z1(s)exp (-int\_0\textasciicircum{}lambda V\_lambda'(s) dlambda'), and the total multiplier is Pi\_rm tr(s):=exp (-int\_0\textasciicircum{}1 V\_lambda'(s) dlambda')=(P\_Z1(s))/(P\_RVM\textasciicircum{}crit(s)), zeta...}
\expandafter\gdef\csname CGTip@eq:integrated-zeta\endcsname{Eq. (11.13), p. 51: zeta\_lambda(s)=zeta\_Z1(s)exp (int\_0\textasciicircum{}lambda V\_lambda'(s) dlambda'),}
\expandafter\gdef\csname CGTip@eq:integrated-inverse\endcsname{Eq. (11.14), p. 51: F\_lambda(s)=F\_Z1(s)exp (-int\_0\textasciicircum{}lambda V\_lambda'(s) dlambda'),}
\expandafter\gdef\csname CGTip@eq:transport-multiplier\endcsname{Eq. (11.15), p. 51: Pi\_rm tr(s):=exp (-int\_0\textasciicircum{}1 V\_lambda'(s) dlambda')=(P\_Z1(s))/(P\_RVM\textasciicircum{}crit(s)), zeta\_Z1=zeta\_RVM Pi\_rm tr .}
\expandafter\gdef\csname CGTip@eq:mode-integral\endcsname{Eq. (11.16), p. 51: int\_0\textasciicircum{}1d\_lambda'log (1+(q)/(A\_k,lambda'))dlambda'==log(1+q/((1)/(4)+tau\_k\textasciicircum{}2))/(1+q/((1)/(4)+gamma\_k\textasciicircum{}2))mod 2pi i.}
\expandafter\gdef\csname CGTip@cor:endpoint-expansion\endcsname{Corollary 11.5 (Expansion at the endpoint) (p. 51): Let Kin C be compact, disjoint from the zeros of P\_Z1 and from s=1 (equivalently one may state the expansion for the product ratio P\_lambda/P\_Z1 without this exclusion). Uniformly for sin K, as lambda decreasing to 0, zeta\_lambda(s)=zeta\_Z1(s)(1+lambda V\_0(s)+O\_K(lambda\textasciicircum{}2)), F\_lambda(s)=F\_Z1(s)(1-lambda V\_0(s)+O\_K(l...}
\expandafter\gdef\csname CGTip@eq:endpoint-zeta-expansion\endcsname{Eq. (11.17), p. 51: zeta\_lambda(s)=zeta\_Z1(s)(1+lambda V\_0(s)+O\_K(lambda\textasciicircum{}2)),}
\expandafter\gdef\csname CGTip@eq:endpoint-inverse-expansion\endcsname{Eq. (11.18), p. 51: F\_lambda(s)=F\_Z1(s)(1-lambda V\_0(s)+O\_K(lambda\textasciicircum{}2)).}
\expandafter\gdef\csname CGTip@prop:residues\endcsname{Proposition 11.6 (Residues along the family) (p. 52): Suppose gamma\_k is a simple ordinate (it occurs once in the list). Put rho\_k,lambda=(1)/(2)+igamma\_k(lambda) and, for 0<lambda <=1, r\_k,lambda:=Res\_s=rho\_k,lambda(1)/(s zeta\_lambda(s))=(G(rho\_k,lambda))/(rho\_k,lambdaP\_lambda'(rho\_k,lambda)). Then r\_k,lambda->(G((1)/(2)+igamma\_k))/(((1)/(2)+igamma\_k)P\_Z1'((1)/(2)+iga...}
\expandafter\gdef\csname CGTip@eq:residue-transport\endcsname{Eq. (11.19), p. 52: r\_k,lambda:=Res\_s=rho\_k,lambda(1)/(s zeta\_lambda(s))=(G(rho\_k,lambda))/(rho\_k,lambdaP\_lambda'(rho\_k,lambda)).}
\expandafter\gdef\csname CGTip@eq:residue-limit\endcsname{Eq. (11.20), p. 52: r\_k,lambda->(G((1)/(2)+igamma\_k))/(((1)/(2)+igamma\_k)P\_Z1'((1)/(2)+igamma\_k))(lambda decreasing to 0).}
\expandafter\gdef\csname CGTip@rem:multiplicity\endcsname{Remark 11.7 (Multiplicity) (p. 52): For 0<lambda <=1 every zero of zeta\_lambda is simple. At lambda=0 a repeated ordinate of zeta appears in zeta\_Z1 with the same multiplicity, and the residues of Proposition 11.6 are replaced by the full principal part. Multiplicity plays no role in Step 3: Theorem 12.8 uses only that zeta does not vanish in Re s>(1)...}
\expandafter\gdef\csname CGTip@sec:Z1-step3\endcsname{Section 12: The endpoint zeta\_Z1, Hamburger's theorem, and Step 3 (p. 52)}
\expandafter\gdef\csname CGTip@subsec:Z1-class\endcsname{Section 12.1: The analytic class of the endpoint (p. 52)}
\expandafter\gdef\csname CGTip@prop:Z1-class\endcsname{Proposition 12.1 (Analytic class of zeta\_Z1) (p. 52): The endpoint P\_Z1(s)=prod\_k>=1(1+(q(s))/((1)/(4)+gamma\_k\textasciicircum{}2)), zeta\_Z1(s)=(P\_Z1(s))/(G(s)), has the following properties. item P\_Z1 is entire of order exactly one, P\_Z1(1-s)=P\_Z1(s), P\_Z1(0)=P\_Z1(1)=1; item the nontrivial zeros of zeta\_Z1 are exactly (1)/(2)+- igamma\_k, k>=1, with the multiplicities of the ordinate l...}
\expandafter\gdef\csname CGTip@eq:Z1-product\endcsname{Eq. (12.1), p. 52: P\_Z1(s)=prod\_k>=1(1+(q(s))/((1)/(4)+gamma\_k\textasciicircum{}2)), zeta\_Z1(s)=(P\_Z1(s))/(G(s)),}
\expandafter\gdef\csname CGTip@eq:Z1-divisor\endcsname{Eq. (12.2), p. 52: (1)/(2)+- igamma\_k, k>=1,}
\expandafter\gdef\csname CGTip@eq:Z1-FE\endcsname{Eq. (12.3), p. 52: pi\textasciicircum{}-s/2Gamma(s/2)zeta\_Z1(s)=pi\textasciicircum{}-(1-s)/2Gamma((1-s)/2)zeta\_Z1(1-s);}
\expandafter\gdef\csname CGTip@eq:G-near-one\endcsname{Eq. (12.4), p. 52: G(s)=(s-1)+O((s-1)\textasciicircum{}2)(s->1).}
\expandafter\gdef\csname CGTip@eq:Z1-FE-proof\endcsname{Eq. (12.5), p. 52: pi\textasciicircum{}-s/2Gamma(s/2)zeta\_Z1(s)=(P\_Z1(s))/(q(s)),}
\expandafter\gdef\csname CGTip@eq:Z1-entire-removal\endcsname{Eq. (12.6), p. 52: (s-1)zeta\_Z1(s)=P\_Z1(s) (pi\textasciicircum{}s/2)/(s Gamma(s/2)),}
\expandafter\gdef\csname CGTip@eq:Z1-defect-ratio\endcsname{Eq. (12.7), p. 53: R\_Z1(s):=(zeta\_Z1(s))/(zeta(s))=(P\_Z1(s))/(P\_zeta(s)),}
\expandafter\gdef\csname CGTip@subsec:Hamburger\endcsname{Section 12.2: Hamburger's theorem at the endpoint (p. 53)}
\expandafter\gdef\csname CGTip@thm:Hamburger-Z1\endcsname{Theorem 12.2 (Hamburger closure) (p. 53): Consider the condition (D)zeta\_Z1(s)=sum\_n>=1(a\_n)/(n\textasciicircum{}s)absolutely for Re s>1 . Then (D)<=> R\_Z1==1<=>zeta\_Z1=zeta <=>RH.}
\expandafter\gdef\csname CGTip@eq:Z1-D\endcsname{Eq. (12.8), p. 53: (D)zeta\_Z1(s)=sum\_n>=1(a\_n)/(n\textasciicircum{}s)absolutely for Re s>1 .}
\expandafter\gdef\csname CGTip@eq:Hamburger-equivalence\endcsname{Eq. (12.9), p. 53: (D)<=> R\_Z1==1<=>zeta\_Z1=zeta <=>RH.}
\expandafter\gdef\csname CGTip@eq:Hamburger-conclusion\endcsname{Eq. (12.10), p. 53: zeta\_Z1=C\_H zeta}
\expandafter\gdef\csname CGTip@eq:RH-product\endcsname{Eq. (12.11), p. 53: P\_zeta(s)=prod\_k>=1(1+(q(s))/((1)/(4)+gamma\_k\textasciicircum{}2))=P\_Z1(s)}
\expandafter\gdef\csname CGTip@subsec:D-multiplier\endcsname{Section 12.3: Condition (D) and the integer-weighted zeta (p. 53)}
\expandafter\gdef\csname CGTip@eq:K-multiplier\endcsname{Eq. (12.12), p. 53: K(s):=(zeta\_Z1(s))/(zeta\_Z,RVM(s)) .}
\expandafter\gdef\csname CGTip@thm:D-multiplier\endcsname{Theorem 12.3 (Reduction of (D) to the multiplier) (p. 53): item If K has an expansion K(s)=sum\_n>=1k\_nn\textasciicircum{}-s that converges absolutely in Re s>1, then (12.8) holds. item If (12.8) holds, then K(s)=sum\_n>=1k\_nn\textasciicircum{}-s absolutely on the whole half-plane Re s>1. Thus (D) holds if and only if K is an absolutely convergent Dirichlet series in Re s>1. Hence, by Theorem 12.2, an absolut...}
\expandafter\gdef\csname CGTip@eq:K-explicit\endcsname{Eq. (12.13), p. 53: K(s)=Pi\_rm tr(s) e\textasciicircum{}-h\_RVM(s)=(P\_Z1(s))/(P\_RVM\textasciicircum{}crit(s)) e\textasciicircum{}-h\_RVM(s) ,}
\expandafter\gdef\csname CGTip@rem:D-atomic\endcsname{Remark 12.4 (Continuous-to-atomic meaning of (D)) (p. 54): The prime-power count J\_Z1\textasciicircum{}rm M of zeta\_Z1 (Theorem 13.5) is an L\textasciicircum{}2\_rm loc count built from the zeros, and (D) makes its distributional derivative the atomic measure sum\_p,mm\textasciicircum{}-1delta\_p\textasciicircum{}m. The final step of the method is therefore the passage of zeta\_Z1 into the discrete Dirichlet category. Theorem 13.7 states it exa...}
\expandafter\gdef\csname CGTip@subsec:step3\endcsname{Section 12.4: Step 3: the Mertens bound from zeta\_Z1 = zeta (p. 54)}
\expandafter\gdef\csname CGTip@lem:zeta-growth\endcsname{Lemma 12.5 (Vertical growth of zeta) (p. 54): For (1)/(4)<=sigma <=4 and |t|>=2, |zeta(sigma+it)|<=15|t| .}
\expandafter\gdef\csname CGTip@eq:zeta-growth\endcsname{Eq. (12.14), p. 54: |zeta(sigma+it)|<=15|t| .}
\expandafter\gdef\csname CGTip@eq:zeta-integral-rep\endcsname{Eq. (12.15), p. 54: zeta(s)=(s)/(s-1)-sint\_1\textasciicircum{}inf yy\textasciicircum{}-s-1 dy .}
\expandafter\gdef\csname CGTip@lem:log-zeta\endcsname{Lemma 12.6 (Logarithm of zeta off the critical line) (p. 54): Assume zeta(s)!=0 for Re s>(1)/(2). Let 0<delta <=(1)/(4) and omega(delta):=(log(3-2delta))/(log(3-delta))<1 . Then, with the branch of log zeta equal to sum\_n>=2Lambda(n)n\textasciicircum{}-s/log n on Re s>1, |log zeta(sigma+it)|<=((48)/(delta)log|t|)\textasciicircum{}omega(delta)(sigma >=(1)/(2)+delta, |t|>=4). Consequently, for every eps>0 there...}
\expandafter\gdef\csname CGTip@eq:three-circles-exponent\endcsname{Eq. (12.16), p. 54: omega(delta):=(log(3-2delta))/(log(3-delta))<1 .}
\expandafter\gdef\csname CGTip@eq:log-zeta-bound\endcsname{Eq. (12.17), p. 54: |log zeta(sigma+it)|<=((48)/(delta)log|t|)\textasciicircum{}omega(delta)(sigma >=(1)/(2)+delta, |t|>=4).}
\expandafter\gdef\csname CGTip@eq:inv-zeta-bound\endcsname{Eq. (12.18), p. 54: |(1)/(zeta(sigma+it))|<=|t|\textasciicircum{}eps(sigma >=(1)/(2)+delta, |t|>= t\_2).}
\expandafter\gdef\csname CGTip@eq:borel-caratheodory\endcsname{Eq. (12.19), p. 54: max\_|s-s\_0|<= r\_3|log zeta(s)|<=(2r\_3)/(R-r\_3)max\_|s-s\_0|=RRe log zeta(s)+(R+r\_3)/(R-r\_3)|log zeta(s\_0)|,}
\expandafter\gdef\csname CGTip@eq:BC-bound\endcsname{Eq. (12.20), p. 54: max\_|s-s\_0|<= r\_3|log zeta(s)|<=(12)/(delta)(log(30|t|)+log zeta(2))<=(48)/(delta)log|t| .}
\expandafter\gdef\csname CGTip@eq:three-circles\endcsname{Eq. (12.21), p. 54: M\_2<= M\_1\textasciicircum{}1-omegaM\_3\textasciicircum{}omega, omega=(log(r\_2/r\_1))/(log(r\_3/r\_1))<=(log(3-2delta))/(log(3-delta))=omega(delta).}
\expandafter\gdef\csname CGTip@lem:perron-kernel\endcsname{Lemma 12.7 (Perron kernel) (p. 55): For y>0, y!=1, c>0, T>0 [ref, Ch. 17], |(1)/(2pi i)int\_c-iT\textasciicircum{}c+iT(y\textasciicircum{}s)/(s) ds- 1\_y>1|<= y\textasciicircum{}cmin (1,(1)/(T|log y|)).}
\expandafter\gdef\csname CGTip@eq:perron-kernel\endcsname{Eq. (12.22), p. 55: |(1)/(2pi i)int\_c-iT\textasciicircum{}c+iT(y\textasciicircum{}s)/(s) ds- 1\_y>1|<= y\textasciicircum{}cmin (1,(1)/(T|log y|)).}
\expandafter\gdef\csname CGTip@thm:step3\endcsname{Theorem 12.8 (Step 3) (p. 55): If zeta\_Z1=zeta, then for every eps>0 M(x)=O\_eps (x\textasciicircum{}1/2+eps)(x>=1).}
\expandafter\gdef\csname CGTip@eq:step3-bound\endcsname{Eq. (12.23), p. 55: M(x)=O\_eps (x\textasciicircum{}1/2+eps)(x>=1).}
\expandafter\gdef\csname CGTip@eq:perron-truncated\endcsname{Eq. (12.24), p. 55: M(x)=(1)/(2pi i)int\_c-iT\textasciicircum{}c+iT(x\textasciicircum{}s)/(s zeta(s)) ds+O(sum\_n>=1((x)/(n))\textasciicircum{}cmin (1,(1)/(T|log(x/n)|))).}
\expandafter\gdef\csname CGTip@eq:perron-error\endcsname{Eq. (12.25), p. 55: M(x)=(1)/(2pi i)int\_c-iT\textasciicircum{}c+iT(x\textasciicircum{}s)/(s zeta(s)) ds+O(log x).}
\expandafter\gdef\csname CGTip@eq:contour-shift\endcsname{Eq. (12.26), p. 55: int\_c-iT\textasciicircum{}c+iT=int\_(1)/(2)+delta-iT\textasciicircum{}(1)/(2)+delta+iT+int\_(1)/(2)+delta+iT\textasciicircum{}c+iT-int\_(1)/(2)+delta-iT\textasciicircum{}c-iT.}
\expandafter\gdef\csname CGTip@eq:horizontal-bound\endcsname{Eq. (12.27), p. 55: (1)/(2pi)int\_(1)/(2)+delta\textasciicircum{}c(x\textasciicircum{}sigma T\textasciicircum{}eps)/(T) dsigma <= x\textasciicircum{}cT\textasciicircum{}eps-1=e x\textasciicircum{}eps.}
\expandafter\gdef\csname CGTip@eq:left-bound\endcsname{Eq. (12.28), p. 55: (1)/(2pi)int\_-T\textasciicircum{}T(x\textasciicircum{}1/2+delta)/(|(1)/(2)+delta+it| |zeta((1)/(2)+delta+it)|) dt<= x\textasciicircum{}1/2+delta(C'(delta,eps)+int\_t\_2\textasciicircum{}Tt\textasciicircum{}eps-1dt)<<\_delta,epsx\textasciicircum{}1/2+delta+eps.}
\expandafter\gdef\csname CGTip@thm:transfer\endcsname{Theorem 12.9 (Transfer of the synthetic Mertens bound to M) (p. 55): Assume zeta\_Z1=zeta and let 0<delta <=(1)/(4), eps>0. Then: item on Re s>(1)/(2), (1)/(zeta(s))=F\_RVM(s) Pi\_rm tr(s)\textasciicircum{}-1, Pi\_rm tr=(P\_Z1)/(P\_RVM\textasciicircum{}crit)=(zeta)/(zeta\_RVM), where Pi\_rm tr is holomorphic and zero-free in Re s>(1)/(2) and Pi\_rm tr(1)=1; item there is t\_3=t\_3(delta,eps) such that, for (1)/(2)+delta <=sigma...}
\expandafter\gdef\csname CGTip@eq:transfer-factorization\endcsname{Eq. (12.29), p. 55: (1)/(zeta(s))=F\_RVM(s) Pi\_rm tr(s)\textasciicircum{}-1, Pi\_rm tr=(P\_Z1)/(P\_RVM\textasciicircum{}crit)=(zeta)/(zeta\_RVM),}
\expandafter\gdef\csname CGTip@eq:transfer-multiplier-bound\endcsname{Eq. (12.30), p. 56: |F\_RVM(sigma+it)|<=4L, |Pi\_rm tr(sigma+it)\textasciicircum{}+-1|<=|t|\textasciicircum{}eps;}
\expandafter\gdef\csname CGTip@eq:transfer-conclusion\endcsname{Eq. (12.31), p. 56: M(x)=O\_eps (x\textasciicircum{}1/2+eps).}
\expandafter\gdef\csname CGTip@eq:log-zetaRVM-bounded\endcsname{Eq. (12.32), p. 56: log zeta\_RVM=log varpi\_RVM+Log(1+w)+Log(1+ E\_0)=O(1).}
\expandafter\gdef\csname CGTip@eq:log-Pi-tr\endcsname{Eq. (12.33), p. 56: |log Pi\_rm tr(sigma+it)|=|log zeta-log zeta\_RVM|<=((48)/(delta)log|t|)\textasciicircum{}omega(delta)+O(1)<=eps log|t|}
\expandafter\gdef\csname CGTip@rem:transfer-mellin\endcsname{Remark 12.10 (The transfer in Mellin form) (p. 56): Formally, (12.29) reads M=(Pi\_rm tr\textasciicircum{}-1)\textasciicircum{}v*Mel M\_RVM\textasciicircum{}crit: the Mertens function is the synthetic Mertens function convolved with the inverse Mellin transform of the inverse transport multiplier. The synthetic factor F\_RVM contributes the exponent (1)/(2) and is bounded (Theorem 9.1); the multiplier contributes at mos...}
\expandafter\gdef\csname CGTip@cor:step3-D\endcsname{Corollary 12.11 (Step 3 through (D)) (p. 56): Condition (12.8) implies M(x)=O\_eps(x\textasciicircum{}1/2+eps) for every eps>0.}
\expandafter\gdef\csname CGTip@rem:Step3-status\endcsname{Remark 12.12 (Status of Step 3) (p. 56): The implication zeta\_Z1=zeta => M(x)=O\_eps(x\textasciicircum{}1/2+eps) is the classical consequence RH => M(x)=O\_eps(x\textasciicircum{}1/2+eps) (see [ref, S 14.25]) rewritten and reproved in the endpoint notation. It requires nothing beyond the identity: no separate statement about the zeros of zeta enters, because zeta=zeta\_Z1 already carries the...}
\expandafter\gdef\csname CGTip@subsec:reciprocal\endcsname{Section 12.5: Reciprocal endpoint and finite-height Perron integrals (p. 56)}
\expandafter\gdef\csname CGTip@eq:FZ1\endcsname{Eq. (12.34), p. 56: F\_Z1(s):=(1)/(zeta\_Z1(s))=F\_RVM(s) Pi\_rm tr(s)\textasciicircum{}-1,}
\expandafter\gdef\csname CGTip@eq:truncated-Perron\endcsname{Eq. (12.35), p. 56: M\_lambda,U(x):=(1)/(2pi i)int\_c-iU\textasciicircum{}c+iUF\_lambda(s)(x\textasciicircum{}s)/(s) ds, M\_Z1,U:=M\_0,U.}
\expandafter\gdef\csname CGTip@prop:finite-Perron\endcsname{Proposition 12.13 (Finite-height Perron integrals along the family) (p. 57): For fixed c>1, U<inf and x>0, M\_lambda,U(x)-> M\_Z1,U(x)(lambda decreasing to 0).}
\expandafter\gdef\csname CGTip@eq:finite-Perron-limit\endcsname{Eq. (12.36), p. 57: M\_lambda,U(x)-> M\_Z1,U(x)(lambda decreasing to 0).}
\expandafter\gdef\csname CGTip@thm:coefficient-criterion\endcsname{Theorem 12.14 (Reciprocal-coefficient criterion) (p. 57): Assume that in Re s>1 F\_Z1(s)=sum\_n>=1(mu\_Z1(n))/(n\textasciicircum{}s) absolutely, and let M\_Z1(x)=sum\_n<= xmu\_Z1(n), so that mu\_Z1(n)=M\_Z1(n)-M\_Z1(n-1). Then the following are equivalent: (i) M\_Z1(n)=M(n) for every n>=1; (ii) mu\_Z1=mu; (iii) F\_Z1=1/zeta in Re s>1; (iv) zeta\_Z1=zeta; (v) RH.}
\expandafter\gdef\csname CGTip@eq:Z1-reciprocal-Dirichlet\endcsname{Eq. (12.37), p. 57: F\_Z1(s)=sum\_n>=1(mu\_Z1(n))/(n\textasciicircum{}s)}
\expandafter\gdef\csname CGTip@eq:Z1-increments\endcsname{Eq. (12.38), p. 57: mu\_Z1(n)=M\_Z1(n)-M\_Z1(n-1).}
\expandafter\gdef\csname CGTip@rem:negative-convolution\endcsname{Remark 12.15 (A route that does not close Step 3) (p. 57): The following approach does not work. One may try to expand both factors of F\_Z1=F\_RVM Pi\_rm tr\textasciicircum{}-1 in (12.34) as absolutely convergent Dirichlet series and identify their Dirichlet convolution with mu. By Theorem 9.13 the expansion of F\_RVM does not exist in any half-plane, so there are no coefficients to convolve....}
\expandafter\gdef\csname CGTip@sec:euler-products\endcsname{Section 13: The endpoint at infinity and the Euler products (p. 57)}
\expandafter\gdef\csname CGTip@subsec:value-infinity\endcsname{Section 13.1: The value of the endpoint at infinity (p. 57)}
\expandafter\gdef\csname CGTip@eq:frak-e\endcsname{Eq. (13.1), p. 57: e(tau):=N\_zeta(tau)-N\_RVM(tau), e\_1(tau):=int\_tau\_0\textasciicircum{}tau e(v) dv(tau >=tau\_0=4pi).}
\expandafter\gdef\csname CGTip@eq:frak-e1-bound\endcsname{Eq. (13.2), p. 57: | e\_1(tau)|<= C\_ elog tau (tau >=tau\_0).}
\expandafter\gdef\csname CGTip@thm:Z1-clock\endcsname{Theorem 13.1 (Clock decomposition of the endpoint) (p. 57): Let s=(1)/(2)+z, z=sigma'+it, sigma'>0. There is an absolute constant R\_rm end>0 such that log zeta\_Z1(s)=log varpi\_Z1-int\_tau\_0\textasciicircum{}inf e\_1(tau)((1)/((tau-iz)\textasciicircum{}2)+(1)/((tau+iz)\textasciicircum{}2))dtau+O((1)/(|z|))(|z|>= R\_rm end), where varpi\_Z1:=P\_Z1((1)/(2))exp (nu\_3+2int\_tau\_0\textasciicircum{}inf(e\_1(tau))/(tau\textasciicircum{}2) dtau)>0 . Consequently: item for s...}
\expandafter\gdef\csname CGTip@eq:Z1-clock\endcsname{Eq. (13.3), p. 57: log zeta\_Z1(s)=log varpi\_Z1-int\_tau\_0\textasciicircum{}inf e\_1(tau)((1)/((tau-iz)\textasciicircum{}2)+(1)/((tau+iz)\textasciicircum{}2))dtau+O((1)/(|z|))(|z|>= R\_rm end),}
\expandafter\gdef\csname CGTip@eq:varpi-Z1\endcsname{Eq. (13.4), p. 57: varpi\_Z1:=P\_Z1((1)/(2))exp (nu\_3+2int\_tau\_0\textasciicircum{}inf(e\_1(tau))/(tau\textasciicircum{}2) dtau)>0 .}
\expandafter\gdef\csname CGTip@eq:Z1-growth\endcsname{Eq. (13.5), p. 57: |log zeta\_Z1(s)-log varpi\_Z1|<=(2(pi+2)C\_ e)/(sigma')(1+log (2|t|+2sigma'+tau\_0))+(C)/(|z|);}
\expandafter\gdef\csname CGTip@eq:varpi-Z1-product\endcsname{Eq. (13.6), p. 58: varpi\_Z1=varpi\_RVMprod\_k>=1((1)/(4)+tau\_k\textasciicircum{}2)/((1)/(4)+gamma\_k\textasciicircum{}2):=varpi\_RVMlim\_K->infprod\_k<= K((1)/(4)+tau\_k\textasciicircum{}2)/((1)/(4)+gamma\_k\textasciicircum{}2).}
\expandafter\gdef\csname CGTip@eq:Z1-split\endcsname{Eq. (13.7), p. 58: sum\_kf(gamma\_k)=Lambda\_rm sm(z)+int\_tau\_0\textasciicircum{}inf e\_1f'' dtau,}
\expandafter\gdef\csname CGTip@eq:log-kernel\endcsname{Eq. (13.8), p. 58: int\_tau\_0\textasciicircum{}inf(log tau)/((tau-a)\textasciicircum{}2+sigma'\textasciicircum{}2) dtau <=(pi)/(sigma')log X+4 (1+log X)/(X)<=(pi+2)/(sigma')(1+log X).}
\expandafter\gdef\csname CGTip@eq:varpi-ratio\endcsname{Eq. (13.9), p. 58: (varpi\_Z1)/(varpi\_RVM)=(P\_Z1((1)/(2)))/(P\_RVM\textasciicircum{}crit((1)/(2)))exp (2int\_tau\_0\textasciicircum{}inf(e\_1(tau)+ s\_1(tau))/(tau\textasciicircum{}2) dtau).}
\expandafter\gdef\csname CGTip@eq:count-difference-integral\endcsname{Eq. (13.10), p. 58: int\_tau\_0\textasciicircum{}inf(e\_1+ s\_1)/(tau\textasciicircum{}2) dtau=lim\_X->infint\_tau\_0\textasciicircum{}X(N\_zeta(tau)-n\_RVM(tau))/(tau) dtau=lim\_K->infsum\_k<= Klog(tau\_k)/(gamma\_k).}
\expandafter\gdef\csname CGTip@thm:varpi-Z1\endcsname{Theorem 13.2 (The value at infinity decides the endpoint identity) (p. 58): varpi\_Z1>=1, and varpi\_Z1=1 if and only if zeta\_Z1=zeta.}
\expandafter\gdef\csname CGTip@eq:pair-ratio\endcsname{Eq. (13.11), p. 59: (W\_rho\textasciicircum{}2)/((W\_rho-d\_rho\textasciicircum{}2)\textasciicircum{}2+4d\_rho\textasciicircum{}2gamma\textasciicircum{}2)* c\_gamma, c\_gamma=(|rho(1-rho)|\textasciicircum{}2)/(((1)/(4)+gamma\textasciicircum{}2)\textasciicircum{}2)=1+(2d\_rho\textasciicircum{}2(gamma\textasciicircum{}2-(1)/(4))+d\_rho\textasciicircum{}4)/(((1)/(4)+gamma\textasciicircum{}2)\textasciicircum{}2)>1 .}
\expandafter\gdef\csname CGTip@cor:product-identity\endcsname{Corollary 13.3 (The endpoint identity as an identity of the displacements) (p. 59): The following are equivalent: (D); zeta\_Z1=zeta; varpi\_Z1=1; lim\_K->infprod\_k<= K((1)/(4)+gamma\_k\textasciicircum{}2)/((1)/(4)+tau\_k\textasciicircum{}2)=varpi\_RVM (ordered product). Unconditionally lim\_sigma ->+infPi\_rm tr(sigma)=varpi\_Z1/varpi\_RVM=varpi\_Z1L>= L, and the equality lim\_sigma ->+infPi\_rm tr(sigma)=L=1/varpi\_RVM is a fourth condition eq...}
\expandafter\gdef\csname CGTip@eq:product-identity\endcsname{Eq. (13.12), p. 59: lim\_K->infprod\_k<= K((1)/(4)+gamma\_k\textasciicircum{}2)/((1)/(4)+tau\_k\textasciicircum{}2)=varpi\_RVM}
\expandafter\gdef\csname CGTip@prop:Z1-bounded\endcsname{Proposition 13.4 (Boundedness of the endpoint) (p. 59): Let sigma\_0>1. Then |zeta\_Z1| is bounded above and below on Re s>=sigma\_0 if and only if sup\_sigma'>=sigma\_0-(1)/(2), t>=0|Re int\_tau\_0\textasciicircum{}inf(e\_1(tau)+ s\_1(tau))/((tau-t+isigma')\textasciicircum{}2) dtau |<inf .}
\expandafter\gdef\csname CGTip@eq:conjugate-poisson\endcsname{Eq. (13.13), p. 59: sup\_sigma'>=sigma\_0-(1)/(2), t>=0|Re int\_tau\_0\textasciicircum{}inf(e\_1(tau)+ s\_1(tau))/((tau-t+isigma')\textasciicircum{}2) dtau |<inf .}
\expandafter\gdef\csname CGTip@subsec:euler-products\endcsname{Section 13.2: The Euler products of zeta\_RVM and zeta\_Z1 (p. 59)}
\expandafter\gdef\csname CGTip@thm:mellin-counts\endcsname{Theorem 13.5 (Prime-power counts as Mellin transforms) (p. 59): Let Xin RVM,Z1. There is a unique function J\_X\textasciicircum{}rm M on (1,inf) with J\_X\textasciicircum{}rm M(x)x\textasciicircum{}-cin L\textasciicircum{}2((1,inf),dx/x) for every c>1 and log zeta\_X(s)=sint\_1\textasciicircum{}inf J\_X\textasciicircum{}rm M(x) x\textasciicircum{}-s-1 dx(Re s>1), the branch of log zeta\_X being real on (1,inf). We use the notation int x\textasciicircum{}-s dJ\_X\textasciicircum{}rm M:=sint\_1\textasciicircum{}inf J\_X\textasciicircum{}rm M(x)x\textasciicircum{}-s-1dx distributionally; no...}
\expandafter\gdef\csname CGTip@eq:mellin-count\endcsname{Eq. (13.14), p. 59: log zeta\_X(s)=sint\_1\textasciicircum{}inf J\_X\textasciicircum{}rm M(x) x\textasciicircum{}-s-1 dx(Re s>1),}
\expandafter\gdef\csname CGTip@eq:mellin-count-laplace\endcsname{Eq. (13.15), p. 60: (log zeta\_X(s))/(s)=int\_0\textasciicircum{}inf j(u) e\textasciicircum{}-su du(Re s>c), e\textasciicircum{}-cujin L\textasciicircum{}2(0,inf).}
\expandafter\gdef\csname CGTip@thm:distributional-EF\endcsname{Theorem 13.6 (Distributional explicit formula) (p. 60): Let Xin RVM,Z1, with ordinates kappa\_k=tau\_k for X=RVM and kappa\_k=gamma\_k for X=Z1. In D'(1,inf), (log x) (dJ\_X\textasciicircum{}rm M)/(dx)=1-(2)/(sqrt(x)) Re sum\_k>=1x\textasciicircum{}ikappa\_k-(1)/(x(x\textasciicircum{}2-1)), the series converging in D'(1,inf). The explicit-formula count J\_X(x):=li(x)-2Re sum\_k>=1Ei (((1)/(2)+ikappa\_k)log x)-log 2+int\_x\textasciicircum{}inf(dy)/(...}
\expandafter\gdef\csname CGTip@eq:distributional-EF\endcsname{Eq. (13.16), p. 60: (log x) (dJ\_X\textasciicircum{}rm M)/(dx)=1-(2)/(sqrt(x)) Re sum\_k>=1x\textasciicircum{}ikappa\_k-(1)/(x(x\textasciicircum{}2-1)),}
\expandafter\gdef\csname CGTip@eq:J-X\endcsname{Eq. (13.17), p. 60: J\_X(x):=li(x)-2Re sum\_k>=1Ei (((1)/(2)+ikappa\_k)log x)-log 2+int\_x\textasciicircum{}inf(dy)/(y(y\textasciicircum{}2-1)log y)}
\expandafter\gdef\csname CGTip@eq:log-derivative-X\endcsname{Eq. (13.18), p. 60: -(zeta\_X')/(zeta\_X)(s)=(1)/(s-1)-sum\_k>=1((1)/(s-rho\_k)+(1)/(s-conj rho\_k))+sum\_n>=1((1)/(2n)-(1)/(s+2n))-(gamma\_E)/(2)-(1)/(2)log pi .}
\expandafter\gdef\csname CGTip@thm:cell-form\endcsname{Theorem 13.7 (Cell form of the endpoint identity) (p. 61): The following are equivalent. item (D), that is, zeta\_Z1=zeta. item J\_Z1\textasciicircum{}rm M is constant on every interval (n,n+1), n>=1, and its jumps b\_n:=v\_n-v\_n-1, where v\_n is its value on (n,n+1) and v\_0:=0, satisfy sum\_n<= X| b\_n|=O\_eps(X\textasciicircum{}1+eps) for every eps>0. item For every n>=1, in D'(n,n+1), 2 Re sum\_k>=1(x\textasciicircum{}igamma\_k-x\textasciicircum{}...}
\expandafter\gdef\csname CGTip@eq:cell-identity\endcsname{Eq. (13.19), p. 61: 2 Re sum\_k>=1(x\textasciicircum{}igamma\_k-x\textasciicircum{}itau\_k)=sqrt(x) log x J\_RVM'(x),}
\expandafter\gdef\csname CGTip@eq:cell-difference\endcsname{Eq. (13.20), p. 61: (log x)((dJ\_Z1\textasciicircum{}rm M)/(dx)-J\_RVM'(x))=-(2)/(sqrt(x)) Re sum\_k>=1(x\textasciicircum{}igamma\_k-x\textasciicircum{}itau\_k),}
\expandafter\gdef\csname CGTip@rem:cell-form\endcsname{Remark 13.8 (The final step in cell form) (p. 61): Theorem 13.7 is the exact form of the continuous-to-atomic passage of Remark 12.4. Both sides of (13.19) are built from the midpoint construction: the right side from the clock alone (J\_RVM' is bounded on [2,inf) by Theorem 7.1, equal to 1/log x plus the unit-periodic first harmonic and the harmonics with stationary...}
\expandafter\gdef\csname CGTip@sec:atomized\endcsname{Section 14: The atomized family and the programme for the final step (p. 61)}
\expandafter\gdef\csname CGTip@subsec:family-infinity\endcsname{Section 14.1: The family at infinity (p. 62)}
\expandafter\gdef\csname CGTip@prop:family-clock\endcsname{Proposition 14.1 (Clock decomposition along the family) (p. 62): There is an absolute constant C with | e\_1\textasciicircum{}(lambda)(tau)|<= Clog tau for tau >=tau\_0 and all lambda in[0,1]. Consequently Theorem 13.1 holds for every zeta\_lambda with e\_1 replaced by e\_1\textasciicircum{}(lambda), uniformly in lambda; in particular varpi\_lambda:=lim\_sigma ->+infzeta\_lambda(sigma) exists and varpi\_lambda=varpi\_RVMli...}
\expandafter\gdef\csname CGTip@eq:varpi-lambda\endcsname{Eq. (14.1), p. 62: varpi\_lambda=varpi\_RVMlim\_K->infprod\_k<= K((1)/(4)+tau\_k\textasciicircum{}2)/((1)/(4)+gamma\_k(lambda)\textasciicircum{}2), varpi\_1=varpi\_RVM,varpi\_0=varpi\_Z1.}
\expandafter\gdef\csname CGTip@eq:family-count-primitive\endcsname{Eq. (14.2), p. 62: int\_tau\_0\textasciicircum{}X(N\_lambda-n\_RVM)=sum\_k((X-gamma\_k(lambda))\textasciicircum{}+-(X-tau\_k)\textasciicircum{}+)=(1-lambda)sum\_gamma\_k(lambda),tau\_k<= Xeta\_k+O(log X),}
\expandafter\gdef\csname CGTip@eq:Upsilon\endcsname{Eq. (14.3), p. 62: Upsilon(lambda):=sum\_k>=1(2gamma\_k(lambda) eta\_k)/((1)/(4)+gamma\_k(lambda)\textasciicircum{}2)(summed in the order of k).}
\expandafter\gdef\csname CGTip@prop:family-flow\endcsname{Proposition 14.2 (The value at infinity along the family) (p. 62): The series (14.3) converges uniformly on [0,1], lambda ->log varpi\_lambda is C\textasciicircum{}1 on [0,1], and (d)/(dlambda)log varpi\_lambda=-Upsilon(lambda)=lim\_sigma ->+inf V\_lambda(sigma), Upsilon'(lambda)=sum\_k>=1(2eta\_k\textasciicircum{}2((1)/(4)-gamma\_k(lambda)\textasciicircum{}2))/(((1)/(4)+gamma\_k(lambda)\textasciicircum{}2)\textasciicircum{}2)<0 , with V\_lambda the generator of Theorem 11....}
\expandafter\gdef\csname CGTip@eq:varpi-derivative\endcsname{Eq. (14.4), p. 62: (d)/(dlambda)log varpi\_lambda=-Upsilon(lambda)=lim\_sigma ->+inf V\_lambda(sigma), Upsilon'(lambda)=sum\_k>=1(2eta\_k\textasciicircum{}2((1)/(4)-gamma\_k(lambda)\textasciicircum{}2))/(((1)/(4)+gamma\_k(lambda)\textasciicircum{}2)\textasciicircum{}2)<0 ,}
\expandafter\gdef\csname CGTip@eq:flow-integral\endcsname{Eq. (14.5), p. 62: int\_0\textasciicircum{}1Upsilon(lambda) dlambda=sum\_k>=1log((1)/(4)+tau\_k\textasciicircum{}2)/((1)/(4)+gamma\_k\textasciicircum{}2)=log L+log varpi\_Z1,}
\expandafter\gdef\csname CGTip@prop:varpi-verified\endcsname{Proposition 14.3 (The value at infinity and the verified zeros) (p. 63): Let T\_rm ver>= e\textasciicircum{}4. If every zero beta+igamma with 0<gamma <= T\_rm ver has beta=(1)/(2), then 0<=log varpi\_Z1<=0.255sum\_gamma>T\_rm ver(1)/(gamma\textasciicircum{}2)<=0.255((log(T\_rm ver/2pi)+1)/(2pi T\_rm ver)+(2log T\_rm ver+1)/(T\_rm ver\textasciicircum{}2)), the sum running over the zeros with gamma>T\_rm ver, with multiplicity. With T\_rm ver=3*10\textasciicircum{}12...}
\expandafter\gdef\csname CGTip@eq:varpi-verified\endcsname{Eq. (14.6), p. 63: 0<=log varpi\_Z1<=0.255sum\_gamma>T\_rm ver(1)/(gamma\textasciicircum{}2)<=0.255((log(T\_rm ver/2pi)+1)/(2pi T\_rm ver)+(2log T\_rm ver+1)/(T\_rm ver\textasciicircum{}2)),}
\expandafter\gdef\csname CGTip@subsec:smooth-cells\endcsname{Section 14.2: Atomization by smooth cells (p. 63)}
\expandafter\gdef\csname CGTip@def:atomized\endcsname{Definition 14.4 (Atomized family) (p. 63): The cell masses, the atomized zeta function and its coefficients are tilde c\_n(lambda):=< dpi\_lambda,varphi\_n>, tilde zeta\_Z,lambda(s):=prod\_n>=2(1-n\textasciicircum{}-s)\textasciicircum{}-tilde c\_n(lambda)=sum\_N>=1(tilde a\_N(lambda))/(N\textasciicircum{}s), tilde a\_N(lambda)=sum\_prod\_nn\textasciicircum{}r\_n=N prod\_n: r\_n>=1binom(tilde c\_n(lambda)+r\_n-1,r\_n), tilde a\_1(lambda)=1 .}
\expandafter\gdef\csname CGTip@eq:atomized-family\endcsname{Eq. (14.7), p. 63: tilde c\_n(lambda):=< dpi\_lambda,varphi\_n>, tilde zeta\_Z,lambda(s):=prod\_n>=2(1-n\textasciicircum{}-s)\textasciicircum{}-tilde c\_n(lambda)=sum\_N>=1(tilde a\_N(lambda))/(N\textasciicircum{}s),}
\expandafter\gdef\csname CGTip@eq:atomized-coefficients\endcsname{Eq. (14.8), p. 63: tilde a\_N(lambda)=sum\_prod\_nn\textasciicircum{}r\_n=N prod\_n: r\_n>=1binom(tilde c\_n(lambda)+r\_n-1,r\_n), tilde a\_1(lambda)=1 .}
\expandafter\gdef\csname CGTip@eq:J-lambda-derivative\endcsname{Eq. (14.9), p. 63: (d J\_lambda)/(d lambda)(x)=2sum\_k>=1eta\_k Im(x\textasciicircum{}rho\_k,lambda)/(rho\_k,lambda), (d)/(d lambda)((log x)(dJ\_lambda)/(dx))=(2log x)/(sqrt(x))sum\_k>=1eta\_ksin (gamma\_k(lambda)log x).}
\expandafter\gdef\csname CGTip@subsec:joining-factor\endcsname{Section 14.3: The joining factor cell by cell (p. 64)}
\expandafter\gdef\csname CGTip@eq:atomization-kernel\endcsname{Eq. (14.10), p. 64: a\_s(x):=sum\_n>=2varphi\_n(x) n\textasciicircum{}-s-x\textasciicircum{}-s, h\_n(lambda;s):=< dpi\_lambda, varphi\_n(x)(n\textasciicircum{}-s-x\textasciicircum{}-s)> .}
\expandafter\gdef\csname CGTip@eq:first-cell\endcsname{Eq. (14.11), p. 64: u\_lambda(s):=C\_lambda 2\textasciicircum{}-s+sint\_1\textasciicircum{}2J\textasciicircum{}rm M\_lambda(x) x\textasciicircum{}-s-1 dx, C\_lambda:=J\textasciicircum{}rm M\_lambda-J\_lambda ,}
\expandafter\gdef\csname CGTip@prop:first-cell\endcsname{Proposition 14.5 (The first cell) (p. 64): Let lambda in[0,1]. item u\_lambda is entire, and for Re s>1 log zeta\_lambda(s)= u\_lambda(s)+< dJ\textasciicircum{}*\_lambda,x\textasciicircum{}-s>, < dJ\textasciicircum{}*\_lambda,x\textasciicircum{}-s>:=sint\_2\textasciicircum{}inf J\_lambda(x) x\textasciicircum{}-s-1 dx . item u\_lambda(sigma+it)=o(|t|) as |t|->inf, uniformly for sigma in compact sets. item u\_lambda is constant if and only if J\_lambda=0 on (1,2); in th...}
\expandafter\gdef\csname CGTip@eq:first-cell-split\endcsname{Eq. (14.12), p. 64: log zeta\_lambda(s)= u\_lambda(s)+< dJ\textasciicircum{}*\_lambda,x\textasciicircum{}-s>, < dJ\textasciicircum{}*\_lambda,x\textasciicircum{}-s>:=sint\_2\textasciicircum{}inf J\_lambda(x) x\textasciicircum{}-s-1 dx .}
\expandafter\gdef\csname CGTip@lem:L2-partial-sums\endcsname{Lemma 14.6 (Mean-square convergence of the ordinate series) (p. 64): Let kappa be a list with |kappa\_k-tau\_k|<= C\_eta (in particular kappa=(gamma\_k(lambda)) for any lambda in[0,1]), rho\_k:=(1)/(2)+ikappa\_k, and let Iin(1,inf) be compact. The partial sums S\_K(x):=sum\_k<= K(x\textasciicircum{}rho\_k/rho\_k+x\textasciicircum{}conj rho\_k/conj rho\_k) and E\_K(x):=sum\_k<= K(Ei(rho\_klog x)+Ei(conj rho\_klog x)) converge in L\textasciicircum{}2(...}
\expandafter\gdef\csname CGTip@eq:L2-partial-sums\endcsname{Eq. (14.13), p. 64: ||S\_K'-S\_K||\_L\textasciicircum{}2(I)\textasciicircum{}2+||E\_K'-E\_K||\_L\textasciicircum{}2(I)\textasciicircum{}2 <= C\_Isum\_K<k<= K'(log(kappa\_k+2))/(kappa\_k\textasciicircum{}2)(K<K'),}
\expandafter\gdef\csname CGTip@eq:L2-partial-sums-proof\endcsname{Eq. (14.14), p. 65: int omega(u)|sum\_K<k<= K'(e\textasciicircum{}ikappa\_ku)/(rho\_k)|\textasciicircum{}2du <<sum\_K<k,k'<= K'((1+|kappa\_k-kappa\_k'|)\textasciicircum{}-2)/(kappa\_kkappa\_k'); <<sum\_K<k<= K'(1)/(kappa\_k\textasciicircum{}2)sum\_j>=0(no.k':j<=|kappa\_k'-kappa\_k|<j+1)/((1+j)\textasciicircum{}2) <<sum\_K<k<= K'(log(kappa\_k+2))/(kappa\_k\textasciicircum{}2),}
\expandafter\gdef\csname CGTip@thm:joining-factor\endcsname{Theorem 14.7 (The joining factor) (p. 65): Let tilde h\_lambda:=log tilde zeta\_Z,lambda-log zeta\_lambda on the half-plane where tilde zeta\_Z,lambda converges absolutely (Theorem 14.13). Then tilde zeta\_Z,lambda=zeta\_lambda e\textasciicircum{}tilde h\_lambda, tilde h\_lambda(s)=- u\_lambda(s)+sum\_m>=1(1)/(m)sum\_n>=2 h\_n(lambda;ms), and, explicitly, tilde h\_lambda(s)=- u\_lambda(s)...}
\expandafter\gdef\csname CGTip@eq:joining-factor\endcsname{Eq. (14.15), p. 65: tilde zeta\_Z,lambda=zeta\_lambda e\textasciicircum{}tilde h\_lambda, tilde h\_lambda(s)=- u\_lambda(s)+sum\_m>=1(1)/(m)sum\_n>=2 h\_n(lambda;ms),}
\expandafter\gdef\csname CGTip@eq:joining-explicit\endcsname{Eq. (14.16), p. 65: tilde h\_lambda(s)=- u\_lambda(s)+sum\_m>=1(1)/(m)sum\_j>=1(mu(j))/(j)int\_2\textasciicircum{}inf a\_ms(y\textasciicircum{}j) (1)/(log y)(1-(2)/(sqrt(y))Re sum\_k>=1y\textasciicircum{}igamma\_k(lambda)-(1)/(y(y\textasciicircum{}2-1)))dy,}
\expandafter\gdef\csname CGTip@eq:joining-derivative\endcsname{Eq. (14.17), p. 65: d\_lambda tilde h\_lambda(s)=-d\_lambda u\_lambda(s)+sum\_m>=1(1)/(m)sum\_j>=1(mu(j))/(j)int\_2\textasciicircum{}inf a\_ms(y\textasciicircum{}j) (2)/(sqrt(y))sum\_k>=1eta\_ksin (gamma\_k(lambda)log y) dy .}
\expandafter\gdef\csname CGTip@prop:cell-explicit\endcsname{Proposition 14.8 (Explicit cell defects) (p. 66): For every n>=2, lambda in[0,1] and Re s>1, h\_n(lambda;s)=sum\_j>=1(mu(j))/(j)int\_2\textasciicircum{}inf varphi\_n(y\textasciicircum{}j)(n\textasciicircum{}-s-y\textasciicircum{}-js) (1)/(log y)(1-(2)/(sqrt(y))Re sum\_k>=1y\textasciicircum{}igamma\_k(lambda)-(1)/(y(y\textasciicircum{}2-1)))dy, d\_lambda h\_n(lambda;s)=sum\_j>=1(mu(j))/(j)int\_2\textasciicircum{}inf varphi\_n(y\textasciicircum{}j)(n\textasciicircum{}-s-y\textasciicircum{}-js) (2)/(sqrt(y))sum\_k>=1eta\_ksin (gamma\_k(lambda)log y...}
\expandafter\gdef\csname CGTip@eq:cell-explicit\endcsname{Eq. (14.18), p. 66: h\_n(lambda;s)=sum\_j>=1(mu(j))/(j)int\_2\textasciicircum{}inf varphi\_n(y\textasciicircum{}j)(n\textasciicircum{}-s-y\textasciicircum{}-js) (1)/(log y)(1-(2)/(sqrt(y))Re sum\_k>=1y\textasciicircum{}igamma\_k(lambda)-(1)/(y(y\textasciicircum{}2-1)))dy,}
\expandafter\gdef\csname CGTip@eq:cell-explicit-derivative\endcsname{Eq. (14.19), p. 66: d\_lambda h\_n(lambda;s)=sum\_j>=1(mu(j))/(j)int\_2\textasciicircum{}inf varphi\_n(y\textasciicircum{}j)(n\textasciicircum{}-s-y\textasciicircum{}-js) (2)/(sqrt(y))sum\_k>=1eta\_ksin (gamma\_k(lambda)log y) dy ,}
\expandafter\gdef\csname CGTip@prop:cell-limit\endcsname{Proposition 14.9 (The limit lambda ->0 and the meaning of a vanishing cell) (p. 66): For every n>=2 and Re s>1, h\_n(lambda;s)-> h\_n(0;s) and tilde c\_n(lambda)->tilde c\_n(0) as lambda ->0\textasciicircum{}+, and h\_n(0;s)= h\_n(1;s)-int\_0\textasciicircum{}1d\_lambda h\_n(lambda;s) dlambda . Moreover h\_n(0;s)=0 for all Re s>1 if and only if varphi\_n dpi\_Z1=tilde c\_n(0) delta\_n.}
\expandafter\gdef\csname CGTip@eq:cell-flow\endcsname{Eq. (14.20), p. 66: h\_n(0;s)= h\_n(1;s)-int\_0\textasciicircum{}1d\_lambda h\_n(lambda;s) dlambda .}
\expandafter\gdef\csname CGTip@subsec:mean-square\endcsname{Section 14.4: Mean square of the cell masses (p. 67)}
\expandafter\gdef\csname CGTip@eq:VW\endcsname{Eq. (14.21), p. 67: V\_n(gamma):=sum\_j>=1(mu(j))/(j)int\_2\textasciicircum{}inf varphi\_n(y\textasciicircum{}j) y\textasciicircum{}-1/2e\textasciicircum{}igamma log y dy, W\_n(gamma):=sum\_j>=1(mu(j))/(j)int\_2\textasciicircum{}inf varphi\_n(y\textasciicircum{}j) (y\textasciicircum{}-1/2e\textasciicircum{}igamma log y)/(log y) dy ,}
\expandafter\gdef\csname CGTip@eq:cell-mass-W\endcsname{Eq. (14.22), p. 67: tilde c\_n(lambda)= m\_n-2 Re sum\_k>=1W\_n(gamma\_k(lambda)), d\_lambda tilde c\_n(lambda)=2sum\_k>=1eta\_k Im V\_n(gamma\_k(lambda)),}
\expandafter\gdef\csname CGTip@eq:cell-mass-power-two\endcsname{Eq. (14.23), p. 67: -(mu(j))/(j)int\_2\textasciicircum{}3 J\_lambda(y) (d)/(dy)varphi\_n(y\textasciicircum{}j) dy,}
\expandafter\gdef\csname CGTip@lem:mean-value-ordinates\endcsname{Lemma 14.10 (Mean value over the moving ordinates) (p. 67): There is C such that for all lambda in[0,1], T>=2, u\_0in R and complex a\_k with |a\_k|<=1, int\_u\_0-1\textasciicircum{}u\_0+1|sum\_gamma\_k(lambda)<= Ta\_k e\textasciicircum{}igamma\_k(lambda)u|\textasciicircum{}2du<= C Tlog\textasciicircum{}2T .}
\expandafter\gdef\csname CGTip@eq:mean-value-ordinates\endcsname{Eq. (14.24), p. 67: int\_u\_0-1\textasciicircum{}u\_0+1|sum\_gamma\_k(lambda)<= Ta\_k e\textasciicircum{}igamma\_k(lambda)u|\textasciicircum{}2du<= C Tlog\textasciicircum{}2T .}
\expandafter\gdef\csname CGTip@eq:fejer-pairs\endcsname{Eq. (14.25), p. 67: <<no.(k,k'):gamma\_k(lambda)<= T, |gamma\_k(lambda)-gamma\_k'(lambda)|<=1.}
\expandafter\gdef\csname CGTip@thm:mean-square-cells\endcsname{Theorem 14.11 (Mean square of the cell masses along the family) (p. 67): For every eps>0 there is C\_eps such that for all X>=2 and all lambda in[0,1] sum\_X<n<=2X|tilde c\_n(lambda)|\textasciicircum{}2<= C\_eps X\textasciicircum{}1+eps, sum\_X<n<=2X|d\_lambda tilde c\_n(lambda)|\textasciicircum{}2<= C\_eps X\textasciicircum{}1+eps. In particular the same bound already includes lambda=0: sum\_X<n<=2X|tilde c\_n(0)|\textasciicircum{}2<= C\_eps X\textasciicircum{}1+eps.}
\expandafter\gdef\csname CGTip@eq:mean-square-cells\endcsname{Eq. (14.26), p. 68: sum\_X<n<=2X|tilde c\_n(lambda)|\textasciicircum{}2<= C\_eps X\textasciicircum{}1+eps, sum\_X<n<=2X|d\_lambda tilde c\_n(lambda)|\textasciicircum{}2<= C\_eps X\textasciicircum{}1+eps.}
\expandafter\gdef\csname CGTip@eq:A-nj\endcsname{Eq. (14.27), p. 68: A\_n,j=sum\_kint\_2\textasciicircum{}inf varphi\_n(y\textasciicircum{}j) (y\textasciicircum{}-1/2+igamma\_k(lambda))/(log y) dy , 1<= j<= J:=floor((log(2X+1))/(log 2)) .}
\expandafter\gdef\csname CGTip@eq:cell-CS\endcsname{Eq. (14.28), p. 68: |int varphi\_n(y\textasciicircum{}j) S\_n,j,lambda(y) dy|\textasciicircum{}2<=||varphi||\_inf\textasciicircum{}2 |I\_n,j|int\_I\_n,j| S\_n,j,lambda|\textasciicircum{}2dy .}
\expandafter\gdef\csname CGTip@eq:cell-block\endcsname{Eq. (14.29), p. 68: sum\_X<n<=2X|int varphi\_n(y\textasciicircum{}j)S dy|\textasciicircum{}2<< X\textasciicircum{}1/j-1((j)/(log X))\textasciicircum{}2int\_u\_0-1\textasciicircum{}u\_0+1|sum\_gamma\_k(lambda)<= Te\textasciicircum{}igamma\_k(lambda)u|\textasciicircum{}2du<< j\textasciicircum{}2X\textasciicircum{}1/j+2eps.}
\expandafter\gdef\csname CGTip@cor:abscissa-one\endcsname{Corollary 14.12 (Abscissa 1 for every member) (p. 68): For every lambda in[0,1] and sigma>1, sum\_n>=2|tilde c\_n(lambda)| n\textasciicircum{}-sigma<inf ; hence tilde zeta\_Z,lambda and 1/tilde zeta\_Z,lambda are absolutely convergent Dirichlet series in Re s>1, for every lambda, including the endpoint lambda=0. Moreover, uniformly in lambda, no.n<= X:|tilde c\_n(lambda)|>n\textasciicircum{}delta<<\_eps X\textasciicircum{}1-2de...}
\expandafter\gdef\csname CGTip@eq:abscissa-one\endcsname{Eq. (14.30), p. 68: sum\_n>=2|tilde c\_n(lambda)| n\textasciicircum{}-sigma<inf ;}
\expandafter\gdef\csname CGTip@subsec:atomized-continuation\endcsname{Section 14.5: Analytic continuation of the atomized family (p. 69)}
\expandafter\gdef\csname CGTip@thm:atomized-continuation\endcsname{Theorem 14.13 (Continuation of the atomized family) (p. 69): Let lambda in[0,1]. item For every c>(1)/(2), (pi\_lambda- L) x\textasciicircum{}-cin L\textasciicircum{}2((2,inf),dx/x). item |tilde c\_n(lambda)|<= C(1+I\_n\textasciicircum{}1/2) with I\_n as below, and sum\_n|tilde c\_n(lambda)|n\textasciicircum{}-sigma<inf for sigma>(3)/(2); by Corollary 14.12 this holds for sigma>1, and tilde zeta\_Z,lambda is an absolutely convergent Dirichlet series...}
\expandafter\gdef\csname CGTip@eq:cell-defect-bound\endcsname{Eq. (14.31), p. 69: | h\_n(lambda;s)|<= C|s| n\textasciicircum{}-sigma-1(1+I\_n\textasciicircum{}1/2), sum\_nI\_n\textasciicircum{}1/2n\textasciicircum{}-sigma-1<inf (sigma>c).}
\expandafter\gdef\csname CGTip@thm:closing\endcsname{Theorem 14.14 (Closing criterion; via Hamburger's theorem [ref]) (p. 69): The following are equivalent. item (D), that is, zeta\_Z1=zeta. item tilde h\_0==0. item J\_Z1=0 on (1,2), and for every n>=2 and every Re s>1, int\_0\textasciicircum{}1d\_lambda h\_n(lambda;s) dlambda= h\_n(1;s). In this case tilde c\_n(0)= 1\_rm prime(n), tilde a\_N(0)=1 for all N, and varpi\_Z1=1.}
\expandafter\gdef\csname CGTip@eq:closing-cells\endcsname{Eq. (14.32), p. 69: int\_0\textasciicircum{}1d\_lambda h\_n(lambda;s) dlambda= h\_n(1;s).}
\expandafter\gdef\csname CGTip@rem:A6-scale\endcsname{Remark 14.15 (The pointwise size of the cell masses) (p. 70): Theorem 14.14 needs no pointwise bound on tilde c\_n(0): its proof uses the cell masses only through (14.30). The bound tilde c\_n(0)=O(n\textasciicircum{}eps) for every n is not proved directly; Corollary 14.12 restricts the exceptions to a set of n<= X of size O(X\textasciicircum{}1-2delta+eps) for the threshold n\textasciicircum{}delta, uniformly in lambda, and und...}
\expandafter\gdef\csname CGTip@subsec:atomized-endpoints\endcsname{Section 14.6: The two endpoints of the atomized family (p. 70)}
\expandafter\gdef\csname CGTip@prop:atomized-clock\endcsname{Proposition 14.16 (Atomization at the clock endpoint) (p. 70): At lambda=1: item tilde c\_n(1)=O(1), and tilde zeta\_Z,1 is an absolutely convergent Dirichlet series in Re s>1; item tilde h\_1(sigma)->-log varpi\_RVM=log L as sigma ->+inf; this limit is !=0 exactly when varpi\_RVM!=1, and it is >0 by Theorem 9.11; item tilde h\_1 is not constant: for every sigma>1, zeta\_RVM(sigma+it)...}
\expandafter\gdef\csname CGTip@eq:duality\endcsname{Eq. (14.33), p. 70: underbrace zeta\_lambda\_substack functional equation zeros (1)/(2)+- igamma\_k(lambda)* e\textasciicircum{}tilde h\_lambda = underbrace tilde zeta\_Z,lambda\_substack Dirichlet series Euler product, tilde h\_1 not constant, (D)<=>tilde h\_0==0 .}
\expandafter\gdef\csname CGTip@eq:atomized-mobius\endcsname{Eq. (14.34), p. 70: (1)/(tilde zeta\_Z,lambda(s))=prod\_n>=2(1-n\textasciicircum{}-s)\textasciicircum{}tilde c\_n(lambda)=sum\_N>=1(tilde mu\_lambda(N))/(N\textasciicircum{}s), tilde mu\_lambda(N)=sum\_prod\_nn\textasciicircum{}r\_n=N prod\_n: r\_n>=1(-1)\textasciicircum{}r\_nbinom(tilde c\_n(lambda),r\_n),}
\expandafter\gdef\csname CGTip@eq:atomized-mertens\endcsname{Eq. (14.35), p. 70: tilde M\_lambda(x):=sum\_N<= xtilde mu\_lambda(N), (1)/(tilde zeta\_Z,lambda)=e\textasciicircum{}-tilde h\_lambdaF\_lambda ,}
\expandafter\gdef\csname CGTip@prop:Mellin-Mertens\endcsname{Proposition 14.17 (Mellin form of the endpoint Mertens function) (p. 70): There are c\_1>=2 and a unique measurable B\_Z1 on (0,inf), vanishing on (0,1), with x\textasciicircum{}-cB\_Z1in L\textasciicircum{}2((1,inf),dx/x) and F\_Z1(s)=sint\_1\textasciicircum{}inf B\_Z1(x) x\textasciicircum{}-s-1 dx=exp (-int\_[1,inf)x\textasciicircum{}-s dJ\textasciicircum{}rm M\_Z1(x))(Re s>c\_1), for every c>c\_1. Moreover (D) holds if and only if B\_Z1=M almost everywhere; and if F\_Z1 has an absolutely convergen...}
\expandafter\gdef\csname CGTip@eq:Mellin-Mertens\endcsname{Eq. (14.36), p. 71: F\_Z1(s)=sint\_1\textasciicircum{}inf B\_Z1(x) x\textasciicircum{}-s-1 dx=exp (-int\_[1,inf)x\textasciicircum{}-s dJ\textasciicircum{}rm M\_Z1(x))(Re s>c\_1),}
\expandafter\gdef\csname CGTip@subsec:hamburger-table\endcsname{Section 14.7: Hamburger's hypotheses at the two endpoints (p. 71)}
\expandafter\gdef\csname CGTip@eq:complement\endcsname{Eq. (14.37), p. 72: underbrace zeta\_lambda\_substack keeps the functional equation and the zeros on Re s=(1)/(2)longleftrightarrow underbrace tilde zeta\_Z,lambda\_substack keeps the Dirichlet series and the Euler product, tilde zeta\_Z,lambda=zeta\_lambda e\textasciicircum{}tilde h\_lambda.}
\expandafter\gdef\csname CGTip@eq:closure-common\endcsname{Eq. (14.38), p. 72: tilde h\_0==0<=> J\_Z1=0 on (1,2), h\_n(0;s)=0 for all n>=2, Re s>1, and varpi\_Z1=1 .}
\expandafter\gdef\csname CGTip@subsec:programme\endcsname{Section 14.8: Programme (p. 73)}
\expandafter\gdef\csname CGTip@rem:input-A7\endcsname{Remark 14.18 (Input for A7) (p. 74): Steps A1--A6 use only two properties of an ordinate list: bounded displacement from the clock and a primitive of the count difference of size O(log tau). The list (tau\_k) has both properties, and for it the corresponding function zeta\_RVM is not a Dirichlet series (Theorem 9.13). Step A7 must therefore use a propert...}
\expandafter\gdef\csname CGTip@sec:criteria\endcsname{Section 15: Criteria for the final step and the clock end of the family (p. 74)}
\expandafter\gdef\csname CGTip@eq:cell-moments\endcsname{Eq. (15.1), p. 74: hat varphi(xi):=int\_Rvarphi(t)e\textasciicircum{}-2pi ixi t dt, M\_r:=int\_Rt\textasciicircum{}rvarphi(t) dt, C\_r:=int\_Rt\textasciicircum{}rvarphi(t)cos(2pi t) dt .}
\expandafter\gdef\csname CGTip@eq:partition-fourier\endcsname{Eq. (15.2), p. 74: hat varphi(q)=int\_0\textasciicircum{}1sum\_min Zvarphi(t-m) e\textasciicircum{}-2pi iqt dt=delta\_q0(qin Z),}
\expandafter\gdef\csname CGTip@subsec:open-path\endcsname{Section 15.1: The open path (p. 75)}
\expandafter\gdef\csname CGTip@prop:nonconstant\endcsname{Proposition 15.1 (No member of the open path is a Dirichlet series) (p. 75): Let 0<lambda <=1. Then zeta\_lambda is not a constant multiple of a Dirichlet series sum\_na\_nn\textasciicircum{}-s converging absolutely in Re s>1. In particular the joining factor tilde h\_lambda is not constant.}
\expandafter\gdef\csname CGTip@subsec:convexity\endcsname{Section 15.2: Convexity along the family (p. 75)}
\expandafter\gdef\csname CGTip@prop:convexity\endcsname{Proposition 15.2 (Convexity on the real axis) (p. 75): For every real sigma>1 the function lambda ->log zeta\_lambda(sigma) is strictly convex on [0,1]: d\_lambda\textasciicircum{}2log zeta\_lambda(sigma)=sum\_k>=1eta\_k\textasciicircum{}2(C(sigma-(1)/(2),gamma\_k(lambda))- C((1)/(2),gamma\_k(lambda)))>0, C(a,gamma):=(2(a\textasciicircum{}2-gamma\textasciicircum{}2))/((a\textasciicircum{}2+gamma\textasciicircum{}2)\textasciicircum{}2). Consequently, with the generator V\_lambda of Theorem 11.3,...}
\expandafter\gdef\csname CGTip@eq:convexity-second\endcsname{Eq. (15.3), p. 75: d\_lambda\textasciicircum{}2log zeta\_lambda(sigma)=sum\_k>=1eta\_k\textasciicircum{}2(C(sigma-(1)/(2),gamma\_k(lambda))- C((1)/(2),gamma\_k(lambda)))>0, C(a,gamma):=(2(a\textasciicircum{}2-gamma\textasciicircum{}2))/((a\textasciicircum{}2+gamma\textasciicircum{}2)\textasciicircum{}2).}
\expandafter\gdef\csname CGTip@eq:convexity-sandwich\endcsname{Eq. (15.4), p. 75: log zeta\_RVM(sigma)- V\_1(sigma)<log zeta\_Z1(sigma)<log zeta\_RVM(sigma)- V\_0(sigma)(sigma>1).}
\expandafter\gdef\csname CGTip@subsec:pair\endcsname{Section 15.3: Hamburger's theorem for the pair (p. 76)}
\expandafter\gdef\csname CGTip@thm:pair-criterion\endcsname{Theorem 15.3 (Hamburger's theorem for the pair (zeta\_lambda,tilde zeta\_Z,lambda)) (p. 76): Let lambda in[0,1]. Suppose that the joining factor e\textasciicircum{}tilde h\_lambda extends to a meromorphic function on C of finite order with finitely many poles, and that e\textasciicircum{}tilde h\_lambda(s)=e\textasciicircum{}tilde h\_lambda(1-s) . Then tilde zeta\_Z,lambda=zeta. For lambda=0, moreover, zeta\_Z1=zeta and tilde h\_0==0.}
\expandafter\gdef\csname CGTip@eq:pair-symmetry\endcsname{Eq. (15.5), p. 76: e\textasciicircum{}tilde h\_lambda(s)=e\textasciicircum{}tilde h\_lambda(1-s) .}
\expandafter\gdef\csname CGTip@thm:double-hamburger\endcsname{Theorem 15.4 (Real-variable and vertical forms of Hamburger's theorem for the pair) (p. 76): For lambda in[0,1], y>0 and sigma>1 put T\_lambda(y):=sum\_N>=1tilde a\_N(lambda) e\textasciicircum{}-pi N\textasciicircum{}2y-y\textasciicircum{}-1/2sum\_N>=1tilde a\_N(lambda) e\textasciicircum{}-pi N\textasciicircum{}2/y-(1)/(2)(y\textasciicircum{}-1/2-1), k\_n(lambda):=sum\_d| ntilde a\_d(lambda) mu(n/d), V\_sigma(lambda):=sum\_n>=2|k\_n(lambda)|\textasciicircum{}2 n\textasciicircum{}-2sigma. item tilde zeta\_Z,lambda/zeta=sum\_n>=1k\_n(lambda)n\textasciicircum{}-s converges...}
\expandafter\gdef\csname CGTip@eq:theta-defect\endcsname{Eq. (15.6), p. 76: T\_lambda(y):=sum\_N>=1tilde a\_N(lambda) e\textasciicircum{}-pi N\textasciicircum{}2y-y\textasciicircum{}-1/2sum\_N>=1tilde a\_N(lambda) e\textasciicircum{}-pi N\textasciicircum{}2/y-(1)/(2)(y\textasciicircum{}-1/2-1),}
\expandafter\gdef\csname CGTip@eq:vertical-defect\endcsname{Eq. (15.7), p. 76: k\_n(lambda):=sum\_d| ntilde a\_d(lambda) mu(n/d), V\_sigma(lambda):=sum\_n>=2|k\_n(lambda)|\textasciicircum{}2 n\textasciicircum{}-2sigma.}
\expandafter\gdef\csname CGTip@eq:vertical-mean\endcsname{Eq. (15.8), p. 76: lim\_T->inf(1)/(2T)int\_-T\textasciicircum{}T|(tilde zeta\_Z,lambda(sigma+it))/(zeta(sigma+it))|\textasciicircum{}2dt=1+ V\_sigma(lambda) >=1 .}
\expandafter\gdef\csname CGTip@eq:double-hamburger-limit\endcsname{Eq. (15.9), p. 76: (D)<=> T\_0==0<=> V\_sigma(0)=0<=>lim\_lambda ->0\textasciicircum{}+T\_lambda(y)=0 (y>0)<=>lim\_lambda ->0\textasciicircum{}+ V\_sigma(lambda)=0 .}
\expandafter\gdef\csname CGTip@eq:quotient-three\endcsname{Eq. (15.10), p. 77: (tilde zeta\_Z,lambda)/(zeta)=e\textasciicircum{}tilde h\_lambda (zeta\_lambda)/(zeta)}
\expandafter\gdef\csname CGTip@prop:clock-asymmetry\endcsname{Proposition 15.5 (At the clock end the symmetry fails) (p. 77): The joining factor e\textasciicircum{}tilde h\_1 has no meromorphic continuation to C that satisfies (15.5).}
\expandafter\gdef\csname CGTip@subsec:positivity\endcsname{Section 15.4: Positivity and the value at infinity (p. 77)}
\expandafter\gdef\csname CGTip@lem:summation-positive\endcsname{Lemma 15.6 (Summation formula for a representing measure) (p. 77): Let lambda in[0,1] and let A be a positive measure on [1,inf) with zeta\_lambda(s)=int\_[1,inf)u\textasciicircum{}-s d A(u)(Re s>1). Then for every even Schwartz function g g(0)+2int\_[1,inf) g d A=hat g(0)+2int\_[1,inf)hat g d A .}
\expandafter\gdef\csname CGTip@eq:representing-measure\endcsname{Eq. (15.11), p. 77: zeta\_lambda(s)=int\_[1,inf)u\textasciicircum{}-s d A(u)(Re s>1).}
\expandafter\gdef\csname CGTip@eq:summation-positive\endcsname{Eq. (15.12), p. 77: g(0)+2int\_[1,inf) g d A=hat g(0)+2int\_[1,inf)hat g d A .}
\expandafter\gdef\csname CGTip@eq:tate-local\endcsname{Eq. (15.13), p. 77: pi\textasciicircum{}-s/2Gamma((s)/(2)) hat g\textasciicircum{}sharp(1-s)=pi\textasciicircum{}-(1-s)/2Gamma((1-s)/(2)) g\textasciicircum{}sharp(s)}
\expandafter\gdef\csname CGTip@eq:summation-shift\endcsname{Eq. (15.14), p. 77: int g d A=(1)/(2)hat g(0)-(1)/(2) g(0)+(1)/(2pi i)int\_(1-c)zeta\_lambda(s) g\textasciicircum{}sharp(s) ds .}
\expandafter\gdef\csname CGTip@thm:positivity\endcsname{Theorem 15.7 (Positivity criterion) (p. 78): Let lambda in[0,1]. If zeta\_lambda has a representation (15.11) with a positive measure A, then varpi\_lambda= A(1)<=1 . In particular, if zeta\_Z1 has such a representation, then varpi\_Z1=1 and zeta\_Z1=zeta.}
\expandafter\gdef\csname CGTip@eq:positivity-bound\endcsname{Eq. (15.15), p. 78: varpi\_lambda= A(1)<=1 .}
\expandafter\gdef\csname CGTip@eq:extremal-function\endcsname{Eq. (15.16), p. 78: g(u):=((sin pi u)/(pi u))\textasciicircum{}4(1-u\textasciicircum{}2), hat g(xi)= B(xi)+(1)/(4pi\textasciicircum{}2) B''(xi) .}
\expandafter\gdef\csname CGTip@eq:lp-inequality\endcsname{Eq. (15.17), p. 78: 1 >= g\_epsilon(0)+2int g\_epsilon d A=hat g\_epsilon(0)+2int hat g\_epsilon d A >= hat g\_epsilon(0)+2varpi\_lambda hat g\_epsilon(1) .}
\expandafter\gdef\csname CGTip@eq:lp-value\endcsname{Eq. (15.18), p. 78: varpi\_lambda <=(1-hat g(0))/(2hat g(1))=((1)/(3)+(1)/(2pi\textasciicircum{}2))/((1)/(3)+(1)/(2pi\textasciicircum{}2))=1 .}
\expandafter\gdef\csname CGTip@cor:positivity\endcsname{Corollary 15.8 (Positive form of the final step) (p. 78): The following are equivalent: (i) (D); (ii) zeta\_Z1 is the Mellin--Stieltjes transform (15.11) of a positive measure on [1,inf); (iii) J\textasciicircum{}rm M\_Z1, extended by 0 to (0,1], is non-decreasing on (0,inf).}
\expandafter\gdef\csname CGTip@lem:summation-all\endcsname{Lemma 15.9 (Mellin form and summation formula for every member) (p. 78): There is c\_1>=2 such that for every lambda in[0,1] there is a unique measurable function A\_lambda on (0,inf), vanishing on (0,1), with x\textasciicircum{}-c A\_lambda in L\textasciicircum{}2((1,inf),dx/x) for every c>c\_1 and zeta\_lambda(s)=sint\_1\textasciicircum{}inf A\_lambda(x) x\textasciicircum{}-s-1 dx(Re s>c\_1). If zeta\_lambda has a representation (15.11), then A\_lambda(x)= A([1,...}
\expandafter\gdef\csname CGTip@eq:mellin-A\endcsname{Eq. (15.19), p. 78: zeta\_lambda(s)=sint\_1\textasciicircum{}inf A\_lambda(x) x\textasciicircum{}-s-1 dx(Re s>c\_1).}
\expandafter\gdef\csname CGTip@eq:summation-all\endcsname{Eq. (15.20), p. 79: g(0)+2< A\_lambda, g>=hat g(0)+2< A\_lambda,hat g> .}
\expandafter\gdef\csname CGTip@eq:parseval-A\endcsname{Eq. (15.21), p. 79: < A\_lambda, g>=(1)/(2pi i)int\_(c)zeta\_lambda(s) g\textasciicircum{}sharp(s) ds .}
\expandafter\gdef\csname CGTip@prop:one-number\endcsname{Proposition 15.10 (The value at infinity as one number) (p. 79): Let g\_epsilon be as in the proof of Theorem 15.7, w\_epsilon:=hat g\_epsilon- g\_epsilon, which is non-negative on |u|>=1, and A'\_lambda(x):= A\_lambda(x)-varpi\_lambda for x>=1. For every lambda in[0,1] and epsilon>0, varpi\_lambda w\_epsilon(1)=(1)/(2)(1-hat g\_epsilon(0))-< A'\_lambda, w\_epsilon >, < A'\_lambda, w\_epsilon...}
\expandafter\gdef\csname CGTip@eq:one-number-eps\endcsname{Eq. (15.22), p. 79: varpi\_lambda w\_epsilon(1)=(1)/(2)(1-hat g\_epsilon(0))-< A'\_lambda, w\_epsilon >, < A'\_lambda, w\_epsilon >:=-int\_1\textasciicircum{}inf A'\_lambda(x) w\_epsilon'(x) dx,}
\expandafter\gdef\csname CGTip@eq:one-number\endcsname{Eq. (15.23), p. 79: varpi\_lambda-1=-(W(lambda))/(hat g(1)), hat g(1)=(1)/(6)+(1)/(4pi\textasciicircum{}2).}
\expandafter\gdef\csname CGTip@thm:real-axis\endcsname{Theorem 15.11 (The endpoint dominates zeta on the real axis) (p. 79): For every real sigma>1, zeta\_Z1(sigma) >= zeta(sigma), and equality for one sigma>1 holds if and only if zeta\_Z1=zeta. In particular prod\_k>=1(1+(2)/((1)/(4)+gamma\_k\textasciicircum{}2)) >= (pi)/(3),with equality if and only if (D).}
\expandafter\gdef\csname CGTip@eq:real-axis\endcsname{Eq. (15.24), p. 79: zeta\_Z1(sigma) >= zeta(sigma),}
\expandafter\gdef\csname CGTip@eq:sigma-two\endcsname{Eq. (15.25), p. 79: prod\_k>=1(1+(2)/((1)/(4)+gamma\_k\textasciicircum{}2)) >= (pi)/(3),with equality if and only if (D).}
\expandafter\gdef\csname CGTip@eq:pair-monotone\endcsname{Eq. (15.26), p. 80: (2)/(W)-(2(W-d\_rho\textasciicircum{}2))/((W-d\_rho\textasciicircum{}2)\textasciicircum{}2+4d\_rho\textasciicircum{}2gamma\textasciicircum{}2)=(2d\_rho\textasciicircum{}2 (4gamma\textasciicircum{}2+d\_rho\textasciicircum{}2-W))/(W((W-d\_rho\textasciicircum{}2)\textasciicircum{}2+4d\_rho\textasciicircum{}2gamma\textasciicircum{}2)),}
\expandafter\gdef\csname CGTip@prop:open-path-varpi\endcsname{Proposition 15.12 (The final step on the open path) (p. 80): The following are equivalent: (i) (D); (ii) varpi\_lambda<1 for every lambda in(0,1]; (iii) varpi\_lambda\_m<=1 for some sequence lambda\_m->0\textasciicircum{}+; (iv) W(lambda)>0 for every lambda in(0,1]. The implication (i)=>(ii) uses varpi\_RVM<1 (Theorem 9.11), which is the case lambda=1 of (ii); the implications (ii)=>(iii)=>(i) and...}
\expandafter\gdef\csname CGTip@subsec:crossing\endcsname{Section 15.5: The crossing point of the family and what convexity can give (p. 80)}
\expandafter\gdef\csname CGTip@prop:crossing\endcsname{Proposition 15.13 (The crossing point) (p. 80): item Upsilon(0)>-log varpi\_RVM>0.020615, without hypothesis. item If RH holds, then for every lambda in(0,1] f(lambda)<=lambda f(1)and f(lambda)+lambda Upsilon(lambda)<=0 . item If RH fails, there is exactly one lambda\textasciicircum{}*in(0,1) with varpi\_lambda\textasciicircum{}*=1; f>0 on [0,lambda\textasciicircum{}*), f<0 on (lambda\textasciicircum{}*,1], and (f(0))/(Upsilon(0))<...}
\expandafter\gdef\csname CGTip@eq:crossing-refuters\endcsname{Eq. (15.27), p. 80: f(lambda)<=lambda f(1)and f(lambda)+lambda Upsilon(lambda)<=0 .}
\expandafter\gdef\csname CGTip@eq:crossing-bounds\endcsname{Eq. (15.28), p. 80: (f(0))/(Upsilon(0))<=lambda\textasciicircum{}*<=(f(0))/(f(0)+|f(1)|).}
\expandafter\gdef\csname CGTip@rem:convexity-limits\endcsname{Remark 15.14 (What convexity cannot give) (p. 81): For a convex function, values in the interior bound the interior from the end points (chord), and the end points only from below (tangents): f(0)>= f(lambda)+lambda Upsilon(lambda). The only upper bound for f(0) obtained from convexity on [0,1] is f(0)<f(1)+Upsilon(0), which cannot give f(0)<=0 because Upsilon(0)>|f...}
\expandafter\gdef\csname CGTip@prop:kappa-family\endcsname{Proposition 15.15 (A family with the same properties and the opposite endpoint) (p. 81): Let kappa\_1=12.6 and, for 0<=lambda <=1, let kappa(lambda) be the list ((1-lambda)kappa\_1+lambda tau\_1,tau\_2,tau\_3,...), zeta\textasciicircum{}kappa\_lambda=P\_kappa(lambda)/G with P\_kappa(s)=prod\_k(1+q(s)/((1)/(4)+kappa\_k\textasciicircum{}2)), and varpi\textasciicircum{}kappa\_lambda=lim\_sigma ->+infzeta\textasciicircum{}kappa\_lambda(sigma). Then every zeta\textasciicircum{}kappa\_lambda has order one,...}
\expandafter\gdef\csname CGTip@rem:hypothesis-X\endcsname{Remark 15.16 (The hypothesis added to the functional equation) (p. 81): The functional equation, the order and the zeros on the critical line do not characterize zeta: every zeta\_lambda has them. A converse theorem for zeta\_Z1 needs one further hypothesis on the function A\_Z1 of Lemma 15.9, and each of the following suffices and is equivalent to (D): A\_Z1 is a step function with jumps b...}
\expandafter\gdef\csname CGTip@subsec:clock-density\endcsname{Section 15.6: The density of the clock (p. 82)}
\expandafter\gdef\csname CGTip@eq:U-def\endcsname{Eq. (15.29), p. 82: U(M,t\_0):=int\_t\_0\textasciicircum{}infe\textasciicircum{}-iMvarXi(t) dt:=lim\_T->infint\_t\_0\textasciicircum{}Te\textasciicircum{}-iMvarXi(t) dt .}
\expandafter\gdef\csname CGTip@lem:stationary-U\endcsname{Lemma 15.17 (Stationary phase for varXi) (p. 82): The limit (15.29) exists, and there is an absolute constant C\_ U with | U(M,t\_0)-sqrt((2pi)/(M)) e\textasciicircum{}i(M-pi/4)|<=(C\_ U)/(M)(M>=1, 0<t\_0<=(1)/(2)).}
\expandafter\gdef\csname CGTip@eq:U-asymptotic\endcsname{Eq. (15.30), p. 82: | U(M,t\_0)-sqrt((2pi)/(M)) e\textasciicircum{}i(M-pi/4)|<=(C\_ U)/(M)(M>=1, 0<t\_0<=(1)/(2)).}
\expandafter\gdef\csname CGTip@eq:U-substitution\endcsname{Eq. (15.31), p. 82: p(w):=(dt)/(d w)=(w)/(log t), p(0)=1, q(w):=(p(w)-1)/(w)=(1)/(log t)-(1)/(w).}
\expandafter\gdef\csname CGTip@eq:U-split\endcsname{Eq. (15.32), p. 82: e\textasciicircum{}-iM U(M,t\_0)=int\_ w\_0\textasciicircum{}inf e\textasciicircum{}-iM w\textasciicircum{}2/2 d w+int\_ w\_0\textasciicircum{}inf e\textasciicircum{}-iM w\textasciicircum{}2/2 w q(w) d w .}
\expandafter\gdef\csname CGTip@eq:clock-harmonic\endcsname{Eq. (15.33), p. 82: R\_m(y):=-2(-1)\textasciicircum{}m m\textasciicircum{}-3/2 y\textasciicircum{}(1)/(2m)-(1)/(2)cos (2pi m y\textasciicircum{}1/m-((7m+1)pi)/(4)),}
\expandafter\gdef\csname CGTip@thm:clock-density\endcsname{Theorem 15.18 (The density of the clock) (p. 82): J\_RVM is continuously differentiable on (2,inf), and for y>=16 J\_RVM'(y)=(1)/(log y)+2cos(2pi y)-(y\textasciicircum{}-1/4)/(sqrt(2))cos (4pisqrt(y)+(pi)/(4))+O(y\textasciicircum{}-1/3). On [2,inf) the measure dpi\_RVM is the sum of the atoms (mu(j))/(j)J\_RVM(2) delta\_2\textasciicircum{}j, j>=1, and of a density that is continuous off the points 2\textasciicircum{}j and equals J\_RVM'(...}
\expandafter\gdef\csname CGTip@eq:clock-density\endcsname{Eq. (15.34), p. 82: J\_RVM'(y)=(1)/(log y)+2cos(2pi y)-(y\textasciicircum{}-1/4)/(sqrt(2))cos (4pisqrt(y)+(pi)/(4))+O(y\textasciicircum{}-1/3).}
\expandafter\gdef\csname CGTip@eq:H-harmonic\endcsname{Eq. (15.35), p. 82: H\_m\textasciicircum{}+-(v):=int\_tau\_0\textasciicircum{}inf e\textasciicircum{}iPsi\_m\textasciicircum{}+-(tau) dtau ,}
\expandafter\gdef\csname CGTip@eq:S-zero\endcsname{Eq. (15.36), p. 83: sum\_kD\_k=int\_u\_1\textasciicircum{}u\_2(e\textasciicircum{}v/2)/(v) S\_0(v) dv, S\_0(v):=S\_rm sm(v)+sum\_m>=1((-1)\textasciicircum{}m v)/(2pi m)(H\_m\textasciicircum{}-(v)- H\_m\textasciicircum{}+(v)),}
\expandafter\gdef\csname CGTip@eq:J-prime-exact\endcsname{Eq. (15.37), p. 83: J\_RVM'(y)=(1)/(log y)(1-(1)/(y(y\textasciicircum{}2-1))-(2)/(sqrt(y)) Re S\_0(log y)).}
\expandafter\gdef\csname CGTip@eq:H-by-U\endcsname{Eq. (15.38), p. 83: H\_m\textasciicircum{}-(v)= t\_m e\textasciicircum{}-7pi im/4 U(m t\_m,tau\_0/ t\_m), tau\_0/ t\_m=2e\textasciicircum{}-v/m<=(1)/(2)(y>=4\textasciicircum{}m).}
\expandafter\gdef\csname CGTip@eq:H-asymptotic\endcsname{Eq. (15.39), p. 83: ((-1)\textasciicircum{}mv)/(2pi m) H\_m\textasciicircum{}-(v)=(-1)\textasciicircum{}m (v)/(m\textasciicircum{}3/2) y\textasciicircum{}(1)/(2m)exp (i(2pi m y\textasciicircum{}1/m-((7m+1)pi)/(4)))+O((v)/(m\textasciicircum{}2)).}
\expandafter\gdef\csname CGTip@cor:clock-cells\endcsname{Corollary 15.19 (Cell masses and cell defects at the clock end) (p. 83): For n>=17, with 1\_Pow\_2(n)=1 if n=2\textasciicircum{}j for an integer j>=1 and 0 otherwise (the product below is 0 when n is not a power of 2, and log\_2n=j when it is), c\_n\textasciicircum{}RVM=int\_n-1\textasciicircum{}n((1)/(log x)+2cos(2pi x))dx+ 1\_Pow\_2(n) (mu(log\_2n))/(log\_2n) J\_RVM(2)+O(n\textasciicircum{}-1/4), tilde c\_n(1)=(1)/(log n)+ 1\_Pow\_2(n) (mu(log\_2n))/(log\_2n) J\_RVM(2...}
\expandafter\gdef\csname CGTip@eq:clock-cell-sharp\endcsname{Eq. (15.40), p. 83: c\_n\textasciicircum{}RVM=int\_n-1\textasciicircum{}n((1)/(log x)+2cos(2pi x))dx+ 1\_Pow\_2(n) (mu(log\_2n))/(log\_2n) J\_RVM(2)+O(n\textasciicircum{}-1/4),}
\expandafter\gdef\csname CGTip@eq:clock-cell-smooth\endcsname{Eq. (15.41), p. 83: tilde c\_n(1)=(1)/(log n)+ 1\_Pow\_2(n) (mu(log\_2n))/(log\_2n) J\_RVM(2)+O(n\textasciicircum{}-1/4),}
\expandafter\gdef\csname CGTip@eq:clock-defect\endcsname{Eq. (15.42), p. 83: | h\_n(1;s)|<= C\_sigma\_1(|s| n\textasciicircum{}-sigma-4/3+|s|(1+|s|) n\textasciicircum{}-sigma-2).}
\expandafter\gdef\csname CGTip@subsec:height-kernels\endcsname{Section 15.7: The cell masses in the heights of the ordinates (p. 84)}
\expandafter\gdef\csname CGTip@prop:height-kernels\endcsname{Proposition 15.20 (Height kernel of the cell masses) (p. 84): For every eps>0, lambda in[0,1] and n>=3 that is not a power of 2, tilde c\_n(lambda)= m\_n-(2)/(sqrt(n) log n)sum\_k>=1hat varphi((gamma\_k(lambda))/(2pi n))cos (gamma\_k(lambda)log n)+E\_n(lambda), where sum\_X<n<=2X|E\_n(lambda)|\textasciicircum{}2<= C\_eps X\textasciicircum{}1/2+eps uniformly in lambda (sum over the n that are not powers of 2). The kerne...}
\expandafter\gdef\csname CGTip@eq:height-kernel\endcsname{Eq. (15.43), p. 84: tilde c\_n(lambda)= m\_n-(2)/(sqrt(n) log n)sum\_k>=1hat varphi((gamma\_k(lambda))/(2pi n))cos (gamma\_k(lambda)log n)+E\_n(lambda),}
\expandafter\gdef\csname CGTip@eq:Xi-rotation\endcsname{Eq. (15.44), p. 84: Xi\_n(Omega):=sum\_k>=1hat varphi((tau\_k)/(2pi n))n\textasciicircum{}itau\_ke\textasciicircum{}-iOmega eta\_k(0<=Omega <=log n).}
\expandafter\gdef\csname CGTip@prop:flow-count\endcsname{Proposition 15.21 (The flow of a cell mass and the count difference) (p. 84): For every n>=3 that is not a power of 2, tilde c\_n(0)-tilde c\_n(1)=-2int\_0\textasciicircum{}1sum\_k>=1eta\_k Im V\_n(gamma\_k(lambda)) dlambda=-2 Im int\_tau\_0\textasciicircum{}inf V\_n(t)(N\_zeta(t)-n\_RVM(t)) dt .}
\expandafter\gdef\csname CGTip@eq:flow-count\endcsname{Eq. (15.45), p. 84: tilde c\_n(0)-tilde c\_n(1)=-2int\_0\textasciicircum{}1sum\_k>=1eta\_k Im V\_n(gamma\_k(lambda)) dlambda=-2 Im int\_tau\_0\textasciicircum{}inf V\_n(t)(N\_zeta(t)-n\_RVM(t)) dt .}
\expandafter\gdef\csname CGTip@prop:defect-kernel\endcsname{Proposition 15.22 (Height kernel of the cell defects) (p. 85): For every eps>0, lambda in[0,1], Re s>0 and n>=3 that is not a power of 2, with an implied constant uniform for s in compact subsets of Re s>0 and for lambda in[0,1], h\_n(lambda;s)=-(s n\textasciicircum{}-s-3/2)/(pi log n)sum\_k>=1hat varphi'((gamma\_k(lambda))/(2pi n))sin (gamma\_k(lambda)log n)+O\_s,eps(n\textasciicircum{}-sigma-3/4+eps).}
\expandafter\gdef\csname CGTip@eq:defect-kernel\endcsname{Eq. (15.46), p. 85: h\_n(lambda;s)=-(s n\textasciicircum{}-s-3/2)/(pi log n)sum\_k>=1hat varphi'((gamma\_k(lambda))/(2pi n))sin (gamma\_k(lambda)log n)+O\_s,eps(n\textasciicircum{}-sigma-3/4+eps).}
\expandafter\gdef\csname CGTip@rem:scale-test\endcsname{Remark 15.23 (The scale of the displacements) (p. 85): By Proposition 15.1, a property of the sequence (eta\_k) that yields (D) is false for every sequence (c eta\_k) with 0<= c<1. Properties that are invariant under eta\_k-> c eta\_k with c>0 (the signs of the eta\_k, the non-vanishing of the eta\_k, their boundedness) hold along the whole open path and do not suffice; the p...}
\expandafter\gdef\csname CGTip@eq:eta-exact\endcsname{Eq. (15.47), p. 85: eta\_k N\_RVM'(xi\_k)=N\_zeta\textasciicircum{}*(gamma\_k)-N\_RVM(gamma\_k)for some xi\_k between gamma\_k and tau\_k .}
\expandafter\gdef\csname CGTip@subsec:gaussian-defect\endcsname{Section 15.8: The Gaussian defect of the atomization (p. 85)}
\expandafter\gdef\csname CGTip@lem:vertical-growth\endcsname{Lemma 15.24 (Vertical growth of the joining factor) (p. 85): For every lambda in[0,1] and (1)/(2)<sigma\_0<sigma\_1, tilde h\_lambda(sigma+it)=o(|t|) as |t|->inf, uniformly for sigma\_0<=sigma <=sigma\_1.}
\expandafter\gdef\csname CGTip@eq:cell-RL\endcsname{Eq. (15.48), p. 86: h\_n(lambda;ms)= int varphi\_n(x)(n\textasciicircum{}-ms-x\textasciicircum{}-ms)d L(x)-int E(x)varphi\_n'(x)(n\textasciicircum{}-ms-x\textasciicircum{}-ms)dx; -msint E(x)varphi\_n(x) x\textasciicircum{}-ms-1dx .}
\expandafter\gdef\csname CGTip@eq:gaussian-defect\endcsname{Eq. (15.49), p. 86: Theta\_lambda(y):=(1)/(4pi i)int\_(c)zeta\_lambda(s) Gamma((s)/(2)) pi\textasciicircum{}-s/2 y\textasciicircum{}-s/2 ds, G\_lambda(y):=sum\_N>=1tilde a\_N(lambda) e\textasciicircum{}-pi N\textasciicircum{}2y-Theta\_lambda(y).}
\expandafter\gdef\csname CGTip@prop:gaussian-defect\endcsname{Proposition 15.25 (Gaussian defect) (p. 86): Let lambda in[0,1]. item Theta\_lambda(y)=y\textasciicircum{}-1/2 Theta\_lambda(1/y)+(1)/(2)(y\textasciicircum{}-1/2-1) for y>0, and Theta\_lambda(y)=O\_A(y\textasciicircum{}-A) as y->inf for every A. item The coefficients tilde a\_N(lambda) satisfy the theta relation of zeta, sum\_N>=1tilde a\_N(lambda) e\textasciicircum{}-pi N\textasciicircum{}2y=y\textasciicircum{}-1/2sum\_N>=1tilde a\_N(lambda) e\textasciicircum{}-pi N\textasciicircum{}2/y+(1)/(2)(y\textasciicircum{}-1/2...}
\expandafter\gdef\csname CGTip@eq:theta-atomized\endcsname{Eq. (15.50), p. 86: sum\_N>=1tilde a\_N(lambda) e\textasciicircum{}-pi N\textasciicircum{}2y=y\textasciicircum{}-1/2sum\_N>=1tilde a\_N(lambda) e\textasciicircum{}-pi N\textasciicircum{}2/y+(1)/(2)(y\textasciicircum{}-1/2-1)(y>0),}
\expandafter\gdef\csname CGTip@eq:defect-selfdual\endcsname{Eq. (15.51), p. 86: G\_lambda(y)=y\textasciicircum{}-1/2 G\_lambda(1/y)(y>0).}
\expandafter\gdef\csname CGTip@eq:defect-expansion\endcsname{Eq. (15.52), p. 86: G\_lambda(y)=(1)/(2)(e\textasciicircum{}tilde h\_lambda(1)-1) y\textasciicircum{}-1/2+O\_delta (y\textasciicircum{}-1/4-delta),}
\expandafter\gdef\csname CGTip@cor:residue-condition\endcsname{Corollary 15.26 (The density of the atomized integers) (p. 87): sum\_N<= xtilde a\_N(lambda) has Gaussian mean density e\textasciicircum{}tilde h\_lambda(1) in the sense of (15.52) (a statement about the Gaussian weights e\textasciicircum{}-pi N\textasciicircum{}2y as y->0\textasciicircum{}+, not a Tauberian statement about the partial sums), and (D) implies tilde h\_0(1)=0,that is, u\_0(1)=sum\_m>=1(1)/(m)sum\_n>=2 h\_n(0;m).}
\expandafter\gdef\csname CGTip@eq:residue-condition\endcsname{Eq. (15.53), p. 87: tilde h\_0(1)=0,that is, u\_0(1)=sum\_m>=1(1)/(m)sum\_n>=2 h\_n(0;m).}
\expandafter\gdef\csname CGTip@eq:defect-expansion-clock\endcsname{Eq. (15.54), p. 87: G\_1(y)=(1)/(2)(e\textasciicircum{}tilde h\_1(1)-1) y\textasciicircum{}-1/2+O\_delta (y\textasciicircum{}-delta)(y->0\textasciicircum{}+),}
\expandafter\gdef\csname CGTip@subsec:entire-lambda\endcsname{Section 15.9: Entire dependence on lambda and the exceptional set (p. 87)}
\expandafter\gdef\csname CGTip@prop:entire-lambda\endcsname{Proposition 15.27 (The cell masses as entire functions of lambda) (p. 87): Let n>=3 not be a power of 2. item The real-axis function lambda ->tilde c\_n(lambda) has the entire extension [ tilde C\_n(lambda):= m\_n-sum\_k(W\_n(gamma\_k+lambda eta\_k)+conj W\_n(gamma\_k+conj lambda eta\_k)), ] which is of exponential type at most C\_eta log(n+1) and agrees with tilde c\_n(lambda) for real lambda. item F...}
\expandafter\gdef\csname CGTip@eq:lambda-derivatives\endcsname{Eq. (15.55), p. 87: d\_lambda\textasciicircum{}rtilde c\_n(lambda)=-2 Re sum\_k>=1eta\_k\textasciicircum{} r W\_n\textasciicircum{}(r)(gamma\_k(lambda)), |d\_lambda\textasciicircum{}rtilde c\_n(lambda)|<= C\_r,eps(C\_eta log(n+1))\textasciicircum{}r n\textasciicircum{}1/2+eps.}
\expandafter\gdef\csname CGTip@prop:exceptional-set\endcsname{Proposition 15.28 (The exceptional set) (p. 88): The set E:=lambda in[0,1]: tilde zeta\_Z,lambda=zeta is finite, 1not in E, and (D) holds if and only if 0in E. For every lambda not in E the hypotheses of Theorem 15.3 do not hold together, T\_lambda not ==0 and V\_sigma(lambda)>0.}
\expandafter\gdef\csname CGTip@eq:exceptional-set\endcsname{Eq. (15.56), p. 88: E:=lambda in[0,1]: tilde zeta\_Z,lambda=zeta}
\expandafter\gdef\csname CGTip@subsec:cell-mass-form\endcsname{Section 15.10: The final step in the cell masses (p. 88)}
\expandafter\gdef\csname CGTip@cor:cell-mass-form\endcsname{Corollary 15.29 (The final step is one number per cell) (p. 88): The following are equivalent: item (D); item tilde c\_n(0)= 1\_rm prime(n) for every n>=2; item there is a list kappain O (Section 16) with tilde c\_n(kappa)= 1\_rm prime(n) for every n>=2.}
\expandafter\gdef\csname CGTip@eq:cell-mass-euler\endcsname{Eq. (15.57), p. 88: tilde zeta\_Z,kappa(s)=prod\_n>=2(1-n\textasciicircum{}-s)\textasciicircum{}-tilde c\_n(kappa)=prod\_p(1-p\textasciicircum{}-s)\textasciicircum{}-1=zeta(s)(Re s>1),}
\expandafter\gdef\csname CGTip@prop:witness\endcsname{Proposition 15.30 (Every point of the exceptional set is a witness) (p. 88): For the set E of (15.56), E!=empty <=>(D),so that either E=empty or 0in E . If (D) holds, then E0=lambda in(0,1]: D\_n(lambda)=0 for every n>=2, D\_n(lambda):=tilde c\_n(lambda)-tilde c\_n(0), and for n>=3 not a power of 2, the real-axis function D\_n(lambda)=-2Re sum\_k>=1(W\_n(gamma\_k+lambda eta\_k)-W\_n(gamma\_k)) has the...}
\expandafter\gdef\csname CGTip@eq:witness\endcsname{Eq. (15.58), p. 88: E!=empty <=>(D),so that either E=empty or 0in E .}
\expandafter\gdef\csname CGTip@eq:witness-zeros\endcsname{Eq. (15.59), p. 88: E0=lambda in(0,1]: D\_n(lambda)=0 for every n>=2, D\_n(lambda):=tilde c\_n(lambda)-tilde c\_n(0),}
\expandafter\gdef\csname CGTip@rem:dual-function\endcsname{Remark 15.31 (The dual function of the atomized family) (p. 89): For Re s<(1)/(2) put B\textasciicircum{}v\_lambda(s):=zeta\_lambda(s) e\textasciicircum{}tilde h\_lambda(1-s). By the functional equation of zeta\_lambda, pi\textasciicircum{}-s/2Gamma ((s)/(2))B\textasciicircum{}v\_lambda(s)=pi\textasciicircum{}-(1-s)/2Gamma ((1-s)/(2)) tilde zeta\_Z,lambda(1-s)(Re s<(1)/(2)), so B\textasciicircum{}v\_lambda is the function attached to tilde zeta\_Z,lambda by the functional equation. To as...}
\expandafter\gdef\csname CGTip@eq:dual-function\endcsname{Eq. (15.60), p. 89: pi\textasciicircum{}-s/2Gamma ((s)/(2))B\textasciicircum{}v\_lambda(s)=pi\textasciicircum{}-(1-s)/2Gamma ((1-s)/(2)) tilde zeta\_Z,lambda(1-s)(Re s<(1)/(2)),}
\expandafter\gdef\csname CGTip@prop:per-cell\endcsname{Proposition 15.32 (The identity that remains, cell by cell) (p. 89): Let n>=17 not be a power of 2. Then tilde c\_n(0)= 1\_rm prime(n) if and only if 1\_rm prime(n)-tilde c\_n(1)=-2 Im int\_tau\_0\textasciicircum{}inf V\_n(t)(N\_zeta(t)-n\_RVM(t)) dt, tilde c\_n(1)=(1)/(log n)+O(n\textasciicircum{}-1/4).}
\expandafter\gdef\csname CGTip@eq:per-cell\endcsname{Eq. (15.61), p. 89: 1\_rm prime(n)-tilde c\_n(1)=-2 Im int\_tau\_0\textasciicircum{}inf V\_n(t)(N\_zeta(t)-n\_RVM(t)) dt, tilde c\_n(1)=(1)/(log n)+O(n\textasciicircum{}-1/4).}
\expandafter\gdef\csname CGTip@subsec:symmetric-factors\endcsname{Section 15.11: Symmetric joining factors (p. 89)}
\expandafter\gdef\csname CGTip@rem:pair-no-slack\endcsname{Remark 15.33 (The symmetry criterion has no slack at lambda=0) (p. 89): At lambda=0 the hypothesis of Theorem 15.3 is equivalent to its conclusion. If (D) holds, then tilde h\_0==0 (Theorem 14.14), and e\textasciicircum{}tilde h\_0=1 is symmetric and of finite order; hence e\textasciicircum{}tilde h\_0 symmetric of finite order<=>tilde h\_0==0<=>(D), and in the same way T\_0==0 and V\_sigma(0)=0 are equivalent forms of (D), n...}
\expandafter\gdef\csname CGTip@eq:pair-no-slack\endcsname{Eq. (15.62), p. 89: e\textasciicircum{}tilde h\_0 symmetric of finite order<=>tilde h\_0==0<=>(D),}
\expandafter\gdef\csname CGTip@prop:dichotomy\endcsname{Proposition 15.34 (Symmetric joining factors in general) (p. 89): Let f be meromorphic in C with (s-1)f(s) entire of order at most one, and let f satisfy the functional equation (12.3). Let g=f e\textasciicircum{}h, where e\textasciicircum{}h is entire, without zeros, of finite order, and e\textasciicircum{}h(s)=e\textasciicircum{}h(1-s). If gnot ==0 and g(s)=sum\_n>=1a\_nn\textasciicircum{}-s converges absolutely in Re s>1, then g=C\_Hzeta with C\_H!=0 and e\textasciicircum{}h is con...}
\expandafter\gdef\csname CGTip@rem:dichotomy-counterexample\endcsname{Remark 15.35 (Order one cannot be dropped) (p. 89): For real nu\_5!=0 put f(s)=zeta(s)exp (-nu\_5(s-(1)/(2))\textasciicircum{}2) and e\textasciicircum{}h(s)=exp (nu\_5(s-(1)/(2))\textasciicircum{}2). Then f is meromorphic of order two with the functional equation, the pole and the zeros of zeta; e\textasciicircum{}h is entire, without zeros, of finite order and symmetric; g=f e\textasciicircum{}h=zeta is an absolutely convergent Dirichlet series in Re s...}
\expandafter\gdef\csname CGTip@subsec:halfline\endcsname{Section 15.12: The half-line form of Hamburger's theorem for a pair (p. 90)}
\expandafter\gdef\csname CGTip@thm:ladder\endcsname{Theorem 15.36 (The ladder) (p. 90): Let dcolon[1,inf)->R be measurable with int\_1\textasciicircum{}inf| d(x)| x\textasciicircum{}- e\_1-1 dx<inf for some e\_1>=0, and put h(sigma):=sigma int\_1\textasciicircum{}inf d(x) x\textasciicircum{}-sigma-1 dx(sigma> e\_1). For every integer N>=2, h(sigma)=O(sigma N\textasciicircum{}-sigma)(sigma ->+inf)<=> d(x)=0 for almost every xin[1,N).}
\expandafter\gdef\csname CGTip@eq:ladder-h\endcsname{Eq. (15.63), p. 90: h(sigma):=sigma int\_1\textasciicircum{}inf d(x) x\textasciicircum{}-sigma-1 dx(sigma> e\_1).}
\expandafter\gdef\csname CGTip@eq:ladder\endcsname{Eq. (15.64), p. 90: h(sigma)=O(sigma N\textasciicircum{}-sigma)(sigma ->+inf)<=> d(x)=0 for almost every xin[1,N).}
\expandafter\gdef\csname CGTip@eq:ladder-Y\endcsname{Eq. (15.65), p. 90: Y(s):=int\_0\textasciicircum{}log Nd(u) e\textasciicircum{}s(log N-u) du .}
\expandafter\gdef\csname CGTip@eq:ladder-Y-real\endcsname{Eq. (15.66), p. 90: Y(sigma)=(N\textasciicircum{}sigma h(sigma))/(sigma)-int\_log N\textasciicircum{}inf d(u) e\textasciicircum{}-sigma(u-log N) du ,}
\expandafter\gdef\csname CGTip@thm:halfline-hamburger\endcsname{Theorem 15.37 (Hamburger's theorem for a pair, half-line form) (p. 90): Let f satisfy (H0) and the functional equation (12.3), and suppose that for Re s>1 log f(s)=sint\_1\textasciicircum{}inf J\_f(x) x\textasciicircum{}-s-1 dx,int\_1\textasciicircum{}inf|J\_f(x)| x\textasciicircum{}-sigma-1 dx<inf (sigma>1), with J\_f real and measurable. Let g(s)=sum\_n>=1a\_nn\textasciicircum{}-s converge absolutely in Re s>1, with a\_1=1, and let J\_g(x):=sum\_2<= n<= xb\_n, where log g(s)=sum...}
\expandafter\gdef\csname CGTip@eq:halfline-f\endcsname{Eq. (15.67), p. 90: log f(s)=sint\_1\textasciicircum{}inf J\_f(x) x\textasciicircum{}-s-1 dx,int\_1\textasciicircum{}inf|J\_f(x)| x\textasciicircum{}-sigma-1 dx<inf (sigma>1),}
\expandafter\gdef\csname CGTip@cor:ladder-pair\endcsname{Corollary 15.38 (The ladder for the pair (zeta\_lambda,tilde zeta\_Z,lambda)) (p. 91): Let lambda in[0,1] and N>=2. Then tilde h\_lambda(sigma)=O(sigma N\textasciicircum{}-sigma) as sigma ->+inf if and only if the Mellin count J\textasciicircum{}rm M\_lambda of zeta\_lambda coincides almost everywhere on [1,N) with the step function sum\_n\textasciicircum{}m<= xtilde c\_n(lambda)/m; in particular J\textasciicircum{}rm M\_lambda then vanishes on (1,2) and is constant between...}
\expandafter\gdef\csname CGTip@eq:ladder-D\endcsname{Eq. (15.68), p. 91: (D)<=>tilde h\_0(sigma)=O\_N(sigma N\textasciicircum{}-sigma) for every N>=2 ,}
\expandafter\gdef\csname CGTip@subsec:bridge\endcsname{Section 15.13: The bridge theorem: the phase of the joining factor on the critical line (p. 91)}
\expandafter\gdef\csname CGTip@thm:bridge\endcsname{Theorem 15.39 (Bridge theorem, phase form) (p. 91): Let f be meromorphic in C, real on the real axis, with (H0), the functional equation (12.3), and lim\_sigma ->+inff(sigma)=varpi\_f>0. Let g(s)=sum\_n>=1a\_nn\textasciicircum{}-s converge absolutely in Re s>1, with real a\_n and a\_1=1. Suppose that g=f e\textasciicircum{}h in Re s>(1)/(2) with h holomorphic there and real on (1,inf), and that item[(B1)]...}
\expandafter\gdef\csname CGTip@eq:bridge-B2\endcsname{Eq. (15.69), p. 91: |h(s)|<= C (1+|s|)\textasciicircum{} e\_1(1+(sigma-(1)/(2))\textasciicircum{}- e\_2)(Re s>(1)/(2)).}
\expandafter\gdef\csname CGTip@eq:bridge-mean\endcsname{Eq. (15.70), p. 91: log|h(s)|<=(1)/(2pi)int\_0\textasciicircum{}2pilog |h(s+e\textasciicircum{}iphi)| dphi <=log C+ e\_1log(2+|s|)+log 2+(e\_2)/(2pi)int\_0\textasciicircum{}2pilog\textasciicircum{}+(1)/(|sigma-(1)/(2)+cos phi|) dphi .}
\expandafter\gdef\csname CGTip@cor:bridge-family\endcsname{Corollary 15.40 (The bridge theorem for the family) (p. 92): item For 0<lambda <=1 the joining factor tilde h\_lambda does not satisfy both (B1) and (B2). item (D) holds if and only if tilde h\_0 satisfies (B1) and (B2).}
\expandafter\gdef\csname CGTip@eq:phase-defect\endcsname{Eq. (15.71), p. 92: v\_lambda(t):=lim\_sigma ->(1)/(2)\textasciicircum{}+Im tilde h\_lambda(sigma+it)=lim\_sigma ->(1)/(2)\textasciicircum{}+(Im log tilde zeta\_Z,lambda(sigma+it)-Im log zeta\_lambda(sigma+it)).}
\expandafter\gdef\csname CGTip@prop:continuous-phase\endcsname{Proposition 15.41 (The phase of the continuous side) (p. 92): Let lambda in[0,1] and t>0 not an ordinate gamma\_k(lambda). With the branch of log zeta\_lambda that is real on (1,inf), continued along any path in the open quadrant Re s>(1)/(2), Im s>0 (the continuation is well defined there, since zeta\_lambda has no zeros or poles in the open quadrant: its zeros lie on the line R...}
\expandafter\gdef\csname CGTip@eq:continuous-phase\endcsname{Eq. (15.72), p. 92: lim\_sigma ->(1)/(2)\textasciicircum{}+Im log zeta\_lambda(sigma+it)=-vartheta(t)+pi (N\_lambda(t)-1), N\_lambda(t)=no.k:gamma\_k(lambda)<= t,}
\expandafter\gdef\csname CGTip@eq:generator-phase\endcsname{Eq. (15.73), p. 92: lim\_sigma ->(1)/(2)\textasciicircum{}+Im V\_lambda(sigma+it)=-pi sum\_k>=1eta\_k delta (t-gamma\_k(lambda)).}
\expandafter\gdef\csname CGTip@prop:one-ordinate\endcsname{Proposition 15.42 (The cell masses are sums of a one-ordinate function) (p. 93): For n>=2 there is a real function o\_n(kappa,kappa') of two real variables, depending only on n and the cell varphi\_n, such that for all lists kappa,kappa'in O tilde c\_n(kappa)-tilde c\_n(kappa')=sum\_k>=1 o\_n(kappa\_k,kappa'\_k), o\_n(kappa,kappa')=-2 Re (W\_n(kappa)-W\_n(kappa')) if n is not a power of 2, and, with L\_n(s)...}
\expandafter\gdef\csname CGTip@eq:one-ordinate\endcsname{Eq. (15.74), p. 93: tilde c\_n(kappa)-tilde c\_n(kappa')=sum\_k>=1 o\_n(kappa\_k,kappa'\_k), o\_n(kappa,kappa')=-2 Re (W\_n(kappa)-W\_n(kappa')) if n is not a power of 2,}
\expandafter\gdef\csname CGTip@eq:one-ordinate-log\endcsname{Eq. (15.75), p. 93: log tilde zeta\_Z,kappa(s)-log tilde zeta\_Z,kappa'(s)=sum\_n>=2(sum\_k>=1 o\_n(kappa\_k,kappa'\_k)) L\_n(s),}
\expandafter\gdef\csname CGTip@eq:one-ordinate-logzeta\endcsname{Eq. (15.76), p. 93: log zeta\_kappa(s)-log zeta\_kappa'(s)=sum\_k>=1log(1+q(s)/((1)/(4)+kappa\_k\textasciicircum{}2))/(1+q(s)/((1)/(4)+kappa'\_k\textasciicircum{}2)) .}
\expandafter\gdef\csname CGTip@eq:one-ordinate-joining\endcsname{Eq. (15.77), p. 93: j(kappa,kappa';s):=sum\_n>=2 o\_n(kappa,kappa') L\_n(s)-log(1+q(s)/((1)/(4)+kappa\textasciicircum{}2))/(1+q(s)/((1)/(4)+kappa'\textasciicircum{}2)), tilde h\_kappa-tilde h\_kappa'=sum\_k>=1 j(kappa\_k,kappa'\_k;s) formally.}
\expandafter\gdef\csname CGTip@heur:phase-rotation\endcsname{Heuristic 15.43 (What happens to the phase defect along the path) (p. 93): The following is a computation by stationary phase without remainder terms; it is not proved. For n much larger than kappa the cell resolves the oscillation, Re W\_n(kappa)\textasciitilde{} n\textasciicircum{}-1/2cos(kappa log n)/log n, and for n much smaller than kappa/2pi the kernel W\_n(kappa) is negligible because varphi is smooth (Proposition 15...}
\expandafter\gdef\csname CGTip@eq:phase-rotation\endcsname{Eq. (15.78), p. 93: Im j(kappa,* ;(1)/(2)+it)\textasciitilde{}amplitudex cos (2pi N\_RVM(kappa)+phase(t)),| j(kappa,* ;sigma+it)|\textasciitilde{}kappa\textasciicircum{}-1/2-sigma,}
\expandafter\gdef\csname CGTip@subsec:transport\endcsname{Section 15.14: Transport from the clock on the rungs (p. 93)}
\expandafter\gdef\csname CGTip@eq:transport-defs\endcsname{Eq. (15.79), p. 93: c\_ f(t):=int\_0\textasciicircum{}inf f(u)cos(tu) du, k(u):=(1)/(2)(e\textasciicircum{}u/2-(e\textasciicircum{}-u/2)/(e\textasciicircum{}2u-1)), k\_ f:=int\_0\textasciicircum{}inf f(u) k(u) du .}
\expandafter\gdef\csname CGTip@prop:rung-functional\endcsname{Proposition 15.44 (The rung functional) (p. 93): For every list kappain O (Section 16), R\_ f(kappa):=(1)/(2)int\_1\textasciicircum{}inf f(log x) x\textasciicircum{}-1/2log x dJ\textasciicircum{}rm M\_kappa(x)= k\_ f-sum\_k>=1 c\_ f(kappa\_k), the series converging absolutely. Along the family, with R\_ f(lambda):= R\_ f((gamma\_k(lambda))), R\_ f(lambda)= R\_ f(1)- S\_ f(lambda), S\_ f(lambda):=sum\_k>=1(c\_ f(gamma\_k(lambda))-...}
\expandafter\gdef\csname CGTip@eq:rung-functional\endcsname{Eq. (15.80), p. 93: R\_ f(kappa):=(1)/(2)int\_1\textasciicircum{}inf f(log x) x\textasciicircum{}-1/2log x dJ\textasciicircum{}rm M\_kappa(x)= k\_ f-sum\_k>=1 c\_ f(kappa\_k),}
\expandafter\gdef\csname CGTip@eq:transport\endcsname{Eq. (15.81), p. 94: R\_ f(lambda)= R\_ f(1)- S\_ f(lambda), S\_ f(lambda):=sum\_k>=1(c\_ f(gamma\_k(lambda))- c\_ f(tau\_k)), S\_ f(1)=0 .}
\expandafter\gdef\csname CGTip@eq:transport-D\endcsname{Eq. (15.82), p. 94: (D) => S\_ f(0)= R\_ f(1),}
\expandafter\gdef\csname CGTip@thm:transport-band\endcsname{Theorem 15.45 (The transport is band-limited in lambda) (p. 94): Let fin C\_c\textasciicircum{}inf(0,log N). The function lambda -> S\_ f(lambda) extends to an entire function with | S\_ f(lambda)|<= C\_ f(1+|lambda|log(2+|lambda|)) e\textasciicircum{}C\_eta log N |Im lambda|(lambda in C), and for every j>=1 d\_lambda\textasciicircum{}j S\_ f(1)= m\_j(f):=sum\_k>=1eta\_k\textasciicircum{} j c\_ f\textasciicircum{}(j)(tau\_k), S\_ f(0)=sum\_j>=1((-1)\textasciicircum{}j)/(j!) m\_j(f). Moreover |...}
\expandafter\gdef\csname CGTip@eq:transport-band\endcsname{Eq. (15.83), p. 94: | S\_ f(lambda)|<= C\_ f(1+|lambda|log(2+|lambda|)) e\textasciicircum{}C\_eta log N |Im lambda|(lambda in C),}
\expandafter\gdef\csname CGTip@eq:transport-moments\endcsname{Eq. (15.84), p. 94: d\_lambda\textasciicircum{}j S\_ f(1)= m\_j(f):=sum\_k>=1eta\_k\textasciicircum{} j c\_ f\textasciicircum{}(j)(tau\_k), S\_ f(0)=sum\_j>=1((-1)\textasciicircum{}j)/(j!) m\_j(f).}
\expandafter\gdef\csname CGTip@prop:clock-defect-poisson\endcsname{Proposition 15.46 (The defect of the clock on a rung) (p. 94): Let t\_0in(2pi,tau\_0) be the point with N\_RVM(t\_0)=0, and let fin C\_c\textasciicircum{}inf(0,inf). Then sum\_k>=1 c\_ f(tau\_k)= I\_0+2sum\_m>=1(-1)\textasciicircum{}m I\_m, I\_m:=int\_ t\_0\textasciicircum{}inf c\_ f(t)cos (2pi mN\_RVM(t))N\_RVM'(t) dt=O\_ f(m\textasciicircum{}-2), I\_0=-int\_0\textasciicircum{}inf(f(u))/(2pi u)(sin(u t\_0)log(t\_0)/(2pi)+(pi)/(2)-int\_0\textasciicircum{}u t\_0(sin y)/(y) dy)du, and therefore R\_ f(1)=...}
\expandafter\gdef\csname CGTip@eq:clock-defect-poisson\endcsname{Eq. (15.85), p. 94: sum\_k>=1 c\_ f(tau\_k)= I\_0+2sum\_m>=1(-1)\textasciicircum{}m I\_m, I\_m:=int\_ t\_0\textasciicircum{}inf c\_ f(t)cos (2pi mN\_RVM(t))N\_RVM'(t) dt=O\_ f(m\textasciicircum{}-2),}
\expandafter\gdef\csname CGTip@eq:clock-defect-I0\endcsname{Eq. (15.86), p. 94: I\_0=-int\_0\textasciicircum{}inf(f(u))/(2pi u)(sin(u t\_0)log(t\_0)/(2pi)+(pi)/(2)-int\_0\textasciicircum{}u t\_0(sin y)/(y) dy)du,}
\expandafter\gdef\csname CGTip@subsec:multiscale-jacobian\endcsname{Section 15.15: The multiscale Jacobian of the midpoint clock (p. 95)}
\expandafter\gdef\csname CGTip@eq:clock-jacobian\endcsname{Eq. (15.87), p. 95: (J\_T v)(u):=-usum\_tau\_k<= Tv\_ksin(tau\_k u).}
\expandafter\gdef\csname CGTip@thm:clock-jacobian\endcsname{Theorem 15.47 (Multiscale Jacobian of the midpoint clock) (p. 95): The following statements hold. item For every T>tau\_2, Delta\_T>= (2pi)/(log(T/2pi)). item Fix a>0 and varkappa\_rm win>0, and put [ L\_T:=(1+varkappa\_rm win)log(T/2pi), I\_T=[a,a+L\_T]. ] There is c\_a,varkappa\_rm win>0, independent of T, such that for all finitely supported v=(v\_k) on k:tau\_k<= T, || J\_T v||\_L\textasciicircum{}2(I\_T)\textasciicircum{}2...}
\expandafter\gdef\csname CGTip@eq:clock-gap-lower\endcsname{Eq. (15.88), p. 95: Delta\_T>= (2pi)/(log(T/2pi)).}
\expandafter\gdef\csname CGTip@eq:multiscale-frame\endcsname{Eq. (15.89), p. 95: || J\_T v||\_L\textasciicircum{}2(I\_T)\textasciicircum{}2 >= c\_a,varkappa\_rm winL\_Tsum\_tau\_k<= T|v\_k|\textasciicircum{}2.}
\expandafter\gdef\csname CGTip@eq:fixed-window-pair\endcsname{Eq. (15.90), p. 95: || J v\textasciicircum{}(k)||\_L\textasciicircum{}2(0,U) <= (U\textasciicircum{}5/2)/(sqrt(10))(tau\_k+1-tau\_k)->0.}
\expandafter\gdef\csname CGTip@eq:smooth-tail-clock\endcsname{Eq. (15.91), p. 96: sum\_tau\_k>T| c\_ f(kappa\_k)- c\_ f(tau\_k)| <<\_A, f,C\_eta T\textasciicircum{}1-Alog T.}
\expandafter\gdef\csname CGTip@eq:multiscale-residual\endcsname{Eq. (15.92), p. 97: varepsilon\_r(T,U):= inf\_h,b|| D\_RVM+usum\_tau\_k<= Th\_ksin(tau\_k u) -sum\_log n<= Ub\_ndelta(u-log n) ||\_H\textasciicircum{}-r(0,U).}
\expandafter\gdef\csname CGTip@sec:lists\endcsname{Section 16: The space of ordinate lists and the existence form of the final step (p. 97)}
\expandafter\gdef\csname CGTip@subsec:lists\endcsname{Section 16.1: Lists, their zeta functions and their atomizations (p. 97)}
\expandafter\gdef\csname CGTip@eq:list-space\endcsname{Eq. (16.1), p. 97: tau\_0\textasciicircum{}+<=kappa\_1<=kappa\_2<=...,|kappa\_k-tau\_k|<= C\_eta,|int\_tau\_0\textasciicircum{}X(N\_kappa(tau)-N\_RVM(tau))dtau |<= Clog X(X>=tau\_0),}
\expandafter\gdef\csname CGTip@eq:list-objects\endcsname{Eq. (16.2), p. 97: zeta\_kappa(s):=(1)/(G(s))prod\_k>=1(1+(q(s))/((1)/(4)+kappa\_k\textasciicircum{}2)), tilde c\_n(kappa), tilde a\_N(kappa), k\_N(kappa), tilde zeta\_Z,kappa=zeta\_kappa e\textasciicircum{}tilde h\_kappa, V\_sigma(kappa),}
\expandafter\gdef\csname CGTip@thm:transfer-O\endcsname{Theorem 16.1 (Transfer to O) (p. 97): Let kappain O. The following statements, proved above for the lists (gamma\_k(lambda))\_k>=1, hold for kappa with the same proofs, and their constants depend only on tau\_0\textasciicircum{}+, C\_eta and C in (16.1) and on the parameters of the statement (sigma, eps, n, N), never on the list: C=C(tau\_0\textasciicircum{}+,C\_eta,C;sigma,eps,...), hence th...}
\expandafter\gdef\csname CGTip@thm:rigidity\endcsname{Theorem 16.2 (Rigidity of the coincidence) (p. 98): The set C:=kappain O: tilde h\_kappa is constant of the lists whose continuous and atomized zeta functions coincide up to a constant factor has at most one element. If kappain C, then zeta\_kappa=tilde zeta\_Z,kappa=zeta, kappa is the list of the ordinates of zeta, and (D) holds. Conversely, (D) implies C=(gamma\_k).}
\expandafter\gdef\csname CGTip@eq:coincidence-set\endcsname{Eq. (16.3), p. 98: C:=kappain O: tilde h\_kappa is constant}
\expandafter\gdef\csname CGTip@prop:lists-compact\endcsname{Proposition 16.3 (Compactness and continuity) (p. 98): The space O and the functions (16.2) have the following properties. item O is compact in the product topology of prod\_k[tau\_k-C\_eta,tau\_k+C\_eta]. item For every n>=2, N>=1 and sigma>1, the functions kappa->tilde c\_n(kappa), tilde a\_N(kappa), k\_N(kappa) and V\_sigma(kappa) are continuous on O. item For every real sigm...}
\expandafter\gdef\csname CGTip@eq:lists-compact-dominant\endcsname{Eq. (16.4), p. 99: | F\_kappa(t)|<= H(t):=(C'\_sigma (1+log(2+|t|)))/(1+|t|), Hin L\textasciicircum{}2(R),}
\expandafter\gdef\csname CGTip@cor:existence-form\endcsname{Corollary 16.4 (Existence form of the final step) (p. 99): For every sigma>1 the minimum min\_ O V\_sigma is attained, and the following are equivalent: (i) (D); (ii) C!=empty; (iii) there is kappain O (Section 16) with tilde zeta\_Z,kappa=zeta; (iv) min\_ O V\_sigma=0.}
\expandafter\gdef\csname CGTip@rem:existence-input\endcsname{Remark 16.5 (What an existence proof uses) (p. 99): Compactness, continuity and convexity do not decide whether min\_ O V\_sigma is 0: they do not involve the form of the functional equation. The existence statement uses the functional equation with the factor Gamma((s)/(2))pi\textasciicircum{}-s/2, for which a function with both an absolutely convergent Dirichlet series and the functi...}
\expandafter\gdef\csname CGTip@eq:first-variation\endcsname{Eq. (16.5), p. 100: (d tilde c\_n(kappa))/(d kappa\_k)=2 Im V\_n(kappa\_k),}
\expandafter\gdef\csname CGTip@subsec:free-masses\endcsname{Section 16.2: The rung system with free masses (p. 100)}
\expandafter\gdef\csname CGTip@eq:rung-system\endcsname{Eq. (16.6), p. 100: sum\_k>=1cos(kappa\_ku)+sum\_n>=2 a\_n delta(u-log n)= k(u)in D'(0,inf).}
\expandafter\gdef\csname CGTip@thm:free-masses\endcsname{Theorem 16.6 (Existence form with free masses) (p. 100): (D) holds if and only if (16.6) has a solution (kappa,(a\_n)) with kappain O (Section 16). The solution is then unique: kappa=(gamma\_k), and a\_n=(log p)/(2sqrt(n)) if n=p\textasciicircum{}m is a prime power, a\_n=0 otherwise.}
\expandafter\gdef\csname CGTip@eq:free-masses\endcsname{Eq. (16.7), p. 100: a\_n=(log p)/(2sqrt(n)) if n=p\textasciicircum{}m is a prime power, a\_n=0 otherwise.}
\expandafter\gdef\csname CGTip@eq:free-masses-dJ\endcsname{Eq. (16.8), p. 100: (log x) dJ\textasciicircum{}rm M\_kappa(x)=sum\_n>=22sqrt(n) a\_n delta\_nin D'(1,inf).}
\expandafter\gdef\csname CGTip@eq:free-masses-factor\endcsname{Eq. (16.9), p. 100: zeta\_kappa(s)=exp (nu\_6+nu\_7 2\textasciicircum{}-s) tilde zeta\_Z,kappa(s)(Re s>1)}
\expandafter\gdef\csname CGTip@thm:nondegenerate\endcsname{Theorem 16.7 (Nondegeneracy of the solution) (p. 101): Let (kappa,(a\_n)) be a solution of (16.6) with kappain O (Section 16). Let (dot kappa\_k) be a bounded real sequence and (dot a\_n)\_n>=2 real numbers with dot a\_n=O\_eps(n\textasciicircum{}-1/2+eps) for every eps>0, such that sum\_k>=1dot kappa\_k usin(kappa\_ku)=sum\_n>=2dot a\_n delta(u-log n)in D'(0,inf), the derivative of (16.6) in the...}
\expandafter\gdef\csname CGTip@eq:rung-linearized\endcsname{Eq. (16.10), p. 101: sum\_k>=1dot kappa\_k usin(kappa\_ku)=sum\_n>=2dot a\_n delta(u-log n)in D'(0,inf),}
\expandafter\gdef\csname CGTip@eq:nondeg-H1\endcsname{Eq. (16.11), p. 101: H\_1(z):=sum\_k>=1dot kappa\_kint\_0\textasciicircum{}inf (1- x\_0(u)) usin(kappa\_ku) e\textasciicircum{}-zu du=(1)/(2)< J,(1- x\_0) e\textasciicircum{}-z|u|> .}
\expandafter\gdef\csname CGTip@eq:nondeg-split\endcsname{Eq. (16.12), p. 101: H\_1(z)=sum\_k>=1dot kappa\_k (2kappa\_kz)/((z\textasciicircum{}2+kappa\_k\textasciicircum{}2)\textasciicircum{}2)- y(z), y(z):=sum\_k>=1dot kappa\_kint\_0\textasciicircum{}log 2 x\_0(u) usin(kappa\_ku) e\textasciicircum{}-zu du .}
\expandafter\gdef\csname CGTip@eq:nondeg-H\endcsname{Eq. (16.13), p. 102: H(z)=nu\_8+sum\_k>=1dot kappa\_k((kappa\_k)/(z\textasciicircum{}2+kappa\_k\textasciicircum{}2)-(1)/(kappa\_k)),}
\expandafter\gdef\csname CGTip@eq:nondeg-order\endcsname{Eq. (16.14), p. 102: max\_|s-(1)/(2)|<= r\_j|(s-1) Z(s)|<= C\_ Hlog\textasciicircum{}2r\_j*max\_|s-(1)/(2)|=r\_j|(s-1)zeta(s)|<=exp (2r\_j\textasciicircum{}1+eps)}
\expandafter\gdef\csname CGTip@rem:nondegenerate-scope\endcsname{Remark 16.8 (Scope of Theorem 16.7) (p. 102): The theorem is conditional on the existence of a solution; it shows that a solution is isolated to first order, the Jacobian of (16.6) being injective there on the directions considered, except for the splitting of multiple ordinates. This is a necessary endpoint nondegeneracy condition for any continuation scheme;...}
\expandafter\gdef\csname CGTip@subsec:gradient\endcsname{Section 16.3: The gradient of the vertical defect (p. 103)}
\expandafter\gdef\csname CGTip@eq:K-kappa\endcsname{Eq. (16.15), p. 103: K\_kappa:=(tilde zeta\_Z,kappa)/(zeta)=sum\_n>=1k\_n(kappa) n\textasciicircum{}-s, V\_sigma(kappa)=< K\_kappa, K\_kappa>\_sigma-1, r\_m(kappa):=Re < K\_kappa L\_m, K\_kappa>\_sigma ,}
\expandafter\gdef\csname CGTip@prop:critical-points\endcsname{Proposition 16.9 (Critical points of the vertical defect) (p. 103): Let sigma>1 and kappain O (Section 16). Then | r\_m(kappa)|<=2m\textasciicircum{}-sigma(1+ V\_sigma(kappa)), | u\_m(t)|<= Cm\textasciicircum{}-1/2 off the powers of 2 and | u\_m(t)|<= C on them, and (d V\_sigma)/(d kappa\_k)(kappa)=2 Q\_kappa(kappa\_k), Q\_kappa(t):=sum\_m>=2 r\_m(kappa) u\_m(t), the series converging absolutely and uniformly in t. If V\_sigma(k...}
\expandafter\gdef\csname CGTip@eq:gradient\endcsname{Eq. (16.16), p. 103: (d V\_sigma)/(d kappa\_k)(kappa)=2 Q\_kappa(kappa\_k), Q\_kappa(t):=sum\_m>=2 r\_m(kappa) u\_m(t),}
\expandafter\gdef\csname CGTip@sec:models\endcsname{Section 17: Acceptance criteria and models in which the final step is a theorem (p. 104)}
\expandafter\gdef\csname CGTip@subsec:acceptance\endcsname{Section 17.1: The acceptance criterion: window and frequencies (p. 104)}
\expandafter\gdef\csname CGTip@eq:model-product\endcsname{Eq. (17.1), p. 104: C\_kappa(z):= zprod\_k>=1(1-(z\textasciicircum{}2)/(kappa\_k\textasciicircum{}2)),}
\expandafter\gdef\csname CGTip@eq:model-logder\endcsname{Eq. (17.2), p. 104: (C\_kappa')/(C\_kappa)(z)=(1)/(z)+sum\_k>=1(2 z)/(z\textasciicircum{}2-kappa\_k\textasciicircum{}2)=-iint\_0\textasciicircum{}inf e\textasciicircum{}i zu(1+2sum\_k>=1cos(kappa\_ku))du ,}
\expandafter\gdef\csname CGTip@thm:acceptance\endcsname{Theorem 17.1 (Acceptance criterion) (p. 104): For a list kappa as above the following are equivalent: item the spectral density is atomic on the frequencies: there are real w\_ l with sum\_ lin F| w\_ l|e\textasciicircum{}-nu\_10 l<inf for some admissible nu\_10 and 1+2sum\_k>=1cos(kappa\_ku)=sum\_ lin F w\_ l delta(u- l)in D'(0,inf); item the continuous function has an Euler form: for...}
\expandafter\gdef\csname CGTip@eq:acceptance-density\endcsname{Eq. (17.3), p. 104: 1+2sum\_k>=1cos(kappa\_ku)=sum\_ lin F w\_ l delta(u- l)in D'(0,inf);}
\expandafter\gdef\csname CGTip@eq:acceptance-euler\endcsname{Eq. (17.4), p. 104: e\textasciicircum{}i k z C\_kappa(z)= c\_0exp (sum\_ lin F b\_ l e\textasciicircum{}i l z), b\_ l=-(w\_ l)/(l);}
\expandafter\gdef\csname CGTip@eq:acceptance-window\endcsname{Eq. (17.5), p. 104: e\textasciicircum{}i k z C\_kappa(z)=sum\_ lin F\textasciicircum{}+n[0,2 k] c\_ l e\textasciicircum{}i l z, c\_0!=0 .}
\expandafter\gdef\csname CGTip@subsec:graph-model\endcsname{Section 17.2: A non-periodic accepted list: one unitary matrix (p. 105)}
\expandafter\gdef\csname CGTip@eq:graph-matrices\endcsname{Eq. (17.6), p. 105: U:= v sqrt(1- v\textasciicircum{}2) sqrt(1- v\textasciicircum{}2) - v}
\expandafter\gdef\csname CGTip@eq:graph-det\endcsname{Eq. (17.7), p. 105: D(z):=det (I- U X(z))=1- v e\textasciicircum{}2i p z+ v e\textasciicircum{}2i q z-e\textasciicircum{}2i(p+ q) z=-2i e\textasciicircum{}i(p+ q) z S(z),}
\expandafter\gdef\csname CGTip@eq:graph-sine\endcsname{Eq. (17.8), p. 105: S(z):=sin ((p+ q) z)+ vsin ((p- q) z).}
\expandafter\gdef\csname CGTip@prop:graph-model\endcsname{Proposition 17.2 (The two-edge model) (p. 105): item All zeros of S are real. Each interval between consecutive points (m+(1)/(2))pi/ k, min Z, contains a zero, and the number of zeros in [0,T], with multiplicity, is kT/pi+O(1); hence the positive zeros form a list kappa with |kappa\_k-kpi/ k| bounded. The set of zeros is not periodic. item S(- z)=- S(z), and D(z)...}
\expandafter\gdef\csname CGTip@eq:graph-euler\endcsname{Eq. (17.9), p. 105: log D(z)=-sum\_r>=1(1)/(r)tr(U X(z))\textasciicircum{}r=sum\_ lin F b\_ l e\textasciicircum{}i l z, F=2(m p+m' q):m,m'>=00,}
\expandafter\gdef\csname CGTip@eq:graph-eigenphase\endcsname{Eq. (17.10), p. 106: theta\_j'(t)=2< Dv\_j,v\_j>=2(p|v\_j,1|\textasciicircum{}2+ q|v\_j,2|\textasciicircum{}2) >= 2min(p, q)>0 .}
\expandafter\gdef\csname CGTip@subsec:finite-euler-model\endcsname{Section 17.3: A finite Euler product: the constructions of this paper with the final step proved (p. 106)}
\expandafter\gdef\csname CGTip@eq:finite-euler\endcsname{Eq. (17.11), p. 106: F(s):=(1-2\textasciicircum{}-s)(1-3\textasciicircum{}-s)=1-2\textasciicircum{}-s-3\textasciicircum{}-s+6\textasciicircum{}-s=sum\_n|6mu(n) n\textasciicircum{}-s.}
\expandafter\gdef\csname CGTip@eq:finite-euler-zeros\endcsname{Eq. (17.12), p. 106: y\_k\_k>=1=(2pi m)/(log 2)\_m>=1u (2pi m)/(log 3)\_m>=1in increasing order,}
\expandafter\gdef\csname CGTip@eq:finite-euler-hadamard\endcsname{Eq. (17.13), p. 107: F(s)=6\textasciicircum{}-s/2 (log 2log 3) s\textasciicircum{}2prod\_k>=1(1+(s\textasciicircum{}2)/(y\_k\textasciicircum{}2)).}
\expandafter\gdef\csname CGTip@eq:finite-euler-list\endcsname{Eq. (17.14), p. 107: F\_kappa(s):=6\textasciicircum{}-s/2 (log 2log 3) s\textasciicircum{}2prod\_k>=1(1+(s\textasciicircum{}2)/(kappa\_k\textasciicircum{}2)), F\_kappa(-s)=6\textasciicircum{}s F\_kappa(s),}
\expandafter\gdef\csname CGTip@eq:finite-euler-clock\endcsname{Eq. (17.15), p. 107: F\_1(s)=(log 2log 3)/(2)*(s\textasciicircum{}2 (1+6\textasciicircum{}-s))/(1+s\textasciicircum{}2log\textasciicircum{}26/pi\textasciicircum{}2).}
\expandafter\gdef\csname CGTip@thm:finite-euler\endcsname{Theorem 17.3 (The final step in the model) (p. 107): item (Poisson, prime by prime.) In D'(R), 2+2sum\_k>=1cos(y\_ku)=sum\_pin 2,3 sum\_min Zcos(2pi mu)/(log p)=sum\_pin 2,3log psum\_min Zdelta(u-mlog p). The spectral density of the list (y\_k) is atomic on the points mlog p: the continuous function built from the list is the Dirichlet series with Euler product, F\_0= F. item...}
\expandafter\gdef\csname CGTip@eq:finite-euler-poisson\endcsname{Eq. (17.16), p. 107: 2+2sum\_k>=1cos(y\_ku)=sum\_pin 2,3 sum\_min Zcos(2pi mu)/(log p)=sum\_pin 2,3log psum\_min Zdelta(u-mlog p).}
\expandafter\gdef\csname CGTip@eq:finite-euler-coefficients\endcsname{Eq. (17.17), p. 107: (d)/(dsigma) A\_lambda,T(sigma)=(-i)/(2T)[f(sigma+it)e\textasciicircum{}ilambda t]\_t=-T\textasciicircum{}t=T-lambda A\_lambda,T(sigma), so |(d)/(dsigma)(e\textasciicircum{}lambda sigma A\_lambda,T(sigma))|<=(M)/(T)e\textasciicircum{}lambda sigma,}
\expandafter\gdef\csname CGTip@eq:finite-euler-spectrum\endcsname{Eq. (17.18), p. 107: e\textasciicircum{}lambda sigma\_1 A\_lambda(sigma\_1)=e\textasciicircum{}lambda sigma\_2 A\_lambda(sigma\_2)(lambda in R).}
\expandafter\gdef\csname CGTip@eq:finite-euler-superposition\endcsname{Eq. (17.19), p. 108: y\_k=(2pi)/(log 2) N u (2pi)/(log 3) N ,}
\expandafter\gdef\csname CGTip@subsec:three-abscissas\endcsname{Section 17.4: The diagonal operator over all primes and the three abscissas (p. 108)}
\expandafter\gdef\csname CGTip@eq:prime-det\endcsname{Eq. (17.20), p. 108: det (I- X\_ P(s))=prod\_p(1-p\textasciicircum{}-s)=(1)/(zeta(s))(Re s>1).}
\expandafter\gdef\csname CGTip@eq:prime-det2\endcsname{Eq. (17.21), p. 108: D\_2(s):=det \_2(I- X\_ P(s))=prod\_p(1-p\textasciicircum{}-s)e\textasciicircum{}p\textasciicircum{}-s=exp (-sum\_m>=2(1)/(m) P(ms))(Re s>(1)/(2)),}
\expandafter\gdef\csname CGTip@prop:no-finite-rank\endcsname{Proposition 17.4 (No common object of finite rank in the window of zeta) (p. 108): Let X>=2. There is no function fnot ==0 that satisfies (H0) and the functional equation (12.3) and is an absolutely convergent Dirichlet series in Re s>1 supported on the integers all of whose prime factors are at most X.}
\expandafter\gdef\csname CGTip@eq:three-abscissas\endcsname{Eq. (17.22), p. 109: contraction: sigma>0;convergence of the determinant: sigma>1;axis of symmetry: sigma=(1)/(2) .}
\expandafter\gdef\csname CGTip@sec:routeB-finite\endcsname{Section 18: Route B for finite sections: compactness, uniqueness, certificates, and a Kadec--Avdonin criterion with blocks of variable length (p. 109)}
\expandafter\gdef\csname CGTip@subsec:rb-intro\endcsname{Section 18.1: Introduction to the section (p. 109)}
\expandafter\gdef\csname CGTip@subsec:rb-setting\endcsname{Section 18.2: Setting (p. 110)}
\expandafter\gdef\csname CGTip@eq:rb-NRVM\endcsname{Eq. (18.1), p. 110: N\_RVM(t):=(t)/(2pi)log(t)/(2pi e)+(7)/(8), N\_RVM'(t)=(1)/(2pi)log(t)/(2pi)(t>0).}
\expandafter\gdef\csname CGTip@eq:rb-Ntau-floor\endcsname{Eq. (18.2), p. 110: n\_RVM(t)=floor(N\_RVM(t)+(1)/(2)), |n\_RVM(t)-N\_RVM(t)|<=(1)/(2)(t>=2pi).}
\expandafter\gdef\csname CGTip@def:rb-lists\endcsname{Definition 18.1 (Admissible lists) (p. 110): For C>0, a list kappa=(kappa\_k)\_k>=1 is C-admissible if it is real and nondecreasing, kappa\_1>=2pi, and |kappa\_k-tau\_k|<= C for all k. Put N\_kappa(t):=no.k:kappa\_k<= t (with multiplicity) and P\_kappa(s):=prod\_k>=1(1+(q(s))/((1)/(4)+kappa\_k\textasciicircum{}2)), zeta\_kappa:=(P\_kappa)/(G). The clock list tau=(tau\_k) gives zeta\_tau=zeta\_...}
\expandafter\gdef\csname CGTip@eq:rb-Pk\endcsname{Eq. (18.3), p. 110: P\_kappa(s):=prod\_k>=1(1+(q(s))/((1)/(4)+kappa\_k\textasciicircum{}2)), zeta\_kappa:=(P\_kappa)/(G).}
\expandafter\gdef\csname CGTip@eq:rb-rung\endcsname{Eq. (18.4), p. 110: sum\_k>=1cos(kappa\_ku)+sum\_n>=2 a\_n delta(u-log n)= k(u)in D'(0,inf),}
\expandafter\gdef\csname CGTip@eq:rb-residual\endcsname{Eq. (18.5), p. 110: < R(kappa, a),varphi >:=sum\_k>=1 c\_varphi(kappa\_k)+sum\_n>=2 a\_nvarphi(log n)-< k,varphi >, c\_varphi(t):=int\_0\textasciicircum{}inf varphi(u)cos(tu) du .}
\expandafter\gdef\csname CGTip@subsec:rb-routeB\endcsname{Section 18.3: Route B: compactness, selection and the canonical box (p. 110)}
\expandafter\gdef\csname CGTip@eq:rb-rho\endcsname{Eq. (18.6), p. 110: s\_m(kappa, a):=max\_varphi in T\_m|< R(kappa, a),varphi >| .}
\expandafter\gdef\csname CGTip@def:rb-B\endcsname{Definition 18.2 (Statement B) (p. 110): Statement B(C, M) says that there are finite sections (kappa\textasciicircum{}(m), a\textasciicircum{}(m))in K(C, M) with s\_m(kappa\textasciicircum{}(m), a\textasciicircum{}(m))->0.}
\expandafter\gdef\csname CGTip@lem:rb-canonical-box\endcsname{Lemma 18.3 (Canonical box of the exact solution) (p. 111): Let M\_*:=max\_n>=2Lambda(n)/(2sqrt(n)). Then M\_*=log 7/(2sqrt(7))=0.367742...<1/e. If condition (D) holds, the pair (gamma, a\textasciicircum{}zeta), with a\textasciicircum{}zeta\_n:=Lambda(n)/(2sqrt(n)), solves (18.4) and lies in K(C\_eta, M\_*), where C\_eta is the constant of 2.2.}
\expandafter\gdef\csname CGTip@cor:rb-D-implies-B\endcsname{Corollary 18.4 ((D) implies B) (p. 111): If (D) holds, then B(C, M) holds for every C>= C\_eta and M>= M\_*.}
\expandafter\gdef\csname CGTip@prop:rb-compactness\endcsname{Proposition 18.5 (Compactness and passage to the limit) (p. 111): Let (kappa\textasciicircum{}(m), a\textasciicircum{}(m))in K(C, M) with s\_m(kappa\textasciicircum{}(m), a\textasciicircum{}(m))->0. Then a subsequence converges coordinatewise to some (kappa, a)in K(C, M), and every such limit satisfies R(kappa, a)=0 in D'(0,inf).}
\expandafter\gdef\csname CGTip@lem:rb-mass-ms\endcsname{Lemma 18.6 (Mean square of the masses of an exact solution) (p. 111): Let kappa be C-admissible for some C>0, and let a=(a\_n)\_n>=2 be real numbers, with no bound assumed, such that R(kappa, a)=0 in D'(0,inf). Then for every epsilon>0 there is C\_epsilon, depending only on epsilon and C, such that sum\_X<n<=2X a\_n\textasciicircum{}2<= C\_epsilon X\textasciicircum{}epsilon(X>=2).}
\expandafter\gdef\csname CGTip@eq:rb-mass-ms\endcsname{Eq. (18.7), p. 111: sum\_X<n<=2X a\_n\textasciicircum{}2<= C\_epsilon X\textasciicircum{}epsilon(X>=2).}
\expandafter\gdef\csname CGTip@eq:rb-mass-isolated\endcsname{Eq. (18.8), p. 112: a\_n=< k,varphi\_n>-Re S(log n)- E\_n, y\_k:=(hat psi(kappa\_k/X))/(X),; S(u):=sum\_kappa\_k<= T y\_ke\textasciicircum{}ikappa\_ku, E\_n:=sum\_kappa\_k>T y\_kcos(kappa\_klog n).}
\expandafter\gdef\csname CGTip@thm:rb-uniqueness\endcsname{Theorem 18.7 (Uniqueness in the tube, with free masses) (p. 112): Let kappa be C-admissible for some C>0, and let a=(a\_n)\_n>=2 be real numbers, with no bound assumed. If R(kappa, a)=0 in D'(0,inf), then all nontrivial zeros of zeta lie on Re s=(1)/(2), kappa=gamma, a= a\textasciicircum{}zeta, and (D) holds.}
\expandafter\gdef\csname CGTip@eq:rb-uniq-factor\endcsname{Eq. (18.9), p. 112: zeta\_kappa(s)=(P\_kappa(s))/(G(s))=C\_1exp (sum\_n>=2 c\_nn\textasciicircum{}-s), c\_n:=(2sqrt(n) a\_n)/(log n)(Re s>(3)/(2)),}
\expandafter\gdef\csname CGTip@rem:rb-uniqueness-input\endcsname{Remark 18.8 (What Theorem 18.7 uses) (p. 113): Theorem 18.7 contains the uniqueness half of Theorem 16.6, for every list at bounded distance from the clock and with no assumption on the masses: neither the third condition of (16.1) nor the threshold kappa\_1>=tau\_0\textasciicircum{}+ is used. In the proof of Theorem 16.6 the third condition enters through Theorem 16.1(ii), (iii)...}
\expandafter\gdef\csname CGTip@cor:rb-U\endcsname{Corollary 18.9 (Uniqueness in the box) (p. 113): For all C, M>0, the only possible zero of R in K(C, M) is (gamma, a\textasciicircum{}zeta), and R has a zero in K(C, M) only if the Riemann hypothesis holds.}
\expandafter\gdef\csname CGTip@cor:rb-selection\endcsname{Corollary 18.10 (Selection of the zeta list) (p. 113): Assume B(C, M), and let (kappa\textasciicircum{}(m), a\textasciicircum{}(m)) be any sequence in K(C, M) with s\_m->0. Then the whole sequence converges coordinatewise to (gamma, a\textasciicircum{}zeta), and (D) and the Riemann hypothesis hold.}
\expandafter\gdef\csname CGTip@cor:rb-equivalence\endcsname{Corollary 18.11 (Route B is equivalent to (D)) (p. 114): (D) holds if and only if B(C, M) holds for some C, M>0; and then B(C, M) holds for all C>= C\_eta and M>= M\_*.}
\expandafter\gdef\csname CGTip@rem:rb-B-status\endcsname{Remark 18.12 (Logical status of route B) (p. 114): By Corollary 18.11, route B with free masses is an equivalence, B(C, M)<=>(D) for C>= C\_eta and M>= M\_*, with a single open piece, the existence statement B, which is therefore not weaker than (D). Compactness is not an issue in the box: Proposition 18.5 uses only the imposed bounds; and uniqueness holds in every tu...}
\expandafter\gdef\csname CGTip@subsec:rb-prescribed\endcsname{Section 18.3.1: Route B with prescribed von Mangoldt masses (p. 114)}
\expandafter\gdef\csname CGTip@eq:rb-laplace-def\endcsname{Eq. (18.10), p. 114: L T(s):=< T, j(u) e\textasciicircum{}-zu>,}
\expandafter\gdef\csname CGTip@lem:rb-archimedean\endcsname{Lemma 18.13 (The Archimedean distribution) (p. 114): Let g(u):=(e\textasciicircum{}-u/2)/(e\textasciicircum{}2u-1)-(1)/(2u) for u>0; g is bounded, g(0\textasciicircum{}+)=-(3)/(4), and g(u)=-(1)/(2u)+O(e\textasciicircum{}-5u/2) as u->inf, in particular bounded. Put c\_ K:=-(1)/(4)log(2pi e\textasciicircum{}gamma\_E), gamma\_E Euler's constant, and K:=(1)/(2)e\textasciicircum{}u/2 1\_+-(1)/(2) g 1\_+-(1)/(4) fp (1\_+)/(u)+ c\_ K delta\_0 . Then K is a distribution of order 1 s...}
\expandafter\gdef\csname CGTip@eq:rb-K-dist\endcsname{Eq. (18.11), p. 114: K:=(1)/(2)e\textasciicircum{}u/2 1\_+-(1)/(2) g 1\_+-(1)/(4) fp (1\_+)/(u)+ c\_ K delta\_0 .}
\expandafter\gdef\csname CGTip@eq:rb-LK\endcsname{Eq. (18.12), p. 114: L K(s)=(1)/(2) (G')/(G)(s)(Re s>1).}
\expandafter\gdef\csname CGTip@eq:rb-aux-12\endcsname{Eq. (18.13), p. 115: (1)/(2)(Gamma')/(Gamma)((s)/(2))=int\_0\textasciicircum{}inf ((e\textasciicircum{}-2u)/(2u)-(e\textasciicircum{}-su)/(1-e\textasciicircum{}-2u))du =int\_0\textasciicircum{}inf ((e\textasciicircum{}-2u)/(2u)-e\textasciicircum{}-zu(e\textasciicircum{}-u/2)/(e\textasciicircum{}2u-1))du-(1)/(s),}
\expandafter\gdef\csname CGTip@eq:rb-aux-13\endcsname{Eq. (18.14), p. 115: (1)/(2)(G')/(G)(s)=(1)/(2s)+(1)/(2(s-1))-(1)/(4)log pi+(1)/(4)(Gamma')/(Gamma)((s)/(2)) = L((1)/(2)e\textasciicircum{}u/2 1\_+)(s)-(1)/(2) L(g 1\_+)(s)+(1)/(4)log(z)/(2)-(1)/(4)log pi .}
\expandafter\gdef\csname CGTip@lem:rb-laplace\endcsname{Lemma 18.14 (Laplace transform of the residual) (p. 115): Let kappa be C-admissible and put H\_kappa(u):=sum\_kkappa\_k\textasciicircum{}-2cos(kappa\_ku), which converges absolutely and uniformly on R since H\_kappa(0)=sum\_kkappa\_k\textasciicircum{}-2<inf, and S\_kappa:=-d\_u\textasciicircum{}2(H\_kappa 1\_+)+ H\_kappa(0) delta\_0', A\_zeta:=sum\_n>=2 a\textasciicircum{}zeta\_n delta\_log n, E\_kappa:= S\_kappa+ A\_zeta- K . Then E\_kappa is a distribution o...}
\expandafter\gdef\csname CGTip@eq:rb-S-A-dist\endcsname{Eq. (18.15), p. 115: S\_kappa:=-d\_u\textasciicircum{}2(H\_kappa 1\_+)+ H\_kappa(0) delta\_0', A\_zeta:=sum\_n>=2 a\textasciicircum{}zeta\_n delta\_log n, E\_kappa:= S\_kappa+ A\_zeta- K .}
\expandafter\gdef\csname CGTip@eq:rb-LR\endcsname{Eq. (18.16), p. 115: L E\_kappa(s)=hat R(s)=(1)/(2)((P\_kappa')/(P\_kappa)-(xi')/(xi))(s),xi:=(1)/(2)Gzeta ,}
\expandafter\gdef\csname CGTip@eq:rb-aux-14\endcsname{Eq. (18.17), p. 115: L S\_kappa(s)=-z\textasciicircum{}2sum\_k(z)/(kappa\_k\textasciicircum{}2(z\textasciicircum{}2+kappa\_k\textasciicircum{}2))+zsum\_k(1)/(kappa\_k\textasciicircum{}2) =sum\_k(z)/(z\textasciicircum{}2+kappa\_k\textasciicircum{}2)=(1)/(2)(P\_kappa')/(P\_kappa)(s),}
\expandafter\gdef\csname CGTip@rem:rb-boundary-term\endcsname{Remark 18.15 (The boundary term) (p. 116): Without the correction H\_kappa(0)delta\_0' in (18.15), the transform of -d\_u\textasciicircum{}2(H\_kappa 1\_+) is (1)/(2)P'\_kappa/P\_kappa-(s-(1)/(2)) H\_kappa(0): by (18.17), L(-d\_u\textasciicircum{}2(H\_kappa 1\_+))=-z\textasciicircum{}2 L(H\_kappa 1\_+), which differs from (1)/(2)P'\_kappa/P\_kappa by -z H\_kappa(0)=- L(H\_kappa(0)delta\_0'), the transform of the boundary term...}
\expandafter\gdef\csname CGTip@prop:rb-prescribed-uniqueness\endcsname{Proposition 18.16 (Uniqueness with prescribed masses) (p. 116): Let kappa be C-admissible for some C, and suppose R(kappa, a\textasciicircum{}zeta)=0 in D'(0,inf). Then all nontrivial zeros of zeta lie on Re s=(1)/(2), and kappa=gamma.}
\expandafter\gdef\csname CGTip@def:rb-Bprime\endcsname{Definition 18.17 (Statement B') (p. 116): Statement B'(C) says that there are C-admissible lists kappa\textasciicircum{}(m), each equal to the clock beyond some finite index, with s\_m(kappa\textasciicircum{}(m), a\textasciicircum{}zeta)->0.}
\expandafter\gdef\csname CGTip@cor:rb-Bprime\endcsname{Corollary 18.18 (Route B') (p. 116): item If B'(C) holds for some C, then every sequence as in Definition 18.17 converges coordinatewise to gamma, and all nontrivial zeros of zeta lie on Re s=(1)/(2). item If all nontrivial zeros of zeta lie on Re s=(1)/(2), then B'(C) holds for every C>= C\_eta. Hence the Riemann hypothesis is equivalent to B'(C) for s...}
\expandafter\gdef\csname CGTip@rem:rb-prescribed\endcsname{Remark 18.19 (What prescribing the masses changes, and what it does not) (p. 116): With prescribed masses, uniqueness for every list at bounded distance from the clock follows from Hadamard's theorem and the functional equation; with free masses it needs in addition Lemma 18.6 and Hamburger's theorem (Theorem 18.7). Routes B and B' have the same single open piece, existence, and each is equivalent...}
\expandafter\gdef\csname CGTip@subsec:rb-nonfinitary\endcsname{Section 18.3.2: Finite stages do not decide route B (p. 117)}
\expandafter\gdef\csname CGTip@eq:rb-Wr\endcsname{Eq. (18.18), p. 117: W\_r(f):=(1+u\textasciicircum{}*2)cosh(u\textasciicircum{}*)/(2)sum\_j=0\textasciicircum{}rbinom(r,j)|| f\textasciicircum{}(r-j)||\_L\textasciicircum{}1.}
\expandafter\gdef\csname CGTip@prop:rb-nonfinitary\endcsname{Proposition 18.20 (The zeta list solves every finite stage up to the verified height) (p. 117): Let H\_0:=3*10\textasciicircum{}12. For every r>=2 and every real f as above, |< R(gamma, a\textasciicircum{}zeta), f>|<=(1)/(4) W\_r(f)sum\_k: d\_k!=0gamma\_k\textasciicircum{}-r<=(r)/(4(r-1)) W\_r(f) H\_0\textasciicircum{}1-rlog H\_0 . The list gamma is C\_eta-admissible. Consequently, for every finite set T of such test functions (with the same r) and every delta>0, there is a finite sect...}
\expandafter\gdef\csname CGTip@eq:rb-nonfinitary\endcsname{Eq. (18.19), p. 117: |< R(gamma, a\textasciicircum{}zeta), f>|<=(1)/(4) W\_r(f)sum\_k: d\_k!=0gamma\_k\textasciicircum{}-r<=(r)/(4(r-1)) W\_r(f) H\_0\textasciicircum{}1-rlog H\_0 .}
\expandafter\gdef\csname CGTip@eq:rb-nonfinitary-section\endcsname{Eq. (18.20), p. 117: max\_ fin T|< R(kappa, a), f>|<=max\_ fin T(r)/(4(r-1)) W\_r(f) H\_0\textasciicircum{}1-rlog H\_0+delta ,}
\expandafter\gdef\csname CGTip@rem:rb-nonfinitary\endcsname{Remark 18.21 (Reading of Proposition 18.20) (p. 118): If the Riemann hypothesis holds, the residual in (18.19) vanishes. If it fails, B and B' fail (Corollaries 18.11 and 18.18), and still every finite stage has sections in the box with residual at most (r)/(4(r-1))max\_ T W\_r H\_0\textasciicircum{}1-rlog H\_0+delta. No computation on finitely many tests at a precision coarser than this b...}
\expandafter\gdef\csname CGTip@subsec:rb-staged\endcsname{Section 18.3.3: Staged construction: a certificate theorem and a graded extension (p. 118)}
\expandafter\gdef\csname CGTip@def:rb-staged\endcsname{Definition 18.22 (Staged sections) (p. 118): Fix u\textasciicircum{}*\_1<u\textasciicircum{}*\_2<...->inf, finite test sets T\_1in T\_2in ... supported in (0,u\textasciicircum{}*\_m) with dense union as in Section 18.3, t\_m:=| T\_m|, and k\textasciicircum{}*\_m->inf. Stage m starts from a list kappa\textasciicircum{}(m,0) that equals the clock beyond some finite index, moves the coordinates k<= k\textasciicircum{}*\_m by xin R\textasciicircum{}k\textasciicircum{}*\_m, and has masses a\textasciicircum{}zeta. Its residua...}
\expandafter\gdef\csname CGTip@eq:rb-Fm\endcsname{Eq. (18.21), p. 118: F\_m: R\textasciicircum{}k\textasciicircum{}*\_m-> R\textasciicircum{} t\_m, F\_m(x)\_varphi:=< R(kappa\textasciicircum{}(m,0)+x, a\textasciicircum{}zeta),varphi >(varphi in T\_m),}
\expandafter\gdef\csname CGTip@eq:rb-jacobian\endcsname{Eq. (18.22), p. 118: (J\_m)\_varphi,k=d\_x\_k F\_m(0)\_varphi=-int\_0\textasciicircum{}inf varphi(u) usin(kappa\textasciicircum{}(m,0)\_ku) du(varphi in T\_m, k<= k\textasciicircum{}*\_m).}
\expandafter\gdef\csname CGTip@eq:rb-eta-H\endcsname{Eq. (18.23), p. 118: h\_m:=|| F\_m(0)||\_inf, H\_m:=k\textasciicircum{}*\_mmax\_varphi in T\_mint\_0\textasciicircum{}inf|varphi(u)| u\textasciicircum{}2 du .}
\expandafter\gdef\csname CGTip@lem:rb-matching\endcsname{Lemma 18.23 (Monotone matching) (p. 118): Let (a\_k)\_k>=1 be real with | a\_k-tau\_k|<= C for all k, and let (a\_(k)) be its nondecreasing rearrangement. Then | a\_(k)-tau\_k|<= C for all k. If moreover a\_k=tau\_k for k>k\textasciicircum{}*, then a\_(k)=tau\_k for all k>k\textasciicircum{}*+j, where j:=no.k>k\textasciicircum{}*:tau\_k<=tau\_k\textasciicircum{}*+C<inf.}
\expandafter\gdef\csname CGTip@lem:rb-lipschitz\endcsname{Lemma 18.24 (Lipschitz bound for the Jacobian) (p. 119): For all x,x'in R\textasciicircum{}k\textasciicircum{}*\_m, ||D F\_m(x)-D F\_m(x')||\_inf ->inf<= H\_m||x-x'||\_inf, and || F\_m(x)- F\_m(0)- J\_mx||\_inf <=(1)/(2) H\_m||x||\_inf\textasciicircum{}2 .}
\expandafter\gdef\csname CGTip@eq:rb-taylor\endcsname{Eq. (18.24), p. 119: || F\_m(x)- F\_m(0)- J\_mx||\_inf <=(1)/(2) H\_m||x||\_inf\textasciicircum{}2 .}
\expandafter\gdef\csname CGTip@thm:rb-certificate\endcsname{Theorem 18.25 (Certificate for one stage) (p. 119): Suppose J\_m has a right inverse Q\_m: R\textasciicircum{} t\_m-> R\textasciicircum{}k\textasciicircum{}*\_m, J\_m Q\_m=I\_ t\_m, with || Q\_m||\_inf ->inf<= L\_m, and l\_m:=2 L\_m\textasciicircum{}2 H\_m h\_m<1 .A\_m Then the map T\_m(y):=y- F\_m(Q\_my) is a contraction of the ball ||y||\_inf <=2 h\_min R\textasciicircum{} t\_m with constant l\_m, and x\_*:= Q\_my\_*, y\_* its unique fixed point in that ball, satisfies F\_m(x...}
\expandafter\gdef\csname CGTip@eq:rb-Am\endcsname{Eq. (A_m), p. 119: l\_m:=2 L\_m\textasciicircum{}2 H\_m h\_m<1 .A\_m}
\expandafter\gdef\csname CGTip@eq:rb-cert-conclusion\endcsname{Eq. (18.25), p. 119: F\_m(x\_*)=0,||x\_*||\_inf <=2 L\_m h\_m .}
\expandafter\gdef\csname CGTip@rem:rb-certificate-errors\endcsname{Remark 18.26 (Evaluation errors and order of the ordinates) (p. 119): Theorem 18.25 is stated for exact values of h\_m, L\_m and H\_m. If the residual is only known up to an error, ||tilde F\_m(0)- F\_m(0)||\_inf <=epsilon\_m, the theorem applies with h\_m replaced by ||tilde F\_m(0)||\_inf+epsilon\_m; if the computed matrix tilde Q\_m satisfies only ||I- J\_mtilde Q\_m||\_inf ->inf<=delta\_m<1, then...}
\expandafter\gdef\csname CGTip@thm:rb-staged\endcsname{Theorem 18.27 (Staged certificate implies route B') (p. 120): Let C<tau\_1-2pi. Start at kappa\textasciicircum{}(1,0):=tau and, for m>=1, let kappa\textasciicircum{}(m) be the nondecreasing rearrangement of kappa\textasciicircum{}(m,0)+x\_*\textasciicircum{}(m) given by Theorem 18.25, and kappa\textasciicircum{}(m+1,0):=kappa\textasciicircum{}(m). Suppose (A\_m) holds for every m and sum\_j>=12 L\_j h\_j=sup\_m>=1 sum\_j<= m2 L\_j h\_j<= C .B Then every kappa\textasciicircum{}(m) is C-admissible and equ...}
\expandafter\gdef\csname CGTip@eq:rb-B\endcsname{Eq. (B), p. 120: sum\_j>=12 L\_j h\_j=sup\_m>=1 sum\_j<= m2 L\_j h\_j<= C .B}
\expandafter\gdef\csname CGTip@rem:rb-certificate-status\endcsname{Remark 18.28 (What the certificate theorem does and does not do) (p. 120): Theorem 18.27 reduces the extension from one stage to the next to two computable quantities per stage, l\_m=2 L\_m\textasciicircum{}2 H\_m h\_m and 2 L\_m h\_m. It requires no contraction factor q, no separate control of the perturbation of old tests, and no prior ``entry into the basin'': old tests are kept in T\_m+1 and are solved again...}
\expandafter\gdef\csname CGTip@prop:rb-graded\endcsname{Proposition 18.29 (Graded extension) (p. 120): Let s\_m>=0 satisfy s\_m+1<= q\_m s\_m+ h\_m with 0<= q\_m<=1 and sum\_m h\_m<inf, and assume prod\_j=L\textasciicircum{}m q\_j->0(m->inf)for every L>=1 . Then s\_m->0. Condition (18.26) holds if q\_m<= q<1 for all large m, and more generally if sum\_m(1- q\_m)=inf. If the lists kappa\textasciicircum{}(m) producing s\_m= s\_m(kappa\textasciicircum{}(m), a\textasciicircum{}zeta) are C-admissible sec...}
\expandafter\gdef\csname CGTip@eq:rb-graded-tail\endcsname{Eq. (18.26), p. 120: prod\_j=L\textasciicircum{}m q\_j->0(m->inf)for every L>=1 .}
\expandafter\gdef\csname CGTip@rem:rb-graded-status\endcsname{Remark 18.30 (Tube budget, and the two geometries) (p. 121): In Proposition 18.29 the tube |kappa\textasciicircum{}(m)\_k-tau\_k|<= C, uniform in m, is a separate hypothesis: summability of the coordinate increments sum\_m|kappa\textasciicircum{}(m+1)\_k-kappa\textasciicircum{}(m)\_k| gives convergence of each coordinate but not that the limit stays in a prescribed tube, so a construction must carry a displacement budget sum\_m>= m...}
\expandafter\gdef\csname CGTip@heur:rb-frame\endcsname{Heuristic 18.31 (Frame bound from exponential sums of the clock) (p. 121): In the opposite regime k\textasciicircum{}*\_m>> t\_m one may take Q\_m:= J\_m\textasciicircum{} T(J\_m J\_m\textasciicircum{} T)\textasciicircum{}-1. If ||(1)/(k\textasciicircum{}*\_m) J\_m J\_m\textasciicircum{} T-D||\_inf ->inf<=epsilon<d\_- for a diagonal D>= d\_-I, and |(J\_m)\_varphi,k|<= a\_J, then L\_m<= t\_m a\_J/(k\textasciicircum{}*\_m(d\_--epsilon))=O(t\_m/k\textasciicircum{}*\_m) and L\_m\textasciicircum{}2 H\_m=O(t\_m\textasciicircum{}2 B\_0/k\textasciicircum{}*\_m) with B\_0:=max\_varphi int|varphi|u\textasciicircum{}2, so (A\_m)...}
\expandafter\gdef\csname CGTip@subsec:rb-gonek\endcsname{Section 18.3.4: Explicit lists in the tube: Gonek's symmetrised Euler products (p. 121)}
\expandafter\gdef\csname CGTip@eq:rb-alphaX\endcsname{Eq. (18.27), p. 121: p\_ X(t):=arg P\_ X((1)/(2)+it)=-sum\_n<= X\textasciicircum{}2(Lambda\_ X(n)sin(tlog n))/(sqrt(n)log n),}
\expandafter\gdef\csname CGTip@eq:rb-aux-16\endcsname{Eq. (18.28), p. 121: m\_ X:=sum\_n<= X\textasciicircum{}2(Lambda\_ X(n))/(sqrt(n)log n), n\_ X:=sum\_n<= X\textasciicircum{}2(Lambda\_ X(n))/(sqrt(n)),}
\expandafter\gdef\csname CGTip@eq:rb-gonek-phase\endcsname{Eq. (18.29), p. 121: varTheta\_ X(t):=vartheta(t)+ p\_ X(t)in(pi)/(2)+pi Z .}
\expandafter\gdef\csname CGTip@lem:rb-gonek-tube\endcsname{Lemma 18.32 (Gonek lists lie in a tube around the clock) (p. 122): Fix X>=2 and put t\_ X:=max (2pi e\textasciicircum{}2 n\_ X+2, t\_1), where t\_1 is such that |vartheta'(t)-(1)/(2)log(t)/(2pi)|<=(1)/(2) for t>= t\_1. Then varTheta\_ X is strictly increasing on [ t\_ X,inf) with varTheta\_ X'>=(1)/(2), the solutions of (18.29) in [ t\_ X,inf) form a strictly increasing sequence t\_1\textasciicircum{}(X)< t\_2\textasciicircum{}(X)<..., each a...}
\expandafter\gdef\csname CGTip@eq:rb-gonek-list\endcsname{Eq. (18.30), p. 122: kappa\textasciicircum{}(X)\_k:= tau\_k, k<= j\_ X, t\textasciicircum{}(X)\_k- j\_ X, k> j\_ X,}
\expandafter\gdef\csname CGTip@eq:rb-gonek-tube\endcsname{Eq. (18.31), p. 122: |kappa\textasciicircum{}(X)\_k-tau\_k|<=(2 m\_ X+o(1))/(log(tau\_k/2pi))(k->inf).}
\expandafter\gdef\csname CGTip@eq:rb-aux-15\endcsname{Eq. (18.32), p. 122: N\_RVM (kappa\textasciicircum{}(X)\_k)=k-(1)/(2)-(p\_ X(kappa\textasciicircum{}(X)\_k))/(pi)+ r\_vartheta (kappa\textasciicircum{}(X)\_k)(k> j\_ X),}
\expandafter\gdef\csname CGTip@eq:rb-aux-17\endcsname{Eq. (18.33), p. 122: |N\_RVM(kappa\textasciicircum{}(X)\_k)-N\_RVM(tau\_k)|<=(m\_ X)/(pi)+(c\_r)/(kappa\textasciicircum{}(X)\_k)}
\expandafter\gdef\csname CGTip@rem:rb-gonek-status\endcsname{Remark 18.33 (What the Gonek lists give and what remains) (p. 122): The lists kappa\textasciicircum{}(X) use the primes up to X\textasciicircum{}2 and the clock, and nothing about the zeros of zeta. Their prime side is built in: formally, sum\_kdelta(t-kappa\textasciicircum{}(X)\_k)=(varTheta\_ X'(t))/(pi)[1+2sum\_j>=1(-1)\textasciicircum{}jcos(2jvarTheta\_ X(t))] where varTheta\_ X'>0, and the j=0 term carries p\_ X'/pi=-(1)/(pi)sum\_n<= X\textasciicircum{}2Lambda\_ X(n)n\textasciicircum{}-...}
\expandafter\gdef\csname CGTip@subsec:rb-window\endcsname{Section 18.4: Route B on a window: unknown lists, a Bernstein form, and minimizers (p. 123)}
\expandafter\gdef\csname CGTip@lem:rb-peel\endcsname{Lemma 18.34 (Laplace transforms against the zeros) (p. 123): Let mu be a positive Borel measure on R with int e\textasciicircum{}-yt dmu(t)<inf for y>y\_1, and suppose that int\_ Re\textasciicircum{}-yt dmu(t)=sum\_k>=1e\textasciicircum{}iyd\_ke\textasciicircum{}-ygamma\_k(y>y\_1). Then all nontrivial zeros of zeta lie on Re s=(1)/(2), and mu=sum\_kdelta\_gamma\_k.}
\expandafter\gdef\csname CGTip@eq:rb-peel\endcsname{Eq. (18.34), p. 123: int\_ Re\textasciicircum{}-yt dmu(t)=sum\_k>=1e\textasciicircum{}iyd\_ke\textasciicircum{}-ygamma\_k(y>y\_1).}
\expandafter\gdef\csname CGTip@thm:rb-window-unknown\endcsname{Theorem 18.35 (Route B on one window, with an unknown list) (p. 123): Let Jin(0,inf) be an open interval containing no point log p\textasciicircum{}m, for instance J=(log 2,log 3). Let kappa=(kappa\_k)\_k>=1 be a nondecreasing sequence of real numbers with kappa\_1>0 and N\_kappa(T)<<(1+T)\textasciicircum{}A for some A. If sum\_k>=1hat f(kappa\_k)=int\_J f(u) e\textasciicircum{}-u/2V\_+(u) dufor every real fin C\_c\textasciicircum{}inf(J), then all nontrivial...}
\expandafter\gdef\csname CGTip@eq:rb-window-unknown\endcsname{Eq. (18.35), p. 123: sum\_k>=1hat f(kappa\_k)=int\_J f(u) e\textasciicircum{}-u/2V\_+(u) dufor every real fin C\_c\textasciicircum{}inf(J),}
\expandafter\gdef\csname CGTip@thm:rb-bernstein\endcsname{Theorem 18.36 (A Bernstein form of the Riemann hypothesis) (p. 124): The function y-> Z\_xi(iy)=e\textasciicircum{}-iy/2V\_+(iy)=sum\_ke\textasciicircum{}iyd\_ke\textasciicircum{}-ygamma\_k is real on (0,inf) and has the zero-free representation (5.11) (iyin Omega\_theta). The following are equivalent: item RH; item y-> Z\_xi(iy) is completely monotone on (0,inf): (-1)\textasciicircum{}n(d\textasciicircum{}n)/(dy\textasciicircum{}n) Z\_xi(iy)>=0 for all n>=0 and y>0; item y-> Z\_xi(iy) is com...}
\expandafter\gdef\csname CGTip@rem:rb-window-positivity\endcsname{Remark 18.37 (Where the positivity sits) (p. 124): In Proposition 5.2 one imaginary frequency suffices, because every term of the defect X(iy\_0) is nonnegative. For an unknown list there is no such sign, and Theorems 18.35 and 18.36 use all frequencies of a half-line; the positivity is that of the unknown measure mu, which Lemma 18.34 turns, group by group from the...}
\expandafter\gdef\csname CGTip@eq:rb-box-window\endcsname{Eq. (18.36), p. 124: T\_C:=prod\_k>=1[max(2pi,tau\_k-C), tau\_k+C],}
\expandafter\gdef\csname CGTip@eq:rb-EW\endcsname{Eq. (18.37), p. 124: E\_W(kappa):=sum\_j>=1 w\_j < R(kappa, a\textasciicircum{}zeta),varphi\_j>\textasciicircum{}2, P\_kappa:=sum\_j>=1 w\_j < R(kappa, a\textasciicircum{}zeta),varphi\_j> varphi\_j .}
\expandafter\gdef\csname CGTip@lem:rb-minimizer\endcsname{Lemma 18.38 (Minimizers on a window) (p. 124): item E\_W is continuous on T\_C and attains its minimum. item Let kappa\textasciicircum{}* be a minimizer, and put G(z):=int\_u\_1\textasciicircum{}u\_2u P\_kappa\textasciicircum{}*(u)sin(zu) du(zin C). Then E\_W(kappa\textasciicircum{}*)=< R(kappa\textasciicircum{}*, a\textasciicircum{}zeta), P\_kappa\textasciicircum{}*>, and G(kappa\textasciicircum{}*\_k)=0 for every k for which kappa\textasciicircum{}*\_k is an interior point of its interval in (18.36). item Dichotomy. Eit...}
\expandafter\gdef\csname CGTip@eq:rb-G-minimizer\endcsname{Eq. (18.38), p. 124: G(z):=int\_u\_1\textasciicircum{}u\_2u P\_kappa\textasciicircum{}*(u)sin(zu) du(zin C).}
\expandafter\gdef\csname CGTip@thm:rb-window-equivalence\endcsname{Theorem 18.39 (Window solvability) (p. 125): Let C>= C\_eta. The following are equivalent: item RH; item for every U>1 there is kappa in T\_C with R(kappa, a\textasciicircum{}zeta)=0 in D'(U\textasciicircum{}-1,U); item for every U>1, min\_ T\_C E\_(U\textasciicircum{}-1,U)=0, for one (equivalently every) admissible choice of tests and weights. In particular RH holds if, for every U>1, some minimizer of E\_(U\textasciicircum{}-1,U)...}
\expandafter\gdef\csname CGTip@rem:rb-minimizer-open\endcsname{Remark 18.40 (What the dichotomy leaves open) (p. 125): Lemma 18.38 uses no stability estimate: the first-order condition at a minimizer is a completeness statement for usin(kappa\_ku) on a bounded interval, and the clock density (log r)/2pi per unit length is far above the density u\_2/log 2 that a nonzero function of exponential type u\_2 can absorb (Jensen's formula; the...}
\expandafter\gdef\csname CGTip@heur:rb-window-zerofree\endcsname{Heuristic 18.41 (Window size and zero-free regions) (p. 126): Theorem 5.7 gives, without hypothesis, k-sum\_n a\textasciicircum{}zeta\_ndelta\_log n=sum\_kcosh(d\_ku)cos(gamma\_ku) in D'(0,inf). So a solution kappa in T\_C on the window (U\textasciicircum{}-1,U) satisfies sum\_kcos(kappa\_ku)=sum\_kcosh(d\_ku)cos(gamma\_ku) there. Test it with psi(u)cos(gamma\_0u), psi >=0 supported in [U- w,U), where (1)/(2)+d\_0+igamma\_0...}
\expandafter\gdef\csname CGTip@eq:rb-window-zerofree\endcsname{Eq. (18.39), p. 126: cosh (d\_0(U- w))<\textasciitilde{}(C+(1)/(w))(log gamma\_0)/(2pi),that is,|d\_0|<\textasciitilde{}(log(Clog gamma\_0))/(U) :}
\expandafter\gdef\csname CGTip@subsec:rb-apriori\endcsname{Section 18.5: An a priori bound for quasi-solutions (p. 126)}
\expandafter\gdef\csname CGTip@def:rb-laplace-residual\endcsname{Definition 18.42 (Laplace residual) (p. 126): For Re s>1 put hat R(s):=(1)/(2) (P\_kappa')/(P\_kappa)(s)-(1)/(2) (G')/(G)(s)+sum\_n=2\textasciicircum{}n\textasciicircum{}* a\_n n\textasciicircum{}(1)/(2)-s. For T>=0 put E(T):=int\_-T\textasciicircum{}T|hat R(2+iv)| dv, upsilon:=sup\_T>=2(E(T))/(log T)in[0,inf]. Since int\_0\textasciicircum{}inf e\textasciicircum{}-(s-(1)/(2))ucos(kappa u) du=(s-(1)/(2))/((s-(1)/(2))\textasciicircum{}2+kappa\textasciicircum{}2)=(1)/(2)(d)/(ds)log (1+(q(s))/((1)/(4)+kap...}
\expandafter\gdef\csname CGTip@eq:rb-Rhat\endcsname{Eq. (18.40), p. 126: hat R(s):=(1)/(2) (P\_kappa')/(P\_kappa)(s)-(1)/(2) (G')/(G)(s)+sum\_n=2\textasciicircum{}n\textasciicircum{}* a\_n n\textasciicircum{}(1)/(2)-s.}
\expandafter\gdef\csname CGTip@eq:rb-E\endcsname{Eq. (18.41), p. 126: E(T):=int\_-T\textasciicircum{}T|hat R(2+iv)| dv, upsilon:=sup\_T>=2(E(T))/(log T)in[0,inf].}
\expandafter\gdef\csname CGTip@lem:rb-upsilon-infinite\endcsname{Lemma 18.43 (When the Laplace residual grows logarithmically) (p. 126): For a finite section, upsilon<inf if and only if a\_n=0 for every n. More precisely, with D(v):=sum\_n=2\textasciicircum{}n\textasciicircum{}* a\_nn\textasciicircum{}-3/2-iv, E(T)=int\_-T\textasciicircum{}T| D(v)| dv+O\_C,k\textasciicircum{}*(log T), liminf\_T->inf(1)/(2T)int\_-T\textasciicircum{}T| D(v)| dv>=(sum\_n<= n\textasciicircum{}*| a\_n|\textasciicircum{}2n\textasciicircum{}-3)/(sum\_n<= n\textasciicircum{}*| a\_n|n\textasciicircum{}-3/2)>0 (a!=0).}
\expandafter\gdef\csname CGTip@eq:rb-upsilon-split\endcsname{Eq. (18.42), p. 126: E(T)=int\_-T\textasciicircum{}T| D(v)| dv+O\_C,k\textasciicircum{}*(log T), liminf\_T->inf(1)/(2T)int\_-T\textasciicircum{}T| D(v)| dv>=(sum\_n<= n\textasciicircum{}*| a\_n|\textasciicircum{}2n\textasciicircum{}-3)/(sum\_n<= n\textasciicircum{}*| a\_n|n\textasciicircum{}-3/2)>0 (a!=0).}
\expandafter\gdef\csname CGTip@rem:rb-count-scope\endcsname{Remark 18.44 (Scope of the counting bound) (p. 127): Lemma 18.43 restricts Theorem 18.51 and Corollary 18.52, as stated for finite sections, to the sections without masses, for which |N\_kappa-n\_RVM|<= k\textasciicircum{}*+1 is immediate. The reason is that on a vertical line the Laplace transform of the prime side is almost periodic and does not decay, so an L\textasciicircum{}1 residual of logarithmi...}
\expandafter\gdef\csname CGTip@eq:rb-beta\endcsname{Eq. (18.43), p. 127: B(C):=sum\_k>=1 sup\_xin[max(tau\_k-C,2pi), tau\_k+C] sup\_sin J| f(x;s)- f(tau\_k;s)| .}
\expandafter\gdef\csname CGTip@lem:rb-PL\endcsname{Lemma 18.45 (Phragmen--Lindelof, two forms) (p. 127): (a) Let S:= l\_1<=Re s<= l\_2, Im s>= t\_1 have width a:= l\_2- l\_1. Let u be upper semicontinuous on S and subharmonic in its interior, with u<= M\_0 on d S and u(s)<= A\_0e\textasciicircum{}cIm s for some c<pi/a. Then u<= M\_0 on S.; (b) Let g be holomorphic in Re s>1, continuous on Re s>=2, with |g|<=1 on Re s=2 and | g(s)|<= A\_0(2+|s|)...}
\expandafter\gdef\csname CGTip@lem:rb-factor\endcsname{Lemma 18.46 (The Dirichlet factor) (p. 127): Put D\_ a(s):=sum\_n=2\textasciicircum{}n\textasciicircum{}*(2 a\_n)/(log n)n\textasciicircum{}(1)/(2)-s and Y:=zeta\_kappae\textasciicircum{}- D\_ a. Then Y is holomorphic and zero-free in Re s>1, and | D\_ a|<= M c\_l on Re s>=2. For the branch of log Y that is real on (1,inf), (log Y)'=2 hat R, |log Y(2)|<= B\_0:=|log zeta\_RVM(2)|+ B(C)+ M c\_l, and, for all real t, |log Y(2+it)|<= B\_0+ E...}
\expandafter\gdef\csname CGTip@eq:rb-logY-deriv\endcsname{Eq. (18.44), p. 127: (log Y)'=2 hat R, |log Y(2)|<= B\_0:=|log zeta\_RVM(2)|+ B(C)+ M c\_l,}
\expandafter\gdef\csname CGTip@eq:rb-line\endcsname{Eq. (18.45), p. 128: |log Y(2+it)|<= B\_0+ E(|t|)<= B\_0+upsilon log(2+|t|).}
\expandafter\gdef\csname CGTip@lem:rb-halfplane\endcsname{Lemma 18.47 (Bound on the half-plane) (p. 128): With W:=log|zeta\_kappa|=log|P\_kappa|-log|G| and B\_1:= B\_0+upsilon log(2sqrt(2)), B\_2:= B\_1+ M c\_l, for Re s>=2, |log| Y(s)||<= B\_1+upsilon log|s-1|, | W(s)|<= B\_2+upsilon log|s-1| .}
\expandafter\gdef\csname CGTip@eq:rb-halfplane\endcsname{Eq. (18.46), p. 128: |log| Y(s)||<= B\_1+upsilon log|s-1|, | W(s)|<= B\_2+upsilon log|s-1| .}
\expandafter\gdef\csname CGTip@lem:rb-strip\endcsname{Lemma 18.48 (Bound on the half-strip) (p. 128): Let S:=-1<=Re s<=2, Im s>=1 and B\_3:=max (B\_2+(3)/(4)log 2, c\_tau+ B(C)). Then W(s)<= B\_3+((3)/(2)+upsilon)log|s+2|(sin S).}
\expandafter\gdef\csname CGTip@eq:rb-strip\endcsname{Eq. (18.47), p. 128: W(s)<= B\_3+((3)/(2)+upsilon)log|s+2|(sin S).}
\expandafter\gdef\csname CGTip@eq:rb-G-ratio\endcsname{Eq. (18.48), p. 128: |(G(2+it))/(G(-1+it))|=(|1+it|)/(2) pi\textasciicircum{}-3/2((t)/(2)coth(pi t)/(2))\textasciicircum{}1/2<=(1+t)\textasciicircum{}3/2(t>=1).}
\expandafter\gdef\csname CGTip@lem:rb-backlund\endcsname{Lemma 18.49 (Backlund--Jensen on a horizontal segment) (p. 128): There are absolute constants T\_0>=3 and c\_5 such that, for T>= T\_0 with Tnot in kappa\_k, the variation of arg P\_kappa along the segment from 2+iT to (1)/(2)+iT is at most pi(n\_T+1), where n\_T<=(((5)/(2)+upsilon)log(T+8)+ E(T)+ B\_0+ B\_3+ M c\_l+2 c\_5)/(log(4/3)) .}
\expandafter\gdef\csname CGTip@eq:rb-nu\endcsname{Eq. (18.49), p. 129: n\_T<=(((5)/(2)+upsilon)log(T+8)+ E(T)+ B\_0+ B\_3+ M c\_l+2 c\_5)/(log(4/3)) .}
\expandafter\gdef\csname CGTip@lem:rb-argument\endcsname{Lemma 18.50 (Argument principle) (p. 129): For T>0, Tnot in kappa\_k, let G\_T be the path 2->2+iT->(1)/(2)+iT. Then pi N\_kappa(T)=Delta\_ G\_Targ P\_kappa =pi+vartheta(T)+Delta\_ G\_Targ zeta\_kappa, where vartheta(T)=arg (pi\textasciicircum{}-iT/2Gamma((1)/(4)+(iT)/(2))) is the Riemann--Siegel theta function, defined by continuous variation from vartheta(0)=0.}
\expandafter\gdef\csname CGTip@eq:rb-arg-principle\endcsname{Eq. (18.50), p. 129: pi N\_kappa(T)=Delta\_ G\_Targ P\_kappa =pi+vartheta(T)+Delta\_ G\_Targ zeta\_kappa,}
\expandafter\gdef\csname CGTip@thm:rb-count\endcsname{Theorem 18.51 (Counting bound for quasi-solutions) (p. 129): There are absolute constants T\_0 and c\_5 such that, for every finite section in K(C, M) and every T>= T\_0, |N\_kappa(T)-n\_RVM(T)|<=(c\_1+ c\_2upsilon)log(T+8)+ c\_3 B+ c\_4, with c\_1=(5)/(2log(4/3))=8.690..., c\_2=(2)/(log(4/3))+(1)/(pi)=7.270..., c\_3=(1)/(log(4/3))+(1)/(pi)=3.794..., c\_4:=(3)/(2)+(c\_6)/(pi)+(1)/(T\_0)+(2...}
\expandafter\gdef\csname CGTip@eq:rb-count\endcsname{Eq. (18.51), p. 129: |N\_kappa(T)-n\_RVM(T)|<=(c\_1+ c\_2upsilon)log(T+8)+ c\_3 B+ c\_4,}
\expandafter\gdef\csname CGTip@eq:rb-a-constants\endcsname{Eq. (18.52), p. 129: c\_1=(5)/(2log(4/3))=8.690..., c\_2=(2)/(log(4/3))+(1)/(pi)=7.270..., c\_3=(1)/(log(4/3))+(1)/(pi)=3.794...,}
\expandafter\gdef\csname CGTip@eq:rb-aux-1\endcsname{Eq. (18.53), p. 130: |N\_kappa(T)-n\_RVM(T)|<= n\_T+1+(E(T)+ M c\_l+ c\_6)/(pi)+(1)/(T)+(1)/(2) .}
\expandafter\gdef\csname CGTip@cor:rb-displacement\endcsname{Corollary 18.52 (Displacement bound) (p. 130): Assume T\_0>=2pi e. For every k with tau\_k-C>= T\_0, |kappa\_k-tau\_k|<=(2pi [(c\_1+ c\_2upsilon)log(tau\_k+C+8)+ c\_3 B+ c\_4])/(log ((tau\_k-C)/2pi)) . In particular limsup\_k->inf|kappa\_k-tau\_k|<=2pi(c\_1+ c\_2upsilon) uniformly over sections with given C, M,upsilon. For every p>0 there is k\_0(C, M,upsilon, p) such that |kapp...}
\expandafter\gdef\csname CGTip@eq:rb-disp\endcsname{Eq. (18.54), p. 130: |kappa\_k-tau\_k|<=(2pi [(c\_1+ c\_2upsilon)log(tau\_k+C+8)+ c\_3 B+ c\_4])/(log ((tau\_k-C)/2pi)) .}
\expandafter\gdef\csname CGTip@rem:rb-tube\endcsname{Remark 18.53 (What Corollary 18.52 does not give) (p. 130): The proof uses the imposed tube |kappa\_k-tau\_k|<= C twice: through B(C) in B, and to replace D(tau\_k+ q) by D(tau\_k+C). The conclusion is that, inside any finite tube, the coordinates of sufficiently high index lie in the smaller tube of radius 2pi(c\_1+ c\_2upsilon)+ p. So the constraint is inactive there once C>2pi(...}
\expandafter\gdef\csname CGTip@heur:rb-window\endcsname{Heuristic 18.54 (What a finite section controls) (p. 130): The quantity in Theorem 18.51 is E(T), the L\textasciicircum{}1 norm on Re s=2 of the Laplace transform of the residual. It is not the Euclidean residual that a numerical section minimises over finitely many tests. Let the tests be j\_0(u)varphi(u) with a fixed window j\_0in C\_c\textasciicircum{}inf(0,inf), j\_0==1 on [u\_0,u\textasciicircum{}*], and varphi band-limited...}
\expandafter\gdef\csname CGTip@heur:rb-informative\endcsname{Heuristic 18.55 (Informative part of a section) (p. 130): Where the atoms are denser than the resolution of the tests, e\textasciicircum{}-u<<pi/ w\textasciicircum{}sharp, the masses alone can approximate a smooth target by a Riemann sum a\_n\textasciitilde{} g(log n)/n without moving any ordinate. That part of the window then carries no information on the ordinates. Counting unknowns against equations in the remaining win...}
\expandafter\gdef\csname CGTip@subsec:rb-riesz\endcsname{Section 18.6: A Kadec--Avdonin criterion with blocks of variable length (p. 131)}
\expandafter\gdef\csname CGTip@subsec:rb-generating\endcsname{Section 18.6.1: Discrepancy and generating function (p. 131)}
\expandafter\gdef\csname CGTip@def:rb-discrepancy\endcsname{Definition 18.56 (Discrepancy, Poisson and conjugate operators) (p. 131): Let X=(x\_k)\_kin Z be a real nondecreasing sequence with | x\_k-(k-(1)/(2))|<= c\_ Xlog(2+|k|)(kin Z). Its discrepancy with respect to the half-integers is varPsi\_ X(x):=sum\_kin Z(1\_[ x\_k,inf)(x)- 1\_[k-(1)/(2),inf)(x)). The k-th term is +-1 on the interval between x\_k and k-(1)/(2) and 0 elsewhere. By (18.55), at most...}
\expandafter\gdef\csname CGTip@eq:rb-H\endcsname{Eq. (18.55), p. 131: | x\_k-(k-(1)/(2))|<= c\_ Xlog(2+|k|)(kin Z).}
\expandafter\gdef\csname CGTip@eq:rb-psi\endcsname{Eq. (18.56), p. 131: varPsi\_ X(x):=sum\_kin Z(1\_[ x\_k,inf)(x)- 1\_[k-(1)/(2),inf)(x)).}
\expandafter\gdef\csname CGTip@eq:rb-PQ\endcsname{Eq. (18.57), p. 131: P\_ y g(x):=(1)/(pi)int(y g(s) d s)/((x- s)\textasciicircum{}2+ y\textasciicircum{}2), Y\_ y g(x):=(1)/(pi)int g(s)[(x- s)/((x- s)\textasciicircum{}2+ y\textasciicircum{}2)+(s)/(1+ s\textasciicircum{}2)]d s .}
\expandafter\gdef\csname CGTip@lem:rb-generating\endcsname{Lemma 18.57 (Generating function) (p. 131): Let X=(x\_k)\_kin Z satisfy (18.55), understood as a sequence indexed by Z in which a point of multiplicity m occurs m times. Let F\_0(z):=cos pi z, and fix z\_0in R(Xu(Z+(1)/(2))). The series G\_ X(z):=sum\_kin Z((1)/(z- x\_k)-(1)/(z-k+(1)/(2))) converges absolutely and locally uniformly off Xu(Z+(1)/(2)), and F\_ X(z):= F...}
\expandafter\gdef\csname CGTip@eq:rb-aux-2\endcsname{Eq. (18.58), p. 131: G\_ X(z):=sum\_kin Z((1)/(z- x\_k)-(1)/(z-k+(1)/(2)))}
\expandafter\gdef\csname CGTip@eq:rb-aux-3\endcsname{Eq. (18.59), p. 131: F\_ X(z):= F\_0(z)exp int\_ z\_0\textasciicircum{} z G\_ X(z) d z}
\expandafter\gdef\csname CGTip@eq:rb-F-formula\endcsname{Eq. (18.60), p. 131: log| F\_ X(x+iy)|=log| F\_0(x+iy)|+pi Y\_ yvarPsi\_ X(x)+ c\_ X'' ;}
\expandafter\gdef\csname CGTip@eq:rb-Q-growth\endcsname{Eq. (18.61), p. 132: | Y\_ yvarPsi\_ X(x)|<= c\_ X'''(1+log(1+ y\textasciicircum{}-2))log\textasciicircum{}2(4+| z|)(z= x+iy, y>0).}
\expandafter\gdef\csname CGTip@subsec:rb-criterion\endcsname{Section 18.6.2: The criterion (p. 132)}
\expandafter\gdef\csname CGTip@eq:rb-aux-4\endcsname{Eq. (18.62), p. 132: [ w]\_A\_2:=sup\_I((1)/(|I|)int\_Iw)((1)/(|I|)int\_Iw\textasciicircum{}-1)<inf ,}
\expandafter\gdef\csname CGTip@thm:rb-criterion\endcsname{Theorem 18.58 (Block (1)/(4)-criterion) (p. 132): Let X satisfy (18.55) and suppose varPsi\_ X= A+ B with measurable A, B such that: item || A||\_inf <= d\_*<(1)/(4); item sup\_x| Y\_1 B(x)- c\_ B|<= U for some constant c\_ B. Then w(x):=| F\_ X(x+i)|\textasciicircum{}2 is in (A\_2) and [ w]\_A\_2<=(pi\textasciicircum{}2coth\textasciicircum{}2(pi) e\textasciicircum{}4pi U)/(cos\textasciicircum{}2(2pi d\_*)) .}
\expandafter\gdef\csname CGTip@eq:rb-A2const\endcsname{Eq. (18.63), p. 132: [ w]\_A\_2<=(pi\textasciicircum{}2coth\textasciicircum{}2(pi) e\textasciicircum{}4pi U)/(cos\textasciicircum{}2(2pi d\_*)) .}
\expandafter\gdef\csname CGTip@eq:rb-aux-5\endcsname{Eq. (18.64), p. 132: | W(o)|=e\textasciicircum{}2pi Y\_1+ y\_1 A(x\_1),arg W(o)=-2pi P\_1+ y\_1 A(x\_1)in[-2pi d\_*,2pi d\_*].}
\expandafter\gdef\csname CGTip@eq:rb-aux-6\endcsname{Eq. (18.65), p. 132: (1)/(|I|)int\_I| W|<=(1)/(cos b\_*) (1)/(|I|)int\_IRe W<=(pi)/(cos b\_*)Re W(o\_I)<=(pi)/(cos b\_*)| W(o\_I)| .}
\expandafter\gdef\csname CGTip@cor:rb-riesz\endcsname{Corollary 18.59 (Riesz bases) (p. 133): Under the hypotheses of Theorem 18.58: item if X is uniformly discrete (in particular without repeated points), e\textasciicircum{}i x\_k v\_kin Z is a Riesz basis of L\textasciicircum{}2(-pi,pi); item if X is a finite union of uniformly discrete sets, the generalised divided differences of e\textasciicircum{}i x\_k v in the sense of [ref] form a Riesz basis of L\textasciicircum{}2(-pi...}
\expandafter\gdef\csname CGTip@subsec:rb-variable\endcsname{Section 18.6.3: Kadec, Avdonin, and variable block length (p. 133)}
\expandafter\gdef\csname CGTip@prop:rb-variable-blocks\endcsname{Proposition 18.60 (Bounded conjugate for block-supported oscillation) (p. 133): Let J\_m\_min Z be a partition of R into consecutive intervals of lengths n\_m>=1. Let B=sum\_m B\_m with supp B\_min J\_m, int B\_m=0 and || B\_m||\_L\textasciicircum{}inf<= D'\_m, sup\_m D'\_m<inf. Put U\_rm loc:=sup\_m sup\_ xin J\_m-1u J\_mu J\_m+1|int\_ J\_m((x- s) B\_m(s))/(pi((x- s)\textasciicircum{}2+1)) d s|, O:=sup\_m\_0sum\_|m-m\_0|>=2(n\_m\textasciicircum{}2)/(dist(J\_m\_0, J\_m)\textasciicircum{}2)....}
\expandafter\gdef\csname CGTip@eq:rb-L0\endcsname{Eq. (18.66), p. 133: U\_rm loc:=sup\_m sup\_ xin J\_m-1u J\_mu J\_m+1|int\_ J\_m((x- s) B\_m(s))/(pi((x- s)\textasciicircum{}2+1)) d s|, O:=sup\_m\_0sum\_|m-m\_0|>=2(n\_m\textasciicircum{}2)/(dist(J\_m\_0, J\_m)\textasciicircum{}2).}
\expandafter\gdef\csname CGTip@cor:rb-kadec-avdonin\endcsname{Corollary 18.61 (Kadec and Avdonin, with the constant (1)/(4)) (p. 133): (a) [ref] If | d\_k|<= d\_*<(1)/(4) for all k, then e\textasciicircum{}i x\_k v is a Riesz basis of L\textasciicircum{}2(-pi,pi).; (b) [ref] If X is uniformly discrete, sup\_k| d\_k|<= D, and for a fixed n the means of d\_k over consecutive blocks of n indices are at most d\_*<(1)/(4) in modulus, then e\textasciicircum{}i x\_k v is a Riesz basis of L\textasciicircum{}2(-pi,pi).}
\expandafter\gdef\csname CGTip@thm:rb-variable-N\endcsname{Theorem 18.62 (Avdonin's theorem with variable block length) (p. 133): Let X be uniformly discrete with sup\_k| d\_k|<= D. Let B\_m be a partition of Z into consecutive blocks of n\_m indices, and J\_m:=bigcup\_kin B\_mI\_k. Let conj d\_m be the mean of d\_k over B\_m, D\_m:=max\_kin B\_m| d\_k-conj d\_m|, and n\_m\textasciicircum{}*:= n\_m-1+ n\_m+ n\_m+1. Assume item sup\_m|conj d\_m|<= d\_*<(1)/(4); item O<inf, with O as...}
\expandafter\gdef\csname CGTip@eq:rb-U-variable\endcsname{Eq. (18.67), p. 134: U=16(1+ D)\textasciicircum{}3+(6 V)/(pi)+(V O)/(2pi log 2) .}
\expandafter\gdef\csname CGTip@rem:rb-price\endcsname{Remark 18.63 (The price of variable length) (p. 134): For fixed n, conditions (b) and (c) hold automatically, and Theorem 18.62 is Avdonin's theorem. If n\_m->inf, the averaging condition (a) is still only |conj d\_m|<= d\_*<(1)/(4). However, (c) then requires the oscillation inside the block to decay like 1/log n\_m. In Theorem 18.58 the oscillation enters only through th...}
\expandafter\gdef\csname CGTip@subsec:rb-application\endcsname{Section 18.7: Application to lists near the RvM clock (p. 134)}
\expandafter\gdef\csname CGTip@subsec:rb-unfolding\endcsname{Section 18.7.1: Unfolding (p. 134)}
\expandafter\gdef\csname CGTip@def:rb-unfolding\endcsname{Definition 18.64 (Unfolded list) (p. 134): Let t\_0:=9.6769... be the zero of N\_RVM in (2pi,inf). For a C-admissible list kappa with kappa\_1> t\_0 (true for tau and gamma), put x\_k:=N\_RVM(kappa\_k)>0 for k>=1 and x\_1-k:=- x\_k. The resulting symmetric sequence X\_kappa has discrepancy varPsi\_kappa:=varPsi\_ X\_kappa, and F\_kappa:=F\_ X\_kappa with z\_0=0. For the cloc...}
\expandafter\gdef\csname CGTip@lem:rb-psi-bound\endcsname{Lemma 18.65 (Unfolded discrepancy) (p. 134): For a C-admissible list with kappa\_1> t\_0, | x\_k-(k-(1)/(2))|<=(C)/(2pi)log(tau\_k+C)/(2pi) for k>=1, so (18.55) holds. Moreover, for x=N\_RVM(t) with x>0, varPsi\_kappa(x)=N\_kappa(t)-n\_RVM(t), |N\_kappa(t)-n\_RVM(t)|<=(C)/(pi)log(t+C)/(2pi)+2, and varPsi\_kappa(- x)=-varPsi\_kappa(x) for almost every x.}
\expandafter\gdef\csname CGTip@eq:rb-psi-t\endcsname{Eq. (18.68), p. 134: varPsi\_kappa(x)=N\_kappa(t)-n\_RVM(t), |N\_kappa(t)-n\_RVM(t)|<=(C)/(pi)log(t+C)/(2pi)+2,}
\expandafter\gdef\csname CGTip@prop:rb-O-blocks\endcsname{Proposition 18.66 (The condition of O gives the averaging hypothesis with blocks \textasciitilde{}log\textasciicircum{}2T) (p. 135): Let kappa be C-admissible with kappa\_1> t\_0, and assume the third condition of O (16.1) in the form |int\_2pi\textasciicircum{}R(N\_kappa(t)-N\_RVM(t))dt|<= C\_2log R(R>=3), with a constant C\_2>0; this is no restriction, since (18.69) with one constant holds with every larger one, and it is needed because the block length H(T):=16C\_2log...}
\expandafter\gdef\csname CGTip@eq:rb-O3\endcsname{Eq. (18.69), p. 135: |int\_2pi\textasciicircum{}R(N\_kappa(t)-N\_RVM(t))dt|<= C\_2log R(R>=3),}
\expandafter\gdef\csname CGTip@eq:rb-block-mean\endcsname{Eq. (18.70), p. 135: || A\_kappa||\_inf <=(3)/(16), | J\_m|=(8C\_2)/(pi)log\textasciicircum{}2T\_m (1+o(1))(m->inf).}
\expandafter\gdef\csname CGTip@eq:rb-aux-7\endcsname{Eq. (18.71), p. 135: int\_ JvarPsi\_kappa d x =int\_T\textasciicircum{}T+ H(N\_kappa-n\_RVM)(t) N\_RVM'(t) dt; =N\_RVM'(T) int\_T\textasciicircum{}T+ H (N\_kappa-n\_RVM)+int\_T\textasciicircum{}T+ H (N\_kappa-n\_RVM)(N\_RVM'(t)-N\_RVM'(T)) dt .}
\expandafter\gdef\csname CGTip@eq:rb-aux-8\endcsname{Eq. (18.72), p. 135: |(1)/(| J|)int\_ JvarPsi\_kappa|<=(2C\_2log(T+ H))/(H)+(5)/(4N\_RVM'(T) H)+(H)/(2pi TN\_RVM'(T))((C)/(pi)log(T+ H+C)/(2pi)+2).}
\expandafter\gdef\csname CGTip@rem:rb-not-equivalent\endcsname{Remark 18.67 (Proposition 18.66 is not an equivalence) (p. 135): Proposition 18.66 is an implication. The averaging condition on blocks of length \textasciitilde{}log\textasciicircum{}2T does not imply (18.69): summing |int\_ J\_mvarPsi|<=(3)/(16)| J\_m| over the blocks up to height R gives only O(Rlog R), not O(log R). The two conditions do not coincide.}
\expandafter\gdef\csname CGTip@cor:rb-one-condition\endcsname{Corollary 18.68 (Reduction to one condition) (p. 136): For every kappa as in Proposition 18.66: if Y\_1 B\_kappa is bounded up to a constant, then | F\_kappa(x+i)|\textasciicircum{}2in(A\_2), with constant (18.63) for d\_*=(3)/(16). Corollary 18.59(a) then holds for X\_kappa if it is uniformly discrete, and Corollary 18.59(b) if it is a finite union of uniformly discrete sets.}
\expandafter\gdef\csname CGTip@subsec:rb-weight\endcsname{Section 18.7.2: The weight is a quotient of synthetic zetas (p. 136)}
\expandafter\gdef\csname CGTip@prop:rb-weight\endcsname{Proposition 18.69 (Weight identity) (p. 136): Let kappa be C-admissible and satisfy (18.69). For 0<hbar <=1 and t>=2pi, log |(zeta\_kappa)/(zeta\_RVM)((1)/(2)+hbar+it)|=int\_2pi\textasciicircum{}inf((t- s) (N\_kappa-n\_RVM)(s))/((t- s)\textasciicircum{}2+hbar\textasciicircum{}2) d s+ c\_kappa+ m\_kappa,hbar(t), where the integral converges as an improper integral and c\_kappa is a constant. The remainder satisfies | m\_...}
\expandafter\gdef\csname CGTip@eq:rb-weight-identity\endcsname{Eq. (18.73), p. 136: log |(zeta\_kappa)/(zeta\_RVM)((1)/(2)+hbar+it)|=int\_2pi\textasciicircum{}inf((t- s) (N\_kappa-n\_RVM)(s))/((t- s)\textasciicircum{}2+hbar\textasciicircum{}2) d s+ c\_kappa+ m\_kappa,hbar(t),}
\expandafter\gdef\csname CGTip@eq:rb-aux-9\endcsname{Eq. (18.74), p. 136: log |(P\_kappa)/(P\_tau)(s)|=sum\_k[ l\_hbar(t-kappa\_k)- l\_hbar(t-tau\_k)]+sum\_k[ l\_hbar(t+kappa\_k)- l\_hbar(t+tau\_k)]-sum\_klog((1)/(4)+kappa\_k\textasciicircum{}2)/((1)/(4)+tau\_k\textasciicircum{}2).}
\expandafter\gdef\csname CGTip@rem:rb-local-scale\endcsname{Remark 18.70 (Local scale) (p. 136): With hbar=1 and hbar\_T:=hbar/N\_RVM'(T)=2pi/log(T/2pi) (the notation of the proof of Theorem 18.71), the substitution x=N\_RVM(t) maps the kernel ((t- s) d s)/((t- s)\textasciicircum{}2+hbar\_T\textasciicircum{}2) to ((x- x') d x')/((x- x')\textasciicircum{}2+1) to first order near t=T. Comparing (18.73) with (18.60) at y=1, the unfolded weight | F\_kappa(x+i)|\textasciicircum{}2 near x...}
\expandafter\gdef\csname CGTip@subsec:rb-obstruction\endcsname{Section 18.7.3: Obstruction for the zeta list (p. 136)}
\expandafter\gdef\csname CGTip@thm:rb-obstruction\endcsname{Theorem 18.71 (No (A\_2) for the unfolded zeta ordinates) (p. 136): Let X\_gamma be the unfolded list of the ordinates gamma\_k>0 of the nontrivial zeros of zeta, counted with multiplicity, and let F be any entire function of exponential type whose zero set is X\_gamma, with multiplicities. Then | F(x+ihbar)|\textasciicircum{}2not in(A\_2) for every real hbar !=0. The proof does not assume the Riemann h...}
\expandafter\gdef\csname CGTip@eq:rb-aux-10\endcsname{Eq. (18.75), p. 137: (1)/(|I|)int\_I|f- f\_I|\textasciicircum{}2<=4[ w\_0]\_A\_2for every interval I .}
\expandafter\gdef\csname CGTip@eq:rb-BMO-L2\endcsname{Eq. (18.76), p. 137: (1)/(|I|)int\_I| P\_hbar varPsi-(P\_hbar varPsi)\_I|\textasciicircum{}2<= c\_rm JNfor every interval I .}
\expandafter\gdef\csname CGTip@eq:rb-aux-11\endcsname{Eq. (18.77), p. 138: int\_T\textasciicircum{}2T|S- c|\textasciicircum{}2>=int\_T\textasciicircum{}2TS\textasciicircum{}2-(1)/(T)(int\_T\textasciicircum{}2TS)\textasciicircum{}2=(1)/(2pi\textasciicircum{}2)Tlog log T (1+o(1))}
\expandafter\gdef\csname CGTip@lem:rb-smoothing\endcsname{Lemma 18.72 (Smoothing S at scale 1/log T) (p. 138): Let hbar>0 be fixed, hbar\_T:=2pi hbar/log(T/2pi), and S:=0 on (-inf, t\_0). Then || P\_hbar\_T* S-S||\_L\textasciicircum{}2[T,2T]=O\_hbar(sqrt(T)), unconditionally.}
\expandafter\gdef\csname CGTip@eq:rb-fujii\endcsname{Eq. (18.78), p. 138: int\_T\textasciicircum{}2T(S(t+u)-S(t))\textasciicircum{}2dt<= c\_FTlog(2+ulog T),}
\expandafter\gdef\csname CGTip@rem:rb-fujii\endcsname{Remark 18.73 (A second proof) (p. 138): Lemma 18.72 also follows from Selberg's mean-square approximation of S(t) by S\_ x\_0(t):=-(1)/(pi)sum\_p<= xp\textasciicircum{}-1/2sin(tlog p), x\_0=T\textasciicircum{}delta [ref], and the Montgomery--Vaughan mean value theorem [ref, Cor. 3]. Indeed P\_hbar\_T* S\_ x\_0-S\_ x\_0=-(1)/(pi)sum\_p<= xp\textasciicircum{}-1/2(e\textasciicircum{}-hbar\_Tlog p-1)sin(tlog p) has mean square << T hbar\_...}
\expandafter\gdef\csname CGTip@cor:rb-obstruction\endcsname{Corollary 18.74 (Consequences for the zeta list) (p. 138): (a) e\textasciicircum{}ix\_kxi\_x\_kin X\_gamma is not a Riesz basis of L\textasciicircum{}2(-pi,pi).; (b) If X\_gamma is a finite union of uniformly discrete sets, its generalised divided differences are not a Riesz basis of L\textasciicircum{}2(-pi,pi).; (c) For every decomposition varPsi\_gamma=A+B with ||A||\_inf<(1)/(4), Y\_1B is unbounded. In particular this holds for...}
\expandafter\gdef\csname CGTip@heur:rb-graduated\endcsname{Heuristic 18.75 (Graduated version) (p. 138): The proof of Theorem 18.71 uses intervals I\_T of length \textasciitilde{} Tlog T. On a window of L unfolded points at height T, the mean square oscillation of P\_hbar varPsi is governed by the variance of increments of S over t-intervals of length L/N\_RVM'(T). By Fujii's asymptotics this is expected to be of size (1)/(pi\textasciicircum{}2)log L. Th...}
\expandafter\gdef\csname CGTip@rem:rb-routeB-stability\endcsname{Remark 18.76 (Bearing on route B) (p. 139): Theorem 18.71 rules out a uniform Riesz-basis (Fredholm) structure for the unfolded zeta list on the whole line, and hence a linear stability argument with constants independent of the height. It says nothing about statement B or B'. The compactness step of route B (Proposition 18.5) uses no stability, and uniquenes...}
\expandafter\gdef\csname CGTip@rem:rb-obstruction-literature\endcsname{Remark 18.77 (Literature on the zeta ordinates as an exponential system) (p. 139): We have not found any study of Riesz bases of exponentials, or of the condition (A\_2), for the ordinates of zeta; the search was not exhaustive. The closest work found is the Fourier interpolation theorem of Bondarenko, Radchenko and Seip [ref, Thm. 1.1]: an even function f analytic in a strip |Im z|<(1)/(2)+epsilon...}
\expandafter\gdef\csname CGTip@subsec:rb-numerics\endcsname{Section 18.8: Numerical certificates for finite sections (p. 139)}
\expandafter\gdef\csname CGTip@subsec:rb-num-setup\endcsname{Section 18.8.1: Set-up (p. 139)}
\expandafter\gdef\csname CGTip@subsec:rb-num-init\endcsname{Section 18.8.2: Initialisation: the clock is outside the basin, the Gonek lists are inside (p. 139)}
\expandafter\gdef\csname CGTip@fig:rb-gonek\endcsname{Figure 1 (p. 141): Distance to the zeta ordinates, in units of the local mean spacing 2pi/log(gamma\_k/2pi), of the midpoint clock and of the Gonek lists kappa\textasciicircum{}(X) for X=3,8,20 (k<=160 for X=3,20; k<=400 for X=8). The dashed line is the observed radius of the Newton basin. The ordinates gamma\_k are used here only to measure the distance.}
\expandafter\gdef\csname CGTip@subsec:rb-num-kantorovich\endcsname{Section 18.8.3: Certificates at the Gonek list: the direct bounds fail (p. 140)}
\expandafter\gdef\csname CGTip@subsec:rb-num-certified\endcsname{Section 18.8.4: Floating-point candidate certificates for exact square sections (p. 140)}
\expandafter\gdef\csname CGTip@eq:rb-num-banach\endcsname{Eq. (18.79), p. 141: l(r):= d\_0+|| Q J'||\_inf ->inf r+(1)/(2) ||| Q| A\textasciicircum{}(3)(r)||\_inf ->inf r\textasciicircum{}2<1, h<=(1- l(r)) r,}
\expandafter\gdef\csname CGTip@eq:rb-num-scaling\endcsname{Eq. (18.80), p. 142: L\_m\textasciitilde{} T\textasciicircum{}1.81+-0.10,|| Q J'||\_inf ->inf\textasciitilde{} T\textasciicircum{}2.10+-0.06,sampled margin\textasciitilde{} T\textasciicircum{}-2.08+-0.06(T<=487);}
\expandafter\gdef\csname CGTip@subsec:rb-num-verified\endcsname{Section 18.8.5: Verified certificates with compactly supported tests (p. 142)}
\expandafter\gdef\csname CGTip@eq:rb-bspline-transform\endcsname{Eq. (18.81), p. 142: c\_varphi\textasciicircum{} N\_c,h(t)=hcos(ct)((sin(ht/2))/(ht/2))\textasciicircum{}8,| c\_varphi\textasciicircum{} N\_c,h(t)|<= h((2)/(ht))\textasciicircum{}8 .}
\expandafter\gdef\csname CGTip@eq:rb-W-bspline\endcsname{Eq. (18.82), p. 142: W\_r(varphi\textasciicircum{} N\_c,h)<=(1+u\textasciicircum{}*2)cosh(u\textasciicircum{}*)/(2) h(1+(2)/(h))\textasciicircum{}r .}
\expandafter\gdef\csname CGTip@eq:rb-clock-tail\endcsname{Eq. (18.83), p. 143: sum\_k>K| c\_varphi\textasciicircum{} N\_c,h(tau\_k)|<= h(pi h(1-(11)/(8K)))\textasciicircum{}-8 ((log K)\textasciicircum{}8K\textasciicircum{}-7)/(7(1-(8)/(7log K))),}
\expandafter\gdef\csname CGTip@eq:rb-krawczyk\endcsname{Eq. (18.84), p. 143: l(rho):=sup\_xin B(x\_1,rho)||I- Q D F\_m(x)||\_inf ->inf<1,|| Q F\_m(x\_1)||\_inf <=(1- l(rho)) rho ;}
\expandafter\gdef\csname CGTip@thm:rb-verified\endcsname{Theorem 18.78 (Exact square sections up to height 487; computer-assisted) (p. 143): For each m=1,...,15 the square system F\_m described above has exactly one zero x\_*\textasciicircum{}(m) in the ball ||x-x\_1\textasciicircum{}(m)||\_inf <=rho\_m; the centres x\_1\textasciicircum{}(m) (exact binary numbers) and the radii rho\_m (double-precision numbers) are recorded in the file certV\_centres.json of the repository, and Table 7 gives rho\_m rounded up. Th...}
\expandafter\gdef\csname CGTip@tab:rb-verified\endcsname{Table 7 (p. 143): Verified square sections with B-spline tests (program certV). rho\_m: radius of the ball of Theorem 18.78; upper bounds are rounded up and lower bounds down; l(rho\_m): upper bound of the contraction constant (18.84); L\_m=|| Q||\_inf ->inf; margin: largest (1- l(rho))rho over 36 radii in [10\textasciicircum{}-8,10\textasciicircum{}-1], each value a rig...}
\expandafter\gdef\csname CGTip@rem:rb-verified\endcsname{Remark 18.79 (What Theorem 18.78 says and does not say) (p. 144): For these fifteen stages the candidate certificates of Section 18.8.4 are replaced by verified ones: the inequalities hold in ball arithmetic, for tests in C\_c\textasciicircum{}6(0,inf), with every evaluation error included. The radii rho\_m, from 6.5*10\textasciicircum{}-21 to 1.1*10\textasciicircum{}-12, are set by the bound (18.83) for the clock tail, which more f...}
\expandafter\gdef\csname CGTip@subsec:rb-num-entry\endcsname{Section 18.8.6: The entry step (p. 144)}
\expandafter\gdef\csname CGTip@fig:rb-scaling\endcsname{Figure 2 (p. 145): Constants of the Banach certificate (18.79) for the square sections of Section 18.8.4 (solid lines, with fitted slopes in log T), and the band-only entry test of Section 18.8.6 (dashed): the reciprocal of the band margin and the reciprocal of the residual of the new band at its Gonek values. The entry test passes, i...}
\expandafter\gdef\csname CGTip@subsec:rb-num-reading\endcsname{Section 18.8.7: Reading of the computations (p. 145)}
\expandafter\gdef\csname CGTip@subsec:rb-status\endcsname{Section 18.9: Status of the statements (p. 146)}
\expandafter\gdef\csname CGTip@sec:numerics\endcsname{Section 19: Numerical illustration with 2001052 zeros (p. 146)}
\expandafter\gdef\csname CGTip@subsec:num-clock\endcsname{Section 19.1: The clock and the zeros (p. 147)}
\expandafter\gdef\csname CGTip@exp:clock\endcsname{Experiment 19.1 (Clock, midpoints and displacements) (p. 147): Object: N\_RVM, N\_zeta, tau\_k, eta\_k=tau\_k-gamma\_k. Range: k<= N. Expected: (2.4) to rounding, and |eta\_k|<= C\_eta (Lemma 2.2). Observed: max\_k|N\_RVM(tau\_k)-(k-(1)/(2))|=6.985*10\textasciicircum{}-10; max\_k|eta\_k|=0.9937 (at k=34), mean |eta\_k|=0.1512, root mean square 0.1881, mean eta\_k=1.7*10\textasciicircum{}-7. By blocks of k, max|eta\_k| is 0.840...}
\expandafter\gdef\csname CGTip@fig:clock-zoom\endcsname{Figure 3 (p. 147): The midpoint clock N\_RVM (red), the zero staircase N\_zeta (black), the midpoint ordinates tau\_k (red points) and the points (gamma\_k,k-(1)/(2)) (squares). Left: the first nine zeros. Right: the last twelve zeros of the file.}
\expandafter\gdef\csname CGTip@fig:remainder\endcsname{Figure 4 (p. 148): The computed difference N\_rm file(T)-N\_RVM(T) for T<=200, where N\_rm file counts the supplied ordinates. The file is complete on this range, so the curve is the remainder N\_zeta(T)-N\_RVM(T)=S(T)+O(1/T) of (2.2).}
\expandafter\gdef\csname CGTip@subsec:num-family\endcsname{Section 19.2: The lambda-family and the endpoint (p. 147)}
\expandafter\gdef\csname CGTip@fig:family\endcsname{Figure 5 (p. 148): Top: the observed displacements eta\_k for every 50th k<= N of the supplied range and their empirical distribution. Bottom: moving ordinates in clock units, k<=30, and the functions Z\_lambda(t) on the critical line.}
\expandafter\gdef\csname CGTip@exp:Z1\endcsname{Experiment 19.2 (The endpoint zeta\_Z1 against zeta) (p. 148): Object: zeta\_Z1=P\_Z1/G computed from the N ordinates gamma\_k, and zeta computed by mpmath. Range: sigma in 2,1.25,0.75, 1<= t<=2000 (499 points per line), and the critical line 0.5<= t<=60. Expected: if zeta\_Z1=zeta, the ratio is 1 up to the truncation of the product. Observed: max|zeta\_Z1/zeta-1|=8.7*10\textasciicircum{}-10, 1.1*10...}
\expandafter\gdef\csname CGTip@fig:Z1\endcsname{Figure 6 (p. 149): Left: |zeta\_Z1/zeta-1| on three vertical lines, with zeta\_Z1 computed from the supplied ordinates and the clock tail beyond them. Right: Z\_Z1 and Hardy's Z on the critical line.}
\expandafter\gdef\csname CGTip@exp:product-identity\endcsname{Experiment 19.3 (The value at infinity: the product identity) (p. 148): Object: the ordered partial products P\_K=prod\_k<= K((1)/(4)+gamma\_k\textasciicircum{}2)/((1)/(4)+tau\_k\textasciicircum{}2) of (13.12). Range: K<= N. Expected: by Corollary 13.3, zeta\_Z1=zeta is the statement P\_K->varpi\_RVM. Observed: P\_K=0.96256, 0.96248, 0.96229, 0.961951, 0.961930, 0.9619418189, 0.9619415390 for K=10,10\textasciicircum{}2,...,10\textasciicircum{}6,N, against varpi...}
\expandafter\gdef\csname CGTip@fig:product-identity\endcsname{Figure 7 (p. 149): The partial products P\_K of (13.12) against varpi\_RVM (dashed). Left: K<=10\textasciicircum{}4. Right: (P\_K-varpi\_RVM)*10\textasciicircum{}6 for 10\textasciicircum{}5<= K<= N. The value of the endpoint at +inf is varpi\_Z1=varpi\_RVM/lim\_KP\_K (Theorem 13.1).}
\expandafter\gdef\csname CGTip@subsec:num-transition\endcsname{Section 19.3: The transition theorem and zeta\_RVM against zeta (p. 149)}
\expandafter\gdef\csname CGTip@exp:transition\endcsname{Experiment 19.4 (Transition theorem) (p. 149): Object: zeta\_RVM, w(s), varpi\_RVM, L. Range: sigma'in 0.25,0.5,1.5, 10<= t<=10\textasciicircum{}4. Expected: Theorem 9.1 and Corollary 9.8: t |zeta\_RVM/(varpi\_RVM(1+w))-1| bounded, |w|=(t/2pi)\textasciicircum{}-sigma'(1+O(1/t)). Observed: Richardson extrapolation on Re s=2 (which presupposes an expansion in 1/T and 1/T\textasciicircum{}2 finer than the O(1/T) proved...}
\expandafter\gdef\csname CGTip@exp:varpi-bound\endcsname{Experiment 19.5 (The transition constant: the inequalities of Theorem 9.11) (p. 150): Object: nu\_3, the enclosure of tau\_1, the lower bound M(tau\_1) of (9.48), and, as a cross-check, nu\_4, P\textasciicircum{}crit\_RVM((1)/(2)) and varpi\_RVM. Method: Part A in interval arithmetic (mpmath.iv, 40 digits, outward rounding); Part B in floating point (25 digits for nu\_4, double precision for the product over k<=2*10\textasciicircum{}6). The...}
\expandafter\gdef\csname CGTip@fig:transition\endcsname{Figure 8 (p. 151): Top: the computed residual of the transition formula multiplied by t (it contains the truncation of the product, the tail approximation, the extrapolated value of varpi\_RVM and rounding), and the rescaled vertical limit t\textasciicircum{}min(sigma',1)|1/zeta\_RVM-L|. Bottom: |zeta| and |zeta\_RVM| on sigma=2 and sigma=0.75.}
\expandafter\gdef\csname CGTip@subsec:num-expsums\endcsname{Section 19.4: Exponential sums (p. 150)}
\expandafter\gdef\csname CGTip@exp:expsums\endcsname{Experiment 19.6 (Exponential sums over the two spectra) (p. 150): Object: S(N;log x)=sum\_k<= Nx\textasciicircum{}io\_k for o\_k=tau\_k and o\_k=gamma\_k. Range: N<=2 001 052, xin 2,3,4,6,10. Expected: |S\_RVM|<<1+log N (Theorem 4.1); for gamma\_k, the main term -(T)/(2pi)(Lambda(x))/(sqrt(x)) of Landau's formula (Theorem 1.3), whose error term O\_x(log T) makes the comparison at a finite height illustrati...}
\expandafter\gdef\csname CGTip@fig:expsums\endcsname{Figure 9 (p. 151): Observed size of the exponential sums on the computed range N<=2 001 052, compared with log N: the midpoint spectrum shows no prime frequencies (left); the spectrum of zeta resonates at prime powers (right).}
\expandafter\gdef\csname CGTip@subsec:num-centroids\endcsname{Section 19.5: Centroids, walks and Cramer's function (p. 150)}
\expandafter\gdef\csname CGTip@exp:cramer\endcsname{Experiment 19.7 (Cramer's function and the centroids of the two spectra) (p. 150): Object: e\textasciicircum{}-1/2V\_+(1) from the zero-free representation (5.11); the triangular centroids C\textasciicircum{}pathc\_zeta(N;1) and C\textasciicircum{}pathc\_RVM(N;1); the limit C\textasciicircum{}pathc\_RVM(inf;1) of Theorem 5.13. Method: program s17 (Appendix B.18), mpmath with 30 digits, rays at theta=pi/3, prime powers up to 2*10\textasciicircum{}6, a=1 and a=2. Expected: independence...}
\expandafter\gdef\csname CGTip@exp:walks\endcsname{Experiment 19.8 (Walks of e\textasciicircum{}igamma\_k, e\textasciicircum{}itau\_k and e\textasciicircum{}i(gamma\_k+tau\_k)/2) (p. 152): Object: the walks S\_k=sum\_j<= ke\textasciicircum{}ix\_j for x=gamma, tau and theta=(gamma+tau)/2, k<=1000, and their triangular centroids. Method: program s16 (Appendix B.17). Observed: C\_gamma=0.7720742579+0.7626754110 i, C\_tau=0.0401443298+1.7565288414 i, C\_theta=0.6133423097+1.3939939459 i; the midpoint of the chord is (1)/(2)(C\_g...}
\expandafter\gdef\csname CGTip@fig:walks3d\endcsname{Figure 10 (p. 152): The walks (Re S\_k,Im S\_k,k), k<=1000, with k on the vertical axis, for e\textasciicircum{}igamma\_k (left), e\textasciicircum{}itau\_k (centre) and e\textasciicircum{}i(gamma\_k+tau\_k)/2 (right); the marker is the centroid at height (N+1)/2.}
\expandafter\gdef\csname CGTip@fig:walks-proj\endcsname{Figure 11 (p. 153): Left: the three walks of Figure 10 together. Right: their projection on the complex plane with the three centroids and the midpoint of the chord (1)/(2)(C\_gamma+C\_tau) (cross).}
\expandafter\gdef\csname CGTip@subsec:num-primes\endcsname{Section 19.6: Prime counts: pi and pi\_RVM (p. 152)}
\expandafter\gdef\csname CGTip@exp:primes\endcsname{Experiment 19.9 (Prime counts and RVM primes) (p. 152): Object: J\_RVM, pi\_RVM, c\_n\textasciicircum{}RVM, p\_j\textasciicircum{}RVM against J, pi, the primes. Range: 2<= x<=100 with all N ordinates of each list (3000 points); 2<= x<=1000 with the first K=17 061 ordinates and a smooth cut (40 000 points). Expected: the explicit formula with the gamma\_k reproduces pi(x); pi\_RVM is continuous and pi\_RVM(x)-li...}
\expandafter\gdef\csname CGTip@fig:prime-counts\endcsname{Figure 12 (p. 153): pi(x) (black staircase), the explicit formula with the zeros of zeta (blue), pi\_RVM(x) (red) and li(x)-li(2) (green). The distance between the red and green curves is the quantity bounded in Theorem 6.5; the ratios are in Figure 14.}
\expandafter\gdef\csname CGTip@fig:rvm-primes\endcsname{Figure 13 (p. 153): Left: cell masses c\_n\textasciicircum{}RVM, the prime indicator and 1/log n. Right: real RVM prime levels against the primes up to 1000.}
\expandafter\gdef\csname CGTip@exp:pi-ratios\endcsname{Experiment 19.10 (The three ratios and the square-root scale) (p. 152): Object: li/pi\_RVM, pi/li, pi/pi\_RVM, and (pi\_RVM-li)/B, (pi-li)/B with B(x)=sqrt(x)(log log x)\textasciicircum{}2/log x. Range: the grid of Experiment 19.9 on [8,1000] (K=17 061 ordinates). Expected: the three ratios tend to 1 (Corollary 6.6) and (pi\_RVM-li)/B stays bounded (Theorem 6.5). Observed: at x=100,500,1000, li/pi\_RVM=1.155...}
\expandafter\gdef\csname CGTip@fig:pi-ratios\endcsname{Figure 14 (p. 154): Left: the three ratios of Corollary 6.6. Right: pi\_RVM-li and pi-li in units of the scale B(x) of Theorem 6.5. The ratio pi/pi\_RVM (black) is the closest of the three to 1; the ratios with li approach 1 at the rate (1)/(2)li(sqrt(x))/li(x) of the bias J-pi.}
\expandafter\gdef\csname CGTip@subsec:num-mertens\endcsname{Section 19.7: The Mertens functions and the ternary division (p. 154)}
\expandafter\gdef\csname CGTip@eq:ternary-update\endcsname{Eq. (19.1), p. 154: M\_h(n)=M\_h(n-1)+mu(n)\textasciicircum{}2 varsigma\_n, |M\_h(n)- M(n)| small,}
\expandafter\gdef\csname CGTip@tab:ternary\endcsname{Table 9 (p. 155): The ternary division for n<=15 (the first 40 values are in Appendix C.6). The zero columns of mu and mu\_h coincide; the signs of mu\_h are chosen by comparing M\_h(n-1) with M(n). For n<=12 the two walks agree; at n=13 the target -1.85 lies above M\_h(12)=-2, so mu\_h(13)=+1 while mu(13)=-1.}
\expandafter\gdef\csname CGTip@exp:mertens\endcsname{Experiment 19.11 (Mertens functions and the ternary walk) (p. 154): Object: M, M\_RVM\textasciicircum{}crit, M\_h; the explicit formulas. Range: n<=10\textasciicircum{}6; explicit formulas on [2,2000] with all N ordinates of each list. Expected: the explicit formula with the zeros of zeta reproduces M; M\_RVM\textasciicircum{}crit(x)=O(x\textasciicircum{}1/2+eps) (Theorem 9.14); M\_h tracks M\_RVM\textasciicircum{}crit with the densities of mu. Observed: item explicit fo...}
\expandafter\gdef\csname CGTip@fig:mertens\endcsname{Figure 15 (p. 156): Left: M (black), the tapered-and-interpolated target M for M\_RVM\textasciicircum{}crit (red) and M\_h (green) at every integer n<=2000; the residues are computed from the product for k<=1000 and modelled beyond. Right: range of M and mean of M in blocks of 500 integers up to 10\textasciicircum{}6, with pmsqrt(n).}
\expandafter\gdef\csname CGTip@fig:ternary\endcsname{Figure 16 (p. 156): Left: mu(n) and mu\_h(n) for n<=60: same zeros, different signs. Right: M, the tapered-and-interpolated target M for M\_RVM\textasciicircum{}crit and M\_h for n<=300.}
\expandafter\gdef\csname CGTip@fig:explicit\endcsname{Figure 17 (p. 156): Left: M(x), its explicit formula from the zeros of zeta and M\_RVM\textasciicircum{}crit(x) from the tau\_k on [2,300]. Right: normalized derivatives at the zeros: erratic for zeta, constant varpi\_RVM for zeta\_RVM.}
\expandafter\gdef\csname CGTip@rem:anticorrelation\endcsname{Remark 19.12 (The opposite phase of M and M\_RVM\textasciicircum{}crit between 10\textasciicircum{}3 and 2*10\textasciicircum{}3) (p. 155): In the left panel of Figure 15 the curves of M and M\_RVM\textasciicircum{}crit are in opposite phase between n\textasciitilde{}10\textasciicircum{}3 and n\textasciitilde{}2*10\textasciicircum{}3. Heuristically, the visible anticorrelation is largely explained by the first mode, 2sqrt(x) Re (x\textasciicircum{}igamma\_1/(rho\_1zeta'(rho\_1))) and 2sqrt(x) Re (x\textasciicircum{}itau\_1/(rho\_1\textasciicircum{}RVMzeta\_RVM'(rho\_1\textasciicircum{}RVM))), two oscillations...}
\expandafter\gdef\csname CGTip@subsec:num-first-1000\endcsname{Section 19.8: The first 1000 values (p. 157)}
\expandafter\gdef\csname CGTip@exp:first-1000\endcsname{Experiment 19.13 (The first 1000 values) (p. 157): Object: tau\_k against gamma\_k, p\_j\textasciicircum{}RVM against p\_j, pi\_RVM(n) against pi(n), and M\_RVM\textasciicircum{}crit(n) against M(n), for k,j,n<=1000. Method: program s15 (Appendix B.16). For pi\_RVM on [2,8200] the zero sum of (6.2) uses the first K=183 029 ordinates (tau\_K=128 805.3=2.5*2pi *8200) with the smooth cut of Experiment 19.9, on...}
\expandafter\gdef\csname CGTip@fig:clock-1000\endcsname{Figure 18 (p. 158): The clock against the zeros, k<=1000: gamma\_k and tau\_k; the displacements eta\_k=tau\_k-gamma\_k; the two spectra for k<=60; the distribution of eta\_k (Table 13).}
\expandafter\gdef\csname CGTip@fig:primes-1000\endcsname{Figure 19 (p. 158): The RVM prime levels p\_j\textasciicircum{}RVM against the primes p\_j, j<=1000, and the difference p\_j\textasciicircum{}RVM-p\_j (Table 14).}
\expandafter\gdef\csname CGTip@fig:pi-1000\endcsname{Figure 20 (p. 159): pi(n), pi\_RVM(n) and li(n)-li(2) at the integers n<=1000, and the difference pi\_RVM(n)-pi(n) (Table 15).}
\expandafter\gdef\csname CGTip@fig:mertens-1000\endcsname{Figure 21 (p. 159): M(n) and M\_RVM\textasciicircum{}crit(n) at the integers n<=1000 with pmsqrt(n), and the difference M\_RVM\textasciicircum{}crit(n)-M(n) (Table 16).}
\expandafter\gdef\csname CGTip@subsec:num-rungs\endcsname{Section 19.9: The rung equations: transport, the defect of the clock, and the solution from the clock (p. 157)}
\expandafter\gdef\csname CGTip@exp:transport\endcsname{Experiment 19.14 (Transport on the first rung) (p. 157): Object: R\_ f(1), S\_ f(lambda) and the moments m\_j(f) of (15.80)--(15.84). Range: the first 20 000 ordinates of each list; max\_k>=1000| c\_ f(tau\_k)|=2.5*10\textasciicircum{}-10 and max\_k>=3000| c\_ f(tau\_k)|=2.1*10\textasciicircum{}-13. Expected: if zeta\_Z1=zeta, then S\_ f(0)= R\_ f(1) (Proposition 15.44); the series (15.84) converges to S\_ f(0) with t...}
\expandafter\gdef\csname CGTip@tab:transport\endcsname{Table 10 (p. 160): Experiment 19.14. Left: the moments m\_j(f)=sum\_keta\_k\textasciicircum{} j c\textasciicircum{}(j)\_ f(tau\_k) and the partial sums of the series (15.84) for S\_ f(0), compared with the defect R\_ f(1)=-0.042463282161 of the clock. Right: the rung functional along the family.}
\expandafter\gdef\csname CGTip@exp:clock-defect\endcsname{Experiment 19.15 (The defect of the clock by Poisson summation) (p. 159): Object: the two parts of R\_ f(1) in Proposition 15.46. Expected: R\_ f(1)=(k\_ f- I\_0)-2sum\_m>=1(-1)\textasciicircum{}m I\_m. Observed: t\_0=9.676906787; I\_0=0.063533498702 from (15.86); the mismatch of the smooth densities is k\_ f- I\_0=+0.070110032214, and the oscillating part, obtained as the difference R\_ f(1)-(k\_ f- I\_0) and not by...}
\expandafter\gdef\csname CGTip@exp:clock-to-zeros\endcsname{Experiment 19.16 (From the clock to the ordinates and to the masses of the primes) (p. 160): Object: a finite section of the system (16.6): the unknowns are the ordinates kappa\_k below a height and the masses a\_n, 2<= n<=11; the equations are the pairings of (16.6) with omega\_rm test(u)cos(omega u) and omega\_rm test(u)sin(omega u), where omega\_rm test is a fixed smooth function supported in (0,log 13.5) and...}
\expandafter\gdef\csname CGTip@tab:clock-to-zeros\endcsname{Table 11 (p. 161): Experiment 19.16. Solution of a finite section of the rung system (16.6), started at the clock. Left: the ordinates. Right: the masses at the points log n, which are not assumed, against log p/(2sqrt(n)) at prime powers n=p\textasciicircum{}m.}
\expandafter\gdef\csname CGTip@exp:multiscale-jacobian\endcsname{Experiment 19.17 (Multiscale Jacobian of the clock) (p. 160): Object: the minimum midpoint-clock gap, its bound (15.88), the Ingham critical window 2pi/Delta\_T, and the adjacent-pair bound (15.90). Input: none of the genuine zeta ordinates; the program s12\_multiscale\_jacobian.py generates the tau\_k directly from the Lambert formula up to K=2 001 052. Observed: rrrr toprule K T...}
\expandafter\gdef\csname CGTip@fig:multiscale-gap\endcsname{Figure 22 (p. 162): Minimum gap of the midpoint clock up to height T and the explicit lower bound 2pi/log(T/2pi). Their ratio is already 1+2*10\textasciicircum{}-8 at K=2 001 052.}
\expandafter\gdef\csname CGTip@fig:fixed-window-degeneracy\endcsname{Figure 23 (p. 162): Upper bound in (15.90) for adjacent clock frequencies on four fixed u-windows. The decay is slow, of order 1/log tau\_k, but it is enough to exclude a height-uniform lower frame bound on any fixed window.}
\expandafter\gdef\csname CGTip@sec:conclusions\endcsname{Section 20: Conclusions and open problems (p. 163)}
\expandafter\gdef\csname CGTip@rem:not-proved-list\endcsname{Remark 20.1 (Statements derived in part or conjectural, not proved here) (p. 168): The following statements arose in the work on the final step; none of them is used in this paper. item Rates in lambda. Off the powers of 2 the formulas of Proposition 15.27 suggest tilde h\_lambda(s)-tilde h\_0(s)=O\_s(lambda),| h\_n(lambda;s)- h\_n(0;s)|<<\_s,epslambda n\textasciicircum{}-sigma-(1)/(2)+eps; the cells n=2\textasciicircum{}j and the unifo...}
\expandafter\gdef\csname CGTip@eq:rates-lambda\endcsname{Eq. (20.1), p. 168: tilde h\_lambda(s)-tilde h\_0(s)=O\_s(lambda),| h\_n(lambda;s)- h\_n(0;s)|<<\_s,epslambda n\textasciicircum{}-sigma-(1)/(2)+eps;}
\expandafter\gdef\csname CGTip@rem:open-problems\endcsname{Remark 20.2 (Open problems) (p. 168): item Prove the identity zeta\_Z1=zeta. Each of the following is proved to be equivalent to it, unconditionally and within the class of objects stated: (a) condition (D), an absolutely convergent Dirichlet series for zeta\_Z1 in Re s>1 (Theorem 12.2); (b) an absolutely convergent Dirichlet expansion of K=zeta\_Z1/zeta\_Z...}
\expandafter\gdef\csname CGTip@app:tables-1000\endcsname{Section A: Tables of the first 1000 values (p. 169)}
\expandafter\gdef\csname CGTip@tab:clock-1000\endcsname{Table 13 (p. 170): The ordinates gamma\_k of the zeros of zeta, the midpoint ordinates tau\_k of the clock and the displacements eta\_k=tau\_k-gamma\_k, k<=1000. The displacements are computed from the full-precision ordinates and rounded independently, so they may differ in the last digit from the difference of the displayed columns.}
\expandafter\gdef\csname CGTip@tab:primes-1000\endcsname{Table 14 (p. 179): The primes p\_j and the RVM prime levels p\_j\textasciicircum{}RVM, j<=1000. Here p\_j\textasciicircum{}RVM is a discrete version of inf x:pi\_RVM(x)>= j: the first point of the computational grid in u=log x (step 5.9*10\textasciicircum{}-6) at which pi\_RVM>= j, refined by linear interpolation from the previous grid point. It is not the inverse of a monotone function, a...}
\expandafter\gdef\csname CGTip@tab:pi-1000\endcsname{Table 15 (p. 184): The prime count pi(n), exact, and the synthetic prime count pi\_RVM(n), a real-valued function computed from truncated sums (not an integer count), at the integers n<=1000, with the difference pi\_RVM(n)-pi(n).}
\expandafter\gdef\csname CGTip@tab:mertens-1000\endcsname{Table 16 (p. 190): The Mobius function mu(n) and the Mertens function M(n), exact (sieve), and the critical RVM Mertens function M\_RVM\textasciicircum{}crit(n), approximate (the tapered explicit formula of program s05 with the residues computed from the product for k<=1000 and modelled beyond), at the integers n<=1000, with the difference M\_RVM\textasciicircum{}crit(n...}
\expandafter\gdef\csname CGTip@app:code\endcsname{Section B: Python programs (p. 196)}
\expandafter\gdef\csname CGTip@app:code-rvm-common\endcsname{Section B.1: rvm\_common.py (p. 196)}
\expandafter\gdef\csname CGTip@app:code-s01-clock-and-family\endcsname{Section B.2: s01\_clock\_and\_family.py (p. 200)}
\expandafter\gdef\csname CGTip@app:code-s02-transition\endcsname{Section B.3: s02\_transition.py (p. 203)}
\expandafter\gdef\csname CGTip@app:code-s03-exponential-sums\endcsname{Section B.4: s03\_exponential\_sums.py (p. 205)}
\expandafter\gdef\csname CGTip@app:code-s04-primes\endcsname{Section B.5: s04\_primes.py (p. 206)}
\expandafter\gdef\csname CGTip@app:code-s05-mertens\endcsname{Section B.6: s05\_mertens.py (p. 210)}
\expandafter\gdef\csname CGTip@app:code-s06-prime-ratios\endcsname{Section B.7: s06\_prime\_ratios.py (p. 215)}
\expandafter\gdef\csname CGTip@app:code-s07-product-identity\endcsname{Section B.8: s07\_product\_identity.py (p. 216)}
\expandafter\gdef\csname CGTip@app:code-s08-transport-moments\endcsname{Section B.9: s08\_transport\_moments.py (p. 217)}
\expandafter\gdef\csname CGTip@app:code-s09-clock-defect-split\endcsname{Section B.10: s09\_clock\_defect\_split.py (p. 219)}
\expandafter\gdef\csname CGTip@app:code-s10-clock-to-zeros\endcsname{Section B.11: s10\_clock\_to\_zeros.py (p. 220)}
\expandafter\gdef\csname CGTip@app:code-s11-rvm-z1\endcsname{Section B.12: s11\_rvm\_z1\_difference\_analysis.py (p. 222)}
\expandafter\gdef\csname CGTip@app:code-s12-multiscale\endcsname{Section B.13: s12\_multiscale\_jacobian.py (p. 227)}
\expandafter\gdef\csname CGTip@app:code-s13-bounded\endcsname{Section B.14: s13\_nested\_bounded\_endpoint.py (p. 230)}
\expandafter\gdef\csname CGTip@app:code-s14\endcsname{Section B.15: s14\_varpi\_bound.py (p. 237)}
\expandafter\gdef\csname CGTip@app:code-s15\endcsname{Section B.16: s15\_tables\_1000.py (p. 240)}
\expandafter\gdef\csname CGTip@app:code-s16\endcsname{Section B.17: s16\_walks\_3d\_1000.py (p. 245)}
\expandafter\gdef\csname CGTip@app:code-s17\endcsname{Section B.18: s17\_cramer\_centroids.py (p. 247)}
\expandafter\gdef\csname CGTip@app:output\endcsname{Section C: Console output of the programs (p. 249)}
\expandafter\gdef\csname CGTip@app:out-s01-clock-and-family\endcsname{Section C.1: s01\_clock\_and\_family (p. 249)}
\expandafter\gdef\csname CGTip@app:out-s01b-clock-zoom\endcsname{Section C.2: s01b\_clock\_zoom (p. 250)}
\expandafter\gdef\csname CGTip@app:out-s02-transition\endcsname{Section C.3: s02\_transition (p. 250)}
\expandafter\gdef\csname CGTip@app:out-s03-exponential-sums\endcsname{Section C.4: s03\_exponential\_sums (p. 250)}
\expandafter\gdef\csname CGTip@app:out-s04-primes\endcsname{Section C.5: s04\_primes (p. 250)}
\expandafter\gdef\csname CGTip@app:out-s05-mertens\endcsname{Section C.6: s05\_mertens (p. 252)}
\expandafter\gdef\csname CGTip@app:out-s06-prime-ratios\endcsname{Section C.7: s06\_prime\_ratios (p. 253)}
\expandafter\gdef\csname CGTip@app:out-s07-product-identity\endcsname{Section C.8: s07\_product\_identity (p. 254)}
\expandafter\gdef\csname CGTip@app:out-s08-transport-moments\endcsname{Section C.9: s08\_transport\_moments (p. 254)}
\expandafter\gdef\csname CGTip@app:out-s09-clock-defect-split\endcsname{Section C.10: s09\_clock\_defect\_split (p. 255)}
\expandafter\gdef\csname CGTip@app:out-s10-clock-to-zeros\endcsname{Section C.11: s10\_clock\_to\_zeros (p. 255)}
\expandafter\gdef\csname CGTip@app:out-s12-multiscale\endcsname{Section C.12: s12\_multiscale\_jacobian (p. 256)}
\expandafter\gdef\csname CGTip@app:out-s13-bounded\endcsname{Section C.13: s13\_nested\_bounded\_endpoint (p. 256)}
\expandafter\gdef\csname CGTip@app:output-s14\endcsname{Section C.14: s14\_varpi\_bound (p. 258)}
\expandafter\gdef\csname CGTip@app:output-s15\endcsname{Section C.15: s15\_tables\_1000 (p. 258)}
\expandafter\gdef\csname CGTip@app:output-s16\endcsname{Section C.16: s16\_walks\_3d\_1000 (p. 259)}
\expandafter\gdef\csname CGTip@app:output-s17\endcsname{Section C.17: s17\_cramer\_centroids (p. 259)}
\expandafter\gdef\csname CGTip@app:code-routeB\endcsname{Section D: Python programs of Section 18 (p. 264)}
\expandafter\gdef\csname CGTip@app:code-routeB-core\endcsname{Section D.1: core.py (p. 264)}
\expandafter\gdef\csname CGTip@app:code-routeB-zeros\endcsname{Section D.2: zeros.py (p. 265)}
\expandafter\gdef\csname CGTip@app:code-routeB-gonek-fast\endcsname{Section D.3: gonek\_fast.py (p. 265)}
\expandafter\gdef\csname CGTip@app:code-routeB-horizonG\endcsname{Section D.4: horizonG.py (p. 265)}
\expandafter\gdef\csname CGTip@app:code-routeB-cert7\endcsname{Section D.5: cert7.py (p. 266)}
\expandafter\gdef\csname CGTip@app:code-routeB-entry\endcsname{Section D.6: entry.py (p. 267)}
\expandafter\gdef\csname CGTip@app:code-routeB-entry2\endcsname{Section D.7: entry2.py (p. 268)}
\expandafter\gdef\csname CGTip@app:code-routeB-certV\endcsname{Section D.8: certV.py (p. 269)}
\expandafter\gdef\csname CGTip@app:code-routeB-laplace-check\endcsname{Section D.9: laplace\_check.py (p. 274)}
\expandafter\gdef\csname CGTip@app:code-routeB-diag\endcsname{Section D.10: diag\_cert7\_gamma.py (p. 274)}
\expandafter\gdef\csname CGTip@app:code-routeB-figs\endcsname{Section D.11: figs.py (p. 275)}
\expandafter\gdef\csname CGTip@app:code-routeB-summary\endcsname{Section D.12: results\_summary.py (p. 275)}
\expandafter\gdef\csname CGTip@app:out-routeB\endcsname{Section E: Console output of the programs of Section 18 (p. 277)}
\expandafter\gdef\csname CGTip@app:out-routeB-results\endcsname{Section E.1: RESULTS.txt (p. 277)}
\expandafter\gdef\csname CGTip@app:out-routeB-horizonG\endcsname{Section E.2: out\_horizonG.txt (p. 279)}
\expandafter\gdef\csname CGTip@app:out-routeB-cert7\endcsname{Section E.3: out\_cert7.txt (p. 280)}
\expandafter\gdef\csname CGTip@app:out-routeB-cert7-gamma\endcsname{Section E.4: out\_cert7\_gamma.txt (p. 280)}
\expandafter\gdef\csname CGTip@app:out-routeB-entry\endcsname{Section E.5: out\_entry.txt (p. 280)}
\expandafter\gdef\csname CGTip@app:out-routeB-entry2\endcsname{Section E.6: out\_entry2.txt (p. 281)}
\expandafter\gdef\csname CGTip@app:out-routeB-certV\endcsname{Section E.7: out\_certV.txt (p. 281)}
\expandafter\gdef\csname CGTip@app:out-routeB-certG\endcsname{Section E.8: out\_certG.txt (p. 282)}
\expandafter\gdef\csname CGTip@app:out-routeB-cert3\endcsname{Section E.9: out\_cert3.txt (p. 282)}
\expandafter\gdef\csname CGTip@app:out-routeB-gonek\endcsname{Section E.10: out\_gonek.txt and out\_gonek\_fast.txt (p. 283)}
\expandafter\gdef\csname CGTip@app:out-routeB-laplace-check\endcsname{Section E.11: out\_laplace\_check.txt (p. 283)}
\expandafter\gdef\csname CGTip@sec:appendix-endpoint\endcsname{Section F: Endpoint-only closure, interpolation tests, and autonomous residual questions (p. 284)}
\expandafter\gdef\csname CGTip@subsec:three-interpolations\endcsname{Section F.1: Three interpolations and their status (p. 284)}
\expandafter\gdef\csname CGTip@eq:additive-family\endcsname{Eq. (F.1), p. 284: G\_varepsilon(s):=(1-varepsilon)zeta\_RVM(s)+varepsilon zeta(s), 0<=varepsilon <=1.}
\expandafter\gdef\csname CGTip@prop:additive-level\endcsname{Proposition F.1 (The additive family is a level-set reformulation) (p. 284): Away from the zeros of zeta\_RVM, G\_varepsilon(s)=zeta\_RVM(s)((1-varepsilon)+varepsilon Q\_RVM(s)). Hence for 0<varepsilon<1 every zero of G\_varepsilon which is not a zero of zeta\_RVM lies on the level set Q\_RVM(s)=-(1-varepsilon)/(varepsilon). On the critical line e\textasciicircum{}ivartheta(t) G\_varepsilon((1)/(2)+it) is real. At v...}
\expandafter\gdef\csname CGTip@eq:additive-level\endcsname{Eq. (F.2), p. 284: G\_varepsilon(s)=zeta\_RVM(s)((1-varepsilon)+varepsilon Q\_RVM(s)).}
\expandafter\gdef\csname CGTip@eq:Q-negative-level\endcsname{Eq. (F.3), p. 284: Q\_RVM(s)=-(1-varepsilon)/(varepsilon).}
\expandafter\gdef\csname CGTip@subsec:first-transition\endcsname{Section F.2: The clock as the first transition term (p. 284)}
\expandafter\gdef\csname CGTip@eq:rvm-first-rs\endcsname{Eq. (F.4), p. 284: zeta\_RVM(s)=varpi\_RVM(1+chi(s))(1+O(|t|\textasciicircum{}-1)).}
\expandafter\gdef\csname CGTip@eq:first-rs-hardy\endcsname{Eq. (F.5), p. 284: e\textasciicircum{}ivartheta(t)(1+chi((1)/(2)+it))=2cos vartheta(t).}
\expandafter\gdef\csname CGTip@eq:finite-additive\endcsname{Eq. (F.6), p. 285: P\_varepsilon,M(s):=1+varepsilon sum\_2<= n<= Mn\textasciicircum{}-s, R\_varepsilon,M(s):= P\_varepsilon,M(s)+chi(s) P\_varepsilon,M(1-s),}
\expandafter\gdef\csname CGTip@eq:finite-additive-hardy\endcsname{Eq. (F.7), p. 285: e\textasciicircum{}ivartheta(t) R\_varepsilon,M((1)/(2)+it) =2Re (e\textasciicircum{}ivartheta(t) P\_varepsilon,M((1)/(2)+it)).}
\expandafter\gdef\csname CGTip@eq:dominance\endcsname{Eq. (F.8), p. 285: varepsilon sum\_2<= n<= Mn\textasciicircum{}-1/2<1,}
\expandafter\gdef\csname CGTip@subsec:live\endcsname{Section F.3: What remains genuinely live (p. 285)}
\expandafter\gdef\csname CGTip@eq:clock-injectivity\endcsname{Eq. (F.9), p. 285: H(z)=nu\_8+sum\_k>=1dot kappa\_k((tau\_k)/(z\textasciicircum{}2+tau\_k\textasciicircum{}2)-(1)/(tau\_k)) =sum\_n>=2(dot a\_n)/(log n) n\textasciicircum{}-z}
\expandafter\gdef\csname CGTip@subsec:rung-apriori\endcsname{Section F.4: A priori bounded displacement for exact rung solutions (p. 285)}
\expandafter\gdef\csname CGTip@prop:rung-apriori\endcsname{Proposition F.2 (A priori displacement bound for an exact rung solution) (p. 286): Fix tin[0,1]. Let kappa=(kappa\_k)\_k>=1 satisfy 0<kappa\_1<=kappa\_2<=... and kappa\_k->inf, and let (a\_n)\_n>=2 be real masses satisfying, in D'(0,inf), sum\_k>=1cos(kappa\_k u) -(1- t)sum\_k>=1cos(tau\_k u)- t k(u) =-sum\_n>=2 a\_n delta(u-log n). Assume that for every epsilon\_0>0 the local count and the masses satisfy N\_kap...}
\expandafter\gdef\csname CGTip@eq:rung-apriori-system\endcsname{Eq. (F.10), p. 286: sum\_k>=1cos(kappa\_k u) -(1- t)sum\_k>=1cos(tau\_k u)- t k(u) =-sum\_n>=2 a\_n delta(u-log n).}
\expandafter\gdef\csname CGTip@eq:rung-apriori-density\endcsname{Eq. (F.11), p. 286: N\_kappa(T+1)-N\_kappa(T)<<\_epsilon\_0(1+T)\textasciicircum{}epsilon\_0(T>=0),}
\expandafter\gdef\csname CGTip@eq:rung-apriori-mass\endcsname{Eq. (F.12), p. 286: sum\_X<n<=2X| a\_n|\textasciicircum{}2<<\_epsilon\_0X\textasciicircum{}1+epsilon\_0(X>=1).}
\expandafter\gdef\csname CGTip@eq:rung-apriori-b\endcsname{Eq. (F.13), p. 286: b:=lim\_sigma ->+inf[sum\_k>=1(sigma-(1)/(2))/((sigma-(1)/(2))\textasciicircum{}2+kappa\_k\textasciicircum{}2)-(1- t)sum\_k>=1(sigma-(1)/(2))/((sigma-(1)/(2))\textasciicircum{}2+tau\_k\textasciicircum{}2)- t (G')/(2G)(sigma)]}
\expandafter\gdef\csname CGTip@eq:rung-apriori-count\endcsname{Eq. (F.14), p. 286: |N\_kappa(T)-n\_RVM(T)-(2 b)/(pi)T|<= C\textasciicircum{}*log(2+T)(T>=0)}
\expandafter\gdef\csname CGTip@eq:rung-apriori-bound\endcsname{Eq. (F.15), p. 286: |kappa\_k-tau\_k+(4 b kappa\_k)/(log(s\_k/2pi))|<= C\textasciicircum{}*, |kappa\_k-tau\textasciicircum{}(b)\_k|<= C\textasciicircum{}*,}
\expandafter\gdef\csname CGTip@eq:rung-laplace-S\endcsname{Eq. (F.16), p. 287: L S\_kappa(s)=-z\textasciicircum{}2sum\_k(z)/(kappa\_k\textasciicircum{}2(z\textasciicircum{}2+kappa\_k\textasciicircum{}2))+zsum\_k(1)/(kappa\_k\textasciicircum{}2)=sum\_k>=1(z)/(z\textasciicircum{}2+kappa\_k\textasciicircum{}2)=(1)/(2) (P\_kappa')/(P\_kappa)(s)(Re s>(1)/(2)):}
\expandafter\gdef\csname CGTip@eq:rung-E\endcsname{Eq. (F.17), p. 287: E:= S\_kappa-(1- t) S\_tau- t K+ A}
\expandafter\gdef\csname CGTip@eq:rung-laplace-identity\endcsname{Eq. (F.18), p. 287: (1)/(2)(P\_kappa')/(P\_kappa)(s)-(1- t) (1)/(2)((P\textasciicircum{}crit\_RVM)')/(P\textasciicircum{}crit\_RVM)(s)- t (1)/(2)(G')/(G)(s)+sum\_n>=2 a\_nn\textasciicircum{}(1)/(2)-s=p(s-(1)/(2))(Re s>(3)/(2)).}
\expandafter\gdef\csname CGTip@eq:rung-Y\endcsname{Eq. (F.19), p. 287: Y\_ t(s)=exp (sum\_n>=2(2 a\_n)/(log n) n\textasciicircum{}1/2-s)}
\expandafter\gdef\csname CGTip@eq:rung-three-factor\endcsname{Eq. (F.20), p. 288: P\_kappa(s) =C\_ t(P\_RVM\textasciicircum{}crit(s))\textasciicircum{}1- t G(s)\textasciicircum{} t Y\_ t(s) e\textasciicircum{}2 b(s-2) =C\_ t G(s) zeta\_RVM(s)\textasciicircum{}1- t Y\_ t(s) e\textasciicircum{}2 b(s-2)(Re s>(3)/(2)),}
\expandafter\gdef\csname CGTip@eq:rung-B\endcsname{Eq. (F.21), p. 288: sup\_Re s>=2|log Y\_ t(s)|<= B\_2:=sum\_n>=2(2| a\_n|)/(log n)n\textasciicircum{}-3/2<inf,}
\expandafter\gdef\csname CGTip@eq:rung-line\endcsname{Eq. (F.22), p. 288: | W\_kappa(s)|<= K\_rm E:=|log C\_ t|+ c\_RVM+B\_2+4| b|(2<=Re s<=4),}
\expandafter\gdef\csname CGTip@eq:rung-strip\endcsname{Eq. (F.23), p. 288: W\_kappa(s)<= K\_rm E''+(3)/(2)log|s+2|(sin S), K\_rm E'':=max (K\_rm E+(3)/(4)log 2, 5sum\_kkappa\_k\textasciicircum{}-2+ c\_7).}
\expandafter\gdef\csname CGTip@eq:rung-nT\endcsname{Eq. (F.24), p. 289: n\_T<=(K\_rm E+ K\_rm E''+(5)/(2)log(T+8)+2 c\_5)/(log(4/3)),}
\expandafter\gdef\csname CGTip@eq:rung-count\endcsname{Eq. (F.25), p. 289: |N\_kappa(T)-n\_RVM(T)-(2 b)/(pi) T| <= n\_T+1+(2 c\_RVM+2B\_2+ c\_6)/(pi)+(1)/(T)+(1)/(2)(T>= T\_0, Tnot in kappa\_k),}
\expandafter\gdef\csname CGTip@eq:rung-quantile\endcsname{Eq. (F.26), p. 290: | F\_ b(kappa\_k)-k|<= C\textasciicircum{}*log(2+kappa\_k)+2(k>=1).}
\expandafter\gdef\csname CGTip@cor:rung-apriori-b0\endcsname{Corollary F.3 (Exact solutions without mass at the origin) (p. 290): Under the hypotheses of Proposition F.2, if the extension of the left side of (F.10) constructed in the proof has no mass at u=0, that is, b=0, then |N\_kappa(T)-n\_RVM(T)|<= C\textasciicircum{}*log(2+T) and |kappa\_k-tau\_k|<= C\textasciicircum{}* for all T>=0 and k>=1, with C\textasciicircum{}* depending only on the implied constants of (F.11)--(F.12) for the fixed ep...}
\expandafter\gdef\csname CGTip@rem:rung-apriori-scope\endcsname{Remark F.4 (p. 290): Proposition F.2 is an a priori theorem, not an existence theorem. The exact-atomic interpolation has no established interior solution for 0< t<1; the proposition only says that any exact solution satisfying its hypotheses lies in a uniform drift-corrected class, with ordinates within C\textasciicircum{}* of the drift-corrected clock...}
\expandafter\gdef\csname CGTip@subsec:finite-diagnostics\endcsname{Section F.5: Additional finite-section diagnostics (p. 291)}
\expandafter\gdef\csname CGTip@eq:bounded-section\endcsname{Eq. (F.27), p. 291: |kappa\_k-tau\_k|<=1.2, |a\_n|<=0.5,}
\expandafter\gdef\csname CGTip@fig:nested-bounded\endcsname{Figure 24 (p. 292): Nested bounded endpoint sections N=50,70,100: root mean square residual of each section on every permanent test block introduced so far; the horizontal axis is the size N of the section that introduced the block. The legend gives, for each section, the number of equations and blocks, the numbers of ordinates and mas...}
\expandafter\gdef\csname CGTip@fig:bounded-masses\endcsname{Figure 25 (p. 293): Recovered free masses a\_n, n<=15, of the single bounded sections N=20,34,50, each solved alone from the clock. All masses with n<=15 are shown, and none of them is at the box |a\_n|=0.5 (dashed lines; a mass at the box would be circled in red). The masses at the box are edge masses, listed in the legend: n=20 for N=2...}
\makeatother


\hypersetup{
  colorlinks=false,
  pdfborder={0 0 1},
  linkbordercolor={1 0 0},
  citebordercolor={0 0.3 1},
  urlbordercolor={0 0.6 0.6},
  filebordercolor={0 0.6 0.6},
  linktoc=all,
  bookmarksnumbered=true,
  bookmarksopen=true,
  bookmarksopenlevel=1,
  pdfstartview=FitH,
  pdftitle={The midpoint Riemann-von Mangoldt synthetic zeta system},
  pdfauthor={Rotkiv Ayn\"asac},
  pdfsubject={Midpoint RVM zeta function, exponential sums over the zeta ordinates and Cramer's function, window forms of the Riemann hypothesis, the lambda-family to the endpoint Z1, Hamburger closure, and the Mertens bound},
  pdfkeywords={Riemann zeta function, Riemann-von Mangoldt clock, exponential sums, Cramer function, Riemann hypothesis, Hadamard product, Hamburger theorem, Mertens function}
}
% Back-references in the bibliography (placed after the DOI).
\renewcommand*{\backref}[1]{}
\renewcommand*{\backrefalt}[4]{%
  \ifcase #1 %
    \relax
  \or
    \ (Cited on p.~#2.)%
  \else
    \ (Cited on pp.~#2.)%
  \fi}

% In \cite[note]{key} the note (e.g. a section number) is linked too.
\makeatletter
\let\CG@origcite\cite
\renewcommand*{\cite}[2][]{\def\CG@citekey{#2}%
  \ifx\relax#1\relax\CG@origcite{#2}\else\CG@origcite[#1]{#2}\fi}
\def\@cite#1#2{{\@citestyle[\citeform{#1}\if@tempswa, \hyper@link{cite}{cite.\CG@citekey}{#2}\fi]}}
\makeatother

\geometry{margin=1in}
\numberwithin{equation}{section}
\setcounter{tocdepth}{3}
\setcounter{secnumdepth}{3}
\setlength{\emergencystretch}{1.5em}

\theoremstyle{plain}
\newtheorem{theorem}{Theorem}[section]
\newtheorem{proposition}[theorem]{Proposition}
\newtheorem{lemma}[theorem]{Lemma}
\newtheorem{corollary}[theorem]{Corollary}
\newtheorem{conjecture}[theorem]{Conjecture}
\theoremstyle{definition}
\newtheorem{definition}[theorem]{Definition}
\newtheorem{hypothesis}[theorem]{Hypothesis}
\newtheorem{heuristic}[theorem]{Heuristic}
\newtheorem{experiment}[theorem]{Experiment}
\theoremstyle{remark}
\newtheorem{remark}[theorem]{Remark}
\newtheorem{warning}[theorem]{Warning}

\newcommand{\R}{\mathbb R}
\newcommand{\C}{\mathbb C}
\newcommand{\N}{\mathbb N}
\newcommand{\Z}{\mathbb Z}
\newcommand{\eps}{\varepsilon}
\newcommand{\li}{\operatorname{li}}
\newcommand{\Log}{\operatorname{Log}}
\newcommand{\Ei}{\operatorname{Ei}}
\newcommand{\Res}{\operatorname*{Res}}
\newcommand{\Mell}{\ast_{\rm Mel}}
\newcommand{\RVM}{\mathrm{RVM}}
\newcommand{\crit}{\mathrm{crit}}
\newcommand{\pts}{\mathrm{pts}}
\newcommand{\pathc}{\mathrm{path}}
\newcommand{\NR}{\mathcal N_{\RVM}}
\newcommand{\vpi}{\varpi_{\RVM}}
\newcommand{\Zo}{Z1}
\providecommand{\NoCaseChange}[1]{\lowercase{#1}}

\begin{document}

\title[The midpoint RVM synthetic zeta system]{The midpoint Riemann--von Mangoldt synthetic zeta system:\\
canonical products, exponential sums, prime counts, Mertens transforms,\\
and the transformation to $\zeta_{Z1}$}
\author{Rotkiv Ayn\"asac\\[2pt]\mbox{}}
\date{15 October 2026 -- version 25}

% The abstract is typeset after \maketitle as ordinary text so that it may run over several pages.
\makeatletter
\renewenvironment{abstract}{\par\begingroup\small\noindent{\scshape Abstract.}\enspace\ignorespaces}{\par\endgroup\medskip}
\makeatother


\maketitle
\markboth{ROTKIV AYN\"ASAC}{THE MIDPOINT RVM SYNTHETIC ZETA SYSTEM}

\begin{abstract}
We study the zeta-like system generated by the midpoint inverse of the smooth Riemann--von Mangoldt counting law.  The midpoint ordinates $\tau_k$, defined by $\frac{\tau_k}{2\pi}\log\frac{\tau_k}{2\pi e}+\frac78=k-\frac12$ and given by the Lambert formula $\tau_k=2\pi(k-\frac{11}8)/W_0((k-\frac{11}8)/e)$, are those of Fran\c{c}a and LeClair; the constructions built on them are, to the author's knowledge, new.  The synthetic zeta function $\zeta_{\RVM}$ has all its nontrivial zeros at $\frac12\pm i\tau_k$.  We prove normal convergence of its canonical product, a uniform logarithmic bound for the exponential sums of the $\tau_k$, the explicit-formula prime count $\pi_{\RVM}$ with the unconditional estimate $\pi_{\RVM}(x)-\li(x)=O(\sqrt x(\log\log x)^2/\log x)$, the transition theorem $\zeta_{\RVM}=\varpi_{\RVM}(1+w)(1+O(1/|t|))$ with $\varpi_{\RVM}<0.9796$, and the synthetic Mertens bound $M_{\RVM}^{\crit}(x)=L+O_\eps(x^{1/2+\eps})$, $L=1/\varpi_{\RVM}$.

\emph{Exponential sums over the ordinates.}  We compare $\sum_ke^{iw\gamma_k}$, over the ordinates $\gamma_k$ of the zeros of $\zeta$, with Cram\'er's function $V_+(w)=\sum_{\Im\rho>0}e^{\rho w}$, for which we give a representation that involves no sum over the zeros, valid on a half-plane containing $(0,\infty)$.  This yields equivalent forms of the Riemann hypothesis written with the ordinates alone, which we have not found in the literature consulted: RH holds if and only if $\sum_ke^{-y_0\gamma_k}=e^{-iy_0/2}V_+(iy_0)$ for one $y_0\in(0,4\pi]$; and RH holds if and only if $\sum_k\widehat{\mathsf f}(\gamma_k)=\int\mathsf f(u)e^{-u/2}V_+(u)\,du$ for every test function $\mathsf f$ supported in the window $(\log2,\log3)$, where the explicit formula has no prime atom, that is, if and only if the triangular centroids $\sum_{k\le N}(1-\frac{k-1}N)e^{iu\gamma_k}$ converge there to $e^{-u/2}V_+(u)$ in the sense of distributions.  Without hypothesis the defect is an explicit absolutely convergent sum over the zeros off the line.  The rung functional of the endpoint is that same sum: $(\log x)(dJ^{\rm M}_{Z1}-dJ)/dx=\frac2{\sqrt x}\sum_k(\cosh(d_k\log x)-1)\cos(\gamma_k\log x)$, $d_k=\Re\rho_k-\frac12$, so that the real part of the window identity is the single cell $(2,3)$, and the whole ladder of cells, equivalently RH, is that cell together with a sine identity on the window.  If only finitely many zeros were off the line, $J^{\rm M}_{Z1}$ would be constant on no cell $(n,n+1)$.  A Frullani identity writes $\log L+\log\varpi_{Z1}$ as $\lim_K2\int_0^\infty u^{-1}e^{-u/2}\sum_{k\le K}(\cos\gamma_ku-\cos\tau_ku)\,du$.  Under RH the triangular centroid of $e^{i\gamma_k}$ converges to $e^{-1/2}V_+(1)$, and that of $e^{i\tau_k}$ converges without hypothesis; the two limits are, numerically, $0.77689426079\ldots+0.76785566886\ldots\,i$, whose real part is $\frac1{2\sqrt e}(e-\frac1{e^2-1})$, and $0.041569+1.757541\,i$.

\emph{The family and its endpoint.}  The ordinates move along $\gamma_k(\lambda)=(1-\lambda)\gamma_k+\lambda\tau_k$; every member $\zeta_\lambda$ has its nontrivial zeros on $\Re s=\frac12$ by construction, $\zeta_1=\zeta_{\RVM}$, and $\zeta_0=\zeta_{Z1}$.  The final step is the identity $\zeta_{Z1}=\zeta$, which by Hamburger's theorem is equivalent to the condition (D) that $\zeta_{Z1}$ is an absolutely convergent Dirichlet series in $\Re s>1$, and to RH.  We prove that $\varpi_{Z1}=\lim_{\sigma\to+\infty}\zeta_{Z1}(\sigma)=\varpi_{\RVM}\lim_K\prod_{k\le K}\frac{1/4+\tau_k^2}{1/4+\gamma_k^2}\ge1$, with equality exactly under (D); that $\log\varpi_\lambda$ is strictly convex in $\lambda$, and that if RH fails the family crosses $\varpi_\lambda=1$ at exactly one $\lambda^*$, and $\lambda^*<2\cdot10^{-11}$ by the verified zeros; that every member admits an atomization by smooth cells which is an absolutely convergent Dirichlet series; that (D) is equivalent to the vanishing of every cell defect, to a positivity criterion, to Hamburger's theorem for the pair (continuous, atomized) in several forms, and to the existence of a list in a compact space of ordinate lists whose continuous and atomized zeta functions coincide; and, for completeness, that $\zeta_{Z1}=\zeta$ implies $M(x)=O_\eps(x^{1/2+\eps})$.  A last section treats route B for finite sections of the rung system: an exact solution whose ordinates stay at bounded distance from the clock, with arbitrary real masses, is the zeta list with the von Mangoldt masses, so that route B is equivalent to RH; since RH is verified up to height $3\cdot10^{12}$, no finite stage decides it; on one window, RH is equivalent to the existence of a list of ordinates reproducing Cram\'er's function, and to the complete monotonicity of $e^{-iy/2}V_+(iy)$; and fifteen finite sections up to height $487$ are solved exactly, with a proof in ball arithmetic.

\emph{Tables.}  We tabulate and draw the first $1000$ values of $\tau_k$ against $\gamma_k$, of the RVM prime levels $p_j^{\RVM}$ against $p_j$, of $\pi_{\RVM}(n)$ against $\pi(n)$ and of $M_{\RVM}^{\crit}(n)$ against $M(n)$, and the walks of $e^{i\gamma_k}$, $e^{i\tau_k}$ and $e^{i(\gamma_k+\tau_k)/2}$ in three dimensions, computed from $2\,001\,052$ zeros.

\emph{What is proved and what is not.}  The mathematical statements of the first three paragraphs are theorems of the paper, with the hypotheses (RH or (D)) indicated where they are used; the decimal values of the two centroid limits are computed, the bound on $\lambda^*$ uses the zeros verified by Platt and Trudgian, and the paragraph on tables describes computations.  The identity $\zeta_{Z1}=\zeta$, equivalently (D), equivalently RH, equivalently each of the imaginary-frequency, window and cell criteria above, is not proved.  The computations are experiments and are not used in any proof, except for the interval evaluation of a few explicit inequalities in the proof of $\varpi_{\RVM}<1$ and the ball-arithmetic verification of the finite sections of Theorem~\CGref{thm:rb-verified}.
\end{abstract}

\tableofcontents

\section*{Notation}
\label{sec:notation}

Every symbol below has a single meaning throughout the paper.  The last column links to the defining equation; hovering over the link shows the formula.  Symbols introduced inside a single proof (indices $j,k_0,m$, auxiliary variables $v,y$, thresholds $t_0,t_1,t_2$, radii $r_1,r_2,r_3$) are local to that proof and do not appear elsewhere; so are the following local symbols, which coincide in shape with global ones but never with their meaning: $E_k$ and $T_m$ in the proofs of Theorems~\CGref{thm:exp-sum} and~\CGref{thm:unit-cell}, $\phi$ in Lemmas~\CGref{lem:A-RVM} and~\CGref{lem:TR-clock}, $K_m^\pm$ and $r$ in Section~\CGref{sec:transition}, $J$, $S$, $C_0$, $E$ and $A$ in the proofs of Section~\CGref{sec:atomized}, $U$ as a window length in Subsection~\CGref{subsec:multiscale-jacobian}, $\mathsf X_{\mathfrak P}$ for the prime diagonal operator, $\mathfrak g_1$ in Subsection~\CGref{subsec:rb-application}, and $c_1$ in Propositions~\CGref{prop:Mellin-Mertens} and~\CGref{prop:one-number}.

\medskip
{\small
\begin{longtable}{>{\raggedright\arraybackslash}p{0.27\textwidth}>{\raggedright\arraybackslash}p{0.53\textwidth}>{\raggedright\arraybackslash}p{0.12\textwidth}}
\toprule
Symbol & Meaning & Defined\\
\midrule
\endfirsthead
\toprule
Symbol & Meaning & Defined\\
\midrule
\endhead
\bottomrule
\endfoot
\multicolumn{3}{l}{\emph{Variables and the Riemann zeta function}}\\
$s=\sigma+it$ & complex variable, $\sigma=\Re s$, $t=\Im s$ & --\\
$z=s-\tfrac12=\sigma'+it$ & centred variable, $\sigma'=\sigma-\tfrac12$ & \CGeqref{eq:transition-zw}\\
$q(s)=s(s-1)$ & symmetric quadratic, $q(1-s)=q(s)$ & \CGeqref{eq:q-def}\\
$\zeta$, $\rho$ & Riemann zeta function; $\rho$ runs over its nontrivial zeros with multiplicity & \S\CGref{sec:clock}\\
$\gamma_k$ & $k$th positive ordinate of the nontrivial zeros of $\zeta$, with multiplicity, nondecreasing & \CGeqref{eq:zero-count}\\
$N_\zeta(T)$, $N_\zeta^{*}(T)$ & zero count up to height $T$; its symmetric version & \CGeqref{eq:zero-count}, \CGeqref{eq:symmetric-count}\\
$S(T)$ & remainder in the Riemann--von Mangoldt formula & \CGeqref{eq:RvM-formula}\\
$G(s)$ & $s(s-1)\Gamma(s/2)\pi^{-s/2}$ & \CGeqref{eq:Gdef}\\
$P_\zeta=2\xi$ & completed product $G\zeta$ & \CGeqref{eq:Pzeta}\\
$\mu$, $M(x)$ & M\"obius function; Mertens function $\sum_{n\le x}\mu(n)$ & \CGeqref{eq:mertens-def}\\
$\pi(x)$, $\li$ & prime-counting function; $\li(x):=\operatorname{PV}\!\int_0^x dt/\log t$ & --\\
$\chi(s)$, $\vartheta$ & functional-equation factor; Riemann--Siegel theta function & \CGeqref{eq:chi}, \CGeqref{eq:theta-def}\\
\midrule
\multicolumn{3}{l}{\emph{The midpoint clock}}\\
$\NR(T)$ & midpoint Riemann--von Mangoldt clock & \CGeqref{eq:clock}\\
$\tau_k$ & midpoint ordinates: the unique solutions in $[2\pi,\infty)$ of $\NR(\tau_k)=k-\tfrac12$ & \CGeqref{eq:midpoint-condition}\\
$W_0$ & principal branch of Lambert's $W$-function & \CGeqref{eq:lambert}\\
$\mathcal T$ & inverse function of the clock & \CGeqref{eq:inverse-clock}\\
$\ell_k$ & $\log(\tau_k/2\pi)$ & text after \CGeqref{eq:inverse-clock}\\
$n_{\RVM}(H)$ & number of $\tau_k\le H$ & \CGeqref{eq:n-RVM}\\
$\eta_k$, $C_\eta$ & $\eta_k=\tau_k-\gamma_k$; a bound $|\eta_k|\le C_\eta$ & \CGeqref{eq:eta-def}, \CGeqref{eq:bounded-displacement}\\
\midrule
\multicolumn{3}{l}{\emph{The synthetic zeta function and its sums}}\\
$P_{\RVM}^{\crit}$, $\xi_{\RVM}$, $\zeta_{\RVM}$ & canonical product, completed and synthetic zeta & \CGeqref{eq:Pcrit}, \CGeqref{eq:zetaRVM}\\
$\rho_k^{\RVM}$ & synthetic zero $\tfrac12+i\tau_k$ & \CGeqref{eq:synthetic-divisor}\\
$S_{\RVM}(N;u)$, $C_{\RVM}^{\pts}$, $C_{\RVM}^{\pathc}$ & exponential sum; point and path centroids & \CGeqref{eq:S-RVM}, \CGeqref{eq:path-centroid}\\
$u_1,u_2$ & a frequency interval $[u_1,u_2]\subset(0,\infty)$ & \CGeqref{eq:exp-sum-bound}\\
$M_H(x)$, $M_H^{\RVM}(x)$ & finite-height phase averages (not Mertens functions) & \CGeqref{eq:MH-genuine}, \CGeqref{eq:MH-RVM}\\
$\Ei$ & exponential integral on $\Im w\ne0$ & \CGeqref{eq:Ei-def}\\
$J_{\RVM}$, $J_{\RVM}^{*}$, $\pi_{\RVM}$ & explicit-formula prime-power count; truncation; synthetic prime count (real-valued, not an integer count) & \CGeqref{eq:J-RVM}, \CGeqref{eq:J-star}, \CGeqref{eq:pi-RVM}\\
$\mathcal A_{\RVM}(u)$, $\ell(t)$ & weighted clock sum $\sum_ke^{iu\tau_k}/\tau_k$; $\log(t/2\pi)$ & \CGeqref{eq:A-RVM}\\
$\mathfrak s_1$; $\mathfrak e$, $\mathfrak e_1$ & primitive of the midpoint sawtooth; $N_\zeta-\NR$ and its primitive & \CGeqref{eq:saw-primitive}, \CGeqref{eq:frak-e}\\
$\varpi_{Z1}$ & value of $\zeta_{Z1}$ at $+\infty$ & \CGeqref{eq:varpi-Z1}\\
$J_X^{\rm M}$, $J_X$, $C_X$ & Mellin and explicit-formula prime-power counts of $\zeta_X$, $X\in\{\RVM,Z1\}$; their difference & \CGeqref{eq:mellin-count}, \CGeqref{eq:J-X}\\
$\mathsf b_n$ & jumps of $J_{Z1}^{\rm M}$ at the integers & Thm.~\CGref{thm:cell-form}\\
$\varpi_\lambda$; $\Upsilon(\lambda)$ & value of $\zeta_\lambda$ at $+\infty$; flow at infinity & \CGeqref{eq:varpi-lambda}, \CGeqref{eq:Upsilon}\\
$\varphi_n$, $\tilde c_n(\lambda)$, $\tilde a_N(\lambda)$ & smooth cells; atomized masses and coefficients & \CGeqref{eq:atomized-family}, \CGeqref{eq:atomized-coefficients}\\
$\tilde\zeta_{\Z,\lambda}$, $\tilde h_\lambda$ & atomized family; joining factor & \CGeqref{eq:atomized-family}, \CGeqref{eq:joining-factor}\\
$\mathfrak a_s$, $\mathfrak h_n(\lambda;s)$, $\mathfrak L$ & atomization kernel; cell defects; $\int_1^x(1-y^{-1})dy/\log y$ & \CGeqref{eq:atomization-kernel}\\
$V_n$, $W_n$ & ordinate kernels of the cell $n$ & \CGeqref{eq:VW}\\
$\tilde\mu_\lambda$, $\tilde M_\lambda$ & atomized M\"obius coefficients; summatory function & \CGeqref{eq:atomized-mobius}\\
$B_{Z1}$ & Mellin form of the endpoint Mertens function & \CGeqref{eq:Mellin-Mertens}\\
$c_n^{\RVM}$, $\zeta_{\Z,\RVM}$ & integer weights; weighted Euler product & \CGeqref{eq:cn}, \CGeqref{eq:weighted-zeta}\\
$J_{\Z,\RVM}$, $J_{\RVM}^{*,\Z}$, $h_{\RVM}$ & integer-step counts; Mellin correction & \CGeqref{eq:JZ}, \CGeqref{eq:J-star-Z}, \CGeqref{eq:hRVM}\\
$\theta$ & unit-cell step, $0<\theta\le1$ & \CGeqref{eq:unit-cell}\\
$p_j^{\RVM}$ & real RVM prime levels $\pi_{\RVM}^{-1}(j)$ on monotone branches & \CGeqref{eq:intro-real-primes}\\
$\mathcal M(n)$, $\varsigma_n$, $\mu_h$, $M_h$ & target samples of $M_{\RVM}^{\crit}$; signs; ternary walk and its sum (companion object; $h$ is a label) & \CGeqref{eq:intro-muh}, \CGeqref{eq:intro-quantizer}\\
$a_N^{\Z}$ & Dirichlet coefficients of $\zeta_{\Z,\RVM}$ & \CGeqref{eq:ZRVM-Dirichlet}\\
$C_J$ & Lipschitz constant of $J_{\RVM}$ & \CGeqref{eq:unit-cell}\\
$\mathfrak d_n^{\RVM}$ & prime defect $\mathbf 1_{\rm prime}(n)-c_n^{\RVM}$ & \CGeqref{eq:prime-defect}\\
\midrule
\multicolumn{3}{l}{\emph{Mertens transforms and the transition theorem}}\\
$F_{\RVM}=\widehat M_{\RVM}^{\crit}$ & $1/\zeta_{\RVM}$ & \CGeqref{eq:F-RVM}\\
$c$, $U$ & Perron abscissa; truncation height & \CGeqref{eq:Mcrit-symmetric}\\
$M_{\RVM}^{\crit}$, $M_{\RVM,G}^{\crit}$ & symmetric Perron inverse; Gaussian-smoothed inverse & \CGeqref{eq:Mcrit-symmetric}, \CGeqref{eq:MG}\\
$\tau_*$, $w(s)$ & transition point $-iz=t-i\sigma'$; transition variable $e^{-2\pi i\NR(\tau_*)}$ & \CGeqref{eq:transition-zw}\\
$B$; $\sigma'_1,\sigma'_2$ & upper bound for $\sigma'$; a range $[\sigma'_1,\sigma'_2]\subset(0,\tfrac12)$ & \CGeqref{eq:transition-w}, \CGeqref{eq:transition-chi}\\
$\mathcal E_0,\mathcal E$ & error factors of the transition theorem & \CGeqref{eq:transition-w}, \CGeqref{eq:transition-chi}\\
$\varpi_{\RVM}$, $L$ & transition constant, $\varpi_{\RVM}<0.9796$ (Theorem~\CGref{thm:varpi-less-one}), floating-point value $0.9619415789$; $L=1/\varpi_{\RVM}>1$ & \CGeqref{eq:varpi-def}, \CGeqref{eq:vertical-limit}\\
$\mathscr W$, $\mathscr I_0,\mathscr I_n$, $\mathfrak M$, $\mathbf B_2,\mathbf B_3$ & weight $d\tau/\tau=\mathscr W(u)\,du$ of the clock; pieces of $\nu_4$; lower bound in the proof of Theorem~\CGref{thm:varpi-less-one}; Bernoulli polynomials & \CGeqref{eq:nu4-weight}, \CGeqref{eq:nu4-bound}\\
$f,\Lambda,\Lambda_{\rm sm},\Lambda_{\rm saw}$ & logarithm of the product and its smooth and sawtooth parts (used in the transition and related clock estimates) & \CGeqref{eq:transition-product}, \CGeqref{eq:transition-split}\\
$\psi,\psi_r,\Phi_r,\nu_1,\dots,\nu_4$ & sawtooth, Abel means, ray remainder, real constants (used in the transition and related clock estimates; $\nu_3$ in closed form \CGeqref{eq:nu3-closed}, $\nu_4$ by \CGeqref{eq:nu4-closed}) & \CGeqref{eq:sawtooth-count}, \CGeqref{eq:psi-r}\\
$\alpha_c$ & $\min(c-\tfrac12,1)$ & \CGeqref{eq:FminusL-all}\\
\midrule
\multicolumn{3}{l}{\emph{Comparison with $\zeta$}}\\
$\mathcal Q_{\RVM}$, $\mathcal Q_{\RVM}^{\vee}$ & quotient $\zeta/\zeta_{\RVM}$; its inverse Mellin transform & \CGeqref{eq:Q-RVM}, \CGeqref{eq:M-convolution}\\
\midrule
\multicolumn{3}{l}{\emph{The $\lambda$-family and the endpoint $\zeta_{Z1}$}}\\
$\lambda$ & family parameter; $\lambda=1$: $\zeta_{\RVM}$, $\lambda=0$: $\zeta_{Z1}$ & \CGeqref{eq:gamma-lambda}\\
$\gamma_k(\lambda)$ & moving ordinates $(1-\lambda)\gamma_k+\lambda\tau_k$ & \CGeqref{eq:gamma-lambda}\\
$P_\lambda$, $\zeta_\lambda$, $F_\lambda$ & family product, family zeta, $1/\zeta_\lambda$ & \CGeqref{eq:P-lambda}, \CGeqref{eq:zeta-lambda}\\
$A_{k,\lambda}$ & $\tfrac14+\gamma_k(\lambda)^2$ & \CGeqref{eq:generator}\\
$\mathcal V_\lambda$ & infinitesimal generator $\partial_\lambda\log P_\lambda$ & \CGeqref{eq:generator}\\
$\Pi_{\rm tr}$ & total multiplier $P_{Z1}/P_{\RVM}^{\crit}$ & \CGeqref{eq:transport-multiplier}\\
$\rho_{k,\lambda}$, $\mathfrak r_{k,\lambda}$ & moving zero $\tfrac12+i\gamma_k(\lambda)$; residue of $1/(s\zeta_\lambda)$ & \CGeqref{eq:residue-transport}\\
$P_{Z1}$, $\zeta_{Z1}$, $F_{Z1}$ & endpoint product, endpoint zeta, $1/\zeta_{Z1}$ & \CGeqref{eq:Z1-product}, \CGeqref{eq:FZ1}\\
$\mathscr R_{Z1}$ & ratio $\zeta_{Z1}/\zeta$ & \CGeqref{eq:Z1-defect-ratio}\\
$\mathcal K$ & multiplier $\zeta_{Z1}/\zeta_{\Z,\RVM}$ & \CGeqref{eq:K-multiplier}\\
(D), $a_n$ & Dirichlet condition for $\zeta_{Z1}$ and its coefficients & \CGeqref{eq:Z1-D}\\
$C_H$ & constant in Hamburger's theorem & \CGeqref{eq:Hamburger-conclusion}\\
$\mu_{Z1}$, $M_{Z1}$ & coefficients of $F_{Z1}$; summatory function & \CGeqref{eq:Z1-reciprocal-Dirichlet}\\
$M_{\lambda,U}$, $M_{Z1,U}$ & finite-height Perron integrals & \CGeqref{eq:truncated-Perron}\\
$\delta$, $\eps$, $\omega(\delta)$ & margins $\sigma\ge\tfrac12+\delta$, exponent $\eps>0$, three-circle exponent & \CGeqref{eq:three-circles-exponent}\\
\midrule
\multicolumn{3}{l}{\emph{Criteria in the half-plane, the rung system and the models}}\\
$\mathsf D_n(\lambda)$, $B^\vee_\lambda$ & $\tilde c_n(\lambda)-\tilde c_n(0)$; dual function of the atomized family & \CGeqref{eq:witness-zeros}, \CGeqref{eq:dual-function}\\
$\mathsf d$, $\mathsf Y$, $\mathsf e_1,\mathsf e_2$ & difference of two counts; auxiliary entire function; exponents & \CGeqref{eq:ladder-h}, \CGeqref{eq:ladder-Y}, \CGeqref{eq:bridge-B2}\\
$J_f$, $J_g$, $b_n$, $\varpi_f$ & counts and coefficients of a general pair $(f,g)$; value of $f$ at $+\infty$ & \CGeqref{eq:halfline-f}, Thm.~\CGref{thm:bridge}\\
\textup{(B1)}, \textup{(B2)} & the two hypotheses of the bridge theorem & Thm.~\CGref{thm:bridge}\\
$\mathsf v_\lambda$ & phase defect of the family on the critical line & \CGeqref{eq:phase-defect}\\
$\mathsf o_n$, $\mathsf L_n$, $\mathfrak j$ & one-ordinate difference of cell masses; $-\log(1-n^{-s})$; one-ordinate joining factor & \CGeqref{eq:one-ordinate}, \CGeqref{eq:one-ordinate-joining}\\
$\mathsf f$, $\mathsf c_{\mathsf f}$, $\mathsf k$, $\mathsf k_{\mathsf f}$ & test function in $u=\log x$; its cosine transform; the part of $G$ and the pole & \CGeqref{eq:transport-defs}\\
$\mathsf R_{\mathsf f}$, $\mathsf S_{\mathsf f}$, $\mathsf m_j(\mathsf f)$ & rung functional; transport; moments sampled on the clock & \CGeqref{eq:rung-functional}, \CGeqref{eq:transport}, \CGeqref{eq:transport-moments}\\
$\mathfrak t_0$, $\mathsf I_m$ & zero of the clock, $\NR(\mathfrak t_0)=0$; Poisson integrals & \CGeqref{eq:clock-defect-poisson}\\
$\mathsf a_n$, $\dot\kappa_k$, $\dot{\mathsf a}_n$ & free masses of the rung system; a direction of variation & \CGeqref{eq:rung-system}, \CGeqref{eq:rung-linearized}\\
$\mathsf J$, $\mathsf x_0$, $\mathsf H_1$, $\mathsf y$, $\mathsf H$, $\mathsf Z$ & objects of the proof of Theorem~\CGref{thm:nondegenerate} & \CGeqref{eq:nondeg-H1}--\CGeqref{eq:nondeg-H}\\
$\langle\cdot,\cdot\rangle_\sigma$, $\mathsf K_{\boldsymbol{\kappa}}$, $\mathsf r_m$, $\mathsf u_m$, $\mathsf Q_{\boldsymbol{\kappa}}$ & pairing of Dirichlet series; $\tilde\zeta_{\Z,\boldsymbol{\kappa}}/\zeta$; gradient data of $\mathsf V_\sigma$ & \CGeqref{eq:K-kappa}, \CGeqref{eq:gradient}\\
$\mathsf z$, $\mathfrak k$, $\mathfrak C_{\boldsymbol{\kappa}}$, $\mathsf N$ & variable, half-type and canonical product of the models; $e^{i\mathfrak k\mathsf z}\mathfrak C_{\boldsymbol{\kappa}}$ & \CGeqref{eq:model-product}\\
$\mathscr F$, $\mathscr F^{+}$, $\mathfrak l$, $\mathfrak w_{\mathfrak l}$, $\mathfrak b_{\mathfrak l}$, $\mathfrak c_{\mathfrak l}$ & frequencies of a discrete category; a frequency; masses, exponents, coefficients & \CGeqref{eq:acceptance-density}--\CGeqref{eq:acceptance-window}\\
$\mathfrak p,\mathfrak q,\mathfrak v$, $\mathsf U$, $\mathsf X$, $\mathfrak D$, $\mathfrak S$ & data, unitary matrix, diagonal factor, determinant and sine function of the two-edge model & \CGeqref{eq:graph-matrices}--\CGeqref{eq:graph-sine}\\
$\mathsf F$, $\mathsf F_{\boldsymbol{\kappa}}$, $\mathsf F_\lambda$, $\mathfrak y_k$, $\mathfrak x_k$ & finite Euler product, its list functions and family; its ordinates and its clock & \CGeqref{eq:finite-euler}--\CGeqref{eq:finite-euler-list}\\
$\mathscr L$, $\mathfrak P$ & $\ell^2$ over the primes; $\sum_pp^{-s}$ & \CGeqref{eq:prime-det}, \CGeqref{eq:prime-det2}\\
$\theta_{\RVM}(s)$, $\tau_0$ & bounded half-plane phase of the clock; fixed transition threshold $\tau_0=4\pi$ (not a clock ordinate) & \CGeqref{eq:halfplane-RVM}, \S\CGref{subsec:transition-proof}\\
$\mathfrak m_n$, $\mathfrak u_\lambda$, $N_\lambda$, $\mathfrak e^{(\lambda)}$, $C_\lambda$ & cell baseline; first-cell term; moving count and count defect; Mellin/explicit-count constant & \CGeqref{eq:cell-mass-W}, \CGeqref{eq:joining-factor}, \CGeqref{eq:family-count-primitive}\\
$\mathsf A_\lambda$, $\mathsf W(\lambda)$, $T_\lambda$, $k_n(\lambda)$, $\mathsf V_\sigma$ & representing function; one-number defect; theta defect; quotient coefficients; vertical defect & \CGeqref{eq:one-number}, \CGeqref{eq:theta-defect}, \CGeqref{eq:vertical-defect}\\
$\Theta_\lambda$, $\mathfrak G_\lambda$, $\mathcal J_T$, $\Delta_T$, $K(T)$, $L_T$ & auxiliary criteria and multiscale clock-Jacobian quantities & \S\CGref{sec:criteria}\\
$E_n$, $\Xi_n$, $\varepsilon_r$ & cell-error and defect-kernel quantities; local regularization parameter where displayed & \S\CGref{sec:criteria}\\
$\boldsymbol{\kappa}$, $N_{\boldsymbol{\kappa}}$, $\tau_0^+$ & generic ordinate list, its count, and lower list threshold & \CGeqref{eq:list-space}\\
$\mathcal G_\varepsilon$, $\mathsf P_{\varepsilon,M}$, $\mathsf R_{\varepsilon,M}$ & additive interpolation and finite completed Dirichlet polynomials & \CGeqref{eq:additive-family}, \CGeqref{eq:finite-additive}\\
$\nu_5,\dots,\nu_{12}$ & real constants, each local to the statement where it occurs & --\\
\midrule
\multicolumn{3}{l}{\emph{Route B for finite sections (Section~\CGref{sec:routeB-finite})}}\\
$\mathbb K(C,\mathbb M)$, $\mathbb M_*$, $\mathsf a^\zeta_n$ & box of $C$-admissible lists with masses $\le\mathbb M$; $\log7/(2\sqrt7)$; prescribed masses $\Lambda(n)/(2\sqrt n)$ & \S\CGref{subsec:rb-routeB}\\
$\mathfrak R(\boldsymbol\kappa,\mathsf a)$, $\mathsf c_\varphi$, $\mathscr T_m$, $\mathsf s_m$ & residual distribution of the rung system; cosine transform of a test; nested test sets; section residual & \CGeqref{eq:rb-residual}, \CGeqref{eq:rb-rho}\\
$\mathcal L$, $\mathscr K$, $\mathscr S_{\boldsymbol\kappa}$, $\mathscr A_\zeta$, $\mathscr E_{\boldsymbol\kappa}$, $\mathfrak H_{\boldsymbol\kappa}$, $\mathfrak g$, $\mathrm c_{\mathscr K}$ & Laplace transform; Archimedean, spectral, arithmetic and total residual distributions; $\sum_k\kappa_k^{-2}\cos(\kappa_ku)$; auxiliary function and constant & \CGeqref{eq:rb-laplace-def}--\CGeqref{eq:rb-S-A-dist}\\
$\mathbb F_m$, $\mathbb J_m$, $\mathbb Q_m$, $\mathbb L_m$, $\mathbb H_m$, $\mathsf h_m$, $\mathsf l_m$, $k^*_m$, $\mathrm t_m$, $\mathbb T_m$ & residual map, Jacobian, right inverse and its norm, Lipschitz constant, start residual and contraction index of stage $m$; numbers of unknowns and tests; contraction map & \CGeqref{eq:rb-Fm}--\CGeqref{eq:rb-Am}\\
$\mathrm q_m$ & graded contraction factors & Prop.~\CGref{prop:rb-graded}\\
$P_{\mathrm X}$, $\zeta_{\mathrm X}$, $\Lambda_{\mathrm X}$, $\mathrm p_{\mathrm X}$, $\mathrm m_{\mathrm X}$, $\mathrm n_{\mathrm X}$, $\varTheta_{\mathrm X}$, $\boldsymbol\kappa^{(\mathrm X)}$ & Gonek's truncated Euler product and symmetrised function with parameter $\mathrm X$; its prime phase and two constants; the Gonek phase; the Gonek list & \CGeqref{eq:rb-alphaX}--\CGeqref{eq:rb-gonek-list}\\
$\widehat{\mathsf R}$, $\mathbb E(T)$, $\upsilon$, $\mathfrak B(C)$, $\mathbb B,\mathbb B_0,\dots,\mathbb B_3$, $\mathrm c_1,\dots,\mathrm c_6$, $\mathrm c_\ell$ & Laplace residual of a section; its $L^1$ size on $\Re s=2$ and growth ratio; constants of the a priori bound & \CGeqref{eq:rb-Rhat}--\CGeqref{eq:rb-a-constants}\\
$\mathcal Y$, $\mathcal D_{\mathsf a}$, $\mathbb W$, $\mathbb S$, $\mathbb J$, $\mathfrak f$, $\mathfrak F_T$, $\mathfrak R_T$, $\mathbb G_T$, $n_T$, $\mathbb D(T)$ & Dirichlet factor and polynomial; $\log|\zeta_{\boldsymbol\kappa}|$; half-strip; compact set; log-factor; Backlund function; rectangle, path, zero count and bound of the counting argument & \S\CGref{subsec:rb-apriori}\\
$\mathscr X$, $\mathrm x_k$, $\varPsi_{\mathscr X}$, $\mathbb P_{\mathrm y}$, $\mathbb Y_{\mathrm y}$, $\mathfrak G_{\mathscr X}$, $\mathcal F_{\mathscr X}$, $\mathcal F_0$, $\mathfrak z=\mathrm x+i\mathrm y$ & unfolded list and its points; discrepancy; Poisson and conjugate operators; partial-fraction series; generating function; $\cos\pi\mathfrak z$; unfolded complex variable & \CGeqref{eq:rb-H}--\CGeqref{eq:rb-F-formula}\\
$\mathrm A$, $\mathrm B$, $\mathrm d_\star$, $\mathbb U$, $\mathcal W$, $\mathrm w$, $\mathrm d_k$, $\mathfrak Z_m$, $\mathfrak J_m$, $\mathfrak n_m$, $\mathrm D_m$, $\mathbb O$, $\mathbb V$, $\mathbb U_{\rm loc}$ & mean and oscillating parts of the discrepancy and their bounds; sector function; weight; displacements, index blocks, intervals, lengths, oscillations and the constants of the variable-length theorem & \S\CGref{subsec:rb-riesz}\\
$\hbar$, $\hbar_T$, $\mathbb P^{\NR}_t$, $\mathbb P^{\circ}_t$, $\mathrm c_{\rm JN}$, $\mathfrak I_T$, $\mathrm c_F$ & line height and its $t$-scale; the two kernels, the John--Nirenberg constant, the interval and Fujii's constant in the obstruction theorem & \S\CGref{subsec:rb-obstruction}\\
$\mathrm w_m$, $T_m$, $u^*_m$, $\mathbb J'$, $\mathbb A^{(3)}$, $\mathsf l(\mathrm r)$ & test width, height and window of stage $m$; second- and third-derivative matrices; Banach index of the computed certificates & \S\CGref{subsec:rb-numerics}\\
$\mathfrak W_r(\mathsf f)$, $H_0$ & constant of a test in the bound for the zeta list; height of the verified zeros, $3\cdot10^{12}$ & \CGeqref{eq:rb-Wr}, Prop.~\CGref{prop:rb-nonfinitary}\\
$\mathscr Z_{\boldsymbol\kappa}$, $\mathfrak T_C$, $\mathfrak E_W$, $\mathscr P_{\boldsymbol\kappa}$, $\mathscr G$ & ordinate transform of an unknown list; box of lists; window residual, its gradient function and the function of its first-order condition & \S\CGref{subsec:rb-window}\\
$\mathrm N_8$, $\varphi^{\mathrm N}_{c,h}$, $h_m$, $\rho_m$ & cardinal B-spline of order $8$, B-spline test, its width at stage $m$, radius of the verified ball & \S\CGref{subsec:rb-num-verified}\\
\midrule
\multicolumn{3}{l}{\emph{Exponential sums over the ordinates (Section~\CGref{sec:ordinate-sums})}}\\
$\rho_k=\frac12+d_k+i\gamma_k$ & the zero with positive ordinate of index $k$; $d_k$ its distance to the critical line & Sec.~\CGref{sec:ordinate-sums}\\
$\widehat{\mathsf f}$ & $\int\mathsf f(u)e^{iuz}du$ & \CGeqref{eq:f-hat}\\
$V_+$ & Cram\'er's function $\sum_{\Im\rho>0}e^{\rho z}$ & \CGeqref{eq:cramer-V}\\
$\mathscr Z_\gamma$, $\mathscr Z_\tau$, $\mathscr Z_\xi$, $\mathfrak X$ & ordinate, clock and zero transforms; their defect $\mathscr Z_\gamma-\mathscr Z_\xi$ & \CGeqref{eq:ordinate-transforms}, \CGeqref{eq:ordinate-defect}\\
$g$; $\mathscr V_{\rm b},\mathscr V_{\rm R},\mathscr V_{\rm L},\mathscr V_{\rm p}$; $\Omega_\theta$ & archimedean part of $\xi'/\xi$; pieces of the zero-free representation; its half-plane & \CGeqref{eq:g-arch}, \CGeqref{eq:cramer-pieces}\\
$\mathscr Y$ & off-line cosine distribution $\sum_k(\cosh(d_ku)-1)\cos(\gamma_ku)$ & \CGeqref{eq:Y-offline}\\
$\mathscr Q_K$, $\mathscr Q_\infty$ & Frullani integrals of the cosine difference & \CGeqref{eq:frullani-QK}\\
$C^{\pathc}_\zeta$, $C^{\pathc}_{\RVM}(\infty;u)$ & triangular centroid of the ordinates; limit of the clock centroid & \CGeqref{eq:triangular-centroid}, \CGeqref{eq:clock-centroid-limit}\\
$b_k$, $d^{\rm b}_k$ & $(e^{iu(\tau_{k+1}-\tau_k)}-1)^{-1}$, $b_{k-1}-b_k$ & Thm.~\CGref{thm:clock-centroid}\\
$f(\lambda)$, $\lambda^*$ & $\log\varpi_\lambda$; crossing point & Prop.~\CGref{prop:crossing}\\
\end{longtable}
}

\medskip\noindent\textbf{The three counts.}  Three counting symbols occur and are never interchanged: $\NR(T)=\mathcal N_{\RVM}(T)$ is the \emph{continuous} midpoint clock \CGeqref{eq:clock}, a real-analytic increasing function; $n_{\RVM}(T)=\#\{k:\tau_k\le T\}=\lfloor\NR(T)+\tfrac12\rfloor$ \CGeqref{eq:n-RVM} is the \emph{integer} count of the clock ordinates; and $N_{\boldsymbol\kappa}(T)=\#\{k:\kappa_k\le T\}$ is the integer count of a general list (with $N_\zeta$ for the ordinates of $\zeta$ and $N_\lambda$ for the list $(\gamma_k(\lambda))$).  Differences such as $N_\zeta-\NR$ compare an integer staircase with the continuous clock; differences such as $N_{\boldsymbol\kappa}-n_{\RVM}$ compare two staircases.  When a mass $\Lambda(n)/(2\sqrt n)$ is called a ``von Mangoldt mass'' it is this scaled mass that is meant.  A distributional identity is stated with its space: $\mathcal D'(0,\infty)$ (the origin excluded), $\mathcal S'(\mathbb R)$, or distributions on $\mathbb R$ supported in $[0,\infty)$ (the origin included); the constants $\mathrm b$ of Appendix~\CGref{subsec:rung-apriori} and the polynomials of Lemma~\CGref{lem:rb-laplace} live exactly in the difference between the first and the third.

\medskip\noindent\textbf{Collision check.}
Global symbols in the table have one meaning throughout the paper.  The interpolation parameter of Appendix~\CGref{subsec:rung-apriori} is written $\mathsf t\in[0,1]$ ($\theta$ is the unit-cell step), and the constants of that appendix are $\mathbb K_{\rm E},\mathbb K_{\rm E}'$ (not $\mathbb K(C,\mathbb M)$, the box of Section~\CGref{sec:routeB-finite}), $B_2$ (not $\mathbb B_2$ or $\mathrm B_2$ of that section) and $C^*$; the letters $K$, $N$, $T$, $L$ and $C$ with different fonts or indices are different symbols, listed separately in the table.  Local proof variables do not collide with global notation: in particular $c_0$ is the de la Vall\'ee Poussin constant, $\varkappa_j$ and $\mathcal P_j$ are local van-der-Corput parameters, $\omega_0$ and $\omega_{\rm test}$ are test weights, $d_\rho$ is a horizontal zero displacement, $T_{\rm ver}$ is the verified-zero height, $\mathcal H^\pm$ are transition half-planes, and $\Psi_m^\pm$ are local stationary phases; $R_0$ is reserved for the clock half-plane threshold and $R_{\rm end}$ for the endpoint clock-decomposition threshold.  The symbol $c$ is reserved for a Perron abscissa, $\eps$ for an exponent/margin, $\varepsilon$ for the additive-family parameter, $\theta$ for the unit-cell step, and $\vartheta$ for the Riemann--Siegel function.  The fixed number $\tau_0=4\pi$ is a transition threshold and is not a member of the sequence $(\tau_k)$.  Dotted variables $\dot\kappa_k,\dot{\mathsf a}_n$ denote a direction of variation.

\section{Introduction}
\label{sec:introduction}

\subsection{Background}
\label{subsec:background}

The theory of the Riemann zeta function joins four structures that are usually introduced in different languages.  On the spectral side there are the nontrivial zeros and their Riemann--von Mangoldt counting law; on the analytic side the completed function $\xi$, its functional equation and its Hadamard product; on the prime side the prime-power measure and the explicit formula; and on the inverse side the reciprocal Dirichlet series $1/\zeta$ and the Mertens function
\begin{equation}\label{eq:mertens-def}
 M(x):=\sum_{n\le x}\mu(n).
\end{equation}
Riemann's memoir made the first three structures inseparable, and the work of von Mangoldt, Hadamard and Hamburger supplied, respectively, the rigorous zero count and explicit formula, the entire-product factorization, and a converse characterization of $\zeta$ from its Dirichlet-series and functional-equation data \cite{Riemann1859,vonMangoldt1895,Hadamard1896,Hamburger1921,Titchmarsh1986,Edwards1974}.  This paper asks what remains of that architecture when the zero staircase is replaced, at the outset, by the exact midpoint inverse of its smooth Riemann--von Mangoldt clock, and how the resulting object is carried back to $\zeta$.

\subsection{Known results on the Riemann--von Mangoldt law, the ordinates and the clock}
\label{subsec:known}

This subsection collects the known theorems that the paper uses or compares with.  They are quoted, not proved here; the theorems of the later sections are new unless a citation is given.  Let $\vartheta$ be the Riemann--Siegel theta function, $\vartheta(T)=\arg\Gamma(\frac14+\frac{iT}2)-\frac T2\log\pi$, and $S(T)=\frac1\pi\arg\zeta(\frac12+iT)$ as in \CGeqref{eq:RvM-full}.

\begin{theorem}[Riemann--von Mangoldt; \cite{vonMangoldt1895}, {\cite[Thms.~9.3--9.4]{Titchmarsh1986}}]\label{thm:known-RvM}
For $T>0$ not an ordinate, $N_\zeta(T)=1+\frac1\pi\vartheta(T)+S(T)$, with $1+\frac1\pi\vartheta(T)=\frac T{2\pi}\log\frac T{2\pi e}+\frac78+O(T^{-1})$ and $S(T)=O(\log T)$.
\end{theorem}

The smooth part is the midpoint clock $\NR$ of \CGeqref{eq:clock}.  Explicit forms of the bound for $S(T)$ were given by Trudgian \cite{Trudgian2014}; Hasanalizade, Shen and Wong \cite{HasanalizadeShenWong2022} proved $\bigl|N_\zeta(T)-\frac T{2\pi}\log\frac T{2\pi e}\bigr|\le0.1038\log T+0.2573\log\log T+9.3675$ for $T\ge e$, and Bellotti and Wong \cite{BellottiWong2025} improved the constants to $0.10076$, $0.24460$, $8.08344$.

\begin{theorem}[Littlewood; Selberg; \cite{Littlewood1924}, {\cite[Thm.~9.9]{Titchmarsh1986}}, \cite{Selberg1946}]\label{thm:known-S}
$\int_0^TS(t)\,dt=O(\log T)$, and $\int_0^TS(t)^2\,dt=\frac T{2\pi^2}\log\log T+O\bigl(T(\log\log T)^{1/2}\bigr)$.
\end{theorem}

\emph{Gram points.}  The Gram points $g_n$ are defined by $\vartheta(g_n)=n\pi$; Gram's law is the statement $(-1)^nZ(g_n)>0$ for Hardy's function $Z$.  The name and the first failure are due to Hutchinson \cite{Hutchinson1925}; the law fails infinitely often (Titchmarsh) and, by Trudgian \cite{Trudgian2011}, a positive proportion of the time; Brent's computations \cite{Brent1979} used Gram blocks and Rosser's rule.  By Theorem~\CGref{thm:known-RvM}, $\NR(g_{k-1})=k+O(g_{k-1}^{-1})$, while $\NR(\tau_k)=k-\frac12$: the midpoint ordinates lie half-way between consecutive Gram points in clock units.

\emph{Computations.}  The Riemann hypothesis has been verified for all zeros up to height $3\cdot10^{12}$ (Platt and Trudgian \cite{PlattTrudgian2021}); Platt \cite{Platt2017} isolated the zeros up to height about $3.06\cdot10^{10}$ in interval arithmetic.

\begin{theorem}[Landau's formula; \cite{Landau1912}, uniform version \cite{Gonek1993}]\label{thm:known-landau}
For fixed $x>1$, $\sum_{0<\gamma\le T}x^\rho=-\frac T{2\pi}\Lambda(x)+O_x(\log T)$, where $\Lambda(x)=0$ unless $x$ is a prime power.
\end{theorem}

\begin{theorem}[Hlawka; \cite{Hlawka1975}, see also \cite{FordZaharescu2005}]\label{thm:known-hlawka}
For every real $\alpha\ne0$ the sequence $(\alpha\gamma_k)_{k\ge1}$ is uniformly distributed modulo $1$; in particular $\frac1N\sum_{k\le N}e^{iu\gamma_k}\to0$ for every $u\ne0$.
\end{theorem}

\begin{theorem}[Cram\'er; \cite{Cramer1919}]\label{thm:known-cramer}
The function $V_+(z)=\sum_{\Im\rho>0}e^{\rho z}$, holomorphic in $\Im z>0$, has a meromorphic continuation to the plane cut along the non-positive imaginary axis; its poles are simple and lie at the points $\pm\log p^m$.
\end{theorem}

The explicit formulas of Riemann and von Mangoldt, in the distributional form of Guinand and Weil \cite{Riemann1859,vonMangoldt1895,Guinand1948,Weil1952}, are used throughout.  For the Mertens function, $M(x)=O_\eps(x^{1/2+\eps})$ is equivalent to RH \cite[\S14.25]{Titchmarsh1986}; the Mertens conjecture $|M(x)|<\sqrt x$ is false \cite{OdlyzkoTeRiele1985}; $M(x)$ has been computed to $10^{14}$ by Kotnik and van de Lune \cite{KotnikvandeLune2004} and to $10^{16}$ by Hurst \cite{Hurst2018}, who proved $\limsup M(x)/\sqrt x\ge1.826054$ and $\liminf M(x)/\sqrt x\le-1.837625$; under RH and further hypotheses, Ng \cite{Ng2004} determined the limiting distribution of $e^{-y/2}M(e^y)$.

\emph{The clock in the literature.}  The midpoint clock and its Lambert inversion are due to Fran\c{c}a and LeClair: equation (24) of \cite{FrancaLeClair2014}, $\frac{\tilde y_n}{2\pi}\log\frac{\tilde y_n}{2\pi e}=n-\frac{11}8$, is the midpoint condition \CGeqref{eq:midpoint-condition}, and its solution (25), $\tilde y_n=2\pi(n-\frac{11}8)/W_0((n-\frac{11}8)/e)$, is \CGeqref{eq:lambert}; thus $\tau_k=\tilde y_k$.  Fran\c{c}a and LeClair obtain $\tilde y_n$ as the leading-order solution of their transcendental equation for the $n$th zero, in which the term $\arg\zeta$ is dropped, use it as an approximation to $\gamma_n$, study the statistics of the zeros with it, and reconstruct $\pi(x)$ numerically from the explicit formula with these approximations in place of the zeros \cite{FrancaLeClair2014,FrancaLeClair2015}.  In the literature consulted we have not found a zeta function built on the clock ordinates, nor the prime count, the prime levels and the Mertens function constructed from it below, nor the family joining it to $\zeta$.

\emph{Beurling systems with a Riemann hypothesis.}  Zhang \cite{Zhang2013} constructed Beurling generalized primes whose zeta function has no zeros in $\Re s>\frac12$ and whose integer counting function is $x+O(x^{1/2}e^{c(\log x)^{2/3}})$; Broucke, Debruyne and Vindas \cite{BrouckeDebruyneVindas2020} proved the existence, by a probabilistic construction, of a Beurling system with $\pi_B(x)=\operatorname{Li}(x)+O(\sqrt x)$ and large oscillation of the integers; see also \cite{DiamondMontgomeryVorhauer2006,HilberdinkLapidus2006,BrouckeDebruyneRevesz2025}.  These are genuine Beurling systems, with non-decreasing counting functions, obtained by existence arguments.  The synthetic count $\pi_{\RVM}$ of this paper is deterministic and is built from the clock, with $\pi_{\RVM}(x)-\li(x)=O(\sqrt x(\log\log x)^2/\log x)$ without hypothesis (Theorem~\CGref{thm:Lambert-estimate}); it is not the counting function of a Beurling system, since its density $\frac1{\log y}+2\cos(2\pi y)+\cdots$ (Theorem~\CGref{thm:clock-density}) changes sign.  The comparison is therefore one of size of the error term, not of the structure.

\subsection{The midpoint RVM system}
\label{subsec:midpoint-system}

The point of departure is the smooth counting function \CGeqref{eq:clock},
\begin{equation}\label{eq:intro-clock}
 \NR(T)=\frac{T}{2\pi}\log\frac{T}{2\pi e}+\frac78,
\end{equation}
and the real ordinates $\tau_k$ determined by the midpoint condition $\NR(\tau_k)=k-\tfrac12$.  The half-integer convention is the midpoint convention for a staircase and leads to the exact Lambert inversion \CGeqref{eq:lambert},
\begin{equation}\label{eq:intro-lambert}
 \tau_k=\frac{2\pi(k-11/8)}{W_0\bigl((k-11/8)/e\bigr)} .
\end{equation}
The clock, the midpoint convention (the constant $\frac{11}8$) and the Lambert inversion are those of Fran\c{c}a and LeClair \cite[Sec.~III.A, Eqs.~(24)--(25)]{FrancaLeClair2014}, who obtained \CGeqref{eq:intro-lambert} as the leading-order solution of their equation for the $n$th zero and used it as an approximation to $\gamma_n$ (Subsection~\CGref{subsec:known}); for the function $W$ see \cite{Corless1996,DLMF2024}.  What is new here is what is built on the clock: the midpoint inverse is treated as the exact spectrum of a deterministic system and is propagated, without the fluctuating term $S(T)$, through canonical products, explicit prime counting and Mertens inversion.

The first object is the canonical product and the synthetic zeta function
\begin{equation}\label{eq:intro-zeta-rvm}
 P_{\RVM}^{\crit}(s)=\prod_{k\ge1}\left(1+\frac{s(s-1)}{\tfrac14+\tau_k^2}\right),
 \qquad
 \zeta_{\RVM}(s)=\frac{P_{\RVM}^{\crit}(s)}{G(s)},
 \qquad G(s)=s(s-1)\Gamma(s/2)\pi^{-s/2}.
\end{equation}
Because $\tau_k\asymp k/\log k$, the factors form a normally convergent product in the symmetric variable $s(s-1)$; $P_{\RVM}^{\crit}$ is entire and invariant under $s\mapsto1-s$, and $\zeta_{\RVM}$ has a residue-one pole at $1$, the trivial zeros, and the simple nontrivial zeros $\rho_k^{\RVM}=\tfrac12\pm i\tau_k$ (Theorem~\CGref{thm:canonical}).  The real part $\tfrac12$ is part of the definition.

\emph{Generalized prime systems.}  In Beurling's theory the generalized primes are arbitrary reals $1<p_1\le p_2\le\cdots\to\infty$; their multiplicative semigroup defines generalized integers and a zeta function by a Dirichlet series, an Euler product or a Mellin--Stieltjes transform, separating prime distribution from the additive structure of the integers \cite{Beurling1937,BatemanDiamond1969,Diamond1970}.  Beurling systems with prescribed analytic behaviour, critical zeros and large oscillation show that zeta-like zero data alone do not determine the arithmetic realization \cite{DiamondMontgomeryVorhauer2006,BrouckeDebruyneVindas2020}.  The RVM system runs in the opposite direction: it starts from a prescribed zero clock, constructs an entire completed product first, obtains a continuous prime-power count from its explicit formula, and only afterwards asks for an atomic realization on the integers.

\emph{Continuous and discrete prime systems.}  The passage from a continuous prime system to a discrete one is classical in the theory of Beurling generalized primes \cite{Beurling1937,DiamondMontgomeryVorhauer2006,BrouckeVindas2024,DiamondZhang2016}: the atoms are placed at suitably chosen real points, and the discretization changes the zeta function by a factor holomorphic and zero-free in $\Re s>\frac12$ (compare Theorem~\CGref{thm:atomized-continuation}(iii)).  Measures with a continuous part and non-integer multiplicities also appear in the theory of complex dimensions \cite[Ch.~3]{LapidusvanFrankenhuysen2000}.  Grosswald and Schnitzer \cite{GrosswaldSchnitzer1978} replaced each prime $p_n$ by a real $q_n$ with $p_n\le q_n\le p_{n+1}$ and proved that the modified Euler product is $\zeta$ times a factor holomorphic and zero-free in $\Re s>0$, so that it has the same zeros as $\zeta$ there; Corollary~\CGref{cor:weighted-factorization} and Theorem~\CGref{thm:atomized-continuation}(iv) are of the same type for $\zeta_{\RVM}$.  Here the masses of $d\pi_\lambda$ are placed at the integers, with real exponents $\tilde c_n(\lambda)$; this yields an ordinary Dirichlet series for every $\lambda$, absolutely convergent in $\Re s>1$ (Corollary~\CGref{cor:abscissa-one}), and makes Hamburger's theorem \cite{Hamburger1921,Siegel1922} available at the endpoint.  Hamburger's theorem and its extensions \cite{BochnerChandrasekharan1956,KaczorowskiPerelli1999} need absolute convergence in $\Re s>1$, and Knopp \cite{Knopp1994} showed that a half-plane of absolute convergence further to the right does not suffice.  We are not aware of an earlier use of this combination.

\emph{Prime coordinate.}  From the RVM zeros we define the prime-power count $J_{\RVM}$ by \CGeqref{eq:J-RVM}, with the branch \CGeqref{eq:Ei-def} of $\Ei$.  Finite M\"obius inversion in the exponent variable, which relates Riemann's prime-power count to $\pi(x)$ \cite{Riemann1859,Titchmarsh1986,Edwards1974}, is used here as the \emph{definition} of the synthetic prime count \CGeqref{eq:pi-RVM},
\begin{equation}\label{eq:intro-pi}
 \pi_{\RVM}(x)=\sum_{1\le m\le\log_2x}\frac{\mu(m)}m\,J_{\RVM}^{*}(x^{1/m}),
 \qquad
 J_{\RVM}^{*}(x)=\sum_{1\le m\le\log_2x}\frac1m\,\pi_{\RVM}(x^{1/m}).
\end{equation}
The function $\pi_{\RVM}$ is continuous away from the points $2^j$ with $\mu(j)\ne0$, where it has the jump $\mu(j)J_{\RVM}(2)/j$; in particular it is not a staircase on the integers.  It supports two notions of RVM prime: the continuous measure $d\pi_{\RVM}$, and, on every interval where $\pi_{\RVM}$ is continuous and strictly increasing, the real level crossings
\begin{equation}\label{eq:intro-real-primes}
 p_j^{\RVM}:=\pi_{\RVM}^{-1}(j),\qquad \pi_{\RVM}(p_j^{\RVM})=j,
\end{equation}
which need not be integers.  In the computations of Section~\CGref{sec:numerics} the levels are the first crossings $\inf\{x:\pi_{\RVM}(x)\ge j\}$, which agree with \CGeqref{eq:intro-real-primes} when the crossing lies on such an interval, evaluated on a grid.  The synthetic count has square-root accuracy without any hypothesis on $\zeta$: $J_{\RVM}(x)-\li(x)$ and $\pi_{\RVM}(x)-\li(x)$ are $O(\sqrt x(\log\log x)^2/\log x)$ (Theorem~\CGref{thm:Lambert-estimate}).  The proof reduces the zero sum to $\sum_ke^{iu\tau_k}/\tau_k$, which is $O(\log^2u)$ by the second-derivative test on dyadic blocks of the clock (Lemma~\CGref{lem:A-RVM}).  Hence $\li(x)/\pi_{\RVM}(x)\to1$, and with the prime number theorem $\pi(x)/\li(x)\to1$ and $\pi(x)/\pi_{\RVM}(x)\to1$ (Corollary~\CGref{cor:three-limits}, Figure~\CGref{fig:pi-ratios}).  Sampling on unit cells gives the masses $c_n^{\RVM}=\pi_{\RVM}(n)-\pi_{\RVM}(n-1)$ and the weighted Euler object $\zeta_{\Z,\RVM}$ of \CGeqref{eq:weighted-zeta}.  The masses are not $0$ or $1$ and need not be positive, but they are bounded: $J_{\RVM}$ is Lipschitz on $[2,\infty)$ (Theorem~\CGref{thm:unit-cell}), so $c_n^{\RVM}=O(1)$ and $\zeta_{\Z,\RVM}$ is an absolutely convergent Dirichlet series in $\Re s>1$ (Corollaries~\CGref{cor:cn-bounded} and~\CGref{cor:ZRVM-Dirichlet}).  The unit-periodic first harmonic of the clock integrates to zero on every integer cell (Lemma~\CGref{lem:cell-cancel}); under the cell hypothesis one obtains $c_n^{\RVM}\sim1/\log n$ (Proposition~\CGref{prop:eventual-positive}).

\emph{Mertens coordinate.}  The Mertens coordinate is the symmetric inverse Mellin object of $F_{\RVM}=1/\zeta_{\RVM}$.  A smooth/sawtooth Stieltjes decomposition of the clock product with Abel regularization of the sawtooth gives on every band to the right of $\tfrac12$
\begin{equation}\label{eq:intro-transition}
 \zeta_{\RVM}(s)=\varpi_{\RVM}\bigl(1+w(s)\bigr)\bigl(1+O(1/|t|)\bigr),
 \qquad w(s)=\exp\bigl\{-2\pi i\,\NR\bigl(-i(s-\tfrac12)\bigr)\bigr\},
\end{equation}
and $\zeta_{\RVM}=\varpi_{\RVM}(1+\chi)(1+O(1/|t|))$ on closed substrips of $\tfrac12<\Re s<1$ (Theorem~\CGref{thm:transition}).  Hence $F_{\RVM}(\sigma+it)\to L=1/\varpi_{\RVM}$ for every fixed $\sigma>\tfrac12$; neither $\zeta_{\RVM}$ nor $F_{\RVM}$ is an absolutely convergent Dirichlet series in any half-plane (Theorem~\CGref{thm:no-Dirichlet-RVM}); and after subtracting the vertical limit the Perron inverse converges absolutely, with
\begin{equation}\label{eq:intro-Mbound}
 M_{\RVM}^{\crit}(x)=L+O_\eps(x^{1/2+\eps})
\end{equation}
(Theorem~\CGref{thm:unsmoothed-RVM}).  The RVM Mertens function therefore exists as a continuous spectral inverse although no coefficient sequence for $F_{\RVM}$ exists.  Since the primitive of the midpoint sawtooth is bounded, the same decomposition without contours gives $C^{-1}\le|\zeta_{\RVM}|\le C$ on every half-plane $\Re s\ge\sigma_0>1$ and $\zeta_{\RVM}(s)\to\varpi_{\RVM}$ as $\Re s\to+\infty$ (Theorem~\CGref{thm:halfplane-RVM}).

\emph{Ternary discretization.}  An integer-valued representative of $M_{\RVM}^{\crit}$ is obtained by sampling the target $\mathcal M(n):=M_{\RVM}^{\crit}(n)$ and choosing signs on the squarefree support:
\begin{equation}\label{eq:intro-muh}
 \mu_h(n)=\mu(n)^2\varsigma_n,\qquad\varsigma_n\in\{-1,+1\},\qquad M_h(n):=\sum_{m\le n}\mu_h(m),
\end{equation}
\begin{equation}\label{eq:intro-quantizer}
 \varsigma_n=\begin{cases}+1,&M_h(n-1)<\mathcal M(n),\\-1,&M_h(n-1)\ge\mathcal M(n),\end{cases}\qquad(\mu(n)^2=1).
\end{equation}
This is a greedy sigma--delta feedback quantization; feedback quantization and noise shaping are standard themes of sigma--delta theory \cite{DaubechiesDeVore2003,GunturkThao2003,BenedettoPowellYilmaz2006}.  Here the target is a zero-generated Mertens inverse and the admissible times are the squarefree integers.  The walk $\mu_h$ is a comparison object: it is not claimed to be multiplicative or equal to $\mu$.  The two discretizations are distinct: the cell masses $c_n^{\RVM}$ come from $\pi_{\RVM}$, the signs $\varsigma_n$ track $M_{\RVM}^{\crit}$.

\emph{Comparison with $\zeta$.}  At the level of transforms the defect is the exact quotient $\mathcal Q_{\RVM}=\zeta/\zeta_{\RVM}$ of \CGeqref{eq:Q-RVM}, with $F_{\RVM}=\mathcal Q_{\RVM}\cdot(1/\zeta)$.  Since $|\zeta_{\RVM}|$ is bounded above and below on $\frac12+\eps\le\sigma\le c'$, $|t|\ge t_0$ (Corollary~\CGref{cor:Q-control}), on each bounded band $\frac12+\eps\le\sigma\le c'$, every bound for $\zeta$ or $1/\zeta$ is equivalent to the corresponding bound for $\mathcal Q_{\RVM}$ by Corollary~\CGref{cor:Q-control}; for $\sigma>c'$ one uses Theorem~\CGref{thm:halfplane-RVM}.  For $M$ itself the bound $M(x)=O_\eps(x^{1/2+\eps})$ is equivalent to the Riemann hypothesis \cite[\S14.25]{Titchmarsh1986}, and the disproof of the Mertens conjecture shows that the $\eps$ cannot simply be dropped \cite{OdlyzkoTeRiele1985}.

\subsection{The core of the method}
\label{subsec:core}

The transfer from $\zeta_{\RVM}$ to $\zeta$ is a transformation of $\zeta_{\RVM}$ depending on a parameter $\lambda\in[0,1]$.  The ordinates move along
\begin{equation}\label{eq:intro-family}
 \gamma_k(\lambda)=(1-\lambda)\gamma_k+\lambda\tau_k,
 \qquad
 \zeta_\lambda(s)=\frac1{G(s)}\prod_{k\ge1}\left(1+\frac{s(s-1)}{\tfrac14+\gamma_k(\lambda)^2}\right),
\end{equation}
where $\gamma_k$ are the ordinates of the nontrivial zeros of $\zeta$; thus $\zeta_1=\zeta_{\RVM}$ and $\zeta_0=\zeta_{Z1}$ (Section~\CGref{sec:lambda-family}).
\begin{enumerate}[label=\textup{(\roman*)}]
\item Every member $\zeta_\lambda$ has its nontrivial zeros exactly at $\tfrac12\pm i\gamma_k(\lambda)$.  The real part $\tfrac12$ is part of the definition of the family, like the gamma factor, and never changes along it: it is impossible for $\zeta_{Z1}$, or for any $\zeta_\lambda$, to have a zero off the critical line other than the trivial zeros.  For $0<\lambda\le1$ all zeros are simple.
\item All differences between the moving zetas lie in the numbers $\gamma_k-\tau_k$, which are bounded (Lemma~\CGref{lem:bounded-displacement}); the quotient of two members is an explicit product in $\gamma_k$ and $\lambda(\tau_k-\gamma_k)$ (Remark~\CGref{rem:only-gamma-tau}).
\item The final step is the identity $\zeta_{Z1}=\zeta$.  It is the single and complete test: while $\zeta_{Z1}=\zeta$ there is nothing outside $\Re s=\tfrac12$ to speak of, since the identity itself places every nontrivial zero of $\zeta$ at $\tfrac12\pm i\gamma_k$, and $\zeta_{Z1}\ne\zeta$ is the only possible obstruction.
\end{enumerate}
The endpoint has the functional equation, the pole, the residue and the order of $\zeta$ (Proposition~\CGref{prop:Z1-class}).  Hamburger's theorem belongs to the family of converse theorems that recover Dirichlet series from analytic continuation, functional equation and growth \cite{Hamburger1921,Perelli2016}; applied to $\zeta_{Z1}$ it gives (Theorem~\CGref{thm:Hamburger-Z1})
\begin{equation}\label{eq:intro-hamburger}
 \zeta_{Z1}=\zeta
 \quad\Longleftrightarrow\quad
 \text{(D): }\zeta_{Z1}\text{ is an absolutely convergent Dirichlet series in }\Re s>1 .
\end{equation}
With the integer-weighted zeta $\zeta_{\Z,\RVM}$, which is a Dirichlet series, (D) is equivalent to an absolutely convergent Dirichlet expansion of the explicit multiplier $\mathcal K=\zeta_{Z1}/\zeta_{\Z,\RVM}=\Pi_{\rm tr}e^{-h_{\RVM}}$ (Theorem~\CGref{thm:D-multiplier}); since $1/\zeta_{\Z,\RVM}$ is also an absolutely convergent Dirichlet series in $\Re s>1$, the equivalence holds on the whole half-plane.  Section~\CGref{sec:euler-products} gives two sharper forms of (D).  First, the clock decomposition applied to the ordinates $\gamma_k$ gives the value of the endpoint at infinity,
\begin{equation}\label{eq:intro-varpi-Z1}
 \varpi_{Z1}=\lim_{\sigma\to+\infty}\zeta_{Z1}(\sigma)=\varpi_{\RVM}\lim_{K\to\infty}\prod_{k\le K}\frac{\frac14+\tau_k^2}{\frac14+\gamma_k^2}\ \ge1,
\end{equation}
and $\zeta_{Z1}=\zeta$ holds if and only if $\varpi_{Z1}=1$ (Theorems~\CGref{thm:Z1-clock} and~\CGref{thm:varpi-Z1}, Corollary~\CGref{cor:product-identity}): the final step is the single ordered-limit identity $\lim_{K\to\infty}\prod_{k\le K}(\frac14+\gamma_k^2)/(\frac14+\tau_k^2)=\varpi_{\RVM}$ between the displacements and the transition constant.  Second, the prime-power counts $J_{\RVM}^{\rm M}$ and $J_{Z1}^{\rm M}$ are the Mellin transforms of $\log\zeta_{\RVM}$ and $\log\zeta_{Z1}$ and satisfy a distributional explicit formula (Theorems~\CGref{thm:mellin-counts} and~\CGref{thm:distributional-EF}); this proves the Mellin--Stieltjes identification of the integer weights (Theorem~\CGref{thm:MS}), and (D) holds if and only if $J_{Z1}^{\rm M}$ is constant between consecutive integers and its jumps satisfy the growth hypothesis of Theorem~\CGref{thm:cell-form}.  Third, along the family $\frac d{d\lambda}\log\varpi_\lambda=-\Upsilon(\lambda)$, the generator at infinity, and (D) is the identity $\int_0^1\Upsilon(\lambda)\,d\lambda=\log L$ (Proposition~\CGref{prop:family-flow}).  Section~\CGref{sec:atomized} atomizes every member: $\tilde\zeta_{\Z,\lambda}=\zeta_\lambda e^{\tilde h_\lambda}$ keeps the Euler product and the Dirichlet series while $\zeta_\lambda$ keeps the functional equation; the joining factor $\tilde h_\lambda$ is written cell by cell and is holomorphic in $\Re s>\frac12$ (Theorems~\CGref{thm:joining-factor} and~\CGref{thm:atomized-continuation}); the cell masses have mean square $X^{1+\eps}$ on $(X,2X]$ uniformly in $\lambda$, so every $\tilde\zeta_{\Z,\lambda}$, the endpoint included, is an absolutely convergent Dirichlet series in $\Re s>1$ (Theorem~\CGref{thm:mean-square-cells}, Corollary~\CGref{cor:abscissa-one}); and (D) is equivalent to the vanishing of every cell defect at $\lambda=0$ (Theorem~\CGref{thm:closing}).  At the clock endpoint the joining factor $\tilde h_1$ is not constant, and the endpoint Mertens function has a Mellin form $B_{Z1}$ with $B_{Z1}=M$ exactly when (D) holds (Propositions~\CGref{prop:atomized-clock} and~\CGref{prop:Mellin-Mertens}).  The remaining steps form the programme of Subsection~\CGref{subsec:programme}.  In RVM language, the final step asks whether the zero-generated analytic object $\zeta_{Z1}$ belongs to the discrete Dirichlet category; if it does, Hamburger uniqueness makes it $\zeta$, and the prime and Mertens coordinates become those of the integers.  Step~3 is Theorem~\CGref{thm:step3}:
\begin{equation}\label{eq:intro-step3}
 \zeta_{Z1}=\zeta
 \quad\Longrightarrow\quad
 M(x)=O_\eps\bigl(x^{1/2+\eps}\bigr)\qquad(\eps>0),
\end{equation}
with a self-contained proof by the Borel--Carath\'eodory lemma, Hadamard's three-circle theorem and the truncated Perron formula.  Theorem~\CGref{thm:transfer} rewrites the same endpoint consequence through the synthetic factorization.  Its estimate for $M$ is a corollary of Theorem~\CGref{thm:step3}; the synthetic Mertens bound is not an independent input to the classical implication.  The identity $\zeta_{Z1}=\zeta$ itself, equivalently (D), is not proved in this paper; it is the open step of the method (Remark~\CGref{rem:Step3-status}).

The relations among the principal objects are
\begin{equation}\label{eq:intro-prime-chain}
 \{\tau_k\}\longrightarrow\zeta_{\RVM}\longrightarrow J_{\RVM}\longleftrightarrow\pi_{\RVM}\longrightarrow\bigl(d\pi_{\RVM},\,p_j^{\RVM},\,c_n^{\RVM}\bigr),
\end{equation}
\begin{equation}\label{eq:intro-M-chain}
 F_{\RVM}=\zeta_{\RVM}^{-1}\longrightarrow M_{\RVM}^{\crit}\longrightarrow(\mathcal M,\mu_h,M_h),
\end{equation}
\begin{equation}\label{eq:intro-family-chain}
 \zeta_{\RVM}=\zeta_1\xrightarrow{\ \lambda\downarrow0\ }\zeta_0=\zeta_{Z1}\overset{?}{=}\zeta\ \Longrightarrow\ M(x)=O_\eps(x^{1/2+\eps}).
\end{equation}
The first chain is the prime coordinate, the second the Mertens coordinate, the third the $\lambda$-family; they share the midpoint spectrum and solve different realization problems.

\subsection{Exponential sums over the ordinates and the Riemann hypothesis}
\label{subsec:intro-expsums}

The exponential sums of the two spectra behave differently.  For the clock, $\sum_{k\le N}e^{iu\tau_k}\ll1+\log N$ for every $u>0$ (Theorem~\CGref{thm:exp-sum}), and the triangular centroid
\begin{equation}\label{eq:intro-centroid}
 C^{\pathc}(N;u)=\frac1N\sum_{n\le N}\sum_{k\le n}e^{iuo_k}=\sum_{k\le N}\Bigl(1-\frac{k-1}N\Bigr)e^{iuo_k}
\end{equation}
of the list $o_k=\tau_k$ converges, without hypothesis, to an explicit series (Theorem~\CGref{thm:clock-centroid}); at $u=1$ the limit is $0.041569+1.757541\,i$.  For the ordinates $o_k=\gamma_k$ the sums resonate at the frequencies $\log p^m$ (Landau), and the natural comparison is with Cram\'er's function $V_+(z)=\sum_{\Im\rho>0}e^{\rho z}$, which sees the complete zeros.  Section~\CGref{sec:ordinate-sums} proves the following statements, all but the last without any hypothesis.
\begin{enumerate}[label=\textup{(\alph*)}]
\item \emph{Zero-free representation.}  $V_+$ is given on the half-plane $\Omega_\theta$, which contains $(0,\infty)$, by an explicit formula in terms of $\Lambda$, $\Gamma$ and $\zeta$ on a real segment, with no sum over the zeros (Proposition~\CGref{prop:cramer-zero-free}).
\item \emph{One imaginary frequency.}  For any $y_0\in(0,4\pi]$,
\begin{equation}\label{eq:intro-imaginary}
 \mathrm{RH}\iff\sum_{k\ge1}e^{-y_0\gamma_k}=e^{-iy_0/2}V_+(iy_0)
\end{equation}
(Proposition~\CGref{prop:imaginary-frequency}); the defect is $\sum_k(1-\cos(y_0d_k))e^{-y_0\gamma_k}\ge0$, $d_k=\Re\rho_k-\frac12$.
\item \emph{The window form.}  With $\widehat{\mathsf f}(z)=\int\mathsf f(u)e^{iuz}du$ and $J=(\log2,\log3)$, where the prime side of the explicit formula has no atom,
\begin{equation}\label{eq:intro-window}
 \mathrm{RH}\iff\sum_{k\ge1}\widehat{\mathsf f}(\gamma_k)=\int_J\mathsf f(u)\,e^{-u/2}V_+(u)\,du\quad\text{for every }\mathsf f\in C_c^\infty(J),
\end{equation}
equivalently the triangular centroids $C^{\pathc}_\zeta(N;u)$ converge to $e^{-u/2}V_+(u)$ in $\mathscr D'(J)$ (Theorem~\CGref{thm:window-RH}).  Without hypothesis the difference of the two sides is $\sum_k\bigl(\widehat{\mathsf f}(\gamma_k)-\frac12\widehat{\mathsf f}(\gamma_k-id_k)-\frac12\widehat{\mathsf f}(\gamma_k+id_k)\bigr)$, an absolutely convergent sum over the zeros off the line.  After inserting the representation of (a), both sides are expressed without the abscissas of the zeros.  The proof uses a one-variable edge-of-the-wedge lemma and positivity at an imaginary frequency.
\item \emph{The endpoint defect is a sum over the zeros off the line.}  For every test function, $\mathsf R_{\mathsf f}(0)-\frac12\sum_n\Lambda(n)n^{-1/2}\mathsf f(\log n)=\sum_k\int\mathsf f(u)(\cosh(d_ku)-1)\cos(\gamma_ku)\,du$; equivalently $(\log x)(dJ^{\rm M}_{Z1}-dJ)/dx=\frac2{\sqrt x}\sum_k(\cosh(d_k\log x)-1)\cos(\gamma_k\log x)$; and the cosine difference $\sum_k(\mathsf c_{\mathsf f}(\gamma_k)-\mathsf c_{\mathsf f}(\tau_k))$ equals the defect of the clock minus that sum (Theorem~\CGref{thm:rung-zero-sum}).  If only finitely many zeros were off the line, $J^{\rm M}_{Z1}$ would be constant on no cell $(n,n+1)$ (Proposition~\CGref{prop:finite-offline}).
\item \emph{One window is one cell.}  The real part of \CGeqref{eq:intro-window} is the constancy of $J^{\rm M}_{Z1}$ on the single cell $(2,3)$, and (D), equivalently RH, is equivalent to that cell together with the sine identity $\sum_k\int\mathsf f(u)\sin(\gamma_ku)du=\int\mathsf f(u)e^{-u/2}\Im V_+(u)du$ on $J$ (Corollary~\CGref{cor:one-window}).  The whole ladder of cells is thereby compressed into one spectral window.
\item \emph{Frullani form.}  $2\int_0^\infty\frac{e^{-u/2}}u\sum_{k\le K}(\cos\gamma_ku-\cos\tau_ku)\,du\to\log L+\log\varpi_{Z1}\ge\log L$, with equality exactly under RH (Theorem~\CGref{thm:frullani}).
\item \emph{The centroid of $e^{i\gamma_k}$ under RH.}  If RH holds, $C^{\pathc}_\zeta(N;u)\to e^{-u/2}V_+(u)$ for every $u>0$ that is not the logarithm of a prime power; at $u=1$ the limit is
\begin{equation}\label{eq:intro-centroid-value}
 e^{-1/2}V_+(1)=\frac1{2\sqrt e}\Bigl(e-\frac1{e^2-1}\Bigr)+i\,0.76785566886\ldots=0.77689426079\ldots+0.76785566886\ldots\,i
\end{equation}
(Theorem~\CGref{thm:centroid-RH}).  The proof uses Landau's formula, Selberg's mean value of $S$, the zero-free representation and a local Tauberian lemma; the converse implication is not proved.
\end{enumerate}
The numerical side of these statements (walks of $e^{i\gamma_k}$, $e^{i\tau_k}$ and $e^{i(\gamma_k+\tau_k)/2}$ in three dimensions, the centroids with $2\,001\,052$ zeros, and $V_+(1)$ from the zero-free formula) is in Subsection~\CGref{subsec:num-centroids}.  Subsection~\CGref{subsec:num-first-1000} and Appendix~\CGref{app:tables-1000} give the first $1000$ values of $\tau_k$ against $\gamma_k$, of the RVM prime levels $p_j^{\RVM}$ against $p_j$, of $\pi_{\RVM}(n)$ against $\pi(n)$ and of $M_{\RVM}^{\crit}(n)$ against $M(n)$, with graphs.

\subsection{Originality, scope and status}
\label{subsec:status}

\emph{What this paper does not prove.}  Before the list of results: this paper does not prove the Riemann hypothesis; it does not prove (D), the identity $\zeta_{Z1}=\zeta$, or any of the statements shown below to be equivalent to it; it does not prove the existence of a solution of the rung system \CGeqref{eq:rung-system} in $\mathscr O$, nor of an interior exact-atomic path, nor the convergence of the finite sections of Section~\CGref{sec:routeB-finite} to a distributional limit; and, apart from the interval evaluation of a few explicit elementary expressions in the proof of Theorem~\CGref{thm:varpi-less-one} (which the margins allow to be checked by hand with five-decimal values of $\log2$, $\log\pi$ and $\log\tau_1$), it contains no computer-assisted proof, the other computations being floating-point experiments.  The only points left open by the method are these: the endpoint identity (D) in any of its equivalent forms; the existence of a solution of the infinite rung system; the injectivity of the clock-only linearization (Theorem~\CGref{thm:nondegenerate} proves injectivity at a solution, not at the clock); and the passage to the limit from finite sections.  Every ``equivalent form'' below is a reduction to (D), not a resolution of it.

The ingredients are standard or belong to neighbouring theories: the Riemann--von Mangoldt law, Lambert inversion, Hadamard products, explicit formulas, M\"obius and Perron inversion, Beurling generalized primes, sigma--delta feedback, and Hamburger's converse theorem.  The contribution is their assembly into one midpoint RVM space with four coordinated levels: (i) an exact midpoint zero clock and its canonical zeta-like product; (ii) a continuous prime-power and prime coordinate, with real prime levels and non-integral cell masses; (iii) a continuous synthetic Mertens inverse with its ternary feedback discretization; (iv) the $\lambda$-family from the midpoint spectrum to the ordinates of $\zeta$, ending in the single identity $\zeta_{Z1}=\zeta$ and the Mertens bound that follows from it.  To the author's knowledge the cited literature contains the individual mechanisms but not this chain.  Standard material is cited at the point of use; statements without citation are proved here.

\begingroup\small
\begin{longtable}{p{0.50\textwidth}p{0.40\textwidth}}
\toprule
Statement & Status\\
\midrule
\endfirsthead
\toprule
Statement & Status\\
\midrule
\endhead
\bottomrule
\endfoot
Midpoint clock and Lambert inversion (due to Fran\c{c}a--LeClair \cite{FrancaLeClair2014}); bounded displacement (from $S(T)=O(\log T)$) & theorem\\
Canonical product $P_{\RVM}^{\crit}$ and $\zeta_{\RVM}$ & theorem\\
$S_{\RVM}(N;u)\ll1+\log N$; centroid and $M_H^{\RVM}$ decay & theorem\\
Clock path centroid converges for every $u>0$, with the limit as an explicit series & theorem\\
Zero-free representation of Cram\'er's function $V_+$ on $\Omega_\theta$ & theorem (continuation due to Cram\'er)\\
RH $\Leftrightarrow\sum_ke^{-y_0\gamma_k}=e^{-iy_0/2}V_+(iy_0)$ for one $y_0\in(0,4\pi]$ & theorem\\
RH $\Leftrightarrow$ window identity on $(\log2,\log3)$; exact defect as a sum over zeros off the line & theorem\\
$\mathsf R_{\mathsf f}(0)=\langle\mathscr Y,\mathsf f\rangle$; $(\log x)(dJ^{\rm M}_{Z1}-dJ)/dx=2x^{-1/2}\mathscr Y(\log x)$; cosine difference & theorem, unconditional\\
(D) $\Leftrightarrow$ the cell $(2,3)$ together with the sine identity on the window & theorem\\
Finitely many zeros off the line $\Rightarrow$ every cell fails & theorem\\
Frullani form: RH $\Leftrightarrow\mathscr Q_\infty=\log L$ & theorem\\
Local Tauberian lemma & theorem (self-contained case of known results)\\
Triangular centroid of $e^{iu\gamma_k}$ tends to $e^{-u/2}V_+(u)$ & theorem under RH; converse open\\
Compact-uniform existence of $J_{\RVM}$; finite inversion & theorem\\
Real prime levels $p_j^{\RVM}$ & defined on monotone branches of $\pi_{\RVM}$\\
Square-root estimate for $J_{\RVM},\pi_{\RVM}$ & theorem\\
$\li/\pi_{\RVM}\to1$; $\pi/\pi_{\RVM}\to1$ & theorem (with the prime number theorem)\\
$J_{\RVM}$ Lipschitz, $c_n^{\RVM}=O(1)$ & theorem\\
$\zeta_{\Z,\RVM}$ absolutely convergent Dirichlet series in $\Re s>1$ & theorem\\
Mellin--Stieltjes identification; $\zeta_{\Z,\RVM}=\zeta_{\RVM}e^{h_{\RVM}}$ & theorem\\
$|\zeta|\asymp|\mathcal Q_{\RVM}|$ for $\sigma\ge\frac12+\eps$, $|t|\ge t_0$ & theorem\\
Transition theorem $\varpi_{\RVM}(1+w)$, $\varpi_{\RVM}(1+\chi)$ & theorem\\
$\varpi_{\RVM}<0.9796$; closed form $\nu_3=8-\frac{33}4\log2-2\log\pi$ & theorem (explicit inequalities checked in interval arithmetic, program \texttt{s14})\\
$C^{-1}\le|\zeta_{\RVM}|\le C$ on $\Re s\ge\sigma_0>1$; $\zeta_{\RVM}\to\varpi_{\RVM}$ at $+\infty$ & theorem\\
Unsmoothed $M_{\RVM}^{\crit}=L+O_\eps(x^{1/2+\eps})$ & theorem\\
Dirichlet expansion of $\zeta_{\RVM}^{\pm1}$ & impossible (theorem)\\
Ternary walk $\mu_h$, $M_h$ & definition; comparison object\\
Exact Hadamard product of $2\xi$ & theorem\\
$\lambda$-family, generator, multiplier, simplicity for $\lambda>0$ & theorem\\
Analytic class of $\zeta_{Z1}$ & theorem\\
$\zeta_{Z1}=\zeta\Leftrightarrow$ (D) $\Leftrightarrow$ RH & theorem\\
(D) $\Leftrightarrow$ Dirichlet expansion of $\mathcal K=\zeta_{Z1}/\zeta_{\Z,\RVM}$ in $\Re s>1$ & theorem\\
$\varpi_{Z1}=\varpi_{\RVM}\lim_{K\to\infty}\prod_{k\le K}\frac{1/4+\tau_k^2}{1/4+\gamma_k^2}$ (ordered product); (D) $\Leftrightarrow\varpi_{Z1}=1$ & theorem\\
Mellin prime-power counts; distributional explicit formula & theorem\\
(D) $\Leftrightarrow$ $J_{Z1}^{\rm M}$ constant between integers with the jump-growth bound of Theorem~\CGref{thm:cell-form} & theorem\\
$\frac d{d\lambda}\log\varpi_\lambda=-\Upsilon$; (D) $\Leftrightarrow\int_0^1\Upsilon=\log L$ & theorem\\
Atomized family; joining factor holomorphic in $\Re s>\frac12$ & theorem\\
Atomized family: Dirichlet series in $\Re s>1$ for every $\lambda$ & theorem\\
(D) $\Leftrightarrow$ vanishing cell defects at $\lambda=0$ & theorem\\
$\zeta_\lambda$ not a Dirichlet series in $\Re s>1$ for $0<\lambda\le1$ & theorem\\
$\lambda\mapsto\log\zeta_\lambda(\sigma)$ strictly convex for real $\sigma>1$ & theorem\\
Hamburger's theorem for the pair $(\zeta_\lambda,\tilde\zeta_{\Z,\lambda})$ with symmetric joining factor & theorem\\
(D) $\Leftrightarrow$ $\zeta_{Z1}$ has a positive representing measure; then $\varpi_{Z1}\le1$ & theorem\\
Density of $d\pi_{\RVM}$; cell masses and defects at $\lambda=1$ & theorem\\
Height kernel of the cell masses; Gaussian defect of the atomization & theorem\\
$\zeta_{Z1}(\sigma)\ge\zeta(\sigma)$ for real $\sigma>1$; (D) $\Leftrightarrow$ equality at one $\sigma$ & theorem\\
Hamburger's theorem for the pair in theta form and in vertical mean-square form & theorem\\
Symmetry of the joining factor fails at $\lambda=1$; $\tilde\zeta_{\Z,\lambda}=\zeta$ for finitely many $\lambda$ at most & theorem\\
Summation formula for every $\zeta_\lambda$; $\varpi_\lambda-1$ as one number $\mathsf W(\lambda)$ & theorem\\
(D) $\Leftrightarrow$ $\varpi_\lambda<1$ on the open path $0<\lambda\le1$ & theorem (uses $\varpi_{\RVM}<1$, Theorem~\CGref{thm:varpi-less-one})\\
Crossing point: unique, $f(0)/\Upsilon(0)\le\lambda^*\le f(0)/(f(0)+|f(1)|)$, $\lambda^*<2\cdot10^{-11}$ if RH fails & theorem (the last bound uses the verified zeros)\\
A family with all the listed properties and $\varpi_0>1$ & theorem\\
Space $\mathscr O$ of ordinate lists: rigidity of the coincidence; (D) $\Leftrightarrow$ existence of a coincidence list & theorem\\
(D) $\Leftrightarrow$ $\tilde c_n(0)=\mathbf 1_{\rm prime}(n)$ for every $n$; every point of $\mathsf E$ is a witness of (D) & theorem\\
Symmetric joining factors: order one is needed; the symmetry criterion has no slack at $\lambda=0$ & theorem\\
Ladder: $\tilde h_\lambda(\sigma)=O(\sigma N^{-\sigma})$ $\Leftrightarrow$ no mass of $dJ^{\rm M}_\lambda$ outside the integers below $N$ & theorem\\
Bridge theorem in phase form; (D) $\Leftrightarrow$ \textup{(B1)}, \textup{(B2)} for $\tilde h_0$; impossible for $\lambda>0$ & theorem\\
Phase of $\zeta_\lambda$ on the critical line; cell masses as sums of a one-ordinate function & theorem\\
Transport $\mathsf S_{\mathsf f}(\lambda)$: entire of exponential type; moment series on the clock; defect of the clock by Poisson summation & theorem\\
Multiscale clock Jacobian: lower frame bound on $U\asymp\log T$, no fixed-window bound, dense range, smooth tails & theorem; quantitative nonlinear continuation open\\
(D) $\Leftrightarrow$ solvability of the rung system with free masses; the solution is nondegenerate & theorem\\
Gradient of $\mathsf V_\sigma$ on $\mathscr O$ & theorem\\
Acceptance criterion (window and frequencies); two-edge model; the finite Euler product over $2,3$ with the final step proved & theorem\\
No object of finite rank in the window of $\zeta$; the three abscissas of the diagonal operator over the primes & theorem\\
Phase defect along the path (rotation of the one-ordinate terms) & heuristic\\
Endpoint-only interpolation audit; clock-only injectivity problem; a priori rung displacement bound & Appendix~\CGref{sec:appendix-endpoint}: theorem/proposition/open problem as labelled\\
Route B: compactness in a box; uniqueness in every tube with free masses (an exact solution is the zeta list with the von Mangoldt masses, and RH holds); B $\Leftrightarrow$ (D) & theorem (Section~\CGref{sec:routeB-finite})\\
Finite stages do not decide route B: the zeta list solves every finite set of tests up to $(3\cdot10^{12})^{1-r}$ & theorem (uses the zeros verified up to $3\cdot10^{12}$)\\
Route B on one window: RH $\Leftrightarrow$ an unknown list reproduces Cram\'er's function on $(\log2,\log3)$ $\Leftrightarrow$ $e^{-iy/2}V_+(iy)$ completely monotone; minimizers on windows: exact or degenerate; RH $\Leftrightarrow$ solvability on every window & theorem; the exclusion of degenerate minimizers is open\\
Laplace transform of the residual with prescribed masses; uniqueness; RH $\Leftrightarrow$ B$'(C)$ & theorem; classical in substance\\
Certificate theorem for staged sections; graded extension; Gonek lists in a tube & theorem; the certificate hypotheses (A$_m$), (B) and (G1), (G2) are open and imply RH\\
Counting bound $|N_{\boldsymbol\kappa}-n_{\RVM}|\le(\mathrm c_1+\mathrm c_2\upsilon)\log(T+8)+\mathrm c_3\mathbb B+\mathrm c_4$ for quasi-solutions & theorem; void for finite sections with masses, since then $\upsilon=\infty$ (Lemma~\CGref{lem:rb-upsilon-infinite})\\
Block Kadec--Avdonin criterion; variable-length Avdonin theorem & theorem; the block criterion is close to the sufficiency half of the Hru\v{s}\v{c}ev--Nikol'skii--Pavlov theorem; the variable-length version was not found in the literature consulted\\
No $(A_2)$ for the unfolded zeta ordinates on any horizontal line & theorem, unconditional\\
Candidate Banach certificates of exact square sections up to $T=487$; basin $\sim T^{-2}$ & computation (floating point, Gaussian tests; not verified)\\
Exact solutions of fifteen square sections up to $T=487$ with B-spline tests & theorem, computer-assisted (ball arithmetic)\\
First $1000$ values of $\tau_k$, $p_j^{\RVM}$, $\pi_{\RVM}(n)$, $M_{\RVM}^{\crit}(n)$; walks and centroids & computation\\
$\zeta_{Z1}=\zeta\Rightarrow M(x)=O_\eps(x^{1/2+\eps})$ & theorem (Step~3)\\
$\zeta_{Z1}=\zeta$ & open\\
\end{longtable}
\endgroup

\emph{What is proved and what is not.}  Four groups of statements are proved in full.  \emph{(1) The synthetic system.}  The clock and its Lambert inversion, the canonical product and $\zeta_{\RVM}$, the logarithmic bound for the exponential sums, the prime coordinate $J_{\RVM}$, $\pi_{\RVM}$ with its square-root estimate and its density (Theorem~\CGref{thm:clock-density}), the transition theorem, the impossibility of a Dirichlet expansion of $\zeta_{\RVM}^{\pm1}$ and the synthetic Mertens bound.  \emph{(2) The family and its endpoint.}  Every $\zeta_\lambda$ has the functional equation, order one and all its nontrivial zeros on $\Re s=\frac12$; $\zeta_{Z1}=\zeta$ is equivalent to (D), to $\varpi_{Z1}=1$ and to the constancy of $J^{\rm M}_{Z1}$ between consecutive integers together with the jump-growth bound of Theorem~\CGref{thm:cell-form}; $\varpi_{Z1}\ge1$ and $\zeta_{Z1}(\sigma)\ge\zeta(\sigma)$ for real $\sigma>1$; $\log\zeta_\lambda(\sigma)$ is strictly convex in $\lambda$; and $\zeta_{Z1}=\zeta$ implies $M(x)=O_\eps(x^{1/2+\eps})$ (Step~3).  \emph{(3) The atomized family and the criteria.}  Every $\tilde\zeta_{\Z,\lambda}$ is an absolutely convergent Dirichlet series in $\Re s>1$, zero-free in $\Re s>\frac12$, joined to $\zeta_\lambda$ by an explicit factor; no $\zeta_\lambda$ with $\lambda>0$ is a Dirichlet series; Hamburger's theorem holds for the pair in three forms (Theorems~\CGref{thm:pair-criterion}, \CGref{thm:double-hamburger}); the positivity criterion (Theorem~\CGref{thm:positivity}); and the rigidity of the coincidence in the space of ordinate lists (Theorem~\CGref{thm:rigidity}).  \emph{(4) Criteria written where the joining factor is defined, and models.}  The reduction of (D) to the cell masses (Corollary~\CGref{cor:cell-mass-form}); the ladder and the half-line form of Hamburger's theorem for a pair (Theorems~\CGref{thm:ladder}, \CGref{thm:halfline-hamburger}); the bridge theorem in phase form (Theorem~\CGref{thm:bridge}, Corollary~\CGref{cor:bridge-family}); the transport from the clock and the exact defect of the clock (Theorem~\CGref{thm:transport-band}, Proposition~\CGref{prop:clock-defect-poisson}); the existence form with free masses and the nondegeneracy of the solution (Theorems~\CGref{thm:free-masses}, \CGref{thm:nondegenerate}); the acceptance criterion and the models in which the final step is a theorem (Theorems~\CGref{thm:acceptance}, \CGref{thm:finite-euler} for the primes $2,3$, Proposition~\CGref{prop:graph-model}).

\emph{Not proved:} the identity $\zeta_{Z1}=\zeta$.  Each of the following statements is proved to be equivalent to it, and none of them is proved: (D); $\varpi_{Z1}\le1$; $\zeta_{Z1}(\sigma)\le\zeta(\sigma)$ for one $\sigma>1$; $\varpi_\lambda<1$ for all $\lambda\in(0,1]$ (this equivalence uses Theorem~\CGref{thm:varpi-less-one}); $J^{\rm M}_{Z1}$ non-decreasing; $\mathsf W(0)\ge0$ \CGeqref{eq:one-number}; $\tilde h_0$ constant; the theta relation $T_0\equiv0$ \CGeqref{eq:theta-defect}; $\mathsf V_\sigma(0)=0$ \CGeqref{eq:vertical-defect}; the existence of a list in $\mathscr O$ whose continuous and atomized zeta functions coincide (Corollary~\CGref{cor:existence-form}); $\tilde c_n(0)=\mathbf 1_{\rm prime}(n)$ for every $n\ge2$ (Corollary~\CGref{cor:cell-mass-form}); $\tilde h_0(\sigma)=O_N(\sigma N^{-\sigma})$ for every $N$ \CGeqref{eq:ladder-D}; the conditions \textup{(B1)} and \textup{(B2)} for $\tilde h_0$ (Corollary~\CGref{cor:bridge-family}); $\mathsf S_{\mathsf f}(0)=\mathsf R_{\mathsf f}(1)$ for every $\mathsf f$ whose support avoids the points $\log n$, together with the cell jump-growth condition (Proposition~\CGref{prop:rung-functional}); the solvability of the rung system \CGeqref{eq:rung-system} in $\mathscr O$ (Theorem~\CGref{thm:free-masses}); the identity \CGeqref{eq:imaginary-frequency} at one imaginary frequency; the window identity of Theorem~\CGref{thm:window-RH}(ii); the constancy of $J^{\rm M}_{Z1}$ on $(2,3)$ together with the sine identity \CGeqref{eq:sine-identity}; $\mathscr Q_\infty=\log L$ (Theorem~\CGref{thm:frullani}).  The symmetry of the joining factor $e^{\tilde h_0}$ with finite order is a sufficient condition (Theorem~\CGref{thm:pair-criterion}), and it is also necessary (Remark~\CGref{rem:pair-no-slack}).

\subsection{Organization}
\label{subsec:organization}

Section~\CGref{sec:clock} treats the clock, its Lambert inversion and the bounded displacement; Section~\CGref{sec:zeta-rvm} the canonical product; Section~\CGref{sec:exp-sums} the exponential sums of the clock; Section~\CGref{sec:ordinate-sums} the exponential sums over the ordinates of $\zeta$, Cram\'er's function, the window form of RH, the endpoint defect as a sum over the zeros off the line, and the centroids; Sections~\CGref{sec:explicit-formula} and~\CGref{sec:integer-weights} the prime coordinate and the integer weights; Sections~\CGref{sec:mertens-transform} and~\CGref{sec:transition} the Mertens transform and the transition theorem; Section~\CGref{sec:comparison} the exact Hadamard product and the comparison with $\zeta$; Section~\CGref{sec:lambda-family} the $\lambda$-family; Section~\CGref{sec:Z1-step3} the endpoint, Hamburger's theorem and Step~3; Section~\CGref{sec:euler-products} the value of the endpoint at infinity and the Euler products of $\zeta_{\RVM}$ and $\zeta_{Z1}$; Section~\CGref{sec:atomized} the atomized family, Hamburger's hypotheses at the two endpoints and the programme for the final step; Section~\CGref{sec:criteria} the open path, convexity, Hamburger's theorem for the pair, the positivity criterion, the density of the clock, the Gaussian defect, the reduction to the cell masses, the half-line form, the bridge theorem and the transport from the clock; Section~\CGref{sec:lists} the space of ordinate lists, the rigidity of the coincidence, the existence form of the final step, the rung system with free masses and the gradient of the vertical defect; Section~\CGref{sec:models} the acceptance criterion and the models in which the final step is a theorem; Section~\CGref{sec:routeB-finite} route B for finite sections, the prescribed-mass uniqueness, the staged certificates, the Gonek lists, the a priori counting bound, the block Kadec--Avdonin criterion and its failure for the zeta list, with the certificate computations; Section~\CGref{sec:numerics} the numerical illustration; Section~\CGref{sec:conclusions} the conclusions and open problems; Appendix~\CGref{sec:appendix-endpoint} records the endpoint-only interpolation audit, the autonomous clock-rigidity question, the a priori rung bound, and the documented bounded finite-section diagnostics; Appendix~\CGref{app:tables-1000} gives the tables of the first $1000$ values.

\section{The midpoint Riemann--von Mangoldt clock}
\label{sec:clock}

Let $\rho$ run over the nontrivial zeros of $\zeta$, counted with multiplicity.  No nontrivial zero is real, so the positive ordinates can be listed with multiplicity as $0<\gamma_1\le\gamma_2\le\cdots$.  Put
\begin{equation}\label{eq:zero-count}
 N_\zeta(T):=\#\{\rho:\ 0<\Im\rho\le T\}=\#\{k:\gamma_k\le T\}.
\end{equation}
The Riemann--von Mangoldt formula \cite[Thms.~9.3--9.4]{Titchmarsh1986}, \cite{vonMangoldt1895} states that for $T\ge2$, $T\ne\gamma_k$,
\begin{equation}\label{eq:RvM-full}
 N_\zeta(T)=\frac{T}{2\pi}\log\frac{T}{2\pi e}+\frac78+S(T)+O\Bigl(\frac1T\Bigr),
 \qquad
 S(T):=\frac1\pi\arg\zeta\bigl(\tfrac12+iT\bigr)=O(\log T),
\end{equation}
where $\arg\zeta(\frac12+iT)$ is obtained by continuous variation of the argument along the segments from $2$ to $2+iT$ and from $2+iT$ to $\frac12+iT$, starting from $\arg\zeta(2)=0$.  The $O(1/T)$ term comes from Stirling's formula for $\Gamma$ and is explicit; $S(T)$ is the fluctuating part, with a jump $+m$ at an ordinate of multiplicity $m$.  The \emph{midpoint RVM clock} is the smooth part of \CGeqref{eq:RvM-full}:
\begin{equation}\label{eq:clock}
 \NR(T):=\frac{T}{2\pi}\log\frac{T}{2\pi e}+\frac78 .
\end{equation}
It is strictly increasing on $[2\pi,\infty)$.  The midpoint ordinates $\tau_k$ are the unique solutions in $[2\pi,\infty)$ of
\begin{equation}\label{eq:midpoint-condition}
 \NR(\tau_k)=k-\frac12,
 \qquad k\ge1,
\end{equation}
that is, the abscissae on the increasing branch where the smooth count crosses the midpoints $k-\frac12$ of the steps of a unit staircase.  For $k=1$ the same level has a second positive solution in $(0,2\pi)$ on the $W_{-1}$ branch; it is not part of the midpoint clock used here.  Thus $\NR$ replaces the staircase $N_\zeta$ by its smooth part, and $\tau_k$ is the point where the $k$th step would be crossed half-way (Figure~\CGref{fig:clock-zoom}).

The Lambert $W$-function is the multivalued inverse of $v\mapsto ve^v$: $W(z)e^{W(z)}=z$.  On $[-1/e,\infty)$ the equation $ve^v=z$ has exactly one solution $v\ge-1$; this defines the principal branch $W_0:[-1/e,\infty)\to[-1,\infty)$, which is real, increasing, with $W_0(0)=0$ and $W_0(z)\sim\log z-\log\log z$ as $z\to\infty$ \cite[\S4.13]{DLMF2024}, \cite{Corless1996}.

\begin{proposition}[Lambert inversion]\label{prop:lambert}
For every $k\ge1$,
\begin{equation}\label{eq:lambert}
 \tau_k=\frac{2\pi(k-11/8)}{W_0\bigl((k-11/8)/e\bigr)},
\end{equation}
The quotient is positive also for $k=1$.
\end{proposition}

\begin{proof}
Put $y=\tau/(2\pi)$.  Equation \CGeqref{eq:midpoint-condition} reads
\begin{equation}\label{eq:lambert-proof-y}
 y(\log y-1)=k-\frac{11}{8}.
\end{equation}
The function $y\mapsto y(\log y-1)$ decreases on $(0,1)$, increases on $(1,\infty)$ and has minimum $-1$ at $y=1$; the clock is increasing exactly where $y>1$, and $\tau_k$ is the solution with $y>1$.  With $v=\log y-1$ one has $y=e^{v+1}$ and $ve^v=(k-11/8)/e$; the solution with $y>1$ has $v>-1$, so $v=W_0((k-11/8)/e)$ and $y=(k-11/8)/W_0((k-11/8)/e)$.  For $k=1$, $(k-11/8)/e=-3/(8e)\in(-1/e,0)$, where $W_0$ is real with $W_0(-3/(8e))\in(-1,0)$; numerator and denominator are both negative, and $y=e^{1+W_0(-3/(8e))}\in(1,e)$, so $2\pi<\tau_1<2\pi e$.  (The other solution of \CGeqref{eq:lambert-proof-y} for $k=1$, with $y<1$, is given by the branch $W_{-1}$ and is not a clock point.)
\end{proof}

Equation \CGeqref{eq:lambert-proof-y} and its solution \CGeqref{eq:lambert} are equations (24) and (25) of Fran\c{c}a and LeClair \cite[Sec.~III.A]{FrancaLeClair2014}, who derived them as the leading-order approximation to the ordinate of the $n$th zero; here they define the clock exactly, and nothing is approximated.

The clock derivative and the resulting growth are
\begin{equation}\label{eq:clock-derivative}
 \NR'(T)=\frac1{2\pi}\log\frac{T}{2\pi},
\end{equation}
\begin{equation}\label{eq:clock-growth}
 \tau_k\asymp\frac{k}{\log k}\qquad(k\ge2).
\end{equation}
In terms of the clock, \CGeqref{eq:RvM-full} reads
\begin{equation}\label{eq:RvM-formula}
 N_\zeta(T)=\NR(T)+S(T)+O(T^{-1}),\qquad S(T)=O(\log T)\qquad(T\ge2).
\end{equation}
Let
\begin{equation}\label{eq:symmetric-count}
 N_\zeta^{*}(T):=\tfrac12\bigl(N_\zeta(T+0)+N_\zeta(T-0)\bigr)
\end{equation}
be the symmetric count, and define the displacements
\begin{equation}\label{eq:eta-def}
 \eta_k:=\tau_k-\gamma_k .
\end{equation}

\begin{lemma}[Bounded displacement]\label{lem:bounded-displacement}
There is a constant $C_\eta$ such that
\begin{equation}\label{eq:bounded-displacement}
 |\eta_k|\le C_\eta\qquad(k\ge1).
\end{equation}
\end{lemma}

\begin{proof}
Let $\gamma=\gamma_k$ and let $m=m(\gamma)$ be the number of indices $j$ with $\gamma_j=\gamma$; these indices form a block $j_0+1,\dots,j_0+m$ containing $k$.  Then $N_\zeta(\gamma-0)=j_0$ and $N_\zeta(\gamma+0)=j_0+m$, so $N_\zeta^{*}(\gamma)=j_0+m/2$ and
\begin{equation}\label{eq:block-count}
 \bigl|N_\zeta^{*}(\gamma_k)-(k-\tfrac12)\bigr|\le\frac{m(\gamma_k)}2 .
\end{equation}
By \CGeqref{eq:RvM-formula}, applied at the one-sided limits $N_\zeta(\gamma+1+0)$ and $N_\zeta(\gamma-1-0)$ if $\gamma\pm1$ is itself an ordinate, $m(\gamma)\le N_\zeta(\gamma+1+0)-N_\zeta(\gamma-1-0)=O(\log\gamma)$, and likewise $N_\zeta^{*}(\gamma)=\NR(\gamma)+O(\log\gamma)$.  With \CGeqref{eq:midpoint-condition},
\begin{equation}\label{eq:clock-discrepancy}
 \bigl|\NR(\gamma_k)-\NR(\tau_k)\bigr|=O(\log\gamma_k).
\end{equation}
Explicitly, $\NR(\tau_k)=k-\frac12$ and $\NR(\gamma_k)=N_\zeta^*(\gamma_k)+O(\log\gamma_k)=k+O(\log\gamma_k)$; since $\NR(T)=\frac T{2\pi}\log\frac T{2\pi e}+O(1)$ is increasing with $\NR(T)\asymp T\log T$, both relations give $\gamma_k\asymp\tau_k\asymp k/\log k$ and hence $\log\gamma_k\asymp\log\tau_k\asymp\log k$.  Hence there is $k_1$ such that for $k\ge k_1$ both $\gamma_k$ and $\tau_k$ exceed $2\pi e$, where $\NR$ is increasing with $\NR'\ge\frac1{2\pi}$.  For $k\ge k_1$, by the mean-value theorem, $\NR(\gamma_k)-\NR(\tau_k)=\NR'(\xi_k)(\gamma_k-\tau_k)$ with $\xi_k$ between $\gamma_k$ and $\tau_k$, and by \CGeqref{eq:clock-derivative} $\NR'(\xi_k)=\frac1{2\pi}\log\frac{\xi_k}{2\pi}\gg\log k$ because $\xi_k\ge\min(\gamma_k,\tau_k)\gg k/\log k$; with \CGeqref{eq:clock-discrepancy} this gives $|\gamma_k-\tau_k|=O(1)$.  The finitely many indices $k<k_1$ are absorbed into $C_\eta$.
\end{proof}

\section{The pure critical RVM zeta function}
\label{sec:zeta-rvm}

Put
\begin{equation}\label{eq:q-def}
 q(s):=s(s-1),
\end{equation}
\begin{equation}\label{eq:Pcrit}
 P_{\RVM}^{\crit}(s):=\prod_{k\ge1}\left(1+\frac{q(s)}{\tfrac14+\tau_k^2}\right),
\end{equation}
\begin{equation}\label{eq:Gdef}
 G(s):=s(s-1)\Gamma(s/2)\pi^{-s/2},
\end{equation}
and
\begin{equation}\label{eq:zetaRVM}
 \xi_{\RVM}(s):=\frac12P_{\RVM}^{\crit}(s),
 \qquad
 \zeta_{\RVM}(s):=\frac{P_{\RVM}^{\crit}(s)}{G(s)}.
\end{equation}

\begin{theorem}[Canonical characterization of $\zeta_{\RVM}$]\label{thm:canonical}
The product \CGeqref{eq:Pcrit} converges absolutely and locally uniformly in $s\in\C$; as a function of the symmetric variable $q=q(s)$ it is a canonical product of genus zero.  It defines an entire function with
\begin{equation}\label{eq:P-symmetry}
 P_{\RVM}^{\crit}(1-s)=P_{\RVM}^{\crit}(s),
 \qquad
 P_{\RVM}^{\crit}(\bar s)=\overline{P_{\RVM}^{\crit}(s)},
 \qquad
 P_{\RVM}^{\crit}(0)=P_{\RVM}^{\crit}(1)=1 .
\end{equation}
Its zeros are exactly the simple zeros
\begin{equation}\label{eq:synthetic-divisor}
 \rho_k^{\RVM}=\frac12+i\tau_k,\qquad \overline{\rho_k^{\RVM}}=\frac12-i\tau_k,
 \qquad k\ge1 .
\end{equation}
Consequently $\zeta_{\RVM}$ has a simple pole at $s=1$ with residue $1$, simple trivial zeros at $s=-2,-4,-6,\dots$, and nontrivial zeros exactly at \CGeqref{eq:synthetic-divisor}; moreover $\xi_{\RVM}(s)=\xi_{\RVM}(1-s)$.
\end{theorem}

\begin{proof}
By \CGeqref{eq:clock-growth}, $\sum_k\tau_k^{-2}<\infty$.  On a compact set $K$,
\begin{equation}\label{eq:canonical-normal}
 \sup_{s\in K}\left|\frac{q(s)}{\tfrac14+\tau_k^2}\right|\ll_K\tau_k^{-2},
\end{equation}
so the Weierstrass criterion gives normal convergence.  Each factor depends on $s$ only through $q(s)$ and equals $1$ at $s=0,1$, which gives \CGeqref{eq:P-symmetry}; real coefficients give the conjugation symmetry.  The $k$th factor vanishes exactly when
\begin{equation}\label{eq:factor-zero}
 \Bigl(s-\frac12\Bigr)^2+\tau_k^2=0 .
\end{equation}
The clock is strictly increasing on $[2\pi,\infty)$ and $\tau_1>2\pi$, so the $\tau_k$ are distinct; the derivative $2s-1=\pm2i\tau_k\ne0$ of the vanishing factor and the non-vanishing of the other factors make the zeros simple.  At $s=1$, $P_{\RVM}^{\crit}(1)=1$ and $G(s)=(s-1)(1+O(s-1))$, so the residue is $1$.  At $s=-2n$ ($n\ge1$) the simple pole of $\Gamma(s/2)$ gives a simple zero of $1/G$, while $P_{\RVM}^{\crit}(-2n)>0$ because $q(-2n)>0$, so there is no cancellation.  At $s=0$, $s\Gamma(s/2)=2\Gamma(1+s/2)\to2$, so $G(0)=-2\ne0$ and $s=0$ is not a zero.
\end{proof}

\section{RVM exponential sums, centroids, and the height transform}
\label{sec:exp-sums}

For $u\in\R$ and $N\ge1$ define
\begin{equation}\label{eq:S-RVM}
 S_{\RVM}(N;u):=\sum_{k\le N}e^{iu\tau_k},
 \qquad
 C_{\RVM}^{\pts}(N;u):=\frac{S_{\RVM}(N;u)}{N},
\end{equation}
and
\begin{equation}\label{eq:path-centroid}
 C_{\RVM}^{\pathc}(N;u):=\frac1N\sum_{n\le N}S_{\RVM}(n;u)
 =\sum_{k\le N}\frac{N-k+1}{N}e^{iu\tau_k}.
\end{equation}

\begin{theorem}[Uniform logarithmic exponential-sum bound]\label{thm:exp-sum}
Let $0<u_1\le u_2<\infty$.  Then
\begin{equation}\label{eq:exp-sum-bound}
 \sup_{u\in[u_1,u_2]}|S_{\RVM}(N;u)|\ll_{u_1,u_2}1+\log N .
\end{equation}
\end{theorem}

\begin{proof}
Let $\mathcal T:[-\tfrac18,\infty)\to[2\pi,\infty)$ be the inverse of the clock, so $\tau_k=\mathcal T(k-\tfrac12)$, and put $\ell_k=\log(\tau_k/2\pi)>0$.  Differentiating $\NR(\mathcal T(y))=y$,
\begin{equation}\label{eq:inverse-clock}
 \mathcal T'(y)=\frac{2\pi}{\log(\mathcal T(y)/2\pi)},
 \qquad
 \mathcal T''(y)=-\frac{4\pi^2}{\mathcal T(y)\log^3(\mathcal T(y)/2\pi)} .
\end{equation}
Since $\mathcal T'$ decreases and $|\mathcal T''|$ decreases,
\begin{equation}\label{eq:rvm-gap-bounds}
 \frac{2\pi}{\ell_{k+1}}\le\tau_{k+1}-\tau_k\le\frac{2\pi}{\ell_k},
\end{equation}
\begin{equation}\label{eq:rvm-second-difference}
 |\tau_{k+2}-2\tau_{k+1}+\tau_k|\le\frac{4\pi^2}{\tau_k\ell_k^3}.
\end{equation}
From \CGeqref{eq:midpoint-condition}, $\tau_k\ell_k=2\pi(k-\tfrac{11}8)+\tau_k$, hence
\begin{equation}\label{eq:tau-ell-k}
 \tau_k\ell_k\asymp k\qquad(k\ge2).
\end{equation}
Fix $u\in[u_1,u_2]$ and put $\Delta_k=u(\tau_{k+1}-\tau_k)$, $E_k=e^{iu\tau_k}$.  Choose $k_0=k_0(u_2)$ with $\ell_{k_0}\ge2u_2$; then $0<\Delta_k\le\pi$ for $k\ge k_0$.  With $\mathfrak b_k^{\rm sbp}:=(e^{i\Delta_k}-1)^{-1}$ one has $E_{k+1}=e^{i\Delta_k}E_k$, hence
\begin{equation}\label{eq:Ek-identity}
 E_k=\mathfrak b_k^{\rm sbp}(E_{k+1}-E_k).
\end{equation}
Summing \CGeqref{eq:Ek-identity} over $k_0\le k\le N$ and shifting the index in the first part gives the exact identity
\begin{equation}\label{eq:summation-by-parts}
 \sum_{k=k_0}^{N}E_k=\mathfrak b_N^{\rm sbp}E_{N+1}-\mathfrak b_{k_0}^{\rm sbp}E_{k_0}+\sum_{k=k_0+1}^{N}(\mathfrak b_{k-1}^{\rm sbp}-\mathfrak b_k^{\rm sbp})E_k .
\end{equation}
For $0<x\le\pi$, $|e^{ix}-1|=2\sin(x/2)\ge2x/\pi$.  Using \CGeqref{eq:rvm-gap-bounds} and \CGeqref{eq:rvm-second-difference},
\begin{equation}\label{eq:kappa-bounds}
 |\mathfrak b_k^{\rm sbp}|\le\frac{\pi}{2\Delta_k}\le\frac{\ell_{k+1}}{4u_1},
 \quad
 |\mathfrak b_{k-1}^{\rm sbp}-\mathfrak b_k^{\rm sbp}|\le|\Delta_k-\Delta_{k-1}|\,|\mathfrak b_{k-1}^{\rm sbp}||\mathfrak b_k^{\rm sbp}|
 \le\frac{\pi^2u_2}{4u_1^2}\,\frac{\ell_{k+1}^2}{\tau_{k-1}\ell_{k-1}^3}\ll_{u_1,u_2}\frac1{\tau_k\ell_k}\ll\frac1k,
\end{equation}
where $\tau_{k\pm1}\asymp\tau_k$ and $\ell_{k\pm1}\asymp\ell_k$ uniformly for $k\ge k_0$ by \CGeqref{eq:rvm-gap-bounds}, and the last step is \CGeqref{eq:tau-ell-k}.  Inserting \CGeqref{eq:kappa-bounds} into \CGeqref{eq:summation-by-parts} gives $|\sum_{k_0}^{N}E_k|\ll\ell_{N+1}+\ell_{k_0+1}+\sum_{k\le N}k^{-1}\ll1+\log N$.  The indices $k<k_0$ contribute at most $k_0$.  All constants depend only on $u_1,u_2$.
\end{proof}

\begin{corollary}[Point centroid]\label{cor:point-centroid}
Uniformly for $u\in[u_1,u_2]\subset(0,\infty)$,
\begin{equation}\label{eq:point-centroid}
 C_{\RVM}^{\pts}(N;u)=O_{u_1,u_2}\!\left(\frac{1+\log N}{N}\right)\longrightarrow0 .
\end{equation}
By direct averaging, Theorem~\CGref{thm:exp-sum} gives only $C_{\RVM}^{\pathc}(N;u)=O(\log N)$ for the path centroid \CGeqref{eq:path-centroid}, not even boundedness; the path centroid converges by a second summation by parts (Theorem~\CGref{thm:clock-centroid}).
\end{corollary}

\begin{proof}
The first assertion is \CGeqref{eq:exp-sum-bound} divided by $N$.  For the second, the identity $C_{\RVM}^{\pathc}=N^{-1}\sum_{n\le N}S_{\RVM}(n;u)$ with $S_{\RVM}(n;u)=O(\log n)$ gives $O(\log N)$ and nothing better.
\end{proof}

Define the midpoint count
\begin{equation}\label{eq:n-RVM}
 n_{\RVM}(H):=\#\{k:\tau_k\le H\}.
\end{equation}
For $x>0$ and $n_{\RVM}(H)\ge1$ define
\begin{equation}\label{eq:MH-RVM}
 M_H^{\RVM}(x):=\frac1{n_{\RVM}(H)}\sum_{\tau_k\le H}x^{i\tau_k},
\end{equation}
and for the ordinates of $\zeta$
\begin{equation}\label{eq:MH-genuine}
 M_H(x):=\frac1{N_\zeta(H)}\sum_{\gamma_k\le H}x^{i\gamma_k}.
\end{equation}

\begin{corollary}[Synthetic height-transform decay]\label{cor:MH-decay}
For every fixed $x>1$, with $N=n_{\RVM}(H)$,
\begin{equation}\label{eq:MH-decay}
 M_H^{\RVM}(x)=O_x\!\left(\frac{1+\log N}{N}\right)\longrightarrow0\qquad(H\to\infty).
\end{equation}
\end{corollary}

\begin{proof}
Take $u=\log x>0$ in Theorem~\CGref{thm:exp-sum}; \CGeqref{eq:MH-RVM} is the point centroid at $N=n_{\RVM}(H)$.
\end{proof}

\begin{remark}\label{rem:MH-not-mertens}
Neither \CGeqref{eq:MH-RVM} nor \CGeqref{eq:MH-genuine} is the Mertens function; both are normalized phase averages.
\end{remark}

\section{Exponential sums over the ordinates of \texorpdfstring{$\zeta$}{zeta}: centroids, Cram\'er's function and the window form of the Riemann hypothesis}
\label{sec:ordinate-sums}

Section~\CGref{sec:exp-sums} bounded the exponential sums of the clock.  This section studies the same sums over the ordinates $\gamma_k$ of $\zeta$ and compares them with the sums over the complete zeros.  The comparison gives three equivalent forms of the Riemann hypothesis written with the ordinates alone (Proposition~\CGref{prop:imaginary-frequency}, Theorem~\CGref{thm:window-RH}, Corollary~\CGref{cor:one-window}), which we have not found in the literature consulted (Remark~\CGref{rem:window}(d)), an exact formula for the defect between the two sums (Proposition~\CGref{prop:pair-decomposition}, \CGeqref{eq:window-residual}), an unconditional identity between the rung functional of $\zeta_{Z1}$ and the zeros off the critical line (Theorem~\CGref{thm:rung-zero-sum}), and the value of the triangular centroid of $e^{i\gamma_k}$ under RH (Theorem~\CGref{thm:centroid-RH}).  The clock centroid converges without any hypothesis (Theorem~\CGref{thm:clock-centroid}).

Throughout, $\rho$ runs over the nontrivial zeros of $\zeta$ with multiplicity, and the zeros with positive ordinate are listed as $\rho_k=\frac12+d_k+i\gamma_k$, $k\ge1$, so that $(\gamma_k)$ is the list of Section~\CGref{sec:clock} and $d_k=\Re\rho_k-\frac12$; zeros with equal ordinate are listed in any order.  By the functional equation the zeros with $d_k\ne0$ come in pairs $\{\rho,1-\bar\rho\}$ with the same ordinate, opposite $d$ and equal multiplicity, and $|d_k|<\frac12$ because $\zeta$ has no zeros on $\Re s=0$ and $\Re s=1$ \cite{Titchmarsh1986}.  RH is the statement that $d_k=0$ for every $k$.  For a function $\mathsf f\in C_c^\infty(\R)$ we write
\begin{equation}\label{eq:f-hat}
 \widehat{\mathsf f}(z):=\int_{\R}\mathsf f(u)\,e^{iuz}\,du\qquad(z\in\C);
\end{equation}
for real $\mathsf f$ and real $t$, $\Re\widehat{\mathsf f}(t)=\mathsf c_{\mathsf f}(t)$ is the cosine transform of \CGeqref{eq:transport-defs}.  If $\operatorname{supp}\mathsf f\subset[u_1,u_2]$, $M$ integrations by parts give
\begin{equation}\label{eq:f-hat-decay}
 |\widehat{\mathsf f}(t+iy)|\le C_{\mathsf f,M}\,e^{u_2|y|}(1+|t|)^{-M}\qquad(t,y\in\R,\ M\ge0).
\end{equation}

\subsection{The ordinate transform and Cram\'er's function}
\label{subsec:cramer}

For $\Im w>0$ define
\begin{equation}\label{eq:ordinate-transforms}
 \mathscr Z_\gamma(w):=\sum_{k\ge1}e^{iw\gamma_k},\qquad
 \mathscr Z_\tau(w):=\sum_{k\ge1}e^{iw\tau_k},\qquad
 \mathscr Z_\xi(w):=\sum_{k\ge1}e^{w(\rho_k-\frac12)}=e^{-w/2}\,V_+(w),
\end{equation}
where
\begin{equation}\label{eq:cramer-V}
 V_+(z):=\sum_{\Im\rho>0}e^{\rho z}\qquad(\Im z>0)
\end{equation}
is Cram\'er's function \cite{Cramer1919}.  Since $|e^{iw\gamma}|=e^{-\gamma\Im w}$, $|e^{w(\rho-\frac12)}|\le e^{|\Re w|/2}e^{-\gamma\Im w}$ and $N_\zeta(T)\ll T\log T$, the three series converge absolutely and locally uniformly on $\Im w>0$ and define holomorphic functions there; summation by parts gives
\begin{equation}\label{eq:Z-gamma-growth}
 |\mathscr Z_\gamma(u+iy)|\le\sum_ke^{-y\gamma_k}\le C\,y^{-2}\qquad(u\in\R,\ 0<y\le1).
\end{equation}
The defect between the ordinate transform and the zero transform is
\begin{equation}\label{eq:ordinate-defect}
 \mathfrak X(w):=\mathscr Z_\gamma(w)-\mathscr Z_\xi(w)\qquad(\Im w>0).
\end{equation}

\begin{proposition}[Pair decomposition of the defect]\label{prop:pair-decomposition}
For $\Im w>0$,
\begin{equation}\label{eq:pair-decomposition}
 \mathfrak X(w)=-\sum_{k\ge1}\bigl(\cosh(wd_k)-1\bigr)e^{iw\gamma_k},
\end{equation}
the series converging absolutely; here $k$ runs over both members of each pair $\{\frac12\pm d+i\gamma\}$, with multiplicity, and only the zeros with $d_k\ne0$ contribute.  In particular $\mathfrak X\equiv0$ under RH, and for $y>0$
\begin{equation}\label{eq:defect-imaginary}
 \mathfrak X(iy)=\sum_{k\ge1}\bigl(1-\cos(yd_k)\bigr)e^{-y\gamma_k}\ \ge0 .
\end{equation}
\end{proposition}

\begin{proof}
$e^{w(\rho_k-\frac12)}=e^{wd_k}e^{iw\gamma_k}$.  The zeros with $d_k\ne0$ form pairs with ordinates $\gamma$ and abscissas $\pm d$, so $\sum_ke^{wd_k}e^{iw\gamma_k}=\sum_k\cosh(wd_k)e^{iw\gamma_k}$, the rearrangement being allowed by absolute convergence.  Subtracting from $\mathscr Z_\gamma$ gives \CGeqref{eq:pair-decomposition}, and $\cosh(iyd)=\cos(yd)$ gives \CGeqref{eq:defect-imaginary}.
\end{proof}

\begin{proposition}[RH at one imaginary frequency]\label{prop:imaginary-frequency}
For every $y_0\in(0,4\pi]$,
\begin{equation}\label{eq:imaginary-frequency}
 \mathrm{RH}\iff\sum_{k\ge1}e^{-y_0\gamma_k}=e^{-iy_0/2}\,V_+(iy_0) .
\end{equation}
\end{proposition}

\begin{proof}
The right side of \CGeqref{eq:imaginary-frequency} is $\mathscr Z_\xi(iy_0)$ and the left side is $\mathscr Z_\gamma(iy_0)$.  Under RH, $\mathfrak X\equiv0$.  Conversely, if $\mathfrak X(iy_0)=0$, every term of \CGeqref{eq:defect-imaginary} vanishes, so $y_0d_k\in2\pi\Z$ for every $k$; since $|y_0d_k|<4\pi\cdot\frac12=2\pi$, $d_k=0$.
\end{proof}

The right side of \CGeqref{eq:imaginary-frequency} has a representation that does not use the zeros (Proposition~\CGref{prop:cramer-zero-free}; $iy_0\in\Omega_\theta$).  The identity is invisible numerically: if every zero up to height $T_{\rm ver}$ is on the critical line, $0\le\mathfrak X(iy)\le2\sum_{\gamma_k>T_{\rm ver}}e^{-y\gamma_k}$, which is of size $T_{\rm ver}\log T_{\rm ver}\,e^{-yT_{\rm ver}}$.  At a real frequency the coefficients $1-\cosh(ud_k)$ of \CGeqref{eq:pair-decomposition} keep a sign but the phases $e^{iu\gamma_k}$ do not, and the argument of the proof does not apply; Theorem~\CGref{thm:window-RH} replaces one real point by an interval.

Here and below, a \emph{zero-free representation} of $V_+$ is a representation that involves no sum over the zeros of $\zeta$; the term does not refer to zeros of $V_+$.

\begin{proposition}[Zero-free representation of Cram\'er's function]\label{prop:cramer-zero-free}
Let $a>0$ and $0<\theta<\frac\pi2$, and put
\begin{equation}\label{eq:g-arch}
 g(s):=\frac1s+\frac1{s-1}-\frac12\log\pi+\frac12\,\frac{\Gamma'}{\Gamma}\Bigl(\frac s2\Bigr),\qquad\text{so that}\qquad \frac{\xi'}{\xi}(s)=g(s)+\frac{\zeta'}{\zeta}(s).
\end{equation}
For $\Im z>0$,
\begin{equation}\label{eq:cramer-zero-free}
 2\pi i\,V_+(z)=\mathscr V_{\rm b}(z)+\mathscr V_{\rm R}(z)+\mathscr V_{\rm L}(z)+\mathscr V_{\rm p}(z),
\end{equation}
\begin{equation}\label{eq:cramer-pieces}
\begin{aligned}
 \mathscr V_{\rm b}(z)&=\int_{-a}^{1+a}e^{\sigma z}\,\frac{\xi'}{\xi}(\sigma)\,d\sigma,\qquad
 \mathscr V_{\rm R}(z)=i\,e^{(1+a)z}\int_0^{\infty e^{i\theta}}e^{itz}\,g(1+a+it)\,dt,\\
 \mathscr V_{\rm L}(z)&=i\,e^{-az}\int_0^{\infty e^{i\theta}}e^{itz}\,g(1+a-it)\,dt,\qquad
 \mathscr V_{\rm p}(z)=\sum_{n\ge2}\Lambda(n)\Bigl(\frac{e^{(1+a)(z-\log n)}}{z-\log n}+\frac{n^{-1-a}e^{-az}}{z+\log n}\Bigr),
\end{aligned}
\end{equation}
the two $t$-integrals being taken along the ray $\arg t=\theta$.  The right side of \CGeqref{eq:cramer-zero-free} does not depend on $a$ and $\theta$, and it continues $V_+$ meromorphically to the half-plane $\Omega_\theta:=\{z=x+iy:\ x\sin\theta+y\cos\theta>0\}$, which contains the positive real axis; in $\Omega_\theta$ the only singularities are simple poles at the points $z=\log n$, $n$ a prime power, with residues $\Lambda(n)/2\pi i$.  In particular $\mathscr Z_\xi(u)=e^{-u/2}V_+(u)$ is real-analytic on $(0,\infty)\setminus\{\log p^m\}$, and
\begin{equation}\label{eq:cramer-real-window}
 2\Re V_+(u)=e^u-\frac1{e^{2u}-1}\qquad(\log2<u<\log3).
\end{equation}
\end{proposition}

\begin{proof}
\emph{Contour.}  $\xi$ is entire, its zeros are the $\rho$, and $\xi(\sigma)>0$ for real $\sigma$, so $\xi'/\xi$ is meromorphic with simple poles at the $\rho$ (residue the multiplicity) and no pole on $\R$.  Fix $z=x+iy$ with $y>0$ and let $R_T$ be the rectangle with vertices $-a$, $1+a$, $1+a+iT$, $-a+iT$, with $T$ not an ordinate, oriented positively: the bottom side $s=\sigma$, $-a\le\sigma\le1+a$; the right side $s=1+a+it$, $t$ from $0$ to $T$; the top side from $1+a+iT$ to $-a+iT$; the left side $s=-a+it$, $t$ from $T$ to $0$.  The residue theorem gives $\sum_{0<\gamma<T}e^{\rho z}=\frac1{2\pi i}\oint_{\partial R_T}e^{sz}\frac{\xi'}{\xi}(s)\,ds$.  There are $T_j\in[j,j+1]$ with $\frac{\zeta'}{\zeta}(\sigma+iT_j)\ll\log^2T_j$ uniformly for $-1\le\sigma\le2$ \cite[Ch.~17]{Davenport1980}; with $\xi'/\xi(s)=-\xi'/\xi(1-s)$, the boundedness of $\zeta'/\zeta$ on $\Re s\ge2$ and Stirling's formula, $\frac{\xi'}{\xi}(\sigma+iT_j)\ll\log^2T_j$ for $-a\le\sigma\le1+a$.  Since $|e^{sz}|\le e^{(1+a)|x|}e^{-T_jy}$ on the top side, its contribution tends to $0$, and the left side tends to $V_+(z)$.  On the right side $s=1+a+it$, $\frac{\xi'}{\xi}=g+\frac{\zeta'}{\zeta}$ with $\frac{\zeta'}{\zeta}(s)=-\sum_n\Lambda(n)n^{-s}$; on the left side $s=-a+it$, $\frac{\xi'}{\xi}(s)=-g(1+a-it)-\frac{\zeta'}{\zeta}(1+a-it)$.  The integrals converge absolutely because $|e^{(c+it)z}|=e^{cx-ty}$ and $g(1+a\pm it)\ll\log(2+t)$, and the Dirichlet series may be integrated termwise; with $\int_0^\infty e^{it(z\mp\log n)}\,dt=i/(z\mp\log n)$ for $\Im z>0$ this gives \CGeqref{eq:cramer-zero-free} with the rays on the positive real axis ($\theta=0$).

\emph{Rotation.}  For $0\le\varphi\le\theta$ and $t=re^{i\varphi}$, $\Im(tz)=r(x\sin\varphi+y\cos\varphi)$; for $z\in\Omega_\theta$ with $y>0$ the sinusoid $\varphi\mapsto x\sin\varphi+y\cos\varphi$ is positive at $\varphi=0$ and $\varphi=\theta$, hence on $[0,\theta]$, since $\theta<\pi$.  The singularities of $t\mapsto g(1+a+it)$ lie on the positive imaginary axis and those of $t\mapsto g(1+a-it)$ on the negative imaginary axis, and $g(1+a\pm it)\ll\log(2+|t|)$ in the closed sector $0\le\arg t\le\theta$.  Cauchy's theorem in the sector moves the rays to $\arg t=\theta$.  The rotated integrals converge absolutely and locally uniformly on $\Omega_\theta$, where they are holomorphic; $\mathscr V_{\rm b}$ is entire; and $\mathscr V_{\rm p}$ converges absolutely away from $\{\pm\log n\}$ and is meromorphic on $\C$ with simple poles at $\pm\log n$, of residue $\Lambda(n)$ at $\log n$; the points $-\log n$ are not in $\Omega_\theta$.  For two admissible choices $(a,\theta)$ and $(a',\theta')$ the two expressions agree with $V_+$ on $\Omega_\theta\cap\Omega_{\theta'}\cap\{\Im z>0\}$, which is non-empty and open; $\Omega_\theta\cap\Omega_{\theta'}$ is a half-plane, hence connected, and the identity theorem gives agreement on all of it.  The continuation is therefore unique and independent of $a$ and $\theta$.

\emph{The window.}  Von Mangoldt's explicit formula $\psi_0(x)=x-\sum_\rho x^\rho/\rho-\log2\pi-\frac12\log(1-x^{-2})$ \cite[Ch.~3]{Edwards1974}, \cite{Guinand1948,Weil1952}, differentiated in $\mathscr D'(1,\infty)$ and written in $u=\log x$, is $\sum_\rho e^{\rho u}=e^u-(e^{2u}-1)^{-1}-\sum_n\Lambda(n)\delta_{\log n}$ in $\mathscr D'(0,\infty)$, the sum over $\rho$ being taken in symmetric order.  Pairing $V_+(\cdot+iy)$ with $\mathsf f\in C_c^\infty(0,\infty)$ gives $\sum_{\Im\rho>0}e^{iy\rho}\int\mathsf f(u)e^{\rho u}du$, and by \CGeqref{eq:f-hat-decay} and dominated convergence its limit as $y\downarrow0$ is $\sum_{\Im\rho>0}\int\mathsf f(u)e^{\rho u}du$.  The zeros with negative ordinate are the conjugates, so $2\Re\,\mathrm{bv}\,V_+=e^u-(e^{2u}-1)^{-1}-\sum_n\Lambda(n)\delta_{\log n}$ in $\mathscr D'(0,\infty)$, where $\mathrm{bv}$ denotes the distributional boundary value from $\Im z>0$.  On $(\log2,\log3)$ there is no atom and $V_+$ is analytic, which gives \CGeqref{eq:cramer-real-window}.
\end{proof}

The full meromorphic continuation of $V_+$, with a logarithmic singularity at $0$, is Cram\'er's theorem \cite{Cramer1919}; only the part in $\Omega_\theta$ is used here.  The representation \CGeqref{eq:cramer-zero-free} is the unilateral form of the explicit formula: on $(\log2,\log3)$ there are no prime atoms, but the prime terms $\Lambda(n)e^{(1+a)(u-\log n)}/(u-\log n)$ of $\mathscr V_{\rm p}$ do not vanish there.  Numerically (Experiment~\CGref{exp:cramer}) $e^{-1/2}V_+(1)=0.77689426079325\ldots+0.76785566886\ldots\,i$, the same for $a=1$ and $a=2$ up to the truncation of the prime sum.

\subsection{The window form of the Riemann hypothesis}
\label{subsec:window}

\begin{lemma}[Vanishing boundary values]\label{lem:edge}
Let $J\subset\R$ be an open interval, $h>0$, and $F$ holomorphic on $J\times(0,h)$ with $|F(u+iy)|\le Cy^{-m}$ there for some $m\ge0$.  If $\lim_{y\downarrow0}\int F(u+iy)\mathsf f(u)\,du=0$ for every $\mathsf f\in C_c^\infty(J)$, then $F\equiv0$.
\end{lemma}

\begin{proof}
Fix $[a',b']\subset J$, $w_0\in(a',b')\times(0,h)$, and let $F_1(w)=\int_{w_0}^wF$, $F_{j+1}(w)=\int_{w_0}^wF_j$, the integrals taken along paths in the rectangle $R=(a',b')\times(0,h)$; $R$ is convex, so by Cauchy's theorem they do not depend on the path, and $F_{j+1}'=F_j$.  Integrating along vertical segments, $|F_j(u+iy)|\le C_jy^{-m+j}$ for $j<m$, $|F_m|\le C_m(1+|\log y|)$, and $F_{m+1}$ is bounded on $(a',b')\times(0,h/2]$.  For $z,w$ in this convex set, $F_{m+2}(z)-F_{m+2}(w)=\int_{[w,z]}F_{m+1}(\zeta)\,d\zeta$, so $F_{m+2}$ is uniformly Lipschitz there and extends continuously to $[a',b']$, with boundary function $\beta$.  For $\mathsf f\in C_c^\infty((a',b'))$, $m+2$ integrations by parts in $u$ (using $\partial_uF_{j+1}=F_j$) give $\int F(u+iy)\mathsf f(u)du=(-1)^{m+2}\int F_{m+2}(u+iy)\mathsf f^{(m+2)}(u)du$, and letting $y\downarrow0$, by the uniform convergence of $F_{m+2}(\cdot+iy)$ to $\beta$, $\int\beta\,\mathsf f^{(m+2)}=0$, that is, $\beta^{(m+2)}=0$ in $\mathscr D'((a',b'))$.  A distribution on an interval with vanishing derivative is a constant, because every $\psi\in C_c^\infty((a',b'))$ with $\int\psi=0$ is the derivative of a test function; applied $m+2$ times, $\beta$ coincides in $\mathscr D'$ with a polynomial $P$ of degree at most $m+1$, and, $\beta$ being continuous, $\beta=P$ at every point of $(a',b')$.  The function $E=F_{m+2}-P$ is holomorphic in $R$, continuous up to $(a',b')$, and vanishes identically there.  By the Schwarz reflection principle \cite[Thm.~11.14]{Rudin1987} (whose hypothesis $\Im E(z_n)\to0$ as $z_n\to u\in(a',b')$ holds), the function equal to $E$ on $R$, to $0$ on $(a',b')$ and to $\overline{E(\bar z)}$ on the reflected rectangle is holomorphic on $(a',b')\times(-h,h)$; it vanishes on $(a',b')$, hence $E\equiv0$, $F_{m+2}=P$ on $R$ and $F=P^{(m+2)}=0$.  Since $J\times(0,h)$ is connected, $F\equiv0$.
\end{proof}

\begin{theorem}[The window form of the Riemann hypothesis]\label{thm:window-RH}
Let $J\subset(0,\infty)$ be an open interval containing no point $\log p^m$, for instance $J=(\log2,\log3)$.  Without any hypothesis, for every $\mathsf f\in C_c^\infty(J)$,
\begin{equation}\label{eq:window-residual}
 \sum_{k\ge1}\widehat{\mathsf f}(\gamma_k)-\int_J\mathsf f(u)\,e^{-u/2}V_+(u)\,du
 =\sum_{k\ge1}\Bigl(\widehat{\mathsf f}(\gamma_k)-\tfrac12\widehat{\mathsf f}(\gamma_k-id_k)-\tfrac12\widehat{\mathsf f}(\gamma_k+id_k)\Bigr),
\end{equation}
both series converging absolutely; only the zeros with $d_k\ne0$ contribute to the right side.  The following are equivalent:
\begin{enumerate}[label=\textup{(\roman*)}]
\item RH;
\item for every real $\mathsf f\in C_c^\infty(J)$, $\displaystyle\sum_{k\ge1}\widehat{\mathsf f}(\gamma_k)=\int_J\mathsf f(u)\,e^{-u/2}V_+(u)\,du$;
\item the triangular centroids $C^{\pathc}_\zeta(N;u):=\sum_{k\le N}\bigl(1-\frac{k-1}N\bigr)e^{iu\gamma_k}$ converge in $\mathscr D'(J)$ to $e^{-u/2}V_+(u)$.
\end{enumerate}
\end{theorem}

\begin{proof}
For $y>0$, $\int\mathscr Z_\gamma(u+iy)\mathsf f(u)du=\sum_ke^{-y\gamma_k}\widehat{\mathsf f}(\gamma_k)$ and $\int\mathscr Z_\xi(u+iy)\mathsf f(u)du=\sum_ke^{iy(d_k+i\gamma_k)}\widehat{\mathsf f}(\gamma_k-id_k)$.  By \CGeqref{eq:f-hat-decay} with $|d_k|<\frac12$ and dominated convergence, the limits as $y\downarrow0$ are $\sum_k\widehat{\mathsf f}(\gamma_k)$ and $\sum_k\widehat{\mathsf f}(\gamma_k-id_k)$; the second equals $\int\mathsf f(u)e^{-u/2}V_+(u)du$ because $\mathscr Z_\xi$ is analytic on a neighbourhood of $\operatorname{supp}\mathsf f$ (Proposition~\CGref{prop:cramer-zero-free}).  The multiset $\{(\gamma_k,d_k)\}$ is invariant under $d_k\mapsto-d_k$ (the functional equation pairs $\rho$ with $1-\bar\rho$, with equal multiplicity), so $\sum_k\widehat{\mathsf f}(\gamma_k-id_k)=\sum_k\widehat{\mathsf f}(\gamma_k+id_k)$, for real or complex $\mathsf f$; this gives \CGeqref{eq:window-residual}.  (i)$\Rightarrow$(ii): every term of the right side of \CGeqref{eq:window-residual} vanishes.  For (ii)$\Leftrightarrow$(iii): $\int C^{\pathc}_\zeta(N;u)\mathsf f(u)du=\sum_{k\le N}(1-\frac{k-1}N)\widehat{\mathsf f}(\gamma_k)$ is a Ces\`aro mean of the absolutely convergent series $\sum_k\widehat{\mathsf f}(\gamma_k)$, so for every $\mathsf f\in C_c^\infty(J)$ it converges to $\sum_k\widehat{\mathsf f}(\gamma_k)$, with no hypothesis.  (ii)$\Rightarrow$(iii): by (ii) and linearity this limit is $\int\mathsf f(u)e^{-u/2}V_+(u)du$ for every complex $\mathsf f$, which is the convergence in $\mathscr D'(J)$.  (iii)$\Rightarrow$(ii): the limit of the pairings is then both $\int\mathsf f e^{-u/2}V_+$ and $\sum_k\widehat{\mathsf f}(\gamma_k)$.  (ii)$\Rightarrow$(i): $\mathfrak X=\mathscr Z_\gamma-\mathscr Z_\xi$ is holomorphic on the half-plane $\Im w>0$, which is connected.  Fix a compact interval $J'\subset J$.  By \CGeqref{eq:Z-gamma-growth}, and because $\mathscr Z_\xi$ is holomorphic, hence bounded, on a neighbourhood of $J'\times[0,1]$ (its only singularities in $\Omega_\theta$ are the points $\log p^m\notin J$, Proposition~\CGref{prop:cramer-zero-free}), $|\mathfrak X(u+iy)|\le C_{J'}y^{-2}$ on $J'\times(0,1]$.  By (ii), linearity and the first part, $\lim_{y\downarrow0}\int\mathfrak X(u+iy)\mathsf f(u)du=0$ for every $\mathsf f\in C_c^\infty(J)$; this is the vanishing of the distributional boundary value.  Lemma~\CGref{lem:edge}, applied on the interior of $J'$ with $m=2$ and $h=1$, gives $\mathfrak X=0$ there, and the identity theorem gives $\mathfrak X\equiv0$ on $\Im w>0$.  In particular $\mathfrak X(iy_0)=0$, and Proposition~\CGref{prop:imaginary-frequency}, through the positivity of every term of \CGeqref{eq:defect-imaginary}, gives $d_k=0$ for every $k$, that is, RH.
\end{proof}

\begin{remark}[What the window form needs and what it does not]\label{rem:window}
\begin{enumerate}[label=\textup{(\alph*)}]
\item After inserting \CGeqref{eq:cramer-zero-free}, both sides of (ii) are expressed without the abscissas of the zeros: the left side uses the ordinates only, the right side is given in terms of $\Lambda$, $\Gamma$ and the values of $\zeta$ on a real segment.  Identity (ii) is the explicit formula of Guinand and Weil \cite{Guinand1948,Weil1952} written for the measure of the ordinates instead of the measure of the zeros, and required only in one gap between prime powers, where its prime side has no atom.
\item The interval cannot be replaced by the argument at one real point: at $u=1$ the defect $\mathfrak X(1+i0)$ is a sum with positive coefficients $\cosh d_k-1$ and free phases $e^{i\gamma_k}$, and the triangular centroid at $u=1$ (Theorem~\CGref{thm:centroid-RH}) is implied by RH, not shown equivalent to it.  Nor can the one-sided transform be replaced by the symmetric one: the cosine part of (ii) controls only $\Re\,\mathrm{bv}\,\mathfrak X$ on $J$ (Corollary~\CGref{cor:one-window}), and Lemma~\CGref{lem:edge} needs the full boundary value.  The window itself is used only to make $\mathscr Z_\xi$ holomorphic across $J$, so that the right side of (ii) is an ordinary integral.  Cancellation on $J$ between the terms of different zeros off the line is excluded by the half-plane: $\mathfrak X$ is holomorphic in $\Im w>0$ because the ordinates are positive, and the boundary value of such a function of polynomial growth cannot vanish on an interval unless the function vanishes (Lemma~\CGref{lem:edge}); positivity at an imaginary frequency then gives $d_k=0$ for every $k$.
\item The Davenport--Heilbronn function \cite{DavenportHeilbronn1936a,DavenportHeilbronn1936b} has a functional equation of the same type and a Dirichlet series, and zeros off its critical line; an argument proving (ii) must therefore use the Euler product.
\item Literature.  Cram\'er \cite{Cramer1919} continued $V_+$; Guillera \cite{Guillera2013} proved height-regularized forms of the real part of such sums at non-prime-power arguments and stated Abel summation in the ordinal index (his Corollary~4.4), leaving a version with ordinal weights depending on $N$ as a conjecture (his Conjecture~4.5); Arias de Reyna \cite{AriasdeReyna2024} proved that RH is equivalent to the temperedness of a measure built from the primes, which then equals the Fourier transform of the symmetric measure of the ordinates.  The form (ii) differs in three respects: it is local (one window without prime atoms), one-sided ($\sum_ke^{iu\gamma_k}$ against $V_+$), and its mechanism is Lemma~\CGref{lem:edge} with the positivity \CGeqref{eq:defect-imaginary} at an imaginary frequency.  We have not found it in the literature; the search included the Fourier-quasicrystal literature \cite{LevOlevskii2015,KurasovSarnak2020} but was not exhaustive.
\end{enumerate}
\end{remark}

\subsection{One window is one cell: the cosine and sine identities}
\label{subsec:one-window}

Define the distribution
\begin{equation}\label{eq:Y-offline}
 \langle\mathscr Y,\mathsf f\rangle:=\sum_{k\ge1}\int_0^\infty\mathsf f(u)\bigl(\cosh(d_ku)-1\bigr)\cos(\gamma_ku)\,du\qquad(\mathsf f\in C_c^\infty(0,\infty)),
\end{equation}
which converges absolutely by \CGeqref{eq:f-hat-decay}; only zeros off the critical line contribute, and $\mathscr Y=0$ under RH.  By \CGeqref{eq:pair-decomposition}, $\Re\,\mathrm{bv}\,\mathfrak X=-\mathscr Y$ on every interval of $(0,\infty)$ free of the points $\log p^m$.  Let $J(x)=\sum_{p^m\le x}\frac1m$ be Riemann's prime-power count.

\begin{theorem}[The rung functional of the endpoint is a sum over the zeros off the line]\label{thm:rung-zero-sum}
For every real $\mathsf f\in C_c^\infty(0,\infty)$,
\begin{equation}\label{eq:rung-zero-sum}
 \mathsf R_{\mathsf f}(0)-\frac12\sum_{n\ge2}\frac{\Lambda(n)}{\sqrt n}\,\mathsf f(\log n)=\langle\mathscr Y,\mathsf f\rangle .
\end{equation}
Equivalently, in $\mathscr D'(1,\infty)$,
\begin{equation}\label{eq:cell-defect-identity}
 (\log x)\Bigl(\frac{dJ^{\rm M}_{Z1}}{dx}-\frac{dJ}{dx}\Bigr)=\frac2{\sqrt x}\,\mathscr Y(\log x).
\end{equation}
Here the derivatives are distributional derivatives on $(1,\infty)$ ($dJ/dx$ is the measure $\sum_{p^m}m^{-1}\delta_{p^m}$), and $\mathscr Y(\log x)$ denotes the distribution $\sum_k(\cosh(d_k\log x)-1)\cos(\gamma_k\log x)$ on $(1,\infty)$, the series converging in $\mathscr D'(1,\infty)$.  In particular, if the support of $\mathsf f$ contains no point $\log n$, then $\mathsf R_{\mathsf f}(0)=\langle\mathscr Y,\mathsf f\rangle$ and the cosine difference of Proposition~\CGref{prop:rung-functional} is
\begin{equation}\label{eq:cosine-difference}
 \mathsf S_{\mathsf f}(0)=\sum_{k\ge1}\bigl(\mathsf c_{\mathsf f}(\gamma_k)-\mathsf c_{\mathsf f}(\tau_k)\bigr)=\mathsf R_{\mathsf f}(1)-\langle\mathscr Y,\mathsf f\rangle .
\end{equation}
\end{theorem}

\begin{proof}
By \CGeqref{eq:rung-functional}, $\mathsf R_{\mathsf f}(0)=\mathsf k_{\mathsf f}-\sum_k\mathsf c_{\mathsf f}(\gamma_k)$.  Pair the explicit formula $\sum_n\Lambda(n)\delta_n=1-\sum_\rho x^{\rho-1}-x^{-1}(x^2-1)^{-1}$ in $\mathscr D'(1,\infty)$ (proof of Proposition~\CGref{prop:cramer-zero-free}) with $\frac12\mathsf f(\log x)x^{-1/2}$ and put $u=\log x$: the term $1$ and the last term give $\mathsf k_{\mathsf f}$ as in the proof of Proposition~\CGref{prop:rung-functional}, and the zeros give $\frac12\sum_\rho\int\mathsf f(u)e^{(\rho-\frac12)u}du=\sum_k\int\mathsf f(u)\cosh(d_ku)\cos(\gamma_ku)du$ (pairs $\rho,\bar\rho$ and the pair symmetry in $d$).  Hence $\frac12\sum_n\Lambda(n)n^{-1/2}\mathsf f(\log n)=\mathsf k_{\mathsf f}-\sum_k\int\mathsf f(u)\cosh(d_ku)\cos(\gamma_ku)du$, and subtraction gives \CGeqref{eq:rung-zero-sum}.  Theorem~\CGref{thm:distributional-EF} for $Z1$ and the same explicit formula for $J$, whose derivative is $(\log x)^{-1}\sum_n\Lambda(n)\delta_n$, give \CGeqref{eq:cell-defect-identity} by subtraction.  The last statement is \CGeqref{eq:rung-zero-sum} with the prime sum empty, together with \CGeqref{eq:transport}.
\end{proof}

Theorem~\CGref{thm:rung-zero-sum} makes the defect of the endpoint explicit: on every rung, the defect $\mathsf R_{\mathsf f}(0)$, which (D) requires to vanish (Proposition~\CGref{prop:rung-functional}), is the pairing of $\mathsf f$ with the zeros off the critical line, and the transport from the clock is $\mathsf R_{\mathsf f}(1)-\langle\mathscr Y,\mathsf f\rangle$.  Any lemma that proves $\mathsf S_{\mathsf f}(0)=\mathsf R_{\mathsf f}(1)$ on a rung proves $\langle\mathscr Y,\mathsf f\rangle=0$ there.

\begin{corollary}[One window is one cell]\label{cor:one-window}
Let $J=(\log2,\log3)$.
\begin{enumerate}[label=\textup{(\alph*)}]
\item $\Re\,\mathrm{bv}\,\mathfrak X=0$ on $J$ $\iff$ $\langle\mathscr Y,\mathsf f\rangle=0$ for every $\mathsf f\in C_c^\infty(J)$ $\iff$ $\mathsf R_{\mathsf f}(0)=0$ for every $\mathsf f\in C_c^\infty(J)$ $\iff$ $J^{\rm M}_{Z1}$ is constant on $(2,3)$, the rung $N=3$ of Theorem~\CGref{thm:ladder}.
\item RH holds if and only if $J^{\rm M}_{Z1}$ is constant on $(2,3)$ and, for every real $\mathsf f\in C_c^\infty(J)$, the sine identity
\begin{equation}\label{eq:sine-identity}
 \sum_{k\ge1}\int_J\mathsf f(u)\sin(\gamma_ku)\,du=\int_J\mathsf f(u)\,e^{-u/2}\,\Im V_+(u)\,du
\end{equation}
holds, with $V_+$ given on $J$ by \CGeqref{eq:cramer-zero-free}.
\item Consequently (D), which is equivalent to the absence of mass of $dJ^{\rm M}_{Z1}$ outside the integers on every rung (Theorem~\CGref{thm:ladder}, Corollary~\CGref{cor:ladder-pair}), is equivalent to the single rung $(2,3)$ together with \CGeqref{eq:sine-identity} on $J$.
\end{enumerate}
\end{corollary}

\begin{proof}
(a) On $(2,3)$, $dJ=0$, so by \CGeqref{eq:cell-defect-identity} $J^{\rm M}_{Z1}$ is constant there if and only if $\mathscr Y=0$ on $J$; the other equivalences are \CGeqref{eq:rung-zero-sum} and $\Re\,\mathrm{bv}\,\mathfrak X=-\mathscr Y$.  (b) $\mathrm{bv}\,\mathfrak X=0$ on $J$ if and only if its real and imaginary parts vanish; by Theorem~\CGref{thm:window-RH} and \CGeqref{eq:cramer-real-window} the real part is (a) and the imaginary part is \CGeqref{eq:sine-identity}.  (c) is (b) with the equivalence of (D) and RH (Theorem~\CGref{thm:Hamburger-Z1}).
\end{proof}

\begin{remark}[The conjugate part is not local]\label{rem:conjugate-nonlocal}
Put $Q(u):=2\mathscr Y(u)=u\,e^{u/2}\bigl(\frac{dJ^{\rm M}_{Z1}}{dx}-\frac{dJ}{dx}\bigr)(e^u)$.  For a holomorphic function on the upper half-plane with suitable decay, the boundary values satisfy $\mathrm{bv}=-\frac12(Q+i\mathcal HQ)$ with $\mathcal H$ the Hilbert transform; $\mathfrak X$ grows like $e^{\max|d_k|\,|\Re w|}$ and the relation holds only after a regularization, so the following is a heuristic description, not used in any proof.  If $Q=0$ on $J$, then on $J$, $\mathcal HQ(u)=\frac1\pi\,\mathrm{pv}\!\int_{\R\setminus J}\frac{Q(v)}{u-v}dv$ is produced by the values of $Q$ outside the window: a mass at the end point, $Q=\delta_{\log2}$, vanishes on $J$ and has $\mathcal HQ(u)=1/(\pi(u-\log2))\ne0$ there.  The terms $\Lambda(n)e^{(1+a)(u-\log n)}/(u-\log n)$ of \CGeqref{eq:cramer-pieces} are of this kind.  Thus the cell $(2,3)$ controls the cosine part of the window, and the sine identity \CGeqref{eq:sine-identity} carries the information of all other cells into the window.
\end{remark}

\begin{proposition}[Finitely many zeros off the line]\label{prop:finite-offline}
Suppose that the set of zeros with $d_k\ne0$ is finite and non-empty.  Then for every open interval $I\subset(0,\infty)$ there is a real $\mathsf f\in C_c^\infty(I)$ with $\langle\mathscr Y,\mathsf f\rangle>0$.  Consequently, if $J^{\rm M}_{Z1}$ is constant on one interval $(n,n+1)$, $n\ge1$, then either RH holds or infinitely many zeros lie off the critical line.
\end{proposition}

\begin{proof}
$\mathscr Y$ is then the finite exponential polynomial $\mathscr Y(u)=\sum_{k\in O}(\cosh(d_ku)-1)\cos(\gamma_ku)$, $O$ the finite set of indices with $d_k\ne0$.  Writing $\cosh(du)\cos(\gamma u)=\frac14\sum_{\epsilon,\epsilon'=\pm1}e^{(\epsilon d+i\epsilon'\gamma)u}$, let $A=\max_{k\in O}|d_k|>0$ and $\gamma_0$ the ordinate of an index where it is attained.  The exponent $A+i\gamma_0$ occurs only for the $m\ge1$ indices with $|d_k|=A$, $\gamma_k=\gamma_0$, each with coefficient $\frac14$, so its total coefficient is $m/4>0$, also when several indices share the ordinate $\gamma_0$; the terms $-\cos(\gamma_ku)$ have exponents of real part $0$.  Distinct exponentials are linearly independent on every interval, so $\mathscr Y\not\equiv0$ on $I$, and $\mathscr Y$ is real-analytic.  Take $\mathsf f=\chi\mathscr Y$ with $\chi\ge0$ a bump in $C_c^\infty(I)$ positive on a subinterval; then $\langle\mathscr Y,\mathsf f\rangle=\int\chi\mathscr Y^2>0$.  The last statement follows from \CGeqref{eq:cell-defect-identity}, because $dJ=0$ on $(n,n+1)$.
\end{proof}

The proposition is the contrapositive of the easy direction for a finite exceptional set; with infinitely many zeros off the line $\mathscr Y$ is only a distribution, and a test function supported in one cell does not separate ordinates closer than about $1/\log2$, so infinitely many terms might in principle conspire.

\begin{theorem}[Frullani form of the value at infinity]\label{thm:frullani}
For $K\ge1$ put
\begin{equation}\label{eq:frullani-QK}
 \mathscr Q_K:=2\int_0^\infty\frac{e^{-u/2}}u\sum_{k\le K}\bigl(\cos(\gamma_ku)-\cos(\tau_ku)\bigr)\,du .
\end{equation}
Then $\mathscr Q_K=\sum_{k\le K}\log\frac{\frac14+\tau_k^2}{\frac14+\gamma_k^2}$, the limit $\mathscr Q_\infty=\lim_{K\to\infty}\mathscr Q_K$ exists, and
\begin{equation}\label{eq:frullani-limit}
 \mathscr Q_\infty=\log L+\log\varpi_{Z1}=\log L+\sum_{\text{pairs}}\log c_\gamma\ \ge\ \log L>0.0206 ,
\end{equation}
with $c_\gamma$ the pair factor of \CGeqref{eq:pair-ratio}.  Hence RH holds if and only if $\mathscr Q_\infty=\log L=-\log\varpi_{\RVM}$.
\end{theorem}

\begin{proof}
For real $t$, Frullani's integral gives $\int_0^\infty(1-\cos tu)e^{-u/2}u^{-1}du=\frac12\log(1+4t^2)$, so $\log(\frac14+t^2)=-\log4+2\int_0^\infty(1-\cos tu)e^{-u/2}u^{-1}du$ and $\log\frac{\frac14+\tau^2}{\frac14+\gamma^2}=2\int_0^\infty(\cos\gamma u-\cos\tau u)e^{-u/2}u^{-1}du$, an absolutely convergent integral.  Summing over $k\le K$ gives the first statement.  By Theorem~\CGref{thm:Z1-clock}, $\lim_K\prod_{k\le K}\frac{1/4+\tau_k^2}{1/4+\gamma_k^2}=\varpi_{Z1}/\varpi_{\RVM}=\varpi_{Z1}L$; Theorem~\CGref{thm:varpi-Z1} gives $\log\varpi_{Z1}=\sum\log c_\gamma\ge0$ with equality exactly under RH, and Theorem~\CGref{thm:varpi-less-one} gives $\log L>0.020615$.
\end{proof}

\begin{remark}[The arithmetic of both sides and the size of the defect]\label{rem:frullani-arithmetic}
\begin{enumerate}[label=\textup{(\alph*)}]
\item Inserting the explicit formula, $\sum_k\cos(\gamma_ku)=\mathsf k(u)-\frac12\sum_n\Lambda(n)n^{-1/2}\delta_{\log n}-\mathscr Y(u)$ (proof of Theorem~\CGref{thm:rung-zero-sum}), into \CGeqref{eq:frullani-QK} returns the same identity: by Frullani, $\int_0^\infty e^{-u/2}(\cosh(du)-1)\cos(\gamma u)\frac{du}u=\frac14\log\frac{(\frac14+\gamma^2)^2}{|\rho|^2|1-\rho|^2}=-\frac14\log c_\gamma$ for each zero $\rho=\frac12+d+i\gamma$, so the off-line part of \CGeqref{eq:frullani-QK} is exactly $-\log\varpi_{Z1}$, and what remains is $\zeta(\sigma)\to1$ and the value $\varpi_{\RVM}$ already used in Theorem~\CGref{thm:varpi-Z1}.  With this kernel the arithmetic of the two sides contracts with ratio exactly $1$; a proof of $\log\varpi_{Z1}\le0$ along these lines needs a kernel and a sign on the prime side giving a ratio $<1$, the structure of Weil's positivity and Li's criterion.
\item Every estimate gives $\log\varpi_{Z1}\le\eps$ with $\eps>0$: by Proposition~\CGref{prop:varpi-verified}, $0\le\log\varpi_{Z1}<4\cdot10^{-13}$, and zero-density estimates improve the constant only.  There is no gap below: one pair with $d\to0$ or $\gamma\to\infty$ contributes $\log c_\gamma\approx2d^2/\gamma^2\to0$.
\item The criterion $\log\varpi_{Z1}=\sum\log c_\gamma\ge0$, with equality exactly under RH, is of the type of the integral criteria of Volchkov \cite{Volchkov1995}, Balazard--Saias--Yor \cite{BalazardSaiasYor1999} and Sekatskii--Beltraminelli--Merlini \cite{SekatskiiBeltraminelliMerlini2012}, which also express a sum over the zeros off the line, with weights of first order in $d$, as an explicit quantity; here the weights are of second order, $\log c_\gamma=\frac{2d^2(\gamma^2-\frac14)+d^4}{(\frac14+\gamma^2)^2}+O(d^4\gamma^{-4})$.
\end{enumerate}
\end{remark}

\subsection{Centroids of the two spectra}
\label{subsec:centroids}

The triangular centroid of a list $(o_k)$ at the frequency $u$ is the mean of the partial sums of its exponential walk,
\begin{equation}\label{eq:triangular-centroid}
 C^{\pathc}(N;u)=\frac1N\sum_{n\le N}\sum_{k\le n}e^{iuo_k}=\sum_{k\le N}\Bigl(1-\frac{k-1}N\Bigr)e^{iuo_k};
\end{equation}
for $o_k=\tau_k$ it is the path centroid $C^{\pathc}_{\RVM}(N;u)$ of \CGeqref{eq:path-centroid}, for $o_k=\gamma_k$ it is $C^{\pathc}_\zeta(N;u)$ of Theorem~\CGref{thm:window-RH}.  In three dimensions, with $k$ on the vertical axis, the walk $(\Re S_k,\Im S_k,k)$ has its centroid at $(\Re C^{\pathc},\Im C^{\pathc},\frac{N+1}2)$ (Figures~\CGref{fig:walks3d}, \CGref{fig:walks-proj}).  The normalized sums $S_N/N$ tend to $0$ for both lists (Corollary~\CGref{cor:point-centroid} and Hlawka's theorem \cite{Hlawka1975}); the centroid is the mean of the partial sums, and it does not tend to $0$.

\begin{theorem}[The clock centroid converges]\label{thm:clock-centroid}
Let $u>0$, $E_k=e^{iu\tau_k}$, $S_n=\sum_{k\le n}E_k$, $b_k=(e^{iu(\tau_{k+1}-\tau_k)}-1)^{-1}$ (this is $\mathfrak b^{\rm sbp}_k$ of the proof of Theorem~\CGref{thm:exp-sum}), and let $k_0=k_0(u)$ be as in that proof.  The series
\begin{equation}\label{eq:clock-centroid-limit}
 C^{\pathc}_{\RVM}(\infty;u):=\sum_{k<k_0}E_k-b_{k_0}E_{k_0}+\sum_{k>k_0}(b_{k-1}-b_k)E_k
\end{equation}
converges, $S_n-b_nE_{n+1}=C^{\pathc}_{\RVM}(\infty;u)+O_u(\log n/n)$, and
\begin{equation}\label{eq:clock-centroid-rate}
 C^{\pathc}_{\RVM}(N;u)=C^{\pathc}_{\RVM}(\infty;u)+O_u\Bigl(\frac{(\log N)^2}N\Bigr).
\end{equation}
\end{theorem}

\begin{proof}
Put $d_k^{\rm b}:=b_{k-1}-b_k$.  By \CGeqref{eq:summation-by-parts}, for $n\ge k_0$, $S_n=\sum_{k<k_0}E_k-b_{k_0}E_{k_0}+\sum_{k_0<k\le n}d^{\rm b}_kE_k+b_nE_{n+1}$.  Write $b_k=\beta(D(k))$ with $\beta(x)=(e^{ix}-1)^{-1}$ and $D(y)=u\bigl(\mathcal T(y+\frac12)-\mathcal T(y-\frac12)\bigr)$, $\mathcal T$ the inverse clock.  On $0<x\le\pi$, $|\beta^{(j)}(x)|\ll_jx^{-1-j}$; by \CGeqref{eq:inverse-clock}, differentiating $\mathcal T''=-4\pi^2/(\mathcal T\ell^3)$ with $\ell=\log(\mathcal T/2\pi)$ and $\ell'=\mathcal T'/\mathcal T$ gives $\mathcal T'''=4\pi^2\mathcal T'(\ell^3+3\ell^2)/(\mathcal T^2\ell^6)=8\pi^3(\ell+3)/(\mathcal T^2\ell^5)$, so that $\mathcal T'\asymp\ell^{-1}$, $|\mathcal T''|\asymp1/(\mathcal T\ell^3)$ and $\mathcal T'''\ll1/(\mathcal T^2\ell^4)$; hence $D\asymp1/\ell$, $D'\ll1/(\mathcal T\ell^3)$ and $D''\ll1/(\mathcal T^2\ell^4)$, so $(\beta\circ D)'\ll1/(\mathcal T\ell)$ and $(\beta\circ D)''\ll1/(\mathcal T\ell)^2$.  With $\tau_k\ell_k\asymp k$ \CGeqref{eq:tau-ell-k}: $|d^{\rm b}_k|\ll1/k$ and $|d^{\rm b}_k-d^{\rm b}_{k+1}|\ll1/k^2$.  Abel summation with $S_k=O(\log k)$ (Theorem~\CGref{thm:exp-sum}) gives, for $M>n$, $\sum_{n<k\le M}d^{\rm b}_kE_k=d^{\rm b}_MS_M-d^{\rm b}_{n+1}S_n+\sum_{n<k<M}(d^{\rm b}_k-d^{\rm b}_{k+1})S_k=O(\log n/n)$ uniformly in $M$.  Hence the series \CGeqref{eq:clock-centroid-limit} converges and $S_n=C^{\pathc}_{\RVM}(\infty;u)+b_nE_{n+1}+O(\log n/n)$.  Averaging over $n\le N$, the error terms contribute $O((\log N)^2/N)$ and the indices $n<k_0$ contribute $O(k_0\log k_0/N)$.  Finally $E_{n+1}=b_{n+1}(E_{n+2}-E_{n+1})$ by \CGeqref{eq:Ek-identity}; with $c_n:=b_nb_{n+1}$ the exact identity
\[
 \sum_{n=k_0}^Nb_nE_{n+1}=\sum_{n=k_0}^Nc_n(E_{n+2}-E_{n+1})=c_NE_{N+2}-c_{k_0}E_{k_0+1}-\sum_{n=k_0+1}^N(c_n-c_{n-1})E_{n+1}
\]
holds, and $|c_n|\ll\log^2n$, $|c_n-c_{n-1}|=|b_n|\,|b_{n+1}-b_{n-1}|\le|b_n|(|d^{\rm b}_n|+|d^{\rm b}_{n+1}|)\ll\log n/n$ give $\sum_{k_0\le n\le N}b_nE_{n+1}=O((\log N)^2)$.  This is \CGeqref{eq:clock-centroid-rate}.
\end{proof}

This answers the question left open after Corollary~\CGref{cor:point-centroid}: Theorem~\CGref{thm:exp-sum} alone does not give the convergence of the path centroid, and the second summation by parts does.  Numerically $C^{\pathc}_{\RVM}(\infty;1)=0.041569+1.757541\,i$ (Experiment~\CGref{exp:cramer}).

For the ordinates the same quantity is controlled by Cram\'er's function.  We use the following local Tauberian lemma; Tauberian theorems with local boundary behaviour are developed in \cite{Korevaar2004,DebruyneVindas2019}, and we give a self-contained proof of the case needed.

\begin{lemma}[A local Tauberian lemma]\label{lem:local-tauberian}
Let $F:[0,\infty)\to\C$ be locally absolutely continuous, $F=0$ on $[0,x_0]$, with $|F(x)|\le C\log(2+x)$ and $|F'(x)|\le C\log(2+x)/(1+x)$ almost everywhere.  Suppose that for $\Re s>0$
\begin{equation}\label{eq:tauberian-hyp}
 \int_0^\infty F(x)e^{-sx}\,dx=\frac cs+c'\Log s+\Phi(s),
\end{equation}
with constants $c,c'$, $\Log$ the principal branch on $\Re s>0$, and $\Phi$ holomorphic in a disc $|s|<r$.  Then $F(x)\to c$ as $x\to\infty$.
\end{lemma}

\begin{proof}
Let $K(y)=\kappa_0(\sin(\alpha y)/(\alpha y))^6$ with $\alpha=r/12$ and $\kappa_0$ chosen so that $\int K=1$; its Fourier transform $\widehat K(\xi)=\int K(y)e^{-i\xi y}dy$ is continuous and supported in $[-r/2,r/2]$, and $K(y)\ll y^{-6}$.  Put $f=F-c$ on $(0,\infty)$ and $f=0$ on $(-\infty,0]$; the jump at $0$ is harmless, and the transform used below is the unilateral Laplace transform of $F-c$, which is $c'\Log s+\Phi(s)$ by \CGeqref{eq:tauberian-hyp}.  For $\sigma>0$, $f_\sigma(x)=f(x)e^{-\sigma x}$ is integrable with Fourier transform $c'\Log(\sigma+i\xi)+\Phi(\sigma+i\xi)$ on $|\xi|<r$, so $(f_\sigma*K)(x)=\frac1{2\pi}\int\bigl(c'\Log(\sigma+i\xi)+\Phi(\sigma+i\xi)\bigr)\widehat K(\xi)e^{i\xi x}d\xi$.  As $\sigma\downarrow0$ the left side tends to $(f*K)(x)$ by dominated convergence, and the right side tends to the same integral with $\sigma=0$, because $|\Log(\sigma+i\xi)|\le|\log|\xi||+C$ on $|\xi|\le r/2$; the boundary values are $\Log(i\xi)=\log|\xi|+\frac{i\pi}2\operatorname{sgn}\xi$.  The function $(c'\Log(i\xi)+\Phi(i\xi))\widehat K(\xi)$ is integrable, so $(f*K)(x)\to0$ by the Riemann--Lebesgue lemma.  Finally $f(x)-(f*K)(x)=\int(f(x)-f(x-y))K(y)dy$ (using $\int K=1$); let $x\ge4x_0+4$.  For $|y|\le\sqrt x$, $x-y\ge x-\sqrt x>x_0$, and $|f(x)-f(x-y)|\le|y|\sup_{[x-\sqrt x,x+\sqrt x]}|F'|\ll|y|\log x/x$, which integrates against $K$ to $O(\log x/x)$.  For $|y|>\sqrt x$ with $x-y>0$, $|f(x)-f(x-y)|\le|F(x)|+|F(x-y)|+2|c|\ll\log(2+x)+\log(2+|y|)$; for $|y|>\sqrt x$ with $x-y\le0$, $f(x-y)=0$ and $|f(x)-f(x-y)|=|f(x)|\ll\log(2+x)$.  Both cases integrate against $K(y)\ll y^{-6}$ over $|y|>\sqrt x$ to $O(\log(2+x)\,x^{-5/2})$.  Hence $f(x)\to0$.
\end{proof}

\begin{theorem}[The ordinate centroid under RH]\label{thm:centroid-RH}
Assume RH.  For every $u\in(0,\infty)\setminus\{\log p^m\}$,
\begin{equation}\label{eq:centroid-RH}
 \lim_{N\to\infty}C^{\pathc}_\zeta(N;u)=\lim_{N\to\infty}\sum_{k\le N}\Bigl(1-\frac{k-1}N\Bigr)e^{iu\gamma_k}=e^{-u/2}\,V_+(u).
\end{equation}
For $u=1$ the limit is $e^{-1/2}V_+(1)$, whose real part is $\frac1{2\sqrt e}\bigl(e-\frac1{e^2-1}\bigr)=0.7768942607932\ldots$ by \CGeqref{eq:cramer-real-window}, and whose numerical value is $0.77689426079325\ldots+0.76785566886\ldots\,i$ (Experiment~\CGref{exp:cramer}).
\end{theorem}

\begin{proof}
Let $B_u(t)=\sum_{\gamma_k\le t}e^{iu\gamma_k}$ and $\mathfrak e=N_\zeta-\NR$ as in \CGeqref{eq:frak-e}.  Under RH, $x^\rho=\sqrt x\,x^{i\gamma}$ and Landau's formula $\sum_{0<\gamma\le T}x^\rho=-\frac T{2\pi}\Lambda(x)+O_x(\log T)$ \cite{Landau1912,Gonek1993} at $x=e^u$, which is not a prime power, give $B_u(T)=O_u(\log T)$.

\emph{Step 1 (index form).}  It suffices to take $N$ at the end of a block of equal ordinates.  A block has $m(\gamma)\ll\log\gamma$ elements (proof of Lemma~\CGref{lem:bounded-displacement}), and $|S_n|\le|B_u(\gamma_n)|+m(\gamma_n)\ll\log\gamma_n$; if $N'$ is the end of the block containing $N$, then $C^{\pathc}_\zeta(N')-C^{\pathc}_\zeta(N)=\frac1{N'}\sum_{N<n\le N'}S_n+(\frac1{N'}-\frac1N)\sum_{n\le N}S_n$, and both terms are $O(m(\gamma_N)\log\gamma_N/N)=O(\log^2\gamma_N/N)$.  Put $T=\gamma_N$.  Summing over the block of each ordinate, $\sum_{k\le N}(N-k+1)e^{iu\gamma_k}=\int_{(0,T]}(N+\frac12-N_\zeta^*(t))e^{iut}dN_\zeta(t)$ exactly, with the symmetric count \CGeqref{eq:symmetric-count}: for a block $k=j+1,\dots,j+m$ at the ordinate $\gamma$, $N_\zeta^*(\gamma)=j+\frac m2$ and $\sum_{k=j+1}^{j+m}(N-k+1)=m(N+1)-mj-\frac{m(m+1)}2=m\bigl(N+\frac12-N_\zeta^*(\gamma)\bigr)$, which is the mass of the right side at $\gamma$.  Write $N^*_\zeta=\NR+\mathfrak e^*$, $\mathfrak e^*=\frac12(\mathfrak e(t+0)+\mathfrak e(t-0))$, and $dN_\zeta=\NR'dt+d\mathfrak e$.  With $G_N=N+\frac12-\NR$, continuous, and $G_N(T)=\mathfrak e(T)+\frac12=O(\log T)$,
\[
 \int_{(0,T]}G_Ne^{iut}dN_\zeta=G_N(T)B_u(T)+\int_0^TB_u(t)\NR'(t)\,dt .
\]
The remaining term is $\int\mathfrak e^*e^{iut}dN_\zeta=\int_{t_*}^T\mathfrak e\,e^{iut}\NR'dt+\frac12\int_{(t_*,T]}e^{iut}d(\mathfrak e^2)$, $t_*=\tau_0=4\pi<\gamma_1$ (below $t_*$ the measure $dN_\zeta$ vanishes), because $\mathfrak e^*d\mathfrak e=\frac12d(\mathfrak e^2)$: at a jump from $a$ to $b$, $\frac{a+b}2(b-a)=\frac{b^2-a^2}2$, and between the jumps $\mathfrak e$ is smooth and $\mathfrak e^*=\mathfrak e$.  Integration by parts and Selberg's bound $\int_0^TS(t)^2dt\ll T\log\log T$ \cite{Selberg1946} give $\frac12\int e^{iut}d(\mathfrak e^2)=O(\log^2T)+O(T\log\log T)$.  Let $F_u(t)=\int_{t_*}^t\mathfrak e(v)e^{iuv}dv$.  Since $B_u(t)=\int_{t_*}^te^{iuv}\NR'(v)dv+[e^{iuv}\mathfrak e(v)]_{t_*}^t-iuF_u(t)$ and the first integral is $O_u(\log t)$, $F_u(t)=O_u(\log t)$, so $\int_{t_*}^T\mathfrak e\,e^{iut}\NR'dt=\NR'(T)F_u(T)-\int_{t_*}^TF_u\NR''dt=O(\log^2T)$.  As $N\asymp T\log T$,
\[
 C^{\pathc}_\zeta(N;u)=\frac1N\int_0^TB_u(t)\NR'(t)\,dt+o(1).
\]
\emph{Step 2.}  Let $A_u(x)=\int_0^xB_u$ and $\mathcal M_u(x)=A_u(x)/x$.  If $\mathcal M_u(x)\to c$, write $A_u(t)=ct+r(t)$; for every $\epsilon>0$ there is $t_\epsilon$ with $|r(t)|\le\epsilon t$ for $t\ge t_\epsilon$.  Then $\int_0^TB_u\NR'=\NR'(T)A_u(T)-\int_{t_*}^TA_u\NR''$, with $\NR''(t)=1/(2\pi t)$, $\NR'(T)A_u(T)=c\,T\NR'(T)+o(T\log T)$ and $\int_{t_*}^TA_u\NR''=\frac{c(T-t_*)}{2\pi}+O_\epsilon(1)+O(\epsilon T)$; since $T\NR'(T)-\frac T{2\pi}=\NR(T)-\frac78$, this is $c\,\NR(T)+o(T\log T)$.  With $N=\NR(T)+O(\log T)$ this gives $C^{\pathc}_\zeta(N;u)\to c$.

\emph{Step 3.}  For $\Re s>0$, $\int_0^\infty e^{-sx}dB_u(x)=\mathscr Z_\gamma(u+is)=:h(s)$.  Under RH $h(s)=\mathscr Z_\xi(u+is)=e^{-(u+is)/2}V_+(u+is)$, which by Proposition~\CGref{prop:cramer-zero-free} is holomorphic in $|s|<r_u=\min(u,\min_n|u-\log n|)$.  Since $|B_u(x)|\ll\log(2+x)$ and $|A_u(x)|\ll x\log(2+x)$, two integrations by parts give $\int_0^\infty e^{-sx}A_u(x)dx=h(s)/s^2$ for $\Re s>0$.  For real $s>0$, Fubini's theorem (the double integral of $|A_u(x)|e^{-\sigma x}$ over $x>0$, $\sigma>s$ is $\int|A_u|e^{-sx}x^{-1}dx<\infty$) gives $\int_0^\infty e^{-sx}\mathcal M_u(x)dx=\int_0^\infty A_u(x)\int_s^\infty e^{-\sigma x}d\sigma\,dx=\int_s^\infty h(\sigma)\sigma^{-2}d\sigma$.  Writing $h(\sigma)=h(0)+h'(0)\sigma+(h(\sigma)-h(0)-h'(0)\sigma)$ on $(s,r_u/2)$ gives, for real $0<s<r_u/2$, $\int_s^\infty h(\sigma)\sigma^{-2}d\sigma=\frac{h(0)}s-h'(0)\Log s+\Phi(s)$, where $\Phi'(s)=-(h(s)-h(0)-h'(0)s)/s^2$ is holomorphic in $|s|<r_u$.  Both sides of this identity are holomorphic in $\{\Re s>0\}\cap\{|s|<r_u\}$ (the left side as a Laplace transform), so it holds there by analytic continuation, which is the form \CGeqref{eq:tauberian-hyp}.  Moreover $|\mathcal M_u|\ll\log x$ and $\mathcal M_u'=(B_u-\mathcal M_u)/x\ll\log x/x$.  Lemma~\CGref{lem:local-tauberian} gives $\mathcal M_u(x)\to h(0)=e^{-u/2}V_+(u)$, and Step~2 gives \CGeqref{eq:centroid-RH}.
\end{proof}

\begin{remark}[Status of the centroid value]\label{rem:centroid-status}
Without RH neither the existence of $\lim_NC^{\pathc}_\zeta(N;1)$ nor its value is proved: Step~3 uses RH to identify $\mathscr Z_\gamma$ with the zero-free function $\mathscr Z_\xi$, and Step~1 uses $B_u(T)=O(\log T)$.  The equality \CGeqref{eq:centroid-RH} at $u=1$ is implied by RH; the converse does not follow from the argument of Proposition~\CGref{prop:imaginary-frequency} (Remark~\CGref{rem:window}(b)).  The real part of the value is Cram\'er's and Guillera's \cite{Cramer1919,Guillera2013}; we have not found the imaginary part $0.76785566886\ldots$, or the complex limit of the triangular centroid, in the literature consulted.  At $u=\log p^m$ the centroid diverges.  Under RH, Landau's formula gives $B_u(T)=-\frac T{2\pi}\frac{\log p}{p^{m/2}}+O(\log T)$; with $\gamma_n\sim2\pi n/\log n$ and $|S_n-B_u(\gamma_n)|\le m(\gamma_n)$ this is $S_n=-\frac{\log p}{p^{m/2}}\frac n{\log n}(1+o(1))$, and averaging, $C^{\pathc}_\zeta(N;u)=-\frac{\log p}{2p^{m/2}}\frac N{\log N}(1+o(1))=-\frac{\log p}{4\pi p^{m/2}}\,\gamma_N(1+o(1))$.  For the clock there is no resonance (Theorem~\CGref{thm:clock-centroid} holds for every $u>0$).
\end{remark}

\section{The RVM explicit formula, \texorpdfstring{$J_{\RVM}$}{J\_RVM}, and \texorpdfstring{$\pi_{\RVM}$}{pi\_RVM}}
\label{sec:explicit-formula}

For $w=a+ib$ with $b>0$ define
\begin{equation}\label{eq:Ei-def}
 \Ei(w):=\int_{-\infty}^{a}\frac{e^{v+ib}}{v+ib}\,dv,
 \qquad
 \Ei(\bar w):=\overline{\Ei(w)} .
\end{equation}
The path is the horizontal line $\Im=b$, which avoids $0$; this is the branch of $\Ei(\rho\log x)$ used in Riemann's explicit formula \cite{Edwards1974}.  For $x>1$ put
\begin{equation}\label{eq:J-RVM}
 J_{\RVM}(x):=\li(x)-2\Re\sum_{k\ge1}\Ei\!\left(\Bigl(\frac12+i\tau_k\Bigr)\log x\right)-\log2+\int_x^\infty\frac{dy}{y(y^2-1)\log y},
\end{equation}
the zero sum being taken in the order of $k$.

\begin{lemma}[Asymptotics of $\Ei$ off the real axis]\label{lem:Ei}
For $w=a+ib$, $b>0$,
\begin{equation}\label{eq:Ei-remainder}
 \Ei(w)=\frac{e^w}{w}+\varrho(w),\qquad|\varrho(w)|\le\frac{e^a}{b^2}.
\end{equation}
\end{lemma}

\begin{proof}
Integrate \CGeqref{eq:Ei-def} by parts: $\int_{-\infty}^ae^{v+ib}(v+ib)^{-1}dv=e^w/w+\int_{-\infty}^ae^{v+ib}(v+ib)^{-2}dv$, the boundary term at $-\infty$ being $0$.  The last integral is bounded by $\int_{-\infty}^ae^v b^{-2}dv=e^a/b^2$.
\end{proof}

\begin{theorem}[Compact-uniform existence of the RVM zero sum]\label{thm:J-exists}
The series $\sum_k\Re\Ei\bigl((\tfrac12+i\tau_k)u\bigr)$ converges uniformly for $u$ in every compact interval $[u_1,u_2]\subset(0,\infty)$.  Hence $J_{\RVM}$ is a continuous function on $(1,\infty)$.
\end{theorem}

\begin{proof}
Apply Lemma~\CGref{lem:Ei} with $w=(\tfrac12+i\tau_k)u$, $a=u/2$, $b=\tau_ku$.  Since $\Re\bigl(e^{iu\tau_k}/(\tfrac12+i\tau_k)\bigr)=(\tfrac12\cos u\tau_k+\tau_k\sin u\tau_k)/(\tfrac14+\tau_k^2)$,
\begin{equation}\label{eq:Ei-asymptotic}
 \Re\Ei\bigl((\tfrac12+i\tau_k)u\bigr)=\frac{e^{u/2}}u\,\frac{\sin(u\tau_k)}{\tau_k}+O_{u_1,u_2}(\tau_k^{-2}),
\end{equation}
uniformly for $u\in[u_1,u_2]$.  The error series converges absolutely and uniformly.  For the main series, Abel summation with $S_{\RVM}(K;u)=O_{u_1,u_2}(1+\log K)$ (Theorem~\CGref{thm:exp-sum}) and
\begin{equation}\label{eq:abel-weights}
 \frac1{\tau_K}-\frac1{\tau_{K+1}}\le\frac{\tau_{K+1}-\tau_K}{\tau_K^2}\le\frac{2\pi}{\ell_K\tau_K^2}\ll\frac{\log K}{K^2}
\end{equation}
reduces the uniform Cauchy estimate to the convergent majorant $\sum_K(\log K)^2K^{-2}$ (a majorant for the differences produced by Abel summation, not for the original terms), while the boundary term $S_{\RVM}(K;u)/\tau_K\to0$ uniformly.  Hence $\sum_ke^{iu\tau_k}/\tau_k$ converges uniformly, and so does its imaginary part.
\end{proof}

Define the support-truncated version
\begin{equation}\label{eq:J-star}
 J_{\RVM}^{*}(x):=\begin{cases}J_{\RVM}(x),&x\ge2,\\0,&0<x<2,\end{cases}
\end{equation}
set $\pi_{\RVM}(x)=0$ for $0<x<2$, and for $x\ge2$
\begin{equation}\label{eq:pi-RVM}
 \pi_{\RVM}(x):=\sum_{1\le m\le\log_2x}\frac{\mu(m)}m\,J_{\RVM}^{*}(x^{1/m}).
\end{equation}

\begin{proposition}[Exact finite prime-power inversion]\label{prop:finite-inversion}
For every $x\ge2$,
\begin{equation}\label{eq:J-pi-inverse}
 J_{\RVM}^{*}(x)=\sum_{1\le m\le\log_2x}\frac1m\,\pi_{\RVM}(x^{1/m}).
\end{equation}
\end{proposition}

\begin{proof}
All sums are finite because $x^{1/m}<2$ for $m>\log_2x$.  Substituting \CGeqref{eq:pi-RVM} into the right side of \CGeqref{eq:J-pi-inverse}, the coefficient of $J_{\RVM}^{*}(x^{1/r})$ is
\begin{equation}\label{eq:mobius-coefficient}
 \frac1r\sum_{d\mid r}\mu(d),
\end{equation}
which is $1$ for $r=1$ and $0$ otherwise.
\end{proof}

\subsection{The square-root estimate}
\label{subsec:square-root}

For $u>0$ put
\begin{equation}\label{eq:A-RVM}
 \mathcal A_{\RVM}(u):=\sum_{k\ge1}\frac{e^{iu\tau_k}}{\tau_k},
\end{equation}
summed in the order of $k$, and write $\ell(t)=\log(t/2\pi)$, so that $\NR'(t)=\ell(t)/2\pi$ and $\ell(\tau_k)=\ell_k$.

\begin{lemma}[The weighted clock sum]\label{lem:A-RVM}
For every $u\ge\log2$ the series \CGeqref{eq:A-RVM} converges, and
\begin{equation}\label{eq:A-RVM-bound}
 |\mathcal A_{\RVM}(u)|\le\frac{\log^2(2+2u)}{4\pi}+\frac{\log(2+2u)}{2\pi}+C_{\mathcal A},
\end{equation}
with an absolute constant $C_{\mathcal A}$.
\end{lemma}

\begin{proof}
Put $Y_j=2^j\cdot4\pi(1+u)$ for $j\ge0$.

\emph{Head $\tau_k\le Y_0$.}  Let $n_{\RVM}(t)=\#\{k:\tau_k\le t\}$.  By \CGeqref{eq:midpoint-condition}, $\tau_k\le t$ if and only if $k\le\NR(t)+\tfrac12$, so $n_{\RVM}(t)\le\NR(t)+\tfrac12=\frac t{2\pi}(\ell(t)-1)+\frac{11}8\le\frac{t\,\ell(t)}{2\pi}$ for $t\ge\tau_1>\frac{11\pi}4$.  Stieltjes integration gives
\begin{equation}\label{eq:A-head}
 \sum_{\tau_k\le Y_0}\frac1{\tau_k}=\frac{n_{\RVM}(Y_0)}{Y_0}+\int_{\tau_1}^{Y_0}\frac{n_{\RVM}(t)}{t^2}\,dt\le\frac{\ell(Y_0)}{2\pi}+\int_{\tau_1}^{Y_0}\frac{\ell(t)}{2\pi t}\,dt\le\frac{\ell(Y_0)}{2\pi}+\frac{\ell(Y_0)^2}{4\pi},
\end{equation}
and $\ell(Y_0)=\log(2+2u)$.

\emph{Blocks.}  Let $\mathcal B_j=\{k:Y_j<\tau_k\le Y_{j+1}\}$, let $\mathcal T$ be the inverse clock of \CGeqref{eq:inverse-clock}, and put $\phi(y)=u\,\mathcal T(y-\tfrac12)/2\pi$, so that $e^{2\pi i\phi(k)}=e^{iu\tau_k}$.  By \CGeqref{eq:inverse-clock},
\begin{equation}\label{eq:phi-second}
 \phi''(y)=-\frac{2\pi u}{\mathcal T\,\ell(\mathcal T)^3},\qquad \mathcal T=\mathcal T(y-\tfrac12).
\end{equation}
On $\mathcal B_j$ put $\varkappa_j=2\pi u/(Y_{j+1}\ell(Y_{j+1})^3)$.  Since $Y_{j+1}=2Y_j$ and $\ell(Y_{j+1})\le2\ell(Y_j)$ (because $\ell(Y_j)\ge\ell(Y_0)\ge\log2$), we have $\varkappa_j\le-\phi''\le16\varkappa_j$ at the integers of the block, and $\varkappa_j\le-\phi''\le256\varkappa_j$ on the whole real interval $[a,b]$ over which a run of $\mathcal B_j$ is summed, since that interval may begin inside the previous block, across which $-\phi''$ changes by at most another factor $16$; in particular $\phi''$ keeps a constant sign on $[a,b]$, as the second-derivative test requires.  The number of indices in $\mathcal B_j$ is at most $\NR(Y_{j+1})-\NR(Y_j)+1\le Y_j\ell(Y_{j+1})/2\pi+1$.  The second-derivative test for exponential sums \cite[Thm.~5.9]{Titchmarsh1986} gives, for every run $a<k\le b$ of consecutive indices of $\mathcal B_j$,
\begin{equation}\label{eq:block-vdc}
 \Bigl|\sum_{a<k\le b}e^{iu\tau_k}\Bigr|\le A_{\rm vdC}\Bigl(256\Bigl(\frac{Y_j\ell(Y_{j+1})}{2\pi}+1\Bigr)\varkappa_j^{1/2}+\varkappa_j^{-1/2}\Bigr)=:\mathcal P_j,
\end{equation}
with an absolute constant $A_{\rm vdC}$; if the real interval $[a,b]$ contains at most one integer, the sum has at most one term and is bounded trivially by $1$, and $1\le\varkappa_j^{-1/2}$ because $\varkappa_j\le u/(4(1+u)\log^32)<1$.  The weights $1/\tau_k$ are positive and decreasing, and $1/\tau_k<1/Y_j$ on $\mathcal B_j$, so Abel's inequality bounds the contribution of any run in $\mathcal B_j$ by $\mathcal P_j/Y_j$:
\begin{equation}\label{eq:block-weighted}
 \frac{\mathcal P_j}{Y_j}\le A_{\rm vdC}\left(\frac{16}{2\pi}\Bigl(\frac{\pi u}{Y_j\ell(Y_{j+1})}\Bigr)^{1/2}+\frac{16}{Y_j}+\Bigl(\frac{\ell(Y_{j+1})^3}{\pi uY_j}\Bigr)^{1/2}\right).
\end{equation}
With $Y_j=2^j\cdot4\pi(1+u)$ and $\ell(Y_{j+1})=\log(2^{j+2}(1+u))\le j+2+\log(1+u)$, each of the three terms is at most an absolute constant times $(j+2)^{3/2}2^{-j/2}$, uniformly in $u\ge\log2$, because $\sup_{u\ge\log2}(1+\log(1+u))^3/(u(1+u))<\infty$.  Hence $\sum_{j\ge0}\mathcal P_j/Y_j\le C_{\mathcal A}$ with $C_{\mathcal A}$ absolute.

\emph{Convergence and bound.}  Any finite range $K'<k\le K''$ of indices, with $\tau_{K'+1}\in\mathcal B_J$, consists of at most two partial runs (the initial segment of $\mathcal B_J$ and the terminal segment of the block containing $\tau_{K''}$) and of complete blocks in between; each run, partial or complete, contributes at most $\mathcal P_j/Y_j$ by \CGeqref{eq:block-vdc} and Abel's inequality, so $|\sum_{K'<k\le K''}e^{iu\tau_k}/\tau_k|\le\sum_{j\ge J}\mathcal P_j/Y_j\to0$ as $K'\to\infty$ (then $J\to\infty$), uniformly in $K''$.  The series converges, and \CGeqref{eq:A-head} together with $\sum_j\mathcal P_j/Y_j\le C_{\mathcal A}$ gives \CGeqref{eq:A-RVM-bound}.
\end{proof}

\begin{theorem}[Square-root estimate for $J_{\RVM}$ and $\pi_{\RVM}$]\label{thm:Lambert-estimate}
For $x\ge8$,
\begin{equation}\label{eq:J-baseline}
 J_{\RVM}(x)-\li(x)=O\!\left(\frac{\sqrt x(\log\log x)^2}{\log x}\right),
\end{equation}
\begin{equation}\label{eq:pi-baseline}
 \pi_{\RVM}(x)-\li(x)=O\!\left(\frac{\sqrt x(\log\log x)^2}{\log x}\right),
\end{equation}
with absolute implied constants.  No hypothesis on $\zeta$ is used.
\end{theorem}

\begin{proof}
Let $u=\log x\ge\log8$ and $\rho_k^{\RVM}=\tfrac12+i\tau_k$.  Lemma~\CGref{lem:Ei} with $w=\rho_k^{\RVM}u$, $a=u/2$, $b=\tau_ku$ gives
\begin{equation}\label{eq:Ei-split}
 \Ei(\rho_k^{\RVM}u)=\frac{e^{\rho_k^{\RVM}u}}{\rho_k^{\RVM}u}+\varrho_k,\qquad|\varrho_k|\le\frac{e^{u/2}}{\tau_k^2u^2}.
\end{equation}
Moreover $\frac1{\rho_k^{\RVM}}=-\frac i{\tau_k}+r_k$ with $r_k=\frac{i/2}{\tau_k\rho_k^{\RVM}}$, so $|r_k|\le\frac1{2\tau_k^2}$.  Since $\sum_k\tau_k^{-2}<\infty$, Lemma~\CGref{lem:A-RVM} shows that the series in \CGeqref{eq:J-RVM} converges in the order of $k$ and
\begin{equation}\label{eq:Ei-sum-A}
 \sum_{k\ge1}\Ei(\rho_k^{\RVM}u)=\frac{e^{u/2}}u\Bigl(-i\,\mathcal A_{\RVM}(u)+\sum_kr_ke^{iu\tau_k}\Bigr)+\sum_k\varrho_k
 =\frac{e^{u/2}}u\bigl(-i\,\mathcal A_{\RVM}(u)+O(1)\bigr).
\end{equation}
For $u\ge2$ one has $\log(2+2u)\le\log(2u)+\log2\le3\log u$, so \CGeqref{eq:A-RVM-bound} gives $\mathcal A_{\RVM}(u)=O(\log^2u)$.  The remaining terms $-\log2+\int_x^\infty dy/(y(y^2-1)\log y)$ of \CGeqref{eq:J-RVM} are bounded for $x\ge2$.  With $e^{u/2}=\sqrt x$ and $\log u=\log\log x$ this proves \CGeqref{eq:J-baseline}.

For \CGeqref{eq:pi-baseline}, \CGeqref{eq:pi-RVM} gives
\begin{equation}\label{eq:pi-minus-li}
 \pi_{\RVM}(x)-\li(x)=\bigl(J_{\RVM}(x)-\li(x)\bigr)+\sum_{2\le m\le\log_2x}\frac{\mu(m)}m\,J_{\RVM}^{*}(x^{1/m}).
\end{equation}
By \CGeqref{eq:J-baseline}, $\li(y)\ll y/\log y$, and the continuity of $J_{\RVM}$ on $[2,8]$ (Theorem~\CGref{thm:J-exists}), $|J^*_{\RVM}(y)|\le C\,y/\log y$ for all $y\ge2$.  The term $m=2$ is therefore $O(\sqrt x/\log x)$.  For $3\le m\le\log_2x$, $x^{1/m}\ge2$ and $\frac1mJ^*_{\RVM}(x^{1/m})\le\frac Cm\cdot\frac{x^{1/m}}{\log(x^{1/m})}=\frac{C\,x^{1/m}}{\log x}\le\frac{C\,x^{1/3}}{\log x}$, so these terms add up to at most $C\,x^{1/3}\log_2x/\log x=O(x^{1/3})$.  Both are absorbed in \CGeqref{eq:pi-baseline}.
\end{proof}

\begin{corollary}[The three limits]\label{cor:three-limits}
As $x\to\infty$,
\begin{equation}\label{eq:three-limits}
 \frac{\li(x)}{\pi_{\RVM}(x)}\to1,\qquad
 \frac{\pi(x)}{\li(x)}\to1,\qquad
 \frac{\pi(x)}{\pi_{\RVM}(x)}\to1 .
\end{equation}
More precisely, for some absolute $c_0>0$,
\begin{equation}\label{eq:pi-minus-piRVM}
 \pi(x)-\pi_{\RVM}(x)=O\!\left(x\,e^{-c_0\sqrt{\log x}}\right).
\end{equation}
\end{corollary}

\begin{proof}
Since $\li(x)\sim x/\log x$, \CGeqref{eq:pi-baseline} gives $\pi_{\RVM}(x)=\li(x)\bigl(1+O(x^{-1/2}(\log\log x)^2)\bigr)$; in particular $\pi_{\RVM}(x)>0$ for large $x$, and the first limit follows.  The second is the prime number theorem \cite{Hadamard1896,delaValleePoussin1896}, in the form $\pi(x)=\li(x)+O(x\,e^{-c_0\sqrt{\log x}})$ \cite[Ch.~18]{Davenport1980}.  Subtracting \CGeqref{eq:pi-baseline} gives \CGeqref{eq:pi-minus-piRVM}, and dividing by $\pi_{\RVM}(x)\sim x/\log x$ gives the third limit.
\end{proof}

The first limit and \CGeqref{eq:pi-baseline} concern only the synthetic system; the second and third compare it with the primes and use the prime number theorem.  Figures~\CGref{fig:prime-counts} and~\CGref{fig:pi-ratios} show them numerically.

\section{Integer weights and the continuous-to-atomic problem}
\label{sec:integer-weights}

Define
\begin{equation}\label{eq:cn}
 c_n^{\RVM}:=\pi_{\RVM}(n)-\pi_{\RVM}(n-1),\qquad n\ge2,
\end{equation}
and, initially for $\Re s>1$,
\begin{equation}\label{eq:weighted-zeta}
 \zeta_{\Z,\RVM}(s):=\prod_{n\ge2}(1-n^{-s})^{-c_n^{\RVM}}.
\end{equation}
No positivity or $0/1$ property is built into \CGeqref{eq:cn}.  Define the integer-step prime-power count and the right-continuous extension
\begin{equation}\label{eq:JZ}
 J_{\Z,\RVM}(x):=\sum_{1\le m\le\log_2x}\frac1m\,\pi_{\RVM}\bigl(\lfloor x^{1/m}\rfloor\bigr),
\end{equation}
\begin{equation}\label{eq:J-star-Z}
 J_{\RVM}^{*,\Z}(x):=J_{\RVM}^{*}(\lfloor x\rfloor).
\end{equation}

\begin{theorem}[Lipschitz control of $J_{\RVM}$]\label{thm:unit-cell}
There is an absolute constant $C_J$ such that
\begin{equation}\label{eq:unit-cell}
 |J_{\RVM}(x+\theta)-J_{\RVM}(x)|\le C_J\,\theta\qquad(x\ge2,\ 0<\theta\le1).
\end{equation}
\end{theorem}

\begin{proof}
Put $u_1=\log x$, $u_2=\log(x+\theta)$, so $0<u_2-u_1=\log(1+\theta/x)\le\theta/x$ and $u_2-u_1\le\log2$.  The terms $\li$ and $\int_x^\infty dy/(y(y^2-1)\log y)$ of \CGeqref{eq:J-RVM} have derivatives bounded by $1/\log2$ and $3/\log2$ on $[2,\infty)$, so they change by $O(\theta)$.  It remains to bound $\sum_kD_k$, where, by \CGeqref{eq:Ei-def} and $\frac{d}{dv}\Ei(\rho v)=e^{\rho v}/v$,
\begin{equation}\label{eq:Dk-def}
 D_k:=\Ei(\rho_k^{\RVM}u_2)-\Ei(\rho_k^{\RVM}u_1)=\int_{u_1}^{u_2}\frac{e^{v/2}e^{iv\tau_k}}{v}\,dv .
\end{equation}
The series $\sum_kD_k$ converges (Theorem~\CGref{thm:J-exists} and its proof).  By Abel's theorem for Dirichlet series with the increasing exponents $\tau_k$ (not the index $k$), $\sum_kD_k=\lim_{\alpha\downarrow0}\sum_ke^{-\alpha\tau_k}D_k$: with $R_K:=\sum_{k>K}D_k$, summation by parts gives $\sum_{k>K}e^{-\alpha\tau_k}D_k=e^{-\alpha\tau_{K+1}}R_K-\sum_{k>K}(e^{-\alpha\tau_k}-e^{-\alpha\tau_{k+1}})R_k$, of modulus at most $2\sup_{k\ge K}|R_k|$ uniformly in $\alpha>0$ because $e^{-\alpha\tau_k}$ decreases in $k$, while the finitely many terms $k\le K$ converge as $\alpha\downarrow0$.  Moreover, for $\alpha>0$ Fubini gives
\begin{equation}\label{eq:damped-sum}
 \sum_ke^{-\alpha\tau_k}D_k=\int_{u_1}^{u_2}\frac{e^{v/2}}v\,S_\alpha(v)\,dv,\qquad S_\alpha(v):=\sum_{k\ge1}g(\tau_k),\quad g(\tau):=e^{(iv-\alpha)\tau}.
\end{equation}
\emph{Clock decomposition.}  With $\tau_0=4\pi$ and the sawtooth count \CGeqref{eq:sawtooth-count} (proved later in Section~\CGref{sec:transition}), Stieltjes integration exactly as in the forward reference \CGeqref{eq:transition-split} ($g(\infty)=0$, $g'\in L^1$) gives
\begin{equation}\label{eq:S-eps-split}
 S_\alpha(v)=\int_{\tau_0}^\infty g\,d\NR+g(\tau_0)\NR(\tau_0)+\int_{\tau_0}^\infty\psi\bigl(\NR(\tau)+\tfrac12\bigr)g'(\tau)\,d\tau .
\end{equation}
The middle term is at most $\tfrac12$.  For the first, two integrations by parts with $\NR'(\tau)=\frac1{2\pi}\log\frac\tau{2\pi}$, $\NR''(\tau)=\frac1{2\pi\tau}$ and $v\ge\log2$ give $|\int_{\tau_0}^\infty g\,d\NR|\le C/v$, uniformly in $\alpha\in(0,1]$.

\emph{Harmonics.}  The partial sums of $\psi(y+\frac12)=-\sum_{m\ge1}(-1)^m\sin(2\pi my)/(\pi m)$ are uniformly bounded, $\sup_{M\ge1}\sup_{y\in\mathbb R}\bigl|\sum_{m\le M}\sin(2\pi my)/m\bigr|<\infty$ (for $0<y\le\frac12$ split at $M_0:=\lfloor1/y\rfloor$: the terms $m\le M_0$ contribute at most $\sum_{m\le M_0}2\pi y\le2\pi$, and Abel summation with $|\sum_{M_0<m\le M'}\sin(2\pi my)|\le1/\sin(\pi y)\le1/(2y)$ and $1/m\le1/M_0\le2y$ bounds the rest by $2$; the sum is odd and $1$-periodic in $y$), and the series converges to $\psi(y+\frac12)$ for every $y\notin\frac12+\mathbb Z$; the exceptional set $\{\tau:\NR(\tau)+\frac12\in\mathbb Z\}$ is discrete, hence Lebesgue-null.  Since $g'\in L^1(\tau_0,\infty)$ for $\alpha>0$, dominated convergence allows termwise integration:
\begin{equation}\label{eq:K-harmonics}
 \int_{\tau_0}^\infty\psi\bigl(\NR+\tfrac12\bigr)g'\,d\tau=-\sum_{m\ge1}\frac{(-1)^m}{2\pi i\,m}\bigl(K_m^+-K_m^-\bigr),\qquad
 K_m^\pm:=\int_{\tau_0}^\infty(iv-\alpha)e^{-\alpha\tau}e^{i\Psi_m^\pm(\tau)}\,d\tau,
\end{equation}
with $\Psi_m^\pm(\tau)=v\tau\pm2\pi m\NR(\tau)$.  The weight $e^{-\alpha\tau}$ is positive and decreasing and $|iv-\alpha|\le v+1$ for $0<\alpha\le1$; the second mean-value theorem (Bonnet's form), applied to the real and to the imaginary part of $e^{i\Psi_m^\pm}$ separately, bounds $|\int_a^be^{-\alpha\tau}e^{i\Psi_m^\pm}d\tau|$ by $2e^{-\alpha a}\sup_{\xi\in[a,b]}|\int_a^\xi e^{i\Psi_m^\pm}d\tau|$, so each integral is at most $2(v+1)\sup_{\xi}|\int_a^\xi e^{i\Psi_m^\pm}d\tau|$ over the relevant interval $[a,b]$.  We use the first- and second-derivative tests \cite[Lemmas~4.2 and~4.4]{Titchmarsh1986} for a real phase $\varTheta$: $|\int_a^b e^{i\varTheta}|\le4/\mu$ if $\varTheta'$ is monotone with $|\varTheta'|\ge\mu$, and $\le8/\sqrt\lambda$ if $|\varTheta''|\ge\lambda$.

For $K_m^+$: $(\Psi_m^+)'=v+m\log\frac\tau{2\pi}\ge v+m\log2$ is increasing, so $|K_m^+|\le8(v+1)/(v+m\log2)$.

For $K_m^-$: $(\Psi_m^-)'=v-m\log\frac\tau{2\pi}$ is decreasing and vanishes at $T_m:=2\pi e^{v/m}>2\pi$, and $(\Psi_m^-)''=-m/\tau$.  On $[\tau_0,T_m/2]$ (if nonempty) $(\Psi_m^-)'\ge m\log2$; on $[2T_m,\infty)$, $(\Psi_m^-)'\le-m\log2$; on $[\max(\tau_0,T_m/2),2T_m]$, $|(\Psi_m^-)''|\ge m/(2T_m)$.  Hence
\begin{equation}\label{eq:K-minus-bound}
 |K_m^-|\le2(v+1)\Bigl(\frac{8}{m\log2}+8\sqrt{\frac{2T_m}m}\Bigr)
 =2(v+1)\Bigl(\frac{8}{m\log2}+16\sqrt\pi\,\frac{e^{v/(2m)}}{\sqrt m}\Bigr).
\end{equation}
Summing with the weights $1/(2\pi m)$ and using $\sum_mm^{-3/2}e^{v/(2m)}\le e^{v/2}+\zeta(\tfrac32)e^{v/4}$,
\begin{equation}\label{eq:S-eps-bound}
 |S_\alpha(v)|\le C\,(v+1)\,e^{v/2}\qquad(v\ge\log2,\ 0<\alpha\le1),
\end{equation}
with an absolute constant $C$.

\emph{Conclusion.}  Inserting \CGeqref{eq:S-eps-bound} into \CGeqref{eq:damped-sum}, for $v\in[u_1,u_2]$ we have $e^{v/2}\le\sqrt{2x}$ and $(v+1)/v\le1+1/\log2$, so
\begin{equation}\label{eq:Dk-sum-bound}
 \Bigl|\sum_ke^{-\alpha\tau_k}D_k\Bigr|\le C'(u_2-u_1)\,x\le C'\theta,
\end{equation}
uniformly in $\alpha$.  Letting $\alpha\downarrow0$ and taking $2\Re$ proves \CGeqref{eq:unit-cell}.
\end{proof}

The $m=1$ harmonic, with stationary point $T_1=2\pi x$, produces the unit-periodic ripple of $J_{\RVM}$ (Figure~\CGref{fig:prime-counts}); the harmonics $m\ge2$ have stationary points $2\pi x^{1/m}$ and contribute $O(x^{-1/2+1/(2m)}m^{-3/2})$ to a unit increment.

\begin{corollary}[Bounded cell masses]\label{cor:cn-bounded}
\begin{equation}\label{eq:unit-cell-consequence}
 \pi_{\RVM}(y)-\pi_{\RVM}(\lfloor y\rfloor)=O(1)\quad(y\ge2),\qquad c_n^{\RVM}=O(1)\quad(n\ge2),
\end{equation}
with absolute implied constants, uniform in $y\ge2$ and $n\ge2$.
\end{corollary}

\begin{proof}
By \CGeqref{eq:pi-RVM}, $c_n^{\RVM}=\sum_{m\le\log_2n}\frac{\mu(m)}m\bigl(J_{\RVM}^{*}(n^{1/m})-J_{\RVM}^{*}((n-1)^{1/m})\bigr)$.  The term $m=1$ is $O(1)$ by \CGeqref{eq:unit-cell}.  For $m\ge2$ and $(n-1)^{1/m}\ge2$, \CGeqref{eq:unit-cell} with $\theta=n^{1/m}-(n-1)^{1/m}\le n^{1/m-1}$ bounds the term by $C_Jn^{1/m-1}\le C_Jn^{-1/2}$, and there are at most $\log_2n$ such terms.  The jump of $J_{\RVM}^{*}$ at $2$ enters only for $n=2^m$, once, with size $|J_{\RVM}(2)|/m$.  The same argument applies to $\pi_{\RVM}(y)-\pi_{\RVM}(\lfloor y\rfloor)$.
\end{proof}

\begin{corollary}[The integer-weighted zeta is a Dirichlet series]\label{cor:ZRVM-Dirichlet}
The product \CGeqref{eq:weighted-zeta} converges absolutely for $\Re s>1$ and
\begin{equation}\label{eq:ZRVM-Dirichlet}
 \zeta_{\Z,\RVM}(s)=\sum_{N\ge1}\frac{a_N^{\Z}}{N^s},\qquad a_1^{\Z}=1,\qquad \sum_{N\ge1}\frac{|a_N^{\Z}|}{N^\sigma}<\infty\quad(\sigma>1).
\end{equation}
The inverse is also an absolutely convergent Dirichlet series on the whole half-plane:
\begin{equation}\label{eq:ZRVM-inverse-Dirichlet}
 \frac1{\zeta_{\Z,\RVM}(s)}=\sum_{N\ge1}\frac{\bar a_N^{\Z}}{N^s},\qquad \bar a_1^{\Z}=1,\qquad \sum_{N\ge1}\frac{|\bar a_N^{\Z}|}{N^\sigma}<\infty\quad(\sigma>1).
\end{equation}
\end{corollary}

\begin{proof}
By Corollary~\CGref{cor:cn-bounded},
\begin{equation}\label{eq:log-ZRVM}
 \log\zeta_{\Z,\RVM}(s)=\sum_{n\ge2}c_n^{\RVM}\sum_{m\ge1}\frac{n^{-ms}}m=\sum_{N\ge2}\frac{b_N}{N^s},\qquad b_N:=\sum_{n^m=N}\frac{c_n^{\RVM}}m,
\end{equation}
with $\sum_N|b_N|N^{-\sigma}\le C\sum_{n\ge2}\sum_m n^{-m\sigma}/m<\infty$ for $\sigma>1$.  The exponential of an absolutely convergent Dirichlet series without constant term is an absolutely convergent Dirichlet series with $a_1^{\Z}=1$, and $\sum_N|a_N^{\Z}|N^{-\sigma}\le\exp\bigl(\sum_N|b_N|N^{-\sigma}\bigr)$.  The same applies to $1/\zeta_{\Z,\RVM}=\exp\bigl(-\sum_Nb_NN^{-s}\bigr)$, which gives \CGeqref{eq:ZRVM-inverse-Dirichlet}.
\end{proof}

Numerically the cell masses lie in $[-0.38,1.04]$ for $n\le1000$ (Experiment~\CGref{exp:primes}).

\begin{theorem}[Mellin--Stieltjes identification]\label{thm:MS}
There is a function $J_{\RVM}^{\rm M}$ on $(1,\infty)$ with $J_{\RVM}^{\rm M}(x)x^{-c}\in L^2((1,\infty),dx/x)$ for every $c>1$ such that
\begin{equation}\label{eq:MS-RVM}
 \log\zeta_{\RVM}(s)=s\int_1^\infty J_{\RVM}^{\rm M}(x)\,x^{-s-1}\,dx\qquad(\Re s>1),
\end{equation}
and $J_{\RVM}^{\rm M}-J_{\RVM}$ is constant on $(1,\infty)$.  In particular $J_{\RVM}^{\rm M}=J_{\RVM}^{*}+C_{\RVM}$ on $[2,\infty)$ for a constant $C_{\RVM}$.
\end{theorem}

The proof is given in Section~\CGref{sec:euler-products} (Theorems~\CGref{thm:mellin-counts} and~\CGref{thm:distributional-EF}); it uses the half-plane bound of Theorem~\CGref{thm:halfplane-RVM}.

\begin{corollary}[Weighted factorization]\label{cor:weighted-factorization}
\begin{equation}\label{eq:JZ-diff}
 J_{\Z,\RVM}(x)-J_{\RVM}^{*}(x)=O(1+\log\log x),
\end{equation}
and
\begin{equation}\label{eq:weighted-factorization}
 \zeta_{\Z,\RVM}(s)=\zeta_{\RVM}(s)\,e^{h_{\RVM}(s)},
\end{equation}
where
\begin{equation}\label{eq:hRVM}
 h_{\RVM}(s):=s\int_1^\infty\bigl(J_{\Z,\RVM}(x)-J_{\RVM}^{\rm M}(x)\bigr)x^{-s-1}\,dx
\end{equation}
is holomorphic for $\Re s>0$.  The two functions have the same zeros in $\Re s>0$, and $\zeta_{\Z,\RVM}$ inherits through $e^{h_{\RVM}}$ the pole of $\zeta_{\RVM}$ at $s=1$; compare the modified Euler products of Grosswald and Schnitzer \cite{GrosswaldSchnitzer1978}.
\end{corollary}

\begin{proof}
By \CGeqref{eq:unit-cell-consequence} replacing $x^{1/m}$ by its integer part costs $O(1/m)$ in the $m$th summand; summing $m^{-1}$ up to $\log_2x$ gives \CGeqref{eq:JZ-diff}.  The product \CGeqref{eq:weighted-zeta} is initially defined for $\Re s>1$.  For non-integral exponents we use the principal logarithmic branch, equivalently $-\log(1-n^{-s})=\sum_{m\ge1}n^{-ms}/m$, which is unambiguous because $|n^{-s}|<1$ there.  It gives $\log\zeta_{\Z,\RVM}(s)=s\int_1^\infty J_{\Z,\RVM}(x)x^{-s-1}dx$ for $\Re s>1$, because $\sum_{n\le y}c_n^{\RVM}=\pi_{\RVM}(\lfloor y\rfloor)$.  Subtracting \CGeqref{eq:MS-RVM} gives \CGeqref{eq:weighted-factorization} on $\Re s>1$.  On $[2,\infty)$, $J_{\Z,\RVM}-J_{\RVM}^{\rm M}=(J_{\Z,\RVM}-J_{\RVM}^{*})-C_{\RVM}=O(1+\log\log x)$ by \CGeqref{eq:JZ-diff}; on $(1,2)$, $J_{\Z,\RVM}=0$ and $J_{\RVM}^{\rm M}$ is square-integrable.  Hence \CGeqref{eq:hRVM} converges locally normally for $\Re s>0$, and \CGeqref{eq:weighted-factorization} defines the holomorphic continuation of $\zeta_{\Z,\RVM}$ from $\Re s>1$ to $\Re s>0$; exponentiation introduces no zeros.
\end{proof}

\begin{lemma}[Exact first-harmonic cell cancellation]\label{lem:cell-cancel}
For every integer $n$,
\begin{equation}\label{eq:cell-cancel}
 \int_{n-1}^{n}2\cos(2\pi x)\,dx=0 .
\end{equation}
\end{lemma}

\begin{hypothesis}[Uniform cell remainder away from the exceptional powers of $2$; proved in Corollary~\CGref{cor:clock-cells}]\label{hyp:cell-remainder}
For $n\ge3$ with $n\notin\{2^j:\mu(j)\ne0\}$,
\begin{equation}\label{eq:cell-remainder}
 c_n^{\RVM}-\int_{n-1}^{n}\left(\frac1{\log x}+2\cos(2\pi x)\right)dx=o\!\left(\frac1{\log n}\right).
\end{equation}
\end{hypothesis}

Hypothesis~\CGref{hyp:cell-remainder} is proved below, in Corollary~\CGref{cor:clock-cells}, for exactly the set of $n$ in its statement; it is stated here as a hypothesis only because its proof needs the density theorem of Section~\CGref{sec:criteria}, and the proposition that follows is therefore unconditional.

\begin{proposition}[Eventual positivity]\label{prop:eventual-positive}
Hypothesis~\CGref{hyp:cell-remainder} implies
\begin{equation}\label{eq:eventual-positive}
 c_n^{\RVM}\sim\frac1{\log n},
\end{equation}
hence $0<c_n^{\RVM}<1$ for all sufficiently large $n$ in that set.
\end{proposition}

\begin{proof}
The main-density integral is $1/\log n+O(1/(n\log^2n))$; the first harmonic vanishes by Lemma~\CGref{lem:cell-cancel}.
\end{proof}

\section{The critical RVM Mertens transform}
\label{sec:mertens-transform}

Define
\begin{equation}\label{eq:F-RVM}
 F_{\RVM}(s)=\widehat M_{\RVM}^{\crit}(s):=\frac1{\zeta_{\RVM}(s)}=\frac{G(s)}{P_{\RVM}^{\crit}(s)} .
\end{equation}
Formally the associated summatory object is the inverse Mellin transform
\begin{equation}\label{eq:Mcrit-formal}
 M_{\RVM}^{\crit}(x)=\frac1{2\pi i}\int_{(c)}F_{\RVM}(s)\frac{x^s}{s}\,ds ;
\end{equation}
its existence as a symmetric limit, for every $c>\tfrac12$, is proved in Theorem~\CGref{thm:unsmoothed-RVM}.

\begin{proposition}[Synthetic pole abscissa]\label{prop:Mcrit-poles}
$F_{\RVM}=G/P_{\RVM}^{\crit}$ is meromorphic on $\C$ and holomorphic in $\Re s>\tfrac12$.  The poles of this meromorphic function are the simple poles
\begin{equation}\label{eq:Mcrit-poles}
 \frac12\pm i\tau_k\quad(k\ge1)\qquad\text{and}\qquad s=-2,-4,-6,\dots
\end{equation}
and $s=1$ is a simple zero of $F_{\RVM}$.
\end{proposition}

\begin{proof}
This is Theorem~\CGref{thm:canonical} applied to $F_{\RVM}=G/P_{\RVM}^{\crit}$: the zeros of $P_{\RVM}^{\crit}$ are \CGeqref{eq:synthetic-divisor}, and the poles of $G$ are those of $\Gamma(s/2)$ except $s=0$, where $s\Gamma(s/2)\to2$.
\end{proof}

\begin{proposition}[Gaussian square-root bound under vertical growth]\label{prop:Gaussian}
Fix $c>1$ and $0<\eps<c-\tfrac12$.  Assume, uniformly for $\tfrac12+\eps\le\sigma\le c$,
\begin{equation}\label{eq:subgaussian}
 F_{\RVM}(\sigma+it)=\exp(o(t^2))\qquad(|t|\to\infty).
\end{equation}
Define
\begin{equation}\label{eq:MG}
 M_{\RVM,G}^{\crit}(x):=\frac1{2\pi i}\int_{(c)}e^{s^2}F_{\RVM}(s)\frac{x^s}{s}\,ds .
\end{equation}
Then
\begin{equation}\label{eq:MG-bound}
 M_{\RVM,G}^{\crit}(x)=O_\eps(x^{1/2+\eps}).
\end{equation}
\end{proposition}

\begin{proof}
On $s=\sigma+it$, $|e^{s^2}|=e^{\sigma^2-t^2}$, so \CGeqref{eq:subgaussian} gives absolute convergence on vertical lines and makes the horizontal integrals tend to $0$.  The strip $\tfrac12+\eps\le\Re s\le c$ contains no pole of $F_{\RVM}$ (Proposition~\CGref{prop:Mcrit-poles}) and does not contain $s=0$, so Cauchy's theorem moves the line to $\Re s=\tfrac12+\eps$, where $|x^s|=x^{1/2+\eps}$.
\end{proof}

Theorem~\CGref{thm:unsmoothed-RVM} below removes both the assumption \CGeqref{eq:subgaussian} and the smoothing.

\section{The transition theorem and the unsmoothed synthetic Mertens bound}
\label{sec:transition}

Let
\begin{equation}\label{eq:chi}
 \chi(s):=2^s\pi^{s-1}\sin\!\Bigl(\frac{\pi s}2\Bigr)\Gamma(1-s)=\pi^{s-1/2}\frac{\Gamma((1-s)/2)}{\Gamma(s/2)},
\end{equation}
the factor of the functional equation $\zeta(s)=\chi(s)\zeta(1-s)$ \cite[\S2.1]{Titchmarsh1986}.  Write
\begin{equation}\label{eq:transition-zw}
 s=\tfrac12+z,\quad z=\sigma'+it,\qquad
 \tau_*:=-iz=t-i\sigma',\qquad
 w(s):=\exp\bigl\{-2\pi i\,\NR(\tau_*)\bigr\}.
\end{equation}
Stirling's formula gives, uniformly on compact $\sigma$-sets,
\begin{equation}\label{eq:chi-size}
 |\chi(\sigma+it)|=\Bigl(\frac{|t|}{2\pi}\Bigr)^{1/2-\sigma}\bigl(1+O(|t|^{-1})\bigr).
\end{equation}

\begin{theorem}[Transition evaluation of the midpoint product]\label{thm:transition}
There is a constant $\varpi_{\RVM}>0$ with the following properties.
\begin{enumerate}[label=\textup{(\alph*)}]
\item For every $B>0$ there are $t_0(B)$ and $C_B$ such that, for $0<\sigma'\le B$ and $t\ge t_0(B)$,
\begin{equation}\label{eq:transition-w}
 \zeta_{\RVM}(s)=\varpi_{\RVM}\bigl(1+w(s)\bigr)\bigl(1+\mathcal E_0(s)\bigr),\qquad|\mathcal E_0(s)|\le\frac{C_B}t .
\end{equation}
\item For every $0<\sigma'_1<\sigma'_2<\tfrac12$ there are $t_0$ and $C$ such that, for $\sigma'_1\le\sigma'\le\sigma'_2$ and $t\ge t_0$,
\begin{equation}\label{eq:transition-chi}
 \zeta_{\RVM}(s)=\varpi_{\RVM}\bigl(1+\chi(s)\bigr)\bigl(1+\mathcal E(s)\bigr),\qquad|\mathcal E(s)|\le\frac Ct .
\end{equation}
\end{enumerate}
For $t<0$ the same statements hold by $\zeta_{\RVM}(\bar s)=\overline{\zeta_{\RVM}(s)}$.
\end{theorem}

The phase identity $w=\chi(s)(1+O(1/t))$ alone does not give \CGeqref{eq:transition-w}; the product has to be evaluated across the transition region $\tau_k\asymp t$.  $\Log$ denotes the principal logarithm and $\Log\Gamma$ the branch of $\log\Gamma$ on $\Re v>0$ that is real on $(0,\infty)$.  Nothing about $\zeta$ is used.

\subsection{Proof of Theorem~\ref{thm:transition}}
\label{subsec:transition-proof}

Throughout, $t>0$, $0<\sigma'\le B$, $\tau_0:=4\pi$, $t^2\ge2(\tau_0^2+B^2)$, and implied constants depend only on $B$ (and on $\sigma'_1,\sigma'_2$ in part (b)).

\subsubsection*{Step 0: reduction to a sum over the clock}
Since $q(s)=z^2-\tfrac14$, each factor of \CGeqref{eq:Pcrit} equals $(\tau_k^2+z^2)/(\tau_k^2+\tfrac14)$, so
\begin{equation}\label{eq:transition-product}
 \frac{P_{\RVM}^{\crit}(s)}{P_{\RVM}^{\crit}(\tfrac12)}=\prod_{k\ge1}\Bigl(1+\frac{z^2}{\tau_k^2}\Bigr)=\exp\Lambda(z),\qquad
 \Lambda(z):=\sum_{k\ge1}f(\tau_k),\quad f(\tau):=\Log\Bigl(1+\frac{z^2}{\tau^2}\Bigr).
\end{equation}
For $\tau>0$, $\Im(1+z^2/\tau^2)=2\sigma't/\tau^2>0$, so $f$ is smooth on $(0,\infty)$, $f(\tau)=O(|z|^2\tau^{-2})$, the series converges absolutely, and
\begin{equation}\label{eq:fprime}
 f'(\tau)=-\frac{2z^2}{\tau(\tau^2+z^2)}=-\frac2\tau+\frac1{\tau-iz}+\frac1{\tau+iz}.
\end{equation}

\subsubsection*{Step 1: the midpoint sawtooth}
The clock increases on $[2\pi,\infty)$, $\NR(\tau_0)=2\log\frac2e+\frac78\in(0,\tfrac12)$, and $\tau_1>\tau_0$.  With $\psi(y):=y-\lfloor y\rfloor-\tfrac12$, the count \CGeqref{eq:n-RVM} satisfies, for $\tau\ge\tau_0$,
\begin{equation}\label{eq:sawtooth-count}
 n_{\RVM}(\tau)=\bigl\lfloor\NR(\tau)+\tfrac12\bigr\rfloor=\NR(\tau)-\psi\bigl(\NR(\tau)+\tfrac12\bigr),\qquad n_{\RVM}(\tau_0)=0 .
\end{equation}
Stieltjes integration of $f$ against \CGeqref{eq:sawtooth-count}, followed by integration by parts in the $\psi$-part ($f(\infty)=0$, $|\psi|\le\tfrac12$, $\psi(\NR(\tau_0)+\tfrac12)=\NR(\tau_0)$), gives
\begin{equation}\label{eq:transition-split}
\begin{split}
 &\Lambda(z)=\Lambda_{\rm sm}(z)+\Lambda_{\rm saw}(z),\qquad
 \Lambda_{\rm sm}(z):=\int_{\tau_0}^\infty f\,d\NR+f(\tau_0)\NR(\tau_0),\\
 &\Lambda_{\rm saw}(z):=\int_{\tau_0}^{\infty}\psi\bigl(\NR(\tau)+\tfrac12\bigr)f'(\tau)\,d\tau .
\end{split}
\end{equation}

\begin{lemma}[Smooth part]\label{lem:TR-smooth}
With $\log G(s):=\Log(z^2-\tfrac14)+\Log\Gamma(\tfrac s2)-\tfrac s2\log\pi$,
\begin{equation}\label{eq:smooth-difference}
 \Lambda_{\rm sm}(z)-\log G(s)=\nu_3+O(t^{-1}),\qquad\nu_3\in\R\ \text{independent of }z .
\end{equation}
\end{lemma}

\begin{proof}
A second integration by parts ($f(\tau)=O(|z|^2\tau^{-2})$ and $\NR(\tau)=O(\tau\log\tau)$, so $f\NR\to0$ at $\infty$) cancels the boundary term:
\begin{equation}\label{eq:smooth-J-ibp}
 \Lambda_{\rm sm}(z)=-\int_{\tau_0}^\infty\NR f'\,d\tau=\int_{\tau_0}^\infty\NR(\tau)\Bigl(\frac2\tau-\frac{2\tau}{\tau^2+z^2}\Bigr)d\tau .
\end{equation}
The constant $\tfrac78$ of the clock contributes $\tfrac78\Log(1+z^2/\tau_0^2)$, because $\Log\frac{\tau^2}{\tau^2+z^2}$ is an antiderivative of the bracket on $(0,\infty)$ tending to $0$ at $\infty$.  As $0<\arg z<\frac\pi2$, $\arg(1+\tau_0^2/z^2)=O(t^{-2})$ and
\begin{equation}\label{eq:J2}
 \tfrac78\Log\Bigl(1+\frac{z^2}{\tau_0^2}\Bigr)=\tfrac74\Log z-\tfrac74\log\tau_0+O(t^{-2}).
\end{equation}
The term $\frac\tau{2\pi}\log\frac\tau{2\pi e}$ contributes $\frac1\pi\int_{\tau_0}^\infty\log\frac{\tau}{2\pi e}\frac{z^2}{\tau^2+z^2}\,d\tau$.  Put
\begin{equation}\label{eq:I-def}
 \mathcal I(z):=\frac1\pi\int_0^\infty\log\frac{\tau}{2\pi e}\,\frac{z^2}{\tau^2+z^2}\,d\tau .
\end{equation}
The integrand is holomorphic in $z$ on $\Re z>0$ (its singularities are at $z=\pm i\tau$), $O(|z|^2\tau^{-2}\log\tau)$ at $\infty$ and $O(|\log\tau|)$ at $0$, locally uniformly; so $\mathcal I$ is holomorphic on $\Re z>0$.  For real $z=x>0$ the substitution $\tau=xv$ and $\int_0^\infty\frac{dv}{1+v^2}=\frac\pi2$, $\int_0^\infty\frac{\log v}{1+v^2}\,dv=0$ give $\mathcal I(x)=\frac x2\log\frac x{2\pi e}$; by the identity theorem $\mathcal I(z)=\frac z2\Log\frac z{2\pi e}$ on $\Re z>0$.  On $[0,\tau_0]$, $|\tau^2+z^2|\ge t^2/2$, so $\frac{z^2}{\tau^2+z^2}=1+O(t^{-2})$ there.  Hence
\begin{equation}\label{eq:J-asymp}
 \Lambda_{\rm sm}(z)=\frac z2\Log\frac z{2\pi e}+\frac74\Log z+\nu_1+O(t^{-2}),\qquad
 \nu_1=-\tfrac74\log\tau_0-\tfrac1\pi\int_0^{\tau_0}\log\tfrac\tau{2\pi e}\,d\tau .
\end{equation}
Stirling's formula $\Log\Gamma(v)=(v-\tfrac12)\Log v-v+\tfrac12\log2\pi+O(|v|^{-1})$ \cite[\S5.11]{DLMF2024} with $v=\frac s2=\frac{z+1/2}2$ and $\Log v=\Log z-\log2+\frac1{2z}+O(t^{-2})$ gives
\begin{equation}\label{eq:stirling-loggamma}
 \Log\Gamma(\tfrac s2)=\tfrac z2\Log\tfrac z2-\tfrac14\Log z+\tfrac14\log2-\tfrac z2+\tfrac12\log2\pi+O(t^{-1}),
\end{equation}
and with $\Log(z^2-\tfrac14)=2\Log z+O(t^{-2})$, $-\frac s2\log\pi=-\frac z2\log\pi-\frac14\log\pi$,
\begin{equation}\label{eq:G-asymp}
 \log G(s)=\frac z2\Log\frac z{2\pi e}+\frac74\Log z+\nu_2+O(t^{-1}),\qquad
 \nu_2=\tfrac14\log2+\tfrac12\log2\pi-\tfrac14\log\pi .
\end{equation}
Subtracting gives \CGeqref{eq:smooth-difference} with $\nu_3=\nu_1-\nu_2$.
\end{proof}

\begin{lemma}[The clock as a holomorphic phase]\label{lem:TR-clock}
$\NR(\tau)=\frac\tau{2\pi}\bigl(\Log\frac\tau{2\pi}-1\bigr)+\frac78$ is holomorphic on $\Re\tau>0$, $\NR(\bar\tau)=\overline{\NR(\tau)}$, and for $\Re\tau\ge\tau_0$,
\begin{equation}\label{eq:clock-halfplane-bound}
 2\pi\,\Im\NR(\tau)\ge(\log2)\,\Im\tau\quad(\Im\tau\ge0),\qquad
 2\pi\,\Im\NR(\tau)\le(\log2)\,\Im\tau\quad(\Im\tau\le0).
\end{equation}
\end{lemma}

\begin{proof}
For $\tau=X+iY$, $X\ge\tau_0$, $Y>0$, put $v=Y/X$.  Then
\begin{equation}\label{eq:clock-im-identity}
 2\pi\Im\NR(\tau)=Y\log\frac{|\tau|}{2\pi e}+X\arg\tau
 =Y\Bigl[\log\frac{X}{2\pi e}+\tfrac12\log(1+v^2)+\frac{\arctan v}v\Bigr]\ge Y\bigl[\log2+\varphi(v)\bigr],
\end{equation}
using $X\ge4\pi$, where $\varphi(v)=\frac12\log(1+v^2)+\frac{\arctan v}v-1$.  Since $\varphi(0^+)=0$ and $\varphi'(v)=(v-\arctan v)/v^2\ge0$, $\varphi\ge0$.  The case $Y<0$ follows by conjugation.
\end{proof}

In particular $|w(s)|=e^{2\pi\Im\NR(\tau_*)}\le2^{-\sigma'}<1$, and on the vertical rays $|e^{\pm2\pi im\NR(\tau_0\pm iy)}|\le2^{-my}$.

\begin{lemma}[Abel-regularized harmonic sum]\label{lem:TR-harmonic}
\begin{equation}\label{eq:Psi-asymp}
 \Lambda_{\rm saw}(z)=\Log\bigl(1+w(s)\bigr)+\nu_4+O(t^{-2}),\qquad\nu_4\in\R\ \text{independent of }z .
\end{equation}
\end{lemma}

\begin{proof}
\emph{Regularization.}  For $0<r<1$ let
\begin{equation}\label{eq:psi-r}
 \psi_r(y):=-\sum_{m\ge1}r^m\frac{\sin2\pi my}{\pi m}=\frac1\pi\Im\Log\bigl(1-re^{2\pi iy}\bigr).
\end{equation}
Since $\Re(1-re^{2\pi iy})>0$, $|\psi_r|\le\tfrac12$, and $\psi_r(y)\to\psi(y)$ as $r\uparrow1$ for $y\notin\Z$ (Abel summation of the convergent Fourier series of $\psi$).  The set $\{\tau:\NR(\tau)+\tfrac12\in\Z\}$ is discrete and $f'\in L^1(\tau_0,\infty)$, so by dominated convergence
\begin{equation}\label{eq:Psi-r-limit}
 \Lambda_{\rm saw}(z)=\lim_{r\uparrow1}\Lambda_r(z),\qquad
 \Lambda_r(z):=\int_{\tau_0}^\infty\psi_r\bigl(\NR+\tfrac12\bigr)f'\,d\tau .
\end{equation}
\emph{Termwise integration.}  For fixed $r<1$ the series for $\psi_r(\NR+\tfrac12)$, with $\sin2\pi m(y+\tfrac12)=(-1)^m\sin2\pi my$, is dominated by $\sum_mr^m/m$; by Fubini,
\begin{equation}\label{eq:Psi-r-series}
 \Lambda_r(z)=-\sum_{m\ge1}\frac{(-r)^m}{2\pi i\,m}\bigl(I_m^+-I_m^-\bigr),\qquad
 I_m^\pm:=\int_{\tau_0}^\infty e^{\pm2\pi im\NR(\tau)}f'(\tau)\,d\tau .
\end{equation}
\emph{Contours.}  Let $\mathcal H^\pm=\{\Re\tau\ge\tau_0,\ \pm\Im\tau\ge0\}$.  On $\mathcal H^\pm$, $|e^{\pm2\pi im\NR}|\le1$ by Lemma~\CGref{lem:TR-clock}, and $|f'(\tau)|\le4|z|^2|\tau|^{-3}$ for $|\tau|\ge2|z|$, so the quarter-circle arcs of radius $R$ contribute $O(|z|^2R^{-2})\to0$.  The poles of $f'$ are $0$, $iz=-t+i\sigma'\notin\mathcal H^\pm$, and $\tau_*=-iz=t-i\sigma'\in\mathcal H^-$, the last with residue $1$ (from $\frac1{\tau+iz}$).  Define the ray integrals
\begin{equation}\label{eq:ray-integrals}
 R_m^+:=i\int_0^\infty e^{2\pi im\NR(\tau_0+iy)}f'(\tau_0+iy)\,dy,\qquad
 R_m^-:=-i\int_0^\infty e^{-2\pi im\NR(\tau_0-iy)}f'(\tau_0-iy)\,dy .
\end{equation}
On $\mathcal H^+$ Cauchy's theorem gives $I_m^+=R_m^+$.  On $\mathcal H^-$ the positively oriented boundary of the quarter disc runs down the ray, along the arc, and back along $[\tau_0,\tau_0+R]$ from right to left; letting $R\to\infty$,
\begin{equation}\label{eq:lower-contour}
 R_m^--I_m^-=2\pi i\,e^{-2\pi im\NR(\tau_*)}=2\pi i\,w^m .
\end{equation}
Therefore, using $|rw|<1$,
\begin{equation}\label{eq:Psi-r-closed}
 \Lambda_r(z)=-\sum_{m\ge1}\frac{(-rw)^m}m+\Phi_r(z)=\Log(1+rw)+\Phi_r(z),\qquad
 \Phi_r(z):=-\sum_{m\ge1}\frac{(-r)^m}{2\pi i\,m}\bigl(R_m^+-R_m^-\bigr).
\end{equation}
\emph{The rays.}  Split $f'(\tau)=-\frac2\tau+g_z(\tau)$, $g_z(\tau)=\frac{2\tau}{\tau^2+z^2}$, and correspondingly $R_m^\pm=R_m^{\pm,(0)}+R_m^{\pm,(z)}$, $\Phi_r=\Phi_r^{(0)}+\Phi_r^{(z)}$.  The ray integrands carry the factor $2^{-my}$.  Since $|2/\tau|\le2/\tau_0$ on the rays, $|R^{\pm,(0)}_m|\le2/(\tau_0m\log2)$, so the series for $\Phi_r^{(0)}$ is dominated by $\sum_mC/m^2$ uniformly in $r$, and $\Phi_r^{(0)}\to\nu_4:=\Phi_1^{(0)}$, a constant depending only on $\tau_0$.  It is real: $\NR(\bar\tau)=\overline{\NR(\tau)}$ gives $R_m^{-,(0)}=\overline{R_m^{+,(0)}}$, so each term equals $-(-r)^m\Im R_m^{+,(0)}/(\pi m)$.  For $\tau=\tau_0\pm iy$, $\Re(\tau^2+z^2)=\tau_0^2+\sigma'^2-y^2-t^2\le-\frac12(y^2+t^2)$ because $t^2\ge2(\tau_0^2+B^2)$; hence $|g_z(\tau)|\le4(\tau_0+y)t^{-2}$ and
\begin{equation}\label{eq:ray-bound}
 |R^{\pm,(z)}_m|\le\frac4{t^2}\int_0^\infty2^{-my}(\tau_0+y)\,dy=\frac4{t^2}\Bigl(\frac{\tau_0}{m\log2}+\frac1{(m\log2)^2}\Bigr).
\end{equation}
So $|\Phi_r^{(z)}(z)|\le C\,t^{-2}\sum_mm^{-2}$, uniformly in $0<r<1$ and $0<\sigma'\le B$, and $\Phi_r^{(z)}$ converges as $r\uparrow1$ by dominated convergence of the series.  Since $\Log(1+rw)\to\Log(1+w)$, letting $r\uparrow1$ in \CGeqref{eq:Psi-r-closed} and \CGeqref{eq:Psi-r-limit} gives \CGeqref{eq:Psi-asymp}.  \emph{Uniformity.}  The three limits $r\uparrow1$ taken in this proof are dominated by $r$-independent majorants: $|\psi_r|\le\frac12$ against $|f'|\in L^1$ in \CGeqref{eq:Psi-r-limit}; after multiplication by the exterior factor $r^m/m\le1/m$, the ray bounds $|R_m^{\pm,(0)}|\le2/(\tau_0m\log2)$ and \CGeqref{eq:ray-bound} give the summable, $r$-independent majorants $C/m^2$ and $Ct^{-2}/m^2$ for the series $\Phi_r^{(0)}$ and $\Phi_r^{(z)}$, so dominated convergence applies to these series as $r\uparrow1$; and $|rw|\le|w|<1$.  Hence the remainder $O(t^{-2})$ in \CGeqref{eq:Psi-asymp} is uniform in $0<\sigma'\le B$ and $t\ge t_0(B)$, and does not depend on the regularization.
\end{proof}

\subsubsection*{Step 2: assembly}
By \CGeqref{eq:transition-product}, \CGeqref{eq:transition-split} and Lemmas~\CGref{lem:TR-smooth} and~\CGref{lem:TR-harmonic},
\begin{equation}\label{eq:assembly}
\begin{split}
 \zeta_{\RVM}(s)&=\frac{P_{\RVM}^{\crit}(s)}{G(s)}=P_{\RVM}^{\crit}(\tfrac12)\exp\bigl(\Lambda_{\rm sm}(z)-\log G(s)+\Lambda_{\rm saw}(z)\bigr)\\
 &=P_{\RVM}^{\crit}(\tfrac12)\,e^{\nu_3+\nu_4}(1+w)\bigl(1+O(t^{-1})\bigr),
\end{split}
\end{equation}
uniformly for $0<\sigma'\le B$.  This is \CGeqref{eq:transition-w} with
\begin{equation}\label{eq:varpi-def}
 \varpi_{\RVM}:=P_{\RVM}^{\crit}(\tfrac12)\,e^{\nu_3+\nu_4}>0,
\end{equation}
which depends only on the clock; its numerical value $0.9619415789\ldots$ plays no role in the proof.  Theorem~\CGref{thm:varpi-less-one} below proves $\varpi_{\RVM}<1$.

\subsubsection*{Step 3: from the clock phase to \texorpdfstring{$\chi$}{chi}}
Let
\begin{equation}\label{eq:theta-def}
 \vartheta(v):=\frac1{2i}\Bigl(\Log\Gamma\bigl(\tfrac14+\tfrac{iv}2\bigr)-\Log\Gamma\bigl(\tfrac14-\tfrac{iv}2\bigr)\Bigr)-\frac v2\log\pi,
\end{equation}
holomorphic on $|\Im v|<\tfrac12$.  The identity $\chi(\tfrac12+iv)=e^{-2i\vartheta(v)}$ holds for real $v$, hence on the strip, and $s=\tfrac12+i\tau_*$ with $|\Im\tau_*|=\sigma'<\tfrac12$.  Stirling's formula in the strip gives $\vartheta(v)=\frac v2\Log\frac v{2\pi}-\frac v2-\frac\pi8+O(|v|^{-1})$, i.e.
\begin{equation}\label{eq:clock-theta}
 \NR(v)=\frac{\vartheta(v)}\pi+1+O(|v|^{-1}).
\end{equation}
Therefore
\begin{equation}\label{eq:w-chi}
 w=e^{-2i\vartheta(\tau_*)-2\pi i+O(1/t)}=\chi(s)\bigl(1+O(t^{-1})\bigr).
\end{equation}
For $\sigma'\ge\sigma'_1$, \CGeqref{eq:chi-size} gives $|\chi(s)|\le\tfrac12$ when $t\ge t_0(\sigma'_1)$, so $1+w=(1+\chi)\bigl(1+(w-\chi)/(1+\chi)\bigr)=(1+\chi)(1+O(t^{-1}))$, which with \CGeqref{eq:transition-w} gives \CGeqref{eq:transition-chi}.  This completes the proof of Theorem~\CGref{thm:transition}. \qed

\begin{lemma}[Extension to the critical line]\label{lem:TR-boundary}
The function $\mathcal E_0=\zeta_{\RVM}/(\varpi_{\RVM}(1+w))-1$ extends holomorphically to $\{|\sigma'|<\tfrac12,\ t>t_0\}$, where $t_0=t_0(\tfrac14)$ is the threshold of Theorem~\CGref{thm:transition}(a) with $B=\tfrac14$, and $|\mathcal E_0|\le C/t$ for $0\le\sigma'\le\tfrac14$.
\end{lemma}

\begin{proof}
By Lemma~\CGref{lem:TR-clock}, $\Im\NR(\tau_*)\ne0$ for $\sigma'\ne0$, so $1+w$ vanishes in the strip only on $\sigma'=0$, exactly at $t=\tau_k$; these zeros are simple because $\frac{d}{dt}w=-i\log(\tau_*/2\pi)\,w\ne0$.  The zeros of $\zeta_{\RVM}$ in the strip are the same simple points.  Hence $\mathcal E_0$ has only removable singularities, and the bound extends to $\sigma'=0$ by continuity.
\end{proof}

\begin{corollary}[Derivative at the zeros]\label{cor:derivative-zeros}
For $k$ with $\tau_k>t_0$,
\begin{equation}\label{eq:derivative-zeros}
 \zeta_{\RVM}'(\rho_k^{\RVM})=\varpi_{\RVM}\log\frac{\tau_k}{2\pi}\Bigl(1+O\Bigl(\frac1{\tau_k}\Bigr)\Bigr).
\end{equation}
\end{corollary}

\begin{proof}
By Lemma~\CGref{lem:TR-boundary}, $\zeta_{\RVM}=\varpi_{\RVM}(1+w)(1+\mathcal E_0)$ holds holomorphically near $\sigma'=0$.  At $s=\rho_k^{\RVM}$, $1+w=0$, so $\zeta_{\RVM}'(\rho_k^{\RVM})=\varpi_{\RVM}w'(\rho_k^{\RVM})(1+\mathcal E_0(\rho_k^{\RVM}))$.  From \CGeqref{eq:transition-zw}, $w'(s)=-2\pi i\,\NR'(\tau_*)\cdot(-i)\,w(s)=-\log(\tau_*/2\pi)\,w(s)$, and at the zero $w=-1$, $\tau_*=\tau_k$.  Finally $|\mathcal E_0|\le C/\tau_k$.
\end{proof}

\begin{remark}[Structure of the proof]\label{rem:TR-audit}
The delicate points are handled as follows: the smooth part is computed with principal branches by the identity theorem (Lemma~\CGref{lem:TR-smooth}); the clock estimate is quantified (Lemma~\CGref{lem:TR-clock}); the sawtooth series is replaced by its Abel means, termwise integration is justified for $r<1$, and $r\uparrow1$ is taken after the contour argument uniformly in $r$ and $z$ (Lemma~\CGref{lem:TR-harmonic}); the orientation of the lower contour and the sign of the residue are written out in \CGeqref{eq:lower-contour}; the $z$-dependence of the rays is isolated in $g_z$ and is $O(t^{-2})$, the rest is the explicit real constant $\nu_4$; Stirling is used uniformly for $0<\sigma'\le B$, and the only place where $\sigma'<\tfrac12$ enters is Step~3.
\end{remark}

\begin{corollary}[Vertical constant limit]\label{cor:vertical-limit}
For every fixed $\sigma>\tfrac12$,
\begin{equation}\label{eq:vertical-limit}
 \zeta_{\RVM}(\sigma+it)\longrightarrow\varpi_{\RVM},\qquad
 F_{\RVM}(\sigma+it)\longrightarrow L:=\frac1{\varpi_{\RVM}}\qquad(|t|\to\infty).
\end{equation}
More precisely, with $\sigma'=\sigma-\tfrac12$,
\begin{equation}\label{eq:w-decay}
 |w(\sigma+it)|\le (|t|/2\pi)^{-\sigma'}(1+O_B(|t|^{-1}))\qquad(0<\sigma'\le B),
\end{equation}
and, with $\alpha_c:=\min(c-\tfrac12,1)$ for $c>\tfrac12$,
\begin{equation}\label{eq:FminusL-all}
 F_{\RVM}(c+it)-L=O_c(|t|^{-\alpha_c}).
\end{equation}
\end{corollary}

\begin{proof}
Expanding the clock at $\tau_*=t-i\sigma'$ gives $\Im\NR(\tau_*)=-\frac{\sigma'}{2\pi}\log\frac t{2\pi}+O_\sigma(t^{-1})$, hence $|w|=e^{2\pi\Im\NR(\tau_*)}=O_\sigma(t^{-\sigma'})$.  Inverting \CGeqref{eq:transition-w} for large $t$ gives $F_{\RVM}=L(1+O(|w|)+O(t^{-1}))$.
\end{proof}

\subsection{Bounds on the whole half-plane}
\label{subsec:halfplane}

The constants of Theorem~\CGref{thm:transition} depend on the width $B$ of the band.  The following estimate holds on the whole half-plane and does not use contours.  Put
\begin{equation}\label{eq:saw-primitive}
 \mathfrak s_1(\tau):=\int_{\tau_0}^{\tau}\psi\bigl(\NR(v)+\tfrac12\bigr)\,dv\qquad(\tau\ge\tau_0).
\end{equation}

\begin{lemma}[Bounded primitive of the midpoint sawtooth]\label{lem:saw-primitive}
$|\mathfrak s_1(\tau)|\le\dfrac{\pi}{4\log2}$ for all $\tau\ge\tau_0$.
\end{lemma}

\begin{proof}
Substitute $y=\NR(v)$, $dv=\mathcal T'(y)\,dy$ with $\mathcal T'(y)=2\pi/\log(\mathcal T(y)/2\pi)$ from \CGeqref{eq:inverse-clock}.  The weight $\mathcal T'$ is positive and decreasing, with $\mathcal T'\le2\pi/\log2$ on $y\ge\NR(\tau_0)$ because $\log(\tau_0/2\pi)=\log2$.  The function $\Phi(y):=\frac12\bigl(\{y+\frac12\}-\frac12\bigr)^2-\frac18$ is continuous, $1$-periodic, with values in $[-\tfrac18,0]$, and $\Phi'(y)=\{y+\frac12\}-\frac12=\psi(y+\tfrac12)$ for $y\notin\frac12+\Z$; it is therefore a primitive of $y\mapsto\psi(y+\tfrac12)$, so every integral of $\psi(y+\frac12)$ over an interval has modulus at most $\tfrac18$.  The second mean-value theorem gives $|\mathfrak s_1(\tau)|\le\frac{2\pi}{\log2}\cdot\frac18$.
\end{proof}

\begin{theorem}[Two-sided bound and limit at $+\infty$]\label{thm:halfplane-RVM}
There is an absolute constant $R_0$ such that, for $s=\tfrac12+z$, $z=\sigma'+it$, $\sigma'>\tfrac12$ and $|z|\ge R_0$,
\begin{equation}\label{eq:halfplane-RVM}
 \zeta_{\RVM}(s)=\varpi_{\RVM}\,e^{\theta_{\RVM}(s)},\qquad |\theta_{\RVM}(s)|\le\frac{\pi^2}{2\log2}\cdot\frac1{\sigma'}+\frac{C}{|z|}.
\end{equation}
Consequently, for every $\sigma_0>1$ there is $C(\sigma_0)$ with
\begin{equation}\label{eq:halfplane-two-sided}
 C(\sigma_0)^{-1}\le|\zeta_{\RVM}(s)|\le C(\sigma_0)\qquad(\Re s\ge\sigma_0),
\end{equation}
and $\zeta_{\RVM}(s)\to\varpi_{\RVM}$ uniformly in $t$ as $\Re s\to+\infty$.  Moreover
\begin{equation}\label{eq:nu4-closed}
 \nu_4=-2\int_{\tau_0}^\infty\frac{\mathfrak s_1(\tau)}{\tau^2}\,d\tau .
\end{equation}
\end{theorem}

\begin{proof}
\emph{Smooth part.}  For $\Re z>0$ the computation in the proof of Lemma~\CGref{lem:TR-smooth} is exact:
\begin{equation}\label{eq:smooth-exact}
 \Lambda_{\rm sm}(z)=\frac z2\Log\frac z{2\pi e}+\frac78\Log\Bigl(1+\frac{z^2}{\tau_0^2}\Bigr)-\frac1\pi\int_0^{\tau_0}\log\frac\tau{2\pi e}\,\frac{z^2}{\tau^2+z^2}\,d\tau .
\end{equation}
For $|z|\ge2\tau_0$ one has $|\tau^2+z^2|\ge\frac34|z|^2$ on $[0,\tau_0]$, and Stirling's formula is uniform in $|\arg(s/2)|<\frac\pi2$; the expansions \CGeqref{eq:J-asymp} and \CGeqref{eq:G-asymp} were derived for $t>0$, but every step (the branches of $\mathrm{Log}(z^2)$, of $\mathrm{Log}(1+z^2/\tau_0^2)$, of $\mathrm{Log}(z^2-\frac14)$ and of $f(\tau)=\mathrm{Log}(1+z^2/\tau^2)$ on $[\tau_0,\infty)$) only uses $|\arg z|<\frac\pi2$, which holds on $\Re z>0$ for either sign of $t$; they therefore hold with $O(|z|^{-1})$ in place of $O(t^{-1})$, uniformly on $\Re z>\tfrac12$, and $\Lambda_{\rm sm}(z)-\log G(s)=\nu_3+O(|z|^{-1})$.

\emph{Sawtooth part.}  Integrating \CGeqref{eq:transition-split} by parts with \CGeqref{eq:saw-primitive} ($\mathfrak s_1(\tau_0)=0$, $f'(\tau)=O(|z|^2\tau^{-3})$),
\begin{equation}\label{eq:saw-ibp}
 \Lambda_{\rm saw}(z)=-\int_{\tau_0}^\infty\mathfrak s_1f''\,d\tau,\qquad f''(\tau)=\frac2{\tau^2}-\frac1{(\tau-iz)^2}-\frac1{(\tau+iz)^2}.
\end{equation}
Since $|\tau\mp iz|^2=(\tau\pm t)^2+\sigma'^2$ and $\int_{\R}\frac{d\tau}{(\tau-a)^2+\sigma'^2}=\frac\pi{\sigma'}$, Lemma~\CGref{lem:saw-primitive} gives
\begin{equation}\label{eq:saw-bound}
 \Bigl|\Lambda_{\rm saw}(z)+2\int_{\tau_0}^\infty\frac{\mathfrak s_1(\tau)}{\tau^2}\,d\tau\Bigr|\le\frac{\pi}{4\log2}\cdot\frac{2\pi}{\sigma'} .
\end{equation}

\emph{Assembly.}  By \CGeqref{eq:assembly}, $\log\zeta_{\RVM}(s)=\log\varpi_\infty+\theta_{\RVM}(s)$ with $\varpi_\infty:=P_{\RVM}^{\crit}(\tfrac12)\exp(\nu_3-2\int_{\tau_0}^\infty\mathfrak s_1\tau^{-2}d\tau)>0$ and $\theta_{\RVM}$ bounded as in \CGeqref{eq:halfplane-RVM}.  Both $\varpi_\infty$ and $\varpi_{\RVM}$ are positive real numbers, and we use their real logarithms.  Taking moduli, $|\zeta_{\RVM}(s)|=\varpi_\infty e^{\Re\theta_{\RVM}(s)}$; for fixed $\sigma'$ and $t\to\infty$, Corollary~\CGref{cor:vertical-limit} gives $|\zeta_{\RVM}(s)|\to\varpi_{\RVM}$, hence $|\log\varpi_{\RVM}-\log\varpi_\infty|\le\pi^2/(2\sigma'\log2)$, with no $2\pi i$-ambiguity; letting $\sigma'\to\infty$ gives $\varpi_\infty=\varpi_{\RVM}$, which is \CGeqref{eq:nu4-closed} by \CGeqref{eq:varpi-def}.  On the compact set $\{\Re s\ge\sigma_0,\ |z|\le R_0\}$ the function $\zeta_{\RVM}$ is continuous and has no zeros or poles, which gives \CGeqref{eq:halfplane-two-sided}.
\end{proof}

\begin{theorem}[The transition constant is less than one]\label{thm:varpi-less-one}
The smooth constant of Lemma~\CGref{lem:TR-smooth} is
\begin{equation}\label{eq:nu3-closed}
 \nu_3=8-\frac{33}4\log2-2\log\pi=-0.0079240113\ldots,
\end{equation}
the sawtooth constant of Lemma~\CGref{lem:TR-harmonic} satisfies $\nu_4<-0.012691$, and
\begin{equation}\label{eq:varpi-less-one}
 \varpi_{\RVM}=P^{\crit}_{\RVM}(\tfrac12)\,e^{\nu_3+\nu_4}<P^{\crit}_{\RVM}(\tfrac12)\,e^{-0.020615}<0.97960 .
\end{equation}
In particular $\varpi_{\RVM}<1$ and $L=1/\varpi_{\RVM}>1$.
\end{theorem}

\begin{proof}
\emph{(a) The canonical factor.}  $q(\frac12)=-\frac14$, so $P^{\crit}_{\RVM}(\frac12)=\prod_k\bigl(1-\frac{1/4}{1/4+\tau_k^2}\bigr)=\prod_k\frac{\tau_k^2}{\frac14+\tau_k^2}$; every factor lies in $(0,1)$ and $\sum_k\tau_k^{-2}<\infty$, so $0<P^{\crit}_{\RVM}(\frac12)<1$.  No numerical value of this product is used.

\emph{(b) The smooth constant.}  In \CGeqref{eq:J-asymp}, with $\tau_0=4\pi$, $\int_0^{\tau_0}\log\frac\tau{2\pi e}\,d\tau=\tau_0\log\frac{\tau_0}{2\pi e}-\tau_0=4\pi\log2-8\pi$, so $\nu_1=8-4\log2-\frac74\log(4\pi)=8-\frac{15}2\log2-\frac74\log\pi$; in \CGeqref{eq:G-asymp}, $\nu_2=\frac14\log2+\frac12\log2\pi-\frac14\log\pi=\frac34\log2+\frac14\log\pi$.  Hence $\nu_3=\nu_1-\nu_2$ is \CGeqref{eq:nu3-closed}.  Its sign is the inequality $e^8<2^{33/4}\pi^2$: with $\log2>0.69314$ and $\log\pi>1.14472$, $\frac{33}4\log2+2\log\pi>5.718405+2.28944>8.0078$, so $\nu_3<-0.0078$.

\emph{(c) The sawtooth constant as a sum over the periods of the clock.}  By \CGeqref{eq:nu4-closed} and one integration by parts ($\mathfrak s_1'=\psi(\NR+\frac12)$, $\mathfrak s_1(\tau_0)=0$, $\mathfrak s_1$ bounded by Lemma~\CGref{lem:saw-primitive}),
\begin{equation}\label{eq:nu4-limit}
 \nu_4=-2\lim_{X\to\infty}\int_{\tau_0}^X\psi\bigl(\NR(\tau)+\tfrac12\bigr)\frac{d\tau}\tau ,
\end{equation}
and the limit may be taken along $X=\tau_{M+1}$, $M\to\infty$, because the left side of the integration by parts converges absolutely.  Since $\NR$ increases on $[2\pi,\infty)$ and $0<\NR(\tau_0)=2\log\frac2e+\frac78<\frac12$, on $[\tau_0,\tau_1)$ one has $\NR+\frac12\in(\frac12,1)$ and $\psi(\NR+\frac12)=\NR$; so the first piece is
\begin{equation}\label{eq:nu4-I0}
 \mathscr I_0(\tau_1):=\int_{\tau_0}^{\tau_1}\NR(\tau)\frac{d\tau}\tau=\frac{\tau_1\ell_1-2\tau_1-\tau_0\log2+2\tau_0}{2\pi}+\frac78\log\frac{\tau_1}{\tau_0},\qquad \ell_1:=\log\frac{\tau_1}{2\pi}.
\end{equation}
For $n\ge1$, on $[\tau_n,\tau_{n+1})$ one has $\NR+\frac12\in[n,n+1)$ and $\psi(\NR+\frac12)=\NR+\frac12-n-\frac12$.  Substitute $u=\NR(\tau)+\frac12$, i.e.\ $\tau=\mathcal T(u-\frac12)$ with the inverse clock of \CGeqref{eq:inverse-clock}; then
\begin{equation}\label{eq:nu4-weight}
 \frac{d\tau}\tau=\mathscr W(u)\,du,\qquad \mathscr W(u):=\frac{2\pi}{\mathcal T\,\log(\mathcal T/2\pi)},\qquad \mathcal T=\mathcal T(u-\tfrac12),
\end{equation}
and $\mathscr I_n:=\int_{\tau_n}^{\tau_{n+1}}\psi(\NR+\frac12)\frac{d\tau}\tau=\int_n^{n+1}(u-n-\frac12)\mathscr W(u)\,du$.  With $\ell=\log(\mathcal T/2\pi)>0$ for $u\ge1$,
\[
 \mathscr W'(u)=-\frac{4\pi^2(\ell+1)}{\mathcal T^2\ell^3}<0,\qquad
 \frac d{d\mathcal T}\log\frac{\ell+1}{\mathcal T^2\ell^3}=-\frac{2\ell^2+4\ell+3}{\mathcal T\ell(\ell+1)}<0,
\]
so $-\mathscr W'$ is decreasing in $u$: $\mathscr W$ is decreasing and convex on $[1,\infty)$, and $\mathscr W(u),\mathscr W'(u)\to0$ as $u\to\infty$.  Let $\mathbf B_2(x)=x^2-x+\frac16$ and $\mathbf B_3(x)=x^3-\frac32x^2+\frac12x$ be the Bernoulli polynomials, $\bar{\mathbf B}_j(u):=\mathbf B_j(u-\lfloor u\rfloor)$; since $\frac12\mathbf B_2'(x)=x-\frac12$, $\frac16\mathbf B_3'=\frac12\mathbf B_2$, $\mathbf B_2(0)=\mathbf B_2(1)=\frac16$ and $\mathbf B_3(0)=\mathbf B_3(1)=0$, two integrations by parts (the Euler--Maclaurin formula on one period) give
\[
 \mathscr I_n=\frac{\mathscr W(n+1)-\mathscr W(n)}{12}+\frac16\int_n^{n+1}\bar{\mathbf B}_3(u)\,\mathscr W''(u)\,du .
\]
Summing over $1\le n\le M$ and letting $M\to\infty$,
\[
 \sum_{n\ge1}\mathscr I_n=-\frac{\mathscr W(1)}{12}+\frac16\int_1^\infty\bar{\mathbf B}_3\,\mathscr W''\,du,\qquad
 \Bigl|\frac16\int_1^\infty\bar{\mathbf B}_3\,\mathscr W''\,du\Bigr|\le\frac{\sqrt3}{216}\int_1^\infty\mathscr W''\,du=\frac{\sqrt3}{216}\,|\mathscr W'(1)|,
\]
because $\max_{[0,1]}|\mathbf B_3|=\mathbf B_3\bigl(\frac{3-\sqrt3}6\bigr)=\frac{\sqrt3}{36}$ and $\mathscr W''\ge0$.  Since $\mathcal T(\frac12)=\tau_1$, $\mathscr W(1)=2\pi/(\tau_1\ell_1)$ and $|\mathscr W'(1)|=4\pi^2(\ell_1+1)/(\tau_1^2\ell_1^3)$.  With \CGeqref{eq:nu4-limit},
\begin{equation}\label{eq:nu4-bound}
 \nu_4=-2\Bigl(\mathscr I_0(\tau_1)+\sum_{n\ge1}\mathscr I_n\Bigr)\le-2\,\mathfrak M(\tau_1),\qquad
 \mathfrak M(t):=\mathscr I_0(t)-\frac{\pi}{6t\ell}-\frac{\sqrt3}{216}\cdot\frac{4\pi^2(\ell+1)}{t^2\ell^3},\quad \ell=\log\frac t{2\pi}.
\end{equation}

\emph{(d) The numerical inequalities.}  Since
\[
 \NR(14.52134)=0.4999991\ldots<\tfrac12<0.5000004\ldots=\NR(14.52135),
\]
$\tau_1\in[14.52134,14.52135]$ (the Lambert formula \CGeqref{eq:lambert} gives $\tau_1=14.5213469530\ldots$).  On this interval $\mathscr I_0\in[0.054079,0.054086]$, $\pi/(6t\ell)\in[0.0430408,0.0430410]$, the last term of $\mathfrak M$ lies in $[0.0046925,0.0046926]$, and therefore $\mathfrak M(\tau_1)\in[0.0063457,0.0063526]$; hence $\nu_4\le-0.0126914$.  These enclosures, and those of (b), are computed in interval arithmetic with outward rounding by the program \texttt{s14\_varpi\_bound.py} (Appendix~\CGref{app:code-s14}, output in Appendix~\CGref{app:output-s14}); the margins ($\nu_3<-0.0079$, $\mathfrak M>0.0063$ against terms of size $0.05$) are large enough for a hand computation with five decimals.

\emph{(e) Conclusion.}  By \CGeqref{eq:varpi-def}, $\log\varpi_{\RVM}=\log P^{\crit}_{\RVM}(\frac12)+\nu_3+\nu_4<0-0.0079240-0.0126914<-0.020615$, and $e^{-0.020615}<0.97960$.
\end{proof}

\begin{remark}[High-precision values; what the proof uses]\label{rem:varpi-values}
The proof uses only $P^{\crit}_{\RVM}(\frac12)<1$ and the two upper bounds $\nu_3<0$, $\nu_4\le-2\mathfrak M(\tau_1)$.  The floating-point values computed by the same program (Part~B of its output, not part of the proof) are $\nu_4=-0.0251569083$ (period-by-period quadrature of $\mathscr I_n$ for $n\le2000$ and the Euler--Maclaurin tail), $e^{\nu_3+\nu_4}=0.9674602699$, $P^{\crit}_{\RVM}(\frac12)=0.9942956924$ (Lambert clock for $k\le2\cdot10^6$ and the tail integral), and $\varpi_{\RVM}=0.9619415789$.  The Richardson value $0.9619415789$ of Experiment~\CGref{exp:transition} is obtained from $\zeta_{\RVM}=P^{\crit}_{\RVM}/G$ on $\Re s=2$ without using $\nu_3$ or $\nu_4$; the agreement to $2\cdot10^{-11}$ is an independent check of \CGeqref{eq:varpi-def} and \CGeqref{eq:nu4-closed}.  The Euler--Maclaurin remainder $\frac16\int\bar{\mathbf B}_3\mathscr W''$ is $+0.00154$, about a third of the bound used, which is why the true $\nu_4$ is about twice the bound.
\end{remark}

Theorem~\CGref{thm:halfplane-RVM} extends the two-sided bound \CGeqref{eq:zetaRVM-two-sided} of Corollary~\CGref{cor:Q-control} from bands to the whole half-plane $\Re s\ge\sigma_0>1$, and identifies $\varpi_{\RVM}$ also as the limit of $\zeta_{\RVM}(\sigma)$ as $\sigma\to+\infty$.

\begin{theorem}[No absolutely convergent Dirichlet expansion]\label{thm:no-Dirichlet-RVM}
Neither $\zeta_{\RVM}$ nor $F_{\RVM}=1/\zeta_{\RVM}$ is represented by a Dirichlet series $\sum_na_nn^{-s}$ converging absolutely in some half-plane $\Re s>\sigma_0$.
\end{theorem}

\begin{proof}
Suppose $F_{\RVM}(s)=\sum_{n\ge1}a_nn^{-s}$ absolutely for $\Re s>\sigma_0$.  Fix $\sigma>\max(\sigma_0,\tfrac12)$ and put $\tilde a_1=a_1-L$, $\tilde a_n=a_n$ ($n\ge2$), so $F_{\RVM}(\sigma+it)-L=\sum_n\tilde a_nn^{-\sigma-it}$ with $\sum_n|\tilde a_n|n^{-\sigma}<\infty$.  Then
\begin{equation}\label{eq:mean-square}
 \frac1{2T}\int_{-T}^{T}\bigl|F_{\RVM}(\sigma+it)-L\bigr|^2dt
 =\sum_{m,n}\tilde a_m\overline{\tilde a_n}(mn)^{-\sigma}\,\frac1{2T}\int_{-T}^{T}\Bigl(\frac nm\Bigr)^{it}dt ,
\end{equation}
the interchange being justified by absolute convergence of the double series, uniformly in $t$.  The inner average is $1$ for $m=n$ and is bounded by $\min\bigl(1,1/(T|\log(n/m)|)\bigr)\to0$ for $m\ne n$; each term is bounded by $|\tilde a_m\tilde a_n|(mn)^{-\sigma}$, which is summable, so dominated convergence gives
\begin{equation}\label{eq:mean-square-limit}
 \lim_{T\to\infty}\frac1{2T}\int_{-T}^{T}\bigl|F_{\RVM}(\sigma+it)-L\bigr|^2dt=\sum_{n\ge1}|\tilde a_n|^2n^{-2\sigma}.
\end{equation}
By \CGeqref{eq:vertical-limit} the integrand tends to $0$ and is bounded on the line (continuous, with a finite limit), so the left side of \CGeqref{eq:mean-square-limit} is $0$.  Hence every $\tilde a_n=0$, $F_{\RVM}\equiv L$ in the half-plane and, by analytic continuation, everywhere, which contradicts $F_{\RVM}(1)=0$.  The same argument applied to $\zeta_{\RVM}-\varpi_{\RVM}$ forces an absolutely convergent Dirichlet series for $\zeta_{\RVM}$ to be the constant $\varpi_{\RVM}$, which contradicts its pole at $s=1$.
\end{proof}

\begin{theorem}[Unsmoothed synthetic Mertens bound]\label{thm:unsmoothed-RVM}
For every $x>1$ and every $c>\tfrac12$ the symmetric Perron integral
\begin{equation}\label{eq:Mcrit-symmetric}
 M_{\RVM}^{\crit}(x):=\lim_{U\to\infty}\frac1{2\pi i}\int_{c-iU}^{c+iU}F_{\RVM}(s)\frac{x^s}{s}\,ds
\end{equation}
exists, is independent of $c>\tfrac12$, and equals the absolutely convergent expression
\begin{equation}\label{eq:Mcrit-renorm}
 M_{\RVM}^{\crit}(x)=L+\frac1{2\pi i}\int_{(c)}\bigl(F_{\RVM}(s)-L\bigr)\frac{x^s}{s}\,ds .
\end{equation}
For every $\eps\in(0,\tfrac12)$,
\begin{equation}\label{eq:Mcrit-unsmoothed-bound}
 M_{\RVM}^{\crit}(x)=L+O_\eps(x^{1/2+\eps}).
\end{equation}
\end{theorem}

\begin{proof}
Fix $c>\tfrac12$.  By \CGeqref{eq:FminusL-all} and continuity of $F_{\RVM}$ on $\Re s=c$ (no poles there by Proposition~\CGref{prop:Mcrit-poles}), the integral in \CGeqref{eq:Mcrit-renorm} converges absolutely.  The Perron kernel gives, for $x>1$ and $c>0$ \cite[Ch.~17]{Davenport1980} (for $0<x<1$ the limit is $0$),
\begin{equation}\label{eq:bromwich}
 \lim_{U\to\infty}\frac1{2\pi i}\int_{c-iU}^{c+iU}\frac{x^s}{s}\,ds=1 .
\end{equation}
Splitting $F_{\RVM}=L+(F_{\RVM}-L)$ proves existence of \CGeqref{eq:Mcrit-symmetric} and \CGeqref{eq:Mcrit-renorm}.  If $\tfrac12<c<c'$, the integrand of \CGeqref{eq:Mcrit-symmetric} is holomorphic on the closed strip $c\le\Re s\le c'$ (Proposition~\CGref{prop:Mcrit-poles}); by Theorem~\CGref{thm:transition}(a) with $B=c'-\tfrac12$, $F_{\RVM}$ is bounded on the horizontal segments $\Im s=\pm U$, $U\ge t_0(B)$, so these contribute $O(x^{c'}U^{-1})\to0$, and Cauchy's theorem gives independence of $c$.  Take $c=\tfrac12+\eps$ in \CGeqref{eq:Mcrit-renorm}:
\begin{equation}\label{eq:Mcrit-final-estimate}
 |M_{\RVM}^{\crit}(x)-L|\le\frac{x^{1/2+\eps}}{2\pi}\int_{-\infty}^{\infty}\frac{|F_{\RVM}(\tfrac12+\eps+it)-L|}{|\tfrac12+\eps+it|}\,dt,
\end{equation}
and the integral is finite by \CGeqref{eq:FminusL-all}, the integrand being $O(|t|^{-1-\alpha_c})$.
\end{proof}

\begin{remark}[Scope]\label{rem:unsmoothed-scope}
Theorem~\CGref{thm:unsmoothed-RVM} concerns $1/\zeta_{\RVM}$.  The transfer to $M$ goes through the endpoint of the $\lambda$-family and is Theorem~\CGref{thm:step3}.
\end{remark}

\section{The exact Hadamard product and the comparison with \texorpdfstring{$\zeta$}{zeta}}
\label{sec:comparison}

Put
\begin{equation}\label{eq:Pzeta}
 P_\zeta(s):=G(s)\zeta(s)=2\xi(s),\qquad \xi(s)=\tfrac12s(s-1)\pi^{-s/2}\Gamma(s/2)\zeta(s).
\end{equation}

\begin{theorem}[Exact Hadamard product]\label{thm:Hadamard}
Let $\rho$ run over the nontrivial zeros of $\zeta$ with $\Im\rho>0$, with multiplicity, ordered by $\Im\rho$.  The product
\begin{equation}\label{eq:Hadamard-exact}
 P_\zeta(s)=\prod_{\Im\rho>0}\left(1+\frac{q(s)}{\rho(1-\rho)}\right)
\end{equation}
converges absolutely and locally uniformly in $s\in\C$, and the identity holds with no exponential or constant factor.
\end{theorem}

\begin{proof}
$\xi$ is entire of order one with $\xi(0)=\tfrac12$, so Hadamard's factorization \cite{Hadamard1896}, \cite[\S2.12]{Titchmarsh1986}, \cite[Ch.~12]{Davenport1980} gives
\begin{equation}\label{eq:hadamard-general}
 2\xi(s)=e^{Bs}\prod_{\rho}\Bigl(1-\frac s\rho\Bigr)e^{s/\rho},
\end{equation}
the product running over all nontrivial zeros and converging absolutely and locally uniformly: for $|s|\le K$ and $|\rho|\ge2K$ the logarithm of the factor $(1-s/\rho)e^{s/\rho}$ is $O_K(|\rho|^{-2})$, and $\sum_\rho|\rho|^{-2}<\infty$.  The functional equation $\xi(s)=\xi(1-s)$ makes $\rho\mapsto1-\rho$ an involution of the zero multiset; no nontrivial zero is real, since $\eta(\sigma):=(1-2^{1-\sigma})\zeta(\sigma)=\sum_{n\ge1}(-1)^{n-1}n^{-\sigma}>0$ for $0<\sigma<1$ (an alternating series with decreasing terms) while $1-2^{1-\sigma}<0$ there, so $\zeta(\sigma)<0$; so this involution exchanges $\Im\rho>0$ and $\Im\rho<0$.  Grouping $\rho$ with $1-\rho$, which absolute convergence permits,
\begin{equation}\label{eq:pair-factor}
 \Bigl(1-\frac s\rho\Bigr)\Bigl(1-\frac s{1-\rho}\Bigr)e^{s/\rho+s/(1-\rho)}
 =\Bigl(1+\frac{q(s)}{\rho(1-\rho)}\Bigr)e^{s/(\rho(1-\rho))}.
\end{equation}
Since $|\rho(1-\rho)|\ge(\Im\rho)^2$ and $\sum_{\Im\rho>0}(\Im\rho)^{-2}<\infty$, the numbers $1/(\rho(1-\rho))$ are absolutely summable, so the exponentials collect into a single factor:
\begin{equation}\label{eq:hadamard-exp}
 2\xi(s)=e^{B's}\,\Pi_\zeta(s),\qquad \Pi_\zeta(s):=\prod_{\Im\rho>0}\Bigl(1+\frac{q(s)}{\rho(1-\rho)}\Bigr),\qquad B'=B+\sum_{\Im\rho>0}\frac1{\rho(1-\rho)} .
\end{equation}
$\Pi_\zeta$ depends on $s$ only through $q(s)$, hence $\Pi_\zeta(1-s)=\Pi_\zeta(s)$.  From $\xi(1-s)=\xi(s)$ we get $e^{B's}=e^{B'(1-s)}$ on the open set where $\Pi_\zeta\ne0$, so $B'=0$. No constant factor is left either: the general factorization \CGeqref{eq:hadamard-general} carries the factor $e^{A+Bs}$ with $e^{A}=2\xi(0)=1$, which is why it was written without $A$; equivalently, at $s=0$ one has $q(0)=0$, $\Pi_\zeta(0)=1$ and $2\xi(0)=1$, so \CGeqref{eq:hadamard-exp} with $B'=0$ is \CGeqref{eq:Hadamard-exact}.
\end{proof}

Define the exact quotient
\begin{equation}\label{eq:Q-RVM}
 \mathcal Q_{\RVM}(s):=\frac{\zeta(s)}{\zeta_{\RVM}(s)}=\frac{P_\zeta(s)}{P_{\RVM}^{\crit}(s)}.
\end{equation}

\begin{proposition}[Exact transform relation]\label{prop:M-transform-relation}
For $\Re s>1$, and by meromorphic continuation,
\begin{equation}\label{eq:M-transform-relation}
 F_{\RVM}(s)=\mathcal Q_{\RVM}(s)\,\frac1{\zeta(s)} .
\end{equation}
If the inverse Mellin transform $\mathcal Q_{\RVM}^{\vee}$ of $\mathcal Q_{\RVM}$ and the inverse defining $M_{\RVM}^{\crit}$ exist in a common distribution space with Mellin convolution, then
\begin{equation}\label{eq:M-convolution}
 M_{\RVM}^{\crit}=\mathcal Q_{\RVM}^{\vee}\Mell M .
\end{equation}
\end{proposition}

\CGeqref{eq:M-transform-relation} is the identity $1/\zeta_{\RVM}=(\zeta/\zeta_{\RVM})(1/\zeta)$.  Equation~\CGeqref{eq:M-convolution} is only notation for the Mellin-convolution theorem when both inverse transforms belong to a specified common convolution space; no such space is constructed here.

\begin{corollary}[All growth of $\zeta$ to the right of $\frac12$ lies in $\mathcal Q_{\RVM}$]\label{cor:Q-control}
For every $c'>\tfrac12$ and $0<\eps<c'-\tfrac12$ there are $t_0=t_0(\eps,c')$ and $0<c_-(\eps,c')\le c_+(\eps,c')$ such that
\begin{equation}\label{eq:zetaRVM-two-sided}
 c_-\le|\zeta_{\RVM}(\sigma+it)|\le c_+\qquad(\tfrac12+\eps\le\sigma\le c',\ |t|\ge t_0),
\end{equation}
and therefore, in the same region,
\begin{equation}\label{eq:zeta-asymp-Q}
 c_-|\mathcal Q_{\RVM}(s)|\le|\zeta(s)|\le c_+|\mathcal Q_{\RVM}(s)|,\qquad
 \frac1{c_+|\mathcal Q_{\RVM}(s)|}\le\Bigl|\frac1{\zeta(s)}\Bigr|\le\frac1{c_-|\mathcal Q_{\RVM}(s)|}.
\end{equation}
Consequently:
\begin{enumerate}[label=\textup{(\roman*)}]
\item the Lindel\"of hypothesis holds if and only if $\mathcal Q_{\RVM}(\sigma+it)=O_{\sigma,\eps}(|t|^{\eps})$ for every $\sigma>\tfrac12$;
\item every bound for $\zeta$ or $1/\zeta$ on $\tfrac12+\eps\le\sigma\le c'$, $|t|\ge t_0$, is equivalent, up to the constants $c_\pm$, to the same bound for $\mathcal Q_{\RVM}$ or $1/\mathcal Q_{\RVM}$.
\end{enumerate}
\end{corollary}

\begin{proof}
By Theorem~\CGref{thm:transition}(a) with $B=c'-\tfrac12$, $\zeta_{\RVM}=\varpi_{\RVM}(1+w)(1+\mathcal E_0)$ with $|\mathcal E_0|\le C_B/|t|$ and, by the uniform form \CGeqref{eq:w-decay}, $|w|\le C_B(|t|/2\pi)^{-\eps}$; for $|t|\ge t_0$ this gives $\tfrac12\le|1+w|\le2$ and $\tfrac12\le|1+\mathcal E_0|\le2$, hence \CGeqref{eq:zetaRVM-two-sided} with $c_-=\varpi_{\RVM}/4$, $c_+=4\varpi_{\RVM}$.  Then \CGeqref{eq:zeta-asymp-Q} follows from $\zeta=\mathcal Q_{\RVM}\zeta_{\RVM}$.  For (i): the Lindel\"of hypothesis is equivalent to $\zeta(\sigma+it)=O(|t|^\eps)$ for every $\sigma>\tfrac12$ \cite[Ch.~XIII]{Titchmarsh1986}, and \CGeqref{eq:zeta-asymp-Q} transfers this to $\mathcal Q_{\RVM}$.  (ii) is \CGeqref{eq:zeta-asymp-Q}.
\end{proof}

On the prime side define
\begin{equation}\label{eq:prime-defect}
 \mathfrak d_n^{\RVM}:=\mathbf 1_{\{n\ \mathrm{prime}\}}-c_n^{\RVM}.
\end{equation}
Since $\pi_{\RVM}(1)=0$, for every integer $N\ge2$
\begin{equation}\label{eq:prime-defect-cumulative}
 \sum_{2\le n\le N}\mathfrak d_n^{\RVM}=\pi(N)-\pi_{\RVM}(N).
\end{equation}

\section{The \texorpdfstring{$\lambda$}{lambda}-family from \texorpdfstring{$\zeta_{\RVM}$}{zeta\_RVM} to \texorpdfstring{$\zeta_{Z1}$}{zeta\_Z1}}
\label{sec:lambda-family}

For $\lambda\in[0,1]$ define the moving ordinates, the family product and the family zeta function
\begin{equation}\label{eq:gamma-lambda}
 \gamma_k(\lambda):=(1-\lambda)\gamma_k+\lambda\tau_k=\gamma_k+\lambda\eta_k,
\end{equation}
\begin{equation}\label{eq:P-lambda}
 P_\lambda(s):=\prod_{k\ge1}\left(1+\frac{q(s)}{\tfrac14+\gamma_k(\lambda)^2}\right),
\end{equation}
\begin{equation}\label{eq:zeta-lambda}
 \zeta_\lambda(s):=\frac{P_\lambda(s)}{G(s)},\qquad F_\lambda(s):=\frac1{\zeta_\lambda(s)},
\end{equation}
and the endpoint
\begin{equation}\label{eq:family-endpoints}
 P_{Z1}:=P_0,\qquad \zeta_{Z1}:=\zeta_0,\qquad \zeta_1=\zeta_{\RVM}.
\end{equation}

\begin{theorem}[The $\lambda$-family]\label{thm:path}
The product \CGeqref{eq:P-lambda} converges absolutely and locally uniformly jointly on $[0,1]\times\C$.  For every $\lambda$, $P_\lambda$ is entire of order exactly one, $P_\lambda(1-s)=P_\lambda(s)$, $P_\lambda(\bar s)=\overline{P_\lambda(s)}$, $P_\lambda(0)=P_\lambda(1)=1$, and its zeros, with multiplicity, are exactly
\begin{equation}\label{eq:family-divisor}
 \frac12\pm i\gamma_k(\lambda),\qquad k\ge1 .
\end{equation}
For $0<\lambda\le1$ the ordinates $\gamma_k(\lambda)$ are strictly increasing in $k$ and all zeros of $P_\lambda$ are simple; at $\lambda=0$ the multiplicities are those of the ordinate list $(\gamma_k)$.  For every $j\ge0$, $P_\lambda^{(j)}\to P_{\lambda_0}^{(j)}$ locally uniformly as $\lambda\to\lambda_0$.  In particular no member $\zeta_\lambda$ of the family, $\zeta_{Z1}=\zeta_0$ included, can have a zero off the critical line other than $-2,-4,\dots$.
\end{theorem}

\begin{proof}
By Lemma~\CGref{lem:bounded-displacement}, $|\gamma_k(\lambda)-\gamma_k|\le C_\eta$, and $\gamma_k\asymp k/\log k$; so $\gamma_k(\lambda)\ge\tfrac12\gamma_k$ for $k\ge k_0$, uniformly in $\lambda$.  On a compact $K\subset\C$,
\begin{equation}\label{eq:lambda-normal}
 \sum_{k\ge1}\sup_{s\in K,\ 0\le\lambda\le1}\left|\frac{q(s)}{\tfrac14+\gamma_k(\lambda)^2}\right|<\infty,
\end{equation}
which gives joint normal convergence and, by Cauchy's formula, convergence of all $s$-derivatives.  For $|s|\le R_{\rm fam}$, uniformly in $\lambda$,
\begin{equation}\label{eq:order-one}
 \log\max_{|s|\le R_{\rm fam}}|P_\lambda(s)|\le\sum_{k\ge1}\log\Bigl(1+\frac{R_{\rm fam}^2+R_{\rm fam}}{\tfrac14+\gamma_k(\lambda)^2}\Bigr)=O(R_{\rm fam}\log R_{\rm fam}),
\end{equation}
so the order is at most one; the zero count is $\asymp T\log T$, so the exponent of convergence is one and the order is exactly one.  Symmetries, normalization and divisor follow factor by factor as in Theorem~\CGref{thm:canonical}.  For $0<\lambda\le1$,
\begin{equation}\label{eq:ordinate-gaps}
 \gamma_{k+1}(\lambda)-\gamma_k(\lambda)=(1-\lambda)(\gamma_{k+1}-\gamma_k)+\lambda(\tau_{k+1}-\tau_k)\ge\lambda(\tau_{k+1}-\tau_k)>0,
\end{equation}
so distinct factors vanish at distinct points and each zero is simple; at $\lambda=0$ equal ordinates give equal factors.
\end{proof}

\begin{remark}[All differences lie in $\gamma_k-\tau_k$]\label{rem:only-gamma-tau}
For $\lambda,\lambda'\in[0,1]$ and $s$ off the divisors,
\begin{equation}\label{eq:ratio-members}
 \frac{\zeta_\lambda(s)}{\zeta_{\lambda'}(s)}=\prod_{k\ge1}\frac{\bigl(\tfrac14+(\gamma_k+\lambda'\eta_k)^2\bigr)\bigl(\tfrac14+(\gamma_k+\lambda\eta_k)^2+q(s)\bigr)}{\bigl(\tfrac14+(\gamma_k+\lambda\eta_k)^2\bigr)\bigl(\tfrac14+(\gamma_k+\lambda'\eta_k)^2+q(s)\bigr)} .
\end{equation}
The gamma factor cancels; the members differ only through $\eta_k=\tau_k-\gamma_k$, and all of them coincide if every $\eta_k$ vanishes.
\end{remark}

\begin{theorem}[Infinitesimal generator]\label{thm:generator}
Put $A_{k,\lambda}:=\tfrac14+\gamma_k(\lambda)^2$.  Off the divisor of $P_\lambda$,
\begin{equation}\label{eq:generator}
 \mathcal V_\lambda(s):=\partial_\lambda\log P_\lambda(s)=-2q(s)\sum_{k\ge1}\frac{\gamma_k(\lambda)\,\eta_k}{A_{k,\lambda}\bigl(A_{k,\lambda}+q(s)\bigr)},
\end{equation}
the series converging locally uniformly on compact sets separated from the divisor.  Since $G$ does not depend on $\lambda$,
\begin{equation}\label{eq:generator-zeta}
 \partial_\lambda\log\zeta_\lambda=\mathcal V_\lambda,\qquad \partial_\lambda\log F_\lambda=-\mathcal V_\lambda .
\end{equation}
\end{theorem}

\begin{proof}
Since $\partial_\lambda A_{k,\lambda}=2\gamma_k(\lambda)\eta_k$,
\begin{equation}\label{eq:mode-derivative}
 \partial_\lambda\log\Bigl(1+\frac q{A_{k,\lambda}}\Bigr)=-\frac{2q\,\gamma_k(\lambda)\eta_k}{A_{k,\lambda}(A_{k,\lambda}+q)} .
\end{equation}
On compact sets separated from the divisor the summand is $O_K(|\eta_k|\gamma_k^{-3})$, and $\sum_k\gamma_k^{-3}<\infty$, so termwise differentiation is justified.
\end{proof}

\begin{proposition}[Integrated multiplier]\label{prop:integrated-multiplier}
On every component of $\C\setminus\bigcup_{\lambda'\in[0,1]}\{\tfrac12\pm i\gamma_k(\lambda'):k\ge1\}$,
\begin{equation}\label{eq:integrated-zeta}
 \zeta_\lambda(s)=\zeta_{Z1}(s)\exp\Bigl(\int_0^\lambda\mathcal V_{\lambda'}(s)\,d\lambda'\Bigr),
\end{equation}
\begin{equation}\label{eq:integrated-inverse}
 F_\lambda(s)=F_{Z1}(s)\exp\Bigl(-\int_0^\lambda\mathcal V_{\lambda'}(s)\,d\lambda'\Bigr),
\end{equation}
and the total multiplier is
\begin{equation}\label{eq:transport-multiplier}
 \Pi_{\rm tr}(s):=\exp\Bigl(-\int_0^1\mathcal V_{\lambda'}(s)\,d\lambda'\Bigr)=\frac{P_{Z1}(s)}{P_{\RVM}^{\crit}(s)},\qquad \zeta_{Z1}=\zeta_{\RVM}\,\Pi_{\rm tr} .
\end{equation}
The identities extend meromorphically to $\C$ through the product form.
\end{proposition}

\begin{proof}
\CGeqref{eq:integrated-zeta} and \CGeqref{eq:integrated-inverse} are the fundamental theorem of calculus applied to \CGeqref{eq:generator-zeta}.  For \CGeqref{eq:transport-multiplier}, integrate each mode with $r=\gamma_k(\lambda')$, $dr=\eta_k\,d\lambda'$:
\begin{equation}\label{eq:mode-integral}
 \int_0^1\partial_{\lambda'}\log\Bigl(1+\frac q{A_{k,\lambda'}}\Bigr)d\lambda'\equiv\log\frac{1+q/(\tfrac14+\tau_k^2)}{1+q/(\tfrac14+\gamma_k^2)}\pmod{2\pi i}.
\end{equation}
Normal convergence permits summation, and exponentiation removes the branch ambiguity.
\end{proof}

\begin{corollary}[Expansion at the endpoint]\label{cor:endpoint-expansion}
Let $K\subset\C$ be compact, disjoint from the zeros of $P_{Z1}$ and from $s=1$ (equivalently one may state the expansion for the product ratio $P_\lambda/P_{Z1}$ without this exclusion).  Uniformly for $s\in K$, as $\lambda\downarrow0$,
\begin{equation}\label{eq:endpoint-zeta-expansion}
 \zeta_\lambda(s)=\zeta_{Z1}(s)\bigl(1+\lambda\mathcal V_0(s)+O_K(\lambda^2)\bigr),
\end{equation}
\begin{equation}\label{eq:endpoint-inverse-expansion}
 F_\lambda(s)=F_{Z1}(s)\bigl(1-\lambda\mathcal V_0(s)+O_K(\lambda^2)\bigr).
\end{equation}
\end{corollary}

\begin{proof}
By Theorem~\CGref{thm:path} the zeros of $P_\lambda$ converge to those of $P_{Z1}$, so $K$ avoids the moving divisor for small $\lambda$.  The series for $\mathcal V_\lambda$ and its $\lambda$-derivative converge normally on $K$, uniformly for $\lambda$ near $0$: the summands of \CGeqref{eq:generator} are $O_K(|\eta_k|\gamma_k^{-3})$, and, since $\partial_\lambda\gamma_k(\lambda)=\eta_k$ and $\partial_\lambda A_{k,\lambda}=2\gamma_k(\lambda)\eta_k$,
\[
 \partial_\lambda\mathcal V_\lambda(s)=-2q(s)\sum_{k\ge1}\eta_k^2\Bigl(\frac1{A_{k,\lambda}(A_{k,\lambda}+q(s))}-\frac{2\gamma_k(\lambda)^2\,(2A_{k,\lambda}+q(s))}{A_{k,\lambda}^2(A_{k,\lambda}+q(s))^2}\Bigr),
\]
whose summands are $O_K(\eta_k^2\gamma_k^{-4})$.  Hence $\mathcal V_\lambda=\mathcal V_0+O_K(\lambda)$, and \CGeqref{eq:integrated-zeta} gives $\log(\zeta_\lambda/\zeta_{Z1})=\lambda\mathcal V_0+O_K(\lambda^2)$.
\end{proof}

\begin{proposition}[Residues along the family]\label{prop:residues}
Suppose $\gamma_k$ is a simple ordinate (it occurs once in the list).  Put $\rho_{k,\lambda}=\tfrac12+i\gamma_k(\lambda)$ and, for $0<\lambda\le1$,
\begin{equation}\label{eq:residue-transport}
 \mathfrak r_{k,\lambda}:=\Res_{s=\rho_{k,\lambda}}\frac1{s\,\zeta_\lambda(s)}=\frac{G(\rho_{k,\lambda})}{\rho_{k,\lambda}P_\lambda'(\rho_{k,\lambda})}.
\end{equation}
Then
\begin{equation}\label{eq:residue-limit}
 \mathfrak r_{k,\lambda}\longrightarrow\frac{G(\tfrac12+i\gamma_k)}{(\tfrac12+i\gamma_k)P_{Z1}'(\tfrac12+i\gamma_k)}\qquad(\lambda\downarrow0).
\end{equation}
\end{proposition}

\begin{proof}
Theorem~\CGref{thm:path} gives convergence of the zero locations and local uniform convergence of $P_\lambda'$; the limit denominator is nonzero because the ordinate $\gamma_k$ occurs once in the list, so the corresponding endpoint zero $\frac12+i\gamma_k$ of $P_{Z1}$ is simple; and $G$ is continuous.
\end{proof}

\begin{remark}[Multiplicity]\label{rem:multiplicity}
For $0<\lambda\le1$ every zero of $\zeta_\lambda$ is simple.  At $\lambda=0$ a repeated ordinate of $\zeta$ appears in $\zeta_{Z1}$ with the same multiplicity, and the residues of Proposition~\CGref{prop:residues} are replaced by the full principal part.  Multiplicity plays no role in Step~3: Theorem~\CGref{thm:step3} uses only that $\zeta$ does not vanish in $\Re s>\tfrac12$.
\end{remark}

\section{The endpoint \texorpdfstring{$\zeta_{Z1}$}{zeta\_Z1}, Hamburger's theorem, and Step 3}
\label{sec:Z1-step3}

\subsection{The analytic class of the endpoint}
\label{subsec:Z1-class}

\begin{proposition}[Analytic class of $\zeta_{Z1}$]\label{prop:Z1-class}
The endpoint
\begin{equation}\label{eq:Z1-product}
 P_{Z1}(s)=\prod_{k\ge1}\left(1+\frac{q(s)}{\tfrac14+\gamma_k^2}\right),\qquad \zeta_{Z1}(s)=\frac{P_{Z1}(s)}{G(s)},
\end{equation}
has the following properties.
\begin{enumerate}[label=\textup{(\roman*)}]
\item $P_{Z1}$ is entire of order exactly one, $P_{Z1}(1-s)=P_{Z1}(s)$, $P_{Z1}(0)=P_{Z1}(1)=1$;
\item the nontrivial zeros of $\zeta_{Z1}$ are exactly
\begin{equation}\label{eq:Z1-divisor}
 \frac12\pm i\gamma_k,\qquad k\ge1,
\end{equation}
with the multiplicities of the ordinate list; its other zeros are $s=-2,-4,\dots$.  It is impossible for $\zeta_{Z1}$ to vanish at any point off the critical line other than $-2,-4,\dots$: every factor of $P_{Z1}$ vanishes only at $\tfrac12\pm i\gamma_k$;
\item $\zeta_{Z1}$ has a simple pole at $s=1$ with residue $1$;
\item $\zeta_{Z1}$ satisfies the functional equation of $\zeta$:
\begin{equation}\label{eq:Z1-FE}
 \pi^{-s/2}\Gamma(s/2)\zeta_{Z1}(s)=\pi^{-(1-s)/2}\Gamma((1-s)/2)\zeta_{Z1}(1-s);
\end{equation}
\item $(s-1)\zeta_{Z1}(s)$ is entire of order exactly one.
\end{enumerate}
\end{proposition}

\begin{proof}
(i) and the divisor in (ii) are the case $\lambda=0$ of Theorem~\CGref{thm:path}; the trivial zeros and the absence of other zeros come from $1/G$ as in Theorem~\CGref{thm:canonical}.  For (iii), $P_{Z1}(1)=1$ and
\begin{equation}\label{eq:G-near-one}
 G(s)=(s-1)+O\bigl((s-1)^2\bigr)\qquad(s\to1).
\end{equation}
For (iv),
\begin{equation}\label{eq:Z1-FE-proof}
 \pi^{-s/2}\Gamma(s/2)\zeta_{Z1}(s)=\frac{P_{Z1}(s)}{q(s)},
\end{equation}
and both $P_{Z1}$ and $q$ are invariant under $s\mapsto1-s$.  For (v),
\begin{equation}\label{eq:Z1-entire-removal}
 (s-1)\zeta_{Z1}(s)=P_{Z1}(s)\,\frac{\pi^{s/2}}{s\,\Gamma(s/2)},
\end{equation}
where $s\Gamma(s/2)\to2$ at $s=0$ and $1/\Gamma$ is entire of order one; the product has order at most one, and the zero count $\asymp T\log T$ gives order at least one.
\end{proof}

Define the ratio
\begin{equation}\label{eq:Z1-defect-ratio}
 \mathscr R_{Z1}(s):=\frac{\zeta_{Z1}(s)}{\zeta(s)}=\frac{P_{Z1}(s)}{P_\zeta(s)},
\end{equation}
a meromorphic function; $\mathscr R_{Z1}\equiv1$ is the identity $\zeta_{Z1}=\zeta$.

\subsection{Hamburger's theorem at the endpoint}
\label{subsec:Hamburger}

\begin{theorem}[Hamburger closure]\label{thm:Hamburger-Z1}
Consider the condition
\begin{equation}\label{eq:Z1-D}
 \text{(D)}\qquad\zeta_{Z1}(s)=\sum_{n\ge1}\frac{a_n}{n^s}\qquad\text{absolutely for }\Re s>1 .
\end{equation}
Then
\begin{equation}\label{eq:Hamburger-equivalence}
 (\mathrm D)\quad\Longleftrightarrow\quad\mathscr R_{Z1}\equiv1\quad\Longleftrightarrow\quad\zeta_{Z1}=\zeta\quad\Longleftrightarrow\quad\mathrm{RH}.
\end{equation}
\end{theorem}

\begin{proof}
Assume (D).  Then $\zeta_{Z1}(1-s)=\sum_na_nn^{-(1-s)}$ converges absolutely for $\Re s<0$, $(s-1)\zeta_{Z1}$ is entire of finite order (Proposition~\CGref{prop:Z1-class}(v)), and \CGeqref{eq:Z1-FE} holds.  These are the hypotheses of Hamburger's theorem \cite{Hamburger1921}, \cite[\S2.13]{Titchmarsh1986}, which gives
\begin{equation}\label{eq:Hamburger-conclusion}
 \zeta_{Z1}=C_H\,\zeta
\end{equation}
for a constant $C_H$.  Both functions have residue $1$ at $s=1$, so $C_H=1$.  If $\zeta_{Z1}=\zeta$, the nontrivial zeros of $\zeta$ are those of $\zeta_{Z1}$, namely \CGeqref{eq:Z1-divisor}, which is RH.  If RH holds, every $\rho$ with $\Im\rho>0$ is $\tfrac12+i\gamma_k$, so $\rho(1-\rho)=\tfrac14+\gamma_k^2$, and Theorem~\CGref{thm:Hadamard} gives
\begin{equation}\label{eq:RH-product}
 P_\zeta(s)=\prod_{k\ge1}\Bigl(1+\frac{q(s)}{\tfrac14+\gamma_k^2}\Bigr)=P_{Z1}(s)
\end{equation}
exactly, hence $\zeta=\zeta_{Z1}$.  Finally $\zeta=\sum_nn^{-s}$ converges absolutely for $\Re s>1$, so $\zeta_{Z1}=\zeta$ gives (D).
\end{proof}

\subsection{Condition (D) and the integer-weighted zeta}
\label{subsec:D-multiplier}

By Corollary~\CGref{cor:ZRVM-Dirichlet}, $\zeta_{\Z,\RVM}$ is an absolutely convergent Dirichlet series in $\Re s>1$ with first coefficient $1$, unconditionally.  Define the multiplier
\begin{equation}\label{eq:K-multiplier}
 \mathcal K(s):=\frac{\zeta_{Z1}(s)}{\zeta_{\Z,\RVM}(s)} .
\end{equation}

\begin{theorem}[Reduction of (D) to the multiplier]\label{thm:D-multiplier}
\begin{enumerate}[label=\textup{(\roman*)}]
\item If $\mathcal K$ has an expansion $\mathcal K(s)=\sum_{n\ge1}k_nn^{-s}$ that converges absolutely in $\Re s>1$, then \CGeqref{eq:Z1-D} holds.
\item If \CGeqref{eq:Z1-D} holds, then $\mathcal K(s)=\sum_{n\ge1}k_nn^{-s}$ absolutely on the whole half-plane $\Re s>1$.
\end{enumerate}
Thus (D) holds if and only if $\mathcal K$ is an absolutely convergent Dirichlet series in $\Re s>1$.
Hence, by Theorem~\CGref{thm:Hamburger-Z1}, an absolutely convergent Dirichlet expansion of $\mathcal K$ in $\Re s>1$ gives $\zeta_{Z1}=\zeta$, RH, and, by Theorem~\CGref{thm:step3}, $M(x)=O_\eps(x^{1/2+\eps})$.
\end{theorem}

\begin{proof}
(i) $\zeta_{Z1}=\zeta_{\Z,\RVM}\,\mathcal K$ is the Dirichlet convolution of two absolutely convergent series in $\Re s>1$, and $\sum_N|(a^{\Z}*k)_N|N^{-\sigma}\le\sum_N|a_N^{\Z}|N^{-\sigma}\sum_n|k_n|n^{-\sigma}<\infty$.  (ii) By \CGeqref{eq:ZRVM-inverse-Dirichlet}, $1/\zeta_{\Z,\RVM}$ is an absolutely convergent Dirichlet series in $\Re s>1$, and so is its product $\mathcal K$ with \CGeqref{eq:Z1-D}.
\end{proof}

By \CGeqref{eq:transport-multiplier} and Corollary~\CGref{cor:weighted-factorization},
\begin{equation}\label{eq:K-explicit}
 \mathcal K(s)=\Pi_{\rm tr}(s)\,e^{-h_{\RVM}(s)}=\frac{P_{Z1}(s)}{P_{\RVM}^{\crit}(s)}\,e^{-h_{\RVM}(s)} ,
\end{equation}
an explicit function of the displacements $\gamma_k-\tau_k$ and of the cell masses $c_n^{\RVM}$.  Condition (D) is thereby located in a single object built from the two constructions of this paper: the $\lambda$-family and the integer weights.

\begin{remark}[Continuous-to-atomic meaning of (D)]\label{rem:D-atomic}
The prime-power count $J_{Z1}^{\rm M}$ of $\zeta_{Z1}$ (Theorem~\CGref{thm:mellin-counts}) is an $L^2_{\rm loc}$ count built from the zeros, and (D) makes its distributional derivative the atomic measure $\sum_{p,m}m^{-1}\delta_{p^m}$.  The final step of the method is therefore the passage of $\zeta_{Z1}$ into the discrete Dirichlet category.  Theorem~\CGref{thm:cell-form} states it exactly: (D) holds if and only if $J_{Z1}^{\rm M}$ is constant between consecutive integers and its jumps obey $\sum_{n\le X}|\mathsf b_n|=O_\eps(X^{1+\eps})$, and Corollary~\CGref{cor:product-identity} reduces it to the single identity \CGeqref{eq:product-identity} between the displacements $\gamma_k-\tau_k$ and $\varpi_{\RVM}$.
\end{remark}

\subsection{Step 3: the Mertens bound from \texorpdfstring{$\zeta_{Z1}=\zeta$}{zeta\_Z1 = zeta}}
\label{subsec:step3}

\begin{lemma}[Vertical growth of $\zeta$]\label{lem:zeta-growth}
For $\tfrac14\le\sigma\le4$ and $|t|\ge2$,
\begin{equation}\label{eq:zeta-growth}
 |\zeta(\sigma+it)|\le15|t| .
\end{equation}
\end{lemma}

\begin{proof}
For $\sigma>0$ \cite[(2.1.5)]{Titchmarsh1986},
\begin{equation}\label{eq:zeta-integral-rep}
 \zeta(s)=\frac s{s-1}-s\int_1^\infty\{y\}y^{-s-1}\,dy .
\end{equation}
For $|t|\ge2$ and $\sigma\le4$ one has $|s|\le3|t|$, $|s/(s-1)|\le|s|/|t|\le3$, and the integral is at most $1/\sigma\le4$; hence $|\zeta(s)|\le3+12|t|\le15|t|$.
\end{proof}

\begin{lemma}[Logarithm of $\zeta$ off the critical line]\label{lem:log-zeta}
Assume $\zeta(s)\ne0$ for $\Re s>\tfrac12$.  Let $0<\delta\le\tfrac14$ and
\begin{equation}\label{eq:three-circles-exponent}
 \omega(\delta):=\frac{\log(3-2\delta)}{\log(3-\delta)}<1 .
\end{equation}
Then, with the branch of $\log\zeta$ equal to $\sum_{n\ge2}\Lambda(n)n^{-s}/\log n$ on $\Re s>1$,
\begin{equation}\label{eq:log-zeta-bound}
 |\log\zeta(\sigma+it)|\le\Bigl(\frac{48}\delta\log|t|\Bigr)^{\omega(\delta)}\qquad(\sigma\ge\tfrac12+\delta,\ |t|\ge4).
\end{equation}
Consequently, for every $\eps>0$ there is $t_2=t_2(\delta,\eps)$ with
\begin{equation}\label{eq:inv-zeta-bound}
 \Bigl|\frac1{\zeta(\sigma+it)}\Bigr|\le|t|^{\eps}\qquad(\sigma\ge\tfrac12+\delta,\ |t|\ge t_2).
\end{equation}
\end{lemma}

\begin{proof}
Fix $t$ with $|t|\ge4$ and put $s_0=2+it$, $R=\tfrac32-\tfrac\delta4$, $r_3=\tfrac32-\tfrac\delta2$.  The closed disc $|s-s_0|\le R$ lies in $\tfrac12+\tfrac\delta4\le\Re s\le\tfrac72$, $|\Im s-t|\le\tfrac32$; it avoids $s=1$ and, by hypothesis, every zero of $\zeta$.  Hence $\log\zeta$ has a holomorphic branch on a neighbourhood of this disc, chosen to agree with the Dirichlet-series branch on the connected overlap with $\Re s>1$.  Such branches are defined component by component on zero-free regions; no global branch on $\Re s>\tfrac12$ is asserted.  On the disc $|\Im s|\ge|t|/2\ge2$ and $|\Im s|\le2|t|$, so Lemma~\CGref{lem:zeta-growth} gives $\Re\log\zeta=\log|\zeta|\le\log(30|t|)$ there, while $|\log\zeta(s_0)|\le\log\zeta(2)$.  The Borel--Carath\'eodory lemma \cite[\S5.5]{Titchmarsh1939}, applied with radii $r_3<R$,
\begin{equation}\label{eq:borel-caratheodory}
 \max_{|s-s_0|\le r_3}|\log\zeta(s)|\le\frac{2r_3}{R-r_3}\max_{|s-s_0|=R}\Re\log\zeta(s)+\frac{R+r_3}{R-r_3}|\log\zeta(s_0)|,
\end{equation}
together with $R-r_3=\delta/4$, $2r_3\le3$, $R+r_3\le3$ and $\log30+\log\zeta(2)<4\le3\log|t|$, gives
\begin{equation}\label{eq:BC-bound}
 \max_{|s-s_0|\le r_3}|\log\zeta(s)|\le\frac{12}\delta\bigl(\log(30|t|)+\log\zeta(2)\bigr)\le\frac{48}\delta\log|t| .
\end{equation}
Now let $\tfrac12+\delta\le\sigma<\tfrac32$ and set $r_1=\tfrac12$, $r_2=2-\sigma\in(\tfrac12,\tfrac32-\delta]$.  On $|s-s_0|=r_1$ one has $\Re s\ge\tfrac32$, so $|\log\zeta(s)|\le\log\zeta(\tfrac32)<1$.  Hadamard's three-circle theorem \cite[\S5.3]{Titchmarsh1939} gives, with $M_j=\max_{|s-s_0|=r_j}|\log\zeta|$,
\begin{equation}\label{eq:three-circles}
 M_2\le M_1^{1-\omega}M_3^{\omega},\qquad \omega=\frac{\log(r_2/r_1)}{\log(r_3/r_1)}\le\frac{\log(3-2\delta)}{\log(3-\delta)}=\omega(\delta).
\end{equation}
Since $M_1<1$ and $M_3\le\frac{48}\delta\log|t|$ with $\frac{48}\delta\log|t|>1$, \CGeqref{eq:three-circles} gives \CGeqref{eq:log-zeta-bound} at $s=\sigma+it$, which lies on the circle $|s-s_0|=r_2$ because $|(\sigma+it)-(2+it)|=2-\sigma=r_2$.  For $\sigma\ge\tfrac32$, $|\log\zeta|\le\log\zeta(\tfrac32)<1$.  Finally $|1/\zeta|=\exp(-\Re\log\zeta)\le\exp|\log\zeta|$, and $(\frac{48}\delta\log|t|)^{\omega(\delta)}\le\eps\log|t|$ for $|t|\ge t_2(\delta,\eps)$ because $\omega(\delta)<1$.
\end{proof}

\begin{lemma}[Perron kernel]\label{lem:perron-kernel}
For $y>0$, $y\ne1$, $c>0$, $T>0$ \cite[Ch.~17]{Davenport1980},
\begin{equation}\label{eq:perron-kernel}
 \left|\frac1{2\pi i}\int_{c-iT}^{c+iT}\frac{y^s}{s}\,ds-\mathbf 1_{\{y>1\}}\right|\le y^c\min\Bigl(1,\frac1{T|\log y|}\Bigr).
\end{equation}
\end{lemma}

\begin{theorem}[Step 3]\label{thm:step3}
If $\zeta_{Z1}=\zeta$, then for every $\eps>0$
\begin{equation}\label{eq:step3-bound}
 M(x)=O_\eps\bigl(x^{1/2+\eps}\bigr)\qquad(x\ge1).
\end{equation}
\end{theorem}

\begin{proof}
\emph{Non-vanishing.}  By Proposition~\CGref{prop:Z1-class}(ii), $\zeta_{Z1}$ has no zeros in $\Re s>\tfrac12$; since $\zeta=\zeta_{Z1}$, neither has $\zeta$, and Lemma~\CGref{lem:log-zeta} applies.

\emph{Truncation.}  Let $x=N+\tfrac12$ with $N\ge3$ an integer, $c=1+1/\log x$ and $T=x$.  Applying \CGeqref{eq:perron-kernel} with $y=x/n$ to each term of $\sum_n\mu(n)n^{-s}=1/\zeta(s)$, which converges absolutely on $\Re s=c$, gives
\begin{equation}\label{eq:perron-truncated}
 M(x)=\frac1{2\pi i}\int_{c-iT}^{c+iT}\frac{x^s}{s\,\zeta(s)}\,ds+O\Bigl(\sum_{n\ge1}\Bigl(\frac xn\Bigr)^c\min\Bigl(1,\frac1{T|\log(x/n)|}\Bigr)\Bigr).
\end{equation}
For $n\le x/2$ or $n\ge2x$, $|\log(x/n)|\ge\log2$, and these terms contribute at most $x^c\zeta(c)/(T\log2)\le e\,x(1+\log x)/(T\log2)$.  For $x/2<n<2x$, $(x/n)^c\le4$ and $|\log(x/n)|\ge|n-x|/(2x)$; as $|n-x|\in\{j-\tfrac12:j\ge1\}$, each value taken at most twice, these terms contribute at most $8\sum_{1\le j\le2x}\min(1,4x/(T(2j-1)))\ll1+(x/T)\log x$.  Hence
\begin{equation}\label{eq:perron-error}
 M(x)=\frac1{2\pi i}\int_{c-iT}^{c+iT}\frac{x^s}{s\,\zeta(s)}\,ds+O(\log x).
\end{equation}
\emph{Contour shift.}  Fix $0<\delta\le\tfrac14$.  The integrand is holomorphic on the closed rectangle with vertices $c\pm iT$, $\tfrac12+\delta\pm iT$: $\zeta$ has no zeros there, $s=1$ is a zero of $1/\zeta$, and $s=0$ lies outside.  By Cauchy's theorem
\begin{equation}\label{eq:contour-shift}
 \int_{c-iT}^{c+iT}=\int_{\frac12+\delta-iT}^{\frac12+\delta+iT}+\int_{\frac12+\delta+iT}^{c+iT}-\int_{\frac12+\delta-iT}^{c-iT}.
\end{equation}
For $x\ge t_2(\delta,\eps)$, \CGeqref{eq:inv-zeta-bound} on $|t|=T=x$ bounds the horizontal integrals by
\begin{equation}\label{eq:horizontal-bound}
 \frac1{2\pi}\int_{\frac12+\delta}^{c}\frac{x^\sigma T^{\eps}}{T}\,d\sigma\le x^cT^{\eps-1}=e\,x^{\eps}.
\end{equation}
On the left side, $1/\zeta$ is continuous on the compact segment $\{\tfrac12+\delta+it:|t|\le t_2\}$, bounded by a constant $C(\delta,\eps)$, and \CGeqref{eq:inv-zeta-bound} holds for $|t|\ge t_2$; hence
\begin{equation}\label{eq:left-bound}
 \frac1{2\pi}\int_{-T}^{T}\frac{x^{1/2+\delta}}{|\tfrac12+\delta+it|\,|\zeta(\tfrac12+\delta+it)|}\,dt\le x^{1/2+\delta}\Bigl(C'(\delta,\eps)+\int_{t_2}^{T}t^{\eps-1}dt\Bigr)\ll_{\delta,\eps}x^{1/2+\delta+\eps}.
\end{equation}
\emph{Conclusion.}  \CGeqref{eq:perron-error}--\CGeqref{eq:left-bound} give $M(x)\ll_{\delta,\eps}x^{1/2+\delta+\eps}$ for $x=N+\tfrac12$; it suffices to prove the assertion for $0<\eps_0\le\tfrac12$, since larger exponents follow from smaller ones, and for such $\eps_0$ the choice $\delta=\eps=\eps_0/2$ satisfies $\delta\le\tfrac14$ and gives $O_{\eps_0}(x^{1/2+\eps_0})$.  For arbitrary $x\ge3$, $M(x)=M(\lfloor x\rfloor+\tfrac12)$ and $\lfloor x\rfloor+\tfrac12\le2x$; for $1\le x<3$, $|M(x)|\le2$.
\end{proof}

\begin{theorem}[Transfer of the synthetic Mertens bound to $M$]\label{thm:transfer}
Assume $\zeta_{Z1}=\zeta$ and let $0<\delta\le\tfrac14$, $\eps>0$.  Then:
\begin{enumerate}[label=\textup{(\roman*)}]
\item on $\Re s>\tfrac12$,
\begin{equation}\label{eq:transfer-factorization}
 \frac1{\zeta(s)}=F_{\RVM}(s)\,\Pi_{\rm tr}(s)^{-1},
 \qquad
 \Pi_{\rm tr}=\frac{P_{Z1}}{P_{\RVM}^{\crit}}=\frac{\zeta}{\zeta_{\RVM}},
\end{equation}
where $\Pi_{\rm tr}$ is holomorphic and zero-free in $\Re s>\tfrac12$ and $\Pi_{\rm tr}(1)=1$;
\item there is $t_3=t_3(\delta,\eps)$ such that, for $\tfrac12+\delta\le\sigma\le2$ and $|t|\ge t_3$,
\begin{equation}\label{eq:transfer-multiplier-bound}
 |F_{\RVM}(\sigma+it)|\le4L,\qquad |\Pi_{\rm tr}(\sigma+it)^{\pm1}|\le|t|^{\eps};
\end{equation}
\item the following classical consequence of RH is recovered at the endpoint:
\begin{equation}\label{eq:transfer-conclusion}
 M(x)=O_\eps\bigl(x^{1/2+\eps}\bigr).
\end{equation}
\end{enumerate}
\end{theorem}

\begin{proof}
(i) By \CGeqref{eq:transport-multiplier}, $\zeta_{Z1}=\zeta_{\RVM}\Pi_{\rm tr}$; with $\zeta_{Z1}=\zeta$ this is \CGeqref{eq:transfer-factorization}.  In $\Re s>\tfrac12$, $\zeta=\zeta_{Z1}$ has no zeros (Proposition~\CGref{prop:Z1-class}(ii)) and $\zeta_{\RVM}$ has none (Theorem~\CGref{thm:canonical}); both have a simple pole of residue $1$ at $s=1$, so $\Pi_{\rm tr}$ is holomorphic there with $\Pi_{\rm tr}(1)=1$, and it is zero-free in $\Re s>\tfrac12$.

(ii) By Theorem~\CGref{thm:transition}(a) with $B=\tfrac32$, for $\delta\le\sigma'\le\tfrac32$ and $|t|\ge t_0$ we have $\zeta_{\RVM}=\varpi_{\RVM}(1+w)(1+\mathcal E_0)$ with $|w|\le C|t|^{-\delta}\le\tfrac12$ and $|\mathcal E_0|\le\tfrac12$ for $|t|$ large; hence $|F_{\RVM}|\le4L$ and
\begin{equation}\label{eq:log-zetaRVM-bounded}
 \log\zeta_{\RVM}=\log\varpi_{\RVM}+\Log(1+w)+\Log(1+\mathcal E_0)=O(1).
\end{equation}
By Lemma~\CGref{lem:log-zeta}, which applies because $\zeta=\zeta_{Z1}$ has no zeros in $\Re s>\tfrac12$, $|\log\zeta(\sigma+it)|\le(\frac{48}\delta\log|t|)^{\omega(\delta)}$ with $\omega(\delta)<1$.  Therefore
\begin{equation}\label{eq:log-Pi-tr}
 |\log\Pi_{\rm tr}(\sigma+it)|=|\log\zeta-\log\zeta_{\RVM}|\le\Bigl(\frac{48}\delta\log|t|\Bigr)^{\omega(\delta)}+O(1)\le\eps\log|t|
\end{equation}
for $|t|\ge t_3(\delta,\eps)$, which gives \CGeqref{eq:transfer-multiplier-bound}.

(iii) The required estimate for $1/\zeta$ already follows from the zero-free hypothesis $\zeta=\zeta_{Z1}$ through Lemma~\CGref{lem:log-zeta} and Theorem~\CGref{thm:step3}.  The factorization \CGeqref{eq:transfer-factorization} is therefore a reformulation of the same estimate, not an independent transfer from the synthetic Mertens bound.  By (i)--(ii), $\Pi_{\rm tr}^{-1}$ is holomorphic in $\Re s>\tfrac12$ and grows at most like $|t|^{\eps}$.  Hence $|1/\zeta(\sigma+it)|\le4L|t|^{\eps}$ for $\tfrac12+\delta\le\sigma\le2$, $|t|\ge t_3$, and the truncated Perron formula, contour shift and estimates \CGeqref{eq:perron-error}--\CGeqref{eq:left-bound} of Theorem~\CGref{thm:step3}, applied with this bound, give $M(x)\ll_{\delta,\eps}x^{1/2+\delta+\eps}$.  Choosing $\delta=\eps=\eps_0/2$ gives \CGeqref{eq:transfer-conclusion}.
\end{proof}

\begin{remark}[The transfer in Mellin form]\label{rem:transfer-mellin}
Formally, \CGeqref{eq:transfer-factorization} reads $M=(\Pi_{\rm tr}^{-1})^{\vee}\Mell M_{\RVM}^{\crit}$: the Mertens function is the synthetic Mertens function convolved with the inverse Mellin transform of the inverse transport multiplier.  The synthetic factor $F_{\RVM}$ contributes the exponent $\tfrac12$ and is bounded (Theorem~\CGref{thm:transition}); the multiplier contributes at most $|t|^{\eps}$.  All the information of the $\lambda$-family enters through $\Pi_{\rm tr}=P_{Z1}/P_{\RVM}^{\crit}$, that is, through the displacements $\gamma_k-\tau_k$.
\end{remark}

\begin{corollary}[Step 3 through (D)]\label{cor:step3-D}
Condition \CGeqref{eq:Z1-D} implies $M(x)=O_\eps(x^{1/2+\eps})$ for every $\eps>0$.
\end{corollary}

\begin{proof}
Theorem~\CGref{thm:Hamburger-Z1} and Theorem~\CGref{thm:step3}.
\end{proof}

\begin{remark}[Status of Step 3]\label{rem:Step3-status}
The implication $\zeta_{Z1}=\zeta\Rightarrow M(x)=O_\eps(x^{1/2+\eps})$ is the classical consequence RH $\Rightarrow M(x)=O_\eps(x^{1/2+\eps})$ (see \cite[\S14.25]{Titchmarsh1986}) rewritten and reproved in the endpoint notation.  It requires nothing beyond the identity: no separate statement about the zeros of $\zeta$ enters, because $\zeta=\zeta_{Z1}$ already carries the divisor \CGeqref{eq:Z1-divisor}.  The identity $\zeta_{Z1}=\zeta$ itself, equivalently (D), equivalently an absolutely convergent Dirichlet expansion of the multiplier $\mathcal K$ of \CGeqref{eq:K-multiplier} in $\Re s>1$ (Theorem~\CGref{thm:D-multiplier}), is not proved in this paper.  It is the open step of the method.
\end{remark}

\subsection{Reciprocal endpoint and finite-height Perron integrals}
\label{subsec:reciprocal}

Put
\begin{equation}\label{eq:FZ1}
 F_{Z1}(s):=\frac1{\zeta_{Z1}(s)}=F_{\RVM}(s)\,\Pi_{\rm tr}(s)^{-1},
\end{equation}
the second identity by \CGeqref{eq:transport-multiplier}.  For $c>1$ and finite $U$ define
\begin{equation}\label{eq:truncated-Perron}
 M_{\lambda,U}(x):=\frac1{2\pi i}\int_{c-iU}^{c+iU}F_\lambda(s)\frac{x^s}{s}\,ds,\qquad M_{Z1,U}:=M_{0,U}.
\end{equation}

\begin{proposition}[Finite-height Perron integrals along the family]\label{prop:finite-Perron}
For fixed $c>1$, $U<\infty$ and $x>0$,
\begin{equation}\label{eq:finite-Perron-limit}
 M_{\lambda,U}(x)\longrightarrow M_{Z1,U}(x)\qquad(\lambda\downarrow0).
\end{equation}
\end{proposition}

\begin{proof}
The segment $\{c+it:|t|\le U\}$ is compact and disjoint from every divisor \CGeqref{eq:family-divisor}, and $P_\lambda\to P_{Z1}$ uniformly on it, so $F_\lambda\to F_{Z1}$ uniformly there.
\end{proof}

\begin{theorem}[Reciprocal-coefficient criterion]\label{thm:coefficient-criterion}
Assume that in $\Re s>1$
\begin{equation}\label{eq:Z1-reciprocal-Dirichlet}
 F_{Z1}(s)=\sum_{n\ge1}\frac{\mu_{Z1}(n)}{n^s}
\end{equation}
absolutely, and let $M_{Z1}(x)=\sum_{n\le x}\mu_{Z1}(n)$, so that
\begin{equation}\label{eq:Z1-increments}
 \mu_{Z1}(n)=M_{Z1}(n)-M_{Z1}(n-1).
\end{equation}
Then the following are equivalent: (i) $M_{Z1}(n)=M(n)$ for every $n\ge1$; (ii) $\mu_{Z1}=\mu$; (iii) $F_{Z1}=1/\zeta$ in $\Re s>1$; (iv) $\zeta_{Z1}=\zeta$; (v) RH.
\end{theorem}

\begin{proof}
(i)$\Leftrightarrow$(ii) by \CGeqref{eq:Z1-increments}; (ii)$\Leftrightarrow$(iii) by uniqueness of absolutely convergent Dirichlet coefficients; (iii)$\Leftrightarrow$(iv) by analytic continuation, both functions being nonzero in $\Re s>1$; (iv)$\Leftrightarrow$(v) by Theorem~\CGref{thm:Hamburger-Z1}.
\end{proof}

\begin{remark}[A route that does not close Step 3]\label{rem:negative-convolution}
The following approach does not work.  One may try to expand both factors of $F_{Z1}=F_{\RVM}\,\Pi_{\rm tr}^{-1}$ in \CGeqref{eq:FZ1} as absolutely convergent Dirichlet series and identify their Dirichlet convolution with $\mu$.  By Theorem~\CGref{thm:no-Dirichlet-RVM} the expansion of $F_{\RVM}$ does not exist in any half-plane, so there are no coefficients to convolve.  The endpoint criterion concerns $\zeta_{Z1}$ itself, as in Theorem~\CGref{thm:Hamburger-Z1}.
\end{remark}

\section{The endpoint at infinity and the Euler products}
\label{sec:euler-products}

This section applies the clock decomposition of Theorem~\CGref{thm:halfplane-RVM} to the ordinates $\gamma_k$ of the endpoint, and constructs the prime-power counts of $\zeta_{\RVM}$ and $\zeta_{Z1}$ as Mellin transforms of their logarithms.  No hypothesis is used.

\subsection{The value of the endpoint at infinity}
\label{subsec:value-infinity}

With the zero count $N_\zeta$ of \CGeqref{eq:RvM-full} put
\begin{equation}\label{eq:frak-e}
 \mathfrak e(\tau):=N_\zeta(\tau)-\NR(\tau),\qquad \mathfrak e_1(\tau):=\int_{\tau_0}^{\tau}\mathfrak e(v)\,dv\qquad(\tau\ge\tau_0=4\pi).
\end{equation}
By \CGeqref{eq:RvM-full}, $\mathfrak e(\tau)=S(\tau)+O(1/\tau)$, and Littlewood's bound $\int_0^TS(t)\,dt=O(\log T)$ \cite[Thm.~9.9]{Titchmarsh1986} gives an absolute constant $C_{\mathfrak e}$ with
\begin{equation}\label{eq:frak-e1-bound}
 |\mathfrak e_1(\tau)|\le C_{\mathfrak e}\log\tau\qquad(\tau\ge\tau_0).
\end{equation}
For the clock the corresponding primitive is $-\mathfrak s_1$, bounded by Lemma~\CGref{lem:saw-primitive}.

\begin{theorem}[Clock decomposition of the endpoint]\label{thm:Z1-clock}
Let $s=\tfrac12+z$, $z=\sigma'+it$, $\sigma'>0$.  There is an absolute constant $R_{\rm end}>0$ such that
\begin{equation}\label{eq:Z1-clock}
 \log\zeta_{Z1}(s)=\log\varpi_{Z1}-\int_{\tau_0}^\infty\mathfrak e_1(\tau)\Bigl(\frac1{(\tau-iz)^2}+\frac1{(\tau+iz)^2}\Bigr)d\tau+O\Bigl(\frac1{|z|}\Bigr)\qquad(|z|\ge R_{\rm end}),
\end{equation}
where
\begin{equation}\label{eq:varpi-Z1}
 \varpi_{Z1}:=P_{Z1}(\tfrac12)\exp\Bigl(\nu_3+2\int_{\tau_0}^\infty\frac{\mathfrak e_1(\tau)}{\tau^2}\,d\tau\Bigr)>0 .
\end{equation}
Consequently:
\begin{enumerate}[label=\textup{(\roman*)}]
\item for $\sigma'>0$ and $|z|\ge R_{\rm end}$,
\begin{equation}\label{eq:Z1-growth}
 \bigl|\log\zeta_{Z1}(s)-\log\varpi_{Z1}\bigr|\le\frac{2(\pi+2)C_{\mathfrak e}}{\sigma'}\Bigl(1+\log\bigl(2|t|+2\sigma'+\tau_0\bigr)\Bigr)+\frac C{|z|};
\end{equation}
\item $\zeta_{Z1}(\sigma)\to\varpi_{Z1}$ as $\sigma\to+\infty$;
\item the value at infinity is the ordered product
\begin{equation}\label{eq:varpi-Z1-product}
 \varpi_{Z1}=\varpi_{\RVM}\prod_{k\ge1}\frac{\tfrac14+\tau_k^2}{\tfrac14+\gamma_k^2}:=\varpi_{\RVM}\lim_{K\to\infty}\prod_{k\le K}\frac{\tfrac14+\tau_k^2}{\tfrac14+\gamma_k^2}.
\end{equation}
\end{enumerate}
\end{theorem}

\begin{proof}
As in \CGeqref{eq:transition-product}, $\log P_{Z1}(s)-\log P_{Z1}(\frac12)=\sum_kf(\gamma_k)$ with $f(\tau)=\Log(1+z^2/\tau^2)$.  Since $\gamma_1>\tau_0$, $N_\zeta(\tau_0)=0$ and $\mathfrak e(\tau_0)=-\NR(\tau_0)$.  Stieltjes integration against $N_\zeta=\NR+\mathfrak e$ and two integrations by parts ($f\mathfrak e\to0$ and $\mathfrak e_1f'\to0$ at $\infty$, because $f(\tau)=O(|z|^2\tau^{-2})$, $f'(\tau)=O(|z|^2\tau^{-3})$, $\mathfrak e(\tau)=O(\log\tau)$ and $\mathfrak e_1(\tau)=O(\log\tau)$ by \CGeqref{eq:frak-e1-bound}; and $\mathfrak e_1(\tau_0)=0$) give
\begin{equation}\label{eq:Z1-split}
 \sum_kf(\gamma_k)=\Lambda_{\rm sm}(z)+\int_{\tau_0}^\infty\mathfrak e_1f''\,d\tau,
\end{equation}
with the same $\Lambda_{\rm sm}$ as in \CGeqref{eq:transition-split}.  For $\sigma'>\tfrac12$ the proof of Theorem~\CGref{thm:halfplane-RVM} gives $\Lambda_{\rm sm}-\log G=\nu_3+O(|z|^{-1})$; for $0<\sigma'\le\tfrac12$ the same estimate holds because \CGeqref{eq:smooth-exact} is an exact formula valid on all of $\Re z>0$ (the singularities of $f'$ at $\tau=\pm iz$ lie off the positive axis), $|\tau^2+z^2|\ge\frac34|z|^2$ on $[0,\tau_0]$ for $|z|\ge2\tau_0$ whatever $\sigma'>0$, Stirling's formula with remainder $O(|v|^{-1})$ is uniform in $|\arg v|\le\frac\pi2$, and $\mathrm{Log}(z^2-\frac14)=2\,\mathrm{Log}\,z+O(|z|^{-2})$ since $\Im(z^2-\frac14)=2\sigma't$ has the sign of $t$, so no cut is crossed; an $O$-estimate is not transported by analytic continuation, and none is needed.  Thus, inserting $f''$ from \CGeqref{eq:saw-ibp} gives \CGeqref{eq:Z1-clock}.

For (i), $|\tau\mp iz|^2=(\tau\pm t)^2+\sigma'^2$.  For real $a$, splitting at $X=2|a|+2\sigma'+\tau_0$ and using $(\tau-a)^2\ge\tau^2/4$ for $\tau\ge X$,
\begin{equation}\label{eq:log-kernel}
 \int_{\tau_0}^\infty\frac{\log\tau}{(\tau-a)^2+\sigma'^2}\,d\tau\le\frac\pi{\sigma'}\log X+4\,\frac{1+\log X}X\le\frac{\pi+2}{\sigma'}(1+\log X).
\end{equation}
With \CGeqref{eq:frak-e1-bound} this gives \CGeqref{eq:Z1-growth}.  For (ii) take $t=0$: the right side of \CGeqref{eq:Z1-growth} tends to $0$ as $\sigma'\to\infty$.

For (iii), Theorem~\CGref{thm:halfplane-RVM} gives $\varpi_{\RVM}=P_{\RVM}^{\crit}(\frac12)\exp(\nu_3-2\int\mathfrak s_1\tau^{-2}d\tau)$, so
\begin{equation}\label{eq:varpi-ratio}
 \frac{\varpi_{Z1}}{\varpi_{\RVM}}=\frac{P_{Z1}(\frac12)}{P_{\RVM}^{\crit}(\frac12)}\exp\Bigl(2\int_{\tau_0}^\infty\frac{\mathfrak e_1(\tau)+\mathfrak s_1(\tau)}{\tau^2}\,d\tau\Bigr).
\end{equation}
Since $\mathfrak e+\psi(\NR+\frac12)=N_\zeta-n_{\RVM}$ and $(\mathfrak e_1+\mathfrak s_1)(X)/X\to0$,
\begin{equation}\label{eq:count-difference-integral}
 \int_{\tau_0}^\infty\frac{\mathfrak e_1+\mathfrak s_1}{\tau^2}\,d\tau=\lim_{X\to\infty}\int_{\tau_0}^X\frac{N_\zeta(\tau)-n_{\RVM}(\tau)}\tau\,d\tau=\lim_{K\to\infty}\sum_{k\le K}\log\frac{\tau_k}{\gamma_k}.
\end{equation}
In the last step, $\int_{\tau_0}^X(N_\zeta-n_{\RVM})\tau^{-1}d\tau=\sum_{\gamma_k\le X}\log\frac X{\gamma_k}-\sum_{\tau_k\le X}\log\frac X{\tau_k}$; the indices counted by only one of the two sums are $O(\log X)$ in number and, by Lemma~\CGref{lem:bounded-displacement}, have ordinates in $[X-C_\eta,X]$, where $0\le\log(X/\gamma_k)\le-\log(1-C_\eta/X)=O(1/X)$, so they contribute $O(\log X/X)$.  Finally $P_{Z1}(\frac12)/P_{\RVM}^{\crit}(\frac12)=\prod_k\frac{\gamma_k^2}{\frac14+\gamma_k^2}\cdot\frac{\frac14+\tau_k^2}{\tau_k^2}$, an absolutely convergent product, and multiplying by $\prod_{k\le K}\tau_k^2/\gamma_k^2$ gives \CGeqref{eq:varpi-Z1-product}.
\end{proof}

The ordinates enter $\varpi_{Z1}$ only through the displacements $\gamma_k-\tau_k$, in agreement with Remark~\CGref{rem:only-gamma-tau}.  The product in \CGeqref{eq:varpi-Z1-product} is an \emph{ordered} product: its factors are $1+O(|\eta_k|/\tau_k)$, and $\sum_k|\eta_k|/\tau_k$ is not known to converge ($\sum_k\tau_k^{-1}$ diverges), so absolute convergence is not asserted and no rearrangement is permitted.  The limit is that of the partial products over $k\le K$ in increasing height, which is the order in which \CGeqref{eq:count-difference-integral} is evaluated, and this order is part of the definition of $\varpi_{Z1}$; the equality of the ordered limit with $\varpi_{Z1}$ is what \CGeqref{eq:count-difference-integral} proves, and the value is not claimed to be independent of other groupings.

\begin{theorem}[The value at infinity decides the endpoint identity]\label{thm:varpi-Z1}
$\varpi_{Z1}\ge1$, and $\varpi_{Z1}=1$ if and only if $\zeta_{Z1}=\zeta$.
\end{theorem}

\begin{proof}
Write $\zeta_{Z1}=\zeta\,\mathscr R_{Z1}$ with $\mathscr R_{Z1}=P_{Z1}/P_\zeta$ as in \CGeqref{eq:Z1-defect-ratio}, and compare the products factor by factor, with the zeros $\rho_k=\beta_k+i\gamma_k$ listed with the ordinates (Theorem~\CGref{thm:Hadamard}).  Factors with $\rho_k(1-\rho_k)=\frac14+\gamma_k^2$ cancel.  The remaining factors of $P_\zeta$ come in pairs $\{\rho,1-\bar\rho\}$ with the same ordinate $\gamma$; let $\mathcal P$ be a set of representatives, one zero $\rho$ with $\Im\rho>0$ and $\Re\rho\ne\frac12$ from each pair, repeated according to multiplicity.  With $d_\rho=\Re\rho-\frac12$ and $W_\rho=(s-\frac12)^2+\gamma^2$, the corresponding factor of $\mathscr R_{Z1}$ is
\begin{equation}\label{eq:pair-ratio}
 \frac{W_\rho^2}{(W_\rho-d_\rho^2)^2+4d_\rho^2\gamma^2}\cdot c_\gamma,\qquad
 c_\gamma=\frac{|\rho(1-\rho)|^2}{(\frac14+\gamma^2)^2}=1+\frac{2d_\rho^2(\gamma^2-\frac14)+d_\rho^4}{(\frac14+\gamma^2)^2}>1 .
\end{equation}
For real $\sigma>1$, $W_\rho\ge\gamma^2\ge\gamma_1^2$, and the logarithm $\ell_\rho(\sigma)$ of the first factor satisfies $|\ell_\rho(\sigma)|\le14d_\rho^2/W_\rho\le14d_\rho^2/\gamma^2$, while $0<\log c_\gamma\le c_\gamma-1\le3d_\rho^2/\gamma^2$.  Hence $\log\mathscr R_{Z1}(\sigma)=\sum_{\rho\in\mathcal P}\bigl(\ell_\rho(\sigma)+\log c_\gamma\bigr)$ is a series dominated, uniformly in $\sigma>1$, by $\sum_{\rho\in\mathcal P}17d_\rho^2/\gamma^2\le\frac{17}4\sum_{\rho\in\mathcal P}\gamma^{-2}<\infty$; each term tends to $\log c_\gamma$ as $\sigma\to\infty$ because $W_\rho\to\infty$; so, by dominated convergence for series, $\log\mathscr R_{Z1}(\sigma)\to\sum_{\rho\in\mathcal P}\log c_\gamma$.  With $\zeta(\sigma)\to1$ this gives $\varpi_{Z1}=\lim_{\sigma\to\infty}\zeta_{Z1}(\sigma)=\prod_{\rho\in\mathcal P}c_\gamma\ge1$ (an absolutely convergent product of factors $>1$), with equality exactly when there are no such pairs, that is, when $\mathscr R_{Z1}\equiv1$.
\end{proof}

\begin{corollary}[The endpoint identity as an identity of the displacements]\label{cor:product-identity}
The following are equivalent: (D); $\zeta_{Z1}=\zeta$; $\varpi_{Z1}=1$;
\begin{equation}\label{eq:product-identity}
 \lim_{K\to\infty}\prod_{k\le K}\frac{\tfrac14+\gamma_k^2}{\tfrac14+\tau_k^2}=\varpi_{\RVM}
\end{equation}
(ordered product). Unconditionally $\lim_{\sigma\to+\infty}\Pi_{\rm tr}(\sigma)=\varpi_{Z1}/\varpi_{\RVM}=\varpi_{Z1}L\ge L$, and the equality $\lim_{\sigma\to+\infty}\Pi_{\rm tr}(\sigma)=L=1/\varpi_{\RVM}$ is a fourth condition equivalent to the three above.
\end{corollary}

\begin{proof}
Theorems~\CGref{thm:Hamburger-Z1}, \CGref{thm:Z1-clock} and~\CGref{thm:varpi-Z1}, and $\Pi_{\rm tr}=\zeta_{Z1}/\zeta_{\RVM}\to\varpi_{Z1}/\varpi_{\RVM}$ by Theorems~\CGref{thm:halfplane-RVM} and~\CGref{thm:Z1-clock}(ii).
\end{proof}

\begin{proposition}[Boundedness of the endpoint]\label{prop:Z1-bounded}
Let $\sigma_0>1$.  Then $|\zeta_{Z1}|$ is bounded above and below on $\Re s\ge\sigma_0$ if and only if
\begin{equation}\label{eq:conjugate-poisson}
 \sup_{\sigma'\ge\sigma_0-\frac12,\ t\ge0}\Bigl|\Re\int_{\tau_0}^\infty\frac{\mathfrak e_1(\tau)+\mathfrak s_1(\tau)}{(\tau-t+i\sigma')^2}\,d\tau\Bigr|<\infty .
\end{equation}
\end{proposition}

\begin{proof}
$\tau+iz=\tau-t+i\sigma'$.  For $t\ge0$ the term with $(\tau-iz)^{-2}$ in \CGeqref{eq:Z1-clock} is bounded by \CGeqref{eq:frak-e1-bound}, since $|\tau-iz|\ge\tau+t$; the separate integral $\int_{\tau_0}^\infty\mathfrak s_1(\tau)(\tau+iz)^{-2}d\tau$ is bounded directly by Lemma~\CGref{lem:saw-primitive} and $\int_{\tau_0}^\infty|\tau+iz|^{-2}d\tau\ll1/\sigma'$.  By \CGeqref{eq:Z1-clock}, $\log|\zeta_{Z1}|$ is bounded on $\Re s\ge\sigma_0$, $t\ge0$, exactly when \CGeqref{eq:conjugate-poisson} holds, and $\zeta_{Z1}(\bar s)=\overline{\zeta_{Z1}(s)}$ covers $t<0$.
\end{proof}

The function in \CGeqref{eq:conjugate-poisson} is, after one integration by parts, the conjugate Poisson integral at distance $\sigma'$ of the count difference $N_\zeta-n_{\RVM}$; it depends only on the displacements $\gamma_k-\tau_k$.

\subsection{The Euler products of \texorpdfstring{$\zeta_{\RVM}$}{zeta\_RVM} and \texorpdfstring{$\zeta_{Z1}$}{zeta\_Z1}}
\label{subsec:euler-products}

\begin{theorem}[Prime-power counts as Mellin transforms]\label{thm:mellin-counts}
Let $X\in\{\RVM,Z1\}$.  There is a unique function $J_X^{\rm M}$ on $(1,\infty)$ with $J_X^{\rm M}(x)x^{-c}\in L^2((1,\infty),dx/x)$ for every $c>1$ and
\begin{equation}\label{eq:mellin-count}
 \log\zeta_X(s)=s\int_1^\infty J_X^{\rm M}(x)\,x^{-s-1}\,dx\qquad(\Re s>1),
\end{equation}
the branch of $\log\zeta_X$ being real on $(1,\infty)$.  We use the notation $\int x^{-s}\,dJ_X^{\rm M}:=s\int_1^\infty J_X^{\rm M}(x)x^{-s-1}dx$ distributionally; no bounded-variation representative is asserted.  With this convention $\zeta_X(s)=\exp\bigl(\int x^{-s}\,dJ_X^{\rm M}(x)\bigr)$.
\end{theorem}

\begin{proof}
$\zeta_X$ has no zeros or poles in $\Re s>1$, so $\log\zeta_X$ is holomorphic there.  For $c>1$, Theorem~\CGref{thm:halfplane-RVM} ($X=\RVM$) and Theorem~\CGref{thm:Z1-clock} ($X=Z1$), together with continuity on the compact part, give $|\log\zeta_X(s)|\le C_c(1+\log(|t|+2))$ on $\Re s\ge c$.  Hence $\mathsf F(s):=\log\zeta_X(s)/s$ is holomorphic on $\Re s>c$ and satisfies $\sup_{\sigma>c}\int_{\R}|\mathsf F(\sigma+it)|^2dt\le B_c^2<\infty$: $\mathsf F(\,\cdot+c)$ belongs to the Hardy space $H^2$ of the right half-plane.  By the Paley--Wiener theorem for $H^2$ of a half-plane \cite[Thm.~19.2]{Rudin1987} there is $g\in L^2(0,\infty)$ with $\mathsf F(s)=\int_0^\infty g(u)e^{-(s-c)u}du$ for $\Re s>c$, the integral converging absolutely.  Put $j(u):=g(u)e^{cu}$ for $u>0$; then
\begin{equation}\label{eq:mellin-count-laplace}
 \frac{\log\zeta_X(s)}s=\int_0^\infty j(u)\,e^{-su}\,du\qquad(\Re s>c),\qquad e^{-cu}j\in L^2(0,\infty).
\end{equation}
The function $j$ does not depend on $c$: if $c'>c$, the representation obtained from $c'$ agrees with \CGeqref{eq:mellin-count-laplace} on $\Re s>c'$, and the Laplace transform is injective on functions with $e^{-c'u}j\in L^2$ (the Fourier transform of $e^{-\sigma u}j(u)\mathbf 1_{u>0}$ is $\mathsf F(\sigma+i\cdot)$, and the Fourier--Plancherel transform is injective).  Define $J_X^{\rm M}(e^u):=j(u)$ for $u>0$, i.e.\ $J_X^{\rm M}(x)=j(\log x)$ on $(1,\infty)$; the substitution $x=e^u$ turns \CGeqref{eq:mellin-count-laplace} into \CGeqref{eq:mellin-count}, and $e^{-cu}j\in L^2(0,\infty)$ is the stated condition $J_X^{\rm M}(x)x^{-c}\in L^2((1,\infty),dx/x)$ for every $c>1$.  Uniqueness is the injectivity just used.
\end{proof}

\begin{theorem}[Distributional explicit formula]\label{thm:distributional-EF}
Let $X\in\{\RVM,Z1\}$, with ordinates $\kappa_k=\tau_k$ for $X=\RVM$ and $\kappa_k=\gamma_k$ for $X=Z1$.  In $\mathcal D'(1,\infty)$,
\begin{equation}\label{eq:distributional-EF}
 (\log x)\,\frac{dJ_X^{\rm M}}{dx}=1-\frac2{\sqrt x}\,\Re\sum_{k\ge1}x^{i\kappa_k}-\frac1{x(x^2-1)},
\end{equation}
the series converging in $\mathcal D'(1,\infty)$.  The explicit-formula count
\begin{equation}\label{eq:J-X}
 J_X(x):=\li(x)-2\Re\sum_{k\ge1}\Ei\Bigl(\bigl(\tfrac12+i\kappa_k\bigr)\log x\Bigr)-\log2+\int_x^\infty\frac{dy}{y(y^2-1)\log y}
\end{equation}
converges in $\mathcal D'(1,\infty)$ ($J_{\RVM}$ is the function \CGeqref{eq:J-RVM}), and $J_X^{\rm M}-J_X$ is a constant $C_X$ on $(1,\infty)$.
\end{theorem}

\begin{proof}
With $\rho_k=\frac12+i\kappa_k$, the factorization $\zeta_X=P_X/G$, $P_X'/P_X=\sum_k\bigl(\frac1{s-\rho_k}+\frac1{s-\bar\rho_k}\bigr)$, and the digamma expansion $\frac12\frac{\Gamma'}{\Gamma}(\frac s2)+\frac1s=-\frac{\gamma_E}2+\sum_{n\ge1}\bigl(\frac1{2n}-\frac1{s+2n}\bigr)$ give, for $\Re s>1$,
\begin{equation}\label{eq:log-derivative-X}
 -\frac{\zeta_X'}{\zeta_X}(s)=\frac1{s-1}-\sum_{k\ge1}\Bigl(\frac1{s-\rho_k}+\frac1{s-\bar\rho_k}\Bigr)+\sum_{n\ge1}\Bigl(\frac1{2n}-\frac1{s+2n}\Bigr)-\frac{\gamma_E}2-\frac12\log\pi .
\end{equation}
Each term is a Mellin--Stieltjes transform $\int_{[1,\infty)}x^{-s}d\mu$: $dx\mapsto\frac1{s-1}$, $x^{\rho-1}dx\mapsto\frac1{s-\rho}$, and the sum over the trivial zeros is kept as a whole, $\sum_n\bigl(\frac1{2n}-\frac1{s+2n}\bigr)=\int_1^\infty\frac{(1-x^{-s})\,dx}{x(x^2-1)}$, a tempered distribution in $u=\log x$ whose restriction to $(1,\infty)$ is the density $-1/(x(x^2-1))$ (the two pieces $\sum\frac1{2n}$ and $x^{-2n-1}dx$ are not separately summable, nor is $1/(x(x^2-1))$ integrable at $x=1$, so the split is made only on $(1,\infty)$); $\delta_1\mapsto1$ accounts for the constant term.  After multiplication by $x^{-c}$ and the change $u=\log x$, the Fourier transforms of the partial sums over $k$ and over $n$ converge locally uniformly with polynomial bounds, so the series converge in $\mathcal S'$, and on $(1,\infty)$ the limit measure is the right side of \CGeqref{eq:distributional-EF} times $dx$.  Differentiating \CGeqref{eq:mellin-count} in $s$ shows that $-\zeta_X'/\zeta_X$ is the distributional Mellin transform of $(\log x)\,dJ_X^{\rm M}$; injectivity gives \CGeqref{eq:distributional-EF}.  The series over $k$ converges in $\mathcal D'(1,\infty)$ because $\int\varphi(x)x^{-1/2+i\kappa}dx$ is rapidly decreasing in $\kappa$ for $\varphi\in C_c^\infty(1,\infty)$ and $\kappa_k\asymp k/\log k$.  By Lemma~\CGref{lem:Ei}, $\Ei(\rho_ku)=e^{\rho_ku}/(\rho_ku)+\varrho_k$ with $\rho_k=\tfrac12+i\kappa_k$ and $\sum_k|\varrho_k|$ locally uniformly convergent, so \CGeqref{eq:J-X} converges in $\mathcal D'(1,\infty)$ for the same reason; differentiating termwise with $\frac d{dx}\Ei(\rho\log x)=x^{\rho-1}/\log x$, $(\log x)J_X'$ equals the right side of \CGeqref{eq:distributional-EF}.  A distribution on an interval with zero derivative is constant.
\end{proof}

\begin{proof}[Proof of Theorem~\CGref{thm:MS}]
Take $X=\RVM$ in Theorems~\CGref{thm:mellin-counts} and~\CGref{thm:distributional-EF}.  By Theorem~\CGref{thm:J-exists}, $J_{\RVM}$ is continuous on $(1,\infty)$, so $J_{\RVM}^{\rm M}=J_{\RVM}+C_{\RVM}$ pointwise, and $J_{\RVM}^{*}=J_{\RVM}$ on $[2,\infty)$.
\end{proof}

\begin{theorem}[Cell form of the endpoint identity]\label{thm:cell-form}
The following are equivalent.
\begin{enumerate}[label=\textup{(\roman*)}]
\item (D), that is, $\zeta_{Z1}=\zeta$.
\item $J_{Z1}^{\rm M}$ is constant on every interval $(n,n+1)$, $n\ge1$, and its jumps $\mathsf b_n:=v_n-v_{n-1}$, where $v_n$ is its value on $(n,n+1)$ and $v_0:=0$, satisfy $\sum_{n\le X}|\mathsf b_n|=O_\eps(X^{1+\eps})$ for every $\eps>0$.
\item For every $n\ge1$, in $\mathcal D'(n,n+1)$,
\begin{equation}\label{eq:cell-identity}
 2\,\Re\sum_{k\ge1}\bigl(x^{i\gamma_k}-x^{i\tau_k}\bigr)=\sqrt x\,\log x\;J_{\RVM}'(x),
\end{equation}
and the jumps of $J_{Z1}^{\rm M}$ satisfy the bound in (ii).
\end{enumerate}
In this case $J_{Z1}^{\rm M}$ is Riemann's prime-power count $J(x)=\sum_{p^m\le x}1/m$.
\end{theorem}

\begin{proof}
(ii)$\Rightarrow$(i): $|v_n|\le\sum_{m\le n}|\mathsf b_m|\ll n^{1+\eps}$, so for $\Re s>1$ summation by parts in \CGeqref{eq:mellin-count} gives $\log\zeta_{Z1}(s)=\sum_nv_n(n^{-s}-(n+1)^{-s})=\sum_{n\ge1}\mathsf b_nn^{-s}$, absolutely convergent by partial summation.  The exponential of an absolutely convergent Dirichlet series is again an absolutely convergent Dirichlet series, so $\zeta_{Z1}$ satisfies \CGeqref{eq:Z1-D}.

(i)$\Rightarrow$(ii): by Theorem~\CGref{thm:Hamburger-Z1}, $\zeta_{Z1}=\zeta$, and $\log\zeta(s)=s\int_1^\infty J(x)x^{-s-1}dx$ for $\Re s>1$; by uniqueness in Theorem~\CGref{thm:mellin-counts}, $J_{Z1}^{\rm M}=J$, which is constant between consecutive integers, with jumps $1/m$ at $p^m$ and $\sum_{n\le X}|\mathsf b_n|=J(X)\ll X$.

(ii)$\Leftrightarrow$(iii): Theorem~\CGref{thm:distributional-EF} for $X=Z1$ and $X=\RVM$, with $dJ_{\RVM}^{\rm M}/dx=J_{\RVM}'$, gives in $\mathcal D'(1,\infty)$
\begin{equation}\label{eq:cell-difference}
 (\log x)\Bigl(\frac{dJ_{Z1}^{\rm M}}{dx}-J_{\RVM}'(x)\Bigr)=-\frac2{\sqrt x}\,\Re\sum_{k\ge1}\bigl(x^{i\gamma_k}-x^{i\tau_k}\bigr),
\end{equation}
so $dJ_{Z1}^{\rm M}/dx=0$ on $(n,n+1)$ is \CGeqref{eq:cell-identity}.
\end{proof}

\begin{remark}[The final step in cell form]\label{rem:cell-form}
Theorem~\CGref{thm:cell-form} is the exact form of the continuous-to-atomic passage of Remark~\CGref{rem:D-atomic}.  Both sides of \CGeqref{eq:cell-identity} are built from the midpoint construction: the right side from the clock alone ($J_{\RVM}'$ is bounded on $[2,\infty)$ by Theorem~\CGref{thm:unit-cell}, equal to $1/\log x$ plus the unit-periodic first harmonic and the harmonics with stationary points $2\pi x^{1/m}$, Section~\CGref{sec:integer-weights}), the left side from the displacements $\gamma_k-\tau_k$.  The final step of the method is the identity \CGeqref{eq:cell-identity} on every cell, equivalently \CGeqref{eq:product-identity}.  Numerically, the explicit formula \CGeqref{eq:J-X} with the $\gamma_k$ is constant between consecutive integers up to the truncation of the list (Experiment~\CGref{exp:primes}, Figure~\CGref{fig:prime-counts}), and the partial products of \CGeqref{eq:product-identity} approach $\varpi_{\RVM}$ (Experiment~\CGref{exp:product-identity}).
\end{remark}

\section{The atomized family and the programme for the final step}
\label{sec:atomized}

The continuous family $\zeta_\lambda$ keeps, for every $\lambda$, the functional equation and the zeros on the critical line.  Its atomization $\tilde\zeta_{\Z,\lambda}$, built from the cell masses of the prime coordinate, keeps the Euler product and the Dirichlet series.  The two are joined by a factor $e^{\tilde h_\lambda}$, and both families have a limit at $\lambda=0$.  Hamburger's theorem needs the functional equation and the Dirichlet series on one and the same function; at $\lambda=0$ this is the vanishing of the joining factor.  This section constructs the atomized family for every $\lambda$, proves its analytic properties, writes the joining factor cell by cell, and lists in Subsection~\CGref{subsec:programme} the steps that remain, in the order in which they will be followed.  The passage from a continuous prime system to a discrete one is classical for Beurling generalized primes \cite{Beurling1937,DiamondMontgomeryVorhauer2006,BrouckeVindas2024,DiamondZhang2016}; the atomization used here places the masses at the integers, which produces ordinary Dirichlet series.

Throughout, $\eta_k=\tau_k-\gamma_k$, $\gamma_k(\lambda)=(1-\lambda)\gamma_k+\lambda\tau_k$, $\rho_{k,\lambda}=\frac12+i\gamma_k(\lambda)$, and $\lambda'$ denotes an integration variable in $[0,1]$.

\subsection{The family at infinity}
\label{subsec:family-infinity}

Put $N_\lambda(\tau)=\#\{k:\gamma_k(\lambda)\le\tau\}$, $\mathfrak e^{(\lambda)}=N_\lambda-\NR$ and $\mathfrak e_1^{(\lambda)}(\tau)=\int_{\tau_0}^\tau\mathfrak e^{(\lambda)}$, so that $\mathfrak e_1^{(0)}=\mathfrak e_1$ and $\mathfrak e_1^{(1)}=-\mathfrak s_1$.

\begin{proposition}[Clock decomposition along the family]\label{prop:family-clock}
There is an absolute constant $C$ with $|\mathfrak e_1^{(\lambda)}(\tau)|\le C\log\tau$ for $\tau\ge\tau_0$ and all $\lambda\in[0,1]$.  Consequently Theorem~\CGref{thm:Z1-clock} holds for every $\zeta_\lambda$ with $\mathfrak e_1$ replaced by $\mathfrak e_1^{(\lambda)}$, uniformly in $\lambda$; in particular $\varpi_\lambda:=\lim_{\sigma\to+\infty}\zeta_\lambda(\sigma)$ exists and
\begin{equation}\label{eq:varpi-lambda}
 \varpi_\lambda=\varpi_{\RVM}\lim_{K\to\infty}\prod_{k\le K}\frac{\frac14+\tau_k^2}{\frac14+\gamma_k(\lambda)^2},\qquad \varpi_1=\varpi_{\RVM},\quad\varpi_0=\varpi_{Z1}.
\end{equation}
\end{proposition}

\begin{proof}
$\gamma_k(\lambda)$ lies between $\gamma_k$ and $\tau_k$, so $\gamma_1(\lambda)>\tau_0$.  For $X\ge\tau_0$,
\begin{equation}\label{eq:family-count-primitive}
 \int_{\tau_0}^X(N_\lambda-n_{\RVM})=\sum_k\bigl((X-\gamma_k(\lambda))^+-(X-\tau_k)^+\bigr)=(1-\lambda)\sum_{\gamma_k(\lambda),\tau_k\le X}\eta_k+O(\log X),
\end{equation}
the indices counted by only one of the two sums being $O(\log X)$ in number with ordinates in $[X-C_\eta,X]$ (Lemma~\CGref{lem:bounded-displacement}).  The case $\lambda=0$ of \CGeqref{eq:family-count-primitive} identifies the last sum with $\mathfrak e_1+\mathfrak s_1+O(\log X)=O(\log X)$, by \CGeqref{eq:frak-e1-bound} and Lemma~\CGref{lem:saw-primitive}.  Since $\mathfrak e_1^{(\lambda)}=\int(N_\lambda-n_{\RVM})-\mathfrak s_1$, the bound follows, and the proof of Theorem~\CGref{thm:Z1-clock} applies verbatim.
\end{proof}

Define the \emph{flow at infinity}
\begin{equation}\label{eq:Upsilon}
 \Upsilon(\lambda):=\sum_{k\ge1}\frac{2\gamma_k(\lambda)\,\eta_k}{\frac14+\gamma_k(\lambda)^2}\qquad(\text{summed in the order of }k).
\end{equation}

\begin{proposition}[The value at infinity along the family]\label{prop:family-flow}
The series \CGeqref{eq:Upsilon} converges uniformly on $[0,1]$, $\lambda\mapsto\log\varpi_\lambda$ is $C^1$ on $[0,1]$, and
\begin{equation}\label{eq:varpi-derivative}
 \frac{d}{d\lambda}\log\varpi_\lambda=-\Upsilon(\lambda)=\lim_{\sigma\to+\infty}\mathcal V_\lambda(\sigma),\qquad
 \Upsilon'(\lambda)=\sum_{k\ge1}\frac{2\eta_k^2\bigl(\frac14-\gamma_k(\lambda)^2\bigr)}{\bigl(\frac14+\gamma_k(\lambda)^2\bigr)^2}<0 ,
\end{equation}
with $\mathcal V_\lambda$ the generator of Theorem~\CGref{thm:generator}.  Moreover
\begin{equation}\label{eq:flow-integral}
 \int_0^1\Upsilon(\lambda)\,d\lambda=\sum_{k\ge1}\log\frac{\frac14+\tau_k^2}{\frac14+\gamma_k^2}=\log L+\log\varpi_{Z1},
\end{equation}
so that (D) holds if and only if $\int_0^1\Upsilon(\lambda)\,d\lambda=\log L$, and $\Upsilon(1)<\log L+\log\varpi_{Z1}<\Upsilon(0)$.
\end{proposition}

\begin{proof}
Put $g(u)=2u/(\frac14+u^2)$.  The series $\Upsilon(1)=\sum_kg(\tau_k)\eta_k$ converges by summation by parts against $N_\zeta-n_{\RVM}$, exactly as in \CGeqref{eq:count-difference-integral}.  Since $|g(\gamma_k(\lambda))-g(\tau_k)|\le C|\eta_k|\gamma_k^{-2}$ and $\sum_k\eta_k^2\gamma_k^{-2}<\infty$, the series $\Upsilon(\lambda)-\Upsilon(1)$ converges absolutely and uniformly; the same holds for the termwise derivative, which is the series in \CGeqref{eq:varpi-derivative}. Each term of that series is $\le0$ because $\gamma_k(\lambda)>\frac12$, and the term $k=1$ is $<0$ because $\eta_1\ne0$: the first ordinate satisfies $\gamma_1<14.14$ (its rigorous location $\gamma_1=14.1347\ldots$ is classical, \cite[Ch.~XV]{Titchmarsh1986}), whereas $\NR(14.14)=0.4499\ldots<\frac12=\NR(\tau_1)$ by \CGeqref{eq:clock}, so $\tau_1>14.14>\gamma_1$ by the monotonicity of $\NR$ on $[2\pi,\infty)$, and $\eta_1=\tau_1-\gamma_1>0$. Hence $\Upsilon'<0$; the formula alone gives only $\Upsilon'\le0$, with strict inequality exactly when $\boldsymbol\gamma\ne\boldsymbol\tau$.  Differentiating \CGeqref{eq:varpi-lambda} termwise gives $\frac d{d\lambda}\log\varpi_\lambda=-\Upsilon(\lambda)$, and the limit $\sigma\to+\infty$ in the generator \CGeqref{eq:generator} is justified by Abel summation of the ordered ordinate sum: its partial sums are controlled by the count-difference integral, while the multiplier $q/(A+q)$ is bounded and monotone in $A$.  This gives the same ordered series.  Integrating from $0$ to $1$ and using $\varpi_1=\varpi_{\RVM}$, $\varpi_0=\varpi_{Z1}$ gives \CGeqref{eq:flow-integral}; the last statement follows from Theorem~\CGref{thm:varpi-Z1} and the monotonicity of $\Upsilon$.
\end{proof}

\begin{proposition}[The value at infinity and the verified zeros]\label{prop:varpi-verified}
Let $T_{\rm ver}\ge e^4$.  If every zero $\beta+i\gamma$ with $0<\gamma\le T_{\rm ver}$ has $\beta=\frac12$, then
\begin{equation}\label{eq:varpi-verified}
 0\le\log\varpi_{Z1}\le0.255\sum_{\gamma>T_{\rm ver}}\frac1{\gamma^2}\le0.255\Bigl(\frac{\log(T_{\rm ver}/2\pi)+1}{2\pi T_{\rm ver}}+\frac{2\log T_{\rm ver}+1}{T_{\rm ver}^2}\Bigr),
\end{equation}
the sum running over the zeros with $\gamma>T_{\rm ver}$, with multiplicity.  With $T_{\rm ver}=3\cdot10^{12}$ \cite{PlattTrudgian2021} the right side is less than $3.78\cdot10^{-13}$; in particular $0\le\log\varpi_{Z1}<4\cdot10^{-13}$.
\end{proposition}

\begin{proof}
By the proof of Theorem~\CGref{thm:varpi-Z1}, $\log\varpi_{Z1}=\sum_{\rho\in\mathcal P}\log c_\gamma$, and every $\rho\in\mathcal P$ has $\gamma>T_{\rm ver}$.  For $|d_\rho|\le\frac12$ and $\gamma\ge T_{\rm ver}$, \CGeqref{eq:pair-ratio} gives $\log c_\gamma\le c_\gamma-1\le(2d_\rho^2\gamma^2+d_\rho^4)\gamma^{-4}\le0.51\,\gamma^{-2}$.  Each $\rho\in\mathcal P$ stands for the two distinct zeros $\rho$ and $1-\bar\rho$ with ordinate $\gamma$, so $\sum_{\rho\in\mathcal P}0.51\gamma^{-2}\le0.255\sum_{\gamma>T_{\rm ver}}\gamma^{-2}$, the sum over all zeros; this is the first inequality.  For the second, Trudgian's explicit bound \cite[Cor.~1]{Trudgian2014} gives $N_\zeta(t)=\mathcal M(t)+\mathcal R(t)$ with $\mathcal M(t):=\frac t{2\pi}\log\frac t{2\pi e}+\frac78$ and $|\mathcal R(t)|\le0.112\log t+0.278\log\log t+2.51+0.2/t\le\log t$ for $t\ge e^4$.  Stieltjes integration and one integration by parts (the boundary term at $\infty$ vanishes because $N_\zeta(t)=O(t\log t)$) give
\[
 \sum_{\gamma>T}\frac1{\gamma^2}=\int_T^\infty\frac{\mathcal M'(t)}{t^2}\,dt-\frac{\mathcal R(T)}{T^2}+2\int_T^\infty\frac{\mathcal R(t)}{t^3}\,dt
 \le\frac{\log(T/2\pi)+1}{2\pi T}+\frac{\log T}{T^2}+\frac{\log T+\frac12}{T^2}
\]
for $T\ge e^4$, since $\mathcal M'(t)=\frac1{2\pi}\log\frac t{2\pi}$ and $2\int_T^\infty t^{-3}\log t\,dt=(\log T+\frac12)/T^2$; this is the second inequality.  For $T_{\rm ver}=3\cdot10^{12}$, $\log(T_{\rm ver}/2\pi)+1=27.8917\ldots$, the first term is $1.4797\ldots\cdot10^{-12}$, the second is below $10^{-23}$, and $0.255\cdot1.4798\cdot10^{-12}<3.78\cdot10^{-13}$.
\end{proof}

\subsection{Atomization by smooth cells}
\label{subsec:smooth-cells}

Fix $\varphi_n\in C_c^\infty(n-1,n+1)$, $n\ge2$, with $\sum_{n\ge2}\varphi_n=1$ on $[2,\infty)$ and $\varphi_n(m)=\delta_{nm}$ for integers $m$.  For $\lambda\in[0,1]$ let $J_\lambda$ be the explicit-formula count \CGeqref{eq:J-X} with the ordinates $\gamma_k(\lambda)$ (a distribution on $(1,\infty)$ by Theorem~\CGref{thm:distributional-EF}, which applies to every $\zeta_\lambda$ by Theorem~\CGref{thm:Z1-clock}), $J^*_\lambda=J_\lambda\mathbf 1_{[2,\infty)}$, and $\pi_\lambda=\sum_{j\ge1}\frac{\mu(j)}jJ^*_\lambda(x^{1/j})$.

\begin{definition}[Atomized family]\label{def:atomized}
The cell masses, the atomized zeta function and its coefficients are
\begin{equation}\label{eq:atomized-family}
 \tilde c_n(\lambda):=\langle d\pi_\lambda,\varphi_n\rangle,\qquad
 \tilde\zeta_{\Z,\lambda}(s):=\prod_{n\ge2}\bigl(1-n^{-s}\bigr)^{-\tilde c_n(\lambda)}=\sum_{N\ge1}\frac{\tilde a_N(\lambda)}{N^s},
\end{equation}
\begin{equation}\label{eq:atomized-coefficients}
 \tilde a_N(\lambda)=\sum_{\prod_nn^{r_n}=N}\ \prod_{n:\,r_n\ge1}\binom{\tilde c_n(\lambda)+r_n-1}{r_n},\qquad \tilde a_1(\lambda)=1 .
\end{equation}
\end{definition}

Unlike the discretizations of Beurling systems, where the atoms are placed at chosen real points \cite{DiamondMontgomeryVorhauer2006,BrouckeVindas2024,DiamondZhang2016}, the masses of $d\pi_\lambda$ are placed at the integers with real exponents; measures with continuous part and non-integer multiplicities of this kind appear also in \cite[Ch.~3]{LapidusvanFrankenhuysen2000}.  In \CGeqref{eq:atomized-coefficients}, $\binom{c+r-1}r:=c(c+1)\cdots(c+r-1)/r!$ for real $c$ and integers $r\ge1$.  Thus $\tilde a_p(\lambda)=\tilde c_p(\lambda)$ for primes $p$, $\tilde a_4=\tilde c_4+\frac12\tilde c_2(\tilde c_2+1)$, $\tilde a_6=\tilde c_6+\tilde c_2\tilde c_3$.  The pairing with a fixed smooth partition defines $\tilde c_n(\lambda)$ for every $\lambda$ without point values of $\pi_\lambda$, and $\lambda\mapsto\tilde c_n(\lambda)$ is continuous.  The individual masses depend on the partition $(\varphi_n)$, which is fixed throughout; the closure statements below are made for this fixed partition.  If $d\pi_0=\sum_p\delta_p$, then $\tilde c_n(0)=\mathbf 1_{\rm prime}(n)$ and $\tilde a_N(0)=1$ for all $N$.  Along the family, in $\mathcal D'(1,\infty)$,
\begin{equation}\label{eq:J-lambda-derivative}
 \frac{\partial J_\lambda}{\partial\lambda}(x)=2\sum_{k\ge1}\eta_k\,\Im\frac{x^{\rho_{k,\lambda}}}{\rho_{k,\lambda}},\qquad
 \frac{\partial}{\partial\lambda}\Bigl((\log x)\frac{dJ_\lambda}{dx}\Bigr)=\frac{2\log x}{\sqrt x}\sum_{k\ge1}\eta_k\sin\bigl(\gamma_k(\lambda)\log x\bigr).
\end{equation}

\subsection{The joining factor cell by cell}
\label{subsec:joining-factor}

Define the \emph{atomization kernel} and the \emph{cell defects}
\begin{equation}\label{eq:atomization-kernel}
 \mathfrak a_s(x):=\sum_{n\ge2}\varphi_n(x)\,n^{-s}-x^{-s},\qquad
 \mathfrak h_n(\lambda;s):=\bigl\langle d\pi_\lambda,\ \varphi_n(x)\bigl(n^{-s}-x^{-s}\bigr)\bigr\rangle .
\end{equation}
$\mathfrak a_s$ vanishes at every integer $\ge2$, and $\mathfrak h_n(\lambda;s)$ measures how far the mass of the cell $n$ is from being concentrated at $n$.

The explicit-formula count $J_\lambda$ need not vanish on $(1,2)$, whereas $\pi_\lambda$ is built from $J^*_\lambda=J_\lambda\mathbf 1_{[2,\infty)}$.  The first cell $(1,2)$ enters the joining factor through
\begin{equation}\label{eq:first-cell}
 \mathfrak u_\lambda(s):=C_\lambda\,2^{-s}+s\int_1^2J^{\rm M}_\lambda(x)\,x^{-s-1}\,dx,\qquad C_\lambda:=J^{\rm M}_\lambda-J_\lambda ,
\end{equation}
where $J^{\rm M}_\lambda$ is the Mellin count of $\zeta_\lambda$ (Theorem~\CGref{thm:mellin-counts}, which applies to every $\zeta_\lambda$ by Proposition~\CGref{prop:family-clock}) and $C_\lambda$ is the constant of Theorem~\CGref{thm:distributional-EF}.

\begin{proposition}[The first cell]\label{prop:first-cell}
Let $\lambda\in[0,1]$.
\begin{enumerate}[label=\textup{(\roman*)}]
\item $\mathfrak u_\lambda$ is entire, and for $\Re s>1$
\begin{equation}\label{eq:first-cell-split}
 \log\zeta_\lambda(s)=\mathfrak u_\lambda(s)+\langle dJ^*_\lambda,x^{-s}\rangle,\qquad \langle dJ^*_\lambda,x^{-s}\rangle:=s\int_2^\infty J_\lambda(x)\,x^{-s-1}\,dx .
\end{equation}
\item $\mathfrak u_\lambda(\sigma+it)=o(|t|)$ as $|t|\to\infty$, uniformly for $\sigma$ in compact sets.
\item $\mathfrak u_\lambda$ is constant if and only if $J_\lambda=0$ on $(1,2)$; in that case $\mathfrak u_\lambda\equiv C_\lambda$.
\end{enumerate}
For Riemann's count $J(x)=\sum_{p^m\le x}1/m$, which vanishes on $(1,2)$, the corresponding term is $0$.
\end{proposition}

\begin{proof}
(i) $J^{\rm M}_\lambda(x)x^{-c}\in L^2((1,\infty),dx/x)$ for $c>1$, so $J^{\rm M}_\lambda\in L^1(1,2)$ and the integral in \CGeqref{eq:first-cell} is entire in $s$.  By \CGeqref{eq:mellin-count}, $\log\zeta_\lambda(s)=s\int_1^2J^{\rm M}_\lambda x^{-s-1}dx+s\int_2^\infty(J_\lambda+C_\lambda)x^{-s-1}dx$, and $s\int_2^\infty C_\lambda x^{-s-1}dx=C_\lambda2^{-s}$.  The second integral is the pairing of $x^{-s}$ with the distributional derivative of $J^*_\lambda$, which vanishes on $(0,2)$.
(ii) $C_\lambda2^{-s}$ is bounded on vertical strips, and $\int_1^2J^{\rm M}_\lambda(x)x^{-s-1}dx=\int_0^{\log2}J^{\rm M}_\lambda(e^u)e^{-\sigma u}e^{-itu}du\to0$ as $|t|\to\infty$ by the Riemann--Lebesgue lemma, uniformly for $\sigma$ in compact sets.
(iii) For a constant $c$ one has $c=c\,2^{-s}+s\int_1^2c\,x^{-s-1}dx$, so $\mathfrak u_\lambda(s)-c=s\int_1^\infty H(x)x^{-s-1}dx$ with $H=(J^{\rm M}_\lambda-c)\mathbf 1_{(1,2)}+(C_\lambda-c)\mathbf 1_{[2,\infty)}$.  By the injectivity used in Theorem~\CGref{thm:mellin-counts}, $\mathfrak u_\lambda\equiv c$ if and only if $H=0$ almost everywhere, that is, $C_\lambda=c$ and $J^{\rm M}_\lambda=C_\lambda$ on $(1,2)$, which is $J_\lambda=0$ on $(1,2)$.
\end{proof}

\begin{lemma}[Mean-square convergence of the ordinate series]\label{lem:L2-partial-sums}
Let $\boldsymbol{\kappa}$ be a list with $|\kappa_k-\tau_k|\le C_\eta$ (in particular $\boldsymbol{\kappa}=(\gamma_k(\lambda))$ for any $\lambda\in[0,1]$), $\rho_k:=\frac12+i\kappa_k$, and let $I\subset(1,\infty)$ be compact. The partial sums $S_K(x):=\sum_{k\le K}\bigl(x^{\rho_k}/\rho_k+x^{\bar\rho_k}/\bar\rho_k\bigr)$ and $E_K(x):=\sum_{k\le K}\bigl(\Ei(\rho_k\log x)+\Ei(\bar\rho_k\log x)\bigr)$ converge in $L^2(I)$ as $K\to\infty$, with
\begin{equation}\label{eq:L2-partial-sums}
 \|S_{K'}-S_K\|_{L^2(I)}^2+\|E_{K'}-E_K\|_{L^2(I)}^2\ \le\ C_I\sum_{K<k\le K'}\frac{\log(\kappa_k+2)}{\kappa_k^2}\qquad(K<K'),
\end{equation}
where $C_I$ depends only on $I$ and $C_\eta$. Consequently: \textup{(a)} for every Lipschitz $\psi$ with compact support in $(1,\infty)$, the pairing $\langle dJ_X^{\rm M},\psi\rangle:=-\int J_X^{\rm M}\psi'\,dx$ equals the series obtained by integrating \CGeqref{eq:distributional-EF} term by term, which converges; if moreover $\psi'$ is of bounded variation, the series converges absolutely; \textup{(b)} $\lambda\mapsto J_\lambda$ and $\lambda\mapsto\partial_\lambda J_\lambda$ are continuous from $[0,1]$ to $L^2(I)$; \textup{(c)} the constant $C_\lambda:=J^{\rm M}_\lambda-J_\lambda$ is continuous on $[0,1]$, hence $\sup_\lambda|C_\lambda|<\infty$.
\end{lemma}

\begin{proof}
In $u=\log x$, $u\in\log I=:[u_1,u_2]$, the difference $S_{K'}-S_K$ is $e^{u/2}\sum_{K<k\le K'}\bigl(e^{i\kappa_ku}/\rho_k+e^{-i\kappa_ku}/\bar\rho_k\bigr)$. Let $\omega\in C_c^\infty(\mathbb R)$ be $\ge1$ on $[u_1,u_2]$. Expanding the square and using $|\hat\omega(\xi)|\ll_N(1+|\xi|)^{-N}$,
\begin{equation}\label{eq:L2-partial-sums-proof}
\begin{split}
 \int\omega(u)\Bigl|\sum_{K<k\le K'}\frac{e^{i\kappa_ku}}{\rho_k}\Bigr|^2du
 &\ll\sum_{K<k,k'\le K'}\frac{(1+|\kappa_k-\kappa_{k'}|)^{-2}}{\kappa_k\kappa_{k'}}\\
 &\ll\sum_{K<k\le K'}\frac1{\kappa_k^2}\sum_{j\ge0}\frac{\#\{k':j\le|\kappa_{k'}-\kappa_k|<j+1\}}{(1+j)^2}
 \ \ll\sum_{K<k\le K'}\frac{\log(\kappa_k+2)}{\kappa_k^2},
\end{split}
\end{equation}
because the tube gives $\#\{k':|\kappa_{k'}-t|\le1\}\ll\log(t+2)$ uniformly (the $\kappa_{k'}$ in $[t-1,t+1]$ have their clock points in $[t-1-C_\eta,t+1+C_\eta]$, an interval of length $2+2C_\eta$, which contains $O(\log(t+2))$ clock points), $\kappa_{k'}\asymp\kappa_k$ on the pairs that matter, and the pairs with $|\kappa_k-\kappa_{k'}|>1$ are summed with the weight $(1+|\kappa_k-\kappa_{k'}|)^{-2}$ against the same local count. The same bound holds for the $\Ei$ sums by Lemma~\CGref{lem:Ei}, since $\Ei(\rho u)=e^{\rho u}/(\rho u)+\varrho_k$ with $\sum_k|\varrho_k|$ locally uniformly convergent and $u\ge u_1>0$. This is \CGeqref{eq:L2-partial-sums}, and the right side tends to $0$ as $K\to\infty$ because $\kappa_k\asymp k/\log k$.

(a) For $\psi\in C_c^\infty(1,\infty)$ the identity is the convergence in $\mathcal D'(1,\infty)$ established in Theorem~\CGref{thm:distributional-EF}. For Lipschitz $\psi$ with compact support in $(1,\infty)$, $\psi'\in L^\infty\subset L^2(I)$, and both sides are continuous in the $L^2(I)$ topology of the partial sums: the left side is the pairing of an $L^2(I)$ limit with $\psi'$, and the terms of the right side are the pairings of the summands with $\psi'$, whose partial sums converge by \CGeqref{eq:L2-partial-sums} and the Cauchy--Schwarz inequality.  Lipschitz continuity alone gives convergence, not absolute convergence: the $k$-th term is $\int\psi'(x)\,\Re(x^{\rho_k}/\rho_k)\,dx$, of size $O(1/\kappa_k)$ in general.  When $\psi'$ is of bounded variation, one integration by parts in $u=\log x$ (against $e^{i\kappa_ku}$, with $\psi'$ absorbing the boundary terms) gains a factor $1/\kappa_k$ and the terms become $O(\kappa_k^{-2}\,\mathrm{TV}(\psi'))$, so the series converges absolutely; this is the case for the test functions $\mathfrak a_{ms}(y^j)\mathbf 1_{y\ge2}$ used below, whose derivatives are piecewise smooth.

(b) With $\boldsymbol{\kappa}=(\gamma_k(\lambda))$ and $\boldsymbol{\kappa}'=(\gamma_k(\lambda'))$, $|\gamma_k(\lambda)-\gamma_k(\lambda')|=|\lambda-\lambda'||\eta_k|\le C_\eta|\lambda-\lambda'|$; the difference of the $k$-th summands is $O(|\lambda-\lambda'|\,\log x\,x^{1/2}/\kappa_k)$ termwise, and the tail beyond $K$ is controlled uniformly in $\lambda$ by \CGeqref{eq:L2-partial-sums}. Hence $J_\lambda\to J_{\lambda'}$ in $L^2(I)$ as $\lambda\to\lambda'$.  For $\partial_\lambda J_\lambda$ the summands are $\eta_k\,\partial_\kappa(x^{\rho}/\rho)=i\eta_k(\log x-1/\rho_k)x^{\rho_k}/\rho_k$, of modulus at most $C_I/\kappa_k$ on $I$ because $|\eta_k|\le C_\eta$; the part $-i\eta_kx^{\rho_k}/\rho_k^2$ converges absolutely, and for the part $i\eta_k(\log x)x^{\rho_k}/\rho_k$ the proof of \CGeqref{eq:L2-partial-sums} applies verbatim with $\omega(u)u^2$ in place of $\omega(u)$ and the coefficients bounded by $C_\eta$, which gives the same uniform tail bound and hence the continuity of $\partial_\lambda J_\lambda$.

(c) By Plancherel on a fixed line $\Re s=c>1$ (proof of Proposition~\CGref{prop:cell-limit}), $\lambda\mapsto J^{\rm M}_\lambda$ is continuous into the weighted $L^2$, hence into $L^2(I)$. Pairing $J^{\rm M}_\lambda-J_\lambda$ with a fixed $\phi\in C_c^\infty(2,3)$, $\int\phi=1$, gives $C_\lambda$ as a difference of two continuous functions of $\lambda$.
\end{proof}

\begin{theorem}[The joining factor]\label{thm:joining-factor}
Let $\tilde h_\lambda:=\log\tilde\zeta_{\Z,\lambda}-\log\zeta_\lambda$ on the half-plane where $\tilde\zeta_{\Z,\lambda}$ converges absolutely (Theorem~\CGref{thm:atomized-continuation}).  Then
\begin{equation}\label{eq:joining-factor}
 \tilde\zeta_{\Z,\lambda}=\zeta_\lambda\,e^{\tilde h_\lambda},\qquad
 \tilde h_\lambda(s)=-\mathfrak u_\lambda(s)+\sum_{m\ge1}\frac1m\sum_{n\ge2}\mathfrak h_n(\lambda;ms),
\end{equation}
and, explicitly,
\begin{equation}\label{eq:joining-explicit}
 \tilde h_\lambda(s)=-\mathfrak u_\lambda(s)+\sum_{m\ge1}\frac1m\sum_{j\ge1}\frac{\mu(j)}j\int_2^\infty\mathfrak a_{ms}(y^j)\,\frac1{\log y}\Bigl(1-\frac2{\sqrt y}\Re\sum_{k\ge1}y^{i\gamma_k(\lambda)}-\frac1{y(y^2-1)}\Bigr)dy,
\end{equation}
the series of the integrated terms over $k$ converging absolutely.  Its derivative along the family is
\begin{equation}\label{eq:joining-derivative}
 \partial_\lambda\tilde h_\lambda(s)=-\partial_\lambda\mathfrak u_\lambda(s)+\sum_{m\ge1}\frac1m\sum_{j\ge1}\frac{\mu(j)}j\int_2^\infty\mathfrak a_{ms}(y^j)\,\frac2{\sqrt y}\sum_{k\ge1}\eta_k\sin\bigl(\gamma_k(\lambda)\log y\bigr)\,dy .
\end{equation}
Here $\partial_\lambda\mathfrak u_\lambda(s)=\mathcal V_\lambda(s)-\partial_\lambda\langle dJ^*_\lambda,x^{-s}\rangle$ by \CGeqref{eq:first-cell-split} and \CGeqref{eq:generator-zeta}.  Moreover $\mathfrak u_\lambda(\sigma)\to\log\varpi_\lambda$ as $\sigma\to+\infty$, and $\mathfrak u_\lambda\equiv\log\varpi_\lambda$ if $J_\lambda$ vanishes on $(1,2)$.
\end{theorem}

\begin{proof}
By the Euler product, $\log\tilde\zeta_{\Z,\lambda}(s)=\sum_{m\ge1}\frac1m\sum_n\tilde c_n(\lambda)n^{-ms}$.  By Proposition~\CGref{prop:first-cell}, $\log\zeta_\lambda(s)=\mathfrak u_\lambda(s)+\langle dJ^*_\lambda,x^{-s}\rangle$; finite inversion (Proposition~\CGref{prop:finite-inversion}, whose proof applies to $J^*_\lambda$) gives $dJ^*_\lambda=\sum_m\frac1m\,d[\pi_\lambda(x^{1/m})]$, and $\langle d[\pi_\lambda(x^{1/m})],x^{-s}\rangle=\langle d\pi_\lambda,y^{-ms}\rangle$.  Inserting $1=\sum_n\varphi_n$ on $[2,\infty)$ gives $\langle d\pi_\lambda,y^{-ms}\rangle=\sum_n\bigl(\tilde c_n(\lambda)n^{-ms}-\mathfrak h_n(\lambda;ms)\bigr)$, hence \CGeqref{eq:joining-factor}.  Letting $\sigma\to+\infty$, $\log\tilde\zeta_{\Z,\lambda}\to0$ (because $\tilde a_1=1$), $\log\zeta_\lambda\to\log\varpi_\lambda$ (Proposition~\CGref{prop:family-clock}) and the double series tends to $0$ by the bounds of Theorem~\CGref{thm:atomized-continuation}; so $\mathfrak u_\lambda(\sigma)\to\log\varpi_\lambda$, and a constant $\mathfrak u_\lambda$ (Proposition~\CGref{prop:first-cell}(iii)) equals $\log\varpi_\lambda$.  Atoms of $d\pi_\lambda$ at integers contribute nothing because the test functions vanish there.  Formula \CGeqref{eq:joining-explicit} is \CGeqref{eq:joining-factor} with $d\pi_\lambda$ written by \CGeqref{eq:distributional-EF} and the substitution $x=y^j$; the integrated terms in $k$ are $O_{m,j,s}(\gamma_k(\lambda)^{-2})$ because $\mathfrak a_{ms}(2)=0$, and the termwise integration against the Lipschitz test $\mathfrak a_{ms}(y^j)\mathbf 1_{y\ge2}$ is justified by Lemma~\CGref{lem:L2-partial-sums}(a).  The remaining double summation is absolutely controlled by $\sum_{m\ge1}m^{-1}\sum_{j\ge1}j^{-1}(mj|s|)^2 2^{-jm\sigma}<\infty$, so the exchanges of the $m$- and $j$-sums with the integrated ordinate sum are justified.  Differentiating in $\lambda$ with \CGeqref{eq:J-lambda-derivative} gives \CGeqref{eq:joining-derivative}.  The differentiability of $\mathfrak u_\lambda(s)$ in $\lambda$ rests on three facts, for a fixed $\sigma>1$.
(1) $|\mathcal V_\lambda(\sigma+it)|\ll_\sigma\log^2(|t|+2)$ uniformly in $\lambda\in[0,1]$: in \CGeqref{eq:generator}, $|A_{k,\lambda}+q(s)|=|\gamma_k(\lambda)-t+i(\sigma-\frac12)|\,|\gamma_k(\lambda)+t-i(\sigma-\frac12)|$, so each of the $O(\log(|t|+2))$ ordinates with $|\gamma_k(\lambda)-t|\le1$ contributes $O_\sigma(1)$, through $|\eta_k|\le C_\eta$, and the remaining ordinates contribute $O(\log^2(|t|+2))$ in total by summation against the count.
(2) $\lambda\mapsto\log\zeta_\lambda(\sigma+it)/(\sigma+it)$ is $C^1$ into $L^2(\R,dt)$ with derivative $\mathcal V_\lambda(\sigma+it)/(\sigma+it)$: by \CGeqref{eq:generator-zeta} and the mean-value theorem the difference quotients are dominated by $C_\sigma\log^2(|t|+2)/|\sigma+it|\in L^2(\R)$, and $\lambda\mapsto\mathcal V_\lambda(\sigma+it)$ is continuous for each $t$ by dominated convergence on the ordinate sum.
(3) By Plancherel, $J^{\rm M}_\lambda$ corresponds to $\log\zeta_\lambda(\sigma+it)/(\sigma+it)$ isometrically (Theorem~\CGref{thm:mellin-counts}), so $\lambda\mapsto J^{\rm M}_\lambda$ is $C^1$ in the corresponding weighted $L^2$ space; the constant $C_\lambda=J^{\rm M}_\lambda-J_\lambda$ is obtained by pairing with one fixed test function and is $C^1$ as well, by Lemma~\CGref{lem:L2-partial-sums}(b).
In particular $\|\partial_\lambda J_\lambda\|_{L^2(2,3)}$ is locally uniformly bounded.
\end{proof}

\begin{proposition}[Explicit cell defects]\label{prop:cell-explicit}
For every $n\ge2$, $\lambda\in[0,1]$ and $\Re s>1$,
\begin{equation}\label{eq:cell-explicit}
 \mathfrak h_n(\lambda;s)=\sum_{j\ge1}\frac{\mu(j)}j\int_2^\infty\varphi_n(y^j)\bigl(n^{-s}-y^{-js}\bigr)\,\frac1{\log y}\Bigl(1-\frac2{\sqrt y}\Re\sum_{k\ge1}y^{i\gamma_k(\lambda)}-\frac1{y(y^2-1)}\Bigr)dy,
\end{equation}
\begin{equation}\label{eq:cell-explicit-derivative}
 \partial_\lambda\mathfrak h_n(\lambda;s)=\sum_{j\ge1}\frac{\mu(j)}j\int_2^\infty\varphi_n(y^j)\bigl(n^{-s}-y^{-js}\bigr)\,\frac2{\sqrt y}\sum_{k\ge1}\eta_k\sin\bigl(\gamma_k(\lambda)\log y\bigr)\,dy ,
\end{equation}
the sums over $j$ being finite and the series of the integrated terms over $k$ converging absolutely, because the test functions are smooth with compact support.  At $\lambda=1$ the cell defect $\mathfrak h_n(1;\cdot)$ is the dispersion of the cell $n$ for $\zeta_{\RVM}$, computed from the clock alone; the absolutely continuous part of $d\pi_{\RVM}$ has bounded density there by Theorem~\CGref{thm:unit-cell}, while the atoms at $2^j$ are treated separately, so $\mathfrak h_n(1;\cdot)\equiv0$ holds only if this density vanishes on the support of $\varphi_n$.
\end{proposition}

\begin{proof}
Insert \CGeqref{eq:distributional-EF} for $d\pi_\lambda=\sum_j\frac{\mu(j)}j\,d[J^*_\lambda(x^{1/j})]$ into the definition \CGeqref{eq:atomization-kernel} and substitute $x=y^j$; the atoms of $d\pi_\lambda$ at integers, for instance the jump of $J^*_\lambda$ at $2$, contribute nothing because $n^{-s}-x^{-s}$ vanishes at $x=n$ and $\varphi_n$ vanishes at the other integers.  Differentiation in $\lambda$ uses \CGeqref{eq:J-lambda-derivative}; the factor $\log y$ cancels.  For the last statement, $\mathfrak h_n(1;\cdot)\equiv0$ forces $\varphi_n\,d\pi_{\RVM}=\tilde c_n(1)\delta_n$ by the injectivity used in Proposition~\CGref{prop:cell-limit}, and, away from the explicit atoms at $2^j$, an absolutely continuous measure equals a multiple of $\delta_n$ only if both vanish; at the powers of $2$ the boundary atom is handled by \CGeqref{eq:cell-mass-power-two}.
\end{proof}

In the form of the path, the final step requires that in every cell the flow of the displacements cancel exactly the dispersion of the clock, \CGeqref{eq:closing-cells}.

\begin{proposition}[The limit $\lambda\to0$ and the meaning of a vanishing cell]\label{prop:cell-limit}
For every $n\ge2$ and $\Re s>1$, $\mathfrak h_n(\lambda;s)\to\mathfrak h_n(0;s)$ and $\tilde c_n(\lambda)\to\tilde c_n(0)$ as $\lambda\to0^+$, and
\begin{equation}\label{eq:cell-flow}
 \mathfrak h_n(0;s)=\mathfrak h_n(1;s)-\int_0^1\partial_\lambda\mathfrak h_n(\lambda;s)\,d\lambda .
\end{equation}
Moreover $\mathfrak h_n(0;s)=0$ for all $\Re s>1$ if and only if $\varphi_n\,d\pi_{Z1}=\tilde c_n(0)\,\delta_n$.
\end{proposition}

\begin{proof}
The test function $\varphi_n(x)(n^{-s}-x^{-s})$ is smooth with compact support, and the integrated terms over $k$ are dominated uniformly in $\lambda$ because $|\gamma_k(\lambda)-\gamma_k|\le C_\eta$; dominated convergence gives the limits for cells away from the powers of $2$.  For $n=2^j$, the boundary term \CGeqref{eq:cell-mass-power-two} also converges: by Plancherel, $\log\zeta_\lambda(\sigma+it)/(\sigma+it)\to\log\zeta_{Z1}(\sigma+it)/(\sigma+it)$ in $L^2(dt)$ on a fixed line $\sigma>1$, hence $J^{\rm M}_\lambda\to J^{\rm M}_{Z1}$ in the corresponding weighted $L^2$ space, and therefore $J_\lambda\to J_{Z1}$ in $L^2(2,3)$ by Lemma~\CGref{lem:L2-partial-sums}(b),(c).  This proves all stated limits and \CGeqref{eq:cell-flow}.  Since $\mathfrak h_n(0;s)=\tilde c_n(0)n^{-s}-\langle\varphi_n\,d\pi_{Z1},x^{-s}\rangle$, its vanishing for all $s$ says that the compactly supported distribution $\varphi_n\,d\pi_{Z1}$ has the Mellin transform of $\tilde c_n(0)\delta_n$; the Mellin transform is injective on such distributions.
\end{proof}

\subsection{Mean square of the cell masses}
\label{subsec:mean-square}

From now on the cells are translates of one function: $\varphi_n(x)=\varphi(x-n)$, where $\varphi=\beta_0/\sum_{m\in\Z}\beta_0(\cdot-m)$ for a fixed bump $\beta_0\in C_c^\infty(-1,1)$ positive on $[-\frac12,\frac12]$.  Then $\sum_{n\ge2}\varphi_n=1$ on $[2,\infty)$, $\varphi_n(m)=\delta_{nm}$, and $\|\varphi_n^{(A)}\|_\infty=\|\varphi^{(A)}\|_\infty=:C_A$ for every $A\ge0$, uniformly in $n$.  For $n\ge2$ put
\begin{equation}\label{eq:VW}
 V_n(\gamma):=\sum_{j\ge1}\frac{\mu(j)}j\int_2^\infty\varphi_n(y^j)\,y^{-1/2}e^{i\gamma\log y}\,dy,\qquad
 W_n(\gamma):=\sum_{j\ge1}\frac{\mu(j)}j\int_2^\infty\varphi_n(y^j)\,\frac{y^{-1/2}e^{i\gamma\log y}}{\log y}\,dy ,
\end{equation}
so that $W_n'=iV_n$; both are rapidly decreasing in $\gamma$ when $n$ is not a power of $2$ (for $n=2^j$ the term $j$ carries the cut of $J^*_\lambda$ at $y=2$ and is treated in the proof of Theorem~\CGref{thm:mean-square-cells}).  For $n$ not a power of $2$, the density of $d\pi_\lambda$ used in \CGeqref{eq:cell-explicit} and the derivative \CGeqref{eq:J-lambda-derivative} give, with $\eta_k=\tau_k-\gamma_k$:
\begin{equation}\label{eq:cell-mass-W}
 \tilde c_n(\lambda)=\mathfrak m_n-2\,\Re\sum_{k\ge1}W_n\bigl(\gamma_k(\lambda)\bigr),\qquad
 \partial_\lambda\tilde c_n(\lambda)=2\sum_{k\ge1}\eta_k\,\Im V_n\bigl(\gamma_k(\lambda)\bigr),
\end{equation}
where $\mathfrak m_n$ collects the terms that do not involve the ordinates and is $O(1)$.  If $n=2^j$, the $j$th M\"obius-inversion term crosses the cut at $y=2$ and contributes instead the additional boundary term
\begin{equation}\label{eq:cell-mass-power-two}
 -\frac{\mu(j)}j\int_2^3 J_\lambda(y)\,\frac d{dy}\varphi_n(y^j)\,dy,
\end{equation}
whose $\lambda$-derivative is obtained by replacing $J_\lambda$ by $\partial_\lambda J_\lambda$.

\begin{lemma}[Mean value over the moving ordinates]\label{lem:mean-value-ordinates}
There is $C$ such that for all $\lambda\in[0,1]$, $T\ge2$, $u_0\in\R$ and complex $a_k$ with $|a_k|\le1$,
\begin{equation}\label{eq:mean-value-ordinates}
 \int_{u_0-1}^{u_0+1}\Bigl|\sum_{\gamma_k(\lambda)\le T}a_k\,e^{i\gamma_k(\lambda)u}\Bigr|^2du\le C\,T\log^2T .
\end{equation}
\end{lemma}

\begin{proof}
Let $\omega_0(u)=C_0\bigl(\sin(\frac{u-u_0}2)/\frac{u-u_0}2\bigr)^2$ with $C_0$ such that $\omega_0\ge1$ on $|u-u_0|\le1$; the Fourier transform of $\omega_0$ is bounded and supported in $[-1,1]$.  Expanding the square,
\begin{align}
\int \omega_0(u)\Bigl|\sum_ka_ke^{i\gamma_k(\lambda)u}\Bigr|^2du&=\sum_{k,k'}a_k\bar a_{k'}\,e^{i(\gamma_k(\lambda)-\gamma_{k'}(\lambda))u_0}\,\widehat{\omega_0}\bigl(\gamma_k(\lambda)-\gamma_{k'}(\lambda)\bigr)\nonumber\\
&\ll\#\{(k,k'):\gamma_k(\lambda)\le T,\ |\gamma_k(\lambda)-\gamma_{k'}(\lambda)|\le1\}.\label{eq:fejer-pairs}
\end{align}
Since $|\gamma_k(\lambda)-\tau_k|\le C_\eta$, every interval $[t,t+1]$ contains $O(\log(t+2))$ of the $\gamma_k(\lambda)$, uniformly in $\lambda$, because the corresponding $\tau_k$ lie in $[t-C_\eta,t+1+C_\eta]$, an interval of length $1+2C_\eta$ containing $O(\log(t+2))$ clock points.  The number of pairs is therefore $O(T\log^2T)$.
\end{proof}

\begin{theorem}[Mean square of the cell masses along the family]\label{thm:mean-square-cells}
For every $\eps>0$ there is $C_\eps$ such that for all $X\ge2$ and all $\lambda\in[0,1]$
\begin{equation}\label{eq:mean-square-cells}
 \sum_{X<n\le2X}|\tilde c_n(\lambda)|^2\le C_\eps X^{1+\eps},\qquad
 \sum_{X<n\le2X}|\partial_\lambda\tilde c_n(\lambda)|^2\le C_\eps X^{1+\eps}.
\end{equation}
In particular the same bound already includes $\lambda=0$: $\sum_{X<n\le2X}|\tilde c_n(0)|^2\le C_\eps X^{1+\eps}$.
\end{theorem}

\begin{proof}
Fix $\eps$, $X$, $\lambda$ and put $T=X^{1+\eps}$.  For $n$ not a power of $2$, \CGeqref{eq:cell-mass-W} gives $\tilde c_n(\lambda)-\mathfrak m_n=-2\Re\sum_{j}\frac{\mu(j)}jA_{n,j}$, with
\begin{equation}\label{eq:A-nj}
 A_{n,j}=\sum_k\int_2^\infty\varphi_n(y^j)\,\frac{y^{-1/2+i\gamma_k(\lambda)}}{\log y}\,dy ,\qquad 1\le j\le J:=\Bigl\lfloor\frac{\log(2X+1)}{\log2}\Bigr\rfloor .
\end{equation}
\emph{Large ordinates.}  In $u=\log y$ the function $\varphi_n(e^{ju})e^{u/2}/u$ has support of length $O((jn)^{-1})$, size $O(n^{1/(2j)})$ and $r$-th derivative $O_r((jn)^rn^{1/(2j)})$: by Fa\`a di Bruno's formula the $r$-th derivative of $\varphi_n(e^{ju})$ is a sum of terms $\varphi_n^{(m)}(e^{ju})(je^{ju})^mj^{r-m}$, $m\le r$, and $\|\varphi_n^{(m)}\|_\infty=C_m$ does not depend on $n$; the derivatives of $e^{u/2}/u$ are $O_r(e^{u/2})$ because $u\ge\log2$.  Integrating by parts $A$ times, each term with $\gamma_k(\lambda)>T$ is $O_A\bigl((jn/\gamma_k(\lambda))^An^{1/(2j)-1}\bigr)$, and their sum over $k$ is $O(X^{-1})$ once $A>(3+\eps)/\eps$.
\emph{The cell $n=2^j$.}  Each block contains exactly one $n=2^j$; for it the term $j$ carries the cut of $J^*_\lambda$ at $y=2$ and is treated directly.  It equals $-\frac{\mu(j)}j\int_2^3J_\lambda(y)\,\frac d{dy}\varphi_n(y^j)\,dy$; the weight has $L^2$ norm $O((jn)^{1/2})$, and $\|J_\lambda\|_{L^2(2,3)}$ is bounded uniformly in $\lambda$ (Theorem~\CGref{thm:mellin-counts} with the uniform bound of Proposition~\CGref{prop:family-clock} for $J^{\rm M}_\lambda$, and Lemma~\CGref{lem:L2-partial-sums}(c) for $\sup_\lambda|C_\lambda|$).  This term is therefore $O((n/j)^{1/2})$ and contributes $O(X)$ to the sum of squares; the same holds for $\partial_\lambda\tilde c_n$, with $\partial_\lambda J_\lambda$ in place of $J_\lambda$.
\emph{Small ordinates.}  Let $\mathcal S_{n,j,\lambda}(y)=y^{-1/2}(\log y)^{-1}\sum_{\gamma_k(\lambda)\le T}y^{i\gamma_k(\lambda)}$ and let $I_{n,j}$ be the support of $y\mapsto\varphi_n(y^j)$, of length $O(X^{1/j-1})$.  By Cauchy--Schwarz in the cell,
\begin{equation}\label{eq:cell-CS}
 \Bigl|\int\varphi_n(y^j)\mathcal S_{n,j,\lambda}(y)\,dy\Bigr|^2\le\|\varphi\|_\infty^2\,|I_{n,j}|\int_{I_{n,j}}|\mathcal S_{n,j,\lambda}|^2dy .
\end{equation}
For fixed $j$ the intervals $I_{n,j}$, $X<n\le2X$, cover each point at most twice and lie in $[(X-1)^{1/j},(2X+1)^{1/j}]$, an interval of length at most $2$ in $u=\log y$.  With $dy/y=du$ and Lemma~\CGref{lem:mean-value-ordinates} ($a_k=1$),
\begin{equation}\label{eq:cell-block}
 \sum_{X<n\le2X}\Bigl|\int\varphi_n(y^j)S\,dy\Bigr|^2\ll X^{1/j-1}\Bigl(\frac{j}{\log X}\Bigr)^2\int_{u_0-1}^{u_0+1}\Bigl|\sum_{\gamma_k(\lambda)\le T}e^{i\gamma_k(\lambda)u}\Bigr|^2du\ll j^2X^{1/j+2\eps}.
\end{equation}
\emph{Assembly.}  By Minkowski's inequality over $j$, the $\ell^2$ norm over $X<n\le2X$ of $\tilde c_n(\lambda)-\mathfrak m_n$ is $\ll\sum_{j\le J}\bigl(X^{1/(2j)+\eps}+X^{-1/2}\bigr)\ll X^{1/2+\eps}$, and $\sum_{X<n\le2X}\mathfrak m_n^2\ll X$.  This is the first bound with $3\eps$ in place of $\eps$.  Away from powers of $2$, the argument for $\partial_\lambda\tilde c_n$ is the same, by \CGeqref{eq:cell-mass-W}, with coefficients $a_k=\eta_k/C_\eta$ in Lemma~\CGref{lem:mean-value-ordinates} and without the factor $(\log y)^{-1}$.
\end{proof}

\begin{corollary}[Abscissa $1$ for every member]\label{cor:abscissa-one}
For every $\lambda\in[0,1]$ and $\sigma>1$,
\begin{equation}\label{eq:abscissa-one}
 \sum_{n\ge2}|\tilde c_n(\lambda)|\,n^{-\sigma}<\infty ;
\end{equation}
hence $\tilde\zeta_{\Z,\lambda}$ and $1/\tilde\zeta_{\Z,\lambda}$ are absolutely convergent Dirichlet series in $\Re s>1$, for every $\lambda$, including the endpoint $\lambda=0$.  Moreover, uniformly in $\lambda$, $\#\{n\le X:|\tilde c_n(\lambda)|>n^{\delta}\}\ll_\eps X^{1-2\delta+\eps}$ for every $\delta>0$.
\end{corollary}

\begin{proof}
By Cauchy--Schwarz on dyadic blocks, $\sum_{X<n\le2X}|\tilde c_n(\lambda)|\le(X\cdot C_\eps X^{1+\eps})^{1/2}$, and summing over $X=2^i$ gives \CGeqref{eq:abscissa-one} for $\sigma>1+\eps/2$, with $\eps$ arbitrary.  Then $\log\tilde\zeta_{\Z,\lambda}=\sum_m\frac1m\sum_n\tilde c_n(\lambda)n^{-ms}$ converges absolutely in $\Re s>1$, and the exponentials of $\pm$ this series are absolutely convergent Dirichlet series there.  The last statement is Chebyshev's inequality on dyadic blocks.
\end{proof}

The abscissa $\frac32$ obtained below from the primitive $\pi_\lambda-\mathfrak L$ (Theorem~\CGref{thm:atomized-continuation}(ii)) is not the true one: a bound on the size of $\pi_\lambda-\mathfrak L$ does not see the cancellation of the density of $d\pi_\lambda$ in mean square, which Theorem~\CGref{thm:mean-square-cells} captures.  Absolute convergence in $\Re s>1$ is exactly what Hamburger's theorem requires, and Knopp's obstruction (Subsection~\CGref{subsec:hamburger-table}) does not arise.

\subsection{Analytic continuation of the atomized family}
\label{subsec:atomized-continuation}

Let $\mathfrak L(x)=\int_1^x(1-y^{-1})\,dy/\log y$; then $\log\frac s{s-1}=\int_1^\infty x^{-s}\,d\mathfrak L(x)$ for $\Re s>1$ (both sides have derivative $\frac1s-\frac1{s-1}$ and tend to $0$ as $s\to+\infty$), and $\mathfrak L(x)=\li(x)-\log\log x-\gamma_E$.

\begin{theorem}[Continuation of the atomized family]\label{thm:atomized-continuation}
Let $\lambda\in[0,1]$.
\begin{enumerate}[label=\textup{(\roman*)}]
\item For every $c>\frac12$, $(\pi_\lambda-\mathfrak L)\,x^{-c}\in L^2\bigl((2,\infty),dx/x\bigr)$.
\item $|\tilde c_n(\lambda)|\le C(1+I_n^{1/2})$ with $I_n$ as below, and $\sum_n|\tilde c_n(\lambda)|n^{-\sigma}<\infty$ for $\sigma>\frac32$; by Corollary~\CGref{cor:abscissa-one} this holds for $\sigma>1$, and $\tilde\zeta_{\Z,\lambda}$ is an absolutely convergent Dirichlet series in $\Re s>1$.
\item $\tilde h_\lambda$ extends holomorphically to $\Re s>\frac12$.  Hence $\tilde\zeta_{\Z,\lambda}=\zeta_\lambda e^{\tilde h_\lambda}$ is meromorphic in $\Re s>\frac12$, with a single pole at $s=1$ and no zeros.
\item For $\lambda=1$, $\tilde h_1$ extends holomorphically to $\Re s>0$, and the zeros of $\tilde\zeta_{\Z,1}$ in $\Re s>0$ are exactly $\frac12\pm i\tau_k$.
\end{enumerate}
\end{theorem}

\begin{proof}
(i) By Proposition~\CGref{prop:family-clock}, for $c>\frac12$ there is $C_c$ with $|\log\zeta_\lambda(s)-\log\frac s{s-1}|\le C_c(1+\log(|t|+2))$ on $\Re s\ge c$; the function $\log\zeta_\lambda-\log\frac s{s-1}$ is holomorphic on $\Re s>\frac12$, where $\zeta_\lambda$ has no zeros, and its pole at $1$ is removed.  It equals $s\int_1^\infty(J_\lambda^{\rm M}-\mathfrak L)x^{-s-1}dx$ for $\Re s>1$.  The argument of Theorem~\CGref{thm:mellin-counts}, applied on the lines $\Re s=c>\frac12$, gives $(J^{\rm M}_\lambda-\mathfrak L)x^{-c}\in L^2$.  The constant $C_\lambda=J^{\rm M}_\lambda-J_\lambda$, the truncation at $2$, the terms $j\ge2$ of the M\"obius inversion and $\mathfrak L(x)-\mathfrak L(x^{1/j})$ preserve the property for $c>\frac12$: with the substitution $y=x^{1/j}$, $\int_1^\infty|(J^{\rm M}_\lambda-\mathfrak L)(x^{1/j})|^2x^{-2c}\frac{dx}x=j\int_1^\infty|(J^{\rm M}_\lambda-\mathfrak L)(y)|^2y^{-2cj}\frac{dy}y<\infty$ because $cj>\frac12$, and $\mathfrak L(x^{1/j})\ll x^{1/j}/\log x$ lies in the weighted space for $2c>1$.

(ii) Put $E=\pi_\lambda-\mathfrak L$ and $I_n=\int_{n-1}^{n+1}|E|^2dx$.  Since $\int\varphi_n'=0$, $\tilde c_n(\lambda)=\int\varphi_n\,d\mathfrak L-\int E\varphi_n'\,dx$, so $|\tilde c_n(\lambda)|\le C(1+I_n^{1/2})$.  By (i), $\sum_nI_nn^{-2c-1}<\infty$, and Cauchy--Schwarz gives $\sum_nI_n^{1/2}n^{-\sigma}<\infty$ for $\sigma>c+1$.

(iii) With $\psi_n=\varphi_n(x)(n^{-s}-x^{-s})$ one has $|\psi_n'|\le C|s|n^{-\sigma-1}$ and $\psi_n$ vanishes at the ends of the cell, so
\begin{equation}\label{eq:cell-defect-bound}
 |\mathfrak h_n(\lambda;s)|\le C|s|\,n^{-\sigma-1}\bigl(1+I_n^{1/2}\bigr),\qquad \sum_nI_n^{1/2}n^{-\sigma-1}<\infty\quad(\sigma>c).
\end{equation}
Choosing $c\in(\frac12,\sigma)$, the series \CGeqref{eq:joining-factor} converges locally uniformly on $\Re s>\frac12$; the terms $m\ge2$ are evaluated at $\Re(ms)\ge2\sigma$.  Since $\mathfrak u_\lambda$ is entire (Proposition~\CGref{prop:first-cell}), $\tilde h_\lambda$ is holomorphic on $\Re s>\frac12$, and $\zeta_\lambda$ has no zeros there.

(iv) By Corollary~\CGref{cor:cn-bounded}, $\pi_{\RVM}$ oscillates by $O(1)$ on every unit cell; integrating by parts, $|\mathfrak h_n(1;s)|\le C|s|n^{-\sigma-1}$, and the series converges for $\sigma>0$.  The factor $e^{\tilde h_1}$ has no zeros, so the zeros of $\tilde\zeta_{\Z,1}$ in $\Re s>0$ are those of $\zeta_{\RVM}$.
\end{proof}

\begin{theorem}[Closing criterion; via Hamburger's theorem \cite{Hamburger1921}]\label{thm:closing}
The following are equivalent.
\begin{enumerate}[label=\textup{(\roman*)}]
\item (D), that is, $\zeta_{Z1}=\zeta$.
\item $\tilde h_0\equiv0$.
\item $J_{Z1}=0$ on $(1,2)$, and for every $n\ge2$ and every $\Re s>1$,
\begin{equation}\label{eq:closing-cells}
 \int_0^1\partial_\lambda\mathfrak h_n(\lambda;s)\,d\lambda=\mathfrak h_n(1;s).
\end{equation}
\end{enumerate}
In this case $\tilde c_n(0)=\mathbf 1_{\rm prime}(n)$, $\tilde a_N(0)=1$ for all $N$, and $\varpi_{Z1}=1$.
\end{theorem}

\begin{proof}
(ii)$\Rightarrow$(i): $\tilde h_0\equiv0$ gives $\zeta_{Z1}=\tilde\zeta_{\Z,0}=\exp\bigl(\sum_m\frac1m\sum_n\tilde c_n(0)n^{-ms}\bigr)$, an absolutely convergent Dirichlet series in $\Re s>1$ by Corollary~\CGref{cor:abscissa-one}.  (iii)$\Rightarrow$(ii): by \CGeqref{eq:cell-flow} every $\mathfrak h_n(0;\cdot)$ vanishes, and $\mathfrak u_0\equiv\log\varpi_{Z1}$ (Theorem~\CGref{thm:joining-factor}), so $\tilde h_0\equiv-\log\varpi_{Z1}$ is constant; then $\zeta_{Z1}=\varpi_{Z1}\tilde\zeta_{\Z,0}$ is an absolutely convergent Dirichlet series in $\Re s>1$ (Corollary~\CGref{cor:abscissa-one}), Theorem~\CGref{thm:Hamburger-Z1} gives $\zeta_{Z1}=\zeta$, and comparing the first coefficients, $\varpi_{Z1}=1$.  (i)$\Rightarrow$(iii): if $\zeta_{Z1}=\zeta$, Theorem~\CGref{thm:cell-form} gives $J^{\rm M}_{Z1}=J$; and the explicit-formula count $J_{Z1}$ built from the ordinates (Theorem~\CGref{thm:distributional-EF}) also equals $J$, because every zero of $\zeta$ is then $\frac12\pm i\gamma_k$ and the series defining $J_{Z1}$ is von Mangoldt's explicit formula for $J$ \cite{Edwards1974}; so $J_{Z1}$ vanishes on $(1,2)$, $d\pi_{Z1}=\sum_p\delta_p$, $\tilde c_n(0)=\mathbf 1_{\rm prime}(n)$, every $\mathfrak h_n(0;\cdot)$ vanishes, and \CGeqref{eq:cell-flow} gives \CGeqref{eq:closing-cells}.
\end{proof}

\begin{remark}[The pointwise size of the cell masses]\label{rem:A6-scale}
Theorem~\CGref{thm:closing} needs no pointwise bound on $\tilde c_n(0)$: its proof uses the cell masses only through \CGeqref{eq:abscissa-one}.  The bound $\tilde c_n(0)=O(n^{\eps})$ for every $n$ is not proved directly; Corollary~\CGref{cor:abscissa-one} restricts the exceptions to a set of $n\le X$ of size $O(X^{1-2\delta+\eps})$ for the threshold $n^{\delta}$, uniformly in $\lambda$, and under the closure every cell mass is bounded, $\tilde c_n(0)=\mathbf 1_{\rm prime}(n)$.  For $n$ not a power of $2$, \CGeqref{eq:cell-mass-W} gives the trivial pointwise bound $\tilde c_n(\lambda)=O(n^{1/2+\eps})$; the mean square \CGeqref{eq:mean-square-cells} shows that the typical cell mass is $O(n^{\eps})$.
\end{remark}

\subsection{The two endpoints of the atomized family}
\label{subsec:atomized-endpoints}

\begin{proposition}[Atomization at the clock endpoint]\label{prop:atomized-clock}
At $\lambda=1$:
\begin{enumerate}[label=\textup{(\roman*)}]
\item $\tilde c_n(1)=O(1)$, and $\tilde\zeta_{\Z,1}$ is an absolutely convergent Dirichlet series in $\Re s>1$;
\item $\tilde h_1(\sigma)\to-\log\varpi_{\RVM}=\log L$ as $\sigma\to+\infty$; this limit is $\ne0$ exactly when $\varpi_{\RVM}\ne1$, and it is $>0$ by Theorem~\CGref{thm:varpi-less-one};
\item $\tilde h_1$ is not constant: for every $\sigma>1$, $\zeta_{\RVM}(\sigma+it)\to\varpi_{\RVM}$ as $|t|\to\infty$, whereas $\tilde\zeta_{\Z,1}(\sigma+it)$ has no limit as $|t|\to\infty$.
\end{enumerate}
\end{proposition}

\begin{proof}
(i) By Corollary~\CGref{cor:cn-bounded}, $\pi_{\RVM}$ oscillates by $O(1)$ on every unit cell, and $\tilde c_n(1)=-\int\pi_{\RVM}\varphi_n'\,dx$ with $\int\varphi_n'=0$.  (ii) is \CGeqref{eq:joining-factor} with $\sigma\to+\infty$; the limit is $\ne0$ if and only if $\varpi_{\RVM}\ne1$.  (iii) The first limit is \CGeqref{eq:vertical-limit}.  If $\tilde\zeta_{\Z,1}(\sigma+it)$ had a limit $\ell$ as $|t|\to\infty$, the mean-square argument \CGeqref{eq:mean-square}--\CGeqref{eq:mean-square-limit} applied to $\tilde\zeta_{\Z,1}-\ell$ on the line $\Re s=\sigma$ would give $\tilde a_N(1)=0$ for $N\ge2$, so $\tilde\zeta_{\Z,1}\equiv1$, which contradicts its pole at $s=1$ (Theorem~\CGref{thm:atomized-continuation}(iii)).  If $\tilde h_1$ were constant, $\zeta_{\RVM}$ would be a constant multiple of an absolutely convergent Dirichlet series, which Theorem~\CGref{thm:no-Dirichlet-RVM} excludes.
\end{proof}

The two families thus keep complementary structures and are joined by a zero-free factor:
\begin{equation}\label{eq:duality}
 \underbrace{\zeta_\lambda}_{\substack{\text{functional equation}\\ \text{zeros }\frac12\pm i\gamma_k(\lambda)}}\cdot\ e^{\tilde h_\lambda}
 \ =\ \underbrace{\tilde\zeta_{\Z,\lambda}}_{\substack{\text{Dirichlet series}\\ \text{Euler product}}},\qquad
 \tilde h_1\ \text{not constant},\qquad (D)\iff\tilde h_0\equiv0 .
\end{equation}
At $\lambda=1$ the joining factor carries the whole difference between the clock and an Euler product; Theorem~\CGref{thm:closing} states that this difference disappears at $\lambda=0$.

The reciprocal of the atomized family is again an Euler product:
\begin{equation}\label{eq:atomized-mobius}
 \frac1{\tilde\zeta_{\Z,\lambda}(s)}=\prod_{n\ge2}\bigl(1-n^{-s}\bigr)^{\tilde c_n(\lambda)}=\sum_{N\ge1}\frac{\tilde\mu_\lambda(N)}{N^s},\qquad
 \tilde\mu_\lambda(N)=\sum_{\prod_nn^{r_n}=N}\ \prod_{n:\,r_n\ge1}(-1)^{r_n}\binom{\tilde c_n(\lambda)}{r_n},
\end{equation}
absolutely convergent in $\Re s>1$ for every $\lambda$ (Corollary~\CGref{cor:abscissa-one}), because it is the exponential of the absolutely convergent series $-\sum_m\frac1m\sum_n\tilde c_n(\lambda)n^{-ms}$.  Thus $\tilde\mu_\lambda(p)=-\tilde c_p(\lambda)$, and the atomized Mertens functions
\begin{equation}\label{eq:atomized-mertens}
 \tilde M_\lambda(x):=\sum_{N\le x}\tilde\mu_\lambda(N),\qquad \frac1{\tilde\zeta_{\Z,\lambda}}=e^{-\tilde h_\lambda}F_\lambda ,
\end{equation}
exist for every $\lambda\in[0,1]$.  Under the conditions of Theorem~\CGref{thm:closing}, $\tilde c_n(0)=\mathbf 1_{\rm prime}(n)$, so $\tilde\mu_0=\mu$ and $\tilde M_0=M$.

\begin{proposition}[Mellin form of the endpoint Mertens function]\label{prop:Mellin-Mertens}
There are $c_1\ge2$ and a unique measurable $B_{Z1}$ on $(0,\infty)$, vanishing on $(0,1)$, with $x^{-c}B_{Z1}\in L^2\bigl((1,\infty),dx/x\bigr)$ and
\begin{equation}\label{eq:Mellin-Mertens}
 F_{Z1}(s)=s\int_1^\infty B_{Z1}(x)\,x^{-s-1}\,dx=\exp\Bigl(-\int_{[1,\infty)}x^{-s}\,dJ^{\rm M}_{Z1}(x)\Bigr)\qquad(\Re s>c_1),
\end{equation}
for every $c>c_1$.  Moreover $(D)$ holds if and only if $B_{Z1}=M$ almost everywhere; and if $F_{Z1}$ has an absolutely convergent Dirichlet expansion \CGeqref{eq:Z1-reciprocal-Dirichlet}, then $B_{Z1}=M_{Z1}$.
\end{proposition}

\begin{proof}
Put $A=2(\pi+2)C_{\mathfrak e}$.  By \CGeqref{eq:Z1-growth}, for $\sigma'\ge\max(4A,R_{\rm end})$, $|F_{Z1}(s)|\le C(|t|+\sigma)^{1/4}$; take $c_1=\max(2,\frac12+4A,\frac12+R_{\rm end})$.  Then $\int_{\R}|F_{Z1}(\sigma+it)|^2|\sigma+it|^{-2}dt\le C'\sigma^{-1/2}$ for $\sigma>c_1$, and the argument of Theorem~\CGref{thm:mellin-counts} gives $B_{Z1}$, its vanishing on $(0,1)$, its independence of $c$ and its uniqueness; the second equality in \CGeqref{eq:Mellin-Mertens} is Theorem~\CGref{thm:mellin-counts}.  If $\zeta_{Z1}=\zeta$, then $F_{Z1}(s)=1/\zeta(s)=s\int_1^\infty M(x)x^{-s-1}dx$ and uniqueness gives $B_{Z1}=M$.  Conversely, $B_{Z1}=M$ gives $F_{Z1}=1/\zeta$ on $\Re s>c_1$, hence $\zeta_{Z1}=\zeta$ by analytic continuation.  Under \CGeqref{eq:Z1-reciprocal-Dirichlet}, partial summation gives $F_{Z1}(s)=s\int_1^\infty M_{Z1}(x)x^{-s-1}dx$, and uniqueness gives $B_{Z1}=M_{Z1}$.
\end{proof}

\subsection{Hamburger's hypotheses at the two endpoints}
\label{subsec:hamburger-table}

Hamburger's theorem \cite{Hamburger1921}, in the form of \cite[\S2.13]{Titchmarsh1986} (see also \cite{Siegel1922,BochnerChandrasekharan1956}), has four hypotheses: (H0) $f=G/P$ with $G$ entire of finite order and $P$ a polynomial; (H1) $f(s)=\sum_na_nn^{-s}$ absolutely for $\sigma>1$, with arbitrary complex $a_n$; (H2) $f(s)\Gamma(\frac s2)\pi^{-s/2}=g(1-s)\Gamma(\frac{1-s}2)\pi^{-(1-s)/2}$; (H3) $g(1-s)=\sum_nb_nn^{s-1}$ absolutely for $\sigma<-\alpha<0$.  It contains no hypothesis on the location of zeros, and its conclusion is $f=C\zeta$.  Absolute convergence for $\sigma>1$ in (H1) cannot be weakened to some half-plane: Knopp constructed infinitely many linearly independent solutions of (H0), (H2), (H3) that are Dirichlet series absolutely convergent only in some half-plane \cite{Knopp1994,Nakamura2020}.  Within the Selberg class, the functions of degree $1$ are classified in \cite{KaczorowskiPerelli1999}.  The table records the status at the two endpoints and is updated as the programme advances.

\begin{center}
\small
\begin{tabular}{p{0.06\textwidth}p{0.40\textwidth}p{0.44\textwidth}}
\toprule
 & $\zeta_{Z1}$ (continuous limit) & $\tilde\zeta_{\Z,0}$ (atomized limit)\\
\midrule
H0 & proved: $(s-1)\zeta_{Z1}$ entire of order $1$ (Proposition~\CGref{prop:Z1-class}) & open: meromorphic in $\Re s>\frac12$ only (Theorem~\CGref{thm:atomized-continuation})\\
H1 & open: this is (D) & proved in $\Re s>1$ (Corollary~\CGref{cor:abscissa-one})\\
H2 & proved, with $g=f$ (Proposition~\CGref{prop:Z1-class}) & open\\
H3 & follows from H1 & open\\
zeros & no hypothesis & no hypothesis\\
\bottomrule
\end{tabular}
\end{center}

\noindent The two columns meet when $\tilde h_0\equiv0$: then $\zeta_{Z1}=\tilde\zeta_{\Z,0}$ has H0--H3, and Hamburger's theorem gives $\zeta_{Z1}=\zeta$ (Theorem~\CGref{thm:closing}).  It suffices that the joining factor be symmetric (Theorem~\CGref{thm:pair-criterion}), or that $\zeta_{Z1}$ have a positive representing measure (Theorem~\CGref{thm:positivity}).

The parallel between the two families, topic by topic of the paper, is the following.

\begingroup\small
\begin{longtable}{>{\raggedright\arraybackslash}p{0.20\textwidth}>{\raggedright\arraybackslash}p{0.355\textwidth}>{\raggedright\arraybackslash}p{0.355\textwidth}}
\toprule
Topic of the paper & continuous family $\zeta_\lambda$ & atomized family $\tilde\zeta_{\Z,\lambda}$\\
\midrule
\endfirsthead
\toprule
Topic of the paper & continuous family $\zeta_\lambda$ & atomized family $\tilde\zeta_{\Z,\lambda}$\\
\midrule
\endhead
\bottomrule
\endfoot
\addlinespace[2pt]
clock, $\tau_k$, bounded displacement & same & same\\
definition & Hadamard product over $\frac12\pm i\gamma_k(\lambda)$ & Euler product over $n\ge2$ with masses $\tilde c_n(\lambda)=\langle d\pi_\lambda,\varphi_n\rangle$\\
zeros & exactly $\frac12\pm i\gamma_k(\lambda)$; impossible elsewhere (Theorem~\CGref{thm:path}) & none in $\Re s>\frac12$ for every $\lambda$, since $\tilde h_\lambda$ is holomorphic there (Theorem~\CGref{thm:atomized-continuation}(iii)); for $\lambda=1$, exactly $\frac12\pm i\tau_k$ in $\Re s>0$ (Theorem~\CGref{thm:atomized-continuation}(iv))\\
functional equation & proved for every $\lambda$ & open; the factor $e^{\tilde h_\lambda}$ carries no functional equation\\
Dirichlet series & not for $\lambda=1$ in any half-plane (Theorem~\CGref{thm:no-Dirichlet-RVM}); not in $\Re s>1$ for $0<\lambda<1$ (Proposition~\CGref{prop:nonconstant}); at $\lambda=0$ this is (D) & proved in $\Re s>1$ for every $\lambda$, including $\lambda=0$ (Corollary~\CGref{cor:abscissa-one})\\
first coefficient ($\sigma\to+\infty$) & $\varpi_\lambda=\varpi_{\RVM}\prod_k\frac{1/4+\tau_k^2}{1/4+\gamma_k(\lambda)^2}$ \CGeqref{eq:varpi-lambda} & $\tilde a_1(\lambda)=1$ for every $\lambda$\\
vertical behaviour & $\lambda=1$: $\zeta_{\RVM}(\sigma+it)\to\varpi_{\RVM}$ as $|t|\to\infty$ \CGeqref{eq:vertical-limit} & almost periodic on every line of absolute convergence; $\lambda=1$: no limit as $|t|\to\infty$ (Proposition~\CGref{prop:atomized-clock})\\
half-plane bounds & $\lambda=1$: $C^{-1}\le|\zeta_{\RVM}|\le C$ on $\Re s\ge\sigma_0>1$ (Theorem~\CGref{thm:halfplane-RVM}); every $\lambda$: $|\log\zeta_\lambda-\log\frac s{s-1}|\ll_c1+\log(|t|+2)$ on $\Re s\ge c>\frac12$ & $|\tilde\zeta_{\Z,\lambda}|\le\sum_N|\tilde a_N(\lambda)|N^{-\sigma_0}$ and $|\tilde\zeta_{\Z,\lambda}^{-1}|\le\sum_N|\tilde\mu_\lambda(N)|N^{-\sigma_0}$ on $\Re s\ge\sigma_0>1$, for every $\lambda$\\
prime coordinate & counts $J_\lambda$, $\pi_\lambda$ (distributions, \CGeqref{eq:J-X}) & masses $\tilde c_n(\lambda)$ at the integers; $\tilde c_n(1)=O(1)$; mean square $X^{1+\eps}$ on $(X,2X]$ for every $\lambda$ (Theorem~\CGref{thm:mean-square-cells})\\
reciprocal & $F_\lambda=1/\zeta_\lambda$; $F_{Z1}$ is the Mellin transform of $B_{Z1}$ (Proposition~\CGref{prop:Mellin-Mertens}) & Dirichlet series with coefficients $\tilde\mu_\lambda(N)$ \CGeqref{eq:atomized-mobius}\\
Mertens & $M^{\crit}_{\RVM}=L+O_\eps(x^{1/2+\eps})$ (Theorem~\CGref{thm:unsmoothed-RVM}); $B_{Z1}=M\iff(D)$ & $\tilde M_\lambda=\sum_{N\le x}\tilde\mu_\lambda(N)$ exists for every $\lambda$ \CGeqref{eq:atomized-mertens}; the bound $x^{1/2+\eps}$ is open (it needs control of $e^{-\tilde h_\lambda}$)\\
flow in $\lambda$ & generator $\mathcal V_\lambda$ on the Hadamard product; $\Upsilon(\lambda)$ at $+\infty$ (Proposition~\CGref{prop:family-flow}) & $\log\tilde\zeta_{\Z,\lambda}=\sum_n\tilde c_n(\lambda)\bigl(-\log(1-n^{-s})\bigr)$, each mass differentiable in $\lambda$ by \CGeqref{eq:J-lambda-derivative}; joining factor \CGeqref{eq:joining-derivative}\\
limit $\lambda\to0^+$ & $\zeta_\lambda\to\zeta_{Z1}$ & $\tilde c_n(\lambda)\to\tilde c_n(0)$, $\mathfrak h_n(\lambda;\cdot)\to\mathfrak h_n(0;\cdot)$ (Proposition~\CGref{prop:cell-limit})\\
endpoint: Hamburger's hypotheses & H0 and H2 proved; H1 is (D), open & H1 proved in $\Re s>1$ (Corollary~\CGref{cor:abscissa-one}); H0 and H2 open\\
relation between the endpoints & $\tilde\zeta_{\Z,0}=\zeta_{Z1}\,e^{\tilde h_0}$ & same\\
what closes the route & $\zeta_{Z1}$ satisfies H1 & $\tilde\zeta_{\Z,0}$ satisfies H0 and H2\\
common form of the closure & $\tilde h_0\equiv0$ & $\tilde h_0\equiv0$\\
equivalent forms of the closure & $\zeta_{Z1}$ has a positive representing measure (Theorem~\CGref{thm:positivity}) & $e^{\tilde h_0}$ of finite order and symmetric under $s\mapsto1-s$ (Theorem~\CGref{thm:pair-criterion})\\
consequence of the closure & $\zeta_{Z1}=\zeta$; Step~3 gives $M(x)=O_\eps(x^{1/2+\eps})$ (Theorem~\CGref{thm:step3}) & $\tilde a_N(0)=1$, $\tilde c_n(0)=\mathbf 1_{\rm prime}(n)$, $\tilde\mu_0=\mu$, $\tilde M_0=M$; $\tilde\zeta_{\Z,0}=\zeta=\zeta_{Z1}$\\
\end{longtable}
\endgroup

\noindent\emph{How the two routes complement each other.}  Along $\lambda$ each family keeps exactly the half of Hamburger's hypotheses that the other one lacks:
\begin{equation}\label{eq:complement}
 \underbrace{\zeta_\lambda}_{\substack{\text{keeps the functional equation}\\ \text{and the zeros on }\Re s=\frac12}}\qquad\longleftrightarrow\qquad\underbrace{\tilde\zeta_{\Z,\lambda}}_{\substack{\text{keeps the Dirichlet series}\\ \text{and the Euler product}}},\qquad \tilde\zeta_{\Z,\lambda}=\zeta_\lambda\,e^{\tilde h_\lambda}.
\end{equation}
The closure is the same in both columns, the vanishing of the joining factor at the endpoint: by \CGeqref{eq:joining-factor} and Theorem~\CGref{thm:closing},
\begin{equation}\label{eq:closure-common}
 \tilde h_0\equiv0\iff J_{Z1}=0\text{ on }(1,2),\ \ \mathfrak h_n(0;s)=0\ \ \text{for all }n\ge2,\ \Re s>1,\ \text{and}\ \varpi_{Z1}=1 .
\end{equation}
No further condition is needed in the atomized column: the cell masses have mean square $X^{1+\eps}$ uniformly in $\lambda$ (Theorem~\CGref{thm:mean-square-cells}), so $\tilde\zeta_{\Z,0}$ satisfies H1 (Corollary~\CGref{cor:abscissa-one}).

\subsection{Programme}
\label{subsec:programme}

The steps are numbered in the order in which they are followed.  The status column records what is proved in this paper.

\begingroup\small
\begin{longtable}{p{0.05\textwidth}p{0.55\textwidth}p{0.30\textwidth}}
\toprule
Step & Statement & Status\\
\midrule
\endfirsthead
\toprule
Step & Statement & Status\\
\midrule
\endhead
\bottomrule
\endfoot
A1 & Atomization by smooth cells, coefficients $\tilde a_N(\lambda)$ & proved (Definition~\CGref{def:atomized})\\
A2 & Joining factor and explicit cell defects, \CGeqref{eq:joining-factor}--\CGeqref{eq:cell-explicit-derivative} & proved (Thm.~\CGref{thm:joining-factor}, Prop.~\CGref{prop:cell-explicit})\\
A3 & $L^2$ square-root estimate for every $\lambda$; Dirichlet series in $\Re s>\frac32$ & proved (Theorem~\CGref{thm:atomized-continuation}); superseded by A6\\
A4 & $\tilde h_\lambda$ holomorphic in $\Re s>\frac12$; in $\Re s>0$ for $\lambda=1$ & proved (Theorem~\CGref{thm:atomized-continuation})\\
A5 & Limits $\lambda\to0^+$ of $\tilde c_n$, $\mathfrak h_n$, $\varpi_\lambda$; flow at infinity $\Upsilon$ & proved (Propositions~\CGref{prop:family-flow}, \CGref{prop:cell-limit})\\
A5$'$ & Endpoints: $\tilde h_1$ not constant; atomized M\"obius coefficients $\tilde\mu_\lambda$; Mellin form $B_{Z1}$, with $B_{Z1}=M_{Z1}$ only under \CGeqref{eq:Z1-reciprocal-Dirichlet} & proved (Propositions~\CGref{prop:atomized-clock}, \CGref{prop:Mellin-Mertens})\\
A6 & Mean square of the cell masses uniformly in $\lambda$; Dirichlet series in $\Re s>1$ for every $\lambda$, including $\lambda=0$ & proved (Theorem~\CGref{thm:mean-square-cells}, Corollary~\CGref{cor:abscissa-one})\\
A7 & Cell identity \CGeqref{eq:closing-cells} for every $n\ge2$, and $J_{Z1}=0$ on the first cell $(1,2)$ & open; reduced to one number for each cell in A20\\
A8 & Equivalent scalar form: $\int_0^1\Upsilon\,d\lambda=\log L$, i.e.\ \CGeqref{eq:product-identity} & open; follows from A7\\
A9 & $\zeta_{Z1}=\zeta$ by Hamburger; $M(x)=O_\eps(x^{1/2+\eps})$ by Step~3 & proved as implication from A7 (Thms.~\CGref{thm:closing}, \CGref{thm:step3})\\
A10 & Continuation of $\tilde h_\lambda$ to $0<\Re s\le\frac12$ for $\lambda<1$ & open; not needed for A9; at $\lambda=0$ it follows from \textup{(B1)}, \textup{(B2)} (A23)\\
A11 & $\zeta_\lambda$ is not a Dirichlet series and $\tilde h_\lambda$ is not constant for $0<\lambda\le1$; convexity of $\log\zeta_\lambda(\sigma)$ in $\lambda$ & proved (Propositions~\CGref{prop:nonconstant}, \CGref{prop:convexity})\\
A12 & Hamburger's theorem for the pair: A9 follows if $e^{\tilde h_0}$ is of finite order and symmetric under $s\mapsto1-s$ & proved as implication (Theorem~\CGref{thm:pair-criterion}); at $\lambda=0$ the hypothesis is equivalent to (D) (Remark~\CGref{rem:pair-no-slack})\\
A13 & Positivity: A9 follows if $\zeta_{Z1}$ has a positive representing measure; then $\varpi_{Z1}\le1$ & proved as implication (Theorem~\CGref{thm:positivity}, Corollary~\CGref{cor:positivity})\\
A14 & Clock end: density plus explicit atoms of $d\pi_{\RVM}$, cell masses $\frac1{\log n}+\mathbf1_{\mathrm{Pow}_2}(n)\frac{\mu(\log_2n)}{\log_2n}J_{\RVM}(2)+O(n^{-1/4})$, defects $O(|s|n^{-\sigma-4/3})$ & proved (Theorem~\CGref{thm:clock-density}, Corollary~\CGref{cor:clock-cells})\\
A15 & Height kernel of the cell masses; flow as count difference; Gaussian defect and the density condition $\tilde h_0(1)=0$ & proved (Propositions~\CGref{prop:height-kernels}, \CGref{prop:flow-count}, \CGref{prop:gaussian-defect})\\
A16 & Hamburger's theorem for the pair without continuation: A9 follows if $T_0\equiv0$, or if $\mathsf V_\sigma(0)=0$; symmetry fails at $\lambda=1$; exceptional set finite & proved (Theorem~\CGref{thm:double-hamburger}, Propositions~\CGref{prop:clock-asymmetry}, \CGref{prop:exceptional-set})\\
A17 & $\zeta_{Z1}(\sigma)\ge\zeta(\sigma)$ for $\sigma>1$; A9 follows from the reverse inequality at one $\sigma$, or from $\varpi_\lambda\le1$ along $\lambda\to0^+$, or from $\mathsf W(0)\ge0$ & proved (Theorem~\CGref{thm:real-axis}, Propositions~\CGref{prop:one-number}, \CGref{prop:open-path-varpi})\\
A18 & Space $\mathscr O$ of ordinate lists: at most one coincidence list; A9 is equivalent to its existence & proved (Theorem~\CGref{thm:rigidity}, Corollary~\CGref{cor:existence-form})\\
A19 & Existence of a list in $\mathscr O$ with $\tilde h_{\boldsymbol{\kappa}}$ constant, equivalently $\min_{\mathscr O}\mathsf V_\sigma=0$ & open; equivalent to A7\\
A20 & A9 holds if and only if $\tilde c_n(0)=\mathbf 1_{\rm prime}(n)$ for every $n$; every $\lambda\in\mathsf E$ is a witness; the identity \CGeqref{eq:per-cell} cell by cell & proved (Corollary~\CGref{cor:cell-mass-form}, Propositions~\CGref{prop:witness}, \CGref{prop:per-cell}); the identity \CGeqref{eq:per-cell} itself is open\\
A21 & Symmetric joining factors: order one is needed, and A12 has no slack at $\lambda=0$ & proved (Proposition~\CGref{prop:dichotomy}, Remarks~\CGref{rem:pair-no-slack}, \CGref{rem:dichotomy-counterexample})\\
A22 & Half-line form: A9 follows if $\tilde h_0(\sigma)=O_N(\sigma N^{-\sigma})$ for every $N$; the $N$th rung is the constancy of $J^{\rm M}_{Z1}$ on $(N-1,N)$ & proved (Theorems~\CGref{thm:ladder}, \CGref{thm:halfline-hamburger}, Corollary~\CGref{cor:ladder-pair})\\
A23 & Bridge theorem: A9 follows if $\Im\tilde h_0$ vanishes on $\Re s=\frac12$ and $\tilde h_0$ is temperate; the two conditions hold for no $\lambda>0$ & proved (Theorem~\CGref{thm:bridge}, Corollary~\CGref{cor:bridge-family}); \textup{(B1)}, \textup{(B2)} at $\lambda=0$ open\\
A24 & Transport from the clock: $\mathsf S_{\mathsf f}$ entire of type $C_\eta\log N$, moment series on the clock, defect of the clock by Poisson summation & proved (Proposition~\CGref{prop:rung-functional}, Theorem~\CGref{thm:transport-band}, Proposition~\CGref{prop:clock-defect-poisson})\\
A24b & Multiscale Jacobian at the clock: stable finite ordinate block on windows $U\asymp\log T$, fixed-window obstruction, dense range, and rapidly summable smooth tails & proved (Theorem~\CGref{thm:clock-jacobian}); quantitative multiscale continuation to a solution open\\
A25 & Rung system with free masses: A9 holds if and only if it has a solution in $\mathscr O$; the solution is unique and nondegenerate & proved (Theorems~\CGref{thm:free-masses}, \CGref{thm:nondegenerate}); existence open\\
A26 & Coordinate gradient of $\mathsf V_\sigma$ on $\mathscr O$; A9 would follow from a boundary-aware coordinate/KKT criterion forcing every minimizer to have zero defect & proved (Proposition~\CGref{prop:critical-points}); the forcing statement is open\\
A27 & Acceptance criterion and models: the final step is a theorem for the finite Euler product over the primes $2,3$ (not for more primes, see after Theorem~\CGref{thm:finite-euler}); no object of finite rank in the window of $\zeta$ & proved (Theorems~\CGref{thm:acceptance}, \CGref{thm:finite-euler}, Propositions~\CGref{prop:graph-model}, \CGref{prop:no-finite-rank})\\
\end{longtable}
\endgroup

\begin{remark}[Input for A7]\label{rem:input-A7}
Steps A1--A6 use only two properties of an ordinate list: bounded displacement from the clock and a primitive of the count difference of size $O(\log\tau)$.  The list $(\tau_k)$ has both properties, and for it the corresponding function $\zeta_{\RVM}$ is not a Dirichlet series (Theorem~\CGref{thm:no-Dirichlet-RVM}).  Step A7 must therefore use a property of the ordinates $\gamma_k$ that the clock ordinates do not have; the natural candidate is the arithmetic of $\zeta$ carried by the $\gamma_k$, for instance Landau's formula $\sum_{0<\gamma\le T}x^\rho=-\frac T{2\pi}\Lambda(x)+O(\log T)$, in which the prime powers appear as resonances of the ordinates.  The multiscale analysis below separates a different issue: the linearized clock has enough directions, but only on windows whose length grows like $\log T$; thus the obstruction is not fixed-height injectivity but quantitative nonlinear continuation across scales.
\end{remark}

\section{Criteria for the final step and the clock end of the family}
\label{sec:criteria}

This section collects the consequences of Corollary~\CGref{cor:abscissa-one} for the final step.  Along the open path $0<\lambda\le1$ no member $\zeta_\lambda$ is a Dirichlet series (Proposition~\CGref{prop:nonconstant}); the family is convex in $\lambda$ on the real axis (Proposition~\CGref{prop:convexity}); Hamburger's theorem applies to the pair $(\zeta_\lambda,\tilde\zeta_{\Z,\lambda})$ as soon as the joining factor is symmetric (Theorem~\CGref{thm:pair-criterion}); positivity of the measure representing $\zeta_{Z1}$ gives $\varpi_{Z1}\le1$ and closes the route (Theorem~\CGref{thm:positivity}).  At the clock end the density of $d\pi_{\RVM}$ is computed explicitly (Theorem~\CGref{thm:clock-density}), and the flow of the cell masses is written in the heights of the ordinates (Propositions~\CGref{prop:height-kernels} and~\CGref{prop:flow-count}).  Subsections~\CGref{subsec:cell-mass-form}--\CGref{subsec:transport} reduce the final step to the cell masses (Corollary~\CGref{cor:cell-mass-form}), state two criteria whose hypotheses are written where the joining factor is defined, the half-line form (Theorem~\CGref{thm:halfline-hamburger}) and the bridge theorem (Theorem~\CGref{thm:bridge}), and compute the transport from the clock on each cell (Theorem~\CGref{thm:transport-band}).

Throughout this section the function $\chi$ of Subsection~\CGref{subsec:mean-square} is even, so that $\varphi$ is even, and
\begin{equation}\label{eq:cell-moments}
 \hat{\varphi}(\xi):=\int_{\R}\varphi(t)e^{-2\pi i\xi t}\,dt,\qquad \mathsf M_r:=\int_{\R}t^r\varphi(t)\,dt,\qquad \mathsf C_r:=\int_{\R}t^r\varphi(t)\cos(2\pi t)\,dt .
\end{equation}
$\hat{\varphi}$ is real and even, and $\mathsf M_1=\mathsf C_1=0$.  Since $\sum_{m\in\Z}\varphi(t-m)=1$,
\begin{equation}\label{eq:partition-fourier}
 \hat{\varphi}(q)=\int_0^1\sum_{m\in\Z}\varphi(t-m)\,e^{-2\pi iqt}\,dt=\delta_{q0}\qquad(q\in\Z),
\end{equation}
so $\mathsf M_0=1$ and $\mathsf C_0=\hat{\varphi}(1)=0$.

\subsection{The open path}
\label{subsec:open-path}

\begin{proposition}[No member of the open path is a Dirichlet series]\label{prop:nonconstant}
Let $0<\lambda\le1$.  Then $\zeta_\lambda$ is not a constant multiple of a Dirichlet series $\sum_na_nn^{-s}$ converging absolutely in $\Re s>1$.  In particular the joining factor $\tilde h_\lambda$ is not constant.
\end{proposition}

\begin{proof}
Suppose $\zeta_\lambda(s)=\sum_na_nn^{-s}$ absolutely for $\Re s>1$.  By Theorem~\CGref{thm:path}, exactly as in Proposition~\CGref{prop:Z1-class}, $(s-1)\zeta_\lambda(s)$ is entire of order one and $\zeta_\lambda$ satisfies \CGeqref{eq:Z1-FE}; these are the hypotheses of Hamburger's theorem, so $\zeta_\lambda=C\zeta$ with $C\ne0$.  The nontrivial zeros of $\zeta$ with positive ordinate are then $\frac12+i\gamma_k(\lambda)$, $k\ge1$, with multiplicity.  The ordinates of $\zeta$, listed in non-decreasing order, are the $\gamma_k$, and the $\gamma_k(\lambda)$ are increasing; hence $\gamma_k(\lambda)=\gamma_k$, that is $\lambda\eta_k=0$, for every $k$.  Since $\lambda>0$, every $\eta_k$ vanishes, all members of the family coincide (Remark~\CGref{rem:only-gamma-tau}), and $\zeta_{\RVM}=\zeta_\lambda$ is an absolutely convergent Dirichlet series, against Theorem~\CGref{thm:no-Dirichlet-RVM}.  If $\tilde h_\lambda$ were a constant $c$, then $\zeta_\lambda=e^{-c}\tilde\zeta_{\Z,\lambda}$ would be such a series by Corollary~\CGref{cor:abscissa-one}.
\end{proof}

The two columns of \CGeqref{eq:complement} exist for every $\lambda$, and for every $\lambda\in(0,1]$ they are different functions.  The joining factor converges, $\tilde h_\lambda\to\tilde h_0$ as $\lambda\to0^+$, through non-constant functions; (D) is the statement that the limit is constant.  Any property of the displacement sequence that would imply (D) must distinguish the full sequence $(\eta_k)$ from every scaled sequence $(c\,\eta_k)$ with $0\le c<1$: the latter correspond to members $\lambda>0$ of the family, for which Proposition~\CGref{prop:nonconstant} excludes closure.  Thus this observation is a necessary discrimination test for a proposed criterion, not an assertion that the actual displacement sequence already satisfies such a criterion.

\subsection{Convexity along the family}
\label{subsec:convexity}

\begin{proposition}[Convexity on the real axis]\label{prop:convexity}
For every real $\sigma>1$ the function $\lambda\mapsto\log\zeta_\lambda(\sigma)$ is strictly convex on $[0,1]$:
\begin{equation}\label{eq:convexity-second}
 \partial_\lambda^2\log\zeta_\lambda(\sigma)=\sum_{k\ge1}\eta_k^2\Bigl(\mathcal C\bigl(\sigma-\tfrac12,\gamma_k(\lambda)\bigr)-\mathcal C\bigl(\tfrac12,\gamma_k(\lambda)\bigr)\Bigr)>0,\qquad \mathcal C(a,\gamma):=\frac{2(a^2-\gamma^2)}{(a^2+\gamma^2)^2}.
\end{equation}
Consequently, with the generator $\mathcal V_\lambda$ of Theorem~\CGref{thm:generator},
\begin{equation}\label{eq:convexity-sandwich}
 \log\zeta_{\RVM}(\sigma)-\mathcal V_1(\sigma)<\log\zeta_{Z1}(\sigma)<\log\zeta_{\RVM}(\sigma)-\mathcal V_0(\sigma)\qquad(\sigma>1).
\end{equation}
\end{proposition}

\begin{proof}
With $a=\sigma-\frac12>\frac12$ one has $q(\sigma)=a^2-\frac14$, so by \CGeqref{eq:P-lambda} $\log\zeta_\lambda(\sigma)=-\log G(\sigma)+\sum_k\bigl(\log(a^2+\gamma_k(\lambda)^2)-\log(\frac14+\gamma_k(\lambda)^2)\bigr)$.  Since $\partial_\gamma^2\log(a^2+\gamma^2)=\mathcal C(a,\gamma)$ and $\partial_\lambda\gamma_k(\lambda)=\eta_k$, termwise differentiation gives the series in \CGeqref{eq:convexity-second}; its terms are $O(\eta_k^2\gamma_k^{-2})$, so it converges uniformly on $[0,1]$.  Fix $\gamma>\frac12$.  Then $\mathcal C(\frac12,\gamma)<0$; $\partial_a\mathcal C(a,\gamma)=4a(3\gamma^2-a^2)(a^2+\gamma^2)^{-3}>0$ for $0<a<\sqrt3\,\gamma$; and $\mathcal C(a,\gamma)>0$ for $a>\gamma$.  Hence $\mathcal C(a,\gamma)>\mathcal C(\frac12,\gamma)$ for every $a>\frac12$.  All ordinates satisfy $\gamma_k(\lambda)>\tau_0>\frac12$, so every term of the series is $\ge0$; the sum is $>0$ exactly when some $\eta_k\ne0$, that is, when $\boldsymbol\gamma\ne\boldsymbol\tau$, and this holds with $k=1$: $\eta_1=\tau_1-\gamma_1>0$ by the comparison $\gamma_1<14.14<\tau_1$ made in the proof of Proposition~\CGref{prop:family-flow}. So the sum is positive.  A strictly convex function $F$ on $[0,1]$ satisfies $F(1)-F'(1)<F(0)<F(1)-F'(0)$; with $F(\lambda)=\log\zeta_\lambda(\sigma)$ and $F'=\mathcal V_\lambda(\sigma)$ \CGeqref{eq:generator-zeta} this is \CGeqref{eq:convexity-sandwich}.
\end{proof}

As $\sigma\to+\infty$, \CGeqref{eq:convexity-second} becomes $-\Upsilon'(\lambda)>0$ and \CGeqref{eq:convexity-sandwich} becomes the last statement of Proposition~\CGref{prop:family-flow}.  The sign in \CGeqref{eq:convexity-second} comes from the test function $x^{-\sigma}$, whose transform is not oscillatory in the ordinate; for the cell functions $\varphi_n$ the corresponding second derivative $\partial_\lambda^2\tilde c_n(\lambda)=-2\Re\sum_k\eta_k^2W_n''(\gamma_k(\lambda))$ has terms of both signs.

\subsection{Hamburger's theorem for the pair}
\label{subsec:pair}

\begin{theorem}[Hamburger's theorem for the pair $(\zeta_\lambda,\tilde\zeta_{\Z,\lambda})$]\label{thm:pair-criterion}
Let $\lambda\in[0,1]$.  Suppose that the joining factor $e^{\tilde h_\lambda}$ extends to a meromorphic function on $\C$ of finite order with finitely many poles, and that
\begin{equation}\label{eq:pair-symmetry}
 e^{\tilde h_\lambda(s)}=e^{\tilde h_\lambda(1-s)} .
\end{equation}
Then $\tilde\zeta_{\Z,\lambda}=\zeta$.  For $\lambda=0$, moreover, $\zeta_{Z1}=\zeta$ and $\tilde h_0\equiv0$.
\end{theorem}

\begin{proof}
Put $f=\tilde\zeta_{\Z,\lambda}=\zeta_\lambda e^{\tilde h_\lambda}$.  By Theorem~\CGref{thm:path}, $(s-1)\zeta_\lambda$ is entire of order one, so $f=\mathsf G/\mathsf P$ with $\mathsf G$ entire of finite order and $\mathsf P$ a polynomial: this is (H0).  By \CGeqref{eq:Z1-FE} for $\zeta_\lambda$ and \CGeqref{eq:pair-symmetry}, $f(s)\Gamma(\frac s2)\pi^{-s/2}=f(1-s)\Gamma(\frac{1-s}2)\pi^{-(1-s)/2}$: this is (H2) with $g=f$.  (H1) is Corollary~\CGref{cor:abscissa-one}, and it gives (H3) for $g=f$.  Hamburger's theorem gives $f=C\zeta$, and $C=1$ because $\tilde a_1(\lambda)=1$.  By Theorem~\CGref{thm:atomized-continuation}(iii), $f=\zeta$ has no zeros in $\Re s>\frac12$, and by the functional equation none in $\Re s<\frac12$ apart from the trivial ones: the Riemann hypothesis holds.  For $\lambda=0$, Theorem~\CGref{thm:Hamburger-Z1} (RH $\Leftrightarrow\zeta_{Z1}=\zeta$) then gives $\zeta_{Z1}=\zeta=\tilde\zeta_{\Z,0}$, and $e^{\tilde h_0}=1$.
\end{proof}

Theorem~\CGref{thm:pair-criterion} uses the functional equation of $\zeta_\lambda$ and the Dirichlet series of $\tilde\zeta_{\Z,\lambda}$ as they stand; the only hypothesis is on the factor that joins them, and it asks for symmetry instead of constancy.  At $\lambda=1$ the two hypotheses do not hold together: the conclusion $\tilde\zeta_{\Z,1}=\zeta$ would place the zeros of $\zeta$ at $\frac12\pm i\tau_k$ (Theorem~\CGref{thm:atomized-continuation}(iv)), so $\zeta=\zeta_{\RVM}$ by Theorem~\CGref{thm:Hadamard}, against Theorem~\CGref{thm:no-Dirichlet-RVM}.

\begin{theorem}[Real-variable and vertical forms of Hamburger's theorem for the pair]\label{thm:double-hamburger}
For $\lambda\in[0,1]$, $y>0$ and $\sigma>1$ put
\begin{equation}\label{eq:theta-defect}
 T_\lambda(y):=\sum_{N\ge1}\tilde a_N(\lambda)\,e^{-\pi N^2y}-y^{-1/2}\sum_{N\ge1}\tilde a_N(\lambda)\,e^{-\pi N^2/y}-\tfrac12\bigl(y^{-1/2}-1\bigr),
\end{equation}
\begin{equation}\label{eq:vertical-defect}
 k_n(\lambda):=\sum_{d\mid n}\tilde a_d(\lambda)\,\mu(n/d),\qquad \mathsf V_\sigma(\lambda):=\sum_{n\ge2}|k_n(\lambda)|^2\,n^{-2\sigma}.
\end{equation}
\begin{enumerate}[label=\textup{(\roman*)}]
\item $\tilde\zeta_{\Z,\lambda}/\zeta=\sum_{n\ge1}k_n(\lambda)n^{-s}$ converges absolutely in $\Re s>1$, $k_1(\lambda)=1$, and
\begin{equation}\label{eq:vertical-mean}
 \lim_{T\to\infty}\frac1{2T}\int_{-T}^{T}\left|\frac{\tilde\zeta_{\Z,\lambda}(\sigma+it)}{\zeta(\sigma+it)}\right|^2dt=1+\mathsf V_\sigma(\lambda)\ \ge1 .
\end{equation}
\item For each $\lambda$ the following are equivalent: \textup{(a)} $\tilde\zeta_{\Z,\lambda}=\zeta$; \textup{(b)} $T_\lambda(y)=0$ for every $y>0$; \textup{(c)} $\mathsf V_\sigma(\lambda)=0$ for some, equivalently every, $\sigma>1$.  Each of them implies (D).
\item $\lambda\mapsto T_\lambda(y)$ and $\lambda\mapsto\mathsf V_\sigma(\lambda)$ are continuous on $[0,1]$, and
\begin{equation}\label{eq:double-hamburger-limit}
 (\mathrm D)\iff T_0\equiv0\iff\mathsf V_\sigma(0)=0\iff\lim_{\lambda\to0^+}T_\lambda(y)=0\ \ (y>0)\iff\lim_{\lambda\to0^+}\mathsf V_\sigma(\lambda)=0 .
\end{equation}
At the clock end $T_1\not\equiv0$ and $\mathsf V_\sigma(1)>0$.
\end{enumerate}
\end{theorem}

\begin{proof}
(i) $1/\zeta=\sum_n\mu(n)n^{-s}$ and $\tilde\zeta_{\Z,\lambda}$ (Corollary~\CGref{cor:abscissa-one}) converge absolutely in $\Re s>1$, and so does their product, with $k_1=\tilde a_1=1$.  The mean square of an absolutely convergent Dirichlet series on a vertical line is $\sum_n|k_n|^2n^{-2\sigma}$, as in \CGeqref{eq:mean-square}--\CGeqref{eq:mean-square-limit}.
(ii) (a)$\Leftrightarrow$(c) because $\mathsf V_\sigma(\lambda)=0$ means $k_n(\lambda)=0$ for $n\ge2$.  (a)$\Rightarrow$(b) is the theta relation $\sum_{n\ge1}e^{-\pi n^2y}=y^{-1/2}\sum_{n\ge1}e^{-\pi n^2/y}+\frac12(y^{-1/2}-1)$.  (b)$\Rightarrow$(a): the relation $T_\lambda\equiv0$ is the identity reached in the first half of the proof of Hamburger's theorem in \cite[\S2.13]{Titchmarsh1986}, with $b_n=a_n=\tilde a_n(\lambda)$; the second half of that proof uses only this identity and the absolute convergence of $\sum_na_nn^{-\sigma'}$ for $\sigma'>1$, and gives $a_k=a_{k+1}$ for every $k$, hence $\tilde a_N(\lambda)=\tilde a_1(\lambda)=1$.  If (a) holds, then by Theorem~\CGref{thm:atomized-continuation}(iii) and the functional equation the nontrivial zeros of $\zeta$ lie on $\Re s=\frac12$; they are the points $\frac12\pm i\gamma_k$, and Theorem~\CGref{thm:Hamburger-Z1} gives (D).
(iii) By Theorem~\CGref{thm:mean-square-cells}, for every $\sigma'>1$ the sums $\sum_n|\tilde c_n(\lambda)|n^{-\sigma'}$ are bounded uniformly in $\lambda$; hence so are $\sum_N|\tilde a_N(\lambda)|N^{-\sigma'}\le\exp\bigl(\sum_m\frac1m\sum_n|\tilde c_n(\lambda)|n^{-m\sigma'}\bigr)$ and $\sum_n|k_n(\lambda)|n^{-\sigma'}$.  Each $\tilde a_N(\lambda)$ and each $k_n(\lambda)$ is a polynomial in finitely many cell masses, which are continuous in $\lambda$.  The tails $\sum_{N>N_0}$ of \CGeqref{eq:theta-defect} and of \CGeqref{eq:vertical-defect} (with $1<\sigma'<\sigma$) are small uniformly in $\lambda$, which gives the continuity.  Under (D), $\tilde a_N(0)=1$ (Theorem~\CGref{thm:closing}), so $T_0\equiv0$ and $\mathsf V_\sigma(0)=0$; the converse is (ii) at $\lambda=0$, and the two limits are the values at $\lambda=0$ by continuity.  Finally $\tilde\zeta_{\Z,1}\ne\zeta$ by the remark after Theorem~\CGref{thm:pair-criterion}.
\end{proof}

Theorem~\CGref{thm:double-hamburger} needs no continuation of the joining factor: \CGeqref{eq:theta-defect} is a real function of $y$ built from the atomized coefficients alone, and \CGeqref{eq:vertical-mean} lives in $\Re s>1$, where
\begin{equation}\label{eq:quotient-three}
 \frac{\tilde\zeta_{\Z,\lambda}}{\zeta}=e^{\tilde h_\lambda}\,\frac{\zeta_\lambda}{\zeta}
\end{equation}
is a quotient of three explicit functions.  Its vertical mean is $1$ and its vertical mean square is at least $1$ for every $\lambda$; (D) is the statement that at $\lambda=0$ the mean square does not exceed $1$.

\begin{proposition}[At the clock end the symmetry fails]\label{prop:clock-asymmetry}
The joining factor $e^{\tilde h_1}$ has no meromorphic continuation to $\C$ that satisfies \CGeqref{eq:pair-symmetry}.
\end{proposition}

\begin{proof}
$e^{\tilde h_1}$ is holomorphic and zero-free in $\Re s>0$ (Theorem~\CGref{thm:atomized-continuation}(iv)).  A meromorphic continuation with \CGeqref{eq:pair-symmetry} is holomorphic and zero-free also in $\Re s<1$, hence entire and zero-free, so it equals $e^{g}$ with $g$ entire, $g=\tilde h_1$ on $\Re s>0$, and $g(s)-g(1-s)$ constant.  For $\Re s\ge\frac14$, the bound $|\mathfrak h_n(1;ms)|\le C\,m|s|\,n^{-m\sigma-1}$ of Theorem~\CGref{thm:atomized-continuation}(iv) gives $\sum_m\frac1m\sum_n|\mathfrak h_n(1;ms)|\le C|s|\sum_{n\ge2}n^{-1}(n^{1/4}-1)^{-1}$, and $|\mathfrak u_1(s)|\le C(1+|s|)$ by \CGeqref{eq:first-cell}; so $|g(s)|\le C(1+|s|)$ on $\Re s\ge\frac14$, and by the symmetry also on $\Re s\le\frac34$.  An entire function of linear growth is a polynomial of degree at most $1$, and $g(s)-g(1-s)$ constant then forces $g$ to be constant.  This contradicts Proposition~\CGref{prop:atomized-clock}(iii).
\end{proof}

At $\lambda=1$ it is therefore the symmetry \CGeqref{eq:pair-symmetry} itself that fails in Theorem~\CGref{thm:pair-criterion}, whatever the order of the continuation.

\subsection{Positivity and the value at infinity}
\label{subsec:positivity}

For an even Schwartz function $\mathsf g$ on $\R$ put $\hat{\mathsf g}(\xi)=\int_{\R}\mathsf g(u)e^{-2\pi i\xi u}du$ and $\mathsf g^{\sharp}(s)=\int_0^\infty\mathsf g(u)u^{s-1}du$.

\begin{lemma}[Summation formula for a representing measure]\label{lem:summation-positive}
Let $\lambda\in[0,1]$ and let $\mathsf A$ be a positive measure on $[1,\infty)$ with
\begin{equation}\label{eq:representing-measure}
 \zeta_\lambda(s)=\int_{[1,\infty)}u^{-s}\,d\mathsf A(u)\qquad(\Re s>1).
\end{equation}
Then for every even Schwartz function $\mathsf g$
\begin{equation}\label{eq:summation-positive}
 \mathsf g(0)+2\int_{[1,\infty)}\mathsf g\,d\mathsf A=\hat{\mathsf g}(0)+2\int_{[1,\infty)}\hat{\mathsf g}\,d\mathsf A .
\end{equation}
\end{lemma}

\begin{proof}
$\mathsf g^{\sharp}$ is meromorphic with simple poles at $s=0,-2,-4,\dots$, residue $\mathsf g(0)$ at $s=0$, and decays faster than any power of $|t|$ on vertical strips; $\mathsf g^{\sharp}(1)=\frac12\hat{\mathsf g}(0)$; and
\begin{equation}\label{eq:tate-local}
 \pi^{-s/2}\Gamma(\tfrac s2)\,\hat{\mathsf g}^{\sharp}(1-s)=\pi^{-(1-s)/2}\Gamma(\tfrac{1-s}2)\,\mathsf g^{\sharp}(s)
\end{equation}
for $0<\Re s<1$, and then by meromorphic continuation: the Fourier transform of the even homogeneous distribution $|\xi|^{-s}$ is a multiple of $|u|^{s-1}$, and the multiple is fixed by $\mathsf g(u)=e^{-\pi u^2}$, for which $\hat{\mathsf g}=\mathsf g$ and $\mathsf g^{\sharp}(s)=\frac12\pi^{-s/2}\Gamma(\frac s2)$.  Fix $1<c<2$.  By Mellin inversion and Fubini, $\int\mathsf g\,d\mathsf A=\frac1{2\pi i}\int_{(c)}\zeta_\lambda(s)\mathsf g^{\sharp}(s)\,ds$, since $\int u^{-c}d\mathsf A<\infty$.  On $\Re s\ge c$, $|\zeta_\lambda(s)|\le\zeta_\lambda(c)$; on $\Re s=1-c$ the functional equation \CGeqref{eq:Z1-FE}, valid for every $\zeta_\lambda$ by Theorem~\CGref{thm:path}, gives $\zeta_\lambda(s)=\zeta_\lambda(1-s)G(1-s)/G(s)$, a power of $|t|$ times $\zeta_\lambda(c)$.  The Phragm\'en--Lindel\"of principle, applied to the entire function $(s-1)\zeta_\lambda(s)$ of finite order on the strip $1-c\le\Re s\le c$, bounds it by a power of $|t|$ there, and dividing by $s-1$, with $|s-1|\ge1$ for $|t|\ge1$, gives the same for $\zeta_\lambda$ on $1-c\le\Re s\le c$, $|t|\ge1$.  Moving the line to $\Re s=1-c$ crosses $s=1$, where $\zeta_\lambda$ has residue $1$, and $s=0$, where $\zeta_\lambda(0)=P_\lambda(0)/G(0)=-\frac12$:
\begin{equation}\label{eq:summation-shift}
 \int\mathsf g\,d\mathsf A=\tfrac12\hat{\mathsf g}(0)-\tfrac12\mathsf g(0)+\frac1{2\pi i}\int_{(1-c)}\zeta_\lambda(s)\mathsf g^{\sharp}(s)\,ds .
\end{equation}
By \CGeqref{eq:Z1-FE} and \CGeqref{eq:tate-local}, $\zeta_\lambda(s)\mathsf g^{\sharp}(s)=\zeta_\lambda(1-s)\hat{\mathsf g}^{\sharp}(1-s)$, and the last integral is $\frac1{2\pi i}\int_{(c)}\zeta_\lambda(s)\hat{\mathsf g}^{\sharp}(s)\,ds=\int\hat{\mathsf g}\,d\mathsf A$.
\end{proof}

\begin{theorem}[Positivity criterion]\label{thm:positivity}
Let $\lambda\in[0,1]$.  If $\zeta_\lambda$ has a representation \CGeqref{eq:representing-measure} with a positive measure $\mathsf A$, then
\begin{equation}\label{eq:positivity-bound}
 \varpi_\lambda=\mathsf A(\{1\})\le1 .
\end{equation}
In particular, if $\zeta_{Z1}$ has such a representation, then $\varpi_{Z1}=1$ and $\zeta_{Z1}=\zeta$.
\end{theorem}

\begin{proof}
By monotone convergence $\zeta_\lambda(\sigma)=\int u^{-\sigma}d\mathsf A\to\mathsf A(\{1\})$ as $\sigma\to+\infty$, and the limit is $\varpi_\lambda$ (Proposition~\CGref{prop:family-clock}).  Let $\mathsf T(\xi)=(1-|\xi|)_+$ and $\mathsf B=\mathsf T*\mathsf T$, so that $\mathsf B(\xi)=\frac23-\xi^2+\frac12|\xi|^3$ for $|\xi|\le1$, $\mathsf B(\xi)=\frac16(2-|\xi|)^3$ for $1\le|\xi|\le2$, $\mathsf B=0$ elsewhere, and $\mathsf B''(\xi)=\mathsf T(\xi+1)-2\mathsf T(\xi)+\mathsf T(\xi-1)$.  Put
\begin{equation}\label{eq:extremal-function}
 \mathsf g(u):=\Bigl(\frac{\sin\pi u}{\pi u}\Bigr)^4(1-u^2),\qquad \hat{\mathsf g}(\xi)=\mathsf B(\xi)+\frac1{4\pi^2}\mathsf B''(\xi) .
\end{equation}
Then $\mathsf g(0)=1$; $\mathsf g(u)\le0$ for $|u|\ge1$; $\hat{\mathsf g}(0)=\frac23-\frac1{2\pi^2}$ and $\hat{\mathsf g}(1)=\frac16+\frac1{4\pi^2}$; and $\hat{\mathsf g}\ge0$ on $\R$: for $0\le\xi\le1$, $\hat{\mathsf g}(\xi)=\frac23-\xi^2+\frac12\xi^3+\frac{3\xi-2}{4\pi^2}$ is positive at $\xi=0$ and $\xi=1$, and its derivative $-2\xi+\frac32\xi^2+\frac3{4\pi^2}$ has the roots $\xi_\pm=\frac13\bigl(2\pm\sqrt{4-9/(2\pi^2)}\bigr)$, $\xi_-=0.0391\ldots$, $\xi_+=1.294\ldots$, so on $[0,1]$ it is positive on $[0,\xi_-)$ and negative on $(\xi_-,1]$, and the minimum of $\hat{\mathsf g}$ on $[0,1]$ is $\min(\hat{\mathsf g}(0),\hat{\mathsf g}(1))=\frac16+\frac1{4\pi^2}>0$; and for $1\le\xi\le2$, $\hat{\mathsf g}(\xi)=\frac16(2-\xi)^3+\frac{2-\xi}{4\pi^2}\ge0$.  For $\epsilon>0$ let $\mathsf g_\epsilon(u)=\mathsf g(u)e^{-\epsilon\pi u^2}$, an even Schwartz function with $\mathsf g_\epsilon(0)=1$, $\mathsf g_\epsilon\le0$ on $|u|\ge1$, and $\hat{\mathsf g}_\epsilon=\hat{\mathsf g}*(\epsilon^{-1/2}e^{-\pi\xi^2/\epsilon})\ge0$.  By Lemma~\CGref{lem:summation-positive}, since $\mathsf A\ge0$ is carried by $[1,\infty)$ and has mass $\varpi_\lambda$ at $u=1$,
\begin{equation}\label{eq:lp-inequality}
 1\ \ge\ \mathsf g_\epsilon(0)+2\int\mathsf g_\epsilon\,d\mathsf A=\hat{\mathsf g}_\epsilon(0)+2\int\hat{\mathsf g}_\epsilon\,d\mathsf A\ \ge\ \hat{\mathsf g}_\epsilon(0)+2\varpi_\lambda\,\hat{\mathsf g}_\epsilon(1) .
\end{equation}
Letting $\epsilon\to0$, $\hat{\mathsf g}_\epsilon\to\hat{\mathsf g}$ pointwise because $\hat{\mathsf g}$ is continuous, and
\begin{equation}\label{eq:lp-value}
 \varpi_\lambda\le\frac{1-\hat{\mathsf g}(0)}{2\hat{\mathsf g}(1)}=\frac{\frac13+\frac1{2\pi^2}}{\frac13+\frac1{2\pi^2}}=1 .
\end{equation}
The last statement follows from Theorem~\CGref{thm:varpi-Z1}.
\end{proof}

\begin{corollary}[Positive form of the final step]\label{cor:positivity}
The following are equivalent: \textup{(i)} (D); \textup{(ii)} $\zeta_{Z1}$ is the Mellin--Stieltjes transform \CGeqref{eq:representing-measure} of a positive measure on $[1,\infty)$; \textup{(iii)} $J^{\rm M}_{Z1}$, extended by $0$ to $(0,1]$, is non-decreasing on $(0,\infty)$.
\end{corollary}

\begin{proof}
(i)$\Rightarrow$(iii): then $J^{\rm M}_{Z1}=J$ (Theorem~\CGref{thm:cell-form}).  (iii)$\Rightarrow$(ii): $J^{\rm M}_{Z1}(x)x^{-c}\in L^2((1,\infty),dx/x)$ for every $c>1$, and a non-decreasing $J^{\rm M}_{Z1}$, which is $\ge0$, satisfies $J^{\rm M}_{Z1}(x)^2\,(2x)^{-2c}\log2\le\int_x^{2x}J^{\rm M}_{Z1}(y)^2y^{-2c}\,dy/y$ for $x\ge1$; so $J^{\rm M}_{Z1}$ is $O_c(x^{c})$ and $dJ^{\rm M}_{Z1}$ is a positive measure on $[1,\infty)$ with $\int x^{-\sigma}dJ^{\rm M}_{Z1}<\infty$ for $\sigma>1$.  By Theorem~\CGref{thm:mellin-counts}, $\zeta_{Z1}(s)=\exp\int x^{-s}dJ^{\rm M}_{Z1}=\int u^{-s}d\mathsf A(u)$, where $\mathsf A=\sum_{r\ge0}\frac1{r!}(dJ^{\rm M}_{Z1})^{*r}$ (multiplicative convolution powers) is a positive measure on $[1,\infty)$.  (ii)$\Rightarrow$(i) is Theorem~\CGref{thm:positivity}.
\end{proof}

In Theorem~\CGref{thm:positivity} the Dirichlet series is not used: (H1) is replaced by the positivity of the representing measure, the functional equation gives $\varpi_{Z1}\le1$ through \CGeqref{eq:lp-inequality}, and the Hadamard product gives $\varpi_{Z1}\ge1$ (Theorem~\CGref{thm:varpi-Z1}).  The function \CGeqref{eq:extremal-function} is extremal for the integers: for $\mathsf A=\sum_{n\ge1}\delta_n$ both inequalities in \CGeqref{eq:lp-inequality} are equalities in the limit, since $\mathsf g(n)=0$ for $n\ge1$ and $\hat{\mathsf g}(n)=0$ for $n\ge2$.  At the clock end $\varpi_{\RVM}<1$ (Theorem~\CGref{thm:varpi-less-one}); unconditionally, the density of $dJ_{\RVM}$ takes negative values in every cell $(n,n+1)$ with $n$ large (Theorem~\CGref{thm:clock-density}).

\begin{lemma}[Mellin form and summation formula for every member]\label{lem:summation-all}
There is $c_1\ge2$ such that for every $\lambda\in[0,1]$ there is a unique measurable function $\mathsf A_\lambda$ on $(0,\infty)$, vanishing on $(0,1)$, with $x^{-c}\mathsf A_\lambda\in L^2\bigl((1,\infty),dx/x\bigr)$ for every $c>c_1$ and
\begin{equation}\label{eq:mellin-A}
 \zeta_\lambda(s)=s\int_1^\infty\mathsf A_\lambda(x)\,x^{-s-1}\,dx\qquad(\Re s>c_1).
\end{equation}
If $\zeta_\lambda$ has a representation \CGeqref{eq:representing-measure}, then $\mathsf A_\lambda(x)=\mathsf A([1,x])$ almost everywhere.  With the pairing $\langle\mathsf A_\lambda,\mathsf g\rangle:=-\int_1^\infty\mathsf A_\lambda(x)\,\mathsf g'(x)\,dx$, every even Schwartz function $\mathsf g$ satisfies
\begin{equation}\label{eq:summation-all}
 \mathsf g(0)+2\langle\mathsf A_\lambda,\mathsf g\rangle=\hat{\mathsf g}(0)+2\langle\mathsf A_\lambda,\hat{\mathsf g}\rangle .
\end{equation}
Equivalently, the even tempered distribution $\delta_0+d\mathsf A_\lambda(|u|)$ on $\R$ is its own Fourier transform.
\end{lemma}

\begin{proof}
By Proposition~\CGref{prop:family-clock} the bound \CGeqref{eq:Z1-growth} holds for every $\zeta_\lambda$ with constants independent of $\lambda$, so, as in the proof of Proposition~\CGref{prop:Mellin-Mertens}, $|\zeta_\lambda(s)|\le C(|t|+\sigma)^{1/4}$ for $\Re s>c_1$, and the argument of Theorem~\CGref{thm:mellin-counts} applied to $\zeta_\lambda(s)/s$ gives $\mathsf A_\lambda$, its vanishing on $(0,1)$ and its uniqueness.  For a positive measure $\mathsf A$ as in \CGeqref{eq:representing-measure}, partial integration gives \CGeqref{eq:mellin-A} with $\mathsf A([1,x])$, and uniqueness applies.

Let $c>c_1$.  In the variable $u=\log x$, the function $\mathsf A_\lambda(e^u)e^{-cu}$ has Fourier transform $\zeta_\lambda(c+it)/(c+it)$, and $-\mathsf g'(e^u)e^{(c+1)u}$ has Mellin transform $\int_0^\infty(-\mathsf g'(x))x^{s}dx=s\,\mathsf g^{\sharp}(s)$ on $\Re s=c$; Parseval's formula gives
\begin{equation}\label{eq:parseval-A}
 \langle\mathsf A_\lambda,\mathsf g\rangle=\frac1{2\pi i}\int_{(c)}\zeta_\lambda(s)\,\mathsf g^{\sharp}(s)\,ds .
\end{equation}
$\zeta_\lambda$ is bounded by a power of $|t|$ on every vertical strip (finite order, the functional equation and the bound on $\Re s\ge c$), and $\mathsf g^{\sharp}$ decays faster than any power.  Move the line to $\Re s=1-c$.  The poles of $\mathsf g^{\sharp}$ at $s=-2,-4,\dots$ are cancelled by the trivial zeros of $\zeta_\lambda$; the remaining poles are $s=1$ and $s=0$, exactly as in \CGeqref{eq:summation-shift}.  On the new line \CGeqref{eq:Z1-FE} and \CGeqref{eq:tate-local} give $\zeta_\lambda(s)\mathsf g^{\sharp}(s)=\zeta_\lambda(1-s)\hat{\mathsf g}^{\sharp}(1-s)$, and \CGeqref{eq:parseval-A} for $\hat{\mathsf g}$ gives \CGeqref{eq:summation-all}.
\end{proof}

\begin{proposition}[The value at infinity as one number]\label{prop:one-number}
Let $\mathsf g_\epsilon$ be as in the proof of Theorem~\CGref{thm:positivity}, $\mathsf w_\epsilon:=\hat{\mathsf g}_\epsilon-\mathsf g_\epsilon$, which is non-negative on $|u|\ge1$, and $\mathsf A'_\lambda(x):=\mathsf A_\lambda(x)-\varpi_\lambda$ for $x\ge1$.  For every $\lambda\in[0,1]$ and $\epsilon>0$,
\begin{equation}\label{eq:one-number-eps}
 \varpi_\lambda\,\mathsf w_\epsilon(1)=\tfrac12\bigl(1-\hat{\mathsf g}_\epsilon(0)\bigr)-\langle\mathsf A'_\lambda,\mathsf w_\epsilon\rangle,\qquad \langle\mathsf A'_\lambda,\mathsf w_\epsilon\rangle:=-\int_1^\infty\mathsf A'_\lambda(x)\,\mathsf w_\epsilon'(x)\,dx,
\end{equation}
the limit $\mathsf W(\lambda):=\lim_{\epsilon\to0}\langle\mathsf A'_\lambda,\mathsf w_\epsilon\rangle$ exists, and
\begin{equation}\label{eq:one-number}
 \varpi_\lambda-1=-\frac{\mathsf W(\lambda)}{\hat{\mathsf g}(1)},\qquad \hat{\mathsf g}(1)=\frac16+\frac1{4\pi^2}.
\end{equation}
Consequently $\mathsf W(0)\le0$; $\mathsf W(1)>0$, since this is equivalent to $\varpi_{\RVM}<1$ (Theorem~\CGref{thm:varpi-less-one}); and (D) holds if and only if $\mathsf W(0)\ge0$.  If $\mathsf A_\lambda$ is non-decreasing, then $\langle\mathsf A'_\lambda,\mathsf w_\epsilon\rangle=\int_{(1,\infty)}\mathsf w_\epsilon\,d\mathsf A_\lambda\ge0$, which is Theorem~\CGref{thm:positivity}.
\end{proposition}

\begin{proof}
By \CGeqref{eq:summation-all} for $\mathsf g_\epsilon$, with $\mathsf g_\epsilon(0)=1$, $2\langle\mathsf A_\lambda,\mathsf w_\epsilon\rangle=1-\hat{\mathsf g}_\epsilon(0)$.  Since $\mathsf A_\lambda=\varpi_\lambda+\mathsf A'_\lambda$ on $[1,\infty)$ and $-\int_1^\infty\mathsf w_\epsilon'=\mathsf w_\epsilon(1)$, $\langle\mathsf A_\lambda,\mathsf w_\epsilon\rangle=\varpi_\lambda\mathsf w_\epsilon(1)+\langle\mathsf A'_\lambda,\mathsf w_\epsilon\rangle$; this is \CGeqref{eq:one-number-eps}.  As $\epsilon\to0$, $\mathsf w_\epsilon(1)=\hat{\mathsf g}_\epsilon(1)\to\hat{\mathsf g}(1)$, because $\mathsf g_\epsilon(1)=0$, and $\frac12(1-\hat{\mathsf g}_\epsilon(0))\to\frac12(1-\hat{\mathsf g}(0))=\hat{\mathsf g}(1)$ by \CGeqref{eq:lp-value}; hence the limit exists and \CGeqref{eq:one-number} holds.  The sign statements follow from \CGeqref{eq:one-number}: at $\lambda=0$ from $\varpi_{Z1}\ge1$ (Theorem~\CGref{thm:varpi-Z1}), and at $\lambda=1$ from $\varpi_1=\varpi_{\RVM}$, so that $\mathsf W(1)>0$ is the same statement as $\varpi_{\RVM}<1$.
\end{proof}

\begin{theorem}[The endpoint dominates $\zeta$ on the real axis]\label{thm:real-axis}
For every real $\sigma>1$,
\begin{equation}\label{eq:real-axis}
 \zeta_{Z1}(\sigma)\ \ge\ \zeta(\sigma),
\end{equation}
and equality for one $\sigma>1$ holds if and only if $\zeta_{Z1}=\zeta$.  In particular
\begin{equation}\label{eq:sigma-two}
 \prod_{k\ge1}\Bigl(1+\frac2{\frac14+\gamma_k^2}\Bigr)\ \ge\ \frac{\pi}{3},\qquad\text{with equality if and only if }(\mathrm D).
\end{equation}
\end{theorem}

\begin{proof}
By the proof of Theorem~\CGref{thm:varpi-Z1}, $\zeta_{Z1}(\sigma)/\zeta(\sigma)=\mathscr R_{Z1}(\sigma)$ is the convergent product of the factors \CGeqref{eq:pair-ratio}, with $W=(\sigma-\frac12)^2+\gamma^2$, which is strictly increasing in $\sigma>\frac12$, hence on $\sigma\ge1$.  At $\sigma=1$, $W=\frac14+\gamma^2$ and the factor equals $1$ by the definition of $c_\gamma$.  The logarithmic derivative in $W$ of $W_\rho^2/((W_\rho-d_\rho^2)^2+4d_\rho^2\gamma^2)$ is
\begin{equation}\label{eq:pair-monotone}
 \frac2W-\frac{2(W-d_\rho^2)}{(W-d_\rho^2)^2+4d_\rho^2\gamma^2}=\frac{2d_\rho^2\,(4\gamma^2+d_\rho^2-W)}{W\bigl((W-d_\rho^2)^2+4d_\rho^2\gamma^2\bigr)},
\end{equation}
positive for $W<4\gamma^2+d_\rho^2$ and negative afterwards, and the factor tends to $c_\gamma>1$ as $W\to\infty$.  A function that equals $1$ at $\sigma=1$, increases and then decreases to a limit greater than $1$ is greater than $1$ for every $\sigma>1$.  Hence every factor of the convergent product $\mathscr R_{Z1}(\sigma)$ is $\ge1$, and every factor attached to a pair with $d_\rho\ne0$ is $>1$; so $\mathscr R_{Z1}(\sigma)\ge1$, with equality if and only if no such factor exists, that is, if and only if the product is empty and $\mathscr R_{Z1}\equiv1$.  For $\sigma=2$: $q(2)=2$, $G(2)=2/\pi$ and $P_\zeta(2)=G(2)\zeta(2)=\pi/3$, so \CGeqref{eq:real-axis} reads $P_{Z1}(2)\ge\pi/3$.
\end{proof}

\begin{proposition}[The final step on the open path]\label{prop:open-path-varpi}
The following are equivalent: \textup{(i)} (D); \textup{(ii)} $\varpi_\lambda<1$ for every $\lambda\in(0,1]$; \textup{(iii)} $\varpi_{\lambda_m}\le1$ for some sequence $\lambda_m\to0^+$; \textup{(iv)} $\mathsf W(\lambda)>0$ for every $\lambda\in(0,1]$. The implication (i)$\Rightarrow$(ii) uses $\varpi_{\RVM}<1$ (Theorem~\CGref{thm:varpi-less-one}), which is the case $\lambda=1$ of (ii); the implications (ii)$\Rightarrow$(iii)$\Rightarrow$(i) and (ii)$\Leftrightarrow$(iv) do not use it.
\end{proposition}

\begin{proof}
$f(\lambda)=\log\varpi_\lambda$ is continuous and strictly convex on $[0,1]$ (Proposition~\CGref{prop:family-flow}: $f'=-\Upsilon$, $\Upsilon'<0$), $f(1)=\log\varpi_{\RVM}<0$ by Theorem~\CGref{thm:varpi-less-one}, and $f(0)\ge0$ (Theorem~\CGref{thm:varpi-Z1}).  If $f(0)=0$, convexity gives $f(\lambda)<(1-\lambda)f(0)+\lambda f(1)<0$ for $0<\lambda<1$.  If $f(\lambda_m)\le0$ along $\lambda_m\to0^+$, then $f(0)\le0$ by continuity, so $f(0)=0$, which is (D) by Theorem~\CGref{thm:varpi-Z1}.  (ii)$\Leftrightarrow$(iv) is \CGeqref{eq:one-number}.
\end{proof}

\subsection{The crossing point of the family and what convexity can give}
\label{subsec:crossing}

Write $f(\lambda):=\log\varpi_\lambda$ for $0\le\lambda\le1$.  By Proposition~\CGref{prop:family-flow}, $f$ is $C^1$ with $f'=-\Upsilon$ and $\Upsilon'<0$, so $f$ is strictly convex; $f(1)=\log\varpi_{\RVM}<-0.020615$ by Theorem~\CGref{thm:varpi-less-one}, and $f(0)=\log\varpi_{Z1}\ge0$, with equality exactly under RH, by Theorem~\CGref{thm:varpi-Z1}.

\begin{proposition}[The crossing point]\label{prop:crossing}
\begin{enumerate}[label=\textup{(\roman*)}]
\item $\Upsilon(0)>-\log\varpi_{\RVM}>0.020615$, without hypothesis.
\item If RH holds, then for every $\lambda\in(0,1]$
\begin{equation}\label{eq:crossing-refuters}
 f(\lambda)\le\lambda f(1)\qquad\text{and}\qquad f(\lambda)+\lambda\Upsilon(\lambda)\le0 .
\end{equation}
\item If RH fails, there is exactly one $\lambda^*\in(0,1)$ with $\varpi_{\lambda^*}=1$; $f>0$ on $[0,\lambda^*)$, $f<0$ on $(\lambda^*,1]$, and
\begin{equation}\label{eq:crossing-bounds}
 \frac{f(0)}{\Upsilon(0)}\le\lambda^*\le\frac{f(0)}{f(0)+|f(1)|}.
\end{equation}
With the verified zeros (Proposition~\CGref{prop:varpi-verified}), $\lambda^*<4\cdot10^{-13}/0.020615<1.95\cdot10^{-11}$.
\end{enumerate}
\end{proposition}

\begin{proof}
(i) Proposition~\CGref{prop:family-flow} gives $\Upsilon(0)>\log L+\log\varpi_{Z1}\ge\log L=-\log\varpi_{\RVM}$.  (ii) Under RH $f(0)=0$; the chord of the convex function $f$ gives $f(\lambda)\le(1-\lambda)f(0)+\lambda f(1)=\lambda f(1)$, and the tangent at $\lambda$ gives $0=f(0)\ge f(\lambda)-\lambda f'(\lambda)=f(\lambda)+\lambda\Upsilon(\lambda)$.  (iii) If $f(0)>0>f(1)$, $f$ has a zero; if it had two, $\lambda_1<\lambda_2$, convexity on $[\lambda_1,1]$ would give $f(\lambda_2)\le\frac{1-\lambda_2}{1-\lambda_1}f(\lambda_1)+\frac{\lambda_2-\lambda_1}{1-\lambda_1}f(1)<0$.  The chord at $\lambda^*$ gives $0\le(1-\lambda^*)f(0)+\lambda^*f(1)$, which is the upper bound, and the tangent at $0$ gives $0=f(\lambda^*)\ge f(0)-\lambda^*\Upsilon(0)$, which is the lower bound.
\end{proof}

For a fixed $\lambda$ a violation of either inequality \CGeqref{eq:crossing-refuters} would refute RH, and their validity is necessary, not sufficient.  Along a sequence $\lambda_m\to0^+$ the first one is sufficient: it gives $f(0)=\lim f(\lambda_m)\le0$, which is Proposition~\CGref{prop:open-path-varpi}(iii).  If RH fails, the set on which $f>0$ is the single interval $[0,\lambda^*)$, of length $\lambda^*<2\cdot10^{-11}$ (Proposition~\CGref{prop:crossing}(iii)); the information that decides RH is therefore asymptotic in $\lambda\to0^+$ and of the same nature as the information at $\lambda=0$.

\begin{remark}[What convexity cannot give]\label{rem:convexity-limits}
For a convex function, values in the interior bound the interior from the end points (chord), and the end points only from below (tangents): $f(0)\ge f(\lambda)+\lambda\Upsilon(\lambda)$.  The only upper bound for $f(0)$ obtained from convexity on $[0,1]$ is $f(0)<f(1)+\Upsilon(0)$, which cannot give $f(0)\le0$ because $\Upsilon(0)>|f(1)|$ by (i).  If $f$ were extended as a strictly convex function to $[-c,1]$, the chord through $-c$ and $1$ would give $f(0)<U_c:=\frac{cf(1)+f(-c)}{1+c}$, but $U_c>f(0)\ge0$ for every $c>0$, whether $f(0)=0$ or $f(0)>0$; and $U_c\to f(0)$ as $c\to0^+$.  The same holds for the other convexity principles that bound interior values by boundary values (three lines, Phragm\'en--Lindel\"of, Hardy's convexity, Lyapunov's inequality).  The properties used in Proposition~\CGref{prop:crossing} are also satisfied by a family whose endpoint has $f(0)>0$; Proposition~\CGref{prop:kappa-family} gives one.
\end{remark}

\begin{proposition}[A family with the same properties and the opposite endpoint]\label{prop:kappa-family}
Let $\kappa_1=12.6$ and, for $0\le\lambda\le1$, let $\boldsymbol\kappa(\lambda)$ be the list $\bigl((1-\lambda)\kappa_1+\lambda\tau_1,\tau_2,\tau_3,\dots\bigr)$, $\zeta^{\boldsymbol\kappa}_\lambda=P_{\boldsymbol\kappa(\lambda)}/G$ with $P_{\boldsymbol\kappa}(s)=\prod_k\bigl(1+q(s)/(\frac14+\kappa_k^2)\bigr)$, and $\varpi^{\boldsymbol\kappa}_\lambda=\lim_{\sigma\to+\infty}\zeta^{\boldsymbol\kappa}_\lambda(\sigma)$.  Then every $\zeta^{\boldsymbol\kappa}_\lambda$ has order one, the functional equation \CGeqref{eq:Z1-FE} and its nontrivial zeros exactly at $\frac12\pm i\kappa_k(\lambda)$; $\zeta^{\boldsymbol\kappa}_1=\zeta_{\RVM}$ and $\varpi^{\boldsymbol\kappa}_1=\varpi_{\RVM}<1$; $\lambda\mapsto\log\varpi^{\boldsymbol\kappa}_\lambda$ is strictly convex; and $\varpi^{\boldsymbol\kappa}_0>1$.
\end{proposition}

\begin{proof}
Only the first factor differs from $P^{\crit}_{\RVM}$, so the first three properties are those of $\zeta_{\RVM}$ (Theorem~\CGref{thm:canonical}), and $\zeta^{\boldsymbol\kappa}_\lambda/\zeta_{\RVM}=\frac{1+q(s)/(\frac14+\kappa_1(\lambda)^2)}{1+q(s)/(\frac14+\tau_1^2)}\to\frac{\frac14+\tau_1^2}{\frac14+\kappa_1(\lambda)^2}$ as $\sigma\to+\infty$, with $\kappa_1(\lambda)=(1-\lambda)\kappa_1+\lambda\tau_1$.  With $\zeta_{\RVM}(\sigma)\to\varpi_{\RVM}$ (Theorem~\CGref{thm:halfplane-RVM}), $\log\varpi^{\boldsymbol\kappa}_\lambda=\log\varpi_{\RVM}+\log(\frac14+\tau_1^2)-\log(\frac14+\kappa_1(\lambda)^2)$, whose second derivative in $\lambda$ is $2(\tau_1-\kappa_1)^2(\kappa_1(\lambda)^2-\frac14)/(\frac14+\kappa_1(\lambda)^2)^2>0$.  For the endpoint we bound $\varpi_{\RVM}$ from below.  By \CGeqref{eq:varpi-def}, $\log\varpi_{\RVM}=\log P^{\crit}_{\RVM}(\frac12)+\nu_3+\nu_4$.  First, $\log P^{\crit}_{\RVM}(\frac12)=-\sum_k\log(1+\frac1{4\tau_k^2})\ge-\frac14\sum_k\tau_k^{-2}$; from $k-\frac12=\NR(\tau_k)\le\frac{\tau_k}{2\pi}\log\frac{\tau_k}{2\pi}\le(\frac{\tau_k}{2\pi})^{3/2}$ (because $\tau_k\ge\tau_1>\frac{7\pi}4$ and $\log y\le\sqrt y$), $\tau_k\ge2\pi(k-\frac12)^{2/3}$ and $\sum_k\tau_k^{-2}\le\frac1{4\pi^2}\bigl(2^{4/3}+3\cdot2^{1/3}\bigr)<0.1596$.  Second, $\nu_3>-0.007925$ by \CGeqref{eq:nu3-closed}.  Third, by \CGeqref{eq:nu4-closed} and Lemma~\CGref{lem:saw-primitive}, $|\nu_4|\le\frac{2\pi}{4\log2}\int_{\tau_0}^\infty\tau^{-2}d\tau=\frac1{8\log2}<0.18034$.  Hence $\log\varpi_{\RVM}>-0.0399-0.007925-0.18034>-0.2283$ and $\varpi_{\RVM}>0.7958$.  Since $\tau_1>14.52134$, $\frac{\frac14+\tau_1^2}{\frac14+\kappa_1^2}>\frac{211.119}{159.01}>1.3277$, and $\varpi^{\boldsymbol\kappa}_0>0.7958\cdot1.3277>1.056$.
\end{proof}

With the floating-point value $\varpi_{\RVM}=0.96194$, already $\kappa_1=14$ gives $\varpi^{\boldsymbol\kappa}_0\approx1.035$.  The list $\boldsymbol\kappa(0)$ differs from the clock in one entry, by $1.93$, and lies at bounded distance from it.  Proposition~\CGref{prop:kappa-family} does not concern $\zeta$: it shows that the properties ``zeros on the critical line, functional equation, order one, bounded displacement, $C^1$ strictly convex $\log\varpi_\lambda$, the endpoint $\lambda=1$ equal to the clock with $\varpi<1$'' do not imply $\varpi_0\le1$, equivalently $\mathsf W(0)\ge0$ in Proposition~\CGref{prop:one-number}.  A transport lemma from the clock end to $\lambda=0$ must therefore use that the list at $\lambda=0$ is the list of the ordinates of $\zeta$; in the form of Theorem~\CGref{thm:rung-zero-sum} it must show that the sum over the zeros off the line vanishes.  The positivity transported along the family is the number $\mathsf W(1)>0$; the monotonicity of the measure $\mathsf A_\lambda$ is a hypothesis of Proposition~\CGref{prop:one-number} and, at $\lambda=0$, is equivalent to (D) (Corollary~\CGref{cor:positivity}).

A statement that holds as an identity only at the endpoint cannot be reached as a limit of the same identity along the path, where it is false (Proposition~\CGref{prop:nonconstant}); an inequality can.  Proposition~\CGref{prop:open-path-varpi} is of this form: the inequality $\varpi_\lambda\le1$ on the open path passes to the limit, and the opposite inequality $\varpi_{Z1}\ge1$ holds at the endpoint.  By Theorem~\CGref{thm:real-axis} the same scheme is available at every real $\sigma>1$: (D) holds if and only if $\zeta_{Z1}(\sigma)\le\zeta(\sigma)$ for one $\sigma>1$.

\begin{remark}[The hypothesis added to the functional equation]\label{rem:hypothesis-X}
The functional equation, the order and the zeros on the critical line do not characterize $\zeta$: every $\zeta_\lambda$ has them.  A converse theorem for $\zeta_{Z1}$ needs one further hypothesis on the function $\mathsf A_{Z1}$ of Lemma~\CGref{lem:summation-all}, and each of the following suffices and is equivalent to (D): $\mathsf A_{Z1}$ is a step function with jumps $\mathsf b_n$ at the integers satisfying $\sum_{n\le X}|\mathsf b_n|=O_\eps(X^{1+\eps})$, so that $\zeta_{Z1}$ is a Dirichlet series absolutely convergent in $\Re s>1$ (Hamburger's theorem, in the form of Theorem~\CGref{thm:cell-form}; the growth condition on the jumps is part of the hypothesis, since Hamburger's theorem requires the absolute convergence); $\mathsf A_{Z1}$ is non-decreasing (Theorem~\CGref{thm:positivity}); $\mathsf W(0)\ge0$ (Proposition~\CGref{prop:one-number}), which is the weakest of the three, being equivalent to $\varpi_{Z1}\le1$.  The inequality \CGeqref{eq:lp-inequality} is the one-dimensional case of the linear programming bound of Cohn and Elkies \cite{CohnElkies2003}.  By Lemma~\CGref{lem:summation-all} the measure $\delta_0+d\mathsf A_\lambda(|u|)$ is self-dual for every $\lambda$; measures on $\R$ whose support and spectrum are both uniformly discrete are classified in \cite{LevOlevskii2015}.
\end{remark}

\subsection{The density of the clock}
\label{subsec:clock-density}

Let $\varXi(t):=t\log t-t$ for $t>0$, and for $M\ge1$, $0<t_0\le\frac12$ let
\begin{equation}\label{eq:U-def}
 \mathcal U(M,t_0):=\int_{t_0}^{\infty}e^{-iM\varXi(t)}\,dt:=\lim_{T\to\infty}\int_{t_0}^{T}e^{-iM\varXi(t)}\,dt .
\end{equation}

\begin{lemma}[Stationary phase for $\varXi$]\label{lem:stationary-U}
The limit \CGeqref{eq:U-def} exists, and there is an absolute constant $C_{\mathcal U}$ with
\begin{equation}\label{eq:U-asymptotic}
 \Bigl|\mathcal U(M,t_0)-\sqrt{\frac{2\pi}M}\,e^{i(M-\pi/4)}\Bigr|\le\frac{C_{\mathcal U}}M\qquad(M\ge1,\ 0<t_0\le\tfrac12).
\end{equation}
\end{lemma}

\begin{proof}
$\varXi'(t)=\log t$ is increasing and unbounded, so the limit exists by the first-derivative test.  $\varXi$ is convex with minimum $\varXi(1)=-1$.  Put $\mathsf w:=\operatorname{sgn}(t-1)\sqrt{2(\varXi(t)+1)}$.  Near $t=1$, $\varXi(t)+1=\frac12(t-1)^2(1+O(t-1))$ with an analytic factor, so $t\mapsto\mathsf w$ is analytic with derivative $1$ at $t=1$; elsewhere $\mathsf w\,d\mathsf w=\log t\,dt$.  Hence $t\mapsto\mathsf w$ is an increasing analytic bijection of $(0,\infty)$ onto $(-\sqrt2,\infty)$, with
\begin{equation}\label{eq:U-substitution}
 \mathsf p(\mathsf w):=\frac{dt}{d\mathsf w}=\frac{\mathsf w}{\log t},\qquad \mathsf p(0)=1,\qquad \mathsf q(\mathsf w):=\frac{\mathsf p(\mathsf w)-1}{\mathsf w}=\frac1{\log t}-\frac1{\mathsf w}.
\end{equation}
$\mathsf q$ is analytic on $(-\sqrt2,\infty)$, bounded, tends to $0$ as $\mathsf w\to\infty$, and $\mathsf q'=-\mathsf w/(t\log^3t)+\mathsf w^{-2}$.  Away from $\mathsf w=0$ one has $\int|\mathsf w|\,|t\log^3t|^{-1}d\mathsf w=\int dt/(t\log^2t)$, which converges at $t=0$ and at $t=\infty$, and $\int\mathsf w^{-2}d\mathsf w$ converges at $\infty$; so $\int_{-\sqrt2}^\infty|\mathsf q'|\,d\mathsf w<\infty$.  With $\mathsf w_0=\mathsf w(t_0)\le\mathsf w(\frac12)=-\sqrt{1-\log2}<-\frac12$ and $e^{-iM\varXi}=e^{iM}e^{-iM\mathsf w^2/2}$,
\begin{equation}\label{eq:U-split}
 e^{-iM}\,\mathcal U(M,t_0)=\int_{\mathsf w_0}^\infty e^{-iM\mathsf w^2/2}\,d\mathsf w+\int_{\mathsf w_0}^\infty e^{-iM\mathsf w^2/2}\,\mathsf w\,\mathsf q(\mathsf w)\,d\mathsf w .
\end{equation}
The first integral is $\int_{\R}e^{-iM\mathsf w^2/2}d\mathsf w=\sqrt{2\pi/M}\,e^{-i\pi/4}$ minus $\int_{-\infty}^{\mathsf w_0}$, and the latter is at most $4/(M|\mathsf w_0|)\le8/M$ by the first-derivative test \cite[Lemma~4.2]{Titchmarsh1986}.  Since $\mathsf w\,e^{-iM\mathsf w^2/2}=\frac iM\frac d{d\mathsf w}e^{-iM\mathsf w^2/2}$, integration by parts bounds the second integral by $M^{-1}\bigl(\sup|\mathsf q|+\int|\mathsf q'|\bigr)$.
\end{proof}

For $m\ge1$ and $y>1$ put
\begin{equation}\label{eq:clock-harmonic}
 \mathcal R_m(y):=-2(-1)^m\,m^{-3/2}\,y^{\frac1{2m}-\frac12}\cos\Bigl(2\pi m\,y^{1/m}-\frac{(7m+1)\pi}4\Bigr),
\end{equation}
so that $\mathcal R_1(y)=2\cos(2\pi y)$ and $\mathcal R_2(y)=-2^{-1/2}y^{-1/4}\cos(4\pi\sqrt y+\frac\pi4)$.

\begin{theorem}[The density of the clock]\label{thm:clock-density}
$J_{\RVM}$ is continuously differentiable on $(2,\infty)$, and for $y\ge16$
\begin{equation}\label{eq:clock-density}
 J_{\RVM}'(y)=\frac1{\log y}+2\cos(2\pi y)-\frac{y^{-1/4}}{\sqrt2}\cos\Bigl(4\pi\sqrt y+\frac\pi4\Bigr)+O\bigl(y^{-1/3}\bigr).
\end{equation}
On $[2,\infty)$ the measure $d\pi_{\RVM}$ is the sum of the atoms $\frac{\mu(j)}jJ_{\RVM}(2)\,\delta_{2^j}$, $j\ge1$, and of a density that is continuous off the points $2^j$ and equals $J_{\RVM}'(y)+O(y^{-1/2})$.
\end{theorem}

\begin{proof}
\emph{An exact series.}  We use the notation of the proof of Theorem~\CGref{thm:unit-cell}: $v=\log y$, $S_\alpha(v)=\sum_ke^{(iv-\alpha)\tau_k}$, the decomposition \CGeqref{eq:S-eps-split} and the harmonics $K_m^\pm$ of \CGeqref{eq:K-harmonics}, with phases $\Psi_m^\pm(\tau)=v\tau\pm2\pi m\NR(\tau)$.  Put
\begin{equation}\label{eq:H-harmonic}
 \mathcal H_m^\pm(v):=\int_{\tau_0}^\infty e^{i\Psi_m^\pm(\tau)}\,d\tau ,
\end{equation}
improper integrals that converge, locally uniformly in $v\ge\log2$, by the first-derivative test, because $(\Psi_m^\pm)'$ is monotone and unbounded.  As $\alpha\downarrow0$, $K_m^\pm\to iv\,\mathcal H_m^\pm(v)$, and the first two terms of \CGeqref{eq:S-eps-split} tend to a continuous function $S_{\rm sm}(v)$ with $|S_{\rm sm}(v)|\le C/v+\frac12$.  The bounds for $K_m^\pm$ in that proof are uniform in the Abel parameter and, for $K_m^-$, include the factor displayed in \CGeqref{eq:K-minus-bound}; only the finite range $m\le v/\log2$ can contain the stationary point.  After this finite range is separated, the remaining bounds are summable against $1/(2\pi m)$, so dominated convergence in \CGeqref{eq:damped-sum} and \CGeqref{eq:K-harmonics}
\begin{equation}\label{eq:S-zero}
 \sum_kD_k=\int_{u_1}^{u_2}\frac{e^{v/2}}v\,S_0(v)\,dv,\qquad S_0(v):=S_{\rm sm}(v)+\sum_{m\ge1}\frac{(-1)^m\,v}{2\pi m}\bigl(\mathcal H_m^-(v)-\mathcal H_m^+(v)\bigr),
\end{equation}
the series converging uniformly on compact sets.  Thus $S_0$ is continuous, $J_{\RVM}$ is continuously differentiable on $(2,\infty)$, and by \CGeqref{eq:J-RVM}
\begin{equation}\label{eq:J-prime-exact}
 J_{\RVM}'(y)=\frac1{\log y}\Bigl(1-\frac1{y(y^2-1)}-\frac2{\sqrt y}\,\Re S_0(\log y)\Bigr).
\end{equation}

\emph{The terms without a stationary point.}  $(\Psi_m^+)'\ge v+m\log2$, so $|\mathcal H_m^+|\le4/(v+m\log2)$ and $\sum_m\frac v{2\pi m}|\mathcal H_m^+|\ll1+\log v$.  For $m>v/\log2$, $(\Psi_m^-)'\le v-m\log2<0$ on $[\tau_0,\infty)$; the first-derivative test, and the second-derivative test on $[\tau_0,2\tau_0]$ with $|(\Psi_m^-)''|=m/\tau$, give $|\mathcal H_m^-|\ll\min\bigl((m\log2-v)^{-1},m^{-1/2}\bigr)$, and $\sum_{m>v/\log2}\frac v{2\pi m}|\mathcal H_m^-|\ll1+\log v$.

\emph{The terms $3\le m\le v/\log2$.}  By the argument for \CGeqref{eq:K-minus-bound}, $|\mathcal H_m^-|\le\frac8{m\log2}+16\sqrt\pi\,e^{v/(2m)}m^{-1/2}$, so these terms contribute to $S_0$ at most $C\,v\sum_{m\ge3}\bigl(m^{-2}+m^{-3/2}e^{v/(2m)}\bigr)\ll v\,e^{v/6}$.

\emph{The terms $m=1,2$.}  The stationary point of $\Psi_m^-$ is $\mathsf t_m:=2\pi e^{v/m}=2\pi y^{1/m}$.  With $\tau=\mathsf t_mt$ and \CGeqref{eq:clock}, $\Psi_m^-(\mathsf t_mt)=-m\mathsf t_m\varXi(t)-\frac{7\pi m}4$, hence
\begin{equation}\label{eq:H-by-U}
 \mathcal H_m^-(v)=\mathsf t_m\,e^{-7\pi im/4}\,\mathcal U\bigl(m\mathsf t_m,\tau_0/\mathsf t_m\bigr),\qquad \tau_0/\mathsf t_m=2e^{-v/m}\le\tfrac12\quad(y\ge4^m).
\end{equation}
By Lemma~\CGref{lem:stationary-U} with $M=m\mathsf t_m$,
\begin{equation}\label{eq:H-asymptotic}
 \frac{(-1)^mv}{2\pi m}\,\mathcal H_m^-(v)=(-1)^m\,\frac v{m^{3/2}}\,y^{\frac1{2m}}\exp\Bigl(i\Bigl(2\pi m\,y^{1/m}-\frac{(7m+1)\pi}4\Bigr)\Bigr)+O\Bigl(\frac v{m^2}\Bigr).
\end{equation}
Inserting the four estimates into \CGeqref{eq:J-prime-exact} gives $J_{\RVM}'(y)=\frac1{\log y}+\mathcal R_1(y)+\mathcal R_2(y)+O(y^{-1/3})$, which is \CGeqref{eq:clock-density}.

\emph{The prime coordinate.}  By \CGeqref{eq:pi-RVM}, on each interval $(2^r,2^{r+1})$ the function $\pi_{\RVM}$ is the finite sum of the continuously differentiable functions $\frac{\mu(j)}jJ_{\RVM}(x^{1/j})$, $j\le r$; the terms $j\ge2$ have derivative $\frac{\mu(j)}{j^2}x^{1/j-1}J_{\RVM}'(x^{1/j})$, and $|J_{\RVM}'|\le C_J$ by \CGeqref{eq:unit-cell}, so their sum is $O(x^{-1/2})$.  At $x=2^j$ the term $j$ jumps by $\frac{\mu(j)}jJ_{\RVM}(2)$.
\end{proof}

\begin{corollary}[Cell masses and cell defects at the clock end]\label{cor:clock-cells}
For $n\ge17$, with $\mathbf 1_{\mathrm{Pow}_2}(n)=1$ if $n=2^j$ for an integer $j\ge1$ and $0$ otherwise (the product below is $0$ when $n$ is not a power of $2$, and $\log_2n=j$ when it is),
\begin{equation}\label{eq:clock-cell-sharp}
 c_n^{\RVM}=\int_{n-1}^n\Bigl(\frac1{\log x}+2\cos(2\pi x)\Bigr)dx+\mathbf 1_{\mathrm{Pow}_2}(n)\,\frac{\mu(\log_2n)}{\log_2n}\,J_{\RVM}(2)+O\bigl(n^{-1/4}\bigr),
\end{equation}
\begin{equation}\label{eq:clock-cell-smooth}
 \tilde c_n(1)=\frac1{\log n}+\mathbf 1_{\mathrm{Pow}_2}(n)\,\frac{\mu(\log_2n)}{\log_2n}\,J_{\RVM}(2)+O\bigl(n^{-1/4}\bigr),
\end{equation}
and for $0<\Re s\le\sigma_1$
\begin{equation}\label{eq:clock-defect}
 |\mathfrak h_n(1;s)|\le C_{\sigma_1}\bigl(|s|\,n^{-\sigma-4/3}+|s|(1+|s|)\,n^{-\sigma-2}\bigr).
\end{equation}
In particular Hypothesis~\CGref{hyp:cell-remainder} is proved (it holds along the integers $n$ that are not of the form $2^j$ with $\mu(j)\ne0$, which is its whole range), so Proposition~\CGref{prop:eventual-positive} is unconditional, and $c_n^{\RVM}\sim1/\log n$, $0<c_n^{\RVM}<1$ for all large such $n$; at $n=2^j$, $c_n^{\RVM}=\bigl(1+\mu(j)J_{\RVM}(2)\log2\bigr)/\log n+O(n^{-1/4})$.
\end{corollary}

\begin{proof}
\CGeqref{eq:clock-cell-sharp}: $c_n^{\RVM}=\pi_{\RVM}(n)-\pi_{\RVM}(n-1)$ is the integral of the density over $(n-1,n)$ plus the atom at $n$, and $\int_{n-1}^n(|\mathcal R_2|+O(x^{-1/3}))dx=O(n^{-1/4})$.  \CGeqref{eq:clock-cell-smooth}: $\varphi_n(2^j)=1$ if $n=2^j$ and $0$ otherwise; $\int\varphi(t)\,dt/\log(n+t)=\frac1{\log n}+O(\frac1{n\log^2n})$ because $\mathsf M_0=1$; and $\int\varphi(t)\,2\cos(2\pi(n+t))\,dt=2\mathsf C_0=0$ by \CGeqref{eq:partition-fourier}.  \CGeqref{eq:clock-defect}: the test function $\varphi_n(x)(n^{-s}-x^{-s})$ vanishes at the atoms, and $n^{-s}-(n+t)^{-s}=s\,n^{-s-1}t+O_{\sigma_1}(|s|(1+|s|)n^{-\sigma-2})$ for $|t|\le1$.  The first moment of the density is $\int t\varphi(t)\pi_{\RVM}'(n+t)\,dt=\frac{\mathsf M_1}{\log n}+2\mathsf C_1+O(n^{-1/3})=O(n^{-1/3})$, because $\mathsf M_1=\mathsf C_1=0$, $\mathcal R_2(n+t)=\mathcal R_2(n)+O(n^{-3/4})$ and the density is bounded.  The last statements follow from \CGeqref{eq:clock-cell-sharp} and Lemma~\CGref{lem:cell-cancel}.
\end{proof}

By \CGeqref{eq:clock-density}, $\frac{d\pi_{\RVM}}{dy}=\bigl(\frac1{\log y}-1\bigr)+\bigl(1+2\cos(2\pi y)\bigr)+O(y^{-1/4})$: the clock carries the first harmonic of the comb $\sum_n\delta_n=1+2\sum_{q\ge1}\cos(2\pi qy)$ of all integers, with weight $1$, on a smooth background.  This harmonic comes from the ordinates $\tau_k$ near the height $2\pi y$, the stationary point of $\Phi_1^-$; by \CGeqref{eq:partition-fourier} it does not enter the cell masses \CGeqref{eq:clock-cell-smooth}, which are $1/\log n$ to leading order, and it enters the cell defects only through the moments $\mathsf C_r$, $r\ge2$.

\subsection{The cell masses in the heights of the ordinates}
\label{subsec:height-kernels}

\begin{proposition}[Height kernel of the cell masses]\label{prop:height-kernels}
For every $\eps>0$, $\lambda\in[0,1]$ and $n\ge3$ that is not a power of $2$,
\begin{equation}\label{eq:height-kernel}
 \tilde c_n(\lambda)=\mathfrak m_n-\frac2{\sqrt n\,\log n}\sum_{k\ge1}\hat{\varphi}\Bigl(\frac{\gamma_k(\lambda)}{2\pi n}\Bigr)\cos\bigl(\gamma_k(\lambda)\log n\bigr)+E_n(\lambda),
\end{equation}
where $\sum_{X<n\le2X}|E_n(\lambda)|^2\le C_\eps X^{1/2+\eps}$ uniformly in $\lambda$ (sum over the $n$ that are not powers of $2$).  The kernel $\hat{\varphi}(\gamma/2\pi n)$ vanishes at the heights $\gamma=2\pi qn$, $q=1,2,\dots$
\end{proposition}

\begin{proof}
The term $j=1$ of $W_n$ in \CGeqref{eq:VW} is $\int\varphi(t)(n+t)^{-1/2+i\gamma}\,dt/\log(n+t)$.  Uniformly for $|t|\le1$ and $0<\gamma\le n^{1+\eps}$, $(1+t/n)^{-1/2+i\gamma}=e^{i\gamma t/n}\bigl(1+O((1+\gamma/n)/n)\bigr)$ and $\frac1{\log(n+t)}=\frac1{\log n}(1+O(\frac1n))$ on $|t|\le1$, so this term equals $\frac{n^{-1/2+i\gamma}}{\log n}\hat{\varphi}(\frac\gamma{2\pi n})+O\bigl(n^{-3/2}(1+\gamma/n)\bigr)$, because $\int\varphi(t)e^{i\gamma t/n}dt=\hat{\varphi}(\gamma/2\pi n)$ for even $\varphi$.  Summed over the $O(n^{1+\eps}\log n)$ ordinates $\gamma_k(\lambda)\le n^{1+\eps}$ the error is at most $O(n^{1+\eps}\log n)\cdot n^{-3/2}(1+n^{\eps})=O(n^{-1/2+2\eps}\log n)=O(n^{-1/2+3\eps})$.  For $\gamma>n^{1+\eps}$ both the term and $\hat{\varphi}(\gamma/2\pi n)$ are $O_A((n/\gamma)^A)$, as in the proof of Theorem~\CGref{thm:mean-square-cells}.  The terms $j\ge2$ have sum of squares $O(X^{1/2+3\eps})$ over $X<n\le2X$ by \CGeqref{eq:cell-block}.  The kernel vanishes at $\gamma=2\pi qn$, $q\in\Z\setminus\{0\}$, because the partition of unity gives $\hat{\varphi}(q)=0$ for $q\in\Z\setminus\{0\}$, \CGeqref{eq:partition-fourier}.
\end{proof}

Since $\gamma_k(\lambda)\log n=\tau_k\log n-(1-\lambda)\eta_k\log n$, the passage along the family rotates each clock phase $n^{i\tau_k}$ by the angle $-(1-\lambda)\eta_k\log n$: up to the change of $\hat{\varphi}$ over a distance $O(1/n)$, the sum in \CGeqref{eq:height-kernel} is $\Re\,\Xi_n\bigl((1-\lambda)\log n\bigr)$ with
\begin{equation}\label{eq:Xi-rotation}
 \Xi_n(\Omega):=\sum_{k\ge1}\hat{\varphi}\Bigl(\frac{\tau_k}{2\pi n}\Bigr)n^{i\tau_k}e^{-i\Omega\eta_k}\qquad(0\le\Omega\le\log n).
\end{equation}
At $\Omega=0$ this is the clock sum, whose real part is $O(n^{1/4}\log n)$: the kernel of the term $j=1$ is, up to the $O(1/n)$ change of kernel, the ordinate part of $\int\varphi(t)J_{\RVM}'(n+t)\,dt$ in \CGeqref{eq:J-prime-exact}, and Theorem~\CGref{thm:clock-density} gives $\int\varphi(t)\bigl[2\cos(2\pi(n+t))+\mathcal R_2(n+t)+O((n+t)^{-1/3})\bigr]dt=O(n^{-1/4})$ pointwise because $C_0=0$ (the bound is not obtained from $\tilde c_n(1)$, which carries the terms $j\ge2$ whose remainder is controlled only in mean square), because the resonant heights $\tau_k\approx2\pi n$ are zeros of the kernel; the endpoint of the family is $\Omega=\log n$.

\begin{proposition}[The flow of a cell mass and the count difference]\label{prop:flow-count}
For every $n\ge3$ that is not a power of $2$,
\begin{equation}\label{eq:flow-count}
 \tilde c_n(0)-\tilde c_n(1)=-2\int_0^1\sum_{k\ge1}\eta_k\,\Im V_n\bigl(\gamma_k(\lambda)\bigr)\,d\lambda=-2\,\Im\int_{\tau_0}^\infty V_n(t)\bigl(N_\zeta(t)-n_{\RVM}(t)\bigr)\,dt .
\end{equation}
\end{proposition}

\begin{proof}
The first equality is \CGeqref{eq:cell-mass-W} and the continuity of $\tilde c_n$ at $\lambda=0$.  By \CGeqref{eq:cell-mass-W}, $\tilde c_n(0)-\tilde c_n(1)=2\Re\sum_k\bigl(W_n(\tau_k)-W_n(\gamma_k)\bigr)=2\Re\int_{\tau_0}^\infty W_n\,d(n_{\RVM}-N_\zeta)$.  Both counts vanish at $\tau_0$, their difference is $O(\log t)$ by Lemma~\CGref{lem:bounded-displacement}, and $W_n$ is rapidly decreasing; integration by parts with $W_n'=iV_n$ gives $2\Re\bigl(i\int V_n(N_\zeta-n_{\RVM})\,dt\bigr)$.
\end{proof}

\begin{proposition}[Height kernel of the cell defects]\label{prop:defect-kernel}
For every $\eps>0$, $\lambda\in[0,1]$, $\Re s>0$ and $n\ge3$ that is not a power of $2$, with an implied constant uniform for $s$ in compact subsets of $\Re s>0$ and for $\lambda\in[0,1]$,
\begin{equation}\label{eq:defect-kernel}
 \mathfrak h_n(\lambda;s)=-\frac{s\,n^{-s-3/2}}{\pi\log n}\sum_{k\ge1}\hat{\varphi}'\Bigl(\frac{\gamma_k(\lambda)}{2\pi n}\Bigr)\sin\bigl(\gamma_k(\lambda)\log n\bigr)+O_{s,\eps}\bigl(n^{-\sigma-3/4+\eps}\bigr).
\end{equation}
\end{proposition}

\begin{proof}
For $|t|\le1$, $n^{-s}-(n+t)^{-s}=s\,n^{-s-1}t+\psi_2(t)$, where $\psi_2$ and all its derivatives in $t$ are $O_s(n^{-\sigma-2})$.  \emph{The term $j=1$.}  The part of the density that does not involve the ordinates contributes $O_s(n^{-\sigma-2})$, because $\int t\varphi(t)\,dt/\log(n+t)=O(1/(n\log^2n))$ ($\mathsf M_1=0$).  The pairing of $\varphi\psi_2$ with the ordinate sum is $O_s(n^{-\sigma-2})$ times $O(n^{1/2+\eps})$, as in Remark~\CGref{rem:A6-scale}.  For the real function $t\varphi(t)$, the computation in the proof of Proposition~\CGref{prop:height-kernels}, with $\int t\varphi(t)e^{2\pi i\xi t}dt=-\frac i{2\pi}\hat{\varphi}'(\xi)$ for even $\varphi$ and $-2\Re\bigl(-\frac i{2\pi}n^{i\gamma}\bigr)=-\frac1\pi\sin(\gamma\log n)$, gives the sum in \CGeqref{eq:defect-kernel} with an error $O_s(n^{-\sigma-3/2+3\eps})$.  \emph{The terms $j\ge2$.}  The test function $\varphi_n(y^j)(n^{-s}-y^{-js})$ is $O(|s|n^{-\sigma-1})$ on a support of length $O(n^{1/j-1})$; using the uniform upper bound $O(n^{1+\eps}\log n)$ for the number of relevant ordinates, which is more than sufficient for every $j\ge2$, and $y^{-1/2}=O(n^{-1/(2j)})$, these terms are $O_s(n^{-\sigma-1+1/(2j)+2\eps})$, at most $O_s(n^{-\sigma-3/4+2\eps})$.  The bounds used depend on $s$ only through $|s|$ and $\sigma$, hence are uniform on compact subsets of $\Re s>0$; replacing $\eps$ by $\eps/3$ gives the stated error.
\end{proof}

The kernel $\hat{\varphi}$ of the masses vanishes at the heights $2\pi qn$ and the kernel $\hat{\varphi}'$ of the defects in general does not: the masses and the defects of a cell read complementary heights of the ordinates.  At $\lambda=1$ the resonant heights $\tau_k\approx2\pi n$ contribute to $\sum_kn^{i\tau_k}$ the real quantity $-\sqrt n\log n\,e^{2\pi in}$ (proof of Theorem~\CGref{thm:clock-density}), which the sine in \CGeqref{eq:defect-kernel} does not see: the first harmonic of the clock is centred at the integers.

\begin{remark}[The scale of the displacements]\label{rem:scale-test}
By Proposition~\CGref{prop:nonconstant}, a property of the sequence $(\eta_k)$ that yields (D) is false for every sequence $(c\,\eta_k)$ with $0\le c<1$.  Properties that are invariant under $\eta_k\mapsto c\,\eta_k$ with $c>0$ (the signs of the $\eta_k$, the non-vanishing of the $\eta_k$, their boundedness) hold along the whole open path and do not suffice; the property fixes the size of the displacements.  In the form \CGeqref{eq:Xi-rotation} it is a condition on the angles $\eta_k\log n$ modulo $2\pi$, jointly in $k$ and $n$.  The construction contains one exact relation of this kind: for a simple ordinate $\gamma_k$, the midpoint value of the count is $N_\zeta^*(\gamma_k)=\frac12\bigl(N_\zeta(\gamma_k+0)+N_\zeta(\gamma_k-0)\bigr)=k-\frac12$, and \CGeqref{eq:midpoint-condition} and \CGeqref{eq:block-count} give $\NR(\tau_k)-\NR(\gamma_k)=N_\zeta^{*}(\gamma_k)-\NR(\gamma_k)$, that is,
\begin{equation}\label{eq:eta-exact}
 \eta_k\,\NR'(\xi_k)=N_\zeta^{*}(\gamma_k)-\NR(\gamma_k)\qquad\text{for some }\xi_k\text{ between }\gamma_k\text{ and }\tau_k .
\end{equation}
\end{remark}

\subsection{The Gaussian defect of the atomization}
\label{subsec:gaussian-defect}

\begin{lemma}[Vertical growth of the joining factor]\label{lem:vertical-growth}
For every $\lambda\in[0,1]$ and $\frac12<\sigma_0<\sigma_1$, $\tilde h_\lambda(\sigma+it)=o(|t|)$ as $|t|\to\infty$, uniformly for $\sigma_0\le\sigma\le\sigma_1$.
\end{lemma}

\begin{proof}
By \CGeqref{eq:joining-factor} and Proposition~\CGref{prop:first-cell}(ii) it suffices to treat the double series.  Fix $c\in(\frac12,\sigma_0)$.  With $I_n:=\int_{n-1}^{n+1}|\pi_\lambda-\mathfrak L|^2dx$ as in the proof of Theorem~\CGref{thm:atomized-continuation}, \CGeqref{eq:cell-defect-bound} gives $\sum_{m\ge1}\frac1m\sum_{n>N}|\mathfrak h_n(\lambda;ms)|\le C|s|\sum_{n>N}n^{-\sigma_0-1}(1+I_n^{1/2})$ and $\sum_{m>N}\frac1m\sum_{n\ge2}|\mathfrak h_n(\lambda;ms)|\le C|s|\sum_{n\ge2}n^{-N\sigma_0-1}(1+I_n^{1/2})$; both are at most $\epsilon_N|s|$ with $\epsilon_N\to0$ as $N\to\infty$.  For the finitely many pairs $m,n\le N$, with $E=\pi_\lambda-\mathfrak L$ as in the proof of Theorem~\CGref{thm:atomized-continuation},
\begin{equation}\label{eq:cell-RL}
\begin{split}
 \mathfrak h_n(\lambda;ms)={}&\int\varphi_n(x)\bigl(n^{-ms}-x^{-ms}\bigr)d\mathfrak L(x)-\int E(x)\varphi_n'(x)\bigl(n^{-ms}-x^{-ms}\bigr)dx\\
 &-ms\int E(x)\varphi_n(x)\,x^{-ms-1}dx .
\end{split}
\end{equation}
The first two terms are bounded in $t$, uniformly for $\sigma_0\le\sigma\le\sigma_1$.  In $u=\log x$ the last integral is the Fourier transform at $mt$ of the amplitude $E(e^u)\varphi_n(e^u)e^{-m\sigma u}$; as $\sigma$ ranges over $[\sigma_0,\sigma_1]$ these amplitudes form a compact subset of $L^1$ (a continuous image of a compact interval, because $E\in L^1(n-1,n+1)$), so the Riemann--Lebesgue convergence to $0$ as $|t|\to\infty$ is uniform in $\sigma$.  Hence $\limsup_{|t|\to\infty}|t|^{-1}\sup_\sigma|\tilde h_\lambda(\sigma+it)|\le2\epsilon_N$ for every $N$.
\end{proof}

For $y>0$ and $c>1$ put
\begin{equation}\label{eq:gaussian-defect}
 \Theta_\lambda(y):=\frac1{4\pi i}\int_{(c)}\zeta_\lambda(s)\,\Gamma(\tfrac s2)\,\pi^{-s/2}\,y^{-s/2}\,ds,\qquad
 \mathfrak G_\lambda(y):=\sum_{N\ge1}\tilde a_N(\lambda)\,e^{-\pi N^2y}-\Theta_\lambda(y).
\end{equation}
For $\zeta$ itself $\Theta(y)=\sum_{n\ge1}e^{-\pi n^2y}$, and if $\zeta_\lambda$ has a representation \CGeqref{eq:representing-measure}, then $\Theta_\lambda(y)=\int e^{-\pi u^2y}\,d\mathsf A(u)$.  $\mathfrak G_\lambda$ is the change produced by the atomization in the Gaussian count of the integers of the system.

\begin{proposition}[Gaussian defect]\label{prop:gaussian-defect}
Let $\lambda\in[0,1]$.
\begin{enumerate}[label=\textup{(\roman*)}]
\item $\Theta_\lambda(y)=y^{-1/2}\,\Theta_\lambda(1/y)+\frac12\bigl(y^{-1/2}-1\bigr)$ for $y>0$, and $\Theta_\lambda(y)=O_A(y^{-A})$ as $y\to\infty$ for every $A$.
\item The coefficients $\tilde a_N(\lambda)$ satisfy the theta relation of $\zeta$,
\begin{equation}\label{eq:theta-atomized}
 \sum_{N\ge1}\tilde a_N(\lambda)\,e^{-\pi N^2y}=y^{-1/2}\sum_{N\ge1}\tilde a_N(\lambda)\,e^{-\pi N^2/y}+\tfrac12\bigl(y^{-1/2}-1\bigr)\qquad(y>0),
\end{equation}
if and only if
\begin{equation}\label{eq:defect-selfdual}
 \mathfrak G_\lambda(y)=y^{-1/2}\,\mathfrak G_\lambda(1/y)\qquad(y>0).
\end{equation}
Under the hypotheses of Theorem~\CGref{thm:pair-criterion}, \CGeqref{eq:defect-selfdual} holds.
\item For every $\delta>0$, as $y\to0^+$,
\begin{equation}\label{eq:defect-expansion}
 \mathfrak G_\lambda(y)=\tfrac12\bigl(e^{\tilde h_\lambda(1)}-1\bigr)\,y^{-1/2}+O_\delta\bigl(y^{-1/4-\delta}\bigr),
\end{equation}
and $e^{\tilde h_\lambda(1)}=\Res_{s=1}\tilde\zeta_{\Z,\lambda}(s)$.
\item If \CGeqref{eq:defect-selfdual} holds, then $\tilde h_\lambda(1)=0$.
\end{enumerate}
\end{proposition}

\begin{proof}
(i) $\zeta_\lambda$ is bounded by a power of $|t|$ on vertical strips (proof of Lemma~\CGref{lem:summation-positive}; on $\Re s\ge c$ use Proposition~\CGref{prop:family-clock}), and $\Gamma(\frac s2)$ decays exponentially.  Moving the line to $\Re s=1-c$ crosses $s=1$, where the integrand has residue $\Gamma(\frac12)\pi^{-1/2}y^{-1/2}=y^{-1/2}$, and $s=0$, with residue $2\zeta_\lambda(0)=-1$; the prefactor $\frac1{4\pi i}=\frac12\cdot\frac1{2\pi i}$ turns these residues into the term $\frac12(y^{-1/2}-1)$.  On the new line \CGeqref{eq:Z1-FE} and $y^{-s/2}=y^{-1/2}(1/y)^{-(1-s)/2}$ give $y^{-1/2}\Theta_\lambda(1/y)$.  The decay follows by moving the line to the right, to $\Re s=2A$ with $2A>c$: no pole is crossed, since $\zeta_\lambda(s)\Gamma(\frac s2)$ is holomorphic on $\Re s>1$, the integrals converge uniformly by the same bounds, and $|y^{-s/2}|=y^{-A}$ on the new line.
(ii) By $e^{-x}=\frac1{2\pi i}\int_{(c/2)}\Gamma(w)x^{-w}dw$ and Corollary~\CGref{cor:abscissa-one}, $\sum_N\tilde a_N(\lambda)e^{-\pi N^2y}=\frac1{4\pi i}\int_{(c)}\tilde\zeta_{\Z,\lambda}(s)\Gamma(\frac s2)\pi^{-s/2}y^{-s/2}ds$; subtracting (i) from the theta relation gives the equivalence.  Under the hypotheses of Theorem~\CGref{thm:pair-criterion}, $\tilde\zeta_{\Z,\lambda}=\zeta$.
(iii) By (ii), $2\mathfrak G_\lambda(y)=\frac1{2\pi i}\int_{(c)}\zeta_\lambda(s)\bigl(e^{\tilde h_\lambda(s)}-1\bigr)\Gamma(\frac s2)\pi^{-s/2}y^{-s/2}ds$.  By Lemma~\CGref{lem:vertical-growth}, $|e^{\tilde h_\lambda(s)}|\le e^{o(|t|)}$ on $\frac12+2\delta\le\Re s\le c$, so the line can be moved to $\Re s=\frac12+2\delta$.  The only pole crossed is $s=1$, where $\zeta_\lambda$ has residue $1$ and $\tilde h_\lambda$ is holomorphic; the residue is $(e^{\tilde h_\lambda(1)}-1)y^{-1/2}$, and the remaining integral is $O_\delta(y^{-1/4-\delta})$.
(iv) $\sum_N\tilde a_N(\lambda)e^{-\pi N^2y}=O(e^{-\pi y})$ as $y\to\infty$, so by (i) $\mathfrak G_\lambda(y)=O_A(y^{-A})$ there; then \CGeqref{eq:defect-selfdual} gives $\mathfrak G_\lambda(y)=O_A(y^{A})$ as $y\to0^+$ for every $A$, and \CGeqref{eq:defect-expansion} forces $e^{\tilde h_\lambda(1)}=1$; since $\tilde h_\lambda$ is real on the real axis, $\tilde h_\lambda(1)$ is real, so $\tilde h_\lambda(1)=0$ and not $2\pi im$ with $m\ne0$.
\end{proof}

\begin{corollary}[The density of the atomized integers]\label{cor:residue-condition}
$\sum_{N\le x}\tilde a_N(\lambda)$ has Gaussian mean density $e^{\tilde h_\lambda(1)}$ in the sense of \CGeqref{eq:defect-expansion} (a statement about the Gaussian weights $e^{-\pi N^2y}$ as $y\to0^+$, not a Tauberian statement about the partial sums), and (D) implies
\begin{equation}\label{eq:residue-condition}
 \tilde h_0(1)=0,\qquad\text{that is,}\qquad \mathfrak u_0(1)=\sum_{m\ge1}\frac1m\sum_{n\ge2}\mathfrak h_n(0;m).
\end{equation}
\end{corollary}

\begin{proof}
Under (D), $\tilde h_0\equiv0$ (Theorem~\CGref{thm:closing}); the second form is \CGeqref{eq:joining-factor} at $s=1$.
\end{proof}

The identity $\tilde h_0(\sigma)=0$ for real $\sigma$ has two distinguished members: $\sigma=+\infty$ is $\varpi_{Z1}=1$, the first coefficient (Corollary~\CGref{cor:product-identity}), and $\sigma=1$ is \CGeqref{eq:residue-condition}, the density of the integers.

By Proposition~\CGref{prop:gaussian-defect}(i), the function \CGeqref{eq:theta-defect} is $T_\lambda(y)=\mathfrak G_\lambda(y)-y^{-1/2}\mathfrak G_\lambda(1/y)$, so \CGeqref{eq:defect-selfdual} is $T_\lambda\equiv0$.  At the clock end the expansion \CGeqref{eq:defect-expansion} holds with a smaller remainder: for every $\delta>0$,
\begin{equation}\label{eq:defect-expansion-clock}
 \mathfrak G_1(y)=\tfrac12\bigl(e^{\tilde h_1(1)}-1\bigr)\,y^{-1/2}+O_\delta\bigl(y^{-\delta}\bigr)\qquad(y\to0^+),
\end{equation}
because $\tilde h_1$ is holomorphic in $\Re s>0$ with $\tilde h_1(\sigma+it)=o(|t|)$ uniformly on $\sigma_0\le\sigma\le\sigma_1$ for every $\sigma_0>0$ (the proof of Lemma~\CGref{lem:vertical-growth} with the bound $|\mathfrak h_n(1;s)|\le C|s|n^{-\sigma-1}$ of Theorem~\CGref{thm:atomized-continuation}(iv)), and the line of integration can be moved to $\Re s=2\delta$.

\subsection{Entire dependence on \texorpdfstring{$\lambda$}{lambda} and the exceptional set}
\label{subsec:entire-lambda}

\begin{proposition}[The cell masses as entire functions of $\lambda$]\label{prop:entire-lambda}
Let $n\ge3$ not be a power of $2$.
\begin{enumerate}[label=\textup{(\roman*)}]
\item The real-axis function $\lambda\mapsto\tilde c_n(\lambda)$ has the entire extension
\[
 \widetilde C_n(\lambda):=\mathfrak m_n-\sum_k\Bigl(W_n(\gamma_k+\lambda\eta_k)+\overline{W_n(\gamma_k+\bar\lambda\eta_k)}\Bigr),
\]
which is of exponential type at most $C_\eta\log(n+1)$ and agrees with $\tilde c_n(\lambda)$ for real $\lambda$.
\item For every $r\ge0$, $\eps>0$ and $\lambda\in[0,1]$,
\begin{equation}\label{eq:lambda-derivatives}
 \partial_\lambda^r\tilde c_n(\lambda)=-2\,\Re\sum_{k\ge1}\eta_k^{\,r}\,W_n^{(r)}\bigl(\gamma_k(\lambda)\bigr),\qquad \bigl|\partial_\lambda^r\tilde c_n(\lambda)\bigr|\le C_{r,\eps}\bigl(C_\eta\log(n+1)\bigr)^r\,n^{1/2+\eps}.
\end{equation}
\item $|\tilde c_n(\lambda)-\tilde c_n(\lambda')|\le C_\eps|\lambda-\lambda'|\,n^{1/2+\eps}$ on $[0,1]$.  Hence, for sets of integers $n$ that are not powers of $2$, $\tilde c_n(0)=O_\eps(n^\eps)$ if and only if $\tilde c_n(\lambda)=O_\eps(n^\eps)$ uniformly for $0\le\lambda\le n^{-1/2}$, and $\tilde c_n(\lambda)=O_\eps(n^{\eps})$ uniformly for $1-n^{-1/2}\le\lambda\le1$.
\end{enumerate}
\end{proposition}

\begin{proof}
For complex $\gamma$, $|y^{i\gamma}|\le e^{|\Im\gamma|\log y}$ with $\log y\le\log(n+1)$ on the support of $\varphi_n(y^j)$, and repeated integration by parts in \CGeqref{eq:VW} gives $|W_n^{(r)}(\gamma)|\le C_{A,r,n}\,e^{|\Im\gamma|\log(n+1)}(1+|\Re\gamma|)^{-A}$ for every $A$.  With $\gamma=\gamma_k+\lambda\eta_k$, $|\Im\gamma|\le C_\eta|\Im\lambda|$, the two reflected series defining $\widetilde C_n$ converge locally uniformly in $\lambda\in\C$, so $\widetilde C_n$ is entire; on the real axis it equals $\tilde c_n$.  Termwise differentiation on the real axis gives the first formula in \CGeqref{eq:lambda-derivatives}.  For real $\lambda$, $W_n^{(r)}$ carries the factor $(i\log y)^r$ with $\log y\le\log(n+1)$, the ordinates $\gamma_k(\lambda)>n^{1+\eps}$ contribute $O(n^{-1})$, and each of the $O(n^{1+\eps}\log n)$ remaining ordinates contributes $O\bigl((\log(n+1))^rn^{-1+1/(2j)}\bigr)$ to the term $j$, as in Remark~\CGref{rem:A6-scale}; this is the bound in \CGeqref{eq:lambda-derivatives}.  (iii) is the case $r=1$ and Corollary~\CGref{cor:clock-cells}.
\end{proof}

The bounds \CGeqref{eq:lambda-derivatives} do not by themselves control $\tilde c_n(0)$ from $\tilde c_n(1)$: the function $n^{1/2}\sin\bigl((1-\lambda)C_\eta\log n\bigr)$ is entire of the same type, satisfies the same bounds, vanishes at $\lambda=1$ and has size $n^{1/2}$ at $\lambda=0$ for infinitely many $n$.  The Taylor expansion at the clock end, $\tilde c_n(0)=\sum_{r\ge0}\frac{(-1)^r}{r!}\partial_\lambda^r\tilde c_n(1)$, has the terms $-2\Re\sum_k\eta_k^rW_n^{(r)}(\tau_k)$, clock sums weighted by the powers of the displacements.

\begin{proposition}[The exceptional set]\label{prop:exceptional-set}
The set
\begin{equation}\label{eq:exceptional-set}
 \mathsf E:=\bigl\{\lambda\in[0,1]:\ \tilde\zeta_{\Z,\lambda}=\zeta\bigr\}
\end{equation}
is finite, $1\notin\mathsf E$, and (D) holds if and only if $0\in\mathsf E$.  For every $\lambda\notin\mathsf E$ the hypotheses of Theorem~\CGref{thm:pair-criterion} do not hold together, $T_\lambda\not\equiv0$ and $\mathsf V_\sigma(\lambda)>0$.
\end{proposition}

\begin{proof}
If $\lambda\in\mathsf E$, then $\sum_n\tilde c_n(\lambda)\sum_m\frac{n^{-ms}}m=\log\zeta(s)=\sum_p\sum_m\frac{p^{-ms}}m$ for $\Re s>1$.  Comparing the coefficients of $N^{-s}$ in increasing order of $N$ determines the exponents uniquely, so $\tilde c_n(\lambda)=\mathbf 1_{\rm prime}(n)$ for every $n$.  Choose a composite $n$, not a power of $2$, so large that $\tilde c_n(1)>0$ by \CGeqref{eq:clock-cell-smooth}.  Then $\tilde c_n$ is an entire function of $\lambda$ (Proposition~\CGref{prop:entire-lambda}) that does not vanish identically; the zeros of a nonzero entire function have no finite accumulation point, so only finitely many of them lie in the compact interval $[0,1]$; they contain $\mathsf E$, and $1$ is not one of them.  $0\in\mathsf E$ implies (D) by Theorem~\CGref{thm:double-hamburger}(ii), and (D) implies $\tilde\zeta_{\Z,0}=\zeta_{Z1}=\zeta$ by Theorem~\CGref{thm:closing}.  The last statement is Theorems~\CGref{thm:pair-criterion} and~\CGref{thm:double-hamburger}(ii).
\end{proof}

Proposition~\CGref{prop:nonconstant} excludes that $\zeta_\lambda$ and $\tilde\zeta_{\Z,\lambda}$ coincide for $\lambda>0$; Proposition~\CGref{prop:exceptional-set} excludes, outside a finite set, that $\tilde\zeta_{\Z,\lambda}$ coincides with $\zeta$.  For $0<\lambda<1$ outside $\mathsf E$ it is not decided which of the two hypotheses of Theorem~\CGref{thm:pair-criterion} fails; at $\lambda=1$ it is the symmetry (Proposition~\CGref{prop:clock-asymmetry}).


\subsection{The final step in the cell masses}
\label{subsec:cell-mass-form}

\begin{corollary}[The final step is one number per cell]\label{cor:cell-mass-form}
The following are equivalent:
\begin{enumerate}[label=\textup{(\roman*)}]
\item (D);
\item $\tilde c_n(0)=\mathbf 1_{\rm prime}(n)$ for every $n\ge2$;
\item there is a list $\boldsymbol{\kappa}\in\mathscr O$ (Section~\CGref{sec:lists}) with $\tilde c_n(\boldsymbol{\kappa})=\mathbf 1_{\rm prime}(n)$ for every $n\ge2$.
\end{enumerate}
\end{corollary}

\begin{proof}
(i)$\Rightarrow$(ii) is Theorem~\CGref{thm:closing}, and (ii)$\Rightarrow$(iii) with the list $(\gamma_k)$.  If (iii) holds, then
\begin{equation}\label{eq:cell-mass-euler}
 \tilde\zeta_{\Z,\boldsymbol{\kappa}}(s)=\prod_{n\ge2}\bigl(1-n^{-s}\bigr)^{-\tilde c_n(\boldsymbol{\kappa})}=\prod_p\bigl(1-p^{-s}\bigr)^{-1}=\zeta(s)\qquad(\Re s>1),
\end{equation}
and Corollary~\CGref{cor:existence-form}, (iii)$\Rightarrow$(i), gives (D).
\end{proof}

Step A7 of the programme asks for the vanishing of the functions $\mathfrak h_n(0;\cdot)$, one function of $s$ for each cell, and for $J_{Z1}=0$ on the first cell.  Corollary~\CGref{cor:cell-mass-form} asks for one real number for each cell; the position of the mass inside the cell and the first cell then follow from Hamburger's theorem.

\begin{proposition}[Every point of the exceptional set is a witness]\label{prop:witness}
For the set $\mathsf E$ of \CGeqref{eq:exceptional-set},
\begin{equation}\label{eq:witness}
 \mathsf E\ne\varnothing\iff(\mathrm D),\qquad\text{so that either }\mathsf E=\varnothing\ \text{ or }\ 0\in\mathsf E .
\end{equation}
If (D) holds, then
\begin{equation}\label{eq:witness-zeros}
 \mathsf E\setminus\{0\}=\bigl\{\lambda\in(0,1]:\ \mathsf D_n(\lambda)=0\ \text{for every }n\ge2\bigr\},\qquad \mathsf D_n(\lambda):=\tilde c_n(\lambda)-\tilde c_n(0),
\end{equation}
and for $n\ge3$ not a power of $2$, the real-axis function $\mathsf D_n(\lambda)=-2\Re\sum_{k\ge1}\bigl(W_n(\gamma_k+\lambda\eta_k)-W_n(\gamma_k)\bigr)$ has the entire extension obtained by replacing each real part by the corresponding holomorphic plus reflected-conjugate pair; this extension is of exponential type at most $C_\eta\log(n+1)$ that vanishes at $\lambda=0$.
\end{proposition}

\begin{proof}
If $\lambda_0\in\mathsf E$, the list $(\gamma_k(\lambda_0))$ belongs to $\mathscr O$ and its atomized zeta function is $\zeta$; Corollary~\CGref{cor:existence-form}, (iii)$\Rightarrow$(i), gives (D), and then $0\in\mathsf E$ by Proposition~\CGref{prop:exceptional-set}.  Under (D), $\tilde c_n(0)=\mathbf 1_{\rm prime}(n)$, and $\lambda\in\mathsf E$ means $\tilde c_n(\lambda)=\mathbf 1_{\rm prime}(n)$ for every $n$ (proof of Proposition~\CGref{prop:exceptional-set}).  The last statement is \CGeqref{eq:cell-mass-W} and Proposition~\CGref{prop:entire-lambda}.
\end{proof}

Thus a proof that $\mathsf E\cap(0,1]=\varnothing$ does not approach the final step: if (D) fails, $\mathsf E$ is empty, and if (D) holds, $\mathsf E\cap(0,1]=\varnothing$ is the uniqueness of the witness on the path, which is the statement that the entire functions $\mathsf D_n$ have no common zero in $(0,1]$.  This uniqueness is not proved; Proposition~\CGref{prop:exceptional-set} gives that $\mathsf E$ is finite and $1\notin\mathsf E$.

\begin{remark}[The dual function of the atomized family]\label{rem:dual-function}
For $\Re s<\frac12$ put $B^\vee_\lambda(s):=\zeta_\lambda(s)\,e^{\tilde h_\lambda(1-s)}$.  By the functional equation of $\zeta_\lambda$,
\begin{equation}\label{eq:dual-function}
 \pi^{-s/2}\Gamma\bigl(\tfrac s2\bigr)B^\vee_\lambda(s)=\pi^{-(1-s)/2}\Gamma\bigl(\tfrac{1-s}2\bigr)\,\tilde\zeta_{\Z,\lambda}(1-s)\qquad(\Re s<\tfrac12),
\end{equation}
so $B^\vee_\lambda$ is the function attached to $\tilde\zeta_{\Z,\lambda}$ by the functional equation.  To ask that $B^\vee_\lambda$ continue to $\Re s>1$ as an absolutely convergent Dirichlet series is to ask for the hypotheses of Theorem~\CGref{thm:pair-criterion}; no monotonicity in $\lambda$ is known, and a parameter such as $\inf\{\lambda>0:B^\vee_\lambda\ \text{is not a Dirichlet series}\}$ is not shown to be a transition point.
\end{remark}

\begin{proposition}[The identity that remains, cell by cell]\label{prop:per-cell}
Let $n\ge17$ not be a power of $2$.  Then $\tilde c_n(0)=\mathbf 1_{\rm prime}(n)$ if and only if
\begin{equation}\label{eq:per-cell}
 \mathbf 1_{\rm prime}(n)-\tilde c_n(1)=-2\,\Im\int_{\tau_0}^\infty V_n(t)\bigl(N_\zeta(t)-n_{\RVM}(t)\bigr)\,dt,\qquad \tilde c_n(1)=\frac1{\log n}+O\bigl(n^{-1/4}\bigr).
\end{equation}
\end{proposition}

\begin{proof}
This is \CGeqref{eq:flow-count} and \CGeqref{eq:clock-cell-smooth}.
\end{proof}

The left side of \CGeqref{eq:per-cell} is explicit.  The right side is the pairing of the count difference with the kernel $V_n$, whose term $j=1$ is $n^{-1/2+it}\hat{\varphi}(t/2\pi n)$ up to $O\bigl(n^{-3/2}(1+t/n)\bigr)$ (proof of Proposition~\CGref{prop:defect-kernel}).  For $n=2^j$ the same identity holds with the cut term of Theorem~\CGref{thm:mean-square-cells} added.  No evaluation of the right side of \CGeqref{eq:per-cell} is given in this paper.

\subsection{Symmetric joining factors}
\label{subsec:symmetric-factors}

\begin{remark}[The symmetry criterion has no slack at $\lambda=0$]\label{rem:pair-no-slack}
At $\lambda=0$ the hypothesis of Theorem~\CGref{thm:pair-criterion} is equivalent to its conclusion.  If (D) holds, then $\tilde h_0\equiv0$ (Theorem~\CGref{thm:closing}), and $e^{\tilde h_0}=1$ is symmetric and of finite order; hence
\begin{equation}\label{eq:pair-no-slack}
 e^{\tilde h_0}\ \text{symmetric of finite order}\iff\tilde h_0\equiv0\iff(\mathrm D),
\end{equation}
and in the same way $T_0\equiv0$ and $\mathsf V_\sigma(0)=0$ are equivalent forms of (D), not weaker conditions.  Moreover \CGeqref{eq:pair-symmetry} presupposes a continuation of $\tilde h_0$ across the line $\Re s=\frac12$ (step A10 of the programme): $\tilde h_0$ is holomorphic in $\Re s>\frac12$, and $s\mapsto1-s$ maps this half-plane onto $\Re s<\frac12$.  Theorems~\CGref{thm:halfline-hamburger} and~\CGref{thm:bridge} below are criteria whose hypotheses are written on the half-line $\sigma\to+\infty$ and on the closed half-plane $\Re s\ge\frac12$.
\end{remark}

\begin{proposition}[Symmetric joining factors in general]\label{prop:dichotomy}
Let $f$ be meromorphic in $\C$ with $(s-1)f(s)$ entire of order at most one, and let $f$ satisfy the functional equation \CGeqref{eq:Z1-FE}.  Let $g=f\,e^{h}$, where $e^{h}$ is entire, without zeros, of finite order, and $e^{h(s)}=e^{h(1-s)}$.  If $g\not\equiv0$ and $g(s)=\sum_{n\ge1}a_nn^{-s}$ converges absolutely in $\Re s>1$, then $g=C_H\zeta$ with $C_H\ne0$ and $e^{h}$ is constant.  If moreover $\lim_{\sigma\to+\infty}g(\sigma)=\lim_{\sigma\to+\infty}f(\sigma)\ne0$, then $e^{h}\equiv1$ and $f=g=C_H\zeta$.
\end{proposition}

\begin{proof}
$g$ satisfies (H0)--(H3) of Subsection~\CGref{subsec:hamburger-table} with the second function equal to $g$, so $g=C_H\zeta$ by Hamburger's theorem.  If $C_H=0$ then $f\,e^{h}=0$, which is excluded when $g\not\equiv0$; assume $C_H\ne0$.  Then $e^{h}\,(s-1)f(s)=C_H(s-1)\zeta(s)$ identically, and we compare the Hadamard factorizations of the two entire functions of order at most one, $(s-1)f(s)=e^{A_1+B_1s}\Pi_1(s)$ and $(s-1)\zeta(s)=e^{A_2+B_2s}\Pi_2(s)$, with $\Pi_i$ the canonical products of genus at most one over their zeros.  Since $e^{h}$ has no zeros, the two functions have the same zeros with the same multiplicities, so $\Pi_1=\Pi_2$ and $e^{h}=C_He^{(A_2-A_1)+(B_2-B_1)s}$: $h$ is a polynomial of degree at most one.  The symmetry gives $h(s)-h(1-s)\in2\pi i\Z$, which forces the degree to be zero.  The last statement compares the limits at $+\infty$.
\end{proof}

\begin{remark}[Order one cannot be dropped]\label{rem:dichotomy-counterexample}
For real $\nu_5\ne0$ put $f(s)=\zeta(s)\exp\bigl(-\nu_5(s-\tfrac12)^2\bigr)$ and $e^{h(s)}=\exp\bigl(\nu_5(s-\tfrac12)^2\bigr)$.  Then $f$ is meromorphic of order two with the functional equation, the pole and the zeros of $\zeta$; $e^{h}$ is entire, without zeros, of finite order and symmetric; $g=f\,e^{h}=\zeta$ is an absolutely convergent Dirichlet series in $\Re s>1$; and $e^{h}$ is not constant.  In Theorem~\CGref{thm:pair-criterion} the conclusion $\tilde h_0\equiv0$ uses that $(s-1)\zeta_{Z1}$ has order one (Proposition~\CGref{prop:Z1-class}).  The location of the zeros of $f$ plays no role in Proposition~\CGref{prop:dichotomy}.
\end{remark}

\subsection{The half-line form of Hamburger's theorem for a pair}
\label{subsec:halfline}

The next two theorems are stated for a general pair: a function with the functional equation and a Dirichlet series, joined on the real half-line only.

\begin{theorem}[The ladder]\label{thm:ladder}
Let $\mathsf d\colon[1,\infty)\to\R$ be measurable with $\int_1^\infty|\mathsf d(x)|\,x^{-\mathsf e_1-1}\,dx<\infty$ for some $\mathsf e_1\ge0$, and put
\begin{equation}\label{eq:ladder-h}
 h(\sigma):=\sigma\int_1^\infty\mathsf d(x)\,x^{-\sigma-1}\,dx\qquad(\sigma>\mathsf e_1).
\end{equation}
For every integer $N\ge2$,
\begin{equation}\label{eq:ladder}
 h(\sigma)=O\bigl(\sigma N^{-\sigma}\bigr)\quad(\sigma\to+\infty)\iff\mathsf d(x)=0\ \text{for almost every }x\in[1,N).
\end{equation}
\end{theorem}

\begin{proof}
If $\mathsf d=0$ almost everywhere on $[1,N)$, then $|h(\sigma)|\le\sigma N^{-(\sigma-\mathsf e_1)}\int_N^\infty|\mathsf d|x^{-\mathsf e_1-1}dx$ for $\sigma>\mathsf e_1$.  Conversely, put $d(u):=\mathsf d(e^u)$, so that $d\,e^{-\mathsf e_1u}\in L^1(0,\infty)$, and
\begin{equation}\label{eq:ladder-Y}
 \mathsf Y(s):=\int_0^{\log N}d(u)\,e^{s(\log N-u)}\,du .
\end{equation}
$\mathsf Y$ is entire of exponential type at most $\log N$, bounded on $\Re s\le0$ by $\int_0^{\log N}|d|$, and $\mathsf Y(\sigma)\to0$ as $\sigma\to-\infty$.  For real $\sigma>\mathsf e_1$,
\begin{equation}\label{eq:ladder-Y-real}
 \mathsf Y(\sigma)=\frac{N^\sigma h(\sigma)}\sigma-\int_{\log N}^\infty d(u)\,e^{-\sigma(u-\log N)}\,du ,
\end{equation}
where the first term is bounded by hypothesis and the second by $N^{\mathsf e_1}\int|d|e^{-\mathsf e_1u}du$; so $\mathsf Y$ is bounded on the real axis.  The Phragm\'en--Lindel\"of theorem in the two quadrants $\{\Re s>0,\ \pm\Im s>0\}$ (angle $\pi/2$, order one, bounded on both sides) shows that $\mathsf Y$ is bounded in $\C$, hence constant, hence $0$.  By the uniqueness of the Laplace transform on $(0,\log N)$, $d=0$ almost everywhere there.
\end{proof}

No functional equation, no continuation of $h$ and no uniformity in $t$ enter Theorem~\CGref{thm:ladder}: the hypothesis lives on $\sigma\to+\infty$.

\begin{theorem}[Hamburger's theorem for a pair, half-line form]\label{thm:halfline-hamburger}
Let $f$ satisfy (H0) and the functional equation \CGeqref{eq:Z1-FE}, and suppose that for $\Re s>1$
\begin{equation}\label{eq:halfline-f}
 \log f(s)=s\int_1^\infty J_f(x)\,x^{-s-1}\,dx,\qquad\int_1^\infty|J_f(x)|\,x^{-\sigma-1}\,dx<\infty\quad(\sigma>1),
\end{equation}
with $J_f$ real and measurable.  Let $g(s)=\sum_{n\ge1}a_nn^{-s}$ converge absolutely in $\Re s>1$, with $a_1=1$, and let $J_g(x):=\sum_{2\le n\le x}b_n$, where $\log g(s)=\sum_{n\ge2}b_nn^{-s}$ for large $\sigma$.  Put $h(\sigma):=\log g(\sigma)-\log f(\sigma)$ for large real $\sigma$.  The following are equivalent:
\begin{enumerate}[label=\textup{(\roman*)}]
\item $h(\sigma)=O_N\bigl(\sigma N^{-\sigma}\bigr)$ as $\sigma\to+\infty$, for every integer $N\ge2$;
\item $f=g$;
\item $f=g=\zeta$.
\end{enumerate}
In this case $J_f$ is almost everywhere the step function $\sum_{p^m\le x}\frac1m$.
\end{theorem}

\begin{proof}
For large $\sigma$, $\log g(\sigma)=\sigma\int_1^\infty J_g(x)x^{-\sigma-1}dx$ by partial summation, so $h$ has the form \CGeqref{eq:ladder-h} with $\mathsf d=J_g-J_f$.  (i)$\Rightarrow$(ii): Theorem~\CGref{thm:ladder} gives $\mathsf d=0$ almost everywhere, so $\log f=\log g$ for large $\sigma$ and $f=g$ by analytic continuation.  (ii)$\Rightarrow$(iii): $g$ then satisfies (H0)--(H3), so $g=C_H\zeta$, and $a_1=1$ gives $C_H=1$.  (iii)$\Rightarrow$(i) because $h\equiv0$.
\end{proof}

\begin{corollary}[The ladder for the pair $(\zeta_\lambda,\tilde\zeta_{\Z,\lambda})$]\label{cor:ladder-pair}
Let $\lambda\in[0,1]$ and $N\ge2$.  Then $\tilde h_\lambda(\sigma)=O(\sigma N^{-\sigma})$ as $\sigma\to+\infty$ if and only if the Mellin count $J^{\rm M}_\lambda$ of $\zeta_\lambda$ coincides almost everywhere on $[1,N)$ with the step function $\sum_{n^m\le x}\tilde c_n(\lambda)/m$; in particular $J^{\rm M}_\lambda$ then vanishes on $(1,2)$ and is constant between consecutive integers below $N$.  Consequently
\begin{equation}\label{eq:ladder-D}
 (\mathrm D)\iff\tilde h_0(\sigma)=O_N\bigl(\sigma N^{-\sigma}\bigr)\ \text{for every }N\ge2 ,
\end{equation}
and at the clock end the first rung already fails, by Theorem~\CGref{thm:varpi-less-one}: $\tilde h_1(\sigma)\to\log L=-\log\varpi_{\RVM}$, which is $\ne0$ exactly when $\varpi_{\RVM}\ne1$ (Proposition~\CGref{prop:atomized-clock}(ii)).
\end{corollary}

\begin{proof}
Theorem~\CGref{thm:ladder} with $\mathsf d=J_g-J^{\rm M}_\lambda$, where $J_g$ is the count of $\log\tilde\zeta_{\Z,\lambda}=\sum_m\frac1m\sum_n\tilde c_n(\lambda)n^{-ms}$; the integrability is Theorem~\CGref{thm:mellin-counts}, whose proof applies to every $\zeta_\lambda$, and Corollary~\CGref{cor:abscissa-one}.  \CGeqref{eq:ladder-D} is Theorem~\CGref{thm:halfline-hamburger} with Corollary~\CGref{cor:abscissa-one} and Theorem~\CGref{thm:cell-form}, and the last statement is Proposition~\CGref{prop:atomized-clock}.
\end{proof}

The joining factor therefore measures the mass of $dJ^{\rm M}_\lambda$ outside the integers, and its rate of decay on the real axis is the first integer below which there is such mass.  The $N$th rung at $\lambda=0$ is the statement that $J^{\rm M}_{Z1}$ is constant on $(N-1,N)$; it is (D) restricted to one cell.

\subsection{The bridge theorem: the phase of the joining factor on the critical line}
\label{subsec:bridge}

\begin{theorem}[Bridge theorem, phase form]\label{thm:bridge}
Let $f$ be meromorphic in $\C$, real on the real axis, with (H0), the functional equation \CGeqref{eq:Z1-FE}, and $\lim_{\sigma\to+\infty}f(\sigma)=\varpi_f>0$.  Let $g(s)=\sum_{n\ge1}a_nn^{-s}$ converge absolutely in $\Re s>1$, with real $a_n$ and $a_1=1$.  Suppose that $g=f\,e^{h}$ in $\Re s>\frac12$ with $h$ holomorphic there and real on $(1,\infty)$, and that
\begin{itemize}
\item[\textup{(B1)}] $\Im h$ extends continuously to $\Re s\ge\frac12$ and vanishes on the line $\Re s=\frac12$;
\item[\textup{(B2)}] there are $\mathsf e_1,\mathsf e_2\ge0$ with
\begin{equation}\label{eq:bridge-B2}
 |h(s)|\le C\,(1+|s|)^{\mathsf e_1}\Bigl(1+\bigl(\sigma-\tfrac12\bigr)^{-\mathsf e_2}\Bigr)\qquad(\Re s>\tfrac12).
\end{equation}
\end{itemize}
Then $h\equiv-\log\varpi_f$, $g=\zeta$ and $f=\varpi_f\,\zeta$.
\end{theorem}

\begin{proof}
\emph{Reflection.}  $v:=\Im h$ is harmonic in $\Re s>\frac12$, continuous on the closed half-plane and $0$ on its boundary.  By the reflection principle \cite[Ch.~4, \S6.5]{Ahlfors1979} the function equal to $v(s)$ for $\Re s\ge\frac12$ and to $-v(1-\bar s)$ for $\Re s<\frac12$ is harmonic in $\C$.  It is the imaginary part of an entire function, which can be chosen equal to $h$ in $\Re s>\frac12$; in $\Re s<\frac12$ it differs from $\overline{h(1-\bar s)}$ by a real constant, and this constant is $0$ because the real part of an entire function is continuous across the line.  Thus $h$ is entire and $h(1-\bar s)=\overline{h(s)}$.  Since $h$ is real on $(1,\infty)$, $h(\bar s)=\overline{h(s)}$, and therefore $h(1-s)=h(s)$.

\emph{Growth.}  By \CGeqref{eq:bridge-B2} and the symmetry, $|h(s)|\le C(1+|s|)^{\mathsf e_1}\bigl(1+|\sigma-\frac12|^{-\mathsf e_2}\bigr)$ for $\Re s\ne\frac12$.  The function $\log|h|$ is subharmonic, so for every $s$
\begin{equation}\label{eq:bridge-mean}
 \log|h(s)|\le\frac1{2\pi}\int_0^{2\pi}\log\bigl|h\bigl(s+e^{i\phi}\bigr)\bigr|\,d\phi\le\log C+\mathsf e_1\log(2+|s|)+\log2+\frac{\mathsf e_2}{2\pi}\int_0^{2\pi}\log^+\frac1{|\sigma-\frac12+\cos\phi|}\,d\phi .
\end{equation}
For real $y$, $\frac1{2\pi}\int_0^{2\pi}\log|y+\cos\phi|\,d\phi\ge-\log2$, so the last integral is at most $2\pi\bigl(\log2+\log(2+|s|)\bigr)$.  Hence $|h(s)|\le C'(2+|s|)^{\mathsf e_1+\mathsf e_2}$ in $\C$, and $h$ is a polynomial.  On the real axis $h(\sigma)=\log g(\sigma)-\log f(\sigma)\to-\log\varpi_f$ as $\sigma\to+\infty$, so the polynomial is the constant $-\log\varpi_f$.

\emph{Hamburger.}  $g=f/\varpi_f$ satisfies (H0), (H1), and (H2)--(H3) with the second function equal to $g$; so $g=C_H\zeta$, and $a_1=1$ gives $C_H=1$.
\end{proof}

The symmetry of the joining factor, its continuation to $\Re s\le\frac12$ and its finite order are not hypotheses of Theorem~\CGref{thm:bridge}: they follow from \textup{(B1)} and \textup{(B2)}, which are written on the closed half-plane $\Re s\ge\frac12$ where the joining factor of the family is defined.  Theorem~\CGref{thm:pair-criterion} assumes the continuation and the symmetry instead, and nothing like \textup{(B2)}; Theorem~\CGref{thm:bridge} derives both from \textup{(B1)}--\textup{(B2)}. Neither set of hypotheses contains the other.

\begin{corollary}[The bridge theorem for the family]\label{cor:bridge-family}
\begin{enumerate}[label=\textup{(\roman*)}]
\item For $0<\lambda\le1$ the joining factor $\tilde h_\lambda$ does not satisfy both \textup{(B1)} and \textup{(B2)}.
\item (D) holds if and only if $\tilde h_0$ satisfies \textup{(B1)} and \textup{(B2)}.
\end{enumerate}
\end{corollary}

\begin{proof}
Theorem~\CGref{thm:bridge} applies with $f=\zeta_\lambda$ (Theorem~\CGref{thm:path}, Proposition~\CGref{prop:family-clock}), $g=\tilde\zeta_{\Z,\lambda}$ (Corollary~\CGref{cor:abscissa-one}) and $h=\tilde h_\lambda$ (Theorem~\CGref{thm:atomized-continuation}(iii)).  If \textup{(B1)} and \textup{(B2)} hold, then $\zeta_\lambda=\varpi_\lambda\zeta$ is an absolutely convergent Dirichlet series in $\Re s>1$, which is excluded for $0<\lambda\le1$ by Proposition~\CGref{prop:nonconstant}; for $\lambda=0$ it is (D).  Conversely (D) gives $\tilde h_0\equiv0$ (Theorem~\CGref{thm:closing}).
\end{proof}

Define the \emph{phase defect} of the family, when the limit exists,
\begin{equation}\label{eq:phase-defect}
 \mathsf v_\lambda(t):=\lim_{\sigma\to\frac12^+}\Im\tilde h_\lambda(\sigma+it)=\lim_{\sigma\to\frac12^+}\Bigl(\Im\log\tilde\zeta_{\Z,\lambda}(\sigma+it)-\Im\log\zeta_\lambda(\sigma+it)\Bigr).
\end{equation}

\begin{proposition}[The phase of the continuous side]\label{prop:continuous-phase}
Let $\lambda\in[0,1]$ and $t>0$ not an ordinate $\gamma_k(\lambda)$.  With the branch of $\log\zeta_\lambda$ that is real on $(1,\infty)$, continued along any path in the open quadrant $\{\Re s>\frac12,\ \Im s>0\}$ (the continuation is well defined there, since $\zeta_\lambda$ has no zeros or poles in the open quadrant: its zeros lie on the line $\Re s=\frac12$ and its pole $s=1$ on the real axis, both on the boundary of the open quadrant; the limit below is taken at points $\frac12+it$ that are not zeros),
\begin{equation}\label{eq:continuous-phase}
 \lim_{\sigma\to\frac12^+}\Im\log\zeta_\lambda(\sigma+it)=-\vartheta(t)+\pi\bigl(N_\lambda(t)-1\bigr),\qquad N_\lambda(t)=\#\{k:\gamma_k(\lambda)\le t\},
\end{equation}
and in the sense of distributions in $t>0$
\begin{equation}\label{eq:generator-phase}
 \lim_{\sigma\to\frac12^+}\Im\mathcal V_\lambda(\sigma+it)=-\pi\sum_{k\ge1}\eta_k\,\delta\bigl(t-\gamma_k(\lambda)\bigr).
\end{equation}
Consequently $\mathsf v_\lambda(t)=\lim_{\sigma\to\frac12^+}\Im\log\tilde\zeta_{\Z,\lambda}(\sigma+it)+\vartheta(t)-\pi\bigl(N_\lambda(t)-1\bigr)$.
\end{proposition}

\begin{proof}
$\zeta_\lambda=P_\lambda/G$.  On $s=\frac12+it$, $q(s)=-(\frac14+t^2)$ and $1+q(s)/A_{k,\lambda}=(\gamma_k(\lambda)^2-t^2)/A_{k,\lambda}$ are real, so $P_\lambda(\frac12+it)$ is real and changes sign at each $\gamma_k(\lambda)$; approaching from $\sigma>\frac12$ the argument of the factor $k$ increases by $\pi$ as $t$ passes $\gamma_k(\lambda)$, since $\Im\bigl(\gamma_k(\lambda)^2+(\sigma-\frac12+it)^2\bigr)=2t(\sigma-\frac12)>0$.  Modulo $2\pi$, the argument of $G(\frac12+it)=-(\frac14+t^2)\Gamma(\frac14+\frac{it}2)\pi^{-1/4-it/2}$ is $\pi+\vartheta(t)$ by \CGeqref{eq:theta-def}.  The multiple of $2\pi$ is fixed by continuity from $t=0^+$: for $\frac12<\sigma<1$, $P_\lambda(\sigma)>0$, because $q(\sigma)\ge-\frac14$ makes every factor $\ge1-\frac14/A_{k,\lambda}>0$, and $G(\sigma)=\sigma(\sigma-1)\pi^{-\sigma/2}\Gamma(\frac\sigma2)<0$, so $\zeta_\lambda(\sigma)<0$; the branch continued from $(1,\infty)$ through the upper half-plane around the simple pole $s=1$ has $\Im\log\zeta_\lambda(\sigma)=-\pi$ there.  This gives \CGeqref{eq:continuous-phase}.  For \CGeqref{eq:generator-phase} it is preferable not to interchange a boundary limit and $\partial_\lambda$.  Put $z=\sigma'+it$ and $\gamma=\gamma_k(\lambda)$.  The $k$th term of \CGeqref{eq:generator-zeta} is
\[
 2\gamma\eta_k\left(\frac1{\gamma^2+z^2}-\frac1{\frac14+\gamma^2}\right).
\]
The second summand is real, while
\[
 \frac{2\gamma}{\gamma^2+z^2}=-i\left(\frac1{z-i\gamma}-\frac1{z+i\gamma}\right).
\]
The Poisson-kernel boundary limit on $t>0$ is therefore $-\pi\eta_k\delta(t-\gamma)$; the pole at $-\gamma$ lies outside the positive half-line.  Summing against a compactly supported test function is justified by the bounded displacement and the same mean-value/summation-by-parts estimates used for the generator (equivalently by Cauchy--Schwarz with $\sum_k\eta_k^2\gamma_k^{-2}<\infty$ after one integration by parts).  This proves \CGeqref{eq:generator-phase} directly in distributions.
\end{proof}

By Proposition~\CGref{prop:continuous-phase}, $\mathsf v_\lambda\equiv0$ asks two things of the discrete function on the critical line: jumps of $\pi$ of its argument exactly at the ordinates $\gamma_k(\lambda)$, and the phase $-\vartheta(t)$ between them; equivalently, $e^{i\vartheta(t)}\tilde\zeta_{\Z,\lambda}(\frac12+it)$ is real.  At the clock end $\tilde h_1$ is holomorphic in $\Re s>0$ (Theorem~\CGref{thm:atomized-continuation}(iv)), so $\mathsf v_1(t)=\Im\tilde h_1(\frac12+it)$ is a real-analytic function; it does not vanish identically, by Corollary~\CGref{cor:bridge-family}(i) and the bound $|\tilde h_1(s)|\le C(1+|s|)$ on $\Re s\ge\frac14$ of the proof of Proposition~\CGref{prop:clock-asymmetry}.

\begin{proposition}[The cell masses are sums of a one-ordinate function]\label{prop:one-ordinate}
For $n\ge2$ there is a real function $\mathsf o_n(\kappa,\kappa')$ of two real variables, depending only on $n$ and the cell $\varphi_n$, such that for all lists $\boldsymbol{\kappa},\boldsymbol{\kappa}'\in\mathscr O$
\begin{equation}\label{eq:one-ordinate}
 \tilde c_n(\boldsymbol{\kappa})-\tilde c_n(\boldsymbol{\kappa}')=\sum_{k\ge1}\mathsf o_n(\kappa_k,\kappa'_k),\qquad \mathsf o_n(\kappa,\kappa')=-2\,\Re\bigl(W_n(\kappa)-W_n(\kappa')\bigr)\ \text{ if }n\text{ is not a power of }2,
\end{equation}
and, with $\mathsf L_n(s):=-\log(1-n^{-s})$, for $\Re s>1$
\begin{equation}\label{eq:one-ordinate-log}
 \log\tilde\zeta_{\Z,\boldsymbol{\kappa}}(s)-\log\tilde\zeta_{\Z,\boldsymbol{\kappa}'}(s)=\sum_{n\ge2}\Bigl(\sum_{k\ge1}\mathsf o_n(\kappa_k,\kappa'_k)\Bigr)\mathsf L_n(s),
\end{equation}
\begin{equation}\label{eq:one-ordinate-logzeta}
 \log\zeta_{\boldsymbol{\kappa}}(s)-\log\zeta_{\boldsymbol{\kappa}'}(s)=\sum_{k\ge1}\log\frac{1+q(s)/(\frac14+\kappa_k^2)}{1+q(s)/(\frac14+{\kappa'_k}^{2})} .
\end{equation}
\end{proposition}

\begin{proof}
$\tilde c_n(\boldsymbol{\kappa})=\langle d\pi_{\boldsymbol{\kappa}},\varphi_n\rangle$, where $d\pi_{\boldsymbol{\kappa}}$ is obtained from $J_{\boldsymbol{\kappa}}$ by truncation at $2$ and finite M\"obius inversion, both linear, and by \CGeqref{eq:J-X} the difference $J_{\boldsymbol{\kappa}}-J_{\boldsymbol{\kappa}'}=-2\Re\sum_k\bigl(\Ei((\frac12+i\kappa_k)\log x)-\Ei((\frac12+i\kappa'_k)\log x)\bigr)$ is a sum over $k$ converging in $\mathcal D'(1,\infty)$.  This gives \CGeqref{eq:one-ordinate}, with the value \CGeqref{eq:cell-mass-W} off the powers of $2$.  \CGeqref{eq:one-ordinate-log} is the Euler product with Corollary~\CGref{cor:abscissa-one}, and \CGeqref{eq:one-ordinate-logzeta} is \CGeqref{eq:list-objects}.
\end{proof}

By Proposition~\CGref{prop:one-ordinate}, the difference of two joining factors is, term by term in $n$ and in $k$, a sum over the ordinates of one universal function of one ordinate, the joining factor of a single zero factor:
\begin{equation}\label{eq:one-ordinate-joining}
 \mathfrak j(\kappa,\kappa';s):=\sum_{n\ge2}\mathsf o_n(\kappa,\kappa')\,\mathsf L_n(s)-\log\frac{1+q(s)/(\frac14+\kappa^2)}{1+q(s)/(\frac14+{\kappa'}^{2})},\qquad
 \tilde h_{\boldsymbol{\kappa}}-\tilde h_{\boldsymbol{\kappa}'}=\sum_{k\ge1}\mathfrak j(\kappa_k,\kappa'_k;s)\ \ \text{formally.}
\end{equation}
With $\boldsymbol{\kappa}'=(\tau_k)$ this writes the phase defect as $\mathsf v_{\boldsymbol{\kappa}}=\mathsf v_1+\sum_k\Im\mathfrak j(\kappa_k,\tau_k;\frac12+it)$, the defect of the clock plus a transport term.

\begin{heuristic}[What happens to the phase defect along the path]\label{heur:phase-rotation}
The following is a computation by stationary phase without remainder terms; it is not proved.  For $n$ much larger than $\kappa$ the cell resolves the oscillation, $\Re W_n(\kappa)\approx n^{-1/2}\cos(\kappa\log n)/\log n$, and for $n$ much smaller than $\kappa/2\pi$ the kernel $W_n(\kappa)$ is negligible because $\varphi$ is smooth (Proposition~\CGref{prop:height-kernels}).  The transition is at $n\approx\kappa/2\pi$, where the phase is $\kappa\log\frac\kappa{2\pi}-\kappa=2\pi\bigl(\NR(\kappa)-\frac78\bigr)$.  This suggests
\begin{equation}\label{eq:phase-rotation}
 \Im\mathfrak j\bigl(\kappa,\cdot\,;\tfrac12+it\bigr)\approx\text{amplitude}\times\cos\bigl(2\pi\NR(\kappa)+\text{phase}(t)\bigr),\qquad\bigl|\mathfrak j(\kappa,\cdot\,;\sigma+it)\bigr|\approx\kappa^{-1/2-\sigma},
\end{equation}
which suggests absolute convergence of the sum over the ordinates for $\sigma>\frac12$ and critical behaviour at $\sigma=\frac12$, in agreement with Theorem~\CGref{thm:atomized-continuation}(iii), and is at the edge of convergence on the critical line.  On the clock, $\NR(\tau_k)=k-\frac12$ and the phases are aligned, which is why $\tilde h_1$ crosses the line.  Along the path each term turns by the angle $(1-\lambda)\eta_k\ell_k$, the rotation of \CGeqref{eq:Xi-rotation}.  At $\lambda=0$, \textup{(B1)} is the statement that the rotated sum, with the phases $e^{2\pi i\NR(\gamma_k)}$, converges on the critical line to minus the defect $\mathsf v_1$ of the clock.
\end{heuristic}

\subsection{Transport from the clock on the rungs}
\label{subsec:transport}

For a real function $\mathsf f\in C_c^\infty(0,\infty)$ put
\begin{equation}\label{eq:transport-defs}
 \mathsf c_{\mathsf f}(t):=\int_0^\infty\mathsf f(u)\cos(tu)\,du,\qquad \mathsf k(u):=\frac12\Bigl(e^{u/2}-\frac{e^{-u/2}}{e^{2u}-1}\Bigr),\qquad \mathsf k_{\mathsf f}:=\int_0^\infty\mathsf f(u)\,\mathsf k(u)\,du .
\end{equation}

\begin{proposition}[The rung functional]\label{prop:rung-functional}
For every list $\boldsymbol{\kappa}\in\mathscr O$ (Section~\CGref{sec:lists}),
\begin{equation}\label{eq:rung-functional}
 \mathsf R_{\mathsf f}(\boldsymbol{\kappa}):=\frac12\int_1^\infty\mathsf f(\log x)\,x^{-1/2}\log x\,dJ^{\rm M}_{\boldsymbol{\kappa}}(x)=\mathsf k_{\mathsf f}-\sum_{k\ge1}\mathsf c_{\mathsf f}(\kappa_k),
\end{equation}
the series converging absolutely.  Along the family, with $\mathsf R_{\mathsf f}(\lambda):=\mathsf R_{\mathsf f}\bigl((\gamma_k(\lambda))\bigr)$,
\begin{equation}\label{eq:transport}
 \mathsf R_{\mathsf f}(\lambda)=\mathsf R_{\mathsf f}(1)-\mathsf S_{\mathsf f}(\lambda),\qquad \mathsf S_{\mathsf f}(\lambda):=\sum_{k\ge1}\Bigl(\mathsf c_{\mathsf f}\bigl(\gamma_k(\lambda)\bigr)-\mathsf c_{\mathsf f}(\tau_k)\Bigr),\qquad\mathsf S_{\mathsf f}(1)=0 .
\end{equation}
If the support of $\mathsf f$ contains no point $\log n$, $n\in\N$, then
\begin{equation}\label{eq:transport-D}
 (\mathrm D)\ \Longrightarrow\ \mathsf S_{\mathsf f}(0)=\mathsf R_{\mathsf f}(1),
\end{equation}
and (D) holds if and only if $\mathsf S_{\mathsf f}(0)=\mathsf R_{\mathsf f}(1)$ for every such $\mathsf f$ and the resulting integer jumps satisfy the growth hypothesis of Theorem~\CGref{thm:cell-form}.
\end{proposition}

\begin{proof}
Pair \CGeqref{eq:distributional-EF}, which holds for every list of $\mathscr O$ with the same proof, with $\frac12\mathsf f(\log x)x^{-1/2}$ and put $u=\log x$: the term $1$ gives $\frac12\int\mathsf f(u)e^{u/2}du$, the term $x^{-1}(x^2-1)^{-1}$ gives $\frac12\int\mathsf f(u)e^{-u/2}(e^{2u}-1)^{-1}du$, and the sum gives $\sum_k\mathsf c_{\mathsf f}(\kappa_k)$; $\mathsf c_{\mathsf f}$ is rapidly decreasing.  \CGeqref{eq:transport} is the difference of two cases of \CGeqref{eq:rung-functional}.  By Theorem~\CGref{thm:cell-form}, the vanishing $\mathsf R_{\mathsf f}(0)=0$ for every $\mathsf f$ whose support avoids the points $\log n$ is exactly the constancy part of the cell criterion; (D) follows when the resulting integer jumps also satisfy $\sum_{n\le X}|\mathsf b_n|=O_\eps(X^{1+\eps})$.
\end{proof}

The number $\mathsf R_{\mathsf f}(1)$ belongs to the clock; $\mathsf S_{\mathsf f}(0)$ is the transport from the clock to the ordinates of $\zeta$.  (D) is the statement that the transport cancels the defect of the clock on every rung.

\begin{theorem}[The transport is band-limited in $\lambda$]\label{thm:transport-band}
Let $\mathsf f\in C_c^\infty(0,\log N)$.  The function $\lambda\mapsto\mathsf S_{\mathsf f}(\lambda)$ extends to an entire function with
\begin{equation}\label{eq:transport-band}
 |\mathsf S_{\mathsf f}(\lambda)|\le C_{\mathsf f}\bigl(1+|\lambda|\log(2+|\lambda|)\bigr)\,e^{C_\eta\log N\,|\Im\lambda|}\qquad(\lambda\in\C),
\end{equation}
and for every $j\ge1$
\begin{equation}\label{eq:transport-moments}
 \partial_\lambda^j\mathsf S_{\mathsf f}(1)=\mathsf m_j(\mathsf f):=\sum_{k\ge1}\eta_k^{\,j}\,\mathsf c_{\mathsf f}^{(j)}(\tau_k),\qquad
 \mathsf S_{\mathsf f}(0)=\sum_{j\ge1}\frac{(-1)^j}{j!}\,\mathsf m_j(\mathsf f).
\end{equation}
Moreover $|\mathsf m_j(\mathsf f)|\le C_{\mathsf f}(1+j)^2(C_\eta\log N)^j$ (the support of $\mathsf f$ lies in some $[u_1,\log N]$ with $u_1>0$, and $C_{\mathsf f}$ depends on it).
\end{theorem}

\begin{proof}
By the Paley--Wiener theorem, $|\mathsf c_{\mathsf f}^{(j)}(t)|\le C_{A,j}\,e^{\log N|\Im t|}(1+|t|)^{-A}$ for complex $t$ and every $A$.  With $\gamma_k(\lambda)=\gamma_k+\lambda\eta_k$, $|\Im\gamma_k(\lambda)|\le C_\eta|\Im\lambda|$ and $|\gamma_k(\lambda)|\ge\gamma_k-C_\eta|\lambda|$.  The ordinates with $\gamma_k\le2C_\eta|\lambda|$ are $O(|\lambda|\log(2+|\lambda|))$ in number and each contributes at most $C\,e^{C_\eta\log N|\Im\lambda|}$; for the others $|\mathsf c_{\mathsf f}(\gamma_k(\lambda))|\le C\,e^{C_\eta\log N|\Im\lambda|}(\gamma_k/2)^{-2}$.  This gives \CGeqref{eq:transport-band}, the termwise differentiation and the convergence of the Taylor series at $\lambda=1$ evaluated at $\lambda=0$.  Finally $\mathsf c^{(j)}_{\mathsf f}(t)=\int\mathsf f(u)u^j\cos(tu+\frac{j\pi}2)du$ and two integrations by parts give $(1+t^2)|\mathsf c_{\mathsf f}^{(j)}(t)|\le C_{\mathsf f}(1+j)^2(\log N)^j$.
\end{proof}

The scaled lists $(\tau_k-c'\eta_k)$ are the members $\lambda=1-c'$ of the family, so the scale test of Remark~\CGref{rem:scale-test} is the behaviour of the entire function $\mathsf S_{\mathsf f}$ on the real axis.  In \CGeqref{eq:transport-moments} the nodes $\tau_k$ and the weights $\mathsf c^{(j)}_{\mathsf f}(\tau_k)$ belong to the clock, and the powers $\eta_k^{\,j}$ do not; on the first rung, $N=2$, the term $j$ is at most $C_{\mathsf f}(1+j)^2(C_\eta\log2)^j/j!$.

\begin{proposition}[The defect of the clock on a rung]\label{prop:clock-defect-poisson}
Let $\mathfrak t_0\in(2\pi,\tau_0)$ be the point with $\NR(\mathfrak t_0)=0$, and let $\mathsf f\in C_c^\infty(0,\infty)$.  Then
\begin{equation}\label{eq:clock-defect-poisson}
 \sum_{k\ge1}\mathsf c_{\mathsf f}(\tau_k)=\mathsf I_0+2\sum_{m\ge1}(-1)^m\mathsf I_m,\qquad \mathsf I_m:=\int_{\mathfrak t_0}^\infty\mathsf c_{\mathsf f}(t)\cos\bigl(2\pi m\NR(t)\bigr)\NR'(t)\,dt=O_{\mathsf f}\bigl(m^{-2}\bigr),
\end{equation}
\begin{equation}\label{eq:clock-defect-I0}
 \mathsf I_0=-\int_0^\infty\frac{\mathsf f(u)}{2\pi u}\Bigl(\sin(u\mathfrak t_0)\log\frac{\mathfrak t_0}{2\pi}+\frac\pi2-\int_0^{u\mathfrak t_0}\frac{\sin y}y\,dy\Bigr)du,
\end{equation}
and therefore $\mathsf R_{\mathsf f}(1)=(\mathsf k_{\mathsf f}-\mathsf I_0)-2\sum_{m\ge1}(-1)^m\mathsf I_m$.
\end{proposition}

\begin{proof}
The function $g(v):=\mathsf c_{\mathsf f}(\mathcal T(v))$, $v\ge0$, with $\mathcal T$ the inverse of the clock on $[\mathfrak t_0,\infty)$, is continuous on $[0,\infty)$, $C^\infty$ on $(0,\infty)$ with $g'(0^+)=\mathsf c_{\mathsf f}'(\mathfrak t_0)\mathcal T'(0)$ finite, and rapidly decreasing with its derivatives, and $\tau_k=\mathcal T(k-\frac12)$.  Poisson's summation formula for the even extension of $g$ (continuous, in $L^1$, with a corner at $0$ in general, since $g'(0^+)\ne0$ unless $\mathsf c_{\mathsf f}'(\mathfrak t_0)=0$) at the half-integers gives $2\sum_{k\ge1}g(k-\frac12)=\sum_{m\in\Z}(-1)^m2\int_0^\infty g(v)\cos(2\pi mv)\,dv$, the Fourier coefficients being $O(m^{-2})$ by two integrations by parts, $\int_0^\infty g(v)\cos(2\pi mv)\,dv=-\frac{g'(0^+)}{(2\pi m)^2}-\frac1{(2\pi m)^2}\int_0^\infty g''(v)\cos(2\pi mv)\,dv$, the boundary term coming from the corner (so the decay cannot be improved to $O(m^{-3})$ in general); the substitution $v=\NR(t)$ gives \CGeqref{eq:clock-defect-poisson}.  For $m=0$, $\NR'(t)=\frac1{2\pi}\log\frac t{2\pi}$; integrating by parts with the primitive $\int\mathsf f(u)\sin(tu)u^{-1}du$ of $\mathsf c_{\mathsf f}$ and exchanging the integrals gives \CGeqref{eq:clock-defect-I0}.
\end{proof}

In $\mathsf I_m$ the phases $2\pi m\NR(t)\mp tu$ have the derivative $m\log\frac t{2\pi}\mp u$, which vanishes at $t=2\pi e^{u/m}$: the defect of the clock at the point $x=e^u$ is produced at the heights $2\pi x^{1/m}$, which is the origin of the terms $2\cos(2\pi y)$ and $y^{-1/4}\cos(4\pi\sqrt y+\frac\pi4)$ of \CGeqref{eq:clock-density}.  On the rung $(n,n+1)$ the height for $m=1$ lies in $(2\pi n,2\pi(n+1))$.  On the first rung it lies in $(2\pi,4\pi)$, below $\tau_1$ and on both sides of $\mathfrak t_0$: the defect of the clock on $(1,2)$ is an effect of the start of the clock.  Experiments~\CGref{exp:transport} and~\CGref{exp:clock-defect} evaluate \CGeqref{eq:transport}--\CGeqref{eq:clock-defect-poisson} for one function $\mathsf f$.


\subsection{The multiscale Jacobian of the midpoint clock}
\label{subsec:multiscale-jacobian}

The rung system of Subsection~\CGref{subsec:free-masses} asks for a list and free masses simultaneously.  Before attempting nonlinear continuation from the clock, one must understand the linearization with respect to the ordinates.  The point is genuinely multiscale: the gaps of the midpoint clock tend to zero, so no fixed $u$-window can carry a uniform inverse, while a window of length comparable with $\log T$ is exactly large enough to resolve the ordinates up to height $T$.

For $T>\tau_2$ put
\[
 K(T):=\max\{k:\tau_k\le T\},\qquad
 \Delta_T:=\min_{1\le k<K(T)}(\tau_{k+1}-\tau_k),
\]
and define the finite ordinate block of the linearized rung map by
\begin{equation}\label{eq:clock-jacobian}
 (\mathcal J_T\mathbf v)(u):=-u\sum_{\tau_k\le T}v_k\sin(\tau_k u).
\end{equation}
No genuine zero ordinate and no individual displacement $\eta_k=\tau_k-\gamma_k$ enters this definition.

\begin{theorem}[Multiscale Jacobian of the midpoint clock]\label{thm:clock-jacobian}
The following statements hold.
\begin{enumerate}[label=\textup{(\roman*)}]
\item For every $T>\tau_2$,
\begin{equation}\label{eq:clock-gap-lower}
 \Delta_T\ge \frac{2\pi}{\log(T/2\pi)}.
\end{equation}
\item Fix $a>0$ and $\varkappa_{\rm win}>0$, and put
\[
 L_T:=(1+\varkappa_{\rm win})\log(T/2\pi),\qquad I_T=[a,a+L_T].
\]
There is $c_{a,\varkappa_{\rm win}}>0$, independent of $T$, such that for all finitely supported $\mathbf v=(v_k)$ on $\{k:\tau_k\le T\}$,
\begin{equation}\label{eq:multiscale-frame}
 \|\mathcal J_T\mathbf v\|_{L^2(I_T)}^2
 \ge c_{a,\varkappa_{\rm win}}L_T\sum_{\tau_k\le T}|v_k|^2.
\end{equation}
\item No analogue of \CGeqref{eq:multiscale-frame} with a positive $T$-independent lower constant can hold on a fixed interval $(0,U)$.  Indeed, for the unit vector $\mathbf v^{(k)}=(e_k-e_{k+1})/\sqrt2$,
\begin{equation}\label{eq:fixed-window-pair}
 \|\mathcal J\mathbf v^{(k)}\|_{L^2(0,U)}
 \le \frac{U^{5/2}}{\sqrt{10}}(\tau_{k+1}-\tau_k)\longrightarrow0.
\end{equation}
\item For every finite $U>0$,
\[
 \overline{\operatorname{span}}\{u\sin(\tau_k u):k\ge1\}=L^2(0,U).
\]
Thus the infinite linearized clock has dense range on every finite window, despite the absence of a bounded inverse on a fixed window.
\item Let $\mathsf f\in C_c^\infty(0,U)$ and $\mathsf c_{\mathsf f}(t)=\int_0^\infty\mathsf f(u)\cos(tu)\,du$.  For every list $\boldsymbol\kappa$ satisfying $|\kappa_k-\tau_k|\le C_\eta$ and every $A>1$,
\begin{equation}\label{eq:smooth-tail-clock}
 \sum_{\tau_k>T}\left|\mathsf c_{\mathsf f}(\kappa_k)-\mathsf c_{\mathsf f}(\tau_k)\right|
 \ll_{A,\mathsf f,C_\eta} T^{1-A}\log T.
\end{equation}
\end{enumerate}
\end{theorem}

\begin{proof}
The midpoint equation is $\NR(\tau_k)=k-\frac12$ with
\[
 \NR'(t)=\frac1{2\pi}\log\frac t{2\pi}.
\]
For $\tau_{k+1}\le T$, the mean-value theorem gives
\[
 1=\NR'(\xi_k)(\tau_{k+1}-\tau_k),\qquad \tau_k<\xi_k<\tau_{k+1}\le T,
\]
which proves \CGeqref{eq:clock-gap-lower}.

For (ii), apply Ingham's inequality for non-harmonic exponentials \cite{Ingham1936} to the separated frequency set $\{\pm\tau_k:\tau_k\le T\}$.  By \CGeqref{eq:clock-gap-lower},
\[
 L_T\Delta_T\ge 2\pi(1+\varkappa_{\rm win}),
\]
so after scaling the frequencies by $\Delta_T$ the interval lies a fixed positive distance beyond Ingham's critical length.  The lower Ingham constant is therefore bounded below by a positive constant depending only on $\varkappa_{\rm win}$; rescaling gives a bound of size $L_T$ for the sine sum.  Since $u\ge a$ on $I_T$, multiplication by $u$ gives \CGeqref{eq:multiscale-frame}.

For (iii), $|\sin(\tau_{k+1}u)-\sin(\tau_k u)|\le(\tau_{k+1}-\tau_k)u$, whence
\[
 \|\mathcal J\mathbf v^{(k)}\|_{L^2(0,U)}^2
 \le\frac{(\tau_{k+1}-\tau_k)^2}{2}\int_0^Uu^4\,du
 =\frac{(\tau_{k+1}-\tau_k)^2U^5}{10}.
\]
The same mean-value formula used above shows $\tau_{k+1}-\tau_k\to0$.

For (iv), suppose $F\in L^2(0,U)$ is orthogonal to every $u\sin(\tau_k u)$.  Put $H(u)=uF(u)$ and extend $H$ oddly to $(-U,U)$.  Its Fourier transform
\[
 \widehat H(z)=\int_{-U}^{U}H(u)e^{-izu}\,du
\]
is entire of exponential type at most $U$ and vanishes at every $\pm\tau_k$.  A nonzero entire function of finite exponential type has $O(r)$ zeros in $|z|\le r$ by Jensen's formula, whereas the midpoint clock has $2\NR(r)+O(1)\asymp r\log r$ points $\pm\tau_k$ in that disk.  Hence $\widehat H\equiv0$, so $H=0$ and therefore $F=0$ almost everywhere.

For (v), the Paley--Wiener decay of the cosine transform gives $|\mathsf c'_{\mathsf f}(t)|\ll_{A,\mathsf f}(1+|t|)^{-A}$ for every $A$.  The mean-value theorem and $|\kappa_k-\tau_k|\le C_\eta$ give
\[
 |\mathsf c_{\mathsf f}(\kappa_k)-\mathsf c_{\mathsf f}(\tau_k)|\ll_{A,\mathsf f,C_\eta}\tau_k^{-A}
\]
for large $k$.  Partial summation with $\NR(t)=O(t\log t)$ yields $\sum_{\tau_k>T}\tau_k^{-A}=O(T^{1-A}\log T)$.
\end{proof}

The theorem identifies the correct scale but does not prove existence of a solution of the rung system.  In particular, a fixed-window Banach-space implicit-function argument cannot work because of \CGeqref{eq:fixed-window-pair}.  A quantitative continuation must grow the resolved window together with the spectral height.

One precise residual formulation is the following.  Fix $r>\frac12$, so point masses belong to $H^{-r}(0,U)$, and put
\[
 \mathsf D_{\RVM}(u):=\mathsf k(u)-\sum_{k\ge1}\cos(\tau_k u)
\]
in $\mathcal D'(0,\infty)$.  For $T,U>0$ define
\begin{equation}\label{eq:multiscale-residual}
 \varepsilon_r(T,U):=
 \inf_{h,b}\left\|
 \mathsf D_{\RVM}+u\sum_{\tau_k\le T}h_k\sin(\tau_k u)
 -\sum_{\log n\le U}b_n\delta(u-\log n)
 \right\|_{H^{-r}(0,U)}.
\end{equation}
Theorem~\CGref{thm:clock-jacobian} supplies the finite-height separation, dense-range and tail ingredients.  What is not proved is a bound for \CGeqref{eq:multiscale-residual}, together with control of the minimizing $h$ strong enough to dominate the quadratic remainder in
\[
 \cos((\tau_k+h_k)u)=\cos(\tau_ku)-h_ku\sin(\tau_ku)+O(h_k^2u^2),
\]
uniformly along a nested sequence with $U\asymp\log T$.  Such an estimate would turn the finite-section continuation observed numerically into an existence proof for the infinite rung system.  By Theorem~\CGref{thm:free-masses}, any resulting solution would then be the zeta ordinates with the von Mangoldt masses.  The individual displacements $\eta_k$ are not data in this formulation.

\section{The space of ordinate lists and the existence form of the final step}
\label{sec:lists}

The constructions of Sections~\CGref{sec:lambda-family}--\CGref{sec:criteria} use two properties of the ordinates (Remark~\CGref{rem:input-A7}).  This section applies them to every list with these two properties.  The final step then becomes an existence statement in a compact space, and the list of $\zeta$ is not part of its formulation.

\subsection{Lists, their zeta functions and their atomizations}
\label{subsec:lists}

Let $\mathscr O$ be the set of sequences $\boldsymbol{\kappa}=(\kappa_k)_{k\ge1}$ of real numbers with
\begin{equation}\label{eq:list-space}
 \tau_0^{+}\le\kappa_1\le\kappa_2\le\cdots,\qquad|\kappa_k-\tau_k|\le C_\eta,\qquad\Bigl|\int_{\tau_0}^{X}\bigl(N_{\boldsymbol{\kappa}}(\tau)-\NR(\tau)\bigr)d\tau\Bigr|\le C\log X\quad(X\ge\tau_0),
\end{equation}
where $\tau_0^{+}$ is a fixed number with $\tau_0<\tau_0^{+}<\min(\gamma_1,\tau_1)$, $N_{\boldsymbol{\kappa}}(\tau)=\#\{k:\kappa_k\le\tau\}$, $C_\eta$ is the constant of Lemma~\CGref{lem:bounded-displacement} and $C$ the constant of Proposition~\CGref{prop:family-clock}.  The clock $(\tau_k)$, the ordinates $(\gamma_k)$ of $\zeta$ and every list $(\gamma_k(\lambda))$, $0\le\lambda\le1$, belong to $\mathscr O$.  For $\boldsymbol{\kappa}\in\mathscr O$ (Section~\CGref{sec:lists}) define, exactly as for the family,
\begin{equation}\label{eq:list-objects}
 \zeta_{\boldsymbol{\kappa}}(s):=\frac1{G(s)}\prod_{k\ge1}\Bigl(1+\frac{q(s)}{\frac14+\kappa_k^2}\Bigr),\qquad
 \tilde c_n(\boldsymbol{\kappa}),\quad \tilde a_N(\boldsymbol{\kappa}),\quad k_N(\boldsymbol{\kappa}),\qquad \tilde\zeta_{\Z,\boldsymbol{\kappa}}=\zeta_{\boldsymbol{\kappa}}\,e^{\tilde h_{\boldsymbol{\kappa}}},\qquad \mathsf V_\sigma(\boldsymbol{\kappa}),
\end{equation}
with the cell masses of Definition~\CGref{def:atomized}, the joining factor of Theorem~\CGref{thm:joining-factor} and the function \CGeqref{eq:vertical-defect}.  The results of the preceding sections that are used below for lists are transferred to $\mathscr O$ by the following theorem, which records, result by result, which part of \CGeqref{eq:list-space} each proof uses and which constants are uniform.

\begin{theorem}[Transfer to $\mathscr O$]\label{thm:transfer-O}
Let $\boldsymbol{\kappa}\in\mathscr O$. The following statements, proved above for the lists $(\gamma_k(\lambda))_{k\ge1}$, hold for $\boldsymbol{\kappa}$ with the same proofs, and their constants depend only on $\tau_0^+$, $C_\eta$ and $C$ in \CGeqref{eq:list-space} and on the parameters of the statement ($\sigma$, $\eps$, $n$, $N$), never on the list: $C=C(\tau_0^+,C_\eta,C;\sigma,\eps,\dots)$, hence they are uniform on $\mathscr O$. This uniformity is what Proposition~\CGref{prop:lists-compact} and the compactness arguments of Section~\CGref{sec:routeB-finite} use.
\begin{enumerate}[label=\textup{(\roman*)}]
\item Theorem~\CGref{thm:path}: $\zeta_{\boldsymbol{\kappa}}$ has order one, satisfies the functional equation \CGeqref{eq:Z1-FE}, and its nontrivial zeros are exactly $\frac12\pm i\kappa_k$, with multiplicity. Input: the tube (convergence of the product and $\kappa_k\asymp\tau_k$), the ordering, and $\kappa_1\ge\tau_0^+$.
\item Theorem~\CGref{thm:Z1-clock} in its list form, with $\boldsymbol{\kappa}$ in place of $\boldsymbol\gamma$: the clock decomposition of $\zeta_{\boldsymbol{\kappa}}$, the value $\varpi_{\boldsymbol{\kappa}}$ at $+\infty$ as an ordered product, and the bounds of Proposition~\CGref{prop:family-clock}. Input: the tube and the third condition of \CGeqref{eq:list-space}, which replaces $\int_{\tau_0}^X(N_\zeta-n_{\RVM})=O(\log X)$ in every summation by parts against the count difference.
\item Theorems~\CGref{thm:mellin-counts} and~\CGref{thm:distributional-EF}: the prime-power count $J^{\rm M}_{\boldsymbol{\kappa}}$ as a Mellin transform, its $L^2_{\rm loc}$ bound, and the distributional explicit formula. Input: (i), (ii) and the Perron argument, in which the ordinates enter only through $\log\zeta_{\boldsymbol{\kappa}}$ on vertical lines, controlled by (ii).
\item Proposition~\CGref{prop:one-ordinate}: the cell masses $\tilde c_n(\boldsymbol{\kappa})$ as sums of a one-ordinate function. Input: (iii) and dominated convergence, with the domination supplied by the tube.
\item Lemma~\CGref{lem:mean-value-ordinates} with $\kappa_k$ in place of $\gamma_k(\lambda)$: the tube gives $N_{\boldsymbol{\kappa}}(t+1)-N_{\boldsymbol{\kappa}}(t)\le n_{\RVM}(t+1+C_\eta)-n_{\RVM}(t-C_\eta)\ll\log(t+2)$, with an implied constant depending on $C_\eta$ only.
\item Theorem~\CGref{thm:mean-square-cells} and Corollary~\CGref{cor:abscissa-one}: $\sum_{X<n\le2X}|\tilde c_n(\boldsymbol{\kappa})|^2\le C_\eps X^{1+\eps}$, the absolute convergence of $\tilde\zeta_{\Z,\boldsymbol{\kappa}}^{\pm1}$ in $\Re s>1$, and the exceptional-set bound. Input: (v) for the cells away from the powers of $2$; for the cells $n=2^j$ the term \CGeqref{eq:cell-mass-power-two}, whose bound needs $\|J_{\boldsymbol{\kappa}}\|_{L^2(2,3)}$ bounded uniformly on $\mathscr O$, which is (iii) together with the uniform bound of (ii); and the uniformity in $n$ of the bumps $\varphi_n$. The constant $C_\eps$ depends on $\eps$, $\tau_0^+$, $C_\eta$ and $C$ only. The statement on $\partial_\lambda\tilde c_n$ has no analogue for a single list and is not transferred.
\item Theorem~\CGref{thm:atomized-continuation}(i)--(iii) for $\tilde\zeta_{\Z,\boldsymbol{\kappa}}$. Input: (vi) and (iii).
\end{enumerate}
\end{theorem}

\begin{proof}
Each item names the only places where the cited proofs use the ordinates. They are four: the tube $|\kappa_k-\tau_k|\le C_\eta$, which gives $\kappa_k\asymp\tau_k$, $\sum_k\kappa_k^{-2}<\infty$, the local count of (v), and the domination of (iv); the ordering, used in the monotone rearrangements and in the identification of the zeros of $\zeta_{\boldsymbol{\kappa}}$ with the ordinates; the threshold $\kappa_1\ge\tau_0^+>\tau_0$, used wherever a branch of $\log\bigl(1+q(s)/(\frac14+\kappa^2)\bigr)$ is fixed on $|s-\frac12|<\tau_0$; and the third condition, used in every summation by parts over the ordered ordinates in place of Proposition~\CGref{prop:family-clock}. None of the cited proofs uses that the list is a zero set of $\zeta$ or a convex combination with the clock, except in the statements about $\lambda$-derivatives, which are excluded in (vi). Reading the proofs with these four inputs in place of their sources gives the items, and the constants are those of the inputs.
\end{proof}

\begin{theorem}[Rigidity of the coincidence]\label{thm:rigidity}
The set
\begin{equation}\label{eq:coincidence-set}
 \mathscr C:=\bigl\{\boldsymbol{\kappa}\in\mathscr O:\ \tilde h_{\boldsymbol{\kappa}}\ \text{is constant}\bigr\}
\end{equation}
of the lists whose continuous and atomized zeta functions coincide up to a constant factor has at most one element.  If $\boldsymbol{\kappa}\in\mathscr C$, then $\zeta_{\boldsymbol{\kappa}}=\tilde\zeta_{\Z,\boldsymbol{\kappa}}=\zeta$, $\boldsymbol{\kappa}$ is the list of the ordinates of $\zeta$, and (D) holds.  Conversely, (D) implies $\mathscr C=\{(\gamma_k)\}$.
\end{theorem}

\begin{proof}
If $\tilde h_{\boldsymbol{\kappa}}\equiv c$, then $\zeta_{\boldsymbol{\kappa}}=e^{-c}\tilde\zeta_{\Z,\boldsymbol{\kappa}}$ is an absolutely convergent Dirichlet series in $\Re s>1$ (Corollary~\CGref{cor:abscissa-one}), $(s-1)\zeta_{\boldsymbol{\kappa}}$ is entire of order one and $\zeta_{\boldsymbol{\kappa}}$ satisfies \CGeqref{eq:Z1-FE} (Theorem~\CGref{thm:path}).  Hamburger's theorem gives $\zeta_{\boldsymbol{\kappa}}=C_H\zeta$, and $C_H=1$ because both functions have residue $1$ at $s=1$.  Hence the nontrivial zeros of $\zeta$ with positive ordinate are the points $\frac12+i\kappa_k$, with multiplicity; this determines $\boldsymbol{\kappa}$ as the list $(\gamma_k)$, and $\zeta_{Z1}=\zeta_{\boldsymbol{\kappa}}=\zeta$.  Then $e^{-c}\tilde\zeta_{\Z,\boldsymbol{\kappa}}=\zeta$ and the first coefficients give $e^{-c}=1$; $c$ is real, because $\tilde h_{\boldsymbol\kappa}$ is real on the real axis, so $c=0$.  Conversely, under (D), $\tilde h_0\equiv0$ (Theorem~\CGref{thm:closing}), so $(\gamma_k)\in\mathscr C$.
\end{proof}

\begin{proposition}[Compactness and continuity]\label{prop:lists-compact}
The space $\mathscr O$ and the functions \CGeqref{eq:list-objects} have the following properties.
\begin{enumerate}[label=\textup{(\roman*)}]
\item $\mathscr O$ is compact in the product topology of $\prod_k[\tau_k-C_\eta,\tau_k+C_\eta]$.
\item For every $n\ge2$, $N\ge1$ and $\sigma>1$, the functions $\boldsymbol{\kappa}\mapsto\tilde c_n(\boldsymbol{\kappa})$, $\tilde a_N(\boldsymbol{\kappa})$, $k_N(\boldsymbol{\kappa})$ and $\mathsf V_\sigma(\boldsymbol{\kappa})$ are continuous on $\mathscr O$.
\item For every real $\sigma>1$, $\boldsymbol{\kappa}\mapsto\log\zeta_{\boldsymbol{\kappa}}(\sigma)=-\log G(\sigma)+\sum_k\log\frac{(\sigma-\frac12)^2+\kappa_k^2}{\frac14+\kappa_k^2}$ is a sum of strictly convex functions of the single coordinates $\kappa_k$.
\end{enumerate}
\end{proposition}

\begin{proof}
(i) The product $\mathcal K:=\prod_k[\tau_k-C_\eta,\tau_k+C_\eta]$ is compact by Tychonoff's theorem and metrizable, being a countable product of compact intervals. Each condition in \CGeqref{eq:list-space} defines a closed subset of $\mathcal K$: $\{\kappa_1\ge\tau_0^+\}$ is closed; the ordering is $\bigcap_k\{\kappa_k\le\kappa_{k+1}\}$, an intersection of closed sets; and for fixed $X$ the function $F_X(\boldsymbol{\kappa}):=\sum_k(X-\kappa_k)^+-\int_{\tau_0}^X\NR$ equals $\int_{\tau_0}^X(N_{\boldsymbol{\kappa}}-\NR)$, depends continuously on the finitely many coordinates with $\tau_k\le X+C_\eta$, and the third condition is $\bigcap_{X\ge\tau_0}\{|F_X|\le C\log X\}$, again an intersection of closed sets. So $\mathscr O$ is closed in $\mathcal K$, hence compact.
(ii) For $n$ not a power of $2$, $\tilde c_n(\boldsymbol{\kappa})=\mathfrak m_n-2\Re\sum_kW_n(\kappa_k)$ by \CGeqref{eq:cell-mass-W}, and the series converges uniformly on $\mathscr O$ because $W_n$ is rapidly decreasing and $|\kappa_k-\tau_k|\le C_\eta$.  For $n=2^j$ the term $j$ is $-\frac{\mu(j)}j\int_2^3J_{\boldsymbol{\kappa}}(y)\frac d{dy}\varphi_n(y^j)\,dy$ (proof of Theorem~\CGref{thm:mean-square-cells}), and $\boldsymbol{\kappa}\mapsto J_{\boldsymbol{\kappa}}\in L^2(2,3)$ is continuous, by the following three steps.  (1) Fix $\sigma>1$ and put $\mathsf F_{\boldsymbol{\kappa}}(t):=\log\zeta_{\boldsymbol{\kappa}}(\sigma+it)/(\sigma+it)$.  For fixed $t$, $\log\zeta_{\boldsymbol{\kappa}}(\sigma+it)=-\log G(\sigma+it)+\sum_k\Log\bigl(1+q(\sigma+it)/(\frac14+\kappa_k^2)\bigr)$ is a series converging uniformly on $\mathscr O$ (the $k$-th term is $O(|q|/\tau_k^2)$ on the tube), hence continuous in the product topology; so $\mathsf F_{\boldsymbol{\kappa}}(t)$ is continuous in $\boldsymbol{\kappa}$ for every $t$.  (2) By Theorem~\CGref{thm:transfer-O}(ii) (the bound of Proposition~\CGref{prop:family-clock} on $\mathscr O$), $|\log\zeta_{\boldsymbol{\kappa}}(\sigma+it)-\log\frac{s}{s-1}|\le C_\sigma(1+\log(2+|t|))$ for all $\boldsymbol{\kappa}\in\mathscr O$, so that
\begin{equation}\label{eq:lists-compact-dominant}
 |\mathsf F_{\boldsymbol{\kappa}}(t)|\le\mathsf H(t):=\frac{C'_\sigma\,(1+\log(2+|t|))}{1+|t|},\qquad \mathsf H\in L^2(\R),
\end{equation}
a pointwise majorant independent of $\boldsymbol{\kappa}$; if $\boldsymbol{\kappa}^{(m)}\to\boldsymbol{\kappa}$ in $\mathscr O$, then $|\mathsf F_{\boldsymbol{\kappa}^{(m)}}-\mathsf F_{\boldsymbol{\kappa}}|^2\le4\mathsf H^2$ and dominated convergence gives $\mathsf F_{\boldsymbol{\kappa}^{(m)}}\to\mathsf F_{\boldsymbol{\kappa}}$ in $L^2(\R)$; by Plancherel and Theorem~\CGref{thm:mellin-counts}, $J^{\rm M}_{\boldsymbol{\kappa}^{(m)}}(e^u)e^{-\sigma u}\to J^{\rm M}_{\boldsymbol{\kappa}}(e^u)e^{-\sigma u}$ in $L^2(\R,du)$, hence $J^{\rm M}_{\boldsymbol{\kappa}^{(m)}}\to J^{\rm M}_{\boldsymbol{\kappa}}$ in $L^2(2,3)$.  (3) The constant $C_{\boldsymbol{\kappa}}:=J^{\rm M}_{\boldsymbol{\kappa}}-J_{\boldsymbol{\kappa}}$ equals $\langle J^{\rm M}_{\boldsymbol{\kappa}},\phi\rangle-\langle J_{\boldsymbol{\kappa}},\phi\rangle$ for a fixed $\phi\in C_c^\infty(2,3)$ with $\int\phi=1$: the first pairing is continuous by (2); the second is the sum over $k$ of $\int\bigl(\Ei(\rho_k\log x)+\Ei(\bar\rho_k\log x)\bigr)\phi(x)\,dx$ plus terms independent of the list, a series whose tails are bounded uniformly on $\mathscr O$ by \CGeqref{eq:L2-partial-sums} (whose constant depends on $C_\eta$ only) and whose terms are continuous in $\kappa_k$, hence a continuous function on $\mathscr O$.  Therefore $J_{\boldsymbol{\kappa}}=J^{\rm M}_{\boldsymbol{\kappa}}-C_{\boldsymbol{\kappa}}$ is continuous into $L^2(2,3)$, and the cells $n=2^j$ are continuous.  $\tilde a_N$ and $k_N$ are polynomials in finitely many cell masses.  The sums $\sum_n|k_n(\boldsymbol{\kappa})|n^{-\sigma'}$, $\sigma'>1$, are bounded uniformly on $\mathscr O$ (Theorem~\CGref{thm:mean-square-cells}), so the tails of $\mathsf V_\sigma$ are uniformly small, as in the proof of Theorem~\CGref{thm:double-hamburger}(iii).
(iii) is the computation in the proof of Proposition~\CGref{prop:convexity}: $\partial_\kappa^2\log\frac{a^2+\kappa^2}{1/4+\kappa^2}=\mathcal C(a,\kappa)-\mathcal C(\frac12,\kappa)>0$ for $a>\frac12$, $\kappa>\frac12$.
\end{proof}

\begin{corollary}[Existence form of the final step]\label{cor:existence-form}
For every $\sigma>1$ the minimum $\min_{\mathscr O}\mathsf V_\sigma$ is attained, and the following are equivalent: \textup{(i)} (D); \textup{(ii)} $\mathscr C\ne\varnothing$; \textup{(iii)} there is $\boldsymbol{\kappa}\in\mathscr O$ (Section~\CGref{sec:lists}) with $\tilde\zeta_{\Z,\boldsymbol{\kappa}}=\zeta$; \textup{(iv)} $\min_{\mathscr O}\mathsf V_\sigma=0$.
\end{corollary}

\begin{proof}
A continuous function on a compact space attains its minimum (Proposition~\CGref{prop:lists-compact}).  (i)$\Leftrightarrow$(ii) is Theorem~\CGref{thm:rigidity}, and (ii)$\Rightarrow$(iii) is part of it.  (iii)$\Leftrightarrow$(iv) because $\mathsf V_\sigma(\boldsymbol{\kappa})=\sum_{n\ge2}|k_n(\boldsymbol{\kappa})|^2n^{-2\sigma}\ge0$, where $k_n(\boldsymbol\kappa)$ are the coefficients of $\tilde\zeta_{\Z,\boldsymbol\kappa}/\zeta$ \CGeqref{eq:K-kappa} and $k_1=1$, vanishes exactly when $k_n(\boldsymbol{\kappa})=0$ for all $n\ge2$, that is, when $\tilde\zeta_{\Z,\boldsymbol\kappa}=\zeta$.  (iii)$\Rightarrow$(i): $\tilde\zeta_{\Z,\boldsymbol{\kappa}}=\zeta$ has no zeros in $\Re s>\frac12$ (Theorem~\CGref{thm:atomized-continuation}(iii)), so, with the functional equation, the nontrivial zeros of $\zeta$ lie on $\Re s=\frac12$: the Riemann hypothesis holds, and Theorem~\CGref{thm:Hamburger-Z1} gives (D).  This is not circular: Theorem~\CGref{thm:Hamburger-Z1} rests on the Hadamard product and Hamburger's theorem, not on the present corollary.
\end{proof}

In this form the final step does not ask to identify $\zeta_{Z1}$ with $\zeta$: it asks for the existence of one list in $\mathscr O$ whose continuous and atomized zeta functions coincide, or whose atomized zeta function is $\zeta$; Theorem~\CGref{thm:rigidity} then identifies the list.  The path $\lambda\mapsto(\gamma_k(\lambda))$ is one segment in $\mathscr O$, from the clock to the list of $\zeta$, and along it $\mathsf V_\sigma$ is the function of Theorem~\CGref{thm:double-hamburger}.

\begin{remark}[What an existence proof uses]\label{rem:existence-input}
Compactness, continuity and convexity do not decide whether $\min_{\mathscr O}\mathsf V_\sigma$ is $0$: they do not involve the form of the functional equation.  The existence statement uses the functional equation with the factor $\Gamma(\frac s2)\pi^{-s/2}$, for which a function with both an absolutely convergent Dirichlet series and the functional equation exists, namely $\zeta$.  The first variation of the defect is explicit: by \CGeqref{eq:cell-mass-W}, for $n$ not a power of $2$,
\begin{equation}\label{eq:first-variation}
 \frac{\partial\tilde c_n(\boldsymbol{\kappa})}{\partial\kappa_k}=2\,\Im V_n(\kappa_k),
\end{equation}
so the vanishing of all coordinate derivatives of $\mathsf V_\sigma$ at a list is described by the functions $V_n$ of \CGeqref{eq:VW} alone.  Since $\mathscr O$ has empty interior in the product topology and the displayed derivatives are coordinate derivatives rather than a proved Fr\'echet gradient on a fixed Banach space, boundary minimizers must be treated by one-sided/KKT conditions.
\end{remark}


\subsection{The rung system with free masses}
\label{subsec:free-masses}

With the function $\mathsf k$ of \CGeqref{eq:transport-defs}, consider, for a list $\boldsymbol{\kappa}\in\mathscr O$ (Section~\CGref{sec:lists}) and real numbers $\mathsf a_n$, $n\ge2$, on which nothing is assumed, the system
\begin{equation}\label{eq:rung-system}
 \sum_{k\ge1}\cos(\kappa_ku)+\sum_{n\ge2}\mathsf a_n\,\delta(u-\log n)=\mathsf k(u)\qquad\text{in }\mathcal D'(0,\infty).
\end{equation}
In the variable $x=e^u$ it states that the Mellin count of $\zeta_{\boldsymbol{\kappa}}$ has no mass outside the integers, the masses at the integers being unknown.

\begin{theorem}[Existence form with free masses]\label{thm:free-masses}
(D) holds if and only if \CGeqref{eq:rung-system} has a solution $(\boldsymbol{\kappa},(\mathsf a_n))$ with $\boldsymbol{\kappa}\in\mathscr O$ (Section~\CGref{sec:lists}).  The solution is then unique: $\boldsymbol{\kappa}=(\gamma_k)$, and
\begin{equation}\label{eq:free-masses}
 \mathsf a_n=\frac{\log p}{2\sqrt n}\ \text{ if }n=p^m\text{ is a prime power},\qquad \mathsf a_n=0\ \text{ otherwise.}
\end{equation}
\end{theorem}

\begin{proof}
By \CGeqref{eq:distributional-EF} for the list $\boldsymbol{\kappa}$, with $u=\log x$, \CGeqref{eq:rung-system} is the statement
\begin{equation}\label{eq:free-masses-dJ}
 (\log x)\,dJ^{\rm M}_{\boldsymbol{\kappa}}(x)=\sum_{n\ge2}2\sqrt n\,\mathsf a_n\,\delta_n\qquad\text{in }\mathcal D'(1,\infty).
\end{equation}
If (D) holds, $J^{\rm M}_{Z1}=\sum_{p^m\le x}\frac1m$ (Theorem~\CGref{thm:cell-form}), which is \CGeqref{eq:free-masses-dJ} for the list $(\gamma_k)$ with the masses \CGeqref{eq:free-masses}.  Conversely let \CGeqref{eq:free-masses-dJ} hold.  Then $J^{\rm M}_{\boldsymbol{\kappa}}$ is constant on $(1,2)$ and between consecutive integers; the truncation at $2$ and the finite M\"obius inversion keep the support in the integers, so every cell defect vanishes and, by \CGeqref{eq:first-cell} and the proof of Theorem~\CGref{thm:joining-factor},
\begin{equation}\label{eq:free-masses-factor}
 \zeta_{\boldsymbol{\kappa}}(s)=\exp\bigl(\nu_6+\nu_7\,2^{-s}\bigr)\,\tilde\zeta_{\Z,\boldsymbol{\kappa}}(s)\qquad(\Re s>1)
\end{equation}
with real constants $\nu_6,\nu_7$.  The two factors on the right are the exponential $\exp(\nu_6+\nu_72^{-s})=e^{\nu_6}\sum_{r\ge0}\frac{\nu_7^r}{r!}2^{-rs}$, an absolutely convergent Dirichlet series on all of $\C$ supported on the powers of $2$, and $\tilde\zeta_{\Z,\boldsymbol{\kappa}}$, an absolutely convergent Dirichlet series in $\Re s>1$ by Corollary~\CGref{cor:abscissa-one}, which Theorem~\CGref{thm:transfer-O}(vi) extends to every list of $\mathscr O$; the product of two absolutely convergent Dirichlet series is one.  (Why the support stays in the integers: $J^{\rm M}_{\boldsymbol{\kappa}}$ constant between consecutive integers means $dJ^{\rm M}_{\boldsymbol{\kappa}}$ is a combination of $\delta_n$, $n\ge2$; the truncation $J^*$ at $2$ removes the cell $(1,2)$, on which $J^{\rm M}_{\boldsymbol{\kappa}}$ is a constant, and the finite M\"obius inversion \CGeqref{eq:pi-RVM} maps $\sum_n\mathsf w_n\delta_n$ to $\sum_m\frac{\mu(m)}m\sum_n\mathsf w_n\delta_{n^{1/m}}$, whose support is in the integers exactly when the cell masses vanish off the integers, which is the vanishing of every cell defect $\mathfrak h_n$; the constant on $(1,2)$ is what produces the factor $e^{\nu_6+\nu_72^{-s}}$ through \CGeqref{eq:first-cell}.)  Hence $\zeta_{\boldsymbol{\kappa}}$ satisfies (H1), and it satisfies (H0), (H2) and (H3) by Theorem~\CGref{thm:transfer-O}(i) (order one, functional equation \CGeqref{eq:Z1-FE}, pole at $1$ with residue $1$) for every list of $\mathscr O$; so Hamburger's theorem in the form of Subsection~\CGref{subsec:hamburger-table} gives $\zeta_{\boldsymbol{\kappa}}=C_H\zeta$, and $C_H=1$ by the residues at $s=1$.  The zeros of $\zeta_{\boldsymbol{\kappa}}$ are $\frac12\pm i\kappa_k$, so $\boldsymbol{\kappa}=(\gamma_k)$ and $\zeta_{Z1}=\zeta$.  Then $(\log x)\,dJ^{\rm M}=\sum_{p^m}\log p\,\delta_{p^m}$, which gives \CGeqref{eq:free-masses}.
\end{proof}

In Theorem~\CGref{thm:free-masses} the primes are not part of the data: the data are the function $\mathsf k$, which comes from the factor $G$ and the pole, the space $\mathscr O$, which is defined by the clock, and the condition that the support lie in the points $\log n$.  The masses of von Mangoldt are an output.  What the theorem proves is \emph{uniqueness}: a solution, if it exists, is the zeta list with the von Mangoldt masses.  \emph{Existence} of a solution is exactly (D) and is not proved; the theorem is a reformulation of the final step, not a step towards it.  Experiment~\CGref{exp:clock-to-zeros} solves a finite section of \CGeqref{eq:rung-system} numerically, starting from the clock; nothing proved here says that finite sections converge to a solution of the infinite distributional system.  The uniqueness part of Theorem~\CGref{thm:free-masses} extends to every list at bounded distance from the clock, without the third condition of \CGeqref{eq:list-space} and with no assumption on the masses (Theorem~\CGref{thm:rb-uniqueness}).

\begin{theorem}[Nondegeneracy of the solution]\label{thm:nondegenerate}
Let $(\boldsymbol{\kappa},(\mathsf a_n))$ be a solution of \CGeqref{eq:rung-system} with $\boldsymbol{\kappa}\in\mathscr O$ (Section~\CGref{sec:lists}).  Let $(\dot\kappa_k)$ be a bounded real sequence and $(\dot{\mathsf a}_n)_{n\ge2}$ real numbers with $\dot{\mathsf a}_n=O_\eps(n^{-1/2+\eps})$ for every $\eps>0$, such that
\begin{equation}\label{eq:rung-linearized}
 \sum_{k\ge1}\dot\kappa_k\,u\sin(\kappa_ku)=\sum_{n\ge2}\dot{\mathsf a}_n\,\delta(u-\log n)\qquad\text{in }\mathcal D'(0,\infty),
\end{equation}
the derivative of \CGeqref{eq:rung-system} in the direction $\bigl((\dot\kappa_k),(\dot{\mathsf a}_n)\bigr)$.  Then $\dot{\mathsf a}_n=0$ for every $n$, and $\sum_{k:\,\kappa_k=t}\dot\kappa_k=0$ for every value $t$ of the list.  If the ordinates are simple, $\dot\kappa_k=0$ for every $k$.
\end{theorem}

\begin{proof}
\emph{Step 1.}  The series $\mathsf J:=\sum_k\dot\kappa_k\,u\sin(\kappa_ku)$ converges in $\mathcal S'(\R)$: for $\phi\in\mathcal S$, three integrations by parts give $\int\phi(u)u\sin(\kappa u)\,du=\kappa^{-3}\int(u\phi)'''(u)\cos(\kappa u)\,du$ (the boundary terms vanish since $u\phi$ and its derivatives are Schwartz), so $|\int\phi\,u\sin(\kappa u)|\le\|(u\phi)'''\|_{L^1}\kappa^{-3}$, and $\sum_k|\dot\kappa_k|\kappa_k^{-3}<\infty$ since $\kappa_k\asymp k/\log k$.  Each term $u\sin(\kappa_ku)$ is even, so $\mathsf J$ is even.  By \CGeqref{eq:rung-linearized}, $\mathsf J$ vanishes on $(0,\log2)$, hence by evenness on $(-\log2,0)$, so the restriction of $\mathsf J$ to $(-\log2,\log2)$ is supported at $\{0\}$.  We use the following standard fact.

\emph{Lemma (point support).}  A distribution $\mathsf T\in\mathcal D'(-a,a)$ with $\operatorname{supp}\mathsf T\subset\{0\}$ is a finite combination $\mathsf T=\sum_{j\le m}c_j\delta_0^{(j)}$, with $m$ the order of $\mathsf T$ \cite[Thm.~2.3.4]{Hormander1990}; consequently $\langle\mathsf T,\chi\rangle=\sum_{j\le m}(-1)^jc_j\chi^{(j)}(0)=0$ for every $\chi\in C_c^\infty(-a,a)$ vanishing to order $m$ at $0$, and the same formula defines $\langle\mathsf T,\chi\rangle$ for $\chi\in C^m_c(-a,a)$.  Applied to $\mathsf J|_{(-\log2,\log2)}$ and using evenness, $\mathsf J=\sum_{j\le m,\ j\text{ even}}c_j\delta_0^{(j)}$ on $(-\log2,\log2)$ for some finite $m$.  Fix an even $\mathsf x_0\in C_c^\infty(-\log2,\log2)$ equal to $1$ near $0$, and for $\Re z>0$ put
\begin{equation}\label{eq:nondeg-H1}
 \mathsf H_1(z):=\sum_{k\ge1}\dot\kappa_k\int_0^\infty\bigl(1-\mathsf x_0(u)\bigr)\,u\sin(\kappa_ku)\,e^{-zu}\,du=\tfrac12\bigl\langle\mathsf J,(1-\mathsf x_0)\,e^{-z|u|}\bigr\rangle .
\end{equation}
Approximating by test functions with compact support in $(0,\infty)$ and using \CGeqref{eq:rung-linearized}, $\mathsf H_1(z)=\sum_{n\ge2}\dot{\mathsf a}_nn^{-z}$ for $\Re z>\frac12$, where the series converges absolutely.

\emph{Step 2.}  For $\Re z>0$, $\int_0^\infty u\sin(\kappa u)e^{-zu}du=2\kappa z/(z^2+\kappa^2)^2$, so
\begin{equation}\label{eq:nondeg-split}
 \mathsf H_1(z)=\sum_{k\ge1}\dot\kappa_k\,\frac{2\kappa_kz}{(z^2+\kappa_k^2)^2}-\mathsf y(z),\qquad \mathsf y(z):=\sum_{k\ge1}\dot\kappa_k\int_0^{\log2}\mathsf x_0(u)\,u\sin(\kappa_ku)\,e^{-zu}\,du .
\end{equation}
\emph{$\mathsf y$ is a polynomial.}  Let $m$ be the order of $\mathsf J$ at $0$ from Step~1 and fix an integer $A\ge\max(m,3)$.  Write $e^{-zu}=\sum_{j\le A}\frac{(-z)^j}{j!}u^j+\mathsf R_A(z,u)$ with $\mathsf R_A(z,u)=\frac{(-z)^{A+1}}{A!}\int_0^1(1-\vartheta)^Au^{A+1}e^{-\vartheta zu}\,d\vartheta$, so that $\mathsf R_A(z,\cdot)$ is smooth on $[0,\log2]$ and vanishes to order $A+1$ at $u=0$.  Accordingly $\mathsf y=\mathsf y_{\rm pol}+\mathsf y_{\rm rem}$.  \emph{Polynomial part.}  $\mathsf y_{\rm pol}(z)=\sum_{j\le A}\frac{(-z)^j}{j!}\sum_k\dot\kappa_k\,\mathsf i_{j,k}$ with $\mathsf i_{j,k}:=\int_0^{\log2}\mathsf x_0(u)u^{j+1}\sin(\kappa_ku)\,du$.  For $g\in C^\infty_c[0,\infty)$ two integrations by parts give $\int_0^\infty g(u)\sin(\kappa u)\,du=g(0)/\kappa-\kappa^{-2}\int_0^\infty g''(u)\sin(\kappa u)\,du=g(0)/\kappa-g''(0)/\kappa^3+O_g(\kappa^{-5})$; with $g=\mathsf x_0u^{j+1}$ one has $g(0)=0$ for all $j\ge0$, so $\mathsf i_{j,k}=O_j(\kappa_k^{-3})$ and each inner series converges absolutely since $\sum_k|\dot\kappa_k|\kappa_k^{-3}<\infty$.  Thus $\mathsf y_{\rm pol}$ is a polynomial of degree at most $A$.  \emph{Remainder part.}  Put $\chi_z(u):=\mathsf x_0(u)\mathsf R_A(z,|u|)$, an even function with compact support in $(-\log2,\log2)$, of class $C^A$ (because $|u|^{A+1}$ is $C^A$) and vanishing to order $A+1$ at $0$.  Since $u\sin(\kappa u)$ is even, $\mathsf y_{\rm rem}(z)=\frac12\sum_k\dot\kappa_k\int_{\R}\chi_z(u)\,u\sin(\kappa_ku)\,du$.  Let $\chi_{z,\epsilon}:=\chi_z*\varrho_\epsilon$ with a standard mollifier; for small $\epsilon$ these are in $C_c^\infty(-\log2,\log2)$, $\|(u\chi_{z,\epsilon})'''\|_{L^1}$ is bounded uniformly in $\epsilon$ because $A\ge3$, and $\chi_{z,\epsilon}^{(j)}(0)\to\chi_z^{(j)}(0)=0$ for $j\le A$.  By Step~1 and the point-support lemma, $\sum_k\dot\kappa_k\int\chi_{z,\epsilon}\,u\sin(\kappa_ku)\,du=\langle\mathsf J,\chi_{z,\epsilon}\rangle=\sum_{j\le m}(-1)^jc_j\chi_{z,\epsilon}^{(j)}(0)\to0$ as $\epsilon\to0$; on the other side, each term converges to $\dot\kappa_k\int\chi_z\,u\sin(\kappa_ku)\,du$ and the series is dominated by $\sum_k|\dot\kappa_k|\kappa_k^{-3}\sup_\epsilon\|(u\chi_{z,\epsilon})'''\|_{L^1}<\infty$, so the limit may be taken termwise.  Hence $\mathsf y_{\rm rem}\equiv0$ and $\mathsf y=\mathsf y_{\rm pol}$ is a polynomial.

\emph{$\mathsf y\equiv0$.}  As $z=\sigma\to+\infty$: (a) $\mathsf H_1(\sigma)=\sum_{n\ge2}\dot{\mathsf a}_nn^{-\sigma}\to0$, since for $\sigma\ge1$ the terms are dominated by $|\dot{\mathsf a}_n|n^{-1}$, $\sum_n|\dot{\mathsf a}_n|n^{-1}<\infty$ because $\dot{\mathsf a}_n=O_\eps(n^{-1/2+\eps})$, and each term tends to $0$; (b) the first sum in \CGeqref{eq:nondeg-split} is $O(\sigma^{-1}\log\sigma)$: the terms with $\kappa_k\le\sigma$ are at most $2\kappa_k\sigma/\sigma^4$ and add up to at most $2\sigma^{-3}\sum_{\kappa_k\le\sigma}\kappa_k\le2\sigma^{-2}N_{\boldsymbol{\kappa}}(\sigma)=O(\sigma^{-1}\log\sigma)$, while the terms with $\kappa_k>\sigma$ are at most $2\sigma\kappa_k^{-3}$ and, by partial summation with $N_{\boldsymbol{\kappa}}(t)=O(t\log t)$, $\sum_{\kappa_k>\sigma}\kappa_k^{-3}=\int_\sigma^\infty t^{-3}dN_{\boldsymbol{\kappa}}(t)=O(\sigma^{-2}\log\sigma)$ (the boundary term $N_{\boldsymbol{\kappa}}(\sigma)\sigma^{-3}$ included), which gives $O(\sigma^{-1}\log\sigma)$ again; here $(\dot\kappa_k)$ is bounded.  Hence the polynomial $\mathsf y(\sigma)=\sum_k\dot\kappa_k2\kappa_k\sigma/(\sigma^2+\kappa_k^2)^2-\mathsf H_1(\sigma)\to0$ as $\sigma\to+\infty$, so $\mathsf y\equiv0$.

\emph{Step 3.}  Put $\mathsf H(z):=\sum_{n\ge2}\frac{\dot{\mathsf a}_n}{\log n}n^{-z}$ for $\Re z>\frac12$, so that $\mathsf H'=-\mathsf H_1$.  By Step~2 there is a real constant $\nu_8$ with
\begin{equation}\label{eq:nondeg-H}
 \mathsf H(z)=\nu_8+\sum_{k\ge1}\dot\kappa_k\Bigl(\frac{\kappa_k}{z^2+\kappa_k^2}-\frac1{\kappa_k}\Bigr),
\end{equation}
which continues $\mathsf H$ to a meromorphic function in $\C$, even in $z$, with at most simple poles at $z=\pm i\kappa_k$, of residue $\pm\frac1{2i}\sum_{k':\kappa_{k'}=\kappa_k}\dot\kappa_{k'}$.

\emph{Step 4: finite order.}  By Theorem~\CGref{thm:free-masses}, $\zeta_{\boldsymbol{\kappa}}=\zeta$, with zeros $\frac12\pm i\kappa_k$ of the multiplicity of the list.  Put $\mathsf Z(s):=\mathsf H(s-\frac12)\,\zeta(s)$.  At $s=\frac12\pm i\kappa$ the function $\mathsf H(s-\frac12)$ has at most a simple pole (the residues of the entries of the list with the same value $\kappa$ add up, \CGeqref{eq:nondeg-H}) and $\zeta$ has a zero of order equal to the multiplicity of $\kappa$ in the list, which is at least $1$; so $\mathsf Z$ is holomorphic there and $(s-1)\mathsf Z(s)$ is entire.  \emph{Admissible radii.}  Call $r\ge3$ admissible if $\operatorname{dist}(r,\{\kappa_k\})\ge1/\log r$.  Since $N_{\boldsymbol{\kappa}}(2R+1)-N_{\boldsymbol{\kappa}}(R-1)=(\frac R{2\pi}+o(R))\log R$ (bounded displacement from the clock), the union of the intervals of radius $1/\log R$ about the $\kappa_k\in[R-1,2R+1]$ has total length at most $(\frac1\pi+o(1))R<R$, so every interval $[R,2R]$ with $R\ge R_0$ contains an admissible radius; choose admissible $r_j\in[2^jR_0,2^{j+1}R_0]$, so that $r_j\to\infty$ and $r_{j+1}\le4r_j$.  \emph{Bound on admissible circles.}  Write \CGeqref{eq:nondeg-H} as $\mathsf H(z)=\nu_8-\sum_k\dot\kappa_k\frac{z^2}{\kappa_k(z^2+\kappa_k^2)}$ and let $|z|=r$ be admissible.  Then $|z\mp i\kappa_k|\ge|r-\kappa_k|$ and $|z-i\kappa_k|+|z+i\kappa_k|\ge2\kappa_k$.  Shell $\kappa_k\ge2r$: $|z^2+\kappa_k^2|\ge\kappa_k^2-r^2\ge\frac34\kappa_k^2$, terms $\le\frac43r^2\kappa_k^{-3}\sup|\dot\kappa|$, sum $O(r^2\cdot r^{-2}\log r)=O(\log r)$ by partial summation.  Shell $\kappa_k\le r/2$: $|z^2+\kappa_k^2|\ge\frac34r^2$, terms $\le\frac43\kappa_k^{-1}\sup|\dot\kappa|$, sum $O(\sum_{\kappa_k\le r}\kappa_k^{-1})=O(\log^2r)$.  Shell $r/2<\kappa_k<2r$: one of $|z\pm i\kappa_k|$ is at least $\kappa_k\ge r/2$ and the other is at least $\max(|r-\kappa_k|,1/\log r)$, so the term is at most $\frac{4r}{\kappa_k}\cdot\frac{\sup|\dot\kappa|}{\max(|r-\kappa_k|,1/\log r)}\le\frac{8\sup|\dot\kappa|}{\max(|r-\kappa_k|,1/\log r)}$; the $O(\log r)$ ordinates with $|\kappa_k-r|\le1$ contribute $O(\log^2r)$, and for $i\ge0$ the $O(2^i\log r)$ ordinates with $2^i<|\kappa_k-r|\le2^{i+1}$ contribute $O(\log r)$ each shell, $i\le\log_2r$, so $O(\log^2r)$ in all (multiplicities are counted by $N_{\boldsymbol{\kappa}}$ throughout).  Hence $|\mathsf H(z)|\le C_{\mathsf H}\log^2r$ on every admissible circle $|z|=r\ge R_0$.  \emph{From circles to the plane.}  $(s-1)\zeta(s)$ is entire of order one, so $|(s-1)\zeta(s)|\le\exp(|s|^{1+\eps})$ for $|s|$ large.  By the maximum principle for the entire function $(s-1)\mathsf Z(s)$ on the disc $|s-\frac12|\le r_j$,
\begin{equation}\label{eq:nondeg-order}
 \max_{|s-\frac12|\le r_j}|(s-1)\mathsf Z(s)|\le C_{\mathsf H}\log^2r_j\cdot\max_{|s-\frac12|=r_j}|(s-1)\zeta(s)|\le\exp\bigl(2r_j^{1+\eps}\bigr)
\end{equation}
for $j\ge j_0(\eps)$.  Given $R\ge r_{j_0}$ choose the least $j$ with $r_j\ge R$; then $r_j\le4R$ and \CGeqref{eq:nondeg-order} gives $\max_{|s-\frac12|\le R}|(s-1)\mathsf Z|\le\exp(2(4R)^{1+\eps})$.  So $(s-1)\mathsf Z$ is entire of order at most one: this is (H0).  $\mathsf Z$ is the product of $\sum_n\dot{\mathsf a}_n\sqrt n(\log n)^{-1}n^{-s}$ and $\zeta(s)$, both absolutely convergent in $\Re s>1$: this is (H1).  Since $\mathsf H$ is even, $\mathsf Z(1-s)=\mathsf H(s-\frac12)\zeta(1-s)$, so $\mathsf Z$ satisfies the functional equation of $\zeta$: this is (H2) and (H3) with the second function equal to $\mathsf Z$.  Hamburger's theorem gives $\mathsf Z=C_H\zeta$, that is, $\mathsf H\equiv C_H$.  Since $\mathsf H(\sigma)\to0$ as $\sigma\to+\infty$, $\mathsf H\equiv0$: every $\dot{\mathsf a}_n$ and every residue vanishes.
\end{proof}

\begin{remark}[Scope of Theorem~\CGref{thm:nondegenerate}]\label{rem:nondegenerate-scope}
The theorem is conditional on the existence of a solution; it shows that a solution is isolated to first order, the Jacobian of \CGeqref{eq:rung-system} being injective there on the directions considered, except for the splitting of multiple ordinates.  This is a necessary endpoint nondegeneracy condition for any continuation scheme; it is not, by itself, an implicit-function theorem, because no isomorphism between fixed Banach spaces is proved.  Steps~1--3 do not use that $\boldsymbol{\kappa}$ is a solution: at any list of $\mathscr O$ the kernel of the linearization consists of the Dirichlet series $\mathsf H$ that are even meromorphic functions with simple poles at $\pm i\kappa_k$ of the form \CGeqref{eq:nondeg-H}.  Such series exist for lists in arithmetic progression, for instance $(2^{z}+2^{-z}-2\cos\nu_9)^{-1}$, which is even, a power series in $2^{-z}$ for $\Re z>0$, with simple poles at $z=\pm i(\nu_9+2\pi m)/\log2$, $m\in\Z$; these lists do not belong to $\mathscr O$.  For perturbations $\dot{\mathsf a}_n$ of polynomial growth the same proof applies with Hamburger's theorem for a finite abscissa of absolute convergence \cite{BochnerChandrasekharan1956}.  The injectivity of the Jacobian at the clock and along a path, and with it the existence of a solution, are not proved.
\end{remark}

\subsection{The gradient of the vertical defect}
\label{subsec:gradient}

For Dirichlet series $A=\sum_na_nn^{-s}$, $B=\sum_nb_nn^{-s}$ and $\sigma>1$ put $\langle A,B\rangle_\sigma:=\sum_{n\ge1}a_n\bar b_n\,n^{-2\sigma}$, and for $\boldsymbol{\kappa}\in\mathscr O$ (Section~\CGref{sec:lists})
\begin{equation}\label{eq:K-kappa}
 \mathsf K_{\boldsymbol{\kappa}}:=\frac{\tilde\zeta_{\Z,\boldsymbol{\kappa}}}\zeta=\sum_{n\ge1}k_n(\boldsymbol{\kappa})\,n^{-s},\qquad \mathsf V_\sigma(\boldsymbol{\kappa})=\langle\mathsf K_{\boldsymbol{\kappa}},\mathsf K_{\boldsymbol{\kappa}}\rangle_\sigma-1,\qquad \mathsf r_m(\boldsymbol{\kappa}):=\Re\langle\mathsf K_{\boldsymbol{\kappa}}\mathsf L_m,\mathsf K_{\boldsymbol{\kappa}}\rangle_\sigma ,
\end{equation}
with $\mathsf L_m(s)=-\log(1-m^{-s})$ as in Proposition~\CGref{prop:one-ordinate}.  By Proposition~\CGref{prop:one-ordinate} the derivative $\partial\tilde c_m(\boldsymbol{\kappa})/\partial\kappa_k$ is a function of $m$ and of the value $\kappa_k$ alone; denote it by $\mathsf u_m(\kappa_k)$, so that $\mathsf u_m=2\,\Im V_m$ when $m$ is not a power of $2$ \CGeqref{eq:first-variation}.

\begin{proposition}[Critical points of the vertical defect]\label{prop:critical-points}
Let $\sigma>1$ and $\boldsymbol{\kappa}\in\mathscr O$ (Section~\CGref{sec:lists}).  Then $|\mathsf r_m(\boldsymbol{\kappa})|\le2m^{-\sigma}\bigl(1+\mathsf V_\sigma(\boldsymbol{\kappa})\bigr)$, $|\mathsf u_m(t)|\le Cm^{-1/2}$ off the powers of $2$ and $|\mathsf u_m(t)|\le C$ on them, and
\begin{equation}\label{eq:gradient}
 \frac{\partial\mathsf V_\sigma}{\partial\kappa_k}(\boldsymbol{\kappa})=2\,\mathsf Q_{\boldsymbol{\kappa}}(\kappa_k),\qquad \mathsf Q_{\boldsymbol{\kappa}}(t):=\sum_{m\ge2}\mathsf r_m(\boldsymbol{\kappa})\,\mathsf u_m(t),
\end{equation}
the series converging absolutely and uniformly in $t$.  If $\mathsf V_\sigma(\boldsymbol{\kappa})=0$, then $\mathsf r_m(\boldsymbol{\kappa})=0$ for every $m$ and $\mathsf Q_{\boldsymbol{\kappa}}\equiv0$.
\end{proposition}

\begin{proof}
Since $\log\mathsf K_{\boldsymbol{\kappa}}=\sum_m\bigl(\tilde c_m(\boldsymbol{\kappa})-\mathbf 1_{\rm prime}(m)\bigr)\mathsf L_m$, $\partial\mathsf K_{\boldsymbol{\kappa}}/\partial\tilde c_m=\mathsf K_{\boldsymbol{\kappa}}\mathsf L_m$ and $\partial\mathsf V_\sigma/\partial\tilde c_m=2\mathsf r_m$.  Multiplication by $m^{-rs}$ multiplies the norm $\langle\cdot,\cdot\rangle_\sigma^{1/2}$ by $m^{-r\sigma}$, so $\langle\mathsf K\mathsf L_m,\mathsf K\mathsf L_m\rangle_\sigma^{1/2}\le\langle\mathsf K,\mathsf K\rangle_\sigma^{1/2}\sum_{r\ge1}m^{-r\sigma}/r\le2m^{-\sigma}\langle\mathsf K,\mathsf K\rangle_\sigma^{1/2}$, and the bound for $\mathsf r_m$ follows by the Cauchy--Schwarz inequality.  In \CGeqref{eq:VW} the support of $\varphi_m(y^j)$ has length at most $\frac2jm^{1/j-1}$ and the integrand is at most $m^{-1/(2j)}$ in absolute value, so $|V_m(t)|\le\sum_j\frac2{j^2}m^{-1+1/(2j)}\le Cm^{-1/2}$; for $m=2^j$ the cut term is $-\frac{\mu(j)}j\int_2^3J_{\boldsymbol{\kappa}}(y)\frac d{dy}\varphi_m(y^j)\,dy$ (proof of Theorem~\CGref{thm:mean-square-cells}), and its coordinate derivative needs $\partial_{\kappa_k}J_{\boldsymbol{\kappa}}$ on $(2,3)$: in the series \CGeqref{eq:J-X} only the $k$-th term depends on $\kappa_k$, so the derivative exists and equals $\partial_{\kappa_k}J_{\boldsymbol{\kappa}}(y)=-2\,\Re\bigl(i\,y^{1/2+i\kappa_k}/(\tfrac12+i\kappa_k)\bigr)$, a bounded continuous function of $y\in(2,3)$ with $|\partial_{\kappa_k}J_{\boldsymbol{\kappa}}(y)|\le2\sqrt3/\kappa_k\le2\sqrt3/\tau_0^+$ uniformly on $\mathscr O$ (Lemma~\CGref{lem:L2-partial-sums}(b) is the same computation for the direction $\eta_k$); hence $|\mathsf u_m(t)|\le\frac2j\cdot2\sqrt3/t\cdot\mathrm{TV}(\varphi_m(y^j))\le C$ for $m=2^j$, with $\mathrm{TV}(\varphi_m(y^j))=2\|\varphi\|_\infty$.  This gives the two bounds for $\mathsf u_m$, and they are uniform in $t\ge\tau_0^+$.  The chain rule is justified first for finite Euler products $\prod_{m\le M}$ and then by passage to the limit $M\to\infty$ in the weighted space $\ell^1(n^{-\sigma'})$, $1<\sigma'<\sigma$, of cell masses, in which the coordinate derivatives $\partial_{\kappa_k}\tilde c_m=\mathsf u_m(\kappa_k)$ are bounded uniformly in $m$ and $M$ by the bounds just proved, so that $\sum_nn^{-2\sigma}\partial_{\kappa_k}|k_n|^2$ converges uniformly in the truncation.  It gives \CGeqref{eq:gradient}, with $\sum_m|\mathsf r_m||\mathsf u_m(t)|\le C\sum_mm^{-\sigma-1/2}+C\sum_j2^{-j\sigma}<\infty$ uniformly in $t$, which is the asserted absolute and uniform convergence on all of $t\ge\tau_0^+$, not only on compact sets.  If $\mathsf V_\sigma=0$, then $\mathsf K=1$ and $\mathsf r_m=\Re\langle\mathsf L_m,1\rangle_\sigma=0$, because $\mathsf L_m$ has no constant term.
\end{proof}

A list $\boldsymbol{\kappa}\in\mathscr O$ (Section~\CGref{sec:lists}) at which all coordinate derivatives of $\mathsf V_\sigma$ vanish is therefore a list all of whose ordinates are zeros of the single real-analytic function $\mathsf Q_{\boldsymbol{\kappa}}$.  The minimum of $\mathsf V_\sigma$ on $\mathscr O$ is attained (Corollary~\CGref{cor:existence-form}); if every minimizer satisfied the appropriate coordinate/KKT stationarity conditions and those conditions forced $\mathsf V_\sigma=0$, the minimum would be $0$ and (D) would follow.  This is not proved.  A count of zeros does not exclude critical points with $\mathsf V_\sigma>0$: $\mathsf Q_{\boldsymbol{\kappa}}(t)$ is a sum of terms of frequency $\log m$ with $m\gtrsim t/2\pi$, so it oscillates at the scale $1/\log t$, which is the scale of the gaps between the ordinates.


\section{Acceptance criteria and models in which the final step is a theorem}
\label{sec:models}

In the construction of this paper the continuous function $\zeta_\lambda$ is built from a list of ordinates and the discrete function $\tilde\zeta_{\Z,\lambda}$ from the cells, and the two are joined by $\tilde h_\lambda$.  This section states, without $\zeta$, the general criterion under which a list is accepted by both sides, and gives three models: a non-periodic list accepted by a discrete category with two generators, where the acceptance is proved by one unitary matrix; a finite Euler product over the primes, to which the constructions of Sections~\CGref{sec:clock}--\CGref{sec:lists} apply verbatim and for which the final step is a theorem; and the diagonal operator over all primes, which shows where the argument of the finite models stops.  The complex variable of the first two subsections is denoted by $\mathsf z$, and the half-plane $\Im\mathsf z>0$ plays the role of $\Re s>1$.

\subsection{The acceptance criterion: window and frequencies}
\label{subsec:acceptance}

Let $\mathfrak k>0$ and let $\boldsymbol{\kappa}=(\kappa_k)_{k\ge1}$ be a nondecreasing list of positive numbers with $|\kappa_k-k\pi/\mathfrak k|$ bounded.  The \emph{continuous side} is the canonical product
\begin{equation}\label{eq:model-product}
 \mathfrak C_{\boldsymbol{\kappa}}(\mathsf z):=\mathsf z\prod_{k\ge1}\Bigl(1-\frac{\mathsf z^2}{\kappa_k^2}\Bigr),
\end{equation}
an odd entire function of order one with real zeros; the oddness plays the role of the functional equation.  We assume that $\mathfrak C_{\boldsymbol{\kappa}}$ has exponential type $\mathfrak k$ and polynomial growth on the real axis, as $\sin(\mathfrak k\mathsf z)$ has for the list $\kappa_k=k\pi/\mathfrak k$.  The \emph{discrete side} is given by a closed, locally finite set $\mathscr F\subset(0,\infty)$ of frequencies ($\mathscr F\cap[0,R]$ is finite for every $R$; the set need not be uniformly discrete, and in the two-edge model below it is not) that is closed under addition; put $\mathscr F^{+}:=\mathscr F\cup\{0\}$.  The discrete category consists of the generalized Dirichlet series $\sum_{\mathfrak l\in\mathscr F^{+}}\mathfrak c_{\mathfrak l}\,e^{i\mathfrak l\mathsf z}$, absolutely convergent in some half-plane $\Im\mathsf z>\nu_{10}$ (the number $\nu_{10}$ is not prescribed and may be enlarged), and of their Euler forms $\exp\bigl(\sum_{\mathfrak l\in\mathscr F}\mathfrak b_{\mathfrak l}\,e^{i\mathfrak l\mathsf z}\bigr)$.  For $\Im\mathsf z>0$,
\begin{equation}\label{eq:model-logder}
 \frac{\mathfrak C_{\boldsymbol{\kappa}}'}{\mathfrak C_{\boldsymbol{\kappa}}}(\mathsf z)=\frac1{\mathsf z}+\sum_{k\ge1}\frac{2\mathsf z}{\mathsf z^2-\kappa_k^2}=-i\int_0^\infty e^{i\mathsf zu}\Bigl(1+2\sum_{k\ge1}\cos(\kappa_ku)\Bigr)du ,
\end{equation}
the bracket being a tempered distribution in $u$: it is the spectral density of the list.

\begin{theorem}[Acceptance criterion]\label{thm:acceptance}
For a list $\boldsymbol{\kappa}$ as above the following are equivalent:
\begin{enumerate}[label=\textup{(\roman*)}]
\item the spectral density is atomic on the frequencies: there are real $\mathfrak w_{\mathfrak l}$ with $\sum_{\mathfrak l\in\mathscr F}|\mathfrak w_{\mathfrak l}|e^{-\nu_{10}\mathfrak l}<\infty$ for some admissible $\nu_{10}$ and
\begin{equation}\label{eq:acceptance-density}
 1+2\sum_{k\ge1}\cos(\kappa_ku)=\sum_{\mathfrak l\in\mathscr F}\mathfrak w_{\mathfrak l}\,\delta(u-\mathfrak l)\qquad\text{in }\mathcal D'(0,\infty);
\end{equation}
\item the continuous function has an Euler form: for $\Im\mathsf z>\nu_{10}$, with a constant $\mathfrak c_0\ne0$,
\begin{equation}\label{eq:acceptance-euler}
 e^{i\mathfrak k\mathsf z}\,\mathfrak C_{\boldsymbol{\kappa}}(\mathsf z)=\mathfrak c_0\exp\Bigl(\sum_{\mathfrak l\in\mathscr F}\mathfrak b_{\mathfrak l}\,e^{i\mathfrak l\mathsf z}\Bigr),\qquad\mathfrak b_{\mathfrak l}=-\frac{\mathfrak w_{\mathfrak l}}{\mathfrak l};
\end{equation}
\item the continuous function is a finite sum in the discrete category:
\begin{equation}\label{eq:acceptance-window}
 e^{i\mathfrak k\mathsf z}\,\mathfrak C_{\boldsymbol{\kappa}}(\mathsf z)=\sum_{\mathfrak l\in\mathscr F^{+}\cap[0,2\mathfrak k]}\mathfrak c_{\mathfrak l}\,e^{i\mathfrak l\mathsf z},\qquad\mathfrak c_0\ne0 .
\end{equation}
\end{enumerate}
\end{theorem}

\begin{proof}
(i)$\Leftrightarrow$(ii): the two sides of \CGeqref{eq:acceptance-density} extend to distributions on $\R$ supported in $[0,\infty)$ (the left side as in \CGeqref{eq:model-logder}, the right side as the measure $\sum_{\mathfrak l}\mathfrak w_{\mathfrak l}\delta_{\mathfrak l}$), and (i) says that their difference vanishes on $(0,\infty)$; being supported at $\{0\}$ and of finite order, the difference is a finite combination of derivatives of $\delta_0$ \cite[Thm.~2.3.4]{Hormander1990}, whose transform $u\mapsto e^{i\mathsf zu}$ is a polynomial $\mathsf p(\mathsf z)$.  Hence, by \CGeqref{eq:model-logder}, (i) is the identity $\mathfrak C'_{\boldsymbol{\kappa}}/\mathfrak C_{\boldsymbol{\kappa}}=-i\sum_{\mathfrak l}\mathfrak w_{\mathfrak l}e^{i\mathfrak l\mathsf z}+\mathsf p(\mathsf z)$ in $\Im\mathsf z>\nu_{10}$.  Along $\mathsf z=iy$, $y\to+\infty$, the exponential sum tends to $0$ and $\mathfrak C'_{\boldsymbol{\kappa}}/\mathfrak C_{\boldsymbol{\kappa}}(iy)=-i\bigl(\frac1y+\sum_k\frac{2y}{y^2+\kappa_k^2}\bigr)\to-i\mathfrak k$ (type $\mathfrak k$); so the polynomial $\mathsf p$ has a finite limit along a ray, hence is constant, and the constant is $-i\mathfrak k$.  Integrating, $\log\mathfrak C_{\boldsymbol{\kappa}}=\text{const}-i\mathfrak k\mathsf z-\sum_{\mathfrak l}\frac{\mathfrak w_{\mathfrak l}}{\mathfrak l}e^{i\mathfrak l\mathsf z}$, and conversely.  (ii)$\Rightarrow$(iii): expanding the exponential, $\mathsf N(\mathsf z):=e^{i\mathfrak k\mathsf z}\mathfrak C_{\boldsymbol{\kappa}}(\mathsf z)=\sum_{\mathfrak l\in\mathscr F^{+}}\mathfrak c_{\mathfrak l}e^{i\mathfrak l\mathsf z}$, absolutely convergent in $\Im\mathsf z>\nu_{10}$.  On the other hand $\mathsf N$ satisfies $|\mathsf N(\mathsf z)|\le C(1+|\mathsf z|)^{N}e^{2\mathfrak k\max(0,-\Im\mathsf z)}$: for $\Im\mathsf z\ge0$ it is bounded by a polynomial, since $\mathfrak C_{\boldsymbol{\kappa}}$ has type $\mathfrak k$ and polynomial growth on the real axis and $|e^{i\mathfrak k\mathsf z}|=e^{-\mathfrak k\Im\mathsf z}$ absorbs the exponential growth in the upper half-plane, and for $\Im\mathsf z<0$ the factor $e^{i\mathfrak k\mathsf z}$ contributes $e^{\mathfrak k|\Im\mathsf z|}$ on top of the type $\mathfrak k$ of $\mathfrak C_{\boldsymbol{\kappa}}$. By the Paley--Wiener--Schwartz theorem \cite[Thm.~7.3.1]{Hormander1990}, used with the transform $\mathsf T\mapsto\langle\mathsf T,e^{i\mathsf zu}\rangle$ of \CGeqref{eq:model-logder} (for which a distribution supported in $[0,2\mathfrak k]$ has a transform bounded by $C(1+|\mathsf z|)^Ne^{2\mathfrak k\max(0,-\Im\mathsf z)}$, and conversely), an entire function with this bound is the transform of a distribution supported in $[0,2\mathfrak k]$, so the Fourier transform of the restriction of $\mathsf N$ to a horizontal line is supported in $[0,2\mathfrak k]$.  On the line $\Im\mathsf z=\nu_{10}+1$ this transform is $\sum_{\mathfrak l}\mathfrak c_{\mathfrak l}e^{-\mathfrak l(\nu_{10}+1)}\delta_{\mathfrak l}$, so $\mathfrak c_{\mathfrak l}=0$ for $\mathfrak l>2\mathfrak k$, and the sum is finite because $\mathscr F^{+}\cap[0,2\mathfrak k]$ is finite.  (iii)$\Rightarrow$(ii): write $\mathsf N=\mathfrak c_0(1+\mathsf w)$ with $\mathsf w(\mathsf z)=\sum_{0<\mathfrak l\le2\mathfrak k}(\mathfrak c_{\mathfrak l}/\mathfrak c_0)e^{i\mathfrak l\mathsf z}$; for $\Im\mathsf z\ge\nu'$ with $\sum_{0<\mathfrak l\le2\mathfrak k}|\mathfrak c_{\mathfrak l}/\mathfrak c_0|e^{-\mathfrak l\nu'}\le\frac12$ one has $|\mathsf w|\le\frac12$, so $\log(1+\mathsf w)=\sum_{r\ge1}(-1)^{r+1}\mathsf w^r/r$ converges absolutely, and expanding the powers gives an absolutely convergent series of exponentials whose frequencies are finite sums of elements of $\mathscr F$, hence elements of $\mathscr F$ (closure under addition); this proves (ii) in the half-plane $\Im\mathsf z>\nu'$, after increasing $\nu_{10}$ to $\nu'$ if necessary.
\end{proof}

In model identities of the form $2+2\sum_k\cos(\kappa_k u)$, the leading $2$ records a double zero at the origin; the normalized criterion \CGeqref{eq:acceptance-density} starts with $1$ because $\mathfrak C_{\boldsymbol{\kappa}}$ has a simple zero there.

The continuous side contributes a \emph{window}: by its type and its symmetry, the frequencies of $e^{i\mathfrak k\mathsf z}\mathfrak C_{\boldsymbol{\kappa}}$ lie in $[0,2\mathfrak k]$.  The discrete side contributes the \emph{admissible frequencies} $\mathscr F^{+}$.  A list is accepted by both sides if and only if its function lives in the intersection, and it is then determined by the finitely many coefficients $\mathfrak c_{\mathfrak l}$, $\mathfrak l\in\mathscr F^{+}\cap[0,\mathfrak k)$, the others being fixed by the oddness.  If $\mathscr F=2\mathfrak k\N$, the intersection is $\{0,2\mathfrak k\}$, the only function is a multiple of $\sin(\mathfrak k\mathsf z)$, and the only accepted list is the clock $\kappa_k=k\pi/\mathfrak k$.  Hamburger's theorem is the statement that for the window fixed by the factor $\Gamma(\frac s2)\pi^{-s/2}$ and the frequencies $\log n$ the space of accepted functions has dimension one.  The criterion says nothing about whether a given list, the zeta list in particular, belongs to the intersection; for $\zeta$ that membership is (D).

\subsection{A non-periodic accepted list: one unitary matrix}
\label{subsec:graph-model}

Let $\mathfrak p,\mathfrak q>0$ with $\mathfrak p/\mathfrak q\notin\mathbb Q$, $0<\mathfrak v<1$, $\mathfrak k=\mathfrak p+\mathfrak q$, and
\begin{equation}\label{eq:graph-matrices}
 \mathsf U:=\begin{pmatrix}\mathfrak v&\sqrt{1-\mathfrak v^2}\\ \sqrt{1-\mathfrak v^2}&-\mathfrak v\end{pmatrix},\qquad
 \mathsf X(\mathsf z):=\begin{pmatrix}e^{2i\mathfrak p\mathsf z}&0\\0&e^{2i\mathfrak q\mathsf z}\end{pmatrix},
\end{equation}
\begin{equation}\label{eq:graph-det}
 \mathfrak D(\mathsf z):=\det\bigl(I-\mathsf U\mathsf X(\mathsf z)\bigr)=1-\mathfrak v\,e^{2i\mathfrak p\mathsf z}+\mathfrak v\,e^{2i\mathfrak q\mathsf z}-e^{2i(\mathfrak p+\mathfrak q)\mathsf z}=-2i\,e^{i(\mathfrak p+\mathfrak q)\mathsf z}\,\mathfrak S(\mathsf z),
\end{equation}
\begin{equation}\label{eq:graph-sine}
 \mathfrak S(\mathsf z):=\sin\bigl((\mathfrak p+\mathfrak q)\mathsf z\bigr)+\mathfrak v\sin\bigl((\mathfrak p-\mathfrak q)\mathsf z\bigr).
\end{equation}

\begin{proposition}[The two-edge model]\label{prop:graph-model}
\begin{enumerate}[label=\textup{(\roman*)}]
\item All zeros of $\mathfrak S$ are real.  Each interval between consecutive points $(m+\frac12)\pi/\mathfrak k$, $m\in\Z$, contains a zero, and the number of zeros in $[0,T]$, with multiplicity, is $\mathfrak kT/\pi+O(1)$; hence the positive zeros form a list $\boldsymbol{\kappa}$ with $|\kappa_k-k\pi/\mathfrak k|$ bounded.  The set of zeros is not periodic.
\item $\mathfrak S(-\mathsf z)=-\mathfrak S(\mathsf z)$, and $\mathfrak D(\mathsf z)=-e^{2i\mathfrak k\mathsf z}\,\overline{\mathfrak D(\bar{\mathsf z})}$.
\item For $\Im\mathsf z>0$,
\begin{equation}\label{eq:graph-euler}
 \log\mathfrak D(\mathsf z)=-\sum_{r\ge1}\frac1r\operatorname{tr}\bigl(\mathsf U\mathsf X(\mathsf z)\bigr)^r=\sum_{\mathfrak l\in\mathscr F}\mathfrak b_{\mathfrak l}\,e^{i\mathfrak l\mathsf z},\qquad\mathscr F=\{2(m\mathfrak p+m'\mathfrak q):m,m'\ge0\}\setminus\{0\},
\end{equation}
absolutely convergent, and the spectral density of the zeros of $\mathfrak S$, counted with multiplicity, is atomic on $\mathscr F$ in the sense of \CGeqref{eq:acceptance-density}.
\item $\mathfrak D(\mathsf z)=-2i\,\mathfrak S'(0)\,e^{i\mathfrak k\mathsf z}\,\mathfrak C_{\boldsymbol{\kappa}}(\mathsf z)$ for the list $\boldsymbol{\kappa}$ of all positive zeros of $\mathfrak S$: the continuous and the discrete functions coincide up to the factor $e^{i\mathfrak k\mathsf z}$, and the list is accepted.
\end{enumerate}
\end{proposition}

\begin{proof}
(i) For $\Im\mathsf z>0$ the matrix $\mathsf U\mathsf X(\mathsf z)$ has norm at most $\max(e^{-2\mathfrak p\Im\mathsf z},e^{-2\mathfrak q\Im\mathsf z})<1$, because $\mathsf U$ is unitary; so $I-\mathsf U\mathsf X(\mathsf z)$ is invertible and $\mathfrak D(\mathsf z)\ne0$.  $\mathfrak S$ is odd, so it has no zeros in $\Im\mathsf z<0$ either.  At $\mathsf z_m=(m+\frac12)\pi/\mathfrak k$, $\sin(\mathfrak k\mathsf z_m)=(-1)^m$ and $|\mathfrak v\sin((\mathfrak p-\mathfrak q)\mathsf z_m)|<1$, so $\mathfrak S(\mathsf z_m)$ has the sign $(-1)^m$ and there is a zero in $(\mathsf z_{m-1},\mathsf z_m)$.  For the count, let $\mathsf M(t):=\mathsf U\mathsf X(t)$, unitary and real-analytic in $t\in\R$, with $\mathsf M'(t)=2i\,\mathsf M(t)\mathsf D$, $\mathsf D:=\operatorname{diag}(\mathfrak p,\mathfrak q)$ (diagonal matrices commute).  $\mathsf M(t)$ is never a scalar matrix, since $\mathsf U$ has nonzero off-diagonal entries; so its two eigenvalues $e^{i\theta_1(t)},e^{i\theta_2(t)}$, which can be labelled as continuous functions of $t$ (roots of a quadratic with continuous coefficients), coincide only on a discrete set (the discriminant $\operatorname{tr}^2-4\det=\mathfrak v^2(e^{2i\mathfrak pt}-e^{2i\mathfrak qt})^2+4e^{2i\mathfrak kt}$ is real-analytic and equals $4$ at $t=0$).  Where an eigenvalue $\mu=e^{i\theta_j}$ is simple it is analytic with an analytic unit eigenvector $v_j$, and differentiating $\mathsf Mv_j=\mu v_j$ and pairing with $v_j$ gives $\mu'=\langle\mathsf M'v_j,v_j\rangle$ (the term $\langle\mathsf Mv_j',v_j\rangle=\langle v_j',\mathsf M^*v_j\rangle=\mu\langle v_j',v_j\rangle$ cancels against $\mu\langle v_j',v_j\rangle$), i.e.
\begin{equation}\label{eq:graph-eigenphase}
 \theta_j'(t)=2\langle\mathsf Dv_j,v_j\rangle=2\bigl(\mathfrak p|v_{j,1}|^2+\mathfrak q|v_{j,2}|^2\bigr)\ \ge\ 2\min(\mathfrak p,\mathfrak q)>0 .
\end{equation}
Hence each continuous eigenphase is strictly increasing with $\theta_j(b)-\theta_j(a)\ge2\min(\mathfrak p,\mathfrak q)(b-a)$, and $\theta_1+\theta_2=2\mathfrak kt+\mathrm{const}$ because $\det\mathsf M(t)=-e^{2i\mathfrak kt}$.  For real $t$, $\mathfrak D(t)=(1-e^{i\theta_1(t)})(1-e^{i\theta_2(t)})$, and $1$ is never a double eigenvalue ($\mathsf M(t)=I$ is impossible), so every real zero of $\mathfrak D$ is a simple zero of exactly one factor, at which that eigenphase crosses $2\pi\Z$ with positive speed: the real zeros of $\mathfrak S$ are simple, and their number in $[0,T]$ is $\sum_j\#\{t\in[0,T]:\theta_j(t)\in2\pi\Z\}=\sum_j\bigl(\frac{\theta_j(T)-\theta_j(0)}{2\pi}+O(1)\bigr)=\frac{\mathfrak kT}\pi+O(1)$.  If the zero set had a period $y\ne0$, then $\mathfrak S(\mathsf z+y)$ and $\mathfrak S(\mathsf z)$ would be entire functions of order one with the same zero multiset, so by Hadamard's factorization $\mathfrak S(\mathsf z+y)=e^{a\mathsf z+b}\mathfrak S(\mathsf z)$ for constants $a,b$. As $\Im\mathsf z\to+\infty$, \CGeqref{eq:graph-sine} and $\mathfrak k>|\mathfrak p-\mathfrak q|$ give $\mathfrak S(\mathsf z)=-\frac1{2i}e^{-i\mathfrak k\mathsf z}(1+o(1))$ uniformly in $\Re\mathsf z$, so $e^{a\mathsf z+b}=e^{-i\mathfrak ky}(1+o(1))$ as $\Im\mathsf z\to+\infty$, which forces $a=0$. With $a=0$, comparing the coefficients of $e^{\pm i(\mathfrak p+\mathfrak q)\mathsf z}$ and of $e^{\pm i(\mathfrak p-\mathfrak q)\mathsf z}$ in $\mathfrak S(\mathsf z+y)=e^b\mathfrak S(\mathsf z)$ gives $e^{\pm i(\mathfrak p+\mathfrak q)y}=e^b=e^{\pm i(\mathfrak p-\mathfrak q)y}$, hence $(\mathfrak p+\mathfrak q)y\in\pi\Z$ and $(\mathfrak p-\mathfrak q)y\in\pi\Z$, so $2\mathfrak py\in\pi\Z$ and $2\mathfrak qy\in\pi\Z$, against $\mathfrak p/\mathfrak q\notin\mathbb Q$.  (ii) is immediate; the second identity follows from $\det\mathsf U=-1$.  (iii) The series is the logarithm of the determinant, convergent because the norm is less than $1$; each trace is a combination of exponentials with frequencies in $\mathscr F$.  The statement on the density is Theorem~\CGref{thm:acceptance}, (ii)$\Rightarrow$(i), applied to (iv).  (iv) $\mathfrak S$ is an odd entire function of exponential type $\mathfrak k$, bounded on the real axis, with real zeros; by Hadamard's factorization and the oddness, $\mathfrak S(\mathsf z)=\mathfrak S'(0)\,\mathfrak C_{\boldsymbol{\kappa}}(\mathsf z)$.
\end{proof}

In this model nothing has to be joined: the continuous function and the discrete function are the same determinant.  The Euler form comes from the expansion of $\det(I-\mathsf U\mathsf X)$ in closed paths, and the real zeros and the symmetry come from the unitarity of $\mathsf U$, through the contraction in $\Im\mathsf z>0$.  The acceptance criterion at this level is the existence of the common object $\mathsf U$.  In finite rank this is general: the zeros of $\mathsf P(e^{i\mathfrak l_1\mathsf z},\dots,e^{i\mathfrak l_r\mathsf z})$, for a polynomial $\mathsf P$ without zeros in the product of the open unit discs and in the product of their exteriors, are real and their spectral density is atomic \cite{KurasovSarnak2020}, and conversely the lists with unit masses whose spectral density is atomic with locally finite support are zero sets of exponential polynomials with real zeros \cite{OlevskiiUlanovskii2020}.

\subsection{A finite Euler product: the constructions of this paper with the final step proved}
\label{subsec:finite-euler-model}

Let
\begin{equation}\label{eq:finite-euler}
 \mathsf F(s):=\bigl(1-2^{-s}\bigr)\bigl(1-3^{-s}\bigr)=1-2^{-s}-3^{-s}+6^{-s}=\sum_{n\mid6}\mu(n)\,n^{-s}.
\end{equation}
It is a Dirichlet series over the integers with an Euler product over the primes $2$ and $3$, and $\mathsf F(-s)=6^{s}\,\mathsf F(s)$, with axis of symmetry $\Re s=0$.  Its zeros are $s=0$, double, and $s=\pm i\mathfrak y_k$, where
\begin{equation}\label{eq:finite-euler-zeros}
 \{\mathfrak y_k\}_{k\ge1}=\Bigl\{\frac{2\pi m}{\log2}\Bigr\}_{m\ge1}\cup\Bigl\{\frac{2\pi m}{\log3}\Bigr\}_{m\ge1}\quad\text{in increasing order},
\end{equation}
all simple, a non-periodic list, and, with $2\sinh(\frac{s\log p}2)=s\log p\prod_{m\ge1}\bigl(1+\frac{s^2\log^2p}{4\pi^2m^2}\bigr)$,
\begin{equation}\label{eq:finite-euler-hadamard}
 \mathsf F(s)=6^{-s/2}\,(\log2\log3)\,s^{2}\prod_{k\ge1}\Bigl(1+\frac{s^2}{\mathfrak y_k^2}\Bigr).
\end{equation}
The constructions of this paper apply as follows.  The count of the ordinates is $\lfloor\frac{T\log2}{2\pi}\rfloor+\lfloor\frac{T\log3}{2\pi}\rfloor=\frac{T\log6}{2\pi}-1+O(1)$ with a remainder of absolute value at most $1$; the midpoint clock is $\frac{\mathfrak x_k\log6}{2\pi}-1=k-\frac12$, that is, $\mathfrak x_k=(2k+1)\pi/\log6$, and $\mathfrak x_k-\mathfrak y_k$ is bounded.  For a list $\boldsymbol{\kappa}$ at bounded distance from the clock,
\begin{equation}\label{eq:finite-euler-list}
 \mathsf F_{\boldsymbol{\kappa}}(s):=6^{-s/2}\,(\log2\log3)\,s^{2}\prod_{k\ge1}\Bigl(1+\frac{s^2}{\kappa_k^2}\Bigr),\qquad\mathsf F_{\boldsymbol{\kappa}}(-s)=6^{s}\,\mathsf F_{\boldsymbol{\kappa}}(s),
\end{equation}
and the family is $\mathsf F_\lambda:=\mathsf F_{\boldsymbol{\kappa}(\lambda)}$ with $\kappa_k(\lambda)=(1-\lambda)\mathfrak y_k+\lambda\mathfrak x_k$, so that $\mathsf F_0=\mathsf F$ and, since $\prod_{k\ge0}\bigl(1+s^2/\mathfrak x_k^2\bigr)=\cosh(\frac{s\log6}2)$,
\begin{equation}\label{eq:finite-euler-clock}
 \mathsf F_1(s)=\frac{\log2\log3}2\cdot\frac{s^{2}\,(1+6^{-s})}{1+s^2\log^26/\pi^2}.
\end{equation}

\begin{theorem}[The final step in the model]\label{thm:finite-euler}
\begin{enumerate}[label=\textup{(\roman*)}]
\item \textup{(Poisson, prime by prime.)}  In $\mathcal D'(\R)$,
\begin{equation}\label{eq:finite-euler-poisson}
 2+2\sum_{k\ge1}\cos(\mathfrak y_ku)=\sum_{p\in\{2,3\}}\ \sum_{m\in\Z}\cos\frac{2\pi mu}{\log p}=\sum_{p\in\{2,3\}}\log p\sum_{m\in\Z}\delta(u-m\log p).
\end{equation}
The spectral density of the list $(\mathfrak y_k)$ is atomic on the points $m\log p$: the continuous function built from the list is the Dirichlet series with Euler product, $\mathsf F_0=\mathsf F$.
\item \textup{(Hamburger's theorem of the model.)}  Let $f$ be entire of finite order with $f(-s)=6^{s}f(s)$ and $f(0)=0$, and let $f(s)=\sum_{n\ge1}a_nn^{-s}$ converge absolutely in some half-plane $\Re s>\nu_{12}$.  Then $f=a_1\mathsf F$.
\item \textup{(The family.)}  $\mathsf F_\lambda$ is an absolutely convergent Dirichlet series in a half-plane if and only if $\lambda=0$.  In particular the clock function $\mathsf F_1$ is not a Dirichlet series.
\end{enumerate}
\end{theorem}

\begin{proof}
(i) is the Poisson summation formula $\sum_{m\in\Z}e^{2\pi imu/\log p}=\log p\sum_{m\in\Z}\delta(u-m\log p)$ for $p=2$ and $p=3$; the second statement is \CGeqref{eq:finite-euler-hadamard}.
(ii) Put $\sigma_1:=\nu_{12}+1$ and $\sigma_2:=-\sigma_1$.  On $\Re s\ge\sigma_1$, $f(s)=\sum_na_nn^{-s}$ converges absolutely and $|f|\le C:=\sum_n|a_n|n^{-\sigma_1}$; on $\Re s\le\sigma_2$, the functional equation gives $f(s)=6^{-s}f(-s)=\sum_na_n(n/6)^{s}$, again absolutely convergent, with $|f(s)|\le C\,6^{-\sigma}$.  By the Phragm\'en--Lindel\"of theorem $f$ is bounded in the closed strip $\mathbb S:=\{\sigma_2\le\Re s\le\sigma_1\}$: $f$ has finite order $\rho$, so $|f(s)|\le\exp(|s|^{\rho+1})\le\exp(e^{c|\Im s|})$ for every $c>0$ and $|s|$ large, which is the growth condition of the theorem in a strip of any finite width; let $\mathrm M:=\sup_{\mathbb S}|f|$.  \emph{Comparison of the two Dirichlet expansions.}  For real $\lambda$, $\sigma\in[\sigma_2,\sigma_1]$ and $T>0$ put $\mathsf A_{\lambda,T}(\sigma):=\frac1{2T}\int_{-T}^Tf(\sigma+it)e^{i\lambda t}\,dt$.  Since $f'(\sigma+it)=-i\,\partial_tf(\sigma+it)$, an integration by parts gives
\begin{equation}\label{eq:finite-euler-coefficients}
 \frac{d}{d\sigma}\mathsf A_{\lambda,T}(\sigma)=\frac{-i}{2T}\Bigl[f(\sigma+it)e^{i\lambda t}\Bigr]_{t=-T}^{t=T}-\lambda\,\mathsf A_{\lambda,T}(\sigma),
 \qquad\text{so}\qquad
 \Bigl|\frac{d}{d\sigma}\bigl(e^{\lambda\sigma}\mathsf A_{\lambda,T}(\sigma)\bigr)\Bigr|\le\frac{\mathrm M}{T}e^{\lambda\sigma},
\end{equation}
and integrating over $[\sigma_2,\sigma_1]$, $|e^{\lambda\sigma_1}\mathsf A_{\lambda,T}(\sigma_1)-e^{\lambda\sigma_2}\mathsf A_{\lambda,T}(\sigma_2)|\le\frac{\mathrm M}T\int_{\sigma_2}^{\sigma_1}e^{\lambda\sigma}d\sigma\to0$ as $T\to\infty$.  On the two boundary lines the limits $\mathsf A_\lambda(\sigma_i):=\lim_{T\to\infty}\mathsf A_{\lambda,T}(\sigma_i)$ exist and are computed termwise from the absolutely convergent expansions (dominated convergence; the mean of $e^{i\mu t}$ is $1$ for $\mu=0$ and $0$ otherwise): $\mathsf A_\lambda(\sigma_1)=a_nn^{-\sigma_1}$ if $\lambda=\log n$ and $0$ otherwise; $\mathsf A_\lambda(\sigma_2)=a_m(m/6)^{\sigma_2}$ if $\lambda=\log(6/m)$ and $0$ otherwise.  Hence
\begin{equation}\label{eq:finite-euler-spectrum}
 e^{\lambda\sigma_1}\mathsf A_\lambda(\sigma_1)=e^{\lambda\sigma_2}\mathsf A_\lambda(\sigma_2)\qquad(\lambda\in\R).
\end{equation}
Take $\lambda=\log n$ with $a_n\ne0$: the left side is $a_n\ne0$, so the right side is nonzero, which forces $\log n=\log(6/m)$ for some $m$ with $a_m\ne0$, i.e.\ $n\mid6$ and $m=6/n$, and then $e^{\lambda\sigma_2}\mathsf A_\lambda(\sigma_2)=a_m(m/6)^{\sigma_2}n^{\sigma_2}=a_m$, so $a_n=a_{6/n}$.  Thus $a_n=0$ unless $n\mid6$, $a_6=a_1$ and $a_3=a_2$, and $f(0)=2a_1+2a_2=0$ gives $a_2=-a_1$ and $f=a_1(1-2^{-s}-3^{-s}+6^{-s})=a_1\mathsf F$.  (This is the elementary case of Bohr's theorem that the Dirichlet exponents of a function almost periodic in a strip are the same on every line of the strip \cite{Bohr1926}; \CGeqref{eq:finite-euler-coefficients} is what allows the coefficient formula to be used on the lines $\Re s<0$.)
(iii) $\mathsf F_\lambda$ is entire of order one, satisfies the functional equation and vanishes at $0$.  If it is a Dirichlet series, (ii) gives $\mathsf F_\lambda=a_1\mathsf F$, the lists of zeros coincide, and $\lambda(\mathfrak x_k-\mathfrak y_k)=0$ for every $k$, so $\lambda=0$.
\end{proof}

In the model every object of the programme exists and the analogue of (D) is proved.  What makes $\lambda=0$ close is a property of the list that the clock and the lists with $\lambda\ne0$ do not have,
\begin{equation}\label{eq:finite-euler-superposition}
 \{\mathfrak y_k\}=\frac{2\pi}{\log2}\,\N\ \cup\ \frac{2\pi}{\log3}\,\N ,
\end{equation}
a superposition of one clock for each prime; it passes the scale test of Remark~\CGref{rem:scale-test} and it is the hypothesis of the Poisson formula \CGeqref{eq:finite-euler-poisson}.  The constructions extend to $\prod_{p\le X}(1-p^{-s})=\sum_{n\mid X\#}\mu(n)n^{-s}$, with $X\#$ the product of the primes up to $X$, the functional equation with the factor $(-1)^{\pi(X)}(X\#)^{s}$, a zero of order $\pi(X)$ at the origin and the list $\bigcup_{p\le X}\frac{2\pi}{\log p}\N$: part (i) of Theorem~\CGref{thm:finite-euler} holds prime by prime, and the clock and the family are defined in the same way. The uniqueness statement (ii) does not extend verbatim: for $\pi(X)=3$ the coefficient conditions $a_{X\#/n}=-a_n$ leave four free coefficients, $f(0)=0$ is then automatic, and the two remaining conditions $f'(0)=f''(0)=0$ leave a two-dimensional space, so the analogue of (ii) is false as stated, and for larger $X$ the defect $2^{\pi(X)-1}-\pi(X)$ only grows; since (iii) was proved through (ii), it is proved here only for the two primes $2,3$.  As $X\to\infty$ these functions tend to $1/\zeta$ in $\Re s>1$ only; their axis of symmetry is $\Re s=0$, not $\Re s=\frac12$, and their lists do not converge to the ordinates of $\zeta$, which are not a superposition of one clock for each prime.

\subsection{The diagonal operator over all primes and the three abscissas}
\label{subsec:three-abscissas}

On $\mathscr L:=\ell^2(\{p\ \text{prime}\})$ let $\mathsf X_{\mathfrak P}(s):=\operatorname{diag}(p^{-s})$.  The closed paths of this model are the primes themselves, without scattering between them, and
\begin{equation}\label{eq:prime-det}
 \det\bigl(I-\mathsf X_{\mathfrak P}(s)\bigr)=\prod_p\bigl(1-p^{-s}\bigr)=\frac1{\zeta(s)}\qquad(\Re s>1).
\end{equation}
$\mathsf X_{\mathfrak P}(s)$ has norm $2^{-\sigma}<1$ for $\sigma>0$, it is of trace class if and only if $\sigma>1$, and Hilbert--Schmidt if and only if $\sigma>\frac12$.  With $\mathfrak P(s):=\sum_pp^{-s}$ for $\Re s>1$, the regularized determinant is the function
\begin{equation}\label{eq:prime-det2}
 D_2(s):=\det\nolimits_2\bigl(I-\mathsf X_{\mathfrak P}(s)\bigr)=\prod_p\bigl(1-p^{-s}\bigr)e^{p^{-s}}=\exp\Bigl(-\sum_{m\ge2}\frac1m\,\mathfrak P(ms)\Bigr)\qquad(\Re s>\tfrac12),
\end{equation}
holomorphic and without zeros in $\Re s>\frac12$ because $I-\mathsf X_{\mathfrak P}(s)$ is invertible there; the last expression converges for $\Re s>\frac12$ since $m\sigma>1$ for $m\ge2$. For $\Re s>1$, where $\mathfrak P$ and the Euler product converge, $D_2(s)=e^{\mathfrak P(s)}/\zeta(s)$; in the strip $\frac12<\Re s\le1$ the primary object is $D_2$, and $e^{\mathfrak P}/\zeta$ is defined there only through it, since $\mathfrak P$ itself is not given by its series.

\begin{proposition}[No common object of finite rank in the window of $\zeta$]\label{prop:no-finite-rank}
Let $X\ge2$.  There is no function $f\not\equiv0$ that satisfies (H0) and the functional equation \CGeqref{eq:Z1-FE} and is an absolutely convergent Dirichlet series in $\Re s>1$ supported on the integers all of whose prime factors are at most $X$.
\end{proposition}

\begin{proof}
By Hamburger's theorem $f=C_H\zeta$, and the coefficients of $\zeta$ do not vanish at the integers with a prime factor larger than $X$; so $C_H=0$.
\end{proof}

In the models of Subsections~\CGref{subsec:graph-model} and~\CGref{subsec:finite-euler-model} three boundaries coincide: the boundary of the domain where the matrix is a strict contraction, the boundary of the domain of convergence of the Euler product, and the axis of symmetry.  The real zeros come from this coincidence.  For \CGeqref{eq:prime-det} they are different:
\begin{equation}\label{eq:three-abscissas}
 \text{contraction: }\sigma>0;\qquad\text{convergence of the determinant: }\sigma>1;\qquad\text{axis of symmetry: }\sigma=\tfrac12 .
\end{equation}
The contraction, together with the trace-class determinant identity \CGeqref{eq:prime-det} valid for $\sigma>1$, proves that $\zeta$ has no zeros in $\sigma>1$; the contraction together with the Hilbert--Schmidt property proves that $D_2=\det_2(I-\mathsf X_{\mathfrak P})$ is holomorphic without zeros in $\sigma>\frac12$, and $D_2=e^{\mathfrak P}/\zeta$ there only in the sense that $D_2$ is the continuation of that quotient from $\sigma>1$. Neither statement reaches the axis, and in the strip $\frac12<\sigma\le1$ the regularized determinant moves the question to the continuation of $\mathfrak P$ (equivalently of $\zeta$).  In the language of this paper, $\log\tilde\zeta_{\Z,\lambda}=\sum_n\tilde c_n(\lambda)\mathsf L_n$ is the determinant \CGeqref{eq:prime-det} with real multiplicities $\tilde c_n(\lambda)$ at all integers, and the strip $\frac12<\sigma\le1$ is where the mean square of the cell masses (Theorem~\CGref{thm:mean-square-cells}) and the continuation of the joining factor (Theorem~\CGref{thm:atomized-continuation}) are used.  By Proposition~\CGref{prop:no-finite-rank} the common object of the window of $\zeta$ has no exact finite-rank representative in the discrete category described there (approximation in a topology is not excluded); a unitary object of infinite rank from which both $\zeta_{Z1}$ and $\tilde\zeta_{\Z,0}$ descend, as $\mathsf U$ does in Proposition~\CGref{prop:graph-model}, is not constructed in this paper.  A model without $\zeta$ with the same separation of the three boundaries is a graph with infinitely many edges of lengths $\mathfrak l_j$ such that $\sum_je^{-\mathfrak l_jy}$ diverges for $y\le y_0$, with $y_0>0$.

%======================================================================
\section{Route B for finite sections: compactness, uniqueness, certificates, and a Kadec--Avdonin criterion with blocks of variable length}
\label{sec:routeB-finite}

This section treats route B for the midpoint RvM system in the notation of this paper. Its subsections are self-contained and use the preceding sections only through the results quoted by label. The numerical part (Section~\CGref{subsec:rb-numerics}) is a computation and is not part of any proof.

\subsection{Introduction to the section}\label{subsec:rb-intro}

This paper constructs a synthetic zeta system on the midpoint clock $\tau_k$, defined by $\NR(\tau_k)=k-\frac12$, and reduces the remaining step to condition (D). Route B approaches (D) without constructing the zeta ordinates. One solves finite sections of the rung system \CGeqref{eq:rb-rung}, which imposes simultaneously a real critical spectrum and atomic arithmetic on $\{\log n\}$. If the residuals of these sections tend to zero in a fixed box, compactness yields an exact solution. Theorem~\CGref{thm:rb-uniqueness} identifies every exact solution whose ordinates lie at bounded distance from the clock, with arbitrary real masses, as the list of zeta ordinates with the von Mangoldt masses, and shows that its existence implies the Riemann hypothesis. In this scheme the ordinates $\gamma_k$ are the output of an inverse problem, not an input. Section~\CGref{subsec:rb-routeB} makes this precise: route B is equivalent to (D) (Corollary~\CGref{cor:rb-equivalence}), and its only open piece is the existence of sections with vanishing residual. Section~\CGref{subsec:rb-prescribed} gives the same equivalence with the von Mangoldt masses prescribed: the Laplace transform of the residual is identified (Lemma~\CGref{lem:rb-laplace}), and uniqueness then follows from Hadamard's theorem and the functional equation alone, so that route B$'$ is equivalent to the Riemann hypothesis. Section~\CGref{subsec:rb-nonfinitary} shows that no finite stage decides route B: because the Riemann hypothesis is verified up to height $3\cdot10^{12}$, the zeta list with the von Mangoldt masses solves every finite set of tests up to an explicit error of order $(3\cdot10^{12})^{1-r}$, whether or not the hypothesis holds. Section~\CGref{subsec:rb-staged} gives the existence part a certificate form: two computable quantities per stage, a right inverse of the Jacobian and the residual of the new tests, imply an exact solution of the stage. Section~\CGref{subsec:rb-gonek} shows that Gonek's symmetrised Euler products give explicit lists in a tube around the clock. Section~\CGref{subsec:rb-numerics} reports the computations that measure these quantities, and proves in ball arithmetic, with compactly supported tests, that the square sections of fifteen stages, up to height $487$, have exact solutions (Theorem~\CGref{thm:rb-verified}).

Section~\CGref{subsec:rb-apriori} proves an a priori bound for quasi-solutions. Section~\CGref{subsec:rb-riesz} proves a Kadec--Avdonin criterion with blocks of variable length, which is the tool needed for a Fredholm/stability formulation of the linearised problem. Section~\CGref{subsec:rb-application} applies that criterion to lists near the clock and proves that, for the zeta list, the uniform version fails.

References of the form Theorem~\CGref{thm:free-masses} and the other results of the preceding sections are used as stated there. Except in Proposition~\CGref{prop:rb-nonfinitary} and Section~\CGref{subsec:rb-obstruction}, where the zeta list is the object of study, and in the diagnostic columns of Section~\CGref{subsec:rb-numerics}, nothing below uses the ordinates of $\zeta$.

%======================================================================
\subsection{Setting}\label{subsec:rb-setting}

Throughout, $q(s):=s(s-1)$ and $G(s):=s(s-1)\pi^{-s/2}\Gamma(s/2)$. The smooth clock count is
\begin{equation}\label{eq:rb-NRVM}
 \NR(t):=\frac t{2\pi}\log\frac t{2\pi e}+\frac78,\qquad
 \NR'(t)=\frac1{2\pi}\log\frac t{2\pi}\qquad(t>0).
\end{equation}
It is strictly increasing on $[2\pi,\infty)$, with $\NR(2\pi)=-\frac18$, and the clock points are the solutions $\tau_k>2\pi$ of $\NR(\tau_k)=k-\frac12$, $k\ge1$ ($\tau_1=14.5213\ldots$). We write $n_{\RVM}(t):=\#\{k\ge1:\tau_k\le t\}$. Since $\NR$ is increasing,
\begin{equation}\label{eq:rb-Ntau-floor}
 n_{\RVM}(t)=\Bigl\lfloor\NR(t)+\tfrac12\Bigr\rfloor,\qquad
 \bigl|n_{\RVM}(t)-\NR(t)\bigr|\le\tfrac12\qquad(t\ge2\pi).
\end{equation}

\begin{definition}[Admissible lists]\label{def:rb-lists}
For $C>0$, a list $\boldsymbol\kappa=(\kappa_k)_{k\ge1}$ is \emph{$C$-admissible} if it is real and nondecreasing, $\kappa_1\ge2\pi$, and $|\kappa_k-\tau_k|\le C$ for all $k$. Put $N_{\boldsymbol\kappa}(t):=\#\{k:\kappa_k\le t\}$ (with multiplicity) and
\begin{equation}\label{eq:rb-Pk}
 P_{\boldsymbol\kappa}(s):=\prod_{k\ge1}\Bigl(1+\frac{q(s)}{\frac14+\kappa_k^2}\Bigr),\qquad
 \zeta_{\boldsymbol\kappa}:=\frac{P_{\boldsymbol\kappa}}{G}.
\end{equation}
The clock list $\boldsymbol\tau=(\tau_k)$ gives $\zeta_{\boldsymbol\tau}=\zeta_{\RVM}$.
\end{definition}

$P_{\boldsymbol\kappa}$ is entire and real on $\mathbb R$, and $P_{\boldsymbol\kappa}(1-s)=P_{\boldsymbol\kappa}(s)$. Its zeros are exactly $\frac12\pm i\kappa_k$, with multiplicity, and it is positive on $[-1,2]$, because $q([-1,2])=[-\frac14,2]$, so every factor satisfies $1+q(s)/(\frac14+\kappa_k^2)\ge1-\frac{1/4}{\frac14+\kappa_k^2}>0$ there, and the product converges. Writing $s=\frac12+w$, one has $\frac14+\kappa^2+q(s)=\kappa^2+w^2$.

The \emph{rung system} \CGref{eq:rung-system} for a list $\boldsymbol\kappa$ and real masses $\mathsf a=(\mathsf a_n)_{n\ge2}$ is
\begin{equation}\label{eq:rb-rung}
 \sum_{k\ge1}\cos(\kappa_ku)+\sum_{n\ge2}\mathsf a_n\,\delta(u-\log n)=\mathsf k(u)\qquad\text{in }\mathcal D'(0,\infty),
\end{equation}
with the fixed distribution $\mathsf k$ of \CGref{eq:transport-defs}. Its residual is the distribution $\mathfrak R(\boldsymbol\kappa,\mathsf a)$ defined, for $\varphi\in C_c^\infty(0,\infty)$, by
\begin{equation}\label{eq:rb-residual}
 \langle\mathfrak R(\boldsymbol\kappa,\mathsf a),\varphi\rangle:=\sum_{k\ge1}\mathsf c_\varphi(\kappa_k)+\sum_{n\ge2}\mathsf a_n\varphi(\log n)-\langle\mathsf k,\varphi\rangle,\qquad
 \mathsf c_\varphi(t):=\int_0^\infty\varphi(u)\cos(tu)\,du .
\end{equation}
Two integrations by parts, without boundary terms because $\varphi$ and $\varphi'$ vanish at the ends of its support, give $|\mathsf c_\varphi(t)|\le\|\varphi''\|_{L^1}\,t^{-2}$. Since $\kappa_k\ge\tau_k-C$ and $\tau_k\gg k/\log(k+1)$, the series in \CGeqref{eq:rb-residual} converges absolutely. The mass sum is finite, because $\operatorname{supp}\varphi\subset[u_1,u^*]$ with $0<u_1<u^*$ leaves only $n\le e^{u^*}$.

%======================================================================
\subsection{Route B: compactness, selection and the canonical box}\label{subsec:rb-routeB}

For $C,\mathbb M>0$ let $\mathbb K(C,\mathbb M)$ be the set of pairs $(\boldsymbol\kappa,\mathsf a)$ with $\boldsymbol\kappa$ $C$-admissible and $|\mathsf a_n|\le \mathbb M$ for all $n$. A \emph{finite section} is an element of $\mathbb K(C,\mathbb M)$ with $\kappa_k=\tau_k$ for $k>k^*$ and $\mathsf a_n=0$ for $n>n^*$, for some $k^*,n^*$.

Fix finite sets $\mathscr T_1\subset\mathscr T_2\subset\cdots\subset C_c^\infty(0,\infty)$ whose union $\mathscr T_\infty$ is dense in $C_c^\infty(0,\infty)$. Such sets exist because $C_c^\infty(0,\infty)$ is separable for its usual LF topology: it is the countable strict inductive limit of the separable Fr\'echet spaces $C^\infty_{[1/j,j]}$, and a countable union of countable dense subsets of these is dense. For $(\boldsymbol\kappa,\mathsf a)\in\mathbb K(C,\mathbb M)$ put
\begin{equation}\label{eq:rb-rho}
 \mathsf s_m(\boldsymbol\kappa,\mathsf a):=\max_{\varphi\in\mathscr T_m}\bigl|\langle\mathfrak R(\boldsymbol\kappa,\mathsf a),\varphi\rangle\bigr| .
\end{equation}

\begin{definition}[Statement B]\label{def:rb-B}
\emph{Statement B$(C,\mathbb M)$} says that there are finite sections $(\boldsymbol\kappa^{(m)},\mathsf a^{(m)})\in\mathbb K(C,\mathbb M)$ with $\mathsf s_m(\boldsymbol\kappa^{(m)},\mathsf a^{(m)})\to0$.
\end{definition}

\begin{lemma}[Canonical box of the exact solution]\label{lem:rb-canonical-box}
Let $\mathbb M_*:=\max_{n\ge2}\Lambda(n)/(2\sqrt n)$. Then $\mathbb M_*=\log7/(2\sqrt7)=0.367742\ldots<1/e$. If condition (D) holds, the pair $(\boldsymbol\gamma,\mathsf a^\zeta)$, with $\mathsf a^\zeta_n:=\Lambda(n)/(2\sqrt n)$, solves \CGeqref{eq:rb-rung} and lies in $\mathbb K(C_\eta,\mathbb M_*)$, where $C_\eta$ is the constant of \CGref{lem:bounded-displacement}.
\end{lemma}

\begin{proof}
$x\mapsto\log x/(2\sqrt x)$ increases on $[1,e^2]$ and decreases on $[e^2,\infty)$, with $e^2=7.389\ldots$. Hence among primes its maximum is at $p=7$ or $p=11$; the values are $0.367742\ldots$ and $0.361500\ldots$. A prime power $p^j$ with $j\ge2$ gives $\log p/(2p^{j/2})<\log p/(2\sqrt p)\le \mathbb M_*$. The bound by $1/e$ is the continuous maximum, attained at $x=e^2$. Under (D), $\zeta_{Z1}=\zeta$, so all zeros of $\zeta$ are $\frac12\pm i\gamma_k$. The rung system with list $\boldsymbol\gamma$ and masses $\mathsf a^\zeta$ is then the distributional explicit formula \CGref{thm:distributional-EF}. Finally, $|\gamma_k-\tau_k|\le C_\eta$ by \CGref{lem:bounded-displacement}, and $\gamma_1=14.13\ldots\ge2\pi$.
\end{proof}

\begin{corollary}[(D) implies B]\label{cor:rb-D-implies-B}
If (D) holds, then B$(C,\mathbb M)$ holds for every $C\ge C_\eta$ and $\mathbb M\ge \mathbb M_*$.
\end{corollary}

\begin{proof}
Let $(\boldsymbol\kappa^{(k^*,n^*)},\mathsf a^{(k^*,n^*)})$ be the truncation of $(\boldsymbol\gamma,\mathsf a^\zeta)$. Its list is the nondecreasing rearrangement of the multiset $\{\gamma_1,\dots,\gamma_{k^*}\}\cup\{\tau_{k^*+1},\tau_{k^*+2},\dots\}$, and its masses are $\mathsf a^\zeta_n$ for $n\le n^*$ and $0$ beyond. By Lemma~\CGref{lem:rb-matching} (Section~\CGref{subsec:rb-staged}) the rearranged list satisfies $|\kappa_k-\tau_k|\le C_\eta$, coincides with $(\tau_k)$ beyond some index, and $\mathfrak R$ depends only on the multiset. So this is a finite section in $\mathbb K(C_\eta,\mathbb M_*)$. For fixed $\varphi$, $\langle\mathfrak R(\boldsymbol\kappa^{(k^*,n^*)},\mathsf a^{(k^*,n^*)}),\varphi\rangle-\langle\mathfrak R(\boldsymbol\gamma,\mathsf a^\zeta),\varphi\rangle=\sum_{k>k^*}(\mathsf c_\varphi(\tau_k)-\mathsf c_\varphi(\gamma_k))$ as soon as $e^{u^*}\le n^*$. This tends to $0$ as $k^*\to\infty$. Since $\mathscr T_m$ is finite and $\mathfrak R(\boldsymbol\gamma,\mathsf a^\zeta)=0$ by Lemma~\CGref{lem:rb-canonical-box}, $k^*_m,n^*_m$ can be chosen with $\mathsf s_m\le1/m$.
\end{proof}

\begin{proposition}[Compactness and passage to the limit]\label{prop:rb-compactness}
Let $(\boldsymbol\kappa^{(m)},\mathsf a^{(m)})\in\mathbb K(C,\mathbb M)$ with $\mathsf s_m(\boldsymbol\kappa^{(m)},\mathsf a^{(m)})\to0$. Then a subsequence converges coordinatewise to some $(\boldsymbol\kappa,\mathsf a)\in\mathbb K(C,\mathbb M)$, and every such limit satisfies $\mathfrak R(\boldsymbol\kappa,\mathsf a)=0$ in $\mathcal D'(0,\infty)$.
\end{proposition}

\begin{proof}
Each $\kappa^{(m)}_k$ lies in the compact interval $[\tau_k-C,\tau_k+C]$ and each $\mathsf a^{(m)}_n$ in $[-\mathbb M,\mathbb M]$. A diagonal extraction over the countably many coordinates gives a coordinatewise convergent subsequence. The limit is nondecreasing, satisfies $\kappa_1\ge2\pi$, $|\kappa_k-\tau_k|\le C$ and $|\mathsf a_n|\le \mathbb M$, and therefore lies in $\mathbb K(C,\mathbb M)$.

\emph{Continuity of each pairing.} Fix $\varphi$ with support in $[u_1,u^*]$. The mass part of \CGeqref{eq:rb-residual} involves only $n\le e^{u^*}$ and converges. For the spectral part, $|\mathsf c_\varphi(\kappa^{(m)}_k)|\le\|\varphi''\|_1(\tau_k-C)^{-2}$ when $\tau_k>2C$, and $|\mathsf c_\varphi|\le\|\varphi\|_1$ for the finitely many other $k$. The majorant is summable, so dominated convergence for series gives $\sum_k\mathsf c_\varphi(\kappa_k^{(m)})\to\sum_k\mathsf c_\varphi(\kappa_k)$. Hence $\langle\mathfrak R(\boldsymbol\kappa^{(m)},\mathsf a^{(m)}),\varphi\rangle\to\langle\mathfrak R(\boldsymbol\kappa,\mathsf a),\varphi\rangle$ along the subsequence.

\emph{Vanishing.} If $\varphi\in\mathscr T_\infty$, then $\varphi\in\mathscr T_m$ for all large $m$, so $|\langle\mathfrak R(\boldsymbol\kappa^{(m)},\mathsf a^{(m)}),\varphi\rangle|\le\mathsf s_m\to0$ and $\langle\mathfrak R(\boldsymbol\kappa,\mathsf a),\varphi\rangle=0$. The limit $\mathfrak R(\boldsymbol\kappa,\mathsf a)$ is a single distribution, hence continuous on $C_c^\infty(0,\infty)$. It vanishes on the dense set $\mathscr T_\infty$, so it vanishes identically. No uniform bound on the residuals is needed for this step.
\end{proof}

Theorem~\CGref{thm:free-masses} proves uniqueness for lists in $\mathscr O$. Besides the tube and the ordering, $\mathscr O$ (\CGref{eq:list-space}) carries a third condition, $|\int(N_{\boldsymbol\kappa}-\NR)|\le C\log \mathrm X$, and a coordinatewise limit of sections does not inherit it. The following lemma and theorem prove uniqueness for exact solutions in every tube, without the third condition and with no bound on the masses.

\begin{lemma}[Mean square of the masses of an exact solution]\label{lem:rb-mass-ms}
Let $\boldsymbol\kappa$ be $C$-admissible for some $C>0$, and let $\mathsf a=(\mathsf a_n)_{n\ge2}$ be real numbers, with no bound assumed, such that $\mathfrak R(\boldsymbol\kappa,\mathsf a)=0$ in $\mathcal D'(0,\infty)$. Then for every $\epsilon>0$ there is $C_\epsilon$, depending only on $\epsilon$ and $C$, such that
\begin{equation}\label{eq:rb-mass-ms}
 \sum_{X<n\le2X}\mathsf a_n^2\le C_\epsilon X^{\epsilon}\qquad(X\ge2).
\end{equation}
\end{lemma}

\begin{proof}
Fix an even real $\psi\in C_c^\infty(-\frac15,\frac15)$ with $\psi(0)=1$, and put $\widehat\psi(\xi):=\int\psi(v)\cos(\xi v)\,dv$; then $|\widehat\psi(\xi)|\le\|\psi\|_1$, and $A$ integrations by parts give $|\widehat\psi(\xi)|\le\|\psi^{(A)}\|_1|\xi|^{-A}$ for every integer $A\ge0$. Fix $X\ge2$ and $\delta\in(0,1]$, and put $T:=X^{1+\delta}$. All implied constants below depend only on $\psi$, $\delta$ and $C$.

\emph{One test for each mass.} For $X<n\le2X$ let $\varphi_n(u):=\psi(X(u-\log n))$, supported in $|u-\log n|<1/(5X)$. For an integer $m\ge2$, $m\ne n$, one has $|\log m-\log n|\ge\min\bigl(\log\frac{n+1}n,\log\frac n{n-1}\bigr)\ge\frac1{n+1}\ge\frac1{2X+1}\ge\frac1{5X}$, so $\varphi_n(\log m)=0$, while $\varphi_n(\log n)=1$. Since $\psi$ is even, the substitution $u=\log n+v/X$ gives $\mathsf c_{\varphi_n}(t)=X^{-1}\widehat\psi(t/X)\cos(t\log n)$. Pairing $\mathfrak R(\boldsymbol\kappa,\mathsf a)=0$ with $\varphi_n$,
\begin{equation}\label{eq:rb-mass-isolated}
\begin{aligned}
 &\mathsf a_n=\langle\mathsf k,\varphi_n\rangle-\Re\mathcal S(\log n)-\mathcal E_n,\qquad \mathrm y_k:=\frac{\widehat\psi(\kappa_k/X)}X,\\
 &\mathcal S(u):=\sum_{\kappa_k\le T}\mathrm y_ke^{i\kappa_ku},\qquad
 \mathcal E_n:=\sum_{\kappa_k>T}\mathrm y_k\cos(\kappa_k\log n).
\end{aligned}
\end{equation}

\emph{The term $\mathsf k$.} On the support of $\varphi_n$, $u\ge\log2-\frac1{10}>\frac12$, where $|\mathsf k(u)|\le\cosh\frac u2+\frac{e^{-u/2}}{2(1-e^{-2u})}\le2e^{u/2}$; hence $|\langle\mathsf k,\varphi_n\rangle|\le2e^{1/10}\sqrt n\,\|\psi\|_1X^{-1}\ll X^{-1/2}$.

\emph{Large ordinates.} By the tube, $N_{\boldsymbol\kappa}(t)\le n_{\RVM}(t+C)\ll t\log t$ for $t\ge2$. Let $A:=\lceil2/\delta\rceil+1$, so that $\delta(A-1)\ge2$. Partial summation gives $\sum_{\kappa_k>T}\kappa_k^{-A}\le A\int_T^\infty N_{\boldsymbol\kappa}(t)t^{-A-1}dt\ll T^{1-A}\log T$, hence
\[
 |\mathcal E_n|\le\|\psi^{(A)}\|_1X^{A-1}\sum_{\kappa_k>T}\kappa_k^{-A}\ll X^{A-1}T^{1-A}\log T=X^{-\delta(A-1)}\log T\ll X^{-2}\log X .
\]

\emph{Small ordinates.} The intervals $I_n:=[\log n-\frac1{6X},\log n+\frac1{6X}]$, $X<n\le2X$, are pairwise disjoint (consecutive points $\log n$ differ by at least $\frac1{2X+1}>\frac1{3X}$) and contained in $\mathcal W:=[\log X-1,\log X+2]$. For $g\in C^1(I)$ on an interval $I$ of length $\ell$ and $t_0\in I$, $|g(t_0)|^2\le|g(t)|^2+2\int_I|g||g'|$ for every $t\in I$; averaging over $t\in I$,
\[
 |g(t_0)|^2\le\ell^{-1}\int_I|g|^2+2\int_I|g||g'| .
\]
With $g=\mathcal S$, $I=I_n$, summing over $n$ and using the Cauchy--Schwarz inequality,
\[
 \sum_{X<n\le2X}|\mathcal S(\log n)|^2\le3X\int_{\mathcal W}|\mathcal S|^2+2\Bigl(\int_{\mathcal W}|\mathcal S|^2\Bigr)^{1/2}\Bigl(\int_{\mathcal W}|\mathcal S'|^2\Bigr)^{1/2}.
\]
The proof of Lemma~\CGref{lem:mean-value-ordinates} uses only the tube (Theorem~\CGref{thm:transfer-O}(v)); with $C$ in place of $C_\eta$ it gives $\int_{u_0-1}^{u_0+1}|\sum_{\kappa_k\le T}\mathrm b_ke^{i\kappa_ku}|^2du\ll T\log^2T$ for complex $|\mathrm b_k|\le1$. The interval $\mathcal W$ is covered by two intervals of length $2$. Applied with $\mathrm b_k=X\mathrm y_k/\|\psi\|_1$, and with $\mathrm b_k=X\kappa_k\mathrm y_k/(T\|\psi\|_1)$ for $\mathcal S'(u)=\sum_{\kappa_k\le T}i\kappa_k\mathrm y_ke^{i\kappa_ku}$, this gives $\int_{\mathcal W}|\mathcal S|^2\ll X^{-2}T\log^2T$ and $\int_{\mathcal W}|\mathcal S'|^2\ll X^{-2}T^3\log^2T$. Hence
\[
 \sum_{X<n\le2X}|\mathcal S(\log n)|^2\ll X^{-1}T\log^2T+X^{-2}T^2\log^2T\ll X^{2\delta}\log^2X .
\]

\emph{Conclusion.} By \CGeqref{eq:rb-mass-isolated}, $\mathsf a_n^2\le3\bigl(\langle\mathsf k,\varphi_n\rangle^2+|\mathcal S(\log n)|^2+\mathcal E_n^2\bigr)$, and summing over $X<n\le2X$ gives $\sum_{X<n\le2X}\mathsf a_n^2\ll1+X^{2\delta}\log^2X+X^{-3}\log^2X\ll X^{3\delta}$. Take $\delta=\epsilon/3$.
\end{proof}

\begin{theorem}[Uniqueness in the tube, with free masses]\label{thm:rb-uniqueness}
Let $\boldsymbol\kappa$ be $C$-admissible for some $C>0$, and let $\mathsf a=(\mathsf a_n)_{n\ge2}$ be real numbers, with no bound assumed. If $\mathfrak R(\boldsymbol\kappa,\mathsf a)=0$ in $\mathcal D'(0,\infty)$, then all nontrivial zeros of $\zeta$ lie on $\Re s=\frac12$, $\boldsymbol\kappa=\boldsymbol\gamma$, $\mathsf a=\mathsf a^\zeta$, and (D) holds.
\end{theorem}

\begin{proof}
\emph{Step 1: no drift, and an Euler form.} We apply Proposition~\CGref{prop:rung-apriori} with $\mathsf t=1$, for which \CGeqref{eq:rung-apriori-system} is the equation $\mathfrak R(\boldsymbol\kappa,\mathsf a)=0$. Its hypotheses hold: $0<2\pi\le\kappa_1\le\kappa_2\le\cdots$ and $\kappa_k\ge\tau_k-C\to\infty$; the tube gives $N_{\boldsymbol\kappa}(T+1)-N_{\boldsymbol\kappa}(T)\le n_{\RVM}(T+1+C)-n_{\RVM}(T-C)\ll_C\log(2+T)$, which is \CGeqref{eq:rung-apriori-density}; and Lemma~\CGref{lem:rb-mass-ms} gives \CGeqref{eq:rung-apriori-mass}. Hence the constant $\mathrm b$ of \CGeqref{eq:rung-apriori-b} exists and $\kappa_k-\tau_k=-(4\mathrm b+o(1))\kappa_k/\log\kappa_k$. Since $|\kappa_k-\tau_k|\le C$, $\mathrm b=0$, and \CGeqref{eq:rung-three-factor} with $\mathsf t=1$ and $\mathrm b=0$ reads
\begin{equation}\label{eq:rb-uniq-factor}
 \zeta_{\boldsymbol\kappa}(s)=\frac{P_{\boldsymbol\kappa}(s)}{G(s)}=C_1\exp\Bigl(\sum_{n\ge2}\mathsf c_nn^{-s}\Bigr),\qquad\mathsf c_n:=\frac{2\sqrt n\,\mathsf a_n}{\log n}\qquad(\Re s>\tfrac32),
\end{equation}
with a constant $C_1>0$.

\emph{Step 2: abscissa $1$.} Let $\sigma>1$ and $\epsilon:=\sigma-1$. By the Cauchy--Schwarz inequality and Lemma~\CGref{lem:rb-mass-ms} with $\epsilon/2$ in place of $\epsilon$,
\[
 \sum_{X<n\le2X}|\mathsf c_n|n^{-\sigma}\le\Bigl(\sum_{X<n\le2X}\mathsf a_n^2\Bigr)^{1/2}\Bigl(\sum_{X<n\le2X}\frac{4n^{1-2\sigma}}{\log^2n}\Bigr)^{1/2}\ll_{\sigma,C}X^{\epsilon/4}X^{1-\sigma}=X^{-3\epsilon/4},
\]
which is summable over $X=2^j$. So $\mathsf D(s):=\sum_{n\ge2}\mathsf c_nn^{-s}$ converges absolutely on $\Re s>1$, and so does the Dirichlet series $\exp(\mathsf D(s))=\sum_{j\ge0}\mathsf D(s)^j/j!=\sum_{n\ge1}\mathsf e_nn^{-s}$, since $\sum_n|\mathsf e_n|n^{-\sigma}\le\exp\bigl(\sum_n|\mathsf c_n|n^{-\sigma}\bigr)$. The function $\zeta_{\boldsymbol\kappa}=P_{\boldsymbol\kappa}/G$ is holomorphic on $\Re s>1$, where $G$ has no zeros; by analytic continuation \CGeqref{eq:rb-uniq-factor} holds on $\Re s>1$, and $\zeta_{\boldsymbol\kappa}$ satisfies (H1) of Subsection~\CGref{subsec:hamburger-table}.

\emph{Step 3: Hamburger's theorem.} (H0): $(s-1)\zeta_{\boldsymbol\kappa}(s)=P_{\boldsymbol\kappa}(s)\,\pi^{s/2}/(2\Gamma(1+\frac s2))$ is entire of finite order, because $P_{\boldsymbol\kappa}$ has order at most one (first step of the proof of Proposition~\CGref{prop:rung-apriori}) and $\pi^{s/2}$ and $1/\Gamma$ have order one. (H2), with $g=f=\zeta_{\boldsymbol\kappa}$: $\zeta_{\boldsymbol\kappa}(s)\Gamma(\frac s2)\pi^{-s/2}=P_{\boldsymbol\kappa}(s)/(s(s-1))$ is invariant under $s\mapsto1-s$. (H3) follows from (H1) for $g=f$. Hamburger's theorem gives $\zeta_{\boldsymbol\kappa}=C_H\zeta$. Since $P_{\boldsymbol\kappa}(1)=1$ and $G(s)=(s-1)(1+O(s-1))$ (because $\pi^{-1/2}\Gamma(\frac12)=1$), $\zeta_{\boldsymbol\kappa}$ has residue $1$ at $s=1$, as $\zeta$ has, and $C_H=1$.

\emph{Step 4: identification.} The nontrivial zeros of $\zeta=P_{\boldsymbol\kappa}/G$ are the zeros of $P_{\boldsymbol\kappa}$, namely $\frac12\pm i\kappa_k$ with multiplicity. So they all lie on $\Re s=\frac12$, and the nondecreasing list of their positive ordinates, with multiplicity, is $\boldsymbol\kappa$: $\boldsymbol\kappa=\boldsymbol\gamma$, and $\zeta_{Z1}=\zeta_{\boldsymbol\gamma}=\zeta$, which is (D). Letting $\sigma\to+\infty$ in \CGeqref{eq:rb-uniq-factor} gives $C_1=1$. Then $\mathsf D$ and $\log\zeta=\sum_n\frac{\Lambda(n)}{\log n}n^{-s}$ are continuous real logarithms of $\zeta$ on $(1,\infty)$ that tend to $0$ at $+\infty$, so they coincide there, and the uniqueness of the coefficients of an absolutely convergent Dirichlet series gives $\mathsf c_n=\Lambda(n)/\log n$, that is, $\mathsf a_n=\Lambda(n)/(2\sqrt n)=\mathsf a^\zeta_n$.
\end{proof}

\begin{remark}[What Theorem~\CGref{thm:rb-uniqueness} uses]\label{rem:rb-uniqueness-input}
Theorem~\CGref{thm:rb-uniqueness} contains the uniqueness half of Theorem~\CGref{thm:free-masses}, for every list at bounded distance from the clock and with no assumption on the masses: neither the third condition of \CGeqref{eq:list-space} nor the threshold $\kappa_1\ge\tau_0^+$ is used. In the proof of Theorem~\CGref{thm:free-masses} the third condition enters through Theorem~\CGref{thm:transfer-O}(ii), (iii) and (vi): the bounds of $\log\zeta_{\boldsymbol\kappa}$ on vertical lines, the Perron argument and the mean square of the cell masses. For an exact solution in a tube these inputs are replaced by the Laplace identity \CGeqref{eq:rung-laplace-identity}, which such a solution satisfies (Proposition~\CGref{prop:rung-apriori}, whose mass hypothesis is supplied by Lemma~\CGref{lem:rb-mass-ms}), and by Lemma~\CGref{lem:rb-mass-ms} itself, in which each mass is isolated by one test of width $1/X$ and the remaining exponential sum over the ordinates is controlled by the local mean value of Lemma~\CGref{lem:mean-value-ordinates}. The tube enters twice: it gives the local density used in Lemma~\CGref{lem:rb-mass-ms}, and it forces $\mathrm b=0$. Proposition~\CGref{prop:rung-apriori} alone does not exclude an exact solution with a linear drift, $\kappa_k-\tau_k\sim-4\mathrm b\kappa_k/\log\kappa_k$ (whether such solutions exist is not known); the tube excludes it. The last two steps are the classical route through Hamburger's theorem; what is specific to the rung system is that exactness in $\mathcal D'(0,\infty)$ supplies both the Euler form \CGeqref{eq:rb-uniq-factor} and its abscissa.
\end{remark}

\begin{corollary}[Uniqueness in the box]\label{cor:rb-U}
For all $C,\mathbb M>0$, the only possible zero of $\mathfrak R$ in $\mathbb K(C,\mathbb M)$ is $(\boldsymbol\gamma,\mathsf a^\zeta)$, and $\mathfrak R$ has a zero in $\mathbb K(C,\mathbb M)$ only if the Riemann hypothesis holds.
\end{corollary}

\begin{proof}
Theorem~\CGref{thm:rb-uniqueness}, since $|\mathsf a_n|\le\mathbb M$ is not used there.
\end{proof}

\begin{corollary}[Selection of the zeta list]\label{cor:rb-selection}
Assume B$(C,\mathbb M)$, and let $(\boldsymbol\kappa^{(m)},\mathsf a^{(m)})$ be any sequence in $\mathbb K(C,\mathbb M)$ with $\mathsf s_m\to0$. Then the whole sequence converges coordinatewise to $(\boldsymbol\gamma,\mathsf a^\zeta)$, and (D) and the Riemann hypothesis hold.
\end{corollary}

\begin{proof}
By Proposition~\CGref{prop:rb-compactness}, every subsequence has a further subsequence converging to a zero of $\mathfrak R$ in $\mathbb K(C,\mathbb M)$. By Theorem~\CGref{thm:rb-uniqueness} that zero is $(\boldsymbol\gamma,\mathsf a^\zeta)$, and (D) and the Riemann hypothesis hold. The set $\mathbb K(C,\mathbb M)$ is compact and metrizable in the product topology, and in a compact metric space a sequence all of whose convergent subsequences have the same limit converges to it; so the whole sequence converges to $(\boldsymbol\gamma,\mathsf a^\zeta)$. The last step can also be read directly from the limit: $\mathfrak R(\boldsymbol\gamma,\mathsf a^\zeta)=0$ implies (D), as follows. By Lemma~\CGref{lem:rb-laplace} (Section~\CGref{subsec:rb-prescribed}), $\widehat{\mathsf R}$ of \CGeqref{eq:rb-Rhat} is the Laplace transform of a distribution of finite order on $[0,\infty)$ whose restriction to $(0,\infty)$ is $\mathfrak R$. If $\mathfrak R=0$, that distribution is supported at $\{0\}$, so $\widehat{\mathsf R}$ is a polynomial $p(s)$. Since $\sum_n\mathsf a^\zeta_nn^{\frac12-s}=-\frac12\zeta'/\zeta(s)$, this reads $P_{\boldsymbol\gamma}'/P_{\boldsymbol\gamma}-(G\zeta)'/(G\zeta)=2p$, so $P_{\boldsymbol\gamma}=c\,G\zeta\,e^{Q}$ with $Q'=2p$. Here $G\zeta=2\xi$ and $P_{\boldsymbol\gamma}$ are entire of order $1$, so $e^Q=P_{\boldsymbol\gamma}/(2c\,\xi)$ is entire, zero-free and of order at most $1$, and $\deg Q\le1$. Both $P_{\boldsymbol\gamma}$ and $\xi$ are invariant under $s\mapsto1-s$, so $e^{Q(1-s)}=e^{Q(s)}$, which forces $Q$ to be constant. Hence $\zeta_{\boldsymbol\gamma}=c'\zeta$. Both functions have residue $1$ at $s=1$ ($P_{\boldsymbol\gamma}(1)=1$ and $G(s)\sim s-1$), so $c'=1$. This is (D), because $\zeta_{\boldsymbol\gamma}=\zeta_{Z1}$.
\end{proof}

\begin{corollary}[Route B is equivalent to (D)]\label{cor:rb-equivalence}
(D) holds if and only if B$(C,\mathbb M)$ holds for some $C,\mathbb M>0$; and then B$(C,\mathbb M)$ holds for all $C\ge C_\eta$ and $\mathbb M\ge\mathbb M_*$.
\end{corollary}

\begin{proof}
Corollaries~\CGref{cor:rb-selection} and~\CGref{cor:rb-D-implies-B}.
\end{proof}

\begin{remark}[Logical status of route B]\label{rem:rb-B-status}
By Corollary~\CGref{cor:rb-equivalence}, route B with free masses is an equivalence, B$(C,\mathbb M)\Leftrightarrow$(D) for $C\ge C_\eta$ and $\mathbb M\ge\mathbb M_*$, with a single open piece, the existence statement B, which is therefore not weaker than (D). Compactness is not an issue in the box: Proposition~\CGref{prop:rb-compactness} uses only the imposed bounds; and uniqueness holds in every tube, with no bound on the masses (Theorem~\CGref{thm:rb-uniqueness}). Proposition~\CGref{prop:rb-nonfinitary} shows that no finite stage of B decides it. In this scheme the ordinates $\gamma_k$ never enter the construction; they appear only as the unique admissible zero of $\mathfrak R$. Lemma~\CGref{lem:rb-canonical-box} does \emph{not} show that near-solutions of sections lie in $\mathbb K(C_\eta,\mathbb M_*)$: it locates only the exact solution. Restricting the sections to that box is a choice, justified a posteriori by Corollary~\CGref{cor:rb-D-implies-B}.
\end{remark}

\subsubsection{Route B with prescribed von Mangoldt masses}\label{subsec:rb-prescribed}

The masses $\mathsf a^\zeta_n=\Lambda(n)/(2\sqrt n)$ are known and do not depend on (D) or on the ordinates of $\zeta$. Prescribing them leaves the ordinates as the only unknowns, and uniqueness then follows from Hadamard's factorisation and the functional equation alone, without Lemma~\CGref{lem:rb-mass-ms} and Hamburger's theorem. The tool is the identification of the Laplace transform of the residual.

\emph{Laplace transforms of distributions.} Let $\mathscr T\in\mathcal D'(\mathbb R)$ have support in $[0,\infty)$ and suppose that $e^{-au}\mathscr T$ is a tempered distribution for some $a\in\mathbb R$. For $\Re s>a+\frac12$ put $z:=s-\frac12$ and
\begin{equation}\label{eq:rb-laplace-def}
 \mathcal L\mathscr T(s):=\bigl\langle\mathscr T,\ \mathsf j(u)\,e^{-zu}\bigr\rangle,
\end{equation}
where $\mathsf j\in C^\infty(\mathbb R)$ equals $1$ on $[-1,\infty)$ and $0$ on $(-\infty,-2]$; the value does not depend on $\mathsf j$, and $\mathcal L\mathscr T$ is holomorphic there \cite[Ch.~8]{Zemanian1965}. The shift by $\frac12$ in $z$ is the normalisation of \CGeqref{eq:rb-Rhat}. Since $\mathsf j'$ vanishes on the support of $\mathscr T$, $\mathcal L(\partial_u\mathscr T)(s)=z\,\mathcal L\mathscr T(s)$; also $\mathcal L\delta_0=1$, $\mathcal L\delta_0'=z$, and for a locally integrable $g$ on $[0,\infty)$ with $|g(u)|\le ce^{au}$, $\mathcal L(\mathfrak g\mathbf 1_+)(s)=\int_0^\infty e^{-zu}g(u)\,du$. We write $\mathbf 1_+:=\mathbf 1_{[0,\infty)}$ and $\mathrm{fp}\,\frac{\mathbf 1_+}u$ for the distribution $\varphi\mapsto\int_0^1\frac{\varphi(u)-\varphi(0)}u\,du+\int_1^\infty\frac{\varphi(u)}u\,du$, of order $1$, supported in $[0,\infty)$, with restriction $1/u$ to $(0,\infty)$.

\begin{lemma}[The Archimedean distribution]\label{lem:rb-archimedean}
Let $\mathfrak g(u):=\frac{e^{-u/2}}{e^{2u}-1}-\frac1{2u}$ for $u>0$; $\mathfrak g$ is bounded, $\mathfrak g(0^+)=-\frac34$, and $\mathfrak g(u)=-\frac1{2u}+O(e^{-5u/2})$ as $u\to\infty$, in particular bounded. Put $\mathrm c_{\mathscr K}:=-\frac14\log(2\pi e^{\gamma_E})$, $\gamma_E$ Euler's constant, and
\begin{equation}\label{eq:rb-K-dist}
 \mathscr K:=\tfrac12e^{u/2}\mathbf 1_+-\tfrac12\,\mathfrak g\,\mathbf 1_+-\tfrac14\,\mathrm{fp}\,\frac{\mathbf 1_+}u+\mathrm c_{\mathscr K}\,\delta_0 .
\end{equation}
Then $\mathscr K$ is a distribution of order $1$ supported in $[0,\infty)$, its restriction to $(0,\infty)$ is the function $\mathsf k$ of \CGref{eq:transport-defs}, and
\begin{equation}\label{eq:rb-LK}
 \mathcal L\mathscr K(s)=\tfrac12\,\frac{G'}G(s)\qquad(\Re s>1).
\end{equation}
\end{lemma}

\begin{proof}
On $(0,\infty)$, $\frac12e^{u/2}-\frac12\mathfrak g(u)-\frac1{4u}=\frac12\bigl(e^{u/2}-\frac{e^{-u/2}}{e^{2u}-1}\bigr)=\mathsf k(u)$. The expansion $\frac{e^{-u/2}}{e^{2u}-1}=\frac1{2u}-\frac34+O(u)$ gives the behaviour of $\mathfrak g$ at $0$. With $z=s-\frac12$ and $\Re s>1$: $\mathcal L(\frac12e^{u/2}\mathbf 1_+)=\frac1{2(z-\frac12)}=\frac1{2(s-1)}$. Gauss's integral $\frac{\Gamma'}{\Gamma}(x)=\int_0^\infty\bigl(\frac{e^{-t}}t-\frac{e^{-xt}}{1-e^{-t}}\bigr)dt$ ($\Re x>0$) for the logarithmic derivative of $\Gamma$ \cite[Ch.~XII]{WhittakerWatson1927}, with $t=2u$ and $x=\frac s2$, reads
\begin{equation}\label{eq:rb-aux-12}
 \tfrac12\frac{\Gamma'}{\Gamma}\bigl(\tfrac s2\bigr)=\int_0^\infty\Bigl(\frac{e^{-2u}}{2u}-\frac{e^{-su}}{1-e^{-2u}}\Bigr)du
 =\int_0^\infty\Bigl(\frac{e^{-2u}}{2u}-e^{-zu}\frac{e^{-u/2}}{e^{2u}-1}\Bigr)du-\frac1s,
\end{equation}
because $\frac{e^{-su}}{1-e^{-2u}}=e^{-zu}\bigl(\frac{e^{-u/2}}{e^{2u}-1}+e^{-u/2}\bigr)$ and $\int_0^\infty e^{-zu}e^{-u/2}du=\frac1s$. Writing $\frac{e^{-u/2}}{e^{2u}-1}=\mathfrak g(u)+\frac1{2u}$ and using Frullani's integral $\int_0^\infty\frac{e^{-2u}-e^{-zu}}u\,du=\log\frac z2$ (both sides vanish at $z=2$ and have derivative $\int_0^\infty e^{-zu}du=1/z$ in $z$, for $\Re z>0$), the integral in \CGeqref{eq:rb-aux-12} equals $-\mathcal L(\mathfrak g\mathbf 1_+)(s)+\frac12\log\frac z2$. Hence
\begin{equation}\label{eq:rb-aux-13}
 \tfrac12\frac{G'}G(s)=\frac1{2s}+\frac1{2(s-1)}-\frac14\log\pi+\frac14\frac{\Gamma'}{\Gamma}\bigl(\tfrac s2\bigr)
 =\mathcal L\bigl(\tfrac12e^{u/2}\mathbf 1_+\bigr)(s)-\tfrac12\mathcal L(\mathfrak g\mathbf 1_+)(s)+\tfrac14\log\frac z2-\tfrac14\log\pi .
\end{equation}
Finally $\mathcal L\bigl(\mathrm{fp}\,\frac{\mathbf 1_+}u\bigr)(s)=\int_0^1\frac{e^{-zu}-1}u\,du+\int_1^\infty\frac{e^{-zu}}u\,du=-\gamma_E-\log z$: the left side has derivative $-\int_0^\infty e^{-zu}du=-1/z$ in $z$, and its value at $z=1$ is $\int_0^1\frac{e^{-u}-1}u\,du+\int_1^\infty\frac{e^{-u}}u\,du=-\gamma_E$, the classical integral representation of Euler's constant \cite[Ch.~XII]{WhittakerWatson1927}. Substituting, $\frac14\log\frac z2=-\frac14\mathcal L(\mathrm{fp}\,\frac{\mathbf 1_+}u)-\frac14(\gamma_E+\log2)$, and \CGeqref{eq:rb-aux-13} becomes \CGeqref{eq:rb-LK} with the constant $\mathrm c_{\mathscr K}$. (The identity was also checked numerically to $25$ digits at four points $s$; file \texttt{laplace\_check.py}.)
\end{proof}

\begin{lemma}[Laplace transform of the residual]\label{lem:rb-laplace}
Let $\boldsymbol\kappa$ be $C$-admissible and put $\mathfrak H_{\boldsymbol\kappa}(u):=\sum_k\kappa_k^{-2}\cos(\kappa_ku)$, which converges absolutely and uniformly on $\mathbb R$ since $\mathfrak H_{\boldsymbol\kappa}(0)=\sum_k\kappa_k^{-2}<\infty$, and
\begin{equation}\label{eq:rb-S-A-dist}
 \mathscr S_{\boldsymbol\kappa}:=-\partial_u^2\bigl(\mathfrak H_{\boldsymbol\kappa}\mathbf 1_+\bigr)+\mathfrak H_{\boldsymbol\kappa}(0)\,\delta_0',\qquad
 \mathscr A_\zeta:=\sum_{n\ge2}\mathsf a^\zeta_n\,\delta_{\log n},\qquad
 \mathscr E_{\boldsymbol\kappa}:=\mathscr S_{\boldsymbol\kappa}+\mathscr A_\zeta-\mathscr K .
\end{equation}
Then $\mathscr E_{\boldsymbol\kappa}$ is a distribution of local order at most $2$ (its restriction to every bounded interval is of order $\le2$; the support is not compact, since $\mathscr A_\zeta$ is a locally finite measure on $[\log2,\infty)$), supported in $[0,\infty)$, Laplace-transformable on $\Re s>1$, whose restriction to $(0,\infty)$ is $\mathfrak R(\boldsymbol\kappa,\mathsf a^\zeta)$, and for $\Re s>1$
\begin{equation}\label{eq:rb-LR}
 \mathcal L\mathscr E_{\boldsymbol\kappa}(s)=\widehat{\mathsf R}(s)=\tfrac12\Bigl(\frac{P_{\boldsymbol\kappa}'}{P_{\boldsymbol\kappa}}-\frac{\xi'}{\xi}\Bigr)(s),\qquad\xi:=\tfrac12G\zeta ,
\end{equation}
where $\widehat{\mathsf R}$ is \CGeqref{eq:rb-Rhat} with the masses $\mathsf a^\zeta$ for all $n$. The same holds for a finite section $(\boldsymbol\kappa,\mathsf a)$ with $\mathscr A_\zeta$ replaced by $\sum_{n\le n^*}\mathsf a_n\delta_{\log n}$ and $-\frac12\xi'/\xi$ by $-\frac12G'/G+\sum_{n\le n^*}\mathsf a_nn^{\frac12-s}$.
\end{lemma}

\begin{proof}
\emph{Spectral part.} $\kappa_k\ge\tau_k-C\gg k/\log(k+1)$, so $\mathfrak H_{\boldsymbol\kappa}(0)<\infty$ and $\mathfrak H_{\boldsymbol\kappa}$ is continuous and bounded by $\mathfrak H_{\boldsymbol\kappa}(0)$. In $\mathcal D'(0,\infty)$ a uniformly convergent series may be differentiated termwise, so $-\mathfrak H_{\boldsymbol\kappa}''=\sum_k\cos(\kappa_ku)$ there; the term $\mathfrak H_{\boldsymbol\kappa}(0)\delta_0'$ is supported at $0$. Hence $\mathscr S_{\boldsymbol\kappa}|_{(0,\infty)}=\sum_k\cos(\kappa_ku)$, and $\mathscr S_{\boldsymbol\kappa}$ has order at most $2$. The Laplace transform of a distribution $\mathscr T$ supported in $[0,\infty)$ is $\mathcal L\mathscr T(s):=\langle\mathscr T,\chi e^{-zu}\rangle$, $z=s-\frac12$, with $\chi\in C^\infty(\mathbb R)$ equal to $1$ on a neighbourhood of the support and to $0$ on $(-\infty,-1]$; thus $\mathcal L(\delta_0')(s)=-\frac d{du}e^{-zu}\big|_{u=0}=z$ and $\mathcal L(\mathbf 1_+f)(s)=\int_0^\infty f(u)e^{-zu}du$. For $\Re z>0$, Fubini (the double integral $\sum_k\kappa_k^{-2}\int_0^\infty e^{-\Re z\,u}du$ is finite) gives $\mathcal L(\mathfrak H_{\boldsymbol\kappa}\mathbf 1_+)(s)=\sum_k\kappa_k^{-2}\frac z{z^2+\kappa_k^2}$, so
\begin{equation}\label{eq:rb-aux-14}
 \mathcal L\mathscr S_{\boldsymbol\kappa}(s)=-z^2\sum_k\frac z{\kappa_k^2(z^2+\kappa_k^2)}+z\sum_k\frac1{\kappa_k^2}
 =\sum_k\frac{z}{z^2+\kappa_k^2}=\tfrac12\frac{P_{\boldsymbol\kappa}'}{P_{\boldsymbol\kappa}}(s),
\end{equation}
because $\frac z{\kappa^2}-\frac{z^3}{\kappa^2(z^2+\kappa^2)}=\frac z{z^2+\kappa^2}$, and $P_{\boldsymbol\kappa}=\prod_k\frac{z^2+\kappa_k^2}{\frac14+\kappa_k^2}$ converges locally uniformly, so its logarithmic derivative is the termwise sum $\sum_k\frac{2z}{z^2+\kappa_k^2}$.

\emph{Arithmetic part.} On $[0,u^*]$ only $n\le e^{u^*}$ contribute, so $\mathscr A_\zeta$ is a locally finite measure, of order $0$, and $e^{-au}\mathscr A_\zeta$ is a finite measure for $a>\frac12$. For $\Re s>1$, $\mathcal L\mathscr A_\zeta(s)=\sum_n\frac{\Lambda(n)}{2\sqrt n}n^{-z}=\frac12\sum_n\Lambda(n)n^{-s}=-\frac12\frac{\zeta'}\zeta(s)$.

\emph{Assembly.} By Lemma~\CGref{lem:rb-archimedean}, $\mathcal L\mathscr E_{\boldsymbol\kappa}=\frac12P'_{\boldsymbol\kappa}/P_{\boldsymbol\kappa}-\frac12\zeta'/\zeta-\frac12G'/G$, and $(G\zeta)'/(G\zeta)=\xi'/\xi$ since $\xi=\frac12G\zeta$. The restriction to $(0,\infty)$ is $\sum_k\cos(\kappa_ku)+\sum_n\mathsf a^\zeta_n\delta(u-\log n)-\mathsf k(u)$, which is \CGeqref{eq:rb-residual}. The finite-section case is identical with a finite sum.
\end{proof}

\begin{remark}[The boundary term]\label{rem:rb-boundary-term}
Without the correction $\mathfrak H_{\boldsymbol\kappa}(0)\delta_0'$ in \CGeqref{eq:rb-S-A-dist}, the transform of $-\partial_u^2(\mathfrak H_{\boldsymbol\kappa}\mathbf 1_+)$ is $\frac12P'_{\boldsymbol\kappa}/P_{\boldsymbol\kappa}-(s-\frac12)\mathfrak H_{\boldsymbol\kappa}(0)$: by \CGeqref{eq:rb-aux-14}, $\mathcal L(-\partial_u^2(\mathfrak H_{\boldsymbol\kappa}\mathbf 1_+))=-z^2\mathcal L(\mathfrak H_{\boldsymbol\kappa}\mathbf 1_+)$, which differs from $\frac12P'_{\boldsymbol\kappa}/P_{\boldsymbol\kappa}$ by $-z\mathfrak H_{\boldsymbol\kappa}(0)=-\mathcal L(\mathfrak H_{\boldsymbol\kappa}(0)\delta_0')$, the transform of the boundary term that the second derivative of the cut-off function produces at $u=0$. Either normalisation gives the same restriction to $(0,\infty)$ and differs by a polynomial in $s$, which is harmless in the argument below. Both $\mathscr K$ and $\mathscr S_{\boldsymbol\kappa}$ carry singular terms at $u=0$; this is why the rung system is posed in $\mathcal D'(0,\infty)$ and not on $[0,\infty)$.
\end{remark}

\begin{proposition}[Uniqueness with prescribed masses]\label{prop:rb-prescribed-uniqueness}
Let $\boldsymbol\kappa$ be $C$-admissible for some $C$, and suppose $\mathfrak R(\boldsymbol\kappa,\mathsf a^\zeta)=0$ in $\mathcal D'(0,\infty)$. Then all nontrivial zeros of $\zeta$ lie on $\Re s=\frac12$, and $\boldsymbol\kappa=\boldsymbol\gamma$.
\end{proposition}

\begin{proof}
By Lemma~\CGref{lem:rb-laplace}, $\mathscr E_{\boldsymbol\kappa}$ is a distribution of finite order supported in $[0,\infty)$ whose restriction to $(0,\infty)$ vanishes, so it is supported at $\{0\}$. A distribution of finite order supported at a point is a finite linear combination of $\delta_0$ and its derivatives \cite[Thm.~2.3.4]{Hormander1990}, so $\widehat{\mathsf R}=\mathcal L\mathscr E_{\boldsymbol\kappa}$ is a polynomial $p$ in $z=s-\frac12$. By \CGeqref{eq:rb-LR}, $(P_{\boldsymbol\kappa}e^{-Q}/\xi)'=0$ on $\Re s>1$ with $Q'=2p$, $Q$ a polynomial. Hence $P_{\boldsymbol\kappa}=c\,\xi\,e^{Q}$ on $\Re s>1$ with $c\ne0$, and, both sides being entire, the identity holds on all of $\mathbb C$.  Only now is the quotient discussed: the identity shows that $P_{\boldsymbol\kappa}$ and $\xi$ have the same zeros with the same multiplicities, since $e^Q$ has none.  In the variable $w=s-\frac12$, $P_{\boldsymbol\kappa}$ is a constant multiple of $\prod_k(1+w^2/\kappa_k^2)$, and $N_{\boldsymbol\kappa}(t)\ll t\log t$, so $P_{\boldsymbol\kappa}$ has order $1$; $\xi$ has order $1$. Comparing the Hadamard factorizations $P_{\boldsymbol\kappa}=e^{A_1+B_1s}\Pi(s)$ and $\xi=e^{A_2+B_2s}\Pi(s)$, with the same canonical product $\Pi$ over the common zeros, gives $e^{Q}=c^{-1}e^{(A_1-A_2)+(B_1-B_2)s}$, so $\deg Q\le1$. Since $P_{\boldsymbol\kappa}(1-s)=P_{\boldsymbol\kappa}(s)$ and $\xi(1-s)=\xi(s)$, $e^{Q(1-s)}=e^{Q(s)}$, that is, $Q(1-s)-Q(s)\in2\pi i\mathbb Z$ for every $s$; the left side is continuous, hence constant, and for $Q(s)=a+bs$ it equals $b(1-2s)$, so $b=0$. Therefore $P_{\boldsymbol\kappa}=c'\xi$. The zeros of $P_{\boldsymbol\kappa}$ are $\frac12\pm i\kappa_k$, so every zero of $\xi$, that is, every nontrivial zero of $\zeta$, has real part $\frac12$, and the ordinates counted with multiplicity are the $\kappa_k$. Both lists are nondecreasing, so $\boldsymbol\kappa=\boldsymbol\gamma$.
\end{proof}

\begin{definition}[Statement B$'$]\label{def:rb-Bprime}
\emph{Statement B$'(C)$} says that there are $C$-admissible lists $\boldsymbol\kappa^{(m)}$, each equal to the clock beyond some finite index, with $\mathsf s_m(\boldsymbol\kappa^{(m)},\mathsf a^\zeta)\to0$.
\end{definition}

\begin{corollary}[Route B$'$]\label{cor:rb-Bprime}
\begin{enumerate}[label=\textup{(\alph*)}]
\item If B$'(C)$ holds for some $C$, then every sequence as in Definition~\CGref{def:rb-Bprime} converges coordinatewise to $\boldsymbol\gamma$, and all nontrivial zeros of $\zeta$ lie on $\Re s=\frac12$.
\item If all nontrivial zeros of $\zeta$ lie on $\Re s=\frac12$, then B$'(C)$ holds for every $C\ge C_\eta$.
\end{enumerate}
Hence the Riemann hypothesis is equivalent to B$'(C)$ for some $C<\infty$.
\end{corollary}

\begin{proof}
(a) $|\mathsf a^\zeta_n|\le \mathbb M_*$ by the elementary computation in Lemma~\CGref{lem:rb-canonical-box}, so $(\boldsymbol\kappa^{(m)},\mathsf a^\zeta)\in\mathbb K(C,\mathbb M_*)$. By Proposition~\CGref{prop:rb-compactness}, every subsequence has a further subsequence converging coordinatewise to some $(\boldsymbol\kappa,\mathsf a^\zeta)$ with $\mathfrak R(\boldsymbol\kappa,\mathsf a^\zeta)=0$ and $\boldsymbol\kappa$ $C$-admissible; the masses do not move. By Proposition~\CGref{prop:rb-prescribed-uniqueness}, $\boldsymbol\kappa=\boldsymbol\gamma$ and the zeros lie on the line. The product $\prod_k[\tau_k-C,\tau_k+C]$ is compact and metrizable, for instance by $\mathrm d(\boldsymbol\kappa,\boldsymbol\kappa')=\sum_k2^{-k}\min(1,|\kappa_k-\kappa'_k|)$, which induces the product topology; in a compact metric space a sequence all of whose convergent subsequences have the same limit converges to it. Hence the whole sequence converges. (b) Under this hypothesis the rung system with $(\boldsymbol\gamma,\mathsf a^\zeta)$ is the distributional explicit formula \CGref{thm:distributional-EF}, and $|\gamma_k-\tau_k|\le C_\eta$ by \CGref{lem:bounded-displacement}. The truncations in the proof of Corollary~\CGref{cor:rb-D-implies-B}, with the masses kept equal to $\mathsf a^\zeta$, give the sections.
\end{proof}

\begin{remark}[What prescribing the masses changes, and what it does not]\label{rem:rb-prescribed}
With prescribed masses, uniqueness for every list at bounded distance from the clock follows from Hadamard's theorem and the functional equation; with free masses it needs in addition Lemma~\CGref{lem:rb-mass-ms} and Hamburger's theorem (Theorem~\CGref{thm:rb-uniqueness}). Routes B and B$'$ have the same single open piece, existence, and each is equivalent to the Riemann hypothesis (Corollaries~\CGref{cor:rb-equivalence} and~\CGref{cor:rb-Bprime}). With prescribed masses the whole window $(0,u^*)$ constrains the ordinates; the part above $\log\Omega$ can no longer be absorbed by the masses, as it could in Heuristic~\CGref{heur:rb-informative}.

The equivalence itself is classical in substance. The prime side of the rung system is the prime side of the explicit formula of Riemann and von Mangoldt in the form of Guinand and Weil \cite{Guinand1948,Weil1952}, and the uniqueness is Hadamard's factorisation. Arias de Reyna \cite{AriasdeReyna2024} proves that the Riemann hypothesis is equivalent to the temperedness of the measure $-\sum_n\frac{\Lambda(n)}{\sqrt n}(\delta_{\log n}+\delta_{-\log n})+2\cosh(x/2)\,dx$, the same prime side up to the factor $2$ and the symmetrisation, and identifies its Fourier transform with $\sum_\gamma\delta_{\gamma/2\pi}-2\vartheta'(2\pi t)\,dt$ under that hypothesis; by partial summation his criterion is von Koch's $\sum_{n\le x}\Lambda(n)=x+O(x^{1/2}\log^2x)$. Statement B$'$ is the finite-section form of the same equivalence: it asks for a real spectrum in a tube around the clock that reproduces the prime side against a dense family of tests. It is not a weaker intermediate statement, and the finite sections are not more constrained than the global statement; what they offer is a formulation in which each stage is a finite-dimensional problem with computable certificates (Section~\CGref{subsec:rb-staged}) and explicit candidates (Section~\CGref{subsec:rb-gonek}). Dyson's suggestion that the zeros form a one-dimensional quasicrystal, and the classification of crystalline measures by Kurasov and Sarnak \cite{KurasovSarnak2020} and Olevskii and Ulanovskii \cite{OlevskiiUlanovskii2020}, concern the same object; the zeta list is not in those classes because of its density, as Arias de Reyna points out.
\end{remark}

\subsubsection{Finite stages do not decide route B}\label{subsec:rb-nonfinitary}

Let $r\ge2$. For $\mathsf f$ with compact support in $(0,u^*]$ such that $\mathsf f,\dots,\mathsf f^{(r-1)}$ are absolutely continuous and $\mathsf f^{(r)}\in L^1$, the residual is defined by \CGeqref{eq:rb-residual}, since $|\mathsf c_{\mathsf f}(t)|\le\|\mathsf f''\|_1t^{-2}$, and we put
\begin{equation}\label{eq:rb-Wr}
 \mathfrak W_r(\mathsf f):=(1+u^{*2})\cosh\tfrac{u^*}2\sum_{j=0}^r\binom rj\|\mathsf f^{(r-j)}\|_{L^1}.
\end{equation}
Here $\boldsymbol\gamma$ is the nondecreasing list of the ordinates of all nontrivial zeros of $\zeta$ with positive imaginary part, with multiplicity, and $d_k=\Re\rho_k-\frac12$ as in \CGeqref{eq:Y-offline}; no hypothesis on the zeros is made.

\begin{proposition}[The zeta list solves every finite stage up to the verified height]\label{prop:rb-nonfinitary}
Let $H_0:=3\cdot10^{12}$. For every $r\ge2$ and every real $\mathsf f$ as above,
\begin{equation}\label{eq:rb-nonfinitary}
 \bigl|\langle\mathfrak R(\boldsymbol\gamma,\mathsf a^\zeta),\mathsf f\rangle\bigr|\le\frac14\,\mathfrak W_r(\mathsf f)\sum_{k:\,d_k\ne0}\gamma_k^{-r}\le\frac r{4(r-1)}\,\mathfrak W_r(\mathsf f)\,H_0^{1-r}\log H_0 .
\end{equation}
The list $\boldsymbol\gamma$ is $C_\eta$-admissible. Consequently, for every finite set $\mathscr T$ of such test functions (with the same $r$) and every $\delta>0$, there is a finite section $(\boldsymbol\kappa,\mathsf a)\in\mathbb K(C_\eta,\mathbb M_*)$ whose masses are $\mathsf a^\zeta_n$ for $n\le n^*$ and $0$ beyond, such that
\begin{equation}\label{eq:rb-nonfinitary-section}
 \max_{\mathsf f\in\mathscr T}\bigl|\langle\mathfrak R(\boldsymbol\kappa,\mathsf a),\mathsf f\rangle\bigr|\le\max_{\mathsf f\in\mathscr T}\frac r{4(r-1)}\,\mathfrak W_r(\mathsf f)\,H_0^{1-r}\log H_0+\delta ,
\end{equation}
whether or not the Riemann hypothesis holds.
\end{proposition}

\begin{proof}
\emph{The residual is the sum over the zeros off the line.} For $\mathsf f\in C_c^\infty(0,\infty)$, Theorem~\CGref{thm:rung-zero-sum} and \CGeqref{eq:rung-functional} give $\langle\mathfrak R(\boldsymbol\gamma,\mathsf a^\zeta),\mathsf f\rangle=\sum_k\mathsf c_{\mathsf f}(\gamma_k)+\frac12\sum_n\Lambda(n)n^{-1/2}\mathsf f(\log n)-\mathsf k_{\mathsf f}=-\langle\mathscr Y,\mathsf f\rangle$. Both sides are continuous under convergence of $\mathsf f,\mathsf f',\mathsf f''$ in $L^1$ and of $\mathsf f$ uniformly, with supports in a fixed compact subset of $(0,\infty)$: on the left by $|\mathsf c_{\mathsf g}(t)|\le\|\mathsf g''\|_1t^{-2}$ and the finiteness of the mass sum, on the right by the bound of the next paragraph with $r=2$. Mollification therefore extends the identity to every $\mathsf f$ as in the statement.

\emph{One zero.} Let $d=d_k$, $|d|<\frac12$, and $0\le u\le u^*$. Every derivative of $\mathsf e_d(u):=\cosh(du)-1$ is bounded by $d^2(1+u^2)\cosh\frac u2$: $0\le\mathsf e_d(u)\le\frac{d^2u^2}2\cosh(du)$; $|\mathsf e_d'(u)|=|d\sinh(du)|\le d^2u\cosh(du)$; and $|\mathsf e_d^{(j)}(u)|\le|d|^j\cosh(du)\le d^2\,2^{2-j}\cosh(du)$ for $j\ge2$. Hence $\mathsf g_k:=\mathsf f\,\mathsf e_{d_k}$ has $\|\mathsf g_k^{(r)}\|_1\le d_k^2\,\mathfrak W_r(\mathsf f)$ by the Leibniz rule, and $r$ integrations by parts, without boundary terms, give
\[
 \Bigl|\int_0^\infty\mathsf f(u)\bigl(\cosh(d_ku)-1\bigr)\cos(\gamma_ku)\,du\Bigr|\le\gamma_k^{-r}\|\mathsf g_k^{(r)}\|_1\le\tfrac14\,\mathfrak W_r(\mathsf f)\,\gamma_k^{-r},
\]
and the term vanishes when $d_k=0$. This is the first inequality in \CGeqref{eq:rb-nonfinitary}.

\emph{The verified zeros.} By \cite{PlattTrudgian2021} every zero with $0<\gamma\le H_0$ has $d=0$, so the sum runs over $\gamma_k>H_0$. By \cite{HasanalizadeShenWong2022}, $N_\zeta(t)\le\frac t{2\pi}\log\frac t{2\pi e}+0.1038\log t+0.2573\log\log t+9.3675\le\frac t{2\pi}\log t$ for $t\ge100$, and partial summation gives
\[
\begin{aligned}
 \sum_{\gamma_k>H_0}\gamma_k^{-r}&\le r\int_{H_0}^\infty N_\zeta(t)\,t^{-r-1}dt\le\frac r{2\pi}\int_{H_0}^\infty t^{-r}\log t\,dt\\
 &=\frac{r\,H_0^{1-r}}{2\pi(r-1)}\Bigl(\log H_0+\frac1{r-1}\Bigr)\le\frac r{r-1}H_0^{1-r}\log H_0 .
\end{aligned}
\]

\emph{Sections.} $|\gamma_k-\tau_k|\le C_\eta$ by Lemma~\CGref{lem:bounded-displacement}, which is unconditional, and $\gamma_1>2\pi$, so $\boldsymbol\gamma$ is $C_\eta$-admissible. The truncations of $(\boldsymbol\gamma,\mathsf a^\zeta)$ in the proof of Corollary~\CGref{cor:rb-D-implies-B} are finite sections in $\mathbb K(C_\eta,\mathbb M_*)$, and for $n^*\ge e^{u^*}$ their residual on $\mathsf f$ differs from $\langle\mathfrak R(\boldsymbol\gamma,\mathsf a^\zeta),\mathsf f\rangle$ by $\sum_{k>k^*}(\mathsf c_{\mathsf f}(\tau_k)-\mathsf c_{\mathsf f}(\gamma_k))\to0$ as $k^*\to\infty$; that part of the proof uses no hypothesis. Take $k^*$ large for the finitely many $\mathsf f\in\mathscr T$.
\end{proof}

\begin{remark}[Reading of Proposition~\CGref{prop:rb-nonfinitary}]\label{rem:rb-nonfinitary}
If the Riemann hypothesis holds, the residual in \CGeqref{eq:rb-nonfinitary} vanishes. If it fails, B and B$'$ fail (Corollaries~\CGref{cor:rb-equivalence} and~\CGref{cor:rb-Bprime}), and still every finite stage has sections in the box with residual at most $\frac r{4(r-1)}\max_{\mathscr T}\mathfrak W_r\,H_0^{1-r}\log H_0+\delta$. No computation on finitely many tests at a precision coarser than this bound distinguishes the two cases: route B is not finitary, and finite sections, computed or certified at such a precision, carry no information on the Riemann hypothesis. A test sees a residual of size $\eta$ only if $\mathfrak W_r(\mathsf f)\ge\frac{4(r-1)}r\,\eta\,H_0^{r-1}/\log H_0$ for every admissible $r$. For the B-spline tests $\varphi^{\mathrm N}_{c,h}$ of Section~\CGref{subsec:rb-num-verified}, with $h\le1$ and $r\le7$, $\mathfrak W_r(\varphi^{\mathrm N}_{c,h})\le\mathsf K\,3^rh^{1-r}$ with $\mathsf K:=(1+u^{*2})\cosh\frac{u^*}2$ (by \CGeqref{eq:rb-W-bspline}), so such a test sees $\eta$ only if $h\le3\,(3\mathsf K\log H_0/\eta)^{1/(r-1)}H_0^{-1}$ for $r=2,\dots,7$: the test must resolve frequencies of the order of $H_0$. For the tests used there the bound \CGeqref{eq:rb-nonfinitary} with $r=7$, evaluated in ball arithmetic by program \texttt{certV}, is below $10^{-57}$ at every stage computed (column $\epsilon_\zeta$ of Table~\CGref{tab:rb-verified}).
\end{remark}

\subsubsection{Staged construction: a certificate theorem and a graded extension}\label{subsec:rb-staged}

\begin{definition}[Staged sections]\label{def:rb-staged}
Fix $u^*_1<u^*_2<\cdots\to\infty$, finite test sets $\mathscr T_1\subset\mathscr T_2\subset\cdots$ supported in $(0,u^*_m)$ with dense union as in Section~\CGref{subsec:rb-routeB}, $\mathrm t_m:=|\mathscr T_m|$, and $k^*_m\to\infty$. Stage $m$ starts from a list $\boldsymbol\kappa^{(m,0)}$ that equals the clock beyond some finite index, moves the coordinates $k\le k^*_m$ by $x\in\mathbb R^{k^*_m}$, and has masses $\mathsf a^\zeta$. Its residual map is
\begin{equation}\label{eq:rb-Fm}
 \mathbb F_m:\mathbb R^{k^*_m}\to\mathbb R^{\mathrm t_m},\qquad \mathbb F_m(x)_\varphi:=\bigl\langle\mathfrak R(\boldsymbol\kappa^{(m,0)}+x,\mathsf a^\zeta),\varphi\bigr\rangle\quad(\varphi\in\mathscr T_m),
\end{equation}
where $(\boldsymbol\kappa+x)_k=\kappa_k+x_k$ for $k\le k^*_m$ and $\kappa_k$ otherwise, and its Jacobian at $0$ is
\begin{equation}\label{eq:rb-jacobian}
 (\mathbb J_m)_{\varphi,k}=\partial_{x_k}\mathbb F_m(0)_\varphi=-\int_0^\infty\varphi(u)\,u\sin(\kappa^{(m,0)}_ku)\,du\qquad(\varphi\in\mathscr T_m,\ k\le k^*_m).
\end{equation}
We use the norms $\|x\|_\infty$ on $\mathbb R^{k^*_m}$ and $\mathbb R^{\mathrm t_m}$ and the induced operator norms $\|\cdot\|_{\infty\to\infty}$ (maximal absolute row sum). Put
\begin{equation}\label{eq:rb-eta-H}
 \mathsf h_m:=\|\mathbb F_m(0)\|_\infty,\qquad \mathbb H_m:=k^*_m\max_{\varphi\in\mathscr T_m}\int_0^\infty|\varphi(u)|\,u^2\,du .
\end{equation}
\end{definition}

\begin{lemma}[Monotone matching]\label{lem:rb-matching}
Let $(\mathrm a_k)_{k\ge1}$ be real with $|\mathrm a_k-\tau_k|\le C$ for all $k$, and let $(\mathrm a_{(k)})$ be its nondecreasing rearrangement. Then $|\mathrm a_{(k)}-\tau_k|\le C$ for all $k$. If moreover $\mathrm a_k=\tau_k$ for $k>k^*$, then $\mathrm a_{(k)}=\tau_k$ for all $k>k^*+j$, where $j:=\#\{k>k^*:\tau_k\le\tau_{k^*}+C\}<\infty$.
\end{lemma}

\begin{proof}
At least $k$ of the $\mathrm a_i$ (those with $i\le k$) are $\le\tau_k+C$, so $\mathrm a_{(k)}\le\tau_k+C$; at least $k$ of them (those with $i\ge k$) are $\ge\tau_k-C$, hence at most $k-1$ are $<\tau_k-C$, so $\mathrm a_{(k)}\ge\tau_k-C$. For the second claim, the elements of the multiset not exceeding $\tau_{k^*}+C$ are the $k^*$ moved ones (each is at most $\tau_{k^*}+C$) and the $j$ clock points $\tau_k$, $k^*<k\le k^*+j$; the remaining elements are exactly $\tau_k$, $k>k^*+j$, in increasing order.
\end{proof}

\begin{lemma}[Lipschitz bound for the Jacobian]\label{lem:rb-lipschitz}
For all $x,x'\in\mathbb R^{k^*_m}$, $\|D\mathbb F_m(x)-D\mathbb F_m(x')\|_{\infty\to\infty}\le \mathbb H_m\|x-x'\|_\infty$, and
\begin{equation}\label{eq:rb-taylor}
 \|\mathbb F_m(x)-\mathbb F_m(0)-\mathbb J_mx\|_\infty\le\tfrac12\mathbb H_m\|x\|_\infty^2 .
\end{equation}
\end{lemma}

\begin{proof}
$\mathbb F_m(x)_\varphi=\sum_{k\le k^*_m}\mathsf c_\varphi(\kappa^{(m,0)}_k+x_k)+\text{const}$, so $D^2\mathbb F_m(x)_\varphi$ is diagonal with entries $\partial_{x_k}^2\mathbb F_m(x)_\varphi=-\int\varphi(u)u^2\cos((\kappa_k+x_k)u)\,du$, of modulus at most $\int|\varphi|u^2$. The row $\varphi$ of $D\mathbb F_m(x)-D\mathbb F_m(x')$ has absolute sum at most $\sum_{k\le k^*_m}\int|\varphi|u^2\,|x_k-x'_k|\le \mathbb H_m\|x-x'\|_\infty$. The Taylor bound follows by integrating $D\mathbb F_m$ along the segment.
\end{proof}

\begin{theorem}[Certificate for one stage]\label{thm:rb-certificate}
Suppose $\mathbb J_m$ has a right inverse $\mathbb Q_m:\mathbb R^{\mathrm t_m}\to\mathbb R^{k^*_m}$, $\mathbb J_m\mathbb Q_m=I_{\mathrm t_m}$, with $\|\mathbb Q_m\|_{\infty\to\infty}\le \mathbb L_m$, and
\begin{equation}\label{eq:rb-Am}
 \mathsf l_m:=2\mathbb L_m^2\mathbb H_m\mathsf h_m<1 .\tag{A$_m$}
\end{equation}
Then the map $\mathbb T_m(y):=y-\mathbb F_m(\mathbb Q_my)$ is a contraction of the ball $\{\|y\|_\infty\le2\mathsf h_m\}\subset\mathbb R^{\mathrm t_m}$ with constant $\mathsf l_m$, and $x_*:=\mathbb Q_my_*$, $y_*$ its unique fixed point in that ball, satisfies
\begin{equation}\label{eq:rb-cert-conclusion}
 \mathbb F_m(x_*)=0,\qquad\|x_*\|_\infty\le2\mathbb L_m\mathsf h_m .
\end{equation}
The uniqueness is that of $x_*$ among the zeros of $\mathbb F_m$ in $\mathbb Q_m\{\|y\|_\infty\le2\mathsf h_m\}\subset\operatorname{Ran}\mathbb Q_m$; nothing is asserted about zeros of $\mathbb F_m$ outside $\operatorname{Ran}\mathbb Q_m$.
\end{theorem}

\begin{proof}
$D\mathbb T_m(y)=I-D\mathbb F_m(\mathbb Q_my)\mathbb Q_m=(\mathbb J_m-D\mathbb F_m(\mathbb Q_my))\mathbb Q_m$ because $\mathbb J_m\mathbb Q_m=I$, so by Lemma~\CGref{lem:rb-lipschitz}, for $\|y\|_\infty\le2\mathsf h_m$, $\|D\mathbb T_m(y)\|_{\infty\to\infty}\le \mathbb H_m\|\mathbb Q_my\|_\infty \mathbb L_m\le2\mathbb H_m\mathbb L_m^2\mathsf h_m=\mathsf l_m<1$. Moreover $\mathbb T_m(y)=-\mathbb F_m(0)-\bigl(\mathbb F_m(\mathbb Q_my)-\mathbb F_m(0)-\mathbb J_m\mathbb Q_my\bigr)$, and by \CGeqref{eq:rb-taylor} the bracket has norm at most $\frac12\mathbb H_m\mathbb L_m^2\|y\|_\infty^2\le2\mathbb H_m\mathbb L_m^2\mathsf h_m^2$. Hence $\|\mathbb T_m(y)\|_\infty\le\mathsf h_m(1+\mathsf l_m)<2\mathsf h_m$: the ball is mapped into itself. Banach's fixed point theorem gives a unique fixed point $y_*$ in the ball, $y_*=y_*-\mathbb F_m(\mathbb Q_my_*)$, that is, $\mathbb F_m(x_*)=0$ with $\|x_*\|_\infty\le \mathbb L_m\|y_*\|_\infty\le2\mathbb L_m\mathsf h_m$. The uniqueness is that of $y_*$ in the ball of $\mathbb R^{\mathrm t_m}$, hence of $x_*$ in $\operatorname{Ran}\mathbb Q_m$; when $k^*_m>\mathrm t_m$ the Jacobian has a kernel and $\mathbb F_m$ may have other zeros in the ball of $\mathbb R^{k^*_m}$, about which the theorem says nothing. The iteration $y_{q+1}=\mathbb T_m(y_q)$, $y_0=0$, converges to $y_*$ geometrically with ratio $\mathsf l_m$; in the variable $x$ it is $x_{q+1}=x_q-\mathbb Q_m\mathbb F_m(x_q)$, a \emph{modified} Newton iteration, the right inverse being frozen at $\mathbb Q_m$ instead of being recomputed at each $x_q$.
\end{proof}

\begin{remark}[Evaluation errors and order of the ordinates]\label{rem:rb-certificate-errors}
Theorem~\CGref{thm:rb-certificate} is stated for exact values of $\mathsf h_m$, $\mathbb L_m$ and $\mathbb H_m$.  If the residual is only known up to an error, $\|\tilde{\mathbb F}_m(0)-\mathbb F_m(0)\|_\infty\le\epsilon_m$, the theorem applies with $\mathsf h_m$ replaced by $\|\tilde{\mathbb F}_m(0)\|_\infty+\epsilon_m$; if the computed matrix $\tilde{\mathbb Q}_m$ satisfies only $\|I-\mathbb J_m\tilde{\mathbb Q}_m\|_{\infty\to\infty}\le\delta_m<1$, then $\mathbb Q_m:=\tilde{\mathbb Q}_m(\mathbb J_m\tilde{\mathbb Q}_m)^{-1}$ is an exact right inverse with $\|\mathbb Q_m\|\le\|\tilde{\mathbb Q}_m\|/(1-\delta_m)$, and this is the $\mathbb L_m$ to use.  A rigorous application therefore needs enclosures of $\mathbb F_m(0)$ (quadrature and truncation of the ordinate sums), of $\mathbb J_m\tilde{\mathbb Q}_m$ (rounding) and of $\mathbb H_m$; for the Gaussian tests of Section~\CGref{subsec:rb-numerics} none of these enclosures is produced, and the label \texttt{FLOAT\_PASS} in the console output means only that the floating-point values satisfy the inequalities; for the B-spline tests of Section~\CGref{subsec:rb-num-verified} all of them are produced in ball arithmetic.  The theorem also does not guarantee that the corrected list $\boldsymbol\kappa^{(m,0)}+x_*$ is nondecreasing; this holds when $2\mathbb L_m\mathsf h_m$ is smaller than half the least gap of the moved ordinates, and otherwise the list is rearranged by Lemma~\CGref{lem:rb-matching}, which preserves the tube but changes the section.  Finally, certifying each stage separately gives no compatibility between stages: the solutions $x_*^{(m)}$ need not be restrictions of one another, and the uniform estimates that would be needed to pass to an infinite sequence (a stage-independent bound on $\mathbb L_m\mathbb H_m\mathsf h_m$ together with the compactness of Proposition~\CGref{prop:rb-compactness}) are not available.
\end{remark}

\begin{theorem}[Staged certificate implies route B$'$]\label{thm:rb-staged}
Let $C<\tau_1-2\pi$. Start at $\boldsymbol\kappa^{(1,0)}:=\boldsymbol\tau$ and, for $m\ge1$, let $\boldsymbol\kappa^{(m)}$ be the nondecreasing rearrangement of $\boldsymbol\kappa^{(m,0)}+x_*^{(m)}$ given by Theorem~\CGref{thm:rb-certificate}, and $\boldsymbol\kappa^{(m+1,0)}:=\boldsymbol\kappa^{(m)}$. Suppose \textup{(A$_m$)} holds for every $m$ and
\begin{equation}\label{eq:rb-B}
 \sum_{j\ge1}2\mathbb L_j\mathsf h_j=\sup_{m\ge1}\ \sum_{j\le m}2\mathbb L_j\mathsf h_j\le C .\tag{B}
\end{equation}
Then every $\boldsymbol\kappa^{(m)}$ is $C$-admissible and equals the clock beyond a finite index $\widetilde k^*_m$, $\mathsf s_m(\boldsymbol\kappa^{(m)},\mathsf a^\zeta)=0$ for every $m$, B$'(C)$ holds, and therefore (Corollary~\CGref{cor:rb-Bprime}) $\boldsymbol\kappa^{(m)}\to\boldsymbol\gamma$ coordinatewise and the Riemann hypothesis holds. In \textup{(A$_{m+1}$)} and \textup{(B)} the quantity $\mathsf h_{m+1}$ is $\max_{\varphi\in\mathscr T_{m+1}\setminus\mathscr T_m}|\langle\mathfrak R(\boldsymbol\kappa^{(m)},\mathsf a^\zeta),\varphi\rangle|$, the residual of the new tests only.
\end{theorem}

\begin{proof}
Induction on $m$. Suppose $\boldsymbol\kappa^{(m,0)}$ is $C_{m-1}$-admissible with $C_{m-1}:=\sum_{j<m}2\mathbb L_j\mathsf h_j$ ($C_0=0$) and equals the clock beyond a finite index. By Theorem~\CGref{thm:rb-certificate}, $\boldsymbol\kappa^{(m,0)}+x_*^{(m)}$ satisfies $|(\cdot)_k-\tau_k|\le C_{m-1}+2\mathbb L_m\mathsf h_m=C_m$ for all $k$, and only finitely many coordinates differ from the clock. Lemma~\CGref{lem:rb-matching} shows that its rearrangement $\boldsymbol\kappa^{(m)}$ satisfies $|\kappa^{(m)}_k-\tau_k|\le C_m\le C$ and equals the clock beyond some $\widetilde k^*_m$; it is nondecreasing, and $\kappa^{(m)}_1\ge\tau_1-C>2\pi$. So $\boldsymbol\kappa^{(m)}$ is $C$-admissible. The residual $\mathfrak R$ depends only on the multiset of ordinates, so $\langle\mathfrak R(\boldsymbol\kappa^{(m)},\mathsf a^\zeta),\varphi\rangle=\mathbb F_m(x_*^{(m)})_\varphi=0$ for all $\varphi\in\mathscr T_m$, that is, $\mathsf s_m=0$. Since $\mathscr T_m\subset\mathscr T_{m+1}$, the start of stage $m+1$ has zero residual on $\mathscr T_m$, which gives the stated form of $\mathsf h_{m+1}$. B$'(C)$ holds with $\mathsf s_m=0$, and Corollary~\CGref{cor:rb-Bprime}(a) gives the rest.
\end{proof}

\begin{remark}[What the certificate theorem does and does not do]\label{rem:rb-certificate-status}
Theorem~\CGref{thm:rb-staged} reduces the extension from one stage to the next to two computable quantities per stage, $\mathsf l_m=2\mathbb L_m^2\mathbb H_m\mathsf h_m$ and $2\mathbb L_m\mathsf h_m$. It requires no contraction factor $q$, no separate control of the perturbation of old tests, and no prior ``entry into the basin'': old tests are kept in $\mathscr T_{m+1}$ and are solved again exactly. The two conditions are not independent of the height: $\mathbb H_m\ge k^*_m\int|\varphi|u^2$ grows with the number of unknowns, and $\mathbb L_m$ is the $\ell^\infty$ norm of a right inverse of the Jacobian, whose growth is measured in Section~\CGref{subsec:rb-numerics}. The conditions (A$_m$) for all $m$ and (B) together imply the Riemann hypothesis, so they are not a weaker intermediate step; what changes is that each stage is a statement about one finite matrix and one finite residual vector, verifiable by interval arithmetic; Section~\CGref{subsec:rb-num-verified} does this for square sections. The restriction $C<\tau_1-2\pi=8.24\ldots$ only ensures $\kappa_1\ge2\pi$ and can be removed by treating the first coordinate separately.
\end{remark}

\begin{proposition}[Graded extension]\label{prop:rb-graded}
Let $\mathsf s_m\ge0$ satisfy $\mathsf s_{m+1}\le \mathrm q_m\mathsf s_m+\mathsf h_m$ with $0\le \mathrm q_m\le1$ and $\sum_m\mathsf h_m<\infty$, and assume
\begin{equation}\label{eq:rb-graded-tail}
 \prod_{j=L}^{m}\mathrm q_j\longrightarrow0\quad(m\to\infty)\qquad\text{for every }L\ge1 .
\end{equation}
Then $\mathsf s_m\to0$. Condition \CGeqref{eq:rb-graded-tail} holds if $\mathrm q_m\le\mathrm q<1$ for all large $m$, and more generally if $\sum_m(1-\mathrm q_m)=\infty$. If the lists $\boldsymbol\kappa^{(m)}$ producing $\mathsf s_m=\mathsf s_m(\boldsymbol\kappa^{(m)},\mathsf a^\zeta)$ are $C$-admissible sections for a fixed $C$, then B$'(C)$ holds.
\end{proposition}

\begin{proof}
Iterating, $\mathsf s_m\le\mathsf s_1\prod_{j<m}\mathrm q_j+\sum_{\ell<m}\mathsf h_\ell\prod_{\ell<j<m}\mathrm q_j$. Given $\epsilon>0$ choose $L$ with $\sum_{\ell\ge L}\mathsf h_\ell<\epsilon$. For $m>L$ and $\ell<L$, every $\mathrm q_j$ is at most $1$, so $\prod_{\ell<j<m}\mathrm q_j\le\prod_{j=L}^{m-1}\mathrm q_j$, and the same bound holds for $\prod_{j<m}\mathrm q_j$; hence the first term and the terms with $\ell<L$ are together at most $\bigl(\mathsf s_1+\sum_{\ell<L}\mathsf h_\ell\bigr)\prod_{j=L}^{m-1}\mathrm q_j$, which tends to $0$ by \CGeqref{eq:rb-graded-tail}. The terms with $\ell\ge L$ are at most $\sum_{\ell\ge L}\mathsf h_\ell<\epsilon$. Hence $\limsup_m\mathsf s_m\le\epsilon$, and $\mathsf s_m\to0$. The weaker hypothesis $\prod_{m\ge1}\mathrm q_m=0$ would not suffice: with $\mathrm q_1=0$ and $\mathrm q_m=1$ for $m\ge2$ the product vanishes, while $\mathsf s_{m+1}\le\mathsf s_m+\mathsf h_m$ allows $\mathsf s_m\to\mathsf s_2>0$. For the two sufficient conditions, $\prod_{j=L}^m\mathrm q_j\le\mathrm q^{\,m-\max(L,m_0)}$ when $\mathrm q_j\le\mathrm q$ for $j\ge m_0$, and $\prod_{j=L}^m\mathrm q_j\le\exp\bigl(-\sum_{j=L}^m(1-\mathrm q_j)\bigr)$ in general, since $\mathrm q\le e^{-(1-\mathrm q)}$ on $[0,1]$. The second claim is Definition~\CGref{def:rb-Bprime}.
\end{proof}

\begin{remark}[Tube budget, and the two geometries]\label{rem:rb-graded-status}
In Proposition~\CGref{prop:rb-graded} the tube $|\kappa^{(m)}_k-\tau_k|\le C$, uniform in $m$, is a separate hypothesis: summability of the coordinate increments $\sum_m|\kappa^{(m+1)}_k-\kappa^{(m)}_k|$ gives convergence of each coordinate but not that the limit stays in a prescribed tube, so a construction must carry a displacement budget $\sum_{m\ge m(k)}|\kappa^{(m+1)}_k-\kappa^{(m)}_k|\le C-|\kappa^{(m(k))}_k-\tau_k|$. For B$'(C)$ itself no convergence of the whole sequence is needed: compactness in the product of the intervals $[\tau_k-C,\tau_k+C]$ (Proposition~\CGref{prop:rb-compactness}) extracts a convergent subsequence, and uniqueness identifies the limit. A second point concerns norms. Theorem~\CGref{thm:rb-certificate} is stated in $\ell^\infty$, where $\mathbb H_m$ has the simple form \CGeqref{eq:rb-eta-H} because the Hessian of each $\mathsf c_\varphi$ is diagonal. A bound $\|\mathbb J_m^\dagger\|_{2\to2}\le1/\sigma_{\min}(\mathbb J_m)$ is a Euclidean statement; passing to $\ell^\infty$ costs factors $\sqrt{\mathrm t_m}$ and $\sqrt{K_m}$ unless the right inverse is controlled directly in $\ell^\infty_{\mathrm t_m}\to\ell^\infty_{K_m}$. The numerical constants of Section~\CGref{subsec:rb-numerics} are computed directly in $\ell^\infty$.
\end{remark}

\begin{heuristic}[Frame bound from exponential sums of the clock]\label{heur:rb-frame}
In the opposite regime $k^*_m\gg \mathrm t_m$ one may take $\mathbb Q_m:=\mathbb J_m^{\!T}(\mathbb J_m\mathbb J_m^{\!T})^{-1}$. If $\|\frac1{k^*_m}\mathbb J_m\mathbb J_m^{\!T}-D\|_{\infty\to\infty}\le\epsilon<d_-$ for a diagonal $D\ge d_-I$, and $|(\mathbb J_m)_{\varphi,k}|\le\mathrm a_J$, then $\mathbb L_m\le \mathrm t_m\mathrm a_J/(k^*_m(d_--\epsilon))=O(\mathrm t_m/k^*_m)$ and $\mathbb L_m^2\mathbb H_m=O(\mathrm t_m^2\mathrm B_0/k^*_m)$ with $\mathrm B_0:=\max_\varphi\int|\varphi|u^2$, so (A$_m$) becomes $k^*_m\gg \mathrm t_m^2\mathrm B_0\mathsf h_m$: a wide block compensates a moderate residual. In the point model $\varphi\approx\delta_{u_j}$ the entries of $\frac1{k^*_m}\mathbb J_m\mathbb J_m^{\!T}$ are $u_ju_\ell\cdot\frac1{k^*_m}\sum_{k\in\mathfrak B}\sin(\tau_ku_j)\sin(\tau_ku_\ell)$, where $\mathfrak B$ is the block of the $k^*_m$ interior indices, so the frame property reduces to cancellation in the exponential sums $\frac1{k^*_m}\sum_{k\in\mathfrak B}e^{i\mathrm v\tau_k}$ at the frequencies $\mathrm v=u_j\pm u_\ell\ne0$. On a block of $k^*_m$ consecutive clock points at height $T$ the spacing is $\Delta_T\approx2\pi/\log T$ with curvature $O((k^*_m)^2/(T\log^3T))$, which suggests $|\frac1{k^*_m}\sum_{k\in\mathfrak B}e^{i\mathrm v\tau_k}|\lesssim\frac{\log T}{k^*_m|\mathrm v|}+\frac{|\mathrm v|(k^*_m)^2}{T\log^3T}$ and, for tests on a grid of step $\Delta u$ in $[0,u^*]$ (so that $\Delta u\le|\mathrm v|\le2u^*$), the window $\frac{\log T}{\Delta u}\ll k^*_m\ll(T\log^3T/u^*)^{1/2}$. None of this is proved here, and the point model omits a decisive feature of smooth tests: a test of width $\mathrm w$ in $u$ has $|\partial_\kappa \mathsf c_\varphi(\kappa)|\lesssim\mathrm w^{-1}e^{-(\mathrm w\kappa)^2/2}$ for Gaussian $\varphi$, so ordinates with $\mathrm w\kappa\gg1$ are invisible to it and the corresponding columns of $\mathbb J_m$ are exponentially small. The underdetermined right inverse then has $\mathbb L_m$ dominated by the invisible columns; in the floating-point computations of Section~\CGref{subsec:rb-numerics} this choice produced $\mathbb L_m$ of size $10^{11}$ within four stages, and the stable regime was the opposite one, $\mathrm t_m\ge k^*_m$ with the unknowns restricted to the visible ordinates (a least-squares regime; the certificate theorem applies to its square subsystems only). The quasi-orthogonality heuristic is therefore a statement about tests that are not in $C_c^\infty(0,\infty)$, and it is not pursued further here.
\end{heuristic}

\subsubsection{Explicit lists in the tube: Gonek's symmetrised Euler products}\label{subsec:rb-gonek}

Gonek \cite{Gonek2012} introduced, for $\mathrm X\ge2$, the truncated Euler product $P_{\mathrm X}(s):=\exp\bigl(\sum_{n\le \mathrm X^2}\frac{\Lambda_{\mathrm X}(n)}{n^s\log n}\bigr)$ with $\Lambda_{\mathrm X}(n):=\Lambda(n)$ for $n\le \mathrm X$ and $\Lambda_{\mathrm X}(n):=\Lambda(n)\log(\mathrm X^2/n)/\log \mathrm X$ for $\mathrm X<n\le \mathrm X^2$, and the symmetrised function $\zeta_{\mathrm X}(s):=P_{\mathrm X}(s)+\chi(s)P_{\mathrm X}(1-s)$, where $\chi$ is the factor of the functional equation $\zeta(s)=\chi(s)\zeta(1-s)$. On the critical line $\chi(\frac12+it)=e^{-2i\vartheta(t)}$ and $P_{\mathrm X}(\frac12-it)=\overline{P_{\mathrm X}(\frac12+it)}$, so with
\begin{equation}\label{eq:rb-alphaX}
 \mathrm p_{\mathrm X}(t):=\arg P_{\mathrm X}\bigl(\tfrac12+it\bigr)=-\sum_{n\le \mathrm X^2}\frac{\Lambda_{\mathrm X}(n)\sin(t\log n)}{\sqrt n\log n},
\end{equation}
and the two constants
\begin{equation}\label{eq:rb-aux-16}
 \mathrm m_{\mathrm X}:=\sum_{n\le \mathrm X^2}\frac{\Lambda_{\mathrm X}(n)}{\sqrt n\log n},\qquad \mathrm n_{\mathrm X}:=\sum_{n\le \mathrm X^2}\frac{\Lambda_{\mathrm X}(n)}{\sqrt n},
\end{equation}
one has $\zeta_{\mathrm X}(\frac12+it)=2|P_{\mathrm X}(\frac12+it)|\,e^{-i\vartheta(t)}\cos\bigl(\vartheta(t)+\mathrm p_{\mathrm X}(t)\bigr)$. Since $P_{\mathrm X}$ has no zeros, the zeros of $\zeta_{\mathrm X}$ on the line are exactly the solutions of
\begin{equation}\label{eq:rb-gonek-phase}
 \varTheta_{\mathrm X}(t):=\vartheta(t)+\mathrm p_{\mathrm X}(t)\in\tfrac\pi2+\pi\mathbb Z .
\end{equation}
Gonek proves \cite[Thm.~6.2 of the arXiv version]{Gonek2012} that every zero of $\zeta_{\mathrm X}$ in $0\le\Re s\le1$ with imaginary part at least an absolute constant $C_0$ lies on $\Re s=\frac12$, unconditionally, and that under the Riemann hypothesis the zeros of $\zeta_{\mathrm X}$ approach those of $\zeta$ as $\mathrm X\to\infty$. The phase \CGeqref{eq:rb-gonek-phase} is the clock phase plus a truncated prime sum: $\NR(t)=1+\vartheta(t)/\pi+\mathrm r_\vartheta(t)$ with $|\mathrm r_\vartheta(t)|\le \mathrm c_r/t$ for $t\ge1$, by Stirling's formula \cite[\S12.33]{WhittakerWatson1927}, so \CGeqref{eq:rb-gonek-phase} reads $\NR(t)+\mathrm p_{\mathrm X}(t)/\pi\in\frac12+\mathbb Z+\mathrm r_\vartheta(t)$, to be compared with $\NR(\tau_k)=k-\frac12$.

\begin{lemma}[Gonek lists lie in a tube around the clock]\label{lem:rb-gonek-tube}
Fix $\mathrm X\ge2$ and put $\mathfrak t_{\mathrm X}:=\max\bigl(2\pi e^{2\mathrm n_{\mathrm X}+2},\,t_1\bigr)$, where $t_1$ is such that $|\vartheta'(t)-\frac12\log\frac t{2\pi}|\le\frac12$ for $t\ge t_1$. Then $\varTheta_{\mathrm X}$ is strictly increasing on $[\mathfrak t_{\mathrm X},\infty)$ with $\varTheta_{\mathrm X}'\ge\frac12$, the solutions of \CGeqref{eq:rb-gonek-phase} in $[\mathfrak t_{\mathrm X},\infty)$ form a strictly increasing sequence $\mathrm t_1^{(\mathrm X)}<\mathrm t_2^{(\mathrm X)}<\cdots$, each a simple zero of $\zeta_{\mathrm X}$, and there is an integer $\mathrm j_{\mathrm X}$ such that the list
\begin{equation}\label{eq:rb-gonek-list}
 \kappa^{(\mathrm X)}_k:=\begin{cases}\tau_k,&k\le \mathrm j_{\mathrm X},\\ \mathrm t^{(\mathrm X)}_{k-\mathrm j_{\mathrm X}},&k>\mathrm j_{\mathrm X},\end{cases}
\end{equation}
is $\mathrm c_{\mathrm X}$-admissible for some $\mathrm c_{\mathrm X}<\infty$ and, possibly after a finite permutation, nondecreasing: if the junction at $k=\mathrm j_{\mathrm X}$ breaks monotonicity, the list is replaced by its nondecreasing rearrangement, which keeps the multiset of values, reassigns the indices of finitely many entries (those not exceeding $\tau_{\mathrm j_{\mathrm X}}+\mathrm c_{\mathrm X}$) and preserves the tube bound by Lemma~\CGref{lem:rb-matching}; and
\begin{equation}\label{eq:rb-gonek-tube}
 \bigl|\kappa^{(\mathrm X)}_k-\tau_k\bigr|\le\frac{2\mathrm m_{\mathrm X}+o(1)}{\log(\tau_k/2\pi)}\qquad(k\to\infty).
\end{equation}
\end{lemma}

\begin{proof}
$|\mathrm p_{\mathrm X}'(t)|=|\sum_{n\le \mathrm X^2}\Lambda_{\mathrm X}(n)n^{-1/2}\cos(t\log n)|\le \mathrm n_{\mathrm X}$ and $|\mathrm p_{\mathrm X}|\le \mathrm m_{\mathrm X}$. For $t\ge \mathfrak t_{\mathrm X}$, $\vartheta'(t)\ge\frac12\log\frac t{2\pi}-\frac12\ge \mathrm n_{\mathrm X}+\frac12$, so $\varTheta_{\mathrm X}'\ge\frac12$. Hence $\varTheta_{\mathrm X}$ is a strictly increasing bijection of $[\mathfrak t_{\mathrm X},\infty)$ onto $[\varTheta_{\mathrm X}(\mathfrak t_{\mathrm X}),\infty)$, and the solutions of \CGeqref{eq:rb-gonek-phase} are the preimages of the points $\frac\pi2+\pi j$, $j\ge j_0$, where $j_0$ is the least integer with $\frac\pi2+\pi j_0\ge\varTheta_{\mathrm X}(\mathfrak t_{\mathrm X})$; they are simple zeros because $\frac d{dt}\cos\varTheta_{\mathrm X}=-\varTheta_{\mathrm X}'\sin\varTheta_{\mathrm X}\ne0$ there. The $i$-th solution $\mathrm t^{(\mathrm X)}_i$ satisfies $\varTheta_{\mathrm X}(\mathrm t^{(\mathrm X)}_i)=\frac\pi2+\pi(j_0+i-1)=\pi(j_0+i-\frac12)$; we set $\mathrm j_{\mathrm X}:=j_0+1$, so that $k=\mathrm j_{\mathrm X}+i$ gives $\varTheta_{\mathrm X}(\kappa^{(\mathrm X)}_k)=\pi(k-\frac32)$ for $k>\mathrm j_{\mathrm X}$. Since $\NR=1+\vartheta/\pi+\mathrm r_\vartheta$, the constant $1$ raises the count by one unit and
\begin{equation}\label{eq:rb-aux-15}
 \NR\bigl(\kappa^{(\mathrm X)}_k\bigr)=k-\tfrac12-\frac{\mathrm p_{\mathrm X}(\kappa^{(\mathrm X)}_k)}\pi+\mathrm r_\vartheta\bigl(\kappa^{(\mathrm X)}_k\bigr)\qquad(k>\mathrm j_{\mathrm X}),
\end{equation}
so, since $\NR(\tau_k)=k-\frac12$,
\begin{equation}\label{eq:rb-aux-17}
 \bigl|\NR(\kappa^{(\mathrm X)}_k)-\NR(\tau_k)\bigr|\le\frac{\mathrm m_{\mathrm X}}\pi+\frac{\mathrm c_r}{\kappa^{(\mathrm X)}_k}
\end{equation}
(the choice $\mathrm j_{\mathrm X}=j_0+1$, and not $j_0$, is what makes the two counts agree; in the computations of Section~\CGref{subsec:rb-numerics} the offset was determined by matching). The right-hand side is bounded, and $\NR(2t)-\NR(t)\to\infty$, so $\kappa^{(\mathrm X)}_k/\tau_k\to1$. By the mean value theorem, $|\kappa^{(\mathrm X)}_k-\tau_k|\le(\mathrm m_{\mathrm X}/\pi+o(1))/\NR'(\mathrm s_k)$ with $\mathrm s_k$ between $\kappa^{(\mathrm X)}_k$ and $\tau_k$, and $\NR'(\mathrm s_k)=\frac1{2\pi}\log\frac{\tau_k}{2\pi}(1+o(1))$. This is \CGeqref{eq:rb-gonek-tube}; it gives $|\kappa^{(\mathrm X)}_k-\tau_k|\to0$, so the displacements are bounded, the list is eventually nondecreasing, and finitely many initial indices (where the list is the clock) are harmless. If monotonicity of the concatenated list fails at the junction $k=\mathrm j_{\mathrm X}$, one replaces the prefix $\tau_k$, $k\le \mathrm j_{\mathrm X}$, by the nondecreasing rearrangement of Lemma~\CGref{lem:rb-matching}, which keeps the tube.
\end{proof}

\begin{remark}[What the Gonek lists give and what remains]\label{rem:rb-gonek-status}
The lists $\boldsymbol\kappa^{(\mathrm X)}$ use the primes up to $\mathrm X^2$ and the clock, and nothing about the zeros of $\zeta$. Their prime side is built in: formally, $\sum_k\delta(t-\kappa^{(\mathrm X)}_k)=\frac{\varTheta_{\mathrm X}'(t)}\pi\bigl[1+2\sum_{j\ge1}(-1)^j\cos(2j\varTheta_{\mathrm X}(t))\bigr]$ where $\varTheta_{\mathrm X}'>0$, and the $j=0$ term carries $\mathrm p_{\mathrm X}'/\pi=-\frac1\pi\sum_{n\le \mathrm X^2}\Lambda_{\mathrm X}(n)n^{-1/2}\cos(t\log n)$, whose cosine transform is a sum of masses $\Lambda_{\mathrm X}(n)/(2\sqrt n)$ at $u=\log n$, equal to $\mathsf a^\zeta_n$ for $n\le \mathrm X$. Two statements would make them a solution of B$'$: \textup{(G1)} a tube uniform in $\mathrm X$, $\sup_{\mathrm X\ge\mathrm X_0}\sup_k|\kappa^{(\mathrm X)}_k-\tau_k|\le C$, which by \CGeqref{eq:rb-gonek-tube} concerns the size of $\mathrm p_{\mathrm X}$ at its own crossing points and not its supremum ($\mathrm m_{\mathrm X}\asymp\sqrt{\mathrm X}/\log^2\mathrm X$ by the prime number theorem, and $\mathrm n_{\mathrm X}\asymp\sqrt{\mathrm X}$; the supremum is useless for this); and \textup{(G2)} $\langle\mathfrak R(\boldsymbol\kappa^{(\mathrm X)},\mathsf a^\zeta),\varphi\rangle\to0$ as $\mathrm X\to\infty$ for each test $\varphi$, whose content is the vanishing of the harmonics $j\ge1$ against $\varphi$. Neither is proved here. The monotonicity threshold $\mathfrak t_{\mathrm X}$ of Lemma~\CGref{lem:rb-gonek-tube} grows like $e^{2\mathrm n_{\mathrm X}}$, so the global construction cannot rely on monotonicity of $\varTheta_{\mathrm X}$ at fixed height as $\mathrm X\to\infty$; Gonek's theorem on the location of the zeros does not need it. For fixed $\mathrm X$ the tube constant in \CGeqref{eq:rb-gonek-tube} is $2\mathrm m_{\mathrm X}$, and the proof only uses $\sup_t|\mathrm p_{\mathrm X}(t)|\le\mathrm m_{\mathrm X}$; by hand $\mathrm m_3=0.707+0.577+0.185+0.239+0.087+0.013=1.81$, while the observed suprema of $|\mathrm p_{\mathrm X}|$ on the computed ranges are $1.42$ ($\mathrm X=3$), $3.02$ ($\mathrm X=8$) and $5.98$ ($\mathrm X=20$), which are the constants that matter numerically. Numerically (Section~\CGref{subsec:rb-numerics}) the lists with $3\le \mathrm X\le20$ lie within $0.19$ to $0.035$ mean spacings of the zeta ordinates for $k\le400$, whereas the clock is within $0.66$ spacings.
\end{remark}

%======================================================================
\subsection{Route B on a window: unknown lists, a Bernstein form, and minimizers}\label{subsec:rb-window}

This subsection combines route B with the exponential sums of Section~\CGref{sec:ordinate-sums}. The ordinates are unknowns, as in route B; the prime side is Cram\'er's function on one window, as in Theorem~\CGref{thm:window-RH}. Uniqueness then needs neither the masses nor Hamburger's theorem: it comes from Lemma~\CGref{lem:edge} and from the following lemma, in which the positivity sits on the side of the unknowns.

\begin{lemma}[Laplace transforms against the zeros]\label{lem:rb-peel}
Let $\mu$ be a positive Borel measure on $\mathbb R$ with $\int e^{-yt}\,d\mu(t)<\infty$ for $y>y_1$, and suppose that
\begin{equation}\label{eq:rb-peel}
 \int_{\mathbb R}e^{-yt}\,d\mu(t)=\sum_{k\ge1}e^{iyd_k}e^{-y\gamma_k}\qquad(y>y_1).
\end{equation}
Then all nontrivial zeros of $\zeta$ lie on $\Re s=\frac12$, and $\mu=\sum_k\delta_{\gamma_k}$.
\end{lemma}

\begin{proof}
The right side converges absolutely for $y>0$ and is real, because the zeros with $d_k\ne0$ come in pairs with opposite $d_k$. Let $t_1<t_2<\cdots$ be the distinct values of the $\gamma_k$, $m_j$ the number of $k$ with $\gamma_k=t_j$, and $P_j(y):=\sum_{k:\gamma_k=t_j}e^{iyd_k}=\sum_\omega\mathrm a_{j,\omega}e^{iy\omega}$, a trigonometric polynomial whose coefficients $\mathrm a_{j,\omega}$, the number of $k$ with $\gamma_k=t_j$ and $d_k=\omega$, are nonnegative integers with $\sum_\omega\mathrm a_{j,\omega}=m_j$. For $y\ge y_2>\max(y_1,0)$, $\sum_{\gamma_k>t_j}e^{-y(\gamma_k-t_j)}\le e^{-(y-y_2)(t_{j+1}-t_j)}\sum_ke^{-y_2(\gamma_k-t_j)}$, which tends to $0$ as $y\to\infty$.

We show by induction on $j$ that $\mu_j:=\mu-\sum_{i<j}m_i\delta_{t_i}$ is a positive measure supported in $[t_j,\infty)$ and that $d_k=0$ whenever $\gamma_k<t_j$; for $j=1$ the first claim is about the support of $\mu$. By \CGeqref{eq:rb-peel} and the induction hypothesis, $\int e^{-yt}d\mu_j(t)=\sum_{\gamma_k\ge t_j}e^{iyd_k}e^{-y\gamma_k}$ for $y>y_1$, and the modulus of the right side is at most $Ce^{-yt_j}$ for $y\ge y_2$. \emph{Support.} Let $s:=\inf\operatorname{supp}\mu_j$ and suppose $s<t_j$; for $s<s'<t_j$, $\mu_j((-\infty,s'])>0$, and $\int e^{-yt}d\mu_j\ge e^{-ys'}\mu_j((-\infty,s'])$, which contradicts the bound $Ce^{-yt_j}$ for large $y$. So $\operatorname{supp}\mu_j\subset[t_j,\infty)$. \emph{The lowest group.} Then $e^{yt_j}\int e^{-yt}d\mu_j(t)=\int e^{-y(t-t_j)}d\mu_j(t)$ decreases to $\mu_j(\{t_j\})$ as $y\to\infty$, while $e^{yt_j}\sum_{\gamma_k\ge t_j}e^{iyd_k}e^{-y\gamma_k}=P_j(y)+o(1)$. A trigonometric polynomial with a limit $c$ at infinity has all its coefficients at nonzero frequencies equal to $0$ and constant term $c$, because the mean of $|P_j-c|^2$ over $[Y,2Y]$ tends to $\sum_{\omega\ne0}\mathrm a_{j,\omega}^2+(\mathrm a_{j,0}-c)^2$. Hence $\mathrm a_{j,\omega}=0$ for $\omega\ne0$, that is, $d_k=0$ when $\gamma_k=t_j$, and $\mu_j(\{t_j\})=\mathrm a_{j,0}=m_j$. So $\mu_{j+1}=\mu_j-m_j\delta_{t_j}$ is positive, supported in $[t_j,\infty)$, without atom at $t_j$, and the support argument applied to it with $t_{j+1}$ in place of $t_j$ completes the induction. Since $t_j\to\infty$, every $d_k$ vanishes and $\mu=\sum_jm_j\delta_{t_j}=\sum_k\delta_{\gamma_k}$.
\end{proof}

\begin{theorem}[Route B on one window, with an unknown list]\label{thm:rb-window-unknown}
Let $J\subset(0,\infty)$ be an open interval containing no point $\log p^m$, for instance $J=(\log2,\log3)$. Let $\boldsymbol\kappa=(\kappa_k)_{k\ge1}$ be a nondecreasing sequence of real numbers with $\kappa_1>0$ and $N_{\boldsymbol\kappa}(T)\ll(1+T)^A$ for some $A$. If
\begin{equation}\label{eq:rb-window-unknown}
 \sum_{k\ge1}\widehat{\mathsf f}(\kappa_k)=\int_J\mathsf f(u)\,e^{-u/2}V_+(u)\,du\qquad\text{for every real }\mathsf f\in C_c^\infty(J),
\end{equation}
then all nontrivial zeros of $\zeta$ lie on $\Re s=\frac12$ and $\boldsymbol\kappa=\boldsymbol\gamma$. Conversely, under RH the list $\boldsymbol\gamma$ satisfies \CGeqref{eq:rb-window-unknown} (Theorem~\CGref{thm:window-RH}).
\end{theorem}

\begin{proof}
Put $\mathscr Z_{\boldsymbol\kappa}(w):=\sum_ke^{iw\kappa_k}$ for $\Im w>0$. By partial summation $|\mathscr Z_{\boldsymbol\kappa}(u+iy)|\le\sum_ke^{-y\kappa_k}\ll y^{-A}\le y^{-\lceil A\rceil}$ for $0<y\le1$, and $\int\mathscr Z_{\boldsymbol\kappa}(u+iy)\mathsf f(u)\,du=\sum_ke^{-y\kappa_k}\widehat{\mathsf f}(\kappa_k)\to\sum_k\widehat{\mathsf f}(\kappa_k)$ as $y\downarrow0$, by \CGeqref{eq:f-hat-decay} and dominated convergence. The function $\mathscr Z_\xi=e^{-w/2}V_+$ is holomorphic, hence bounded, on a neighbourhood of $J'\times[0,1]$ for every compact interval $J'\subset J$ (Proposition~\CGref{prop:cramer-zero-free}). By \CGeqref{eq:rb-window-unknown} and linearity, $F:=\mathscr Z_{\boldsymbol\kappa}-\mathscr Z_\xi$ satisfies the hypotheses of Lemma~\CGref{lem:edge} on the interior of $J'$, so $F\equiv0$ there, and, the half-plane being connected, $\mathscr Z_{\boldsymbol\kappa}=\mathscr Z_\xi$ on $\Im w>0$. At $w=iy$ this is \CGeqref{eq:rb-peel} with $\mu=\sum_k\delta_{\kappa_k}$, because $\mathscr Z_\xi(iy)=\sum_ke^{iy(\rho_k-\frac12)}=\sum_ke^{iyd_k}e^{-y\gamma_k}$. Lemma~\CGref{lem:rb-peel} gives RH and $\sum_k\delta_{\kappa_k}=\sum_k\delta_{\gamma_k}$; both lists are nondecreasing, so $\boldsymbol\kappa=\boldsymbol\gamma$.
\end{proof}

\begin{theorem}[A Bernstein form of the Riemann hypothesis]\label{thm:rb-bernstein}
The function $y\mapsto\mathscr Z_\xi(iy)=e^{-iy/2}V_+(iy)=\sum_ke^{iyd_k}e^{-y\gamma_k}$ is real on $(0,\infty)$ and has the zero-free representation \CGeqref{eq:cramer-zero-free} ($iy\in\Omega_\theta$). The following are equivalent:
\begin{enumerate}[label=\textup{(\roman*)}]
\item RH;
\item $y\mapsto\mathscr Z_\xi(iy)$ is completely monotone on $(0,\infty)$: $(-1)^n\frac{d^n}{dy^n}\mathscr Z_\xi(iy)\ge0$ for all $n\ge0$ and $y>0$;
\item $y\mapsto\mathscr Z_\xi(iy)$ is completely monotone on $(a,\infty)$ for some $a\ge0$.
\end{enumerate}
\end{theorem}

\begin{proof}
(i)$\Rightarrow$(ii): under RH, $\mathscr Z_\xi(iy)=\sum_ke^{-y\gamma_k}$ and $(-1)^n\frac{d^n}{dy^n}\mathscr Z_\xi(iy)=\sum_k\gamma_k^ne^{-y\gamma_k}>0$. (ii)$\Rightarrow$(iii) is trivial. (iii)$\Rightarrow$(i): by Bernstein's theorem \cite{Bernstein1929}, \cite[Ch.~IV, \S12]{Widder1941}, $x\mapsto\mathscr Z_\xi(i(a+x))$ is the Laplace transform of a positive measure $\nu$ on $[0,\infty)$, so $\mathscr Z_\xi(iy)=\int e^{-yt}e^{at}d\nu(t)$ for $y>a$, and Lemma~\CGref{lem:rb-peel} with $\mu:=e^{at}\nu$ gives RH.
\end{proof}

\begin{remark}[Where the positivity sits]\label{rem:rb-window-positivity}
In Proposition~\CGref{prop:imaginary-frequency} one imaginary frequency suffices, because every term of the defect $\mathfrak X(iy_0)$ is nonnegative. For an unknown list there is no such sign, and Theorems~\CGref{thm:rb-window-unknown} and~\CGref{thm:rb-bernstein} use all frequencies of a half-line; the positivity is that of the unknown measure $\mu$, which Lemma~\CGref{lem:rb-peel} turns, group by group from the lowest ordinate up, into the vanishing of the abscissas. Theorem~\CGref{thm:rb-window-unknown} is route B on a single window: no masses, no tube (only polynomial counting), no Hamburger theorem, and existence is again equivalent to RH. By the argument of Proposition~\CGref{prop:rb-nonfinitary}, the zeta list satisfies \CGeqref{eq:rb-window-unknown} up to the contribution of the zeros above $H_0$, so no finite computation at a coarser precision decides it either. Theorem~\CGref{thm:rb-bernstein} is a positivity criterion in one real variable whose function has a representation with no sum over the zeros (the primes, $\xi'/\xi$ on a real segment and $\Gamma$, Proposition~\CGref{prop:cramer-zero-free}), in the family of the criteria of Weil and Li; we have not searched the literature for it, and novelty is not claimed.
\end{remark}

\emph{Minimizers of the residual on a window.} The second combination uses compactness, as route B does, and the completeness of exponential systems on bounded intervals in place of stability estimates. Fix $C\ge C_\eta$ and the box
\begin{equation}\label{eq:rb-box-window}
 \mathfrak T_C:=\prod_{k\ge1}\bigl[\max(2\pi,\tau_k-C),\ \tau_k+C\bigr],
\end{equation}
compact and metrizable in the product topology; the residual $\mathfrak R(\boldsymbol\kappa,\mathsf a^\zeta)$ is defined by \CGeqref{eq:rb-residual} for every $\boldsymbol\kappa\in\mathfrak T_C$ (it depends only on the multiset of the coordinates), and $\boldsymbol\gamma\in\mathfrak T_C$. Fix a window $W=(u_1,u_2)$, $0<u_1<u_2$, real test functions $\varphi_j\in C_c^\infty(W)$ whose span is dense in the real space $C_c^\infty(W)$, the bounds $\mathsf B_j:=\sup_{\mathfrak T_C}|\langle\mathfrak R(\boldsymbol\kappa,\mathsf a^\zeta),\varphi_j\rangle|<\infty$, and weights $\mathrm w_j>0$ with $\sum_j\mathrm w_j(1+\mathsf B_j)^2(1+\|\varphi_j\|_{C^2})<\infty$. Put
\begin{equation}\label{eq:rb-EW}
 \mathfrak E_W(\boldsymbol\kappa):=\sum_{j\ge1}\mathrm w_j\,\langle\mathfrak R(\boldsymbol\kappa,\mathsf a^\zeta),\varphi_j\rangle^2,\qquad
 \mathscr P_{\boldsymbol\kappa}:=\sum_{j\ge1}\mathrm w_j\,\langle\mathfrak R(\boldsymbol\kappa,\mathsf a^\zeta),\varphi_j\rangle\,\varphi_j .
\end{equation}

\begin{lemma}[Minimizers on a window]\label{lem:rb-minimizer}
\begin{enumerate}[label=\textup{(\alph*)}]
\item $\mathfrak E_W$ is continuous on $\mathfrak T_C$ and attains its minimum.
\item Let $\boldsymbol\kappa^*$ be a minimizer, and put
\begin{equation}\label{eq:rb-G-minimizer}
 \mathscr G(z):=\int_{u_1}^{u_2}u\,\mathscr P_{\boldsymbol\kappa^*}(u)\sin(zu)\,du\qquad(z\in\mathbb C).
\end{equation}
Then $\mathfrak E_W(\boldsymbol\kappa^*)=\langle\mathfrak R(\boldsymbol\kappa^*,\mathsf a^\zeta),\mathscr P_{\boldsymbol\kappa^*}\rangle$, and $\mathscr G(\kappa^*_k)=0$ for every $k$ for which $\kappa^*_k$ is an interior point of its interval in \CGeqref{eq:rb-box-window}.
\item \emph{Dichotomy.} Either $\mathfrak E_W(\boldsymbol\kappa^*)=0$, that is, $\mathfrak R(\boldsymbol\kappa^*,\mathsf a^\zeta)=0$ in $\mathcal D'(W)$; or there is $c$ such that for every $r\ge1$ the interior coordinates of $\boldsymbol\kappa^*$ take at most $(u_2/\log2)\,r+c$ distinct values in $(0,r]$.
\end{enumerate}
\end{lemma}

\begin{proof}
(a) Each pairing is continuous on $\mathfrak T_C$ (proof of Proposition~\CGref{prop:rb-compactness}) and bounded by $\mathsf B_j$, so the series converges uniformly; a continuous function on a compact space attains its minimum. (b) The series for $\mathscr P_{\boldsymbol\kappa^*}$ converges in $C^2$, so $\mathscr P_{\boldsymbol\kappa^*}\in C^2(\mathbb R)$ is supported in $[u_1,u_2]$ and vanishes with its first two derivatives at the ends. The residual is a distribution of order at most $2$ on $(0,\infty)$, because $\sum_k\cos(\kappa_ku)$ is the second derivative of the uniformly convergent series $-\sum_k\kappa_k^{-2}\cos(\kappa_ku)$, the masses form a measure and $\mathsf k$ is smooth; it is therefore continuous for this convergence, which gives the first claim. For fixed other coordinates, $t\mapsto\langle\mathfrak R,\varphi_j\rangle$ is $C^1$ in the coordinate $\kappa_k=t$ with derivative $\mathsf c'_{\varphi_j}(t)=-\int\varphi_j(u)u\sin(tu)\,du$, bounded by $u_2\|\varphi_j\|_1$; the weights make the differentiated series converge uniformly, so $\partial_{\kappa_k}\mathfrak E_W=-2\int u\mathscr P_{\boldsymbol\kappa}(u)\sin(\kappa_ku)\,du$. At an interior coordinate of a minimizer this derivative vanishes. (c) If $\mathfrak E_W(\boldsymbol\kappa^*)>0$, then $\mathscr P_{\boldsymbol\kappa^*}\ne0$ by (b), so the odd extension of $u\mathscr P_{\boldsymbol\kappa^*}$ is a nonzero function and $\mathscr G\not\equiv0$. $\mathscr G$ is odd, entire, and $|\mathscr G(z)|\le\|u\mathscr P_{\boldsymbol\kappa^*}\|_1e^{u_2|\Im z|}$. Write $\mathscr G(z)=z^m\mathscr H(z)$ with $\mathscr H(0)\ne0$. Jensen's formula gives, for $r\ge1$, $n_{\mathscr H}(r)\log2\le\int_r^{2r}n_{\mathscr H}(t)\frac{dt}t\le\log\max_{|z|=2r}|\mathscr H|-\log|\mathscr H(0)|\le2u_2r+\log\frac{\|u\mathscr P_{\boldsymbol\kappa^*}\|_1}{|\mathscr H(0)|}$. The zeros of $\mathscr G$ are symmetric about $0$, so $\mathscr G$ has at most $\frac12n_{\mathscr H}(r)\le(u_2/\log2)r+c$ zeros in $(0,r]$, and by (b) every interior coordinate of $\boldsymbol\kappa^*$ is one of them. If $\mathfrak E_W(\boldsymbol\kappa^*)=0$, every pairing $\langle\mathfrak R,\varphi_j\rangle$ vanishes, and so does $\mathfrak R$ on $W$, by density.
\end{proof}

\begin{theorem}[Window solvability]\label{thm:rb-window-equivalence}
Let $C\ge C_\eta$. The following are equivalent:
\begin{enumerate}[label=\textup{(\roman*)}]
\item RH;
\item for every $U>1$ there is $\boldsymbol\kappa\in\mathfrak T_C$ with $\mathfrak R(\boldsymbol\kappa,\mathsf a^\zeta)=0$ in $\mathcal D'(U^{-1},U)$;
\item for every $U>1$, $\min_{\mathfrak T_C}\mathfrak E_{(U^{-1},U)}=0$, for one (equivalently every) admissible choice of tests and weights.
\end{enumerate}
In particular RH holds if, for every $U>1$, some minimizer of $\mathfrak E_{(U^{-1},U)}$ has interior coordinates taking more than $(U/\log2+\delta)\,r$ distinct values in $(0,r]$ for some $\delta>0$ and arbitrarily large $r$.
\end{theorem}

\begin{proof}
(i)$\Rightarrow$(ii): take $\boldsymbol\kappa=\boldsymbol\gamma$, which lies in $\mathfrak T_C$ by Lemma~\CGref{lem:bounded-displacement} and has $\mathfrak R(\boldsymbol\gamma,\mathsf a^\zeta)=0$ by Theorem~\CGref{thm:distributional-EF}. (ii)$\Leftrightarrow$(iii): by Lemma~\CGref{lem:rb-minimizer}(a) and density. (ii)$\Rightarrow$(i): let $\boldsymbol\kappa^{(m)}\in\mathfrak T_C$ solve (ii) with $U=m$. A subsequence converges coordinatewise to some $\boldsymbol\kappa^\infty\in\mathfrak T_C$. For $\varphi\in C_c^\infty(0,\infty)$, $\operatorname{supp}\varphi\subset(m^{-1},m)$ for all large $m$, so $\langle\mathfrak R(\boldsymbol\kappa^{(m)},\mathsf a^\zeta),\varphi\rangle=0$ along the subsequence, and continuity of the pairing gives $\mathfrak R(\boldsymbol\kappa^\infty,\mathsf a^\zeta)=0$ in $\mathcal D'(0,\infty)$. The nondecreasing rearrangement of $\boldsymbol\kappa^\infty$ is $C$-admissible (Lemma~\CGref{lem:rb-matching}; all coordinates are at least $2\pi$), and Proposition~\CGref{prop:rb-prescribed-uniqueness} gives RH. The last statement follows from Lemma~\CGref{lem:rb-minimizer}(c) with $u_2=U$.
\end{proof}

\begin{remark}[What the dichotomy leaves open]\label{rem:rb-minimizer-open}
Lemma~\CGref{lem:rb-minimizer} uses no stability estimate: the first-order condition at a minimizer is a completeness statement for $\{u\sin(\kappa_ku)\}$ on a bounded interval, and the clock density $(\log r)/2\pi$ per unit length is far above the density $u_2/\log2$ that a nonzero function of exponential type $u_2$ can absorb (Jensen's formula; the sharp constant is $u_2/\pi$, by the Cartwright--Levinson theorem). Theorem~\CGref{thm:rb-window-equivalence} therefore replaces the existence question of route B by a degeneracy question: whether minimizers can saturate the box, or collapse their interior coordinates onto a set of linear density. If RH fails, then for large $U$ every minimizer is degenerate in this sense, while $\min\mathfrak E_{(U^{-1},U)}\le\mathfrak E_{(U^{-1},U)}(\boldsymbol\gamma)$, which is of the size of Proposition~\CGref{prop:rb-nonfinitary}; at $\boldsymbol\gamma$ all coordinates are interior when $C>\sup_k|\gamma_k-\tau_k|$, and the gradient does not vanish: the identity $\mathfrak E_W=\langle\mathfrak R,\mathscr P_{\boldsymbol\kappa}\rangle$ and the formula for $\partial_{\kappa_k}\mathfrak E_W$ in the proof of Lemma~\CGref{lem:rb-minimizer} hold at every $\boldsymbol\kappa\in\mathfrak T_C$, so if $\mathfrak E_W(\boldsymbol\gamma)>0$ the function \CGeqref{eq:rb-G-minimizer} built at $\boldsymbol\gamma$ vanishes at $O(r)$ points of $(0,r]$, while the ordinates take $\gg r\log r$ distinct values, because a positive proportion of the zeros is simple \cite{Conrey1989}. Under RH and for such $C$, $\boldsymbol\gamma$ is a non-degenerate minimizer, for the same reason. Whether degenerate minimizers can be excluded directly is open; it is not a statement of the form of (D).
\end{remark}

\begin{heuristic}[Window size and zero-free regions]\label{heur:rb-window-zerofree}
Theorem~\CGref{thm:rung-zero-sum} gives, without hypothesis, $\mathsf k-\sum_n\mathsf a^\zeta_n\delta_{\log n}=\sum_k\cosh(d_ku)\cos(\gamma_ku)$ in $\mathcal D'(0,\infty)$. So a solution $\boldsymbol\kappa\in\mathfrak T_C$ on the window $(U^{-1},U)$ satisfies $\sum_k\cos(\kappa_ku)=\sum_k\cosh(d_ku)\cos(\gamma_ku)$ there. Test it with $\psi(u)\cos(\gamma_0u)$, $\psi\ge0$ supported in $[U-\mathrm w,U)$, where $\frac12+d_0+i\gamma_0$ is a zero off the line. The pair of zeros at $\gamma_0$ contributes $\int\psi(u)\cosh(d_0u)\,du$; the left side is at most $\int\psi$ times the number of $\kappa_k$ within about $1/\mathrm w$ of $\gamma_0$, which the box bounds by about $(C+\frac1{\mathrm w})\frac{\log\gamma_0}\pi$. This suggests
\begin{equation}\label{eq:rb-window-zerofree}
 \cosh\bigl(d_0(U-\mathrm w)\bigr)\lesssim\Bigl(C+\frac1{\mathrm w}\Bigr)\frac{\log\gamma_0}{2\pi},\qquad\text{that is,}\qquad|d_0|\lesssim\frac{\log(C\log\gamma_0)}U :
\end{equation}
solvability on the window $(U^{-1},U)$ would exclude zeros at distance more than about $\log(C\log\gamma)/U$ from the critical line; for fixed $U$ this is informative only for $\log(C\log\gamma)<U/2$, and graded by $U$ it would give statements between the verified zeros and RH. Three points are not proved. First, the on-line zeros near $\gamma_0$ and the $\kappa_k$ near $\gamma_0$ contribute terms $\frac12\int\psi(u)\cos((\gamma_k-\gamma_0)u)\,du$ of indefinite sign, since a bump $\psi$ supported near $U$ has no positive Fourier transform. Second, other zeros off the line with ordinates within $1/\mathrm w$ of $\gamma_0$ could cancel the main term; a lower bound of Tur\'an type for exponential sums would be needed, and with it the conclusion weakens to $|d_0|\ll\log\gamma_0/U$. Third, the converse, that such a zero-free region allows a solution on the window, is not constructed. With $H_0=3\cdot10^{12}$, $C=1$ and $\mathrm w=1$, the right side of \CGeqref{eq:rb-window-zerofree} is about $9.1$ at $\gamma_0=H_0$, so the bound says nothing for $|d_0|<\frac12$ unless $U\gtrsim6.8$: window solvability up to $U\approx7$ is consistent with the verified zeros, and beyond it would carry information on the zeros above $H_0$. The stages computed in Section~\CGref{subsec:rb-numerics} use tests supported in $(0,6.4]$.
\end{heuristic}

\subsection{An a priori bound for quasi-solutions}\label{subsec:rb-apriori}

Throughout this section $(\boldsymbol\kappa,\mathsf a)$ is a finite section in $\mathbb K(C,\mathbb M)$ with parameters $k^*,n^*$.

\begin{definition}[Laplace residual]\label{def:rb-laplace-residual}
For $\Re s>1$ put
\begin{equation}\label{eq:rb-Rhat}
 \widehat{\mathsf R}(s):=\tfrac12\,\frac{P_{\boldsymbol\kappa}'}{P_{\boldsymbol\kappa}}(s)-\tfrac12\,\frac{G'}{G}(s)+\sum_{n=2}^{n^*}\mathsf a_n\,n^{\frac12-s}.
\end{equation}
For $T\ge0$ put
\begin{equation}\label{eq:rb-E}
 \mathbb E(T):=\int_{-T}^{T}\bigl|\widehat{\mathsf R}(2+iv)\bigr|\,dv,\qquad
 \upsilon:=\sup_{T\ge2}\frac{\mathbb E(T)}{\log T}\in[0,\infty].
\end{equation}
Since $\int_0^\infty e^{-(s-\frac12)u}\cos(\kappa u)\,du=\frac{s-\frac12}{(s-\frac12)^2+\kappa^2}=\frac12\frac d{ds}\log\bigl(1+\frac{q(s)}{\frac14+\kappa^2}\bigr)$ and $\int_0^\infty e^{-(s-\frac12)u}\delta(u-\log n)\,du=n^{\frac12-s}$, $\widehat{\mathsf R}$ is formally the Laplace transform of the residual distribution of \CGeqref{eq:rb-rung}, with the normalisation of $\mathsf k$ fixed in \CGref{eq:transport-defs}; this is made precise in Lemma~\CGref{lem:rb-laplace}. We take \CGeqref{eq:rb-Rhat} as the definition. Because the coefficients are real, $\widehat{\mathsf R}(\bar s)=\overline{\widehat{\mathsf R}(s)}$, so $\mathbb E(T)=2\int_0^T|\widehat{\mathsf R}(2+iv)|\,dv$. We assume $\upsilon<\infty$ throughout this section; if $\upsilon=\infty$ every bound below is void. By Lemma~\CGref{lem:rb-laplace}, $\widehat{\mathsf R}$ is the Laplace transform of the residual distribution of the section.
\end{definition}

\begin{lemma}[When the Laplace residual grows logarithmically]\label{lem:rb-upsilon-infinite}
For a finite section, $\upsilon<\infty$ if and only if $\mathsf a_n=0$ for every $n$. More precisely, with $\mathcal D(v):=\sum_{n=2}^{n^*}\mathsf a_nn^{-3/2-iv}$,
\begin{equation}\label{eq:rb-upsilon-split}
 \mathbb E(T)=\int_{-T}^{T}|\mathcal D(v)|\,dv+O_{C,k^*}(\log T),\quad
 \liminf_{T\to\infty}\frac1{2T}\int_{-T}^{T}|\mathcal D(v)|\,dv\ge\frac{\sum_{n\le n^*}|\mathsf a_n|^2n^{-3}}{\sum_{n\le n^*}|\mathsf a_n|n^{-3/2}}>0\ \ (\mathsf a\ne0).
\end{equation}
\end{lemma}

\begin{proof}
With $z=\frac32+iv$, the spectral part of $\widehat{\mathsf R}(2+iv)$ is $\frac12\zeta_{\RVM}'/\zeta_{\RVM}(2+iv)+\sum_{k\le k^*}\bigl[\frac z{z^2+\kappa_k^2}-\frac z{z^2+\tau_k^2}\bigr]$, because the list equals the clock beyond $k^*$. Each term of the finite sum is $z(\tau_k^2-\kappa_k^2)/((z^2+\kappa_k^2)(z^2+\tau_k^2))=O_{C,k}(|v|^{-3})$, integrable over $\mathbb R$. By Theorem~\CGref{thm:transition}(a) with $B=\frac32$, $\zeta_{\RVM}=\varpi_{\RVM}(1+w)(1+\mathcal E_0)$ on the band $1\le\sigma'\le2$, $t\ge t_0$, with $|w|\ll t^{-3/2}$, $w'=-\log(\tau_*/2\pi)\,w$ (Corollary~\CGref{cor:derivative-zeros}) and $|\mathcal E_0|\le C/t$; Cauchy's estimate on discs of radius $\frac14$ inside the band gives $|\mathcal E_0'|\le C'/t$, so $|\zeta_{\RVM}'/\zeta_{\RVM}(2+iv)|\ll(1+|v|)^{-1}$ and its integral over $[-T,T]$ is $O(\log T)$. This is the first statement of \CGeqref{eq:rb-upsilon-split}. For the second, $\mathcal D$ is a Dirichlet polynomial with distinct frequencies $\log n$, so $\frac1{2T}\int_{-T}^T|\mathcal D|^2dv\to\sum_n|\mathsf a_n|^2n^{-3}$ (Montgomery--Vaughan, or directly: the off-diagonal terms are $O(1/T)$), while $|\mathcal D|\le\sum_n|\mathsf a_n|n^{-3/2}$; hence $\int_{-T}^T|\mathcal D|\ge\int_{-T}^T|\mathcal D|^2/\sup|\mathcal D|$ gives the stated lower bound. If $\mathsf a\ne0$, then $\mathbb E(T)\asymp T$ and $\upsilon=\infty$; if $\mathsf a=0$, then $\mathbb E(T)=O(\log T)$ and $\upsilon<\infty$.
\end{proof}

\begin{remark}[Scope of the counting bound]\label{rem:rb-count-scope}
Lemma~\CGref{lem:rb-upsilon-infinite} restricts Theorem~\CGref{thm:rb-count} and Corollary~\CGref{cor:rb-displacement}, as stated for finite sections, to the sections without masses, for which $|N_{\boldsymbol\kappa}-n_{\RVM}|\le k^*+1$ is immediate. The reason is that on a vertical line the Laplace transform of the prime side is almost periodic and does not decay, so an $L^1$ residual of logarithmic growth on $\Re s=2$ demands that the spectral side reproduce the prime side at all heights, which no finite section does; the natural examples with $\upsilon<\infty$ are infinite lists with the full mass sequence, such as $(\boldsymbol\gamma,\mathsf a^\zeta)$, for which $\widehat{\mathsf R}\equiv0$ and $\upsilon=0$. The theorem is kept as the template of the count-to-displacement mechanism (Lemmas~\CGref{lem:rb-PL}--\CGref{lem:rb-argument}), which is reused in Appendix~\CGref{sec:appendix-endpoint}; its extension to infinite lists needs a polynomial growth bound for $P_{\boldsymbol\kappa}/P_{\boldsymbol\tau}$ on $\Re s\ge2$, which is not given here (Remark~\CGref{rem:rb-tube}), and Heuristic~\CGref{heur:rb-window} explains why the Euclidean residual of a numerical section is a different quantity.
\end{remark}

Constants. Put $\mathrm c_\ell:=2\sum_{n\ge2}n^{-3/2}/\log n<\infty$, and, for $s\in\mathbb C$ and real $x\ge2\pi$, $\mathfrak f(x;s):=\log\bigl(1+\frac{q(s)}{\frac14+x^2}\bigr)$, with the principal branch, which is defined whenever $|s-\frac12|<x$. Let $\mathbb J:=\{2\}\cup\{\sigma+i:\ -1\le\sigma\le2\}$, a compact set on which $|s-\frac12|\le\frac52<2\pi$, and
\begin{equation}\label{eq:rb-beta}
 \mathfrak B(C):=\sum_{k\ge1}\ \sup_{x\in[\max(\tau_k-C,2\pi),\,\tau_k+C]}\ \sup_{s\in\mathbb J}\bigl|\mathfrak f(x;s)-\mathfrak f(\tau_k;s)\bigr| .
\end{equation}
This is finite. On $\mathbb J$, $|q(s)|\le5$ and $|\partial_xf(x;s)|=\bigl|\frac{2xq(s)}{(\frac14+x^2)(x^2+(s-\frac12)^2)}\bigr|\le\frac{20}{x^3}$ for $x\ge4$. The $k$-th term is therefore at most $160C\tau_k^{-3}$ when $\tau_k\ge2C$, and $\sum_k\tau_k^{-3}<\infty$. Further, $\mathrm c_\tau:=\max_{-1\le\sigma\le2}|\log|\zeta_{\RVM}(\sigma+i)||$ and $\mathrm c_{\RVM}:=\sup_{\Re s\ge2}|\log\zeta_{\RVM}(s)|$. Both are finite: the first because $\zeta_{\RVM}$ has no zeros or poles on that segment, the second by \CGref{thm:halfplane-RVM}.

\begin{lemma}[Phragm\'en--Lindel\"of, two forms]\label{lem:rb-PL}
\textup{(a)} Let $\mathbb S:=\{\mathrm l_1\le\Re s\le\mathrm l_2,\ \Im s\ge t_1\}$ have width $a:=\mathrm l_2-\mathrm l_1$. Let $u$ be upper semicontinuous on $\mathbb S$ and subharmonic in its interior, with $u\le \mathrm M_0$ on $\partial\mathbb S$ and $u(s)\le \mathrm A_0e^{c\Im s}$ for some $c<\pi/a$. Then $u\le \mathrm M_0$ on $\mathbb S$.\\
\textup{(b)} Let $\mathfrak g$ be holomorphic in $\Re s>1$, continuous on $\Re s\ge2$, with $|g|\le1$ on $\Re s=2$ and $|\mathfrak g(s)|\le \mathrm A_0(2+|s|)^{\mathrm L}$ on $\Re s\ge2$. Then $|g|\le1$ on $\Re s\ge2$.
\end{lemma}

\begin{proof}
(a) Choose $c<c'<\pi/a$ and $m:=\frac{\mathrm l_1+\mathrm l_2}2$, and put $\omega(s):=e^{c'\Im s}\cos(c'(\Re s-m))$. This is harmonic, being the real part of $e^{-ic'(s-m)}$. On $\mathbb S$, $\omega\ge\cos(c'a/2)\,e^{c'\Im s}$, and $\cos(c'a/2)>0$. For $\mathrm p>0$, $u-\mathrm M_0-\mathrm p\omega$ is subharmonic, $\le0$ on $\partial\mathbb S$, and tends to $-\infty$ as $\Im s\to\infty$. The maximum principle on $\mathbb S\cap\{\Im s\le t_2\}$, with $t_2$ large, gives $u\le \mathrm M_0+\mathrm p\omega$. Let $\mathrm p\to0$.\\
(b) For $\mathrm p>0$, $\mathfrak g_\eta(s):=\mathfrak g(s)\exp(-\mathrm p(s-1)^{1/2})$, with the principal branch, satisfies $|\mathfrak g_\eta|\le|g|$ on $\Re s\ge2$, because $\Re(s-1)^{1/2}\ge\cos(\pi/4)|s-1|^{1/2}>0$ there. Moreover $\mathfrak g_\eta(s)\to0$ as $|s|\to\infty$ in $\Re s\ge2$. The maximum principle on large half-discs gives $|\mathfrak g_\eta|\le1$. Let $\mathrm p\to0$.
\end{proof}

\begin{lemma}[The Dirichlet factor]\label{lem:rb-factor}
Put $\mathcal D_{\mathsf a}(s):=\sum_{n=2}^{n^*}\frac{2\mathsf a_n}{\log n}n^{\frac12-s}$ and $\mathcal Y:=\zeta_{\boldsymbol\kappa}e^{-\mathcal D_{\mathsf a}}$. Then $\mathcal Y$ is holomorphic and zero-free in $\Re s>1$, and $|\mathcal D_{\mathsf a}|\le \mathbb M\mathrm c_\ell$ on $\Re s\ge2$. For the branch of $\log \mathcal Y$ that is real on $(1,\infty)$,
\begin{equation}\label{eq:rb-logY-deriv}
 (\log \mathcal Y)'=2\,\widehat{\mathsf R},\qquad
 |\log \mathcal Y(2)|\le \mathbb B_0:=|\log\zeta_{\RVM}(2)|+\mathfrak B(C)+\mathbb M\mathrm c_\ell,
\end{equation}
and, for all real $t$,
\begin{equation}\label{eq:rb-line}
 |\log \mathcal Y(2+it)|\le \mathbb B_0+\mathbb E(|t|)\le \mathbb B_0+\upsilon\log(2+|t|).
\end{equation}
\end{lemma}

\begin{proof}
$P_{\boldsymbol\kappa}$ vanishes only on $\Re s=\frac12$, and $G$ is holomorphic and zero-free in $\Re s>1$. Hence $\zeta_{\boldsymbol\kappa}$ and $\mathcal Y$ are holomorphic and zero-free there. $|\mathcal D_{\mathsf a}(s)|\le\sum_n2Mn^{\frac12-\Re s}/\log n\le \mathbb M\mathrm c_\ell$ for $\Re s\ge2$. Differentiating, $(\log \mathcal Y)'=\frac{P'_{\boldsymbol\kappa}}{P_{\boldsymbol\kappa}}-\frac{G'}G+2\sum_n\mathsf a_nn^{\frac12-s}=2\widehat{\mathsf R}$. At $s=2$, $\log\zeta_{\boldsymbol\kappa}(2)=\log\zeta_{\RVM}(2)+\sum_{k\le k^*}(\mathfrak f(\kappa_k;2)-\mathfrak f(\tau_k;2))$, a finite sum of real numbers bounded by $\mathfrak B(C)$, and $|\mathcal D_{\mathsf a}(2)|\le \mathbb M\mathrm c_\ell$. Finally, $\log \mathcal Y(2+it)-\log \mathcal Y(2)=2i\int_0^t\widehat{\mathsf R}(2+iv)\,dv$ has modulus at most $\mathbb E(|t|)$. For $|t|\ge2$, $\mathbb E(|t|)\le\upsilon\log|t|$; for $|t|<2$, $\mathbb E(|t|)\le \mathbb E(2)\le\upsilon\log2$.
\end{proof}

\begin{lemma}[Bound on the half-plane]\label{lem:rb-halfplane}
With $\mathbb W:=\log|\zeta_{\boldsymbol\kappa}|=\log|P_{\boldsymbol\kappa}|-\log|G|$ and $\mathbb B_1:=\mathbb B_0+\upsilon\log(2\sqrt2)$, $\mathbb B_2:=\mathbb B_1+\mathbb M\mathrm c_\ell$, for $\Re s\ge2$,
\begin{equation}\label{eq:rb-halfplane}
 \bigl|\log|\mathcal Y(s)|\bigr|\le \mathbb B_1+\upsilon\log|s-1|,\qquad |\mathbb W(s)|\le \mathbb B_2+\upsilon\log|s-1| .
\end{equation}
\end{lemma}

\begin{proof}
On $\Re s=2$, $|s-1|=\sqrt{1+t^2}\ge(2+|t|)/(2\sqrt2)$, so \CGeqref{eq:rb-line} gives $|\mathcal Y(s)|^{\pm1}\le e^{B_1}|s-1|^{\upsilon}$. A priori growth: on $\Re s\ge2$, $|\log|\zeta_{\RVM}||\le \mathrm c_{\RVM}$. The quotient $\zeta_{\boldsymbol\kappa}/\zeta_{\RVM}$ is a finite product of $k^*$ factors $\frac{(\kappa_k^2+w^2)(\frac14+\tau_k^2)}{(\frac14+\kappa_k^2)(\tau_k^2+w^2)}$ with $w=s-\frac12$. Since $|\kappa^2+w^2|=|\kappa-\Im s+i\Re w|\,|\kappa+\Im s-i\Re w|\ge(\Re w)^2\ge\frac94$, each factor and its inverse are $O((2+|s|)^4)$. Hence $|\mathcal Y|^{\pm1}\le \mathrm A_0(2+|s|)^{4k^*}$. Apply Lemma~\CGref{lem:rb-PL}(b) to $\mathfrak g_+:=\mathcal Y\,(s-1)^{-\upsilon}e^{-\mathbb B_1}$ and to $\mathfrak g_-:=\mathcal Y^{-1}(s-1)^{-\upsilon}e^{-\mathbb B_1}$, with the principal branch of $(s-1)^{-\upsilon}$, holomorphic in $\Re s>1$ with $|(s-1)^{-\upsilon}|=|s-1|^{-\upsilon}$; both are holomorphic in $\Re s>1$ and of polynomial growth on $\Re s\ge2$. Finally $\mathbb W=\log|\mathcal Y|+\Re\mathcal D_{\mathsf a}$.
\end{proof}

\begin{lemma}[Bound on the half-strip]\label{lem:rb-strip}
Let $\mathbb S:=\{-1\le\Re s\le2,\ \Im s\ge1\}$ and $\mathbb B_3:=\max\bigl(\mathbb B_2+\frac34\log2,\ \mathrm c_\tau+\mathfrak B(C)\bigr)$. Then
\begin{equation}\label{eq:rb-strip}
 \mathbb W(s)\le \mathbb B_3+\bigl(\tfrac32+\upsilon\bigr)\log|s+2|\qquad(s\in\mathbb S).
\end{equation}
\end{lemma}

\begin{proof}
$\mathbb W$ is subharmonic in the interior of $\mathbb S$, since $\log|G|$ is harmonic there ($G$ has no zeros or poles with $\Im s\ge1$), and it is upper semicontinuous on $\mathbb S$. The majorant $\mathrm h_0(s):=(\frac32+\upsilon)\log|s+2|$ is harmonic and nonnegative on $\mathbb S$. We check the three sides.

\emph{Right side.} By \CGeqref{eq:rb-halfplane}, $\mathbb W(2+it)\le \mathbb B_2+\upsilon\log|1+it|\le \mathbb B_2+\mathrm h_0(2+it)$.

\emph{Left side.} $P_{\boldsymbol\kappa}$ and $G$ are real, so $\mathbb W(\bar s)=\mathbb W(s)$. Combined with $P_{\boldsymbol\kappa}(1-s)=P_{\boldsymbol\kappa}(s)$, this gives $\mathbb W(-1+it)=\mathbb W(2+it)+\log|G(2+it)/G(-1+it)|$. By $\Gamma(z+1)=z\Gamma(z)$ and the reflection formula, $|\Gamma(\frac12+iy)|^2=\pi/\cosh\pi y$ and $|\Gamma(1+iy)|^2=\pi y/\sinh\pi y$, so
\begin{equation}\label{eq:rb-G-ratio}
 \Bigl|\frac{G(2+it)}{G(-1+it)}\Bigr|=\frac{|1+it|}2\,\pi^{-3/2}\Bigl(\frac t2\coth\frac{\pi t}2\Bigr)^{1/2}\le(1+t)^{3/2}\qquad(t\ge1).
\end{equation}
For the inequality: $|1+it|\le1+t$, $\frac t2\coth\frac{\pi t}2\le\frac{1+t}2\coth\frac\pi2$, and $\frac12\pi^{-3/2}(\frac12\coth\frac\pi2)^{1/2}<1$. Since $(1+t)\le\sqrt2|1+it|$ and $|s+2|=|1+it|$ at $s=-1+it$, we get $\mathbb W(-1+it)\le \mathbb B_2+\frac34\log2+\mathrm h_0(-1+it)$.

\emph{Bottom side.} For $s=\sigma+i$, $\mathbb W(s)=\log|\zeta_{\RVM}(s)|+\sum_{k\le k^*}\Re(\mathfrak f(\kappa_k;s)-\mathfrak f(\tau_k;s))\le \mathrm c_\tau+\mathfrak B(C)\le \mathrm c_\tau+\mathfrak B(C)+\mathrm h_0(s)$.

\emph{Growth.} $\log|P_{\boldsymbol\kappa}(s)|\le\sum_k\log(1+|q(s)|/\kappa_k^2)$. Split at $\kappa_k\le2|s|$ and $\kappa_k>2|s|$: the first part has $N_{\boldsymbol\kappa}(2|s|)=O(|s|\log|s|)$ terms, each at most $\log(1+|q(s)|/4\pi^2)=O(\log|s|)$; in the second, $\log(1+|q(s)|/\kappa_k^2)\le|q(s)|\kappa_k^{-2}$, and $\sum_{\kappa_k>2|s|}\kappa_k^{-2}=O(\log|s|/|s|)$ by partial summation with $N_{\boldsymbol\kappa}(t)=O(t\log t)$. Hence $\log|P_{\boldsymbol\kappa}(s)|=O(|s|\log^2|s|)$, and $-\log|G(s)|\le\frac\pi4|\Im s|+O(\log|s|)$ on $\mathbb S$ by Stirling's formula \cite[\S12.33]{WhittakerWatson1927}. Hence $\mathbb W-\mathrm h_0\le \mathrm A_0e^{c\Im s}$ for any $c>0$. Lemma~\CGref{lem:rb-PL}(a), with width $3$ and $c<\pi/3$, applied to $\mathbb W-\mathrm h_0$ gives the claim.
\end{proof}

\begin{lemma}[Backlund--Jensen on a horizontal segment]\label{lem:rb-backlund}
There are absolute constants $T_0\ge3$ and $\mathrm c_5$ such that, for $T\ge T_0$ with $T\notin\{\kappa_k\}$, the variation of $\arg P_{\boldsymbol\kappa}$ along the segment from $2+iT$ to $\frac12+iT$ is at most $\pi(n_T+1)$, where
\begin{equation}\label{eq:rb-nu}
 n_T\le\frac{(\frac52+\upsilon)\log(T+8)+\mathbb E(T)+\mathbb B_0+\mathbb B_3+\mathbb M\mathrm c_\ell+2\mathrm c_5}{\log(4/3)} .
\end{equation}
\end{lemma}

\begin{proof}
Choose $\mathrm a_T\in\mathbb R$ with $e^{i\mathrm a_T}P_{\boldsymbol\kappa}(2+iT)>0$ and put $\mathfrak F_T(z):=\frac12\bigl(e^{i\mathrm a_T}P_{\boldsymbol\kappa}(z+iT)+e^{-i\mathrm a_T}P_{\boldsymbol\kappa}(z-iT)\bigr)$, which is entire. Since $P_{\boldsymbol\kappa}$ is real, $P_{\boldsymbol\kappa}(x-iT)=\overline{P_{\boldsymbol\kappa}(x+iT)}$ for real $x$, so $\mathfrak F_T(x)=\Re(e^{i\mathrm a_T}P_{\boldsymbol\kappa}(x+iT))$ and $\mathfrak F_T(2)=|P_{\boldsymbol\kappa}(2+iT)|$. If $\mathfrak F_T$ has $n_T$ zeros on $[\frac12,2]$, the segment splits into at most $n_T+1$ pieces on which $\Re(e^{i\mathrm a_T}P_{\boldsymbol\kappa})$ has constant sign. On each piece $\arg$ varies by less than $\pi$, which gives the first claim \cite[\S9.4]{Titchmarsh1986}.

The segment $[\frac12,2]$ lies in the closed disc $|z-2|\le\frac32$. For $\frac32<r<2$, Jensen's formula on $|z-2|\le R=2$ (valid also if $\mathfrak F_T$ has zeros on the circle $|z-2|=2$, which contribute $\log1=0$, since $\mathfrak F_T$ is entire \cite[Thm.~15.18]{Rudin1987}) gives $n_T\log\frac Rr\le\max_{|z-2|=2}\log|\mathfrak F_T(z)|-\log|\mathfrak F_T(2)|$, because the zeros in the closed disc of radius $r$, which include the $n_T$ zeros on the segment, contribute at least $n_T\log(R/r)$; letting $r\downarrow\frac32$ gives the inequality with $\log\frac43$. Moreover $|\mathfrak F_T(z)|\le\max_{|\zeta-2|\le2}|P_{\boldsymbol\kappa}(\zeta+iT)|$, because $|P_{\boldsymbol\kappa}(z-iT)|=|P_{\boldsymbol\kappa}(\bar z+iT)|$. Write $\log|P_{\boldsymbol\kappa}|=\mathbb W+\log|G|$. On the disc $\{|s-(2+iT)|\le2\}$, the part with $\Re s\le2$ lies in $\mathbb S$ and $|s+2|\le T+6$, so Lemma~\CGref{lem:rb-strip} gives $\mathbb W\le \mathbb B_3+(\frac32+\upsilon)\log(T+6)$. The part with $\Re s\ge2$ satisfies $\mathbb W\le \mathbb B_2+\upsilon\log(T+5)$ by Lemma~\CGref{lem:rb-halfplane}, a smaller bound. At the centre, \CGeqref{eq:rb-line} gives $\mathbb W(2+iT)\ge-(\mathbb B_0+\mathbb E(T)+\mathbb M\mathrm c_\ell)$. For the oscillation of $\log|G|$ on the disc: by Stirling's formula for $\Gamma'/\Gamma$ in $|\arg z|\le\frac34\pi$ \cite[\S12.33]{WhittakerWatson1927}, $|G'/G(s)|=|\frac1s+\frac1{s-1}-\frac12\log\pi+\frac12\frac{\Gamma'}{\Gamma}(\frac s2)|\le\frac12\log(T+8)+\mathrm c_5$ on the disc for $T\ge T_0$. So the oscillation is at most $\log(T+8)+2\mathrm c_5$. Collecting terms gives \CGeqref{eq:rb-nu}.
\end{proof}

\begin{lemma}[Argument principle]\label{lem:rb-argument}
For $T>0$, $T\notin\{\kappa_k\}$, let $\mathbb G_T$ be the path $2\to2+iT\to\frac12+iT$. Then
\begin{equation}\label{eq:rb-arg-principle}
 \pi N_{\boldsymbol\kappa}(T)=\Delta_{\mathbb G_T}\arg P_{\boldsymbol\kappa}
 =\pi+\vartheta(T)+\Delta_{\mathbb G_T}\arg\zeta_{\boldsymbol\kappa},
\end{equation}
where $\vartheta(T)=\arg\bigl(\pi^{-iT/2}\Gamma(\frac14+\frac{iT}2)\bigr)$ is the Riemann--Siegel theta function, defined by continuous variation from $\vartheta(0)=0$.
\end{lemma}

\begin{proof}
Take the rectangle $\mathfrak R_T$ with vertices $2\pm iT$, $-1\pm iT$. $P_{\boldsymbol\kappa}$ has no zeros on $\partial\mathfrak R_T$: its zeros lie on $\Re s=\frac12$, and $T\notin\{\kappa_k\}$, so $\frac12\pm iT$ are not zeros. Its zeros inside $\mathfrak R_T$ are $\frac12\pm i\kappa_k$ with $\kappa_k<T$, counted with multiplicity (coincident $\kappa_j=\kappa_k$ give a multiple zero), $2N_{\boldsymbol\kappa}(T)$ in total. Reflection $s\mapsto1-\bar s$ maps $\mathbb G_T$ to the path $-1\to-1+iT\to\frac12+iT$. Since $P_{\boldsymbol\kappa}(1-\bar s)=\overline{P_{\boldsymbol\kappa}(s)}$, the variation of $\arg P_{\boldsymbol\kappa}$ along the reversed image path $\frac12+iT\to-1+iT\to-1$ equals $\Delta_{\mathbb G_T}\arg P_{\boldsymbol\kappa}$. The two lower quarters are handled in the same way with $s\mapsto\bar s$. The four quarter paths make up $\partial\mathfrak R_T$, so $4\Delta_{\mathbb G_T}\arg P_{\boldsymbol\kappa}=2\pi\cdot2N_{\boldsymbol\kappa}(T)$.

For the second equality, write $P_{\boldsymbol\kappa}=G\zeta_{\boldsymbol\kappa}$. Along $\mathbb G_T$, $s(s-1)$ moves from $2$ to $-(\frac14+T^2)$ inside the closed upper half-plane, so its argument varies by $\pi$. $\pi^{-s/2}\Gamma(s/2)$ is holomorphic and zero-free in $\Re s>0$, so its argument variation along $\mathbb G_T$ equals that along $[2,\frac12]\cup[\frac12,\frac12+iT]$, which is $0+\vartheta(T)$.
\end{proof}

\begin{theorem}[Counting bound for quasi-solutions]\label{thm:rb-count}
There are absolute constants $T_0$ and $\mathrm c_5$ such that, for every finite section in $\mathbb K(C,\mathbb M)$ and every $T\ge T_0$,
\begin{equation}\label{eq:rb-count}
 \bigl|N_{\boldsymbol\kappa}(T)-n_{\RVM}(T)\bigr|\le(\mathrm c_1+\mathrm c_2\upsilon)\log(T+8)+\mathrm c_3\mathbb B+\mathrm c_4,
\end{equation}
with
\begin{equation}\label{eq:rb-a-constants}
 \mathrm c_1=\frac5{2\log(4/3)}=8.690\ldots,\quad \mathrm c_2=\frac2{\log(4/3)}+\frac1\pi=7.270\ldots,\quad \mathrm c_3=\frac1{\log(4/3)}+\frac1\pi=3.794\ldots,
\end{equation}
$\mathrm c_4:=\frac32+\frac{\mathrm c_6}\pi+\frac1{T_0}+\frac{2\mathrm c_5}{\log(4/3)}$ with $\mathrm c_5$ the constant of Lemma~\CGref{lem:rb-backlund} and $\mathrm c_6:=\frac32$, and $\mathbb B:=\mathbb B_0+\mathbb B_3+\mathbb M\mathrm c_\ell$. The constants $\mathrm c_1,\dots,\mathrm c_6$ do not depend on $k^*$, $n^*$, $C$, $\mathbb M$ or $\upsilon$; $\mathbb B$ depends on $C$ (through $\mathfrak B(C)$), on $\mathbb M$, and linearly on $\upsilon$ (through $\mathbb B_2$), so the total coefficient of $\upsilon$ in \CGeqref{eq:rb-count} is $\mathrm c_2\log(T+8)+\mathrm c_3\log(2\sqrt2)$. In Corollary~\CGref{cor:rb-displacement} the parameters $C$, $\mathbb M$, $\upsilon$ are fixed before $k\to\infty$.
\end{theorem}

\begin{proof}
By continuity from the right it suffices to take $T\notin\{\kappa_k\}$. By Lemma~\CGref{lem:rb-argument}, $N_{\boldsymbol\kappa}(T)-1-\vartheta(T)/\pi=\pi^{-1}\Delta_{\mathbb G_T}\arg\zeta_{\boldsymbol\kappa}$.

\emph{Vertical piece.} $\arg\zeta_{\boldsymbol\kappa}=\Im\log \mathcal Y+\Im\mathcal D_{\mathsf a}$, with $\Im\log \mathcal Y(2)=\Im\mathcal D_{\mathsf a}(2)=0$. Its variation is at most $\mathbb E(T)+\mathbb M\mathrm c_\ell$ by Lemma~\CGref{lem:rb-factor}.

\emph{Horizontal piece.} The variation is $\Delta\arg P_{\boldsymbol\kappa}-\Delta\arg G$. The first is at most $\pi(n_T+1)$ by Lemma~\CGref{lem:rb-backlund}. The second is $-\Im\int_{1/2}^2(G'/G)(\sigma+iT)\,d\sigma$. Here $|\Im(G'/G)|\le|\frac1s|+|\frac1{s-1}|+\frac12|\arg\frac s2|+O(1/T)\le1$ for $T\ge T_0$, by Stirling's formula for $\Gamma'/\Gamma$ \cite[\S12.33]{WhittakerWatson1927}. So the second variation is at most $\mathrm c_6:=\frac32$.

\emph{Main term.} By the Stirling expansion of $\vartheta$ \cite[\S12.33]{WhittakerWatson1927}, $1+\vartheta(T)/\pi=\NR(T)+O(1/T)$, and this error is at most $1/T$ for $T\ge T_0$. With \CGeqref{eq:rb-Ntau-floor}:
\begin{equation}\label{eq:rb-aux-1}
 |N_{\boldsymbol\kappa}(T)-n_{\RVM}(T)|\le n_T+1+\frac{\mathbb E(T)+\mathbb M\mathrm c_\ell+\mathrm c_6}{\pi}+\frac1T+\frac12 .
\end{equation}
Insert \CGeqref{eq:rb-nu} and $\mathbb E(T)\le\upsilon\log(T+8)$: the coefficient of $\log(T+8)$ is $\frac{5/2+2\upsilon}{\log(4/3)}+\frac\upsilon\pi=\mathrm c_1+\mathrm c_2\upsilon$, the terms in $\mathbb B_0+\mathbb B_3+\mathbb M\mathrm c_\ell$ are at most $(\frac1{\log(4/3)}+\frac1\pi)\mathbb B=\mathrm c_3\mathbb B$, and the remaining constant is $\frac{2\mathrm c_5}{\log(4/3)}+1+\frac{\mathrm c_6}\pi+\frac1{T_0}+\frac12=\mathrm c_4$.
\end{proof}

\begin{corollary}[Displacement bound]\label{cor:rb-displacement}
Assume $T_0\ge2\pi e$. For every $k$ with $\tau_k-C\ge T_0$,
\begin{equation}\label{eq:rb-disp}
 |\kappa_k-\tau_k|\le\frac{2\pi\bigl[(\mathrm c_1+\mathrm c_2\upsilon)\log(\tau_k+C+8)+\mathrm c_3\mathbb B+\mathrm c_4\bigr]}{\log\bigl((\tau_k-C)/2\pi\bigr)} .
\end{equation}
In particular $\limsup_{k\to\infty}|\kappa_k-\tau_k|\le2\pi(\mathrm c_1+\mathrm c_2\upsilon)$ uniformly over sections with given $C,\mathbb M,\upsilon$. For every $\mathrm p>0$ there is $k_0(C,\mathbb M,\upsilon,\mathrm p)$ such that $|\kappa_k-\tau_k|\le2\pi(\mathrm c_1+\mathrm c_2\upsilon)+\mathrm p$ for $k\ge k_0$.
\end{corollary}

\begin{proof}
Let $\mathbb D(T)$ be the right-hand side of \CGeqref{eq:rb-count}; it is nondecreasing. Suppose $\kappa_k>\tau_k+\mathrm q$ with $0<\mathrm q\le C$, and let $T:=\tau_k+\mathrm q-\mathrm h'$ with small $\mathrm h'>0$. The list is nondecreasing, so $N_{\boldsymbol\kappa}(T)\le k-1$. By \CGeqref{eq:rb-Ntau-floor} and because $\NR'$ is increasing, $n_{\RVM}(T)\ge\NR(T)-\frac12\ge k-1+(\mathrm q-\mathrm h')\NR'(\tau_k)$. Hence $(\mathrm q-\mathrm h')\frac1{2\pi}\log\frac{\tau_k}{2\pi}\le\mathbb D(\tau_k+C)$. Now suppose $\kappa_k<\tau_k-\mathrm q$, and let $T:=\tau_k-\mathrm q+\mathrm h'$. Then $N_{\boldsymbol\kappa}(T)\ge k$ and $n_{\RVM}(T)\le\NR(T)+\frac12\le k-(\mathrm q-\mathrm h')\NR'(\tau_k-C)$. Let $\mathrm h'\to0$ in both cases.
\end{proof}

\begin{remark}[What Corollary~\ref{cor:rb-displacement} does not give]\label{rem:rb-tube}
The proof uses the imposed tube $|\kappa_k-\tau_k|\le C$ twice: through $\mathfrak B(C)$ in $\mathbb B$, and to replace $\mathbb D(\tau_k+\mathrm q)$ by $\mathbb D(\tau_k+C)$. The conclusion is that, inside any finite tube, the coordinates of sufficiently high index lie in the smaller tube of radius $2\pi(\mathrm c_1+\mathrm c_2\upsilon)+\mathrm p$. So the constraint is inactive there once $C>2\pi(\mathrm c_1+\mathrm c_2\upsilon)+\mathrm p$. For the indices $k<k_0(C,\mathbb M,\upsilon,\mathrm p)$ nothing is proved, and the corollary does not show that unconstrained near-minimisers lie in a universal tube. The proof uses finite sections: Lemma~\CGref{lem:rb-halfplane} needs the polynomial growth $\mathrm A_0(2+|s|)^{4k^*}$. Extending it to infinite lists requires a separate growth bound for $P_{\boldsymbol\kappa}/P_{\boldsymbol\tau}$ on $\Re s\ge2$, which is not given here. The constant $2\pi \mathrm c_1=54.60\ldots$ is far from optimal.
\end{remark}

\begin{heuristic}[What a finite section controls]\label{heur:rb-window}
The quantity in Theorem~\CGref{thm:rb-count} is $\mathbb E(T)$, the $L^1$ norm on $\Re s=2$ of the Laplace transform of the residual. It is not the Euclidean residual that a numerical section minimises over finitely many tests. Let the tests be $\mathsf j_0(u)\varphi(u)$ with a fixed window $\mathsf j_0\in C_c^\infty(0,\infty)$, $\mathsf j_0\equiv1$ on $[u_0,u^*]$, and $\varphi$ band-limited to $|\omega|\le\mathrm w^\sharp$. Then the kernel $\mathsf j_0(u)e^{-\frac32u}e^{-ivu}$ is, approximately, in the span of the tests only for $|v|\lesssim\mathrm w^\sharp$. The residual on $(0,u_0)$ is then not controlled, and for a list at distance $C$ from the clock it is expected to be of size $C\log(1/u)$ as $u\to0$. Its contribution to $\mathbb E(T)$ is then $O(CTu_0\log(1/u_0))$. Since $u_0\gtrsim1/\mathrm w^\sharp$ in a band-limited section, Theorem~\CGref{thm:rb-count} would apply with a given $\upsilon_0$ only at heights $T\lesssim\upsilon_0\mathrm w^\sharp/(C\log\mathrm w^\sharp)$. None of these size estimates is proved here. They indicate which norm a section must make small. Lemma~\CGref{lem:rb-factor} shows that this norm is the right one.
\end{heuristic}

\begin{heuristic}[Informative part of a section]\label{heur:rb-informative}
Where the atoms are denser than the resolution of the tests, $e^{-u}\ll\pi/\mathrm w^\sharp$, the masses alone can approximate a smooth target by a Riemann sum $\mathsf a_n\approx \mathfrak g(\log n)/n$ without moving any ordinate. That part of the window then carries no information on the ordinates. Counting unknowns against equations in the remaining window, $u^*\le\log\mathrm w^\sharp+O(1)$, gives an overdetermined system. This proves nothing about compatibility. It only explains why B is an existence statement and not a dimension count.
\end{heuristic}

%======================================================================
\subsection{A Kadec--Avdonin criterion with blocks of variable length}\label{subsec:rb-riesz}

\subsubsection{Discrepancy and generating function}\label{subsec:rb-generating}

\begin{definition}[Discrepancy, Poisson and conjugate operators]\label{def:rb-discrepancy}
Let $\mathscr X=(\mathrm x_k)_{k\in\mathbb Z}$ be a real nondecreasing sequence with
\begin{equation}\label{eq:rb-H}
 |\mathrm x_k-(k-\tfrac12)|\le \mathrm c_{\mathscr X}\log(2+|k|)\qquad(k\in\mathbb Z).
\end{equation}
Its \emph{discrepancy} with respect to the half-integers is
\begin{equation}\label{eq:rb-psi}
 \varPsi_{\mathscr X}(\mathrm x):=\sum_{k\in\mathbb Z}\bigl(\mathbf 1_{[\mathrm x_k,\infty)}(\mathrm x)-\mathbf 1_{[k-\frac12,\infty)}(\mathrm x)\bigr).
\end{equation}
The $k$-th term is $\pm1$ on the interval between $\mathrm x_k$ and $k-\frac12$ and $0$ elsewhere. By \CGeqref{eq:rb-H}, at most $O(\log(2+|\mathrm x|))$ terms are nonzero at $\mathrm x$, so $|\varPsi_{\mathscr X}(\mathrm x)|\le \mathrm c'_{\mathscr X}\log(2+|\mathrm x|)$ and $\varPsi_{\mathscr X}\in L^1(d\mathrm x/(1+\mathrm x^2))$. For $\mathrm g\in L^1(d\mathrm x/(1+\mathrm x^2))$ and $\mathrm y>0$ put
\begin{equation}\label{eq:rb-PQ}
 \mathbb P_{\mathrm y}\mathrm g(\mathrm x):=\frac1\pi\int\frac{\mathrm y\,\mathrm g(\mathrm s)\,d\mathrm s}{(\mathrm x-\mathrm s)^2+\mathrm y^2},\qquad
 \mathbb Y_{\mathrm y}\mathrm g(\mathrm x):=\frac1\pi\int \mathrm g(\mathrm s)\Bigl[\frac{\mathrm x-\mathrm s}{(\mathrm x-\mathrm s)^2+\mathrm y^2}+\frac{\mathrm s}{1+\mathrm s^2}\Bigr]d\mathrm s .
\end{equation}
The last term in $\mathbb Y_{\mathrm y}$ only makes the integral converge. Changing it changes $\mathbb Y_{\mathrm y}\mathrm g$ by a constant, which is irrelevant for all statements below.
\end{definition}

\begin{lemma}[Generating function]\label{lem:rb-generating}
Let $\mathscr X=(\mathrm x_k)_{k\in\mathbb Z}$ satisfy \CGeqref{eq:rb-H}, understood as a sequence indexed by $\mathbb Z$ in which a point of multiplicity $m$ occurs $m$ times. Let $\mathcal F_0(\mathfrak z):=\cos\pi \mathfrak z$, and fix $\mathfrak z_0\in\mathbb R\setminus(\mathscr X\cup(\mathbb Z+\frac12))$. The series
\begin{equation}\label{eq:rb-aux-2}
 \mathfrak G_{\mathscr X}(\mathfrak z):=\sum_{k\in\mathbb Z}\Bigl(\frac1{\mathfrak z-\mathrm x_k}-\frac1{\mathfrak z-k+\frac12}\Bigr)
\end{equation}
converges absolutely and locally uniformly off $\mathscr X\cup(\mathbb Z+\frac12)$, and
\begin{equation}\label{eq:rb-aux-3}
 \mathcal F_{\mathscr X}(\mathfrak z):=\mathcal F_0(\mathfrak z)\exp\int_{\mathfrak z_0}^{\mathfrak z}\mathfrak G_{\mathscr X}(\mathrm z)\,d\mathrm z
\end{equation}
is a real entire function whose zero set is $\mathscr X$, with multiplicities. Moreover:
\begin{enumerate}[label=\textup{(\roman*)}]
\item for $\mathrm y>0$ there is a constant $\mathrm c_{\mathscr X}''$ such that
\begin{equation}\label{eq:rb-F-formula}
 \log|\mathcal F_{\mathscr X}(\mathrm x+iy)|=\log|\mathcal F_0(\mathrm x+iy)|+\pi\,\mathbb Y_{\mathrm y}\varPsi_{\mathscr X}(\mathrm x)+\mathrm c_{\mathscr X}'' ;
\end{equation}
\item $\mathcal F_{\mathscr X}$ is of exponential type $\pi$, and its indicator diagram is the segment $[-i\pi,i\pi]$;
\item if $\mathscr X$ is symmetric, $\mathrm x_{1-k}=-\mathrm x_k$, and $\mathfrak z_0=0\notin \mathscr X$, then $\mathcal F_{\mathscr X}$ is even.
\end{enumerate}
\end{lemma}

\begin{proof}
Each term of $\mathfrak G_{\mathscr X}$ equals $\frac{x_k-k+\frac12}{(\mathfrak z-\mathrm x_k)(\mathfrak z-k+\frac12)}=O(\log(2+|k|)/k^2)$ locally uniformly, so the series converges. $\mathfrak G_{\mathscr X}$ is meromorphic with simple poles, with residue $+m$ at a point of $\mathscr X$ of multiplicity $m$ and $-1$ at each half-integer; coincident poles combine their residues. The residues are integers, so $\exp\int_{\mathfrak z_0}^{\mathfrak z}\mathfrak G_{\mathscr X}$ is single-valued. It has a pole of order one at each half-integer not in $\mathscr X$, cancelled by the zero of $\mathcal F_0$. Hence $\mathcal F_{\mathscr X}$ is entire with zero set $\mathscr X$. $\mathfrak G_{\mathscr X}(\bar{\mathfrak z})=\overline{\mathfrak G_{\mathscr X}(\mathfrak z)}$ and $\mathfrak z_0$ is real, so $\mathcal F_{\mathscr X}$ is real on $\mathbb R$. Under the symmetry in (iii), $\mathfrak G_{\mathscr X}$ is odd, so $\int_0^{\mathfrak z}\mathfrak G_{\mathscr X}$ is even, and $\mathcal F_{\mathscr X}$ is even.

(i) For $\Im \mathfrak z\ne0$ put $\mathbb A_{\mathscr X}(\mathfrak z):=\int\varPsi_{\mathscr X}(\mathrm s)\bigl[\frac1{\mathfrak z-\mathrm s}+\frac{\mathrm s}{1+\mathrm s^2}\bigr]d\mathrm s$. For $f(\mathrm s):=\frac1{\mathfrak z-\mathrm s}$, the identity $f(\mathrm x_k)-f(k-\frac12)=-\int f'(\mathrm s)\bigl(\mathbf 1_{[\mathrm x_k,\infty)}-\mathbf 1_{[k-\frac12,\infty)}\bigr)(\mathrm s)\,d\mathrm s$ summed over $k$ (Fubini applies because $\sum_k\int|f'||\cdots|<\infty$ by \CGeqref{eq:rb-H}) gives $\mathfrak G_{\mathscr X}(\mathfrak z)=-\int\varPsi_{\mathscr X}(\mathrm s)(\mathfrak z-\mathrm s)^{-2}d\mathrm s=\mathbb A_{\mathscr X}'(\mathfrak z)$. So $\mathcal F_{\mathscr X}=\mathrm c_+\mathcal F_0e^{\mathbb A_{\mathscr X}}$ on $\Im \mathfrak z>0$, with $\mathrm c_+\ne0$. Since $\Re\frac1{\mathfrak z-\mathrm s}=\frac{\mathrm x-\mathrm s}{(\mathrm x-\mathrm s)^2+\mathrm y^2}$, $\Re\mathbb A_{\mathscr X}(\mathrm x+iy)=\pi\mathbb Y_{\mathrm y}\varPsi_{\mathscr X}(\mathrm x)$ and $\Im\mathbb A_{\mathscr X}(\mathrm x+iy)=-\pi \mathbb P_{\mathrm y}\varPsi_{\mathscr X}(\mathrm x)$.

(ii) Split $\mathbb Y_{\mathrm y}\varPsi_{\mathscr X}(\mathrm x)$ into $|\mathrm s-\mathrm x|\le1$, $1<|\mathrm s-\mathrm x|$ with $|\mathrm s|\le2|\mathfrak z|+2$, and $|\mathrm s|>2|\mathfrak z|+2$. Using $|\varPsi_{\mathscr X}(\mathrm s)|\le \mathrm c'_{\mathscr X}\log(2+|\mathrm s|)$, $\int_{|u|\le1}\frac{|u|\,du}{u^2+\mathrm y^2}=\log(1+\mathrm y^{-2})$, and, for $|\mathrm s|>2|\mathfrak z|+2$, $|\frac{\mathrm x-\mathrm s}{(\mathrm x-\mathrm s)^2+\mathrm y^2}+\frac{\mathrm s}{1+\mathrm s^2}|\le8(1+|\mathfrak z|)/\mathrm s^2$, one gets
\begin{equation}\label{eq:rb-Q-growth}
 |\mathbb Y_{\mathrm y}\varPsi_{\mathscr X}(\mathrm x)|\le \mathrm c_{\mathscr X}'''\bigl(1+\log(1+\mathrm y^{-2})\bigr)\log^2(4+|\mathfrak z|)\qquad(\mathfrak z=\mathrm x+iy,\ \mathrm y>0).
\end{equation}
For $|\mathrm y|\ge1$, $\log|\mathcal F_{\mathscr X}(\mathfrak z)|\le\pi|\mathrm y|+O(\log^2(4+|\mathfrak z|))$ by \CGeqref{eq:rb-F-formula}, \CGeqref{eq:rb-Q-growth} and $|\mathcal F_{\mathscr X}(\bar{\mathfrak z})|=|\mathcal F_{\mathscr X}(\mathfrak z)|$. For $|\mathrm y_1|<1$, apply the sub-mean-value property of $\log|\mathcal F_{\mathscr X}|$ on the circle $|\mathfrak z-\mathfrak z_1|=\frac12$. The factor $\log(1+\mathrm y^{-2})$ is integrable along the circle uniformly in $\mathrm y_1$, since it has at most two logarithmic singularities. This gives $\log|\mathcal F_{\mathscr X}(\mathfrak z_1)|\le\frac{3\pi}2+O(\log^2(4+|\mathfrak z_1|))$ for $|\mathrm y_1|<1$. Hence the type is at most $\pi$ and the indicator satisfies $\mathrm h(\mathrm u)\le\pi|\sin\mathrm u|$. Since \CGeqref{eq:rb-F-formula} is an equality and \CGeqref{eq:rb-Q-growth} is two-sided, $\log|\mathcal F_{\mathscr X}(iy)|=\log|\cos(i\pi \mathrm y)|+\pi\mathbb Y_{\mathrm y}\varPsi_{\mathscr X}(0)+\mathrm c_{\mathscr X}''=\pi|\mathrm y|+O(\log^2(4+|\mathrm y|))$ for $|\mathrm y|\ge1$, so $\mathrm h(\pm\frac\pi2)=\pi$. The indicator diagram lies in $[-i\pi,i\pi]$, whose support function is $\pi|\sin\mathrm u|$, and it contains $\pm i\pi$, so it equals that segment.
\end{proof}

\subsubsection{The criterion}\label{subsec:rb-criterion}

A weight $\mathrm w>0$ on $\mathbb R$ is in $(A_2)$ if
\begin{equation}\label{eq:rb-aux-4}
 [\mathrm w]_{A_2}:=\sup_I\Bigl(\frac1{|I|}\int_Iw\Bigr)\Bigl(\frac1{|I|}\int_Iw^{-1}\Bigr)<\infty ,
\end{equation}
the supremum being over bounded intervals.

\begin{theorem}[Block $\tfrac14$-criterion]\label{thm:rb-criterion}
Let $\mathscr X$ satisfy \CGeqref{eq:rb-H} and suppose $\varPsi_{\mathscr X}=\mathrm A+\mathrm B$ with measurable $\mathrm A,\mathrm B$ such that:
\begin{enumerate}[label=\textup{(\roman*)}]
\item $\|\mathrm A\|_\infty\le \mathrm d_\star<\frac14$;
\item $\sup_x|\mathbb Y_1\mathrm B(\mathrm x)-\mathrm c_{\mathrm B}|\le \mathbb U$ for some constant $\mathrm c_{\mathrm B}$.
\end{enumerate}
Then $\mathrm w(\mathrm x):=|\mathcal F_{\mathscr X}(\mathrm x+i)|^2$ is in $(A_2)$ and
\begin{equation}\label{eq:rb-A2const}
 [\mathrm w]_{A_2}\le\frac{\pi^2\coth^2(\pi)\,e^{4\pi \mathbb U}}{\cos^2(2\pi\mathrm d_\star)} .
\end{equation}
\end{theorem}

\begin{proof}
$\mathrm B=\varPsi_{\mathscr X}-\mathrm A\in L^1(d\mathrm x/(1+\mathrm x^2))$, so $\mathbb Y_1\varPsi_{\mathscr X}=\mathbb Y_1\mathrm A+\mathbb Y_1\mathrm B$. By \CGeqref{eq:rb-F-formula}, $\mathrm w=e^{2\mathrm c''_{\mathscr X}}|\mathcal F_0(\mathrm x+i)|^2e^{2\pi\mathbb Y_1\mathrm A}e^{2\pi\mathbb Y_1\mathrm B}$.

Put $\mathbb A_{\mathrm A}(\mathfrak z):=\int \mathrm A(\mathrm s)\bigl[\frac1{\mathfrak z-\mathrm s}+\frac{\mathrm s}{1+\mathrm s^2}\bigr]d\mathrm s$. It is holomorphic on $\Im \mathfrak z>0$, and $\mathbb A_{\mathrm A}(\mathfrak o+i)$ is holomorphic on $\Im\mathfrak o>-1$. Let $\mathcal W(\mathfrak o):=\exp(2\mathbb A_{\mathrm A}(\mathfrak o+i))$. For $\mathfrak o=\mathrm x_1+i\mathrm y_1$ with $\mathrm y_1\ge0$,
\begin{equation}\label{eq:rb-aux-5}
 |\mathcal W(\mathfrak o)|=e^{2\pi\mathbb Y_{1+\mathrm y_1}\mathrm A(\mathrm x_1)},\qquad\arg\mathcal W(\mathfrak o)=-2\pi \mathbb P_{1+\mathrm y_1}\mathrm A(\mathrm x_1)\in[-2\pi\mathrm d_\star,2\pi\mathrm d_\star].
\end{equation}
So $\mathcal W$ maps the closed upper half-plane into the sector $|\arg|\le \mathrm b_\star:=2\pi\mathrm d_\star<\frac\pi2$, and $|\mathcal W(\mathrm x)|=e^{2\pi\mathbb Y_1\mathrm A(\mathrm x)}$ on $\mathbb R$.

$\Re\mathcal W$ is a positive harmonic function on $\Im\mathfrak o>0$, continuous up to $\mathbb R$. By the Herglotz representation \cite[Ch.~I]{Garnett2007} it equals $\mathrm c\,\mathrm y_1+\int \mathbb P_{\mathfrak o}(t)\,d\mathrm n(t)$ with $\mathrm c\ge0$ and $\mathrm n\ge0$, where $\mathrm n$ is the vague limit of the measures $\Re\mathcal W(t+i\mathrm y)\,dt$ as $\mathrm y\downarrow0$.  Since $\Re\mathcal W$ is continuous on the closed half-plane, $\Re\mathcal W(t+i\mathrm y)\to\Re\mathcal W(t)$ uniformly on compact sets of $t$, so $\mathrm n=\Re\mathcal W(t)\,dt$. Here $\mathbb P_{\mathfrak o}(t):=\frac1\pi\frac{\mathrm y_1}{(t-\mathrm x_1)^2+\mathrm y_1^2}$. Hence $\Re\mathcal W(\mathfrak o)\ge\int \mathbb P_{\mathfrak o}(t)\Re\mathcal W(t)\,dt$. For an interval $I$ with centre $\mathrm x_I$, put $\mathfrak o_I:=\mathrm x_I+i|I|/2$. Then $\mathbb P_{\mathfrak o_I}(t)\ge(\pi|I|)^{-1}$ for $t\in I$, and
\begin{equation}\label{eq:rb-aux-6}
 \frac1{|I|}\int_I|\mathcal W|\le\frac1{\cos\mathrm b_\star}\,\frac1{|I|}\int_I\Re\mathcal W\le\frac\pi{\cos\mathrm b_\star}\Re\mathcal W(\mathfrak o_I)\le\frac\pi{\cos\mathrm b_\star}|\mathcal W(\mathfrak o_I)| .
\end{equation}
$1/\mathcal W$ maps into the same sector, so the same holds for it. Multiplying, $[|\mathcal W|]_{A_2}\le\pi^2/\cos^2\mathrm b_\star$.

Finally $\mathrm w=\mathrm g\,|\mathcal W|$ with $\mathrm g:=e^{2\mathrm c''_{\mathscr X}}|\mathcal F_0(\mathrm x+i)|^2e^{2\pi\mathbb Y_1\mathrm B}$. Since $|\cos\pi(\mathrm x+i)|^2=\cos^2(\pi \mathrm x)\cosh^2\pi+\sin^2(\pi \mathrm x)\sinh^2\pi\in[\sinh^2\pi,\cosh^2\pi]$ and $e^{2\pi\mathbb Y_1\mathrm B}\in[e^{2\pi(\mathrm c_{\mathrm B}-\mathbb U)},e^{2\pi(\mathrm c_{\mathrm B}+\mathbb U)}]$, we have $\sup \mathrm g/\inf \mathrm g\le\coth^2(\pi)e^{4\pi \mathbb U}$. For any weight $\mathrm w_1$, $[\mathrm g\mathrm w_1]_{A_2}\le(\sup \mathrm g/\inf \mathrm g)[\mathrm w_1]_{A_2}$.
\end{proof}

\begin{corollary}[Riesz bases]\label{cor:rb-riesz}
Under the hypotheses of Theorem~\CGref{thm:rb-criterion}:
\begin{enumerate}[label=\textup{(\alph*)}]
\item if $\mathscr X$ is uniformly discrete (in particular without repeated points), $\{e^{i\mathrm x_k\mathrm v}\}_{k\in\mathbb Z}$ is a Riesz basis of $L^2(-\pi,\pi)$;
\item if $\mathscr X$ is a finite union of uniformly discrete sets, the generalised divided differences of $\{e^{i\mathrm x_k\mathrm v}\}$ in the sense of \cite{AvdoninIvanov2002} form a Riesz basis of $L^2(-\pi,\pi)$. Consequently, for every partition of $\mathbb Z$ into consecutive blocks $\mathfrak B_m$ that are unions of the groups of \cite{AvdoninIvanov2002}, the subspaces spanned by $\{\mathrm v^je^{i\mathrm x\mathrm v}\}$, with $\mathrm x=\mathrm x_k$ for $k\in\mathfrak B_m$ and $j$ below the multiplicity of $\mathrm x$, form a Riesz basis of subspaces.
\end{enumerate}
\end{corollary}

\begin{proof}
By Lemma~\CGref{lem:rb-generating}, $\mathcal F_{\mathscr X}$ is entire of exponential type with indicator diagram of width $2\pi$ and zero set $\mathscr X$. By Theorem~\CGref{thm:rb-criterion}, $|\mathcal F_{\mathscr X}(\mathrm x+i)|^2\in(A_2)$, which is equivalent to the Helson--Szeg\H{o} condition \cite[Ch.~VI]{Garnett2007}. Then (a) is Pavlov's criterion \cite[Thm.~1]{AvdoninIvanov2002}, and the first statement of (b) is \cite[Thm.~3]{AvdoninIvanov2002}. Both are stated for $L^2(0,2\pi)$, and translation by $\pi$ multiplies each exponential by a unimodular constant. Grouping the vectors of a Riesz basis into disjoint groups gives a Riesz basis of subspaces.
\end{proof}

\subsubsection{Kadec, Avdonin, and variable block length}\label{subsec:rb-variable}

Write $\mathrm x_k=k-\frac12+\mathrm d_k$, $I_k:=[k-1,k)$, $\mathsf i_k:=\mathbf 1_{[\mathrm x_k,\infty)}-\mathbf 1_{[k-\frac12,\infty)}$, so that $\varPsi_{\mathscr X}=\sum_k\mathsf i_k$ and $\int\mathsf i_k=-\mathrm d_k$.

\begin{proposition}[Bounded conjugate for block-supported oscillation]\label{prop:rb-variable-blocks}
Let $\{\mathfrak J_m\}_{m\in\mathbb Z}$ be a partition of $\mathbb R$ into consecutive intervals of lengths $\mathfrak n_m\ge1$. Let $\mathrm B=\sum_m\mathrm B_m$ with $\operatorname{supp}\mathrm B_m\subset \mathfrak J_m$, $\int \mathrm B_m=0$ and $\|\mathrm B_m\|_{L^\infty}\le \mathrm D'_m$, $\sup_m\mathrm D'_m<\infty$. Put
\begin{equation}\label{eq:rb-L0}
 \mathbb U_{\rm loc}:=\sup_m\ \sup_{\mathrm x\in \mathfrak J_{m-1}\cup \mathfrak J_m\cup \mathfrak J_{m+1}}\Bigl|\int_{\mathfrak J_m}\frac{(\mathrm x-\mathrm s)\,\mathrm B_m(\mathrm s)}{\pi((\mathrm x-\mathrm s)^2+1)}\,d\mathrm s\Bigr|,\qquad
 \mathbb O:=\sup_{m_0}\sum_{|m-m_0|\ge2}\frac{\mathfrak n_m^2}{\operatorname{dist}(\mathfrak J_{m_0},\mathfrak J_m)^2}.
\end{equation}
If $\mathbb U_{\rm loc}<\infty$ and $\mathbb O<\infty$, then $\sup_x|\mathbb Y_1\mathrm B(\mathrm x)-\mathrm c_{\mathrm B}|\le3\mathbb U_{\rm loc}+\frac{\mathbb O}{2\pi}\sup_m\mathrm D'_m$ for a constant $\mathrm c_{\mathrm B}$.
\end{proposition}

\begin{proof}
The term $\frac1\pi\int\frac{\mathrm s}{1+\mathrm s^2}\mathrm B_m(\mathrm s)\,d\mathrm s$ equals $\frac1\pi\int(\frac{\mathrm s}{1+\mathrm s^2}-\frac{\mathrm m_m}{1+\mathrm m_m^2})\mathrm B_m(\mathrm s)\,d\mathrm s$, where $\mathrm m_m$ is the centre of $\mathfrak J_m$. Since $|\frac d{d\mathrm s}\frac{\mathrm s}{1+\mathrm s^2}|\le\frac1{1+\mathrm s^2}$, its modulus is at most $\mathrm D'_m\mathfrak n_m^2/(2\pi(1+\operatorname{dist}(0,\mathfrak J_m)^2))$. If $m_0$ is the index with $0\in \mathfrak J_{m_0}$, then $\operatorname{dist}(0,\mathfrak J_m)\ge\operatorname{dist}(\mathfrak J_{m_0},\mathfrak J_m)$, so these terms are summable by $\mathbb O<\infty$. Their sum is the constant $\mathrm c_{\mathrm B}$. Let $\mathrm x\in \mathfrak J_{m_0}$. The three blocks with $|m-m_0|\le1$ contribute at most $3\mathbb U_{\rm loc}$. For $|m-m_0|\ge2$, $\int \mathrm B_m=0$ gives $\bigl|\int_{\mathfrak J_m}\mathrm k(\mathrm x-\mathrm s)\mathrm B_m(\mathrm s)\,d\mathrm s\bigr|\le \mathrm D'_m\mathfrak n_m\cdot\frac{\mathfrak n_m}2\sup_{\mathrm s\in \mathfrak J_m}|\mathrm k'(\mathrm x-\mathrm s)|$ with $\mathrm k(\mathrm v):=\frac{\mathrm v}{\pi(\mathrm v^2+1)}$. Moreover $|\mathrm k'(\mathrm v)|\le\frac1{\pi\mathrm v^2}$ and $|\mathrm x-\mathrm s|\ge\operatorname{dist}(\mathfrak J_{m_0},\mathfrak J_m)$. Sum over $m$.
\end{proof}

\begin{corollary}[Kadec and Avdonin, with the constant $\tfrac14$]\label{cor:rb-kadec-avdonin}
\textup{(a)} \cite{Kadec1964} If $|\mathrm d_k|\le \mathrm d_\star<\frac14$ for all $k$, then $\{e^{i\mathrm x_k\mathrm v}\}$ is a Riesz basis of $L^2(-\pi,\pi)$.\\
\textup{(b)} \cite{Avdonin1974} If $\mathscr X$ is uniformly discrete, $\sup_k|\mathrm d_k|\le \mathrm D$, and for a fixed $\mathfrak n$ the means of $\mathrm d_k$ over consecutive blocks of $\mathfrak n$ indices are at most $\mathrm d_\star<\frac14$ in modulus, then $\{e^{i\mathrm x_k\mathrm v}\}$ is a Riesz basis of $L^2(-\pi,\pi)$.
\end{corollary}

\begin{proof}
Both statements follow from Theorem~\CGref{thm:rb-variable-N} below: in (a) take $\mathfrak n_m\equiv1$, and in (b) take $\mathfrak n_m\equiv \mathfrak n$. In both cases the blocks at index distance $j\ge2$ from $\mathfrak J_{m_0}$ are at distance $(j-1)\mathfrak n$, so $\mathbb O\le2\sum_{j\ge2}\mathfrak n^2/((j-1)\mathfrak n)^2=\pi^2/3$ and $\mathbb V\le2\mathrm D\log(2+3\mathfrak n)$; for $\mathfrak n=1$, $\mathrm D_m=0$ and $\mathbb V=0$. For (a), $\mathrm d_\star<\frac14$ gives gaps $\ge1-2\mathrm d_\star>\frac12$, so $\mathscr X$ is uniformly discrete.
\end{proof}

\begin{theorem}[Avdonin's theorem with variable block length]\label{thm:rb-variable-N}
Let $\mathscr X$ be uniformly discrete with $\sup_k|\mathrm d_k|\le \mathrm D$. Let $\{\mathfrak B_m\}$ be a partition of $\mathbb Z$ into consecutive blocks of $\mathfrak n_m$ indices, and $\mathfrak J_m:=\bigcup_{k\in\mathfrak B_m}I_k$. Let $\bar{\mathrm d}_m$ be the mean of $\mathrm d_k$ over $\mathfrak B_m$, $\mathrm D_m:=\max_{k\in\mathfrak B_m}|\mathrm d_k-\bar{\mathrm d}_m|$, and $\mathfrak n_m^*:=\mathfrak n_{m-1}+\mathfrak n_m+\mathfrak n_{m+1}$. Assume
\begin{enumerate}[label=\textup{(\alph*)}]
\item $\sup_m|\bar{\mathrm d}_m|\le \mathrm d_\star<\frac14$;
\item $\mathbb O<\infty$, with $\mathbb O$ as in \CGeqref{eq:rb-L0};
\item $\mathbb V:=\sup_m\mathrm D_m\log(2+\mathfrak n_m^*)<\infty$.
\end{enumerate}
Then $\{e^{i\mathrm x_k\mathrm v}\}$ is a Riesz basis of $L^2(-\pi,\pi)$, and \CGeqref{eq:rb-A2const} holds with
\begin{equation}\label{eq:rb-U-variable}
 \mathbb U=16(1+\mathrm D)^3+\frac{6\mathbb V}\pi+\frac{\mathbb V\mathbb O}{2\pi\log2} .
\end{equation}
\end{theorem}

\begin{proof}
Write $\varPsi_{\mathscr X}=\mathrm A+\mathrm B_1+\mathrm B_2$ with $\mathrm A:=-\sum_m\bar{\mathrm d}_m\mathbf 1_{\mathfrak J_m}$, $\mathrm B_1:=\sum_k\mathrm b_k$ where $\mathrm b_k:=\mathsf i_k+\mathrm d_k\mathbf 1_{I_k}$, and $\mathrm B_2:=-\sum_m\sum_{k\in\mathfrak B_m}(\mathrm d_k-\bar{\mathrm d}_m)\mathbf 1_{I_k}$. Then $\|\mathrm A\|_\infty\le \mathrm d_\star$ by (a).

\emph{$\mathrm B_1$.} $\mathrm b_k$ has mean zero, $|\mathrm b_k|\le1+\mathrm D$, and support in an interval of length at most $2\mathrm D+1$ containing $\mathrm m_k:=k-\frac12$. With $\mathrm k$ as above, $\sup|\mathrm k|=\frac1{2\pi}$ and $|\mathrm k'(\mathrm v)|\le\frac1{\pi(1+\mathrm v^2)}$, so $|\int \mathrm k(\mathrm x-\mathrm s)\mathrm b_k(\mathrm s)\,d\mathrm s|=|\int(\mathrm k(\mathrm x-\mathrm s)-\mathrm k(\mathrm x-\mathrm m_k))\mathrm b_k(\mathrm s)\,d\mathrm s|\le(1+\mathrm D)(2\mathrm D+1)\min\bigl(\frac1\pi,(2\mathrm D+1)\sup_{\mathrm s\in\operatorname{supp}\mathrm b_k}|\mathrm k'(\mathrm x-\mathrm s)|\bigr)$. At most $8\mathrm D+5\le8(1+\mathrm D)$ indices have $|\mathrm x-\mathrm m_k|<2(2\mathrm D+1)$, and they contribute at most $16(1+\mathrm D)^3/\pi$. For the others $|\mathrm x-\mathrm s|\ge|\mathrm x-\mathrm m_k|/2$ on the support, so the term is at most $4(1+\mathrm D)^3\cdot\frac4\pi(4+(\mathrm x-\mathrm m_k)^2)^{-1}$. Since $\sum_k(4+(\mathrm x-\mathrm m_k)^2)^{-1}\le\frac12+\frac\pi2$, they contribute at most $\frac{16}\pi(\frac12+\frac\pi2)(1+\mathrm D)^3$. The total is below $16(1+\mathrm D)^3$. The modifying term of $\mathbb Y_1$ contributes a convergent sum of constants, as in Proposition~\CGref{prop:rb-variable-blocks}.

\emph{$\mathrm B_2$.} On $\mathfrak J_m$ it has mean zero and modulus at most $\mathrm D_m$. For $\mathrm x$ in $\mathfrak J_{m-1}\cup \mathfrak J_m\cup \mathfrak J_{m+1}$, $|\int_{\mathfrak J_m}\mathrm k(\mathrm x-\mathrm s)\mathrm B_2(\mathrm s)\,d\mathrm s|\le\frac{\mathrm D_m}\pi\int_{|u|\le \mathfrak n_m^*}\frac{|u|\,du}{1+u^2}\le\frac2\pi \mathrm D_m\log(2+\mathfrak n_m^*)\le\frac{2\mathbb V}\pi$. Moreover $\sup_m\mathrm D_m\le \mathbb V/\log2$. Proposition~\CGref{prop:rb-variable-blocks}, applied with the blocks $\mathfrak J_m$, $\mathrm D'_m=\mathrm D_m$ and $\mathbb U_{\rm loc}\le\frac{2\mathbb V}\pi$, gives $|\mathbb Y_1\mathrm B_2-\mathrm c|\le\frac{6\mathbb V}\pi+\frac{\mathbb V\mathbb O}{2\pi\log2}$.

Apply Theorem~\CGref{thm:rb-criterion} and Corollary~\CGref{cor:rb-riesz}(a).
\end{proof}

\begin{remark}[The price of variable length]\label{rem:rb-price}
For fixed $\mathfrak n$, conditions (b) and (c) hold automatically, and Theorem~\CGref{thm:rb-variable-N} is Avdonin's theorem. If $\mathfrak n_m\to\infty$, the averaging condition (a) is still only $|\bar{\mathrm d}_m|\le \mathrm d_\star<\frac14$. However, (c) then requires the oscillation inside the block to decay like $1/\log \mathfrak n_m$. In Theorem~\CGref{thm:rb-criterion} the oscillation enters only through the bounded-conjugate condition (ii), which allows cancellation inside the blocks and is weaker than (c). Without some condition of this kind the conclusion fails for growing blocks: a block of length $\mathfrak n$ with oscillation of size $\mathrm D_m$ can make $|\mathbb Y_1\mathrm B|$ of size $\mathrm D_m\log \mathfrak n$. Lyubarskii and Seip derive Kadec's theorem from their Muckenhoupt characterisation (\cite[Thm.~2]{LyubarskiiSeip1997} for $p=2$), and indicate that the same technique gives results of Avdonin type. The proof of Theorem~\CGref{thm:rb-criterion} is the sufficiency half of the Helson--Szeg\H{o} argument, applied on the line $\Im \mathfrak z=1$.

\emph{Literature.} The theorem of Hruščev, Nikol'skii and Pavlov \cite[p.~240]{HruscevNikolskiiPavlov1981}, in the form stated by Kozma and Lev \cite[Thm.~5.1]{KozmaLev2011}, characterises the Riesz bases of exponentials $E(\Lambda)$ of $L^2(0,a)$ by three conditions: $\Lambda$ is separated, $n_\Lambda(x)-ax$ belongs to $\mathrm{BMO}(\mathbb R)$, and a decomposition condition of Helson--Szeg\H{o} type on the generating function. Theorem~\CGref{thm:rb-criterion} with Corollary~\CGref{cor:rb-riesz}(a) is close to the sufficiency half of that theorem, written for the discrepancy $\varPsi_{\mathscr X}$, with the bounded-conjugate condition (ii) and the explicit constant \CGeqref{eq:rb-A2const}. Avdonin's theorem appears with a fixed block length in \cite{Avdonin1974} and in \cite[Thm.~3.3]{KozmaLev2011}. In the sources consulted (\cite{KozmaLev2011,LyubarskiiSeip1997,AvdoninIvanov2002}, and \cite{HruscevNikolskiiPavlov1981} through its statement in \cite{KozmaLev2011}) we did not find a version with block lengths that vary from block to block, as in Theorem~\CGref{thm:rb-variable-N}. The search was not exhaustive, so this is weak evidence, and novelty is not claimed.
\end{remark}

%======================================================================
\subsection{Application to lists near the RvM clock}\label{subsec:rb-application}

\subsubsection{Unfolding}\label{subsec:rb-unfolding}

\begin{definition}[Unfolded list]\label{def:rb-unfolding}
Let $\mathfrak t_0:=9.6769\ldots$ be the zero of $\NR$ in $(2\pi,\infty)$. For a $C$-admissible list $\boldsymbol\kappa$ with $\kappa_1>\mathfrak t_0$ (true for $\boldsymbol\tau$ and $\boldsymbol\gamma$), put $\mathrm x_k:=\NR(\kappa_k)>0$ for $k\ge1$ and $\mathrm x_{1-k}:=-\mathrm x_k$. The resulting symmetric sequence $\mathscr X_{\boldsymbol\kappa}$ has discrepancy $\varPsi_{\boldsymbol\kappa}:=\varPsi_{\mathscr X_{\boldsymbol\kappa}}$, and $\mathcal F_{\boldsymbol\kappa}:=F_{\mathscr X_{\boldsymbol\kappa}}$ with $\mathfrak z_0=0$. For the clock, $\mathscr X_{\boldsymbol\tau}=\mathbb Z+\frac12$ and $\mathcal F_{\boldsymbol\tau}=\mathcal F_0$.
\end{definition}

\begin{lemma}[Unfolded discrepancy]\label{lem:rb-psi-bound}
For a $C$-admissible list with $\kappa_1>\mathfrak t_0$, $|\mathrm x_k-(k-\frac12)|\le\frac C{2\pi}\log\frac{\tau_k+C}{2\pi}$ for $k\ge1$, so \CGeqref{eq:rb-H} holds. Moreover, for $\mathrm x=\NR(t)$ with $\mathrm x>0$,
\begin{equation}\label{eq:rb-psi-t}
 \varPsi_{\boldsymbol\kappa}(\mathrm x)=N_{\boldsymbol\kappa}(t)-n_{\RVM}(t),\qquad
 |N_{\boldsymbol\kappa}(t)-n_{\RVM}(t)|\le\frac C\pi\log\frac{t+C}{2\pi}+2,
\end{equation}
and $\varPsi_{\boldsymbol\kappa}(-\mathrm x)=-\varPsi_{\boldsymbol\kappa}(\mathrm x)$ for almost every $\mathrm x$.
\end{lemma}

\begin{proof}
By the mean value theorem and the monotonicity of $\NR'$, $|\NR(\kappa_k)-\NR(\tau_k)|\le C\NR'(\tau_k+C)$. For $\mathrm x>0$ the terms of \CGeqref{eq:rb-psi} with $k\le0$ vanish, because both $\mathrm x_k$ and $k-\frac12$ are negative. For $k\ge1$, $\mathrm x_k\le \mathrm x$ if and only if $\kappa_k\le t$, by monotonicity of $\NR$. If $N_{\boldsymbol\kappa}(t)=j$, then $\tau_j\le t+C$ and $\tau_{j+1}>t-C$, so $n_{\RVM}(t-C)\le j\le n_{\RVM}(t+C)$. By \CGeqref{eq:rb-Ntau-floor}, $n_{\RVM}(t+C)-n_{\RVM}(\max(t-C,2\pi))\le\NR(t+C)-\NR(\max(t-C,2\pi))+1$. This is at most $\frac{2C}{2\pi}\log\frac{t+C}{2\pi}+2$, the extra $1$ covering $t-C<2\pi$, where $\NR\ge-\frac18$ and $n_{\RVM}=0$. Oddness follows from the symmetry of $\mathscr X_{\boldsymbol\kappa}$ and of $\mathbb Z+\frac12$.
\end{proof}

\begin{proposition}[The condition of $\mathscr O$ gives the averaging hypothesis with blocks $\asymp\log^2T$]\label{prop:rb-O-blocks}
Let $\boldsymbol\kappa$ be $C$-admissible with $\kappa_1>\mathfrak t_0$, and assume the third condition of $\mathscr O$ (\CGref{eq:list-space}) in the form
\begin{equation}\label{eq:rb-O3}
 \Bigl|\int_{2\pi}^{R}\bigl(N_{\boldsymbol\kappa}(t)-\NR(t)\bigr)dt\Bigr|\le C_2\log R\qquad(R\ge3),
\end{equation}
with a constant $C_2>0$; this is no restriction, since \CGeqref{eq:rb-O3} with one constant holds with every larger one, and it is needed because the block length $\mathrm H(T):=16C_2\log T$ must be positive. Put $\mathrm H(T):=16C_2\log T$. There is $T_1=T_1(C,C_2)$ such that the following holds for the points $T_{m+1}:=T_m+\mathrm H(T_m)$ ($m\ge1$) and the blocks $\mathfrak J_m:=[\NR(T_m),\NR(T_{m+1}))$ with their mirror images $-\mathfrak J_m$. The step function $\mathrm A_{\boldsymbol\kappa}$, equal to the mean of $\varPsi_{\boldsymbol\kappa}$ on each $\pm \mathfrak J_m$ and to $0$ on $(-\NR(T_1),\NR(T_1))$, satisfies
\begin{equation}\label{eq:rb-block-mean}
 \|\mathrm A_{\boldsymbol\kappa}\|_\infty\le\tfrac3{16},\qquad |\mathfrak J_m|=\frac{8C_2}\pi\log^2T_m\,(1+o(1))\quad(m\to\infty).
\end{equation}
The remainder $\mathrm B_{\boldsymbol\kappa}:=\varPsi_{\boldsymbol\kappa}-\mathrm A_{\boldsymbol\kappa}$ has mean zero on every block. For $\boldsymbol\kappa=\boldsymbol\gamma$, \CGeqref{eq:rb-O3} holds unconditionally for some $C_2$.
\end{proposition}

\begin{proof}
Write $\mathfrak J=[\NR(T),\NR(T+\mathrm H))$. Then
\begin{equation}\label{eq:rb-aux-7}
\begin{split}
 \int_{\mathfrak J}\varPsi_{\boldsymbol\kappa}\,d\mathrm x&=\int_T^{T+\mathrm H}(N_{\boldsymbol\kappa}-n_{\RVM})(t)\,\NR'(t)\,dt\\
 &=\NR'(T)\!\int_T^{T+\mathrm H}\!(N_{\boldsymbol\kappa}-n_{\RVM})+\int_T^{T+\mathrm H}\!(N_{\boldsymbol\kappa}-n_{\RVM})(\NR'(t)-\NR'(T))\,dt .
\end{split}
\end{equation}
The last integral is at most $\mathrm H\cdot\frac{\mathrm H}{2\pi T}\bigl(\frac C\pi\log\frac{T+\mathrm H+C}{2\pi}+1\bigr)$, by \CGeqref{eq:rb-psi-t} and $\NR''(t)=\frac1{2\pi t}$.

Next, $\int_T^{T+\mathrm H}(N_{\boldsymbol\kappa}-n_{\RVM})=\int_T^{T+\mathrm H}(N_{\boldsymbol\kappa}-\NR)+\int_T^{T+\mathrm H}(\NR-n_{\RVM})$. The first piece is at most $2C_2\log(T+\mathrm H)$ by \CGeqref{eq:rb-O3}. For the second, substitute $\mathrm x=\NR(t)$. By \CGeqref{eq:rb-Ntau-floor}, $\NR-n_{\RVM}=\mathrm{saw}(\mathrm x):=\mathrm x-\lfloor \mathrm x+\frac12\rfloor$, which is $1$-periodic with mean zero and $|\mathrm{saw}|\le\frac12$, and $dt=\mathfrak g_1(\mathrm x)\,d\mathrm x$ with $\mathfrak g_1:=1/\NR'(t(\mathrm x))$ positive and decreasing. On a full period $[j-\frac12,j+\frac12)$, $|\int\mathrm{saw}\,\mathfrak g_1|=|\int\mathrm{saw}\,(\mathfrak g_1-\mathfrak g_1(j))|\le\frac14(\mathfrak g_1(j-\frac12)-\mathfrak g_1(j+\frac12))$, which telescopes. The two partial periods contribute at most $\mathfrak g_1(\NR(T))$. Hence $|\int_T^{T+\mathrm H}(\NR-n_{\RVM})|\le\frac5{4\NR'(T)}$.

Since $|\mathfrak J|\ge\NR'(T)\mathrm H$,
\begin{equation}\label{eq:rb-aux-8}
 \Bigl|\frac1{|\mathfrak J|}\int_{\mathfrak J}\varPsi_{\boldsymbol\kappa}\Bigr|\le\frac{2C_2\log(T+\mathrm H)}{\mathrm H}+\frac5{4\NR'(T)\mathrm H}+\frac{\mathrm H}{2\pi T\NR'(T)}\Bigl(\frac C\pi\log\frac{T+\mathrm H+C}{2\pi}+2\Bigr).
\end{equation}
With $\mathrm H=16C_2\log T$, the first term tends to $\frac18$ and the others tend to $0$, so the bound is at most $\frac3{16}$ for $T\ge T_1(C,C_2)$. The length is $|\mathfrak J_m|=\NR'(\mathrm s)\mathrm H(T_m)$ with $\mathrm s\in[T_m,T_{m+1}]$. The mirror blocks follow by oddness. For $\boldsymbol\gamma$: $N_{\boldsymbol\gamma}(t)=1+\vartheta(t)/\pi+S(t)$ \cite[\S9.3]{Titchmarsh1986}, $\int_0^RS=O(\log R)$ \cite[Thm.~9.9]{Titchmarsh1986}, and $1+\vartheta/\pi-\NR=O(1/t)$. Finally, $\boldsymbol\gamma$ is $C_\eta$-admissible, unconditionally, by \CGref{lem:bounded-displacement}. In particular \CGeqref{eq:rb-H} holds for $\mathscr X_{\boldsymbol\gamma}$, by Lemma~\CGref{lem:rb-psi-bound}.
\end{proof}

\begin{remark}[Proposition~\ref{prop:rb-O-blocks} is not an equivalence]\label{rem:rb-not-equivalent}
Proposition~\CGref{prop:rb-O-blocks} is an implication. The averaging condition on blocks of length $\asymp\log^2T$ does not imply \CGeqref{eq:rb-O3}: summing $|\int_{\mathfrak J_m}\varPsi|\le\frac3{16}|\mathfrak J_m|$ over the blocks up to height $R$ gives only $O(R\log R)$, not $O(\log R)$. The two conditions do not coincide.
\end{remark}

\begin{corollary}[Reduction to one condition]\label{cor:rb-one-condition}
For every $\boldsymbol\kappa$ as in Proposition~\CGref{prop:rb-O-blocks}: if $\mathbb Y_1\mathrm B_{\boldsymbol\kappa}$ is bounded up to a constant, then $|\mathcal F_{\boldsymbol\kappa}(\mathrm x+i)|^2\in(A_2)$, with constant \CGeqref{eq:rb-A2const} for $\mathrm d_\star=\frac3{16}$. Corollary~\CGref{cor:rb-riesz}(a) then holds for $\mathscr X_{\boldsymbol\kappa}$ if it is uniformly discrete, and Corollary~\CGref{cor:rb-riesz}(b) if it is a finite union of uniformly discrete sets.
\end{corollary}

\begin{proof}
Lemma~\CGref{lem:rb-psi-bound}, Proposition~\CGref{prop:rb-O-blocks} and Theorem~\CGref{thm:rb-criterion}.
\end{proof}

\subsubsection{The weight is a quotient of synthetic zetas}\label{subsec:rb-weight}

\begin{proposition}[Weight identity]\label{prop:rb-weight}
Let $\boldsymbol\kappa$ be $C$-admissible and satisfy \CGeqref{eq:rb-O3}. For $0<\hbar\le1$ and $t\ge2\pi$,
\begin{equation}\label{eq:rb-weight-identity}
 \log\Bigl|\frac{\zeta_{\boldsymbol\kappa}}{\zeta_{\RVM}}\Bigl(\tfrac12+\hbar+it\Bigr)\Bigr|=\int_{2\pi}^{\infty}\frac{(t-\mathrm s)\,(N_{\boldsymbol\kappa}-n_{\RVM})(\mathrm s)}{(t-\mathrm s)^2+\hbar^2}\,d\mathrm s+\mathrm c_{\boldsymbol\kappa}+\mathrm m_{\boldsymbol\kappa,\hbar}(t),
\end{equation}
where the integral converges as an improper integral and $\mathrm c_{\boldsymbol\kappa}$ is a constant. The remainder satisfies $|\mathrm m_{\boldsymbol\kappa,\hbar}(t)|\le \mathrm c(C,C_2)\log(2+t)/(1+t)$.
\end{proposition}

\begin{proof}
For $s=\frac12+\hbar+it$, $|\frac14+\kappa^2+q(s)|=|\kappa^2+(\hbar+it)^2|=|\kappa-t+i\hbar|\,|\kappa+t-i\hbar|$. Hence, with $\mathrm l_\hbar(u):=\frac12\log(u^2+\hbar^2)$,
\begin{equation}\label{eq:rb-aux-9}
 \log\Bigl|\frac{P_{\boldsymbol\kappa}}{P_{\boldsymbol\tau}}(s)\Bigr|=\sum_k\bigl[\mathrm l_\hbar(t-\kappa_k)-\mathrm l_\hbar(t-\tau_k)\bigr]+\sum_k\bigl[\mathrm l_\hbar(t+\kappa_k)-\mathrm l_\hbar(t+\tau_k)\bigr]-\sum_k\log\frac{\frac14+\kappa_k^2}{\frac14+\tau_k^2}.
\end{equation}
Put $\mathfrak g_1:=N_{\boldsymbol\kappa}-n_{\RVM}$ and $\mathfrak G_1(\mathrm s):=\int_{2\pi}^{\mathrm s}\mathfrak g_1$. By \CGeqref{eq:rb-O3}, by the sawtooth estimate in the proof of Proposition~\CGref{prop:rb-O-blocks} on $[\tau_1,\infty)$, and by $|\NR-n_{\RVM}|\le\frac12$ on $[2\pi,\tau_1]$, $|\mathfrak G_1(\mathrm s)|\le\mathrm c\log(2+\mathrm s)$. For $\mathrm f$ smooth with $\mathrm f'(\mathrm s)=O(1/\mathrm s)$ and $\mathrm f''(\mathrm s)=O(1/\mathrm s^2)$, the sums $\sum_{k\le k^*}[\mathrm f(\kappa_k)-\mathrm f(\tau_k)]=-\int \mathrm f'\mathfrak g_{1,k^*}$ converge as $k^*\to\infty$ to $-\int^\infty \mathrm f'\mathfrak g_1$, an improper integral. Integration by parts against $\mathfrak G_1$ shows both the convergence and that the tail contributions vanish. With $\mathrm f=\mathrm l_\hbar(t-\cdot)$ this gives the integral in \CGeqref{eq:rb-weight-identity}. With $\mathrm f=\log(\frac14+(\cdot)^2)$ it gives the constant $\mathrm c_{\boldsymbol\kappa}$. With $\mathrm f=\mathrm l_\hbar(t+\cdot)$ it gives $\mathrm m_{\boldsymbol\kappa,\hbar}(t)=\int_{2\pi}^\infty \mathrm k_1'(\mathrm s)\mathfrak G_1(\mathrm s)\,d\mathrm s$, where $\mathrm k_1(\mathrm s):=\frac{t+\mathrm s}{(t+\mathrm s)^2+\hbar^2}$ and $|\mathrm k_1'(\mathrm s)|\le(t+\mathrm s)^{-2}$. This yields the bound.
\end{proof}

\begin{remark}[Local scale]\label{rem:rb-local-scale}
With $\hbar=1$ and $\hbar_T:=\hbar/\NR'(T)=2\pi/\log(T/2\pi)$ (the notation of the proof of Theorem~\CGref{thm:rb-obstruction}), the substitution $\mathrm x=\NR(t)$ maps the kernel $\frac{(t-\mathrm s)\,d\mathrm s}{(t-\mathrm s)^2+\hbar_T^2}$ to $\frac{(\mathrm x-\mathrm x')\,d\mathrm x'}{(\mathrm x-\mathrm x')^2+1}$ to first order near $t=T$. Comparing \CGeqref{eq:rb-weight-identity} with \CGeqref{eq:rb-F-formula} at $\mathrm y=1$, the unfolded weight $|\mathcal F_{\boldsymbol\kappa}(\mathrm x+i)|^2$ near $\mathrm x=\NR(T)$ is therefore, to first order, $|\zeta_{\boldsymbol\kappa}/\zeta_{\RVM}|^2$ on the line $\Re s=\frac12+\hbar_T$, up to the bounded factor $|\mathcal F_0(\mathrm x+i)|^2$. The identity \CGeqref{eq:rb-weight-identity} is exact. The identification with the unfolded weight is local, and its quantitative form for the Poisson part is Step~3 of the proof of Theorem~\CGref{thm:rb-obstruction}.
\end{remark}

\subsubsection{Obstruction for the zeta list}\label{subsec:rb-obstruction}

\begin{theorem}[No $(A_2)$ for the unfolded zeta ordinates]\label{thm:rb-obstruction}
Let $\mathscr X_{\boldsymbol\gamma}$ be the unfolded list of the ordinates $\gamma_k>0$ of the nontrivial zeros of $\zeta$, counted with multiplicity, and let $\mathcal F$ be any entire function of exponential type whose zero set is $\mathscr X_{\boldsymbol\gamma}$, with multiplicities. Then $|\mathcal F(\mathrm x+i\hbar)|^2\notin(A_2)$ for every real $\hbar\ne0$. The proof does not assume the Riemann hypothesis.
\end{theorem}

\begin{proof}
Write $\varPsi:=\varPsi_{\boldsymbol\gamma}$ and $\mathcal F_{\boldsymbol\gamma}$ for the function of Definition~\CGref{def:rb-unfolding}. It is even and real by Lemma~\CGref{lem:rb-generating}(iii), because $\mathrm x_1=\NR(\gamma_1)=0.45\ldots>0$, so $0\notin \mathscr X_{\boldsymbol\gamma}$.

\emph{Step 1: reduction to $\mathcal F_{\boldsymbol\gamma}$ and $\hbar>0$.} $\mathcal F$ and $\mathcal F_{\boldsymbol\gamma}$ are entire of order at most $1$ (exponential type) with the same zeros, with multiplicity, so by Hadamard's factorisation theorem both are $e^{\text{linear}}$ times the same canonical product, and $\mathcal F/\mathcal F_{\boldsymbol\gamma}=e^{\mathrm a_1+\mathrm b_1\mathfrak z}$. Hence $|\mathcal F(\mathrm x+i\hbar)|^2=\mathrm c\,e^{2\mathrm b_1'\mathrm x}\mathrm w_0(\mathrm x)$ with $\mathrm b_1'=\Re\mathrm b_1$, a constant $\mathrm c>0$, and $\mathrm w_0(\mathrm x):=|\mathcal F_{\boldsymbol\gamma}(\mathrm x+i\hbar)|^2$, which is even in $\mathrm x$. Moreover $|\mathcal F_{\boldsymbol\gamma}(\mathrm x-i\hbar)|=|\mathcal F_{\boldsymbol\gamma}(\mathrm x+i\hbar)|$. Suppose $e^{2\mathrm b_1'\mathrm x}\mathrm w_0\in(A_2)$. Reflection preserves $(A_2)$, so $e^{-2\mathrm b_1'\mathrm x}\mathrm w_0\in(A_2)$ as well. By Step~2 below both logarithms are in BMO, so $4\mathrm b_1'\mathrm x\in\mathrm{BMO}$, which forces $\mathrm b_1'=0$. So we may assume $\mathcal F=\mathcal F_{\boldsymbol\gamma}$ and $\hbar>0$, and that $\mathrm w_0\in(A_2)$.

\emph{Step 2: BMO.} Let $\mathrm f:=\log \mathrm w_0$ and $\mathrm f_I$ its mean on $I$. By Jensen's inequality, $\frac1{|I|}\int_Ie^{\pm(f-\mathrm f_I)}\ge1$. The product of the two averages is at most $[\mathrm w_0]_{A_2}$, so each is at most $[\mathrm w_0]_{A_2}$. Since $u^2\le2e^{|u|}\le2(e^u+e^{-u})$,
\begin{equation}\label{eq:rb-aux-10}
 \frac1{|I|}\int_I|f-\mathrm f_I|^2\le4[\mathrm w_0]_{A_2}\qquad\text{for every interval }I .
\end{equation}
By \CGeqref{eq:rb-F-formula} and the boundedness of $\log|\mathcal F_0(\mathrm x+i\hbar)|$, $\mathrm g_0:=\pi\mathbb Y_\hbar\varPsi\in\mathrm{BMO}$. Now take $\mathbb A_\hbar(\mathfrak z):=\mathbb A_{\mathscr X}(\mathfrak z+i\hbar)$, in the notation of the proof of Lemma~\CGref{lem:rb-generating}, which is holomorphic on $\Im \mathfrak z>-\hbar$. Its real part on $\mathbb R$ is $\mathrm g_0$, and its imaginary part on $\mathbb R$ is $-\pi \mathbb P_\hbar\varPsi$. By \CGeqref{eq:rb-Q-growth}, $\Re\mathbb A_\hbar(\mathfrak z)=O(\log^2(4+|\mathfrak z|))$ on $\Im \mathfrak z\ge0$, and so is the Poisson integral $\mathbb P[\mathrm g_0]$. Hence $\Re\mathbb A_\hbar-\mathbb P[\mathrm g_0]$ is harmonic on $\Im \mathfrak z>0$, continuous with zero boundary values, and $O(\log^2|\mathfrak z|)$. After odd reflection it is harmonic on $\mathbb C$ and $o(|\mathfrak z|)$, hence constant, hence zero. Therefore $\mathbb A_\hbar=\mathbb P[\mathrm g_0]+i\mathbb Y[\mathrm g_0]+ic$ on $\Im \mathfrak z>0$, where $\mathbb P[\mathrm g_0]+i\mathbb Y[\mathrm g_0]$ is the holomorphic function with real part $\mathbb P[\mathrm g_0]$ and $\mathbb Y[\mathrm g_0](\mathrm x+iy)=\mathbb Y_{\mathrm y}\mathrm g_0(\mathrm x)$. Letting $\mathrm y\to0$, $-\pi \mathbb P_\hbar\varPsi=\tilde{\mathrm g}_0+c$ almost everywhere on $\mathbb R$, where $\tilde{\mathrm g}_0$ is the conjugate function of $\mathrm g_0$, the a.e.\ boundary value of $\mathbb Y_{\mathrm y}u$ \cite[Ch.~III]{Garnett2007}. Conjugation is bounded on BMO \cite[Ch.~VI]{Garnett2007}, so $\mathbb P_\hbar\varPsi\in\mathrm{BMO}$. By the John--Nirenberg inequality \cite[Ch.~VI]{Garnett2007}, there is $\mathrm c_{\rm JN}<\infty$ with
\begin{equation}\label{eq:rb-BMO-L2}
 \frac1{|I|}\int_I\bigl|\mathbb P_\hbar\varPsi-(\mathbb P_\hbar\varPsi)_I\bigr|^2\le\mathrm c_{\rm JN}\qquad\text{for every interval }I .
\end{equation}

\emph{Step 3: back to the ordinates.} By \cite[\S9.3]{Titchmarsh1986}, $N_{\boldsymbol\gamma}(t)=1+\vartheta(t)/\pi+S(t)$, with $S(t)=O(\log t)$ \cite[Thm.~9.4]{Titchmarsh1986}. Lemma~\CGref{lem:rb-psi-bound} and \CGeqref{eq:rb-Ntau-floor} then give, for $\mathrm x=\NR(t)>0$, $\varPsi(\mathrm x)=S(t)+\mathrm r(\mathrm x)$ with $|\mathrm r|\le \mathrm c_e$. Here $\mathrm c_e$ bounds $|1+\vartheta/\pi-\NR|+\frac12$ on $t\ge2\pi$. Let $\mathfrak t_0:=\NR^{-1}(0)$ and $T\to\infty$, with $t\in[T,2T]$ and $\hbar_t:=\hbar/\NR'(t)$, $\hbar_T:=\hbar/\NR'(T)$.
\begin{enumerate}[label=(\alph*)]
\item The part of $\mathbb P_\hbar\varPsi(\NR(t))$ coming from $\mathrm x'<0$ is $O(\log T/T)$, since $|\varPsi(\mathrm x')|\le \mathrm c\log(2+|\mathrm x'|)$. The part coming from $\mathrm r$ is at most $\mathrm c_e$. Substituting $\mathrm x'=\NR(\mathrm s)$, the rest is $\int_{\mathfrak t_0}^\infty \mathbb P^{\NR}_t(\mathrm s)S(\mathrm s)\,d\mathrm s$ with $\mathbb P^{\NR}_t(\mathrm s):=\mathbb P_\hbar(\NR(t)-\NR(\mathrm s))\NR'(\mathrm s)$, where $\mathbb P_\hbar(\mathrm v):=\frac \hbar{\pi(\mathrm v^2+\hbar^2)}$.
\item Put $\mathbb P^{\circ}_t(\mathrm s):=\mathbb P_{\hbar_t}(t-\mathrm s)=\mathbb P_\hbar(\NR'(t)(t-\mathrm s))\NR'(t)$. For $|\mathrm s-t|\le T^{1/3}$, Taylor's formula with $\NR''=\frac1{2\pi \mathrm s}$ gives $|\NR(t)-\NR(\mathrm s)-\NR'(t)(t-\mathrm s)|\le\mathrm d_T:=T^{-1/3}/\pi$. Hence $|\mathbb P^{\NR}_t(\mathrm s)-\mathbb P^{\circ}_t(\mathrm s)|\le(5\mathrm d_T/\hbar)\mathbb P^{\circ}_t(\mathrm s)$ for $T$ large. With $|S|\le \mathrm c\log(3T)$ on that range, the core difference is $O(T^{-1/3}\log T)$. For $|\mathrm s-t|>T^{1/3}$ we use $\mathbb P_\hbar(\mathrm v)\le\frac \hbar{\pi \mathrm v^2}$, the convexity of $\NR$ and $|S(\mathrm s)|\le \mathrm c\log(2+\mathrm s)$. If $\mathrm s>t+T^{1/3}$, then $\NR(\mathrm s)-\NR(t)\ge\NR'(t)(\mathrm s-t)\ge \mathrm s-t$ for $T$ large, so $\int_{t+T^{1/3}}^\infty \mathbb P^{\NR}_t(\mathrm s)|S(\mathrm s)|\,d\mathrm s\le\frac{\mathrm c\hbar}\pi\int_{T^{1/3}}^\infty\frac{\NR'(t+\mathrm v)\log(2+t+\mathrm v)}{\mathrm v^2}\,d\mathrm v=O(T^{-1/3}\log^2T)$, splitting the $\mathrm v$-integral at $\mathrm v=T$. If $\mathfrak t_0\le \mathrm s<t-T^{1/3}$, then $\NR(t)-\NR(\mathrm s)\ge\NR'(\mathrm s)(t-\mathrm s)$ and $\NR'(\mathrm s)\ge\NR'(\mathfrak t_0)>0$, so $\mathbb P^{\NR}_t(\mathrm s)\le\frac{h}{\pi\NR'(\mathrm s)(t-\mathrm s)^2}\le\frac{h}{\pi\NR'(\mathfrak t_0)(t-\mathrm s)^2}$ and $\int_{\mathfrak t_0}^{t-T^{1/3}}\mathbb P^{\NR}_t|S|\le\frac{\mathrm c\hbar\log(3T)}{\pi\NR'(\mathfrak t_0)}\int_{T^{1/3}}^\infty \mathrm v^{-2}\,d\mathrm v=O(T^{-1/3}\log T)$. The same bounds hold for $\mathbb P^{\circ}_t$ with $\hbar_t$ in place of $\hbar$ and $\NR'(t)$ in place of $\NR'(\mathrm s)$, the region $\mathrm s\gg T$ being handled by the split at $\mathrm v=T$. So that part tends to $0$ as well.
\item $\|\mathbb P_a-\mathbb P_b\|_{L^1}\le\frac2\pi|\log(a/b)|$, because $\int|\partial_c\mathbb P_c|=\frac2{\pi c}$. Since $\hbar_t/\hbar_T\in[1-\frac{2\log2}{\log(T/2\pi)},1]$, $|\int(\mathbb P_{\hbar_t}-\mathbb P_{\hbar_T})(t-\mathrm s)S(\mathrm s)\,d\mathrm s|=O(\log T/\log T)=O(1)$, again with the tails controlled by $|S(\mathrm s)|\le \mathrm c\log(2+\mathrm s)$.
\item $\|\mathbb P_{\hbar_T}\ast S-S\|_{L^2[T,2T]}=O(\sqrt T)$, with $S:=0$ on $(-\infty,\mathfrak t_0)$; this is Lemma~\CGref{lem:rb-smoothing} below.
\end{enumerate}
Together, (a)--(d) give $\|\mathbb P_\hbar\varPsi\circ\NR-S\|_{L^2[T,2T]}=O(\sqrt T)$.

\emph{Step 4: contradiction.} Let $\mathfrak I_T:=[\NR(T),\NR(2T)]$ and $\bar{\mathrm p}_T:=(\mathbb P_\hbar\varPsi)_{\mathfrak I_T}$. By \CGeqref{eq:rb-BMO-L2}, $\int_T^{2T}|\mathbb P_\hbar\varPsi(\NR(t))-\bar{\mathrm p}_T|^2\NR'(t)\,dt\le\mathrm c_{\rm JN}|\mathfrak I_T|\le\mathrm c_{\rm JN}T\NR'(2T)$, so $\int_T^{2T}|\mathbb P_\hbar\varPsi(\NR(t))-\bar{\mathrm p}_T|^2\,dt\le2\mathrm c_{\rm JN}T$. With Step~3, $\int_T^{2T}|S-\bar{\mathrm p}_T|^2=O(T)$. But
\begin{equation}\label{eq:rb-aux-11}
 \int_T^{2T}|S-\mathrm c|^2\ge\int_T^{2T}S^2-\frac1T\Bigl(\int_T^{2T}S\Bigr)^2=\frac1{2\pi^2}T\log\log T\,(1+o(1))
\end{equation}
for every constant $\mathrm c$, by Selberg's theorem $\int_0^TS^2=\frac{T}{2\pi^2}\log\log T+O(T(\log\log T)^{1/2})$ \cite{Selberg1946}, with the normalisation $S(t)=\pi^{-1}\arg\zeta(\frac12+it)$ used here and in \cite[p.~225]{Fujii1974}, and $\int_T^{2T}S=O(\log T)$ \cite[Thm.~9.9]{Titchmarsh1986}. This is a contradiction.
\end{proof}

\begin{lemma}[Smoothing $S$ at scale $1/\log T$]\label{lem:rb-smoothing}
Let $\hbar>0$ be fixed, $\hbar_T:=2\pi \hbar/\log(T/2\pi)$, and $S:=0$ on $(-\infty,\mathfrak t_0)$. Then $\|\mathbb P_{\hbar_T}\ast S-S\|_{L^2[T,2T]}=O_\hbar(\sqrt T)$, unconditionally.
\end{lemma}

\begin{proof}
Fujii's Main Theorem \cite[Main Theorem]{Fujii1974}, proved unconditionally by modifying Selberg's method, states for $q=1$, $k=1$, $T^{1/2+a}\le H\le T$ and $0<u\le H-(H/\sqrt T)^{1/8}$ that $\int_T^{T+H}(S(t+u)-S(t))^2dt=c_1H\cdot2\log(2+u\log T)\,(1+O((\log(2+u\log T))^{-1/2}))$ with an absolute $c_1$. The relative error is $O(1)$ when $u\log T$ is bounded, so the statement carries no asymptotic information there, and we use it only as the upper bound it contains: for $0<u\le1$, $T\ge T_0$ and $H=T$,
\begin{equation}\label{eq:rb-fujii}
 \int_T^{2T}\bigl(S(t+u)-S(t)\bigr)^2dt\le \mathrm c_FT\log(2+u\log T),
\end{equation}
with an absolute constant $\mathrm c_F$.
The range $0<u\le H-(H/\sqrt T)^{1/8}=T-T^{1/16}$ of \cite{Fujii1974} contains $0<u\le1$, and the same bound holds for $-1\le u<0$: the substitution $t\mapsto t-u$ moves the integral to $[T-|u|,2T-|u|]\subset[T/2,2T]$, covered by \CGeqref{eq:rb-fujii} at the base points $T/2$ (with $H=T/2$) and $T$ (with $H=T$). By Minkowski's inequality, $\|\mathbb P_{\hbar_T}\ast S-S\|_{L^2[T,2T]}\le\int \mathbb P_{\hbar_T}(u)\|S(\cdot-u)-S\|_{L^2[T,2T]}\,du$. For $0<|u|\le1$ the norm is at most $(2\mathrm c_FT\log(2+|u|\log T))^{1/2}$. For $1<|u|\le T/2$, Selberg's theorem $\int_0^TS^2=\frac1{2\pi^2}T\log\log T+O(T(\log\log T)^{1/2})$ \cite{Selberg1946} bounds it by $\mathrm c(T\log\log T)^{1/2}$. For $|u|>T/2$, $|S(t)|\le \mathrm c\log(2+t)$ \cite[Thm.~9.4]{Titchmarsh1986} bounds it by $\mathrm c\sqrt T\log(3T+|u|)$. Substituting $u=\hbar_Tv$ in the first range, where $\hbar_T\log T\to2\pi \hbar$, and using $\int_{|u|>1}\mathbb P_{\hbar_T}=O(\hbar_T)$ and $\int_{|u|>T/2}\mathbb P_{\hbar_T}(u)\log(3T+|u|)\,du=O(\hbar_T\log T/T)$, all three contributions are $O(\sqrt T)$.
\end{proof}

\begin{remark}[A second proof]\label{rem:rb-fujii}
Lemma~\CGref{lem:rb-smoothing} also follows from Selberg's mean-square approximation of $S(t)$ by $S_{\mathrm x_0}(t):=-\frac1\pi\sum_{p\le \mathrm x}p^{-1/2}\sin(t\log p)$, $\mathrm x_0=T^\delta$ \cite{Selberg1946}, and the Montgomery--Vaughan mean value theorem \cite[Cor.~3]{MontgomeryVaughan1974}. Indeed $\mathbb P_{\hbar_T}\ast S_{\mathrm x_0}-S_{\mathrm x_0}=-\frac1\pi\sum_{p\le \mathrm x}p^{-1/2}(e^{-\hbar_T\log p}-1)\sin(t\log p)$ has mean square $\ll T\,\hbar_T^2\log^2\mathrm x\ll T$ on $[T,2T]$, and the remainder $S-S_{\mathrm x_0}$ is handled by Minkowski's inequality. Fujii's announcement \cite{Fujii1975} states the moment lemma of \cite{Fujii1974} in the form used in \cite[Lemma~1]{Fujii1975b}.
\end{remark}

\begin{corollary}[Consequences for the zeta list]\label{cor:rb-obstruction}
\textup{(a)} $\{e^{ix_k\xi}\}_{x_k\in \mathscr X_{\boldsymbol\gamma}}$ is not a Riesz basis of $L^2(-\pi,\pi)$.\\
\textup{(b)} If $\mathscr X_{\boldsymbol\gamma}$ is a finite union of uniformly discrete sets, its generalised divided differences are not a Riesz basis of $L^2(-\pi,\pi)$.\\
\textup{(c)} For every decomposition $\varPsi_{\boldsymbol\gamma}=A+B$ with $\|A\|_\infty<\frac14$, $\mathbb Y_1B$ is unbounded. In particular this holds for the blocks of Proposition~\CGref{prop:rb-O-blocks}, where the averaging hypothesis is satisfied.
\end{corollary}

\begin{proof}
(a) and (b): the criteria \cite[Thm.~1, Thm.~3]{AvdoninIvanov2002} are necessary and sufficient, and require a generating function satisfying Helson--Szeg\H{o}, equivalently $(A_2)$, on some line $\Im \mathfrak z=\hbar$. The line $\hbar=0$ is excluded, since $|\mathcal F(\mathrm x)|^{-2}$ is not locally integrable at a real zero. All lines with $\hbar\ne0$ are excluded by Theorem~\CGref{thm:rb-obstruction}. (c): otherwise Theorem~\CGref{thm:rb-criterion} would give $|\mathcal F_{\boldsymbol\gamma}(\mathrm x+i)|^2\in(A_2)$.
\end{proof}

\begin{heuristic}[Graduated version]\label{heur:rb-graduated}
The proof of Theorem~\CGref{thm:rb-obstruction} uses intervals $\mathfrak I_T$ of length $\asymp T\log T$. On a window of $L$ unfolded points at height $T$, the mean square oscillation of $\mathbb P_\hbar\varPsi$ is governed by the variance of increments of $S$ over $t$-intervals of length $L/\NR'(T)$. By Fujii's asymptotics this is expected to be of size $\frac1{\pi^2}\log L$. This suggests that local $(A_2)$ constants on windows of $L$ points grow with $L$ but stay bounded for $L$ bounded. That is the graduated stability that a Fredholm formulation would need. None of this is proved here. The computations of Section~\CGref{subsec:rb-numerics} measure the $\ell^\infty$ norm of the inverse of square finite sections directly and find growth like $T^{1.8}$ up to $T=487$.
\end{heuristic}

\begin{remark}[Bearing on route B]\label{rem:rb-routeB-stability}
Theorem~\CGref{thm:rb-obstruction} rules out a uniform Riesz-basis (Fredholm) structure for the unfolded zeta list on the whole line, and hence a linear stability argument with constants independent of the height. It says nothing about statement B or B$'$. The compactness step of route B (Proposition~\CGref{prop:rb-compactness}) uses no stability, and uniqueness holds in every tube (Theorem~\CGref{thm:rb-uniqueness}); the open part of route B and of route B$'$ is existence alone (Corollaries~\CGref{cor:rb-equivalence} and~\CGref{cor:rb-Bprime}).
\end{remark}

\begin{remark}[Literature on the zeta ordinates as an exponential system]\label{rem:rb-obstruction-literature}
We have not found any study of Riesz bases of exponentials, or of the condition $(A_2)$, for the ordinates of $\zeta$; the search was not exhaustive. The closest work found is the Fourier interpolation theorem of Bondarenko, Radchenko and Seip \cite[Thm.~1.1]{BondarenkoRadchenkoSeip2023}: an even function $f$ analytic in a strip $|\Im z|<\frac12+\epsilon$, with an integrability condition, is recovered by an explicit series from the values $f^{(j)}((\rho-\frac12)/i)$, $0\le j<m(\rho)$, at the nontrivial zeros $\rho$ with positive imaginary part, and the values of $\widehat f$ at the points $\log n/(4\pi)$. The theorem is unconditional (zeros off the line enter as complex nodes), it rests on a strengthened form of Knopp's abundance principle for Dirichlet series with a functional equation \cite{Knopp1994}, and removing a single node breaks it. There the zeros and the logarithms of the integers are dual families of nodes, as the ordinates and the points $\log n$ are in the rung system \CGeqref{eq:rb-rung}; Theorem~\CGref{thm:rb-obstruction} concerns instead the Riesz-basis property of the exponentials of the unfolded ordinates, which is not used there.
\end{remark}

%======================================================================
\subsection{Numerical certificates for finite sections}\label{subsec:rb-numerics}

This section reports computations. Except for Theorem~\CGref{thm:rb-verified}, which is proved by the ball-arithmetic computation of Section~\CGref{subsec:rb-num-verified}, they are not part of any proof; their role is to measure the quantities of Theorems~\CGref{thm:rb-certificate} and~\CGref{thm:rb-staged} on actual sections and to locate where the staged construction fails. The programs and their outputs are listed in Appendix~\CGref{app:code-routeB} and in \texttt{code/routeB/} of the repository. Except in Section~\CGref{subsec:rb-num-verified}, floating-point arithmetic was used and no interval arithmetic was run.

\subsubsection{Set-up}\label{subsec:rb-num-setup}

The tests are Gaussians $\varphi_{c,\mathrm w}(u):=e^{-(u-c)^2/2\mathrm w^2}$ with centre $c\ge9\mathrm w$, so that their restriction to $u>0$ is, to working precision, supported in $(0,\infty)$; their cosine transform is $\mathsf c_{\varphi_{c,\mathrm w}}(t)=\mathrm w\sqrt{2\pi}\,e^{-(\mathrm wt)^2/2}\cos(ct)$. A Gaussian is not compactly supported, so these tests are not elements of $C_c^\infty(0,\infty)$ and the computed quantities are not literally those of Theorem~\CGref{thm:rb-certificate}, whose tests are in that space. The pairings against a compactly supported bump $\varphi_{c,\mathrm w}\chi$, with $\chi\in C_c^\infty(0,\infty)$ equal to $1$ on $[c/2,3c/2]$, differ from the computed ones by the pairings against the tail $\varphi_{c,\mathrm w}(1-\chi)$, which were not computed; the alternative, an extension of Theorem~\CGref{thm:rb-certificate} to Schwartz tests, is not carried out here. For this reason, and because of the floating-point arithmetic, the certificates reported below are \emph{candidate certificates}: floating-point evidence for the inequalities of the theorem evaluated on tests outside its hypothesis, not verified statements. Stage $m$ uses the width $\mathrm w_m:=0.08\cdot0.85^{m-1}$, the height $T_m:=4/\mathrm w_m$, and the new tests $\varphi_{c,\mathrm w_m}$ with centres $c\in[9\mathrm w_m,\,\log(\max(T_m,7)/2\pi)+2]$ in steps of $\mathrm w_m$; the test set $\mathscr T_m$ is the union of the levels $\le m$. The unknowns of stage $m$ are the ordinates with $\tau_k\le2.4/\mathrm w_m$ (``interior''); the factor $e^{-(\mathrm w_m\kappa)^2/2}$ makes the others invisible to the tests of level $m$. The masses are $\mathsf a^\zeta_n$ for all $n\le e^{c+9\mathrm w}$, and $\mathsf k$ is \CGref{eq:transport-defs}. The pairings $\langle\mathsf k,\varphi\rangle$ were computed by adaptive quadrature. The zeta ordinates $\gamma_k$, $k\le400$ ($\gamma_{400}=679.74$), computed with \texttt{mpmath}, enter only in the diagnostic columns; no construction uses them.

Two checks fix the normalisation: $\langle\mathfrak R(\boldsymbol\gamma,\mathsf a^\zeta),\varphi\rangle\approx10^{-16}$ for all tests tried, and $\langle\mathfrak R(\boldsymbol\tau,\mathsf a^\zeta),\varphi\rangle$ between $10^{-3}$ and $10^{-1}$.

\subsubsection{Initialisation: the clock is outside the basin, the Gonek lists are inside}\label{subsec:rb-num-init}

\emph{Receding horizon from the clock.} Starting at $\boldsymbol\kappa=\boldsymbol\tau$, each stage solved the overdetermined least-squares problem for the interior ordinates with all tests of levels $\le m$, with bounds $|\kappa_k-\tau_k|\le3$ and the non-interior ordinates frozen at their current values; every residual reported for these runs is therefore the residual of a frozen-tail finite model, not of the infinite system. This is the overdetermined geometry $\mathrm t_m\ge k^*_m$: the Jacobian has a left inverse (the pseudo-inverse) and Gauss--Newton is a least-squares iteration. Theorem~\CGref{thm:rb-certificate} requires a right inverse, $\mathbb J_m\mathbb Q_m=I_{\mathrm t_m}$, which forces $\mathrm t_m\le k^*_m$; the two geometries meet only in the square case $\mathrm t_m=k^*_m$, which is the one used for the certificates of Section~\CGref{subsec:rb-num-certified}. The receding-horizon runs are therefore diagnostics of the basin, not evaluations of the theorem. The iteration stalls at the first Gram-type failure of the clock: the fit returns $\kappa_{34}=109.75$ where $\gamma_{34}=111.03$ ($\tau_{34}=110.04$), and $\kappa_{41}=125.50$ where $\gamma_{41}=124.26$; from $m=9$ on the residual on the lagged levels $\le m-4$ grows from $2\cdot10^{-5}$ to $4\cdot10^{-2}$ and the front of ordinates within $10^{-2}$ of $\boldsymbol\gamma$ stays at $k\approx31$.

\emph{Receding horizon from a Gonek list.} The same iteration started at $\boldsymbol\kappa^{(\mathrm X)}$ with $\mathrm X=8$ ($983$ roots up to $t=1400$, that is, every solution of \CGeqref{eq:rb-gonek-phase} from $t\approx14$ on; this is not the list of Lemma~\CGref{lem:rb-gonek-tube}, whose prefix is the clock below $\mathfrak t_8\approx2\pi e^{2\mathrm n_8+2}$ with $\mathrm n_8\approx7$, so that the lemma's tube does not cover the list used here, whose monotonicity was checked directly; the roots were computed from the Riemann--Siegel expansion of $\vartheta$ and the sum \CGeqref{eq:rb-alphaX}) converged through $m=16$, $T_{16}=572$, $165$ interior ordinates and $4985$ tests: all interior ordinates within $10^{-2}$ of $\gamma_k$, $90$ of them within $10^{-4}$, the first twenty within $5\cdot10^{-6}$, residual on the lagged levels between $10^{-6}$ and $9\cdot10^{-6}$ at every stage, and $\max_k|\kappa_k-\tau_k|=0.99$ throughout. The residual over all tests of level $\le m$ stays at $2\cdot10^{-3}$, which is the contribution of the frozen edge.

\emph{Why.} The basin of the exact solution, measured by random perturbations of $\boldsymbol\gamma$ in the interior block, has radius $0.1$ to $0.25$ mean spacings. The clock lies at up to $0.66$ spacings from $\boldsymbol\gamma$ for $k\le400$, with $50$ indices beyond $0.25$; the lists $\boldsymbol\kappa^{(\mathrm X)}$ lie within $0.184$ ($\mathrm X=3$), $0.081$ ($\mathrm X=5$), $0.053$ ($\mathrm X=8$), $0.047$ ($\mathrm X=12$) and $0.035$ ($\mathrm X=20$) spacings for $k\le160$, and within $0.103$ spacings for $k\le400$ at $\mathrm X=8$ (Figure~\CGref{fig:rb-gonek}). The $\ell^2$ conditioning of the interior block at $\boldsymbol\gamma$ is moderate and roughly uniform up to $T=299$: condition numbers $17$ to $75$, smallest singular value $0.08$ to $0.15$. The coupling of the interior to the frozen edge, $\|\mathbb J_I^+\mathbb J_E\|_{\infty\to\infty}$, is between $0.25$ and $0.47$, and the shift it induces on the low ordinates decays from $2\cdot10^{-2}$ at $m=2$ to $10^{-5}$ for $m\ge10$.

\begin{figure}[tbp]
\centering
\includegraphics[width=0.9\textwidth]{fig18_gonek_vs_clock.png}
\caption{Distance to the zeta ordinates, in units of the local mean spacing $2\pi/\log(\gamma_k/2\pi)$, of the midpoint clock and of the Gonek lists $\boldsymbol\kappa^{(\mathrm X)}$ for $\mathrm X=3,8,20$ ($k\le160$ for $\mathrm X=3,20$; $k\le400$ for $\mathrm X=8$). The dashed line is the observed radius of the Newton basin. The ordinates $\gamma_k$ are used here only to measure the distance.}
\label{fig:rb-gonek}
\end{figure}

\subsubsection{Certificates at the Gonek list: the direct bounds fail}\label{subsec:rb-num-kantorovich}

At $\boldsymbol\kappa^{(8)}$, with the interior unknowns of stage $m$ and three choices of the right inverse, the quantities of Theorem~\CGref{thm:rb-certificate} were computed for $m=1,\dots,14$ ($T$ from $50$ to $414$).

\begingroup\small
\begin{longtable}{p{0.40\textwidth}p{0.22\textwidth}p{0.30\textwidth}}
\toprule
Right inverse and norm & Certificate index & Range over $m=1,\dots,14$\\
\midrule
square subsystem (tests chosen by pivoted QR), $\ell^\infty$ & $2\mathbb L_m^2\mathbb H_m\mathsf h_m$ & $56\to1.2\cdot10^{7}$\\
pseudo-inverse of the full system, $\ell^2$, $\beta_2=1/\sigma_{\min}$ & $\beta_2\mathrm K_2\delta_2$ (Kantorovich) & $0.47\to233$\\
pseudo-inverse of the full system, $\ell^\infty$, diagonal Hessian & $2\mathrm c_{\rm diag}\mathsf h$ & $1.3\to1.6\cdot10^{3}$\\
\bottomrule
\end{longtable}
\endgroup

None of the certificates is satisfied, not even at $T=50$. Yet one Gauss--Newton step from $\boldsymbol\kappa^{(8)}$ reduces $\max_k|\kappa_k-\gamma_k|$ by a factor between $4$ and $15$ at every height ($0.088\to0.016$ at $T=572$), and $\sigma_{\min}$ in $\ell^2$ stays between $1/17$ and $1/5$. The worst-case bounds lose a factor $10^3$ to $10^4$ through cancellation: $\mathbb H_m$ counts every unknown with the full weight $\int|\varphi|u^2$, while the actual second-order term of $\mathbb F_m$ along the Newton step is a sum of $k^*_m$ oscillating terms.

\subsubsection{Floating-point candidate certificates for exact square sections}\label{subsec:rb-num-certified}

The iterate of Section~\CGref{subsec:rb-num-init} was polished, stage by stage, by Newton's method on the square system ($n$ tests selected by pivoted QR on $\mathbb J^{\!T}$, $n$ interior unknowns) to residual $10^{-13}$. At the polished point $x_0$ the certificate of Theorem~\CGref{thm:rb-certificate} was evaluated in the sharper form that uses the separable structure of $F_m$: with $\mathbb Q$ an approximate inverse of the square Jacobian $\mathbb J$ at $x_0$, computed by solving $\mathbb J\mathbb Q=I$ with an LU factorization, $\mathsf d_0:=\|I-\mathbb Q\mathbb J\|_{\infty\to\infty}$ its defect, $\mathbb J'$ the matrix of second derivatives $\partial^2_{\kappa_k}\mathsf c_\varphi$ at $x_0$, and $\mathbb A^{(3)}(\mathrm r)$ an entrywise bound for the third derivatives on $|x-x_0|\le\mathrm r$, recomputed for each $\mathrm r$, the map $\mathbb T(x)=x-\mathbb Q\mathbb F(x)$ contracts the ball of radius $\mathrm r$ in $\ell^\infty$ if
\begin{equation}\label{eq:rb-num-banach}
 \mathsf l(\mathrm r):=\mathsf d_0+\|\mathbb Q\mathbb J'\|_{\infty\to\infty}\,\mathrm r+\tfrac12\,\bigl\||\mathbb Q|\mathbb A^{(3)}(\mathrm r)\bigr\|_{\infty\to\infty}\,\mathrm r^2<1,\qquad \mathsf h\le(1-\mathsf l(\mathrm r))\,\mathrm r,
\end{equation}
in which case the ball contains an exact solution of the square system; here $\mathrm t_m=k^*_m$, so the right inverse of Theorem~\CGref{thm:rb-certificate} is an inverse.  The defect enters $\mathsf l$ because $D\mathbb T(x)=(I-\mathbb Q\mathbb J)+\mathbb Q(\mathbb J-D\mathbb F(x))$, and a fixed point of $\mathbb T$ is a zero of $\mathbb F$ because $\mathsf d_0<1$ makes $\mathbb Q\mathbb J$, hence $\mathbb Q$, invertible. The sampled margin of a stage is the largest value of $(1-\mathsf l(\mathrm r))\mathrm r$ over $60$ radii $\mathrm r$ spaced logarithmically in $[10^{-12},10^{-1}]$: a sampled estimate, not a maximum over all $\mathrm r$, of the largest residual that \CGeqref{eq:rb-num-banach} could accept at that point.

\begingroup\small
\begin{longtable}{rrrrrrr}
\toprule
$m$ & $T_m$ & $n$ & $\mathbb L_m=\|Q\|_{\infty\to\infty}$ & sampled margin & $\mathsf s_m$ on all tests $\le m$ & $\max_k|\kappa_k-\gamma_k|$\\
\midrule
1 & 50 & 4 & $41$ & $7.0\cdot10^{-3}$ & $7.7\cdot10^{-8}$ & $7\cdot10^{-7}$\\
4 & 81 & 9 & $50$ & $3.1\cdot10^{-3}$ & $2.6\cdot10^{-7}$ & $1\cdot10^{-6}$\\
8 & 156 & 26 & $150$ & $7.5\cdot10^{-4}$ & $9.6\cdot10^{-7}$ & $1\cdot10^{-5}$\\
12 & 299 & 68 & $460$ & $2.2\cdot10^{-4}$ & $5.5\cdot10^{-6}$ & $4\cdot10^{-5}$\\
15 & 487 & 133 & $1600$ & $6.7\cdot10^{-5}$ & $3.3\cdot10^{-3}$ & $4\cdot10^{-2}$\\
16 & 572 & 165 & -- & square Jacobian singular & -- & --\\
\bottomrule
\end{longtable}
\endgroup

For stages $1$ to $15$ the floating-point evaluation of \CGeqref{eq:rb-num-banach} holds with $\mathrm r$ equal to the polishing residual ($10^{-15}$ to $5\cdot10^{-13}$), with defect $\mathsf d_0\le6.3\cdot10^{-13}$ and $\mathsf l(\mathrm r)\le1.9\cdot10^{-9}$.  Radii of this size are at the level of rounding and are numerically fragile as individual numbers; what measures the stability of the test is the sampled margin, which shows that at the same point every residual up to $7\cdot10^{-3}$ ($m=1$) to $7\cdot10^{-5}$ ($m=15$) would pass, eight to twelve orders of magnitude above the polishing residual. If the inequalities held in exact arithmetic for compactly supported tests, Theorem~\CGref{thm:rb-certificate} would give an exact solution of the square section within $\mathrm r$ of $x_0$, inside the tube with $C=1$; what has been established is the floating-point inequality, without rounding bounds and with Gaussian tests, so these stages carry candidate certificates and not verified ones. For compactly supported tests this is settled in Section~\CGref{subsec:rb-num-verified}: with B-spline tests the same fifteen stages pass in ball arithmetic. The candidate solutions are not unique: at $m=15$ the polished point is $4\cdot10^{-2}$ from $\boldsymbol\gamma$. At $m=16$ the square Jacobian is singular to working precision and the stage has no candidate certificate. Least-squares log-log fits over the computed stages $m\ge4$ (Figure~\CGref{fig:rb-scaling}) give the empirical exponents
\begin{equation}\label{eq:rb-num-scaling}
 \mathbb L_m\sim T^{1.81\pm0.10},\qquad\|\mathbb Q\mathbb J'\|_{\infty\to\infty}\sim T^{2.10\pm0.06},\qquad\text{sampled margin}\sim T^{-2.08\pm0.06}\qquad(T\le487);
\end{equation}
these are regressions over the $12$ computed stages $4\le m\le15$, with one standard error of the fitted slope after $\pm$, not asymptotic laws.

\subsubsection{Verified certificates with compactly supported tests}\label{subsec:rb-num-verified}

\emph{Tests.} Let $\mathrm N_8$ be the cardinal B-spline of order $8$ on $[0,8]$, the eightfold convolution of $\mathbf 1_{[0,1]}$: a piecewise polynomial of degree $7$ and class $C^6$ with $\int\mathrm N_8=1$. For $c,h>0$ put $\varphi^{\mathrm N}_{c,h}(u):=\mathrm N_8\bigl((u-c)/h+4\bigr)$, supported in $[c-4h,c+4h]$. Then
\begin{equation}\label{eq:rb-bspline-transform}
 \mathsf c_{\varphi^{\mathrm N}_{c,h}}(t)=h\cos(ct)\Bigl(\frac{\sin(ht/2)}{ht/2}\Bigr)^8,\qquad|\mathsf c_{\varphi^{\mathrm N}_{c,h}}(t)|\le h\Bigl(\frac2{ht}\Bigr)^8 .
\end{equation}
Since $\mathrm N_8^{(i)}=\nabla^i\mathrm N_{8-i}$ with $\nabla g(x)=g(x)-g(x-1)$ and $\|\mathrm N_q\|_1=1$, $\|(\varphi^{\mathrm N}_{c,h})^{(i)}\|_1\le2^ih^{1-i}$ for $i\le7$, and the constant \CGeqref{eq:rb-Wr} satisfies, for $r\le7$ and $u^*=c+4h$,
\begin{equation}\label{eq:rb-W-bspline}
 \mathfrak W_r(\varphi^{\mathrm N}_{c,h})\le(1+u^{*2})\cosh\tfrac{u^*}2\;h\Bigl(1+\frac2h\Bigr)^r .
\end{equation}
The residual $\mathfrak R(\boldsymbol\kappa,\mathsf a)$ is a distribution of order at most $2$ on $(0,\infty)$, because $\sum_k\cos(\kappa_ku)$ is the second derivative of the uniformly convergent series $-\sum_k\kappa_k^{-2}\cos(\kappa_ku)$, the masses form a measure and $\mathsf k$ is smooth. So \CGeqref{eq:rb-residual} defines $\langle\mathfrak R,\varphi\rangle$ for every $\varphi\in C_c^2(0,\infty)$, in particular for $\varphi^{\mathrm N}_{c,h}$, and Proposition~\CGref{prop:rb-compactness}, Theorems~\CGref{thm:rb-certificate} and~\CGref{thm:rb-staged} and Corollaries~\CGref{cor:rb-equivalence} and~\CGref{cor:rb-Bprime} hold verbatim for test sets in $C_c^2(0,\infty)$ whose union has a span that approximates every $\varphi\in C_c^\infty(0,\infty)$ in $C^2$ with supports in a fixed compact set: the proofs use the tests only through $|\mathsf c_\varphi(t)|\le\|\varphi''\|_1t^{-2}$, the finiteness of the mass sum, and density. B-splines with dyadic widths provide such sets: for $f\in C_c^\infty(0,\infty)$ and $h=2^{-i}$ the splines $\sum_jf(jh)\,\varphi^{\mathrm N}_{jh,h}$ converge to $f$ in $C^2$, because $\frac d{du}\sum_j\mathsf c_j\mathrm N_q(u/h-j)=h^{-1}\sum_j(\mathsf c_j-\mathsf c_{j-1})\mathrm N_{q-1}(u/h-j)$, $\sum_j\mathrm N_q(x-j)=1$, and the divided differences of $f$ converge uniformly to its derivatives. The computations below use other widths and centres; they concern finite square sections, for which density plays no role.

\emph{Sections.} Stage $m$, $1\le m\le16$, uses $\mathrm w_m=0.08\cdot0.85^{m-1}$ and $T_m=4/\mathrm w_m$ as in Section~\CGref{subsec:rb-num-setup}, the width $h_m:=\mathrm w_m\sqrt{12/8}$ (so that $\varphi^{\mathrm N}_{c,h_m}$ has the variance of the Gaussian of width $\mathrm w_m$), and the tests $\varphi^{\mathrm N}_{c,h_j}$ with centres $c\in[9\mathrm w_j,\log(\max(T_j,7)/2\pi)+2]$ in steps of $\mathrm w_j$, for all levels $j\le m$; for $m\le15$ their supports lie in $[4.1\,\mathrm w_j,\,6.4]$. The unknowns are the $n_m$ ordinates with $\tau_k\le2.4/\mathrm w_m$. The other ordinates are frozen: at the receding-horizon iterate of Section~\CGref{subsec:rb-num-init} for $n_m<k\le60\,000$ (double-precision numbers, taken as exact; they coincide with the clock to $10^{-12}$ beyond $k=1000$), and at the clock for $k>60\,000$. The masses are $\mathsf a^\zeta$. The square system $\mathbb F_m$ consists of $n_m$ tests chosen by pivoted QR on $\mathbb J^{\!T}$ at the iterate.

\emph{Computation.} Program \texttt{certV} (Appendix~\CGref{app:code-routeB-certV}) polishes $\mathbb F_m$ by Newton's method in double precision and then works in ball arithmetic (Arb \cite{Johansson2017} through python-flint, $128$ bits). It computes enclosures of: the cosine transforms by \CGeqref{eq:rb-bspline-transform}; the frozen ordinates $k\le60\,000$, summed explicitly; the clock tail $k>60\,000$, bounded by
\begin{equation}\label{eq:rb-clock-tail}
 \sum_{k>K}|\mathsf c_{\varphi^{\mathrm N}_{c,h}}(\tau_k)|\le h\Bigl(\pi h\Bigl(1-\frac{11}{8K}\Bigr)\Bigr)^{-8}\,\frac{(\log K)^8K^{-7}}{7\bigl(1-\frac8{7\log K}\bigr)},
\end{equation}
which follows from \CGeqref{eq:rb-bspline-transform}, from $\tau_k\ge2\pi(1-\frac{11}{8K})k/\log k$ for $k>K$ ($W_0(x)\le\log x$ for $x\ge e$) and from the integral test; the masses, exactly; and $\langle\mathsf k,\varphi^{\mathrm N}_{c,h}\rangle$, by rigorous numerical integration over the eight polynomial pieces of $\varphi^{\mathrm N}_{c,h}$ \cite{Johansson2018}. Two Newton steps in ball arithmetic give the centre $x_1$. With $\mathbb Q$ the double-precision inverse of the Jacobian at the polished point, taken as an exact matrix, the map $\mathbb T(x)=x-\mathbb Q\mathbb F_m(x)$ maps the ball $B(x_1,\rho)$ of $\ell^\infty$ into itself and contracts if
\begin{equation}\label{eq:rb-krawczyk}
 \mathsf l(\rho):=\sup_{x\in B(x_1,\rho)}\bigl\|I-\mathbb Q\,D\mathbb F_m(x)\bigr\|_{\infty\to\infty}<1,\qquad\|\mathbb Q\mathbb F_m(x_1)\|_\infty\le(1-\mathsf l(\rho))\,\rho ;
\end{equation}
an upper bound for $\mathsf l(\rho)$ is obtained by evaluating the Jacobian on the ball, entry by entry, and multiplying interval matrices. As in Section~\CGref{subsec:rb-num-certified}, $\mathsf l<1$ makes $\mathbb Q$ invertible, so the fixed point is the only zero of $\mathbb F_m$ in the ball. The program also checks, in ball arithmetic, that the frozen list is increasing and lies within $0.9938$ of the clock (with $\tau_k$ enclosed through Arb's Lambert function), it bounds the distance of the solution to the zeta ordinates using Arb's rigorous enclosures of the zeros (\texttt{acb.zeta\_zero}), and it evaluates the bound \CGeqref{eq:rb-nonfinitary} with $r=7$ and \CGeqref{eq:rb-W-bspline} on every test of the stage.

\begin{theorem}[Exact square sections up to height $487$; computer-assisted]\label{thm:rb-verified}
For each $m=1,\dots,15$ the square system $\mathbb F_m$ described above has exactly one zero $x_*^{(m)}$ in the ball $\|x-x_1^{(m)}\|_\infty\le\rho_m$; the centres $x_1^{(m)}$ (exact binary numbers) and the radii $\rho_m$ (double-precision numbers) are recorded in the file \texttt{certV\_centres.json} of the repository, and Table~\CGref{tab:rb-verified} gives $\rho_m$ rounded up. The list formed by $x_*^{(m)}$ and the frozen ordinates is strictly increasing and $0.995$-admissible, and with the von Mangoldt masses it satisfies exactly the $n_m$ selected equations of the rung system. Its first $n_m$ ordinates are within the distance given in Table~\CGref{tab:rb-verified} of $\gamma_1,\dots,\gamma_{n_m}$.
\end{theorem}

\begin{proof}
For every stage the ball-arithmetic evaluation of program \texttt{certV} gives enclosures that satisfy \CGeqref{eq:rb-krawczyk} with the values of Table~\CGref{tab:rb-verified}, and Banach's fixed point theorem gives the zero and its uniqueness in the ball. The gaps between consecutive ordinates of $x_1^{(m)}$, and between the last of them and the first frozen ordinate, exceed $2\rho_m$ (the smallest gap is 0.71); $\max_k|x^{(m)}_{1,k}-\tau_k|+\rho_m\le0.9937$, and the frozen ordinates are within $0.9938$ of the clock. The distances to $\gamma_k$ are bounded with the enclosures of the zeros computed by Arb.
\end{proof}

\begingroup\small
\setlength{\tabcolsep}{3.5pt}
\begin{longtable}{rrrrrrrrr}
\caption{Verified square sections with B-spline tests (program \texttt{certV}). $\rho_m$: radius of the ball of Theorem~\CGref{thm:rb-verified}; upper bounds are rounded up and lower bounds down; $\mathsf l(\rho_m)$: upper bound of the contraction constant \CGeqref{eq:rb-krawczyk}; $\mathbb L_m=\|\mathbb Q\|_{\infty\to\infty}$; margin: largest $(1-\mathsf l(\rho))\rho$ over $36$ radii in $[10^{-8},10^{-1}]$, each value a rigorous lower bound for the residual that \CGeqref{eq:rb-krawczyk} would accept; $\max_{k\le n}|x_*-\gamma|$: rigorous bound of the distance of the solution to $\gamma_1,\dots,\gamma_{n_m}$; $\epsilon_\zeta$: the bound of Proposition~\CGref{prop:rb-nonfinitary} ($r=7$) for the residual of the zeta list on every test of level $\le m$.}\label{tab:rb-verified}\\
\toprule
$m$ & $T_m$ & $n_m$ & $\rho_m$ & $\mathsf l(\rho_m)$ & $\mathbb L_m$ & margin & $\max_{k\le n}|x_*-\gamma|$ & $\epsilon_\zeta$\\
\midrule
\endfirsthead
\toprule
$m$ & $T_m$ & $n_m$ & $\rho_m$ & $\mathsf l(\rho_m)$ & $\mathbb L_m$ & margin & $\max_{k\le n}|x_*-\gamma|$ & $\epsilon_\zeta$\\
\midrule
\endhead
1 & 50 & 4 & $6.5\cdot10^{-21}$ & $1.3\cdot10^{-11}$ & $1.1\cdot10^{2}$ & $2.6\cdot10^{-3}$ & $7.7\cdot10^{-7}$ & $2.2\cdot10^{-64}$\\
2 & 59 & 5 & $1.3\cdot10^{-20}$ & $2.4\cdot10^{-12}$ & $6.9\cdot10^{1}$ & $4.3\cdot10^{-3}$ & $9.0\cdot10^{-8}$ & $6.3\cdot10^{-64}$\\
3 & 69 & 7 & $6.1\cdot10^{-20}$ & $9.7\cdot10^{-12}$ & $1.2\cdot10^{2}$ & $1.2\cdot10^{-3}$ & $1.1\cdot10^{-6}$ & $1.8\cdot10^{-63}$\\
4 & 81 & 9 & $2.9\cdot10^{-19}$ & $1.3\cdot10^{-11}$ & $1.7\cdot10^{2}$ & $8.5\cdot10^{-4}$ & $1.2\cdot10^{-6}$ & $5.0\cdot10^{-63}$\\
5 & 96 & 12 & $9.8\cdot10^{-19}$ & $1.5\cdot10^{-11}$ & $1.7\cdot10^{2}$ & $5.2\cdot10^{-4}$ & $4.8\cdot10^{-7}$ & $1.5\cdot10^{-62}$\\
6 & 113 & 16 & $3.9\cdot10^{-18}$ & $2.5\cdot10^{-11}$ & $2.2\cdot10^{2}$ & $3.8\cdot10^{-4}$ & $2.6\cdot10^{-6}$ & $4.2\cdot10^{-62}$\\
7 & 133 & 20 & $2.5\cdot10^{-17}$ & $4.9\cdot10^{-11}$ & $4.6\cdot10^{2}$ & $1.3\cdot10^{-4}$ & $1.8\cdot10^{-5}$ & $1.3\cdot10^{-61}$\\
8 & 156 & 26 & $6.0\cdot10^{-17}$ & $3.0\cdot10^{-11}$ & $4.2\cdot10^{2}$ & $1.7\cdot10^{-4}$ & $3.9\cdot10^{-6}$ & $3.7\cdot10^{-61}$\\
9 & 183 & 34 & $3.1\cdot10^{-16}$ & $1.7\cdot10^{-10}$ & $6.6\cdot10^{2}$ & $8.3\cdot10^{-5}$ & $1.7\cdot10^{-5}$ & $1.1\cdot10^{-60}$\\
10 & 216 & 43 & $1.3\cdot10^{-15}$ & $1.3\cdot10^{-10}$ & $9.1\cdot10^{2}$ & $5.1\cdot10^{-5}$ & $2.1\cdot10^{-5}$ & $3.2\cdot10^{-60}$\\
11 & 254 & 54 & $5.5\cdot10^{-15}$ & $1.3\cdot10^{-10}$ & $1.2\cdot10^{3}$ & $3.6\cdot10^{-5}$ & $3.2\cdot10^{-5}$ & $9.4\cdot10^{-60}$\\
12 & 299 & 68 & $1.7\cdot10^{-14}$ & $3.4\cdot10^{-10}$ & $1.4\cdot10^{3}$ & $2.7\cdot10^{-5}$ & $5.3\cdot10^{-5}$ & $2.8\cdot10^{-59}$\\
13 & 352 & 85 & $9.4\cdot10^{-14}$ & $1.8\cdot10^{-9}$ & $2.1\cdot10^{3}$ & $1.5\cdot10^{-5}$ & $3.1\cdot10^{-4}$ & $8.4\cdot10^{-59}$\\
14 & 414 & 107 & $3.3\cdot10^{-13}$ & $7.6\cdot10^{-9}$ & $2.5\cdot10^{3}$ & $1.1\cdot10^{-5}$ & $3.9\cdot10^{-3}$ & $2.5\cdot10^{-58}$\\
15 & 487 & 133 & $1.1\cdot10^{-12}$ & $3.0\cdot10^{-8}$ & $2.8\cdot10^{3}$ & $8.9\cdot10^{-6}$ & $3.7\cdot10^{-2}$ & $7.5\cdot10^{-58}$\\
16 & 572 & 165 & \multicolumn{6}{l}{Newton polishing fails: the square Jacobian is singular to working precision}\\
\bottomrule
\end{longtable}
\endgroup

\begin{remark}[What Theorem~\CGref{thm:rb-verified} says and does not say]\label{rem:rb-verified}
For these fifteen stages the candidate certificates of Section~\CGref{subsec:rb-num-certified} are replaced by verified ones: the inequalities hold in ball arithmetic, for tests in $C_c^6(0,\infty)$, with every evaluation error included. The radii $\rho_m$, from $6.5\cdot10^{-21}$ to $1.1\cdot10^{-12}$, are set by the bound \CGeqref{eq:rb-clock-tail} for the clock tail, which more frozen ordinates would reduce; $\mathsf l(\rho_m)\le3.0\cdot10^{-8}$. The verified margins are smaller than the sampled margins of the Gaussian candidates, because the B-spline transform changes sign and decays only like $t^{-8}$. At $m=16$, $T=572$, the square Jacobian is singular to working precision, as with Gaussian tests. Three things are not given. The selected equations are $n_m$ of the tests of level $\le m$; on all of them the floating-point residual of the solution is between $2.4\cdot10^{-8}$ and $7.2\cdot10^{-4}$. The solutions of different stages are not restrictions of one another. And by Proposition~\CGref{prop:rb-nonfinitary} no statement of this kind bears on the Riemann hypothesis: whether or not RH holds, the zeta list with the von Mangoldt masses has residual at most $7.5\cdot10^{-58}$ on every test used here (column $\epsilon_\zeta$). The distances to $\gamma_k$ (at most $7.6\cdot10^{-4}$ for $k\le20$ at every stage) are diagnostics.
\end{remark}

\subsubsection{The entry step}\label{subsec:rb-num-entry}

The open point of the staged construction, in the form that Theorem~\CGref{thm:rb-staged} makes precise, is whether the residual $\mathsf h_{m+1}$ of the new tests at the start of stage $m+1$ is below the margin of that stage. This was tested as follows: start from the polished point of stage $m-1$ (Section~\CGref{subsec:rb-num-certified}), put the new band of unknowns (those with $\tau_k\in(2.4/\mathrm w_{m-1},2.4/\mathrm w_m]$) and everything beyond at their Gonek values, and evaluate \CGeqref{eq:rb-num-banach} (i) for the band alone, with the new tests of level $m$, and (ii) for the full square system of stage $m$, without polishing.

The band-only test passes (in the same floating-point sense) for $T\le96$ and fails from $T=113$ on; its margin falls from $9\cdot10^{-2}$ to $3\cdot10^{-3}$ while the Gonek residual of the band stays between $10^{-2}$ and $3\cdot10^{-1}$ (the Gonek error in the band is $0.02$ to $0.09$). The full-system test passes only at $T=59$. The margin decreases as a power of $T$ by \CGeqref{eq:rb-num-scaling}, whereas the error of the Gonek lists with fixed $\mathrm X$ does not decrease with $T$, and with $\mathrm X=T^{\mathrm x}$ it would decrease at most like $1/\log T$. Within the range computed, no worst-case certificate in $\ell^\infty$ closes the entry step, and the trend of the constants says that none will at larger $T$. This is a numerical statement, not a theorem.

\emph{The same test in $\ell^2$.} Program \texttt{entry2} (Appendix~\CGref{app:code-routeB-entry2}) evaluates the entry test with the norms $\|\cdot\|_{2\to2}$ and $\|\cdot\|_2$ in place of $\|\cdot\|_{\infty\to\infty}$ and $\|\cdot\|_\infty$ in \CGeqref{eq:rb-num-banach}. This is legitimate: the second-order term is $\mathbb Q\mathbb J'\operatorname{diag}(\mathrm e)\mathrm e$, and $\|\operatorname{diag}(\mathrm e)\|_{2\to2}=\|\mathrm e\|_\infty\le\|\mathrm e\|_2$; a zero in the $\ell^2$ ball lies in the $\ell^\infty$ ball of the same radius. In $\ell^2$ the inverse is smaller ($\|\mathbb Q\|_{2\to2}$ is $1.0$ to $2.3$ times below $\|\mathbb Q\|_{\infty\to\infty}$ for the band and $1.7$ to $3.9$ times for the full system), the $\ell^2$ norm of the preconditioned residual of the new tests is up to $1.9$ times larger than its $\ell^\infty$ norm, and the margins grow by factors $1.0$ to $2.0$ (band) and $1.2$ to $2.8$ (full system). The pattern does not change: the band test passes for $T\le96$ and fails from $T=113$ on, and the full-system test passes only at $T=59$, in both norms. The moderate $\ell^2$ conditioning at the zeta list (Section~\CGref{subsec:rb-num-init}) does not close the entry step; what fails is the size of the error of the Gonek lists relative to the margin, not the choice of norm.

\begin{figure}[tbp]
\centering
\includegraphics[width=0.9\textwidth]{fig19_cert_scaling.png}
\caption{Constants of the Banach certificate \CGeqref{eq:rb-num-banach} for the square sections of Section~\CGref{subsec:rb-num-certified} (solid lines, with fitted slopes in $\log T$), and the band-only entry test of Section~\CGref{subsec:rb-num-entry} (dashed): the reciprocal of the band margin and the reciprocal of the residual of the new band at its Gonek values. The entry test passes, in floating point, where the triangles lie below the squares.}
\label{fig:rb-scaling}
\end{figure}

\subsubsection{Reading of the computations}\label{subsec:rb-num-reading}

\begin{enumerate}[label=(\roman*)]
\item The constructions that use only the clock and the primes (the Gonek lists, the square sections, the certificates) behave as predicted by the tube lemma and the certificate theorem: the candidate certificates hold in floating point up to $T=487$, and fifteen square sections have verified exact solutions (Theorem~\CGref{thm:rb-verified}). By Proposition~\CGref{prop:rb-nonfinitary} none of this bears on the Riemann hypothesis: if it were false, the zeta list with the von Mangoldt masses would still solve every test used here up to $10^{-57}$, and sections in the box with residuals of that size would exist at every stage. A zero $\frac12+d_0+i\gamma_0$ off the line contributes $2\cosh(d_0u)\cos(\gamma_0u)$ to the prime side; a stage detects it only through tests whose constant $\mathfrak W_r$ is of order $H_0^{r-1}$, that is, tests that resolve frequencies of the order of $H_0$ (Remark~\CGref{rem:rb-nonfinitary}); Heuristic~\CGref{heur:rb-window-zerofree} describes how the size of the window enters.
\item The quantities $\mathbb L_m$ and $\|QB\|$ grow polynomially in $T$. This is the quantitative face of the obstruction theorem of Section~\CGref{subsec:rb-obstruction}: no height-uniform linear stability, hence no height-uniform Newton basin. A certificate for the infinite construction needs a basin that shrinks no faster than the initial error of a $\gamma$-free list decreases, and the computed trend is $T^{-2}$ against $1/\log T$.
\item The observed contraction of Gauss--Newton (factor $4$ to $15$ per step at every height) is a mean-square phenomenon that the worst-case constants do not see. A proof of the entry step would have to bound $\|\mathbb Q\mathbb J'\,\mathrm{diag}(\mathrm e)\,\mathrm e\|_\infty$ for the actual error vector $\mathrm e$ of the Gonek list, not for an arbitrary vector of the same norm; that error vector is $\gamma$-free, and its residual $\mathbb F_m(\boldsymbol\kappa^{(\mathrm X)})$ is controlled through the explicit formula by the moments of $S-S_{\mathrm X}$, which are unconditional (Selberg, Fujii). Whether this closes is not known; any proof along this line would contain a proof of the Riemann hypothesis, by Corollary~\CGref{cor:rb-Bprime}.
\end{enumerate}

%======================================================================
\subsection{Status of the statements}\label{subsec:rb-status}

\begingroup\small
\begin{longtable}{p{0.42\textwidth}p{0.5\textwidth}}
\toprule
Statement & Status\\
\midrule
Lemma~\CGref{lem:rb-canonical-box}, Corollary~\CGref{cor:rb-D-implies-B} & proved; conditional on (D) as stated\\
Proposition~\CGref{prop:rb-compactness}, Lemma~\CGref{lem:rb-matching} & proved, unconditional\\
Lemma~\CGref{lem:rb-mass-ms}, Theorem~\CGref{thm:rb-uniqueness}, Corollaries~\CGref{cor:rb-U}, \CGref{cor:rb-selection}, \CGref{cor:rb-equivalence} & proved, unconditional: an exact solution in a tube, with arbitrary real masses, is $(\boldsymbol\gamma,\mathsf a^\zeta)$ and implies RH; B$(C,\mathbb M)\Leftrightarrow$(D); uses Proposition~\CGref{prop:rung-apriori}, Lemma~\CGref{lem:mean-value-ordinates} and Hamburger's theorem\\
Proposition~\CGref{prop:rb-nonfinitary} & proved, unconditional; uses \cite{PlattTrudgian2021,HasanalizadeShenWong2022}; no finite stage decides B or B$'$\\
Lemmas~\CGref{lem:rb-archimedean}, \CGref{lem:rb-laplace} & proved, unconditional\\
Proposition~\CGref{prop:rb-prescribed-uniqueness}, Corollary~\CGref{cor:rb-Bprime} & proved, unconditional: RH $\Leftrightarrow$ B$'(C)$; classical in substance\\
Lemma~\CGref{lem:rb-lipschitz}, Theorems~\CGref{thm:rb-certificate}, \CGref{thm:rb-staged}, Proposition~\CGref{prop:rb-graded} & proved; their hypotheses (A$_m$) for all $m$ and (B) are open and imply RH\\
Heuristic~\CGref{heur:rb-frame} & not proved; its point model is not a test of $C_c^\infty(0,\infty)$\\
Lemma~\CGref{lem:rb-gonek-tube} & proved, unconditional; uses \cite{Gonek2012} only for the interpretation of the list\\
(G1), (G2) of Remark~\CGref{rem:rb-gonek-status} & open; together they imply B$'$ and RH\\
Lemmas~\CGref{lem:rb-PL}--\CGref{lem:rb-argument}, Theorem~\CGref{thm:rb-count}, Corollary~\CGref{cor:rb-displacement} & proved, unconditional; uses \CGref{thm:halfplane-RVM}; by Lemma~\CGref{lem:rb-upsilon-infinite} the hypothesis $\upsilon<\infty$ holds for a finite section only when all masses vanish, so the bound is a template, not a bound for sections with primes (Remark~\CGref{rem:rb-count-scope})\\
Lemma~\CGref{lem:rb-peel}, Theorems~\CGref{thm:rb-window-unknown}, \CGref{thm:rb-bernstein} & proved, unconditional: RH $\Leftrightarrow$ an unknown list solves the window identity $\Leftrightarrow$ $e^{-iy/2}V_+(iy)$ is completely monotone; uses \cite{Bernstein1929}; literature not searched for Theorem~\CGref{thm:rb-bernstein}\\
Lemma~\CGref{lem:rb-minimizer}, Theorem~\CGref{thm:rb-window-equivalence} & proved, unconditional: dichotomy for minimizers; RH $\Leftrightarrow$ solvability on every window in a fixed box; the exclusion of degenerate minimizers is open\\
Heuristics~\CGref{heur:rb-window}, \CGref{heur:rb-informative}, \CGref{heur:rb-graduated}, \CGref{heur:rb-window-zerofree} & not proved\\
Lemma~\CGref{lem:rb-generating}, Theorem~\CGref{thm:rb-criterion}, Corollary~\CGref{cor:rb-riesz} & proved; Corollary~\CGref{cor:rb-riesz} uses \cite{AvdoninIvanov2002}\\
Corollary~\CGref{cor:rb-kadec-avdonin}, Theorem~\CGref{thm:rb-variable-N}, Proposition~\CGref{prop:rb-variable-blocks} & proved; the variable-length version was not found in the literature consulted (Remark~\CGref{rem:rb-price})\\
Lemma~\CGref{lem:rb-psi-bound}, Proposition~\CGref{prop:rb-O-blocks}, Proposition~\CGref{prop:rb-weight}, Corollary~\CGref{cor:rb-one-condition} & proved, unconditional\\
Lemma~\CGref{lem:rb-smoothing}, Theorem~\CGref{thm:rb-obstruction}, Corollary~\CGref{cor:rb-obstruction} & proved, unconditional; uses the upper bound contained in \cite[Main Theorem]{Fujii1974}, \cite{Selberg1946,Titchmarsh1986,Garnett2007}; no study of $(A_2)$ for the zeta ordinates was found (Remark~\CGref{rem:rb-obstruction-literature})\\
Theorem~\CGref{thm:rb-verified} & proved with computer assistance (ball arithmetic, B-spline tests): exact solutions of fifteen square sections up to $T=487$; no bearing on RH (Proposition~\CGref{prop:rb-nonfinitary})\\
Other computations of Section~\CGref{subsec:rb-numerics} & floating point; Gaussian candidate certificates and the entry step in $\ell^\infty$ and $\ell^2$; not part of any proof\\
\bottomrule
\end{longtable}
\endgroup

%======================================================================

\section{Numerical illustration with \texorpdfstring{$2\,001\,052$}{2001052} zeros}
\label{sec:numerics}

This section shows what is being constructed and checks that the theorems are consistent with computation.  It uses the ordinates $\gamma_1,\dots,\gamma_N$ of the first $N=2\,001\,052$ zeros of $\zeta$ (file \texttt{zeross.bin}, $\gamma_N=1\,132\,490.6587$) and the midpoint ordinates $\tau_1,\dots,\tau_N$ computed from \CGeqref{eq:lambert} (file \texttt{zLamb.bin}, $\tau_N=1\,132\,490.7226$).  The programs are listed in Appendix~\CGref{app:code} and their complete console output in Appendix~\CGref{app:output}; all figures have linear axes.  Each item is an \emph{experiment}: it reports a computation and does not replace a proof.

Products over the ordinates are evaluated as follows, in double precision throughout: the first $10^5$ factors one by one, the factors $10^5<k\le N$ through power sums of $(\tfrac14+o_k^2)^{-1}$, and the factors $k>N$ by the midpoint rule for the clock, which gives the tail $\frac1{2\pi}\bigl(q\frac{L_U+1}U-\frac{q^2}2\frac{3L_U+1}{9U^3}+\frac{q^3}3\frac{5L_U+1}{25U^5}\bigr)$ with $\NR(U)=N$, $L_U=\log(U/2\pi)$.  For the products over the $\gamma_k$ ($P_{Z1}$, $\zeta_{Z1}$) this means that the unknown tail of the zeta ordinates beyond the file is \emph{replaced by the clock tail}; every agreement reported for $\zeta_{Z1}$ against $\zeta$ includes that substitution and is an agreement of the construction with itself beyond height $\gamma_N$.  ``All $N$ ordinates'' or ``all zeros'' below always means all zeros in the supplied file, $\gamma_k\le\gamma_N$.  No quantity in this section is certified: the words ``exact'' and ``certified'' are not used for floating-point values, and the label \texttt{FLOAT\_PASS} of the console outputs (Appendix~\CGref{app:out-routeB}) marks a floating-point test, as explained in Remark~\CGref{rem:rb-certificate-errors}.

\subsection{The clock and the zeros}
\label{subsec:num-clock}

\begin{experiment}[Clock, midpoints and displacements]\label{exp:clock}
\emph{Object:} $\NR$, $N_\zeta$, $\tau_k$, $\eta_k=\tau_k-\gamma_k$.  \emph{Range:} $k\le N$.  \emph{Expected:} \CGeqref{eq:midpoint-condition} to rounding, and $|\eta_k|\le C_\eta$ (Lemma~\CGref{lem:bounded-displacement}).  \emph{Observed:} $\max_k|\NR(\tau_k)-(k-\tfrac12)|=6.985\cdot10^{-10}$; $\max_k|\eta_k|=0.9937$ (at $k=34$), mean $|\eta_k|=0.1512$, root mean square $0.1881$, mean $\eta_k=1.7\cdot10^{-7}$.  By blocks of $k$, $\max|\eta_k|$ is $0.840,0.994,0.915,0.834,0.786,0.763,0.717$ on $[1,10]$, $[11,10^2]$, $[10^2,10^3]$, \dots, $[10^6,N]$.
\end{experiment}

Figure~\CGref{fig:clock-zoom} shows how the clock places the midpoint ordinates.  The staircase $N_\zeta$ jumps by one at each $\gamma_k$; the smooth clock $\NR$ crosses the half-integer level $k-\tfrac12$ at $\tau_k$.  Since $N_\zeta-\NR=S(T)+O(1/T)$ (Figure~\CGref{fig:remainder}), the horizontal distance between the open square $(\gamma_k,k-\tfrac12)$ and the red point $(\tau_k,k-\tfrac12)$ is the remainder $S$ converted into ordinate units by the slope $\NR'=\frac1{2\pi}\log\frac T{2\pi}$.  The right panel repeats the picture at height $1.13\cdot10^6$, where the mean spacing is $0.519$: every $\tau_k$ still lies next to its $\gamma_k$.

\begin{figure}[htbp]
\centering
\includegraphics[width=0.86\textwidth]{fig00_clock_zoom.png}
\caption{The midpoint clock $\NR$ (red), the zero staircase $N_\zeta$ (black), the midpoint ordinates $\tau_k$ (red points) and the points $(\gamma_k,k-\frac12)$ (squares).  Left: the first nine zeros.  Right: the last twelve zeros of the file.}
\label{fig:clock-zoom}
\end{figure}

\begin{figure}[htbp]
\centering
\includegraphics[width=0.86\textwidth]{fig00b_remainder.png}
\caption{The computed difference $N_{\rm file}(T)-\NR(T)$ for $T\le200$, where $N_{\rm file}$ counts the supplied ordinates.  The file is complete on this range, so the curve is the remainder $N_\zeta(T)-\NR(T)=S(T)+O(1/T)$ of \CGeqref{eq:RvM-full}.}
\label{fig:remainder}
\end{figure}

\subsection{The \texorpdfstring{$\lambda$}{lambda}-family and the endpoint}
\label{subsec:num-family}

Figure~\CGref{fig:family} shows the core of the method.  The lower left panel draws the moving ordinates $\gamma_k(\lambda)$ in clock units, $\NR(\gamma_k(\lambda))-(k-\frac12)$: at $\lambda=1$ all lines meet at $0$ (the $\tau_k$), at $\lambda=0$ they sit at the ordinates of $\zeta$; for $\lambda>0$ the ordinates $\gamma_k(\lambda)$ remain strictly increasing in $k$ (Theorem~\CGref{thm:path}).  The lower right panel draws the real functions $Z_\lambda(t)=e^{i\vartheta(t)}\zeta_\lambda(\frac12+it)$ for $\lambda=0,0.25,0.5,0.75,1$ together with Hardy's function $Z(t)=e^{i\vartheta(t)}\zeta(\frac12+it)$: the zeros of every $Z_\lambda$ are real, they move continuously with $\lambda$, and the curve for $\lambda=0$ lies on top of $Z$ at the scale of the figure, because it is computed from the supplied ordinates with the clock tail substituted beyond them; this is a property of the computed curves, not evidence for $\zeta_{Z1}=\zeta$.

\begin{figure}[htbp]
\centering
\includegraphics[width=0.80\textwidth]{fig01_clock_family.png}
\caption{Top: the observed displacements $\eta_k$ for every 50th $k\le N$ of the supplied range and their empirical distribution.  Bottom: moving ordinates in clock units, $k\le30$, and the functions $Z_\lambda(t)$ on the critical line.}
\label{fig:family}
\end{figure}

\begin{experiment}[The endpoint $\zeta_{Z1}$ against $\zeta$]\label{exp:Z1}
\emph{Object:} $\zeta_{Z1}=P_{Z1}/G$ computed from the $N$ ordinates $\gamma_k$, and $\zeta$ computed by mpmath.  \emph{Range:} $\sigma\in\{2,1.25,0.75\}$, $1\le t\le2000$ ($499$ points per line), and the critical line $0.5\le t\le60$.  \emph{Expected:} if $\zeta_{Z1}=\zeta$, the ratio is $1$ up to the truncation of the product.  \emph{Observed:} $\max|\zeta_{Z1}/\zeta-1|=8.7\cdot10^{-10}$, $1.1\cdot10^{-9}$, $2.6\cdot10^{-9}$ on the three lines, and $\max|Z_{Z1}-Z|/\max(1,|Z|)=8.4\cdot10^{-10}$ on the critical line; the balanced error $|\zeta_{Z1}-\zeta|/\max(1,|\zeta_{Z1}|,|\zeta|)$, which is not amplified near the zeros, is at most $7.5\cdot10^{-10}$, $1.1\cdot10^{-9}$, $1.9\cdot10^{-9}$ on the three lines, and the imaginary parts discarded in $Z_\lambda$ are below $3\cdot10^{-14}$ (Figure~\CGref{fig:Z1}); the tail beyond $\gamma_N$ is the clock tail, so the $10^{-9}$ figure is consistent with, and is expected to be dominated by, the file truncation and the clock-tail substitution; it does not test (D).
\end{experiment}

\begin{figure}[htbp]
\centering
\includegraphics[width=0.86\textwidth]{fig02_zetaZ1_vs_zeta.png}
\caption{Left: $|\zeta_{Z1}/\zeta-1|$ on three vertical lines, with $\zeta_{Z1}$ computed from the supplied ordinates and the clock tail beyond them.  Right: $Z_{Z1}$ and Hardy's $Z$ on the critical line.}
\label{fig:Z1}
\end{figure}

Experiment~\CGref{exp:Z1} shows numerical agreement at the displayed truncation precision.  This does not test (D): Proposition~\CGref{prop:varpi-verified} already forces $|\log\varpi_{Z1}|<4\cdot10^{-13}$ from verified zeros, while the observed $10^{-9}$ scale is dominated by truncating the ordered clock tail.

\begin{experiment}[The value at infinity: the product identity]\label{exp:product-identity}
\emph{Object:} the ordered partial products $P_K=\prod_{k\le K}(\frac14+\gamma_k^2)/(\frac14+\tau_k^2)$ of \CGeqref{eq:product-identity}.  \emph{Range:} $K\le N$.  \emph{Expected:} by Corollary~\CGref{cor:product-identity}, $\zeta_{Z1}=\zeta$ is the statement $P_K\to\varpi_{\RVM}$.  \emph{Observed:} $P_K=0.96256$, $0.96248$, $0.96229$, $0.961951$, $0.961930$, $0.9619418189$, $0.9619415390$ for $K=10,10^2,\dots,10^6,N$, against $\varpi_{\RVM}=0.9619415789$ from Experiment~\CGref{exp:transition}; for $K\ge10^6$, $P_K$ stays in $[0.9619389,0.9619440]$ with mean $0.9619416633$ (Figure~\CGref{fig:product-identity}).  Proposition~\CGref{prop:varpi-verified} already bounds the endpoint discrepancy by $4\cdot10^{-13}$ under the verified-zero input, so the $10^{-9}$--$10^{-6}$ discrepancies in these finite products are truncation/tail errors and are not numerical evidence for (D).  The flow at infinity \CGeqref{eq:Upsilon}, truncated to the supplied ordinates $k\le N$, is $\Upsilon(\lambda)=0.0436950$, $0.0412445$, $0.0387987$, $0.0363572$, $0.0339198$ for $\lambda=0,\frac14,\frac12,\frac34,1$, decreasing as in \CGeqref{eq:varpi-derivative} with slope close to the auxiliary approximation $-\sum_{k\le N}2\eta_k^2/\gamma_k^2=-0.0098211$ (which drops the $\frac14$ of the denominators and the terms of higher order in $\eta_k/\gamma_k$), and the truncated flow integral is $\int_0^1\Upsilon\,d\lambda=\sum_{k\le N}\log\frac{1/4+\tau_k^2}{1/4+\gamma_k^2}=0.0388016004$, against $\log L=0.0388015590$; the agreement of $P_N$ with $\varpi_{\RVM}$ is an internal check of two computations built on the same clock, not a certification of the infinite product.
\end{experiment}

\begin{figure}[htbp]
\centering
\includegraphics[width=0.86\textwidth]{fig11_product_identity.png}
\caption{The partial products $P_K$ of \CGeqref{eq:product-identity} against $\varpi_{\RVM}$ (dashed).  Left: $K\le10^4$.  Right: $(P_K-\varpi_{\RVM})\cdot10^6$ for $10^5\le K\le N$.  The value of the endpoint at $+\infty$ is $\varpi_{Z1}=\varpi_{\RVM}/\lim_KP_K$ (Theorem~\CGref{thm:Z1-clock}).}
\label{fig:product-identity}
\end{figure}

\subsection{The transition theorem and \texorpdfstring{$\zeta_{\RVM}$}{zeta\_RVM} against \texorpdfstring{$\zeta$}{zeta}}
\label{subsec:num-transition}

\begin{experiment}[Transition theorem]\label{exp:transition}
\emph{Object:} $\zeta_{\RVM}$, $w(s)$, $\varpi_{\RVM}$, $L$.  \emph{Range:} $\sigma'\in\{0.25,0.5,1.5\}$, $10\le t\le10^4$.  \emph{Expected:} Theorem~\CGref{thm:transition} and Corollary~\CGref{cor:vertical-limit}: $t\,|\zeta_{\RVM}/(\varpi_{\RVM}(1+w))-1|$ bounded, $|w|=(t/2\pi)^{-\sigma'}(1+O(1/t))$.  \emph{Observed:} Richardson extrapolation on $\Re s=2$ (which presupposes an expansion in $1/T$ and $1/T^2$ finer than the $O(1/T)$ proved) gives $\varpi_{\RVM}=0.9619415789$ and $L=1.0395641710$, values without error bound; for this extrapolation the product is evaluated with $10^6$ exact factors and eight power sums, because the default evaluation of the programs ($10^5$ exact factors, six power sums) has relative error $7\cdot10^{-10}$ at $T=10^4$ and $1.2\cdot10^{-5}$ at $T=2\cdot10^4$; the heights $4000,7000,10000$ change the extrapolated value by $9\cdot10^{-11}$, doubling all heights by $1.8\cdot10^{-9}$, and the linear model on $5000,10000$ by $1.1\cdot10^{-7}$; for the computed residual, which contains the $O(1/t)$ term of the theorem together with the truncation of the product, the clock tail, the extrapolated $\varpi_{\RVM}$ and rounding, $\max_{t>100}t\,|\zeta_{\RVM}/(\varpi_{\RVM}(1+w))-1|=0.0617,0.0618,0.0620$ for the three values of $\sigma'$, and $0.0790$ for $\varpi_{\RVM}(1+\chi)$ at $\sigma'=0.25$; $|w|\,t^{\sigma'}=1.5832$, $2.5066$, $15.66$--$15.75$, i.e.\ $(2\pi)^{\sigma'}$.  The vertical limit of Corollary~\CGref{cor:vertical-limit} is approached at the rate of $|w|$: at $t=10^4$, $|1/\zeta_{\RVM}-L|=0.143$, $0.025$, $1.6\cdot10^{-5}$ for $\sigma'=0.25,0.5,1.5$, so on $\sigma'=0.25$ the convergence is slow and the limit is not reached numerically in the computed range.
\end{experiment}

\begin{experiment}[The transition constant: the inequalities of Theorem~\CGref{thm:varpi-less-one}]\label{exp:varpi-bound}
\emph{Object:} $\nu_3$, the enclosure of $\tau_1$, the lower bound $\mathfrak M(\tau_1)$ of \CGeqref{eq:nu4-bound}, and, as a cross-check, $\nu_4$, $P^{\crit}_{\RVM}(\frac12)$ and $\varpi_{\RVM}$.  \emph{Method:} Part~A in interval arithmetic (\texttt{mpmath.iv}, $40$ digits, outward rounding); Part~B in floating point ($25$ digits for $\nu_4$, double precision for the product over $k\le2\cdot10^6$).  The program generates the clock from the Lambert formula and reads no data file (it reads \texttt{output/s02\_varpi.txt} only for the comparison in Part~B).  \emph{Observed (Part~A, the proof):} $\nu_3\in[-0.00792401131834915,-0.00792401131834915]$; $\NR(\tau_0)=0.261294\ldots\in(0,\frac12)$; $\NR(14.52134)=0.49999907<\frac12<0.50000041=\NR(14.52135)$; $\mathfrak M(\tau_1)\in[0.00634570,0.00635251]$; hence $\nu_4\le-0.0126914$, $\nu_3+\nu_4\le-0.0206154$ and $\varpi_{\RVM}<0.979596$.  \emph{Observed (Part~B, not part of the proof):} $\nu_4=-0.0251569083$, $e^{\nu_3+\nu_4}=0.9674602699$, $P^{\crit}_{\RVM}(\frac12)=0.9942956924$, $\varpi_{\RVM}=0.9619415789$, which agrees with the Richardson value of Experiment~\CGref{exp:transition} to $2.1\cdot10^{-11}$; the Euler--Maclaurin remainder $\sum_n\mathscr I_n+\mathscr W(1)/12=0.00154$ lies within its bound $0.00469$.  Console output in Appendix~\CGref{app:output-s14}.
\end{experiment}

\begin{figure}[htbp]
\centering
\includegraphics[width=0.80\textwidth]{fig03_transition.png}
\caption{Top: the computed residual of the transition formula multiplied by $t$ (it contains the truncation of the product, the tail approximation, the extrapolated value of $\varpi_{\RVM}$ and rounding), and the rescaled vertical limit $t^{\min(\sigma',1)}|1/\zeta_{\RVM}-L|$.  Bottom: $|\zeta|$ and $|\zeta_{\RVM}|$ on $\sigma=2$ and $\sigma=0.75$.}
\label{fig:transition}
\end{figure}

The bottom panels of Figure~\CGref{fig:transition} compare $|\zeta|$ and $|\zeta_{\RVM}|$ on two vertical lines: $\zeta$ carries the Euler product and oscillates with the primes, whereas $\zeta_{\RVM}$ approaches the constant $\varpi_{\RVM}$ on $\sigma=2$ and $\varpi_{\RVM}|1+\chi|$ on $\sigma=0.75$.  This is the analytic content of Theorem~\CGref{thm:no-Dirichlet-RVM}.

\subsection{Exponential sums}
\label{subsec:num-expsums}

\begin{experiment}[Exponential sums over the two spectra]\label{exp:expsums}
\emph{Object:} $S(N;\log x)=\sum_{k\le N}x^{io_k}$ for $o_k=\tau_k$ and $o_k=\gamma_k$.  \emph{Range:} $N\le2\,001\,052$, $x\in\{2,3,4,6,10\}$.  \emph{Expected:} $|S_{\RVM}|\ll1+\log N$ (Theorem~\CGref{thm:exp-sum}); for $\gamma_k$, the main term $-\frac T{2\pi}\frac{\Lambda(x)}{\sqrt x}$ of Landau's formula (Theorem~\CGref{thm:known-landau}), whose error term $O_x(\log T)$ makes the comparison at a finite height illustrative.  \emph{Observed:} the maximum over all $N\le2\,001\,052$ of $|S_{\RVM}(N;\log x)|$ is $3.58$, $3.98$, $3.96$, $4.53$, $10.34$ (the figure samples every $2000$th $N$); at $N=2\,001\,052$, $|\sum x^{i\gamma_k}|=88\,301$, $114\,268$, $62\,435$, $0.26$, $2.42$, against the main terms $88\,342$, $114\,324$, $62\,467$, $0$, $0$.
\end{experiment}

\begin{figure}[htbp]
\centering
\includegraphics[width=0.86\textwidth]{fig04_exponential_sums.png}
\caption{Observed size of the exponential sums on the computed range $N\le2\,001\,052$, compared with $\log N$: the midpoint spectrum shows no prime frequencies (left); the spectrum of $\zeta$ resonates at prime powers (right).}
\label{fig:expsums}
\end{figure}

\subsection{Centroids, walks and Cram\'er's function}
\label{subsec:num-centroids}

\begin{experiment}[Cram\'er's function and the centroids of the two spectra]\label{exp:cramer}
\emph{Object:} $e^{-1/2}V_+(1)$ from the zero-free representation \CGeqref{eq:cramer-zero-free}; the triangular centroids $C^{\pathc}_\zeta(N;1)$ and $C^{\pathc}_{\RVM}(N;1)$; the limit $C^{\pathc}_{\RVM}(\infty;1)$ of Theorem~\CGref{thm:clock-centroid}.  \emph{Method:} program \texttt{s17} (Appendix~\CGref{app:code-s17}), mpmath with $30$ digits, rays at $\theta=\pi/3$, prime powers up to $2\cdot10^6$, $a=1$ and $a=2$.  \emph{Expected:} independence of $a$; $2\Re V_+(u)=e^u-(e^{2u}-1)^{-1}$ on $(\log2,\log3)$ \CGeqref{eq:cramer-real-window}; under RH, $C^{\pathc}_\zeta(N;1)\to e^{-1/2}V_+(1)$ (Theorem~\CGref{thm:centroid-RH}).  \emph{Observed:}
\begin{itemize}
\item for $a=2$ the value of $e^{-1/2}V_+(1)$ is $0.7768942607932536+0.7678556688644353\,i$, and for $a=1$ it is $0.7768942607932536+0.7678556452589\,i$; the difference $2.4\cdot10^{-8}$ is the truncation of the prime sum, of order $\sum_{n>2\cdot10^6}\Lambda(n)n^{-2}e^2/\log n$ for $a=1$ and below $10^{-12}$ for $a=2$; the real part equals $\frac1{2\sqrt e}(e-\frac1{e^2-1})=0.7768942607932536$;
\item $2\Re V_+(u)$ agrees with $e^u-(e^{2u}-1)^{-1}$ to all printed digits at $u=0.75$, $0.90$, $1.05$;
\item $C^{\pathc}_\zeta(N;1)=0.77776961+0.76786173\,i$, $0.77702246+0.76802715\,i$, $0.77689988+0.76784866\,i$, $0.7769071196+0.7678445908\,i$ for $N=10^4,10^5,10^6,2\,001\,052$: the distance to $e^{-1/2}V_+(1)$ is $1.7\cdot10^{-5}$ at $N=2\,001\,052$;
\item $C^{\pathc}_{\RVM}(N;1)=0.0415702193+1.7575417496\,i$ at $N=2\,001\,052$, and $S_n-b_nE_{n+1}=0.0415691734+1.7575409824\,i$, $0.0415688005+1.7575410722\,i$ at $n=10^6$ and $n=2\,001\,051$, so $C^{\pathc}_{\RVM}(\infty;1)=0.041569+1.757541\,i$.
\end{itemize}
\end{experiment}

\begin{experiment}[Walks of $e^{i\gamma_k}$, $e^{i\tau_k}$ and $e^{i(\gamma_k+\tau_k)/2}$]\label{exp:walks}
\emph{Object:} the walks $S_k=\sum_{j\le k}e^{ix_j}$ for $x=\gamma$, $\tau$ and $\theta=(\gamma+\tau)/2$, $k\le1000$, and their triangular centroids.  \emph{Method:} program \texttt{s16} (Appendix~\CGref{app:code-s16}).  \emph{Observed:} $C_\gamma=0.7720742579+0.7626754110\,i$, $C_\tau=0.0401443298+1.7565288414\,i$, $C_\theta=0.6133423097+1.3939939459\,i$; the midpoint of the chord is $\frac12(C_\gamma+C_\tau)=0.4061092939+1.2596021262\,i$, and $C_\theta$ lies $0.2470$ away from it (difference $0.2072+0.1344\,i$); $\max_{k\le1000}|S_k|=3.418$, $2.668$, $3.310$; $\gamma_k-\tau_k$ has mean $-2.2\cdot10^{-5}$ and standard deviation $0.3587$ on $k\le1000$.
\end{experiment}

The clock walk is a regular spiral around its centroid, the walk of the zeros an irregular cloud around $C_\gamma$, and the walk of the midpoints $\theta_k$ is intermediate in shape and position (Figures~\CGref{fig:walks3d} and~\CGref{fig:walks-proj}).  Since $e^{i\theta_k}=e^{i\gamma_k}e^{i\eta_k/2}$, the centroid of the midpoints is not the midpoint of the centroids: the passage from the clock to the zeros acts on the walk through the phases $e^{i\eta_k/2}$, not linearly.

\begin{figure}[htbp]
\centering
\includegraphics[width=\textwidth]{fig24_walks_3d_1000.png}
\caption{The walks $(\Re S_k,\Im S_k,k)$, $k\le1000$, with $k$ on the vertical axis, for $e^{i\gamma_k}$ (left), $e^{i\tau_k}$ (centre) and $e^{i(\gamma_k+\tau_k)/2}$ (right); the marker is the centroid at height $(N+1)/2$.}
\label{fig:walks3d}
\end{figure}

\begin{figure}[htbp]
\centering
\includegraphics[width=\textwidth]{fig25_walks_projection_1000.png}
\caption{Left: the three walks of Figure~\CGref{fig:walks3d} together.  Right: their projection on the complex plane with the three centroids and the midpoint of the chord $\frac12(C_\gamma+C_\tau)$ (cross).}
\label{fig:walks-proj}
\end{figure}

\subsection{Prime counts: \texorpdfstring{$\pi$}{pi} and \texorpdfstring{$\pi_{\RVM}$}{pi\_RVM}}
\label{subsec:num-primes}

\begin{experiment}[Prime counts and RVM primes]\label{exp:primes}
\emph{Object:} $J_{\RVM}$, $\pi_{\RVM}$, $c_n^{\RVM}$, $p_j^{\RVM}$ against $J$, $\pi$, the primes.  \emph{Range:} $2\le x\le100$ with all $N$ ordinates of each list ($3000$ points); $2\le x\le1000$ with the first $K=17\,061$ ordinates and a smooth cut ($40\,000$ points).  \emph{Expected:} the explicit formula with the $\gamma_k$ reproduces $\pi(x)$; $\pi_{\RVM}$ is continuous and $\pi_{\RVM}(x)-\li(x)=O(\sqrt x(\log\log x)^2/\log x)$ (Theorem~\CGref{thm:Lambert-estimate}).  \emph{Observed:} away from the primes, $|\pi_{\gamma}-\pi|\le0.40$ on $[2,100]$; $\pi_{\RVM}(100)=26.08$, $\pi(100)=25$, $\operatorname{li}(100)=30.13$; $\pi_{\RVM}(1000)=169.98$, $\pi(1000)=168$, $\operatorname{li}(1000)=177.61$.  $\pi_{\RVM}$ is not monotone: it carries a ripple of unit scale, with minimum grid increment $-0.25$ on $[2,100]$.  The linear interpolation of $J_{\RVM}(x^{1/m})$ in the M\"obius inversion changes $\pi_{\RVM}$ by at most $5.2\cdot10^{-5}$ when the step in $u$ is doubled, an interpolation error of about $2\cdot10^{-5}$ at the step used.  The first grid crossings of the levels $j=1,\dots,8$, refined by linear interpolation, are $p_j^{\RVM}=2.21,3.15,5.19,7.16,11.14,14.02,17.26,19.18$ against $p_j=2,3,5,7,11,13,17,19$.
\end{experiment}

\begin{figure}[htbp]
\centering
\includegraphics[width=0.86\textwidth]{fig05_prime_counts.png}
\caption{$\pi(x)$ (black staircase), the explicit formula with the zeros of $\zeta$ (blue), $\pi_{\RVM}(x)$ (red) and $\operatorname{li}(x)-\operatorname{li}(2)$ (green).  The distance between the red and green curves is the quantity bounded in Theorem~\CGref{thm:Lambert-estimate}; the ratios are in Figure~\CGref{fig:pi-ratios}.}
\label{fig:prime-counts}
\end{figure}

\begin{figure}[htbp]
\centering
\includegraphics[width=0.86\textwidth]{fig06_rvm_primes.png}
\caption{Left: cell masses $c_n^{\RVM}$, the prime indicator and $1/\log n$.  Right: real RVM prime levels against the primes up to $1000$.}
\label{fig:rvm-primes}
\end{figure}

Figure~\CGref{fig:prime-counts} shows $\pi$ and $\pi_{\RVM}$ in the usual format.  The RVM count follows the prime staircase much more closely than $\operatorname{li}$: replacing each $\gamma_k$ by $\tau_k$ moves every zero by less than one unit (Experiment~\CGref{exp:clock}), and the low ordinates, which shape $\pi(x)$ for small $x$, are moved least in relative terms.  The steps are replaced by a continuous ripple, so the cell masses $c_n^{\RVM}$ of Figure~\CGref{fig:rvm-primes} are real numbers, occasionally negative, and are not a prime indicator.

\begin{experiment}[The three ratios and the square-root scale]\label{exp:pi-ratios}
\emph{Object:} $\li/\pi_{\RVM}$, $\pi/\li$, $\pi/\pi_{\RVM}$, and $(\pi_{\RVM}-\li)/B$, $(\pi-\li)/B$ with $B(x)=\sqrt x(\log\log x)^2/\log x$.  \emph{Range:} the grid of Experiment~\CGref{exp:primes} on $[8,1000]$ ($K=17\,061$ ordinates).  \emph{Expected:} the three ratios tend to $1$ (Corollary~\CGref{cor:three-limits}) and $(\pi_{\RVM}-\li)/B$ stays bounded (Theorem~\CGref{thm:Lambert-estimate}).  \emph{Observed:} at $x=100,500,1000$, $\li/\pi_{\RVM}=1.155,1.086,1.045$, $\pi/\li=0.830,0.933,0.946$, $\pi/\pi_{\RVM}=0.959,1.013,0.988$; $\max|\pi_{\RVM}-\li|/B=2.72$ on $[8,1000]$, $1.35$ on $[50,1000]$ and $0.72$ on $[500,1000]$, while $\max|\pi-\li|/B=2.45$ on $[8,1000]$ (Figure~\CGref{fig:pi-ratios}).
\end{experiment}

\begin{figure}[htbp]
\centering
\includegraphics[width=0.86\textwidth]{fig10_prime_ratios.png}
\caption{Left: the three ratios of Corollary~\CGref{cor:three-limits}.  Right: $\pi_{\RVM}-\li$ and $\pi-\li$ in units of the scale $B(x)$ of Theorem~\CGref{thm:Lambert-estimate}.  The ratio $\pi/\pi_{\RVM}$ (black) is the closest of the three to $1$; the ratios with $\li$ approach $1$ at the rate $\tfrac12\li(\sqrt x)/\li(x)$ of the bias $J-\pi$.}
\label{fig:pi-ratios}
\end{figure}

On $[2,1000]$ the two normalized errors in Figure~\CGref{fig:pi-ratios} have the same size and sign: both counts lie below $\li$ by about $\tfrac12\li(\sqrt x)$, the $m=2$ term of the M\"obius inversion \CGeqref{eq:pi-minus-li}.  This is why $\pi/\pi_{\RVM}$ is much closer to $1$ than either ratio with $\li$.

\subsection{The Mertens functions and the ternary division}
\label{subsec:num-mertens}

The Mertens function is the sum of the ternary sequence $\mu(n)\in\{-1,0,+1\}$.  Its values are specified by two pieces of data:
\begin{enumerate}[label=\textup{(\alph*)}]
\item \emph{the zero set} $\{n:\mu(n)=0\}$, the non-squarefree integers, of density $1-6/\pi^2=0.392073$;
\item \emph{the signs} on the squarefree integers, $\mu(n)=(-1)^{\omega(n)}$, each sign of density $3/\pi^2=0.303964$.
\end{enumerate}
The ternary walk $\mu_h$ of \CGeqref{eq:intro-muh}--\CGeqref{eq:intro-quantizer} keeps (a) exactly, $\mu_h(n)=0\Leftrightarrow\mu(n)=0$, and replaces (b): at a squarefree $n$ it steps up when $M_h(n-1)$ lies below the target $\mathcal M(n)=M_{\RVM}^{\crit}(n)$ and down otherwise.  By construction the walk follows its target within one unit, so its closeness to $M_{\RVM}^{\crit}$ is a property of the quantizer and carries no information about $M$; any coincidence of signs between $\mu_h$ and $\mu$ is reported as an observation without a statistical analysis and is not claimed to have theoretical content.  Thus
\begin{equation}\label{eq:ternary-update}
 M_h(n)=M_h(n-1)+\mu(n)^2\,\varsigma_n,\qquad |M_h(n)-\mathcal M(n)|\ \text{small},
\end{equation}
so that $M_h$ is an integer walk with the zero set of $\mu$ that follows $M_{\RVM}^{\crit}$, in the same way as $M$ is the integer walk with the zero set of $\mu$ whose signs are $(-1)^{\omega(n)}$.

\begin{table}[htbp]
\centering
\small
\begin{tabular}{rrrrrcrrr}
\toprule
$n$ & $\mu(n)$ & $\mu(n)^2$ & $\mathcal M(n)$ & $M_h(n-1)$ & $\varsigma_n$ & $\mu_h(n)$ & $M_h(n)$ & $M(n)$\\
\midrule
1 & 1 & 1 & 0.2925 & 0 & $+1$ & 1 & 1 & 1\\
2 & $-1$ & 1 & 0.4563 & 1 & $-1$ & $-1$ & 0 & 0\\
3 & $-1$ & 1 & $-0.6523$ & 0 & $-1$ & $-1$ & $-1$ & $-1$\\
4 & 0 & 0 & $-0.8720$ & $-1$ & -- & 0 & $-1$ & $-1$\\
5 & $-1$ & 1 & $-1.7255$ & $-1$ & $-1$ & $-1$ & $-2$ & $-2$\\
6 & 1 & 1 & $-1.3063$ & $-2$ & $+1$ & 1 & $-1$ & $-1$\\
7 & $-1$ & 1 & $-1.5931$ & $-1$ & $-1$ & $-1$ & $-2$ & $-2$\\
8 & 0 & 0 & $-2.4029$ & $-2$ & -- & 0 & $-2$ & $-2$\\
9 & 0 & 0 & $-1.3292$ & $-2$ & -- & 0 & $-2$ & $-2$\\
10 & 1 & 1 & $-1.3280$ & $-2$ & $+1$ & 1 & $-1$ & $-1$\\
11 & $-1$ & 1 & $-2.3272$ & $-1$ & $-1$ & $-1$ & $-2$ & $-2$\\
12 & 0 & 0 & $-2.2067$ & $-2$ & -- & 0 & $-2$ & $-2$\\
13 & $-1$ & 1 & $-1.8468$ & $-2$ & $+1$ & 1 & $-1$ & $-3$\\
14 & 1 & 1 & $-1.9324$ & $-1$ & $-1$ & $-1$ & $-2$ & $-2$\\
15 & 1 & 1 & $-1.5435$ & $-2$ & $+1$ & 1 & $-1$ & $-1$\\
\bottomrule
\end{tabular}
\caption{The ternary division for $n\le15$ (the first $40$ values are in Appendix~\CGref{app:out-s05-mertens}).  The zero columns of $\mu$ and $\mu_h$ coincide; the signs of $\mu_h$ are chosen by comparing $M_h(n-1)$ with $\mathcal M(n)$.  For $n\le12$ the two walks agree; at $n=13$ the target $-1.85$ lies above $M_h(12)=-2$, so $\mu_h(13)=+1$ while $\mu(13)=-1$.}
\label{tab:ternary}
\end{table}

\begin{experiment}[Mertens functions and the ternary walk]\label{exp:mertens}
\emph{Object:} $M$, $M_{\RVM}^{\crit}$, $M_h$; the explicit formulas.  \emph{Range:} $n\le10^6$; explicit formulas on $[2,2000]$ with all $N$ ordinates of each list.  \emph{Expected:} the explicit formula with the zeros of $\zeta$ reproduces $M$; $M_{\RVM}^{\crit}(x)=O(x^{1/2+\eps})$ (Theorem~\CGref{thm:unsmoothed-RVM}); $M_h$ tracks $M_{\RVM}^{\crit}$ with the densities of $\mu$.  \emph{Observed:}
\begin{itemize}
\item explicit formula for $M$ with the $\gamma_k$: root mean square error $0.0008$, $0.0005$, $0.0004$ on $[2,2000]$ at distance more than $0.02$, $0.05$, $0.10$ from the integers; the derivatives $\zeta'(\rho_k)$ for $k>1000$ come from the Riemann--Siegel formula with its first correction term, whose relative difference from mpmath at spot checks is at most $2.3\cdot10^{-5}$, and the trivial terms are truncated at six (seven terms change them by $1.5\cdot10^{-5}$ at $x=2$);
\item $\zeta_{\RVM}'(\rho_k^{\RVM})/\log(\tau_k/2\pi)=0.96194$ at $k=1000$, against $\varpi_{\RVM}=0.96194$ (Corollary~\CGref{cor:derivative-zeros}); the values $|\zeta'(\rho_k)|/\log(\gamma_k/2\pi)$ are spread over $(0,3)$ (Figure~\CGref{fig:explicit});
\item truncated spectral sums $2\sum_{k\le N}|a_k|^2$ over the supplied ordinates: $0.029034$ for $\zeta$, $0.027203$ for the midpoint system;
\item the integer target agrees with the full explicit formula for $M_{\RVM}^{\crit}$ to $0.019$ at $100$ integers $\le10^5$ and to $0.018$ at $34$ integers in $[9\cdot10^5,10^6]$; the residues are computed from the product for $k\le1000$ and modelled beyond, so $\mathcal M$ is a tapered-and-interpolated target, not the exact function $M_{\RVM}^{\crit}$; the product formula for $\zeta_{\RVM}'(\rho_k^{\RVM})$, with its sign $(-1)^{k-1}$ (the factors $j<k$ of $P^{\crit}_{\RVM}{}'(\rho^{\RVM}_k)$ are negative), agrees with a central difference of $P^{\crit}_{\RVM}/G$ to $7\cdot10^{-11}$ for $k\le6$, and refitting $c$ on $250<k\le1000$ or adding a term $d/\tau_k^2$ to the model changes the target by at most $5.2\cdot10^{-6}$ at the $134$ check integers;
\item $\max|M_h-\mathcal M|=2.62$ on $[1,10^6]$ and $1.65$ on $[10^4,10^6]$;
\item densities of $(-1,0,+1)$ for $\mu$: $0.303857$, $0.392074$, $0.304069$;
\item densities of $(-1,0,+1)$ for $\mu_h$: $0.304003$, $0.392074$, $0.303923$;
\item $\mu_h(n)=\mu(n)$ at $50.03\%$ of the squarefree $n\le10^6$, and $\sum_{n\le10^6}\mu(n)\mu_h(n)/\#\{n\le10^6:\mu(n)\ne0\}=0.0007$: the empirical normalized correlation of the sign sequences is $0.0007$ on this range (no asymptotic statement is made), although both walks have the same zero set and comparable size;
\item $\max_{10\le n\le10^6}|W(n)|/\sqrt n=0.832,0.754,0.688$ for $W=M,\mathcal M,M_h$;
\item energies $\frac1N\sum_{n\le N}W(n)^2$ at $N=10^6$: $14\,851$, $12\,538$, $12\,538$ for $W=M,\mathcal M,M_h$; the values for $\mathcal M$ carry the truncation, residue-model, taper and interpolation errors of the target.
\end{itemize}
\end{experiment}

\begin{figure}[htbp]
\centering
\includegraphics[width=0.86\textwidth]{fig07_mertens.png}
\caption{Left: $M$ (black), the tapered-and-interpolated target $\mathcal M$ for $M_{\RVM}^{\crit}$ (red) and $M_h$ (green) at every integer $n\le2000$; the residues are computed from the product for $k\le1000$ and modelled beyond.  Right: range of $M$ and mean of $\mathcal M$ in blocks of $500$ integers up to $10^6$, with $\pm\sqrt n$.}
\label{fig:mertens}
\end{figure}

\begin{figure}[htbp]
\centering
\includegraphics[width=0.86\textwidth]{fig08_ternary_walk.png}
\caption{Left: $\mu(n)$ and $\mu_h(n)$ for $n\le60$: same zeros, different signs.  Right: $M$, the tapered-and-interpolated target $\mathcal M$ for $M_{\RVM}^{\crit}$ and $M_h$ for $n\le300$.}
\label{fig:ternary}
\end{figure}

\begin{figure}[htbp]
\centering
\includegraphics[width=0.86\textwidth]{fig09_explicit_weights.png}
\caption{Left: $M(x)$, its explicit formula from the zeros of $\zeta$ and $M_{\RVM}^{\crit}(x)$ from the $\tau_k$ on $[2,300]$.  Right: normalized derivatives at the zeros: erratic for $\zeta$, constant $\varpi_{\RVM}$ for $\zeta_{\RVM}$.}
\label{fig:explicit}
\end{figure}

Figures~\CGref{fig:mertens}--\CGref{fig:explicit} summarize the Mertens side.  $M_{\RVM}^{\crit}$ has the size and the visual texture of $M$, it stays inside $\pm\sqrt n$ in agreement with the square-root-scale behaviour of Theorem~\CGref{thm:unsmoothed-RVM}, and its residues are deterministic; $M_h$ is its ternary integer version with the zero set of $\mu$.  The passage from these synthetic objects to $M$ itself is not numerical: it is the identity $\zeta_{Z1}=\zeta$ of Step~3, which Experiment~\CGref{exp:Z1} is consistent with.


\begin{remark}[{The opposite phase of $M$ and $M_{\RVM}^{\crit}$ between $10^3$ and $2\cdot10^3$}]\label{rem:anticorrelation}
In the left panel of Figure~\CGref{fig:mertens} the curves of $M$ and $M_{\RVM}^{\crit}$ are in opposite phase between $n\approx10^3$ and $n\approx2\cdot10^3$.  Heuristically, the visible anticorrelation is largely explained by the first mode, $2\sqrt x\,\Re\bigl(x^{i\gamma_1}/(\rho_1\zeta'(\rho_1))\bigr)$ and $2\sqrt x\,\Re\bigl(x^{i\tau_1}/(\rho_1^{\RVM}\zeta_{\RVM}'(\rho_1^{\RVM}))\bigr)$, two oscillations in $\log x$ with the frequencies $\gamma_1$ and $\tau_1$.  Their phase difference is $\eta_1\log x$ with $\eta_1=\tau_1-\gamma_1=14.5213-14.1347=0.3866$, up to the constant difference of the arguments of the two residues, which is not computed here; $\eta_1\log x$ equals $\pi$ at $x=e^{\pi/\eta_1}\approx3.4\cdot10^3$ and lies between $2.67$ and $2.94$ on $[10^3,2\cdot10^3]$.  The beat has period $2\pi/\eta_1\approx16.3$ in $\log x$.  This is the rotation \CGeqref{eq:Xi-rotation} seen on the Mertens side: the passage from the clock to $\zeta$ turns the mode $k$ by the angle $\eta_k\log x$.  The higher modes turn at different rates, and at large scale neither a correlation nor an anticorrelation remains, in agreement with the value $0.0007$ of Experiment~\CGref{exp:mertens}.  The effect depends on $\eta_1$ alone and changes sign with $x$; it does not distinguish $\lambda=0$.
\end{remark}

\subsection{The first \texorpdfstring{$1000$}{1000} values}
\label{subsec:num-first-1000}

The four comparisons of the paper are tabulated for their first $1000$ values in Appendix~\CGref{app:tables-1000} (Tables~\CGref{tab:clock-1000}--\CGref{tab:mertens-1000}) and drawn in Figures~\CGref{fig:clock-1000}--\CGref{fig:mertens-1000}.  The functions are those of the preceding subsections: $\tau_k$ from \CGeqref{eq:lambert} (file \texttt{zLamb.bin}), $\gamma_k$ from \texttt{zeross.bin}; $\pi_{\RVM}$ from \CGeqref{eq:pi-RVM}; $p_j^{\RVM}$ the first crossing of the level $j$ by $\pi_{\RVM}$, as in Experiment~\CGref{exp:primes}; $M_{\RVM}^{\crit}(n)$ the integer target of Experiment~\CGref{exp:mertens} and Table~\CGref{tab:ternary}.

\begin{experiment}[The first $1000$ values]\label{exp:first-1000}
\emph{Object:} $\tau_k$ against $\gamma_k$, $p_j^{\RVM}$ against $p_j$, $\pi_{\RVM}(n)$ against $\pi(n)$, and $M_{\RVM}^{\crit}(n)$ against $M(n)$, for $k,j,n\le1000$.  \emph{Method:} program \texttt{s15} (Appendix~\CGref{app:code-s15}).  For $\pi_{\RVM}$ on $[2,8200]$ the zero sum of \CGeqref{eq:J-RVM} uses the first $K=183\,029$ ordinates ($\tau_K=128\,805.3=2.5\cdot2\pi\cdot8200$) with the smooth cut of Experiment~\CGref{exp:primes}, on $1.4\cdot10^6$ points in $u=\log x$; the oscillatory sums are evaluated by a nonuniform fast Fourier transform.  For $M_{\RVM}^{\crit}$ the tapered explicit formula of Experiment~\CGref{exp:mertens} is used; it differs from the ordered sum over all $2\,001\,052$ ordinates by at most $0.028$ on $n\le1000$.  \emph{Observed:}
\begin{itemize}
\item clock: $\max_{k\le1000}|\eta_k|=0.9937$ (at $k=34$), mean $\eta_k=2.2\cdot10^{-5}$, root mean square $0.3587$; $\gamma_{1000}=1419.422481$, $\tau_{1000}=1419.517765$;
\item prime counts: $\pi_{\RVM}(100)=26.0782$, $\pi_{\RVM}(1000)=169.9750$, $\pi_{\RVM}(8200)=1029.688$ against $25$, $168$, $1028$; on $n\le1000$, $\pi_{\RVM}(n)-\pi(n)$ lies in $[-2.957,3.779]$, with mean $0.055$ and root mean square $1.187$; the means of $\pi_{\RVM}(n)-\li(n)$ and $\pi(n)-\li(n)$ are $-6.98$ and $-7.04$;
\item RVM prime levels: $p_1^{\RVM},\dots,p_{12}^{\RVM}=2.2138$, $3.1546$, $5.1880$, $7.1588$, $11.1391$, $14.0201$, $17.2551$, $19.1767$, $26.0524$, $28.0419$, $33.0665$, $38.0878$; $p_{1000}^{\RVM}=7928.195$ against $p_{1000}=7919$; $p_j^{\RVM}-p_j\in[-52.80,65.25]$, with root mean square $19.26$; $\max_{j\le1000}|p_j^{\RVM}-p_j|/\sqrt{p_j}=1.069$;
\item Mertens: $M_{\RVM}^{\crit}(1000)=-8.646$ against $M(1000)=2$; $M_{\RVM}^{\crit}(n)-M(n)\in[-17.08,13.88]$, with root mean square $6.44$; the maxima of $|M(n)|/\sqrt n$ and of $|M_{\RVM}^{\crit}(n)|/\sqrt n$ over $10\le n\le1000$ are $0.832$ and $0.754$; the ternary walk $\mu_h$ agrees in sign with $\mu$ on $51.5\%$ of the squarefree $n\le1000$.
\end{itemize}
\end{experiment}

The proximity of $p_j^{\RVM}$ to $p_j$ is that of $\pi_{\RVM}$ to $\pi$ read horizontally: the error $p_j^{\RVM}-p_j$ is numerically of the scale $\sqrt{p_j}$ on this range, as $\pi_{\RVM}(x)-\pi(x)$ is numerically of the order of $\sqrt x/\log x$ there (Figure~\CGref{fig:pi-1000}); it is not small in units of the prime gap.  The Mertens functions have the same size and the same $\sqrt n$ envelope, and show little sign agreement on the computed range (Remark~\CGref{rem:anticorrelation}).  Odlyzko and te~Riele \cite{OdlyzkoTeRiele1985} used the first $2000$ zeros of $\zeta$ to disprove the Mertens conjecture; here the explicit formula of $M_{\RVM}^{\crit}$ runs over the clock ordinates with the residues $1/(\rho_k^{\RVM}\zeta_{\RVM}'(\rho_k^{\RVM}))$, which are deterministic.

\begin{figure}[htbp]
\centering
\includegraphics[width=0.92\textwidth]{fig20_clock_1000.png}
\caption{The clock against the zeros, $k\le1000$: $\gamma_k$ and $\tau_k$; the displacements $\eta_k=\tau_k-\gamma_k$; the two spectra for $k\le60$; the distribution of $\eta_k$ (Table~\CGref{tab:clock-1000}).}
\label{fig:clock-1000}
\end{figure}

\begin{figure}[htbp]
\centering
\includegraphics[width=0.92\textwidth]{fig21_rvm_primes_1000.png}
\caption{The RVM prime levels $p_j^{\RVM}$ against the primes $p_j$, $j\le1000$, and the difference $p_j^{\RVM}-p_j$ (Table~\CGref{tab:primes-1000}).}
\label{fig:primes-1000}
\end{figure}

\begin{figure}[htbp]
\centering
\includegraphics[width=0.92\textwidth]{fig22_pi_1000.png}
\caption{$\pi(n)$, $\pi_{\RVM}(n)$ and $\li(n)-\li(2)$ at the integers $n\le1000$, and the difference $\pi_{\RVM}(n)-\pi(n)$ (Table~\CGref{tab:pi-1000}).}
\label{fig:pi-1000}
\end{figure}

\begin{figure}[htbp]
\centering
\includegraphics[width=0.92\textwidth]{fig23_mertens_1000.png}
\caption{$M(n)$ and $M_{\RVM}^{\crit}(n)$ at the integers $n\le1000$ with $\pm\sqrt n$, and the difference $M_{\RVM}^{\crit}(n)-M(n)$ (Table~\CGref{tab:mertens-1000}).}
\label{fig:mertens-1000}
\end{figure}

\subsection{The rung equations: transport, the defect of the clock, and the solution from the clock}
\label{subsec:num-rungs}

The four experiments of this subsection evaluate the objects of Subsections~\CGref{subsec:transport} and~\CGref{subsec:free-masses}.  In the first two, $\mathsf f$ is the function $\exp\bigl(-1/(1-y^2)\bigr)$, $y=(u-0.36)/0.26$, on $[0.10,0.62]\subset(0,\log2)$, normalized to $\int\mathsf f=1$; it tests the first rung.  In the programs $\mathsf f$, $\mathsf k_{\mathsf f}$ and $\mathsf c_{\mathsf f}$ are called \texttt{psi}, \texttt{W\_psi} and \texttt{c\_psi}.

\begin{experiment}[Transport on the first rung]\label{exp:transport}
\emph{Object:} $\mathsf R_{\mathsf f}(1)$, $\mathsf S_{\mathsf f}(\lambda)$ and the moments $\mathsf m_j(\mathsf f)$ of \CGeqref{eq:rung-functional}--\CGeqref{eq:transport-moments}.  \emph{Range:} the first $20\,000$ ordinates of each list; $\max_{k\ge1000}|\mathsf c_{\mathsf f}(\tau_k)|=2.5\cdot10^{-10}$ and $\max_{k\ge3000}|\mathsf c_{\mathsf f}(\tau_k)|=2.1\cdot10^{-13}$.  \emph{Expected:} if $\zeta_{Z1}=\zeta$, then $\mathsf S_{\mathsf f}(0)=\mathsf R_{\mathsf f}(1)$ (Proposition~\CGref{prop:rung-functional}); the series \CGeqref{eq:transport-moments} converges to $\mathsf S_{\mathsf f}(0)$ with terms bounded by $C(1+j)^2(C_\eta\log2)^j/j!$ (Theorem~\CGref{thm:transport-band}).  \emph{Observed:}
\begin{itemize}
\item $\mathsf k_{\mathsf f}=0.133643530916$, $\sum_k\mathsf c_{\mathsf f}(\tau_k)=0.176106813076$, $\sum_k\mathsf c_{\mathsf f}(\gamma_k)=0.133643530929$;
\item defect of the clock $\mathsf R_{\mathsf f}(1)=-0.042463282161$; transport $\mathsf S_{\mathsf f}(0)=-0.042463282147$; $\mathsf R_{\mathsf f}(0)=\mathsf R_{\mathsf f}(1)-\mathsf S_{\mathsf f}(0)=-1.4\cdot10^{-11}$, and halving the number of quadrature nodes changes $\mathsf R_{\mathsf f}(1)$ and $\mathsf S_{\mathsf f}(0)$ by $4\cdot10^{-13}$ and $6\cdot10^{-13}$; by Theorem~\CGref{thm:rung-zero-sum}, $\mathsf R_{\mathsf f}(0)=\langle\mathscr Y,\mathsf f\rangle$ is a sum over the zeros off the critical line, all of which have $\gamma>3\cdot10^{12}$ \cite{PlattTrudgian2021}, where $|\widehat{\mathsf f}|$ is far below $10^{-11}$ by \CGeqref{eq:f-hat-decay}; the value $-1.4\cdot10^{-11}$ is the difference of two quantities computed from the same $20\,000$ ordinates with the same quadrature, consistent with the truncation and the quadrature error; it is not a measurement of $\mathscr Y$;
\item moment series (Table~\CGref{tab:transport}): the term $j=1$ is $95.1\%$ of $\mathsf R_{\mathsf f}(1)$, the partial sum with $j\le2$ is $100.8\%$, and the partial sums with $j\le3$, $j\le5$, $j\le8$ differ from $\mathsf R_{\mathsf f}(1)$ by $3.7\cdot10^{-6}$, $6.5\cdot10^{-9}$, $1.0\cdot10^{-11}$;
\item along the family (Table~\CGref{tab:transport}), extended numerically to $\lambda\in[-1,2]$ by the same formula $\gamma_k(\lambda)=\gamma_k+\lambda\eta_k$, $\mathsf R_{\mathsf f}(\lambda)$ is decreasing on the sampled values and almost linear, with slope about $-0.044$ and a single sign change, at $\lambda=0$; outside $[0,1]$ this is only a numerical prolongation of the formula: the products remain defined, since $|\gamma_k(\lambda)-\tau_k|\le(1+|\lambda|)C_\eta$, but the list is not increasing at $\lambda=-1$, $-0.5$ and $2$, and the results of Section~\CGref{sec:lambda-family} are not claimed there;
\item the partial sums $\sum_{k\le K_0}\bigl(\mathsf c_{\mathsf f}(\gamma_k)-\mathsf c_{\mathsf f}(\tau_k)\bigr)$ are $-0.02086$, $-0.04381$, $-0.04649$, $-0.04221$, $-0.04249$, $-0.0424634$ for $K_0=1,5,10,50,100,400$: the first ordinate gives half of the transport and, for this test function, the first-rung sum is numerically stabilised by about fifty ordinates.
\end{itemize}
\end{experiment}

\begin{table}[htbp]
\centering
\small
\begin{tabular}{rrrr@{\qquad\qquad}rr}
\toprule
$j$ & $\mathsf m_j(\mathsf f)$ & $(-1)^j\mathsf m_j(\mathsf f)/j!$ & partial sum $-\ \mathsf R_{\mathsf f}(1)$ & $\lambda$ & $\mathsf R_{\mathsf f}(\lambda)$\\
\midrule
1 & $+4.0395\cdot10^{-2}$ & $-4.0395\cdot10^{-2}$ & $+2.07\cdot10^{-3}$ & $-1.00$ & $+0.045303554710$\\
2 & $-4.8098\cdot10^{-3}$ & $-2.4049\cdot10^{-3}$ & $-3.37\cdot10^{-4}$ & $-0.50$ & $+0.022413444752$\\
3 & $-2.0432\cdot10^{-3}$ & $+3.4053\cdot10^{-4}$ & $+3.69\cdot10^{-6}$ & $-0.25$ & $+0.011133149816$\\
4 & $-5.5184\cdot10^{-5}$ & $-2.2993\cdot10^{-6}$ & $+1.39\cdot10^{-6}$ & $0.00$ & $-0.000000000014$\\
5 & $+1.6734\cdot10^{-4}$ & $-1.3945\cdot10^{-6}$ & $-6.47\cdot10^{-9}$ & $+0.25$ & $-0.010956709511$\\
6 & $+5.5721\cdot10^{-6}$ & $+7.7390\cdot10^{-9}$ & $+1.27\cdot10^{-9}$ & $+0.50$ & $-0.021706821780$\\
7 & $+7.2608\cdot10^{-6}$ & $-1.4406\cdot10^{-9}$ & $-1.69\cdot10^{-10}$ & $+1.00$ & $-0.042463282161$\\
8 & $+6.4039\cdot10^{-6}$ & $+1.5883\cdot10^{-10}$ & $-9.77\cdot10^{-12}$ & $+2.00$ & $-0.080112190380$\\
\bottomrule
\end{tabular}
\caption{Experiment~\CGref{exp:transport}.  Left: the moments $\mathsf m_j(\mathsf f)=\sum_k\eta_k^{\,j}\mathsf c^{(j)}_{\mathsf f}(\tau_k)$ and the partial sums of the series \CGeqref{eq:transport-moments} for $\mathsf S_{\mathsf f}(0)$, compared with the defect $\mathsf R_{\mathsf f}(1)=-0.042463282161$ of the clock.  Right: the rung functional along the family.}
\label{tab:transport}
\end{table}

\begin{experiment}[The defect of the clock by Poisson summation]\label{exp:clock-defect}
\emph{Object:} the two parts of $\mathsf R_{\mathsf f}(1)$ in Proposition~\CGref{prop:clock-defect-poisson}.  \emph{Expected:} $\mathsf R_{\mathsf f}(1)=(\mathsf k_{\mathsf f}-\mathsf I_0)-2\sum_{m\ge1}(-1)^m\mathsf I_m$.  \emph{Observed:} $\mathfrak t_0=9.676906787$; $\mathsf I_0=0.063533498702$ from \CGeqref{eq:clock-defect-I0}; the mismatch of the smooth densities is $\mathsf k_{\mathsf f}-\mathsf I_0=+0.070110032214$, and the oscillating part, obtained as the difference $\mathsf R_{\mathsf f}(1)-(\mathsf k_{\mathsf f}-\mathsf I_0)$ and not by summing the $\mathsf I_m$, is $-0.112573314375$; the program does not verify independently that the oscillatory series $-2\sum_{m\ge1}(-1)^m\mathsf I_m$ converges to this value.  The stationary heights $2\pi e^u$, $u$ in the support of $\mathsf f$, fill $[6.944,11.680]$; only those with $u>0.4319$ lie above $\mathfrak t_0$.
\end{experiment}

\begin{experiment}[From the clock to the ordinates and to the masses of the primes]\label{exp:clock-to-zeros}
\emph{Object:} a finite section of the system \CGeqref{eq:rung-system}: the unknowns are the ordinates $\kappa_k$ below a height and the masses $\mathsf a_n$, $2\le n\le11$; the equations are the pairings of \CGeqref{eq:rung-system} with $\omega_{\rm test}(u)\cos(\omega u)$ and $\omega_{\rm test}(u)\sin(\omega u)$, where $\omega_{\rm test}$ is a fixed smooth function supported in $(0,\log13.5)$ and $\omega$ runs over a grid of step $0.1$ in $[0,\Omega]$.  The ordinates above $\Omega+12$ are kept on the clock; only the first $KT=6000$ clock ordinates are summed numerically.  Although the support window $(0,\log13.5)$ contains $\log12$ and $\log13$, the implementation drops mass columns whose bump weight is below $0.05$, so those two masses are fixed to zero in this run.  These are explicit truncation sources in the residual.  \emph{Method:} Levenberg--Marquardt iteration started at $\kappa_k=\tau_k$; the masses enter linearly and are solved for at each step by ridge regression (the least-squares problem for the masses carries the ridge term $10^{-10}|\mathsf a|^2$, so they are not eliminated exactly; this problem and the damped steps are solved by QR/SVD on stacked matrices, not by normal equations, and a step is accepted only if it decreases the residual), with continuation in height $\Omega=18,22,\dots,50$.  The proximity to $\gamma_k$ is never used inside the solver, only as a diagnostic afterwards.  The file \texttt{zeross.bin} is read only after the iteration, for the comparison.  \emph{Expected:} by Theorem~\CGref{thm:free-masses}, if $\zeta_{Z1}=\zeta$ the system has the unique solution $\kappa_k=\gamma_k$, $\mathsf a_n=\log p/(2\sqrt n)$ at prime powers and $0$ elsewhere.  \emph{Observed:} the iteration reaches small residuals at every stage (at two stages it stops because no damping up to $10^{12}$ decreases the residual further) (residual norms between $4\cdot10^{-4}$ and $5\cdot10^{-3}$, of the size of the contribution of the ordinates kept on the clock); the solution is in Table~\CGref{tab:clock-to-zeros}: the first ten ordinates agree with $\gamma_1,\dots,\gamma_{10}$ within $4.4\cdot10^{-4}$, the eleventh within $2.2\cdot10^{-3}$, and the masses agree with \CGeqref{eq:free-masses} within $8\cdot10^{-4}$, including the values $0$ at $n=6,10$ and the prime-power values at $n=4,8,9$.  The ordinates $k=12,13,14$, which lie above the band $\omega\le50$, are not determined (errors $0.10$, $0.77$, $1.73$).  With masses up to $n\le34$ and the band $\omega\le60$ the section is not resolved and the iteration does not converge: resolving the ordinates up to a height $\Omega$ requires a strict margin $\log N>(1+\varkappa)\log(\Omega/2\pi)$ for some fixed $\varkappa>0$ in a stable finite solve and resolving the masses up to $N$ requires $\Omega$ of the order of $\pi N$, in agreement with the heights $2\pi x^{1/m}$ of Proposition~\CGref{prop:clock-defect-poisson}.  The experiment illustrates the existence form of the final step; it is a computation and not a proof of existence.
\end{experiment}

\begin{table}[htbp]
\centering
\small
\begin{tabular}{rrrr@{\qquad}rrr}
\toprule
$k$ & $\tau_k$ (start) & solution & $\gamma_k$ & $n$ & $\mathsf a_n$ (solution) & \CGeqref{eq:free-masses}\\
\midrule
1 & 14.521346953 & 14.134660056 & 14.134725142 & 2 & $+0.245076213$ & 0.245064536\\
2 & 20.655740356 & 21.022091631 & 21.022039639 & 3 & $+0.317145240$ & 0.317142050\\
3 & 25.492675432 & 25.010827129 & 25.010857580 & 4 & $+0.173268937$ & 0.173286795\\
4 & 29.739411632 & 30.424917668 & 30.424876126 & 5 & $+0.359884977$ & 0.359881258\\
5 & 33.624531889 & 32.935072953 & 32.935061588 & 6 & $-0.000004422$ & 0\\
6 & 37.257370087 & 37.586146548 & 37.586178159 & 7 & $+0.367785301$ & 0.367742452\\
7 & 40.700551604 & 40.918919600 & 40.918719012 & 8 & $+0.122615622$ & 0.122532268\\
8 & 43.994013953 & 43.327322117 & 43.327073281 & 9 & $+0.183143344$ & 0.183102048\\
9 & 47.165128427 & 48.004716753 & 48.005150881 & 10 & $-0.000286111$ & 0\\
10 & 50.233653391 & 49.773657726 & 49.773832478 & 11 & $+0.362276961$ & 0.361496314\\
11 & 53.214422874 & 52.972535790 & 52.970321478 & & &\\
\bottomrule
\end{tabular}
\caption{Experiment~\CGref{exp:clock-to-zeros}.  Solution of a finite section of the rung system \CGeqref{eq:rung-system}, started at the clock.  Left: the ordinates.  Right: the masses at the points $\log n$, which are not assumed, against $\log p/(2\sqrt n)$ at prime powers $n=p^m$.}
\label{tab:clock-to-zeros}
\end{table}


\begin{experiment}[Multiscale Jacobian of the clock]\label{exp:multiscale-jacobian}
\emph{Object:} the minimum midpoint-clock gap, its bound \CGeqref{eq:clock-gap-lower}, the Ingham critical window $2\pi/\Delta_T$, and the adjacent-pair bound \CGeqref{eq:fixed-window-pair}.  \emph{Input:} none of the genuine zeta ordinates; the program \texttt{s12\_multiscale\_jacobian.py} generates the $\tau_k$ directly from the Lambert formula up to $K=2\,001\,052$.  \emph{Observed:}
\begin{center}
\begin{tabular}{rrrr}
\toprule
$K$ & $T=\tau_K$ & $\Delta_T$ & $2\pi/\log(T/2\pi)$\\
\midrule
$10^2$ & $235.987267$ & $1.73462445$ & $1.73286190$\\
$10^3$ & $1419.517765$ & $1.15930485$ & $1.15921748$\\
$10^4$ & $9877.629616$ & $0.85368118$ & $0.85367617$\\
$10^5$ & $74920.891033$ & $0.66939912$ & $0.66939880$\\
$5\cdot10^5$ & $319387.322632$ & $0.57982852$ & $0.57982848$\\
$2\,001\,052$ & $1132490.722588$ & $0.51918344$ & $0.51918343$\\
\bottomrule
\end{tabular}
\end{center}
At the last height, $2\pi/\Delta_T=12.102053$ while $\log(T/2\pi)=12.10205288$.  Thus the rigorous scale $U\asymp\log T$ is numerically almost exact for the midpoint clock.  On every fixed window the adjacent-pair upper bound decreases with the clock gap, illustrating the obstruction in Theorem~\CGref{thm:clock-jacobian}(iii).  These calculations use no $\gamma_k$ and no $\eta_k$.
\end{experiment}

\begin{figure}[htbp]
\centering
\includegraphics[width=0.82\textwidth]{fig14_multiscale_gap.png}
\caption{Minimum gap of the midpoint clock up to height $T$ and the explicit lower bound $2\pi/\log(T/2\pi)$.  Their ratio is already $1+2\cdot10^{-8}$ at $K=2\,001\,052$.}
\label{fig:multiscale-gap}
\end{figure}

\begin{figure}[htbp]
\centering
\includegraphics[width=0.82\textwidth]{fig15_fixed_window_degeneracy.png}
\caption{Upper bound in \CGeqref{eq:fixed-window-pair} for adjacent clock frequencies on four fixed $u$-windows.  The decay is slow, of order $1/\log\tau_k$, but it is enough to exclude a height-uniform lower frame bound on any fixed window.}
\label{fig:fixed-window-degeneracy}
\end{figure}

\clearpage
\section{Conclusions and open problems}
\label{sec:conclusions}

\subsection*{What is new, what is proved, what is assumed, what is open}

\emph{New in this paper.}  The objects and results that, to the author's knowledge, do not appear elsewhere are: the synthetic zeta function $\zeta_{\RVM}$ built on the midpoint clock (Sections~\CGref{sec:clock}--\CGref{sec:zeta-rvm}; the clock and its Lambert inversion are due to Fran\c{c}a and LeClair \cite{FrancaLeClair2014}, Subsection~\CGref{subsec:known}); the study of the exponential sums over the ordinates of $\zeta$ against Cram\'er's function, with the zero-free representation of $V_+$, the forms of RH at one imaginary frequency and on the window $(\log2,\log3)$ with the exact defect, the identification of the endpoint rung functional with the sum over the zeros off the line and the reduction of the ladder to the cell $(2,3)$ with the sine identity, the proposition on finitely many zeros off the line, the Frullani form, the convergence of the clock centroid and the value of the ordinate centroid under RH (Section~\CGref{sec:ordinate-sums}); the crossing point of the family and the family with the opposite endpoint (Subsection~\CGref{subsec:crossing}); the tables of the first $1000$ values (Appendix~\CGref{app:tables-1000}); the unconditional square-root estimate $\pi_{\RVM}(x)-\li(x)=O(\sqrt x(\log\log x)^2/\log x)$ for the synthetic prime count, with its continuous prime levels and cell masses (Sections~\CGref{sec:explicit-formula}--\CGref{sec:integer-weights}); the transition theorem $\zeta_{\RVM}=\varpi_{\RVM}(1+w)(1+O(1/|t|))$ with its constant $\varpi_{\RVM}$, the closed form $\nu_3=8-\frac{33}4\log2-2\log\pi$ of its smooth part and the inequality $\varpi_{\RVM}<1$ (Theorem~\CGref{thm:varpi-less-one}) and the synthetic Mertens bound $M^{\crit}_{\RVM}(x)=L+O_\eps(x^{1/2+\eps})$ (Section~\CGref{sec:transition}); the $\lambda$-family from the clock to the zeta ordinates, with the endpoint $\zeta_{Z1}$, the value $\varpi_{Z1}$ at infinity as an ordered product over the displacements, and the equivalences (D)$\Leftrightarrow\zeta_{Z1}=\zeta\Leftrightarrow\varpi_{Z1}=1\Leftrightarrow$RH (Sections~\CGref{sec:lambda-family}--\CGref{sec:euler-products}); the atomization of every member of the family by smooth cells, the mean-square bound $X^{1+\eps}$ for the cell masses uniformly in $\lambda$, and Hamburger's theorem for the pair (continuous, atomized) in its symmetric, theta and vertical mean-square forms (Sections~\CGref{sec:atomized}--\CGref{sec:criteria}); the compact space $\mathscr O$ of lists with the rigidity, existence-form and free-mass theorems, and the nondegeneracy theorem for the rung system (Section~\CGref{sec:lists}); the acceptance criterion and the two finite models in which the final step is a theorem (Section~\CGref{sec:models}); in Route B for finite sections, the uniqueness of exact solutions in every tube with free masses and the mean-square bound for their masses, the proposition that finite stages do not decide route B, the window forms with an unknown list and the Bernstein form, the dichotomy for minimizers on windows, the identification of the Laplace transform of the residual as $\frac12(P'_{\boldsymbol\kappa}/P_{\boldsymbol\kappa}-\xi'/\xi)$ with the boundary term $\mathfrak H(0)\delta_0'$, the certificate theorem for staged sections, the verified exact solutions of fifteen finite sections, Gonek's lists in a tube around the clock, and the Kadec--Avdonin criterion with blocks of variable length (Section~\CGref{sec:routeB-finite}); and the a priori displacement proposition with its drift constant $\mathrm b$ (Appendix~\CGref{subsec:rung-apriori}).  The equivalence of B$'$ with the Riemann hypothesis is classical in substance (Remark~\CGref{rem:rb-prescribed}), and no novelty is claimed for the variable-block variant beyond its formulation.

\emph{Proved without hypotheses.}  Every numbered statement (1)--(41) and (43)--(51) below is a theorem of this paper, proved with no hypothesis on $\zeta$; statement (42) assumes RH explicitly; including the secondary lemmas and corollaries that carry them (the bounded displacement lemma, the weighted clock sum, the Lipschitz control of $J_{\RVM}$, the mean-square convergence of the ordinate series, the transfer theorem to $\mathscr O$, the compactness and continuity on $\mathscr O$, the point-support and finite-order steps of the nondegeneracy theorem, the Laplace-transform lemma of Route B, the monotone matching lemma); Section~\CGref{sec:routeB-finite} has its own status table (Section~\CGref{subsec:rb-status}).

\emph{No unproved hypothesis.}  The inequality $\varpi_{\RVM}<1$ is Theorem~\CGref{thm:varpi-less-one}, with $\varpi_{\RVM}<0.9796$; the five statements that use it, Proposition~\CGref{prop:atomized-clock}(ii), the sign $\mathsf W(1)>0$ in Proposition~\CGref{prop:one-number}, Proposition~\CGref{prop:open-path-varpi}, Corollary~\CGref{cor:ladder-pair} at the clock end, and the remark after Corollary~\CGref{cor:positivity}, are therefore unconditional.  Hypothesis~\CGref{hyp:cell-remainder} is proved (Corollary~\CGref{cor:clock-cells}) and is a hypothesis only in name.  The jump-growth condition in Theorem~\CGref{thm:cell-form}(ii) is part of a statement equivalent to (D), not an assumption of the paper.  The statements that carry a hypothesis are those that assume (D) explicitly (Theorems~\CGref{thm:step3} and~\CGref{thm:transfer}) and the a priori Proposition~\CGref{prop:rung-apriori}, which assumes an exact solution; they are statements about (D) and about hypothetical solutions, not uses of an unproved lemma.

\emph{Equivalent to (D) and not proved.}  The identity $\zeta_{Z1}=\zeta$, (D), RH, and each of the statements listed in Subsection~\CGref{subsec:status} and after item~(50) below: the identity \CGeqref{eq:imaginary-frequency}; the window identity of Theorem~\CGref{thm:window-RH}(ii) and the distributional convergence of the ordinate centroids on the window; the cell $(2,3)$ together with the sine identity \CGeqref{eq:sine-identity}; $\mathscr Q_\infty=\log L$; $\varpi_{Z1}=1$; the product identity \CGeqref{eq:product-identity}; the cell identity \CGeqref{eq:cell-identity}; $\mathsf W(0)\ge0$; the vanishing of the cell defects at $\lambda=0$; the symmetry of $e^{\tilde h_0}$; $T_0\equiv0$; $\mathsf V_\sigma(0)=0$; the existence of a coincidence list in $\mathscr O$; the solvability of the rung system in $\mathscr O$ or in any tube; statements B and B$'$; the window identity for an unknown list (Theorem~\CGref{thm:rb-window-unknown}); the complete monotonicity of $e^{-iy/2}V_+(iy)$ (Theorem~\CGref{thm:rb-bernstein}); the solvability of the rung system on every window in a fixed box (Theorem~\CGref{thm:rb-window-equivalence}).  These are reductions, each proved as an equivalence; none of them is resolved.

\emph{Open and not claimed.}  The existence of a solution of the rung system; the injectivity of the clock-only linearization and any continuation from the clock; the convergence of finite sections to a distributional limit; a certified enclosure of $\varpi_{\RVM}$ beyond the bounds $0<\varpi_{\RVM}<0.9796$ of Theorem~\CGref{thm:varpi-less-one} (the ten digits $0.9619415789$ are floating-point values); and any computer-assisted statement beyond Theorem~\CGref{thm:rb-verified}, which concerns fifteen finite square sections and, by Proposition~\CGref{prop:rb-nonfinitary}, has no bearing on RH (the Gaussian certificates of Section~\CGref{subsec:rb-numerics} are floating-point candidates); the exclusion of degenerate minimizers in Theorem~\CGref{thm:rb-window-equivalence}; and the relation between window solvability and zero-free regions (Heuristic~\CGref{heur:rb-window-zerofree}).

\emph{Dependence table.}  The following table records what each group of results depends on.

\begingroup\small
\begin{longtable}{p{0.46\textwidth}p{0.46\textwidth}}
\toprule
Results & Depend on\\
\midrule
\endfirsthead
\toprule
Results & Depend on\\
\midrule
\endhead
\bottomrule
\endfoot
Clock, $\zeta_{\RVM}$, exponential sums, $J_{\RVM}$, $\pi_{\RVM}$, unit cell, transition theorem, synthetic Mertens bound (Sections~\CGref{sec:clock}--\CGref{sec:transition}) & the clock only; no property of $\zeta$ beyond the Riemann--von Mangoldt formula (for Lemma~\CGref{lem:bounded-displacement})\\
$\lambda$-family, $\zeta_{Z1}$, $\varpi_{Z1}$, Mellin counts, explicit formula, cell form (Sections~\CGref{sec:lambda-family}--\CGref{sec:euler-products}) & the ordinates $\gamma_k$ as a list in the tube, i.e.\ membership in $\mathscr O$; no hypothesis\\
Atomized family, mean square of cell masses, joining factor, pair criteria (Sections~\CGref{sec:atomized}--\CGref{sec:criteria}) & membership in $\mathscr O$; no hypothesis\\
$\varpi_{\RVM}<1$ (Theorem~\CGref{thm:varpi-less-one}) and the five statements that use it & the clock only; the explicit inequalities of the proof are verified in interval arithmetic (program \texttt{s14}, Appendix~\CGref{app:code-s14})\\
Step~3 and the transfer of the Mertens bound to $M$ (Theorems~\CGref{thm:step3}, \CGref{thm:transfer}) & (D), as an explicit hypothesis of the statement\\
Rigidity, existence form, free masses, nondegeneracy (Section~\CGref{sec:lists}) & membership in $\mathscr O$; nondegeneracy assumes the existence of a solution\\
Route B: compactness, uniqueness with free and with prescribed masses, window forms, minimizers, certificate, Gonek lists, $(A_2)$ criterion (Section~\CGref{sec:routeB-finite}) & $C$-admissibility (tube) only; Theorem~\CGref{thm:rb-uniqueness} also uses Proposition~\CGref{prop:rung-apriori} and Hamburger's theorem, the window forms Lemma~\CGref{lem:edge} and Proposition~\CGref{prop:cramer-zero-free}; B, B$'$ and the window forms are equivalent to RH and not proved\\
Proposition~\CGref{prop:rb-nonfinitary} & the zeros verified up to $3\cdot10^{12}$ \cite{PlattTrudgian2021} and the bound of \cite{HasanalizadeShenWong2022} for $N_\zeta$\\
Theorem~\CGref{thm:rb-verified} & ball-arithmetic computation (program \texttt{certV}; Arb through python-flint)\\
A priori displacement (Appendix~\CGref{subsec:rung-apriori}) & the existence of an exact solution with the local-density and mean-square hypotheses stated there\\
$|\log\varpi_{Z1}|<4\cdot10^{-13}$ (Proposition~\CGref{prop:varpi-verified}) and $\lambda^*<2\cdot10^{-11}$ (Proposition~\CGref{prop:crossing}) & the verified zeros up to $T_{\rm ver}$ (input data)\\
Section~\CGref{sec:ordinate-sums} & Landau's formula, Selberg's mean value of $S$, Cram\'er's function, the explicit formula; Theorem~\CGref{thm:centroid-RH} assumes RH explicitly; all other statements are unconditional\\
Section~\CGref{sec:numerics}, Subsection~\CGref{subsec:rb-numerics} except Theorem~\CGref{thm:rb-verified} & floating-point computations with the supplied files; not part of any proof\\
\end{longtable}
\endgroup

\subsection*{The proved statements}

The following statements are proved (the results of Section~\CGref{sec:routeB-finite} are summarised in its status table, Section~\CGref{subsec:rb-status}):
\begin{enumerate}[label=\textup{(\arabic*)}]
\item the midpoint ordinates $\tau_k$ are given by the Lambert formula \CGeqref{eq:lambert} and stay within a bounded distance of the ordinates $\gamma_k$ of $\zeta$ (Lemma~\CGref{lem:bounded-displacement});
\item $P_{\RVM}^{\crit}$ is entire and symmetric, and $\zeta_{\RVM}$ has exactly the simple zeros $\tfrac12\pm i\tau_k$, the trivial zeros and a residue-one pole at $1$ (Theorem~\CGref{thm:canonical});
\item $S_{\RVM}(N;u)\ll1+\log N$ uniformly on compact positive frequency intervals, so the point centroid and $M_H^{\RVM}(x)$ tend to $0$ (Section~\CGref{sec:exp-sums});
\item $J_{\RVM}$ exists compact-uniformly and is related to $\pi_{\RVM}$ by exact finite M\"obius inversion (Section~\CGref{sec:explicit-formula});
\item $J_{\RVM}(x)-\li(x)$ and $\pi_{\RVM}(x)-\li(x)$ are $O(\sqrt x(\log\log x)^2/\log x)$, so $\li/\pi_{\RVM}\to1$, and with the prime number theorem $\pi/\li\to1$ and $\pi/\pi_{\RVM}\to1$ (Theorem~\CGref{thm:Lambert-estimate}, Corollary~\CGref{cor:three-limits});
\item $J_{\RVM}$ is Lipschitz on $[2,\infty)$, $c_n^{\RVM}=O(1)$, and $\zeta_{\Z,\RVM}$ is an absolutely convergent Dirichlet series in $\Re s>1$ (Section~\CGref{sec:integer-weights});
\item $\zeta_{\RVM}=\varpi_{\RVM}(1+w)(1+O(1/|t|))$ on every band to the right of $\tfrac12$ and $\zeta_{\RVM}=\varpi_{\RVM}(1+\chi)(1+O(1/|t|))$ on closed substrips of $\tfrac12<\Re s<1$ (Theorem~\CGref{thm:transition}); the transition constant is $\varpi_{\RVM}=P^{\crit}_{\RVM}(\frac12)e^{\nu_3+\nu_4}<0.9796$, with $\nu_3=8-\frac{33}4\log2-2\log\pi$ and $\nu_4<-0.0126$ (Theorem~\CGref{thm:varpi-less-one});
\item $1/\zeta_{\RVM}\to L$ on every vertical line to the right of $\tfrac12$, $\zeta_{\RVM}^{\pm1}$ have no absolutely convergent Dirichlet expansion, and $M_{\RVM}^{\crit}(x)=L+O_\eps(x^{1/2+\eps})$ (Section~\CGref{sec:transition});
\item $2\xi$ equals its Hadamard product over $\rho(1-\rho)$ exactly (Theorem~\CGref{thm:Hadamard}), and $|\zeta|\asymp|\mathcal Q_{\RVM}|$ on $\frac12+\eps\le\sigma\le c'$, $|t|\ge t_0$ (Corollary~\CGref{cor:Q-control});
\item the $\lambda$-family from $\zeta_{\RVM}$ to $\zeta_{Z1}$ converges jointly, has simple zeros for $\lambda>0$, an explicit generator and total multiplier, and its members differ only through $\gamma_k-\tau_k$ (Section~\CGref{sec:lambda-family});
\item $\zeta_{Z1}$ has the functional equation, the pole, the residue and the order of $\zeta$, and $\zeta_{Z1}=\zeta\Leftrightarrow$(D)$\Leftrightarrow$RH (Theorem~\CGref{thm:Hamburger-Z1});
\item $\zeta_{\RVM}$ is bounded above and below on every half-plane $\Re s\ge\sigma_0>1$ and tends to $\varpi_{\RVM}$ as $\Re s\to+\infty$ (Theorem~\CGref{thm:halfplane-RVM}); $1/\zeta_{\Z,\RVM}$ is an absolutely convergent Dirichlet series in $\Re s>1$, so (D) is equivalent to an absolutely convergent Dirichlet expansion of $\mathcal K$ in $\Re s>1$ (Theorem~\CGref{thm:D-multiplier});
\item no member of the $\lambda$-family, $\zeta_{Z1}$ included, can have a zero off the critical line other than the trivial zeros (Theorem~\CGref{thm:path});
\item the endpoint has the value $\varpi_{Z1}=\varpi_{\RVM}\lim_{K\to\infty}\prod_{k\le K}(\frac14+\tau_k^2)/(\frac14+\gamma_k^2)\ge1$ at $+\infty$, and $\zeta_{Z1}=\zeta$ if and only if $\varpi_{Z1}=1$ (Theorems~\CGref{thm:Z1-clock}, \CGref{thm:varpi-Z1}, Corollary~\CGref{cor:product-identity});
\item the prime-power counts of $\zeta_{\RVM}$ and $\zeta_{Z1}$ are the Mellin transforms of their logarithms and satisfy a distributional explicit formula; this proves the Mellin--Stieltjes identification and the weighted factorization $\zeta_{\Z,\RVM}=\zeta_{\RVM}e^{h_{\RVM}}$, and, with the jump-growth hypothesis of Theorem~\CGref{thm:cell-form}(ii), (D) holds if and only if the prime-power count of $\zeta_{Z1}$ is constant between consecutive integers (Theorems~\CGref{thm:mellin-counts}, \CGref{thm:distributional-EF}, \CGref{thm:MS}, \CGref{thm:cell-form});
\item along the family $\frac d{d\lambda}\log\varpi_\lambda=-\Upsilon(\lambda)$, and (D) is the identity $\int_0^1\Upsilon\,d\lambda=\log L$; with the verified zeros, $0\le\log\varpi_{Z1}<4\cdot10^{-13}$ (Propositions~\CGref{prop:family-flow}, \CGref{prop:varpi-verified});
\item the atomized family $\tilde\zeta_{\Z,\lambda}=\zeta_\lambda e^{\tilde h_\lambda}$ consists of absolutely convergent Dirichlet series in $\Re s>1$ for every $\lambda$, the endpoint included, because the cell masses have mean square $X^{1+\eps}$ on dyadic blocks uniformly in $\lambda$; the joining factor is explicit cell by cell and holomorphic in $\Re s>\frac12$ (in $\Re s>0$ for $\lambda=1$), and (D) is equivalent to the vanishing of every cell defect at $\lambda=0$ and of $J_{Z1}$ on the first cell $(1,2)$ (Proposition~\CGref{prop:first-cell}, Theorems~\CGref{thm:joining-factor}, \CGref{thm:mean-square-cells}, \CGref{thm:atomized-continuation}, \CGref{thm:closing});
\item at the clock endpoint the joining factor is not constant, $\tilde h_1(\sigma)\to\log L$; the atomized reciprocals have M\"obius coefficients $\tilde\mu_\lambda$, and the endpoint Mertens function has a Mellin form $B_{Z1}$ with $B_{Z1}=M$ if and only if (D) (Propositions~\CGref{prop:atomized-clock}, \CGref{prop:Mellin-Mertens});
\item Step~3: $\zeta_{Z1}=\zeta$ implies $M(x)=O_\eps(x^{1/2+\eps})$ (Theorem~\CGref{thm:step3}); equivalently, the synthetic bound $M_{\RVM}^{\crit}(x)=L+O_\eps(x^{1/2+\eps})$ transfers to $M$ through $1/\zeta=F_{\RVM}\Pi_{\rm tr}^{-1}$ (Theorem~\CGref{thm:transfer});
\item for $0<\lambda\le1$ no $\zeta_\lambda$ is an absolutely convergent Dirichlet series in $\Re s>1$ and no joining factor $\tilde h_\lambda$ is constant, and $\lambda\mapsto\log\zeta_\lambda(\sigma)$ is strictly convex for every real $\sigma>1$ (Propositions~\CGref{prop:nonconstant}, \CGref{prop:convexity});
\item if the joining factor $e^{\tilde h_\lambda}$ is meromorphic of finite order with finitely many poles and symmetric under $s\mapsto1-s$, Hamburger's theorem gives $\tilde\zeta_{\Z,\lambda}=\zeta$, and for $\lambda=0$ also $\zeta_{Z1}=\zeta$ (Theorem~\CGref{thm:pair-criterion});
\item if $\zeta_\lambda$ is the Mellin--Stieltjes transform of a positive measure, then $\varpi_\lambda\le1$; (D) holds if and only if $\zeta_{Z1}$ has such a representation, if and only if $J^{\rm M}_{Z1}$ is non-decreasing (Theorem~\CGref{thm:positivity}, Corollary~\CGref{cor:positivity});
\item $J_{\RVM}$ is continuously differentiable on $(2,\infty)$ with $J_{\RVM}'(y)=\frac1{\log y}+2\cos(2\pi y)-2^{-1/2}y^{-1/4}\cos(4\pi\sqrt y+\frac\pi4)+O(y^{-1/3})$; consequently $\tilde c_n(1)=\frac1{\log n}+O(n^{-1/4})$ off the powers of $2$, and Hypothesis~\CGref{hyp:cell-remainder} holds off the integers $2^j$ with $\mu(j)\ne0$ (Theorem~\CGref{thm:clock-density}, Corollary~\CGref{cor:clock-cells});
\item the cell masses are sums over the ordinates with the height kernel $\hat{\varphi}(\gamma/2\pi n)$, which vanishes at the heights $2\pi qn$; their flow is the pairing of $V_n$ with the count difference $N_\zeta-n_{\RVM}$; and the Gaussian defect of the atomization is $\frac12(e^{\tilde h_\lambda(1)}-1)y^{-1/2}+O(y^{-1/4-\delta})$ (Propositions~\CGref{prop:height-kernels}, \CGref{prop:flow-count}, \CGref{prop:gaussian-defect});
\item $\zeta_{Z1}(\sigma)\ge\zeta(\sigma)$ for every real $\sigma>1$, with equality at one point if and only if $\zeta_{Z1}=\zeta$; in particular $\prod_k\bigl(1+2/(\frac14+\gamma_k^2)\bigr)\ge\pi/3$ (Theorem~\CGref{thm:real-axis});
\item every $\zeta_\lambda$ is the Mellin transform of a function $\mathsf A_\lambda$ that satisfies the summation formula of $\zeta$, and $\varpi_\lambda-1=-\mathsf W(\lambda)/\hat{\mathsf g}(1)$ for one explicit number $\mathsf W(\lambda)$; (D) is equivalent to $\mathsf W(0)\ge0$ and, by Theorem~\CGref{thm:varpi-less-one}, to $\varpi_\lambda<1$ for all $\lambda\in(0,1]$ (Lemma~\CGref{lem:summation-all}, Propositions~\CGref{prop:one-number}, \CGref{prop:open-path-varpi});
\item Hamburger's theorem for the pair holds without continuation of the joining factor: $\tilde\zeta_{\Z,\lambda}=\zeta$ if and only if the atomized coefficients satisfy the theta relation, if and only if the vertical mean square of $\tilde\zeta_{\Z,\lambda}/\zeta$ on a line $\Re s=\sigma>1$ is $1$; both defects are continuous in $\lambda$ (Theorem~\CGref{thm:double-hamburger});
\item at the clock end the joining factor has no symmetric meromorphic continuation, and $\tilde\zeta_{\Z,\lambda}=\zeta$ for at most finitely many $\lambda\in[0,1]$, not for $\lambda=1$ (Propositions~\CGref{prop:clock-asymmetry}, \CGref{prop:exceptional-set});
\item the cell masses are entire functions of $\lambda$ of exponential type $O(\log n)$, and the cell defects have the height kernel $\hat{\varphi}'(\gamma/2\pi n)$ (Propositions~\CGref{prop:entire-lambda}, \CGref{prop:defect-kernel});
\item in the compact space $\mathscr O$ of ordinate lists at bounded distance from the clock, at most one list has coinciding continuous and atomized zeta functions; such a list exists if and only if $\zeta_{Z1}=\zeta$, and it is then the list of the ordinates of $\zeta$ (Theorem~\CGref{thm:rigidity}, Proposition~\CGref{prop:lists-compact}, Corollary~\CGref{cor:existence-form});
\item (D) holds if and only if $\tilde c_n(0)=\mathbf 1_{\rm prime}(n)$ for every $n\ge2$, one real number for each cell; every $\lambda$ with $\tilde\zeta_{\Z,\lambda}=\zeta$ is a witness of (D), so the exceptional set is empty or contains $0$; cell by cell the remaining identity is \CGeqref{eq:per-cell} (Corollary~\CGref{cor:cell-mass-form}, Propositions~\CGref{prop:witness}, \CGref{prop:per-cell});
\item for a pair joined by a symmetric factor Hamburger's theorem gives a constant factor when the continuous function has order one, and not in general; at $\lambda=0$ the symmetry of the joining factor is equivalent to (D) (Proposition~\CGref{prop:dichotomy}, Remarks~\CGref{rem:pair-no-slack}, \CGref{rem:dichotomy-counterexample});
\item the ladder: for every $N\ge2$, $\tilde h_\lambda(\sigma)=O(\sigma N^{-\sigma})$ as $\sigma\to+\infty$ if and only if the Mellin count of $\zeta_\lambda$ has no mass outside the integers below $N$; (D) holds if and only if this holds at $\lambda=0$ for every $N$, and at the clock end the first rung fails (Theorems~\CGref{thm:ladder}, \CGref{thm:halfline-hamburger}, Corollary~\CGref{cor:ladder-pair});
\item the bridge theorem: if the imaginary part of the joining factor of a pair extends continuously to $\Re s\ge\frac12$, vanishes on the critical line, and the factor satisfies the temperate bound \CGeqref{eq:bridge-B2}, then the factor is constant, the discrete function is $\zeta$ and the continuous one is a multiple of $\zeta$; for the family the two conditions hold for no $\lambda\in(0,1]$, and at $\lambda=0$ they are equivalent to (D) (Theorem~\CGref{thm:bridge}, Corollary~\CGref{cor:bridge-family});
\item on the critical line $\Im\log\zeta_\lambda=-\vartheta(t)+\pi(N_\lambda(t)-1)$, and the cell masses of two lists differ by a sum over the ordinates of one function of one ordinate (Propositions~\CGref{prop:continuous-phase}, \CGref{prop:one-ordinate});
\item on each rung the passage from the clock to the ordinates of $\zeta$ is the transport $\mathsf S_{\mathsf f}(\lambda)$, an entire function of $\lambda$ of exponential type at most $C_\eta\log N$ whose value at $\lambda=0$ is the sum of the series of the moments $\sum_k\eta_k^{\,j}\mathsf c_{\mathsf f}^{(j)}(\tau_k)$; (D) holds if and only if the transport equals the defect of the clock on every rung, and this defect is given exactly by Poisson summation, with stationary heights $2\pi x^{1/m}$ (Proposition~\CGref{prop:rung-functional}, Theorem~\CGref{thm:transport-band}, Proposition~\CGref{prop:clock-defect-poisson});
\item (D) holds if and only if the rung system \CGeqref{eq:rung-system}, with free masses at the points $\log n$, has a solution in $\mathscr O$; the solution is then the list of $\zeta$ with the masses of von Mangoldt, and its linearization is injective up to the splitting of multiple ordinates (Theorems~\CGref{thm:free-masses}, \CGref{thm:nondegenerate});
\item at the midpoint clock, the finite ordinate Jacobian up to height $T$ has a uniform lower frame bound on windows of length $(1+\varkappa_{\rm win})\log(T/2\pi)$, no uniform lower bound is possible on a fixed window, the infinite ordinate block has dense range on every finite window, and smooth bounded-displacement transport tails are $O_{A,\mathsf f}(T^{1-A}\log T)$ (Theorem~\CGref{thm:clock-jacobian});
\item the gradient of the vertical defect on $\mathscr O$ is $\partial\mathsf V_\sigma/\partial\kappa_k=2\mathsf Q_{\boldsymbol{\kappa}}(\kappa_k)$ with one real-analytic function $\mathsf Q_{\boldsymbol{\kappa}}$ (Proposition~\CGref{prop:critical-points});
\item a list is accepted by a continuous side and a discrete side if and only if its function lies in the intersection of the window of the first with the frequencies of the second; a non-periodic list is accepted through one unitary matrix; for the finite Euler product $(1-2^{-s})(1-3^{-s})$ the constructions of the paper apply and the endpoint $\lambda=0$ is the only member of the family that is a Dirichlet series; for $\prod_{p\le X}(1-p^{-s})$ with more primes the Poisson part extends and the uniqueness part does not (Theorem~\CGref{thm:acceptance}, Proposition~\CGref{prop:graph-model}, Theorem~\CGref{thm:finite-euler} and the remark after it);
\item there is no function with the functional equation of $\zeta$ and an absolutely convergent Dirichlet series supported on the integers with bounded prime factors, and for the diagonal operator over the primes the contraction, the convergence of the determinant and the symmetry hold in $\sigma>0$, in $\sigma>1$ and on $\sigma=\frac12$ respectively (Proposition~\CGref{prop:no-finite-rank}, Subsection~\CGref{subsec:three-abscissas});
\item under RH, the triangular centroid $\sum_{k\le N}(1-\frac{k-1}N)e^{iu\gamma_k}$ tends to $e^{-u/2}V_+(u)$ for every $u>0$ that is not the logarithm of a prime power; at $u=1$ the limit is $0.77689426079\ldots+0.76785566886\ldots\,i$, with real part $\frac1{2\sqrt e}(e-\frac1{e^2-1})$ (Theorem~\CGref{thm:centroid-RH}, Lemma~\CGref{lem:local-tauberian});
\item the triangular centroid of $e^{iu\tau_k}$ converges for every $u>0$, with the limit given by the series \CGeqref{eq:clock-centroid-limit} and the rate $O((\log N)^2/N)$ (Theorem~\CGref{thm:clock-centroid});
\item Cram\'er's function has the zero-free representation \CGeqref{eq:cramer-zero-free} on the half-plane $\Omega_\theta$, with simple poles at the points $\log p^m$, and $2\Re V_+(u)=e^u-(e^{2u}-1)^{-1}$ on $(\log2,\log3)$ (Proposition~\CGref{prop:cramer-zero-free});
\item the defect between the ordinate transform $\mathscr Z_\gamma$ and the zero transform $\mathscr Z_\xi$ is $-\sum_k(\cosh(wd_k)-1)e^{iw\gamma_k}$, and RH holds if and only if $\sum_ke^{-y_0\gamma_k}=e^{-iy_0/2}V_+(iy_0)$ for one $y_0\in(0,4\pi]$ (Propositions~\CGref{prop:pair-decomposition}, \CGref{prop:imaginary-frequency});
\item on every open interval free of the points $\log p^m$, in particular on $(\log2,\log3)$, RH holds if and only if $\sum_k\widehat{\mathsf f}(\gamma_k)=\int\mathsf f(u)e^{-u/2}V_+(u)du$ for every test function supported there, equivalently if and only if the ordinate centroids converge there to $e^{-u/2}V_+(u)$ in the sense of distributions; without hypothesis the difference is the absolutely convergent sum \CGeqref{eq:window-residual} over the zeros off the line (Lemma~\CGref{lem:edge}, Theorem~\CGref{thm:window-RH});
\item the rung functional of the endpoint is the sum over the zeros off the line, $\mathsf R_{\mathsf f}(0)-\frac12\sum_n\Lambda(n)n^{-1/2}\mathsf f(\log n)=\langle\mathscr Y,\mathsf f\rangle$; equivalently $(\log x)(dJ^{\rm M}_{Z1}-dJ)/dx=2x^{-1/2}\mathscr Y(\log x)$; and the cosine difference is $\mathsf R_{\mathsf f}(1)-\langle\mathscr Y,\mathsf f\rangle$ on every rung (Theorem~\CGref{thm:rung-zero-sum});
\item the real part of the window identity on $(\log2,\log3)$ is the constancy of $J^{\rm M}_{Z1}$ on the cell $(2,3)$, and (D) is equivalent to that cell together with the sine identity \CGeqref{eq:sine-identity} (Corollary~\CGref{cor:one-window});
\item if the set of zeros off the critical line is finite and non-empty, then $\langle\mathscr Y,\mathsf f\rangle>0$ for some test function supported in any given interval, and $J^{\rm M}_{Z1}$ is constant on no interval $(n,n+1)$ (Proposition~\CGref{prop:finite-offline});
\item $2\int_0^\infty u^{-1}e^{-u/2}\sum_{k\le K}(\cos\gamma_ku-\cos\tau_ku)du\to\log L+\log\varpi_{Z1}\ge\log L$, with equality exactly under RH (Theorem~\CGref{thm:frullani});
\item $\log\varpi_\lambda$ satisfies the chord and tangent inequalities \CGeqref{eq:crossing-refuters} under RH, and if RH fails it vanishes at exactly one $\lambda^*$ with $f(0)/\Upsilon(0)\le\lambda^*\le f(0)/(f(0)+|f(1)|)$; the family of Proposition~\CGref{prop:kappa-family} has all the listed properties of the $\lambda$-family and $\varpi_0>1$ (Propositions~\CGref{prop:crossing}, \CGref{prop:kappa-family}).
\end{enumerate}
\emph{Not proved.}  The identity $\zeta_{Z1}=\zeta$ is not proved; neither is the Riemann hypothesis, to which it is equivalent.  The statements proved to be equivalent to it are listed in Subsection~\CGref{subsec:status} and above; in the form of Corollary~\CGref{cor:existence-form} it is the existence of one list in $\mathscr O$ whose continuous and atomized zeta functions coincide, and in the form of Corollary~\CGref{cor:bridge-family} it is the vanishing on the critical line of the imaginary part of the joining factor $\tilde h_0$, with the bound \CGeqref{eq:bridge-B2}.  The theorems of Subsections~\CGref{subsec:cell-mass-form}--\CGref{subsec:transport} and of Sections~\CGref{sec:lists} and~\CGref{sec:models} are formulations of this step and tools for it; none of them evaluates the pairing \CGeqref{eq:per-cell} of the count difference $N_\zeta-n_{\RVM}$, on which every one of them depends at $\lambda=0$.

The computations of Section~\CGref{sec:numerics} with $2\,001\,052$ zeros illustrate only the specific numerical statements explicitly tested there; they do not provide evidence for the unproved equivalent forms of (D).

\begin{remark}[Statements derived in part or conjectural, not proved here]\label{rem:not-proved-list}
The following statements arose in the work on the final step; none of them is used in this paper.
\begin{enumerate}[label=\textup{(\alph*)}]
\item \emph{Rates in $\lambda$.}  Off the powers of $2$ the formulas of Proposition~\CGref{prop:entire-lambda} suggest
\begin{equation}\label{eq:rates-lambda}
 \tilde h_\lambda(s)-\tilde h_0(s)=O_s(\lambda),\qquad\bigl|\mathfrak h_n(\lambda;s)-\mathfrak h_n(0;s)\bigr|\ll_{s,\eps}\lambda\,n^{-\sigma-\frac12+\eps};
\end{equation}
the cells $n=2^j$ and the uniformity in $n$ of the first-cell term are not written.
\item \emph{Flow form of the density condition.}  $\tilde h_0(1)=0$ if and only if $\int_0^1\partial_\lambda\tilde h_\lambda(1)\,d\lambda=\tilde h_1(1)$; the differentiability at $s=1$ uniformly in $\lambda$ is not written.
\item \emph{Pointwise bound.}  $|\tilde c_n(\lambda)|\le C_\eps n^\eps$ uniformly for $\lambda$ in a neighbourhood of $0$ is open; by Corollary~\CGref{cor:cell-mass-form} it is not needed for the final step.
\item \emph{Leading term of the clock defect.}  The form $\mathfrak h_n(1;s)\approx\sqrt2\,\pi\,\mathsf M_2\,s\,n^{-s-7/4}\sin(4\pi\sqrt n+\frac\pi4)$ is not proved; the bound \CGeqref{eq:clock-defect} is.
\item \emph{A heuristic model of the atomized clock.}  By \CGeqref{eq:clock-cell-smooth}, $\tilde\zeta_{\Z,1}(s)$ is close to $\prod_{n\ge2}(1-n^{-s})^{-1/\log n}$, whose logarithm has the leading part $\sum_n n^{-s}/\log n=\int_s^\infty(\zeta(y)-1)\,dy$; no error term is claimed.
\item The comparison between the two columns of Subsection~\CGref{subsec:hamburger-table} and the use of the classification of \cite{LevOlevskii2015} remain as discussed in the open-problem paragraph following Proposition~\CGref{prop:critical-points} and are not proved.
\end{enumerate}
\end{remark}

\begin{remark}[Open problems]\label{rem:open-problems}
\begin{enumerate}[label=\textup{(\roman*)}]
\item Prove the identity $\zeta_{Z1}=\zeta$.  Each of the following is proved to be equivalent to it, unconditionally and within the class of objects stated: (a) condition (D), an absolutely convergent Dirichlet series for $\zeta_{Z1}$ in $\Re s>1$ (Theorem~\CGref{thm:Hamburger-Z1}); (b) an absolutely convergent Dirichlet expansion of $\mathcal K=\zeta_{Z1}/\zeta_{\Z,\RVM}$ in $\Re s>1$; (c) the ordered product identity \CGeqref{eq:product-identity} (Corollary~\CGref{cor:product-identity}); (d) the cell identity \CGeqref{eq:cell-identity} together with the jump bound of Theorem~\CGref{thm:cell-form}; (e) step A7 of the programme in Subsection~\CGref{subsec:programme}; (f) a representation of $\zeta_{Z1}$ as the Mellin--Stieltjes transform \CGeqref{eq:representing-measure} of a positive measure on $[1,\infty)$ (Corollary~\CGref{cor:positivity}); (g) the existence of a list in $\mathscr O$ whose joining factor is constant, equivalently $\min_{\mathscr O}\mathsf V_\sigma=0$ (Corollary~\CGref{cor:existence-form}, step A19).  Conditions that are sufficient but not shown to be equivalent, such as the hypotheses of Theorem~\CGref{thm:pair-criterion} for $0<\lambda<1$, are those of item (iii).  This is the final step of the method.
\item Hypothesis~\CGref{hyp:cell-remainder} is proved for the integers $n$ that are not of the form $2^j$ with $\mu(j)\ne0$ (Corollary~\CGref{cor:clock-cells}); at $n=2^j$ the cell mass is $(1+\mu(j)J_{\RVM}(2)\log2)/\log n+O(n^{-1/4})$.  Determine the sign of $1\pm J_{\RVM}(2)\log2$.
\item Decide, for $0<\lambda<1$ outside the finite set $\mathsf E$, which hypothesis of Theorem~\CGref{thm:pair-criterion} fails, and which of \textup{(B1)}, \textup{(B2)} fails in Theorem~\CGref{thm:bridge}; decide whether the vanishing of all admissible coordinate derivatives (with the corresponding one-sided/KKT conditions at active constraints) forces $\mathsf V_\sigma=0$ (Proposition~\CGref{prop:critical-points}); and decide whether the atomized coefficients $\tilde a_N(0)$ control the sign of $\mathsf W(0)$ in Proposition~\CGref{prop:one-number}.
\item Determine the monotone branches of $\pi_{\RVM}$, on which the real prime levels $p_j^{\RVM}$ are defined, and explain the observed proximity $p_j^{\RVM}\approx p_j$ for small $j$ (Experiment~\CGref{exp:primes}).
\item Derive the bound \textup{(B2)} for $\tilde h_\lambda$ from the bounds of the cell defects, and turn Heuristic~\CGref{heur:phase-rotation} into a lemma with remainder; decide \textup{(B1)} at $\lambda=0$.
\item Evaluate from the construction the right side of \CGeqref{eq:per-cell}, or the moments $\mathsf m_j(\mathsf f)$ of \CGeqref{eq:transport-moments}; on the first rung the moments $j=1,2$ carry the transport up to $4\cdot10^{-6}$ (Experiment~\CGref{exp:transport}).
\item Theorem~\CGref{thm:clock-jacobian} gives the finite-height lower frame bound on the natural scale $U\asymp\log T$, the failure of fixed-window uniform invertibility, dense range on every fixed window, and rapidly summable smooth tails.  The exact-atomic interpolation tested later has no established interior path.  A remaining numerical existence programme is instead to control bounded endpoint finite sections on an exhausting nested family of tests and pass to a distributional limit; proving the required residual-to-zero and compactness statement would already be an existence theorem for the final rung system.
\item Under (D), decide whether $\mathsf E=\{0\}$, that is, whether the functions $\mathsf D_n$ of \CGeqref{eq:witness-zeros} have a common zero in $(0,1]$.
\item Construct an object of infinite rank from which $\zeta_{Z1}$ and $\tilde\zeta_{\Z,0}$ both descend, as the matrix $\mathsf U$ of Proposition~\CGref{prop:graph-model} does for the two-edge model (Subsection~\CGref{subsec:three-abscissas}).
\item Prove the window identity of Theorem~\CGref{thm:window-RH}(ii), or the cell $(2,3)$ together with the sine identity \CGeqref{eq:sine-identity}, directly from $dN_\zeta=d\NR+d\mathfrak e$ and the zero-free representation \CGeqref{eq:cramer-zero-free}, without a contour shift that brings back the terms of the zeros off the line (Remark~\CGref{rem:window}, Remark~\CGref{rem:conjugate-nonlocal}).
\item Decide whether the triangular centroid of $e^{i\gamma_k}$ at $u=1$ converges without RH, and whether its value $e^{-1/2}V_+(1)$, which RH implies (Theorem~\CGref{thm:centroid-RH}), implies anything about the zeros.
\item Find a kernel for which the Frullani form of Theorem~\CGref{thm:frullani} contracts with ratio $<1$ (Remark~\CGref{rem:frullani-arithmetic}).
\end{enumerate}
\end{remark}

\appendix
\section{Tables of the first \texorpdfstring{$1000$}{1000} values}
\label{app:tables-1000}

The tables below list the first $1000$ values of the four comparisons of Experiment~\CGref{exp:first-1000}; they are produced by program \texttt{s15} (Appendix~\CGref{app:code-s15}), which also writes them as CSV files in the folder \texttt{output} (\texttt{s15\_\allowbreak clock\_\allowbreak 1000.csv}, \texttt{s15\_\allowbreak rvm\_\allowbreak primes\_\allowbreak 1000.csv}, \texttt{s15\_\allowbreak pi\_\allowbreak 1000.csv}, \texttt{s15\_\allowbreak mertens\_\allowbreak 1000.csv}).  The ordinates $\gamma_k$ are those of the file \texttt{zeross.bin} and the midpoint ordinates $\tau_k$ those of \texttt{zLamb.bin}, which agree with the Lambert formula \CGeqref{eq:lambert} to $5\cdot10^{-13}$ on $k\le1000$.  Every difference column is computed from the unrounded values and rounded on its own.  After writing each table, \texttt{s15} reads it back and checks that every index $1,\dots,1000$ occurs exactly once (no missing or repeated rows) and that each displayed difference agrees with the difference of the displayed values up to the rounding of the displayed digits (console output in Appendix~\CGref{app:output-s15}).

{\footnotesize\setlength{\tabcolsep}{4pt}
\begin{longtable}{rrrr@{\qquad}rrrr}
\caption{The ordinates $\gamma_k$ of the zeros of $\zeta$, the midpoint ordinates $\tau_k$ of the clock and the displacements $\eta_k=\tau_k-\gamma_k$, $k\le1000$.  The displacements are computed from the full-precision ordinates and rounded independently, so they may differ in the last digit from the difference of the displayed columns.}\label{tab:clock-1000}\\
\toprule
$k$ & $\gamma_k$ & $\tau_k$ & $\eta_k$ & $k$ & $\gamma_k$ & $\tau_k$ & $\eta_k$\\
\midrule
\endfirsthead
\multicolumn{8}{l}{\emph{Table~\thetable\ (continued)}}\\
\toprule
$k$ & $\gamma_k$ & $\tau_k$ & $\eta_k$ & $k$ & $\gamma_k$ & $\tau_k$ & $\eta_k$\\
\midrule
\endhead
\bottomrule
\endfoot
1 & 14.134725 & 14.521347 & $+0.3866$ & 501 & 812.771108 & 812.740087 & $-0.0310$ \\
2 & 21.022040 & 20.655740 & $-0.3663$ & 502 & 814.045914 & 814.032038 & $-0.0139$ \\
3 & 25.010858 & 25.492675 & $+0.4818$ & 503 & 814.870540 & 815.323569 & $+0.4530$ \\
4 & 30.424876 & 29.739412 & $-0.6855$ & 504 & 816.727738 & 816.614679 & $-0.1131$ \\
5 & 32.935062 & 33.624532 & $+0.6895$ & 505 & 818.380669 & 817.905369 & $-0.4753$ \\
6 & 37.586178 & 37.257370 & $-0.3288$ & 506 & 819.204642 & 819.195642 & $-0.0090$ \\
7 & 40.918719 & 40.700552 & $-0.2182$ & 507 & 820.721898 & 820.485497 & $-0.2364$ \\
8 & 43.327073 & 43.994014 & $+0.6669$ & 508 & 821.713454 & 821.774936 & $+0.0615$ \\
9 & 48.005151 & 47.165128 & $-0.8400$ & 509 & 822.197757 & 823.063960 & $+0.8662$ \\
10 & 49.773832 & 50.233653 & $+0.4598$ & 510 & 824.526294 & 824.352571 & $-0.1737$ \\
11 & 52.970321 & 53.214423 & $+0.2441$ & 511 & 826.039287 & 825.640768 & $-0.3985$ \\
12 & 56.446248 & 56.118922 & $-0.3273$ & 512 & 826.905811 & 826.928553 & $+0.0227$ \\
13 & 59.347044 & 58.956266 & $-0.3908$ & 513 & 828.340174 & 828.215928 & $-0.1242$ \\
14 & 60.831779 & 61.733839 & $+0.9021$ & 514 & 829.437011 & 829.502893 & $+0.0659$ \\
15 & 65.112544 & 64.457725 & $-0.6548$ & 515 & 830.895884 & 830.789449 & $-0.1064$ \\
16 & 67.079811 & 67.133010 & $+0.0532$ & 516 & 831.799778 & 832.075597 & $+0.2758$ \\
17 & 69.546402 & 69.764002 & $+0.2176$ & 517 & 833.003641 & 833.361339 & $+0.3577$ \\
18 & 72.067158 & 72.354386 & $+0.2872$ & 518 & 834.651915 & 834.646674 & $-0.0052$ \\
19 & 75.704691 & 74.907346 & $-0.7973$ & 519 & 836.693576 & 835.931605 & $-0.7620$ \\
20 & 77.144840 & 77.425658 & $+0.2808$ & 520 & 837.347335 & 837.216133 & $-0.1312$ \\
21 & 79.337375 & 79.911757 & $+0.5744$ & 521 & 838.249022 & 838.500257 & $+0.2512$ \\
22 & 82.910381 & 82.367798 & $-0.5426$ & 522 & 839.465395 & 839.783980 & $+0.3186$ \\
23 & 84.735493 & 84.795695 & $+0.0602$ & 523 & 841.036390 & 841.067302 & $+0.0309$ \\
24 & 87.425275 & 87.197160 & $-0.2281$ & 524 & 842.041354 & 842.350225 & $+0.3089$ \\
25 & 88.809111 & 89.573734 & $+0.7646$ & 525 & 844.166197 & 843.632748 & $-0.5334$ \\
26 & 92.491899 & 91.926805 & $-0.5651$ & 526 & 844.805994 & 844.914874 & $+0.1089$ \\
27 & 94.651344 & 94.257633 & $-0.3937$ & 527 & 846.194770 & 846.196603 & $+0.0018$ \\
28 & 95.870634 & 96.567366 & $+0.6967$ & 528 & 847.971718 & 847.477936 & $-0.4938$ \\
29 & 98.831194 & 98.857051 & $+0.0259$ & 529 & 848.489281 & 848.758874 & $+0.2696$ \\
30 & 101.317851 & 101.127648 & $-0.1902$ & 530 & 849.862274 & 850.039419 & $+0.1771$ \\
31 & 103.725538 & 103.380040 & $-0.3455$ & 531 & 850.645448 & 851.319570 & $+0.6741$ \\
32 & 105.446623 & 105.615041 & $+0.1684$ & 532 & 853.163113 & 852.599329 & $-0.5638$ \\
33 & 107.168611 & 107.833404 & $+0.6648$ & 533 & 854.095512 & 853.878697 & $-0.2168$ \\
34 & 111.029536 & 110.035825 & $-0.9937$ & 534 & 855.286710 & 855.157676 & $-0.1290$ \\
35 & 111.874659 & 112.222953 & $+0.3483$ & 535 & 856.484117 & 856.436264 & $-0.0479$ \\
36 & 114.320221 & 114.395391 & $+0.0752$ & 536 & 857.310741 & 857.714465 & $+0.4037$ \\
37 & 116.226680 & 116.553701 & $+0.3270$ & 537 & 858.904026 & 858.992278 & $+0.0883$ \\
38 & 118.790783 & 118.698412 & $-0.0924$ & 538 & 860.410671 & 860.269705 & $-0.1410$ \\
39 & 121.370125 & 120.830014 & $-0.5401$ & 539 & 861.171098 & 861.546747 & $+0.3756$ \\
40 & 122.946829 & 122.948971 & $+0.0021$ & 540 & 863.189720 & 862.823404 & $-0.3663$ \\
41 & 124.256819 & 125.055718 & $+0.7989$ & 541 & 864.340824 & 864.099677 & $-0.2411$ \\
42 & 127.516684 & 127.150664 & $-0.3660$ & 542 & 865.594664 & 865.375567 & $-0.2191$ \\
43 & 129.578704 & 129.234193 & $-0.3445$ & 543 & 866.423740 & 866.651076 & $+0.2273$ \\
44 & 131.087689 & 131.306672 & $+0.2190$ & 544 & 867.693123 & 867.926204 & $+0.2331$ \\
45 & 133.497737 & 133.368444 & $-0.1293$ & 545 & 868.670494 & 869.200952 & $+0.5305$ \\
46 & 134.756510 & 135.419835 & $+0.6633$ & 546 & 870.846902 & 870.475320 & $-0.3716$ \\
47 & 138.116042 & 137.461154 & $-0.6549$ & 547 & 872.188751 & 871.749311 & $-0.4394$ \\
48 & 139.736209 & 139.492695 & $-0.2435$ & 548 & 873.098979 & 873.022924 & $-0.0761$ \\
49 & 141.123707 & 141.514736 & $+0.3910$ & 549 & 873.908389 & 874.296160 & $+0.3878$ \\
50 & 143.111846 & 143.527542 & $+0.4157$ & 550 & 875.985285 & 875.569021 & $-0.4163$ \\
51 & 146.000982 & 145.531366 & $-0.4696$ & 551 & 876.600826 & 876.841508 & $+0.2407$ \\
52 & 147.422765 & 147.526447 & $+0.1037$ & 552 & 877.654698 & 878.113620 & $+0.4589$ \\
53 & 150.053520 & 149.513016 & $-0.5405$ & 553 & 879.380952 & 879.385360 & $+0.0044$ \\
54 & 150.925258 & 151.491290 & $+0.5660$ & 554 & 880.834649 & 880.656727 & $-0.1779$ \\
55 & 153.024694 & 153.461481 & $+0.4368$ & 555 & 882.386697 & 881.927723 & $-0.4590$ \\
56 & 156.112909 & 155.423787 & $-0.6891$ & 556 & 883.430332 & 883.198349 & $-0.2320$ \\
57 & 157.597592 & 157.378401 & $-0.2192$ & 557 & 884.198743 & 884.468606 & $+0.2699$ \\
58 & 158.849988 & 159.325507 & $+0.4755$ & 558 & 885.272304 & 885.738493 & $+0.4662$ \\
59 & 161.188964 & 161.265279 & $+0.0763$ & 559 & 886.852802 & 887.008013 & $+0.1552$ \\
60 & 163.030710 & 163.197888 & $+0.1672$ & 560 & 888.475567 & 888.277166 & $-0.1984$ \\
61 & 165.537069 & 165.123495 & $-0.4136$ & 561 & 889.735294 & 889.545953 & $-0.1893$ \\
62 & 167.184440 & 167.042257 & $-0.1422$ & 562 & 890.813132 & 890.814375 & $+0.0012$ \\
63 & 169.094515 & 168.954323 & $-0.1402$ & 563 & 892.386433 & 892.082432 & $-0.3040$ \\
64 & 169.911976 & 170.859837 & $+0.9479$ & 564 & 893.119118 & 893.350125 & $+0.2310$ \\
65 & 173.411537 & 172.758939 & $-0.6526$ & 565 & 894.886292 & 894.617456 & $-0.2688$ \\
66 & 174.754192 & 174.651762 & $-0.1024$ & 566 & 895.397920 & 895.884424 & $+0.4865$ \\
67 & 176.441434 & 176.538435 & $+0.0970$ & 567 & 896.632252 & 897.151032 & $+0.5188$ \\
68 & 178.377408 & 178.419083 & $+0.0417$ & 568 & 899.221523 & 898.417279 & $-0.8042$ \\
69 & 179.916484 & 180.293825 & $+0.3773$ & 569 & 899.858885 & 899.683167 & $-0.1757$ \\
70 & 182.207078 & 182.162777 & $-0.0443$ & 570 & 900.849740 & 900.948696 & $+0.0990$ \\
71 & 184.874468 & 184.026051 & $-0.8484$ & 571 & 902.243208 & 902.213867 & $-0.0293$ \\
72 & 185.598784 & 185.883756 & $+0.2850$ & 572 & 903.099674 & 903.478681 & $+0.3790$ \\
73 & 187.228923 & 187.735995 & $+0.5071$ & 573 & 904.702903 & 904.743139 & $+0.0402$ \\
74 & 189.416159 & 189.582870 & $+0.1667$ & 574 & 905.829941 & 906.007241 & $+0.1773$ \\
75 & 192.026656 & 191.424480 & $-0.6022$ & 575 & 907.656729 & 907.270989 & $-0.3857$ \\
76 & 193.079727 & 193.260918 & $+0.1812$ & 576 & 908.333544 & 908.534382 & $+0.2008$ \\
77 & 195.265397 & 195.092276 & $-0.1731$ & 577 & 910.186334 & 909.797423 & $-0.3889$ \\
78 & 196.876482 & 196.918645 & $+0.0422$ & 578 & 911.234951 & 911.060111 & $-0.1748$ \\
79 & 198.015310 & 198.740109 & $+0.7248$ & 579 & 912.331046 & 912.322448 & $-0.0086$ \\
80 & 201.264752 & 200.556753 & $-0.7080$ & 580 & 912.823999 & 913.584434 & $+0.7604$ \\
81 & 202.493595 & 202.368658 & $-0.1249$ & 581 & 914.730097 & 914.846070 & $+0.1160$ \\
82 & 204.189672 & 204.175902 & $-0.0138$ & 582 & 916.355001 & 916.107357 & $-0.2476$ \\
83 & 205.394697 & 205.978563 & $+0.5839$ & 583 & 917.825378 & 917.368295 & $-0.4571$ \\
84 & 207.906259 & 207.776714 & $-0.1295$ & 584 & 918.836535 & 918.628885 & $-0.2076$ \\
85 & 209.576510 & 209.570428 & $-0.0061$ & 585 & 919.448344 & 919.889129 & $+0.4408$ \\
86 & 211.690863 & 211.359776 & $-0.3311$ & 586 & 921.156396 & 921.149027 & $-0.0074$ \\
87 & 213.347919 & 213.144824 & $-0.2031$ & 587 & 922.500629 & 922.408579 & $-0.0921$ \\
88 & 214.547045 & 214.925641 & $+0.3786$ & 588 & 923.285720 & 923.667786 & $+0.3821$ \\
89 & 216.169539 & 216.702290 & $+0.5328$ & 589 & 924.773484 & 924.926650 & $+0.1532$ \\
90 & 219.067596 & 218.474835 & $-0.5928$ & 590 & 926.551553 & 926.185170 & $-0.3664$ \\
91 & 220.714919 & 220.243336 & $-0.4716$ & 591 & 927.850859 & 927.443348 & $-0.4075$ \\
92 & 221.430706 & 222.007854 & $+0.5771$ & 592 & 928.663659 & 928.701185 & $+0.0375$ \\
93 & 224.007000 & 223.768447 & $-0.2386$ & 593 & 929.874093 & 929.958680 & $+0.0846$ \\
94 & 224.983325 & 225.525170 & $+0.5418$ & 594 & 931.009211 & 931.215836 & $+0.2066$ \\
95 & 227.421444 & 227.278080 & $-0.1434$ & 595 & 931.852741 & 932.472652 & $+0.6199$ \\
96 & 229.337413 & 229.027231 & $-0.3102$ & 596 & 934.385307 & 933.729129 & $-0.6562$ \\
97 & 231.250189 & 230.772674 & $-0.4775$ & 597 & 934.995425 & 934.985268 & $-0.0102$ \\
98 & 231.987235 & 232.514461 & $+0.5272$ & 598 & 936.228649 & 936.241070 & $+0.0124$ \\
99 & 233.693404 & 234.252643 & $+0.5592$ & 599 & 937.532926 & 937.496535 & $-0.0364$ \\
100 & 236.524230 & 235.987267 & $-0.5370$ & 600 & 939.024301 & 938.751665 & $-0.2726$ \\
101 & 237.769820 & 237.718382 & $-0.0514$ & 601 & 939.660941 & 940.006460 & $+0.3455$ \\
102 & 239.555478 & 239.446035 & $-0.1094$ & 602 & 941.157000 & 941.260920 & $+0.1039$ \\
103 & 241.049158 & 241.170269 & $+0.1211$ & 603 & 942.052342 & 942.515046 & $+0.4627$ \\
104 & 242.823272 & 242.891131 & $+0.0679$ & 604 & 944.188036 & 943.768840 & $-0.4192$ \\
105 & 244.070898 & 244.608663 & $+0.5378$ & 605 & 945.333563 & 945.022301 & $-0.3113$ \\
106 & 247.136990 & 246.322908 & $-0.8141$ & 606 & 946.765842 & 946.275431 & $-0.4904$ \\
107 & 248.101990 & 248.033908 & $-0.0681$ & 607 & 947.079183 & 947.528230 & $+0.4490$ \\
108 & 249.573690 & 249.741702 & $+0.1680$ & 608 & 948.346646 & 948.780699 & $+0.4341$ \\
109 & 251.014948 & 251.446331 & $+0.4314$ & 609 & 950.151613 & 950.032838 & $-0.1188$ \\
110 & 253.069987 & 253.147833 & $+0.0778$ & 610 & 951.033249 & 951.284648 & $+0.2514$ \\
111 & 255.306256 & 254.846246 & $-0.4600$ & 611 & 952.727989 & 952.536131 & $-0.1919$ \\
112 & 256.380714 & 256.541608 & $+0.1609$ & 612 & 954.129719 & 953.787286 & $-0.3424$ \\
113 & 258.610439 & 258.233956 & $-0.3765$ & 613 & 954.829309 & 955.038115 & $+0.2088$ \\
114 & 259.874407 & 259.923323 & $+0.0489$ & 614 & 956.675479 & 956.288617 & $-0.3869$ \\
115 & 260.805085 & 261.609746 & $+0.8047$ & 615 & 957.510053 & 957.538794 & $+0.0287$ \\
116 & 263.573894 & 263.293259 & $-0.2806$ & 616 & 958.414593 & 958.788647 & $+0.3741$ \\
117 & 265.557852 & 264.973895 & $-0.5840$ & 617 & 959.459169 & 960.038175 & $+0.5790$ \\
118 & 266.614974 & 266.651687 & $+0.0367$ & 618 & 961.669572 & 961.287380 & $-0.3822$ \\
119 & 267.921915 & 268.326667 & $+0.4048$ & 619 & 963.182087 & 962.536263 & $-0.6458$ \\
120 & 269.970449 & 269.998866 & $+0.0284$ & 620 & 963.567040 & 963.784824 & $+0.2178$ \\
121 & 271.494056 & 271.668316 & $+0.1743$ & 621 & 965.055580 & 965.033063 & $-0.0225$ \\
122 & 273.459609 & 273.335047 & $-0.1246$ & 622 & 966.110755 & 966.280982 & $+0.1702$ \\
123 & 275.587493 & 274.999088 & $-0.5884$ & 623 & 967.371154 & 967.528580 & $+0.1574$ \\
124 & 276.452050 & 276.660469 & $+0.2084$ & 624 & 968.636302 & 968.775860 & $+0.1396$ \\
125 & 278.250744 & 278.319218 & $+0.0685$ & 625 & 970.125611 & 970.022820 & $-0.1028$ \\
126 & 279.229251 & 279.975364 & $+0.7461$ & 626 & 971.071491 & 971.269463 & $+0.1980$ \\
127 & 282.465115 & 281.628933 & $-0.8362$ & 627 & 973.185361 & 972.515789 & $-0.6696$ \\
128 & 283.211186 & 283.279953 & $+0.0688$ & 628 & 973.873079 & 973.761797 & $-0.1113$ \\
129 & 284.835964 & 284.928450 & $+0.0925$ & 629 & 974.774635 & 975.007490 & $+0.2329$ \\
130 & 286.667445 & 286.574450 & $-0.0930$ & 630 & 976.178502 & 976.252867 & $+0.0744$ \\
131 & 287.911921 & 288.217980 & $+0.3061$ & 631 & 976.917202 & 977.497930 & $+0.5807$ \\
132 & 289.579855 & 289.859063 & $+0.2792$ & 632 & 978.766672 & 978.742678 & $-0.0240$ \\
133 & 291.846291 & 291.497725 & $-0.3486$ & 633 & 980.578001 & 979.987113 & $-0.5909$ \\
134 & 293.558434 & 293.133989 & $-0.4244$ & 634 & 981.288615 & 981.231235 & $-0.0574$ \\
135 & 294.965370 & 294.767881 & $-0.1975$ & 635 & 982.396485 & 982.475045 & $+0.0786$ \\
136 & 295.573255 & 296.399422 & $+0.8262$ & 636 & 983.575076 & 983.718543 & $+0.1435$ \\
137 & 297.979277 & 298.028636 & $+0.0494$ & 637 & 985.186929 & 984.961730 & $-0.2252$ \\
138 & 299.840326 & 299.655546 & $-0.1848$ & 638 & 986.130515 & 986.204607 & $+0.0741$ \\
139 & 301.649325 & 301.280173 & $-0.3692$ & 639 & 986.756008 & 987.447174 & $+0.6912$ \\
140 & 302.696750 & 302.902540 & $+0.2058$ & 640 & 988.992622 & 988.689432 & $-0.3032$ \\
141 & 304.864371 & 304.522668 & $-0.3417$ & 641 & 990.223918 & 989.931381 & $-0.2925$ \\
142 & 305.728913 & 306.140577 & $+0.4117$ & 642 & 991.374294 & 991.173023 & $-0.2013$ \\
143 & 307.219496 & 307.756290 & $+0.5368$ & 643 & 992.728696 & 992.414357 & $-0.3143$ \\
144 & 310.109463 & 309.369826 & $-0.7396$ & 644 & 993.214581 & 993.655385 & $+0.4408$ \\
145 & 311.165142 & 310.981204 & $-0.1839$ & 645 & 994.404591 & 994.896107 & $+0.4915$ \\
146 & 312.427801 & 312.590446 & $+0.1626$ & 646 & 996.205336 & 996.136523 & $-0.0688$ \\
147 & 313.985286 & 314.197570 & $+0.2123$ & 647 & 997.511935 & 997.376635 & $-0.1353$ \\
148 & 315.475616 & 315.802596 & $+0.3270$ & 648 & 998.827547 & 998.616442 & $-0.2111$ \\
149 & 317.734806 & 317.405542 & $-0.3293$ & 649 & 999.791572 & 999.855946 & $+0.0644$ \\
150 & 318.853104 & 319.006426 & $+0.1533$ & 650 & 1001.349483 & 1001.095146 & $-0.2543$ \\
151 & 321.160134 & 320.605268 & $-0.5549$ & 651 & 1002.404305 & 1002.334044 & $-0.0703$ \\
152 & 322.144559 & 322.202085 & $+0.0575$ & 652 & 1003.267808 & 1003.572641 & $+0.3048$ \\
153 & 323.466970 & 323.796894 & $+0.3299$ & 653 & 1004.675044 & 1004.810936 & $+0.1359$ \\
154 & 324.862866 & 325.389713 & $+0.5268$ & 654 & 1005.543420 & 1006.048930 & $+0.5055$ \\
155 & 327.443901 & 326.980559 & $-0.4633$ & 655 & 1008.006704 & 1007.286625 & $-0.7201$ \\
156 & 329.033072 & 328.569449 & $-0.4636$ & 656 & 1008.795710 & 1008.524020 & $-0.2717$ \\
157 & 329.953240 & 330.156400 & $+0.2032$ & 657 & 1009.806591 & 1009.761116 & $-0.0455$ \\
158 & 331.474468 & 331.741427 & $+0.2670$ & 658 & 1010.569757 & 1010.997914 & $+0.4282$ \\
159 & 333.645379 & 333.324548 & $-0.3208$ & 659 & 1012.410043 & 1012.234414 & $-0.1756$ \\
160 & 334.211355 & 334.905777 & $+0.6944$ & 660 & 1013.058638 & 1013.470617 & $+0.4120$ \\
161 & 336.841850 & 336.485131 & $-0.3567$ & 661 & 1014.689633 & 1014.706523 & $+0.0169$ \\
162 & 338.339993 & 338.062625 & $-0.2774$ & 662 & 1016.060179 & 1015.942133 & $-0.1180$ \\
163 & 339.858217 & 339.638274 & $-0.2199$ & 663 & 1017.266402 & 1017.177448 & $-0.0890$ \\
164 & 341.042261 & 341.212093 & $+0.1698$ & 664 & 1018.605573 & 1018.412469 & $-0.1931$ \\
165 & 342.054878 & 342.784097 & $+0.7292$ & 665 & 1019.912440 & 1019.647194 & $-0.2652$ \\
166 & 344.661703 & 344.354300 & $-0.3074$ & 666 & 1020.917475 & 1020.881626 & $-0.0358$ \\
167 & 346.347871 & 345.922717 & $-0.4252$ & 667 & 1021.544344 & 1022.115765 & $+0.5714$ \\
168 & 347.272678 & 347.489362 & $+0.2167$ & 668 & 1022.885271 & 1023.349612 & $+0.4643$ \\
169 & 349.316261 & 349.054249 & $-0.2620$ & 669 & 1025.265724 & 1024.583166 & $-0.6826$ \\
170 & 350.408419 & 350.617391 & $+0.2090$ & 670 & 1025.707944 & 1025.816429 & $+0.1085$ \\
171 & 351.878649 & 352.178803 & $+0.3002$ & 671 & 1027.467694 & 1027.049401 & $-0.4183$ \\
172 & 353.488900 & 353.738497 & $+0.2496$ & 672 & 1028.128964 & 1028.282083 & $+0.1531$ \\
173 & 356.017575 & 355.296486 & $-0.7211$ & 673 & 1029.227297 & 1029.514474 & $+0.2872$ \\
174 & 357.151302 & 356.852785 & $-0.2985$ & 674 & 1030.897369 & 1030.746577 & $-0.1508$ \\
175 & 357.952685 & 358.407405 & $+0.4547$ & 675 & 1031.833180 & 1031.978391 & $+0.1452$ \\
176 & 359.743755 & 359.960359 & $+0.2166$ & 676 & 1032.812883 & 1033.209916 & $+0.3970$ \\
177 & 361.289362 & 361.511660 & $+0.2223$ & 677 & 1034.612916 & 1034.441154 & $-0.1718$ \\
178 & 363.331331 & 363.061320 & $-0.2700$ & 678 & 1036.195917 & 1035.672105 & $-0.5238$ \\
179 & 364.736024 & 364.609351 & $-0.1267$ & 679 & 1037.024708 & 1036.902770 & $-0.1219$ \\
180 & 366.212710 & 366.155765 & $-0.0569$ & 680 & 1038.087752 & 1038.133148 & $+0.0454$ \\
181 & 367.993575 & 367.700575 & $-0.2930$ & 681 & 1039.077401 & 1039.363241 & $+0.2858$ \\
182 & 368.968438 & 369.243791 & $+0.2754$ & 682 & 1040.264038 & 1040.593049 & $+0.3290$ \\
183 & 370.050919 & 370.785426 & $+0.7345$ & 683 & 1041.621528 & 1041.822573 & $+0.2010$ \\
184 & 373.061928 & 372.325490 & $-0.7364$ & 684 & 1043.623954 & 1043.051813 & $-0.5721$ \\
185 & 373.864874 & 373.863995 & $-0.0009$ & 685 & 1044.514976 & 1044.280769 & $-0.2342$ \\
186 & 375.825913 & 375.400952 & $-0.4250$ & 686 & 1045.107042 & 1045.509443 & $+0.4024$ \\
187 & 376.324092 & 376.936373 & $+0.6123$ & 687 & 1047.089817 & 1046.737835 & $-0.3520$ \\
188 & 378.436680 & 378.470267 & $+0.0336$ & 688 & 1047.987147 & 1047.965945 & $-0.0212$ \\
189 & 379.872975 & 380.002645 & $+0.1297$ & 689 & 1048.953785 & 1049.193773 & $+0.2400$ \\
190 & 381.484469 & 381.533519 & $+0.0491$ & 690 & 1049.996284 & 1050.421321 & $+0.4250$ \\
191 & 383.443529 & 383.062899 & $-0.3806$ & 691 & 1051.576572 & 1051.648589 & $+0.0720$ \\
192 & 384.956117 & 384.590794 & $-0.3653$ & 692 & 1053.245785 & 1052.875577 & $-0.3702$ \\
193 & 385.861301 & 386.117216 & $+0.2559$ & 693 & 1054.781039 & 1054.102286 & $-0.6788$ \\
194 & 387.222890 & 387.642173 & $+0.4193$ & 694 & 1055.002146 & 1055.328716 & $+0.3266$ \\
195 & 388.846128 & 389.165677 & $+0.3195$ & 695 & 1056.688847 & 1056.554868 & $-0.1340$ \\
196 & 391.456084 & 390.687737 & $-0.7683$ & 696 & 1057.100044 & 1057.780743 & $+0.6807$ \\
197 & 392.245083 & 392.208362 & $-0.0367$ & 697 & 1059.133769 & 1059.006340 & $-0.1274$ \\
198 & 393.427744 & 393.727562 & $+0.2998$ & 698 & 1060.139519 & 1060.231661 & $+0.0921$ \\
199 & 395.582870 & 395.245347 & $-0.3375$ & 699 & 1061.501304 & 1061.456706 & $-0.0446$ \\
200 & 396.381854 & 396.761726 & $+0.3799$ & 700 & 1062.915382 & 1062.681475 & $-0.2339$ \\
201 & 397.918736 & 398.276708 & $+0.3580$ & 701 & 1064.071551 & 1063.905969 & $-0.1656$ \\
202 & 399.985120 & 399.790303 & $-0.1948$ & 702 & 1065.121855 & 1065.130189 & $+0.0083$ \\
203 & 401.839229 & 401.302519 & $-0.5367$ & 703 & 1066.463223 & 1066.354134 & $-0.1091$ \\
204 & 402.861918 & 402.813366 & $-0.0486$ & 704 & 1067.418860 & 1067.577806 & $+0.1589$ \\
205 & 404.236442 & 404.322851 & $+0.0864$ & 705 & 1067.990000 & 1068.801205 & $+0.8112$ \\
206 & 405.134387 & 405.830985 & $+0.6966$ & 706 & 1070.535042 & 1070.024332 & $-0.5107$ \\
207 & 407.581460 & 407.337775 & $-0.2437$ & 707 & 1071.618623 & 1071.247186 & $-0.3714$ \\
208 & 408.947246 & 408.843230 & $-0.1040$ & 708 & 1072.543998 & 1072.469769 & $-0.0742$ \\
209 & 410.513869 & 410.347359 & $-0.1665$ & 709 & 1073.570353 & 1073.692080 & $+0.1217$ \\
210 & 411.972268 & 411.850170 & $-0.1221$ & 710 & 1074.747771 & 1074.914121 & $+0.1664$ \\
211 & 413.262736 & 413.351671 & $+0.0889$ & 711 & 1076.266626 & 1076.135892 & $-0.1307$ \\
212 & 415.018810 & 414.851870 & $-0.1669$ & 712 & 1076.924056 & 1077.357394 & $+0.4333$ \\
213 & 415.455215 & 416.350775 & $+0.8956$ & 713 & 1078.647198 & 1078.578626 & $-0.0686$ \\
214 & 418.387706 & 417.848395 & $-0.5393$ & 714 & 1079.809965 & 1079.799589 & $-0.0104$ \\
215 & 419.861365 & 419.344737 & $-0.5166$ & 715 & 1081.171002 & 1081.020284 & $-0.1507$ \\
216 & 420.643828 & 420.839809 & $+0.1960$ & 716 & 1082.952750 & 1082.240712 & $-0.7120$ \\
217 & 422.076710 & 422.333619 & $+0.2569$ & 717 & 1083.295467 & 1083.460872 & $+0.1654$ \\
218 & 423.716580 & 423.826175 & $+0.1096$ & 718 & 1084.183264 & 1084.680766 & $+0.4975$ \\
219 & 425.069882 & 425.317483 & $+0.2476$ & 719 & 1085.647831 & 1085.900393 & $+0.2526$ \\
220 & 427.208825 & 426.807552 & $-0.4013$ & 720 & 1086.911999 & 1087.119754 & $+0.2078$ \\
221 & 428.127914 & 428.296388 & $+0.1685$ & 721 & 1088.755725 & 1088.338850 & $-0.4169$ \\
222 & 430.328745 & 429.784000 & $-0.5447$ & 722 & 1089.795338 & 1089.557682 & $-0.2377$ \\
223 & 431.301307 & 431.270394 & $-0.0309$ & 723 & 1090.863191 & 1090.776248 & $-0.0869$ \\
224 & 432.138642 & 432.755577 & $+0.6169$ & 724 & 1091.728473 & 1091.994551 & $+0.2661$ \\
225 & 433.889218 & 434.239558 & $+0.3503$ & 725 & 1093.440873 & 1093.212591 & $-0.2283$ \\
226 & 436.161006 & 435.722342 & $-0.4387$ & 726 & 1094.284458 & 1094.430367 & $+0.1459$ \\
227 & 437.581698 & 437.203936 & $-0.3778$ & 727 & 1095.433085 & 1095.647881 & $+0.2148$ \\
228 & 438.621739 & 438.684348 & $+0.0626$ & 728 & 1096.401918 & 1096.865133 & $+0.4632$ \\
229 & 439.918442 & 440.163585 & $+0.2451$ & 729 & 1098.841015 & 1098.082123 & $-0.7589$ \\
230 & 441.683199 & 441.641652 & $-0.0415$ & 730 & 1099.360667 & 1099.298852 & $-0.0618$ \\
231 & 442.904546 & 443.118557 & $+0.2140$ & 731 & 1100.574461 & 1100.515320 & $-0.0591$ \\
232 & 444.319336 & 444.594306 & $+0.2750$ & 732 & 1101.839111 & 1101.731529 & $-0.1076$ \\
233 & 446.860623 & 446.068907 & $-0.7917$ & 733 & 1102.551780 & 1102.947477 & $+0.3957$ \\
234 & 447.441704 & 447.542364 & $+0.1007$ & 734 & 1103.732297 & 1104.163166 & $+0.4309$ \\
235 & 449.148546 & 449.014685 & $-0.1339$ & 735 & 1105.617189 & 1105.378596 & $-0.2386$ \\
236 & 450.126946 & 450.485876 & $+0.3589$ & 736 & 1106.774372 & 1106.593768 & $-0.1806$ \\
237 & 451.403308 & 451.955944 & $+0.5526$ & 737 & 1107.774532 & 1107.808681 & $+0.0341$ \\
238 & 453.986738 & 453.424894 & $-0.5618$ & 738 & 1109.158919 & 1109.023337 & $-0.1356$ \\
239 & 454.974684 & 454.892733 & $-0.0820$ & 739 & 1110.444143 & 1110.237736 & $-0.2064$ \\
240 & 456.328427 & 456.359466 & $+0.0310$ & 740 & 1111.443505 & 1111.451879 & $+0.0084$ \\
241 & 457.903893 & 457.825100 & $-0.0788$ & 741 & 1112.432995 & 1112.665765 & $+0.2328$ \\
242 & 459.513415 & 459.289641 & $-0.2238$ & 742 & 1113.397595 & 1113.879395 & $+0.4818$ \\
243 & 460.087944 & 460.753095 & $+0.6652$ & 743 & 1115.065359 & 1115.092770 & $+0.0274$ \\
244 & 462.065367 & 462.215467 & $+0.1501$ & 744 & 1116.787254 & 1116.305890 & $-0.4814$ \\
245 & 464.057287 & 463.676764 & $-0.3805$ & 745 & 1117.965920 & 1117.518756 & $-0.4472$ \\
246 & 465.671539 & 465.136990 & $-0.5345$ & 746 & 1118.684135 & 1118.731367 & $+0.0472$ \\
247 & 466.570287 & 466.596153 & $+0.0259$ & 747 & 1119.473247 & 1119.943726 & $+0.4705$ \\
248 & 467.439046 & 468.054257 & $+0.6152$ & 748 & 1121.155938 & 1121.155830 & $-0.0001$ \\
249 & 469.536005 & 469.511307 & $-0.0247$ & 749 & 1122.458621 & 1122.367683 & $-0.0909$ \\
250 & 470.773655 & 470.967311 & $+0.1937$ & 750 & 1123.101117 & 1123.579282 & $+0.4782$ \\
251 & 472.799175 & 472.422272 & $-0.3769$ & 751 & 1125.314729 & 1124.790631 & $-0.5241$ \\
252 & 473.835232 & 473.876197 & $+0.0410$ & 752 & 1125.763442 & 1126.001727 & $+0.2383$ \\
253 & 475.600339 & 475.329091 & $-0.2712$ & 753 & 1127.658024 & 1127.212573 & $-0.4455$ \\
254 & 476.769015 & 476.780959 & $+0.0119$ & 754 & 1128.430225 & 1128.423168 & $-0.0071$ \\
255 & 478.075264 & 478.231807 & $+0.1565$ & 755 & 1129.728997 & 1129.633513 & $-0.0955$ \\
256 & 478.942182 & 479.681639 & $+0.7395$ & 756 & 1130.391598 & 1130.843608 & $+0.4520$ \\
257 & 481.830339 & 481.130461 & $-0.6999$ & 757 & 1131.495086 & 1132.053454 & $+0.5584$ \\
258 & 482.834783 & 482.578278 & $-0.2565$ & 758 & 1133.708626 & 1133.263051 & $-0.4456$ \\
259 & 483.851427 & 484.025095 & $+0.1737$ & 759 & 1134.885655 & 1134.472399 & $-0.4133$ \\
260 & 485.539148 & 485.470918 & $-0.0682$ & 760 & 1135.562214 & 1135.681499 & $+0.1193$ \\
261 & 486.528718 & 486.915750 & $+0.3870$ & 761 & 1136.929293 & 1136.890352 & $-0.0389$ \\
262 & 488.380567 & 488.359598 & $-0.0210$ & 762 & 1138.151590 & 1138.098958 & $-0.0526$ \\
263 & 489.661762 & 489.802466 & $+0.1407$ & 763 & 1138.992342 & 1139.307316 & $+0.3150$ \\
264 & 491.398822 & 491.244359 & $-0.1545$ & 764 & 1140.721848 & 1140.515429 & $-0.2064$ \\
265 & 493.314442 & 492.685281 & $-0.6292$ & 765 & 1141.261023 & 1141.723295 & $+0.4623$ \\
266 & 493.957998 & 494.125238 & $+0.1672$ & 766 & 1142.858660 & 1142.930916 & $+0.0723$ \\
267 & 495.358829 & 495.564234 & $+0.2054$ & 767 & 1144.782300 & 1144.138292 & $-0.6440$ \\
268 & 496.429696 & 497.002275 & $+0.5726$ & 768 & 1145.485328 & 1145.345422 & $-0.1399$ \\
269 & 498.580782 & 498.439363 & $-0.1414$ & 769 & 1146.576815 & 1146.552309 & $-0.0245$ \\
270 & 500.309085 & 499.875506 & $-0.4336$ & 770 & 1147.501777 & 1147.758951 & $+0.2572$ \\
271 & 501.604447 & 501.310706 & $-0.2937$ & 771 & 1148.615277 & 1148.965350 & $+0.3501$ \\
272 & 502.276270 & 502.744968 & $+0.4687$ & 772 & 1149.982601 & 1150.171506 & $+0.1889$ \\
273 & 504.499773 & 504.178297 & $-0.3215$ & 773 & 1151.562815 & 1151.377419 & $-0.1854$ \\
274 & 505.415232 & 505.610698 & $+0.1955$ & 774 & 1152.943129 & 1152.583090 & $-0.3600$ \\
275 & 506.464153 & 507.042174 & $+0.5780$ & 775 & 1153.890304 & 1153.788519 & $-0.1018$ \\
276 & 508.800700 & 508.472731 & $-0.3280$ & 776 & 1154.697520 & 1154.993706 & $+0.2962$ \\
277 & 510.264228 & 509.902372 & $-0.3619$ & 777 & 1156.621568 & 1156.198653 & $-0.4229$ \\
278 & 511.562290 & 511.331102 & $-0.2312$ & 778 & 1157.432315 & 1157.403358 & $-0.0290$ \\
279 & 512.623145 & 512.758925 & $+0.1358$ & 779 & 1158.001609 & 1158.607823 & $+0.6062$ \\
280 & 513.668986 & 514.185845 & $+0.5169$ & 780 & 1159.480657 & 1159.812048 & $+0.3314$ \\
281 & 515.435057 & 515.611867 & $+0.1768$ & 781 & 1161.396645 & 1161.016034 & $-0.3806$ \\
282 & 517.589669 & 517.036995 & $-0.5527$ & 782 & 1162.487529 & 1162.219780 & $-0.2677$ \\
283 & 518.234223 & 518.461232 & $+0.2270$ & 783 & 1163.701032 & 1163.423288 & $-0.2777$ \\
284 & 520.106310 & 519.884584 & $-0.2217$ & 784 & 1164.737586 & 1164.626557 & $-0.1110$ \\
285 & 521.525193 & 521.307053 & $-0.2181$ & 785 & 1165.271228 & 1165.829588 & $+0.5584$ \\
286 & 522.456696 & 522.728645 & $+0.2719$ & 786 & 1166.943613 & 1167.032382 & $+0.0888$ \\
287 & 523.960531 & 524.149363 & $+0.1888$ & 787 & 1168.086272 & 1168.234938 & $+0.1487$ \\
288 & 525.077386 & 525.569210 & $+0.4918$ & 788 & 1169.698357 & 1169.437258 & $-0.2611$ \\
289 & 527.903642 & 526.988192 & $-0.9154$ & 789 & 1170.463639 & 1170.639341 & $+0.1757$ \\
290 & 528.406214 & 528.406312 & $+0.0001$ & 790 & 1172.120682 & 1171.841188 & $-0.2795$ \\
291 & 529.806226 & 529.823573 & $+0.0173$ & 791 & 1173.305688 & 1173.042799 & $-0.2629$ \\
292 & 530.866918 & 531.239980 & $+0.3731$ & 792 & 1174.232767 & 1174.244175 & $+0.0114$ \\
293 & 532.688183 & 532.655537 & $-0.0326$ & 793 & 1175.215452 & 1175.445316 & $+0.2299$ \\
294 & 533.779631 & 534.070246 & $+0.2906$ & 794 & 1176.632876 & 1176.646222 & $+0.0133$ \\
295 & 535.664314 & 535.484113 & $-0.1802$ & 795 & 1177.106304 & 1177.846894 & $+0.7406$ \\
296 & 537.069759 & 536.897141 & $-0.1726$ & 796 & 1179.701224 & 1179.047333 & $-0.6539$ \\
297 & 538.428526 & 538.309332 & $-0.1192$ & 797 & 1180.653544 & 1180.247537 & $-0.4060$ \\
298 & 540.213166 & 539.720692 & $-0.4925$ & 798 & 1181.267318 & 1181.447509 & $+0.1802$ \\
299 & 540.631390 & 541.131224 & $+0.4998$ & 799 & 1182.582270 & 1182.647248 & $+0.0650$ \\
300 & 541.847437 & 542.540931 & $+0.6935$ & 800 & 1183.712775 & 1183.846755 & $+0.1340$ \\
301 & 544.323890 & 543.949817 & $-0.3741$ & 801 & 1185.155843 & 1185.046030 & $-0.1098$ \\
302 & 545.636833 & 545.357885 & $-0.2789$ & 802 & 1185.875359 & 1186.245073 & $+0.3697$ \\
303 & 547.010912 & 546.765139 & $-0.2458$ & 803 & 1187.345161 & 1187.443885 & $+0.0987$ \\
304 & 547.931613 & 548.171583 & $+0.2400$ & 804 & 1188.856444 & 1188.642467 & $-0.2140$ \\
305 & 549.497568 & 549.577220 & $+0.0797$ & 805 & 1189.963636 & 1189.840817 & $-0.1228$ \\
306 & 550.970010 & 550.982053 & $+0.0120$ & 806 & 1191.482606 & 1191.038938 & $-0.4437$ \\
307 & 552.049572 & 552.386086 & $+0.3365$ & 807 & 1192.218611 & 1192.236829 & $+0.0182$ \\
308 & 553.764972 & 553.789322 & $+0.0243$ & 808 & 1193.324021 & 1193.434490 & $+0.1105$ \\
309 & 555.792021 & 555.191765 & $-0.6003$ & 809 & 1193.857427 & 1194.631923 & $+0.7745$ \\
310 & 556.899476 & 556.593418 & $-0.3061$ & 810 & 1196.034672 & 1195.829126 & $-0.2055$ \\
311 & 557.564659 & 557.994283 & $+0.4296$ & 811 & 1197.071787 & 1197.026102 & $-0.0457$ \\
312 & 559.316237 & 559.394366 & $+0.0781$ & 812 & 1198.686569 & 1198.222849 & $-0.4637$ \\
313 & 560.240807 & 560.793668 & $+0.5529$ & 813 & 1199.356514 & 1199.419369 & $+0.0629$ \\
314 & 562.559208 & 562.192194 & $-0.3670$ & 814 & 1200.532692 & 1200.615662 & $+0.0830$ \\
315 & 564.160879 & 563.589945 & $-0.5709$ & 815 & 1201.810335 & 1201.811728 & $+0.0014$ \\
316 & 564.506056 & 564.986926 & $+0.4809$ & 816 & 1203.137351 & 1203.007567 & $-0.1298$ \\
317 & 566.698788 & 566.383140 & $-0.3156$ & 817 & 1203.855248 & 1204.203180 & $+0.3479$ \\
318 & 567.731758 & 567.778589 & $+0.0468$ & 818 & 1204.985492 & 1205.398567 & $+0.4131$ \\
319 & 568.923955 & 569.173278 & $+0.2493$ & 819 & 1206.870500 & 1206.593729 & $-0.2768$ \\
320 & 570.051115 & 570.567208 & $+0.5161$ & 820 & 1208.471460 & 1207.788665 & $-0.6828$ \\
321 & 572.419984 & 571.960384 & $-0.4596$ & 821 & 1208.989484 & 1208.983377 & $-0.0061$ \\
322 & 573.614611 & 573.352808 & $-0.2618$ & 822 & 1209.898030 & 1210.177865 & $+0.2798$ \\
323 & 575.093886 & 574.744483 & $-0.3494$ & 823 & 1211.416116 & 1211.372128 & $-0.0440$ \\
324 & 575.807247 & 576.135412 & $+0.3282$ & 824 & 1212.113153 & 1212.566168 & $+0.4530$ \\
325 & 577.039003 & 577.525598 & $+0.4866$ & 825 & 1213.598373 & 1213.759984 & $+0.1616$ \\
326 & 579.098835 & 578.915045 & $-0.1838$ & 826 & 1215.389975 & 1214.953577 & $-0.4364$ \\
327 & 580.136959 & 580.303754 & $+0.1668$ & 827 & 1216.183722 & 1216.146948 & $-0.0368$ \\
328 & 581.946576 & 581.691730 & $-0.2548$ & 828 & 1217.174482 & 1217.340096 & $+0.1656$ \\
329 & 583.236088 & 583.078975 & $-0.1571$ & 829 & 1219.050028 & 1218.533022 & $-0.5170$ \\
330 & 584.561706 & 584.465492 & $-0.0962$ & 830 & 1219.614471 & 1219.725727 & $+0.1113$ \\
331 & 585.984563 & 585.851283 & $-0.1333$ & 831 & 1220.816348 & 1220.918210 & $+0.1019$ \\
332 & 586.742772 & 587.236352 & $+0.4936$ & 832 & 1221.692242 & 1222.110472 & $+0.4182$ \\
333 & 588.139663 & 588.620701 & $+0.4810$ & 833 & 1222.952484 & 1223.302513 & $+0.3500$ \\
334 & 590.660398 & 590.004334 & $-0.6561$ & 834 & 1225.018330 & 1224.494335 & $-0.5240$ \\
335 & 591.725858 & 591.387252 & $-0.3386$ & 835 & 1225.855021 & 1225.685936 & $-0.1691$ \\
336 & 592.571358 & 592.769460 & $+0.1981$ & 836 & 1227.231828 & 1226.877317 & $-0.3545$ \\
337 & 593.974715 & 594.150959 & $+0.1762$ & 837 & 1227.917142 & 1228.068480 & $+0.1513$ \\
338 & 595.728154 & 595.531752 & $-0.1964$ & 838 & 1228.793154 & 1229.259423 & $+0.4663$ \\
339 & 596.362768 & 596.911842 & $+0.5491$ & 839 & 1230.584603 & 1230.450148 & $-0.1345$ \\
340 & 598.493077 & 598.291232 & $-0.2018$ & 840 & 1231.562274 & 1231.640654 & $+0.0784$ \\
341 & 599.545640 & 599.669924 & $+0.1243$ & 841 & 1232.529587 & 1232.830942 & $+0.3014$ \\
342 & 601.602137 & 601.047921 & $-0.5542$ & 842 & 1234.277817 & 1234.021013 & $-0.2568$ \\
343 & 602.579168 & 602.425226 & $-0.1539$ & 843 & 1235.502549 & 1235.210867 & $-0.2917$ \\
344 & 603.625619 & 603.801841 & $+0.1762$ & 844 & 1236.399017 & 1236.400503 & $+0.0015$ \\
345 & 604.616218 & 605.177769 & $+0.5616$ & 845 & 1237.977299 & 1237.589923 & $-0.3874$ \\
346 & 606.383460 & 606.553012 & $+0.1696$ & 846 & 1238.457233 & 1238.779126 & $+0.3219$ \\
347 & 608.413217 & 607.927574 & $-0.4856$ & 847 & 1239.490807 & 1239.968114 & $+0.4773$ \\
348 & 609.389575 & 609.301456 & $-0.0881$ & 848 & 1240.813472 & 1241.156885 & $+0.3434$ \\
349 & 610.839163 & 610.674661 & $-0.1645$ & 849 & 1243.078076 & 1242.345442 & $-0.7326$ \\
350 & 611.774210 & 612.047192 & $+0.2730$ & 850 & 1243.538147 & 1243.533783 & $-0.0044$ \\
351 & 613.599779 & 613.419051 & $-0.1807$ & 851 & 1244.851434 & 1244.721910 & $-0.1295$ \\
352 & 614.646238 & 614.790240 & $+0.1440$ & 852 & 1245.655866 & 1245.909822 & $+0.2540$ \\
353 & 615.538563 & 616.160763 & $+0.6222$ & 853 & 1247.372562 & 1247.097520 & $-0.2750$ \\
354 & 618.112831 & 617.530621 & $-0.5822$ & 854 & 1248.063061 & 1248.285004 & $+0.2219$ \\
355 & 619.184483 & 618.899817 & $-0.2847$ & 855 & 1249.159888 & 1249.472275 & $+0.3124$ \\
356 & 620.272894 & 620.268353 & $-0.0045$ & 856 & 1250.672397 & 1250.659333 & $-0.0131$ \\
357 & 621.709295 & 621.636233 & $-0.0731$ & 857 & 1251.659832 & 1251.846178 & $+0.1863$ \\
358 & 622.375003 & 623.003457 & $+0.6285$ & 858 & 1253.673578 & 1253.032811 & $-0.6408$ \\
359 & 624.269900 & 624.370029 & $+0.1001$ & 859 & 1254.431328 & 1254.219231 & $-0.2121$ \\
360 & 626.019283 & 625.735951 & $-0.2833$ & 860 & 1255.408231 & 1255.405439 & $-0.0028$ \\
361 & 627.268397 & 627.101225 & $-0.1672$ & 861 & 1256.181214 & 1256.591436 & $+0.4102$ \\
362 & 628.325862 & 628.465854 & $+0.1400$ & 862 & 1257.541219 & 1257.777222 & $+0.2360$ \\
363 & 630.473887 & 629.829840 & $-0.6440$ & 863 & 1258.779233 & 1258.962797 & $+0.1836$ \\
364 & 630.805781 & 631.193185 & $+0.3874$ & 864 & 1260.344548 & 1260.148161 & $-0.1964$ \\
365 & 632.225141 & 632.555891 & $+0.3308$ & 865 & 1261.611717 & 1261.333315 & $-0.2784$ \\
366 & 633.546858 & 633.917962 & $+0.3711$ & 866 & 1262.556614 & 1262.518259 & $-0.0384$ \\
367 & 635.523800 & 635.279398 & $-0.2444$ & 867 & 1263.676733 & 1263.702993 & $+0.0263$ \\
368 & 637.397193 & 636.640202 & $-0.7570$ & 868 & 1264.957223 & 1264.887518 & $-0.0697$ \\
369 & 637.925514 & 638.000377 & $+0.0749$ & 869 & 1266.179038 & 1266.071834 & $-0.1072$ \\
370 & 638.927938 & 639.359925 & $+0.4320$ & 870 & 1267.200346 & 1267.255941 & $+0.0556$ \\
371 & 640.694795 & 640.718847 & $+0.0241$ & 871 & 1267.570572 & 1268.439839 & $+0.8693$ \\
372 & 641.945500 & 642.077147 & $+0.1316$ & 872 & 1270.118922 & 1269.623530 & $-0.4954$ \\
373 & 643.278884 & 643.434826 & $+0.1559$ & 873 & 1271.134300 & 1270.807012 & $-0.3273$ \\
374 & 644.990578 & 644.791886 & $-0.1987$ & 874 & 1272.083960 & 1271.990287 & $-0.0937$ \\
375 & 646.348192 & 646.148329 & $-0.1999$ & 875 & 1273.261145 & 1273.173355 & $-0.0878$ \\
376 & 647.761753 & 647.504159 & $-0.2576$ & 876 & 1274.196221 & 1274.356216 & $+0.1600$ \\
377 & 648.786401 & 648.859376 & $+0.0730$ & 877 & 1275.092030 & 1275.538870 & $+0.4468$ \\
378 & 650.197519 & 650.213983 & $+0.0165$ & 878 & 1276.842172 & 1276.721318 & $-0.1209$ \\
379 & 650.668684 & 651.567982 & $+0.8993$ & 879 & 1277.763092 & 1277.903560 & $+0.1405$ \\
380 & 653.649572 & 652.921375 & $-0.7282$ & 880 & 1279.332843 & 1279.085596 & $-0.2472$ \\
381 & 654.301921 & 654.274165 & $-0.0278$ & 881 & 1280.155794 & 1280.267426 & $+0.1116$ \\
382 & 655.709463 & 655.626352 & $-0.0831$ & 882 & 1281.828727 & 1281.449052 & $-0.3797$ \\
383 & 656.964085 & 656.977940 & $+0.0139$ & 883 & 1283.000491 & 1282.630472 & $-0.3700$ \\
384 & 658.175614 & 658.328930 & $+0.1533$ & 884 & 1283.335032 & 1283.811688 & $+0.4767$ \\
385 & 659.663846 & 659.679324 & $+0.0155$ & 885 & 1284.854795 & 1284.992700 & $+0.1379$ \\
386 & 660.716733 & 661.029125 & $+0.3124$ & 886 & 1285.695023 & 1286.173508 & $+0.4785$ \\
387 & 662.296586 & 662.378333 & $+0.0817$ & 887 & 1287.410027 & 1287.354112 & $-0.0559$ \\
388 & 664.244605 & 663.726953 & $-0.5177$ & 888 & 1289.165352 & 1288.534513 & $-0.6308$ \\
389 & 665.342763 & 665.074984 & $-0.2678$ & 889 & 1290.104772 & 1289.714711 & $-0.3901$ \\
390 & 666.515148 & 666.422429 & $-0.0927$ & 890 & 1290.417708 & 1290.894705 & $+0.4770$ \\
391 & 667.148495 & 667.769291 & $+0.6208$ & 891 & 1291.945871 & 1292.074498 & $+0.1286$ \\
392 & 668.975849 & 669.115570 & $+0.1397$ & 892 & 1293.493982 & 1293.254088 & $-0.2399$ \\
393 & 670.323585 & 670.461270 & $+0.1377$ & 893 & 1294.118474 & 1294.433476 & $+0.3150$ \\
394 & 672.458184 & 671.806391 & $-0.6518$ & 894 & 1295.365364 & 1295.612663 & $+0.2473$ \\
395 & 673.043578 & 673.150937 & $+0.1074$ & 895 & 1296.801111 & 1296.791648 & $-0.0095$ \\
396 & 674.355898 & 674.494908 & $+0.1390$ & 896 & 1298.256527 & 1297.970432 & $-0.2861$ \\
397 & 676.139674 & 675.838306 & $-0.3014$ & 897 & 1299.405171 & 1299.149015 & $-0.2562$ \\
398 & 677.230181 & 677.181134 & $-0.0490$ & 898 & 1300.490019 & 1300.327398 & $-0.1626$ \\
399 & 677.800445 & 678.523393 & $+0.7229$ & 899 & 1301.495517 & 1301.505580 & $+0.0101$ \\
400 & 679.742198 & 679.865086 & $+0.1229$ & 900 & 1302.346742 & 1302.683563 & $+0.3368$ \\
401 & 681.894992 & 681.206213 & $-0.6888$ & 901 & 1303.273200 & 1303.861346 & $+0.5881$ \\
402 & 682.602735 & 682.546777 & $-0.0560$ & 902 & 1305.401672 & 1305.038930 & $-0.3627$ \\
403 & 684.013550 & 683.886780 & $-0.1268$ & 903 & 1306.508393 & 1306.216314 & $-0.2921$ \\
404 & 684.972630 & 685.226223 & $+0.2536$ & 904 & 1307.267242 & 1307.393500 & $+0.1263$ \\
405 & 686.163224 & 686.565108 & $+0.4019$ & 905 & 1308.988197 & 1308.570487 & $-0.4177$ \\
406 & 687.961543 & 687.903437 & $-0.0581$ & 906 & 1309.421532 & 1309.747276 & $+0.3257$ \\
407 & 689.368941 & 689.241212 & $-0.1277$ & 907 & 1311.056571 & 1310.923867 & $-0.1327$ \\
408 & 690.474735 & 690.578435 & $+0.1037$ & 908 & 1311.966941 & 1312.100260 & $+0.1333$ \\
409 & 692.451684 & 691.915106 & $-0.5366$ & 909 & 1313.031599 & 1313.276455 & $+0.2449$ \\
410 & 693.176970 & 693.251229 & $+0.0743$ & 910 & 1314.052566 & 1314.452454 & $+0.3999$ \\
411 & 694.533909 & 694.586804 & $+0.0529$ & 911 & 1316.212113 & 1315.628256 & $-0.5839$ \\
412 & 695.726336 & 695.921834 & $+0.1955$ & 912 & 1317.072986 & 1316.803861 & $-0.2691$ \\
413 & 696.626070 & 697.256320 & $+0.6303$ & 913 & 1318.171279 & 1317.979270 & $-0.1920$ \\
414 & 699.132095 & 698.590264 & $-0.5418$ & 914 & 1318.947881 & 1319.154483 & $+0.2066$ \\
415 & 700.296739 & 699.923668 & $-0.3731$ & 915 & 1319.931083 & 1320.329500 & $+0.3984$ \\
416 & 701.301743 & 701.256532 & $-0.0452$ & 916 & 1321.628139 & 1321.504321 & $-0.1238$ \\
417 & 702.227343 & 702.588860 & $+0.3615$ & 917 & 1322.258067 & 1322.678947 & $+0.4209$ \\
418 & 704.033839 & 703.920652 & $-0.1132$ & 918 & 1324.224979 & 1323.853379 & $-0.3716$ \\
419 & 705.125814 & 705.251911 & $+0.1261$ & 919 & 1325.237624 & 1325.027615 & $-0.2100$ \\
420 & 706.184655 & 706.582637 & $+0.3980$ & 920 & 1325.981970 & 1326.201657 & $+0.2197$ \\
421 & 708.269071 & 707.912833 & $-0.3562$ & 921 & 1327.635281 & 1327.375505 & $-0.2598$ \\
422 & 709.229589 & 709.242500 & $+0.0129$ & 922 & 1329.043518 & 1328.549159 & $-0.4944$ \\
423 & 711.130274 & 710.571640 & $-0.5586$ & 923 & 1329.205019 & 1329.722620 & $+0.5176$ \\
424 & 711.900290 & 711.900254 & $-0.0000$ & 924 & 1330.429937 & 1330.895887 & $+0.4659$ \\
425 & 712.749383 & 713.228344 & $+0.4790$ & 925 & 1331.827591 & 1332.068961 & $+0.2414$ \\
426 & 714.082772 & 714.555912 & $+0.4731$ & 926 & 1333.673523 & 1333.241842 & $-0.4317$ \\
427 & 716.112396 & 715.882959 & $-0.2294$ & 927 & 1334.747329 & 1334.414531 & $-0.3328$ \\
428 & 717.482570 & 717.209487 & $-0.2731$ & 928 & 1335.694975 & 1335.587027 & $-0.1079$ \\
429 & 718.742787 & 718.535497 & $-0.2073$ & 929 & 1336.690185 & 1336.759332 & $+0.0691$ \\
430 & 719.697101 & 719.860991 & $+0.1639$ & 930 & 1337.688792 & 1337.931444 & $+0.2427$ \\
431 & 721.351162 & 721.185970 & $-0.1652$ & 931 & 1338.923165 & 1339.103365 & $+0.1802$ \\
432 & 722.277505 & 722.510436 & $+0.2329$ & 932 & 1340.426400 & 1340.275095 & $-0.1513$ \\
433 & 723.845821 & 723.834391 & $-0.0114$ & 933 & 1341.166272 & 1341.446634 & $+0.2804$ \\
434 & 724.562614 & 725.157836 & $+0.5952$ & 934 & 1342.608508 & 1342.617982 & $+0.0095$ \\
435 & 727.056403 & 726.480772 & $-0.5756$ & 935 & 1344.156044 & 1343.789140 & $-0.3669$ \\
436 & 728.405482 & 727.803201 & $-0.6023$ & 936 & 1345.477106 & 1344.960107 & $-0.5170$ \\
437 & 728.758750 & 729.125125 & $+0.3664$ & 937 & 1345.731413 & 1346.130885 & $+0.3995$ \\
438 & 730.416482 & 730.446545 & $+0.0301$ & 938 & 1347.519472 & 1347.301472 & $-0.2180$ \\
439 & 731.417355 & 731.767463 & $+0.3501$ & 939 & 1348.017238 & 1348.471871 & $+0.4546$ \\
440 & 732.818053 & 733.087879 & $+0.2698$ & 940 & 1349.085194 & 1349.642080 & $+0.5569$ \\
441 & 734.789643 & 734.407796 & $-0.3818$ & 941 & 1351.296206 & 1350.812100 & $-0.4841$ \\
442 & 735.765459 & 735.727215 & $-0.0382$ & 942 & 1352.210465 & 1351.981932 & $-0.2285$ \\
443 & 737.052929 & 737.046137 & $-0.0068$ & 943 & 1353.483338 & 1353.151575 & $-0.3318$ \\
444 & 738.580421 & 738.364565 & $-0.2159$ & 944 & 1353.886782 & 1354.321030 & $+0.4342$ \\
445 & 739.909524 & 739.682498 & $-0.2270$ & 945 & 1355.680595 & 1355.490297 & $-0.1903$ \\
446 & 740.573807 & 740.999939 & $+0.4261$ & 946 & 1356.605656 & 1356.659376 & $+0.0537$ \\
447 & 741.757336 & 742.316890 & $+0.5596$ & 947 & 1357.771743 & 1357.828268 & $+0.0565$ \\
448 & 743.895013 & 743.633350 & $-0.2617$ & 948 & 1358.460160 & 1358.996973 & $+0.5368$ \\
449 & 745.344990 & 744.949323 & $-0.3957$ & 949 & 1360.393145 & 1360.165491 & $-0.2277$ \\
450 & 746.499306 & 746.264810 & $-0.2345$ & 950 & 1361.393075 & 1361.333822 & $-0.0593$ \\
451 & 747.674564 & 747.579811 & $-0.0948$ & 951 & 1363.022329 & 1362.501967 & $-0.5204$ \\
452 & 748.242754 & 748.894328 & $+0.6516$ & 952 & 1363.879191 & 1363.669926 & $-0.2093$ \\
453 & 750.655950 & 750.208362 & $-0.4476$ & 953 & 1364.576585 & 1364.837699 & $+0.2611$ \\
454 & 750.966381 & 751.521916 & $+0.5555$ & 954 & 1365.493734 & 1366.005286 & $+0.5116$ \\
455 & 752.887622 & 752.834990 & $-0.0526$ & 955 & 1367.104091 & 1367.172688 & $+0.0686$ \\
456 & 754.322370 & 754.147586 & $-0.1748$ & 956 & 1368.330193 & 1368.339905 & $+0.0097$ \\
457 & 755.839309 & 755.459704 & $-0.3796$ & 957 & 1369.686949 & 1369.506936 & $-0.1800$ \\
458 & 756.768248 & 756.771347 & $+0.0031$ & 958 & 1370.973523 & 1370.673783 & $-0.2997$ \\
459 & 758.101729 & 758.082516 & $-0.0192$ & 959 & 1371.686554 & 1371.840446 & $+0.1539$ \\
460 & 758.900238 & 759.393211 & $+0.4930$ & 960 & 1373.202915 & 1373.006924 & $-0.1960$ \\
461 & 760.282367 & 760.703435 & $+0.4211$ & 961 & 1374.154799 & 1374.173219 & $+0.0184$ \\
462 & 762.700033 & 762.013189 & $-0.6868$ & 962 & 1375.302392 & 1375.339330 & $+0.0369$ \\
463 & 763.593066 & 763.322473 & $-0.2706$ & 963 & 1376.161780 & 1376.505257 & $+0.3435$ \\
464 & 764.307523 & 764.631290 & $+0.3238$ & 964 & 1377.177634 & 1377.671001 & $+0.4934$ \\
465 & 766.087540 & 765.939641 & $-0.1479$ & 965 & 1379.683283 & 1378.836562 & $-0.8467$ \\
466 & 767.218472 & 767.247526 & $+0.0291$ & 966 & 1380.148578 & 1380.001941 & $-0.1466$ \\
467 & 768.281462 & 768.554948 & $+0.2735$ & 967 & 1381.073977 & 1381.167137 & $+0.0932$ \\
468 & 769.693407 & 769.861907 & $+0.1685$ & 968 & 1382.345663 & 1382.332150 & $-0.0135$ \\
469 & 771.070839 & 771.168404 & $+0.0976$ & 969 & 1383.297591 & 1383.496982 & $+0.1994$ \\
470 & 772.961618 & 772.474442 & $-0.4872$ & 970 & 1384.444416 & 1384.661632 & $+0.2172$ \\
471 & 774.117745 & 773.780021 & $-0.3377$ & 971 & 1385.663777 & 1385.826100 & $+0.1623$ \\
472 & 775.047847 & 775.085142 & $+0.0373$ & 972 & 1387.326648 & 1386.990387 & $-0.3363$ \\
473 & 775.999712 & 776.389808 & $+0.3901$ & 973 & 1387.921454 & 1388.154493 & $+0.2330$ \\
474 & 777.299749 & 777.694018 & $+0.3943$ & 974 & 1389.565832 & 1389.318419 & $-0.2474$ \\
475 & 779.157077 & 778.997774 & $-0.1593$ & 975 & 1390.705490 & 1390.482163 & $-0.2233$ \\
476 & 780.348925 & 780.301078 & $-0.0478$ & 976 & 1391.853200 & 1391.645727 & $-0.2075$ \\
477 & 782.137664 & 781.603931 & $-0.5337$ & 977 & 1392.644028 & 1392.809112 & $+0.1651$ \\
478 & 782.597944 & 782.906333 & $+0.3084$ & 978 & 1393.433402 & 1393.972316 & $+0.5389$ \\
479 & 784.288823 & 784.208287 & $-0.0805$ & 979 & 1394.884185 & 1395.135340 & $+0.2512$ \\
480 & 785.739090 & 785.509793 & $-0.2293$ & 980 & 1396.544163 & 1396.298186 & $-0.2460$ \\
481 & 786.461147 & 786.810852 & $+0.3497$ & 981 & 1397.834623 & 1397.460852 & $-0.3738$ \\
482 & 787.468464 & 788.111467 & $+0.6430$ & 982 & 1398.837675 & 1398.623339 & $-0.2143$ \\
483 & 790.059092 & 789.411637 & $-0.6475$ & 983 & 1399.839473 & 1399.785647 & $-0.0538$ \\
484 & 790.831620 & 790.711364 & $-0.1203$ & 984 & 1400.426946 & 1400.947777 & $+0.5208$ \\
485 & 792.427708 & 792.010650 & $-0.4171$ & 985 & 1402.564347 & 1402.109728 & $-0.4546$ \\
486 & 792.888653 & 793.309495 & $+0.4208$ & 986 & 1402.973748 & 1403.271502 & $+0.2978$ \\
487 & 794.483792 & 794.607900 & $+0.1241$ & 987 & 1404.006292 & 1404.433097 & $+0.4268$ \\
488 & 795.606596 & 795.905868 & $+0.2993$ & 988 & 1405.666975 & 1405.594516 & $-0.0725$ \\
489 & 797.263470 & 797.203398 & $-0.0601$ & 989 & 1407.085143 & 1406.755756 & $-0.3294$ \\
490 & 798.707570 & 798.500493 & $-0.2071$ & 990 & 1408.136307 & 1407.916820 & $-0.2195$ \\
491 & 799.654336 & 799.797152 & $+0.1428$ & 991 & 1409.320681 & 1409.077707 & $-0.2430$ \\
492 & 801.604246 & 801.093379 & $-0.5109$ & 992 & 1410.024811 & 1410.238417 & $+0.2136$ \\
493 & 802.541985 & 802.389172 & $-0.1528$ & 993 & 1411.257057 & 1411.398950 & $+0.1419$ \\
494 & 803.243096 & 803.684534 & $+0.4414$ & 994 & 1411.965653 & 1412.559308 & $+0.5937$ \\
495 & 804.762239 & 804.979467 & $+0.2172$ & 995 & 1413.843149 & 1413.719489 & $-0.1237$ \\
496 & 805.861636 & 806.273970 & $+0.4123$ & 996 & 1415.585785 & 1414.879495 & $-0.7063$ \\
497 & 808.151815 & 807.568045 & $-0.5838$ & 997 & 1415.781581 & 1416.039325 & $+0.2577$ \\
498 & 809.197783 & 808.861693 & $-0.3361$ & 998 & 1417.102823 & 1417.198980 & $+0.0962$ \\
499 & 810.081805 & 810.154915 & $+0.0731$ & 999 & 1418.696964 & 1418.358460 & $-0.3385$ \\
500 & 811.184359 & 811.447713 & $+0.2634$ & 1000 & 1419.422481 & 1419.517765 & $+0.0953$ \\
\end{longtable}}

{\scriptsize\setlength{\tabcolsep}{3pt}
\begin{longtable}{rrrr@{\quad}rrrr@{\quad}rrrr}
\caption{The primes $p_j$ and the RVM prime levels $p_j^{\RVM}$, $j\le1000$.  Here $p_j^{\RVM}$ is a discrete version of $\inf\{x:\pi_{\RVM}(x)\ge j\}$: the first point of the computational grid in $u=\log x$ (step $5.9\cdot10^{-6}$) at which $\pi_{\RVM}\ge j$, refined by linear interpolation from the previous grid point.  It is not the inverse of a monotone function, and its accuracy is limited by the grid, the truncation of the ordinate sum and the smooth cut.  The same levels computed by program \texttt{s04} with $K=17\,061$ ordinates on $[2,1000]$ ($40\,000$ grid points), instead of $K=183\,029$ ordinates on $[2,8200]$, differ from these by a median of $1.7\cdot10^{-3}$ over the $168$ levels $j\le168$, by more than $0.01$ at $29$ of them, and by $0.85$, $0.85$, $0.79$ at $j=148$, $118$, $106$: at these three levels $\pi_{\RVM}$ has, before the crossing, a local maximum within $0.006$ of $j$, which one truncation lifts above $j$ and the other does not.}\label{tab:primes-1000}\\
\toprule
$j$ & $p_j$ & $p_j^{\RVM}$ & $p_j^{\RVM}-p_j$ & $j$ & $p_j$ & $p_j^{\RVM}$ & $p_j^{\RVM}-p_j$ & $j$ & $p_j$ & $p_j^{\RVM}$ & $p_j^{\RVM}-p_j$\\
\midrule
\endfirsthead
\multicolumn{12}{l}{\emph{Table~\thetable\ (continued)}}\\
\toprule
$j$ & $p_j$ & $p_j^{\RVM}$ & $p_j^{\RVM}-p_j$ & $j$ & $p_j$ & $p_j^{\RVM}$ & $p_j^{\RVM}-p_j$ & $j$ & $p_j$ & $p_j^{\RVM}$ & $p_j^{\RVM}-p_j$\\
\midrule
\endhead
\bottomrule
\endfoot
1 & 2 & 2.2138 & $+0.2138$ & 335 & 2251 & 2240.2009 & $-10.7991$ & 669 & 4999 & 4995.1090 & $-3.8910$ \\
2 & 3 & 3.1546 & $+0.1546$ & 336 & 2267 & 2264.1226 & $-2.8774$ & 670 & 5003 & 5000.1908 & $-2.8092$ \\
3 & 5 & 5.1880 & $+0.1880$ & 337 & 2269 & 2270.0868 & $+1.0868$ & 671 & 5009 & 5006.2128 & $-2.7872$ \\
4 & 7 & 7.1588 & $+0.1588$ & 338 & 2273 & 2274.1702 & $+1.1702$ & 672 & 5011 & 5014.2390 & $+3.2390$ \\
5 & 11 & 11.1391 & $+0.1391$ & 339 & 2281 & 2279.1152 & $-1.8848$ & 673 & 5021 & 5043.1625 & $+22.1625$ \\
6 & 13 & 14.0201 & $+1.0201$ & 340 & 2287 & 2284.2038 & $-2.7962$ & 674 & 5023 & 5054.1243 & $+31.1243$ \\
7 & 17 & 17.2551 & $+0.2551$ & 341 & 2293 & 2297.2222 & $+4.2222$ & 675 & 5039 & 5060.1530 & $+21.1530$ \\
8 & 19 & 19.1767 & $+0.1767$ & 342 & 2297 & 2311.1364 & $+14.1364$ & 676 & 5051 & 5065.2030 & $+14.2030$ \\
9 & 23 & 26.0524 & $+3.0524$ & 343 & 2309 & 2316.1611 & $+7.1611$ & 677 & 5059 & 5070.1949 & $+11.1949$ \\
10 & 29 & 28.0419 & $-0.9581$ & 344 & 2311 & 2320.1971 & $+9.1971$ & 678 & 5077 & 5075.2101 & $-1.7899$ \\
11 & 31 & 33.0665 & $+2.0665$ & 345 & 2333 & 2324.1726 & $-8.8274$ & 679 & 5081 & 5081.1521 & $+0.1521$ \\
12 & 37 & 38.0878 & $+1.0878$ & 346 & 2339 & 2328.1884 & $-10.8116$ & 680 & 5087 & 5088.1848 & $+1.1848$ \\
13 & 41 & 40.2252 & $-0.7748$ & 347 & 2341 & 2333.1621 & $-7.8379$ & 681 & 5099 & 5100.1882 & $+1.1882$ \\
14 & 43 & 44.1857 & $+1.1857$ & 348 & 2347 & 2341.1412 & $-5.8588$ & 682 & 5101 & 5114.1940 & $+13.1940$ \\
15 & 47 & 46.1156 & $-0.8844$ & 349 & 2351 & 2354.1374 & $+3.1374$ & 683 & 5107 & 5122.1794 & $+15.1794$ \\
16 & 53 & 57.1721 & $+4.1721$ & 350 & 2357 & 2361.1040 & $+4.1040$ & 684 & 5113 & 5128.1328 & $+15.1328$ \\
17 & 59 & 59.1154 & $+0.1154$ & 351 & 2371 & 2365.2375 & $-5.7625$ & 685 & 5119 & 5133.1056 & $+14.1056$ \\
18 & 61 & 61.0717 & $+0.0717$ & 352 & 2377 & 2370.0802 & $-6.9198$ & 686 & 5147 & 5137.1713 & $-9.8287$ \\
19 & 67 & 67.0344 & $+0.0344$ & 353 & 2381 & 2374.1044 & $-6.8956$ & 687 & 5153 & 5142.0815 & $-10.9185$ \\
20 & 71 & 74.0901 & $+3.0901$ & 354 & 2383 & 2379.0944 & $-3.9056$ & 688 & 5167 & 5146.1533 & $-20.8467$ \\
21 & 73 & 76.1129 & $+3.1129$ & 355 & 2389 & 2386.1853 & $-2.8147$ & 689 & 5171 & 5151.1596 & $-19.8404$ \\
22 & 79 & 83.0777 & $+4.0777$ & 356 & 2393 & 2409.1394 & $+16.1394$ & 690 & 5179 & 5157.1915 & $-21.8085$ \\
23 & 83 & 86.0388 & $+3.0388$ & 357 & 2399 & 2415.1474 & $+16.1474$ & 691 & 5189 & 5167.1708 & $-21.8292$ \\
24 & 89 & 92.0775 & $+3.0775$ & 358 & 2411 & 2420.1302 & $+9.1302$ & 692 & 5197 & 5185.1602 & $-11.8398$ \\
25 & 97 & 94.0745 & $-2.9255$ & 359 & 2417 & 2425.1210 & $+8.1210$ & 693 & 5209 & 5195.1193 & $-13.8807$ \\
26 & 101 & 96.1261 & $-4.8739$ & 360 & 2423 & 2431.1450 & $+8.1450$ & 694 & 5227 & 5201.1508 & $-25.8492$ \\
27 & 103 & 103.2182 & $+0.2182$ & 361 & 2437 & 2447.1760 & $+10.1760$ & 695 & 5231 & 5206.2067 & $-24.7933$ \\
28 & 107 & 112.1163 & $+5.1163$ & 362 & 2441 & 2458.1209 & $+17.1209$ & 696 & 5233 & 5211.1935 & $-21.8065$ \\
29 & 109 & 115.0244 & $+6.0244$ & 363 & 2447 & 2463.1388 & $+16.1388$ & 697 & 5237 & 5216.1890 & $-20.8110$ \\
30 & 113 & 118.1253 & $+5.1253$ & 364 & 2459 & 2467.1877 & $+8.1877$ & 698 & 5261 & 5222.1181 & $-38.8819$ \\
31 & 127 & 125.1594 & $-1.8406$ & 365 & 2467 & 2471.2079 & $+4.2079$ & 699 & 5273 & 5228.2500 & $-44.7500$ \\
32 & 131 & 136.0371 & $+5.0371$ & 366 & 2473 & 2476.1039 & $+3.1039$ & 700 & 5279 & 5240.2593 & $-38.7407$ \\
33 & 137 & 138.1768 & $+1.1768$ & 367 & 2477 & 2481.1801 & $+4.1801$ & 701 & 5281 & 5262.1653 & $-18.8347$ \\
34 & 139 & 145.1787 & $+6.1787$ & 368 & 2503 & 2493.1792 & $-9.8208$ & 702 & 5297 & 5271.1112 & $-25.8888$ \\
35 & 149 & 148.0858 & $-0.9142$ & 369 & 2521 & 2506.1738 & $-14.8262$ & 703 & 5303 & 5277.1190 & $-25.8810$ \\
36 & 151 & 150.1093 & $-0.8907$ & 370 & 2531 & 2512.1287 & $-18.8713$ & 704 & 5309 & 5282.1773 & $-26.8227$ \\
37 & 157 & 154.1579 & $-2.8421$ & 371 & 2539 & 2516.2206 & $-22.7794$ & 705 & 5323 & 5288.0995 & $-34.9005$ \\
38 & 163 & 160.1227 & $-2.8773$ & 372 & 2543 & 2521.0787 & $-21.9213$ & 706 & 5333 & 5293.2319 & $-39.7681$ \\
39 & 167 & 163.1171 & $-3.8829$ & 373 & 2549 & 2525.1356 & $-23.8644$ & 707 & 5347 & 5301.1624 & $-45.8376$ \\
40 & 173 & 173.1152 & $+0.1152$ & 374 & 2551 & 2530.2229 & $-20.7771$ & 708 & 5351 & 5323.2036 & $-27.7964$ \\
41 & 179 & 176.0959 & $-2.9041$ & 375 & 2557 & 2543.1879 & $-13.8121$ & 709 & 5381 & 5340.1379 & $-40.8621$ \\
42 & 181 & 186.1175 & $+5.1175$ & 376 & 2579 & 2559.1228 & $-19.8772$ & 710 & 5387 & 5347.1295 & $-39.8705$ \\
43 & 191 & 190.1285 & $-0.8715$ & 377 & 2591 & 2565.1248 & $-25.8752$ & 711 & 5393 & 5352.2287 & $-40.7713$ \\
44 & 193 & 201.0969 & $+8.0969$ & 378 & 2593 & 2570.1536 & $-22.8464$ & 712 & 5399 & 5358.0905 & $-40.9095$ \\
45 & 197 & 203.1883 & $+6.1883$ & 379 & 2609 & 2576.1280 & $-32.8720$ & 713 & 5407 & 5363.1069 & $-43.8931$ \\
46 & 199 & 207.1015 & $+8.1015$ & 380 & 2617 & 2610.1950 & $-6.8050$ & 714 & 5413 & 5368.1864 & $-44.8136$ \\
47 & 211 & 212.1460 & $+1.1460$ & 381 & 2621 & 2618.1134 & $-2.8866$ & 715 & 5417 & 5375.1711 & $-41.8289$ \\
48 & 223 & 215.0763 & $-7.9237$ & 382 & 2633 & 2624.0998 & $-8.9002$ & 716 & 5419 & 5386.1770 & $-32.8230$ \\
49 & 227 & 217.2057 & $-9.7943$ & 383 & 2647 & 2630.1946 & $-16.8054$ & 717 & 5431 & 5403.1571 & $-27.8429$ \\
50 & 229 & 230.1447 & $+1.1447$ & 384 & 2657 & 2654.1582 & $-2.8418$ & 718 & 5437 & 5412.1531 & $-24.8469$ \\
51 & 233 & 234.1155 & $+1.1155$ & 385 & 2659 & 2663.0956 & $+4.0956$ & 719 & 5441 & 5418.1734 & $-22.8266$ \\
52 & 239 & 243.0719 & $+4.0719$ & 386 & 2663 & 2668.0972 & $+5.0972$ & 720 & 5443 & 5423.2029 & $-19.7971$ \\
53 & 241 & 246.0130 & $+5.0130$ & 387 & 2671 & 2672.1596 & $+1.1596$ & 721 & 5449 & 5428.1506 & $-20.8494$ \\
54 & 251 & 248.1162 & $-2.8838$ & 388 & 2677 & 2677.0942 & $+0.0942$ & 722 & 5471 & 5433.1031 & $-37.8969$ \\
55 & 257 & 258.2462 & $+1.2462$ & 389 & 2683 & 2682.1362 & $-0.8638$ & 723 & 5477 & 5438.0936 & $-38.9064$ \\
56 & 263 & 263.0536 & $+0.0536$ & 390 & 2687 & 2691.1547 & $+4.1547$ & 724 & 5479 & 5443.1613 & $-35.8387$ \\
57 & 269 & 267.1805 & $-1.8195$ & 391 & 2689 & 2709.1416 & $+20.1416$ & 725 & 5483 & 5450.2001 & $-32.7999$ \\
58 & 271 & 277.0384 & $+6.0384$ & 392 & 2693 & 2715.2375 & $+22.2375$ & 726 & 5501 & 5465.2375 & $-35.7625$ \\
59 & 277 & 279.1633 & $+2.1633$ & 393 & 2699 & 2720.1824 & $+21.1824$ & 727 & 5503 & 5484.1660 & $-18.8340$ \\
60 & 281 & 282.2079 & $+1.2079$ & 394 & 2707 & 2725.0835 & $+18.0835$ & 728 & 5507 & 5492.1544 & $-14.8456$ \\
61 & 283 & 293.1444 & $+10.1444$ & 395 & 2711 & 2729.1455 & $+18.1455$ & 729 & 5519 & 5498.1613 & $-20.8387$ \\
62 & 293 & 297.0717 & $+4.0717$ & 396 & 2713 & 2734.1897 & $+21.1897$ & 730 & 5521 & 5504.1003 & $-16.8997$ \\
63 & 307 & 311.0820 & $+4.0820$ & 397 & 2719 & 2743.1740 & $+24.1740$ & 731 & 5527 & 5509.1649 & $-17.8351$ \\
64 & 311 & 314.2060 & $+3.2060$ & 398 & 2729 & 2761.1677 & $+32.1677$ & 732 & 5531 & 5515.1948 & $-15.8052$ \\
65 & 313 & 326.1105 & $+13.1105$ & 399 & 2731 & 2768.1501 & $+37.1501$ & 733 & 5557 & 5524.1509 & $-32.8491$ \\
66 & 317 & 329.2472 & $+12.2472$ & 400 & 2741 & 2773.1571 & $+32.1571$ & 734 & 5563 & 5554.2170 & $-8.7830$ \\
67 & 331 & 332.1141 & $+1.1141$ & 401 & 2749 & 2778.1149 & $+29.1149$ & 735 & 5569 & 5566.1453 & $-2.8547$ \\
68 & 337 & 335.1199 & $-1.8801$ & 402 & 2753 & 2783.1305 & $+30.1305$ & 736 & 5573 & 5573.1534 & $+0.1534$ \\
69 & 347 & 343.1291 & $-3.8709$ & 403 & 2767 & 2790.1587 & $+23.1587$ & 737 & 5581 & 5579.2072 & $-1.7928$ \\
70 & 349 & 347.2034 & $-1.7966$ & 404 & 2777 & 2814.2147 & $+37.2147$ & 738 & 5591 & 5586.1546 & $-4.8454$ \\
71 & 353 & 351.0568 & $-1.9432$ & 405 & 2789 & 2822.0996 & $+33.0996$ & 739 & 5623 & 5595.1431 & $-27.8569$ \\
72 & 359 & 365.1411 & $+6.1411$ & 406 & 2791 & 2827.0857 & $+36.0857$ & 740 & 5639 & 5636.1566 & $-2.8434$ \\
73 & 367 & 369.0667 & $+2.0667$ & 407 & 2797 & 2831.1482 & $+34.1482$ & 741 & 5641 & 5645.1130 & $+4.1130$ \\
74 & 373 & 372.1471 & $-0.8529$ & 408 & 2801 & 2836.1031 & $+35.1031$ & 742 & 5647 & 5651.1535 & $+4.1535$ \\
75 & 379 & 380.2049 & $+1.2049$ & 409 & 2803 & 2841.1735 & $+38.1735$ & 743 & 5651 & 5657.1496 & $+6.1496$ \\
76 & 383 & 385.0915 & $+2.0915$ & 410 & 2819 & 2850.1575 & $+31.1575$ & 744 & 5653 & 5663.1924 & $+10.1924$ \\
77 & 389 & 388.0514 & $-0.9486$ & 411 & 2833 & 2863.1914 & $+30.1914$ & 745 & 5657 & 5671.1662 & $+14.1662$ \\
78 & 397 & 391.0724 & $-5.9276$ & 412 & 2837 & 2870.1900 & $+33.1900$ & 746 & 5659 & 5686.1916 & $+27.1916$ \\
79 & 401 & 404.1579 & $+3.1579$ & 413 & 2843 & 2875.1506 & $+32.1506$ & 747 & 5669 & 5706.1647 & $+37.1647$ \\
80 & 409 & 408.1667 & $-0.8333$ & 414 & 2851 & 2879.1584 & $+28.1584$ & 748 & 5683 & 5714.1603 & $+31.1603$ \\
81 & 419 & 413.1228 & $-5.8772$ & 415 & 2857 & 2883.1256 & $+26.1256$ & 749 & 5689 & 5720.1408 & $+31.1408$ \\
82 & 421 & 424.1253 & $+3.1253$ & 416 & 2861 & 2887.1097 & $+26.1097$ & 750 & 5693 & 5725.1545 & $+32.1545$ \\
83 & 431 & 427.1525 & $-3.8475$ & 417 & 2879 & 2891.1585 & $+12.1585$ & 751 & 5701 & 5730.1327 & $+29.1327$ \\
84 & 433 & 430.1187 & $-2.8813$ & 418 & 2887 & 2897.1097 & $+10.1097$ & 752 & 5711 & 5735.1272 & $+24.1272$ \\
85 & 439 & 434.1375 & $-4.8625$ & 419 & 2897 & 2906.2458 & $+9.2458$ & 753 & 5717 & 5740.1864 & $+23.1864$ \\
86 & 443 & 446.0793 & $+3.0793$ & 420 & 2903 & 2921.1375 & $+18.1375$ & 754 & 5737 & 5747.1411 & $+10.1411$ \\
87 & 449 & 449.2083 & $+0.2083$ & 421 & 2909 & 2927.2070 & $+18.2070$ & 755 & 5741 & 5757.1855 & $+16.1855$ \\
88 & 457 & 454.1441 & $-2.8559$ & 422 & 2917 & 2932.1846 & $+15.1846$ & 756 & 5743 & 5777.1652 & $+34.1652$ \\
89 & 461 & 470.1710 & $+9.1710$ & 423 & 2927 & 2937.0967 & $+10.0967$ & 757 & 5749 & 5787.1601 & $+38.1601$ \\
90 & 463 & 475.1873 & $+12.1873$ & 424 & 2939 & 2941.1732 & $+2.1732$ & 758 & 5779 & 5794.1041 & $+15.1041$ \\
91 & 467 & 490.0909 & $+23.0909$ & 425 & 2953 & 2946.2388 & $-6.7612$ & 759 & 5783 & 5799.1482 & $+16.1482$ \\
92 & 479 & 493.0870 & $+14.0870$ & 426 & 2957 & 2955.1707 & $-1.8293$ & 760 & 5791 & 5804.1472 & $+13.1472$ \\
93 & 487 & 496.0928 & $+9.0928$ & 427 & 2963 & 2977.1315 & $+14.1315$ & 761 & 5801 & 5809.1428 & $+8.1428$ \\
94 & 491 & 500.1660 & $+9.1660$ & 428 & 2969 & 2984.1305 & $+15.1305$ & 762 & 5807 & 5814.1876 & $+7.1876$ \\
95 & 499 & 508.1736 & $+9.1736$ & 429 & 2971 & 2989.1881 & $+18.1881$ & 763 & 5813 & 5820.2221 & $+7.2221$ \\
96 & 503 & 512.1244 & $+9.1244$ & 430 & 2999 & 2995.0915 & $-3.9085$ & 764 & 5821 & 5829.1735 & $+8.1735$ \\
97 & 509 & 515.1067 & $+6.1067$ & 431 & 3001 & 3001.1309 & $+0.1309$ & 765 & 5827 & 5848.1921 & $+21.1921$ \\
98 & 521 & 518.1442 & $-2.8558$ & 432 & 3011 & 3028.2015 & $+17.2015$ & 766 & 5839 & 5861.2576 & $+22.2576$ \\
99 & 523 & 531.1694 & $+8.1694$ & 433 & 3019 & 3038.0972 & $+19.0972$ & 767 & 5843 & 5869.1412 & $+26.1412$ \\
100 & 541 & 536.1441 & $-4.8559$ & 434 & 3023 & 3043.1355 & $+20.1355$ & 768 & 5849 & 5875.1237 & $+26.1237$ \\
101 & 547 & 539.2398 & $-7.7602$ & 435 & 3037 & 3048.1091 & $+11.1091$ & 769 & 5851 & 5880.1722 & $+29.1722$ \\
102 & 557 & 543.1977 & $-13.8023$ & 436 & 3041 & 3053.1064 & $+12.1064$ & 770 & 5857 & 5886.0965 & $+29.0965$ \\
103 & 563 & 553.1345 & $-9.8655$ & 437 & 3049 & 3059.1164 & $+10.1164$ & 771 & 5861 & 5891.2185 & $+30.2185$ \\
104 & 569 & 558.0439 & $-10.9561$ & 438 & 3061 & 3069.1605 & $+8.1605$ & 772 & 5867 & 5899.1423 & $+32.1423$ \\
105 & 571 & 561.0337 & $-9.9663$ & 439 & 3067 & 3083.1748 & $+16.1748$ & 773 & 5869 & 5915.1969 & $+46.1969$ \\
106 & 577 & 563.2463 & $-13.7537$ & 440 & 3079 & 3090.1288 & $+11.1288$ & 774 & 5879 & 5939.1321 & $+60.1321$ \\
107 & 587 & 568.1124 & $-18.8876$ & 441 & 3083 & 3095.0947 & $+12.0947$ & 775 & 5881 & 5946.2450 & $+65.2450$ \\
108 & 593 & 583.0774 & $-9.9226$ & 442 & 3089 & 3099.1119 & $+10.1119$ & 776 & 5897 & 5952.2356 & $+55.2356$ \\
109 & 599 & 587.1276 & $-11.8724$ & 443 & 3109 & 3103.1107 & $-5.8893$ & 777 & 5903 & 5958.1194 & $+55.1194$ \\
110 & 601 & 605.1479 & $+4.1479$ & 444 & 3119 & 3107.1361 & $-11.8639$ & 778 & 5923 & 5963.1517 & $+40.1517$ \\
111 & 607 & 609.2115 & $+2.2115$ & 445 & 3121 & 3112.1162 & $-8.8838$ & 779 & 5927 & 5969.1067 & $+42.1067$ \\
112 & 613 & 613.2031 & $+0.2031$ & 446 & 3137 & 3118.2117 & $-18.7883$ & 780 & 5939 & 5975.1703 & $+36.1703$ \\
113 & 617 & 626.1870 & $+9.1870$ & 447 & 3163 & 3134.1748 & $-28.8252$ & 781 & 5953 & 5984.1943 & $+31.1943$ \\
114 & 619 & 631.1994 & $+12.1994$ & 448 & 3167 & 3144.1963 & $-22.8037$ & 782 & 5981 & 6000.2382 & $+19.2382$ \\
115 & 631 & 635.0364 & $+4.0364$ & 449 & 3169 & 3150.1564 & $-18.8436$ & 783 & 5987 & 6012.2388 & $+25.2388$ \\
116 & 641 & 638.0574 & $-2.9426$ & 450 & 3181 & 3155.1151 & $-25.8849$ & 784 & 6007 & 6020.1066 & $+13.1066$ \\
117 & 643 & 642.1105 & $-0.8895$ & 451 & 3187 & 3159.1966 & $-27.8034$ & 785 & 6011 & 6025.1631 & $+14.1631$ \\
118 & 647 & 652.2519 & $+5.2519$ & 452 & 3191 & 3164.1538 & $-26.8462$ & 786 & 6029 & 6030.1357 & $+1.1357$ \\
119 & 653 & 657.1967 & $+4.1967$ & 453 & 3203 & 3170.1508 & $-32.8492$ & 787 & 6037 & 6035.0769 & $-1.9231$ \\
120 & 659 & 661.1242 & $+2.1242$ & 454 & 3209 & 3182.1918 & $-26.8082$ & 788 & 6043 & 6039.1408 & $-3.8592$ \\
121 & 661 & 666.1305 & $+5.1305$ & 455 & 3217 & 3199.2409 & $-17.7591$ & 789 & 6047 & 6044.1026 & $-2.8974$ \\
122 & 673 & 684.1446 & $+11.1446$ & 456 & 3221 & 3206.1529 & $-14.8471$ & 790 & 6053 & 6049.1283 & $-3.8717$ \\
123 & 677 & 688.1253 & $+11.1253$ & 457 & 3229 & 3211.1657 & $-17.8343$ & 791 & 6067 & 6055.1625 & $-11.8375$ \\
124 & 683 & 693.1262 & $+10.1262$ & 458 & 3251 & 3216.1375 & $-34.8625$ & 792 & 6073 & 6064.1640 & $-8.8360$ \\
125 & 691 & 705.1071 & $+14.1071$ & 459 & 3253 & 3221.1711 & $-31.8289$ & 793 & 6079 & 6079.1805 & $+0.1805$ \\
126 & 701 & 709.1331 & $+8.1331$ & 460 & 3257 & 3228.1795 & $-28.8205$ & 794 & 6089 & 6091.1295 & $+2.1295$ \\
127 & 709 & 712.1804 & $+3.1804$ & 461 & 3259 & 3251.1554 & $-7.8446$ & 795 & 6091 & 6098.1337 & $+7.1337$ \\
128 & 719 & 716.0725 & $-2.9275$ & 462 & 3271 & 3260.1233 & $-10.8767$ & 796 & 6101 & 6104.0863 & $+3.0863$ \\
129 & 727 & 721.1769 & $-5.8231$ & 463 & 3299 & 3265.1671 & $-33.8329$ & 797 & 6113 & 6109.0764 & $-3.9236$ \\
130 & 733 & 736.2032 & $+3.2032$ & 464 & 3301 & 3270.1020 & $-30.8980$ & 798 & 6121 & 6113.1661 & $-7.8339$ \\
131 & 739 & 741.0824 & $+2.0824$ & 465 & 3307 & 3274.1585 & $-32.8415$ & 799 & 6131 & 6118.1213 & $-12.8787$ \\
132 & 743 & 746.1513 & $+3.1513$ & 466 & 3313 & 3279.1163 & $-33.8837$ & 800 & 6133 & 6123.1215 & $-9.8785$ \\
133 & 751 & 762.1431 & $+11.1431$ & 467 & 3319 & 3284.2196 & $-34.7804$ & 801 & 6143 & 6128.2314 & $-14.7686$ \\
134 & 757 & 766.1179 & $+9.1179$ & 468 & 3323 & 3293.1799 & $-29.8201$ & 802 & 6151 & 6136.1334 & $-14.8666$ \\
135 & 761 & 770.0602 & $+9.0602$ & 469 & 3329 & 3307.1402 & $-21.8598$ & 803 & 6163 & 6148.1902 & $-14.8098$ \\
136 & 769 & 774.1790 & $+5.1790$ & 470 & 3331 & 3314.2068 & $-16.7932$ & 804 & 6173 & 6166.1645 & $-6.8355$ \\
137 & 773 & 787.1370 & $+14.1370$ & 471 & 3343 & 3320.0863 & $-22.9137$ & 805 & 6197 & 6175.1769 & $-21.8231$ \\
138 & 787 & 791.2035 & $+4.2035$ & 472 & 3347 & 3324.1300 & $-22.8700$ & 806 & 6199 & 6182.1224 & $-16.8776$ \\
139 & 797 & 795.0554 & $-1.9446$ & 473 & 3359 & 3328.1464 & $-30.8536$ & 807 & 6203 & 6187.2377 & $-15.7623$ \\
140 & 809 & 798.0927 & $-10.9073$ & 474 & 3361 & 3332.1809 & $-28.8191$ & 808 & 6211 & 6193.1275 & $-17.8725$ \\
141 & 811 & 802.1355 & $-8.8645$ & 475 & 3371 & 3337.1305 & $-33.8695$ & 809 & 6217 & 6198.2448 & $-18.7552$ \\
142 & 821 & 814.1596 & $-6.8404$ & 476 & 3373 & 3343.1542 & $-29.8458$ & 810 & 6221 & 6205.1467 & $-15.8533$ \\
143 & 823 & 820.1013 & $-2.8987$ & 477 & 3389 & 3358.2128 & $-30.7872$ & 811 & 6229 & 6214.2019 & $-14.7981$ \\
144 & 827 & 824.1010 & $-2.8990$ & 478 & 3391 & 3373.1885 & $-17.8115$ & 812 & 6247 & 6249.1596 & $+2.1596$ \\
145 & 829 & 829.1200 & $+0.1200$ & 479 & 3407 & 3380.1137 & $-26.8863$ & 813 & 6257 & 6259.1317 & $+2.1317$ \\
146 & 839 & 850.0981 & $+11.0981$ & 480 & 3413 & 3385.1452 & $-27.8548$ & 814 & 6263 & 6266.1082 & $+3.1082$ \\
147 & 853 & 854.1110 & $+1.1110$ & 481 & 3433 & 3390.1794 & $-42.8206$ & 815 & 6269 & 6272.1128 & $+3.1128$ \\
148 & 857 & 858.2501 & $+1.2501$ & 482 & 3449 & 3396.1996 & $-52.8004$ & 816 & 6271 & 6278.1248 & $+7.1248$ \\
149 & 859 & 870.1575 & $+11.1575$ & 483 & 3457 & 3409.2101 & $-47.7899$ & 817 & 6277 & 6285.1249 & $+8.1249$ \\
150 & 863 & 876.1066 & $+13.1066$ & 484 & 3461 & 3434.1310 & $-26.8690$ & 818 & 6287 & 6294.2094 & $+7.2094$ \\
151 & 877 & 879.1848 & $+2.1848$ & 485 & 3463 & 3440.1841 & $-22.8159$ & 819 & 6299 & 6319.2203 & $+20.2203$ \\
152 & 881 & 882.1983 & $+1.1983$ & 486 & 3467 & 3446.1293 & $-20.8707$ & 820 & 6301 & 6332.1178 & $+31.1178$ \\
153 & 883 & 886.0679 & $+3.0679$ & 487 & 3469 & 3452.1334 & $-16.8666$ & 821 & 6311 & 6339.1017 & $+28.1017$ \\
154 & 887 & 890.1740 & $+3.1740$ & 488 & 3491 & 3461.1543 & $-29.8457$ & 822 & 6317 & 6344.1800 & $+27.1800$ \\
155 & 907 & 903.1825 & $-3.8175$ & 489 & 3499 & 3490.1483 & $-8.8517$ & 823 & 6323 & 6349.2281 & $+26.2281$ \\
156 & 911 & 908.1821 & $-2.8179$ & 490 & 3511 & 3497.0978 & $-13.9022$ & 824 & 6329 & 6355.0914 & $+26.0914$ \\
157 & 919 & 912.1256 & $-6.8744$ & 491 & 3517 & 3502.1047 & $-14.8953$ & 825 & 6337 & 6360.1410 & $+23.1410$ \\
158 & 929 & 916.1235 & $-12.8765$ & 492 & 3527 & 3507.0822 & $-19.9178$ & 826 & 6343 & 6366.1786 & $+23.1786$ \\
159 & 937 & 923.2156 & $-13.7844$ & 493 & 3529 & 3512.0963 & $-16.9037$ & 827 & 6353 & 6375.1452 & $+22.1452$ \\
160 & 941 & 936.1740 & $-4.8260$ & 494 & 3533 & 3518.1307 & $-14.8693$ & 828 & 6359 & 6392.1884 & $+33.1884$ \\
161 & 947 & 940.1852 & $-6.8148$ & 495 & 3539 & 3528.2055 & $-10.7945$ & 829 & 6361 & 6407.1625 & $+46.1625$ \\
162 & 953 & 944.1096 & $-8.8904$ & 496 & 3541 & 3543.2003 & $+2.2003$ & 830 & 6367 & 6415.1249 & $+48.1249$ \\
163 & 967 & 948.1594 & $-18.8406$ & 497 & 3547 & 3551.0938 & $+4.0938$ & 831 & 6373 & 6421.1049 & $+48.1049$ \\
164 & 971 & 961.1896 & $-9.8104$ & 498 & 3557 & 3556.0958 & $-0.9042$ & 832 & 6379 & 6426.1171 & $+47.1171$ \\
165 & 977 & 968.1073 & $-8.8927$ & 499 & 3559 & 3560.1556 & $+1.1556$ & 833 & 6389 & 6431.1010 & $+42.1010$ \\
166 & 983 & 972.0546 & $-10.9454$ & 500 & 3571 & 3564.2246 & $-6.7754$ & 834 & 6397 & 6436.0893 & $+39.0893$ \\
167 & 991 & 975.1122 & $-15.8878$ & 501 & 3581 & 3569.1259 & $-11.8741$ & 835 & 6421 & 6441.1147 & $+20.1147$ \\
168 & 997 & 979.1111 & $-17.8889$ & 502 & 3583 & 3574.1666 & $-8.8334$ & 836 & 6427 & 6447.1179 & $+20.1179$ \\
169 & 1009 & 987.1885 & $-21.8115$ & 503 & 3593 & 3582.1662 & $-10.8338$ & 837 & 6449 & 6454.2179 & $+5.2179$ \\
170 & 1013 & 999.1110 & $-13.8890$ & 504 & 3607 & 3603.2141 & $-3.7859$ & 838 & 6451 & 6468.1750 & $+17.1750$ \\
171 & 1019 & 1003.2085 & $-15.7915$ & 505 & 3613 & 3612.1946 & $-0.8054$ & 839 & 6469 & 6484.1702 & $+15.1702$ \\
172 & 1021 & 1008.1539 & $-12.8461$ & 506 & 3617 & 3618.1614 & $+1.1614$ & 840 & 6473 & 6493.1410 & $+20.1410$ \\
173 & 1031 & 1035.1474 & $+4.1474$ & 507 & 3623 & 3623.1826 & $+0.1826$ & 841 & 6481 & 6499.1961 & $+18.1961$ \\
174 & 1033 & 1040.2508 & $+7.2508$ & 508 & 3631 & 3629.0905 & $-1.9095$ & 842 & 6491 & 6505.1265 & $+14.1265$ \\
175 & 1039 & 1059.1355 & $+20.1355$ & 509 & 3637 & 3635.1365 & $-1.8635$ & 843 & 6521 & 6510.1525 & $-10.8475$ \\
176 & 1049 & 1064.2205 & $+15.2205$ & 510 & 3643 & 3648.2322 & $+5.2322$ & 844 & 6529 & 6515.1953 & $-13.8047$ \\
177 & 1051 & 1068.1539 & $+17.1539$ & 511 & 3659 & 3670.1553 & $+11.1553$ & 845 & 6547 & 6521.1514 & $-25.8486$ \\
178 & 1061 & 1072.0944 & $+11.0944$ & 512 & 3671 & 3677.1069 & $+6.1069$ & 846 & 6551 & 6528.1567 & $-22.8433$ \\
179 & 1063 & 1076.1849 & $+13.1849$ & 513 & 3673 & 3682.1500 & $+9.1500$ & 847 & 6553 & 6540.1754 & $-12.8246$ \\
180 & 1069 & 1090.1959 & $+21.1959$ & 514 & 3677 & 3687.1971 & $+10.1971$ & 848 & 6563 & 6569.1748 & $+6.1748$ \\
181 & 1087 & 1097.1027 & $+10.1027$ & 515 & 3691 & 3693.2007 & $+2.2007$ & 849 & 6569 & 6578.2150 & $+9.2150$ \\
182 & 1091 & 1101.0906 & $+10.0906$ & 516 & 3697 & 3703.1970 & $+6.1970$ & 850 & 6571 & 6585.1629 & $+14.1629$ \\
183 & 1093 & 1105.0703 & $+12.0703$ & 517 & 3701 & 3728.2204 & $+27.2204$ & 851 & 6577 & 6591.1497 & $+14.1497$ \\
184 & 1097 & 1109.2010 & $+12.2010$ & 518 & 3709 & 3735.2415 & $+26.2415$ & 852 & 6581 & 6597.1282 & $+16.1282$ \\
185 & 1103 & 1125.2160 & $+22.2160$ & 519 & 3719 & 3741.0998 & $+22.0998$ & 853 & 6599 & 6603.1671 & $+4.1671$ \\
186 & 1109 & 1131.1785 & $+22.1785$ & 520 & 3727 & 3745.2187 & $+18.2187$ & 854 & 6607 & 6611.1599 & $+4.1599$ \\
187 & 1117 & 1135.2052 & $+18.2052$ & 521 & 3733 & 3750.1669 & $+17.1669$ & 855 & 6619 & 6625.2230 & $+6.2230$ \\
188 & 1123 & 1140.0986 & $+17.0986$ & 522 & 3739 & 3756.1017 & $+17.1017$ & 856 & 6637 & 6648.1572 & $+11.1572$ \\
189 & 1129 & 1148.2078 & $+19.2078$ & 523 & 3761 & 3763.1526 & $+2.1526$ & 857 & 6653 & 6657.1283 & $+4.1283$ \\
190 & 1151 & 1163.1249 & $+12.1249$ & 524 & 3767 & 3778.1695 & $+11.1695$ & 858 & 6659 & 6663.1601 & $+4.1601$ \\
191 & 1153 & 1167.1264 & $+14.1264$ & 525 & 3769 & 3789.1701 & $+20.1701$ & 859 & 6661 & 6669.0930 & $+8.0930$ \\
192 & 1163 & 1171.0550 & $+8.0550$ & 526 & 3779 & 3795.2208 & $+16.2208$ & 860 & 6673 & 6674.1076 & $+1.1076$ \\
193 & 1171 & 1174.1700 & $+3.1700$ & 527 & 3793 & 3800.2202 & $+7.2202$ & 861 & 6679 & 6679.1290 & $+0.1290$ \\
194 & 1181 & 1179.1805 & $-1.8195$ & 528 & 3797 & 3805.1339 & $+8.1339$ & 862 & 6689 & 6684.2380 & $-4.7620$ \\
195 & 1187 & 1190.1627 & $+3.1627$ & 529 & 3803 & 3810.0886 & $+7.0886$ & 863 & 6691 & 6691.1594 & $+0.1594$ \\
196 & 1193 & 1196.2211 & $+3.2211$ & 530 & 3821 & 3815.1156 & $-5.8844$ & 864 & 6701 & 6700.2104 & $-0.7896$ \\
197 & 1201 & 1200.2294 & $-0.7706$ & 531 & 3823 & 3822.1383 & $-0.8617$ & 865 & 6703 & 6719.1737 & $+16.1737$ \\
198 & 1213 & 1204.1230 & $-8.8770$ & 532 & 3833 & 3847.2074 & $+14.2074$ & 866 & 6709 & 6732.1496 & $+23.1496$ \\
199 & 1217 & 1208.1130 & $-8.8870$ & 533 & 3847 & 3858.1242 & $+11.1242$ & 867 & 6719 & 6739.1971 & $+20.1971$ \\
200 & 1223 & 1214.1813 & $-8.8187$ & 534 & 3851 & 3864.1525 & $+13.1525$ & 868 & 6733 & 6745.1602 & $+12.1602$ \\
201 & 1229 & 1234.1617 & $+5.1617$ & 535 & 3853 & 3870.1187 & $+17.1187$ & 869 & 6737 & 6750.1738 & $+13.1738$ \\
202 & 1231 & 1239.1422 & $+8.1422$ & 536 & 3863 & 3876.1484 & $+13.1484$ & 870 & 6761 & 6755.1420 & $-5.8580$ \\
203 & 1237 & 1244.1696 & $+7.1696$ & 537 & 3877 & 3885.2086 & $+8.2086$ & 871 & 6763 & 6760.1145 & $-2.8855$ \\
204 & 1249 & 1262.1869 & $+13.1869$ & 538 & 3881 & 3917.1363 & $+36.1363$ & 872 & 6779 & 6765.1198 & $-13.8802$ \\
205 & 1259 & 1269.0881 & $+10.0881$ & 539 & 3889 & 3924.1419 & $+35.1419$ & 873 & 6781 & 6770.2180 & $-10.7820$ \\
206 & 1277 & 1273.0736 & $-3.9264$ & 540 & 3907 & 3930.1158 & $+23.1158$ & 874 & 6791 & 6777.1893 & $-13.8107$ \\
207 & 1279 & 1276.1902 & $-2.8098$ & 541 & 3911 & 3936.0980 & $+25.0980$ & 875 & 6793 & 6787.2075 & $-5.7925$ \\
208 & 1283 & 1281.0824 & $-1.9176$ & 542 & 3917 & 3942.2406 & $+25.2406$ & 876 & 6803 & 6803.1936 & $+0.1936$ \\
209 & 1289 & 1288.1618 & $-0.8382$ & 543 & 3919 & 3967.1976 & $+48.1976$ & 877 & 6823 & 6814.1984 & $-8.8016$ \\
210 & 1291 & 1300.1525 & $+9.1525$ & 544 & 3923 & 3980.1423 & $+57.1423$ & 878 & 6827 & 6822.1094 & $-4.8906$ \\
211 & 1297 & 1305.0959 & $+8.0959$ & 545 & 3929 & 3986.1390 & $+57.1390$ & 879 & 6829 & 6828.0874 & $-0.9126$ \\
212 & 1301 & 1308.2039 & $+7.2039$ & 546 & 3931 & 3991.1290 & $+60.1290$ & 880 & 6833 & 6833.1122 & $+0.1122$ \\
213 & 1303 & 1312.0823 & $+9.0823$ & 547 & 3943 & 3996.0869 & $+53.0869$ & 881 & 6841 & 6838.1202 & $-2.8798$ \\
214 & 1307 & 1315.2176 & $+8.2176$ & 548 & 3947 & 4000.2131 & $+53.2131$ & 882 & 6857 & 6843.1460 & $-13.8540$ \\
215 & 1319 & 1320.2297 & $+1.2297$ & 549 & 3967 & 4006.1341 & $+39.1341$ & 883 & 6863 & 6849.1169 & $-13.8831$ \\
216 & 1321 & 1333.1642 & $+12.1642$ & 550 & 3989 & 4013.1954 & $+24.1954$ & 884 & 6869 & 6856.1298 & $-12.8702$ \\
217 & 1327 & 1340.1193 & $+13.1193$ & 551 & 4001 & 4026.2110 & $+25.2110$ & 885 & 6871 & 6867.2183 & $-3.7817$ \\
218 & 1361 & 1344.1685 & $-16.8315$ & 552 & 4003 & 4037.1856 & $+34.1856$ & 886 & 6883 & 6897.1894 & $+14.1894$ \\
219 & 1367 & 1348.1871 & $-18.8129$ & 553 & 4007 & 4044.0995 & $+37.0995$ & 887 & 6899 & 6907.1603 & $+8.1603$ \\
220 & 1373 & 1353.2045 & $-19.7955$ & 554 & 4013 & 4049.0810 & $+36.0810$ & 888 & 6907 & 6914.1568 & $+7.1568$ \\
221 & 1381 & 1376.1714 & $-4.8286$ & 555 & 4019 & 4053.1223 & $+34.1223$ & 889 & 6911 & 6920.1919 & $+9.1919$ \\
222 & 1399 & 1381.2304 & $-17.7696$ & 556 & 4021 & 4057.1494 & $+36.1494$ & 890 & 6917 & 6927.1058 & $+10.1058$ \\
223 & 1409 & 1386.1129 & $-22.8871$ & 557 & 4027 & 4061.2107 & $+34.2107$ & 891 & 6947 & 6934.1205 & $-12.8795$ \\
224 & 1423 & 1391.1422 & $-31.8578$ & 558 & 4049 & 4066.1554 & $+17.1554$ & 892 & 6949 & 6943.2164 & $-5.7836$ \\
225 & 1427 & 1404.2010 & $-22.7990$ & 559 & 4051 & 4072.1499 & $+21.1499$ & 893 & 6959 & 6976.1840 & $+17.1840$ \\
226 & 1429 & 1413.1140 & $-15.8860$ & 560 & 4057 & 4082.1779 & $+25.1779$ & 894 & 6961 & 6988.1276 & $+27.1276$ \\
227 & 1433 & 1417.1484 & $-15.8516$ & 561 & 4073 & 4100.1739 & $+27.1739$ & 895 & 6967 & 6995.1400 & $+28.1400$ \\
228 & 1439 & 1421.0968 & $-17.9032$ & 562 & 4079 & 4108.1728 & $+29.1728$ & 896 & 6971 & 7001.1551 & $+30.1551$ \\
229 & 1447 & 1425.0792 & $-21.9208$ & 563 & 4091 & 4114.1227 & $+23.1227$ & 897 & 6977 & 7007.1397 & $+30.1397$ \\
230 & 1451 & 1430.1040 & $-20.8960$ & 564 & 4093 & 4119.1115 & $+26.1115$ & 898 & 6983 & 7013.1620 & $+30.1620$ \\
231 & 1453 & 1447.1742 & $-5.8258$ & 565 & 4099 & 4124.0913 & $+25.0913$ & 899 & 6991 & 7020.2063 & $+29.2063$ \\
232 & 1459 & 1454.1326 & $-4.8674$ & 566 & 4111 & 4129.1171 & $+18.1171$ & 900 & 6997 & 7032.1631 & $+35.1631$ \\
233 & 1471 & 1459.0779 & $-11.9221$ & 567 & 4127 & 4135.1730 & $+8.1730$ & 901 & 7001 & 7059.2235 & $+58.2235$ \\
234 & 1481 & 1463.1780 & $-17.8220$ & 568 & 4129 & 4147.1853 & $+18.1853$ & 902 & 7013 & 7070.1505 & $+57.1505$ \\
235 & 1483 & 1471.2437 & $-11.7563$ & 569 & 4133 & 4166.1596 & $+33.1596$ & 903 & 7019 & 7077.1204 & $+58.1204$ \\
236 & 1487 & 1491.2253 & $+4.2253$ & 570 & 4139 & 4173.2379 & $+34.2379$ & 904 & 7027 & 7082.2226 & $+55.2226$ \\
237 & 1489 & 1496.2035 & $+7.2035$ & 571 & 4153 & 4179.1365 & $+26.1365$ & 905 & 7039 & 7088.1016 & $+49.1016$ \\
238 & 1493 & 1501.1561 & $+8.1561$ & 572 & 4157 & 4184.1179 & $+27.1179$ & 906 & 7043 & 7093.1208 & $+50.1208$ \\
239 & 1499 & 1508.2039 & $+9.2039$ & 573 & 4159 & 4189.0987 & $+30.0987$ & 907 & 7057 & 7098.1800 & $+41.1800$ \\
240 & 1511 & 1528.1286 & $+17.1286$ & 574 & 4177 & 4194.1336 & $+17.1336$ & 908 & 7069 & 7104.2009 & $+35.2009$ \\
241 & 1523 & 1533.0746 & $+10.0746$ & 575 & 4201 & 4201.1176 & $+0.1176$ & 909 & 7079 & 7112.1650 & $+33.1650$ \\
242 & 1531 & 1536.2254 & $+5.2254$ & 576 & 4211 & 4214.1793 & $+3.1793$ & 910 & 7103 & 7124.1912 & $+21.1912$ \\
243 & 1543 & 1540.1696 & $-2.8304$ & 577 & 4217 & 4231.1522 & $+14.1522$ & 911 & 7109 & 7140.2028 & $+31.2028$ \\
244 & 1549 & 1545.1106 & $-3.8894$ & 578 & 4219 & 4238.1392 & $+19.1392$ & 912 & 7121 & 7150.1656 & $+29.1656$ \\
245 & 1553 & 1552.2090 & $-0.7910$ & 579 & 4229 & 4243.1610 & $+14.1610$ & 913 & 7127 & 7157.1338 & $+30.1338$ \\
246 & 1559 & 1563.2131 & $+4.2131$ & 580 & 4231 & 4248.0954 & $+17.0954$ & 914 & 7129 & 7163.0859 & $+34.0859$ \\
247 & 1567 & 1569.0950 & $+2.0950$ & 581 & 4241 & 4252.1437 & $+11.1437$ & 915 & 7151 & 7168.0901 & $+17.0901$ \\
248 & 1571 & 1573.0829 & $+2.0829$ & 582 & 4243 & 4257.0808 & $+14.0808$ & 916 & 7159 & 7172.2505 & $+13.2505$ \\
249 & 1579 & 1576.2189 & $-2.7811$ & 583 & 4253 & 4262.0881 & $+9.0881$ & 917 & 7177 & 7177.2239 & $+0.2239$ \\
250 & 1583 & 1580.2451 & $-2.7549$ & 584 & 4259 & 4268.1213 & $+9.1213$ & 918 & 7187 & 7183.1121 & $-3.8879$ \\
251 & 1597 & 1587.1309 & $-9.8691$ & 585 & 4261 & 4277.1804 & $+16.1804$ & 919 & 7193 & 7189.1335 & $-3.8665$ \\
252 & 1601 & 1609.1269 & $+8.1269$ & 586 & 4271 & 4290.1888 & $+19.1888$ & 920 & 7207 & 7197.1626 & $-9.8374$ \\
253 & 1607 & 1614.2199 & $+7.2199$ & 587 & 4273 & 4298.2185 & $+25.2185$ & 921 & 7211 & 7214.1783 & $+3.1783$ \\
254 & 1609 & 1620.1133 & $+11.1133$ & 588 & 4283 & 4304.1596 & $+21.1596$ & 922 & 7213 & 7233.1440 & $+20.1440$ \\
255 & 1613 & 1630.2090 & $+17.2090$ & 589 & 4289 & 4309.1209 & $+20.1209$ & 923 & 7219 & 7242.1073 & $+23.1073$ \\
256 & 1619 & 1650.1672 & $+31.1672$ & 590 & 4297 & 4313.1833 & $+16.1833$ & 924 & 7229 & 7248.1581 & $+19.1581$ \\
257 & 1621 & 1655.1315 & $+34.1315$ & 591 & 4327 & 4318.0858 & $-8.9142$ & 925 & 7237 & 7254.1270 & $+17.1270$ \\
258 & 1627 & 1660.0847 & $+33.0847$ & 592 & 4337 & 4322.1738 & $-14.8262$ & 926 & 7243 & 7259.2123 & $+16.2123$ \\
259 & 1637 & 1665.2272 & $+28.2272$ & 593 & 4339 & 4328.1063 & $-10.8937$ & 927 & 7247 & 7265.1601 & $+18.1601$ \\
260 & 1657 & 1685.1699 & $+28.1699$ & 594 & 4349 & 4335.2012 & $-13.7988$ & 928 & 7253 & 7272.1091 & $+19.1091$ \\
261 & 1663 & 1691.1681 & $+28.1681$ & 595 & 4357 & 4357.2365 & $+0.2365$ & 929 & 7283 & 7280.1560 & $-2.8440$ \\
262 & 1667 & 1695.2491 & $+28.2491$ & 596 & 4363 & 4368.2150 & $+5.2150$ & 930 & 7297 & 7297.2521 & $+0.2521$ \\
263 & 1669 & 1699.2422 & $+30.2422$ & 597 & 4373 & 4375.1351 & $+2.1351$ & 931 & 7307 & 7318.1900 & $+11.1900$ \\
264 & 1693 & 1704.1334 & $+11.1334$ & 598 & 4391 & 4381.0877 & $-9.9123$ & 932 & 7309 & 7327.1551 & $+18.1551$ \\
265 & 1697 & 1711.1798 & $+14.1798$ & 599 & 4397 & 4386.1536 & $-10.8464$ & 933 & 7321 & 7334.1023 & $+13.1023$ \\
266 & 1699 & 1727.1269 & $+28.1269$ & 600 & 4409 & 4393.1121 & $-15.8879$ & 934 & 7331 & 7339.2000 & $+8.2000$ \\
267 & 1709 & 1733.0868 & $+24.0868$ & 601 & 4421 & 4404.1742 & $-16.8258$ & 935 & 7333 & 7345.1297 & $+12.1297$ \\
268 & 1721 & 1737.1364 & $+16.1364$ & 602 & 4423 & 4433.1312 & $+10.1312$ & 936 & 7349 & 7351.1041 & $+2.1041$ \\
269 & 1723 & 1741.2102 & $+18.2102$ & 603 & 4441 & 4440.1625 & $-0.8375$ & 937 & 7351 & 7357.1568 & $+6.1568$ \\
270 & 1733 & 1747.1495 & $+14.1495$ & 604 & 4447 & 4446.1503 & $-0.8497$ & 938 & 7369 & 7365.2112 & $-3.7888$ \\
271 & 1741 & 1770.1563 & $+29.1563$ & 605 & 4451 & 4452.1218 & $+1.1218$ & 939 & 7393 & 7383.1842 & $-9.8158$ \\
272 & 1747 & 1776.1354 & $+29.1354$ & 606 & 4457 & 4458.1915 & $+1.1915$ & 940 & 7411 & 7403.1632 & $-7.8368$ \\
273 & 1753 & 1780.2224 & $+27.2224$ & 607 & 4463 & 4469.1808 & $+6.1808$ & 941 & 7417 & 7412.1108 & $-4.8892$ \\
274 & 1759 & 1785.1063 & $+26.1063$ & 608 & 4481 & 4498.1269 & $+17.1269$ & 942 & 7433 & 7418.1327 & $-14.8673$ \\
275 & 1777 & 1790.1557 & $+13.1557$ & 609 & 4483 & 4505.1394 & $+22.1394$ & 943 & 7451 & 7423.1983 & $-27.8017$ \\
276 & 1783 & 1800.2051 & $+17.2051$ & 610 & 4493 & 4510.2300 & $+17.2300$ & 944 & 7457 & 7428.2188 & $-28.7812$ \\
277 & 1787 & 1810.2125 & $+23.2125$ & 611 & 4507 & 4515.2186 & $+8.2186$ & 945 & 7459 & 7433.2203 & $-25.7797$ \\
278 & 1789 & 1815.1773 & $+26.1773$ & 612 & 4513 & 4520.2076 & $+7.2076$ & 946 & 7477 & 7439.0995 & $-37.9005$ \\
279 & 1801 & 1819.1221 & $+18.1221$ & 613 & 4517 & 4526.1275 & $+9.1275$ & 947 & 7481 & 7444.2001 & $-36.7999$ \\
280 & 1811 & 1822.2042 & $+11.2042$ & 614 & 4519 & 4533.1364 & $+14.1364$ & 948 & 7487 & 7451.1815 & $-35.8185$ \\
281 & 1823 & 1826.0840 & $+3.0840$ & 615 & 4523 & 4545.1641 & $+22.1641$ & 949 & 7489 & 7461.1463 & $-27.8537$ \\
282 & 1831 & 1829.2458 & $-1.7542$ & 616 & 4547 & 4559.2047 & $+12.2047$ & 950 & 7499 & 7475.1602 & $-23.8398$ \\
283 & 1847 & 1834.1824 & $-12.8176$ & 617 & 4549 & 4567.1488 & $+18.1488$ & 951 & 7507 & 7486.2426 & $-20.7574$ \\
284 & 1861 & 1843.1978 & $-17.8022$ & 618 & 4561 & 4573.0854 & $+12.0854$ & 952 & 7517 & 7494.1905 & $-22.8095$ \\
285 & 1867 & 1853.2027 & $-13.7973$ & 619 & 4567 & 4577.1858 & $+10.1858$ & 953 & 7523 & 7500.1934 & $-22.8066$ \\
286 & 1871 & 1859.0767 & $-11.9233$ & 620 & 4583 & 4582.1034 & $-0.8966$ & 954 & 7529 & 7506.0829 & $-22.9171$ \\
287 & 1873 & 1863.0681 & $-9.9319$ & 621 & 4591 & 4586.1839 & $-4.8161$ & 955 & 7537 & 7510.1940 & $-26.8060$ \\
288 & 1877 & 1866.1832 & $-10.8168$ & 622 & 4597 & 4591.1700 & $-5.8300$ & 956 & 7541 & 7515.1291 & $-25.8709$ \\
289 & 1879 & 1870.1600 & $-8.8400$ & 623 & 4603 & 4597.1749 & $-5.8251$ & 957 & 7547 & 7520.0902 & $-26.9098$ \\
290 & 1889 & 1875.1388 & $-13.8612$ & 624 & 4621 & 4607.1728 & $-13.8272$ & 958 & 7549 & 7525.0855 & $-23.9145$ \\
291 & 1901 & 1885.2000 & $-15.8000$ & 625 & 4637 & 4628.1755 & $-8.8245$ & 959 & 7559 & 7530.1406 & $-28.8594$ \\
292 & 1907 & 1899.2482 & $-7.7518$ & 626 & 4639 & 4637.1525 & $-1.8475$ & 960 & 7561 & 7537.1119 & $-23.8881$ \\
293 & 1913 & 1905.1448 & $-7.8552$ & 627 & 4643 & 4643.1561 & $+0.1561$ & 961 & 7573 & 7546.1598 & $-26.8402$ \\
294 & 1931 & 1910.1046 & $-20.8954$ & 628 & 4649 & 4648.2062 & $-0.7938$ & 962 & 7577 & 7563.1873 & $-13.8127$ \\
295 & 1933 & 1915.1364 & $-17.8636$ & 629 & 4651 & 4654.0889 & $+3.0889$ & 963 & 7583 & 7577.1558 & $-5.8442$ \\
296 & 1949 & 1940.1808 & $-8.8192$ & 630 & 4657 & 4659.1633 & $+2.1633$ & 964 & 7589 & 7585.1594 & $-3.8406$ \\
297 & 1951 & 1948.1175 & $-2.8825$ & 631 & 4663 & 4666.1874 & $+3.1874$ & 965 & 7591 & 7591.2224 & $+0.2224$ \\
298 & 1973 & 1953.0953 & $-19.9047$ & 632 & 4673 & 4683.2213 & $+10.2213$ & 966 & 7603 & 7597.1469 & $-5.8531$ \\
299 & 1979 & 1957.2030 & $-21.7970$ & 633 & 4679 & 4701.2121 & $+22.2121$ & 967 & 7607 & 7602.2052 & $-4.7948$ \\
300 & 1987 & 1963.1685 & $-23.8315$ & 634 & 4691 & 4709.1040 & $+18.1040$ & 968 & 7621 & 7608.1175 & $-12.8825$ \\
301 & 1993 & 1978.1719 & $-14.8281$ & 635 & 4703 & 4714.1763 & $+11.1763$ & 969 & 7639 & 7614.1082 & $-24.8918$ \\
302 & 1997 & 1987.2099 & $-9.7901$ & 636 & 4721 & 4719.2292 & $-1.7708$ & 970 & 7643 & 7621.1189 & $-21.8811$ \\
303 & 1999 & 1992.1883 & $-6.8117$ & 637 & 4723 & 4725.1228 & $+2.1228$ & 971 & 7649 & 7631.1663 & $-17.8337$ \\
304 & 2003 & 1996.2189 & $-6.7811$ & 638 & 4729 & 4731.1617 & $+2.1617$ & 972 & 7669 & 7656.1998 & $-12.8002$ \\
305 & 2011 & 2000.2526 & $-10.7474$ & 639 & 4733 & 4741.1632 & $+8.1632$ & 973 & 7673 & 7669.2208 & $-3.7792$ \\
306 & 2017 & 2005.1736 & $-11.8264$ & 640 & 4751 & 4765.1530 & $+14.1530$ & 974 & 7681 & 7677.1972 & $-3.8028$ \\
307 & 2027 & 2014.1942 & $-12.8058$ & 641 & 4759 & 4774.1077 & $+15.1077$ & 975 & 7687 & 7684.1274 & $-2.8726$ \\
308 & 2029 & 2034.1186 & $+5.1186$ & 642 & 4783 & 4779.2209 & $-3.7791$ & 976 & 7691 & 7690.1320 & $-0.8680$ \\
309 & 2039 & 2039.1959 & $+0.1959$ & 643 & 4787 & 4784.2094 & $-2.7906$ & 977 & 7699 & 7696.1559 & $-2.8441$ \\
310 & 2053 & 2044.1700 & $-8.8300$ & 644 & 4789 & 4789.1435 & $+0.1435$ & 978 & 7703 & 7703.1624 & $+0.1624$ \\
311 & 2063 & 2050.1539 & $-12.8461$ & 645 & 4793 & 4794.1115 & $+1.1115$ & 979 & 7717 & 7713.1613 & $-3.8387$ \\
312 & 2069 & 2073.1594 & $+4.1594$ & 646 & 4799 & 4799.1402 & $+0.1402$ & 980 & 7723 & 7749.1757 & $+26.1757$ \\
313 & 2081 & 2081.1765 & $+0.1765$ & 647 & 4801 & 4805.1949 & $+4.1949$ & 981 & 7727 & 7761.1384 & $+34.1384$ \\
314 & 2083 & 2086.1956 & $+3.1956$ & 648 & 4813 & 4815.1612 & $+2.1612$ & 982 & 7741 & 7768.1766 & $+27.1766$ \\
315 & 2087 & 2091.1594 & $+4.1594$ & 649 & 4817 & 4829.1706 & $+12.1706$ & 983 & 7753 & 7774.2311 & $+21.2311$ \\
316 & 2089 & 2097.1512 & $+8.1512$ & 650 & 4831 & 4838.1537 & $+7.1537$ & 984 & 7757 & 7780.2035 & $+23.2035$ \\
317 & 2099 & 2117.1977 & $+18.1977$ & 651 & 4861 & 4844.1525 & $-16.8475$ & 985 & 7759 & 7786.2199 & $+27.2199$ \\
318 & 2111 & 2125.2130 & $+14.2130$ & 652 & 4871 & 4849.1438 & $-21.8562$ & 986 & 7789 & 7793.1874 & $+4.1874$ \\
319 & 2113 & 2130.1450 & $+17.1450$ & 653 & 4877 & 4854.0774 & $-22.9226$ & 987 & 7793 & 7802.1813 & $+9.1813$ \\
320 & 2129 & 2134.1472 & $+5.1472$ & 654 & 4889 & 4858.1238 & $-30.8762$ & 988 & 7817 & 7819.1849 & $+2.1849$ \\
321 & 2131 & 2138.1507 & $+7.1507$ & 655 & 4903 & 4862.2267 & $-40.7733$ & 989 & 7823 & 7837.1870 & $+14.1870$ \\
322 & 2137 & 2143.1033 & $+6.1033$ & 656 & 4909 & 4867.2231 & $-41.7769$ & 990 & 7829 & 7846.1774 & $+17.1774$ \\
323 & 2141 & 2149.1816 & $+8.1816$ & 657 & 4919 & 4874.1219 & $-44.8781$ & 991 & 7841 & 7853.1183 & $+12.1183$ \\
324 & 2143 & 2162.1343 & $+19.1343$ & 658 & 4931 & 4884.1951 & $-46.8049$ & 992 & 7853 & 7858.1992 & $+5.1992$ \\
325 & 2153 & 2169.1611 & $+16.1611$ & 659 & 4933 & 4905.1407 & $-27.8593$ & 993 & 7867 & 7864.0833 & $-2.9167$ \\
326 & 2161 & 2174.1224 & $+13.1224$ & 660 & 4937 & 4913.2263 & $-23.7737$ & 994 & 7873 & 7869.0898 & $-3.9102$ \\
327 & 2179 & 2178.1252 & $-0.8748$ & 661 & 4943 & 4920.1034 & $-22.8966$ & 995 & 7877 & 7874.1153 & $-2.8847$ \\
328 & 2203 & 2182.1158 & $-20.8842$ & 662 & 4951 & 4925.1450 & $-25.8550$ & 996 & 7879 & 7879.2346 & $+0.2346$ \\
329 & 2207 & 2186.1720 & $-20.8280$ & 663 & 4957 & 4930.1899 & $-26.8101$ & 997 & 7883 & 7886.1803 & $+3.1803$ \\
330 & 2213 & 2192.1796 & $-20.8204$ & 664 & 4967 & 4936.1552 & $-30.8448$ & 998 & 7901 & 7896.1478 & $-4.8522$ \\
331 & 2221 & 2215.1472 & $-5.8528$ & 665 & 4969 & 4943.2236 & $-25.7764$ & 999 & 7907 & 7914.1644 & $+7.1644$ \\
332 & 2237 & 2222.1383 & $-14.8617$ & 666 & 4973 & 4965.2293 & $-7.7707$ & 1000 & 7919 & 7928.1950 & $+9.1950$ \\
333 & 2239 & 2227.1649 & $-11.8351$ & 667 & 4987 & 4982.1295 & $-4.8705$ &  &  &  &  \\
334 & 2243 & 2232.2362 & $-10.7638$ & 668 & 4993 & 4989.1289 & $-3.8711$ &  &  &  &  \\
\end{longtable}}

{\scriptsize\setlength{\tabcolsep}{3pt}
\begin{longtable}{rrrr@{\quad}rrrr@{\quad}rrrr}
\caption{The prime count $\pi(n)$, exact, and the synthetic prime count $\pi_{\RVM}(n)$, a real-valued function computed from truncated sums (not an integer count), at the integers $n\le1000$, with the difference $\pi_{\RVM}(n)-\pi(n)$.}\label{tab:pi-1000}\\
\toprule
$n$ & $\pi(n)$ & $\pi_{\RVM}(n)$ & $\pi_{\RVM}(n)-\pi(n)$ & $n$ & $\pi(n)$ & $\pi_{\RVM}(n)$ & $\pi_{\RVM}(n)-\pi(n)$ & $n$ & $\pi(n)$ & $\pi_{\RVM}(n)$ & $\pi_{\RVM}(n)-\pi(n)$\\
\midrule
\endfirsthead
\multicolumn{12}{l}{\emph{Table~\thetable\ (continued)}}\\
\toprule
$n$ & $\pi(n)$ & $\pi_{\RVM}(n)$ & $\pi_{\RVM}(n)-\pi(n)$ & $n$ & $\pi(n)$ & $\pi_{\RVM}(n)$ & $\pi_{\RVM}(n)-\pi(n)$ & $n$ & $\pi(n)$ & $\pi_{\RVM}(n)$ & $\pi_{\RVM}(n)-\pi(n)$\\
\midrule
\endhead
\bottomrule
\endfoot
1 & 0 & 0.0000 & $+0.0000$ & 335 & 67 & 67.7510 & $+0.7510$ & 669 & 121 & 120.9426 & $-0.0574$ \\
2 & 1 & 0.5294 & $-0.4706$ & 336 & 67 & 67.9877 & $+0.9877$ & 670 & 121 & 120.9419 & $-0.0581$ \\
3 & 2 & 1.5713 & $-0.4287$ & 337 & 68 & 68.1700 & $+0.1700$ & 671 & 121 & 120.9184 & $-0.0816$ \\
4 & 2 & 1.9484 & $-0.0516$ & 338 & 68 & 68.3041 & $+0.3041$ & 672 & 121 & 120.8793 & $-0.1207$ \\
5 & 3 & 2.5746 & $-0.4254$ & 339 & 68 & 68.4015 & $+0.4015$ & 673 & 122 & 120.8330 & $-1.1670$ \\
6 & 3 & 2.9355 & $-0.0645$ & 340 & 68 & 68.4778 & $+0.4778$ & 674 & 122 & 120.7882 & $-1.2118$ \\
7 & 4 & 3.5681 & $-0.4319$ & 341 & 68 & 68.5504 & $+0.5504$ & 675 & 122 & 120.7532 & $-1.2468$ \\
8 & 4 & 4.1567 & $+0.1567$ & 342 & 68 & 68.6367 & $+0.6367$ & 676 & 122 & 120.7360 & $-1.2640$ \\
9 & 4 & 3.8153 & $-0.1847$ & 343 & 68 & 68.7516 & $+0.7516$ & 677 & 123 & 120.7434 & $-2.2566$ \\
10 & 4 & 3.9868 & $-0.0132$ & 344 & 68 & 68.9061 & $+0.9061$ & 678 & 123 & 120.7808 & $-2.2192$ \\
11 & 5 & 4.6846 & $-0.3154$ & 345 & 68 & 69.1059 & $+1.1059$ & 679 & 123 & 120.8520 & $-2.1480$ \\
12 & 5 & 5.0116 & $+0.0116$ & 346 & 68 & 69.3506 & $+1.3506$ & 680 & 123 & 120.9589 & $-2.0411$ \\
13 & 6 & 5.3739 & $-0.6261$ & 347 & 69 & 69.6339 & $+0.6339$ & 681 & 123 & 121.1014 & $-1.8986$ \\
14 & 6 & 5.9505 & $-0.0495$ & 348 & 69 & 69.9441 & $+0.9441$ & 682 & 123 & 121.2775 & $-1.7225$ \\
15 & 6 & 6.1472 & $+0.1472$ & 349 & 70 & 70.2659 & $+0.2659$ & 683 & 124 & 121.4835 & $-2.5165$ \\
16 & 6 & 6.0944 & $+0.0944$ & 350 & 70 & 70.5815 & $+0.5815$ & 684 & 124 & 121.7141 & $-2.2859$ \\
17 & 7 & 6.4940 & $-0.5060$ & 351 & 70 & 70.8732 & $+0.8732$ & 685 & 124 & 121.9627 & $-2.0373$ \\
18 & 7 & 7.2626 & $+0.2626$ & 352 & 70 & 71.1252 & $+1.1252$ & 686 & 124 & 122.2222 & $-1.7778$ \\
19 & 8 & 7.6976 & $-0.3024$ & 353 & 71 & 71.3251 & $+0.3251$ & 687 & 124 & 122.4848 & $-1.5152$ \\
20 & 8 & 7.6642 & $-0.3358$ & 354 & 71 & 71.4659 & $+0.4659$ & 688 & 124 & 122.7428 & $-1.2572$ \\
21 & 8 & 7.6943 & $-0.3057$ & 355 & 71 & 71.5463 & $+0.5463$ & 689 & 124 & 122.9893 & $-1.0107$ \\
22 & 8 & 8.0689 & $+0.0689$ & 356 & 71 & 71.5711 & $+0.5711$ & 690 & 124 & 123.2180 & $-0.7820$ \\
23 & 9 & 8.4717 & $-0.5283$ & 357 & 71 & 71.5508 & $+0.5508$ & 691 & 125 & 123.4240 & $-1.5760$ \\
24 & 9 & 8.5894 & $-0.4106$ & 358 & 71 & 71.5004 & $+0.5004$ & 692 & 125 & 123.6039 & $-1.3961$ \\
25 & 9 & 8.5995 & $-0.4005$ & 359 & 72 & 71.4378 & $-0.5622$ & 693 & 125 & 123.7561 & $-1.2439$ \\
26 & 9 & 8.8730 & $-0.1270$ & 360 & 72 & 71.3820 & $-0.6180$ & 694 & 125 & 123.8807 & $-1.1193$ \\
27 & 9 & 9.4280 & $+0.4280$ & 361 & 72 & 71.3515 & $-0.6485$ & 695 & 125 & 123.9798 & $-1.0202$ \\
28 & 9 & 9.9043 & $+0.9043$ & 362 & 72 & 71.3617 & $-0.6383$ & 696 & 125 & 124.0569 & $-0.9431$ \\
29 & 10 & 10.0301 & $+0.0301$ & 363 & 72 & 71.4240 & $-0.5760$ & 697 & 125 & 124.1170 & $-0.8830$ \\
30 & 10 & 9.9464 & $-0.0536$ & 364 & 72 & 71.5445 & $-0.4555$ & 698 & 125 & 124.1663 & $-0.8337$ \\
31 & 11 & 10.0197 & $-0.9803$ & 365 & 72 & 71.7236 & $-0.2764$ & 699 & 125 & 124.2118 & $-0.7882$ \\
32 & 11 & 10.3577 & $-0.6423$ & 366 & 72 & 71.9556 & $-0.0444$ & 700 & 125 & 124.2606 & $-0.7394$ \\
33 & 11 & 10.8394 & $-0.1606$ & 367 & 73 & 72.2304 & $-0.7696$ & 701 & 126 & 124.3201 & $-1.6799$ \\
34 & 11 & 11.1855 & $+0.1855$ & 368 & 73 & 72.5338 & $-0.4662$ & 702 & 126 & 124.3968 & $-1.6032$ \\
35 & 11 & 11.2497 & $+0.2497$ & 369 & 73 & 72.8493 & $-0.1507$ & 703 & 126 & 124.4967 & $-1.5033$ \\
36 & 11 & 11.2250 & $+0.2250$ & 370 & 73 & 73.1605 & $+0.1605$ & 704 & 126 & 124.6243 & $-1.3757$ \\
37 & 12 & 11.3795 & $-0.6205$ & 371 & 73 & 73.4518 & $+0.4518$ & 705 & 126 & 124.7827 & $-1.2173$ \\
38 & 12 & 11.7894 & $-0.2106$ & 372 & 73 & 73.7111 & $+0.7111$ & 706 & 126 & 124.9732 & $-1.0268$ \\
39 & 12 & 12.2857 & $+0.2857$ & 373 & 74 & 73.9302 & $-0.0698$ & 707 & 126 & 125.1956 & $-0.8044$ \\
40 & 12 & 12.6379 & $+0.6379$ & 374 & 74 & 74.1063 & $+0.1063$ & 708 & 126 & 125.4475 & $-0.5525$ \\
41 & 13 & 12.7782 & $-0.2218$ & 375 & 74 & 74.2415 & $+0.2415$ & 709 & 127 & 125.7252 & $-1.2748$ \\
42 & 13 & 12.8571 & $-0.1429$ & 376 & 74 & 74.3430 & $+0.3430$ & 710 & 127 & 126.0233 & $-0.9767$ \\
43 & 14 & 13.0980 & $-0.9020$ & 377 & 74 & 74.4222 & $+0.4222$ & 711 & 127 & 126.3354 & $-0.6646$ \\
44 & 14 & 13.5939 & $-0.4061$ & 378 & 74 & 74.4930 & $+0.4930$ & 712 & 127 & 126.6541 & $-0.3459$ \\
45 & 14 & 14.2234 & $+0.2234$ & 379 & 75 & 74.5707 & $-0.4293$ & 713 & 127 & 126.9717 & $-0.0283$ \\
46 & 14 & 14.7450 & $+0.7450$ & 380 & 75 & 74.6702 & $-0.3298$ & 714 & 127 & 127.2803 & $+0.2803$ \\
47 & 15 & 14.9810 & $-0.0190$ & 381 & 75 & 74.8040 & $-0.1960$ & 715 & 127 & 127.5724 & $+0.5724$ \\
48 & 15 & 14.9416 & $-0.0584$ & 382 & 75 & 74.9814 & $-0.0186$ & 716 & 127 & 127.8415 & $+0.8415$ \\
49 & 15 & 14.8016 & $-0.1984$ & 383 & 76 & 75.2070 & $-0.7930$ & 717 & 127 & 128.0819 & $+1.0819$ \\
50 & 15 & 14.7623 & $-0.2377$ & 384 & 76 & 75.4805 & $-0.5195$ & 718 & 127 & 128.2896 & $+1.2896$ \\
51 & 15 & 14.9108 & $-0.0892$ & 385 & 76 & 75.7963 & $-0.2037$ & 719 & 128 & 128.4621 & $+0.4621$ \\
52 & 15 & 15.1776 & $+0.1776$ & 386 & 76 & 76.1444 & $+0.1444$ & 720 & 128 & 128.5986 & $+0.5986$ \\
53 & 16 & 15.4113 & $-0.5887$ & 387 & 76 & 76.5110 & $+0.5110$ & 721 & 128 & 128.7002 & $+0.7002$ \\
54 & 16 & 15.5054 & $-0.4946$ & 388 & 76 & 76.8802 & $+0.8802$ & 722 & 128 & 128.7696 & $+0.7696$ \\
55 & 16 & 15.4828 & $-0.5172$ & 389 & 77 & 77.2352 & $+0.2352$ & 723 & 128 & 128.8111 & $+0.8111$ \\
56 & 16 & 15.4822 & $-0.5178$ & 390 & 77 & 77.5602 & $+0.5602$ & 724 & 128 & 128.8304 & $+0.8304$ \\
57 & 16 & 15.6626 & $-0.3374$ & 391 & 77 & 77.8417 & $+0.8417$ & 725 & 128 & 128.8341 & $+0.8341$ \\
58 & 16 & 16.0949 & $+0.0949$ & 392 & 77 & 78.0702 & $+1.0702$ & 726 & 128 & 128.8296 & $+0.8296$ \\
59 & 17 & 16.7128 & $-0.2872$ & 393 & 77 & 78.2405 & $+1.2405$ & 727 & 129 & 128.8243 & $-0.1757$ \\
60 & 17 & 17.3497 & $+0.3497$ & 394 & 77 & 78.3527 & $+1.3527$ & 728 & 129 & 128.8256 & $-0.1744$ \\
61 & 18 & 17.8352 & $-0.1648$ & 395 & 77 & 78.4121 & $+1.4121$ & 729 & 129 & 128.8406 & $-0.1594$ \\
62 & 18 & 18.0886 & $+0.0886$ & 396 & 77 & 78.4282 & $+1.4282$ & 730 & 129 & 128.8750 & $-0.1250$ \\
63 & 18 & 18.1544 & $+0.1544$ & 397 & 78 & 78.4144 & $+0.4144$ & 731 & 129 & 128.9338 & $-0.0662$ \\
64 & 18 & 18.2061 & $+0.2061$ & 398 & 78 & 78.3862 & $+0.3862$ & 732 & 129 & 129.0202 & $+0.0202$ \\
65 & 18 & 18.3456 & $+0.3456$ & 399 & 78 & 78.3599 & $+0.3599$ & 733 & 130 & 129.1361 & $-0.8639$ \\
66 & 18 & 18.5811 & $+0.5811$ & 400 & 78 & 78.3513 & $+0.3513$ & 734 & 130 & 129.2814 & $-0.7186$ \\
67 & 19 & 18.9195 & $-0.0805$ & 401 & 79 & 78.3736 & $-0.6264$ & 735 & 130 & 129.4546 & $-0.5454$ \\
68 & 19 & 19.2516 & $+0.2516$ & 402 & 79 & 78.4369 & $-0.5631$ & 736 & 130 & 129.6525 & $-0.3475$ \\
69 & 19 & 19.4681 & $+0.4681$ & 403 & 79 & 78.5468 & $-0.4532$ & 737 & 130 & 129.8707 & $-0.1293$ \\
70 & 19 & 19.5266 & $+0.5266$ & 404 & 79 & 78.7044 & $-0.2956$ & 738 & 130 & 130.1034 & $+0.1034$ \\
71 & 20 & 19.4766 & $-0.5234$ & 405 & 79 & 78.9059 & $-0.0941$ & 739 & 131 & 130.3444 & $-0.6556$ \\
72 & 20 & 19.4340 & $-0.5660$ & 406 & 79 & 79.1434 & $+0.1434$ & 740 & 131 & 130.5867 & $-0.4133$ \\
73 & 21 & 19.5193 & $-1.4807$ & 407 & 79 & 79.4054 & $+0.4054$ & 741 & 131 & 130.8234 & $-0.1766$ \\
74 & 21 & 19.7953 & $-1.2047$ & 408 & 79 & 79.6786 & $+0.6786$ & 742 & 131 & 131.0481 & $+0.0481$ \\
75 & 21 & 20.2355 & $-0.7645$ & 409 & 80 & 79.9487 & $-0.0513$ & 743 & 132 & 131.2547 & $-0.7453$ \\
76 & 21 & 20.7379 & $-0.2621$ & 410 & 80 & 80.2022 & $+0.2022$ & 744 & 132 & 131.4384 & $-0.5616$ \\
77 & 21 & 21.1750 & $+0.1750$ & 411 & 80 & 80.4273 & $+0.4273$ & 745 & 132 & 131.5956 & $-0.4044$ \\
78 & 21 & 21.4539 & $+0.4539$ & 412 & 80 & 80.6157 & $+0.6157$ & 746 & 132 & 131.7241 & $-0.2759$ \\
79 & 22 & 21.5573 & $-0.4427$ & 413 & 80 & 80.7629 & $+0.7629$ & 747 & 132 & 131.8234 & $-0.1766$ \\
80 & 22 & 21.5464 & $-0.4536$ & 414 & 80 & 80.8687 & $+0.8687$ & 748 & 132 & 131.8944 & $-0.1056$ \\
81 & 22 & 21.5287 & $-0.4713$ & 415 & 80 & 80.9376 & $+0.9376$ & 749 & 132 & 131.9398 & $-0.0602$ \\
82 & 22 & 21.6057 & $-0.3943$ & 416 & 80 & 80.9776 & $+0.9776$ & 750 & 132 & 131.9636 & $-0.0364$ \\
83 & 23 & 21.8275 & $-1.1725$ & 417 & 80 & 81.0004 & $+1.0004$ & 751 & 133 & 131.9709 & $-1.0291$ \\
84 & 23 & 22.1734 & $-0.8266$ & 418 & 80 & 81.0193 & $+1.0193$ & 752 & 133 & 131.9679 & $-1.0321$ \\
85 & 23 & 22.5664 & $-0.4336$ & 419 & 81 & 81.0486 & $+0.0486$ & 753 & 133 & 131.9614 & $-1.0386$ \\
86 & 23 & 22.9118 & $-0.0882$ & 420 & 81 & 81.1019 & $+0.1019$ & 754 & 133 & 131.9584 & $-1.0416$ \\
87 & 23 & 23.1427 & $+0.1427$ & 421 & 82 & 81.1911 & $-0.8089$ & 755 & 133 & 131.9658 & $-1.0342$ \\
88 & 23 & 23.2500 & $+0.2500$ & 422 & 82 & 81.3250 & $-0.6750$ & 756 & 133 & 131.9901 & $-1.0099$ \\
89 & 24 & 23.2860 & $-0.7140$ & 423 & 82 & 81.5083 & $-0.4917$ & 757 & 134 & 132.0369 & $-1.9631$ \\
90 & 24 & 23.3408 & $-0.6592$ & 424 & 82 & 81.7418 & $-0.2582$ & 758 & 134 & 132.1108 & $-1.8892$ \\
91 & 24 & 23.5027 & $-0.4973$ & 425 & 82 & 82.0216 & $+0.0216$ & 759 & 134 & 132.2151 & $-1.7849$ \\
92 & 24 & 23.8199 & $-0.1801$ & 426 & 82 & 82.3399 & $+0.3399$ & 760 & 134 & 132.3515 & $-1.6485$ \\
93 & 24 & 24.2803 & $+0.2803$ & 427 & 82 & 82.6854 & $+0.6854$ & 761 & 135 & 132.5203 & $-2.4797$ \\
94 & 24 & 24.8160 & $+0.8160$ & 428 & 82 & 83.0446 & $+1.0446$ & 762 & 135 & 132.7199 & $-2.2801$ \\
95 & 24 & 25.3300 & $+1.3300$ & 429 & 82 & 83.4027 & $+1.4027$ & 763 & 135 & 132.9477 & $-2.0523$ \\
96 & 24 & 25.7333 & $+1.7333$ & 430 & 82 & 83.7451 & $+1.7451$ & 764 & 135 & 133.1992 & $-1.8008$ \\
97 & 25 & 25.9782 & $+0.9782$ & 431 & 83 & 84.0586 & $+1.0586$ & 765 & 135 & 133.4692 & $-1.5308$ \\
98 & 25 & 26.0736 & $+1.0736$ & 432 & 83 & 84.3326 & $+1.3326$ & 766 & 135 & 133.7514 & $-1.2486$ \\
99 & 25 & 26.0787 & $+1.0787$ & 433 & 84 & 84.5599 & $+0.5599$ & 767 & 135 & 134.0391 & $-0.9609$ \\
100 & 25 & 26.0782 & $+1.0782$ & 434 & 84 & 84.7374 & $+0.7374$ & 768 & 135 & 134.3254 & $-0.6746$ \\
101 & 26 & 26.1488 & $+0.1488$ & 435 & 84 & 84.8664 & $+0.8664$ & 769 & 136 & 134.6035 & $-1.3965$ \\
102 & 26 & 26.3312 & $+0.3312$ & 436 & 84 & 84.9519 & $+0.9519$ & 770 & 136 & 134.8672 & $-1.1328$ \\
103 & 27 & 26.6166 & $-0.3834$ & 437 & 84 & 85.0029 & $+1.0029$ & 771 & 136 & 135.1109 & $-0.8891$ \\
104 & 27 & 26.9530 & $-0.0470$ & 438 & 84 & 85.0310 & $+1.0310$ & 772 & 136 & 135.3305 & $-0.6695$ \\
105 & 27 & 27.2672 & $+0.2672$ & 439 & 85 & 85.0493 & $+0.0493$ & 773 & 137 & 135.5229 & $-1.4771$ \\
106 & 27 & 27.4935 & $+0.4935$ & 440 & 85 & 85.0715 & $+0.0715$ & 774 & 137 & 135.6866 & $-1.3134$ \\
107 & 28 & 27.5995 & $-0.4005$ & 441 & 85 & 85.1106 & $+0.1106$ & 775 & 137 & 135.8215 & $-1.1785$ \\
108 & 28 & 27.5976 & $-0.4024$ & 442 & 85 & 85.1774 & $+0.1774$ & 776 & 137 & 135.9292 & $-1.0708$ \\
109 & 29 & 27.5415 & $-1.4585$ & 443 & 86 & 85.2801 & $-0.7199$ & 777 & 137 & 136.0126 & $-0.9874$ \\
110 & 29 & 27.5066 & $-1.4934$ & 444 & 86 & 85.4233 & $-0.5767$ & 778 & 137 & 136.0760 & $-0.9240$ \\
111 & 29 & 27.5640 & $-1.4360$ & 445 & 86 & 85.6073 & $-0.3927$ & 779 & 137 & 136.1246 & $-0.8754$ \\
112 & 29 & 27.7561 & $-1.2439$ & 446 & 86 & 85.8289 & $-0.1711$ & 780 & 137 & 136.1647 & $-0.8353$ \\
113 & 30 & 28.0824 & $-1.9176$ & 447 & 86 & 86.0810 & $+0.0810$ & 781 & 137 & 136.2027 & $-0.7973$ \\
114 & 30 & 28.5006 & $-1.4994$ & 448 & 86 & 86.3535 & $+0.3535$ & 782 & 137 & 136.2455 & $-0.7545$ \\
115 & 30 & 28.9409 & $-1.0591$ & 449 & 87 & 86.6343 & $-0.3657$ & 783 & 137 & 136.2996 & $-0.7004$ \\
116 & 30 & 29.3290 & $-0.6710$ & 450 & 87 & 86.9103 & $-0.0897$ & 784 & 137 & 136.3710 & $-0.6290$ \\
117 & 30 & 29.6096 & $-0.3904$ & 451 & 87 & 87.1685 & $+0.1685$ & 785 & 137 & 136.4649 & $-0.5351$ \\
118 & 30 & 29.7635 & $-0.2365$ & 452 & 87 & 87.3973 & $+0.3973$ & 786 & 137 & 136.5855 & $-0.4145$ \\
119 & 30 & 29.8115 & $-0.1885$ & 453 & 87 & 87.5874 & $+0.5874$ & 787 & 138 & 136.7357 & $-1.2643$ \\
120 & 30 & 29.8060 & $-0.1940$ & 454 & 87 & 87.7326 & $+0.7326$ & 788 & 138 & 136.9168 & $-1.0832$ \\
121 & 30 & 29.8131 & $-0.1869$ & 455 & 87 & 87.8302 & $+0.8302$ & 789 & 138 & 137.1288 & $-0.8712$ \\
122 & 30 & 29.8904 & $-0.1096$ & 456 & 87 & 87.8817 & $+0.8817$ & 790 & 138 & 137.3701 & $-0.6299$ \\
123 & 30 & 30.0688 & $+0.0688$ & 457 & 88 & 87.8919 & $-0.1081$ & 791 & 138 & 137.6378 & $-0.3622$ \\
124 & 30 & 30.3437 & $+0.3437$ & 458 & 88 & 87.8693 & $-0.1307$ & 792 & 138 & 137.9277 & $-0.0723$ \\
125 & 30 & 30.6772 & $+0.6772$ & 459 & 88 & 87.8247 & $-0.1753$ & 793 & 138 & 138.2344 & $+0.2344$ \\
126 & 30 & 31.0106 & $+1.0106$ & 460 & 88 & 87.7706 & $-0.2294$ & 794 & 138 & 138.5519 & $+0.5519$ \\
127 & 31 & 31.2839 & $+0.2839$ & 461 & 89 & 87.7202 & $-1.2798$ & 795 & 138 & 138.8737 & $+0.8737$ \\
128 & 31 & 31.4262 & $+0.4262$ & 462 & 89 & 87.6859 & $-1.3141$ & 796 & 138 & 139.1931 & $+1.1931$ \\
129 & 31 & 31.4309 & $+0.4309$ & 463 & 90 & 87.6789 & $-2.3211$ & 797 & 139 & 139.5033 & $+0.5033$ \\
130 & 31 & 31.3888 & $+0.3888$ & 464 & 90 & 87.7075 & $-2.2925$ & 798 & 139 & 139.7985 & $+0.7985$ \\
131 & 32 & 31.3009 & $-0.6991$ & 465 & 90 & 87.7772 & $-2.2228$ & 799 & 139 & 140.0731 & $+1.0731$ \\
132 & 32 & 31.2297 & $-0.7703$ & 466 & 90 & 87.8896 & $-2.1104$ & 800 & 139 & 140.3230 & $+1.3230$ \\
133 & 32 & 31.2332 & $-0.7668$ & 467 & 91 & 88.0428 & $-2.9572$ & 801 & 139 & 140.5448 & $+1.5448$ \\
134 & 32 & 31.3489 & $-0.6511$ & 468 & 91 & 88.2316 & $-2.7684$ & 802 & 139 & 140.7369 & $+1.7369$ \\
135 & 32 & 31.5837 & $-0.4163$ & 469 & 91 & 88.4474 & $-2.5526$ & 803 & 139 & 140.8989 & $+1.8989$ \\
136 & 32 & 31.9130 & $-0.0870$ & 470 & 91 & 88.6797 & $-2.3203$ & 804 & 139 & 141.0319 & $+2.0319$ \\
137 & 33 & 32.2892 & $-0.7108$ & 471 & 91 & 88.9165 & $-2.0835$ & 805 & 139 & 141.1383 & $+2.1383$ \\
138 & 33 & 32.6552 & $-0.3448$ & 472 & 91 & 89.1457 & $-1.8543$ & 806 & 139 & 141.2217 & $+2.2217$ \\
139 & 34 & 32.9609 & $-1.0391$ & 473 & 91 & 89.3559 & $-1.6441$ & 807 & 139 & 141.2868 & $+2.2868$ \\
140 & 34 & 33.1772 & $-0.8228$ & 474 & 91 & 89.5374 & $-1.4626$ & 808 & 139 & 141.3390 & $+2.3390$ \\
141 & 34 & 33.3035 & $-0.6965$ & 475 & 91 & 89.6832 & $-1.3168$ & 809 & 140 & 141.3844 & $+1.3844$ \\
142 & 34 & 33.3672 & $-0.6328$ & 476 & 91 & 89.7892 & $-1.2108$ & 810 & 140 & 141.4291 & $+1.4291$ \\
143 & 34 & 33.4154 & $-0.5846$ & 477 & 91 & 89.8551 & $-1.1449$ & 811 & 141 & 141.4791 & $+0.4791$ \\
144 & 34 & 33.5018 & $-0.4982$ & 478 & 91 & 89.8840 & $-1.1160$ & 812 & 141 & 141.5402 & $+0.5402$ \\
145 & 34 & 33.6717 & $-0.3283$ & 479 & 92 & 89.8825 & $-2.1175$ & 813 & 141 & 141.6172 & $+0.6172$ \\
146 & 34 & 33.9498 & $-0.0502$ & 480 & 92 & 89.8598 & $-2.1402$ & 814 & 141 & 141.7142 & $+0.7142$ \\
147 & 34 & 34.3339 & $+0.3339$ & 481 & 92 & 89.8272 & $-2.1728$ & 815 & 141 & 141.8341 & $+0.8341$ \\
148 & 34 & 34.7950 & $+0.7950$ & 482 & 92 & 89.7972 & $-2.2028$ & 816 & 141 & 141.9785 & $+0.9785$ \\
149 & 35 & 35.2855 & $+0.2855$ & 483 & 92 & 89.7826 & $-2.2174$ & 817 & 141 & 142.1478 & $+1.1478$ \\
150 & 35 & 35.7509 & $+0.7509$ & 484 & 92 & 89.7951 & $-2.2049$ & 818 & 141 & 142.3406 & $+1.3406$ \\
151 & 36 & 36.1439 & $+0.1439$ & 485 & 92 & 89.8448 & $-2.1552$ & 819 & 141 & 142.5547 & $+1.5547$ \\
152 & 36 & 36.4354 & $+0.4354$ & 486 & 92 & 89.9392 & $-2.0608$ & 820 & 141 & 142.7865 & $+1.7865$ \\
153 & 36 & 36.6214 & $+0.6214$ & 487 & 93 & 90.0827 & $-2.9173$ & 821 & 142 & 143.0312 & $+1.0312$ \\
154 & 36 & 36.7228 & $+0.7228$ & 488 & 93 & 90.2763 & $-2.7237$ & 822 & 142 & 143.2835 & $+1.2835$ \\
155 & 36 & 36.7794 & $+0.7794$ & 489 & 93 & 90.5175 & $-2.4825$ & 823 & 143 & 143.5374 & $+0.5374$ \\
156 & 36 & 36.8394 & $+0.8394$ & 490 & 93 & 90.8004 & $-2.1996$ & 824 & 143 & 143.7866 & $+0.7866$ \\
157 & 37 & 36.9470 & $-0.0530$ & 491 & 94 & 91.1164 & $-2.8836$ & 825 & 143 & 144.0251 & $+1.0251$ \\
158 & 37 & 37.1310 & $+0.1310$ & 492 & 94 & 91.4550 & $-2.5450$ & 826 & 143 & 144.2471 & $+1.2471$ \\
159 & 37 & 37.3987 & $+0.3987$ & 493 & 94 & 91.8042 & $-2.1958$ & 827 & 144 & 144.4474 & $+0.4474$ \\
160 & 37 & 37.7339 & $+0.7339$ & 494 & 94 & 92.1517 & $-1.8483$ & 828 & 144 & 144.6216 & $+0.6216$ \\
161 & 37 & 38.1016 & $+1.1016$ & 495 & 94 & 92.4862 & $-1.5138$ & 829 & 145 & 144.7668 & $-0.2332$ \\
162 & 37 & 38.4570 & $+1.4570$ & 496 & 94 & 92.7974 & $-1.2026$ & 830 & 145 & 144.8808 & $-0.1192$ \\
163 & 38 & 38.7565 & $+0.7565$ & 497 & 94 & 93.0777 & $-0.9223$ & 831 & 145 & 144.9633 & $-0.0367$ \\
164 & 38 & 38.9685 & $+0.9685$ & 498 & 94 & 93.3222 & $-0.6778$ & 832 & 145 & 145.0148 & $+0.0148$ \\
165 & 38 & 39.0808 & $+1.0808$ & 499 & 95 & 93.5290 & $-1.4710$ & 833 & 145 & 145.0377 & $+0.0377$ \\
166 & 38 & 39.1038 & $+1.1038$ & 500 & 95 & 93.6998 & $-1.3002$ & 834 & 145 & 145.0353 & $+0.0353$ \\
167 & 39 & 39.0679 & $+0.0679$ & 501 & 95 & 93.8391 & $-1.1609$ & 835 & 145 & 145.0122 & $+0.0122$ \\
168 & 39 & 39.0165 & $+0.0165$ & 502 & 95 & 93.9540 & $-1.0460$ & 836 & 145 & 144.9737 & $-0.0263$ \\
169 & 39 & 38.9962 & $-0.0038$ & 503 & 96 & 94.0538 & $-1.9462$ & 837 & 145 & 144.9260 & $-0.0740$ \\
170 & 39 & 39.0462 & $+0.0462$ & 504 & 96 & 94.1488 & $-1.8512$ & 838 & 145 & 144.8755 & $-0.1245$ \\
171 & 39 & 39.1901 & $+0.1901$ & 505 & 96 & 94.2499 & $-1.7501$ & 839 & 146 & 144.8288 & $-1.1712$ \\
172 & 39 & 39.4304 & $+0.4304$ & 506 & 96 & 94.3672 & $-1.6328$ & 840 & 146 & 144.7922 & $-1.2078$ \\
173 & 40 & 39.7489 & $-0.2511$ & 507 & 96 & 94.5096 & $-1.4904$ & 841 & 146 & 144.7718 & $-1.2282$ \\
174 & 40 & 40.1109 & $+0.1109$ & 508 & 96 & 94.6838 & $-1.3162$ & 842 & 146 & 144.7726 & $-1.2274$ \\
175 & 40 & 40.4729 & $+0.4729$ & 509 & 97 & 94.8939 & $-2.1061$ & 843 & 146 & 144.7989 & $-1.2011$ \\
176 & 40 & 40.7926 & $+0.7926$ & 510 & 97 & 95.1411 & $-1.8589$ & 844 & 146 & 144.8537 & $-1.1463$ \\
177 & 40 & 41.0375 & $+1.0375$ & 511 & 97 & 95.4231 & $-1.5769$ & 845 & 146 & 144.9389 & $-1.0611$ \\
178 & 40 & 41.1922 & $+1.1922$ & 512 & 97 & 95.7351 & $-1.2649$ & 846 & 146 & 145.0550 & $-0.9450$ \\
179 & 41 & 41.2604 & $+0.2604$ & 513 & 97 & 96.0695 & $-0.9305$ & 847 & 146 & 145.2012 & $-0.7988$ \\
180 & 41 & 41.2640 & $+0.2640$ & 514 & 97 & 96.4166 & $-0.5834$ & 848 & 146 & 145.3753 & $-0.6247$ \\
181 & 42 & 41.2376 & $-0.7624$ & 515 & 97 & 96.7654 & $-0.2346$ & 849 & 146 & 145.5742 & $-0.4258$ \\
182 & 42 & 41.2205 & $-0.7795$ & 516 & 97 & 97.1045 & $+0.1045$ & 850 & 146 & 145.7935 & $-0.2065$ \\
183 & 42 & 41.2481 & $-0.7519$ & 517 & 97 & 97.4230 & $+0.4230$ & 851 & 146 & 146.0283 & $+0.0283$ \\
184 & 42 & 41.3444 & $-0.6556$ & 518 & 97 & 97.7110 & $+0.7110$ & 852 & 146 & 146.2729 & $+0.2729$ \\
185 & 42 & 41.5167 & $-0.4833$ & 519 & 97 & 97.9605 & $+0.9605$ & 853 & 147 & 146.5216 & $-0.4784$ \\
186 & 42 & 41.7547 & $-0.2453$ & 520 & 97 & 98.1664 & $+1.1664$ & 854 & 147 & 146.7686 & $-0.2314$ \\
187 & 42 & 42.0325 & $+0.0325$ & 521 & 98 & 98.3262 & $+0.3262$ & 855 & 147 & 147.0082 & $+0.0082$ \\
188 & 42 & 42.3150 & $+0.3150$ & 522 & 98 & 98.4406 & $+0.4406$ & 856 & 147 & 147.2356 & $+0.2356$ \\
189 & 42 & 42.5651 & $+0.5651$ & 523 & 99 & 98.5135 & $-0.4865$ & 857 & 148 & 147.4463 & $-0.5537$ \\
190 & 42 & 42.7517 & $+0.7517$ & 524 & 99 & 98.5515 & $-0.4485$ & 858 & 148 & 147.6373 & $-0.3627$ \\
191 & 43 & 42.8568 & $-0.1432$ & 525 & 99 & 98.5636 & $-0.4364$ & 859 & 149 & 147.8063 & $-1.1937$ \\
192 & 43 & 42.8788 & $-0.1212$ & 526 & 99 & 98.5603 & $-0.4397$ & 860 & 149 & 147.9524 & $-1.0476$ \\
193 & 44 & 42.8337 & $-1.1663$ & 527 & 99 & 98.5533 & $-0.4467$ & 861 & 149 & 148.0759 & $-0.9241$ \\
194 & 44 & 42.7518 & $-1.2482$ & 528 & 99 & 98.5540 & $-0.4460$ & 862 & 149 & 148.1783 & $-0.8217$ \\
195 & 44 & 42.6720 & $-1.3280$ & 529 & 99 & 98.5734 & $-0.4266$ & 863 & 150 & 148.2622 & $-1.7378$ \\
196 & 44 & 42.6344 & $-1.3656$ & 530 & 99 & 98.6209 & $-0.3791$ & 864 & 150 & 148.3314 & $-1.6686$ \\
197 & 45 & 42.6731 & $-2.3269$ & 531 & 99 & 98.7038 & $-0.2962$ & 865 & 150 & 148.3903 & $-1.6097$ \\
198 & 45 & 42.8094 & $-2.1906$ & 532 & 99 & 98.8266 & $-0.1734$ & 866 & 150 & 148.4441 & $-1.5559$ \\
199 & 46 & 43.0489 & $-2.9511$ & 533 & 99 & 98.9911 & $-0.0089$ & 867 & 150 & 148.4983 & $-1.5017$ \\
200 & 46 & 43.3800 & $-2.6200$ & 534 & 99 & 99.1960 & $+0.1960$ & 868 & 150 & 148.5584 & $-1.4416$ \\
201 & 46 & 43.7772 & $-2.2228$ & 535 & 99 & 99.4373 & $+0.4373$ & 869 & 150 & 148.6301 & $-1.3699$ \\
202 & 46 & 44.2058 & $-1.7942$ & 536 & 99 & 99.7085 & $+0.7085$ & 870 & 150 & 148.7183 & $-1.2817$ \\
203 & 46 & 44.6290 & $-1.3710$ & 537 & 99 & 100.0010 & $+1.0010$ & 871 & 150 & 148.8274 & $-1.1726$ \\
204 & 46 & 45.0143 & $-0.9857$ & 538 & 99 & 100.3050 & $+1.3050$ & 872 & 150 & 148.9612 & $-1.0388$ \\
205 & 46 & 45.3398 & $-0.6602$ & 539 & 99 & 100.6101 & $+1.6101$ & 873 & 150 & 149.1220 & $-0.8780$ \\
206 & 46 & 45.5972 & $-0.4028$ & 540 & 99 & 100.9064 & $+1.9064$ & 874 & 150 & 149.3112 & $-0.6888$ \\
207 & 46 & 45.7933 & $-0.2067$ & 541 & 100 & 101.1845 & $+1.1845$ & 875 & 150 & 149.5292 & $-0.4708$ \\
208 & 46 & 45.9475 & $-0.0525$ & 542 & 100 & 101.4371 & $+1.4371$ & 876 & 150 & 149.7748 & $-0.2252$ \\
209 & 46 & 46.0878 & $+0.0878$ & 543 & 100 & 101.6588 & $+1.6588$ & 877 & 151 & 150.0456 & $-0.9544$ \\
210 & 46 & 46.2445 & $+0.2445$ & 544 & 100 & 101.8469 & $+1.8469$ & 878 & 151 & 150.3385 & $-0.6615$ \\
211 & 47 & 46.4441 & $-0.5559$ & 545 & 100 & 102.0013 & $+2.0013$ & 879 & 151 & 150.6491 & $-0.3509$ \\
212 & 47 & 46.7039 & $-0.2961$ & 546 & 100 & 102.1247 & $+2.1247$ & 880 & 151 & 150.9722 & $-0.0278$ \\
213 & 47 & 47.0284 & $+0.0284$ & 547 & 101 & 102.2224 & $+1.2224$ & 881 & 152 & 151.3022 & $-0.6978$ \\
214 & 47 & 47.4081 & $+0.4081$ & 548 & 101 & 102.3015 & $+1.3015$ & 882 & 152 & 151.6331 & $-0.3669$ \\
215 & 47 & 47.8212 & $+0.8212$ & 549 & 101 & 102.3709 & $+1.3709$ & 883 & 153 & 151.9590 & $-1.0410$ \\
216 & 47 & 48.2370 & $+1.2370$ & 550 & 101 & 102.4403 & $+1.4403$ & 884 & 153 & 152.2739 & $-0.7261$ \\
217 & 47 & 48.6217 & $+1.6217$ & 551 & 101 & 102.5195 & $+1.5195$ & 885 & 153 & 152.5726 & $-0.4274$ \\
218 & 47 & 48.9439 & $+1.9439$ & 552 & 101 & 102.6176 & $+1.6176$ & 886 & 153 & 152.8503 & $-0.1497$ \\
219 & 47 & 49.1808 & $+2.1808$ & 553 & 101 & 102.7427 & $+1.7427$ & 887 & 154 & 153.1032 & $-0.8968$ \\
220 & 47 & 49.3215 & $+2.3215$ & 554 & 101 & 102.9009 & $+1.9009$ & 888 & 154 & 153.3286 & $-0.6714$ \\
221 & 47 & 49.3692 & $+2.3692$ & 555 & 101 & 103.0958 & $+2.0958$ & 889 & 154 & 153.5247 & $-0.4753$ \\
222 & 47 & 49.3409 & $+2.3409$ & 556 & 101 & 103.3287 & $+2.3287$ & 890 & 154 & 153.6913 & $-0.3087$ \\
223 & 48 & 49.2643 & $+1.2643$ & 557 & 102 & 103.5980 & $+1.5980$ & 891 & 154 & 153.8292 & $-0.1708$ \\
224 & 48 & 49.1732 & $+1.1732$ & 558 & 102 & 103.8993 & $+1.8993$ & 892 & 154 & 153.9403 & $-0.0597$ \\
225 & 48 & 49.1020 & $+1.1020$ & 559 & 102 & 104.2263 & $+2.2263$ & 893 & 154 & 154.0277 & $+0.0277$ \\
226 & 48 & 49.0805 & $+1.0805$ & 560 & 102 & 104.5702 & $+2.5702$ & 894 & 154 & 154.0955 & $+0.0955$ \\
227 & 49 & 49.1289 & $+0.1289$ & 561 & 102 & 104.9213 & $+2.9213$ & 895 & 154 & 154.1485 & $+0.1485$ \\
228 & 49 & 49.2550 & $+0.2550$ & 562 & 102 & 105.2691 & $+3.2691$ & 896 & 154 & 154.1920 & $+0.1920$ \\
229 & 50 & 49.4539 & $-0.5461$ & 563 & 103 & 105.6029 & $+2.6029$ & 897 & 154 & 154.2316 & $+0.2316$ \\
230 & 50 & 49.7087 & $-0.2913$ & 564 & 103 & 105.9130 & $+2.9130$ & 898 & 154 & 154.2732 & $+0.2732$ \\
231 & 50 & 49.9947 & $-0.0053$ & 565 & 103 & 106.1909 & $+3.1909$ & 899 & 154 & 154.3223 & $+0.3223$ \\
232 & 50 & 50.2833 & $+0.2833$ & 566 & 103 & 106.4302 & $+3.4302$ & 900 & 154 & 154.3839 & $+0.3839$ \\
233 & 51 & 50.5477 & $-0.4523$ & 567 & 103 & 106.6267 & $+3.6267$ & 901 & 154 & 154.4626 & $+0.4626$ \\
234 & 51 & 50.7670 & $-0.2330$ & 568 & 103 & 106.7790 & $+3.7790$ & 902 & 154 & 154.5620 & $+0.5620$ \\
235 & 51 & 50.9301 & $-0.0699$ & 569 & 104 & 106.8882 & $+2.8882$ & 903 & 154 & 154.6848 & $+0.6848$ \\
236 & 51 & 51.0374 & $+0.0374$ & 570 & 104 & 106.9580 & $+2.9580$ & 904 & 154 & 154.8326 & $+0.8326$ \\
237 & 51 & 51.1004 & $+0.1004$ & 571 & 105 & 106.9947 & $+1.9947$ & 905 & 154 & 155.0058 & $+1.0058$ \\
238 & 51 & 51.1399 & $+0.1399$ & 572 & 105 & 107.0061 & $+2.0061$ & 906 & 154 & 155.2036 & $+1.2036$ \\
239 & 52 & 51.1820 & $-0.8180$ & 573 & 105 & 107.0016 & $+2.0016$ & 907 & 155 & 155.4241 & $+0.4241$ \\
240 & 52 & 51.2538 & $-0.7462$ & 574 & 105 & 106.9912 & $+1.9912$ & 908 & 155 & 155.6646 & $+0.6646$ \\
241 & 53 & 51.3786 & $-1.6214$ & 575 & 105 & 106.9849 & $+1.9849$ & 909 & 155 & 155.9211 & $+0.9211$ \\
242 & 53 & 51.5719 & $-1.4281$ & 576 & 105 & 106.9921 & $+1.9921$ & 910 & 155 & 156.1892 & $+1.1892$ \\
243 & 53 & 51.8390 & $-1.1610$ & 577 & 106 & 107.0208 & $+1.0208$ & 911 & 156 & 156.4638 & $+0.4638$ \\
244 & 53 & 52.1735 & $-0.8265$ & 578 & 106 & 107.0774 & $+1.0774$ & 912 & 156 & 156.7396 & $+0.7396$ \\
245 & 53 & 52.5584 & $-0.4416$ & 579 & 106 & 107.1659 & $+1.1659$ & 913 & 156 & 157.0111 & $+1.0111$ \\
246 & 53 & 52.9686 & $-0.0314$ & 580 & 106 & 107.2879 & $+1.2879$ & 914 & 156 & 157.2732 & $+1.2732$ \\
247 & 53 & 53.3748 & $+0.3748$ & 581 & 106 & 107.4425 & $+1.4425$ & 915 & 156 & 157.5210 & $+1.5210$ \\
248 & 53 & 53.7477 & $+0.7477$ & 582 & 106 & 107.6265 & $+1.6265$ & 916 & 156 & 157.7502 & $+1.7502$ \\
249 & 53 & 54.0624 & $+1.0624$ & 583 & 106 & 107.8343 & $+1.8343$ & 917 & 156 & 157.9576 & $+1.9576$ \\
250 & 53 & 54.3022 & $+1.3022$ & 584 & 106 & 108.0587 & $+2.0587$ & 918 & 156 & 158.1405 & $+2.1405$ \\
251 & 54 & 54.4611 & $+0.4611$ & 585 & 106 & 108.2913 & $+2.2913$ & 919 & 157 & 158.2976 & $+1.2976$ \\
252 & 54 & 54.5441 & $+0.5441$ & 586 & 106 & 108.5228 & $+2.5228$ & 920 & 157 & 158.4286 & $+1.4286$ \\
253 & 54 & 54.5667 & $+0.5667$ & 587 & 107 & 108.7442 & $+1.7442$ & 921 & 157 & 158.5342 & $+1.5342$ \\
254 & 54 & 54.5521 & $+0.5521$ & 588 & 107 & 108.9470 & $+1.9470$ & 922 & 157 & 158.6165 & $+1.6165$ \\
255 & 54 & 54.5279 & $+0.5279$ & 589 & 107 & 109.1239 & $+2.1239$ & 923 & 157 & 158.6782 & $+1.6782$ \\
256 & 54 & 54.5215 & $+0.5215$ & 590 & 107 & 109.2693 & $+2.2693$ & 924 & 157 & 158.7231 & $+1.7231$ \\
257 & 55 & 54.5565 & $-0.4435$ & 591 & 107 & 109.3797 & $+2.3797$ & 925 & 157 & 158.7558 & $+1.7558$ \\
258 & 55 & 54.6491 & $-0.3509$ & 592 & 107 & 109.4541 & $+2.4541$ & 926 & 157 & 158.7812 & $+1.7812$ \\
259 & 55 & 54.8061 & $-0.1939$ & 593 & 108 & 109.4938 & $+1.4938$ & 927 & 157 & 158.8048 & $+1.8048$ \\
260 & 55 & 55.0238 & $+0.0238$ & 594 & 108 & 109.5023 & $+1.5023$ & 928 & 157 & 158.8318 & $+1.8318$ \\
261 & 55 & 55.2892 & $+0.2892$ & 595 & 108 & 109.4855 & $+1.4855$ & 929 & 158 & 158.8676 & $+0.8676$ \\
262 & 55 & 55.5823 & $+0.5823$ & 596 & 108 & 109.4508 & $+1.4508$ & 930 & 158 & 158.9170 & $+0.9170$ \\
263 & 56 & 55.8792 & $-0.1208$ & 597 & 108 & 109.4068 & $+1.4068$ & 931 & 158 & 158.9844 & $+0.9844$ \\
264 & 56 & 56.1558 & $+0.1558$ & 598 & 108 & 109.3630 & $+1.3630$ & 932 & 158 & 159.0733 & $+1.0733$ \\
265 & 56 & 56.3917 & $+0.3917$ & 599 & 109 & 109.3287 & $+0.3287$ & 933 & 158 & 159.1864 & $+1.1864$ \\
266 & 56 & 56.5732 & $+0.5732$ & 600 & 109 & 109.3130 & $+0.3130$ & 934 & 158 & 159.3253 & $+1.3253$ \\
267 & 56 & 56.6956 & $+0.6956$ & 601 & 110 & 109.3234 & $-0.6766$ & 935 & 158 & 159.4905 & $+1.4905$ \\
268 & 56 & 56.7634 & $+0.7634$ & 602 & 110 & 109.3662 & $-0.6338$ & 936 & 158 & 159.6815 & $+1.6815$ \\
269 & 57 & 56.7900 & $-0.2100$ & 603 & 110 & 109.4455 & $-0.5545$ & 937 & 159 & 159.8968 & $+0.8968$ \\
270 & 57 & 56.7958 & $-0.2042$ & 604 & 110 & 109.5632 & $-0.4368$ & 938 & 159 & 160.1335 & $+1.1335$ \\
271 & 58 & 56.8046 & $-1.1954$ & 605 & 110 & 109.7189 & $-0.2811$ & 939 & 159 & 160.3883 & $+1.3883$ \\
272 & 58 & 56.8408 & $-1.1592$ & 606 & 110 & 109.9096 & $-0.0904$ & 940 & 159 & 160.6569 & $+1.6569$ \\
273 & 58 & 56.9253 & $-1.0747$ & 607 & 111 & 110.1306 & $-0.8694$ & 941 & 160 & 160.9344 & $+0.9344$ \\
274 & 58 & 57.0729 & $-0.9271$ & 608 & 111 & 110.3750 & $-0.6250$ & 942 & 160 & 161.2157 & $+1.2157$ \\
275 & 58 & 57.2897 & $-0.7103$ & 609 & 111 & 110.6348 & $-0.3652$ & 943 & 160 & 161.4955 & $+1.4955$ \\
276 & 58 & 57.5728 & $-0.4272$ & 610 & 111 & 110.9012 & $-0.0988$ & 944 & 160 & 161.7684 & $+1.7684$ \\
277 & 59 & 57.9101 & $-1.0899$ & 611 & 111 & 111.1651 & $+0.1651$ & 945 & 160 & 162.0296 & $+2.0296$ \\
278 & 59 & 58.2824 & $-0.7176$ & 612 & 111 & 111.4176 & $+0.4176$ & 946 & 160 & 162.2746 & $+2.2746$ \\
279 & 59 & 58.6658 & $-0.3342$ & 613 & 112 & 111.6510 & $-0.3490$ & 947 & 161 & 162.4997 & $+1.4997$ \\
280 & 59 & 59.0351 & $+0.0351$ & 614 & 112 & 111.8587 & $-0.1413$ & 948 & 161 & 162.7019 & $+1.7019$ \\
281 & 60 & 59.3670 & $-0.6330$ & 615 & 112 & 112.0362 & $+0.0362$ & 949 & 161 & 162.8794 & $+1.8794$ \\
282 & 60 & 59.6434 & $-0.3566$ & 616 & 112 & 112.1810 & $+0.1810$ & 950 & 161 & 163.0313 & $+2.0313$ \\
283 & 61 & 59.8537 & $-1.1463$ & 617 & 113 & 112.2929 & $-0.7071$ & 951 & 161 & 163.1576 & $+2.1576$ \\
284 & 61 & 59.9959 & $-1.0041$ & 618 & 113 & 112.3739 & $-0.6261$ & 952 & 161 & 163.2597 & $+2.2597$ \\
285 & 61 & 60.0769 & $-0.9231$ & 619 & 114 & 112.4281 & $-1.5719$ & 953 & 162 & 163.3398 & $+1.3398$ \\
286 & 61 & 60.1114 & $-0.8886$ & 620 & 114 & 112.4616 & $-1.5384$ & 954 & 162 & 163.4011 & $+1.4011$ \\
287 & 61 & 60.1194 & $-0.8806$ & 621 & 114 & 112.4819 & $-1.5181$ & 955 & 162 & 163.4476 & $+1.4476$ \\
288 & 61 & 60.1240 & $-0.8760$ & 622 & 114 & 112.4975 & $-1.5025$ & 956 & 162 & 163.4838 & $+1.4838$ \\
289 & 61 & 60.1476 & $-0.8524$ & 623 & 114 & 112.5174 & $-1.4826$ & 957 & 162 & 163.5149 & $+1.5149$ \\
290 & 61 & 60.2092 & $-0.7908$ & 624 & 114 & 112.5504 & $-1.4496$ & 958 & 162 & 163.5459 & $+1.5459$ \\
291 & 61 & 60.3219 & $-0.6781$ & 625 & 114 & 112.6047 & $-1.3953$ & 959 & 162 & 163.5822 & $+1.5822$ \\
292 & 61 & 60.4908 & $-0.5092$ & 626 & 114 & 112.6872 & $-1.3128$ & 960 & 162 & 163.6289 & $+1.6289$ \\
293 & 62 & 60.7132 & $-1.2868$ & 627 & 114 & 112.8035 & $-1.1965$ & 961 & 162 & 163.6906 & $+1.6906$ \\
294 & 62 & 60.9782 & $-1.0218$ & 628 & 114 & 112.9568 & $-1.0432$ & 962 & 162 & 163.7714 & $+1.7714$ \\
295 & 62 & 61.2689 & $-0.7311$ & 629 & 114 & 113.1486 & $-0.8514$ & 963 & 162 & 163.8746 & $+1.8746$ \\
296 & 62 & 61.5644 & $-0.4356$ & 630 & 114 & 113.3778 & $-0.6222$ & 964 & 162 & 164.0025 & $+2.0025$ \\
297 & 62 & 61.8426 & $-0.1574$ & 631 & 115 & 113.6414 & $-1.3586$ & 965 & 162 & 164.1567 & $+2.1567$ \\
298 & 62 & 62.0837 & $+0.0837$ & 632 & 115 & 113.9343 & $-1.0657$ & 966 & 162 & 164.3375 & $+2.3375$ \\
299 & 62 & 62.2721 & $+0.2721$ & 633 & 115 & 114.2497 & $-0.7503$ & 967 & 163 & 164.5442 & $+1.5442$ \\
300 & 62 & 62.3988 & $+0.3988$ & 634 & 115 & 114.5798 & $-0.4202$ & 968 & 163 & 164.7752 & $+1.7752$ \\
301 & 62 & 62.4627 & $+0.4627$ & 635 & 115 & 114.9157 & $-0.0843$ & 969 & 163 & 165.0277 & $+2.0277$ \\
302 & 62 & 62.4706 & $+0.4706$ & 636 & 115 & 115.2485 & $+0.2485$ & 970 & 163 & 165.2984 & $+2.2984$ \\
303 & 62 & 62.4359 & $+0.4359$ & 637 & 115 & 115.5696 & $+0.5696$ & 971 & 164 & 165.5829 & $+1.5829$ \\
304 & 62 & 62.3775 & $+0.3775$ & 638 & 115 & 115.8712 & $+0.8712$ & 972 & 164 & 165.8767 & $+1.8767$ \\
305 & 62 & 62.3171 & $+0.3171$ & 639 & 115 & 116.1468 & $+1.1468$ & 973 & 164 & 166.1745 & $+2.1745$ \\
306 & 62 & 62.2761 & $+0.2761$ & 640 & 115 & 116.3915 & $+1.3915$ & 974 & 164 & 166.4713 & $+2.4713$ \\
307 & 63 & 62.2737 & $-0.7263$ & 641 & 116 & 116.6025 & $+0.6025$ & 975 & 164 & 166.7618 & $+2.7618$ \\
308 & 63 & 62.3237 & $-0.6763$ & 642 & 116 & 116.7789 & $+0.7789$ & 976 & 164 & 167.0413 & $+3.0413$ \\
309 & 63 & 62.4333 & $-0.5667$ & 643 & 117 & 116.9220 & $-0.0780$ & 977 & 165 & 167.3052 & $+2.3052$ \\
310 & 63 & 62.6023 & $-0.3977$ & 644 & 117 & 117.0351 & $+0.0351$ & 978 & 165 & 167.5498 & $+2.5498$ \\
311 & 64 & 62.8231 & $-1.1769$ & 645 & 117 & 117.1231 & $+0.1231$ & 979 & 165 & 167.7721 & $+2.7721$ \\
312 & 64 & 63.0819 & $-0.9181$ & 646 & 117 & 117.1927 & $+0.1927$ & 980 & 165 & 167.9700 & $+2.9700$ \\
313 & 65 & 63.3604 & $-1.6396$ & 647 & 118 & 117.2512 & $-0.7488$ & 981 & 165 & 168.1423 & $+3.1423$ \\
314 & 65 & 63.6384 & $-1.3616$ & 648 & 118 & 117.3067 & $-0.6933$ & 982 & 165 & 168.2889 & $+3.2889$ \\
315 & 65 & 63.8963 & $-1.1037$ & 649 & 118 & 117.3671 & $-0.6329$ & 983 & 166 & 168.4107 & $+2.4107$ \\
316 & 65 & 64.1173 & $-0.8827$ & 650 & 118 & 117.4400 & $-0.5600$ & 984 & 166 & 168.5093 & $+2.5093$ \\
317 & 66 & 64.2901 & $-1.7099$ & 651 & 118 & 117.5319 & $-0.4681$ & 985 & 166 & 168.5876 & $+2.5876$ \\
318 & 66 & 64.4097 & $-1.5903$ & 652 & 118 & 117.6482 & $-0.3518$ & 986 & 166 & 168.6489 & $+2.6489$ \\
319 & 66 & 64.4786 & $-1.5214$ & 653 & 119 & 117.7922 & $-1.2078$ & 987 & 166 & 168.6972 & $+2.6972$ \\
320 & 66 & 64.5058 & $-1.4942$ & 654 & 119 & 117.9655 & $-1.0345$ & 988 & 166 & 168.7371 & $+2.7371$ \\
321 & 66 & 64.5063 & $-1.4937$ & 655 & 119 & 118.1678 & $-0.8322$ & 989 & 166 & 168.7732 & $+2.7732$ \\
322 & 66 & 64.4990 & $-1.5010$ & 656 & 119 & 118.3967 & $-0.6033$ & 990 & 166 & 168.8103 & $+2.8103$ \\
323 & 66 & 64.5045 & $-1.4955$ & 657 & 119 & 118.6479 & $-0.3521$ & 991 & 167 & 168.8530 & $+1.8530$ \\
324 & 66 & 64.5423 & $-1.4577$ & 658 & 119 & 118.9155 & $-0.0845$ & 992 & 167 & 168.9057 & $+1.9057$ \\
325 & 66 & 64.6291 & $-1.3709$ & 659 & 120 & 119.1926 & $-0.8074$ & 993 & 167 & 168.9722 & $+1.9722$ \\
326 & 66 & 64.7761 & $-1.2239$ & 660 & 120 & 119.4712 & $-0.5288$ & 994 & 167 & 169.0555 & $+2.0555$ \\
327 & 66 & 64.9883 & $-1.0117$ & 661 & 121 & 119.7432 & $-1.2568$ & 995 & 167 & 169.1579 & $+2.1579$ \\
328 & 66 & 65.2637 & $-0.7363$ & 662 & 121 & 120.0004 & $-0.9996$ & 996 & 167 & 169.2810 & $+2.2810$ \\
329 & 66 & 65.5934 & $-0.4066$ & 663 & 121 & 120.2355 & $-0.7645$ & 997 & 168 & 169.4252 & $+1.4252$ \\
330 & 66 & 65.9631 & $-0.0369$ & 664 & 121 & 120.4422 & $-0.5578$ & 998 & 168 & 169.5901 & $+1.5901$ \\
331 & 67 & 66.3543 & $-0.6457$ & 665 & 121 & 120.6154 & $-0.3846$ & 999 & 168 & 169.7741 & $+1.7741$ \\
332 & 67 & 66.7469 & $-0.2531$ & 666 & 121 & 120.7523 & $-0.2477$ & 1000 & 168 & 169.9750 & $+1.9750$ \\
333 & 67 & 67.1211 & $+0.1211$ & 667 & 121 & 120.8515 & $-0.1485$ &  &  &  &  \\
334 & 67 & 67.4600 & $+0.4600$ & 668 & 121 & 120.9140 & $-0.0860$ &  &  &  &  \\
\end{longtable}}

{\scriptsize\setlength{\tabcolsep}{2.5pt}
\begin{longtable}{rrrrr@{\quad}rrrrr@{\quad}rrrrr}
\caption{The M\"obius function $\mu(n)$ and the Mertens function $M(n)$, exact (sieve), and the critical RVM Mertens function $M_{\RVM}^{\crit}(n)$, approximate (the tapered explicit formula of program \texttt{s05} with the residues computed from the product for $k\le1000$ and modelled beyond), at the integers $n\le1000$, with the difference $M_{\RVM}^{\crit}(n)-M(n)$.}\label{tab:mertens-1000}\\
\toprule
$n$ & $\mu$ & $M$ & $M_{\RVM}^{\crit}$ & \shortstack[r]{$M_{\RVM}^{\crit}(n)$\\$-M(n)$} & $n$ & $\mu$ & $M$ & $M_{\RVM}^{\crit}$ & \shortstack[r]{$M_{\RVM}^{\crit}(n)$\\$-M(n)$} & $n$ & $\mu$ & $M$ & $M_{\RVM}^{\crit}$ & \shortstack[r]{$M_{\RVM}^{\crit}(n)$\\$-M(n)$}\\
\midrule
\endfirsthead
\multicolumn{15}{l}{\emph{Table~\thetable\ (continued)}}\\
\toprule
$n$ & $\mu$ & $M$ & $M_{\RVM}^{\crit}$ & \shortstack[r]{$M_{\RVM}^{\crit}(n)$\\$-M(n)$} & $n$ & $\mu$ & $M$ & $M_{\RVM}^{\crit}$ & \shortstack[r]{$M_{\RVM}^{\crit}(n)$\\$-M(n)$} & $n$ & $\mu$ & $M$ & $M_{\RVM}^{\crit}$ & \shortstack[r]{$M_{\RVM}^{\crit}(n)$\\$-M(n)$}\\
\midrule
\endhead
\bottomrule
\endfoot
1 & $1$ & $1$ & $0.2925$ & $-0.7075$ & 335 & $1$ & $2$ & $-1.9107$ & $-3.9107$ & 669 & $1$ & $-10$ & $-1.6301$ & $+8.3699$ \\
2 & $-1$ & $0$ & $0.4563$ & $+0.4563$ & 336 & $0$ & $2$ & $-2.2151$ & $-4.2151$ & 670 & $-1$ & $-11$ & $-1.2832$ & $+9.7168$ \\
3 & $-1$ & $-1$ & $-0.6523$ & $+0.3477$ & 337 & $-1$ & $1$ & $-2.3953$ & $-3.3953$ & 671 & $1$ & $-10$ & $-0.8918$ & $+9.1082$ \\
4 & $0$ & $-1$ & $-0.8720$ & $+0.1280$ & 338 & $0$ & $1$ & $-2.4640$ & $-3.4640$ & 672 & $0$ & $-10$ & $-0.4714$ & $+9.5286$ \\
5 & $-1$ & $-2$ & $-1.7255$ & $+0.2745$ & 339 & $1$ & $2$ & $-2.4456$ & $-4.4456$ & 673 & $-1$ & $-11$ & $-0.0388$ & $+10.9612$ \\
6 & $1$ & $-1$ & $-1.3063$ & $-0.3063$ & 340 & $0$ & $2$ & $-2.3727$ & $-4.3727$ & 674 & $1$ & $-10$ & $0.3882$ & $+10.3882$ \\
7 & $-1$ & $-2$ & $-1.5931$ & $+0.4069$ & 341 & $1$ & $3$ & $-2.2821$ & $-5.2821$ & 675 & $0$ & $-10$ & $0.7923$ & $+10.7923$ \\
8 & $0$ & $-2$ & $-2.4029$ & $-0.4029$ & 342 & $0$ & $3$ & $-2.2108$ & $-5.2108$ & 676 & $0$ & $-10$ & $1.1575$ & $+11.1575$ \\
9 & $0$ & $-2$ & $-1.3292$ & $+0.6708$ & 343 & $0$ & $3$ & $-2.1911$ & $-5.1911$ & 677 & $-1$ & $-11$ & $1.4699$ & $+12.4699$ \\
10 & $1$ & $-1$ & $-1.3280$ & $-0.3280$ & 344 & $0$ & $3$ & $-2.2470$ & $-5.2470$ & 678 & $-1$ & $-12$ & $1.7185$ & $+13.7185$ \\
11 & $-1$ & $-2$ & $-2.3272$ & $-0.3272$ & 345 & $-1$ & $2$ & $-2.3914$ & $-4.3914$ & 679 & $1$ & $-11$ & $1.8958$ & $+12.8958$ \\
12 & $0$ & $-2$ & $-2.2067$ & $-0.2067$ & 346 & $1$ & $3$ & $-2.6249$ & $-5.6249$ & 680 & $0$ & $-11$ & $1.9982$ & $+12.9982$ \\
13 & $-1$ & $-3$ & $-1.8468$ & $+1.1532$ & 347 & $-1$ & $2$ & $-2.9356$ & $-4.9356$ & 681 & $1$ & $-10$ & $2.0263$ & $+12.0263$ \\
14 & $1$ & $-2$ & $-1.9324$ & $+0.0676$ & 348 & $0$ & $2$ & $-3.3007$ & $-5.3007$ & 682 & $-1$ & $-11$ & $1.9844$ & $+12.9844$ \\
15 & $1$ & $-1$ & $-1.5435$ & $-0.5435$ & 349 & $-1$ & $1$ & $-3.6892$ & $-4.6892$ & 683 & $-1$ & $-12$ & $1.8806$ & $+13.8806$ \\
16 & $0$ & $-1$ & $-1.0414$ & $-0.0414$ & 350 & $0$ & $1$ & $-4.0654$ & $-5.0654$ & 684 & $0$ & $-12$ & $1.7261$ & $+13.7261$ \\
17 & $-1$ & $-2$ & $-1.6970$ & $+0.3030$ & 351 & $0$ & $1$ & $-4.3933$ & $-5.3933$ & 685 & $1$ & $-11$ & $1.5347$ & $+12.5347$ \\
18 & $0$ & $-2$ & $-3.0095$ & $-1.0095$ & 352 & $0$ & $1$ & $-4.6406$ & $-5.6406$ & 686 & $0$ & $-11$ & $1.3217$ & $+12.3217$ \\
19 & $-1$ & $-3$ & $-3.2869$ & $-0.2869$ & 353 & $-1$ & $0$ & $-4.7824$ & $-4.7824$ & 687 & $1$ & $-10$ & $1.1032$ & $+11.1032$ \\
20 & $0$ & $-3$ & $-2.2604$ & $+0.7396$ & 354 & $-1$ & $-1$ & $-4.8043$ & $-3.8043$ & 688 & $0$ & $-10$ & $0.8952$ & $+10.8952$ \\
21 & $1$ & $-2$ & $-1.2515$ & $+0.7485$ & 355 & $1$ & $0$ & $-4.7038$ & $-4.7038$ & 689 & $1$ & $-9$ & $0.7125$ & $+9.7125$ \\
22 & $1$ & $-1$ & $-1.1120$ & $-0.1120$ & 356 & $0$ & $0$ & $-4.4909$ & $-4.4909$ & 690 & $1$ & $-8$ & $0.5678$ & $+8.5678$ \\
23 & $-1$ & $-2$ & $-1.3607$ & $+0.6393$ & 357 & $-1$ & $-1$ & $-4.1872$ & $-3.1872$ & 691 & $-1$ & $-9$ & $0.4714$ & $+9.4714$ \\
24 & $0$ & $-2$ & $-1.3669$ & $+0.6331$ & 358 & $1$ & $0$ & $-3.8234$ & $-3.8234$ & 692 & $0$ & $-9$ & $0.4302$ & $+9.4302$ \\
25 & $0$ & $-2$ & $-1.3810$ & $+0.6190$ & 359 & $-1$ & $-1$ & $-3.4364$ & $-2.4364$ & 693 & $0$ & $-9$ & $0.4474$ & $+9.4474$ \\
26 & $1$ & $-1$ & $-1.9805$ & $-0.9805$ & 360 & $0$ & $-1$ & $-3.0652$ & $-2.0652$ & 694 & $1$ & $-8$ & $0.5225$ & $+8.5225$ \\
27 & $0$ & $-1$ & $-3.0141$ & $-2.0141$ & 361 & $0$ & $-1$ & $-2.7470$ & $-1.7470$ & 695 & $1$ & $-7$ & $0.6512$ & $+7.6512$ \\
28 & $0$ & $-1$ & $-3.6155$ & $-2.6155$ & 362 & $1$ & $0$ & $-2.5134$ & $-2.5134$ & 696 & $0$ & $-7$ & $0.8259$ & $+7.8259$ \\
29 & $-1$ & $-2$ & $-3.2066$ & $-1.2066$ & 363 & $0$ & $0$ & $-2.3872$ & $-2.3872$ & 697 & $1$ & $-6$ & $1.0357$ & $+7.0357$ \\
30 & $-1$ & $-3$ & $-2.1599$ & $+0.8401$ & 364 & $0$ & $0$ & $-2.3801$ & $-2.3801$ & 698 & $1$ & $-5$ & $1.2675$ & $+6.2675$ \\
31 & $-1$ & $-4$ & $-1.3649$ & $+2.6351$ & 365 & $1$ & $1$ & $-2.4917$ & $-3.4917$ & 699 & $1$ & $-4$ & $1.5065$ & $+5.5065$ \\
32 & $0$ & $-4$ & $-1.2973$ & $+2.7027$ & 366 & $-1$ & $0$ & $-2.7101$ & $-2.7101$ & 700 & $0$ & $-4$ & $1.7374$ & $+5.7374$ \\
33 & $1$ & $-3$ & $-1.6435$ & $+1.3565$ & 367 & $-1$ & $-1$ & $-3.0129$ & $-2.0129$ & 701 & $-1$ & $-5$ & $1.9448$ & $+6.9448$ \\
34 & $1$ & $-2$ & $-1.7783$ & $+0.2217$ & 368 & $0$ & $-1$ & $-3.3700$ & $-2.3700$ & 702 & $0$ & $-5$ & $2.1143$ & $+7.1143$ \\
35 & $1$ & $-1$ & $-1.4711$ & $-0.4711$ & 369 & $0$ & $-1$ & $-3.7470$ & $-2.7470$ & 703 & $1$ & $-4$ & $2.2335$ & $+6.2335$ \\
36 & $0$ & $-1$ & $-1.0723$ & $-0.0723$ & 370 & $-1$ & $-2$ & $-4.1086$ & $-2.1086$ & 704 & $0$ & $-4$ & $2.2926$ & $+6.2926$ \\
37 & $-1$ & $-2$ & $-1.0774$ & $+0.9226$ & 371 & $1$ & $-1$ & $-4.4228$ & $-3.4228$ & 705 & $-1$ & $-5$ & $2.2849$ & $+7.2849$ \\
38 & $1$ & $-1$ & $-1.5897$ & $-0.5897$ & 372 & $0$ & $-1$ & $-4.6640$ & $-3.6640$ & 706 & $1$ & $-4$ & $2.2071$ & $+6.2071$ \\
39 & $1$ & $0$ & $-2.2227$ & $-2.2227$ & 373 & $-1$ & $-2$ & $-4.8153$ & $-2.8153$ & 707 & $1$ & $-3$ & $2.0598$ & $+5.0598$ \\
40 & $0$ & $0$ & $-2.4881$ & $-2.4881$ & 374 & $-1$ & $-3$ & $-4.8708$ & $-1.8708$ & 708 & $0$ & $-3$ & $1.8473$ & $+4.8473$ \\
41 & $-1$ & $-1$ & $-2.2590$ & $-1.2590$ & 375 & $0$ & $-3$ & $-4.8354$ & $-1.8354$ & 709 & $-1$ & $-4$ & $1.5774$ & $+5.5774$ \\
42 & $-1$ & $-2$ & $-1.8775$ & $+0.1225$ & 376 & $0$ & $-3$ & $-4.7247$ & $-1.7247$ & 710 & $-1$ & $-5$ & $1.2609$ & $+6.2609$ \\
43 & $-1$ & $-3$ & $-1.8384$ & $+1.1616$ & 377 & $1$ & $-2$ & $-4.5625$ & $-2.5625$ & 711 & $0$ & $-5$ & $0.9113$ & $+5.9113$ \\
44 & $0$ & $-3$ & $-2.3559$ & $+0.6441$ & 378 & $0$ & $-2$ & $-4.3787$ & $-2.3787$ & 712 & $0$ & $-5$ & $0.5438$ & $+5.5438$ \\
45 & $0$ & $-3$ & $-3.1822$ & $-0.1822$ & 379 & $-1$ & $-3$ & $-4.2057$ & $-1.2057$ & 713 & $1$ & $-4$ & $0.1745$ & $+4.1745$ \\
46 & $1$ & $-2$ & $-3.8050$ & $-1.8050$ & 380 & $0$ & $-3$ & $-4.0751$ & $-1.0751$ & 714 & $1$ & $-3$ & $-0.1803$ & $+2.8197$ \\
47 & $-1$ & $-3$ & $-3.8344$ & $-0.8344$ & 381 & $1$ & $-2$ & $-4.0139$ & $-2.0139$ & 715 & $-1$ & $-4$ & $-0.5051$ & $+3.4949$ \\
48 & $0$ & $-3$ & $-3.2670$ & $-0.2670$ & 382 & $1$ & $-1$ & $-4.0420$ & $-3.0420$ & 716 & $0$ & $-4$ & $-0.7856$ & $+3.2144$ \\
49 & $0$ & $-3$ & $-2.4450$ & $+0.5550$ & 383 & $-1$ & $-2$ & $-4.1698$ & $-2.1698$ & 717 & $1$ & $-3$ & $-1.0104$ & $+1.9896$ \\
50 & $0$ & $-3$ & $-1.7747$ & $+1.2253$ & 384 & $0$ & $-2$ & $-4.3972$ & $-2.3972$ & 718 & $1$ & $-2$ & $-1.1706$ & $+0.8294$ \\
51 & $1$ & $-2$ & $-1.4374$ & $+0.5626$ & 385 & $-1$ & $-3$ & $-4.7131$ & $-1.7131$ & 719 & $-1$ & $-3$ & $-1.2606$ & $+1.7394$ \\
52 & $0$ & $-2$ & $-1.3016$ & $+0.6984$ & 386 & $1$ & $-2$ & $-5.0972$ & $-3.0972$ & 720 & $0$ & $-3$ & $-1.2785$ & $+1.7215$ \\
53 & $-1$ & $-3$ & $-1.0780$ & $+1.9220$ & 387 & $0$ & $-2$ & $-5.5213$ & $-3.5213$ & 721 & $1$ & $-2$ & $-1.2263$ & $+0.7737$ \\
54 & $0$ & $-3$ & $-0.5793$ & $+2.4207$ & 388 & $0$ & $-2$ & $-5.9522$ & $-3.9522$ & 722 & $0$ & $-2$ & $-1.1091$ & $+0.8909$ \\
55 & $1$ & $-2$ & $0.1096$ & $+2.1096$ & 389 & $-1$ & $-3$ & $-6.3553$ & $-3.3553$ & 723 & $1$ & $-1$ & $-0.9355$ & $+0.0645$ \\
56 & $0$ & $-2$ & $0.6666$ & $+2.6666$ & 390 & $1$ & $-2$ & $-6.6975$ & $-4.6975$ & 724 & $0$ & $-1$ & $-0.7168$ & $+0.2832$ \\
57 & $1$ & $-1$ & $0.7366$ & $+1.7366$ & 391 & $1$ & $-1$ & $-6.9507$ & $-5.9507$ & 725 & $0$ & $-1$ & $-0.4661$ & $+0.5339$ \\
58 & $1$ & $0$ & $0.1610$ & $+0.1610$ & 392 & $0$ & $-1$ & $-7.0945$ & $-6.0945$ & 726 & $0$ & $-1$ & $-0.1982$ & $+0.8018$ \\
59 & $-1$ & $-1$ & $-0.9166$ & $+0.0834$ & 393 & $1$ & $0$ & $-7.1179$ & $-7.1179$ & 727 & $-1$ & $-2$ & $0.0717$ & $+2.0717$ \\
60 & $0$ & $-1$ & $-2.1293$ & $-1.1293$ & 394 & $1$ & $1$ & $-7.0205$ & $-8.0205$ & 728 & $0$ & $-2$ & $0.3288$ & $+2.3288$ \\
61 & $-1$ & $-2$ & $-3.0912$ & $-1.0912$ & 395 & $1$ & $2$ & $-6.8123$ & $-8.8123$ & 729 & $0$ & $-2$ & $0.5592$ & $+2.5592$ \\
62 & $1$ & $-1$ & $-3.5946$ & $-2.5946$ & 396 & $0$ & $2$ & $-6.5126$ & $-8.5126$ & 730 & $-1$ & $-3$ & $0.7512$ & $+3.7512$ \\
63 & $0$ & $-1$ & $-3.6905$ & $-2.6905$ & 397 & $-1$ & $1$ & $-6.1482$ & $-7.1482$ & 731 & $1$ & $-2$ & $0.8951$ & $+2.8951$ \\
64 & $0$ & $-1$ & $-3.6144$ & $-2.6144$ & 398 & $1$ & $2$ & $-5.7506$ & $-7.7506$ & 732 & $0$ & $-2$ & $0.9845$ & $+2.9845$ \\
65 & $1$ & $0$ & $-3.6166$ & $-3.6166$ & 399 & $-1$ & $1$ & $-5.3528$ & $-6.3528$ & 733 & $-1$ & $-3$ & $1.0163$ & $+4.0163$ \\
66 & $-1$ & $-1$ & $-3.8046$ & $-2.8046$ & 400 & $0$ & $1$ & $-4.9865$ & $-5.9865$ & 734 & $1$ & $-2$ & $0.9906$ & $+2.9906$ \\
67 & $-1$ & $-2$ & $-4.0854$ & $-2.0854$ & 401 & $-1$ & $0$ & $-4.6786$ & $-4.6786$ & 735 & $0$ & $-2$ & $0.9111$ & $+2.9111$ \\
68 & $0$ & $-2$ & $-4.2334$ & $-2.2334$ & 402 & $-1$ & $-1$ & $-4.4491$ & $-3.4491$ & 736 & $0$ & $-2$ & $0.7846$ & $+2.7846$ \\
69 & $1$ & $-1$ & $-4.0355$ & $-3.0355$ & 403 & $1$ & $0$ & $-4.3091$ & $-4.3091$ & 737 & $1$ & $-1$ & $0.6209$ & $+1.6209$ \\
70 & $-1$ & $-2$ & $-3.4264$ & $-1.4264$ & 404 & $0$ & $0$ & $-4.2600$ & $-4.2600$ & 738 & $0$ & $-1$ & $0.4316$ & $+1.4316$ \\
71 & $-1$ & $-3$ & $-2.5386$ & $+0.4614$ & 405 & $0$ & $0$ & $-4.2936$ & $-4.2936$ & 739 & $-1$ & $-2$ & $0.2304$ & $+2.2304$ \\
72 & $0$ & $-3$ & $-1.6446$ & $+1.3554$ & 406 & $-1$ & $-1$ & $-4.3929$ & $-3.3929$ & 740 & $0$ & $-2$ & $0.0314$ & $+2.0314$ \\
73 & $-1$ & $-4$ & $-1.0258$ & $+2.9742$ & 407 & $1$ & $0$ & $-4.5343$ & $-4.5343$ & 741 & $-1$ & $-3$ & $-0.1510$ & $+2.8490$ \\
74 & $1$ & $-3$ & $-0.8384$ & $+2.1616$ & 408 & $0$ & $0$ & $-4.6898$ & $-4.6898$ & 742 & $-1$ & $-4$ & $-0.3032$ & $+3.6968$ \\
75 & $0$ & $-3$ & $-1.0461$ & $+1.9539$ & 409 & $-1$ & $-1$ & $-4.8299$ & $-3.8299$ & 743 & $-1$ & $-5$ & $-0.4131$ & $+4.5869$ \\
76 & $0$ & $-3$ & $-1.4474$ & $+1.5526$ & 410 & $-1$ & $-2$ & $-4.9265$ & $-2.9265$ & 744 & $0$ & $-5$ & $-0.4709$ & $+4.5291$ \\
77 & $1$ & $-2$ & $-1.7796$ & $+0.2204$ & 411 & $1$ & $-1$ & $-4.9563$ & $-3.9563$ & 745 & $1$ & $-4$ & $-0.4691$ & $+3.5309$ \\
78 & $-1$ & $-3$ & $-1.8446$ & $+1.1554$ & 412 & $0$ & $-1$ & $-4.9022$ & $-3.9022$ & 746 & $1$ & $-3$ & $-0.4038$ & $+2.5962$ \\
79 & $-1$ & $-4$ & $-1.5955$ & $+2.4045$ & 413 & $1$ & $0$ & $-4.7557$ & $-4.7557$ & 747 & $0$ & $-3$ & $-0.2740$ & $+2.7260$ \\
80 & $0$ & $-4$ & $-1.1456$ & $+2.8544$ & 414 & $0$ & $0$ & $-4.5177$ & $-4.5177$ & 748 & $0$ & $-3$ & $-0.0824$ & $+2.9176$ \\
81 & $0$ & $-4$ & $-0.7025$ & $+3.2975$ & 415 & $1$ & $1$ & $-4.1981$ & $-5.1981$ & 749 & $1$ & $-2$ & $0.1652$ & $+2.1652$ \\
82 & $1$ & $-3$ & $-0.4619$ & $+2.5381$ & 416 & $0$ & $1$ & $-3.8154$ & $-4.8154$ & 750 & $0$ & $-2$ & $0.4599$ & $+2.4599$ \\
83 & $-1$ & $-4$ & $-0.5149$ & $+3.4851$ & 417 & $1$ & $2$ & $-3.3945$ & $-5.3945$ & 751 & $-1$ & $-3$ & $0.7902$ & $+3.7902$ \\
84 & $0$ & $-4$ & $-0.8092$ & $+3.1908$ & 418 & $-1$ & $1$ & $-2.9648$ & $-3.9648$ & 752 & $0$ & $-3$ & $1.1429$ & $+4.1429$ \\
85 & $1$ & $-3$ & $-1.1801$ & $+1.8199$ & 419 & $-1$ & $0$ & $-2.5574$ & $-2.5574$ & 753 & $1$ & $-2$ & $1.5031$ & $+3.5031$ \\
86 & $1$ & $-2$ & $-1.4316$ & $+0.5684$ & 420 & $0$ & $0$ & $-2.2018$ & $-2.2018$ & 754 & $-1$ & $-3$ & $1.8555$ & $+4.8555$ \\
87 & $1$ & $-1$ & $-1.4295$ & $-0.4295$ & 421 & $-1$ & $-1$ & $-1.9238$ & $-0.9238$ & 755 & $1$ & $-2$ & $2.1852$ & $+4.1852$ \\
88 & $0$ & $-1$ & $-1.1643$ & $-0.1643$ & 422 & $1$ & $0$ & $-1.7429$ & $-1.7429$ & 756 & $0$ & $-2$ & $2.4781$ & $+4.4781$ \\
89 & $-1$ & $-2$ & $-0.7568$ & $+1.2432$ & 423 & $0$ & $0$ & $-1.6704$ & $-1.6704$ & 757 & $-1$ & $-3$ & $2.7216$ & $+5.7216$ \\
90 & $0$ & $-2$ & $-0.4069$ & $+1.5931$ & 424 & $0$ & $0$ & $-1.7084$ & $-1.7084$ & 758 & $1$ & $-2$ & $2.9057$ & $+4.9057$ \\
91 & $1$ & $-1$ & $-0.3096$ & $+0.6904$ & 425 & $0$ & $0$ & $-1.8499$ & $-1.8499$ & 759 & $-1$ & $-3$ & $3.0229$ & $+6.0229$ \\
92 & $0$ & $-1$ & $-0.5747$ & $+0.4253$ & 426 & $-1$ & $-1$ & $-2.0792$ & $-1.0792$ & 760 & $0$ & $-3$ & $3.0689$ & $+6.0689$ \\
93 & $1$ & $0$ & $-1.1827$ & $-1.1827$ & 427 & $1$ & $0$ & $-2.3734$ & $-2.3734$ & 761 & $-1$ & $-4$ & $3.0429$ & $+7.0429$ \\
94 & $1$ & $1$ & $-1.9934$ & $-2.9934$ & 428 & $0$ & $0$ & $-2.7045$ & $-2.7045$ & 762 & $-1$ & $-5$ & $2.9470$ & $+7.9470$ \\
95 & $1$ & $2$ & $-2.8017$ & $-4.8017$ & 429 & $-1$ & $-1$ & $-3.0420$ & $-2.0420$ & 763 & $1$ & $-4$ & $2.7868$ & $+6.7868$ \\
96 & $0$ & $2$ & $-3.4154$ & $-5.4154$ & 430 & $-1$ & $-2$ & $-3.3553$ & $-1.3553$ & 764 & $0$ & $-4$ & $2.5708$ & $+6.5708$ \\
97 & $-1$ & $1$ & $-3.7240$ & $-4.7240$ & 431 & $-1$ & $-3$ & $-3.6166$ & $-0.6166$ & 765 & $0$ & $-4$ & $2.3096$ & $+6.3096$ \\
98 & $0$ & $1$ & $-3.7325$ & $-4.7325$ & 432 & $0$ & $-3$ & $-3.8034$ & $-0.8034$ & 766 & $1$ & $-3$ & $2.0159$ & $+5.0159$ \\
99 & $0$ & $1$ & $-3.5497$ & $-4.5497$ & 433 & $-1$ & $-4$ & $-3.9000$ & $+0.1000$ & 767 & $1$ & $-2$ & $1.7034$ & $+3.7034$ \\
100 & $0$ & $1$ & $-3.3365$ & $-4.3365$ & 434 & $-1$ & $-5$ & $-3.8991$ & $+1.1009$ & 768 & $0$ & $-2$ & $1.3864$ & $+3.3864$ \\
101 & $-1$ & $0$ & $-3.2390$ & $-3.2390$ & 435 & $-1$ & $-6$ & $-3.8019$ & $+2.1981$ & 769 & $-1$ & $-3$ & $1.0788$ & $+4.0788$ \\
102 & $-1$ & $-1$ & $-3.3304$ & $-2.3304$ & 436 & $0$ & $-6$ & $-3.6182$ & $+2.3818$ & 770 & $1$ & $-2$ & $0.7935$ & $+2.7935$ \\
103 & $-1$ & $-2$ & $-3.5847$ & $-1.5847$ & 437 & $1$ & $-5$ & $-3.3650$ & $+1.6350$ & 771 & $1$ & $-1$ & $0.5419$ & $+1.5419$ \\
104 & $0$ & $-2$ & $-3.8897$ & $-1.8897$ & 438 & $-1$ & $-6$ & $-3.0650$ & $+2.9350$ & 772 & $0$ & $-1$ & $0.3331$ & $+1.3331$ \\
105 & $-1$ & $-3$ & $-4.0934$ & $-1.0934$ & 439 & $-1$ & $-7$ & $-2.7442$ & $+4.2558$ & 773 & $-1$ & $-2$ & $0.1733$ & $+2.1733$ \\
106 & $1$ & $-2$ & $-4.0637$ & $-2.0637$ & 440 & $0$ & $-7$ & $-2.4294$ & $+4.5706$ & 774 & $0$ & $-2$ & $0.0662$ & $+2.0662$ \\
107 & $-1$ & $-3$ & $-3.7411$ & $-0.7411$ & 441 & $0$ & $-7$ & $-2.1460$ & $+4.8540$ & 775 & $0$ & $-2$ & $0.0121$ & $+2.0121$ \\
108 & $0$ & $-3$ & $-3.1633$ & $-0.1633$ & 442 & $-1$ & $-8$ & $-1.9153$ & $+6.0847$ & 776 & $0$ & $-2$ & $0.0081$ & $+2.0081$ \\
109 & $-1$ & $-4$ & $-2.4549$ & $+1.5451$ & 443 & $-1$ & $-9$ & $-1.7527$ & $+7.2473$ & 777 & $-1$ & $-3$ & $0.0486$ & $+3.0486$ \\
110 & $-1$ & $-5$ & $-1.7869$ & $+3.2131$ & 444 & $0$ & $-9$ & $-1.6662$ & $+7.3338$ & 778 & $1$ & $-2$ & $0.1253$ & $+2.1253$ \\
111 & $1$ & $-4$ & $-1.3206$ & $+2.6794$ & 445 & $1$ & $-8$ & $-1.6557$ & $+6.3443$ & 779 & $1$ & $-1$ & $0.2275$ & $+1.2275$ \\
112 & $0$ & $-4$ & $-1.1564$ & $+2.8436$ & 446 & $1$ & $-7$ & $-1.7130$ & $+5.2870$ & 780 & $0$ & $-1$ & $0.3432$ & $+1.3432$ \\
113 & $-1$ & $-5$ & $-1.3030$ & $+3.6970$ & 447 & $1$ & $-6$ & $-1.8225$ & $+4.1775$ & 781 & $1$ & $0$ & $0.4593$ & $+0.4593$ \\
114 & $-1$ & $-6$ & $-1.6789$ & $+4.3211$ & 448 & $0$ & $-6$ & $-1.9625$ & $+4.0375$ & 782 & $-1$ & $-1$ & $0.5624$ & $+1.5624$ \\
115 & $1$ & $-5$ & $-2.1415$ & $+2.8585$ & 449 & $-1$ & $-7$ & $-2.1074$ & $+4.8926$ & 783 & $0$ & $-1$ & $0.6396$ & $+1.6396$ \\
116 & $0$ & $-5$ & $-2.5354$ & $+2.4646$ & 450 & $0$ & $-7$ & $-2.2295$ & $+4.7705$ & 784 & $0$ & $-1$ & $0.6791$ & $+1.6791$ \\
117 & $0$ & $-5$ & $-2.7417$ & $+2.2583$ & 451 & $1$ & $-6$ & $-2.3017$ & $+3.6983$ & 785 & $1$ & $0$ & $0.6706$ & $+0.6706$ \\
118 & $1$ & $-4$ & $-2.7129$ & $+1.2871$ & 452 & $0$ & $-6$ & $-2.3001$ & $+3.6999$ & 786 & $-1$ & $-1$ & $0.6064$ & $+1.6064$ \\
119 & $1$ & $-3$ & $-2.4825$ & $+0.5175$ & 453 & $1$ & $-5$ & $-2.2054$ & $+2.7946$ & 787 & $-1$ & $-2$ & $0.4812$ & $+2.4812$ \\
120 & $0$ & $-3$ & $-2.1486$ & $+0.8514$ & 454 & $1$ & $-4$ & $-2.0055$ & $+1.9945$ & 788 & $0$ & $-2$ & $0.2928$ & $+2.2928$ \\
121 & $0$ & $-3$ & $-1.8366$ & $+1.1634$ & 455 & $-1$ & $-5$ & $-1.6958$ & $+3.3042$ & 789 & $1$ & $-1$ & $0.0420$ & $+1.0420$ \\
122 & $1$ & $-2$ & $-1.6547$ & $+0.3453$ & 456 & $0$ & $-5$ & $-1.2799$ & $+3.7201$ & 790 & $-1$ & $-2$ & $-0.2672$ & $+1.7328$ \\
123 & $1$ & $-1$ & $-1.6571$ & $-0.6571$ & 457 & $-1$ & $-6$ & $-0.7695$ & $+5.2305$ & 791 & $1$ & $-1$ & $-0.6282$ & $+0.3718$ \\
124 & $0$ & $-1$ & $-1.8256$ & $-0.8256$ & 458 & $1$ & $-5$ & $-0.1834$ & $+4.8166$ & 792 & $0$ & $-1$ & $-1.0314$ & $-0.0314$ \\
125 & $0$ & $-1$ & $-2.0757$ & $-1.0757$ & 459 & $0$ & $-5$ & $0.4544$ & $+5.4544$ & 793 & $1$ & $0$ & $-1.4653$ & $-1.4653$ \\
126 & $0$ & $-1$ & $-2.2829$ & $-1.2829$ & 460 & $0$ & $-5$ & $1.1161$ & $+6.1161$ & 794 & $1$ & $1$ & $-1.9169$ & $-2.9169$ \\
127 & $-1$ & $-2$ & $-2.3227$ & $-0.3227$ & 461 & $-1$ & $-6$ & $1.7729$ & $+7.7729$ & 795 & $-1$ & $0$ & $-2.3720$ & $-2.3720$ \\
128 & $0$ & $-2$ & $-2.1094$ & $-0.1094$ & 462 & $1$ & $-5$ & $2.3966$ & $+7.3966$ & 796 & $0$ & $0$ & $-2.8161$ & $-2.8161$ \\
129 & $1$ & $-1$ & $-1.6233$ & $-0.6233$ & 463 & $-1$ & $-6$ & $2.9627$ & $+8.9627$ & 797 & $-1$ & $-1$ & $-3.2354$ & $-2.2354$ \\
130 & $-1$ & $-2$ & $-0.9186$ & $+1.0814$ & 464 & $0$ & $-6$ & $3.4515$ & $+9.4515$ & 798 & $1$ & $0$ & $-3.6166$ & $-3.6166$ \\
131 & $-1$ & $-3$ & $-0.1091$ & $+2.8909$ & 465 & $-1$ & $-7$ & $3.8501$ & $+10.8501$ & 799 & $1$ & $1$ & $-3.9486$ & $-4.9486$ \\
132 & $0$ & $-3$ & $0.6624$ & $+3.6624$ & 466 & $1$ & $-6$ & $4.1531$ & $+10.1531$ & 800 & $0$ & $1$ & $-4.2220$ & $-5.2220$ \\
133 & $1$ & $-2$ & $1.2621$ & $+3.2621$ & 467 & $-1$ & $-7$ & $4.3627$ & $+11.3627$ & 801 & $0$ & $1$ & $-4.4302$ & $-5.4302$ \\
134 & $1$ & $-1$ & $1.5996$ & $+2.5996$ & 468 & $0$ & $-7$ & $4.4884$ & $+11.4884$ & 802 & $1$ & $2$ & $-4.5695$ & $-6.5695$ \\
135 & $0$ & $-1$ & $1.6493$ & $+2.6493$ & 469 & $1$ & $-6$ & $4.5463$ & $+10.5463$ & 803 & $1$ & $3$ & $-4.6392$ & $-7.6392$ \\
136 & $0$ & $-1$ & $1.4524$ & $+2.4524$ & 470 & $-1$ & $-7$ & $4.5568$ & $+11.5568$ & 804 & $0$ & $3$ & $-4.6414$ & $-7.6414$ \\
137 & $-1$ & $-2$ & $1.1017$ & $+3.1017$ & 471 & $1$ & $-6$ & $4.5436$ & $+10.5436$ & 805 & $-1$ & $2$ & $-4.5815$ & $-6.5815$ \\
138 & $-1$ & $-3$ & $0.7119$ & $+3.7119$ & 472 & $0$ & $-6$ & $4.5309$ & $+10.5309$ & 806 & $-1$ & $1$ & $-4.4670$ & $-5.4670$ \\
139 & $-1$ & $-4$ & $0.3857$ & $+4.3857$ & 473 & $1$ & $-5$ & $4.5416$ & $+9.5416$ & 807 & $1$ & $2$ & $-4.3080$ & $-6.3080$ \\
140 & $0$ & $-4$ & $0.1858$ & $+4.1858$ & 474 & $-1$ & $-6$ & $4.5951$ & $+10.5951$ & 808 & $0$ & $2$ & $-4.1160$ & $-6.1160$ \\
141 & $1$ & $-3$ & $0.1181$ & $+3.1181$ & 475 & $0$ & $-6$ & $4.7056$ & $+10.7056$ & 809 & $-1$ & $1$ & $-3.9038$ & $-4.9038$ \\
142 & $1$ & $-2$ & $0.1329$ & $+2.1329$ & 476 & $0$ & $-6$ & $4.8810$ & $+10.8810$ & 810 & $0$ & $1$ & $-3.6845$ & $-4.6845$ \\
143 & $1$ & $-1$ & $0.1413$ & $+1.1413$ & 477 & $0$ & $-6$ & $5.1221$ & $+11.1221$ & 811 & $-1$ & $0$ & $-3.4712$ & $-3.4712$ \\
144 & $0$ & $-1$ & $0.0424$ & $+1.0424$ & 478 & $1$ & $-5$ & $5.4224$ & $+10.4224$ & 812 & $0$ & $0$ & $-3.2760$ & $-3.2760$ \\
145 & $1$ & $0$ & $-0.2466$ & $-0.2466$ & 479 & $-1$ & $-6$ & $5.7686$ & $+11.7686$ & 813 & $1$ & $1$ & $-3.1096$ & $-4.1096$ \\
146 & $1$ & $1$ & $-0.7654$ & $-1.7654$ & 480 & $0$ & $-6$ & $6.1420$ & $+12.1420$ & 814 & $-1$ & $0$ & $-2.9809$ & $-2.9809$ \\
147 & $0$ & $1$ & $-1.4974$ & $-2.4974$ & 481 & $1$ & $-5$ & $6.5195$ & $+11.5195$ & 815 & $1$ & $1$ & $-2.8965$ & $-3.8965$ \\
148 & $0$ & $1$ & $-2.3724$ & $-3.3724$ & 482 & $1$ & $-4$ & $6.8760$ & $+10.8760$ & 816 & $0$ & $1$ & $-2.8600$ & $-3.8600$ \\
149 & $-1$ & $0$ & $-3.2821$ & $-3.2821$ & 483 & $-1$ & $-5$ & $7.1859$ & $+12.1859$ & 817 & $1$ & $2$ & $-2.8726$ & $-4.8726$ \\
150 & $0$ & $0$ & $-4.1063$ & $-4.1063$ & 484 & $0$ & $-5$ & $7.4257$ & $+12.4257$ & 818 & $1$ & $3$ & $-2.9323$ & $-5.9323$ \\
151 & $-1$ & $-1$ & $-4.7412$ & $-3.7412$ & 485 & $1$ & $-4$ & $7.5754$ & $+11.5754$ & 819 & $0$ & $3$ & $-3.0347$ & $-6.0347$ \\
152 & $0$ & $-1$ & $-5.1236$ & $-4.1236$ & 486 & $0$ & $-4$ & $7.6205$ & $+11.6205$ & 820 & $0$ & $3$ & $-3.1727$ & $-6.1727$ \\
153 & $0$ & $-1$ & $-5.2446$ & $-4.2446$ & 487 & $-1$ & $-5$ & $7.5531$ & $+12.5531$ & 821 & $-1$ & $2$ & $-3.3370$ & $-5.3370$ \\
154 & $-1$ & $-2$ & $-5.1491$ & $-3.1491$ & 488 & $0$ & $-5$ & $7.3724$ & $+12.3724$ & 822 & $-1$ & $1$ & $-3.5168$ & $-4.5168$ \\
155 & $1$ & $-1$ & $-4.9234$ & $-3.9234$ & 489 & $1$ & $-4$ & $7.0847$ & $+11.0847$ & 823 & $-1$ & $0$ & $-3.7000$ & $-3.7000$ \\
156 & $0$ & $-1$ & $-4.6723$ & $-3.6723$ & 490 & $0$ & $-4$ & $6.7033$ & $+10.7033$ & 824 & $0$ & $0$ & $-3.8742$ & $-3.8742$ \\
157 & $-1$ & $-2$ & $-4.4931$ & $-2.4931$ & 491 & $-1$ & $-5$ & $6.2472$ & $+11.2472$ & 825 & $0$ & $0$ & $-4.0268$ & $-4.0268$ \\
158 & $1$ & $-1$ & $-4.4525$ & $-3.4525$ & 492 & $0$ & $-5$ & $5.7396$ & $+10.7396$ & 826 & $-1$ & $-1$ & $-4.1461$ & $-3.1461$ \\
159 & $1$ & $0$ & $-4.5718$ & $-4.5718$ & 493 & $1$ & $-4$ & $5.2061$ & $+9.2061$ & 827 & $-1$ & $-2$ & $-4.2217$ & $-2.2217$ \\
160 & $0$ & $0$ & $-4.8235$ & $-4.8235$ & 494 & $-1$ & $-5$ & $4.6731$ & $+9.6731$ & 828 & $0$ & $-2$ & $-4.2449$ & $-2.2449$ \\
161 & $1$ & $1$ & $-5.1402$ & $-6.1402$ & 495 & $0$ & $-5$ & $4.1654$ & $+9.1654$ & 829 & $-1$ & $-3$ & $-4.2092$ & $-1.2092$ \\
162 & $0$ & $1$ & $-5.4332$ & $-6.4332$ & 496 & $0$ & $-5$ & $3.7045$ & $+8.7045$ & 830 & $-1$ & $-4$ & $-4.1110$ & $-0.1110$ \\
163 & $-1$ & $0$ & $-5.6150$ & $-5.6150$ & 497 & $1$ & $-4$ & $3.3071$ & $+7.3071$ & 831 & $1$ & $-3$ & $-3.9489$ & $-0.9489$ \\
164 & $0$ & $0$ & $-5.6219$ & $-5.6219$ & 498 & $-1$ & $-5$ & $2.9838$ & $+7.9838$ & 832 & $0$ & $-3$ & $-3.7249$ & $-0.7249$ \\
165 & $-1$ & $-1$ & $-5.4291$ & $-4.4291$ & 499 & $-1$ & $-6$ & $2.7387$ & $+8.7387$ & 833 & $0$ & $-3$ & $-3.4434$ & $-0.4434$ \\
166 & $1$ & $0$ & $-5.0574$ & $-5.0574$ & 500 & $0$ & $-6$ & $2.5687$ & $+8.5687$ & 834 & $-1$ & $-4$ & $-3.1115$ & $+0.8885$ \\
167 & $-1$ & $-1$ & $-4.5679$ & $-3.5679$ & 501 & $1$ & $-5$ & $2.4646$ & $+7.4646$ & 835 & $1$ & $-3$ & $-2.7386$ & $+0.2614$ \\
168 & $0$ & $-1$ & $-4.0477$ & $-3.0477$ & 502 & $1$ & $-4$ & $2.4114$ & $+6.4114$ & 836 & $0$ & $-3$ & $-2.3359$ & $+0.6641$ \\
169 & $0$ & $-1$ & $-3.5897$ & $-2.5897$ & 503 & $-1$ & $-5$ & $2.3899$ & $+7.3899$ & 837 & $0$ & $-3$ & $-1.9160$ & $+1.0840$ \\
170 & $-1$ & $-2$ & $-3.2712$ & $-1.2712$ & 504 & $0$ & $-5$ & $2.3783$ & $+7.3783$ & 838 & $1$ & $-2$ & $-1.4922$ & $+0.5078$ \\
171 & $0$ & $-2$ & $-3.1362$ & $-1.1362$ & 505 & $1$ & $-4$ & $2.3539$ & $+6.3539$ & 839 & $-1$ & $-3$ & $-1.0781$ & $+1.9219$ \\
172 & $0$ & $-2$ & $-3.1856$ & $-1.1856$ & 506 & $-1$ & $-5$ & $2.2953$ & $+7.2953$ & 840 & $0$ & $-3$ & $-0.6868$ & $+2.3132$ \\
173 & $-1$ & $-3$ & $-3.3772$ & $-0.3772$ & 507 & $0$ & $-5$ & $2.1836$ & $+7.1836$ & 841 & $0$ & $-3$ & $-0.3304$ & $+2.6696$ \\
174 & $-1$ & $-4$ & $-3.6351$ & $+0.3649$ & 508 & $0$ & $-5$ & $2.0044$ & $+7.0044$ & 842 & $1$ & $-2$ & $-0.0195$ & $+1.9805$ \\
175 & $0$ & $-4$ & $-3.8666$ & $+0.1334$ & 509 & $-1$ & $-6$ & $1.7490$ & $+7.7490$ & 843 & $1$ & $-1$ & $0.2373$ & $+1.2373$ \\
176 & $0$ & $-4$ & $-3.9822$ & $+0.0178$ & 510 & $1$ & $-5$ & $1.4147$ & $+6.4147$ & 844 & $0$ & $-1$ & $0.4338$ & $+1.4338$ \\
177 & $1$ & $-3$ & $-3.9143$ & $-0.9143$ & 511 & $1$ & $-4$ & $1.0052$ & $+5.0052$ & 845 & $0$ & $-1$ & $0.5662$ & $+1.5662$ \\
178 & $1$ & $-2$ & $-3.6318$ & $-1.6318$ & 512 & $0$ & $-4$ & $0.5308$ & $+4.5308$ & 846 & $0$ & $-1$ & $0.6339$ & $+1.6339$ \\
179 & $-1$ & $-3$ & $-3.1450$ & $-0.1450$ & 513 & $0$ & $-4$ & $0.0069$ & $+4.0069$ & 847 & $0$ & $-1$ & $0.6385$ & $+1.6385$ \\
180 & $0$ & $-3$ & $-2.5031$ & $+0.4969$ & 514 & $1$ & $-3$ & $-0.5468$ & $+2.4532$ & 848 & $0$ & $-1$ & $0.5846$ & $+1.5846$ \\
181 & $-1$ & $-4$ & $-1.7827$ & $+2.2173$ & 515 & $1$ & $-2$ & $-1.1076$ & $+0.8924$ & 849 & $1$ & $0$ & $0.4792$ & $+0.4792$ \\
182 & $-1$ & $-5$ & $-1.0713$ & $+3.9287$ & 516 & $0$ & $-2$ & $-1.6520$ & $+0.3480$ & 850 & $0$ & $0$ & $0.3310$ & $+0.3310$ \\
183 & $1$ & $-4$ & $-0.4487$ & $+3.5513$ & 517 & $1$ & $-1$ & $-2.1571$ & $-1.1571$ & 851 & $1$ & $1$ & $0.1508$ & $-0.8492$ \\
184 & $0$ & $-4$ & $0.0289$ & $+4.0289$ & 518 & $-1$ & $-2$ & $-2.6025$ & $-0.6025$ & 852 & $0$ & $1$ & $-0.0500$ & $-1.0500$ \\
185 & $1$ & $-3$ & $0.3394$ & $+3.3394$ & 519 & $1$ & $-1$ & $-2.9716$ & $-1.9716$ & 853 & $-1$ & $0$ & $-0.2591$ & $-0.2591$ \\
186 & $-1$ & $-4$ & $0.4983$ & $+4.4983$ & 520 & $0$ & $-1$ & $-3.2532$ & $-2.2532$ & 854 & $-1$ & $-1$ & $-0.4645$ & $+0.5355$ \\
187 & $1$ & $-3$ & $0.5529$ & $+3.5529$ & 521 & $-1$ & $-2$ & $-3.4422$ & $-1.4422$ & 855 & $0$ & $-1$ & $-0.6546$ & $+0.3454$ \\
188 & $0$ & $-3$ & $0.5711$ & $+3.5711$ & 522 & $0$ & $-2$ & $-3.5397$ & $-1.5397$ & 856 & $0$ & $-1$ & $-0.8191$ & $+0.1809$ \\
189 & $0$ & $-3$ & $0.6257$ & $+3.6257$ & 523 & $-1$ & $-3$ & $-3.5533$ & $-0.5533$ & 857 & $-1$ & $-2$ & $-0.9492$ & $+1.0508$ \\
190 & $-1$ & $-4$ & $0.7773$ & $+4.7773$ & 524 & $0$ & $-3$ & $-3.4962$ & $-0.4962$ & 858 & $1$ & $-1$ & $-1.0384$ & $-0.0384$ \\
191 & $-1$ & $-5$ & $1.0612$ & $+6.0612$ & 525 & $0$ & $-3$ & $-3.3866$ & $-0.3866$ & 859 & $-1$ & $-2$ & $-1.0821$ & $+0.9179$ \\
192 & $0$ & $-5$ & $1.4790$ & $+6.4790$ & 526 & $1$ & $-2$ & $-3.2459$ & $-1.2459$ & 860 & $0$ & $-2$ & $-1.0787$ & $+0.9213$ \\
193 & $-1$ & $-6$ & $1.9971$ & $+7.9971$ & 527 & $1$ & $-1$ & $-3.0974$ & $-2.0974$ & 861 & $-1$ & $-3$ & $-1.0289$ & $+1.9711$ \\
194 & $1$ & $-5$ & $2.5533$ & $+7.5533$ & 528 & $0$ & $-1$ & $-2.9647$ & $-1.9647$ & 862 & $1$ & $-2$ & $-0.9362$ & $+1.0638$ \\
195 & $-1$ & $-6$ & $3.0682$ & $+9.0682$ & 529 & $0$ & $-1$ & $-2.8696$ & $-1.8696$ & 863 & $-1$ & $-3$ & $-0.8061$ & $+2.1939$ \\
196 & $0$ & $-6$ & $3.4608$ & $+9.4608$ & 530 & $-1$ & $-2$ & $-2.8311$ & $-0.8311$ & 864 & $0$ & $-3$ & $-0.6467$ & $+2.3533$ \\
197 & $-1$ & $-7$ & $3.6636$ & $+10.6636$ & 531 & $0$ & $-2$ & $-2.8634$ & $-0.8634$ & 865 & $1$ & $-2$ & $-0.4675$ & $+1.5325$ \\
198 & $0$ & $-7$ & $3.6358$ & $+10.6358$ & 532 & $0$ & $-2$ & $-2.9756$ & $-0.9756$ & 866 & $1$ & $-1$ & $-0.2796$ & $+0.7204$ \\
199 & $-1$ & $-8$ & $3.3703$ & $+11.3703$ & 533 & $1$ & $-1$ & $-3.1705$ & $-2.1705$ & 867 & $0$ & $-1$ & $-0.0946$ & $+0.9054$ \\
200 & $0$ & $-8$ & $2.8952$ & $+10.8952$ & 534 & $-1$ & $-2$ & $-3.4451$ & $-1.4451$ & 868 & $0$ & $-1$ & $0.0755$ & $+1.0755$ \\
201 & $1$ & $-7$ & $2.2680$ & $+9.2680$ & 535 & $1$ & $-1$ & $-3.7903$ & $-2.7903$ & 869 & $1$ & $0$ & $0.2189$ & $+0.2189$ \\
202 & $1$ & $-6$ & $1.5648$ & $+7.5648$ & 536 & $0$ & $-1$ & $-4.1920$ & $-3.1920$ & 870 & $1$ & $1$ & $0.3250$ & $-0.6750$ \\
203 & $1$ & $-5$ & $0.8665$ & $+5.8665$ & 537 & $1$ & $0$ & $-4.6320$ & $-4.6320$ & 871 & $1$ & $2$ & $0.3841$ & $-1.6159$ \\
204 & $0$ & $-5$ & $0.2437$ & $+5.2437$ & 538 & $1$ & $1$ & $-5.0897$ & $-6.0897$ & 872 & $0$ & $2$ & $0.3886$ & $-1.6114$ \\
205 & $1$ & $-4$ & $-0.2547$ & $+3.7453$ & 539 & $0$ & $1$ & $-5.5432$ & $-6.5432$ & 873 & $0$ & $2$ & $0.3331$ & $-1.6669$ \\
206 & $1$ & $-3$ & $-0.6094$ & $+2.3906$ & 540 & $0$ & $1$ & $-5.9713$ & $-6.9713$ & 874 & $-1$ & $1$ & $0.2142$ & $-0.7858$ \\
207 & $0$ & $-3$ & $-0.8324$ & $+2.1676$ & 541 & $-1$ & $0$ & $-6.3550$ & $-6.3550$ & 875 & $0$ & $1$ & $0.0316$ & $-0.9684$ \\
208 & $0$ & $-3$ & $-0.9629$ & $+2.0371$ & 542 & $1$ & $1$ & $-6.6785$ & $-7.6785$ & 876 & $0$ & $1$ & $-0.2129$ & $-1.2129$ \\
209 & $1$ & $-2$ & $-1.0581$ & $+0.9419$ & 543 & $1$ & $2$ & $-6.9309$ & $-8.9309$ & 877 & $-1$ & $0$ & $-0.5146$ & $-0.5146$ \\
210 & $1$ & $-1$ & $-1.1806$ & $-0.1806$ & 544 & $0$ & $2$ & $-7.1067$ & $-9.1067$ & 878 & $1$ & $1$ & $-0.8668$ & $-1.8668$ \\
211 & $-1$ & $-2$ & $-1.3855$ & $+0.6145$ & 545 & $1$ & $3$ & $-7.2060$ & $-10.2060$ & 879 & $1$ & $2$ & $-1.2604$ & $-3.2604$ \\
212 & $0$ & $-2$ & $-1.7091$ & $+0.2909$ & 546 & $1$ & $4$ & $-7.2346$ & $-11.2346$ & 880 & $0$ & $2$ & $-1.6850$ & $-3.6850$ \\
213 & $1$ & $-1$ & $-2.1610$ & $-1.1610$ & 547 & $-1$ & $3$ & $-7.2037$ & $-10.2037$ & 881 & $-1$ & $1$ & $-2.1289$ & $-3.1289$ \\
214 & $1$ & $0$ & $-2.7222$ & $-2.7222$ & 548 & $0$ & $3$ & $-7.1288$ & $-10.1288$ & 882 & $0$ & $1$ & $-2.5795$ & $-3.5795$ \\
215 & $1$ & $1$ & $-3.3474$ & $-4.3474$ & 549 & $0$ & $3$ & $-7.0287$ & $-10.0287$ & 883 & $-1$ & $0$ & $-3.0244$ & $-3.0244$ \\
216 & $0$ & $1$ & $-3.9734$ & $-4.9734$ & 550 & $0$ & $3$ & $-6.9238$ & $-9.9238$ & 884 & $0$ & $0$ & $-3.4512$ & $-3.4512$ \\
217 & $1$ & $2$ & $-4.5299$ & $-6.5299$ & 551 & $1$ & $4$ & $-6.8352$ & $-10.8352$ & 885 & $-1$ & $-1$ & $-3.8488$ & $-2.8488$ \\
218 & $1$ & $3$ & $-4.9520$ & $-7.9520$ & 552 & $0$ & $4$ & $-6.7824$ & $-10.7824$ & 886 & $1$ & $0$ & $-4.2071$ & $-4.2071$ \\
219 & $1$ & $4$ & $-5.1918$ & $-9.1918$ & 553 & $1$ & $5$ & $-6.7824$ & $-11.7824$ & 887 & $-1$ & $-1$ & $-4.5180$ & $-3.5180$ \\
220 & $0$ & $4$ & $-5.2263$ & $-9.2263$ & 554 & $1$ & $6$ & $-6.8485$ & $-12.8485$ & 888 & $0$ & $-1$ & $-4.7756$ & $-3.7756$ \\
221 & $1$ & $5$ & $-5.0620$ & $-10.0620$ & 555 & $-1$ & $5$ & $-6.9888$ & $-11.9888$ & 889 & $1$ & $0$ & $-4.9760$ & $-4.9760$ \\
222 & $-1$ & $4$ & $-4.7331$ & $-8.7331$ & 556 & $0$ & $5$ & $-7.2061$ & $-12.2061$ & 890 & $-1$ & $-1$ & $-5.1183$ & $-4.1183$ \\
223 & $-1$ & $3$ & $-4.2963$ & $-7.2963$ & 557 & $-1$ & $4$ & $-7.4975$ & $-11.4975$ & 891 & $0$ & $-1$ & $-5.2038$ & $-4.2038$ \\
224 & $0$ & $3$ & $-3.8211$ & $-6.8211$ & 558 & $0$ & $4$ & $-7.8544$ & $-11.8544$ & 892 & $0$ & $-1$ & $-5.2361$ & $-4.2361$ \\
225 & $0$ & $3$ & $-3.3782$ & $-6.3782$ & 559 & $1$ & $5$ & $-8.2633$ & $-13.2633$ & 893 & $1$ & $0$ & $-5.2214$ & $-5.2214$ \\
226 & $1$ & $4$ & $-3.0283$ & $-7.0283$ & 560 & $0$ & $5$ & $-8.7067$ & $-13.7067$ & 894 & $-1$ & $-1$ & $-5.1677$ & $-4.1677$ \\
227 & $-1$ & $3$ & $-2.8125$ & $-5.8125$ & 561 & $-1$ & $4$ & $-9.1642$ & $-13.1642$ & 895 & $1$ & $0$ & $-5.0843$ & $-5.0843$ \\
228 & $0$ & $3$ & $-2.7467$ & $-5.7467$ & 562 & $1$ & $5$ & $-9.6138$ & $-14.6138$ & 896 & $0$ & $0$ & $-4.9821$ & $-4.9821$ \\
229 & $-1$ & $2$ & $-2.8197$ & $-4.8197$ & 563 & $-1$ & $4$ & $-10.0335$ & $-14.0335$ & 897 & $-1$ & $-1$ & $-4.8723$ & $-3.8723$ \\
230 & $-1$ & $1$ & $-2.9968$ & $-3.9968$ & 564 & $0$ & $4$ & $-10.4027$ & $-14.4027$ & 898 & $1$ & $0$ & $-4.7664$ & $-4.7664$ \\
231 & $-1$ & $0$ & $-3.2264$ & $-3.2264$ & 565 & $1$ & $5$ & $-10.7036$ & $-15.7036$ & 899 & $1$ & $1$ & $-4.6753$ & $-5.6753$ \\
232 & $0$ & $0$ & $-3.4496$ & $-3.4496$ & 566 & $1$ & $6$ & $-10.9224$ & $-16.9224$ & 900 & $0$ & $1$ & $-4.6094$ & $-5.6094$ \\
233 & $-1$ & $-1$ & $-3.6115$ & $-2.6115$ & 567 & $0$ & $6$ & $-11.0501$ & $-17.0501$ & 901 & $1$ & $2$ & $-4.5773$ & $-6.5773$ \\
234 & $0$ & $-1$ & $-3.6698$ & $-2.6698$ & 568 & $0$ & $6$ & $-11.0831$ & $-17.0831$ & 902 & $-1$ & $1$ & $-4.5864$ & $-5.5864$ \\
235 & $1$ & $0$ & $-3.6032$ & $-3.6032$ & 569 & $-1$ & $5$ & $-11.0232$ & $-16.0232$ & 903 & $-1$ & $0$ & $-4.6416$ & $-4.6416$ \\
236 & $0$ & $0$ & $-3.4141$ & $-3.4141$ & 570 & $1$ & $6$ & $-10.8778$ & $-16.8778$ & 904 & $0$ & $0$ & $-4.7459$ & $-4.7459$ \\
237 & $1$ & $1$ & $-3.1286$ & $-4.1286$ & 571 & $-1$ & $5$ & $-10.6589$ & $-15.6589$ & 905 & $1$ & $1$ & $-4.8997$ & $-5.8997$ \\
238 & $-1$ & $0$ & $-2.7921$ & $-2.7921$ & 572 & $0$ & $5$ & $-10.3824$ & $-15.3824$ & 906 & $-1$ & $0$ & $-5.1011$ & $-5.1011$ \\
239 & $-1$ & $-1$ & $-2.4610$ & $-1.4610$ & 573 & $1$ & $6$ & $-10.0669$ & $-16.0669$ & 907 & $-1$ & $-1$ & $-5.3458$ & $-4.3458$ \\
240 & $0$ & $-1$ & $-2.1939$ & $-1.1939$ & 574 & $-1$ & $5$ & $-9.7326$ & $-14.7326$ & 908 & $0$ & $-1$ & $-5.6276$ & $-4.6276$ \\
241 & $-1$ & $-2$ & $-2.0413$ & $-0.0413$ & 575 & $0$ & $5$ & $-9.3996$ & $-14.3996$ & 909 & $0$ & $-1$ & $-5.9383$ & $-4.9383$ \\
242 & $0$ & $-2$ & $-2.0376$ & $-0.0376$ & 576 & $0$ & $5$ & $-9.0866$ & $-14.0866$ & 910 & $1$ & $0$ & $-6.2682$ & $-6.2682$ \\
243 & $0$ & $-2$ & $-2.1949$ & $-0.1949$ & 577 & $-1$ & $4$ & $-8.8099$ & $-12.8099$ & 911 & $-1$ & $-1$ & $-6.6067$ & $-5.6067$ \\
244 & $0$ & $-2$ & $-2.5015$ & $-0.5015$ & 578 & $0$ & $4$ & $-8.5818$ & $-12.5818$ & 912 & $0$ & $-1$ & $-6.9426$ & $-5.9426$ \\
245 & $0$ & $-2$ & $-2.9228$ & $-0.9228$ & 579 & $1$ & $5$ & $-8.4103$ & $-13.4103$ & 913 & $1$ & $0$ & $-7.2644$ & $-7.2644$ \\
246 & $-1$ & $-3$ & $-3.4071$ & $-0.4071$ & 580 & $0$ & $5$ & $-8.2983$ & $-13.2983$ & 914 & $1$ & $1$ & $-7.5612$ & $-8.5612$ \\
247 & $1$ & $-2$ & $-3.8933$ & $-1.8933$ & 581 & $1$ & $6$ & $-8.2435$ & $-14.2435$ & 915 & $-1$ & $0$ & $-7.8231$ & $-7.8231$ \\
248 & $0$ & $-2$ & $-4.3198$ & $-2.3198$ & 582 & $-1$ & $5$ & $-8.2388$ & $-13.2388$ & 916 & $0$ & $0$ & $-8.0412$ & $-8.0412$ \\
249 & $1$ & $-1$ & $-4.6343$ & $-3.6343$ & 583 & $1$ & $6$ & $-8.2723$ & $-14.2723$ & 917 & $1$ & $1$ & $-8.2087$ & $-9.2087$ \\
250 & $0$ & $-1$ & $-4.8007$ & $-3.8007$ & 584 & $0$ & $6$ & $-8.3291$ & $-14.3291$ & 918 & $0$ & $1$ & $-8.3204$ & $-9.3204$ \\
251 & $-1$ & $-2$ & $-4.8046$ & $-2.8046$ & 585 & $0$ & $6$ & $-8.3913$ & $-14.3913$ & 919 & $-1$ & $0$ & $-8.3735$ & $-8.3735$ \\
252 & $0$ & $-2$ & $-4.6547$ & $-2.6547$ & 586 & $1$ & $7$ & $-8.4402$ & $-15.4402$ & 920 & $0$ & $0$ & $-8.3677$ & $-8.3677$ \\
253 & $1$ & $-1$ & $-4.3808$ & $-3.3808$ & 587 & $-1$ & $6$ & $-8.4570$ & $-14.4570$ & 921 & $1$ & $1$ & $-8.3050$ & $-9.3050$ \\
254 & $1$ & $0$ & $-4.0290$ & $-4.0290$ & 588 & $0$ & $6$ & $-8.4243$ & $-14.4243$ & 922 & $1$ & $2$ & $-8.1894$ & $-10.1894$ \\
255 & $-1$ & $-1$ & $-3.6541$ & $-2.6541$ & 589 & $1$ & $7$ & $-8.3275$ & $-15.3275$ & 923 & $1$ & $3$ & $-8.0274$ & $-11.0274$ \\
256 & $0$ & $-1$ & $-3.3109$ & $-2.3109$ & 590 & $-1$ & $6$ & $-8.1555$ & $-14.1555$ & 924 & $0$ & $3$ & $-7.8273$ & $-10.8273$ \\
257 & $-1$ & $-2$ & $-3.0462$ & $-1.0462$ & 591 & $1$ & $7$ & $-7.9018$ & $-14.9018$ & 925 & $0$ & $3$ & $-7.5986$ & $-10.5986$ \\
258 & $-1$ & $-3$ & $-2.8917$ & $+0.1083$ & 592 & $0$ & $7$ & $-7.5650$ & $-14.5650$ & 926 & $1$ & $4$ & $-7.3524$ & $-11.3524$ \\
259 & $1$ & $-2$ & $-2.8595$ & $-0.8595$ & 593 & $-1$ & $6$ & $-7.1485$ & $-13.1485$ & 927 & $0$ & $4$ & $-7.1000$ & $-11.1000$ \\
260 & $0$ & $-2$ & $-2.9407$ & $-0.9407$ & 594 & $0$ & $6$ & $-6.6609$ & $-12.6609$ & 928 & $0$ & $4$ & $-6.8532$ & $-10.8532$ \\
261 & $0$ & $-2$ & $-3.1075$ & $-1.1075$ & 595 & $-1$ & $5$ & $-6.1151$ & $-11.1151$ & 929 & $-1$ & $3$ & $-6.6234$ & $-9.6234$ \\
262 & $1$ & $-1$ & $-3.3174$ & $-2.3174$ & 596 & $0$ & $5$ & $-5.5277$ & $-10.5277$ & 930 & $1$ & $4$ & $-6.4211$ & $-10.4211$ \\
263 & $-1$ & $-2$ & $-3.5207$ & $-1.5207$ & 597 & $1$ & $6$ & $-4.9181$ & $-10.9181$ & 931 & $0$ & $4$ & $-6.2558$ & $-10.2558$ \\
264 & $0$ & $-2$ & $-3.6678$ & $-1.6678$ & 598 & $-1$ & $5$ & $-4.3068$ & $-9.3068$ & 932 & $0$ & $4$ & $-6.1354$ & $-10.1354$ \\
265 & $1$ & $-1$ & $-3.7172$ & $-2.7172$ & 599 & $-1$ & $4$ & $-3.7148$ & $-7.7148$ & 933 & $1$ & $5$ & $-6.0658$ & $-11.0658$ \\
266 & $-1$ & $-2$ & $-3.6418$ & $-1.6418$ & 600 & $0$ & $4$ & $-3.1615$ & $-7.1615$ & 934 & $1$ & $6$ & $-6.0509$ & $-12.0509$ \\
267 & $1$ & $-1$ & $-3.4334$ & $-2.4334$ & 601 & $-1$ & $3$ & $-2.6645$ & $-5.6645$ & 935 & $-1$ & $5$ & $-6.0921$ & $-11.0921$ \\
268 & $0$ & $-1$ & $-3.1037$ & $-2.1037$ & 602 & $-1$ & $2$ & $-2.2377$ & $-4.2377$ & 936 & $0$ & $5$ & $-6.1889$ & $-11.1889$ \\
269 & $-1$ & $-2$ & $-2.6832$ & $-0.6832$ & 603 & $0$ & $2$ & $-1.8909$ & $-3.8909$ & 937 & $-1$ & $4$ & $-6.3382$ & $-10.3382$ \\
270 & $0$ & $-2$ & $-2.2163$ & $-0.2163$ & 604 & $0$ & $2$ & $-1.6291$ & $-3.6291$ & 938 & $-1$ & $3$ & $-6.5348$ & $-9.5348$ \\
271 & $-1$ & $-3$ & $-1.7558$ & $+1.2442$ & 605 & $0$ & $2$ & $-1.4523$ & $-3.4523$ & 939 & $1$ & $4$ & $-6.7718$ & $-10.7718$ \\
272 & $0$ & $-3$ & $-1.3547$ & $+1.6453$ & 606 & $-1$ & $1$ & $-1.3558$ & $-2.3558$ & 940 & $0$ & $4$ & $-7.0406$ & $-11.0406$ \\
273 & $-1$ & $-4$ & $-1.0594$ & $+2.9406$ & 607 & $-1$ & $0$ & $-1.3302$ & $-1.3302$ & 941 & $-1$ & $3$ & $-7.3314$ & $-10.3314$ \\
274 & $1$ & $-3$ & $-0.9028$ & $+2.0972$ & 608 & $0$ & $0$ & $-1.3623$ & $-1.3623$ & 942 & $-1$ & $2$ & $-7.6336$ & $-9.6336$ \\
275 & $0$ & $-3$ & $-0.9002$ & $+2.0998$ & 609 & $-1$ & $-1$ & $-1.4360$ & $-0.4360$ & 943 & $1$ & $3$ & $-7.9363$ & $-10.9363$ \\
276 & $0$ & $-3$ & $-1.0470$ & $+1.9530$ & 610 & $-1$ & $-2$ & $-1.5337$ & $+0.4663$ & 944 & $0$ & $3$ & $-8.2288$ & $-11.2288$ \\
277 & $-1$ & $-4$ & $-1.3198$ & $+2.6802$ & 611 & $1$ & $-1$ & $-1.6366$ & $-0.6366$ & 945 & $0$ & $3$ & $-8.5009$ & $-11.5009$ \\
278 & $1$ & $-3$ & $-1.6794$ & $+1.3206$ & 612 & $0$ & $-1$ & $-1.7271$ & $-0.7271$ & 946 & $-1$ & $2$ & $-8.7433$ & $-10.7433$ \\
279 & $0$ & $-3$ & $-2.0769$ & $+0.9231$ & 613 & $-1$ & $-2$ & $-1.7890$ & $+0.2110$ & 947 & $-1$ & $1$ & $-8.9483$ & $-9.9483$ \\
280 & $0$ & $-3$ & $-2.4595$ & $+0.5405$ & 614 & $1$ & $-1$ & $-1.8090$ & $-0.8090$ & 948 & $0$ & $1$ & $-9.1098$ & $-10.1098$ \\
281 & $-1$ & $-4$ & $-2.7786$ & $+1.2214$ & 615 & $-1$ & $-2$ & $-1.7776$ & $+0.2224$ & 949 & $1$ & $2$ & $-9.2237$ & $-11.2237$ \\
282 & $-1$ & $-5$ & $-2.9955$ & $+2.0045$ & 616 & $0$ & $-2$ & $-1.6895$ & $+0.3105$ & 950 & $0$ & $2$ & $-9.2879$ & $-11.2879$ \\
283 & $-1$ & $-6$ & $-3.0866$ & $+2.9134$ & 617 & $-1$ & $-3$ & $-1.5438$ & $+1.4562$ & 951 & $1$ & $3$ & $-9.3026$ & $-12.3026$ \\
284 & $0$ & $-6$ & $-3.0459$ & $+2.9541$ & 618 & $-1$ & $-4$ & $-1.3445$ & $+2.6555$ & 952 & $0$ & $3$ & $-9.2702$ & $-12.2702$ \\
285 & $-1$ & $-7$ & $-2.8857$ & $+4.1143$ & 619 & $-1$ & $-5$ & $-1.0998$ & $+3.9002$ & 953 & $-1$ & $2$ & $-9.1951$ & $-11.1951$ \\
286 & $-1$ & $-8$ & $-2.6338$ & $+5.3662$ & 620 & $0$ & $-5$ & $-0.8217$ & $+4.1783$ & 954 & $0$ & $2$ & $-9.0837$ & $-11.0837$ \\
287 & $1$ & $-7$ & $-2.3295$ & $+4.6705$ & 621 & $0$ & $-5$ & $-0.5252$ & $+4.4748$ & 955 & $1$ & $3$ & $-8.9438$ & $-11.9438$ \\
288 & $0$ & $-7$ & $-2.0175$ & $+4.9825$ & 622 & $1$ & $-4$ & $-0.2271$ & $+3.7729$ & 956 & $0$ & $3$ & $-8.7846$ & $-11.7846$ \\
289 & $0$ & $-7$ & $-1.7419$ & $+5.2581$ & 623 & $1$ & $-3$ & $0.0547$ & $+3.0547$ & 957 & $-1$ & $2$ & $-8.6163$ & $-10.6163$ \\
290 & $-1$ & $-8$ & $-1.5394$ & $+6.4606$ & 624 & $0$ & $-3$ & $0.3026$ & $+3.3026$ & 958 & $1$ & $3$ & $-8.4492$ & $-11.4492$ \\
291 & $1$ & $-7$ & $-1.4342$ & $+5.5658$ & 625 & $0$ & $-3$ & $0.5005$ & $+3.5005$ & 959 & $1$ & $4$ & $-8.2939$ & $-12.2939$ \\
292 & $0$ & $-7$ & $-1.4351$ & $+5.5649$ & 626 & $1$ & $-2$ & $0.6348$ & $+2.6348$ & 960 & $0$ & $4$ & $-8.1604$ & $-12.1604$ \\
293 & $-1$ & $-8$ & $-1.5338$ & $+6.4662$ & 627 & $-1$ & $-3$ & $0.6954$ & $+3.6954$ & 961 & $0$ & $4$ & $-8.0580$ & $-12.0580$ \\
294 & $0$ & $-8$ & $-1.7061$ & $+6.2939$ & 628 & $0$ & $-3$ & $0.6758$ & $+3.6758$ & 962 & $-1$ & $3$ & $-7.9947$ & $-10.9947$ \\
295 & $1$ & $-7$ & $-1.9154$ & $+5.0846$ & 629 & $1$ & $-2$ & $0.5745$ & $+2.5745$ & 963 & $0$ & $3$ & $-7.9768$ & $-10.9768$ \\
296 & $0$ & $-7$ & $-2.1174$ & $+4.8826$ & 630 & $0$ & $-2$ & $0.3943$ & $+2.3943$ & 964 & $0$ & $3$ & $-8.0088$ & $-11.0088$ \\
297 & $0$ & $-7$ & $-2.2663$ & $+4.7337$ & 631 & $-1$ & $-3$ & $0.1425$ & $+3.1425$ & 965 & $1$ & $4$ & $-8.0933$ & $-12.0933$ \\
298 & $1$ & $-6$ & $-2.3206$ & $+3.6794$ & 632 & $0$ & $-3$ & $-0.1693$ & $+2.8307$ & 966 & $1$ & $5$ & $-8.2306$ & $-13.2306$ \\
299 & $1$ & $-5$ & $-2.2488$ & $+2.7512$ & 633 & $1$ & $-2$ & $-0.5265$ & $+1.4735$ & 967 & $-1$ & $4$ & $-8.4188$ & $-12.4188$ \\
300 & $0$ & $-5$ & $-2.0336$ & $+2.9664$ & 634 & $1$ & $-1$ & $-0.9116$ & $+0.0884$ & 968 & $0$ & $4$ & $-8.6539$ & $-12.6539$ \\
301 & $1$ & $-4$ & $-1.6743$ & $+2.3257$ & 635 & $1$ & $0$ & $-1.3059$ & $-1.3059$ & 969 & $-1$ & $3$ & $-8.9300$ & $-11.9300$ \\
302 & $1$ & $-3$ & $-1.1871$ & $+1.8129$ & 636 & $0$ & $0$ & $-1.6902$ & $-1.6902$ & 970 & $-1$ & $2$ & $-9.2395$ & $-11.2395$ \\
303 & $1$ & $-2$ & $-0.6029$ & $+1.3971$ & 637 & $0$ & $0$ & $-2.0459$ & $-2.0459$ & 971 & $-1$ & $1$ & $-9.5732$ & $-10.5732$ \\
304 & $0$ & $-2$ & $0.0364$ & $+2.0364$ & 638 & $-1$ & $-1$ & $-2.3564$ & $-1.3564$ & 972 & $0$ & $1$ & $-9.9210$ & $-10.9210$ \\
305 & $1$ & $-1$ & $0.6826$ & $+1.6826$ & 639 & $0$ & $-1$ & $-2.6079$ & $-1.6079$ & 973 & $1$ & $2$ & $-10.2719$ & $-12.2719$ \\
306 & $0$ & $-1$ & $1.2874$ & $+2.2874$ & 640 & $0$ & $-1$ & $-2.7901$ & $-1.7901$ & 974 & $1$ & $3$ & $-10.6150$ & $-13.6150$ \\
307 & $-1$ & $-2$ & $1.8082$ & $+3.8082$ & 641 & $-1$ & $-2$ & $-2.8966$ & $-0.8966$ & 975 & $0$ & $3$ & $-10.9391$ & $-13.9391$ \\
308 & $0$ & $-2$ & $2.2125$ & $+4.2125$ & 642 & $-1$ & $-3$ & $-2.9258$ & $+0.0742$ & 976 & $0$ & $3$ & $-11.2340$ & $-14.2340$ \\
309 & $1$ & $-1$ & $2.4821$ & $+3.4821$ & 643 & $-1$ & $-4$ & $-2.8803$ & $+1.1197$ & 977 & $-1$ & $2$ & $-11.4902$ & $-13.4902$ \\
310 & $-1$ & $-2$ & $2.6144$ & $+4.6144$ & 644 & $0$ & $-4$ & $-2.7669$ & $+1.2331$ & 978 & $-1$ & $1$ & $-11.6999$ & $-12.6999$ \\
311 & $-1$ & $-3$ & $2.6223$ & $+5.6223$ & 645 & $-1$ & $-5$ & $-2.5965$ & $+2.4035$ & 979 & $1$ & $2$ & $-11.8565$ & $-13.8565$ \\
312 & $0$ & $-3$ & $2.5320$ & $+5.5320$ & 646 & $-1$ & $-6$ & $-2.3826$ & $+3.6174$ & 980 & $0$ & $2$ & $-11.9557$ & $-13.9557$ \\
313 & $-1$ & $-4$ & $2.3794$ & $+6.3794$ & 647 & $-1$ & $-7$ & $-2.1412$ & $+4.8588$ & 981 & $0$ & $2$ & $-11.9950$ & $-13.9950$ \\
314 & $1$ & $-3$ & $2.2049$ & $+5.2049$ & 648 & $0$ & $-7$ & $-1.8891$ & $+5.1109$ & 982 & $1$ & $3$ & $-11.9741$ & $-14.9741$ \\
315 & $0$ & $-3$ & $2.0480$ & $+5.0480$ & 649 & $1$ & $-6$ & $-1.6436$ & $+4.3564$ & 983 & $-1$ & $2$ & $-11.8948$ & $-13.8948$ \\
316 & $0$ & $-3$ & $1.9426$ & $+4.9426$ & 650 & $0$ & $-6$ & $-1.4208$ & $+4.5792$ & 984 & $0$ & $2$ & $-11.7609$ & $-13.7609$ \\
317 & $-1$ & $-4$ & $1.9122$ & $+5.9122$ & 651 & $-1$ & $-7$ & $-1.2348$ & $+5.7652$ & 985 & $1$ & $3$ & $-11.5781$ & $-14.5781$ \\
318 & $-1$ & $-5$ & $1.9671$ & $+6.9671$ & 652 & $0$ & $-7$ & $-1.0970$ & $+5.9030$ & 986 & $-1$ & $2$ & $-11.3536$ & $-13.3536$ \\
319 & $1$ & $-4$ & $2.1031$ & $+6.1031$ & 653 & $-1$ & $-8$ & $-1.0152$ & $+6.9848$ & 987 & $-1$ & $1$ & $-11.0959$ & $-12.0959$ \\
320 & $0$ & $-4$ & $2.3022$ & $+6.3022$ & 654 & $-1$ & $-9$ & $-0.9933$ & $+8.0067$ & 988 & $0$ & $1$ & $-10.8147$ & $-11.8147$ \\
321 & $1$ & $-3$ & $2.5349$ & $+5.5349$ & 655 & $1$ & $-8$ & $-1.0311$ & $+6.9689$ & 989 & $1$ & $2$ & $-10.5198$ & $-12.5198$ \\
322 & $-1$ & $-4$ & $2.7637$ & $+6.7637$ & 656 & $0$ & $-8$ & $-1.1241$ & $+6.8759$ & 990 & $0$ & $2$ & $-10.2214$ & $-12.2214$ \\
323 & $1$ & $-3$ & $2.9482$ & $+5.9482$ & 657 & $0$ & $-8$ & $-1.2639$ & $+6.7361$ & 991 & $-1$ & $1$ & $-9.9296$ & $-10.9296$ \\
324 & $0$ & $-3$ & $3.0499$ & $+6.0499$ & 658 & $-1$ & $-9$ & $-1.4391$ & $+7.5609$ & 992 & $0$ & $1$ & $-9.6534$ & $-10.6534$ \\
325 & $0$ & $-3$ & $3.0370$ & $+6.0370$ & 659 & $-1$ & $-10$ & $-1.6354$ & $+8.3646$ & 993 & $1$ & $2$ & $-9.4010$ & $-11.4010$ \\
326 & $1$ & $-2$ & $2.8887$ & $+4.8887$ & 660 & $0$ & $-10$ & $-1.8368$ & $+8.1632$ & 994 & $-1$ & $1$ & $-9.1792$ & $-10.1792$ \\
327 & $1$ & $-1$ & $2.5974$ & $+3.5974$ & 661 & $-1$ & $-11$ & $-2.0268$ & $+8.9732$ & 995 & $1$ & $2$ & $-8.9932$ & $-10.9932$ \\
328 & $0$ & $-1$ & $2.1702$ & $+3.1702$ & 662 & $1$ & $-10$ & $-2.1887$ & $+7.8113$ & 996 & $0$ & $2$ & $-8.8463$ & $-10.8463$ \\
329 & $1$ & $0$ & $1.6280$ & $+1.6280$ & 663 & $-1$ & $-11$ & $-2.3077$ & $+8.6923$ & 997 & $-1$ & $1$ & $-8.7399$ & $-9.7399$ \\
330 & $1$ & $1$ & $1.0035$ & $+0.0035$ & 664 & $0$ & $-11$ & $-2.3706$ & $+8.6294$ & 998 & $1$ & $2$ & $-8.6731$ & $-10.6731$ \\
331 & $-1$ & $0$ & $0.3371$ & $+0.3371$ & 665 & $-1$ & $-12$ & $-2.3676$ & $+9.6324$ & 999 & $0$ & $2$ & $-8.6435$ & $-10.6435$ \\
332 & $0$ & $0$ & $-0.3269$ & $-0.3269$ & 666 & $0$ & $-12$ & $-2.2923$ & $+9.7077$ & 1000 & $0$ & $2$ & $-8.6463$ & $-10.6463$ \\
333 & $0$ & $0$ & $-0.9458$ & $-0.9458$ & 667 & $1$ & $-11$ & $-2.1424$ & $+8.8576$ &  &  &  &  &  \\
334 & $1$ & $1$ & $-1.4828$ & $-2.4828$ & 668 & $0$ & $-11$ & $-1.9197$ & $+9.0803$ &  &  &  &  &  \\
\end{longtable}}

\section{Python programs}
\label{app:code}

The programs below produce every number, table and figure of Section~\CGref{sec:numerics}.  They read the two data files \texttt{data/zeross.bin} (ordinates $\gamma_k$ of the zeros of $\zeta$) and \texttt{data/zLamb.bin} (midpoint ordinates $\tau_k$), both little-endian float64 with $2\,001\,052$ entries, and write their console output, CSV, MAT and binary files to \texttt{output/} and the figures to \texttt{figures/}.  They require Python~3 with NumPy, SciPy, mpmath and Matplotlib (\texttt{s15} also needs \texttt{finufft}).  The outputs distributed with the paper were produced with Python~3.13.16, NumPy~2.5.3, SciPy~1.18.1, mpmath~1.4.1, Matplotlib~3.11.2, finufft~2.5.1 and python-flint~0.9.0 (Arb ball arithmetic, used only by \texttt{code/routeB/certV.py}) on Linux x86\_64, in IEEE double precision except where mpmath or python-flint is called; \texttt{code/requirements.txt} pins these versions.  A double-precision ordinate is resolved to one unit in the last place, $2.3\cdot10^{-10}$ near the top of the files ($1.13\cdot10^6$) and $1.4\cdot10^{-14}$ near $t=100$; differences between computed quantities at or below the level of this representation of the data are not attributed to the algorithms.  The programs are run from the directory \texttt{code/} in the order \texttt{s01}, \texttt{s02} (which writes \texttt{output/s02\_varpi.txt}, read by \texttt{s05}, \texttt{s07} and \texttt{s15}), \texttt{s03}--\texttt{s10}, and then \texttt{s11}--\texttt{s17}; \texttt{s06} reads the CSV file written by \texttt{s04}, and the others depend only on the two data files.  The loader of \texttt{rvm\_common.py} refuses files whose size is not a multiple of $8$ bytes, files of different lengths, non-finite or non-positive values, a list $\tau_k$ that is not strictly increasing or starts below $2\pi$, and a list $\gamma_k$ that decreases; the programs \texttt{s08}--\texttt{s10} and \texttt{s13} use NumPy and SciPy only and check the lengths and the monotonicity themselves.

\subsection{\texttt{rvm\_common.py}}
\label{app:code-rvm-common}

\lstinputlisting[language=Python]{code/rvm_common.py}

\subsection{\texttt{s01\_clock\_and\_family.py}}
\label{app:code-s01-clock-and-family}

\lstinputlisting[language=Python]{code/s01_clock_and_family.py}

\subsection{\texttt{s02\_transition.py}}
\label{app:code-s02-transition}

\lstinputlisting[language=Python]{code/s02_transition.py}

\subsection{\texttt{s03\_exponential\_sums.py}}
\label{app:code-s03-exponential-sums}

\lstinputlisting[language=Python]{code/s03_exponential_sums.py}

\subsection{\texttt{s04\_primes.py}}
\label{app:code-s04-primes}

\lstinputlisting[language=Python]{code/s04_primes.py}

\subsection{\texttt{s05\_mertens.py}}
\label{app:code-s05-mertens}

\lstinputlisting[language=Python]{code/s05_mertens.py}

\subsection{\texttt{s06\_prime\_ratios.py}}
\label{app:code-s06-prime-ratios}

\lstinputlisting[language=Python]{code/s06_prime_ratios.py}

\subsection{\texttt{s07\_product\_identity.py}}
\label{app:code-s07-product-identity}

\lstinputlisting[language=Python]{code/s07_product_identity.py}


\subsection{\texttt{s08\_transport\_moments.py}}
\label{app:code-s08-transport-moments}

\lstinputlisting[language=Python]{code/s08_transport_moments.py}

\subsection{\texttt{s09\_clock\_defect\_split.py}}
\label{app:code-s09-clock-defect-split}

\lstinputlisting[language=Python]{code/s09_clock_defect_split.py}

\subsection{\texttt{s10\_clock\_to\_zeros.py}}
\label{app:code-s10-clock-to-zeros}

\lstinputlisting[language=Python]{code/s10_clock_to_zeros.py}


\subsection{\texttt{s11\_rvm\_z1\_difference\_analysis.py}}
\label{app:code-s11-rvm-z1}
This supplemental diagnostic reproduces the RVM--$Z1$ difference, shell and convergence calculations.  It uses the genuine ordinates only for post-construction diagnostics and is not used in any proof.
\lstinputlisting[language=Python]{code/s11_rvm_z1_difference_analysis.py}

\subsection{\texttt{s12\_multiscale\_jacobian.py}}
\label{app:code-s12-multiscale}
This program produces Experiment~\CGref{exp:multiscale-jacobian} and Figures~\CGref{fig:multiscale-gap} and its companion fixed-window diagnostic.  It generates the clock from the Lambert formula and does not read the genuine ordinates.
\lstinputlisting[language=Python]{code/s12_multiscale_jacobian.py}

\subsection{\texttt{s13\_nested\_bounded\_endpoint.py}}
\label{app:code-s13-bounded}
This program produces the bounded and nested endpoint diagnostics of Appendix~\CGref{sec:appendix-endpoint}.  The solver uses the midpoint clock; genuine ordinates and von Mangoldt masses, when displayed, are read only after the solve as diagnostics.
\lstinputlisting[language=Python]{code/s13_nested_bounded_endpoint.py}

\subsection{\texttt{s14\_varpi\_bound.py}}
\label{app:code-s14}
This program verifies the explicit inequalities in the proof of Theorem~\CGref{thm:varpi-less-one} in interval arithmetic (Part~A) and computes high-precision floating-point values of $\nu_4$, $P^{\crit}_{\RVM}(\frac12)$ and $\varpi_{\RVM}$ as a cross-check (Part~B, not part of the proof).  It uses only the Lambert formula for the clock and reads no data file.  Run time about $30$ seconds.
\lstinputlisting[language=Python]{code/s14_varpi_bound.py}

\subsection{\texttt{s15\_tables\_1000.py}}
\label{app:code-s15}
This program produces Experiment~\CGref{exp:first-1000}, Figures~\CGref{fig:clock-1000}--\CGref{fig:mertens-1000} and the Tables of Appendix~\CGref{app:tables-1000}.  It reads \texttt{data/zeross.bin} and \texttt{data/zLamb.bin}, and needs the package \texttt{finufft} (nonuniform fast Fourier transform) in addition to the packages of the other programs.  Run time about two minutes.  A MATLAB version, \texttt{matlab/s15\_tables\_1000.m}, is distributed with the repository.
\lstinputlisting[language=Python]{code/s15_tables_1000.py}

\subsection{\texttt{s16\_walks\_3d\_1000.py}}
\label{app:code-s16}
This program produces Experiment~\CGref{exp:walks} and Figures~\CGref{fig:walks3d} and~\CGref{fig:walks-proj}.  A MATLAB version, \texttt{matlab/s16\_walks\_3d\_1000.m}, is distributed with the repository.
\lstinputlisting[language=Python]{code/s16_walks_3d_1000.py}

\subsection{\texttt{s17\_cramer\_centroids.py}}
\label{app:code-s17}
This program evaluates Cram\'er's function by the zero-free representation of Proposition~\CGref{prop:cramer-zero-free} and the centroids of Experiment~\CGref{exp:cramer}.  Run time about one minute.
\lstinputlisting[language=Python]{code/s17_cramer_centroids.py}

\section{Console output of the programs}
\label{app:output}

The complete console output of each program, as written to \texttt{output/*\_console.txt}.

{\sloppy The numerical programs that produce the multiscale and bounded-endpoint figures (files \texttt{fig14}--\texttt{fig17}) and the bounded/nested endpoint tables are also distributed with their machine-readable outputs in the repository.  In particular, \texttt{s12\_\allowbreak multiscale\_\allowbreak jacobian.py} writes \texttt{output/\allowbreak s12\_\allowbreak multiscale\_\allowbreak jacobian\_\allowbreak console.txt}, \texttt{output/\allowbreak s12\_\allowbreak multiscale\_\allowbreak summary.csv} and \texttt{output/\allowbreak s12\_\allowbreak fixed\_\allowbreak window\_\allowbreak pairs.csv}; \texttt{s13\_\allowbreak nested\_\allowbreak bounded\_\allowbreak endpoint.py} writes \texttt{output/\allowbreak s13\_\allowbreak bounded\_\allowbreak mass\_\allowbreak summary.csv}, \texttt{output/\allowbreak s13\_\allowbreak bounded\_\allowbreak mass\_\allowbreak ordinates.csv}, \texttt{output/\allowbreak s13\_\allowbreak bounded\_\allowbreak mass\_\allowbreak masses.csv}, \texttt{output/\allowbreak s13\_\allowbreak cross\_\allowbreak validation.csv}, \texttt{output/\allowbreak s13\_\allowbreak nested\_\allowbreak metric\_\allowbreak summary.csv} and the console output \texttt{output/\allowbreak s13\_\allowbreak nested\_\allowbreak bounded\_\allowbreak endpoint\_\allowbreak console.txt}.  Here ``normalized residual'' means the Euclidean residual divided by the square root of the number of test equations.  The bump and permanent-block conventions are specified in the corresponding code and in Appendix~\CGref{sec:appendix-endpoint}.\par}

\subsection{\texttt{s01\_clock\_and\_family}}
\label{app:out-s01-clock-and-family}

\lstinputlisting[style=out]{output/s01_clock_and_family_console.txt}

\subsection{\texttt{s01b\_clock\_zoom}}
\label{app:out-s01b-clock-zoom}

\lstinputlisting[style=out]{output/s01b_clock_zoom_console.txt}

\subsection{\texttt{s02\_transition}}
\label{app:out-s02-transition}

\lstinputlisting[style=out]{output/s02_transition_console.txt}

\subsection{\texttt{s03\_exponential\_sums}}
\label{app:out-s03-exponential-sums}

\lstinputlisting[style=out]{output/s03_exponential_sums_console.txt}

\subsection{\texttt{s04\_primes}}
\label{app:out-s04-primes}

\lstinputlisting[style=out]{output/s04_primes_console.txt}

\subsection{\texttt{s05\_mertens}}
\label{app:out-s05-mertens}

\lstinputlisting[style=out]{output/s05_mertens_console.txt}

\subsection{\texttt{s06\_prime\_ratios}}
\label{app:out-s06-prime-ratios}

\lstinputlisting[style=out]{output/s06_prime_ratios_console.txt}

\subsection{\texttt{s07\_product\_identity}}
\label{app:out-s07-product-identity}

\lstinputlisting[style=out]{output/s07_product_identity_console.txt}

\subsection{\texttt{s08\_transport\_moments}}
\label{app:out-s08-transport-moments}

\lstinputlisting[style=out]{output/s08_transport_console.txt}

\subsection{\texttt{s09\_clock\_defect\_split}}
\label{app:out-s09-clock-defect-split}

\lstinputlisting[style=out]{output/s09_clock_defect_split_console.txt}

\subsection{\texttt{s10\_clock\_to\_zeros}}
\label{app:out-s10-clock-to-zeros}

\lstinputlisting[style=out]{output/s10_clock_to_zeros_console.txt}

\subsection{\texttt{s12\_multiscale\_jacobian}}
\label{app:out-s12-multiscale}

\lstinputlisting[style=out]{output/s12_multiscale_jacobian_console.txt}

\subsection{\texttt{s13\_nested\_bounded\_endpoint}}
\label{app:out-s13-bounded}

\lstinputlisting[style=out]{output/s13_nested_bounded_endpoint_console.txt}

\subsection{\texttt{s14\_varpi\_bound}}
\label{app:output-s14}
The file \texttt{output/s14\_varpi\_bound\_console.txt}; Part~A is the interval verification used in the proof of Theorem~\CGref{thm:varpi-less-one}.
\lstinputlisting[style=out]{output/s14_varpi_bound_console.txt}

\subsection{\texttt{s15\_tables\_1000}}
\label{app:output-s15}
\lstinputlisting[style=out]{output/s15_tables_1000_console.txt}

\subsection{\texttt{s16\_walks\_3d\_1000}}
\label{app:output-s16}
\lstinputlisting[style=out]{output/s16_walks_3d_1000_console.txt}

\subsection{\texttt{s17\_cramer\_centroids}}
\label{app:output-s17}
\lstinputlisting[style=out]{output/s17_cramer_centroids_console.txt}

\begin{thebibliography}{99}

\bibitem{Ahlfors1979}
L.~V. Ahlfors, \emph{Complex Analysis}, 3rd ed., McGraw--Hill, New York, 1979.

\bibitem{AriasdeReyna2024}
J.~Arias de Reyna, Explicit formula and quasicrystal definition, preprint, arXiv:2402.10604 (2024, v2 2025), \href{https://doi.org/10.48550/arXiv.2402.10604}{doi:10.48550/arXiv.2402.10604}.

\bibitem{HruscevNikolskiiPavlov1981}
S.~V. Hru\v{s}\v{c}ev, N.~K. Nikol'skii and B.~S. Pavlov, Unconditional bases of exponentials and of reproducing kernels, in: \emph{Complex Analysis and Spectral Theory (Leningrad, 1979/80)}, Lecture Notes in Math. \textbf{864}, Springer, Berlin, 1981, 214--335, \href{https://doi.org/10.1007/BFb0097000}{doi:10.1007/BFb0097000}.

\bibitem{KozmaLev2011}
G.~Kozma and N.~Lev, Exponential Riesz bases, discrepancy of irrational rotations and BMO, \emph{J. Fourier Anal. Appl.} \textbf{17} (2011), no.~5, 879--898, \href{https://doi.org/10.1007/s00041-011-9178-1}{doi:10.1007/s00041-011-9178-1}; arXiv:1009.2188.

\bibitem{BondarenkoRadchenkoSeip2023}
A.~Bondarenko, D.~Radchenko and K.~Seip, Fourier interpolation with zeros of zeta and $L$-functions, \emph{Constr. Approx.} \textbf{57} (2023), 405--461, \href{https://doi.org/10.1007/s00365-022-09599-w}{doi:10.1007/s00365-022-09599-w}; arXiv:2005.02996.

\bibitem{Johansson2017}
F.~Johansson, Arb: efficient arbitrary-precision midpoint-radius interval arithmetic, \emph{IEEE Trans. Comput.} \textbf{66} (2017), no.~8, 1281--1292, \href{https://doi.org/10.1109/TC.2017.2690633}{doi:10.1109/TC.2017.2690633}; arXiv:1611.02831.

\bibitem{Johansson2018}
F.~Johansson, Numerical integration in arbitrary-precision ball arithmetic, in: \emph{Mathematical Software -- ICMS 2018}, Lecture Notes in Comput. Sci. \textbf{10931}, Springer, Cham, 2018, 255--263, \href{https://doi.org/10.1007/978-3-319-96418-8_30}{doi:10.1007/978-3-319-96418-8\_30}; arXiv:1802.07942.

\bibitem{Bernstein1929}
S.~Bernstein, Sur les fonctions absolument monotones, \emph{Acta Math.} \textbf{52} (1929), 1--66, \href{https://doi.org/10.1007/BF02547400}{doi:10.1007/BF02547400}.

\bibitem{Widder1941}
D.~V. Widder, \emph{The Laplace Transform}, Princeton Mathematical Series 6, Princeton University Press, Princeton, 1941; reprint \href{https://doi.org/10.1515/9781400876457}{doi:10.1515/9781400876457}.

\bibitem{Conrey1989}
J.~B. Conrey, More than two fifths of the zeros of the Riemann zeta function are on the critical line, \emph{J. Reine Angew. Math.} \textbf{399} (1989), 1--26, \href{https://doi.org/10.1515/crll.1989.399.1}{doi:10.1515/crll.1989.399.1}.

\bibitem{Avdonin1974}
S.~A. Avdonin, On the question of Riesz bases of exponential functions in $L^2$, \emph{Vestnik Leningrad. Univ. Mat. Mekh. Astronom.} (1974), no.~13, 5--12. (No DOI.)

\bibitem{AvdoninIvanov2002}
S.~A. Avdonin and S.~A. Ivanov, Exponential Riesz bases of subspaces and divided differences, \emph{St. Petersburg Math. J.} \textbf{13} (2002), no.~3, 339--351, preprint version arXiv:math/0103160, \href{https://doi.org/10.48550/arXiv.math/0103160}{doi:10.48550/arXiv.math/0103160}.

\bibitem{BatemanDiamond1969}
P.~T. Bateman and H.~G. Diamond, Asymptotic distribution of Beurling's generalized prime numbers, in: \emph{Studies in Number Theory}, MAA Studies in Mathematics~6, Mathematical Association of America, 1969.

\bibitem{BenedettoPowellYilmaz2006}
J.~J. Benedetto, A.~M. Powell and \"O.~Y{\i}lmaz, Sigma-delta ($\Sigma\Delta$) quantization and finite frames, \emph{IEEE Trans. Inform. Theory} \textbf{52} (2006), no.~5, 1990--2005, \href{https://doi.org/10.1109/TIT.2006.872849}{doi:10.1109/TIT.2006.872849}.

\bibitem{Fujii1974}
A.~Fujii, On the zeros of Dirichlet $L$-functions. I, \emph{Trans. Amer. Math. Soc.} \textbf{196} (1974), 225--235, \href{https://doi.org/10.1090/S0002-9947-1974-0349603-2}{doi:10.1090/S0002-9947-1974-0349603-2}.

\bibitem{Fujii1975}
A.~Fujii, On the distribution of the zeros of the Riemann zeta function in short intervals, \emph{Bull. Amer. Math. Soc.} \textbf{81} (1975), 139--142, \href{https://doi.org/10.1090/S0002-9904-1975-13674-3}{doi:10.1090/S0002-9904-1975-13674-3}.

\bibitem{Fujii1975b}
A.~Fujii, On the difference between $r$ consecutive ordinates of the zeros of the Riemann zeta function, \emph{Proc. Japan Acad.} \textbf{51} (1975), no.~10, 741--743, \href{https://doi.org/10.3792/pja/1195518466}{doi:10.3792/pja/1195518466}.

\bibitem{Garnett2007}
J.~B. Garnett, \emph{Bounded Analytic Functions}, revised first ed., Graduate Texts in Mathematics 236, Springer, New York, 2007, \href{https://doi.org/10.1007/0-387-49763-3}{doi:10.1007/0-387-49763-3}.

\bibitem{Gonek2012}
S.~M. Gonek, Finite Euler products and the Riemann hypothesis, \emph{Trans. Amer. Math. Soc.} \textbf{364} (2012), 2157--2191, \href{https://doi.org/10.1090/S0002-9947-2011-05546-7}{doi:10.1090/S0002-9947-2011-05546-7}; arXiv:0704.3448.

\bibitem{Guinand1948}
A.~P. Guinand, A summation formula in the theory of prime numbers, \emph{Proc. London Math. Soc. (2)} \textbf{50} (1948), 107--119, \href{https://doi.org/10.1112/plms/s2-50.2.107}{doi:10.1112/plms/s2-50.2.107}.

\bibitem{Hormander1990}
L.~H\"ormander, \emph{The Analysis of Linear Partial Differential Operators I}, 2nd ed., Grundlehren der mathematischen Wissenschaften 256, Springer, Berlin, 1990, \href{https://doi.org/10.1007/978-3-642-61497-2}{doi:10.1007/978-3-642-61497-2}.

\bibitem{Ingham1936}
A.~E. Ingham, Some trigonometrical inequalities with applications to the theory of series, \emph{Math. Z.} \textbf{41} (1936), 367--379.

\bibitem{Beurling1937}
A.~Beurling, Analyse de la loi asymptotique de la distribution des nombres premiers g\'en\'eralis\'es.~I, \emph{Acta Math.} \textbf{68} (1937), 255--291, \href{https://doi.org/10.1007/BF02546666}{doi:10.1007/BF02546666}.

\bibitem{Bohr1926}
H.~Bohr, Zur Theorie der fastperiodischen Funktionen. III. Dirichletentwicklung analytischer Funktionen, \emph{Acta Math.} \textbf{47} (1926), 237--281, \href{https://doi.org/10.1007/BF02543846}{doi:10.1007/BF02543846}.

\bibitem{BochnerChandrasekharan1956}
S.~Bochner and K.~Chandrasekharan, On Riemann's functional equation, \emph{Ann. of Math. (2)} \textbf{63} (1956), 336--360, \href{https://doi.org/10.2307/1969614}{doi:10.2307/1969614}.

\bibitem{BrouckeDebruyneVindas2020}
F.~Broucke, G.~Debruyne and J.~Vindas, Beurling integers with RH and large oscillation, \emph{Adv. Math.} \textbf{370} (2020), 107240, \href{https://doi.org/10.1016/j.aim.2020.107240}{doi:10.1016/j.aim.2020.107240}.

\bibitem{BrouckeVindas2024}
F.~Broucke and J.~Vindas, A new generalized prime random approximation procedure and some of its applications, \emph{Math. Z.} \textbf{307} (2024), Paper No.~62, \href{https://doi.org/10.1007/s00209-024-03526-4}{doi:10.1007/s00209-024-03526-4}.

\bibitem{CohnElkies2003}
H.~Cohn and N.~Elkies, New upper bounds on sphere packings~I, \emph{Ann. of Math. (2)} \textbf{157} (2003), no.~2, 689--714, \href{https://doi.org/10.4007/annals.2003.157.689}{doi:10.4007/annals.2003.157.689}.

\bibitem{Corless1996}
R.~M. Corless, G.~H. Gonnet, D.~E.~G. Hare, D.~J. Jeffrey and D.~E. Knuth, On the Lambert $W$ function, \emph{Adv. Comput. Math.} \textbf{5} (1996), 329--359, \href{https://doi.org/10.1007/BF02124750}{doi:10.1007/BF02124750}.

\bibitem{DaubechiesDeVore2003}
I.~Daubechies and R.~DeVore, Approximating a bandlimited function using very coarsely quantized data: a family of stable sigma-delta modulators of arbitrary order, \emph{Ann. of Math.} \textbf{158} (2003), 679--710, \href{https://doi.org/10.4007/annals.2003.158.679}{doi:10.4007/annals.2003.158.679}.

\bibitem{Davenport1980}
H.~Davenport, \emph{Multiplicative Number Theory}, 2nd ed., revised by H.~L. Montgomery, Graduate Texts in Mathematics~74, Springer, New York, 1980, \href{https://doi.org/10.1007/978-1-4757-5927-3}{doi:10.1007/978-1-4757-5927-3}.

\bibitem{delaValleePoussin1896}
C.-J. de la Vall\'ee Poussin, Recherches analytiques sur la th\'eorie des nombres premiers, \emph{Ann. Soc. Sci. Bruxelles} \textbf{20} (1896), 183--256.

\bibitem{Diamond1970}
H.~G. Diamond, Asymptotic distribution of Beurling's generalized integers, \emph{Illinois J. Math.} \textbf{14} (1970), 12--28, \href{https://doi.org/10.1215/ijm/1256053295}{doi:10.1215/ijm/1256053295}.

\bibitem{DiamondMontgomeryVorhauer2006}
H.~G. Diamond, H.~L. Montgomery and U.~M.~A. Vorhauer, Beurling primes with large oscillation, \emph{Math. Ann.} \textbf{334} (2006), 1--36, \href{https://doi.org/10.1007/s00208-005-0638-2}{doi:10.1007/s00208-005-0638-2}.

\bibitem{DiamondZhang2016}
H.~G. Diamond and W.-B. Zhang, \emph{Beurling Generalized Numbers}, Mathematical Surveys and Monographs~213, American Mathematical Society, Providence, RI, 2016, \href{https://doi.org/10.1090/surv/213}{doi:10.1090/surv/213}.

\bibitem{DLMF2024}
NIST Digital Library of Mathematical Functions, F.~W.~J. Olver, A.~B. Olde Daalhuis, D.~W. Lozier et al. (eds.), Release 1.2.8 of 2026-09-15, \S4.13 (Lambert $W$-function) and \S5.11 (Stirling series), \href{https://dlmf.nist.gov/}{https://dlmf.nist.gov/}.

\bibitem{Edwards1974}
H.~M. Edwards, \emph{Riemann's Zeta Function}, Academic Press, New York, 1974.

\bibitem{FrancaLeClair2014}
G.~Fran\c{c}a and A.~LeClair, Statistical and other properties of Riemann zeros based on an explicit equation for the $n$-th zero on the critical line, preprint, arXiv:1307.8395 (2013), \href{https://doi.org/10.48550/arXiv.1307.8395}{doi:10.48550/arXiv.1307.8395}.

\bibitem{GrosswaldSchnitzer1978}
E.~Grosswald and F.~J. Schnitzer, A class of modified $\zeta$ and $L$-functions, \emph{Pacific J. Math.} \textbf{74} (1978), no.~2, 357--364, \href{https://doi.org/10.2140/pjm.1978.74.357}{doi:10.2140/pjm.1978.74.357}.

\bibitem{GunturkThao2003}
C.~S. G\"unt\"urk and N.~T. Thao, Ergodic dynamics in sigma-delta quantization: tiling invariant sets and spectral analysis of error, preprint, arXiv:math/0308223 (2003), \href{https://doi.org/10.48550/arXiv.math/0308223}{doi:10.48550/arXiv.math/0308223}.

\bibitem{Hadamard1896}
J.~Hadamard, Sur la distribution des z\'eros de la fonction $\zeta(s)$ et ses cons\'equences arithm\'etiques, \emph{Bull. Soc. Math. France} \textbf{24} (1896), 199--220, \href{https://doi.org/10.24033/bsmf.545}{doi:10.24033/bsmf.545}.

\bibitem{Hamburger1921}
H.~Hamburger, \"Uber die Riemannsche Funktionalgleichung der $\zeta$-Funktion, \emph{Math. Z.} \textbf{10} (1921), 240--254, \href{https://doi.org/10.1007/BF01211612}{doi:10.1007/BF01211612}.

\bibitem{KaczorowskiPerelli1999}
J.~Kaczorowski and A.~Perelli, On the structure of the Selberg class, I: $0\le d\le1$, \emph{Acta Math.} \textbf{182} (1999), 207--241, \href{https://doi.org/10.1007/BF02392574}{doi:10.1007/BF02392574}.

\bibitem{Kadec1964}
M.~I. Kadec, The exact value of the Paley--Wiener constant, \emph{Dokl. Akad. Nauk SSSR} \textbf{155} (1964), 1253--1254; English transl. Soviet Math. Dokl. \textbf{5} (1964), 559--561. (No DOI.)

\bibitem{Knopp1994}
M.~I. Knopp, On Dirichlet series satisfying Riemann's functional equation, \emph{Invent. Math.} \textbf{117} (1994), 361--372, \href{https://doi.org/10.1007/BF01232248}{doi:10.1007/BF01232248}.

\bibitem{KurasovSarnak2020}
P.~Kurasov and P.~Sarnak, Stable polynomials and crystalline measures, \emph{J. Math. Phys.} \textbf{61} (2020), no.~8, 083501, \href{https://doi.org/10.1063/5.0012286}{doi:10.1063/5.0012286}.

\bibitem{LapidusvanFrankenhuysen2000}
M.~L. Lapidus and M.~van Frankenhuysen, \emph{Fractal Geometry and Number Theory: Complex Dimensions of Fractal Strings and Zeros of Zeta Functions}, Birkh\"auser, Boston, 2000, \href{https://doi.org/10.1007/978-1-4612-5314-3}{doi:10.1007/978-1-4612-5314-3}.

\bibitem{LevOlevskii2015}
N.~Lev and A.~Olevskii, Quasicrystals and Poisson's summation formula, \emph{Invent. Math.} \textbf{200} (2015), no.~2, 585--606, \href{https://doi.org/10.1007/s00222-014-0542-z}{doi:10.1007/s00222-014-0542-z}.

\bibitem{LyubarskiiSeip1997}
Yu.~I. Lyubarskii and K.~Seip, Complete interpolating sequences for Paley--Wiener spaces and Muckenhoupt's $(A_p)$ condition, \emph{Rev. Mat. Iberoam.} \textbf{13} (1997), no.~2, 361--376, \href{https://doi.org/10.4171/RMI/224}{doi:10.4171/RMI/224}.

\bibitem{MontgomeryVaughan1974}
H.~L. Montgomery and R.~C. Vaughan, Hilbert's inequality, \emph{J. London Math. Soc. (2)} \textbf{8} (1974), 73--82, \href{https://doi.org/10.1112/jlms/s2-8.1.73}{doi:10.1112/jlms/s2-8.1.73}.

\bibitem{OlevskiiUlanovskii2020}
A.~Olevskii and A.~Ulanovskii, Fourier quasicrystals with unit masses, \emph{C. R. Math. Acad. Sci. Paris} \textbf{358} (2020), no.~11--12, 1207--1211, \href{https://doi.org/10.5802/crmath.142}{doi:10.5802/crmath.142}.

\bibitem{OdlyzkoTeRiele1985}
A.~M. Odlyzko and H.~J.~J. te~Riele, Disproof of the Mertens conjecture, \emph{J. reine angew. Math.} \textbf{357} (1985), 138--160, \href{https://doi.org/10.1515/crll.1985.357.138}{doi:10.1515/crll.1985.357.138}.

\bibitem{Nakamura2020}
T.~Nakamura, Dirichlet series with periodic coefficients, Riemann's functional equation and real zeros of Dirichlet $L$-functions, preprint, arXiv:2008.02570 (2020), \href{https://doi.org/10.48550/arXiv.2008.02570}{doi:10.48550/arXiv.2008.02570}.

\bibitem{Perelli2016}
A.~Perelli, Converse theorems: from the Riemann zeta function to the Selberg class, preprint, arXiv:1605.02354 (2016), \href{https://doi.org/10.48550/arXiv.1605.02354}{doi:10.48550/arXiv.1605.02354}.

\bibitem{PlattTrudgian2021}
D.~J. Platt and T.~S. Trudgian, The Riemann hypothesis is true up to $3\cdot10^{12}$, \emph{Bull. London Math. Soc.} \textbf{53} (2021), 792--797, \href{https://doi.org/10.1112/blms.12460}{doi:10.1112/blms.12460}; arXiv:2004.09765.

\bibitem{Riemann1859}
B.~Riemann, \"Uber die Anzahl der Primzahlen unter einer gegebenen Gr\"osse, \emph{Monatsberichte der K\"oniglich Preussischen Akademie der Wissenschaften zu Berlin} (1859), 671--680.

\bibitem{Rudin1987}
W.~Rudin, \emph{Real and Complex Analysis}, 3rd ed., McGraw--Hill, New York, 1987. (No DOI.)

\bibitem{Selberg1946}
A.~Selberg, Contributions to the theory of the Riemann zeta-function, \emph{Arch. Math. Naturvid.} \textbf{48} (1946), no.~5, 89--155. (No DOI.)

\bibitem{Siegel1922}
C.~L. Siegel, Bemerkung zu einem Satz von Hamburger \"uber die Funktionalgleichung der Riemannschen Zetafunktion, \emph{Math. Ann.} \textbf{86} (1922), 276--279, \href{https://doi.org/10.1007/BF01457989}{doi:10.1007/BF01457989}.

\bibitem{Titchmarsh1939}
E.~C. Titchmarsh, \emph{The Theory of Functions}, 2nd ed., Oxford University Press, 1939.

\bibitem{Titchmarsh1986}
E.~C. Titchmarsh, \emph{The Theory of the Riemann Zeta-Function}, 2nd ed., revised by D.~R. Heath-Brown, Clarendon Press, Oxford, 1986.

\bibitem{vonMangoldt1895}
H.~von Mangoldt, Zu Riemanns Abhandlung ``\"Uber die Anzahl der Primzahlen unter einer gegebenen Gr\"osse'', \emph{J. reine angew. Math.} \textbf{114} (1895), 255--305, \href{https://doi.org/10.1515/crll.1895.114.255}{doi:10.1515/crll.1895.114.255}.

\bibitem{Weil1952}
A.~Weil, Sur les ``formules explicites'' de la th\'eorie des nombres premiers, \emph{Comm. S\'em. Math. Univ. Lund [Medd. Lunds Univ. Mat. Sem.]} (1952), Tome Suppl\'ementaire, 252--265. (No DOI.)

\bibitem{WhittakerWatson1927}
E.~T. Whittaker and G.~N. Watson, \emph{A Course of Modern Analysis}, 4th ed., Cambridge University Press, 1927; reissued in the Cambridge Mathematical Library, 1996, \href{https://doi.org/10.1017/CBO9780511608759}{doi:10.1017/CBO9780511608759}.

\bibitem{Zemanian1965}
A.~H. Zemanian, \emph{Distribution Theory and Transform Analysis}, McGraw-Hill, New York, 1965; reprinted by Dover, New York, 1987. (No DOI.)

\bibitem{Cramer1919}
H.~Cram\'er, Studien \"uber die Nullstellen der Riemannschen Zetafunktion, \emph{Math. Z.} \textbf{4} (1919), 104--130, \href{https://doi.org/10.1007/BF01203390}{doi:10.1007/BF01203390}.

\bibitem{Landau1912}
E.~Landau, \"Uber die Nullstellen der Zetafunktion, \emph{Math. Ann.} \textbf{71} (1912), 548--564, \href{https://doi.org/10.1007/BF01456808}{doi:10.1007/BF01456808}.

\bibitem{Gonek1993}
S.~M. Gonek, An explicit formula of Landau and its applications to the theory of the zeta-function, in: \emph{A Tribute to Emil Grosswald: Number Theory and Related Analysis}, Contemp. Math. \textbf{143}, Amer. Math. Soc., Providence, RI, 1993, 395--413, \href{https://doi.org/10.1090/conm/143/1210528}{doi:10.1090/conm/143/1210528}.

\bibitem{DebruyneVindas2019}
G.~Debruyne and J.~Vindas, Complex Tauberian theorems for Laplace transforms with local pseudofunction boundary behavior, \emph{J. Anal. Math.} \textbf{138} (2019), 799--833, \href{https://doi.org/10.1007/s11854-019-0045-3}{doi:10.1007/s11854-019-0045-3}; arXiv:1604.05069.

\bibitem{Korevaar2004}
J.~Korevaar, \emph{Tauberian Theory: A Century of Developments}, Grundlehren der mathematischen Wissenschaften~329, Springer, Berlin, 2004, \href{https://doi.org/10.1007/978-3-662-10225-1}{doi:10.1007/978-3-662-10225-1}.

\bibitem{FrancaLeClair2015}
G.~Fran\c{c}a and A.~LeClair, Transcendental equations satisfied by the individual zeros of Riemann $\zeta$, Dirichlet and modular $L$-functions, \emph{Commun. Number Theory Phys.} \textbf{9} (2015), no.~1, 1--50, \href{https://doi.org/10.4310/CNTP.2015.v9.n1.a1}{doi:10.4310/CNTP.2015.v9.n1.a1}; arXiv:1502.06003.

\bibitem{HasanalizadeShenWong2022}
E.~Hasanalizade, Q.~Shen and P.-J. Wong, Counting zeros of the Riemann zeta function, \emph{J. Number Theory} \textbf{235} (2022), 219--241, \href{https://doi.org/10.1016/j.jnt.2021.06.032}{doi:10.1016/j.jnt.2021.06.032}; arXiv:2107.06506.

\bibitem{BellottiWong2025}
C.~Bellotti and P.-J. Wong, Improved estimates for the argument and zero-counting function of the Riemann zeta-function, \emph{Math. Comp.} (2025, published online), \href{https://doi.org/10.1090/mcom/4133}{doi:10.1090/mcom/4133}; arXiv:2412.15470.

\bibitem{Trudgian2014}
T.~S. Trudgian, An improved upper bound for the argument of the Riemann zeta-function on the critical line II, \emph{J. Number Theory} \textbf{134} (2014), 280--292, \href{https://doi.org/10.1016/j.jnt.2013.07.017}{doi:10.1016/j.jnt.2013.07.017}; arXiv:1208.5846.

\bibitem{Littlewood1924}
J.~E. Littlewood, On the zeros of the Riemann zeta-function, \emph{Proc. Cambridge Philos. Soc.} \textbf{22} (1924), 295--318, \href{https://doi.org/10.1017/S0305004100014225}{doi:10.1017/S0305004100014225}.

\bibitem{Hutchinson1925}
J.~I. Hutchinson, On the roots of the Riemann zeta function, \emph{Trans. Amer. Math. Soc.} \textbf{27} (1925), 49--60, \href{https://doi.org/10.1090/S0002-9947-1925-1501297-5}{doi:10.1090/S0002-9947-1925-1501297-5}.

\bibitem{Trudgian2011}
T.~S. Trudgian, On the success and failure of Gram's Law and the Rosser Rule, \emph{Acta Arith.} \textbf{148} (2011), 225--256, \href{https://doi.org/10.4064/aa148-3-2}{doi:10.4064/aa148-3-2}; arXiv:0811.0883.

\bibitem{Brent1979}
R.~P. Brent, On the zeros of the Riemann zeta function in the critical strip, \emph{Math. Comp.} \textbf{33} (1979), 1361--1372, \href{https://doi.org/10.1090/S0025-5718-1979-0537983-2}{doi:10.1090/S0025-5718-1979-0537983-2}.

\bibitem{Platt2017}
D.~J. Platt, Isolating some non-trivial zeros of zeta, \emph{Math. Comp.} \textbf{86} (2017), 2449--2467, \href{https://doi.org/10.1090/mcom/3198}{doi:10.1090/mcom/3198}.

\bibitem{Hlawka1975}
E.~Hlawka, \"Uber die Gleichverteilung gewisser Folgen, welche mit den Nullstellen der Zetafunktion zusammenh\"angen, \emph{Sitzungsber. \"Osterr. Akad. Wiss. Math.-Naturwiss. Kl. Abt.~II} \textbf{184} (1975), 459--471. (No DOI.)

\bibitem{FordZaharescu2005}
K.~Ford and A.~Zaharescu, On the distribution of imaginary parts of zeros of the Riemann zeta function, \emph{J. reine angew. Math.} \textbf{579} (2005), 145--158, \href{https://doi.org/10.1515/crll.2005.2005.579.145}{doi:10.1515/crll.2005.2005.579.145}; arXiv:math/0405459.

\bibitem{Guillera2013}
J.~Guillera, Some sums over the non-trivial zeros of the Riemann zeta function, preprint, arXiv:1307.5723 (2013), \href{https://doi.org/10.48550/arXiv.1307.5723}{doi:10.48550/arXiv.1307.5723}.

\bibitem{Hurst2018}
G.~Hurst, Computations of the Mertens function and improved bounds on the Mertens conjecture, \emph{Math. Comp.} \textbf{87} (2018), 1013--1028, \href{https://doi.org/10.1090/mcom/3275}{doi:10.1090/mcom/3275}; arXiv:1610.08551.

\bibitem{KotnikvandeLune2004}
T.~Kotnik and J.~van de Lune, On the order of the Mertens function, \emph{Experiment. Math.} \textbf{13} (2004), 473--481, \href{https://doi.org/10.1080/10586458.2004.10504556}{doi:10.1080/10586458.2004.10504556}.

\bibitem{Ng2004}
N.~Ng, The distribution of the summatory function of the M\"obius function, \emph{Proc. London Math. Soc. (3)} \textbf{89} (2004), 361--389, \href{https://doi.org/10.1112/S0024611504014741}{doi:10.1112/S0024611504014741}; arXiv:math/0310381.

\bibitem{Zhang2013}
W.-B. Zhang, Normalization of Beurling generalized primes with Riemann hypothesis, \emph{Ann. Univ. Sci. Budapest. Sect. Comput.} \textbf{39} (2013), 459--469, \href{https://doi.org/10.71352/ac.39.459}{doi:10.71352/ac.39.459}.

\bibitem{BrouckeDebruyneRevesz2025}
F.~Broucke, G.~Debruyne and Sz.~Gy. R\'ev\'esz, Some examples of well-behaved Beurling number systems, \emph{Trans. Amer. Math. Soc.} \textbf{378} (2025), 477--501, \href{https://doi.org/10.1090/tran/9274}{doi:10.1090/tran/9274}; arXiv:2309.01567.

\bibitem{HilberdinkLapidus2006}
T.~W. Hilberdink and M.~L. Lapidus, Beurling zeta functions, generalised primes, and fractal membranes, \emph{Acta Appl. Math.} \textbf{94} (2006), 21--48, \href{https://doi.org/10.1007/s10440-006-9063-0}{doi:10.1007/s10440-006-9063-0}; arXiv:math/0410270.

\bibitem{BalazardSaiasYor1999}
M.~Balazard, \'E.~Saias and M.~Yor, Notes sur la fonction $\zeta$ de Riemann, 2, \emph{Adv. Math.} \textbf{143} (1999), 284--287, \href{https://doi.org/10.1006/aima.1998.1797}{doi:10.1006/aima.1998.1797}.

\bibitem{Volchkov1995}
V.~V. Volchkov, On an equality equivalent to the Riemann hypothesis, \emph{Ukrainian Math. J.} \textbf{47} (1995), 491--493, \href{https://doi.org/10.1007/BF01056314}{doi:10.1007/BF01056314}.

\bibitem{SekatskiiBeltraminelliMerlini2012}
S.~K. Sekatskii, S.~Beltraminelli and D.~Merlini, On equalities involving integrals of the logarithm of the Riemann $\zeta$-function and equivalent to the Riemann hypothesis, \emph{Ukrainian Math. J.} \textbf{64} (2012), 247--261, \href{https://doi.org/10.1007/s11253-012-0642-0}{doi:10.1007/s11253-012-0642-0}.

\bibitem{DavenportHeilbronn1936a}
H.~Davenport and H.~Heilbronn, On the zeros of certain Dirichlet series (I), \emph{J. London Math. Soc.} \textbf{11} (1936), 181--185, \href{https://doi.org/10.1112/jlms/s1-11.3.181}{doi:10.1112/jlms/s1-11.3.181}.

\bibitem{DavenportHeilbronn1936b}
H.~Davenport and H.~Heilbronn, On the zeros of certain Dirichlet series (II), \emph{J. London Math. Soc.} \textbf{11} (1936), 307--312, \href{https://doi.org/10.1112/jlms/s1-11.4.307}{doi:10.1112/jlms/s1-11.4.307}.

\end{thebibliography}


\clearpage
\section{Python programs of Section~\CGref{sec:routeB-finite}}
\label{app:code-routeB}

The programs of \texttt{code/routeB/} produce the numbers of Section~\CGref{subsec:rb-numerics} and Figures~\CGref{fig:rb-gonek} and~\CGref{fig:rb-scaling}. They use only the clock, the primes and the Gonek lists; \texttt{zeros.py} computes the zeta ordinates used in the diagnostic columns with \texttt{mpmath} (at $20$ digits, stored as double-precision numbers).  They are run from \texttt{output/routeB/} in the order \texttt{zeros.py}, \texttt{gonek\_fast.py}, \texttt{horizonG.py}, \texttt{cert7.py}, \texttt{diag\_cert7\_gamma.py}, \texttt{entry.py}, \texttt{entry2.py}, \texttt{certV.py}, \texttt{laplace\_check.py}, \texttt{figs.py} and \texttt{results\_summary.py}; the last one writes \texttt{RESULTS.txt} from the JSON outputs, with the SHA-256 of the programs and of the files read.  All computations are in double precision except \texttt{certV.py}, which works in ball arithmetic (python-flint~0.9.0, Arb, $128$ bits) and proves Theorem~\CGref{thm:rb-verified}; a stage marked \texttt{FLOAT\_PASS} by the other programs satisfies the Banach inequalities in floating point, which is a candidate certificate and not a proof.  \texttt{entry2.py} repeats the entry test of \texttt{entry.py} in $\ell^\infty$ and in $\ell^2$.  The program \texttt{cert7.py} does not read the zeta ordinates; their distance to the polished points is computed separately by \texttt{diag\_cert7\_gamma.py}. The exploratory scripts (\texttt{staged.py}, \texttt{lsq*.py}, \texttt{basin*.py}, \texttt{coupling.py}, \texttt{horizon.py}, \texttt{certG.py}, \texttt{cert2--4.py}) are distributed in the same directory with their logs and are described in \texttt{RESULTS.txt}; \texttt{cert5.py} and \texttt{cert6.py} are variants of \texttt{cert7.py} (\texttt{cert5.py} evaluates the test at the receding-horizon iterate without polishing, \texttt{cert6.py} after the same polishing but without the sampled margin), with the same floating-point labels and defect term, and with the zeta ordinates used only in a final diagnostic block.

\subsection{\texttt{core.py}}
\label{app:code-routeB-core}

\lstinputlisting[language=Python]{code/routeB/core.py}

\subsection{\texttt{zeros.py}}
\label{app:code-routeB-zeros}

\lstinputlisting[language=Python]{code/routeB/zeros.py}

\subsection{\texttt{gonek\_fast.py}}
\label{app:code-routeB-gonek-fast}

\lstinputlisting[language=Python]{code/routeB/gonek_fast.py}

\subsection{\texttt{horizonG.py}}
\label{app:code-routeB-horizonG}

\lstinputlisting[language=Python]{code/routeB/horizonG.py}

\subsection{\texttt{cert7.py}}
\label{app:code-routeB-cert7}

\lstinputlisting[language=Python]{code/routeB/cert7.py}

\subsection{\texttt{entry.py}}
\label{app:code-routeB-entry}

\lstinputlisting[language=Python]{code/routeB/entry.py}

\subsection{\texttt{entry2.py}}
\label{app:code-routeB-entry2}

\lstinputlisting[language=Python]{code/routeB/entry2.py}

\subsection{\texttt{certV.py}}
\label{app:code-routeB-certV}

\lstinputlisting[language=Python]{code/routeB/certV.py}

\subsection{\texttt{laplace\_check.py}}
\label{app:code-routeB-laplace-check}

\lstinputlisting[language=Python]{code/routeB/laplace_check.py}

\subsection{\texttt{diag\_cert7\_gamma.py}}
\label{app:code-routeB-diag}

\lstinputlisting[language=Python]{code/routeB/diag_cert7_gamma.py}

\subsection{\texttt{figs.py}}
\label{app:code-routeB-figs}

\lstinputlisting[language=Python]{code/routeB/figs.py}

\subsection{\texttt{results\_summary.py}}
\label{app:code-routeB-summary}

\lstinputlisting[language=Python]{code/routeB/results_summary.py}

\section{Console output of the programs of Section~\CGref{sec:routeB-finite}}
\label{app:out-routeB}

\subsection{\texttt{RESULTS.txt}}
\label{app:out-routeB-results}

\lstinputlisting[style=out]{output/routeB/RESULTS.txt}

\subsection{\texttt{out\_horizonG.txt}}
\label{app:out-routeB-horizonG}

\lstinputlisting[style=out]{output/routeB/out_horizonG.txt}

\subsection{\texttt{out\_cert7.txt}}
\label{app:out-routeB-cert7}

\lstinputlisting[style=out]{output/routeB/out_cert7.txt}

\subsection{\texttt{out\_cert7\_gamma.txt}}
\label{app:out-routeB-cert7-gamma}

\lstinputlisting[style=out]{output/routeB/out_cert7_gamma.txt}

\subsection{\texttt{out\_entry.txt}}
\label{app:out-routeB-entry}

\lstinputlisting[style=out]{output/routeB/out_entry.txt}

\subsection{\texttt{out\_entry2.txt}}
\label{app:out-routeB-entry2}

\lstinputlisting[style=out]{output/routeB/out_entry2.txt}

\subsection{\texttt{out\_certV.txt}}
\label{app:out-routeB-certV}

\lstinputlisting[style=out]{output/routeB/out_certV.txt}

\subsection{\texttt{out\_certG.txt}}
\label{app:out-routeB-certG}

\lstinputlisting[style=out]{output/routeB/out_certG.txt}

\subsection{\texttt{out\_cert3.txt}}
\label{app:out-routeB-cert3}

\lstinputlisting[style=out]{output/routeB/out_cert3.txt}

\subsection{\texttt{out\_gonek.txt} and \texttt{out\_gonek\_fast.txt}}
\label{app:out-routeB-gonek}

\lstinputlisting[style=out]{output/routeB/out_gonek.txt}
\lstinputlisting[style=out]{output/routeB/out_gonek_fast.txt}

\subsection{\texttt{out\_laplace\_check.txt}}
\label{app:out-routeB-laplace-check}

\lstinputlisting[style=out]{output/routeB/out_laplace_check.txt}

\section{Endpoint-only closure, interpolation tests, and autonomous residual questions}\label{sec:appendix-endpoint}

This appendix records the interpolation and finite-section tests.  It does not alter any theorem above.  Its purpose is to separate the statements that are exact but equivalent to the final step from the one clock-only subproblem that remains autonomous.

\subsection{Three interpolations and their status}\label{subsec:three-interpolations}

There are three natural ways to interpolate between the midpoint clock and the Riemann zeta function.

\emph{The ordinate interpolation.}  The family $\zeta_\lambda$ of Sections~\CGref{sec:lambda-family}--\CGref{sec:criteria} moves the ordinates while keeping every nontrivial zero on $\Re s=\frac12$.  It is an exact analytic family, but for every $0<\lambda\le1$ its continuous prime measure is not in the ordinary Dirichlet category.  Thus the property that is needed at the endpoint is not preserved along the path; this is exactly the role of condition (D).

\emph{The exact-atomic interpolation.}  One may instead ask for real ordinates $\kappa_k$ and masses $\mathsf a_n$ supported at $u=\log n$ throughout an interpolation.  Finite sections of this problem were tested numerically.  The endpoints can be resolved, whereas the linear interpolation of the target distribution develops a persistent interior residual in the tested sections.  This is evidence against that particular interpolation, not a proof that no infinite exact-atomic path exists.  In any case the existence of such a path is not supplied by the results of this paper.

\emph{The additive interpolation of functions.}  There is an exact elementary interpolation
\begin{equation}\label{eq:additive-family}
 \mathcal G_\varepsilon(s):=(1-\varepsilon)\zeta_{\RVM}(s)+\varepsilon\zeta(s),\qquad 0\le\varepsilon\le1.
\end{equation}
Both summands satisfy the same functional equation and reality symmetry, so the same is true of $\mathcal G_\varepsilon$.  This family exists for every $\varepsilon$, but it does not create a new control quantity: it is a level-set representation of the already defined quotient $\mathcal Q_{\RVM}=\zeta/\zeta_{\RVM}$.

\begin{proposition}[The additive family is a level-set reformulation]\label{prop:additive-level}
Away from the zeros of $\zeta_{\RVM}$,
\begin{equation}\label{eq:additive-level}
 \mathcal G_\varepsilon(s)=\zeta_{\RVM}(s)\bigl((1-\varepsilon)+\varepsilon\mathcal Q_{\RVM}(s)\bigr).
\end{equation}
Hence for $0<\varepsilon<1$ every zero of $\mathcal G_\varepsilon$ which is not a zero of $\zeta_{\RVM}$ lies on the level set
\begin{equation}\label{eq:Q-negative-level}
 \mathcal Q_{\RVM}(s)=-\frac{1-\varepsilon}{\varepsilon}.
\end{equation}
On the critical line $e^{i\vartheta(t)}\mathcal G_\varepsilon(\frac12+it)$ is real.  At $\varepsilon=1$ the statement that all its nontrivial zeros lie on the critical line is exactly RH.
\end{proposition}

\begin{proof}
Equation~\CGeqref{eq:additive-level} is the definition of $\mathcal Q_{\RVM}$.  The functional equation and the reality statement follow by linearity from those of the two summands.  The last assertion is $\mathcal G_1=\zeta$.
\end{proof}

Thus the additive family is useful for visualizing the zero motion, but without an additional monotonicity or flow law it is circular as a route to RH: the terminal level is precisely the original zero problem.

\subsection{The clock as the first transition term}\label{subsec:first-transition}

The transition theorem gives the following $1+\chi$ form on every closed subband of $\tfrac12<\Re s<1$:
\begin{equation}\label{eq:rvm-first-rs}
 \zeta_{\RVM}(s)=\varpi_{\RVM}(1+\chi(s))\bigl(1+O(|t|^{-1})\bigr).
\end{equation}
Independently, on $s=\frac12+it$, where $\chi(s)=e^{-2i\vartheta(t)}$, the exact trigonometric identity
\begin{equation}\label{eq:first-rs-hardy}
 e^{i\vartheta(t)}(1+\chi(\tfrac12+it))=2\cos\vartheta(t).
\end{equation}
Consequently the midpoint clock is asymptotically the zero set selected by the $n=1$ term of the Hardy--Littlewood/Riemann--Siegel approximate functional equation.  This is a structural interpretation of the transition theorem, not an exact identification: the $O(1/t)$ term is visible already in the first ordinates.

For a finite Dirichlet polynomial
\begin{equation}\label{eq:finite-additive}
 \mathsf P_{\varepsilon,M}(s):=1+\varepsilon\sum_{2\le n\le M}n^{-s},\qquad
 \mathsf R_{\varepsilon,M}(s):=\mathsf P_{\varepsilon,M}(s)+\chi(s)\mathsf P_{\varepsilon,M}(1-s),
\end{equation}
the functional equation is exact.  On the critical line,
\begin{equation}\label{eq:finite-additive-hardy}
 e^{i\vartheta(t)}\mathsf R_{\varepsilon,M}(\tfrac12+it)
 =2\Re\!\left(e^{i\vartheta(t)}\mathsf P_{\varepsilon,M}(\tfrac12+it)\right).
\end{equation}
At Gram points the $n=1$ term fixes the sign whenever
\begin{equation}\label{eq:dominance}
 \varepsilon\sum_{2\le n\le M}n^{-1/2}<1,
\end{equation}
so in every interval where $\vartheta$ is monotone there is a critical-line zero between consecutive Gram points.  This dominance statement is only sufficient; it does not count all zeros of $\mathsf R_{\varepsilon,M}$ and is not used elsewhere in the paper.

\subsection{What remains genuinely live}\label{subsec:live}

After these tests the status is narrower than the programme suggested by the interpolation pictures.

\begin{enumerate}[label=\textup{(L\arabic*)}]
\item \emph{The endpoint problem remains the endpoint problem.}  Proving that the cell defect vanishes at $\lambda=0$, equivalently proving (D), remains the only step that closes the main route.  The additive family and the exact-atomic interpolation do not simplify this statement.

\item \emph{The clock-only injectivity problem is autonomous.}  The one subproblem that does not contain the genuine ordinates is the injectivity of the linearized rung system at the midpoint clock.  In the meromorphic form of Remark~\CGref{rem:nondegenerate-scope}, the target statement is: if a bounded real sequence $(\dot\kappa_k)$ and coefficients $\dot{\mathsf a}_n=O_\eps(n^{-1/2+\eps})$ give a Dirichlet series $\mathsf H$ satisfying
\begin{equation}\label{eq:clock-injectivity}
 \mathsf H(z)=\nu_8+\sum_{k\ge1}\dot\kappa_k\left(\frac{\tau_k}{z^2+\tau_k^2}-\frac1{\tau_k}\right)
     =\sum_{n\ge2}\frac{\dot{\mathsf a}_n}{\log n}\,n^{-z}
\end{equation}
through a common right half-plane, under the growth hypotheses of the linearization, then $\mathsf H\equiv0$.  Theorem~\CGref{thm:nondegenerate} proves the analogous conclusion at an actual solution, while Remark~\CGref{rem:nondegenerate-scope} explains why it is not automatic for a general list.  Equation~\CGeqref{eq:clock-injectivity} is therefore a closed clock-only rigidity problem, naturally related to Fourier crystalline uniqueness, but by itself it does not prove (D).

\item \emph{Bounded finite sections are a numerical existence programme, not yet an analytic route.}  Direct finite sections of the endpoint rung system can be solved while constraining both the ordinate displacement and the free masses.  If one could prove, for an exhausting nested family of tests, a uniform boundedness and residual-to-zero statement and then pass to a distributional limit, Theorem~\CGref{thm:free-masses} would identify the limit.  Such a theorem would already be an existence theorem for the final system and is therefore not a shortcut around (D).
\end{enumerate}

In particular, the individual displacements $\gamma_k-\tau_k$ are not autonomous data.  They are used below only for post-solution numerical diagnostics.  None of the live statement \CGeqref{eq:clock-injectivity}, nor the definition of the bounded finite-section problem, requires them as input.


\subsection{A priori bounded displacement for exact rung solutions}\label{subsec:rung-apriori}

The following statement is independent of the existence of an interior exact-atomic interpolation.  It says only that, if an exact real solution exists, has the stated local density and masses of the stated mean square, then its count follows the midpoint clock up to a linear drift governed by one constant $\mathrm b$; its ordinates stay within a bounded distance of a drift-corrected clock, and within a bounded distance of the midpoint clock itself exactly when $\mathrm b=0$.  Throughout, the Laplace transform of a distribution $\mathscr T$ supported in $[0,\infty)$ is $\mathcal L\mathscr T(s)=\langle\mathscr T,\chi e^{-(s-\frac12)u}\rangle$, as in the proof of Lemma~\CGref{lem:rb-laplace}; thus $\mathcal L\delta_{\log n}(s)=n^{\frac12-s}$ and $\mathcal L\delta_0'(s)=s-\frac12$.

\begin{proposition}[A priori displacement bound for an exact rung solution]\label{prop:rung-apriori}
Fix $\mathsf t\in[0,1]$.  Let $\boldsymbol{\kappa}=(\kappa_k)_{k\ge1}$ satisfy $0<\kappa_1\le\kappa_2\le\cdots$ and $\kappa_k\to\infty$, and let $(\mathsf a_n)_{n\ge2}$ be real masses satisfying, in $\mathcal D'(0,\infty)$,
\begin{equation}\label{eq:rung-apriori-system}
 \sum_{k\ge1}\cos(\kappa_k u)
 -(1-\mathsf t)\sum_{k\ge1}\cos(\tau_k u)-\mathsf t\,\mathsf k(u)
 =-\sum_{n\ge2}\mathsf a_n\,\delta(u-\log n).
\end{equation}
Assume that for every $\epsilon_0>0$ the local count and the masses satisfy
\begin{equation}\label{eq:rung-apriori-density}
 N_{\boldsymbol{\kappa}}(T+1)-N_{\boldsymbol{\kappa}}(T)\ll_{\epsilon_0}(1+T)^{\epsilon_0}\qquad(T\ge0),
\end{equation}
\begin{equation}\label{eq:rung-apriori-mass}
 \sum_{X<n\le2X}|\mathsf a_n|^2\ll_{\epsilon_0}X^{1+\epsilon_0}\qquad(X\ge1).
\end{equation}
Then the limit
\begin{equation}\label{eq:rung-apriori-b}
 \mathrm b:=\lim_{\sigma\to+\infty}\Bigl[\sum_{k\ge1}\frac{\sigma-\frac12}{(\sigma-\frac12)^2+\kappa_k^2}-(1-\mathsf t)\sum_{k\ge1}\frac{\sigma-\frac12}{(\sigma-\frac12)^2+\tau_k^2}-\mathsf t\,\frac{G'}{2G}(\sigma)\Bigr]
\end{equation}
exists and is real.  Fix $\epsilon_0\in(0,\frac14)$ and a class of such solutions, possibly with different $\mathsf t\in[0,1]$, on which the following four quantities are bounded: the implied constants in \CGeqref{eq:rung-apriori-density} and \CGeqref{eq:rung-apriori-mass} for this $\epsilon_0$, the number $\sum_k\kappa_k^{-2}$ and the number $|\mathrm b|$.  There are $C^*$ and $k_0$, depending only on these four bounds and not on $\mathsf t$, such that every solution of the class satisfies
\begin{equation}\label{eq:rung-apriori-count}
 \Bigl|N_{\boldsymbol{\kappa}}(T)-n_{\RVM}(T)-\frac{2\mathrm b}{\pi}T\Bigr|\le C^*\log(2+T)\qquad(T\ge0)
\end{equation}
and, for every $k\ge k_0$, with some $\mathrm s_k$ between $\kappa_k$ and $\tau_k$,
\begin{equation}\label{eq:rung-apriori-bound}
 \Bigl|\kappa_k-\tau_k+\frac{4\mathrm b\,\kappa_k}{\log(\mathrm s_k/2\pi)}\Bigr|\le C^*,\qquad
 |\kappa_k-\tau^{(\mathrm b)}_k|\le C^*,
\end{equation}
where $\tau^{(\mathrm b)}_k$ is the drift-corrected clock point, the unique solution $x\ge x_{\mathrm b}:=2\pi e^{8|\mathrm b|}$ of $\NR(x)+\frac{2\mathrm b}\pi x=k-\frac12$.  In particular $\kappa_k-\tau_k=-(4\mathrm b+o(1))\kappa_k/\log\kappa_k$ as $k\to\infty$, so the displacement $\kappa_k-\tau_k$ is bounded if and only if $\mathrm b=0$; and if $\mathrm b=0$, then $|\kappa_k-\tau_k|\le C^*$ for all $k\ge1$.  The constant $\mathrm b$ is the mass at $u=0$ of the distribution on $[0,\infty)$ that extends the left side of \CGeqref{eq:rung-apriori-system} with the extensions of the proof, and the identity on the open half-line does not determine it: a list whose count exceeds the clock by $2\mathrm b/\pi$ per unit height satisfies \CGeqref{eq:rung-apriori-b} with that very $\mathrm b$.
\end{proposition}

\begin{proof}
\emph{Order of $P_{\boldsymbol{\kappa}}$.}  By \CGeqref{eq:rung-apriori-density}, $N_{\boldsymbol{\kappa}}(T)\le\sum_{0\le j\le T}\bigl(N_{\boldsymbol{\kappa}}(j+1)-N_{\boldsymbol{\kappa}}(j)\bigr)\ll_{\epsilon_0}(1+T)^{1+\epsilon_0}$ for every $\epsilon_0>0$, and $N_{\boldsymbol{\kappa}}(t)=0$ for $t<\kappa_1$.  Hence $\sum_k\kappa_k^{-2}=\int t^{-2}dN_{\boldsymbol{\kappa}}(t)<\infty$; this number is not controlled by the density constant alone, because $\kappa_1$ may be small, and it is therefore listed separately.  The symmetric canonical product $P_{\boldsymbol{\kappa}}(s)=\prod_k(1+q(s)/(\frac14+\kappa_k^2))$ converges locally uniformly, is entire and real on $\mathbb R$, is positive for real $s>1$ (every factor exceeds $1$ there, since $q(s)>0$), and has order at most $1$.  For $|s|\ge1$, $|q(s)|\le2|s|^2=:R^2$; since $\frac14+\kappa_k^2\ge\kappa_k^2$,
\[
 \log|P_{\boldsymbol{\kappa}}(s)|\le\sum_k\log\Bigl(1+\frac{R^2}{\kappa_k^2}\Bigr)=\int_0^\infty\log\Bigl(1+\frac{R^2}{t^2}\Bigr)dN_{\boldsymbol{\kappa}}(t)=\int_{\kappa_1}^\infty N_{\boldsymbol{\kappa}}(t)\frac{2R^2}{t(t^2+R^2)}\,dt ,
\]
the boundary terms vanishing because $N_{\boldsymbol{\kappa}}(t)=0$ for $t<\kappa_1$ and $N_{\boldsymbol{\kappa}}(t)\log(1+R^2/t^2)\ll t^{1+\epsilon_0}R^2t^{-2}\to0$ as $t\to\infty$.  On $[\kappa_1,1]$ (if $\kappa_1<1$) the integrand is at most $2N_{\boldsymbol{\kappa}}(1)/t$, which contributes $O(1)$; on $[1,\infty)$, with $N_{\boldsymbol{\kappa}}(t)\ll t^{1+\epsilon_0}$ and $t=Rv$, the integral is $\ll R^{1+\epsilon_0}\int_0^\infty\frac{2v^{\epsilon_0}}{1+v^2}dv\ll|s|^{1+\epsilon_0}$.  Since $\epsilon_0>0$ is arbitrary, $\log|P_{\boldsymbol{\kappa}}(s)|=O_\epsilon(|s|^{1+\epsilon})$ for every $\epsilon>0$, with constants that may depend on the solution; only this qualitative growth is used below.

\emph{From $\mathcal D'(0,\infty)$ to the Laplace transform.}  The identity \CGeqref{eq:rung-apriori-system} holds on the open half-line and says nothing about the point $u=0$; the exponential $e^{-(s-1/2)u}$ is not a test function of $\mathcal D(0,\infty)$, so it cannot be paired with \CGeqref{eq:rung-apriori-system} directly.  We first extend each term to a distribution on $\mathbb R$ supported in $[0,\infty)$, as in Lemma~\CGref{lem:rb-laplace}.  Put $\mathfrak H_{\boldsymbol{\kappa}}(u):=\sum_k\kappa_k^{-2}\cos(\kappa_ku)$, continuous and bounded, and $\mathscr S_{\boldsymbol{\kappa}}:=-\partial_u^2(\mathfrak H_{\boldsymbol{\kappa}}\mathbf 1_+)+\mathfrak H_{\boldsymbol{\kappa}}(0)\delta_0'$: a distribution of order at most $2$, supported in $[0,\infty)$, whose restriction to $(0,\infty)$ is $\sum_k\cos(\kappa_ku)$ (termwise differentiation of a uniformly convergent series in $\mathcal D'(0,\infty)$).  Its Laplace transform is computed as in \CGeqref{eq:rb-aux-14}, which uses only $\sum_k\kappa_k^{-2}<\infty$: with $z=s-\frac12$, $\mathcal L(\partial_u\mathscr T)=z\,\mathcal L\mathscr T$ for $\mathscr T$ supported in $[0,\infty)$, $\mathcal L\delta_0'=z$, and $\mathcal L(\mathfrak H_{\boldsymbol{\kappa}}\mathbf 1_+)=\sum_k\kappa_k^{-2}z/(z^2+\kappa_k^2)$ by termwise integration, so
\begin{equation}\label{eq:rung-laplace-S}
 \mathcal L\mathscr S_{\boldsymbol{\kappa}}(s)=-z^2\sum_k\frac{z}{\kappa_k^2(z^2+\kappa_k^2)}+z\sum_k\frac1{\kappa_k^2}=\sum_{k\ge1}\frac{z}{z^2+\kappa_k^2}=\tfrac12\,\frac{P_{\boldsymbol{\kappa}}'}{P_{\boldsymbol{\kappa}}}(s)\qquad(\Re s>\tfrac12):
\end{equation}
the sign of the correction $+\mathfrak H_{\boldsymbol{\kappa}}(0)\delta_0'$ is the one that removes the term $-z\mathfrak H_{\boldsymbol{\kappa}}(0)$, and no multiple of $\delta_0$ arises.  The same extension $\mathscr S_{\boldsymbol\tau}$ serves the clock, with $P^{\crit}_{\RVM}$ in place of $P_{\boldsymbol{\kappa}}$.  The term $\mathsf k$ is extended by the Archimedean distribution $\mathscr K$ of Lemma~\CGref{lem:rb-archimedean}, with $\mathcal L\mathscr K=\frac12G'/G$.  The mass term $\mathscr A:=\sum_{n\ge2}\mathsf a_n\delta_{\log n}$ is a locally finite measure on $\mathbb R$ supported in $[\log2,\infty)$, with no mass at $0$.  For $\sigma_0>\frac32$ choose $\epsilon_0:=\sigma_0-\frac32$, which satisfies $\epsilon_0<2\sigma_0-3$; by \CGeqref{eq:rung-apriori-mass} for this $\epsilon_0$ and Cauchy--Schwarz on dyadic blocks,
\[
 \sum_{X<n\le2X}|\mathsf a_n|n^{1/2-\sigma_0}\le\Bigl(\sum_{X<n\le2X}|\mathsf a_n|^2\Bigr)^{1/2}\Bigl(\sum_{X<n\le2X}n^{1-2\sigma_0}\Bigr)^{1/2}
 \ll_{\epsilon_0}X^{3/2-\sigma_0+\epsilon_0/2},
\]
a negative power of $X$; so $\sum_n|\mathsf a_n|n^{1/2-\sigma}<\infty$ for every $\sigma>\frac32$, and $\mathcal L\mathscr A(s)=\sum_{n\ge2}\mathsf a_nn^{\frac12-s}$ converges absolutely and uniformly on every half-plane $\Re s\ge\sigma_0>\frac32$.  The distribution
\begin{equation}\label{eq:rung-E}
 \mathscr E:=\mathscr S_{\boldsymbol{\kappa}}-(1-\mathsf t)\mathscr S_{\boldsymbol\tau}-\mathsf t\,\mathscr K+\mathscr A
\end{equation}
is supported in $[0,\infty)$, and by \CGeqref{eq:rung-apriori-system} its restriction to $(0,\infty)$ vanishes.  Hence $\mathscr E$ is supported at $\{0\}$; its restriction to a bounded neighbourhood of $0$ has finite order (each term has finite order on bounded intervals), so $\mathscr E=\sum_{j\le J}c_j\delta_0^{(j)}$ with real $c_j$ \cite[Thm.~2.3.4]{Hormander1990}, whose Laplace transform is the real polynomial $p(z)=\sum_{j\le J}c_jz^j$, $z=s-\frac12$.  Taking Laplace transforms of \CGeqref{eq:rung-E} on $\Re s>\frac32$,
\begin{equation}\label{eq:rung-laplace-identity}
 \tfrac12\frac{P_{\boldsymbol{\kappa}}'}{P_{\boldsymbol{\kappa}}}(s)-(1-\mathsf t)\,\tfrac12\frac{(P^{\crit}_{\RVM})'}{P^{\crit}_{\RVM}}(s)-\mathsf t\,\tfrac12\frac{G'}{G}(s)+\sum_{n\ge2}\mathsf a_nn^{\frac12-s}=p\bigl(s-\tfrac12\bigr)\qquad(\Re s>\tfrac32).
\end{equation}
The polynomial $p$ is the price of the passage from the open half-line to the closed one (Remark~\CGref{rem:rb-boundary-term}).

\emph{The polynomial is a constant.}  On the real axis, with $z=\sigma-\frac12\ge1$, each term on the left of \CGeqref{eq:rung-laplace-identity} is $O(z^{\epsilon_0})$ as $z\to+\infty$, with the fixed $\epsilon_0<\frac14$.  Integrating by parts against the count (the boundary term $N_{\boldsymbol{\kappa}}(t)\,z/(z^2+t^2)$ vanishes at $t=0$ and, since $N_{\boldsymbol{\kappa}}(t)\ll(1+t)^{1+\epsilon_0}$, at $t\to\infty$),
\[
 \sum_k\frac{z}{z^2+\kappa_k^2}=\int_0^\infty N_{\boldsymbol{\kappa}}(t)\frac{2zt}{(z^2+t^2)^2}\,dt\ll\int_0^\infty(1+t)^{1+\epsilon_0}\frac{2zt\,dt}{(z^2+t^2)^2}\le z^{\epsilon_0}\int_0^\infty\frac{2(1+v)^{1+\epsilon_0}v\,dv}{(1+v^2)^2}\ll z^{\epsilon_0}
\]
(substitute $t=zv$ and use $1+zv\le z(1+v)$), with the implied constant that of \CGeqref{eq:rung-apriori-density}; the clock term is $O(\log z)$ by the same computation with $n_{\RVM}$; $\frac12G'/G(\sigma)=O(\log\sigma)$ by Stirling's formula; and the mass sum tends to $0$ by absolute convergence.  A polynomial in $z$ that is $O(z^{\epsilon_0})$ as $z\to+\infty$ with $\epsilon_0<1$ has degree zero, so $p\equiv\mathrm b$ with $\mathrm b\in\mathbb R$, and \CGeqref{eq:rung-laplace-identity} read at real $\sigma\to+\infty$ is \CGeqref{eq:rung-apriori-b}: the limit exists and equals the constant $\mathrm b$.

\emph{Integration.}  All terms of \CGeqref{eq:rung-laplace-identity} are real on $(\frac32,\infty)$, where $P_{\boldsymbol{\kappa}}$, $P^{\crit}_{\RVM}$ and $G(\sigma)=\sigma(\sigma-1)\Gamma(\frac\sigma2)\pi^{-\sigma/2}$ are positive.  Since $\frac d{ds}\bigl(\frac{2\mathsf a_n}{\log n}n^{1/2-s}\bigr)=-2\mathsf a_nn^{1/2-s}$, the function
\begin{equation}\label{eq:rung-Y}
 Y_\mathsf t(s)=\exp\!\left(\sum_{n\ge2}\frac{2\mathsf a_n}{\log n}\,n^{1/2-s}\right)
\end{equation}
has logarithmic derivative $-2\sum_n\mathsf a_nn^{1/2-s}$, so twice \CGeqref{eq:rung-laplace-identity} says $(\log P_{\boldsymbol{\kappa}})'=(1-\mathsf t)(\log P^{\crit}_{\RVM})'+\mathsf t(\log G)'+(\log Y_\mathsf t)'+2\mathrm b$.  Integrating from $2$ to $s$ along a path in $\Re s>\frac32$ and exponentiating,
\begin{equation}\label{eq:rung-three-factor}
 P_{\boldsymbol{\kappa}}(s)
 =C_\mathsf t\bigl(P_{\RVM}^{\crit}(s)\bigr)^{1-\mathsf t}\,G(s)^{\mathsf t}\,Y_\mathsf t(s)\,e^{2\mathrm b(s-2)}
 =C_\mathsf t\,G(s)\,\zeta_{\RVM}(s)^{1-\mathsf t}\,Y_\mathsf t(s)\,e^{2\mathrm b(s-2)}\qquad(\Re s>\tfrac32),
\end{equation}
with $C_\mathsf t:=P_{\boldsymbol{\kappa}}(2)\big/\bigl(G(2)\zeta_{\RVM}(2)^{1-\mathsf t}Y_\mathsf t(2)\bigr)>0$, the value that makes both sides agree at $s=2$.  The fractional powers are defined globally, not along paths: the half-plane $\Re s>\frac32$ is simply connected, $P^{\crit}_{\RVM}$ has no zeros there (its zeros lie on $\Re s=\frac12$), $G$ has no zeros there (its only zeros are $s=0$ and $s=1$), and both are positive on $(\frac32,\infty)$; so each has a unique holomorphic logarithm on $\Re s>\frac32$ that is real on the real axis, and $(P^{\crit}_{\RVM})^{1-\mathsf t}:=\exp((1-\mathsf t)\log P^{\crit}_{\RVM})$, $G^{\mathsf t}:=\exp(\mathsf t\log G)$ with these logarithms.  The integration from $2$ to $s$ is then path-independent, and the zero count below is performed on the entire function $P_{\boldsymbol{\kappa}}$ itself, so no branch of a fractional factor enters it.  By the dyadic bound above with $\sigma_0=2$ and the fixed $\epsilon_0<\frac14$ (any $\epsilon_0<1$ would do here), the series in \CGeqref{eq:rung-Y} converges absolutely and uniformly on $\Re s\ge2$ and
\begin{equation}\label{eq:rung-B}
 \sup_{\Re s\ge2}|\log Y_\mathsf t(s)|\le B_2:=\sum_{n\ge2}\frac{2|\mathsf a_n|}{\log n}n^{-3/2}<\infty,
\end{equation}
with $B_2$ bounded in terms of the implied constant in \CGeqref{eq:rung-apriori-mass} for the fixed $\epsilon_0$.  Since $1\le P_{\boldsymbol{\kappa}}(2)=\prod_k(1+2/(\frac14+\kappa_k^2))\le e^{2\sum_k\kappa_k^{-2}}$, $G(2)=2/\pi$ and $|\log\zeta_{\RVM}(2)|\le\mathrm c_{\RVM}$,
\[
 |\log C_\mathsf t|\le2\sum_k\kappa_k^{-2}+\log\tfrac\pi2+\mathrm c_{\RVM}+B_2 ,
\]
a bound in terms of $\sum_k\kappa_k^{-2}$ and $B_2$ only.  No displacement bound is used here.

\emph{The right half-strip.}  Put $\mathbb W_{\boldsymbol{\kappa}}:=\log|P_{\boldsymbol{\kappa}}|-\log|G|$.  On the right half-strip $2\le\Re s\le4$, \CGeqref{eq:rung-three-factor}, Theorem~\CGref{thm:halfplane-RVM} ($|\log\zeta_{\RVM}|\le\mathrm c_{\RVM}$ on $\Re s\ge2$) and \CGeqref{eq:rung-B} give
\begin{equation}\label{eq:rung-line}
 |\mathbb W_{\boldsymbol{\kappa}}(s)|\le\mathbb K_{\rm E}:=|\log C_\mathsf t|+\mathrm c_{\RVM}+B_2+4|\mathrm b|\qquad(2\le\Re s\le4),
\end{equation}
uniformly in $\mathsf t$ and $\Im s$, because $\log|e^{2\mathrm b(s-2)}|=2\mathrm b(\Re s-2)$ lies in $[-4|\mathrm b|,4|\mathrm b|]$ there; on the line $\Re s=2$ itself this term vanishes.  The disc used below lies in $0\le\Re s\le4$; its part in $2\le\Re s\le4$ is covered by \CGeqref{eq:rung-line}.

\emph{The half-strip to the left.}  Let $\mathbb S:=\{-1\le\Re s\le2,\ \Im s\ge1\}$.  $\mathbb W_{\boldsymbol{\kappa}}$ is subharmonic in the interior of $\mathbb S$ and upper semicontinuous on $\mathbb S$, since $\log|G|$ is harmonic there ($G$ has no zeros or poles with $\Im s\ge1$).  The majorant $\mathrm h_0(s):=\frac32\log|s+2|$ is harmonic and nonnegative on $\mathbb S$.  We check the three sides.  \emph{Right side:} $\mathbb W_{\boldsymbol{\kappa}}(2+it)\le\mathbb K_{\rm E}$ by \CGeqref{eq:rung-line}.  \emph{Left side:} the symmetry $P_{\boldsymbol{\kappa}}(1-s)=P_{\boldsymbol{\kappa}}(s)$ and the reality of $P_{\boldsymbol{\kappa}}$ give $P_{\boldsymbol{\kappa}}(-1+it)=P_{\boldsymbol{\kappa}}(2-it)=\overline{P_{\boldsymbol{\kappa}}(2+it)}$, so $\log|P_{\boldsymbol{\kappa}}(-1+it)|=\log|P_{\boldsymbol{\kappa}}(2+it)|$.  The factor $G$ has no such symmetry, and \CGeqref{eq:rb-G-ratio} gives, for $t\ge1$,
\[
 \mathbb W_{\boldsymbol{\kappa}}(-1+it)=\mathbb W_{\boldsymbol{\kappa}}(2+it)+\log\Bigl|\frac{G(2+it)}{G(-1+it)}\Bigr|\le\mathbb K_{\rm E}+\tfrac32\log(1+t)\le\mathbb K_{\rm E}+\tfrac34\log2+\mathrm h_0(-1+it),
\]
because $1+t\le\sqrt2|1+it|$ and $|s+2|=|1+it|$ at $s=-1+it$.  \emph{Bottom side:} for $s=\sigma+i$ with $-1\le\sigma\le2$, $|q(s)|=|s|\,|s-1|\le5$, so $\log|P_{\boldsymbol{\kappa}}(s)|\le\sum_k\log(1+5/\kappa_k^2)\le5\sum_k\kappa_k^{-2}$, and $-\log|G(\sigma+i)|\le\mathrm c_7$, an absolute constant ($G$ is continuous and zero-free on that segment).  \emph{Growth:} $\log|P_{\boldsymbol{\kappa}}(s)|=O(|s|^{1+\epsilon})$ by the first step, and $-\log|G(s)|\le\frac\pi4|\Im s|+O(\log|s|)$ on $\mathbb S$ by Stirling's formula \cite[\S12.33]{WhittakerWatson1927}; hence $\mathbb W_{\boldsymbol{\kappa}}-\mathrm h_0\le\mathrm A_0e^{c\Im s}$ on $\mathbb S$ for every $c>0$, with $\mathrm A_0$ depending on the solution.  Lemma~\CGref{lem:rb-PL}(a), with width $3$ and $c<\pi/3$, applied to $\mathbb W_{\boldsymbol{\kappa}}-\mathrm h_0$, gives
\begin{equation}\label{eq:rung-strip}
 \mathbb W_{\boldsymbol{\kappa}}(s)\le\mathbb K_{\rm E}''+\tfrac32\log|s+2|\qquad(s\in\mathbb S),\qquad
 \mathbb K_{\rm E}'':=\max\Bigl(\mathbb K_{\rm E}+\tfrac34\log2,\ 5\sum_k\kappa_k^{-2}+\mathrm c_7\Bigr).
\end{equation}
The constant $\mathrm A_0$ enters only the applicability of the lemma, not its conclusion, so \CGeqref{eq:rung-strip} is uniform over the class; the compact part of the strip needs no separate bound, because the bottom side $\Im s=1$ is part of $\partial\mathbb S$.

\emph{Horizontal variation of the argument.}  Fix $T\ge T_0$ with $T\notin\{\kappa_k\}$, where $T_0\ge3$ is the absolute constant of Lemma~\CGref{lem:rb-backlund}.  As in that lemma, choose $\mathrm a_T\in\mathbb R$ with $e^{i\mathrm a_T}P_{\boldsymbol{\kappa}}(2+iT)>0$ and put $\mathfrak F_T(z):=\frac12\bigl(e^{i\mathrm a_T}P_{\boldsymbol{\kappa}}(z+iT)+e^{-i\mathrm a_T}P_{\boldsymbol{\kappa}}(z-iT)\bigr)$, an entire function with $\mathfrak F_T(x)=\Re(e^{i\mathrm a_T}P_{\boldsymbol{\kappa}}(x+iT))$ for real $x$ and $\mathfrak F_T(2)=|P_{\boldsymbol{\kappa}}(2+iT)|>0$.  Let $n_T$ be the number of zeros of $\mathfrak F_T$ on $[\frac12,2]$, with multiplicity; the segment splits into at most $n_T+1$ pieces on which $\Re(e^{i\mathrm a_T}P_{\boldsymbol{\kappa}})$ has constant sign, so the variation of $\arg P_{\boldsymbol{\kappa}}$ along the segment from $2+iT$ to $\frac12+iT$ is at most $\pi(n_T+1)$ \cite[\S9.4]{Titchmarsh1986}.  Jensen's formula for $\mathfrak F_T$ on the disc $|z-2|\le2$, valid also if $\mathfrak F_T$ has zeros on the circle (they contribute $\log1=0$) \cite[Thm.~15.18]{Rudin1987}, in which the zeros on $[\frac12,2]\subset\{|z-2|\le\frac32\}$ contribute at least $\log\frac43$ each and the others contribute nonnegatively, together with the bound of the mean over the circle by the maximum, gives
\[
 n_T\log\tfrac43\le\max_{|z-2|=2}\log|\mathfrak F_T(z)|-\log|\mathfrak F_T(2)|\le\max_{|s-(2+iT)|\le2}\log|P_{\boldsymbol{\kappa}}(s)|-\log|P_{\boldsymbol{\kappa}}(2+iT)|,
\]
because $|P_{\boldsymbol{\kappa}}(z-iT)|=|P_{\boldsymbol{\kappa}}(\bar z+iT)|$.  On the disc $\mathbb D_T:=\{|s-(2+iT)|\le2\}$ write $\log|P_{\boldsymbol{\kappa}}|=\mathbb W_{\boldsymbol{\kappa}}+\log|G|$.  The part of $\mathbb D_T$ with $\Re s\le2$ lies in $\mathbb S$ (there $\Im s\ge T-2\ge1$) and has $|s+2|\le T+6$, so \CGeqref{eq:rung-strip} gives $\mathbb W_{\boldsymbol{\kappa}}\le\mathbb K_{\rm E}''+\frac32\log(T+6)$; the part with $\Re s\ge2$ has $\mathbb W_{\boldsymbol{\kappa}}\le\mathbb K_{\rm E}$ by \CGeqref{eq:rung-line}; and at the centre $\mathbb W_{\boldsymbol{\kappa}}(2+iT)\ge-\mathbb K_{\rm E}$.  The oscillation of $\log|G|$ on $\mathbb D_T$ is at most $\log(T+8)+2\mathrm c_5$ (proof of Lemma~\CGref{lem:rb-backlund}).  Hence
\begin{equation}\label{eq:rung-nT}
 n_T\le\frac{\mathbb K_{\rm E}+\mathbb K_{\rm E}''+\frac52\log(T+8)+2\mathrm c_5}{\log(4/3)},
\end{equation}
the numerical constants $6$ and $8$ coming from deliberately coarse bounds for $|s+2|$ and for Stirling's estimate on $\mathbb D_T$.  The variation of $\arg G$ along the same segment is at most $\mathrm c_6=\frac32$ for $T\ge T_0$ (proof of Theorem~\CGref{thm:rb-count}).

\emph{Vertical variation of the argument.}  On the segment from $2$ to $2+iT$, by \CGeqref{eq:rung-three-factor},
\[
 \arg\zeta_{\boldsymbol{\kappa}}(2+it)=\arg\frac{P_{\boldsymbol{\kappa}}}{G}(2+it)=(1-\mathsf t)\Im\log\zeta_{\RVM}(2+it)+\Im\log Y_\mathsf t(2+it)+2\mathrm b\,t ,
\]
with the branches chosen continuous along the segment and real at $t=0$ (the branch of $\log\zeta_{\RVM}$ is the one of Theorem~\CGref{thm:halfplane-RVM}, continuous on $\Re s\ge2$, and $\log Y_{\mathsf t}$ is the series \CGeqref{eq:rung-Y} itself).  What the argument principle uses is the \emph{net} variation of the argument along the segment, that is, the difference of the values of the continuous branch at its two ends, not its total variation; the first two terms are bounded by $\mathrm c_{\RVM}$ and $B_2$ in modulus at both ends, so the net variation along the segment is $2\mathrm bT+O(2\mathrm c_{\RVM}+2B_2)$.  This is where the factor $e^{2\mathrm b(s-2)}$ enters the count: its argument is not bounded, it grows linearly.

\emph{Zero count.}  Lemma~\CGref{lem:rb-argument} applies verbatim to $P_{\boldsymbol{\kappa}}$: it uses only that $P_{\boldsymbol{\kappa}}$ is real on $\mathbb R$, symmetric under $s\mapsto1-s$, with zeros exactly at $\frac12\pm i\kappa_k$ and none on the contour, the last because $T\notin\{\kappa_k\}$.  So $\pi N_{\boldsymbol{\kappa}}(T)=\pi+\vartheta(T)+\Delta_{\mathbb G_T}\arg\zeta_{\boldsymbol{\kappa}}$ along the path $2\to2+iT\to\frac12+iT$, with $\zeta_{\boldsymbol\kappa}:=P_{\boldsymbol\kappa}/G$.  Combining the two variations, and using $|1+\vartheta(T)/\pi-\NR(T)|\le1/T$ for $T\ge T_0$ (Stirling; enlarge $T_0$ if necessary) and $|n_{\RVM}(T)-\NR(T)|\le\frac12$ for $T\ge2\pi$, since $n_{\RVM}(T)=\lfloor\NR(T)+\frac12\rfloor$,
\begin{equation}\label{eq:rung-count}
 \Bigl|N_{\boldsymbol{\kappa}}(T)-n_{\RVM}(T)-\frac{2\mathrm b}\pi T\Bigr|
 \le n_T+1+\frac{2\mathrm c_{\RVM}+2B_2+\mathrm c_6}{\pi}+\frac1{T}+\frac12\qquad(T\ge T_0,\ T\notin\{\kappa_k\}),
\end{equation}
which with \CGeqref{eq:rung-nT} is \CGeqref{eq:rung-apriori-count} for these $T$, with $C^*$ depending only on $\mathbb K_{\rm E}$ and $\mathbb K_{\rm E}''$, hence only on the four bounds of the statement.  It extends to $T\in\{\kappa_k\}$, $T\ge T_0$, by right continuity of $N_{\boldsymbol{\kappa}}$ and $n_{\RVM}$, letting $T$ decrease to $\kappa_k$ through values outside the discrete set $\{\kappa_k\}$; and to $0\le T<T_0$ because there $N_{\boldsymbol{\kappa}}(T)\le N_{\boldsymbol{\kappa}}(T_0)\le n_{\RVM}(T_0)+2|\mathrm b|T_0/\pi+C^*\log(2+T_0)$, while $n_{\RVM}(T)$ and $2|\mathrm b|T/\pi$ are bounded by constants of the same kind; enlarging $C^*$ gives \CGeqref{eq:rung-apriori-count} for all $T\ge0$.  The argument uses $P_{\boldsymbol{\kappa}}$, which is entire, so no branch of a fractional factor enters the zero count.

\emph{From the count to the displacement.}  Put $\mathcal F_{\mathrm b}(x):=\NR(x)+\frac{2\mathrm b}\pi x$.  For $x\ge x_{\mathrm b}=2\pi e^{8|\mathrm b|}$, $\mathcal F_{\mathrm b}'(x)=\frac1{2\pi}\log\frac x{2\pi}+\frac{2\mathrm b}\pi\ge\frac1{4\pi}\log\frac x{2\pi}$, so $\mathcal F_{\mathrm b}$ is strictly increasing on $[x_{\mathrm b},\infty)$ and tends to $\infty$.  Since $N_{\boldsymbol{\kappa}}(\kappa_k)\ge k$ and $N_{\boldsymbol{\kappa}}(\kappa_k-0)\le k-1$ (multiplicities included), \CGeqref{eq:rung-apriori-count} at $T=\kappa_k$ and its left limit, together with $|n_{\RVM}(T\pm0)-\NR(T)|\le1$ for all $T>0$, give
\begin{equation}\label{eq:rung-quantile}
 \bigl|\mathcal F_{\mathrm b}(\kappa_k)-k\bigr|\le C^*\log(2+\kappa_k)+2\qquad(k\ge1).
\end{equation}
The count bound also places the ordinates uniformly: $\kappa_k>T$ whenever $k>n_{\RVM}(T)+2|\mathrm b|T/\pi+C^*\log(2+T)$, and $\kappa_k\le T$ whenever $n_{\RVM}(T)-2|\mathrm b|T/\pi-C^*\log(2+T)\ge k$; so for every $X$ there is $k_0$, and for every $k$ there is a bound for $\kappa_k$, depending only on $X$, $k$, $C^*$ and $|\mathrm b|$.  Fix $X\ge x_{\mathrm b}$, depending only on $C^*$ and $|\mathrm b|$, so large that for $x\ge X$
\[
 \frac x{4\pi}\log\frac x{2\pi}\le\mathcal F_{\mathrm b}(x)\le\frac x\pi\log\frac x{2\pi},\qquad C^*\log(2+x)+\tfrac52\le\frac x{8\pi}\log\frac x{2\pi},
\]
and let $k_0$ be such that $\kappa_k\ge X$, $\tau_k\ge X$ and $k-\frac12\ge\mathcal F_{\mathrm b}(X)$ for $k\ge k_0$.  For $k\ge k_0$, \CGeqref{eq:rung-quantile} and the two-sided bound for $\mathcal F_{\mathrm b}$ give $\kappa_k\log\kappa_k\asymp k\asymp\tau_k\log\tau_k$, hence $\kappa_k\asymp\tau_k$ and $\kappa_k\asymp\tau^{(\mathrm b)}_k$, with absolute implied constants; enlarging $X$, every point $\mathrm s$ between two of $\kappa_k,\tau_k,\tau^{(\mathrm b)}_k$ satisfies $\log(\mathrm s/2\pi)\ge\frac12\log(\kappa_k/2\pi)\ge1$.  The mean-value theorem for $\mathcal F_{\mathrm b}$ between $\kappa_k$ and $\tau^{(\mathrm b)}_k$ gives
\[
 |\kappa_k-\tau^{(\mathrm b)}_k|\le\frac{|\mathcal F_{\mathrm b}(\kappa_k)-(k-\frac12)|}{\min\mathcal F_{\mathrm b}'}\le\frac{4\pi\bigl(C^*\log(2+\kappa_k)+\frac52\bigr)}{\frac12\log(\kappa_k/2\pi)},
\]
which is bounded; this is the second inequality in \CGeqref{eq:rung-apriori-bound}, after enlarging $C^*$.  For the first, write $\NR(\kappa_k)-\NR(\tau_k)=\mathcal F_{\mathrm b}(\kappa_k)-(k-\frac12)-\frac{2\mathrm b}\pi\kappa_k=-\frac{2\mathrm b}\pi\kappa_k+E_k$ with $|E_k|\le C^*\log(2+\kappa_k)+\frac52$; the mean-value theorem for $\NR$ gives
\[
 (\kappa_k-\tau_k)\,\frac1{2\pi}\log\frac{\mathrm s_k}{2\pi}=-\frac{2\mathrm b}\pi\kappa_k+E_k,\qquad\mathrm s_k\ \text{between}\ \kappa_k\ \text{and}\ \tau_k ,
\]
so $\kappa_k-\tau_k+4\mathrm b\kappa_k/\log(\mathrm s_k/2\pi)=2\pi E_k/\log(\mathrm s_k/2\pi)$, of modulus at most $4\pi(C^*\log(2+\kappa_k)+\frac52)/\log(\kappa_k/2\pi)$, which is bounded; this is the first inequality in \CGeqref{eq:rung-apriori-bound}, after enlarging $C^*$.  The intermediate point $\mathrm s_k$ is kept: $\log(\mathrm s_k/2\pi)=\log(\kappa_k/2\pi)+O(1)$ only gives $\kappa_k-\tau_k=-(4\mathrm b+o(1))\kappa_k/\log\kappa_k$, and replacing $\mathrm s_k$ by $\kappa_k$ in the denominator would change the main term by $O(\mathrm b^2\kappa_k/\log^3\kappa_k)$, which is not bounded when $\mathrm b\ne0$.  If $\mathrm b=0$, the first inequality reads $|\kappa_k-\tau_k|\le C^*$ for $k\ge k_0$, and for $k<k_0$, $|\kappa_k-\tau_k|\le\kappa_{k_0}+\tau_{k_0}$, which is bounded by the uniform placement of the ordinates in terms of $k_0$ and $C^*$; enlarging $C^*$ gives the bound for all $k$.  All constants depend only on the implied constants of \CGeqref{eq:rung-apriori-density} and \CGeqref{eq:rung-apriori-mass} for the fixed $\epsilon_0$, on $\sum_k\kappa_k^{-2}$ and on $|\mathrm b|$, and not on $\mathsf t$.
\end{proof}

\begin{corollary}[Exact solutions without mass at the origin]\label{cor:rung-apriori-b0}
Under the hypotheses of Proposition~\CGref{prop:rung-apriori}, if the extension of the left side of \CGeqref{eq:rung-apriori-system} constructed in the proof has no mass at $u=0$, that is, $\mathrm b=0$, then $|N_{\boldsymbol{\kappa}}(T)-n_{\RVM}(T)|\le C^*\log(2+T)$ and $|\kappa_k-\tau_k|\le C^*$ for all $T\ge0$ and $k\ge1$, with $C^*$ depending only on the implied constants of \CGeqref{eq:rung-apriori-density}--\CGeqref{eq:rung-apriori-mass} for the fixed $\epsilon_0$ and on $\sum_k\kappa_k^{-2}$.
\end{corollary}

The mechanism behind the drift is elementary: a mass $\mathrm b\delta_0$ in the extended identity adds the constant $\mathrm b$ to the Laplace residual, hence the factor $e^{2\mathrm b(s-2)}$ to $P_{\boldsymbol{\kappa}}$, whose argument grows linearly along vertical lines; the argument principle converts this into an excess of $2\mathrm bT/\pi$ zeros up to height $T$, and dividing by the clock density $\NR'(\mathrm s_k)=\frac1{2\pi}\log\frac{\mathrm s_k}{2\pi}$ turns the excess into a displacement $-4\mathrm b\kappa_k/\log(\mathrm s_k/2\pi)=-(4\mathrm b+o(1))\kappa_k/\log\kappa_k$, up to a bounded error.  The mean-square hypothesis \CGeqref{eq:rung-apriori-mass} on the masses is independent of the local-density hypothesis and is not implied by it for a general list; both are assumptions on the solution, and the proposition is an a priori statement about hypothetical solutions, not an existence theorem.

\begin{remark}\label{rem:rung-apriori-scope}
Proposition~\CGref{prop:rung-apriori} is an a priori theorem, not an existence theorem.  The exact-atomic interpolation has no established interior solution for $0<\mathsf t<1$; the proposition only says that any exact solution satisfying its hypotheses lies in a uniform drift-corrected class, with ordinates within $C^*$ of the drift-corrected clock points $\tau^{(\mathrm b)}_k$; it has uniformly bounded displacement from the midpoint clock only when $\mathrm b=0$.  The mean-square bound \CGeqref{eq:rung-apriori-mass} on the masses is a hypothesis: for the lists $(\gamma_k(\lambda))$ it is Theorem~\CGref{thm:mean-square-cells}, whose proof uses the tube through Lemma~\CGref{lem:mean-value-ordinates} and the term \CGeqref{eq:cell-mass-power-two}; for a general list satisfying only \CGeqref{eq:rung-apriori-density} it is not derived here, and the proposition assumes it.  The constant $\mathrm b$ of \CGeqref{eq:rung-apriori-b} is the mass at $u=0$ that the identity on the open half-line leaves undetermined; the argument of $e^{2\mathrm b(s-2)}$ along $\Re s=2$ grows linearly and produces the drift $-4\mathrm b\kappa_k/\log(\mathrm s_k/2\pi)$ in \CGeqref{eq:rung-apriori-bound}, so the proposition bounds the displacement from the midpoint clock only for solutions with $\mathrm b=0$, that is, for solutions of the identity extended to $[0,\infty)$ without a mass at the origin.  The zero count follows the mechanism of Lemmas~\CGref{lem:rb-backlund} and~\CGref{lem:rb-argument}: Backlund--Jensen on a horizontal segment, the argument principle on the symmetric rectangle, and Stirling for $\vartheta$; a disc of fixed radius about $2+iT$ controls only the horizontal piece of the argument variation, and the count up to $T$ comes from the vertical piece and the contour.
\end{remark}

\subsection{Additional finite-section diagnostics}\label{subsec:finite-diagnostics}

The calculations of this subsection solve bounded finite sections directly from the midpoint clock with program \texttt{s13} (Appendix~\CGref{app:code-s13-bounded}; console output in Appendix~\CGref{app:out-s13-bounded}).  The unknowns are the ordinates $\kappa_k$ with $\tau_k\le\Omega+14$ and the free masses $a_n$, $2\le n\le N$; the ordinates beyond are kept on the midpoint clock.  Both are bounded,
\begin{equation}\label{eq:bounded-section}
 |\kappa_k-\tau_k|\le1.2,\qquad |a_n|\le0.5,
\end{equation}
as numerical regularizations, not hypotheses of any theorem.  A test is a $C_c^\infty$ bump of unit integral, proportional to $\exp(-1/(1-v^2))$ in the affine coordinate $v\in(-1,1)$ of its support, multiplied by $\cos(\omega u)$ or $\sin(\omega u)$; the equation is the pairing of the rung system with the test.  For a single section of size $N$ the support is $[0,\log(N+\frac12)]$ and $\omega=0,0.25,\dots,\Omega$.  In the nested family, the block of tests introduced by the section of size $N$ has support $[0.08,\log(N+\frac12)]$ and $\omega=0,1,\dots,\Omega_N$, and the rows of each block are divided by the square root of their number, so that every block carries the same weight.  The normalized residual is the Euclidean residual divided by the square root of the number of equations.  The solver is \texttt{scipy.optimize.least\_squares} in double precision (trust-region reflective method with box constraints, analytic Jacobian, scaling by the Jacobian), with tolerances $\texttt{ftol}=\texttt{xtol}=\texttt{gtol}=10^{-11}$ and at most $1200$ evaluations for the single sections, and $2\cdot10^{-10}$ and at most $3000$ evaluations for the nested ones.  For every solve the program prints the termination status, the number of evaluations, the numerical rank of the Jacobian at the solution (singular values above $10^{-12}$ times the largest) with the condition number on that subspace, and the number of variables at their bounds.  The genuine ordinates and the von Mangoldt masses are read only after the solve, for diagnostics.

\emph{Single sections.}  The sections $N=20,34,50$ ($\Omega=70,110,150$; $22,40,60$ ordinates; $19,33,49$ masses; $561,881,1201$ equations), each solved alone from the clock, end with normalized residuals $6.8\cdot10^{-5}$, $1.4\cdot10^{-4}$, $3.9\cdot10^{-4}$, the last one by the \texttt{ftol} test after $127$ evaluations.  The numerical rank of the Jacobian is one or two below the number of unknowns; the missing directions are the masses $a_n$ with $\log n$ at the edge of the test support, which the tests do not see.  On the remaining subspace the condition number is $1.2\cdot10^4$, $2.3\cdot10^8$, $1.2\cdot10^8$.  No ordinate reaches its bound; one or two masses reach $|a_n|=0.5$.  The first ten ordinates are within $7.5\cdot10^{-5}$, $2.1\cdot10^{-4}$, $4.5\cdot10^{-4}$ (root mean square) of $\gamma_k$, and the masses with $n\le15$ follow $\Lambda(n)/(2\sqrt n)$ (Figure~\CGref{fig:bounded-masses}).  Cross-validation with tests not used in the solve (five supports inside $[0,\log(N+\frac12)]$, frequencies on the shifted grid $(j+\frac12)\cdot0.25\le0.8\,\Omega$) gives root mean square residuals between $5\cdot10^{-5}$ and $8\cdot10^{-4}$.

These endpoint calculations are distinct from the linearized clock problem: in the latter, it is the ordinate-direction coefficients $h_k$ that become large while a residual of roughly $15$--$18\%$ persists.  Imposing bounds there does not remove an artifact; it exposes the lack of an exact solution to that finite tested linear system.

\emph{Nested sections.}  In the nested family every test block, once introduced, stays in all later solves, and each section is warm-started from the previous one ($N=13,20,34,50,70,100$; Figure~\CGref{fig:nested-bounded}).  Up to $N=70$ the old blocks are fitted better as the section grows.  At $N=50$ their root mean square residuals are $4.4\cdot10^{-7}$ to $1.7\cdot10^{-6}$, and the newest block has $1.7\cdot10^{-5}$.  At $N=70$ ($93$ ordinates, $69$ masses, $1185$ equations in five blocks) they are $8.7\cdot10^{-7}$ to $2.4\cdot10^{-6}$, with $6.2\cdot10^{-5}$ on the newest block.  The section $N=100$ ($146$ ordinates, $99$ masses, $1786$ equations in six blocks) does not continue this.  The solver stops by the \texttt{ftol} test after $402$ evaluations, with residuals $1.1\cdot10^{-3}$ to $2.4\cdot10^{-3}$ on the old blocks and $2.2\cdot10^{-2}$ on the newest one, and the first ten ordinates move to $3.7\cdot10^{-3}$ (root mean square) from $\gamma_k$.  The rank deficit grows with the section (one to four edge masses), and the condition number on the determined subspace grows from $10^7$ at $N=13$ to $9\cdot10^{10}$ at $N=100$.  In these runs the nested continuation is observed up to $N=70$ and fails at $N=100$.

\emph{Repeatability.}  The section $N=50$ with the same four blocks was solved from three starting points perturbed from the clock (ordinates by up to $\pm0.5$, masses by up to $\pm0.2$).  The solves end with maximal block residuals $2.7\cdot10^{-4}$, $9.5\cdot10^{-6}$ and $7.6\cdot10^{-6}$, against $1.7\cdot10^{-5}$ for the warm-started solution.  Their ordinates differ from those of that solution by up to $1.40$, $0.99$, $0.62$, and their masses by up to $0.92$, $0.62$, $0.11$.  Bounded finite sections therefore have several near-solutions of comparable residual: the residual alone does not determine the computed ordinates and masses.  The closeness to $\gamma_k$ and to $\Lambda(n)/(2\sqrt n)$ is a property of the continuation from the clock, not of every near-solution.

\begin{figure}[htbp]
\centering
\includegraphics[width=0.78\textwidth]{fig16_nested_bounded_endpoint.png}
\caption{Nested bounded endpoint sections $N=50,70,100$: root mean square residual of each section on every permanent test block introduced so far; the horizontal axis is the size $N$ of the section that introduced the block.  The legend gives, for each section, the number of equations and blocks, the numbers of ordinates and masses, and how many masses are at the box $|a_n|=0.5$.  Up to $N=70$ the old blocks are fitted to about $10^{-6}$; at $N=100$ the solver stops with residuals of order $10^{-3}$ on the old blocks and $2\cdot10^{-2}$ on the newest one.  A floating-point computation, not an existence proof for the infinite rung system.}
\label{fig:nested-bounded}
\end{figure}

\begin{figure}[htbp]
\centering
\includegraphics[width=0.78\textwidth]{fig17_bounded_mass_stability.png}
\caption{Recovered free masses $a_n$, $n\le15$, of the single bounded sections $N=20,34,50$, each solved alone from the clock.  All masses with $n\le15$ are shown, and none of them is at the box $|a_n|=0.5$ (dashed lines; a mass at the box would be circled in red).  The masses at the box are edge masses, listed in the legend: $n=20$ for $N=20$, $n=33,34$ for $N=34$, $n=48$ for $N=50$ (where also $a_{47}=0.499994$).  At $\log n$ the test weights are below $3\cdot10^{-6}$ of their peak, so these masses are determined weakly or not at all by the tests; they lie outside the plotted range.  The values $\Lambda(n)/(2\sqrt n)$ (squares) are a post-solution diagnostic and are not supplied to the solver.}
\label{fig:bounded-masses}
\end{figure}

\medskip\noindent\textbf{Status.}

\emph{Proved.}  Proposition~\CGref{prop:rung-apriori}: for any exact rung solution satisfying its hypotheses, the count follows the clock with a linear drift governed by one constant $\mathrm b$; the ordinates lie within a bounded distance of the drift-corrected clock, and within a bounded distance of the midpoint clock exactly when $\mathrm b=0$ (Corollary~\CGref{cor:rung-apriori-b0}).  Theorem~\CGref{thm:clock-jacobian} gives the multiscale linear information at the clock.  The additive interpolation is exact, and is a level-set reformulation of $\mathcal Q_{\RVM}$.

\emph{Numerical evidence.}  The bounded finite sections of this subsection, in floating point: single sections with normalized residuals of order $10^{-4}$; a nested continuation with permanent blocks up to $N=70$ that fails at $N=100$; near-solutions that are not unique.

\emph{Open.}  An interior path for the exact-atomic interpolation; the midpoint-clock injectivity problem \CGeqref{eq:clock-injectivity}, a closed autonomous rigidity question; an infinite solution of the rung system; the endpoint identity (D).  None of the statements above supplies existence.

\end{document}
